% Auto-generated from data/segments.json + layout.json (+ transforms) — do not edit by hand.
% volume: 15An01
% mode: printing
% segments: 3896
% transform_rules: 8
% toc_marks: 576

% Volume layout defaults (layout.json ``layout``).
\csromanlayoutapply{1}{1.5}{21.6pt}{8pt}{8pt}{65pt}{2.5em}%

\setcsromanpitaka{สุตฺตนฺตปิฏก}
\setcsromannikaya{องฺคุตฺตรนิกาย}
\csromanheader{2}{องฺคุตฺตร\-นิกาย}
\gambhira{เอกก\-นิปาต\-ปาฬิ}
\namakkaram{นโม ตสฺส ภควโต อรหโต สมฺมา\-สมฺพุทฺธสฺส.}
\chapterhead{1. \textbf{รูปาทิวคฺค}}
\proseitem{1}{\csromanfolio{1}เอวํ เม สุตํ\csromandash{}เอกํ สมยํ ภควา สาวตฺถิยํ วิหรติ เชตวเน อนาถ\-ปิณฺฑิกสฺส อาราเม. ตตฺร โข ภควา ภิกฺขู อามนฺเตสิ “ภิกฺขโว”ติ. “ภทนฺเต”ติ เต ภิกฺขู ภควโต ปจฺจสฺโสสุํ. ภควา เอตทโวจ\csromandash{}}

\prose{นาหํ ภิกฺขเว อญฺญํ เอกรูปมฺปิ สมนุปสฺสามิ, ยํ เอวํ ปุริสสฺส จิตฺตํ ปริยาทาย ติฏฺฐติ, ยถยิทํ ภิกฺขเว อิตฺถิรูปํ. อิตฺถิรูปํ ภิกฺขเว ปุริสสฺส จิตฺตํ ปริยาทาย ติฏฺฐตีติ.  ปฐมํ.}

\proseitem{2}{นาหํ ภิกฺขเว อญฺญํ เอกสทฺทมฺปิ สมนุปสฺสามิ, ยํ เอวํ ปุริสสฺส จิตฺตํ ปริยาทาย ติฏฺฐติ, ยถยิทํ ภิกฺขเว อิตฺถิสทฺโท. อิตฺถิสทฺโท ภิกฺขเว ปุริสสฺส จิตฺตํ ปริยาทาย ติฏฺฐตีติ.  ทุติยํ.}

\proseitem{3}{นาหํ ภิกฺขเว อญฺญํ เอกคนฺธมฺปิ สมนุปสฺสามิ, ยํ เอวํ ปุริสสฺส จิตฺตํ ปริยาทาย ติฏฺฐติ, ยถยิทํ ภิกฺขเว อิตฺถิคนฺโธ. อิตฺถิคนฺโธ ภิกฺขเว ปุริสสฺส จิตฺตํ ปริยาทาย ติฏฺฐตีติ.  ตติยํ.}

\proseitem{4}{นาหํ ภิกฺขเว อญฺญํ เอกรสมฺปิ สมนุปสฺสามิ, ยํ เอวํ ปุริสสฺส จิตฺตํ ปริยาทาย ติฏฺฐติ, ยถยิทํ ภิกฺขเว อิตฺถิรโส. อิตฺถิรโส ภิกฺขเว ปุริสสฺส จิตฺตํ ปริยาทาย ติฏฺฐตีติ.  จตุตฺถํ.}

\proseitem{5}{\csromanfolio{2}นาหํ ภิกฺขเว อญฺญํ เอก\-โผฏฺฐพฺพมฺปิ สมนุปสฺสามิ, ยํ เอวํ ปุริสสฺส จิตฺตํ ปริยาทาย ติฏฺฐติ, ยถยิทํ ภิกฺขเว อิตฺถิ\-โผฏฺฐพฺโพ. อิตฺถิ\-โผฏฺฐพฺโพ ภิกฺขเว ปุริสสฺส จิตฺตํ ปริยาทาย ติฏฺฐตีติ.  ปญฺจมํ.}

\proseitem{6}{นาหํ ภิกฺขเว อญฺญํ เอกรูปมฺปิ สมนุปสฺสามิ, ยํ เอวํ อิตฺถิยา จิตฺตํ ปริยาทาย ติฏฺฐติ, ยถยิทํ ภิกฺขเว ปุริสรูปํ. ปุริสรูปํ ภิกฺขเว อิตฺถิยา จิตฺตํ ปริยาทาย ติฏฺฐตีติ.  ฉฏฺฐํ.}

\proseitem{7}{นาหํ ภิกฺขเว อญฺญํ เอกสทฺทมฺปิ สมนุปสฺสามิ, ยํ เอวํ อิตฺถิยา จิตฺตํ ปริยาทาย ติฏฺฐติ, ยถยิทํ ภิกฺขเว ปุริสสทฺโท. ปุริสสทฺโท ภิกฺขเว อิตฺถิยา จิตฺตํ ปริยาทาย ติฏฺฐตีติ.  สตฺตมํ.}

\proseitem{8}{นาหํ ภิกฺขเว อญฺญํ เอกคนฺธมฺปิ สมนุปสฺสามิ, ยํ เอวํ อิตฺถิยา จิตฺตํ ปริยาทาย ติฏฺฐติ, ยถยิทํ ภิกฺขเว ปุริสคนฺโธ. ปุริสคนฺโธ ภิกฺขเว อิตฺถิยา จิตฺตํ ปริยาทาย ติฏฺฐตีติ.  อฏฺฐมํ.}

\proseitem{9}{นาหํ ภิกฺขเว อญฺญํ เอกรสมฺปิ สมนุปสฺสามิ, ยํ เอวํ อิตฺถิยา จิตฺตํ ปริยาทาย ติฏฺฐติ, ยถยิทํ ภิกฺขเว ปุริสรโส. ปุริสรโส ภิกฺขเว อิตฺถิยา จิตฺตํ ปริยาทาย ติฏฺฐตีติ.  นวมํ.}

\proseitemwithclosermidruled{10}{นาหํ ภิกฺขเว อญฺญํ เอก\-โผฏฺฐพฺพมฺปิ สมนุปสฺสามิ, ยํ เอวํ อิตฺถิยา จิตฺตํ ปริยาทาย ติฏฺฐติ, ยถยิทํ ภิกฺขเว ปุริส\-โผฏฺฐพฺโพ. ปุริส\-โผฏฺฐพฺโพ ภิกฺขเว อิตฺถิยา จิตฺตํ ปริยาทาย ติฏฺฐตีติ.  ทสมํ.}{รูปาทิวคฺโค ปฐโม.}
\chapterhead{2. นีวรณ\-ปฺปหานวคฺค}
\proseitem{11}{นาหํ ภิกฺขเว อญฺญํ เอกธมฺมมฺปิ สมนุปสฺสามิ, เยน อนุปฺปนฺโน วา กามจฺฉนฺโท อุปฺปชฺชติ, อุปฺปนฺโน วา กามจฺฉนฺโท ภิยฺโยภาวาย เวปุลฺลาย สํวตฺตติ, ยถยิทํ ภิกฺขเว สุภนิมิตฺตํ. สุภนิมิตฺตํ ภิกฺขเว อโยนิโส มนสิ กโรโต อนุปฺปนฺโน เจว กามจฺฉนฺโท อุปฺปชฺชติ, อุปฺปนฺโน จ กามจฺฉนฺโท ภิยฺโยภาวาย เวปุลฺลาย สํวตฺตตีติ.  ปฐมํ.}

\chapterheadpageread{2. นีวรณ\-ปฺปหานวคฺค}
\proseitem{12}{\csromanfolio{3}นาหํ ภิกฺขเว อญฺญํ เอกธมฺมมฺปิ สมนุปสฺสามิ, เยน อนุปฺปนฺโน วา พฺยาปาโท อุปฺปชฺชติ, อุปฺปนฺโน วา พฺยาปาโท ภิยฺโยภาวาย เวปุลฺลาย สํวตฺตติ, ยถยิทํ ภิกฺขเว ปฏิฆ\-นิมิตฺตํ. ปฏิฆ\-นิมิตฺตํ ภิกฺขเว อโยนิโส มนสิ กโรโต อนุปฺปนฺโน เจว พฺยาปาโท อุปฺปชฺชติ, อุปฺปนฺโน จ พฺยาปาโท ภิยฺโยภาวาย เวปุลฺลาย สํวตฺตตีติ.  ทุติยํ.}

\proseitem{13}{นาหํ ภิกฺขเว อญฺญํ เอกธมฺมมฺปิ สมนุปสฺสามิ, เยน อนุปฺปนฺนํ วา ถินมิทฺธํ\csromansharedfootnote{3302f9c593b385b2}{ถีนมิทฺธํ (สี, สฺยา, กํ, อิ)} อุปฺปชฺชติ, อุปฺปนฺนํ วา ถินมิทฺธํ ภิยฺโยภาวาย เวปุลฺลาย สํวตฺตติ, ยถยิทํ ภิกฺขเว อรติ ตนฺที\csromansharedfootnote{b78f366711bc8adf}{ตนฺทิ (ก)} วิชมฺภิตา\csromansharedfootnote{a9de01f3c239097a}{วิชมฺภิกา (สี, สฺยา, กํ, อิ)} ภตฺตสมฺมโท เจตโส จ ลีนตฺตํ. ลีนจิตฺตสฺส ภิกฺขเว อนุปฺปนฺนญฺเจว ถินมิทฺธํ อุปฺปชฺชติ, อุปฺปนฺนญฺจ ถินมิทฺธํ ภิยฺโยภาวาย เวปุลฺลาย สํวตฺตตีติ.  ตติยํ.}

\proseitem{14}{นาหํ ภิกฺขเว อญฺญํ เอกธมฺมมฺปิ สมนุปสฺสามิ, เยน อนุปฺปนฺนํ วา อุทฺธจฺจ\-กุกฺกุจฺจํ อุปฺปชฺชติ, อุปฺปนฺนํ วา อุทฺธจฺจ\-กุกฺกุจฺจํ ภิยฺโยภาวาย เวปุลฺลาย สํวตฺตติ, ยถยิทํ ภิกฺขเว เจตโส อวูปสโม. อวูปสนฺต\-จิตฺตสฺส ภิกฺขเว อนุปฺปนฺนญฺเจว อุทฺธจฺจ\-กุกฺกุจฺจํ อุปฺปชฺชติ, อุปฺปนฺนญฺจ อุทฺธจฺจ\-กุกฺกุจฺจํ ภิยฺโยภาวาย เวปุลฺลาย สํวตฺตตีติ.  จตุตฺถํ.}

\proseitem{15}{นาหํ ภิกฺขเว อญฺญํ เอกธมฺมมฺปิ สมนุปสฺสามิ, เยน อนุปฺปนฺนา วา วิจิกิจฺฉา อุปฺปชฺชติ, อุปฺปนฺนา วา วิจิกิจฺฉา ภิยฺโยภาวาย เวปุลฺลาย สํวตฺตติ, ยถยิทํ ภิกฺขเว อโยนิโส\-มนสิกาโร. อโยนิโส ภิกฺขเว มนสิ กโรโต อนุปฺปนฺนา เจว วิจิกิจฺฉา อุปฺปชฺชติ, อุปฺปนฺนา จ วิจิกิจฺฉา ภิยฺโยภาวาย เวปุลฺลาย สํวตฺตตีติ.  ปญฺจมํ.}

\proseitem{16}{นาหํ ภิกฺขเว อญฺญํ เอกธมฺมมฺปิ สมนุปสฺสามิ, เยน อนุปฺปนฺโน วา กามจฺฉนฺโท นุปฺปชฺชติ, อุปฺปนฺโน วา กามจฺฉนฺโท ปหียติ, ยถยิทํ ภิกฺขเว อสุภนิมิตฺตํ. อสุภนิมิตฺตํ ภิกฺขเว โยนิโส มนสิ กโรโต อนุปฺปนฺโน เจว กามจฺฉนฺโท นุปฺปชฺชติ, อุปฺปนฺโน จ กามจฺฉนฺโท ปหียตีติ.  ฉฏฺฐํ.}

\proseitem{17}{นาหํ ภิกฺขเว อญฺญํ เอกธมฺมมฺปิ สมนุปสฺสามิ, เยน อนุปฺปนฺโน วา พฺยาปาโท นุปฺปชฺชติ, อุปฺปนฺโน วา พฺยาปาโท ปหียติ, ยถยิทํ ภิกฺขเว \csromanfolio{4}เมตฺตา เจโตวิมุตฺติ. เมตฺตํ ภิกฺขเว เจโตวิมุตฺติํ โยนิโส มนสิ กโรโต อนุปฺปนฺโน เจว พฺยาปาโท นุปฺปชฺชติ, อุปฺปนฺโน จ พฺยาปาโท ปหียตีติ.  }

\proseitem{18}{นาหํ ภิกฺขเว อญฺญํ เอกธมฺมมฺปิ สมนุปสฺสามิ, เยน อนุปฺปนฺนํ วา ถินมิทฺธํ นุปฺปชฺชติ, อุปฺปนฺนํ วา ถินมิทฺธํ ปหียติ, ยถยิทํ ภิกฺขเว อารมฺภธาตุ นิกฺกมธาตุ ปรกฺกมธาตุ. อารทฺธ\-วีริยสฺส ภิกฺขเว อนุปฺปนฺนญฺเจว ถินมิทฺธํ นุปฺปชฺชติ, อุปฺปนฺนญฺจ ถินมิทฺธํ ปหียตีติ.  อฏฺฐมํ.}

\proseitem{19}{นาหํ ภิกฺขเว อญฺญํ เอกธมฺมมฺปิ สมนุปสฺสามิ, เยน อนุปฺปนฺนํ วา อุทฺธจฺจ\-กุกฺกุจฺจํ นุปฺปชฺชติ, อุปฺปนฺนํ วา อุทฺธจฺจ\-กุกฺกุจฺจํ ปหียติ, ยถยิทํ ภิกฺขเว เจตโส วูปสโม. วูปสนฺต\-จิตฺตสฺส ภิกฺขเว อนุปฺปนฺนญฺเจว อุทฺธจฺจ\-กุกฺกุจฺจํ นุปฺปชฺชติ, อุปฺปนฺนญฺจ อุทฺธจฺจ\-กุกฺกุจฺจํ ปหียตีติ.  นวมํ.}

\proseitemwithclosermidruled{20}{นาหํ ภิกฺขเว อญฺญํ เอกธมฺมมฺปิ สมนุปสฺสามิ, เยน อนุปฺปนฺนา วา วิจิกิจฺฉา นุปฺปชฺชติ, อุปฺปนฺนา วา วิจิกิจฺฉา ปหียติ, ยถยิทํ ภิกฺขเว โยนิโส\-มนสิกาโร. โยนิโส ภิกฺขเว มนสิ กโรโต อนุปฺปนฺนา เจว วิจิกิจฺฉา นุปฺปชฺชติ, อุปฺปนฺนา จ วิจิกิจฺฉา ปหียตีติ.  ทสมํ.}{นีวรณปฺปหาน\-วคฺโค ทุติโย.}
\chapterhead{3. อกมฺมนิย\-วคฺค}
\proseitem{21}{นาหํ ภิกฺขเว อญฺญํ เอกธมฺมมฺปิ สมนุปสฺสามิ, ยํ เอวํ อภาวิตํ อกมฺมนิยํ โหติ, ยถยิทํ ภิกฺขเว จิตฺตํ\csromansharedfootnote{a5cab43873aeb12e}{ยถยิทํ จิตฺตํ (สี, อิ) เอวมุปริปิ.}. จิตฺตํ ภิกฺขเว อภาวิตํ อกมฺมนิยํ โหตีติ.  ปฐมํ.}

\proseitem{22}{นาหํ ภิกฺขเว อญฺญํ เอกธมฺมมฺปิ สมนุปสฺสามิ, ยํ เอวํ ภาวิตํ กมฺมนิยํ โหติ, ยถยิทํ ภิกฺขเว จิตฺตํ. จิตฺตํ ภิกฺขเว ภาวิตํ กมฺมนิยํ โหตีติ.  ทุติยํ.}

\chapterheadpageread{3. อกมฺมนิย\-วคฺค}
\proseitem{23}{\csromanfolio{5}นาหํ ภิกฺขเว อญฺญํ เอกธมฺมมฺปิ สมนุปสฺสามิ, ยํ เอวํ อภาวิตํ มหโต อนตฺถาย สํวตฺตติ, ยถยิทํ ภิกฺขเว จิตฺตํ. จิตฺตํ ภิกฺขเว อภาวิตํ มหโต อนตฺถาย สํวตฺตตีติ.  ตติยํ.}

\proseitem{24}{นาหํ ภิกฺขเว อญฺญํ เอกธมฺมมฺปิ สมนุปสฺสามิ, ยํ เอวํ ภาวิตํ มหโต อตฺถาย สํวตฺตติ, ยถยิทํ ภิกฺขเว จิตฺตํ. จิตฺตํ ภิกฺขเว ภาวิตํ มหโต อตฺถาย สํวตฺตตีติ.  จตุตฺถํ.}

\proseitem{25}{นาหํ ภิกฺขเว อญฺญํ เอกธมฺมมฺปิ สมนุปสฺสามิ, ยํ เอวํ อภาวิตํ อปาตุภูตํ มหโต อนตฺถาย สํวตฺตติ, ยถยิทํ ภิกฺขเว จิตฺตํ. จิตฺตํ ภิกฺขเว อภาวิตํ อปาตุภูตํ มหโต อนตฺถาย สํวตฺตตีติ.  ปญฺจมํ.}

\proseitem{26}{นาหํ ภิกฺขเว อญฺญํ เอกธมฺมมฺปิ สมนุปสฺสามิ, ยํ เอวํ ภาวิตํ ปาตุภูตํ มหโต อตฺถาย สํวตฺตติ, ยถยิทํ ภิกฺขเว จิตฺตํ. จิตฺตํ ภิกฺขเว ภาวิตํ ปาตุภูตํ มหโต อตฺถาย สํวตฺตตีติ.  ฉฏฺฐํ.}

\proseitem{27}{นาหํ ภิกฺขเว อญฺญํ เอกธมฺมมฺปิ สมนุปสฺสามิ, ยํ เอวํ อภาวิตํ อพหุลีกตํ มหโต อนตฺถาย สํวตฺตติ, ยถยิทํ ภิกฺขเว จิตฺตํ. จิตฺตํ ภิกฺขเว อภาวิตํ อพหุลีกตํ มหโต อนตฺถาย สํวตฺตตีติ.  สตฺตมํ.}

\proseitem{28}{นาหํ ภิกฺขเว อญฺญํ เอกธมฺมมฺปิ สมนุปสฺสามิ, ยํ เอวํ ภาวิตํ พหุลีกตํ มหโต อตฺถาย สํวตฺตติ, ยถยิทํ ภิกฺขเว จิตฺตํ. จิตฺตํ ภิกฺขเว ภาวิตํ พหุลีกตํ มหโต อตฺถาย สํวตฺตตีติ.  อฏฺฐมํ.}

\proseitem{29}{นาหํ ภิกฺขเว อญฺญํ เอกธมฺมมฺปิ สมนุปสฺสามิ, ยํ เอวํ อภาวิตํ อพหุลีกตํ ทุกฺขาธิวหํ โหติ, ยถยิทํ ภิกฺขเว จิตฺตํ. จิตฺตํ ภิกฺขเว อภาวิตํ อพหุลีกตํ ทุกฺขาธิวหํ โหตีติ.  นวมํ.}

\proseitemwithclosermidruled{30}{นาหํ ภิกฺขเว อญฺญํ เอกธมฺมมฺปิ สมนุปสฺสามิ, ยํ เอวํ ภาวิตํ พหุลีกตํ สุขาธิวหํ โหติ, ยถยิทํ ภิกฺขเว จิตฺตํ. จิตฺตํ ภิกฺขเว ภาวิตํ พหุลีกตํ สุขาธิวหํ โหตีติ.  ทสมํ.}{อกมฺมนิย\-วคฺโค ตติโย.}
\chapterheadpageread{4. \textbf{อทนฺตวคฺค}}
\proseitem{31}{\csromanfolio{6}นาหํ ภิกฺขเว อญฺญํ เอกธมฺมมฺปิ สมนุปสฺสามิ, ยํ เอวํ อทนฺตํ มหโต อนตฺถาย สํวตฺตติ, ยถยิทํ ภิกฺขเว จิตฺตํ. จิตฺตํ ภิกฺขเว อทนฺตํ มหโต อนตฺถาย สํวตฺตตีติ.  ปฐมํ.}

\proseitem{32}{นาหํ ภิกฺขเว อญฺญํ เอกธมฺมมฺปิ สมนุปสฺสามิ, ยํ เอวํ ทนฺตํ มหโต อตฺถาย สํวตฺตติ, ยถยิทํ ภิกฺขเว จิตฺตํ. จิตฺตํ ภิกฺขเว ทนฺตํ มหโต อตฺถาย สํวตฺตตีติ.  ทุติยํ.}

\proseitem{33}{นาหํ ภิกฺขเว อญฺญํ เอกธมฺมมฺปิ สมนุปสฺสามิ, ยํ เอวํ อคุตฺตํ มหโต อนตฺถาย สํวตฺตติ, ยถยิทํ ภิกฺขเว จิตฺตํ. จิตฺตํ ภิกฺขเว อคุตฺตํ มหโต อนตฺถาย สํวตฺตตีติ.  ตติยํ.}

\proseitem{34}{นาหํ ภิกฺขเว อญฺญํ เอกธมฺมมฺปิ สมนุปสฺสามิ, ยํ เอวํ คุตฺตํ มหโต อตฺถาย สํวตฺตติ, ยถยิทํ ภิกฺขเว จิตฺตํ. จิตฺตํ ภิกฺขเว คุตฺตํ มหโต อตฺถาย สํวตฺตตีติ.  จตุตฺถํ.}

\proseitem{35}{นาหํ ภิกฺขเว อญฺญํ เอกธมฺมมฺปิ สมนุปสฺสามิ, ยํ เอวํ อรกฺขิตํ มหโต อนตฺถาย สํวตฺตติ, ยถยิทํ ภิกฺขเว จิตฺตํ. จิตฺตํ ภิกฺขเว อรกฺขิตํ มหโต อนตฺถาย สํวตฺตตีติ.  ปญฺจมํ.}

\proseitem{36}{นาหํ ภิกฺขเว อญฺญํ เอกธมฺมมฺปิ สมนุปสฺสามิ, ยํ เอวํ รกฺขิตํ มหโต อตฺถาย สํวตฺตติ, ยถยิทํ ภิกฺขเว จิตฺตํ. จิตฺตํ ภิกฺขเว รกฺขิตํ มหโต อตฺถาย สํวตฺตตีติ.  ฉฏฺฐํ.}

\proseitem{37}{นาหํ ภิกฺขเว อญฺญํ เอกธมฺมมฺปิ สมนุปสฺสามิ, ยํ เอวํ อสํวุตํ มหโต อนตฺถาย สํวตฺตติ, ยถยิทํ ภิกฺขเว จิตฺตํ. จิตฺตํ ภิกฺขเว อสํวุตํ มหโต อนตฺถาย สํวตฺตตีติ.  }

\proseitem{38}{นาหํ ภิกฺขเว อญฺญํ เอกธมฺมมฺปิ สมนุปสฺสามิ, ยํ เอวํ สํวุตํ มหโต อตฺถาย สํวตฺตติ, ยถยิทํ ภิกฺขเว จิตฺตํ. จิตฺตํ ภิกฺขเว สํวุตํ มหโต อตฺถาย สํวตฺตตีติ.  อฏฺฐมํ.}

\chapterheadpageread{5. ปณิหิต\-อจฺฉ\-วคฺค}
\proseitem{39}{\csromanfolio{7}นาหํ ภิกฺขเว อญฺญํ เอกธมฺมมฺปิ สมนุปสฺสามิ, ยํ เอวํ อทนฺตํ อคุตฺตํ อรกฺขิตํ อสํวุตํ มหโต อนตฺถาย สํวตฺตติ, ยถยิทํ ภิกฺขเว จิตฺตํ. จิตฺตํ ภิกฺขเว อทนฺตํ อคุตฺตํ อรกฺขิตํ อสํวุตํ มหโต อนตฺถาย สํวตฺตตีติ.  นวมํ.}

\proseitemwithclosermidruled{40}{นาหํ ภิกฺขเว อญฺญํ เอกธมฺมมฺปิ สมนุปสฺสามิ, ยํ เอวํ ทนฺตํ คุตฺตํ รกฺขิตํ สํวุตํ มหโต อตฺถาย สํวตฺตติ, ยถยิทํ ภิกฺขเว จิตฺตํ. จิตฺตํ ภิกฺขเว ทนฺตํ คุตฺตํ รกฺขิตํ สํวุตํ มหโต อตฺถาย สํวตฺตตีติ.  ทสมํ.}{อทนฺตวคฺโค จตุตฺโถ.}
\chapterhead{5. ปณิหิต\-อจฺฉ\-วคฺค}
\proseitem{41}{เสยฺยถาปิ ภิกฺขเว สาลิสูกํ วา ยวสูกํ วา มิจฺฉา\-ปณิหิตํ หตฺเถน วา ปาเทน วา อกฺกนฺตํ หตฺถํ วา ปาทํ วา เภจฺฉติ\csromansharedfootnote{e06ddd4521012fde}{ภิชฺชิสฺสติ (สฺยา, กํ, ก), เภชฺชติ (สี) โมคฺคลฺลานพฺยากรณํ ปสฺสิตพฺพํ.}, โลหิตํ วา อุปฺปาเทสฺสตีติ เนตํ ฐานํ วิชฺชติ. ตํ กิสฺส เหตุ? มิจฺฉา\-ปณิหิตตฺตา ภิกฺขเว สูกสฺส. เอวเมวํ โข ภิกฺขเว โส วต ภิกฺขุ มิจฺฉา\-ปณิหิเตน จิตฺเตน อวิชฺชํ เภจฺฉติ, วิชฺชํ อุปฺปาเทสฺสติ, นิพฺพานํ สจฺฉิกริสฺสตีติ เนตํ ฐานํ วิชฺชติ. ตํ กิสฺส เหตุ? มิจฺฉา\-ปณิหิตตฺตา ภิกฺขเว จิตฺตสฺสาติ.  ปฐมํ.}

\proseitem{42}{เสยฺยถาปิ ภิกฺขเว สาลิสูกํ วา ยวสูกํ วา สมฺมาปณิหิตํ หตฺเถน วา ปาเทน วา อกฺกนฺตํ หตฺถํ วา ปาทํ วา เภจฺฉติ, โลหิตํ วา อุปฺปาเทสฺสตีติ ฐานเมตํ วิชฺชติ. ตํ กิสฺส เหตุ? สมฺมา\-ปณิหิตตฺตา ภิกฺขเว สูกสฺส. เอวเมวํ โข ภิกฺขเว โส วต ภิกฺขุ สมฺมา\-ปณิหิเตน จิตฺเตน อวิชฺชํ เภจฺฉติ, วิชฺชํ อุปฺปาเทสฺสติ, นิพฺพานํ สจฺฉิกริสฺสตีติ ฐานเมตํ วิชฺชติ. ตํ กิสฺส เหตุ? สมฺมา\-ปณิหิตตฺตา ภิกฺขเว จิตฺตสฺสาติ.  ทุติยํ.}

\proseitem{43}{\csromanfolio{8}อิธาหํ\csromansharedfootnote{f1313266f46720a0}{อิทาหํ (สี)} ภิกฺขเว เอกจฺจํ ปุคฺคลํ ปทุฏฺฐจิตฺตํ เอวํ เจตสา เจโต ปริจฺจ ปชานามิ, อิมมฺหิ เจ อยํ สมเย ปุคฺคโล กาลํ กเรยฺย, ยถาภตํ นิกฺขิตฺโต, เอวํ นิรเย. ตํ กิสฺส เหตุ? จิตฺตํ หิสฺส ภิกฺขเว ปทุฏฺฐํ. เจโตปโทสเหตุ ปน ภิกฺขเว เอวมิเธกจฺเจ สตฺตา กายสฺส เภทา ปรํ มรณา อปายํ ทุคฺคติํ วินิปาตํ นิรยํ อุปปชฺชนฺตีติ.  ตติยํ.}

\proseitem{44}{อิธาหํ ภิกฺขเว เอกจฺจํ ปุคฺคลํ ปสนฺนจิตฺตํ เอวํ เจตสา เจโต ปริจฺจ ปชานามิ, อิมมฺหิ เจ อยํ สมเย ปุคฺคโล กาลํ กเรยฺย, ยถาภตํ นิกฺขิตฺโต, เอวํ สคฺเค. ตํ กิสฺส เหตุ? จิตฺตํ หิสฺส ภิกฺขเว ปสนฺนํ. เจโตปสาทเหตุ ปน ภิกฺขเว เอวมิเธกจฺเจ สตฺตา กายสฺส เภทา ปรํ มรณา สุคติํ สคฺคํ โลกํ อุปปชฺชนฺตีติ.  จตุตฺถํ.}

\proseitem{45}{เสยฺยถาปิ ภิกฺขเว อุทกรหโท อาวิโล ลุฬิโต กลลีภูโต, ตตฺถ จกฺขุมา ปุริโส ตีเร ฐิโต น ปสฺเสยฺย สิปฺปิ\-สมฺพุกมฺปิ\csromansharedfootnote{a5f8077f9631fcbb}{สิปฺปิกสมฺพุกมฺปิ (ก)} สกฺขร\-กฐลมฺปิ มจฺฉคุมฺพมฺปิ จรนฺตมฺปิ ติฏฺฐนฺตมฺปิ. ตํ กิสฺส เหตุ? อาวิลตฺตา ภิกฺขเว อุทกสฺส. เอวเมวํ โข ภิกฺขเว โส วต ภิกฺขุ อาวิเลน จิตฺเตน อตฺตตฺถํ วา ญสฺสติ, ปรตฺถํ วา ญสฺสติ, อุภยตฺถํ วา ญสฺสติ, อุตฺตริํ วา มนุสฺสธมฺมา อลม\-ริย\-ญาณ\-ทสฺสน\-วิเสสํ สจฺฉิกริสฺสตีติ เนตํ ฐานํ วิชฺชติ. ตํ กิสฺส เหตุ? อาวิลตฺตา ภิกฺขเว จิตฺตสฺสาติ.  ปญฺจมํ.}

\proseitem{46}{เสยฺยถาปิ ภิกฺขเว อุทกรหโท อจฺโฉ วิปฺปสนฺโน อนาวิโล, ตตฺถ จกฺขุมา ปุริโส ตีเร ฐิโต ปสฺเสยฺย สิปฺปิ\-สมฺพุกมฺปิ สกฺขร\-กฐลมฺปิ มจฺฉคุมฺพมฺปิ จรนฺตมฺปิ ติฏฺฐนฺตมฺปิ. ตํ กิสฺส เหตุ? อนาวิลตฺตา ภิกฺขเว อุทกสฺส. เอวเมวํ โข ภิกฺขเว โส วต ภิกฺขุ อนาวิเลน จิตฺเตน อตฺตตฺถํ วา ญสฺสติ, ปรตฺถํ วา ญสฺสติ, อุภยตฺถํ วา ญสฺสติ, อุตฺตริํ วา มนุสฺสธมฺมา อลม\-ริย\-ญาณ\-ทสฺสน\-วิเสสํ สจฺฉิกริสฺสตีติ ฐานเมตํ วิชฺชติ. ตํ กิสฺส เหตุ? อนาวิลตฺตา ภิกฺขเว จิตฺตสฺสาติ.  ฉฏฺฐํ.}

\proseitem{47}{เสยฺยถาปิ ภิกฺขเว ยานิ กานิจิ รุกฺขชาตานํ, ผนฺทโน เตสํ อคฺคมกฺขายติ ยทิทํ มุทุตาย เจว กมฺมญฺญตาย จ. เอวเมวํ โข อหํ}

\chapterheadpageread{6. อจฺฉราสงฺฆาต\-วคฺค}
\prose{\csromanfolio{9}ภิกฺขเว นาญฺญํ เอกธมฺมมฺปิ สมนุปสฺสามิ, ยํ เอวํ ภาวิตํ พหุลีกตํ มุทุ จ โหติ กมฺมญฺญญฺจ, ยถยิทํ จิตฺตํ. จิตฺตํ ภิกฺขเว ภาวิตํ พหุลีกตํ มุทุ จ โหติ กมฺมญฺญญฺจ โหตีติ.  }

\proseitem{48}{นาหํ ภิกฺขเว อญฺญํ เอกธมฺมมฺปิ สมนุปสฺสามิ, ยํ เอวํ ลหุปริวตฺตํ, ยถยิทํ จิตฺตํ. ยาวญฺจิทํ ภิกฺขเว อุปมาปิ น สุกรา ยาวลหุปริวตฺตํ จิตฺตนฺติ.  อฏฺฐมํ.}

\proseitem{49}{ปภสฺสรมิทํ ภิกฺขเว จิตฺตํ, ตญฺจ โข อาคนฺตุเกหิ อุปกฺกิเลเสหิ อุปกฺกิลิฏฺฐนฺติ.  นวมํ.}

\proseitemwithclosermidruled{50}{ปภสฺสรมิทํ ภิกฺขเว จิตฺตํ, ตญฺจ โข อาคนฺตุเกหิ อุปกฺกิเลเสหิ วิปฺปมุตฺตนฺติ.  ทสมํ.}{ปณิหิต\-อจฺฉ\-วคฺโค ปญฺจโม.}
\chapterhead{6. อจฺฉราสงฺฆาต\-วคฺค}
\proseitem{51}{ปภสฺสรมิทํ ภิกฺขเว จิตฺตํ, ตญฺจ โข อาคนฺตุเกหิ อุปกฺกิเลเสหิ อุปกฺกิลิฏฺฐํ, ตํ อสฺสุตวา ปุถุชฺชโน ยถาภูตํ นปฺปชานาติ. ตสฺมา “อสฺสุตวโต ปุถุชฺชนสฺส จิตฺตภาวนา นตฺถี”ติ วทามีติ.  ปฐมํ.}

\proseitem{52}{ปภสฺสรมิทํ ภิกฺขเว จิตฺตํ, ตญฺจ โข อาคนฺตุเกหิ อุปกฺกิเลเสหิ วิปฺปมุตฺตํ, ตํ สุตวา อริยสาวโก ยถาภูตํ ปชานาติ. ตสฺมา “สุตวโต อริยสาวกสฺส จิตฺตภาวนา อตฺถี”ติ วทามีติ.  ทุติยํ.}

\proseitem{53}{อจฺฉราสงฺฆาต\-มตฺตมฺปิ เจ ภิกฺขเว ภิกฺขุ เมตฺตาจิตฺตํ\csromansharedfootnote{1084d34ac3f531d6}{เมตฺตํ จิตฺตํ (สี), เมตฺตจิตฺตํ (สฺยา, กํ, อิ, ก)} อาเสวติ, อยํ วุจฺจติ ภิกฺขเว ภิกฺขุ อริตฺตชฺฌาโน วิหรติ สตฺถุ\-สาสน\-กโร โอวาทปติกโร อโมฆํ รฏฺฐปิณฺฑํ ภุญฺชติ, โก ปน วาโท เย นํ พหุลีกโรนฺตีติ.  ตติยํ.}

\proseitem{54}{\csromanfolio{10}อจฺฉราสงฺฆาต\-มตฺตมฺปิ เจ ภิกฺขเว ภิกฺขุ เมตฺตาจิตฺตํ ภาเวติ, อยํ วุจฺจติ ภิกฺขเว ภิกฺขุ อริตฺตชฺฌาโน วิหรติ สตฺถุ\-สาสน\-กโร โอวาทปติกโร อโมฆํ รฏฺฐปิณฺฑํ ภุญฺชติ, โก ปน วาโท เย นํ พหุลีกโรนฺตีติ.  จตุตฺถํ.}

\proseitem{55}{อจฺฉราสงฺฆาต\-มตฺตมฺปิ เจ ภิกฺขเว ภิกฺขุ เมตฺตาจิตฺตํ มนสิ กโรติ, อยํ วุจฺจติ ภิกฺขเว ภิกฺขุ อริตฺตชฺฌาโน วิหรติ สตฺถุ\-สาสน\-กโร โอวาทปติกโร อโมฆํ รฏฺฐปิณฺฑํ ภุญฺชติ, โก ปน วาโท เย นํ พหุลีกโรนฺตีติ.  ปญฺจมํ.}

\proseitem{56}{เย เกจิ ภิกฺขเว ธมฺมา อกุสลา อกุสลภาคิยา อกุสล\-ปกฺขิกา, สพฺเพ เต มโนปุพฺพงฺคมา, มโน เตสํ ธมฺมานํ ปฐมํ อุปฺปชฺชติ, อนฺวเทว อกุสลา ธมฺมาติ.  ฉฏฺฐํ.}

\proseitem{57}{เย เกจิ ภิกฺขเว ธมฺมา กุสลา กุสลภาคิยา กุสลปกฺขิกา, สพฺเพ เต มโนปุพฺพงฺคมา, มโน เตสํ ธมฺมานํ ปฐมํ อุปฺปชฺชติ, อนฺวเทว กุสลา ธมฺมาติ.  สตฺตมํ.}

\proseitem{58}{นาหํ ภิกฺขเว อญฺญํ เอกธมฺมมฺปิ สมนุปสฺสามิ, เยน อนุปฺปนฺนา วา อกุสลา ธมฺมา อุปฺปชฺชนฺติ, อุปฺปนฺนา วา กุสลา ธมฺมา ปริหายนฺติ, ยถยิทํ ภิกฺขเว ปมาโท. ปมตฺตสฺส ภิกฺขเว อนุปฺปนฺนา เจว อกุสลา ธมฺมา อุปฺปชฺชนฺติ, อุปฺปนฺนา จ กุสลา ธมฺมา ปริหายนฺตีติ.  อฏฺฐมํ.}

\proseitem{59}{นาหํ ภิกฺขเว อญฺญํ เอกธมฺมมฺปิ สมนุปสฺสามิ, เยน อนุปฺปนฺนา วา กุสลา ธมฺมา อุปฺปชฺชนฺติ, อุปฺปนฺนา วา อกุสลา ธมฺมา ปริหายนฺติ, ยถยิทํ ภิกฺขเว อปฺปมาโท. อปฺปมตฺตสฺส ภิกฺขเว อนุปฺปนฺนา เจว กุสลา ธมฺมา อุปฺปชฺชนฺติ, อุปฺปนฺนา จ อกุสลา ธมฺมา ปริหายนฺตีติ.  นวมํ.}

\proseitemwithclosermidruled{60}{นาหํ ภิกฺขเว อญฺญํ เอกธมฺมมฺปิ สมนุปสฺสามิ, เยน อนุปฺปนฺนา วา อกุสลา ธมฺมา อุปฺปชฺชนฺติ, อุปฺปนฺนา วา กุสลา ธมฺมา ปริหายนฺติ, ยถยิทํ ภิกฺขเว โกสชฺชํ. กุสีตสฺส ภิกฺขเว อนุปฺปนฺนา เจว อกุสลา ธมฺมา อุปฺปชฺชนฺติ, อุปฺปนฺนา จ กุสลา ธมฺมา ปริหายนฺตีติ.  ทสมํ.}{อจฺฉราสงฺฆาต\-วคฺโค ฉฏฺโฐ.}
\chapterheadpageread{7. วีริยา\-รมฺภา\-ทิ\-วคฺค}
\chapterhead{7. วีริยา\-รมฺภา\-ทิ\-วคฺค}
\proseitem{61}{\csromanfolio{11}นาหํ ภิกฺขเว อญฺญํ เอกธมฺมมฺปิ สมนุปสฺสามิ, เยน อนุปฺปนฺนา วา กุสลา ธมฺมา อุปฺปชฺชนฺติ, อุปฺปนฺนา วา อกุสลา ธมฺมา ปริหายนฺติ, ยถยิทํ ภิกฺขเว วีริยารมฺโภ\csromansharedfootnote{1bf2eb862ccf212e}{วิริยารมฺโภ (สี, สฺยา, กํ, อิ)}. อารทฺธ\-วีริยสฺส ภิกฺขเว อนุปฺปนฺนา เจว กุสลา ธมฺมา อุปฺปชฺชนฺติ, อุปฺปนฺนา จ อกุสลา ธมฺมา ปริหายนฺตีติ.  ปฐมํ.}

\proseitem{62}{นาหํ ภิกฺขเว อญฺญํ เอกธมฺมมฺปิ สมนุปสฺสามิ, เยน อนุปฺปนฺนา วา อกุสลา ธมฺมา อุปฺปชฺชนฺติ, อุปฺปนฺนา วา กุสลา ธมฺมา ปริหายนฺติ, ยถยิทํ ภิกฺขเว มหิจฺฉตา. มหิจฺฉสฺส ภิกฺขเว อนุปฺปนฺนา เจว อกุสลา ธมฺมา อุปฺปชฺชนฺติ, อุปฺปนฺนา จ กุสลา ธมฺมา ปริหายนฺตีติ.  ทุติยํ.}

\proseitem{63}{นาหํ ภิกฺขเว อญฺญํ เอกธมฺมมฺปิ สมนุปสฺสามิ, เยน อนุปฺปนฺนา วา กุสลา ธมฺมา อุปฺปชฺชนฺติ, อุปฺปนฺนา วา อกุสลา ธมฺมา ปริหายนฺติ, ยถยิทํ ภิกฺขเว อปฺปิจฺฉตา. อปฺปิจฺฉสฺส ภิกฺขเว อนุปฺปนฺนา เจว กุสลา ธมฺมา อุปฺปชฺชนฺติ, อุปฺปนฺนา จ อกุสลา ธมฺมา ปริหายนฺตีติ.  ตติยํ.}

\proseitem{64}{นาหํ ภิกฺขเว อญฺญํ เอกธมฺมมฺปิ สมนุปสฺสามิ, เยน อนุปฺปนฺนา วา อกุสลา ธมฺมา อุปฺปชฺชนฺติ, อุปฺปนฺนา วา กุสลา ธมฺมา ปริหายนฺติ, ยถยิทํ ภิกฺขเว อสนฺตุฏฺฐิตา. อสนฺตุฏฺฐสฺส ภิกฺขเว อนุปฺปนฺนา เจว อกุสลา ธมฺมา อุปฺปชฺชนฺติ, อุปฺปนฺนา จ กุสลา ธมฺมา ปริหายนฺตีติ.  จตุตฺถํ.}

\proseitem{65}{นาหํ ภิกฺขเว อญฺญํ เอกธมฺมมฺปิ สมนุปสฺสามิ, เยน อนุปฺปนฺนา วา กุสลา ธมฺมา อุปชฺชนฺติ, อุปฺปนฺนา วา อกุสลา ธมฺมา ปริหายนฺติ, ยถยิทํ ภิกฺขเว สนฺตุฏฺฐิตา. สนฺตุฏฺฐสฺส ภิกฺขเว อนุปฺปนฺนา เจว กุสลา ธมฺมา อุปฺปชฺชนฺติ, อุปฺปนฺนา จ อกุสลา ธมฺมา ปริหายนฺตีติ.  ปญฺจมํ.}

\proseitem{66}{นาหํ ภิกฺขเว อญฺญํ เอกธมฺมมฺปิ สมนุปสฺสามิ, เยน อนุปฺปนฺนา วา อกุสลา ธมฺมา อุปฺปชฺชนฺติ, อุปฺปนฺนา วา กุสลา ธมฺมา ปริหายนฺติ, ยถยิทํ ภิกฺขเว อโยนิโส\-มนสิกาโร. อโยนิโส ภิกฺขเว มนสิ กโรโต อนุปฺปนฺนา เจว อกุสลา ธมฺมา อุปฺปชฺชนฺติ, อุปฺปนฺนา จ กุสลา ธมฺมา ปริหายนฺตีติ.  ฉฏฺฐํ.}

\proseitem{67}{\csromanfolio{12}นาหํ ภิกฺขเว อญฺญํ เอกธมฺมมฺปิ สมนุปสฺสามิ, เยน อนุปฺปนฺนา วา กุสลา ธมฺมา อุปฺปชฺชนฺติ, อุปฺปนฺนา วา อกุสลา ธมฺมา ปริหายนฺติ, ยถยิทํ ภิกฺขเว โยนิโส\-มนสิกาโร. โยนิโส ภิกฺขเว มนสิ กโรโต อนุปฺปนฺนา เจว กุสลา ธมฺมา อุปฺปชฺชนฺติ, อุปฺปนฺนา จ อกุสลา ธมฺมา ปริหายนฺตีติ.  สตฺตมํ.}

\proseitem{68}{นาหํ ภิกฺขเว อญฺญํ เอกธมฺมมฺปิ สมนุปสฺสามิ, เยน อนุปฺปนฺนา วา อกุสลา ธมฺมา อุปฺปชฺชนฺติ, อุปฺปนฺนา วา กุสลา ธมฺมา ปริหายนฺติ, ยถยิทํ ภิกฺขเว อสมฺปชญฺญํ. อสมฺปชานสฺส ภิกฺขเว อนุปฺปนฺนา เจว อกุสลา ธมฺมา อุปฺปชฺชนฺติ, อุปฺปนฺนา จ กุสลา ธมฺมา ปริหายนฺตีติ.  อฏฺฐมํ.}

\proseitem{69}{นาหํ ภิกฺขเว อญฺญํ เอกธมฺมมฺปิ สมนุปสฺสามิ, เยน อนุปฺปนฺนา วา กุสลา ธมฺมา อุปฺปชฺชนฺติ, อุปฺปนฺนา วา อกุสลา ธมฺมา ปริหายนฺติ, ยถยิทํ ภิกฺขเว สมฺปชญฺญํ. สมฺปชานสฺส ภิกฺขเว อนุปฺปนฺนา เจว กุสลา ธมฺมา อุปฺปชฺชนฺติ, อุปฺปนฺนา จ อกุสลา ธมฺมา ปริหายนฺตีติ.  นวมํ.}

\proseitemwithclosermidruled{70}{นาหํ ภิกฺขเว อญฺญํ เอกธมฺมมฺปิ สมนุปสฺสามิ, เยน อนุปฺปนฺนา วา อกุสลา ธมฺมา อุปฺปชฺชนฺติ, อุปฺปนฺนา วา กุสลา ธมฺมา ปริหายนฺติ, ยถยิทํ ภิกฺขเว ปาปมิตฺตตา. ปาปมิตฺตสฺส ภิกฺขเว อนุปฺปนฺนา เจว อกุสลา ธมฺมา อุปฺปชฺชนฺติ, อุปฺปนฺนา จ กุสลา ธมฺมา ปริหายนฺตีติ.  ทสมํ.}{วีริยา\-รมฺภา\-ทิ\-วคฺโค สตฺตโม.}
\chapterhead{8. กลฺยาณมิตฺตา\-ทิ\-วคฺค}
\proseitem{71}{นาหํ ภิกฺขเว อญฺญํ เอกธมฺมมฺปิ สมนุปสฺสามิ, เยน อนุปฺปนฺนา วา กุสลา ธมฺมา อุปฺปชฺชนฺติ, อุปฺปนฺนา วา อกุสลา ธมฺมา ปริหายนฺติ, ยถยิทํ ภิกฺขเว กลฺยาณมิตฺตตา. กลฺยาณ\-มิตฺตสฺส ภิกฺขเว อนุปฺปนฺนา เจว กุสลา ธมฺมา อุปฺปชฺชนฺติ, อุปฺปนฺนา จ อกุสลา ธมฺมา ปริหายนฺตีติ.  ปฐมํ.}

\proseitem{72}{นาหํ ภิกฺขเว อญฺญํ เอกธมฺมมฺปิ สมนุปสฺสามิ, เยน อนุปฺปนฺนา วา อกุสลา ธมฺมา อุปฺปชฺชนฺติ, อุปฺปนฺนา วา กุสลา ธมฺมา ปริหายนฺติ, ยถยิทํ ภิกฺขเว อนุโยโค อกุสลานํ ธมฺมานํ, อนนุโยโค}

\chapterheadpageread{8. กลฺยาณมิตฺตา\-ทิ\-วคฺค}
\prose{\csromanfolio{13}กุสลานํ ธมฺมานํ. อนุโยคา ภิกฺขเว อกุสลานํ ธมฺมานํ, อนนุโยคา กุสลานํ ธมฺมานํ อนุปฺปนฺนา เจว อกุสลา ธมฺมา อุปฺปชฺชนฺติ, อุปฺปนฺนา จ กุสลา ธมฺมา ปริหายนฺตีติ.  ทุติยํ.}

\proseitem{73}{นาหํ ภิกฺขเว อญฺญํ เอกธมฺมมฺปิ สมนุปสฺสามิ, เยน อนุปฺปนฺนา วา กุสลา ธมฺมา อุปฺปชฺชนฺติ, อุปฺปนฺนา วา อกุสลา ธมฺมา ปริหายนฺติ, ยถยิทํ ภิกฺขเว อนุโยโค กุสลานํ ธมฺมานํ, อนนุโยโค อกุสลานํ ธมฺมานํ. อนุโยคา ภิกฺขเว กุสลานํ ธมฺมานํ, อนนุโยคา อกุสลานํ ธมฺมานํ อนุปฺปนฺนา เจว กุสลา ธมฺมา อุปฺปชฺชนฺติ, อุปฺปนฺนา จ อกุสลา ธมฺมา ปริหายนฺตีติ.  ตติยํ.}

\proseitem{74}{นาหํ ภิกฺขเว อญฺญํ เอกธมฺมมฺปิ สมนุปสฺสามิ, เยน อนุปฺปนฺนา วา โพชฺฌงฺคา นุปฺปชฺชนฺติ, อุปฺปนฺนา วา โพชฺฌงฺคา น ภาวนา\-ปาริปูริํ คจฺฉนฺติ, ยถยิทํ ภิกฺขเว อโยนิโส\-มนสิกาโร. อโยนิโส ภิกฺขเว มนสิ กโรโต อนุปฺปนฺนา เจว โพชฺฌงฺคา นุปฺปชฺชนฺติ, อุปฺปนฺนา จ โพชฺฌงฺคา น ภาวนา\-ปาริปูริํ คจฺฉนฺตีติ.  จตุตฺถํ.}

\proseitem{75}{นาหํ ภิกฺขเว อญฺญํ เอกธมฺมมฺปิ สมนุปสฺสามิ, เยน อนุปฺปนฺนา วา โพชฺฌงฺคา อุปฺปชฺชนฺติ, อุปฺปนฺนา วา โพชฺฌงฺคา ภาวนา\-ปาริปูริํ คจฺฉนฺติ, ยถยิทํ ภิกฺขเว โยนิโส มนสิกาโร. โยนิโส ภิกฺขเว มนสิ กโรโต อนุปฺปนฺนา เจว โพชฺฌงฺคา อุปฺปชฺชนฺติ, อุปฺปนฺนา จ โพชฺฌงฺคา ภาวนา\-ปาริปูริํ คจฺฉนฺตีติ.  ปญฺจมํ.}

\proseitem{76}{อปฺปมตฺติกา เอสา ภิกฺขเว ปริหานิ ยทิทํ ญาติปริหานิ. เอตํ ปติกิฏฺฐํ ภิกฺขเว ปริหานีนํ ยทิทํ ปญฺญาปริหานีติ.  ฉฏฺฐํ.}

\proseitem{77}{อปฺปมตฺติกา เอสา ภิกฺขเว วุทฺธิ ยทิทํ ญาติวุทฺธิ. เอตทคฺคํ ภิกฺขเว วุทฺธีนํ ยทิทํ ปญฺญาวุทฺธิ. ตสฺมาติห ภิกฺขเว เอวํ สิกฺขิตพฺพํ “ปญฺญาวุทฺธิยา วทฺธิสฺสามา”ติ. เอวํ หิ โว ภิกฺขเว สิกฺขิตพฺพนฺติ.  สตฺตมํ.}

\proseitem{78}{อปฺปมตฺติกา เอสา ภิกฺขเว ปริหานิ ยทิทํ โภคปริหานิ. เอตํ ปติกิฏฺฐํ ภิกฺขเว ปริหานีนํ ยทิทํ ปญฺญาปริหานีติ.  อฏฺฐมํ.}

\proseitem{79}{\csromanfolio{14}อปฺปมตฺติกา เอสา ภิกฺขเว วุทฺธิ ยทิทํ โภควุทฺธิ. เอตทคฺคํ ภิกฺขเว วุทฺธีนํ ยทิทํ ปญฺญาวุทฺธิ. ตสฺมาติห ภิกฺขเว เอวํ สิกฺขิตพฺพํ “ปญฺญาวุทฺธิยา วทฺธิสฺสามา”ติ. เอวํ หิ โว ภิกฺขเว สิกฺขิตพฺพนฺติ.  นวมํ.}

\proseitemwithclosermidruled{80}{อปฺปมตฺติกา เอสา ภิกฺขเว ปริหานิ ยทิทํ ยโสปริหานิ. เอตํ ปติกิฏฺฐํ ภิกฺขเว ปริหานีนํ ยทิทํ ปญฺญาปริหานีติ.  ทสมํ.}{กลฺยาณมิตฺตา\-ทิ\-วคฺโค อฏฺฐโม.}
\chapterhead{9. \textbf{ปมาทาทิวคฺค}}
\proseitem{81}{อปฺปมตฺติกา เอสา ภิกฺขเว วุทฺธิ ยทิทํ ยโสวุทฺธิ. เอตทคฺคํ ภิกฺขเว วุทฺธีนํ ยทิทํ ปญฺญาวุทฺธิ. ตสฺมาติห ภิกฺขเว เอวํ สิกฺขิตพฺพํ “ปญฺญาวุทฺธิยา วทฺธิสฺสามา”ติ. เอวํ หิ โว ภิกฺขเว สิกฺขิตพฺพนฺติ.  ปฐมํ.}

\proseitem{82}{นาหํ ภิกฺขเว อญฺญํ เอกธมฺมมฺปิ สมนุปสฺสามิ, โย เอวํ มหโต อนตฺถาย สํวตฺตติ, ยถยิทํ ภิกฺขเว ปมาโท. ปมาโท ภิกฺขเว มหโต อนตฺถาย สํวตฺตตีติ.  ทุติยํ.}

\proseitem{83}{นาหํ ภิกฺขเว อญฺญํ เอกธมฺมมฺปิ สมนุปสฺสามิ, โย เอวํ มหโต อตฺถาย สํวตฺตติ, ยถยิทํ ภิกฺขเว อปฺปมาโท. อปฺปมาโท ภิกฺขเว มหโต อตฺถาย สํวตฺตตีติ.  ตติยํ.}

\proseitem{84}{นาหํ ภิกฺขเว อญฺญํ เอกธมฺมมฺปิ สมนุปสฺสามิ, โย เอวํ มหโต อนตฺถาย สํวตฺตติ, ยถยิทํ ภิกฺขเว โกสชฺชํ. โกสชฺชํ ภิกฺขเว มหโต อนตฺถาย สํวตฺตตีติ.  จตุตฺถํ.}

\proseitem{85}{นาหํ ภิกฺขเว อญฺญํ เอกธมฺมมฺปิ สมนุปสฺสามิ, โย เอวํ มหโต อตฺถาย สํวตฺตติ, ยถยิทํ ภิกฺขเว วีริยารมฺโภ. วีริยารมฺโภ ภิกฺขเว มหโต อตฺถาย สํวตฺตตีติ.  ปญฺจมํ.}

\proseitem{86}{นาหํ ภิกฺขเว อญฺญํ เอกธมฺมมฺปิ สมนุปสฺสามิ, โย เอวํ มหโต อนตฺถาย สํวตฺตติ, ยถยิทํ ภิกฺขเว มหิจฺฉตา. มหิจฺฉตา ภิกฺขเว มหโต อนตฺถาย สํวตฺตตีติ.  ฉฏฺฐํ.}

\chapterheadpageread{9. ปมาทาทิวคฺค}
\proseitem{87}{\csromanfolio{15}นาหํ ภิกฺขเว อญฺญํ เอกธมฺมมฺปิ สมนุปสฺสามิ, โย เอวํ มหโต อตฺถาย สํวตฺตติ, ยถยิทํ ภิกฺขเว อปฺปิจฺฉตา. อปฺปิจฺฉตา ภิกฺขเว มหโต อตฺถาย สํวตฺตตีติ.  สตฺตมํ.}

\proseitem{88}{นาหํ ภิกฺขเว อญฺญํ เอกธมฺมมฺปิ สมนุปสฺสามิ, โย เอวํ มหโต อนตฺถาย สํวตฺตติ, ยถยิทํ ภิกฺขเว อสนฺตุฏฺฐิตา. อสนฺตุฏฺฐิตา ภิกฺขเว มหโต อนตฺถาย สํวตฺตตีติ.  อฏฺฐมํ.}

\proseitem{89}{นาหํ ภิกฺขเว อญฺญํ เอกธมฺมมฺปิ สมนุปสฺสามิ, โย เอวํ มหโต อตฺถาย สํวตฺตติ, ยถยิทํ ภิกฺขเว สนฺตุฏฺฐิตา. สนฺตุฏฺฐิตา ภิกฺขเว มหโต อตฺถาย สํวตฺตตีติ.  นวมํ.}

\proseitem{90}{นาหํ ภิกฺขเว อญฺญํ เอกธมฺมมฺปิ สมนุปสฺสามิ, โย เอวํ มหโต อนตฺถาย สํวตฺตติ, ยถยิทํ ภิกฺขเว อโยนิโส\-มนสิกาโร. อโยนิโส\-มนสิกาโร ภิกฺขเว มหโต อนตฺถาย สํวตฺตตีติ.  ทสมํ.}

\proseitem{91}{นาหํ ภิกฺขเว อญฺญํ เอกธมฺมมฺปิ สมนุปสฺสามิ, โย เอวํ มหโต อตฺถาย สํวตฺตติ, ยถยิทํ ภิกฺขเว โยนิโส\-มนสิกาโร. โยนิโส\-มนสิกาโร ภิกฺขเว มหโต อตฺถาย สํวตฺตตีติ.  เอกาทสมํ.}

\proseitem{92}{นาหํ ภิกฺขเว อญฺญํ เอกธมฺมมฺปิ สมนุปสฺสามิ, โย เอวํ มหโต อนตฺถาย สํวตฺตติ, ยถยิทํ ภิกฺขเว อสมฺปชญฺญํ. อสมฺปชญฺญํ ภิกฺขเว มหโต อนตฺถาย สํวตฺตตีติ.  ทฺวาทสมํ.}

\proseitem{93}{นาหํ ภิกฺขเว อญฺญํ เอกธมฺมมฺปิ สมนุปสฺสามิ, โย เอวํ มหโต อตฺถาย สํวตฺตติ, ยถยิทํ ภิกฺขเว สมฺปชญฺญํ. สมฺปชญฺญํ ภิกฺขเว มหโต อตฺถาย สํวตฺตตีติ.  เตรสมํ.}

\proseitem{94}{นาหํ ภิกฺขเว อญฺญํ เอกธมฺมมฺปิ สมนุปสฺสามิ, โย เอวํ มหโต อนตฺถาย สํวตฺตติ, ยถยิทํ ภิกฺขเว ปาปมิตฺตตา. ปาปมิตฺตตา ภิกฺขเว มหโต อนตฺถาย สํวตฺตตีติ.  จุทฺทสมํ.}

\proseitem{95}{\csromanfolio{16}นาหํ ภิกฺขเว อญฺญํ เอกธมฺมมฺปิ สมนุปสฺสามิ, โย เอวํ มหโต อตฺถาย สํวตฺตติ, ยถยิทํ ภิกฺขเว กลฺยาณมิตฺตตา. กลฺยาณมิตฺตตา ภิกฺขเว มหโต อตฺถาย สํวตฺตตีติ.  ปนฺนรสมํ.}

\proseitem{96}{นาหํ ภิกฺขเว อญฺญํ เอกธมฺมมฺปิ สมนุปสฺสามิ, โย เอวํ มหโต อนตฺถาย สํวตฺตติ, ยถยิทํ ภิกฺขเว อนุโยโค อกุสลานํ ธมฺมานํ, อนนุโยโค กุสลานํ ธมฺมานํ. อนุโยโค ภิกฺขเว อกุสลานํ ธมฺมานํ, อนนุโยโค กุสลานํ ธมฺมานํ มหโต อนตฺถาย สํวตฺตตีติ.  โสฬสมํ.}

\proseitemwithclosermidruled{97}{นาหํ ภิกฺขเว อญฺญํ เอกธมฺมมฺปิ สมนุปสฺสามิ, โย เอวํ มหโต อตฺถาย สํวตฺตติ, ยถยิทํ ภิกฺขเว อนุโยโค กุสลานํ ธมฺมานํ, อนนุโยโค อกุสลานํ ธมฺมานํ. อนุโยโค ภิกฺขเว กุสลานํ ธมฺมานํ อนนุโยโค อกุสลานํ ธมฺมานํ มหโต อตฺถาย สํวตฺตตีติ.  สตฺตรสมํ.}{ปมาทาทิวคฺโค นวโม.}
\chapterhead{10. ทุติย\-ปมาทา\-ทิ\-วคฺค}
\proseitem{98}{อชฺฌตฺติกํ ภิกฺขเว องฺคนฺติ กริตฺวา นาญฺญํ เอกงฺคมฺปิ สมนุปสฺสามิ, ยํ เอวํ มหโต อนตฺถาย สํวตฺตติ, ยถยิทํ ภิกฺขเว ปมาโท. ปมาโท ภิกฺขเว มหโต อนตฺถาย สํวตฺตตีติ.  ปฐมํ.}

\proseitem{99}{อชฺฌตฺติกํ ภิกฺขเว องฺคนฺติ กริตฺวา นาญฺญํ เอกงฺคมฺปิ สมนุปสฺสามิ, ยํ เอวํ มหโต อตฺถาย สํวตฺตติ, ยถยิทํ ภิกฺขเว อปฺปมาโท. อปฺปมาโท ภิกฺขเว มหโต อตฺถาย สํวตฺตตีติ.  ทุติยํ.}

\proseitem{100}{อชฺฌตฺติกํ ภิกฺขเว องฺคนฺติ กริตฺวา นาญฺญํ เอกงฺคมฺปิ สมนุปสฺสามิ, ยํ เอวํ มหโต อนตฺถาย สํวตฺตติ, ยถยิทํ ภิกฺขเว โกสชฺชํ. โกสชฺชํ ภิกฺขเว มหโต อนตฺถาย สํวตฺตตีติ.  ตติยํ.}

\proseitem{101}{อชฺฌตฺติกํ ภิกฺขเว องฺคนฺติ กริตฺวา นาญฺญํ เอกงฺคมฺปิ สมนุปสฺสามิ, ยํ เอวํ มหโต อตฺถาย สํวตฺตติ, ยถยิทํ ภิกฺขเว วีริยารมฺโภ. วีริยารมฺโภ ภิกฺขเว มหโต อตฺถาย สํวตฺตตีติ.  จตุตฺถํ.}

\chapterheadpageread{10. ทุติย\-ปมาทา\-ทิ\-วคฺค}
\proseitem{102-109}{\csromanfolio{17}อชฺฌตฺติกํ ภิกฺขเว องฺคนฺติ กริตฺวา นาญฺญํ เอกงฺคมฺปิ สมนุปสฺสามิ, ยํ เอวํ มหโต อนตฺถาย สํวตฺตติ, ยถยิทํ ภิกฺขเว มหิจฺฉตา ฯเปฯ อปฺปิจฺฉตา. อสนฺตุฏฺฐิตา. สนฺตุฏฺฐิตา. อโยนิโส\-มนสิกาโร. โยนิโส\-มนสิกาโร. อสมฺปชญฺญํ. สมฺปชญฺญํ. ทฺวาทสมํ.}

\proseitem{110}{พาหิรํ ภิกฺขเว องฺคนฺติ กริตฺวา นาญฺญํ เอกงฺคมฺปิ สมนุปสฺสามิ, ยํ เอวํ มหโต อนตฺถาย สํวตฺตติ, ยถยิทํ ภิกฺขเว ปาปมิตฺตตา. ปาปมิตฺตตา ภิกฺขเว มหโต อนตฺถาย สํวตฺตตีติ.  เตรสมํ.}

\proseitem{111}{พาหิรํ ภิกฺขเว องฺคนฺติ กริตฺวา นาญฺญํ เอกงฺคมฺปิ สมนุปสฺสามิ, ยํ เอวํ มหโต อตฺถาย สํวตฺตติ, ยถยิทํ ภิกฺขเว กลฺยาณมิตฺตตา. กลฺยาณมิตฺตตา ภิกฺขเว มหโต อตฺถาย สํวตฺตตีติ.  จุทฺทสมํ.}

\proseitem{112}{อชฺฌตฺติกํ ภิกฺขเว องฺคนฺติ กริตฺวา นาญฺญํ เอกงฺคมฺปิ สมนุปสฺสามิ, ยํ เอวํ มหโต อนตฺถาย สํวตฺตติ, ยถยิทํ ภิกฺขเว อนุโยโค อกุสลานํ ธมฺมานํ, อนนุโยโค กุสลานํ ธมฺมานํ. อนุโยโค ภิกฺขเว อกุสลานํ ธมฺมานํ, อนนุโยโค กุสลานํ ธมฺมานํ มหโต อนตฺถาย สํวตฺตตีติ.  ปนฺนรสมํ.}

\proseitem{113}{อชฺฌตฺติกํ ภิกฺขเว องฺคนฺติ กริตฺวา นาญฺญํ เอกงฺคมฺปิ สมนุปสฺสามิ, ยํ เอวํ มหโต อตฺถาย สํวตฺตติ, ยถยิทํ ภิกฺขเว อนุโยโค กุสลานํ ธมฺมานํ, อนนุโยโค อกุสลานํ ธมฺมานํ. อนุโยโค ภิกฺขเว กุสลานํ ธมฺมานํ, อนนุโยโค อกุสลานํ ธมฺมานํ มหโต อตฺถาย สํวตฺตตีติ.  โสฬสมํ.}

\proseitem{114}{นาหํ ภิกฺขเว อญฺญํ เอกธมฺมมฺปิ สมนุปสฺสามิ, โย เอวํ สทฺธมฺมสฺส สมฺโมสาย อนฺตรธานาย สํวตฺตติ, ยถยิทํ ภิกฺขเว ปมาโท. ปมาโท ภิกฺขเว สทฺธมฺมสฺส สมฺโมสาย อนฺตรธานาย สํวตฺตตีติ.  สตฺตรสมํ.}

\proseitem{115}{นาหํ ภิกฺขเว อญฺญํ เอกธมฺมมฺปิ สมนุปสฺสามิ, โย เอวํ สทฺธมฺมสฺส ฐิติยา อสมฺโมสาย อนนฺตรธานาย สํวตฺตติ, ยถยิทํ ภิกฺขเว อปฺปมาโท. อปฺปมาโท ภิกฺขเว สทฺธมฺมสฺส ฐิติยา อสมฺโมสาย อนนฺตรธานาย สํวตฺตตีติ.  อฏฺฐารสมํ.}

\proseitem{116}{\csromanfolio{18}นาหํ ภิกฺขเว อญฺญํ เอกธมฺมมฺปิ สมนุปสฺสามิ, โย เอวํ สทฺธมฺมสฺส สมฺโมสาย อนฺตรธานาย สํวตฺตติ, ยถยิทํ ภิกฺขเว โกสชฺชํ. โกสชฺชํ ภิกฺขเว สทฺธมฺมสฺส สมฺโมสาย อนฺตรธานาย สํวตฺตตีติ.  เอกูนวีสติมํ.}

\proseitem{117}{นาหํ ภิกฺขเว อญฺญํ เอกธมฺมมฺปิ สมนุปสฺสามิ, โย เอวํ สทฺธมฺมสฺส ฐิติยา อสมฺโมสาย อนนฺตรธานาย สํวตฺตติ, ยถยิทํ ภิกฺขเว วีริยารมฺโภ. วีริยารมฺโภ ภิกฺขเว สทฺธมฺมสฺส ฐิติยา อสมฺโมสาย อนนฺตรธานาย สํวตฺตตีติ.  วีสติมํ.}

\proseitem{118-128}{นาหํ ภิกฺขเว อญฺญํ เอกธมฺมมฺปิ สมนุปสฺสามิ, โย เอวํ สทฺธมฺมสฺส สมฺโมสาย อนฺตรธานาย สํวตฺตติ, ยถยิทํ ภิกฺขเว มหิจฺฉตา ฯเปฯ อปฺปิจฺฉตา. อสนฺตุฏฺฐิตา. สนฺตุฏฺฐิตา. อโยนิโส\-มนสิกาโร. โยนิโส\-มนสิกาโร. อสมฺปชญฺญํ. สมฺปชญฺญํ. ปาปมิตฺตตา. กลฺยาณมิตฺตตา. อนุโยโค อกุสลานํ ธมฺมานํ, อนนุโยโค กุสลานํ ธมฺมานํ. อนุโยโค ภิกฺขเว อกุสลานํ ธมฺมานํ, อนนุโยโค กุสลานํ ธมฺมานํ สทฺธมฺมสฺส สมฺโมสาย อนฺตรธานาย สํวตฺตตีติ.  เอกตฺติํสติมํ.}

\proseitem{129}{นาหํ ภิกฺขเว อญฺญํ เอกธมฺมมฺปิ สมนุปสฺสามิ, โย เอวํ สทฺธมฺมสฺส ฐิติยา อสมฺโมสาย อนนฺตรธานาย สํวตฺตติ, ยถยิทํ ภิกฺขเว อนุโยโค กุสลานํ ธมฺมานํ, อนนุโยโค อกุสลานํ ธมฺมานํ. อนุโยโค ภิกฺขเว กุสลานํ ธมฺมานํ, อนนุโยโค อกุสลานํ ธมฺมานํ สทฺธมฺมสฺส ฐิติยา อสมฺโมสาย อนนฺตรธานาย สํวตฺตตีติ.  พาตฺติํสติมํ.  จตุกฺโกฏิกํ นิฏฺฐิตํ.}

\proseitem{130}{เย เต ภิกฺขเว ภิกฺขู อธมฺมํ “ธมฺโม”ติ ทีเปนฺติ, เต ภิกฺขเว ภิกฺขู พหุชน\-อหิตาย ปฏิปนฺนา พหุชน\-อสุขาย พหุโน ชนสฺส อนตฺถาย อหิตาย ทุกฺขาย เทวมนุสฺสานํ. พหุญฺจ เต ภิกฺขเว ภิกฺขู อปุญฺญํ ปสวนฺติ, เต จิมํ\csromansharedfootnote{b7980e615c0a6ac5}{เตปิมํ (สี)} สทฺธมฺมํ อนฺตรธาเปนฺตีติ.  เตตฺติํสติมํ.}

\proseitem{131}{เย เต ภิกฺขเว ภิกฺขู ธมฺมํ “อธมฺโม”ติ ทีเปนฺติ, เต ภิกฺขเว ภิกฺขู พหุชน\-อหิตาย ปฏิปนฺนา พหุชน\-อสุขาย พหุโน ชนสฺส}

\vspace{\tipitakaparskip}
\csromancenter{อธมฺมวคฺค อนตฺถาย อหิตาย ทุกฺขาย เทวมนุสฺสานํ. พหุญฺจ เต ภิกฺขเว ภิกฺขู อปุญฺญํ ปสวนฺติ, เต จิมํ สทฺธมฺมํ อนฺตรธาเปนฺตีติ.  จตุตฺติํสติมํ.}
\proseitemwithclosermidruled{132-139}{\csromanfolio{19}เย เต ภิกฺขเว ภิกฺขูอวินยํ “วินโย”ติ ทีเปนฺติ ฯเปฯ วินยํ “อวินโย”ติ ทีเปนฺติ ฯเปฯ อภาสิตํ อลปิตํ ตถาคเตน “ภาสิตํ ลปิตํ ตถาคเตนา”ติ ทีเปนฺติ ฯเปฯ ภาสิตํ ลปิตํ ตถาคเตน “อภาสิตํ อลปิตํ ตถาคเตนา”ติ ทีเปนฺติ ฯเปฯ อนาจิณฺณํ ตถาคเตน “อาจิณฺณํ ตถาคเตนา”ติ ทีเปนฺติ ฯเปฯ อาจิณฺณํ ตถาคเตน “อนาจิณฺณํ ตถาคเตนา”ติ ทีเปนฺติ ฯเปฯ อปญฺญตฺตํ ตถาคเตน “ปญฺญตฺตํ ตถาคเตนา”ติ ทีเปนฺติ ฯเปฯ ปญฺญตฺตํ ตถาคเตน “อปญฺญตฺตํ ตถาคเตนา”ติ ทีเปนฺติ. เต ภิกฺขเว ภิกฺขู พหุชน\-อหิตาย ปฏิปนฺนา พหุชน\-อสุขาย พหุโน ชนสฺส อนตฺถาย อหิตาย ทุกฺขาย เทวมนุสฺสานํ. พหุญฺจ เต ภิกฺขเว ภิกฺขู อปุญฺญํ ปสวนฺติ, เต จิมํ สทฺธมฺมํ อนฺตรธาเปนฺตีติ.  ทฺวาจตฺตาลีสติมํ.}{ทุติย\-ปมาทา\-ทิ\-วคฺโค ทสโม.}
\chapterhead{11. อธมฺมวคฺค}
\proseitem{140}{เย เต ภิกฺขเว ภิกฺขู อธมฺมํ “อธมฺโม”ติ ทีเปนฺติ, เต ภิกฺขเว ภิกฺขู พหุชนหิตาย ปฏิปนฺนา พหุชน\-สุขาย พหุโน ชนสฺส อตฺถาย หิตาย สุขาย เทวมนุสฺสานํ. พหุญฺจ เต ภิกฺขเว ภิกฺขู ปุญฺญํ ปสวนฺติ, เต จิมํ สทฺธมฺมํ ฐเปนฺตีติ\csromansharedfootnote{907b0628988a8717}{ถเปนฺตีติ (ก)}.  ปฐมํ.}

\proseitem{141}{เย เต ภิกฺขเว ภิกฺขู ธมฺมํ “ธมฺโม”ติ ทีเปนฺติ, เต ภิกฺขเว ภิกฺขู พหุชนหิตาย ปฏิปนฺนา พหุชน\-สุขาย พหุโน ชนสฺส อตฺถาย หิตาย สุขาย เทวมนุสฺสานํ. พหุญฺจ เต ภิกฺขเว ภิกฺขู ปุญฺญํ ปสวนฺติ, เต จิมํ สทฺธมฺมํ ฐเปนฺตีติ.  ทุติยํ.}

\proseitemwithclosermidruled{142-149}{เย เต ภิกฺขเว ภิกฺขู อวินยํ “อวินโย”ติ ทีเปนฺติ ฯเปฯ วินยํ “วินโย”ติ ทีเปนฺติ ฯเปฯ อภาสิตํ อลปิตํ ตถาคเตน “อภาสิตํ อลปิตํ ตถาคเตนา”ติ ทีเปนฺติ ฯเปฯ \csromanfolio{20}ภาสิตํ ลปิตํ ตถาคเตน “ภาสิตํ ลปิตํ ตถาคเตนา”ติ ทีเปนฺติ ฯเปฯ อนาจิณฺณํ ตถาคเตน “อนาจิณฺณํ ตถาคเตนา”ติ ทีเปนฺติ ฯเปฯ อาจิณฺณํ ตถาคเตน “อาจิณฺณํ ตถาคเตนา”ติ ทีเปนฺติ ฯเปฯ อปญฺญตฺตํ ตถาคเตน “อปญฺญตฺตํ ตถาคเตนา”ติ ทีเปนฺติ ฯเปฯ ปญฺญตฺตํ ตถาคเตน “ปญฺญตฺตํ ตถาคเตนา”ติ ทีเปนฺติ. เต ภิกฺขเว ภิกฺขู พหุชนหิตาย ปฏิปนฺนา พหุชน\-สุขาย พหุโน ชนสฺส อตฺถาย หิตาย สุขาย เทวมนุสฺสานํ. พหุญฺจ เต ภิกฺขเว ภิกฺขู ปุญฺญํ ปสวนฺติ, เต จิมํ สทฺธมฺมํ ฐเปนฺตีติ.  ทสมํ.}{อธมฺมวคฺโค เอกาทสโม.}
\chapterhead{12. \textbf{อนาปตฺติวคฺค}}
\proseitem{150}{เย เต ภิกฺขเว ภิกฺขู อนาปตฺติํ “อาปตฺตี”ติ ทีเปนฺติ, เต ภิกฺขเว ภิกฺขู พหุชน\-อหิตาย ปฏิปนฺนา พหุชน\-อสุขาย พหุโน ชนสฺส อนตฺถาย อหิตาย ทุกฺขาย เทวมนุสฺสานํ. พหุญฺจ เต ภิกฺขเว ภิกฺขู อปุญฺญํ ปสวนฺติ, เต จิมํ สทฺธมฺมํ อนฺตรธาเปนฺตีติ.  ปฐมํ.}

\proseitem{151}{เย เต ภิกฺขเว ภิกฺขู อาปตฺติํ “อนาปตฺตี”ติ ทีเปนฺติ, เต ภิกฺขเว ภิกฺขู พหุชน\-อหิตาย ปฏิปนฺนา พหุชน\-อสุขาย พหุโน ชนสฺส อนตฺถาย อหิตาย ทุกฺขาย เทวมนุสฺสานํ. พหุญฺจ เต ภิกฺขเว ภิกฺขู อปุญฺญํ ปสวนฺติ, เต จิมํ สทฺธมฺมํ อนฺตรธาเปนฺตีติ.  ทุติยํ.}

\proseitem{152-159}{เย เต ภิกฺขเว ภิกฺขู ลหุกํ อาปตฺติํ “ครุกา อาปตฺตี”ติ ทีเปนฺติ ฯเปฯ ครุกํ อาปตฺติํ “ลหุกา อาปตฺตี”ติ ทีเปนฺติ ฯเปฯ ทุฏฺฐุลฺลํ อาปตฺติํ “อทุฏฺฐุลฺลา อาปตฺตี”ติ ทีเปนฺติ ฯเปฯ อทุฏฺฐุลฺลํ อาปตฺติํ “ทุฏฺฐุลฺลา อาปตฺตี”ติ เทเปนฺติ ฯเปฯ สาวเสสํ อาปตฺติํ “อนวเสสา อาปตฺตี”ติ ทีเปนฺติ ฯเปฯ อนวเสสํ อาปตฺติํ “สาวเสสา อาปตฺตี”ติ ทีเปนฺติ ฯเปฯ สปฺปฏิกมฺมํ อาปตฺติํ “อปฺปฏิกมฺมา อาปตฺตี”ติ ทีเปนฺติ ฯเปฯ อปฺปฏิกมฺมํ อาปตฺติํ “สปฺปฏิกมฺมา อาปตฺตี”ติ ทีเปนฺติ. เต ภิกฺขเว ภิกฺขู พหุชน\-อหิตาย ปฏิปนฺนา พหุชน\-อสุขาย พหุโน ชนสฺส อนตฺถาย อหิตาย ทุกฺขาย เทวมนุสฺสานํ.}

\vspace{\tipitakaparskip}
\csromancenter{เอกปุคฺคล\-วคฺค พหุญฺจ เต ภิกฺขเว ภิกฺขู อปุญฺญํ ปสวนฺติ, เต จิมํ สทฺธมฺมํ อนฺตรธาเปนฺตีติ.  ทสมํ.}
\proseitem{160}{\csromanfolio{21}เย เต ภิกฺขเว ภิกฺขู อนาปตฺติํ “อนาปตฺตี”ติ ทีเปนฺติ, เต ภิกฺขเว ภิกฺขู พหุชนหิตาย ปฏิปนฺนา พหุชน\-สุขาย พหุโน ชนสฺส อตฺถาย หิตาย สุขาย เทวมนุสฺสานํ. พหุญฺจ เต ภิกฺขเว ภิกฺขู ปุญฺญํ ปสวนฺติ, เต จิมํ สทฺธมฺมํ ฐเปนฺตีติ.  เอกาทสมํ.}

\proseitem{161}{เย เต ภิกฺขเว ภิกฺขู อาปตฺติํ “อาปตฺตี”ติ ทีเปนฺติ, เต ภิกฺขเว ภิกฺขู พหุชนหิตาย ปฏิปนฺนา พหุชน\-สุขาย พหุโน ชนสฺส อตฺถาย หิตาย สุขาย เทวมนุสฺสานํ. พหุญฺจ เต ภิกฺขเว ภิกฺขู ปุญฺญํ ปสวนฺติ, เต จิมํ สทฺธมฺมํ ฐเปนฺตีติ.  ทฺวาทสมํ.}

\proseitemwithclosermidruled{162-169}{เย เต ภิกฺขเว ภิกฺขู ลหุกํ อาปตฺติํ “ลหุกา อาปตฺตี”ติ ทีเปนฺติ. ครุกํ อาปตฺติํ “ครุกา อาปตฺตี”ติ ทีเปนฺติ. ทุฏฺฐุลฺลํ อาปตฺติํ “ทุฏฺฐุลฺลา อาปตฺตี”ติ ทีเปนฺติ. อทุฏฺฐุลฺลํ อาปตฺติํ “อทุฏฺฐุลฺลา อาปตฺตี”ติ ทีเปนฺติ. สาวเสสํ อาปตฺติํ “สาวเสสา อาปตฺตี”ติ ทีเปนฺติ. อนวเสสํ อาปตฺติํ “อนวเสสา อาปตฺตี”ติ ทีเปนฺติ. สปฺปฏิกมฺมํ อาปตฺติํ “สปฺปฏิกมฺมา อาปตฺตี”ติ ทีเปนฺติ. อปฺปฏิกมฺมํ อาปตฺติํ “อปฺปฏิกมฺมา อาปตฺตี”ติ ทีเปนฺติ. เต ภิกฺขเว ภิกฺขู พหุชนหิตาย ปฏิปนฺนา พหุชน\-สุขาย พหุโน ชนสฺส อตฺถาย หิตาย สุขาย เทวมนุสฺสานํ. พหุญฺจ เต ภิกฺขเว ภิกฺขู ปุญฺญํ ปสวนฺติ, เต จิมํ สทฺธมฺมํ ฐเปนฺตีติ.  วีสติมํ.}{อนาปตฺติวคฺโค ทฺวาทสโม.}
\chapterhead{13. เอกปุคฺคล\-วคฺค}
\proseitem{170}{เอกปุคฺคโล ภิกฺขเว โลเก อุปฺปชฺชมาโน อุปฺปชฺชติ พหุชนหิตาย พหุชน\-สุขาย โลกานุกมฺปาย อตฺถายา หิตาย สุขาย เทวมนุสฺสานํ. กตโม เอกปุคฺคโล? ตถาคโต อรหํ สมฺมาสมฺพุทฺโธ, อยํ โข ภิกฺขเว เอกปุคฺคโล โลเก อุปฺปชฺชมาโน อุปฺปชฺชติ \csromanfolio{22}พหุชนหิตาย พหุชน\-สุขาย โลกานุกมฺปาย อตฺถาย หิตาย สุขาย เทว\-มนุสฺสา\-นนฺติ. (1)}

\proseitem{171}{เอกปุคฺคลสฺส ภิกฺขเว ปาตุภาโว ทุลฺลโภ โลกสฺมิํ. กตมสฺส เอกปุคฺคลสฺส? ตถาคตสฺส อรหโต สมฺมา\-สมฺพุทฺธสฺส, อิมสฺส โข ภิกฺขเว เอกปุคฺคลสฺส ปาตุภาโว ทุลฺลโภ โลกสฺมินฺติ. (2)}

\proseitem{172}{เอกปุคฺคโล ภิกฺขเว โลเก อุปฺปชฺชมาโน อุปฺปชฺชติ อจฺฉริย\-มนุสฺโส. กตโม เอกปุคฺคโล? ตถาคโต อรหํ สมฺมาสมฺพุทฺโธ, อยํ โข ภิกฺขเว เอกปุคฺคโล โลเก อุปฺปชฺชมาโน อุปฺปชฺชติ อจฺฉริย\-มนุสฺโสติ. (3)}

\proseitem{173}{เอกปุคฺคลสฺส ภิกฺขเว กาลกิริยา พหุโน ชนสฺส อนุตปฺปา\csromansharedfootnote{b354b494779f9964}{อานุตปฺปา (สี)} โหติ. กตมสฺส เอกปุคฺคลสฺส? ตถาคตสฺส อรหโต สมฺมา\-สมฺพุทฺธสฺส, อิมสฺส โข ภิกฺขเว เอกปุคฺคลสฺส กาลกิริยา พหุโน ชนสฺส อนุตปฺปา โหตีติ. (4)}

\proseitem{174}{เอกปุคฺคโล ภิกฺขเว โลเก อุปฺปชฺชมาโน อุปฺปชฺชติ อทุติโย อสหาโย อปฺปฏิโม อปฺปฏิสโม อปฺปฏิภาโค อปฺปฏิปุคฺคโล อสโม อสมสโม ทฺวิปทานํ อคฺโค. กตโม เอกปุคฺคโล? ตถาคโต อรหํ สมฺมาสมฺพุทฺโธ, อยํ โข ภิกฺขเว เอกปุคฺคโล โลเก อุปฺปชฺชมาโน อุปฺปชฺชติ อทุติโย อสหาโย อปฺปฏิโม อปฺปฏิสโม อปฺปฏิภาโค อปฺปฏิปุคฺคโล อสโม อสมสโม ทฺวิปทานํ อคฺโคติ. (5)}

\proseitem{175-186}{เอกปุคฺคลสฺส ภิกฺขเว ปาตุภาวา มหโต จกฺขุสฺส ปาตุภาโว โหติ, มหโต อาโลกสฺส ปาตุภาโว โหติ, มหโต โอภาสสฺส ปาตุภาโว โหติ, ฉนฺนํ อนุตฺตริยานํ ปาตุภาโว โหติ, จตุนฺนํ ปฏิสมฺภิทานํ สจฺฉิกิริยา โหติ, อเนกธาตุ\-ปฏิเวโธ โหติ, นานาธาตุ\-ปฏิเวโธ โหติ, วิชฺชา\-วิมุตฺติ\-ผล\-สจฺฉิกิริยา โหติ, โสตาปตฺติผล\-สจฺฉิกิริยา โหติ, สกทาคามิ\-ผล\-สจฺฉิกิริยา โหติ, อนาคามิ\-ผล\-สจฺฉิกิริยา โหติ, \csromanfolio{23}อรหตฺตผล\-สจฺฉิกิริยา โหติ. กตมสฺส เอกปุคฺคลสฺส? ตถาคตสฺส อรหโต สมฺมา\-สมฺพุทฺธสฺส, อิมสฺส โข ภิกฺขเว เอกปุคฺคลสฺส ปาตุภาวา มหโต จกฺขุสฺส ปาตุภาโว โหติ, มหโต อาโลกสฺส ปาตุภาโว โหติ, มหโต โอภาสสฺส ปาตุภาโว โหติ, ฉนฺนํ อนุตฺตริยานํ ปาตุภาโว โหติ, จตุนฺนํ ปฏิสมฺภิทานํ สจฺฉิกิริยา โหติ, อเนกธาตุ\-ปฏิเวโธ โหติ, นานาธาตุ\-ปฏิเวโธ โหติ, วิชฺชา\-วิมุตฺติ\-ผล\-สจฺฉิกิริยา โหติ, โสตาปตฺติผล\-สจฺฉิกิริยา โหติ, สกทาคามิ\-ผล\-สจฺฉิกิริยา โหติ, อนาคามิ\-ผล\-สจฺฉิกิริยา โหติ, อรหตฺตผล\-สจฺฉิกิริยา โหตีติ. \mbox{(6-17)}}

\proseitemwithclosermidruled{187}{นาหํ ภิกฺขเว อญฺญํ เอกปุคฺคลมฺปิ สมนุปสฺสามิ, โย เอวํ ตถาคเตน อนุตฺตรํ ธมฺมจกฺกํ ปวตฺติตํ สมฺมเทว อนุปฺปวตฺเตติ, ยถยิทํ ภิกฺขเว สาริปุตฺโต. สาริปุตฺโต ภิกฺขเว ตถาคเตน อนุตฺตรํ ธมฺมจกฺกํ ปวตฺติตํ สมฺมเทว อนุปฺป\-วตฺเตตีติ. (18)}{เอกปุคฺคล\-วคฺโค เตรสโม.}
\chapterhead{14. เอตทคฺควคฺค 1. ปฐมวคฺค}
\proseitem{188}{เอตทคฺคํ ภิกฺขเว มม สาวกานํ ภิกฺขูนํ รตฺตญฺญูนํ ยทิทํ อญฺญาสิโกณฺฑญฺโญ\csromansharedfootnote{45cf490c7b6fbad8}{อญฺญาตโกณฺฑญฺโญติ (ก), อญฺญาโกณฺฑญฺโญ (สี, สฺยา, กํ, อิ)}. (1)}

\prose{189. มหาปญฺญานํ ยทิทํ สาริปุตฺโต. (2)}

\prose{190. อิทฺธิมนฺตานํ ยทิทํ มหาโมคฺคลฺลาโน. (3)}

\prose{191. ธุตวาทานํ\csromansharedfootnote{f687914fa8d084c8}{ธุตงฺคธรานํ (กตฺถจิ)} ยทิทํ มหากสฺสโป. (4)}

\prose{192. ทิพฺพ\-จกฺขุ\-กานํ ยทิทํ อนุรุทฺโธ. (5)}

\prose{193. อุจฺจากุลิกานํ ยทิทํ ภทฺทิโย กาฬิโคธาย\-ปุตฺโต. (6)}

\prose{194. มญฺชุสฺสรานํ ยทิทํ ลกุณฺฑก\csromansharedfootnote{61a80e24ab584fd5}{ลกุณฺฏก (สฺยา, กํ)} ภทฺทิโย. (7)}

\prose{\csromanfolio{24}195. สีหนาทิกานํ ยทิทํ ปิณฺโฑล\-ภารทฺวาโช. (8)}

\prose{196. ธมฺม\-กถิกานํ ยทิทํ ปุณฺโณ มนฺตาณิปุตฺโต. (9)}

\proseitemwithclosermidruled{197}{สํขิตฺเตน ภาสิตสฺส วิตฺถาเรน อตฺถํ วิภชนฺตานํ ยทิทํ มหากจฺจาโนติ. (10)}{วคฺโค ปฐโม.}
\chapterhead{14. \textbf{เอตทคฺควคฺค 2. ทุติยวคฺค}}
\proseitem{198}{เอตทคฺคํ ภิกฺขเว มม สาวกานํ ภิกฺขูนํ มโนมยํ กายํ อภินิ\-มฺมิน\-นฺตานํ ยทิทํ จูฬปนฺถโก\csromansharedfootnote{646c83b901daf1d2}{จุลฺลปนฺถโก (สี, สฺยา, กํ, อิ)}. (1)}

\prose{199. เจโต\-วิวฏฺฏ\-กุสลานํ ยทิทํ จูฬปนฺถโก. (2)}

\prose{200. สญฺญา\-วิวฏฺฏ\-กุสลานํ ยทิทํ มหาปนฺถโก. (3)}

\prose{201. อรณวิหารีนํ ยทิทํ สุภูติ. (4)}

\prose{202. ทกฺขิเณยฺยานํ ยทิทํ สุภูติ. (5)}

\prose{203. อารญฺญกานํ ยทิทํ เรวโต ขทิรวนิโย. (6)}

\prose{204. ฌายีนํ ยทิทํ กงฺขาเรวโต. (7)}

\prose{205. อารทฺธ\-วีริยานํ ยทิทํ โสโณ โกฬิวิโส. (8)}

\prose{206. กลฺยาณ\-วากฺกรณานํ ยทิทํ โสโณ กุฏิกณฺโณ. (9)}

\prose{207. ลาภีนํ ยทิทํ สีวลิ. (10)}

\prosewithclosermidruled{208. สทฺธา\-ธิมุตฺตานํ ยทิทํ วกฺกลีติ. (11)}{วคฺโค ทุติโย.}
\chapterhead{14. \textbf{เอตทคฺควคฺค 3. ตติยวคฺค}}
\proseitem{209}{เอตทคฺคํ ภิกฺขเว มม สาวกานํ ภิกฺขูนํ สิกฺขากามานํ ยทิทํ ราหุโล. (1)}

\chapterheadpageread{14. เอตทคฺควคฺค}
\prose{\csromanfolio{25}210. สทฺธา\-ปพฺพชิตานํ ยทิทํ รฏฺฐปาโล. (2)}

\prose{211. ปฐมํ สลากํ คณฺหนฺตานํ ยทิทํ กุณฺฑธาโน. (3)}

\prose{212. ปฏิภานวนฺตานํ ยทิทํ วงฺคีโส. (4)}

\prose{213. สมนฺต\-ปาสาทิกานํ ยทิทํ อุปเสโน วงฺคนฺตปุตฺโต. (5)}

\prose{214. เสนาสน\-ปญฺญาปกานํ ยทิทํ ทพฺโพ มลฺลปุตฺโต. (6)}

\prose{215. เทวตานํ ปิยมนาปานํ ยทิทํ ปิลินฺทวจฺโฉ. (7)}

\prose{216. ขิปฺปา\-ภิญฺญานํ ยทิทํ พาหิโย ทารุจีริโย. (8)}

\prose{217. จิตฺต\-กถิกานํ ยทิทํ กุมารกสฺสโป. (9)}

\prosewithclosermidruled{218. ปฏิสมฺภิทา\-ปตฺตานํ ยทิทํ มหาโกฏฺฐิโตติ\csromansharedfootnote{e69b0118df4be550}{มหาโกฏฺฐิโกติ (อญฺเญสุ สุตฺเตสุ มรมฺมโปตฺถเก)}. (10)}{วคฺโค ตติโย.}
\chapterhead{14. เอตทคฺควคฺค 4. จตุตฺถวคฺค}
\proseitem{219}{เอตทคฺคํ ภิกฺขเว มม สาวกานํ ภิกฺขูนํ พหุสฺสุตานํ ยทิทํ อานนฺโท. (1)}

\prose{220. สติมนฺตานํ ยทิทํ อานนฺโท. (2)}

\prose{221. คติมนฺตานํ ยทิทํ อานนฺโท. (3)}

\prose{222. ธิติมนฺตานํ ยทิทํ อานนฺโท. (4)}

\prose{223. อุปฏฺฐากานํ ยทิทํ อานนฺโท. (5)}

\prose{224. มหาปริสานํ ยทิทํ อุรุเวลกสฺสโป. (6)}

\prose{225. กุลปฺปสาทกานํ ยทิทํ กาฬุทายี. (7)}

\prose{226. อปฺปาพาธานํ ยทิทํ พากุโล\csromansharedfootnote{9d0ac76b63ea2b71}{พกฺกุโล (สี, สฺยา, กํ, อิ)}. (8)}

\prose{227. ปุพฺเพนิวาสํ อนุสฺสรนฺตานํ ยทิทํ โสภิโต. (9)}

\prose{228. วินยธรานํ ยทิทํ อุปาลิ. (10)}

\prose{\csromanfolio{26}229. ภิกฺขุโน\-วาทกานํ ยทิทํ นนฺทโก. (11)}

\prose{230. อินฺทฺริเยสุ คุตฺตทฺวารานํ ยทิทํ นนฺโท. (12)}

\prose{231. ภิกฺขุ\-โอวาทกานํ ยทิทํ มหากปฺปิโน. (13)}

\prose{232. เตโชธาตุ\-กุสลานํ ยทิทํ สาคโต. (14)}

\prose{233. ปฏิภาเนยฺยกานํ ยทิทํ ราโธ. (15)}

\prosewithclosermidruled{234. ลูขจีวร\-ธรานํ ยทิทํ โมฆราชาติ. (16)}{วคฺโค จตุตฺโถ.}
\chapterhead{14. \textbf{เอตทคฺควคฺค 5. ปญฺจมวคฺค}}
\proseitem{235}{เอตทคฺคํ ภิกฺขเว มม สาวิกานํ ภิกฺขุนีนํ รตฺตญฺญูนํ ยทิทํ มหาปชาปติ\-โคตมี. (1)}

\prose{236. มหาปญฺญานํ ยทิทํ เขมา. (2)}

\prose{237. อิทฺธิมนฺตีนํ ยทิทํ อุปฺปลวณฺณา. (3)}

\prose{238. วินยธรานํ ยทิทํ ปฏาจารา. (4)}

\prose{239. ธมฺม\-กถิกานํ ยทิทํ ธมฺมทินฺนา. (5)}

\prose{240. ฌายีนํ ยทิทํ นนฺทา. (6)}

\prose{241. อารทฺธ\-วีริยานํ ยทิทํ โสณา. (7)}

\prose{242. ทิพฺพ\-จกฺขุ\-กานํ ยทิทํ พกุลา\csromansharedfootnote{152756ce33699a0e}{สกุลา (สี, สฺยา, กํ, อิ)}. (8)}

\prose{243. ขิปฺปา\-ภิญฺญานํ ยทิทํ ภทฺทา กุณฺฑลเกสา. (9)}

\prose{244. ปุพฺเพนิวาสํ อนุสฺสรนฺตีนํ ยทิทํ ภทฺทา กาปิลานี. (10)}

\prose{245. มหาภิญฺญา\-ปตฺตานํ ยทิทํ ภทฺทา กจฺจานา. (11)}

\prose{246. ลูขจีวร\-ธรานํ ยทิทํ กิสาโคตมี. (12)}

\prosewithclosermidruled{247. สทฺธา\-ธิมุตฺตานํ ยทิทํ สิงฺคาลกมาตาติ\csromansharedfootnote{5adbed2e6a6d8fdd}{สิคาลมาตาติ (สี, สฺยา, กํ, อิ)}. (13)}{วคฺโค ปญฺจโม.}
\chapterheadpageread{14. เอตทคฺควคฺค \textbf{14. เอตทคฺควคฺค 6. ฉฏฺฐวคฺค}}
\proseitem{248}{\csromanfolio{27}เอตทคฺคํ ภิกฺขเว มม สาวกานํ อุปาสกานํ ปฐมํ สรณํ คจฺฉนฺตานํ ยทิทํ ตปุสฺส\-ภลฺลิกา\csromansharedfootnote{5efd824c33c81701}{ตปสฺสุภลฺลิกา (สี, อิ)} วาณิชา. (1)}

\prose{249. ทายกานํ ยทิทํ สุทตฺโต คหปติ อนาถปิณฺฑิโก. (2)}

\prose{250. ธมฺม\-กถิกานํ ยทิทํ จิตฺโต คหปติ มจฺฉิกา\-สณฺฑิโก. (3)}

\proseitem{251}{จตูหิ สงฺคห\-วตฺถูหิ ปริสํ สงฺคณฺหนฺตานํ ยทิทํ หตฺถโก อาฬวโก. (4)}

\prose{252. ปณีต\-ทายกานํ ยทิทํ มหานาโม สกฺโก. (5)}

\prose{253. มนาป\-ทายกานํ ยทิทํ อุคฺโค คหปติ เวสาลิโก. (6)}

\prose{254. สํฆุปฏฺฐากานํ ยทิทํ หตฺถิคามโก อุคฺคโต คหปติ. (7)}

\prose{255. อเวจฺจ\-ปฺปสนฺนานํ ยทิทํ สูรมฺพฏฺโฐ\csromansharedfootnote{06701f72770ead53}{สูโร อมฺพฏฺโฐ (สี, สฺยา, กํ, อิ) สุรพนฺโธ (ก)}. (8)}

\prose{256. ปุคฺคล\-ปฺปสนฺนานํ ยทิทํ ชีวโก โกมารภจฺโจ. (9)}

\prosewithclosermidruled{257. วิสฺสาสกานํ ยทิทํ นกุลปิตา คหปตีติ. (10)}{วคฺโค ฉฏฺโฐ.}
\chapterhead{14. เอตทคฺควคฺค 7. สตฺตมวคฺค}
\proseitem{258}{เอตทคฺคํ ภิกฺขเว มม สาวิกานํ อุปาสิกานํ ปฐมํ สรณํ คจฺฉนฺตีนํ ยทิทํ สุชาตา เสนิยธีตา\csromansharedfootnote{6148f09bfa4d0fd7}{เสนานี ธีตา (สี, สฺยา, กํ, อิ)}. (1)}

\prose{259. ทายิกานํ ยทิทํ วิสาขา มิคารมาตา. (2)}

\prose{260. พหุสฺสุตานํ ยทิทํ ขุชฺชุตฺตรา. (3)}

\prose{261. เมตฺตาวิหารีนํ ยทิทํ สามาวตี. (4)}

\prose{262. ฌายีนํ ยทิทํ อุตฺตรา นนฺทมาตา. (5)}

\prose{263. ปณีตทายิกานํ ยทิทํ สุปฺปวาสา โกลิยธีตา. (6)}

\prose{\csromanfolio{28}264. คิลานุ\-ปฏฺฐากีนํ ยทิทํ สุปฺปิยา อุปาสิกา. (7)}

\prose{265. อเวจฺจ\-ปฺปสนฺนานํ ยทิทํ กาติยานี. (8)}

\prose{266. วิสฺสาสิกานํ ยทิทํ นกุลมาตา คหปตานี. (9)}

\proseitemwithclosermid{267}{อนุสฺสว\-ปฺปสนฺนานํ ยทิทํ กาฬี อุปาสิกา กุรรฆริกา\csromansharedfootnote{7404a1498c2e0002}{กุลฆริกา (ก)}ติ. 10)}{วคฺโค สตฺตโม. (เอตทคฺควคฺโค นิฏฺฐิโต.)}
\csromanpalirecto
\gambhira{15. \textbf{อฏฺฐานปาฬิ 1. ปฐมวคฺค}}
\proseitem{268}{อฏฺฐานเมตํ ภิกฺขเว อนวกาโส, ยํ ทิฏฺฐิสมฺปนฺโน ปุคฺคโล กญฺจิ\csromansharedfootnote{a279d1d1d6cac341}{กิญฺจิ (ก)} สงฺขารํ นิจฺจโต อุปคจฺเฉยฺย, เนตํ ฐานํ วิชฺชติ. ฐานญฺจ โข เอตํ ภิกฺขเว วิชฺชติ, ยํ ปุถุชฺชโน กญฺจิ สงฺขารํ นิจฺจโต อุปคจฺเฉยฺย, ฐานเมตํ วิชฺชตีติ. (1)}

\proseitem{269}{อฏฺฐานเมตํ ภิกฺขเว อนวกาโส, ยํ ทิฏฺฐิสมฺปนฺโน ปุคฺคโล กญฺจิ สงฺขารํ สุขโต อุปคจฺเฉยฺย, เนตํ ฐานํ วิชฺชติ. ฐานญฺจ โข เอตํ ภิกฺขเว วิชฺชติ, ยํ ปุถุชฺชโน กญฺจิ สงฺขารํ สุขโต อุปคจฺเฉยฺย, ฐานเมตํ วิชฺชตีติ. (2)}

\proseitem{270}{อฏฺฐานเมตํ ภิกฺขเว อนวกาโส, ยํ ทิฏฺฐิสมฺปนฺโน ปุคฺคโล กญฺจิ ธมฺมํ อตฺตโต อุปคจฺเฉยฺย, เนตํ ฐานํ วิชฺชติ. ฐานญฺจ โข เอตํ ภิกฺขเว วิชฺชติ, ยํ ปุถุชฺชโน กญฺจิ ธมฺมํ อตฺตโต อุปคจฺเฉยฺย, ฐานเมตํ วิชฺชตีติ. (3)}

\proseitem{271}{อฏฺฐานเมตํ ภิกฺขเว อนวกาโส, ยํ ทิฏฺฐิสมฺปนฺโน ปุคฺคโล มาตรํ ชีวิตา โวโรเปยฺย, เนตํ ฐานํ วิชฺชติ. ฐานญฺจ โข ภิกฺขเว วิชฺชติ, ยํ ปุถุชฺชโน มาตรํ ชีวิตา โวโรเปยฺย, ฐานเมตํ วิชฺชตีติ. (4)}

\csromanpalirecto
\gambhira{15. อฏฺฐานปาฬิ}
\proseitem{272}{\csromanfolio{29}อฏฺฐานเมตํ ภิกฺขเว อนวกาโส, ยํ ทิฏฺฐิสมฺปนฺโน ปุคฺคโล ปิตรํ ชีวิตา โวโรเปยฺย, เนตํ ฐานํ วิชฺชติ. ฐานญฺจ โข เอตํ ภิกฺขเว วิชฺชติ, ยํ ปุถุชฺชโน ปิตรํ ชีวิตา โวโรเปยฺย, ฐานเมตํ วิชฺชตีติ. (5)}

\proseitem{273}{อฏฺฐานเมตํ ภิกฺขเว อนวกาโส, ยํ ทิฏฺฐิสมฺปนฺโน ปุคฺคโล อรหนฺตํ ชีวิตา โวโรเปยฺย, เนตํ ฐานํ วิชฺชติ. ฐานญฺจ โข เอตํ ภิกฺขเว วิชฺชติ, ยํ ปุถุชฺชโน อรหนฺตํ ชีวิตา โวโรเปยฺย, ฐานเมตํ วิชฺชตีติ.  (6)}

\proseitem{274}{อฏฺฐานเมตํ ภิกฺขเว อนวกาโส, ยํ ทิฏฺฐิสมฺปนฺโน ปุคฺคโล ตถาคตสฺส ปทุฏฺฐจิตฺโต โลหิตํ อุปฺปาเทยฺย, เนตํ ฐานํ วิชฺชติ. ฐานญฺจ โข เอตํ ภิกฺขเว วิชฺชติ, ยํ ปุถุชฺชโน ตถาคตสฺส ปทุฏฺฐจิตฺโต โลหิตํ อุปฺปาเทยฺย, ฐานเมตํ วิชฺชตีติ. (7)}

\proseitem{275}{อฏฺฐานเมตํ ภิกฺขเว อนวกาโส, ยํ ทิฏฺฐิสมฺปนฺโน ปุคฺคโล สํฆํ ภินฺเทยฺย, เนตํ ฐานํ วิชฺชติ. ฐานญฺจ โข เอตํ ภิกฺขเว วิชฺชติ, ยํ ปุถุชฺชโน สํฆํ ภินฺเทยฺย, ฐานเมตํ วิชฺชตีติ. (8)}

\proseitem{276}{อฏฺฐานเมตํ ภิกฺขเว อนวกาโส, ยํ ทิฏฺฐิสมฺปนฺโน ปุคฺคโล อญฺญํ สตฺถารํ อุทฺทิเสยฺย, เนตํ ฐานํ วิชฺชติ. ฐานญฺจ โข เอตํ ภิกฺขเว วิชฺชติ, ยํ ปุถุชฺชโน อญฺญํ สตฺถารํ อุทฺทิเสยฺย, ฐานเมตํ วิชฺชตีติ. (9)}

\proseitemwithclosermidruled{277}{อฏฺฐานเมตํ ภิกฺขเว อนวกาโส, ยํ เอกิสฺสา โลกธาตุยา เทฺว อรหนฺโต สมฺมาสมฺพุทฺธา อปุพฺพํ อจริมํ อุปฺปชฺเชยฺยุํ, เนตํ ฐานํ วิชฺชติ. ฐานญฺจ โข เอตํ ภิกฺขเว วิชฺชติ, ยํ เอกิสฺสา โลกธาตุยา เอโกว อรหํ สมฺมาสมฺพุทฺโธ อุปฺปชฺเชยฺย, ฐานเมตํ วิชฺชตีติ.}{(10) วคฺโค ปฐโม.}
\csromanpalirecto
\gambhira{15. อฏฺฐานปาฬิ 2. ทุติยวคฺค}
\proseitem{278}{อฏฺฐานเมตํ ภิกฺขเว อนวกาโส, ยํ เอกิสฺสา โลกธาตุยา เทฺว ราชาโน จกฺกวตฺตี อปุพฺพํ อจริมํ อุปฺปชฺเชยฺยุํ, เนตํ ฐานํ \csromanfolio{30}วิชฺชติ. ฐานญฺจ โข เอตํ ภิกฺขเว วิชฺชติ, ยํ เอกิสฺสา โลกธาตุยา เอโก ราชา จกฺกวตฺตี อุปฺปชฺเชยฺย, ฐานเมตํ วิชฺชตีติ. (1)}

\proseitem{279}{อฏฺฐานเมตํ ภิกฺขเว อนวกาโส, ยํ อิตฺถี อรหํ อสฺส สมฺมาสมฺพุทฺโธ, เนตํ ฐานํ วิชฺชติ. ฐานญฺจ โข เอตํ ภิกฺขเว วิชฺชติ, ยํ ปุริโส อรหํ อสฺส สมฺมาสมฺพุทฺโธ, ฐานเมตํ วิชฺชตีติ. (2)}

\proseitem{280}{อฏฺฐานเมตํ ภิกฺขเว อนวกาโส, ยํ อิตฺถี ราชา อสฺส จกฺกวตฺตี, เนตํ ฐานํ วิชฺชติ. ฐานญฺจ โข เอตํ ภิกฺขเว วิชฺชติ, ยํ ปุริโส ราชา อสฺส จกฺกวตฺตี, ฐานเมตํ วิชฺชตีติ. (3)}

\proseitem{281-283}{อฏฺฐานเมตํ ภิกฺขเว อนวกาโส, ยํ อิตฺถี สกฺกตฺตํ กาเรยฺย ฯเปฯ มารตฺตํ กาเรยฺย ฯเปฯ พฺรหฺมตฺตํ กาเรยฺย, เนตํ ฐานํ วิชฺชติ. ฐานญฺจ โข เอตํ ภิกฺขเว วิชฺชติ, ยํ ปุริโส สกฺกตฺตํ กาเรยฺย ฯเปฯ มารตฺตํ กาเรยฺย ฯเปฯ พฺรหฺมตฺตํ กาเรยฺย, ฐนเมตํ วิชฺชตีติ. \mbox{(4-6)}}

\proseitem{284}{อฏฺฐานเมตํ ภิกฺขเว อนวกาโส, ยํ กาย\-ทุจฺจริตสฺส อิฏฺโฐ กนฺโต มนาโป วิปาโก นิพฺพตฺเตยฺย, เนตํ ฐานํ วิชฺชติ. ฐานญฺจ โข เอตํ ภิกฺขเว วิชฺชติ, ยํ กาย\-ทุจฺจริตสฺส อนิฏฺโฐ อกนฺโต อมนาโป วิปาโก นิพฺพตฺเตยฺย, ฐานเมตํ วิชฺชตีติ. (7)}

\proseitemwithclosermidruled{285-286}{อฏฺฐานเมตํ ภิกฺขเว อนวกาโส, ยํ วจีทุจฺจริตสฺส ฯเปฯ ยํ มโน\-ทุจฺจริตสฺส อิฏฺโฐ กนฺโต มนาโป วิปาโก นิพฺพตฺเตยฺย, เนตํ ฐานํ วิชฺชติ. ฐานญฺจ โข เอตํ ภิกฺขเว วิชฺชติ, ยํ มโน\-ทุจฺจริตสฺส อนิฏฺโฐ อกนฺโต อมนาโป วิปาโก นิพฺพตฺเตยฺย, ฐานเมตํ วิชฺชตีติ. \mbox{(8-9)}}{วคฺโค ทุติโย.}
\csromanpalirecto
\gambhira{15. \textbf{อฏฺฐานปาฬิ 3. ตติยวคฺค}}
\proseitem{287}{อฏฺฐานเมตํ ภิกฺขเว อนวกาโส, ยํ กาย\-สุจริตสฺส อนิฏฺโฐ อกนฺโต อมนาโป วิปาโก นิพฺพตฺเตยฺย, เนตํ ฐานํ วิชฺชติ.}

\prose{\csromanfolio{31}ฐานญฺจ โข เอตํ ภิกฺขเว วิชฺชติ, ยํ กาย\-สุจริตสฺส อิฏฺโฐ กนฺโต มนาโป วิปาโก นิพฺพตฺเตยฺย, ฐานเมตํ วิชฺชตีติ. (1)}

\proseitem{288-289}{อฏฺฐานเมตํ ภิกฺขเว อนวกาโส, ยํ วจีสุจริตสฺส ฯเปฯ มโน\-สุจริตสฺส อนิฏฺโฐ อกนฺโต อมนาโป วิปาโก นิพฺพตฺเตยฺย, เนตํ ฐานํ วิชฺชติ. ฐานญฺจ โข เอตํ ภิกฺขเว วิชฺชติ, ยํ มโนสุจิริตสฺส อิฏฺโฐ กนฺโต มนาโป วิปาโก นิพฺพตฺเตยฺย, ฐานเมตํ วิชฺชตีติ. \mbox{(2-3)}}

\proseitem{290}{อฏฺฐานเมตํ ภิกฺขเว อนวกาโส, ยํ กายทุจฺจริต\-สมงฺคี ตนฺนิทานา ตปฺปจฺจยา กายสฺส เภทา ปรํ มรณา สุคติํ สคฺคํ โลกํ อุปปชฺเชยฺย, เนตํ ฐานํ วิชฺชติ. ฐานญฺจ โข เอตํ ภิกฺขเว วิชฺชติ, ยํ กายทุจฺจริต\-สมงฺคี ตนฺนิทานา ตปฺปจฺจยา กายสฺส เภทา ปรํ มรณา อปายํ ทุคฺคติํ วินิปาตํ นิรยํ อุปปชฺเชยฺย, ฐานเมตํ วิชฺชตีติ. (4)}

\proseitem{291-292}{อฏฺฐานเมตํ ภิกฺขเว อนวกาโส, ยํ วจีทุจฺจริต\-สมงฺคี ฯเปฯ ยํ มโนทุจฺจริต\-สมงฺคี ตนฺนิทานา ตปฺปจฺจยา กายสฺส เภทา ปรํ มรณา สุคติํ สคฺคํ โลกํ อุปปชฺเชยฺย, เนตํ ฐานํ วิชฺชติ. ฐานญฺจ โข เอตํ ภิกฺขเว วิชฺชติ, ยํ มโนทุจฺจริต\-สมงฺคี ตนฺนิทานา ตปฺปจฺจยา กายสฺส เภทา ปรํ มรณา อปายํ ทุคฺคติํ วินิปาตํ นิรยํ อุปปชฺเชยฺย, ฐานเมตํ วิชฺชตีติ. \mbox{(5-6)}}

\proseitem{293}{อฏฺฐานเมตํ ภิกฺขเว อนวกาโส, ยํ กายสุจริต\-สมงฺคี ตนฺนิทานา ตปฺปจฺจยา กายสฺส เภทา ปรํ มรณา อปายํ ทุคฺคติํ วินิปาตํ นิรยํ อุปปชฺเชยฺย, เนตํ ฐานํ วิชฺชติ. ฐานญฺจ โข เอตํ ภิกฺขเว วิชฺชติ, ยํ กายสุจริต\-สมงฺคี ตนฺนิทานา ตปฺปจฺจยา กายสฺส เภทา ปรํ มรณา สุคติํ สคฺคํ โลกํ อุปปชฺเชยฺย, ฐานเมตํ วิชฺชตีติ. (7)}

\proseitemwithclosermajor{294-295}{อฏฺฐานเมตํ ภิกฺขเว อนวกาโส, ยํ วจีสุจริต\-สมงฺคี ฯเปฯ ยํ มโนสุจริต\-สมงฺคี ตนฺนิทานา ตปฺปจฺจยา กายสฺส เภทา ปรํ มรณา อปายํ ทุคฺคติํ วินิปาตํ นิรยํ อุปปชฺเชยฺย, เนตํ ฐานํ วิชฺชติ. ฐานญฺจ โข เอตํ ภิกฺขเว วิชฺชติ, ยํ มโนสุจริต\-สมงฺคี ตนฺนิทานา ตปฺปจฺจยา กายสฺส เภทา ปรํ มรณา สุคติํ สคฺคํ โลกํ อุปปชฺเชยฺย, ฐานเมตํ วิชฺชตีติ. \mbox{(8-9)}}{วคฺโค ตติโย. อฏฺฐานปาฬิ นิฏฺฐิตา.}
\csromanpalirecto
\gambhira{16. \textbf{เอกธมฺมปาฬิ 1. ปฐมวคฺค}}
\proseitem{296}{\csromanfolio{32}เอกธมฺโม ภิกฺขเว ภาวิโต พหุลีกโต เอกนฺต\-นิพฺพิทาย วิราคาย นิโรธาย อุปสมาย อภิญฺญาย สมฺโพธาย นิพฺพานาย สํวตฺตติ. กตโม เอกธมฺโม? พุทฺธานุสฺสติ. อยํ โข ภิกฺขเว เอกธมฺโม ภาวิโต พหุลีกโต เอกนฺต\-นิพฺพิทาย วิราคาย นิโรธาย อุปสมาย อภิญฺญาย สมฺโพธาย นิพฺพานาย สํวตฺตตีติ. (1)}

\proseitemwithclosermidruled{297}{เอกธมฺโม ภิกฺขเว ภาวิโต พหุลีกโต เอกนฺต\-นิพฺพิทาย วิราคาย นิโรธาย อุปสมาย อภิญฺญาย สมฺโพธาย นิพฺพานาย สํวตฺตติ. กตโม เอกธมฺโม? ธมฺมานุสฺสติ ฯเปฯ. สํฆานุสฺสติ. สีลานุสฺสติ. จาคานุสฺสติ. เทวตานุสฺสติ. อานาปานสฺสติ. มรณสฺสติ. กายคตาสติ. อุปสมานุสฺสติ. อยํ โข ภิกฺขเว เอกธมฺโม ภาวิโต พหุลีกโต เอกนฺต\-นิพฺพิทาย วิราคาย นิโรธาย อุปสมาย อภิญฺญาย สมฺโพธาย นิพฺพานาย สํวตฺตตีติ. (10)}{วคฺโค ปฐโม.}
\csromanpalirecto
\gambhira{16. \textbf{เอกธมฺมปาฬิ 2. ทุติยวคฺค}}
\proseitem{298}{นาหํ ภิกฺขเว อญฺญํ เอกธมฺมมฺปิ สมนุปสฺสามิ, เยน อนุปฺปนฺนา วา อกุสลา ธมฺมา อุปฺปชฺชนฺติ, อุปฺปนฺนา วา อกุสลา ธมฺมา ภิยฺโยภาวาย เวปุลฺลาย สํวตฺตนฺติ, ยถยิทํ ภิกฺขเว มิจฺฉาทิฏฺฐิ. มิจฺฉา\-ทิฏฺฐิกสฺส ภิกฺขเว อนุปฺปนฺนา เจว อกุสลา ธมฺมา อุปฺปชฺชนฺติ, อุปฺปนฺนา จ อกุสลา ธมฺมา ภิยฺโยภาวาย เวปุลฺลาย สํวตฺตนฺตีติ. (1)}

\proseitem{299}{นาหํ ภิกฺขเว อญฺญํ เอกธมฺมมฺปิ สมนุปสฺสามิ, เยน อนุปฺปนฺนา วา กุสลา ธมฺมา อุปฺปชฺชนฺติ, อุปฺปนฺนา วา กุสลา ธมฺมา ภิยฺโยภาวาย เวปุลฺลาย สํวตฺตนฺติ, ยถยิทํ ภิกฺขเว สมฺมาทิฏฺฐิ. สมฺมา\-ทิฏฺฐิกสฺส ภิกฺขเว อนุปฺปนฺนา เจว กุสลา ธมฺมา อุปฺปชฺชนฺติ, อุปฺปนฺนา จ กุสลา ธมฺมา ภิยฺโยภาวาย เวปุลฺลาย สํวตฺตนฺตีติ. (2)}

\proseitem{300}{นาหํ ภิกฺขเว อญฺญํ เอกธมฺมมฺปิ สมนุปสฺสามิ, เยน อนุปฺปนฺนา วา กุสลา ธมฺมา นุปฺปชฺชนฺติ, อุปฺปนฺนา วา กุสลา ธมฺมา ปริหายนฺติ, \csromanfolio{33}ยถยิทํ ภิกฺขเว มิจฺฉาทิฏฺฐิ. มิจฺฉา\-ทิฏฺฐิกสฺส ภิกฺขเว อนุปฺปนฺนา เจว กุสลา ธมฺมา นุปฺปชฺชนฺติ, อุปฺปนฺนา จ กุสลา ธมฺมา ปริหายนฺตีติ. (3)}

\proseitem{301}{นาหํ ภิกฺขเว อญฺญํ เอกธมฺมมฺปิ สมนุปสฺสามิ, เยน อนุปฺปนฺนา วา อกุสลา ธมฺมา นุปฺปชฺชนฺติ, อุปฺปนฺนา วา อกุสลา ธมฺมา ปริหายนฺติ, ยถยิทํ ภิกฺขเว สมฺมาทิฏฺฐิ. สมฺมา\-ทิฏฺฐิกสฺส ภิกฺขเว อนุปฺปนฺนา เจว อกุสลา ธมฺมา นุปฺปชฺชนฺติ, อุปฺปนฺนา จ อกุสลา ธมฺมา ปริหายนฺตีติ. (4)}

\proseitem{302}{นาหํ ภิกฺขเว อญฺญํ เอกธมฺมมฺปิ สมนุปสฺสามิ, เยน อนุปฺปนฺนา วา มิจฺฉาทิฏฺฐิ อุปฺปชฺชติ, อุปฺปนฺนา วา มิจฺฉาทิฏฺฐิ ปวฑฺฒติ, ยถยิทํ ภิกฺขเว อโยนิโส\-มนสิกาโร. อโยนิโส ภิกฺขเว มนสิ กโรโต อนุปฺปนฺนา เจว มิจฺฉาทิฏฺฐิ อุปฺปชฺชติ, อุปฺปนฺนา จ มิจฺฉาทิฏฺฐิ ปวฑฺฒตีติ. (5)}

\proseitem{303}{นาหํ ภิกฺขเว อญฺญํ เอกธมฺมมฺปิ สมนุปสฺสามิ, เยน อนุปฺปนฺนา วา สมฺมาทิฏฺฐิ อุปฺปชฺชติ, อุปฺปนฺนา วา สมฺมาทิฏฺฐิ ปวฑฺฒติ, ยถยิทํ ภิกฺขเว โยนิโส\-มนสิกาโร. โยนิโส ภิกฺขเว มนสิ กโรโต อนุปฺปนฺนา เจว สมฺมาทิฏฺฐิ อุปฺปชฺชติ, อุปฺปนฺนา จ สมฺมาทิฏฺฐิ ปวฑฺฒตีติ. (6)}

\proseitem{304}{นาหํ ภิกฺขเว อญฺญํ เอกธมฺมมฺปิ สมนุปสฺสามิ, เยน\csromansharedfootnote{de82ecdad78954d2}{เยเนวํ (สี, สฺยา, กํ, อิ)} สตฺตา กายสฺส เภทา ปรํ มรณา อปายํ ทุคฺคติํ วินิปาตํ นิรยํ อุปฺปชฺชนฺติ, ยถยิทํ ภิกฺขเว มิจฺฉาทิฏฺฐิ. มิจฺฉาทิฏฺฐิยา ภิกฺขเว สมนฺนาคตา สตฺตา กายสฺส เภทา ปรํ มรณา อปายํ ทุคฺคติํ วินิปาตํ นิรยํ อุปฺปชฺชนฺตีติ. (7)}

\proseitem{305}{นาหํ ภิกฺขเว อญฺญํ เอกธมฺมมฺปิ สมนุปสฺสามิ, เยน สตฺตา กายสฺส เภทา ปรํ มรณา สุคติํ สคฺคํ โลกํ อุปฺปชฺชนฺติ, ยถยิทํ ภิกฺขเว สมฺมาทิฏฺฐิ. สมฺมาทิฏฺฐิยา ภิกฺขเว สมนฺนาคตา สตฺตา กายสฺส เภทา ปรํ มรณา สุคติํ สคฺคํ โลกํ อุปฺปชฺชนฺตีติ. (8)}

\proseitem{306}{มิจฺฉา\-ทิฏฺฐิกสฺส ภิกฺขเว ปุริส\-ปุคฺคลสฺส ยญฺเจว กายกมฺมํ ยถาทิฏฺฐิ สมตฺตํ สมาทินฺนํ, ยญฺจ วจีกมฺมํ ฯเปฯ ยญฺจ มโนกมฺมํ ยถาทิฏฺฐิ สมตฺตํ สมาทินฺนํ, ยา จ เจตนา, ยา จ ปตฺถนา, โย จ ปณิธิ, เย จ สงฺขารา, สพฺเพ เต ธมฺมา อนิฏฺฐาย อกนฺตาย อมนาปาย อหิตาย ทุกฺขาย สํวตฺตนฺติ. ตํ กิสฺส เหตุ, ทิฏฺฐิ หิสฺส\csromansharedfootnote{e9c7c1457c9644a2}{ทิฏฺฐิ หิ (สี, สฺยา, กํ, อิ)} ภิกฺขเว ปาปิกา. เสยฺยถาปิ \csromanfolio{34}ภิกฺขเว นิมฺพพีชํ วา โกสาตกิพีชํ วา ติตฺตกาลาพุพีชํ วา อลฺลาย ปถวิยา\csromansharedfootnote{70198bb73fb950f6}{ปฐวิยา (สี, สฺยา, กํ, อิ)} นิกฺขิตฺตํ ยญฺเจว ปถวิรสํ อุปาทิยติ, ยญฺจ อาโปรสํ อุปาทิยติ. สพฺพํ ตํ ติตฺตกตฺตาย กฏุกตฺตาย อสาตตฺตาย สํวตฺตติ. ตํ กิสฺส เหตุ? พีชํ หิสฺส\csromansharedfootnote{127664b8dded2c53}{วีชํ (สี, สฺยา, กํ, อิ)} ภิกฺขเว ปาปกํ. เอวเมวํ โข ภิกฺขเว มิจฺฉา\-ทิฏฺฐิกสฺส ปุริส\-ปุคฺคลสฺส ยญฺเจว กายกมฺมํ ยถาทิฏฺฐิ สมตฺตํ สมาทินฺนํ, ยญฺจ วจีกมฺมํ ฯเปฯ ยญฺจ มโนกมฺมํ ยถาทิฏฺฐิ สมตฺตํ สมาทินฺนํ, ยา จ เจตนา, ยา จ ปตฺถนา, โย จ ปณิธิ, เย จ สงฺขารา, สพฺเพ เต ธมฺมา อนิฏฺฐาย อกนฺตาย อมนาปาย อหิตาย ทุกฺขาย สํวตฺตนฺติ. ตํ กิสฺส เหตุ? ทิฏฺฐิ หิสฺส ภิกฺขเว ปาปิกาติ. (9)}

\proseitemwithclosermidruled{307}{สมฺมา\-ทิฏฺฐิกสฺส ภิกฺขเว ปุริส\-ปุคฺคลสฺส ยญฺเจว กายกมฺมํ ยถาทิฏฺฐิ สมตฺตํ สมาทินฺนํ, ยญฺจ วจีกมฺมํ ฯเปฯ ยญฺจ มโนกมฺมํ ยถาทิฏฺฐิ สมตฺตํ สมาทินฺนํ, ยา จ เจตนา, ยา จ ปตฺถนา, โย จ ปณิธิ, เย จ สงฺขารา, สพฺเพ เต ธมฺมา อิฏฺฐาย กนฺตาย มนาปาย หิตาย สุขาย สํวตฺตนฺติ. ตํ กิสฺส เหตุ? ทิฏฺฐิ หิสฺส ภิกฺขเว ภทฺทิกา. เสยฺยถาปิ ภิกฺขเว อุจฺฉุพีชํ วา สาลิพีชํ วา มุทฺทิกาพีชํ วา อลฺลาย ปถวิยา นิกฺขิตฺตํ ยญฺเจว ปถวิรสํ อุปาทิยติ, ยญฺจ อาโปรสํ อุปาทิยติ, สพฺพํ ตํ มธุรตฺตาย สาตตฺตาย อเสจนกตฺตาย สํวตฺตติ. ตํ กิสฺส เหตุ? พีชํ หิสฺส ภิกฺขเว ภทฺทกํ. เอวเมวํ โข ภิกฺขเว สมฺมาทิฏฺฐิกส ปุริส\-ปุคฺคลสฺส ยญฺเจว กายกมฺมํ ยถาทิฏฺฐิ สมตฺตํ สมาทินฺนํ, ยญฺจ วจีกมฺมํ ฯเปฯ ยญฺจ มโนกมฺมํ ยถาทิฏฺฐิ สมตฺตํ สมาทินฺนํ, ยา จ เจตนา, ยา จ ปตฺถนา, โย จ ปณิธิ, เย จ สงฺขารา, สพฺเพ เต ธมฺมา อิฏฺฐาย กนฺตาย มนาปาย หิตาย สุขาย สํวตฺตนฺติ. ตํ กิสฺส เหตุ, ทิฏฺฐิ หิสฺส ภิกฺขเว ภทฺทิกาติ. (10)}{วคฺโค ทุติโย.}
\csromanpalirecto
\gambhira{16. \textbf{เอกธมฺมปาฬิ 3. ตติยวคฺค}}
\proseitem{308}{เอกปุคฺคโล ภิกฺขเว โลเก อุปฺปชฺชมาโน อุปฺปชฺชติ พหุชน\-อหิตาย พหุชน\-อสุขาย พหุโน ชนสฺส อนตฺถาย อหิตาย \csromanfolio{35}ทุกฺขาย เทวมนุสฺสานํ. กตโม เอกปุคฺคโล? มิจฺฉาทิฏฺฐิโก โหติ วิปรีต\-ทสฺสโน, โส พหุชนํ สทฺธมฺมา วุฏฺฐาเปตฺวา อสทฺธมฺเม ปติฏฺฐาเปติ. อยํ โข ภิกฺขเว เอกปุคฺคโล โลเก อุปฺปชฺชมาโน อุปฺปชฺชติ พหุชน\-อหิตาย พหุชน\-อสุขาย พหุโน ชนสฺส อนตฺถาย อหิตาย ทุกฺขาย เทว\-มนุสฺสา\-นนฺติ. (1)}

\proseitem{309}{เอกปุคฺคโล ภิกฺขเว โลเก อุปฺปชฺชมาโน อุปฺปชฺชติ พหุชนหิตาย พหุชน\-สุขาย พหุโน ชนสฺส อตฺถาย หิตาย สุขาย เทวมนุสฺสานํ. กตโม เอกปุคฺคโล? สมฺมาทิฏฺฐิโก โหติ อวิปรีต\-ทสฺสโน, โส พหุชนํ อสทฺธมฺมา วุฏฺฐาเปตฺวา สทฺธมฺเม ปติฏฺฐาเปติ. อยํ โข ภิกฺขเว เอกปุคฺคโล โลเก อุปฺปชฺชมาโน อุปฺปชฺชติ พหุชนหิตาย พหุชน\-สุขาย พหุโน ชนสฺส อตฺถาย หิตาย สุขาย เทว\-มนุสฺสา\-นนฺติ. (2)}

\proseitem{310}{นาหํ ภิกฺขเว อญฺญํ เอกธมฺมมฺปิ สมนุปสฺสามิ, ยํ เอวํ มหาสาวชฺชํ, ยถยิทํ ภิกฺขเว มิจฺฉาทิฏฺฐิ. มิจฺฉาทิฏฺฐิ\-ปรมานิ ภิกฺขเว มหาสาวชฺชานีติ\csromansharedfootnote{2d3f92df046ee171}{วชฺชานีติ (สี, สฺยา, กํ)}. (3)}

\proseitem{311}{นาหํ ภิกฺขเว อญฺญํ เอกปุคฺคลมฺปิ สมนุปสฺสามิ, โย เอวํ พหุชน\-อหิตาย ปฏิปนฺโน พหุชน\-อสุขาย พหุโน ชนสฺส อนตฺถาย อหิตาย ทุกฺขาย เทวมนุสฺสานํ. ยถยิทํ ภิกฺขเว มกฺขลิ โมฆปุริโส. เสยฺยถาปิ ภิกฺขเว นทีมุเข ขิปฺปํ\csromansharedfootnote{c034f8849f783d6c}{ขิปํ (สี, สฺยา, กํ, อิ)} อุฑฺเฑยฺย\csromansharedfootnote{017a8d8e62184d5c}{โอฑฺเฑยฺย (สี), อุชฺเฌยฺย (ก)} พหูนํ มจฺฉานํ อหิตาย ทุกฺขาย อนยาย พฺยสนาย. เอวเมวํ โข ภิกฺขเว มกฺขลิ โมฆปุริโส มนุสฺสขิปฺปํ มญฺเญ โลเก อุปฺปนฺโน พหูนํ สตฺตานํ อหิตาย ทุกฺขาย อนยาย พฺยสนายาติ. (4)}

\proseitem{312}{ทุรกฺขาเต ภิกฺขเว ธมฺมวินเย โย จ สมาทเปติ\csromansharedfootnote{d3f7159cd9cb5fc3}{สมาทาเปติ (?)}, ยญฺจ สมาทเปติ, โย จ สมาทปิโต, ตถตฺตาย ปฏิปชฺชติ, สพฺเพ เต พหุํ อปุญฺญํ ปสวนฺติ. ตํ กิสฺส เหตุ? ทุรกฺขาตตฺตา ภิกฺขเว ธมฺมสฺสาติ. (5)}

\proseitem{313}{\csromanfolio{36}สฺวากฺขาเต ภิกฺขเว ธมฺมวินเย โย จ สมาทเปติ, ยญฺจ สมาทเปติ, โย จ สมาทปิโต ตถตฺตาย ปฏิปชฺชติ, สพฺเพ เต พหุํ ปุญฺญํ ปสวนฺติ. ตํ กิสฺส เหตุ? สฺวากฺขาตตฺตา ภิกฺขเว ธมฺมสฺสาติ. (6)}

\proseitem{314}{ทุรกฺขาเต ภิกฺขเว ธมฺมวินเย ทายเกน มตฺตา ชานิตพฺพา, โน ปฏิคฺคาหเกน. ตํ กิสฺส เหตุ? ทุรกฺขาตตฺตา ภิกฺขเว ธมฺมสฺสาติ. (7)}

\proseitem{315}{สฺวากฺขาเต ภิกฺขเว ธมฺมวินเย ปฏิคฺคาหเกน มตฺตา ชานิตพฺพา, โน ทายเกน. ตํ กิสฺส เหตุ? สฺวากฺขาตตฺตา ภิกฺขเว ธมฺมสฺสาติ. (8)}

\proseitem{316}{ทุรกฺขาเต ภิกฺขเว ธมฺมวินเย โย อารทฺธวีริโย, โส ทุกฺขํ วิหรติ. ตํ กิสฺส เหตุ? ทุรกฺขาตตฺตา ภิกฺขเว ธมฺมสฺสาติ. (9)}

\proseitem{317}{สฺวากฺขาเต ภิกฺขเว ธมฺมวินเย โย กุสีโต, โส ทุกฺขํ วิหรติ. ตํ กิสฺส เหตุ? สฺวากฺขาตตฺตา ภิกฺขเว ธมฺมสฺสาติ. (10)}

\proseitem{318}{ทุรกฺขาเต ภิกฺขเว ธมฺมวินเย โย กุสีโต, โส สุขํ วิหรติ. ตํ กิสฺส เหตุ? ทุรกฺขาตตฺตา ภิกฺขเว ธมฺมสฺสาติ. (11)}

\proseitem{319}{สฺวากฺขาเต ภิกฺขเว ธมฺมวินเย โย อารทฺธวีริโย, โส สุขํ วิหรติ. ตํ กิสฺส เหตุ? สฺวากฺขาตตฺตา ภิกฺขเว ธมฺมสฺสาติ. (12)}

\proseitem{320}{เสยฺยถาปิ ภิกฺขเว อปฺปมตฺตโกปิ คูโถ ทุคฺคนฺโธ โหติ. เอวเมวํ โข อหํ ภิกฺขเว อปฺปมตฺตกมฺปิ ภวํ น วณฺเณมิ อนฺตมโส อจฺฉราสงฺฆาต\-มตฺตมฺปิ. (13)}

\proseitemwithclosermidruled{321}{เสยฺยถาปิ ภิกฺขเว อปฺปมตฺตกมฺปิ มุตฺตํ ทุคฺคนฺธํ โหติ. อปฺปมตฺตโกปิ เขโฬ ทุคฺคนฺโธ โหติ. อปฺปมตฺตโกปิ ปุพฺโพ ทุคฺคนฺโธ โหติ. อปฺปมตฺตกมฺปิ โลหิตํ ทุคฺคนฺธํ โหติ. เอวเมวํ โข อหํ ภิกฺขเว อปฺปมตฺตกมฺปิ ภวํ น วณฺเณมิ อนฺตมโส อจฺฉราสงฺฆาต\-มตฺตมฺปิ. (14)}{วคฺโค ตติโย.}
\chapterheadpageread{16. เอกธมฺมปาฬิ \textbf{16. เอกธมฺมปาฬิ 4. จตุตฺถวคฺค}}
\proseitem{322}{\csromanfolio{37}เสยฺยถาปิ ภิกฺขเว อปฺปมตฺตกํ อิมสฺมิํ ชมฺพุทีเป อาราม\-รามเณยฺยกํ วนรามเณยฺยกํ ภูมิรามเณยฺยกํ โปกฺขรณิ\-รามเณยฺยกํ. อถ โข เอตเทว พหุตรํ, ยทิทํ อุกฺกูลวิกูลํ นทีวิทุคฺคํ ขาณุ\-กณฺฏกฏฺฐานํ\csromansharedfootnote{22778d9db7911a32}{ขาณุกณฺฑกธานํ (สี, อิ)} ปพฺพตวิสมํ. เอวเมวํ โข ภิกฺขเว อปฺปกา เต สตฺตา, เย ถลชา. อถ โข เอเตว สตฺตา พหุตรา, เย โอทกา. (1)}

\proseitem{323}{เอวเมวํ โข ภิกฺขเว อปฺปกา เต สตฺตา, เย มนุสฺเสสุ ปจฺจาชายนฺติ. อถ โข เอเตว สตฺตา พหุตรา, เย อญฺญตฺร มนุสฺเสหิ ปจฺจาชายนฺติ.}

\prose{เอวเมวํ โข ภิกฺขเว อปฺปกา เต สตฺตา, เย มชฺฌิเมสุ ชนปเทสุ ปจฺจาชายนฺติ. อถ โข เอเตว สตฺตา พหุตรา, เย ปจฺจนฺติเมสุ ชนปเทสุ ปจฺจาชายนฺติ อวิญฺญาตาเรสุ มิลกฺเขสุ\csromansharedfootnote{5b4cb73123d2532e}{มิลกฺขูสุ (ก)}. (2)}

\proseitem{324}{เอวเมวํ โข ภิกฺขเว อปฺปกา เต สตฺตา, เย ปญฺญวนฺโต อชฬา อเนฬมูคา ปฏิพลา สุภาสิต\-ทุพฺภาสิตสฺส อตฺถมญฺญาตุํ. อถ โข เอเตว สตฺตา พหุตรา, เย ทุปฺปญฺญา ชฬา เอฬมูคา น ปฏิพลา สุภาสิต\-ทุพฺภาสิตสฺส อตฺถมญฺญาตุํ. (3)}

\proseitem{325}{เอวเมวํ โข ภิกฺขเว อปฺปกา เต สตฺตา, เย อริเยน ปญฺญาจกฺขุนา สมนฺนาคตา, อถ โข เอเตว สตฺตา พหุตรา, เย อวิชฺชาคตา สมฺมูฬฺหา. (4)}

\proseitem{326}{เอวเมวํ โข ภิกฺขเว อปฺปกา เต สตฺตา, เย ลภนฺติ ตถาคตํ ทสฺสนาย. อถ โข เอเตว สตฺตา พหุตรา, เย น ลภนฺติ ตถาคตํ ทสฺสนาย. (5)}

\proseitem{327}{เอวเมวํ โข ภิกฺขเว อปฺปกา เต สตฺตา, เย ลภนฺติ ตถาคต\-ปฺปเวทิตํ ธมฺมวินยํ สวนาย. อถ โข เอเตว สตฺตา พหุตรา, เย น ลภนฺติ ตถาคต\-ปฺปเวทิตํ ธมฺมวินยํ สวนาย. (6)}

\proseitem{328}{\csromanfolio{38}เอวเมวํ โข ภิกฺขเว อปฺปกา เต สตฺตา, เย สุตฺวา ธมฺมํ ธาเรนฺติ. อถ โข เอเตว สตฺตา พหุตรา เย สุตฺวา ธมฺมํ น ธาเรนฺติ. (7)}

\proseitem{329}{เอวเมวํ โข ภิกฺขเว อปฺปกา เต สตฺตา, เย ธาตานํ\csromansharedfootnote{0e167e3d5f245cbf}{ธตานํ (สี, สฺยา, กํ อิ)} ธมฺมานํ อตฺถํ อุปปริกฺขนฺติ. อถ โข เอเตว สตฺตา พหุตรา, เย ธาตานํ ธมฺมานํ อตฺถํ น อุปปริกฺขนฺติ. (8)}

\proseitem{330}{เอวเมวํ โข ภิกฺขเว อปฺปกา เต สตฺตา, เย อตฺถมญฺญาย ธมฺมมญฺญาย ธมฺมา\-นุธมฺมํ ปฏิปชฺชนฺติ. อถ โข เอเตว สตฺตา พหุตรา, เย อตฺถมญฺญาย ธมฺมมญฺญาย ธมฺมา\-นุธมฺมํ น ปฏิปชฺชนฺติ. (9)}

\proseitem{331}{เอวเมวํ โข ภิกฺขเว อปฺปกา เต สตฺตา, เย สํเวชนิเยสุ ฐาเนสุ สํวิชฺชนฺติ. อถ โข เอเตว สตฺตา พหุตรา, เย สํเวชนิเยสุ ฐาเนสุ น สํวิชฺชนฺติ. (10)}

\proseitem{332}{เอวเมวํ โข ภิกฺขเว อปฺปกา เต สตฺตา, เย สํวิคฺคา โยนิโส ปทหนฺติ. อถ โข เอเตว สตฺตา พหุตรา, เย สํวิคฺคา โยนิโส น ปทหนฺติ. (11)}

\proseitem{333}{เอวเมวํ โข ภิกฺขเว อปฺปกา เต สตฺตา, เย ววสฺสคฺคา\-รมฺมณํ กริตฺวา ลภนฺติ สมาธิํ\csromansharedfootnote{5b358b68e3414ca6}{จิตฺตสฺส สมาธิํ (สี)} ลภนฺติ จิตฺตสฺเส\-กคฺคตํ\csromansharedfootnote{aaefcb5739e3c98b}{จิตฺตสฺเสกคฺคํ (สี)}. อถ โข เอเตว สตฺตา พหุตรา, เย ววสฺสคฺคา\-รมฺมณํ กริตฺวา น ลภนฺติ สมาธิํ น ลภนฺติ จิตฺตสฺเส\-กคฺคตํ. (12)}

\proseitem{334}{เอวเมวํ โข ภิกฺขเว อปฺปกา เต สตฺตา, เย อนฺนคฺค\-รสคฺคานํ ลาภิโน. อถ โข เอเตว สตฺตา พหุตรา, เย อนฺนคฺค\-รสคฺคานํ น ลาภิโน อุญฺเฉน กปาลาภเตน ยาเปนฺติ. (13)}

\proseitem{335}{เอวเมวํ โข ภิกฺขเว อปฺปกา เต สตฺตา, เย อตฺถรสสฺส ธมฺมรสสฺส วิมุตฺติรสสฺส ลาภิโน. อถ โข เอเตว สตฺตา พหุตรา, เย อตฺถรสสฺส ธมฺมรสสฺส วิมุตฺติรสสฺส น ลาภิโน. ตสฺมาติห \csromanfolio{39}ภิกฺขเว เอวํ สิกฺขิตพฺพํ “อตฺถรสสฺส ธมฺมรสสฺส วิมุตฺติรสสฺส ลาภิโน ภวิสฺสามา”ติ. เอวํ หิ โว ภิกฺขเว สิกฺขิตพฺพนฺติ. (14)}

\proseitem{336-338}{เสยฺยถาปิ ภิกฺขเว อปฺปมตฺตกํ อิมสฺมิํ ชมฺพุทีเป อาราม\-รามเณยฺยกํ วนรามเณยฺยกํ ภูมิรามเณยฺยกํ โปกฺขรณิ\-รามเณยฺยกํ. อถ โข เอตเทว พหุตรํ, ยทิทํ อุกฺกูลวิกูลํ นทีวิทุคฺคํ ขาณุ\-กณฺฏกฏฺฐานํ ปพฺพตวิสมํ. เอวเมวํ โข ภิกฺขเว อปฺปกา เต สตฺตา, เย มนุสฺสา จุตา มนุสฺเสสุ ปจฺจาชายนฺติ. อถ โข เอเตว สตฺตา พหุตรา, เย มนุสฺสา จุตา นิรเย ปจฺจาชายนฺติ ฯเปฯ ติรจฺฉาน\-โยนิยา ปจฺจาชายนฺติ ฯเปฯ เปตฺติวิสเย ปจฺจาชายนฺติ. \mbox{(15-17)}}

\proseitem{339-341}{เอวเมวํ โข ภิกฺขเว อปฺปกา เต สตฺตา, เย มนุสฺสา จุตา เทเวสุ ปจฺจาชายนฺติ. อถ โข เอเตว สตฺตา พหุตรา, เย มนุสฺสา จุตา นิรเย ปจฺจาชายนฺติ. ติรจฺฉาน\-โยนิยา ปจฺจาชายนฺติ. เปตฺติวิสเย ปจฺจาชายนฺติ. \mbox{(18-20)}}

\proseitem{342-344}{เอวเมวํ โข ภิกฺขเว อปฺปกา เต สตฺตา, เย เทวา จุตา เทเวสุ ปจฺจาชายนฺติ. อถ โข เอเตว สตฺตา พหุตรา, เย เทวา จุตา นิรเย ปจฺจาชายนฺติ. ติรจฺฉาน\-โยนิยา ปจฺจาชายนฺติ. เปตฺติวิสเย ปจฺจาชายนฺติ. \mbox{(21-23)}}

\proseitem{345-347}{เอวเมวํ โข ภิกฺขเว อปฺปกา เต สตฺตา, เย เทวา จุตา มนุสฺเสสุ ปจฺจาชายนฺติ. อถ โข เอเตว สตฺตา พหุตรา, เย เทวา จุตา นิรเย ปจฺจาชายนฺติ. ติรจฺฉาน\-โยนิยา ปจฺจาชายนฺติ. เปตฺติวิสเย ปจฺจาชายนฺติ. \mbox{(24-26)}}

\proseitem{348-350}{เอวเมวํ โข ภิกฺขเว อปฺปกา เต สตฺตา, เย นิรยา จุตา มนุสฺเสสุ ปจฺจาชายนฺติ. อถ โข เอเตว สตฺตา พหุตรา, เย นิรยา จุตา นิรเย ปจฺจาชายนฺติ. ติรจฺฉาน\-โยนิยา ปจฺจาชายนฺติ. เปตฺติวิสเย ปจฺจาชายนฺติ. \mbox{(27-29)}}

\proseitem{351-353}{เอวเมวํ โข ภิกฺขเว อปฺปกา เต สตฺตา, เย นิรยา จุตา เทเวสุ ปจฺจาชายนฺติ. อถ โข เอเตว สตฺตา พหุตรา, เย นิรยา จุตา นิรเย ปจฺจาชายนฺติ. ติรจฺฉาน\-โยนิยา ปจฺจาชายนฺติ. เปตฺติวิสเย ปจฺจาชายนฺติ. \mbox{(30-32)}}

\proseitem{354-356}{\csromanfolio{40}เอวเมวํ โข ภิกฺขเว อปฺปกา เต สตฺตา, เย ติรจฺฉาน\-โยนิยา จุตา มนุสฺเสสุ ปจฺจาชายนฺติ. อถ โข เอเตว สตฺตา พหุตรา, เย ติรจฺฉาน\-โยนิยา จุตา นิรเย ปจฺจาชายนฺติ. ติรจฺฉาน\-โยนิยา ปจฺจาชายนฺติ. เปตฺติวิสเย ปจฺจาชายนฺติ. \mbox{(33-35)}}

\proseitem{357-359}{เอวเมวํ โข ภิกฺขเว อปฺปกา เต สตฺตา, เย ติรจฺฉาน\-โยนิยา จุตา เทเวสุ ปจฺจาชายนฺติ. อถ โข เอเตว สตฺตา พหุตรา, เย ติรจฺฉาน\-โยนิยา จุตา นิรเย ปจฺจาชายนฺติ. ติรจฺฉาน\-โยนิยา ปจฺจาชายนฺติ. เปตฺติวิสเย ปจฺจาชายนฺติ. \mbox{(36-38)}}

\proseitem{360-362}{เอวเมวํ โข ภิกฺขเว อปฺปกา เต สตฺตา, เย เปตฺติวิสยา จุตา มนุสฺเสสุ ปจฺจาชายนฺติ. อถ โข เอเตว สตฺตา พหุตรา, เย เปตฺติวิสยา จุตา นิรเย ปจฺจาชายนฺติ. ติรจฺฉาน\-โยนิยา ปจฺจาชายนฺติ. เปตฺติวิสเย ปจฺจาชายนฺติ. \mbox{(39-41)}}

\proseitemwithclosermid{363-365}{เอวเมวํ โข ภิกฺขเว อปฺปกา เต สตฺตา, เย เปตฺติวิสยา จุตา เทเวสุ ปจฺจาชายนฺติ. อถ โข เอเตว สตฺตา พหุตรา, เย เปตฺติวิสยา จุตา นิรเย ปจฺจาชายนฺติ. ติรจฺฉาน\-โยนิยา ปจฺจาชายนฺติ. เปตฺติวิสเย ปจฺจาชายนฺติ. \mbox{(42-44)}}{วคฺโค จตุตฺโถ. (ชมฺพุทีป\-เปยฺยาโล นิฏฺฐิโต.)}
\chapterhead{17. ปสาทกรธมฺม\-วคฺค}
\proseitemwithcloser{366-381}{อทฺธมิทํ ภิกฺขเว ลาภานํ ยทิทํ อารญฺญิกตฺตํ\csromansharedfootnote{0803b809b943f41d}{อรญฺญกตฺตํ (สพฺพตฺถ)} ฯเปฯ ปิณฺฑ\-ปาติ\-กตฺตํ. ปํสุกูลิ\-กตฺตํ. เตจีวริกตฺตํ. ธมฺม\-กถิกตฺตํ. วินย\-ธรตฺตํ\csromansharedfootnote{31606dd18a89f15e}{วินยธรกตฺตํ (สฺยา, กํ, อิ, ก)}. พาหุสจฺจํ. ถาวเรยฺยํ. อากปฺปสมฺปทา. ปริวาร\-สมฺปทา. มหาปริวารตา. โกลปุตฺติ. วณฺณ\-โปกฺขรตา. กลฺยาณ\-วากฺกรณตา. อปฺปิจฺฉตา. อปฺปาพาธตาติ. (\csromansharedfootnote{0803b809b943f41d}{อรญฺญกตฺตํ (สพฺพตฺถ)}-16)}{(โสฬส ปสาทกร\-ธมฺมา นิฏฺฐิตา.)}
\chapterheadpageread{18. อปร\-อจฺฉราสงฺฆาต\-วคฺค}
\chapterhead{18. อปร\-อจฺฉราสงฺฆาต\-วคฺค}
\proseitem{382}{\csromanfolio{41}อจฺฉราสงฺฆาต\-มตฺตมฺปิ เจ ภิกฺขเว ภิกฺขุ ปฐมํ ฌานํ ภาเวติ, อยํ วุจฺจติ ภิกฺขเว ภิกฺขุ อริตฺตชฺฌาโน วิหรติ สตฺถุ\-สาสน\-กโร โอวาทปติกโร อโมฆํ รฏฺฐปิณฺฑํ ภุญฺชติ, โก ปน วาโท เย นํ พหุลีกโรนฺตีติ. (1)}

\proseitem{383-389}{อจฺฉราสงฺฆาต\-มตฺตมฺปิ เจ ภิกฺขเว ภิกฺขุ ทุติยํ ฌานํ ภาเวติ ฯเปฯ ตติยํ ฌานํ ภาเวติ ฯเปฯ จตุตฺถํ ฌานํ ภาเวติ ฯเปฯ เมตฺตํ เจโตวิมุตฺติํ ภาเวติ ฯเปฯ กรุณํ เจโตวิมุตฺติํ ภาเวติ ฯเปฯ มุทิตํ เจโตวิมุตฺติํ ภาเวติ ฯเปฯ อุเปกฺขํ เจโตวิมุตฺติํ ภาเวติ ฯเปฯ. \mbox{(2-8)}}

\proseitem{390-393}{กาเย กายานุปสฺสี วิหรติ อาตาปี สมฺปชาโน สติมา วิเนยฺย โลเก อภิชฺฌา\-โทมนสฺสํ ฯเปฯ เวทนาสุ เวทนานุปสฺสี วิหรติ. จิตฺเต จิตฺตานุปสฺสี วิหรติ. ธมฺเมสุ ธมฺมานุปสฺสี วิหรติ อาตาปี สมฺปชาโน สติมา วิเนยฺย โลเก อภิชฺฌา\-โทมนสฺสํ. \mbox{(9-12)}}

\proseitem{394-397}{อนุปฺปนฺนานํ ปาปกานํ อกุสลานํ ธมฺมานํ อนุปฺปาทาย ฉนฺทํ ชเนติ วายมติ, วีริยํ\csromansharedfootnote{881fcf767bfdc1fb}{วิริยํ (สี, สฺยา, กํ, อิ)} อารภติ, จิตฺตํ ปคฺคณฺหาติ ปทหติ. อุปฺปนฺนานํ ปาปกานํ อกุสลานํ ธมฺมานํ ปหานาย ฉนฺทํ ชเนติ วายมติ, วีริยํ อารภติ, จิตฺตํ ปคฺคณฺหาติ ปทหติ. อนุปฺปนฺนานํ กุสลานํ ธมฺมานํ อุปฺปาทาย ฉนฺทํ ชเนติ วายมติ, วีริยํ อารภติ, จิตฺตํ ปคฺคณฺหาติ ปทหติ. อุปฺปนฺนานํ กุสลานํ ธมฺมานํ ฐิติยา อสมฺโมสาย ภิยฺโยภาวาย เวปุลฺลาย ภาวนาย ปาริปูริยา ฉนฺทํ ชเนติ วายมติ, วีริยํ อารภติ, จิตฺตํ ปคฺคณฺหาติ ปทหติ. \mbox{(13-16)}}

\proseitem{398-401}{ฉนฺท\-สมาธิ\-ปธาน\-สงฺขาร\-สมนฺนาคตํ อิทฺธิปาทํ ภาเวติ. วีริย\-สมาธิ\-ปธาน\-สงฺขาร\-สมนฺนาคตํ อิทฺธิปาทํ ภาเวติ. จิตฺต\-สมาธิ\-ปธาน\-สงฺขาร\-สมนฺนาคตํ อิทฺธิปาทํ ภาเวติ. วีมํสา\-สมาธิ\-ปธาน\-สงฺขาร\-สมนฺนาคตํ อิทฺธิปาทํ ภาเวติ. \mbox{(17-20)}}

\proseitem{402-406}{สทฺธินฺทฺริยํ ภาเวติ. วีริยินฺทฺริยํ ภาเวติ. สตินฺทฺริยํ ภาเวติ. สมาธินฺทฺริยํ ภาเวติ. ปญฺญินฺทฺริยํ ภาเวติ. \mbox{(21-25)}}

\proseitem{407-411}{\csromanfolio{42}สทฺธาพลํ ภาเวติ. วีริยพลํ ภาเวติ. สติพลํ ภาเวติ. สมาธิพลํ ภาเวติ. ปญฺญาพลํ ภาเวติ. \mbox{(26-30)}}

\proseitem{412-418}{สติสมฺโพชฺฌงฺคํ ภาเวติ. ธมฺมวิจย\-สมฺโพชฺฌงฺคํ ภาเวติ. วีริย\-สมฺโพชฺฌงฺคํ ภาเวติ. ปีติสมฺโพชฺฌงฺคํ ภาเวติ. ปสฺสทฺธิ\-สมฺโพชฺฌงฺคํ ภาเวติ. สมาธิ\-สมฺโพชฺฌงฺคํ ภาเวติ. อุเปกฺขา\-สมฺโพชฺฌงฺคํ ภาเวติ. \mbox{(31-37)}}

\proseitem{419-426}{สมฺมาทิฏฺฐิํ ภาเวติ. สมฺมาสงฺกปฺปํ ภาเวติ. สมฺมาวาจํ ภาเวติ. สมฺมากมฺมนฺตํ ภาเวติ. สมฺมาอาชีวํ ภาเวติ. สมฺมาวายามํ ภาเวติ. สมฺมาสติํ ภาเวติ. สมฺมาสมาธิํ ภาเวติ. \mbox{(38-45)}}

\proseitem{427-434}{1 อชฺฌตฺตํ รูปสญฺญี พหิทฺธา รูปานิ ปสฺสติ ปริตฺตานิ สุวณฺณ\-ทุพฺพณฺณานิ, ตานิ อภิภุยฺย “ชานามิ ปสฺสามี”ติ เอวํสญฺญี โหติ. อชฺฌตฺตํ รูปสญฺญี พหิทฺธา รูปานิ ปสฺสติ อปฺปมาณานิ สุวณฺณ\-ทุพฺพณฺณานิ, ตานิ อภิภุยฺย “ชานามิ ปสฺสามี”ติ เอวํสญฺญี โหติ. อชฺฌตฺตํ อรูปสญฺญี พหิทฺธา รูปานิ ปสฺสติ ปริตฺตานิ สุวณฺณ\-ทุพฺพณฺณานิ, ตานิ อภิภุยฺย “ชานามิ ปสฺสามี”ติ เอวํสญฺญี โหติ. อชฺฌตฺตํ อรูปสญฺญี พหิทฺธา รูปานิ ปสฺสติ อปฺปมาณานิ สุวณฺณ\-ทุพฺพณฺณานิ, ตานิ อภิภุยฺย “ชานามิ ปสฺสามี”ติ เอวํสญฺญี โหติ. อชฺฌตฺตํ อรูปสญฺญี พหิทฺธา รูปานิ ปสฺสติ นีลานิ นีลวณฺณานิ นีลนิทสฺสนานิ นีลนิภาสานิ, ตานิ อภิภุยฺย “ชานามิ ปสฺสามี”ติ เอวํสญฺญี โหติ. อชฺฌตฺตํ อรูปสญฺญี พหิทฺธา รูปานิ ปสฺสติ ปีตานิ ปีตวณฺณานิ ปีตนิทสฺสนานิ ปีตนิภาสานิ, ตานิ อภิภุยฺย “ชานามิ ปสฺสามี”ติ เอวํสญฺญี โหติ. อชฺฌตฺตํ อรูปสญฺญี พหิทฺธา รูปานิ ปสฺสติ โลหิตกานิ โลหิตก\-วณฺณานิ โลหิตก\-นิทสฺสนานิ โลหิตก\-นิภาสานิ, ตานิ อภิภุยฺย “ชานามิ ปสฺสามี”ติ เอวํสญฺญี โหติ. อชฺฌตฺตํ อรูปสญฺญี พหิทฺธา รูปานิ ปสฺสติ โอทาตานิ โอทาตวณฺณานิ โอทาต\-นิทสฺสนานิ โอทาตนิภาสานิ, ตานิ อภิภุยฺย “ชานามิ ปสฺสามี”ติ เอวํสญฺญี โหติ. \mbox{(46-53)}}

\proseitem{435-442}{รูปี รูปานิ ปสฺสติ. อชฺฌตฺตํ อรูปสญฺญี พหิทฺธา รูปานิ ปสฺสติ. “สุภนฺ”เตว อธิมุตฺโต โหติ. สพฺพโส รูปสญฺญานํ สมติกฺกมา ปฏิฆ\-สญฺญานํ อตฺถงฺคมา นานตฺต\-สญฺญานํ อมนสิการา}

\chapterheadpageread{18. อปร\-อจฺฉราสงฺฆาต\-วคฺค}
\prose{\csromanfolio{43}“อนนฺโต อากาโส”ติ อากาสา\-นญฺจา\-ยตนํ อุปสมฺปชฺช วิหรติ. สพฺพโส อากาสา\-นญฺจา\-ยตนํ สมติกฺกมฺม “อนนฺตํ วิญฺญาณนฺ”ติ วิญฺญาณญฺจา\-ยตนํ อุปสมฺปชฺช วิหรติ. สพฺพโส วิญฺญาณญฺจา\-ยตนํ สมติกฺกมฺม “นตฺถิ กิญฺจี”ติ อากิญฺจญฺญา\-ยตนํ อุปสมฺปชฺช วิหรติ. สพฺพโส อากิญฺจญฺญา\-ยตนํ สมติกฺกมฺม เนว\-สญฺญา\-นา\-สญฺญา\-ยตนํ อุปสมฺปชฺช วิหรติ. สพฺพโส เนว\-สญฺญา\-นา\-สญฺญา\-ยตนํ สมติกฺกมฺม สญฺญาเวทยิต\-นิโรธํ อุปสมฺปชฺช วิหรติ. \mbox{(54-61)}}

\proseitem{443-452}{ปถวิกสิณํ ภาเวติ. อาโปกสิณํ ภาเวติ. เตโชกสิณํ ภาเวติ. วาโยกสิณํ ภาเวติ. นีลกสิณํ ภาเวติ. ปีตกสิณํ ภาเวติ. โลหิตกสิณํ ภาเวติ. โอทาตกสิณํ ภาเวติ. อากาสกสิณํ ภาเวติ. วิญฺญาณกสิณํ ภาเวติ. \mbox{(62-71)}}

\proseitem{453-462}{อสุภสญฺญํ ภาเวติ. มรณสญฺญํ ภาเวติ. อาหาเร ปฏิกูลสญฺญํ ภาเวติ. สพฺพโลเก อนภิรติ\-สญฺญํ\csromansharedfootnote{132d1e78e6a06a9d}{อนภิรตสญฺญํ (สี, สฺยา, กํ, อิ)} ภาเวติ. อนิจฺจสญฺญํ ภาเวติ. อนิจฺเจ ทุกฺขสญฺญํ ภาเวติ. ทุกฺเข อนตฺตสญฺญํ ภาเวติ. ปหานสญฺญํ ภาเวติ. วิราคสญฺญํ ภาเวติ. นิโรธสญฺญํ ภาเวติ. \mbox{(72-81)}}

\proseitem{463-472}{อนิจฺจสญฺญํ ภาเวติ. อนตฺตสญฺญํ ภาเวติ. มรณสญฺญํ ภาเวติ. อาหาเร ปฏิกูลสญฺญํ ภาเวติ. สพฺพโลเก อนภิรติ\-สญฺญํ ภาเวติ. อฏฺฐิกสญฺญํ ภาเวติ. ปุฬวกสญฺญํ\csromansharedfootnote{c12fc936704708dc}{ปุฬุวกสญฺญํ (ก)} ภาเวติ. วินีลกสญฺญํ ภาเวติ. วิจฺฉิทฺทก\-สญฺญํ ภาเวติ. อุทฺธุมาตก\-สญฺญํ ภาเวติ. \mbox{(82-91)}}

\proseitem{473-482}{พุทฺธานุสฺสติํ ภาเวติ. ธมฺมานุสฺสติํ ภาเวติ. สํฆานุสฺสติํ ภาเวติ. สีลานุสฺสติํ ภาเวติ. จาคานุสฺสติํ ภาเวติ. เทวตานุสฺสติํ ภาเวติ. อานาปานสฺสติํ ภาเวติ. มรณสฺสติํ ภาเวติ. กายคตาสติํ ภาเวติ. อุปสมา\-นุสฺสติํ ภาเวติ. \mbox{(92-101)}}

\proseitem{483-492}{\csromanfolio{44}ปฐมชฺฌาน\-สหคตํ สทฺธินฺทฺริยํ ภาเวติ. วีริยินฺทฺริยํ ภาเวติ. สตินฺทฺริยํ ภาเวติ. สมาธินฺทฺริยํ ภาเวติ. ปญฺญินฺทฺริยํ ภาเวติ. สทฺธาพลํ ภาเวติ. วีริยพลํ ภาเวติ. สติพลํ ภาเวติ. สมาธิพลํ ภาเวติ. ปญฺญาพลํ ภาเวติ. \mbox{(102-111)}}

\proseitem{493-562}{ทุติยชฺฌาน\-สหคตํ ฯเปฯ ตติยชฺฌาน\-สหคตํ ฯเปฯ จตุตฺถชฺฌาน\-สหคตํ ฯเปฯ เมตฺตาสหคตํ ฯเปฯ กรุณา\-สหคตํ ฯเปฯ มุทิตา\-สหคตํ ฯเปฯ อุเปกฺขา\-สหคตํ สทฺธินฺทฺริยํ ภาเวติ. วีริยินฺทฺริยํ ภาเวติ. สตินฺทฺริยํ ภาเวติ. สมาธินฺทฺริยํ ภาเวติ. ปญฺญินฺทฺริยํ ภาเวติ. สทฺธาพลํ ภาเวติ. วีริยพลํ ภาเวติ. สติพลํ ภาเวติ. สมาธิพลํ ภาเวติ. ปญฺญาพลํ ภาเวติ. อยํ วุจฺจติ ภิกฺขเว ภิกฺขุ อริตฺตชฺฌาโน วิหรติ สตฺถุ\-สาสน\-กโร โอวาทปติกโร อโมฆํ รฏฺฐปิณฺฑํ ภุญฺชติ. โก ปน วาโท เย นํ พหุลีกโรนฺตีติ. \mbox{(112-181)}}

\prose{(อปร\-อจฺฉราสงฺฆาต\-วคฺโค.)}

\chapterhead{19. กายคตาสติ\-วคฺค}
\proseitem{563}{ยสฺส กสฺสจิ ภิกฺขเว มหาสมุทฺโท เจตสา ผุโฏ, อนฺโตคธา ตสฺส กุนฺนทิโย ยา กาจิ สมุทฺทงฺคมา. เอวเมวํ ภิกฺขเว ยสฺส กสฺสจิ กายคตา สติ ภาวิตา พหุลีกตา, อนฺโตคธา ตสฺส กุสลา ธมฺมา เย เกจิ วิชฺชาภาคิยาติ. (1)}

\proseitem{564-570}{เอกธมฺโม ภิกฺขเว ภาวิโต พหุลีกโต มหโต สํเวคาย สํวตฺตติ. มหโต อตฺถาย สํวตฺตติ. มหโต โยคกฺเขมาย สํวตฺตติ. สติสมฺปชญฺญาย สํวตฺตติ. ญาณทสฺสน\-ปฺปฏิลาภาย สํวตฺตติ. ทิฏฺฐ\-ธมฺม\-สุข\-วิหาราย สํวตฺตติ. วิชฺชา\-วิมุตฺติ\-ผล\-สจฺฉิกิริยาย สํวตฺตติ. กตโม เอกธมฺโม? กายคตา สติ. อยํ โข ภิกฺขเว เอกธมฺโม ภาวิโต พหุลีกโต มหโต สํเวคาย สํวตฺตติ. มหโต อตฺถาย สํวตฺตติ. มหโต โยคกฺเขมาย สํวตฺตติ. สติสมฺปชญฺญาย สํวตฺตติ. ญาณทสฺสน\-ปฺปฏิลาภาย สํวตฺตติ. ทิฏฺฐ\-ธมฺม\-สุข\-วิหาราย สํวตฺตติ. วิชฺชา\-วิมุตฺติ\-ผล\-สจฺฉิกิริยาย สํวตฺตตีติ. \mbox{(2-8)}}

\chapterheadpageread{19. กายคตาสติ\-วคฺค}
\proseitem{571}{\csromanfolio{45}เอกธมฺเม ภิกฺขเว ภาวิเต พหุลีกเต กาโยปิ ปสฺสมฺภติ, จิตฺตมฺปิ ปสฺสมฺภติ, วิตกฺก\-วิจาราปิ วูปสมฺมนฺติ, เกวลาปิ วิชฺชาภาคิยา ธมฺมา ภาวนา\-ปาริปูริํ คจฺฉนฺติ. กตมสฺมิํ เอกธมฺเม? กายคตาย สติยา. อิมสฺมิํ โข ภิกฺขเว เอกธมฺเม ภาวิเต พหุลีกเต กาโยปิ ปสฺสมฺภติ, จิตฺตมฺปิ ปสฺสมฺภติ, วิตกฺก\-วิจาราปิ วูปสมฺมนฺติ, เกวลาปิ วิชฺชาภาคิยา ธมฺมา ภาวนา\-ปาริปูริํ คจฺฉนฺตีติ. (9)}

\proseitem{572}{เอกธมฺเม ภิกฺขเว ภาวิเต พหุลีกเต อนุปฺปนฺนา เจว อกุสลา ธมฺมา นุปฺปชฺชนฺติ, อุปฺปนฺนา จ อกุสลา ธมฺมา ปหียนฺติ. กตมสฺมิํ เอกธมฺเม? กายคตาย สติยา. อิมสฺมิํ โข ภิกฺขเว เอกธมฺเม ภาวิเต พหุลีกเต อนุปฺปนฺนา เจว อกุสลา ธมฺมา นุปฺปชฺชนฺติ, อุปฺปนฺนา จ อกุสลา ธมฺมา ปหียนฺตีติ. (10)}

\proseitem{573}{เอกธมฺเม ภิกฺขเว ภาวิเต พหุลีกเต อนุปฺปนฺนา เจว กุสลา ธมฺมา อุปฺปชฺชนฺติ, อุปฺปนฺนา จ กุสลา ธมฺมา ภิยฺโยภาวาย เวปุลฺลาย สํวตฺตนฺติ. กตมสฺมิํ เอกธมฺเม? กายคตาย สติยา. อิมสฺมิํ โข ภิกฺขเว เอกธมฺเม ภาวิเต พหุลีกเต อนุปฺปนฺนา เจว กุสลา ธมฺมา อุปฺปชฺชนฺติ, อุปฺปนฺนา จ กุสลา ธมฺมา ภิยฺโยภาวาย เวปุลฺลาย สํวตฺตนฺตีติ. (11)}

\proseitem{574}{เอกธมฺเม ภิกฺขเว ภาวิเต พหุลีกเต อวิชฺชา ปหียติ, วิชฺชา อุปฺปชฺชติ, อสฺมิมาโน ปหียติ, อนุสยา สมุคฺฆาตํ คจฺฉนฺติ, สํโยชนา ปหียนฺติ. กตมสฺมิํ เอกธมฺเม? กายคตาย สติยา. อิมสฺมิํ โข ภิกฺขเว เอกธมฺเม ภาวิเต พหุลีกเต อวิชฺชา ปหียติ, วิชฺชา อุปฺปชฺชติ, อสฺมิมาโน ปหียติ, อนุสยา สมุคฺฆาตํ คจฺฉนฺติ, สํโยชนา ปหียนฺตีติ. (12)}

\proseitem{575-576}{เอกธมฺโม ภิกฺขเว ภาวิโต พหุลีกโต ปญฺญาปเภทาย สํวตฺตติ. อนุปาทา\-ปรินิพฺพานาย สํวตฺตติ. กตโม เอกธมฺโม? กายคตา สติ. อยํ โข ภิกฺขเว เอกธมฺโม ภาวิโต พหุลีกโต ปญฺญาปเภทาย สํวตฺตติ. อนุปาทา\-ปรินิพฺพานาย สํวตฺตตีติ. \mbox{(13-14)}}

\proseitem{577-579}{\csromanfolio{46}เอกธมฺเม ภิกฺขเว ภาวิเต พหุลีกเต อเนกธาตุ\-ปฏิเวโธ โหติ. นานาธาตุ\-ปฏิเวโธ โหติ. อเนกธาตุ\-ปฏิสมฺภิทา โหติ. กตมสฺมิํ เอกธมฺเม? กายคตาย สติยา. อิมสฺมิํ โข ภิกฺขเว เอกธมฺเม ภาวิเต พหุลีกเต อเนกธาตุ\-ปฏิเวโธ โหติ. นานาธาตุ\-ปฏิเวโธ โหติ. อเนกธาตุ\-ปฏิสมฺภิทา โหตีติ. \mbox{(15-17)}}

\proseitem{580-583}{เอกธมฺโม ภิกฺขเว ภาวิโต พหุลีกโต โสตาปตฺติผล\-สจฺฉิกิริยาย สํวตฺตติ. สกทาคามิ\-ผล\-สจฺฉิกิริยาย สํวตฺตติ. อนาคามิ\-ผล\-สจฺฉิกิริยาย สํวตฺตติ. อรหตฺตผล\-สจฺฉิกิริยาย สํวตฺตติ. กตโม เอกธมฺโม? กายคตา สติ. อยํ โข ภิกฺขเว เอกธมฺโม ภาวิโต พหุลีกโต โสตาปตฺติผล\-สจฺฉิกิริยาย สํวตฺตติ. สกทาคามิ\-ผล\-สจฺฉิกิริยาย สํวตฺตติ. อนาคามิ\-ผล\-สจฺฉิกิริยาย สํวตฺตติ. อรหตฺตผล\-สจฺฉิกิริยาย สํวตฺตตีติ. \mbox{(18-21)}}

\proseitem{584-599}{เอกธมฺโม ภิกฺขเว ภาวิโต พหุลีกโต ปญฺญา\-ปฏิลาภาย สํวตฺตติ. ปญฺญาวุทฺธิยา สํวตฺตติ. ปญฺญาเวปุลฺลาย สํวตฺตติ. มหาปญฺญตาย สํวตฺตติ. ปุถุปญฺญตาย สํวตฺตติ. วิปุล\-ปญฺญตาย สํวตฺตติ. คมฺภีร\-ปญฺญตาย สํวตฺตติ. อสามนฺต\-ปญฺญตาย\csromansharedfootnote{5838b7b651d0ebcf}{อสมตฺถปญฺญตาย (สฺยา, กํ), อสมตฺตปญฺญตาย (ก), อสมนฺตปญฺญตาย (ฏีกา) ปฏิสํฏฺฐ 2. 248 ปิฏฺเฐ ปสฺสิตพฺพํ.} สํวตฺตติ. ภูริปญฺญตาย สํวตฺตติ. ปญฺญาพาหุลฺลาย สํวตฺตติ. สีฆปญฺญตาย สํวตฺตติ. ลหุปญฺญตาย สํวตฺตติ. หาสปญฺญตาย\csromansharedfootnote{ccb151b679aa19f6}{หาสุปญฺญตาย (สี, อิ)} สํวตฺตติ. ชวน\-ปญฺญตาย สํวตฺตติ. ติกฺข\-ปญฺญตาย สํวตฺตติ. นิพฺเพธิก\-ปญฺญตาย สํวตฺตติ. กตโม เอกธมฺโม? กายคตา สติ. อยํ โข ภิกฺขเว เอกธมฺโม ภาวิโต พหุลีกโต ปญฺญา\-ปฏิลาภาย สํวตฺตติ. ปญฺญาวุทฺธิยา สํวตฺตติ. ปญฺญาเวปุลฺลาย สํวตฺตติ. มหาปญฺญตาย สํวตฺตติ. ปุถุปญฺญตาย สํวตฺตติ. วิปุล\-ปญฺญตาย สํวตฺตติ. คมฺภีร\-ปญฺญตาย สํวตฺตติ. อสามนฺต\-ปญฺญตาย สํวตฺตติ. ภูริปญฺญตาย สํวตฺตติ. ปญฺญาพาหุลฺลาย สํวตฺตติ. สีฆปญฺญตาย สํวตฺตติ. ลหุปญฺญตาย สํวตฺตติ. หาสปญฺญตาย สํวตฺตติ. ชวน\-ปญฺญตาย สํวตฺตติ. ติกฺข\-ปญฺญตาย สํวตฺตติ. นิพฺเพธิก\-ปญฺญตาย สํวตฺตตีติ. \mbox{(22-37)}}

\csromanheader{2}{(กายคตาสติ\-วคฺโค.)}
\chapterheadpageread{20. อมตวคฺค}
\chapterhead{20. อมตวคฺค}
\proseitem{600}{\csromanfolio{47}อมตํ เต ภิกฺขเว น ปริภุญฺชนฺติ, เย กายคตาสติํ น ปริภุญฺชนฺติ. อมตํ เต ภิกฺขเว ปริภุญฺชนฺติ, เย กายคตาสติํ ปริภุญฺช\-นฺตีติ. (1)}

\proseitem{601}{อมตํ เตสํ ภิกฺขเว อปริภุตฺตํ, เยสํ กายคตาสติ อปริภุตฺตา. อมตํ เตสํ ภิกฺขเว ปริภุตฺตํ, เยสํ กายคตาสติ ปริภุตฺตาติ. (2)}

\proseitem{602}{อมตํ เตสํ ภิกฺขเว ปริหีนํ, เยสํ กายคตาสติ ปริหีนา. อมตํ เตสํ ภิกฺขเว อปริหีนํ, เยสํ กายคตาสติ อปริหีนาติ. (3)}

\proseitem{603}{อมตํ เตสํ ภิกฺขเว วิรทฺธํ, เยสํ กายคตาสติ วิรทฺธา. อมตํ เตสํ ภิกฺขเว อารทฺธํ\csromansharedfootnote{761a462efb351a8a}{อวิรทฺธํ (ก)}, เยสํ กายคตาสติ อารทฺธาติ. (4)}

\proseitem{604}{อมตํ เต ภิกฺขเว ปมาทิํสุ, เย กายคตาสติํ ปมาทิํสุ. อมตํ เต ภิกฺขเว น ปมาทิํสุ, เย กายคตาสติํ น ปมาทิํสุ. (5)}

\proseitem{605}{อมตํ เตสํ ภิกฺขเว ปมุฏฺฐํ, เยสํ กายคตาสติ ปมุฏฺฐา. อมตํ เตสํ ภิกฺขเว อปฺปมุฏฺฐํ, เยสํ กายคตาสติ อปฺปมุฏฺฐาติ. (6)}

\proseitem{606}{อมตํ เตสํ ภิกฺขเว อนาเสวิตํ, เยสํ กายคตาสติ อนาเสวิตา. อมตํ เตสํ ภิกฺขเว อาเสวิตํ, เยสํ กายคตาสติ อาเสวิตาติ. (7)}

\proseitem{607}{อมตํ เตสํ ภิกฺขเว อภาวิตํ, เยสํ กายคตาสติ อภาวิตา. อมตํ เตสํ ภิกฺขเว ภาวิตํ, เยสํ กายคตาสติ ภาวิตาติ. (8)}

\proseitem{608}{\csromanfolio{48}อมตํ เตสํ ภิกฺขเว อพหุลีกตํ, เยสํ กายคตาสติ อพหุลีกตา. อมตํ เตสํ ภิกฺขเว พหุลีกตํ, เยสํ กายคตาสติ พหุลีกตาติ. (9)}

\proseitem{609}{อมตํ เตสํ ภิกฺขเว อนภิญฺญาตํ, เยสํ กายคตาสติ อนภิญฺญาตา. อมตํ เตสํ ภิกฺขเว อภิญฺญาตํ, เยสํ กายคตาสติ อภิญฺญาตาติ. (10)}

\proseitem{610}{อมตํ เตสํ ภิกฺขเว อปริญฺญาตํ, เยสํ กายคตาสติ อปริญฺญาตา. อมตํ เตสํ ภิกฺขเว ปริญฺญาตํ, เยสํ กายคตาสติ ปริญฺญาตาติ. (11)}

\proseitem{611}{อมตํ เตสํ ภิกฺขเว อสจฺฉิกตํ, เยสํ กายคตาสติ อสจฺฉิกตา. อมตํ เตสํ ภิกฺขเว สจฺฉิกตํ, เยสํ กายคตาสติ สจฺฉิกตาติ. (12) ( )\csromansharedfootnote{74e01824b36d97c2}{(เอกกนิปาตสฺส สุตฺตสหสฺสํ สมตฺตํ.) (สี, สฺยา, กํ, อิ)}}

\prosewithclosermajor{(อิทมโวจ ภควา. อตฺตมนา เต ภิกฺขู ภควโต ภาสิตํ อภินนฺทุนฺติ.)\csromansharedfootnote{58456e94fd777360}{( ) เอตฺถนฺตเร ปาโฐ สีสฺยากํอิโปตฺถเกสุ นตฺถิ.} (อมตวคฺโค.)}{เอกก\-นิปาต\-ปาฬิ นิฏฺฐิตา.}
\csromanpalirecto
\gambhira{\textbf{ทุกนิปาตปาฬิ}}
\namakkaram{นโม ตสฺส ภควโต อรหโต สมฺมา\-สมฺพุทฺธสฺส.}
\csromanmatikaanchor{mtk.0}
\csromanmatikaanchor{mtk.53}
\csromantocmarknopage{chapter}{ติกนิปาตปาฬิ}{mtk.0}
\bookmark[dest={mtk.0},level=0]{ติกนิปาตปาฬิ}
\chapterhead{1. ปฐม\-ปณฺณาสก}
\chapterhead{1. กมฺมกรณ\-วคฺค}
\csromanheader{2}{1. \textbf{วชฺชสุตฺต}}
\proseitem{1}{\csromanfolio{49}เอวํ เม สุตํ\csromandash{}เอกํ สมยํ ภควา สาวตฺถิยํ วิหรติ เชตวเน อนาถ\-ปิณฺฑิกสฺส อาราเม. ตตฺร โข ภควา ภิกฺขู อามนฺเตสิ “ภิกฺขโว”ติ. “ภทนฺเต”ติ เต ภิกฺขู ภควโต ปจฺจสฺโสสุํ. ภควา เอตทโวจ\csromandash{}}

\prose{เทฺวมานิ ภิกฺขเว วชฺชานิ. กตมานิ เทฺว? ทิฏฺฐธมฺมิกญฺจ วชฺชํ สมฺปรายิกญฺจ วชฺชํ. กตมญฺจ ภิกฺขเว ทิฏฺฐ\-ธมฺมิกํ วชฺชํ? อิธ ภิกฺขเว เอกจฺโจ ปสฺสติ โจรํ อาคุจาริํ ราชาโน คเหตฺวา วิวิธา กมฺมการณา\csromansharedfootnote{7989c8fc7b9a6a4f}{วิวิธานิ กมฺมกรณานิ (ก)} กาเรนฺเต กสาหิปิ ตาเฬนฺเต เวตฺเตหิปิ ตาเฬนฺเต อทฺธ\-ทณฺฑเกหิปิ ตาเฬนฺเต หตฺถมฺปิ ฉินฺทนฺเต ปาทมฺปิ ฉินฺทนฺเต หตฺถปาทมฺปิ ฉินฺทนฺเต กณฺณมฺปิ ฉินฺทนฺเต นาสมฺปิ ฉินฺทนฺเต กณฺณนาสมฺปิ ฉินฺทนฺเต พิลงฺค\-ถาลิกมฺปิ กโรนฺเต สงฺข\-มุณฺฑิกมฺปิ กโรนฺเต ราหุมุขมฺปิ กโรนฺเต โชติมาลิกมฺปิ กโรนฺเต หตฺถ\-ปชฺโชติกมฺปิ กโรนฺเต เอรก\-วตฺติกมฺปิ กโรนฺเต จีรก\-วาสิกมฺปิ กโรนฺเต เอเณยฺยกมฺปิ กโรนฺเต พฬิส\-มํสิกมฺปิ กโรนฺเต กหาปณิกมฺปิ กโรนฺเต \csromanfolio{50}ขารา\-ปตจฺฉิกมฺปิ\csromansharedfootnote{40fdca843456e3ac}{ขาราปฏิจฺฉกมฺปิ (สฺยา, กํ, ก)} กโรนฺเต ปลิฆ\-ปริวตฺติกมฺปิ กโรนฺเต ปลาล\-ปีฐกมฺปิ\csromansharedfootnote{1c608027d0f47cb1}{ปลาลปิฏฺฐิกมฺปิ (สี)} กโรนฺเต ตตฺเตนปิ เตเลน โอสิญฺจนฺเต สุนเขหิปิ ขาทาเปนฺเต ชีวนฺตมฺปิ สูเล อุตฺตาเสนฺเต อสินาปิ สีสํ ฉินฺทนฺเต.}

\prose{ตสฺส เอวํ โหติ “ยถารูปานํ โข ปาปกานํ กมฺมานํ เหตุ โจรํ อาคุจาริํ ราชาโน คเหตฺวา วิวิธา กมฺมการณา กาเรนฺติ กสาหิปิ ตาเฬนฺติ, เวตฺเตหิปิ ตาเฬนฺติ, อทฺธ\-ทณฺฑเกหิปิ ตาเฬนฺติ, หตฺถมฺปิ ฉินฺทนฺติ, ปาทมฺปิ ฉินฺทนฺติ, หตฺถปาทมฺปิ ฉินฺทนฺติ, กณฺณมฺปิ ฉินฺทนฺติ, นาสมฺปิ ฉินฺทนฺติ, กณฺณนาสมฺปิ ฉินฺทนฺติ, พิลงฺค\-ถาลิกมฺปิ กโรนฺติ, สงฺข\-มุณฺฑิกมฺปิ กโรนฺติ, ราหุมุขมฺปิ กโรนฺติ, โชติมาลิกมฺปิ กโรนฺติ, หตฺถ\-ปชฺโชติกมฺปิ กโรนฺติ, เอรก\-วตฺติกมฺปิ กโรนฺติ, จีรก\-วาสิกมฺปิ กโรนฺติ, เอเณยฺยกมฺปิ กโรนฺติ, พฬิส\-มํสิกมฺปิ กโรนฺติ, กหาปณิกมฺปิ กโรนฺติ, ขารา\-ปตจฺฉิกมฺปิ กโรนฺติ, ปลิฆ\-ปริวตฺติกมฺปิ กโรนฺติ, ปลาล\-ปีฐกมฺปิ กโรนฺติ, ตตฺเตนปิ เตเลน โอสิญฺจนฺติ, สุนเขหิปิ ขาทาเปนฺติ, ชีวนฺตมฺปิ สูเล อุตฺตาเสนฺติ, อสินาปิ สีสํ ฉินฺทนฺติ. อหญฺเจว\csromansharedfootnote{cf278ffcd2fa33b2}{อหญฺเจ (?)} โข ปน เอวรูปํ ปาปกมฺมํ กเรยฺยํ, มมฺปิ ราชาโน คเหตฺวา เอวรูปา วิวิธา กมฺมการณา กาเรยฺยุํ “กสาหิปิ ตาเฬยฺยุํ ฯเปฯ อสินาปิ สีสํ ฉินฺเทยฺยุนฺ”ติ. โส ทิฏฺฐ\-ธมฺมิกสฺส วชฺชสฺส ภีโต น ปเรสํ ปาภตํ วิลุมฺปนฺโต จรติ. อิทํ วุจฺจติ ภิกฺขเว ทิฏฺฐ\-ธมฺมิกํ วชฺชํ.}

\prose{กตมญฺจ ภิกฺขเว สมฺปรายิกํ วชฺชํ? อิธ ภิกฺขเว เอกจฺโจ อิติ ปฏิสญฺจิกฺขติ “กาย\-ทุจฺจริตสฺส โข ปน ปาปโก ทุกฺโข วิปาโก อภิสมฺปรายํ, วจีทุจฺจริตสฺส ปาปโก ทุกฺโข วิปาโก อภิสมฺปรายํ, มโน\-ทุจฺจริตสฺส ปาปโก ทุกฺโข วิปาโก อภิสมฺปรายํ. อหญฺเจว โข ปน กาเยน ทุจฺจริตํ จเรยฺยํ, วาจาย ทุจฺจริตํ จเรยฺยํ, มนสา ทุจฺจริตํ จเรยฺยํ, กิญฺจ ตํ ยาหํ น กายสฺส เภทา ปรํ มรณา อปายํ ทุคฺคติํ วินิปาตํ นิรยํ อุปปชฺเชยฺยนฺ”ติ. โส สมฺปรายิกสฺส วชฺชสฺส ภีโต กายทุจฺจริตํ ปหาย กายสุจริตํ ภาเวติ, วจีทุจฺจริตํ ปหาย วจีสุจริตํ ภาเวติ, มโนทุจฺจริตํ ปหาย มโนสุจริตํ ภาเวติ, สุทฺธํ อตฺตานํ ปริหรติ. อิทํ วุจฺจติ ภิกฺขเว สมฺปรายิกํ วชฺชํ. อิมานิ โข ภิกฺขเว เทฺว วชฺชานิ. ตสฺมาติห ภิกฺขเว เอวํ สิกฺขิตพฺพํ “ทิฏฺฐ\-ธมฺมิกสฺส วชฺชสฺส ภายิสฺสาม, สมฺปรายิกสฺส วชฺชสฺส ภายิสฺสาม, วชฺชภีรุโน}

\vspace{\tipitakaparskip}
\csromancenter{กมฺมกรณ\-วคฺค ภวิสฺสาม วชฺช\-ภยทสฺสาวิโน”ติ. เอวํ หิ โว ภิกฺขเว สิกฺขิตพฺพํ. วชฺชภีรุโน ภิกฺขเว วชฺช\-ภยทสฺสาวิโน เอตํ ปาฏิกงฺขํ, ยํ ปริมุจฺจิสฺสติ สพฺพวชฺเชหีติ.  ปฐมํ.}
\csromanheader{2}{2. \textbf{ปธานสุตฺต}}
\proseitem{2}{\csromanfolio{51}เทฺวมานิ ภิกฺขเว ปธานานิ ทุรภิสมฺภวานิ โลกสฺมิํ. กตมานิ เทฺว? ยญฺจ คิหีนํ อคารํ อชฺฌาวสตํ จีวร\-ปิณฺฑปาต\-เสนาสน\-คิลานปจฺจย\-เภสชฺช\-ปริกฺขารา\-นุปฺปทานตฺถํ ปธานํ, ยญฺจ อคารสฺมา อนคาริยํ ปพฺพชิตานํ สพฺพูปธิปฏินิสฺสคฺคตฺถํ ปธานํ. อิมานิ โข ภิกฺขเว เทฺว ปธานานิ ทุรภิสมฺภวานิ โลกสฺมิํ.}

\prose{เอตทคฺคํ ภิกฺขเว อิเมสํ ทฺวินฺนํ ปธานานํ ยทิทํ สพฺพูปธิปฏินิสฺสคฺคตฺถํ ปธานํ. ตสฺมาติห ภิกฺขเว เอวํ สิกฺขิตพฺพํ “สพฺพูปธิปฏินิสฺสคฺคตฺถํ ปธานํ ปทหิสฺสามา”ติ. เอวํ หิ โว ภิกฺขเว สิกฺขิตพฺพนฺติ.  ทุติยํ.}

\csromanheader{2}{3. \textbf{ตปนียสุตฺต}}
\proseitem{3}{เทฺวเม ภิกฺขเว ธมฺมา ตปนียา. กตเม เทฺว? อิธ ภิกฺขเว เอกจฺจสฺส กายทุจฺจริตํ กตํ โหติ, อกตํ โหติ กายสุจริตํ. วจีทุจฺจริตํ กตํ โหติ, อกตํ โหติ วจีสุจริตํ. มโนทุจฺจริตํ กตํ โหติ, อกตํ โหติ มโนสุจริตํ. โส “กายทุจฺจริตํ เม กตนฺ”ติ ตปฺปติ, “อกตํ เม กายสุจริตนฺ”ติ ตปฺปติ. “วจีทุจฺจริตํ เม กตนฺ”ติ ตปฺปติ, “อกตํ เม วจีสุจริตนฺ”ติ ตปฺปติ. “มโนทุจฺจริตํ เม กตนฺ”ติ ตปฺปติ, “อกตํ เม มโนสุจริตนฺ”ติ ตปฺปติ. อิเม โข ภิกฺขเว เทฺว ธมฺมา ตปนียาติ.  ตติยํ.}

\csromanheader{2}{4. \textbf{อตปนียสุตฺต}}
\proseitem{4}{เทฺวเม ภิกฺขเว ธมฺมา อตปนียา. กตเม เทฺว? อิธ ภิกฺขเว เอกจฺจสฺส กายสุจริตํ กตํ โหติ, อกตํ โหติ กายทุจฺจริตํ. วจีสุจริตํ กตํ โหติ, อกตํ โหติ วจีทุจฺจริตํ. มโนสุจริตํ กตํ โหติ, อกตํ โหติ มโนทุจฺจริตํ. โส “กายสุจริตํ เม กตนฺ”ติ น ตปฺปติ, “อกตํ เม กายทุจฺจริตนฺ”ติ น ตปฺปติ. “วจีสุจริตํ เม กตนฺ”ติ น ตปฺปติ, “อกตํ เม วจีทุจฺจริตนฺ”ติ น ตปฺปติ. “มโนสุจริตํ เม กตนฺ”ติ น ตปฺปติ, “อกตํ เม มโนทุจฺจริตนฺ”ติ น ตปฺปติ. อิเม โข ภิกฺขเว เทฺว ธมฺมา อตปนียาติ.  จตุตฺถํ.}

\csromanpalirecto
\gambhira{ทุกกนิปาตปาฬิ \textbf{5. อุปญฺญาตสุตฺต}}
\proseitem{5}{\csromanfolio{52}ทฺวินฺนาหํ ภิกฺขเว ธมฺมานํ อุปญฺญาสิํ. ยา จ อสนฺตุฏฺฐิตา กุสเลสุ ธมฺเมสุ, ยา จ อปฺปฏิวานิตา ปธานสฺมิํ. อปฺปฏิวานี สุทาหํ ภิกฺขเว ปทหามิ “กามํ ตโจ จ นฺหารุ\csromansharedfootnote{51a640e57c2c8c9f}{นหารุ (สี, สฺยา, กํ, อิ)} จ อฏฺฐิ จ อวสิสฺสตุ สรีเร อุปสุสฺสตุ มํสโลหิตํ, ยํ ตํ ปุริสถาเมน ปุริสวีริเยน ปุริส\-ปรกฺกเมน ปตฺตพฺพํ, น ตํ อปาปุณิตฺวา วีริยสฺส สณฺฐานํ ภวิสฺสตี”ติ. ตสฺส มยฺหํ ภิกฺขเว อปฺปมาทา\-ธิคตา สมฺโพธิ อปฺปมาทา\-ธิคโต อนุตฺตโร โยคกฺขโม. ตุเมฺห เจปิ ภิกฺขเว อปฺปฏิวานํ ปทเหยฺยาถ “กามํ ตโจ จ นฺหารุ จ อฏฺฐิ จ อวสิสฺสตุ สรีเร อุปสุสฺสตุ มํสโลหิตํ, ยํ ตํ ปุริสถาเมน ปุริสวีริเยน ปุริส\-ปรกฺกเมน ปตฺตพฺพํ, น ตํ อปาปุณิตฺวา วีริยสฺส สณฺฐานํ ภวิสฺสตี”ติ. ตุเมฺหปิ ภิกฺขเว นจิรสฺเสว, ยสฺสตฺถาย กุลปุตฺตา สมฺมเทว อคารสฺมา อนคาริยํ ปพฺพชนฺติ, ตทนุตฺตรํ พฺรหฺมจริย\-ปริโยสานํ ทิฏฺเฐว ธมฺเม สยํ อภิญฺญา สจฺฉิกตฺวา อุปสมฺปชฺช วิหริสฺสถ. ตสฺมาติห ภิกฺขเว เอวํ สิกฺขิตพฺพํ “อปฺปฏิวานํ ปทหิสฺสาม ‘กามํ ตโจ จ นฺหารุ จ อฏฺฐิ จ อวสิสฺสตุ สรีเร อุปสุสฺสตุ มํสโลหิตํ, ยํ ตํ ปุริสถาเมน ปุริสวีริเยน ปุริส\-ปรกฺกเมน ปตฺตพฺพํ, น ตํ อปาปุณิตฺวา วีริยสฺส สณฺฐานํ ภวิสฺสตี’ติ”. เอวํ หิ โว ภิกฺขเว สิกฺขิตพฺพนฺติ.  ปญฺจมํ.}

\csromanheader{2}{6. \textbf{สญฺโญชนสุตฺต}}
\proseitem{6}{เทฺวเม ภิกฺขเว ธมฺมา. กตเม เทฺว? ยา จ สํโยชนิเยสุ ธมฺเมสุ อสฺสาทา\-นุปสฺสิตา, ยา จ สํโยชนิเยสุ ธมฺเมสุ นิพฺพิทา\-นุปสฺสิตา. สํโยชนิเยสุ ภิกฺขเว ธมฺเมสุ อสฺสาทานุปสฺสี วิหรนฺโต ราคํ น ปชหติ, โทสํ น ปชหติ, โมหํ น ปชหติ. ราคํ อปฺปหาย โทสํ อปฺปหาย โมหํ อปฺปหาย น ปริมุจฺจติ ชาติยา ชราย มรเณน โสเกหิ ปริเทเวหิ ทุกฺเขหิ โทมนสฺเสหิ อุปายาเสหิ, น “ปริมุจฺจติ ทุกฺขสฺมา”ติ วทามิ.}

\prose{สํโยชนิเยสุ ภิกฺขเว ธมฺเมสุ นิพฺพิทานุปสฺสี วิหรนฺโต ราคํ ปชหติ, โทสํ ปชหติ, โมหํ ปชหติ. ราคํ ปหาย โทสํ ปหาย}

\vspace{\tipitakaparskip}
\csromancenter{กมฺมกรณ\-วคฺค โมหํ ปหาย ปริมุจฺจติ ชาติยา ชราย มรเณน โสเกหิ ปริเทเวหิ ทุกฺเขหิ โทมนสฺเสหิ อุปายาเสหิ, “ปริมุจฺจติ ทุกฺขสฺมา”ติ วทามิ. อิเม โข ภิกฺขเว เทฺว ธมฺมาติ.  ฉฏฺฐํ.}
\csromanheader{2}{7. \textbf{กณฺหสุตฺต}}
\proseitem{7}{\csromanfolio{53}เทฺวเม ภิกฺขเว ธมฺมา กณฺหา. กตเม เทฺว? อหิริกญฺจ อโนตฺตปฺปญฺจ. อิเม โข ภิกฺขเว เทฺว ธมฺมา กณฺหาติ.  สตฺตมํ.}

\csromanheader{2}{8. \textbf{สุกฺกสุตฺต}}
\proseitem{8}{เทฺวเม ภิกฺขเว ธมฺมา สุกฺกา. กตเม เทฺว? หิรี\csromansharedfootnote{b67bc7960d532281}{หิริ (สี, สฺยา, กํ, อิ)} จ โอตฺตปฺปญฺจ. อิเม โข ภิกฺขเว เทฺว ธมฺมา สุกฺกาติ.  อฏฺฐมํ.}

\csromanheader{2}{9. \textbf{จริยสุตฺต}}
\proseitem{9}{เทฺวเม ภิกฺขเว ธมฺมา สุกฺกา โลกํ ปาเลนฺติ. กตเม เทฺว? หิรี จ โอตฺตปฺปญฺจ. อิเม โข ภิกฺขเว เทฺว สุกฺกา ธมฺมา โลกํ น ปาเลยฺยุํ, นยิธ ปญฺญาเยถ “มาตา”ติ วา “มาตุจฺฉา”ติ วา “มาตุลานี”ติ วา “อาจริยภริยา”ติ วา “ครูนํ ทารา”ติ วา, สมฺเภทํ โลโก อคมิสฺส, ยถา อเชฬกา กุกฺกุฏสูกรา โสณสิงฺคาลา\csromansharedfootnote{3bbd6e811407694c}{โสณสิคาลา (สิ, สฺยา, กํ, อิ)}. ยสฺมา จ โข ภิกฺขเว อิเม เทฺว สุกฺกา ธมฺมา โลกํ ปาเลนฺติ, ตสฺมา ปญฺญายติ\csromansharedfootnote{a64d7ec651404150}{ปญฺญายนฺติ (สี)} “มาตา”ติ วา “มาตุจฺฉา”ติ วา “มาตุลานี”ติ วา “อาจริยภริยา”ติ วา “ครูนํ ทารา”ติ วาติ.  นวมํ.}

\csromanheader{2}{10. วสฺสูปนายิก\-สุตฺต}
\proseitemwithclosermidruled{10}{เทฺวมา ภิกฺขเว วสฺสูปนายิกา. กตมา เทฺว? ปุริมิกา จ ปจฺฉิมิกา จ. อิมา โข ภิกฺขเว เทฺว วสฺสู\-ปนายิกาติ.  ทสมํ.}{กมฺมกรณ\-วคฺโค ปฐโม.}
\csromancenter{\textbf{ตสฺสุทฺทานํ}}
\setlength{\csromangathapagewidth}{0pt}
\setlength{\csromangathawakwidth}{0pt}
\csromangathameasuregroup{วชฺชา ปธานา เทฺว ตปนียา, \\ สํโยชนญฺจ กณฺหญฺจ,}{\csromangathabat{วชฺชา ปธานา เทฺว ตปนียา,}{อุปญฺญาเตน ปญฺจมํ.} \\ \csromangathabat{สํโยชนญฺจ กณฺหญฺจ,}{สุกฺกํ จริยา วสฺสูปนายิเกน วคฺโค.}}
\csromangathasetleft{วชฺชา ปธานา เทฺว ตปนียา, \\ สํโยชนญฺจ กณฺหญฺจ,}
\csromangathagroup{\csromangathastanza{\csromangathabat{วชฺชา ปธานา เทฺว ตปนียา,}{อุปญฺญาเตน ปญฺจมํ.} \\ \csromangathabat{สํโยชนญฺจ กณฺหญฺจ,}{สุกฺกํ จริยา วสฺสูปนายิเกน วคฺโค.}}}

\csromanpalirecto
\gambhira{ทุกกนิปาตปาฬิ 2. อธิกรณ\-วคฺค}
\proseitem{11}{\csromanfolio{54}เทฺวมานิ ภิกฺขเว พลานิ. กตมานิ เทฺว? ปฏิสงฺขานพลญฺจ ภาวนาพลญฺจ. กตมญฺจ ภิกฺขเว ปฏิสงฺขานพลํ? อิธ ภิกฺขเว เอกจฺโจ อิติ ปฏิสญฺจิกฺขติ “กาย\-ทุจฺจริตสฺส โข ปาปโก วิปาโก ทิฏฺเฐ เจว ธมฺเม อภิสมฺปรายญฺจ, วจีทุจฺจริตสฺส ปาปโก วิปาโก ทิฏฺเฐ เจว ธมฺเม อภิสมฺปรายญฺจ, มโน\-ทุจฺจริตสฺส ปาปโก วิปาโก ทิฏฺเฐ เจว ธมฺเม อภิสมฺปรายญฺจา”ติ. โส อิติ ปฏิสงฺขาย กายทุจฺจริตํ ปหาย กายสุจริตํ ภาเวติ, วจีทุจฺจริตํ ปหาย วจีสุจริตํ ภาเวติ, มโนทุจฺจริตํ ปหาย มโนสุจริตํ ภาเวติ, สุทฺธํ อตฺตานํ ปริหรติ. อิทํ วุจฺจติ ภิกฺขเว ปฏิสงฺขานพลํ.}

\prose{กตมญฺจ ภิกฺขเว ภาวนาพลํ? ตตฺร ภิกฺขเว ยมิทํ\csromansharedfootnote{ee166d3bd0f984bc}{ยทิทํ (สี)} ภาวนาพลํ, เสขานเมตํ\csromansharedfootnote{c2504b9c0148f017}{เสขเมตํ (สี, สฺยา, กํ)} พลํ. เสขํ หิ โส ภิกฺขเว พลํ อาคมฺม ราคํ ปชหติ, โทสํ ปชหติ, โมหํ ปชหติ. ราคํ ปหาย โทสํ ปหาย โมหํ ปหาย ยํ อกุสลํ, น ตํ กโรติ. ยํ ปาปํ, น ตํ เสวติ. อิทํ วุจฺจติ ภิกฺขเว ภาวนาพลํ. อิมานิ โข ภิกฺขเว เทฺว พลานีติ. (1)}

\proseitem{12}{เทฺวมานิ ภิกฺขเว พลานิ. กตมานิ เทฺว? ปฏิสงฺขานพลญฺจ ภาวนาพลญฺจ. กตมญฺจ ภิกฺขเว ปฏิสงฺขานพลํ? อิธ ภิกฺขเว เอกจฺโจ อิติ ปฏิสญฺจิกฺขติ “กาย\-ทุจฺจริตสฺส โข ปาปโก วิปาโก ทิฏฺเฐ เจว ธมฺเม อภิสมฺปรายญฺจ, วจีทุจฺจริตสฺส ปาปโก วิปาโก ทิฏฺเฐ เจว ธมฺเม อภิสมฺปรายญฺจ, มโน\-ทุจฺจริตสฺส ปาปโก วิปาโก ทิฏฺเฐ เจว ธมฺเม อภิสมฺปรายญฺจา”ติ. โส อิติ ปฏิสงฺขาย กายทุจฺจริตํ ปหาย กายสุจริตํ ภาเวติ, วจีทุจฺจริตํ ปหาย วจีสุจริตํ ภาเวติ, มโนทุจฺจริตํ ปหาย มโนสุจริตํ ภาเวติ, สุทฺธํ อตฺตานํ ปริหรติ. อิทํ วุจฺจติ ภิกฺขเว ปฏิสงฺขานพลํ.}

\prose{กตมญฺจ ภิกฺขเว ภาวนาพลํ? อิธ ภิกฺขเว ภิกฺขุ สติสมฺโพชฺฌงฺคํ ภาเวติ วิเวกนิสฺสิตํ วิราคนิสฺสิตํ นิโรธ\-นิสฺสิตํ โวสคฺค\-ปริณามิํ. ธมฺมวิจย\-สมฺโพชฺฌงฺคํ ภาเวติ. วีริย\-สมฺโพชฺฌงฺคํ ภาเวติ. ปีติสมฺโพชฺฌงฺคํ ภาเวติ. ปสฺสทฺธิ\-สมฺโพชฺฌงฺคํ ภาเวติ. สมาธิ\-สมฺโพชฺฌงฺคํ ภาเวติ. อุเปกฺขา\-สมฺโพชฺฌงฺคํ ภาเวติ วิเวกนิสฺสิตํ วิราคนิสฺสิตํ นิโรธ\-นิสฺสิตํ \csromanfolio{55}โวสคฺค\-ปริณามิํ. อิทํ วุจฺจติ ภิกฺขเว ภาวนาพลํ. อิมานิ โข ภิกฺขเว เทฺว พลานีติ. (2)}

\proseitem{13}{เทฺวมานิ ภิกฺขเว พลานิ. กตมานิ เทฺว? ปฏิสงฺขานพลญฺจ ภาวนาพลญฺจ. กตมญฺจ ภิกฺขเว ปฏิสงฺขานพลํ? อิธ ภิกฺขเว เอกจฺโจ อิติ ปฏิสญฺจิกฺขติ “กาย\-ทุจฺจริตสฺส โข ปาปโก วิปาโก ทิฏฺเฐ เจว ธมฺเม อภิสมฺปรายญฺจ, วจีทุจฺจริตสฺส โข ปาปโก วิปาโก ทิฏฺเฐ เจว ธมฺเม อภิสมฺปรายญฺจ, มโน\-ทุจฺจริตสฺส โข ปาปโก วิปาโก ทิฏฺเฐ เจว ธมฺเม อภิสมฺปรายญฺจา”ติ. โส อิติ ปฏิสงฺขาย กายทุจฺจริตํ ปหาย กายสุจริตํ ภาเวติ, วจีทุจฺจริตํ ปหาย วจีสุจริตํ ภาเวติ, มโนทุจฺจริตํ ปหาย มโนสุจริตํ ภาเวติ, สุทฺธํ อตฺตานํ ปริหรติ. อิทํ วุจฺจติ ภิกฺขเว ปฏิสงฺขานพลํ.}

\prose{กตมญฺจ ภิกฺขเว ภาวนาพลํ? อิธ ภิกฺขเว ภิกฺขุ วิวิจฺเจว กาเมหิ วิวิจฺจ อกุสเลหิ ธมฺเมหิ สวิตกฺกํ สวิจารํ วิเวกชํ ปีติสุขํ ปฐมํ ฌานํ อุปสมฺปชฺช วิหรติ, วิตกฺก\-วิจารานํ วูปสมา อชฺฌตฺตํ สมฺปสาทนํ เจตโส เอโกทิภาวํ อวิตกฺกํ อวิจารํ สมาธิชํ ปีติสุขํ ทุติยํ ฌานํ อุปสมฺปชฺช วิหรติ, ปีติยา จ วิราคา อุเปกฺขโก จ วิหรติ สโต จ สมฺปชาโน สุขญฺจ กาเยน ปฏิสํเวเทติ, ยํ ตํ อริยา อาจิกฺขนฺติ “อุเปกฺขโก สติมา สุขวิหารี”ติ ตติยํ ฌานํ อุปสมฺปชฺช วิหรติ, สุขสฺส จ ปหานา ทุกฺขสฺส จ ปหานา ปุพฺเพว โสมนสฺส\-โทมนสฺสานํ อตฺถงฺคมา อทุกฺขมสุขํ อุเปกฺขา\-สติ\-ปาริสุทฺธิํ จตุตฺถํ ฌานํ อุปสมฺปชฺช วิหรติ. อิทํ วุจฺจติ ภิกฺขเว ภาวนาพลํ. อิมานิ โข ภิกฺขเว เทฺว พลานีติ. (3)}

\proseitem{12}{เทฺวมา ภิกฺขเว ตถาคตสฺส ธมฺมเทสนา. กตมา เทฺว? สํขิตฺเตน จ วิตฺถาเรน จ. อิมา โข ภิกฺขเว เทฺว ตถาคตสฺส ธมฺมเทสนาติ. (4)}

\proseitem{15}{ยสฺมิํ ภิกฺขเว อธิกรเณ อาปนฺโน\csromansharedfootnote{ff2e36cc75cd8b16}{อาปตฺตาปนฺโน (ก)} จ ภิกฺขุ, โจทโก จ ภิกฺขุ น สาธุกํ อตฺตนาว อตฺตานํ ปจฺจเวกฺขติ, ตสฺเมตํ ภิกฺขเว อธิกรเณ ปาฏิกงฺขํ “ทีฆตฺตาย ขรตฺตาย วาฬตฺตาย สํวตฺติสฺสติ, \csromanfolio{56}ภิกฺขู จ น ผาสุํ\csromansharedfootnote{9e2f4c7ff9f7434e}{น ผาสุ (ก)} วิหริสฺสนฺตี”ติ\csromansharedfootnote{cebecfdffa6445d7}{วิหริสฺสนฺติ (สี, สฺยา, กํ, ก)}. ยสฺมิํ จ โข ภิกฺขเว อธิกรเณ อาปนฺโน จ ภิกฺขุ, โจทโก จ ภิกฺขุ สาธุกํ อตฺตนาว อตฺตานํ ปจฺจเวกฺขติ, ตสฺเมตํ ภิกฺขเว อธิกรเณ ปาฏิกงฺขํ “น ทีฆตฺตาย ขรตฺตาย วาฬตฺตาย สํวตฺติสฺสติ, ภิกฺขู จ ผาสุํ วิหริสฺสนฺตี”ติ.}

\prose{กถญฺจ ภิกฺขเว อาปนฺโน ภิกฺขุ สาธุกํ อตฺตนาว อตฺตานํ ปจฺจเวกฺขติ? อิธ ภิกฺขเว อาปนฺโน ภิกฺขุ อิติ ปฏิสญฺจิกฺขติ\csromandash{}อหํ โข อกุสลํ อาปนฺโน กญฺจิเทว\csromansharedfootnote{5e926657604b5bd9}{กิญฺจิเทว (ก) 4. ตสฺมา มํ โส (สี, สฺยา)} เทสํ กาเยน. มํ โส\csromansharedfootnote{fff476a0436516fc}{อนตฺตมนวาจํ (ก) 6. อนตฺตมนวาจํ นาหํ (ก)} ภิกฺขุ อทฺทส อกุสลํ อาปชฺชมานํ กญฺจิเทว เทสํ กาเยน. โน เจ อหํ อกุสลํ อาปชฺเชยฺยํ กญฺจิเทว เทสํ กาเยน. น มํ โส ภิกฺขุ ปสฺเสยฺย อกุสลํ อาปชฺชมานํ กญฺจิเทว เทสํ กาเยน. ยสฺมา จ โข อหํ อกุสลํ อาปนฺโน กญฺจิเทว เทสํ กาเยน, ตสฺมา มํ โส ภิกฺขุ อทฺทส อกุสลํ อาปชฺชมานํ กญฺจิเทว เทสํ กาเยน. ทิสฺวา จ ปน มํ โส ภิกฺขุ อกุสลํ อาปชฺชมานํ กญฺจิเทว เทสํ กาเยน อนตฺตมโน อโหสิ. อนตฺตมโน สมาโน อนตฺตมน\-วจนํ\csromansharedfootnote{af5e042877daf101}{อตฺตมโน (ก)} มํ โส ภิกฺขุ อวจ. อนตฺตมนวจนาหํ\csromansharedfootnote{2bb0cd1ead8ff7d0}{ตสฺมา มํ โส (สี, สฺยา)} เตน ภิกฺขุนา วุตฺตา สมาโน อนตฺตมโน\csromansharedfootnote{97f99159a927a60c}{อนตฺตมนวาจํ นาหํ (ก)} อโหสิํ. อนตฺตมโน สมาโน ปเรสํ อาโรเจสิํ. อิติ มเมว ตตฺถ อจฺจโย อจฺจคมา สุงฺกทายกํว ภณฺฑสฺมินฺติ. เอวํ โข ภิกฺขเว อาปนฺโน ภิกฺขุ สาธุกํ อตฺตนาว อตฺตานํ ปจฺจเวกฺขติ.}

\prose{กถญฺจ ภิกฺขเว โจทโก ภิกฺขุ สาธุกํ อตฺตนาว อตฺตานํ ปจฺจเวกฺขติ? อิธ ภิกฺขเว โจทโก ภิกฺขุ อิติ ปฏิสญฺจิกฺขติ\csromandash{}อยํ โข ภิกฺขุ อกุสลํ อาปนฺโน กญฺจิเทว เทสํ กาเยน. อหํ อิมํ ภิกฺขุํ อทฺทสํ อกุสลํ อาปชฺชมานํ กญฺจิเทว เทสํ กาเยน. โน เจ อยํ ภิกฺขุ อกุสลํ อาปชฺเชยฺย กญฺจิเทว เทสํ กาเยน. นาหํ อิมํ ภิกฺขุํ ปสฺเสยฺยํ อกุสลํ อาปชฺชมานํ กญฺจิเทว เทสํ กาเยน. ยสฺมา จ โข อยํ ภิกฺขุ อกุสลํ อาปนฺโน กญฺจิเทว เทสํ กาเยน, ตสฺมา อหํ อิมํ ภิกฺขุํ อทฺทสํ อกุสลํ อาปชฺชมานํ กญฺจิเทว เทสํ กาเยน. ทิสฺวา จ ปนาหํ อิมํ ภิกฺขุ อกุสลํ อาปชฺชมานํ กญฺจิเทว เทสํ กาเยน อนตฺตมโน อโหสิํ. \csromanfolio{57}อนตฺตมโน สมาโน อตฺตมนวจนาหํ อิมํ ภิกฺขุํ อวจํ. อนตฺตมนวจนายํ ภิกฺขุ มยา วุตฺโต สมาโน อนตฺตมโน อโหสิ. อนตฺตมโน สมาโน ปเรสํ อาโรเจสิ. อิติ มเมว ตตฺถ อจฺจโย อจฺจคมา สุงฺกทายกํว ภณฺฑสฺมินฺติ. เอวํ โข ภิกฺขเว โจทโก ภิกฺขุ สาธุกํ อตฺตนาว อตฺตานํ ปจฺจเวกฺขติ.}

\prose{ยสฺมิํ ภิกฺขเว อธิกรเณ อาปนฺโน จ ภิกฺขุ, โจทโก จ ภิกฺขุ น สาธุกํ อตฺตนาว อตฺตานํ ปจฺจเวกฺขติ, ตสฺเมตํ ภิกฺขเว อธิกรเณ ปาฏิกงฺขํ “ทีฆตฺตาย ขรตฺตาย วาฬตฺตาย สํวตฺติสฺสติ, ภิกฺขู จ น ผาสุํ วิหริสฺสนฺตี”ติ. ยสฺมิํ จ โข ภิกฺขเว อธิกรเณ อาปนฺโน จ ภิกฺขุ โจทโก จ ภิกฺขุ สาธุกํ อตฺตนาว อตฺตานํ ปจฺจเวกฺขติ, ตสฺเมตํ ภิกฺขเว อธิกรเณ ปาฏิกงฺขํ “น ทีฆตฺตาย ขรตฺตาย วาฬตฺตาย สํวตฺติสฺสติ, ภิกฺขู จ ผาสุํ วิหริสฺสนฺตี”ติ. (5)}

\proseitem{16}{อถ โข อญฺญตโร พฺราหฺมโณ เยน ภควา เตนุปสงฺกมิ, อุปสงฺกมิตฺวา ภควตา สทฺธิํ สมฺโมทิ, สมฺโมทนียํ กถํ สารณียํ วีติสาเรตฺวา เอกมนฺตํ นิสีทิ, เอกมนฺตํ นิสินฺโน โข โส พฺราหฺมโณ ภควนฺตํ เอตทโวจ “โก นุ โข โภ โคตม เหตุ โก ปจฺจโย, เยน มิเธกจฺเจ สตฺตา กายสฺส เภทา ปรํ มรณา อปายํ ทุคฺคติํ วินิปาตํ นิรยํ อุปปชฺชนฺตี”ติ. อธมฺมจริยา\-วิสมจริยา\-เหตุ โข พฺราหฺมณ เอวมิเธกจฺเจ สตฺตา กายสฺส เภทา ปรํ มรณา อปายํ ทุคฺคติํ วินิปาตํ นิรยํ อุปปชฺชนฺตีติ.}

\prose{โก นุ โข โภ โคตม เหตุ โก ปจฺจโย, เยน มิเธกจฺเจ สตฺตา กายสฺส เภทา ปรํ มรณา สุคติํ สคฺคํ โลกํ อุปปชฺชนฺตีติ. ธมฺมจริยา\-สมจริยา\-เหตุ โข พฺราหฺมณ เอวมิเธกจฺเจ สตฺตา กายสฺส เภทา ปรํ มรณา สุคติํ สคฺคํ โลกํ อุปปชฺชนฺตีติ.}

\prose{อภิกฺกนฺตํ โภ โคตม, อภิกฺกนฺตํ โภ โคตม, เสยฺยถาปิ โภ โคตม นิกฺกุชฺชิตํ\csromansharedfootnote{c55eb59ddc1dec73}{นิกุชฺชิตํ (ก)} วา อุกฺกุชฺเชยฺย, ปฏิจฺฉนฺนํ วา วิวเรยฺย, มูฬฺหสฺส วา มคฺคํ อาจิกฺเขยฺย, อนฺธกาเร วา เตลปชฺโชตํ ธาเรยฺย “จกฺขุมนฺโต รูปานิ ทกฺขนฺตี”ติ. เอวเมวํ โภตา โคตเมน อเนก\-ปริยาเยน ธมฺโม ปกาสิโต. เอสาหํ ภวนฺตํ โคตมํ สรณํ คจฺฉามิ ธมฺมญฺจ ภิกฺขุสํฆญฺจ, \csromanfolio{58}อุปาสกํ มํ ภวํ โคตโม ธาเรตุ อชฺชตคฺเค ปาณุเปตํ สรณํ คตนฺติ. (6)}

\proseitem{17}{อถ โข ชาณุสฺโสณิ\csromansharedfootnote{7cb3a9cf3999b608}{ชาณุโสณิ (ก)} พฺราหฺมโณ เยน ภควา เตนุปสงฺกมิ, อุปสงฺกมิตฺวา ภควตา สทฺธิํ สมฺโมทิ, สมฺโมทนียํ กถํ สารณียํ วีติสาเรตฺวา เอกมนฺตํ นิสีทิ, เอกมนฺตํ นิสินฺโน โข ชาณุสฺโสณิ พฺราหฺมโณ ภควนฺตํ เอตทโวจ “โก นุ โข โภ โคตม เหตุ โก ปจฺจโย, เยน มิเธกจฺเจ สตฺตา กายสฺส เภทา ปรํ มรณา อปายํ ทุคฺคติํ วินิปาตํ นิรยํ อุปปชฺชนฺตี”ติ. กตตฺตา จ พฺราหฺมณ อกตตฺตา จ เอวมิเธกจฺเจ สตฺตา กายสฺส เภทา ปรํ มรณา อปายํ ทุคฺคติํ วินิปาตํ นิรยํ อุปปชฺชนฺตีติ. โก ปน โภ โคตม เหตุ โก ปจฺจโย, เยน มิเธกจฺเจ สตฺตา กายสฺส เภทา ปรํ มรณา สุคติํ สคฺคํ โลกํ อุปปชฺชนฺตีติ. กตตฺตา จ พฺราหฺมณ อกตตฺตา จ เอวมิเธกจฺเจ สตฺตา กายสฺส เภทา ปรํ มรณา สุคติํ สคฺคํ โลกํ อุปปชฺชนฺตีติ. น โข อหํ อิมสฺส โภโต โคตมสฺส สํขิตฺเตน ภาสิตสฺส วิตฺถาเรน อตฺถํ อวิภตฺตสฺส วิตฺถาเรน อตฺถํ อาชานามิ, สาธุ เม ภวํ โคตโม ตถา ธมฺมํ เทเสตุ ยถา อหํ อิมสฺส โภโต โคตมสฺส สํขิตฺเตน ภาสิตสฺส วิตฺถาเรน อตฺถํ อวิภตฺตสฺส วิตฺถาเรน อตฺถํ อาชาเนยฺยนฺติ. เตน หิ พฺราหฺมณ สุณาหิ สาธุกํ มนสิ กโรหิ, ภาสิสฺสามีติ. “เอวํ โภ”ติ โข ชาณุสฺโสณิ พฺราหฺมโณ ภควโต ปจฺจสฺโสสิ. ภควา เอตทโวจ\csromandash{}}

\prose{อิธ พฺราหฺมณ เอกจฺจสฺส กายทุจฺจริตํ กตํ โหติ, อกตํ โหติ กายสุจริตํ. วจีทุจฺจริตํ กตํ โหติ, อกตํ โหติ วจีสุจริตํ. มโนทุจฺจริตํ กตํ โหติ, อกตํ โหติ มโนสุจริตํ. เอวํ โข พฺราหฺมณ กตตฺตา จ อกตตฺตา จ เอวมิเธกจฺเจ สตฺตา กายสฺส เภทา ปรํ มรณา อปายํ ทุคฺคติํ วินิปาตํ นิรยํ อุปปชฺชนฺติ. อิธ ปน พฺราหฺมณ เอกจฺจสฺส กายสุจริตํ กตํ โหติ, อกตํ โหติ กายทุจฺจริตํ. วจีสุจริตํ กตํ โหติ, อกตํ โหติ วจีทุจฺจริตํ. มโนสุจริตํ กตํ โหติ, อกตํ โหติ มโนทุจฺจริตํ. เอวํ โข พฺราหฺมณ กตตฺตา จ อกตตฺตา จ เอวมิเธกจฺเจ สตฺตา กายสฺส เภทา ปรํ มรณา สุคติํ สคฺคํ โลกํ อุปปชฺชนฺตีติ.}

\chapterheadpageread{2. อธิกรณ\-วคฺค}
\prose{\csromanfolio{59}อภิกฺกนฺตํ โภ โคตม ฯเปฯ อุปาสกํ มํ ภวํ โคตโม ธาเรตุ อชฺชตคฺเค ปาณุเปตํ สรณํ คตนฺติ. (7)}

\proseitem{18}{อถ โข อายสฺมา อานนฺโท เยน ภควา เตนุปสงฺกมิ, อุปสงฺกมิตฺวา ภควนฺตํ อภิวาเทตฺวา เอกมนฺตํ นิสีทิ, เอกมนฺตํ นิสินฺนํ โข อายสฺมนฺตํ อานนฺทํ ภควา เอตทโวจ “เอกํเสนาหํ อานนฺท อกรณียํ วทามิ กายทุจฺจริตํ วจีทุจฺจริตํ มโนทุจฺจริตนฺ”ติ. ยมิทํ ภนฺเต ภควตา เอกํเสน อกรณียํ อกฺขาตํ กายทุจฺจริตํ วจีทุจฺจริตํ มโนทุจฺจริตํ, ตสฺมิํ อกรณีเย กยิรมาเน โก อาทีนโว ปาฏิกงฺโขติ. ยมิทํ อานนฺท มยา เอกํเสน อกรณียํ อกฺขาตํ กายทุจฺจริตํ วจีทุจฺจริตํ มโนทุจฺจริตํ, ตสฺมิํ อกรณีเย กยิรมาเน อยํ อาทีนโว ปาฏิกงฺโข\csromandash{} อตฺตาปิ อตฺตานํ อุปวทติ, อนุวิจฺจ วิญฺญู ครหนฺติ, ปาปโก กิตฺติสทฺโท อพฺภุคฺคจฺฉติ, สมฺมูโฬฺห กาลํ กโรติ, กายสฺส เภทา ปรํ มรณา อปายํ ทุคฺคติํ วินิปาตํ นิรยํ อุปปชฺชติ. ยมิทํ อานนฺท มยา เอกํเสน อกรณียํ อกฺขาตํ กายทุจฺจริตํ วจีทุจฺจริตํ มโนทุจฺจริตํ, ตสฺมิํ อกรณีเย กยิรมาเน อยํ อาทีนโว ปาฏิกงฺโขติ.}

\prose{เอกํเสนาหํ อานนฺท กรณียํ วทามิ กายสุจริตํ วจีสุจริตํ มโน\-สุจริตนฺติ. ยมิทํ ภนฺเต ภควตา เอกํเสน กรณียํ อกฺขาตํ กายสุจริตํ วจีสุจริตํ มโนสุจริตํ, ตสฺมิํ กรณีเย กยิรมาเน โก อานิสํโส ปาฏิกงฺโขติ. ยมิทํ อานนฺท มยา เอกํเสน กรณียํ อกฺขาตํ กายสุจริตํ วจีสุจริตํ มโนสุจริตํ, ตสฺมิํ กรณีเย กยิรมาเน อยํ อานิสํโส ปาฏิกงฺโข\csromandash{}อตฺตาปิ อตฺตานํ น อุปวทติ, อนุวิจฺจ วิญฺญู ปสํสนฺติ, กลฺยาโณ กิตฺติสทฺโท อพฺภุคฺคจฺฉติ, อสมฺมูโฬฺห กาลํ กโรติ, กายสฺส เภทา ปรํ มรณา สุคติํ สคฺคํ โลกํ อุปปชฺชติ. ยมิทํ อานนฺท มยา เอกํเสน กรณียํ อกฺขาตํ กายสุจริตํ วจีสุจริตํ มโนสุจริตํ, ตสฺมิํ กรณีเย กยิรมาเน อยํ อานิสํโส ปาฏิกงฺโขติ. (8)}

\proseitem{19}{อกุสลํ ภิกฺขเว ปชหถ, สกฺกา ภิกฺขเว อกุสลํ ปชหิตุํ. โน เจทํ\csromansharedfootnote{b266845d257b25b5}{โน เจตํ (สฺยา, กํ, อิ, ก) สํ 2. 25 ปิฏฺเฐ ปสฺสิตพฺพํ.} ภิกฺขเว สกฺกา อภวิสฺส อกุสลํ ปชหิตุํ, นาหํ เอวํ วเทยฺยํ “อกุสลํ ภิกฺขเว ปชหถา”ติ. ยสฺมา จ โข ภิกฺขเว \csromanfolio{60}สกฺกา อกุสลํ ปชหิตุํ, ตสฺมาหํ เอวํ วทามิ “อกุสลํ ภิกฺขเว ปชหถา”ติ. อกุสลญฺจ หิทํ ภิกฺขเว\csromansharedfootnote{5d6e56b586b81391}{อกุสลํ ภิกฺขเว (ก)} ปหีนํ อหิตาย ทุกฺขาย สํวตฺเตยฺย, นาหํ เอวํ วเทยฺยํ “อกุสลํ ภิกฺขเว ปชหถา”ติ. ยสฺมา จ โข ภิกฺขเว อกุสลํ ปหีนํ หิตาย สุขาย สํวตฺตติ, ตสฺมาหํ เอวํ วทามิ “อกุสลํ ภิกฺขเว ปชหถา”ติ.}

\prose{กุสลํ ภิกฺขเว ภาเวถ, สกฺกา ภิกฺขเว กุสลํ ภาเวตุํ. โน เจทํ ภิกฺขเว สกฺกา อภวิสฺส กุสลํ ภาเวตุํ, นาหํ เอวํ วเทยฺยํ “กุสลํ ภิกฺขเว ภาเวถา”ติ. ยสฺมา จ โข ภิกฺขเว สกฺกา กุสลํ ภาเวตุํ, ตสฺมาหํ เอวํ วทามิ “กุสลํ ภิกฺขเว ภาเวถา”ติ. กุสลญฺจ หิทํ ภิกฺขเว ภาวิตํ อหิตาย ทุกฺขาย สํวตฺเตยฺย, นาหํ เอวํ วเทยฺยํ “กุสลํ ภิกฺขเว ภาเวถา”ติ. ยสฺมา จ โข ภิกฺขเว กุสลํ ภาวิตํ หิตาย สุขาย สํวตฺตติ, ตสฺมาหํ เอวํ วทามิ “กุสลํ ภิกฺขเว ภาเวถา”ติ. (9)}

\proseitem{20}{เทฺวเม ภิกฺขเว ธมฺมา สทฺธมฺมสฺส สมฺโมสาย อนฺตรธานาย สํวตฺตนฺติ. กตเม เทฺว? ทุนฺนิกฺขิตฺตญฺจ ปทพฺยญฺชนํ, อตฺโถ จ ทุนฺนีโต. ทุนฺนิกฺขิตฺตสฺส ภิกฺขเว ปทพฺยญฺชนสฺส อตฺโถปิ ทุนฺนโย โหติ. อิเม โข ภิกฺขเว เทฺว ธมฺมา สทฺธมฺมสฺส สมฺโมสาย อนฺตรธานาย สํวตฺตนฺตีติ. (10)}

\proseitemwithclosermidruled{21}{เทฺวเม ภิกฺขเว ธมฺมา สทฺธมฺมสฺส ฐิติยา อสมฺโมสาย อนนฺตรธานาย สํวตฺตนฺติ. กตเม เทฺว? สุนิกฺขิตฺตญฺจ ปทพฺยญฺชนํ, อตฺโถ จ สุนีโต. สุนิกฺขิตฺตสฺส ภิกฺขเว ปทพฺยญฺชนสฺส อตฺโถปิ สุนโย โหติ. อิเม โข ภิกฺขเว เทฺว ธมฺมา สทฺธมฺมสฺส ฐิติยา อสมฺโมสาย อนนฺตรธานาย สํวตฺตนฺตีติ. (11)}{อธิกรณ\-วคฺโค ทุติโย.}
\chapterhead{3. \textbf{พาลวคฺค}}
\proseitem{22}{เทฺวเม ภิกฺขเว พาลา. กตเม เทฺว? โย จ อจฺจยํ อจฺจยโต น ปสฺสติ, โย จ อจฺจยํ เทเสนฺตสฺส ยถาธมฺมํ นปฺปฏิคฺคณฺหาติ. อิเม โข ภิกฺขเว เทฺว พาลาติ.}

\chapterheadpageread{3. พาลวคฺค}
\prose{\csromanfolio{61}เทฺวเม ภิกฺขเว ปณฺฑิตา. กตเม เทฺว? โย จ อจฺจยํ อจฺจยโต ปสฺสติ, โย จ อจฺจยํ เทเสนฺตสฺส ยถาธมฺมํ ปฏิคฺคณฺหาติ. อิเม โข ภิกฺขเว เทฺว ปณฺฑิตาติ. (1)}

\proseitem{23}{เทฺวเม ภิกฺขเว ตถาคตํ อพฺภาจิกฺขนฺติ. กตเม เทฺว? ทุฏฺโฐ วา โทสนฺตโร, สทฺโธ วา ทุคฺคหิเตน\csromansharedfootnote{7187974d876204d2}{ทุคฺคหีเตน (สี)}. อิเม โข ภิกฺขเว เทฺว ตถาคตํ อพฺภาจิกฺขนฺตีติ. (2)}

\proseitem{24}{เทฺวเม ภิกฺขเว ตถาคตํ อพฺภาจิกฺขนฺติ. กตเม เทฺว? โย จ อภาสิตํ อลปิตํ ตถาคเตน “ภาสิตํ ลปิตํ ตถาคเตนา”ติ ทีเปติ, โย จ ภาสิตํ ลปิตํ ตถาคเตน “อภาสิตํ อลปิตํ ตถาคเตนา”ติ ทีเปติ. อิเม โข ภิกฺขเว เทฺว ตถาคตํ อพฺภาจิกฺขนฺตีติ.}

\prose{เทฺวเม ภิกฺขเว ตถาคตํ นาพฺภาจิกฺขนฺติ. กตเม เทฺว? โย จ อภาสิตํ อลปิตํ ตถาคเตน “อภาสิตํ อลปิตํ ตถาคเตนา”ติ ทีเปติ, โย จ ภาสิตํ ลปิตํ ตถาคเตน “ภาสิตํ ลปิตํ ตถาคเตนา”ติ ทีเปติ. อิเม โข ภิกฺขเว เทฺว ตถาคตํ นาพฺภาจิกฺขนฺตีติ. (3)}

\proseitem{25}{เทฺวเม ภิกฺขเว ตถาคตํ อพฺภาจิกฺขนฺติ. กตเม เทฺว? โย จ เนยฺยตฺถํ สุตฺตนฺตํ “นีตตฺโถ สุตฺตนฺโต”ติ ทีเปติ, โย จ นีตตฺถํ สุตฺตนฺตํ “เนยฺยตฺโถ สุตฺตนฺโต”ติ ทีเปติ. อิเม โข ภิกฺขเว เทฺว ตถาคตํ อพฺภาจิกฺขนฺตีติ. (4)}

\proseitem{26}{เทฺวเม ภิกฺขเว ตถาคตํ นาพฺภาจิกฺขนฺติ. กตเม เทฺว? โย จ เนยฺยตฺถํ สุตฺตนฺตํ “เนยฺยตฺโถ สุตฺตนฺโต”ติ ทีเปติ, โย จ นีตตฺถํ สุตฺตนฺตํ “นีตตฺโถ สุตฺตนฺโต”ติ ทีเปติ. อิเม โข ภิกฺขเว เทฺว ตถาคตํ นาพฺภาจิกฺขนฺตีติ. (5)}

\proseitem{27}{ปฏิจฺฉนฺน\-กมฺมนฺตสฺส ภิกฺขเว ทฺวินฺนํ คตีนํ อญฺญตรา คติ ปาฏิกงฺขา นิรโย วา ติรจฺฉานโยนิ วาติ. อปฺปฏิจฺฉนฺน\-กมฺมนฺตสฺส ภิกฺขเว ทฺวินฺนํ คตีนํ อญฺญตรา คติ ปาฏิกงฺขา เทวา วา มนุสฺสา วาติ. (6)}

\csromanpalirecto
\gambhira{ทุกกนิปาตปาฬิ}
\proseitem{28}{\csromanfolio{62}มิจฺฉา\-ทิฏฺฐิกสฺส ภิกฺขเว ทฺวินฺนํ คตีนํ อญฺญตรา คติ ปาฏิกงฺขา นิรโย วา ติรจฺฉานโยนิ วาติ. (7)}

\proseitem{29}{สมฺมา\-ทิฏฺฐิกสฺส ภิกฺขเว ทฺวินฺนํ คตีนํ อญฺญตรา คติ ปาฏิกงฺขา เทวา วา มนุสฺสา วาติ. (8)}

\proseitem{30}{ทุสฺสีลสฺส ภิกฺขเว เทฺว ปฏิคฺคาหา นิรโย วา ติรจฺฉานโยนิ วา, สีลวโต ภิกฺขเว เทฺว ปฏิคฺคาหา เทวา วา มนุสฺสา วาติ\csromansharedfootnote{5237893e1869bae7}{เทโว วา มนุสฺโส วาติ (ก)}. (9)}

\proseitem{31}{ทฺวาหํ ภิกฺขเว อตฺถวเส สมฺปสฺสมาโน อรญฺญ\-วนปตฺถานิ\csromansharedfootnote{cb9f8cf1cdfe365f}{อรญฺเญ ปวนปตฺถานิ (สี, อิ)} ปนฺตานิ เสนาสนานิ ปฏิเสวามิ. กตเม เทฺว? อตฺตโน จ ทิฏฺฐ\-ธมฺม\-สุข\-วิหารํ สมฺปสฺสมาโน, ปจฺฉิมญฺจ ชนตํ อนุกมฺปมาโน. อิเม โข อหํ ภิกฺขเว เทฺว อตฺถวเส สมฺปสฺสมาโน อรญฺญ\-วนปตฺถานิ ปนฺตานิ เสนาสนานิ ปฏิเสวามีติ. (10)}

\proseitemwithclosermidruled{32}{เทฺวเม ภิกฺขเว ธมฺมา วิชฺชาภาคิยา. กตเม เทฺว? สมโถ จ วิปสฺสนา จ. สมโถ ภิกฺขเว ภาวิโต กมตฺถ\csromansharedfootnote{5cd5f34ecb138ca9}{กิมตฺถ (สฺยา, กํ), กตมตฺถ (ก)} มนุโภติ? จิตฺตํ ภาวียติ. จิตฺตํ ภาวิตํ กมตฺถม\-นุโภติ? โย ราโค, โส ปหียติ. วิปสฺสนา ภิกฺขเว ภาวิตา กมตฺถม\-นุโภติ? ปญฺญา ภาวียติ. ปญฺญา ภาวิตา กมตฺถม\-นุโภติ? ยา อวิชฺชา, สา ปหียติ. ราคุ\-ปกฺกิลิฏฺฐํ วา ภิกฺขเว จิตฺตํ น วิมุจฺจติ, อวิชฺชุ\-ปกฺกิลิฏฺฐา วา ปญฺญา น ภาวียติ. อิติ โข ภิกฺขเว ราควิราคา เจโตวิมุตฺติ, อวิชฺชาวิราคา ปญฺญาวิมุตฺตีติ. (11)}{พาลวคฺโค ตติโย.}
\chapterhead{4. \textbf{สมจิตฺตวคฺค}}
\proseitem{33}{อสปฺปุริส\-ภูมิญฺจ โว ภิกฺขเว เทเสสฺสามิ สปฺปุริสภูมิญฺจ, ตํ สุณาถ สาธุกํ มนสิ กโรถ, ภาสิสฺสามีติ. “เอวํ ภนฺเต”ติ โข เต ภิกฺขู ภควโต ปจฺจสฺโสสุํ. ภควา เอตทโวจ\csromandash{}}

\chapterheadpageread{4. สมจิตฺตวคฺค}
\prose{\csromanfolio{63}กตมา จ ภิกฺขเว อสปฺปุริสภูมิ? อสปฺปุริโส ภิกฺขเว อกตญฺญู โหติ อกตเวที. อสพฺภิ เหตํ ภิกฺขเว อุปญฺญาตํ, ยทิทํ อกตญฺญุตา อกตเวทิตา. เกวลา เอสา ภิกฺขเว อสปฺปุริสภูมิ, ยทิทํ อกตญฺญุตา อกตเวทิตา. สปฺปุริโส จ โข ภิกฺขเว กตญฺญู โหติ กตเวที. สพฺภิ เหตํ ภิกฺขเว อุปญฺญาตํ, ยทิทํ กตญฺญุตา กตเวทิตา. เกวลา เอสา ภิกฺขเว สปฺปุริสภูมิ, ยทิทํ กตญฺญุตา กตเวทิตาติ. (1)}

\proseitem{34}{ทฺวินฺนาหํ ภิกฺขเว น สุปฺปติการํ วทามิ. กตเมสํ ทฺวินฺนํ? มาตุ จ ปิตุ จ. เอเกน ภิกฺขเว อํเสน มาตรํ ปริหเรยฺย, เอเกน อํเสน ปิตรํ ปริหเรยฺย วสฺสสตายุโก วสฺสสตชีวี, โส จ เนสํ อุจฺฉาทน\-ปริมทฺทน\-นฺหาปน\-สมฺพาหเนน, เต จ ตตฺเถว มุตฺตกรีสํ จเชยฺยุํ, น เตฺวว ภิกฺขเว มาตาปิตูนํ กตํ วา โหติ ปฏิกตํ วา. อิมิสฺสา จ ภิกฺขเว มหาปถวิยา ปหูตรตฺตรตนาย\csromansharedfootnote{67f043429dc00644}{ปหูตสตฺตรตนาย (สี, สฺยา, กํ, อิ) ติกนิปาเต มหาวคฺเค ทสมสุตฺตฏีกายํ ทสฺสิตปาฬิยา สเมติ.} มาตาปิตโร อิสฺสราธิปจฺเจ รชฺเช ปติฏฺฐาเปยฺย, น เตฺวว ภิกฺขเว มาตาปิตูนํ กตํ วา โหติ ปฏิกตํ วา. ตํ กิสฺส เหตุ? พหุการา\csromansharedfootnote{2979d44c5af72058}{พหูปการา (ก)} ภิกฺขเว มาตาปิตโร ปุตฺตานํ อาปาทกา โปสกา, อิมสฺส โลกสฺส ทสฺเสตาโร. โย จ โข ภิกฺขเว มาตาปิตโร อสฺสทฺเธ สทฺธา\-สมฺปทาย สมาทเปติ นิเวเสติ ปติฏฺฐาเปติ, ทุสฺสีเล สีลสมฺปทาย สมาทเปติ นิเวเสติ ปติฏฺฐาเปติ, มจฺฉรี จาคสมฺปทาย สมาทเปติ นิเวเสติ ปติฏฺฐาเปติ, ทุปฺปญฺเญ ปญฺญาสมฺปทาย สมาทเปติ นิเวเสติ ปติฏฺฐาเปติ. เอตฺตาวตา โข ภิกฺขเว มาตาปิตูนํ กตญฺจ โหติ ปฏิกตญฺจาติ\csromansharedfootnote{a211bae864ae2492}{ปฏิกตญฺจ อติกตญฺจาติ (สี, อิ)}. (2)}

\proseitem{35}{อถ โข อญฺญตโร พฺราหฺมโณ เยน ภควา เตนุปสงฺกมิ, อุปสงฺกมิตฺวา ภควตา สทฺธิํ สมฺโมทิ, สมฺโมทนียํ กถํ ฯเปฯ เอกมนฺตํ นิสินฺโน โข โส พฺราหฺมโณ ภควนฺตํ เอตทโวจ “กิํ วาที ภวํ โคตโม, กิมกฺขายี”ติ. กิริยวาที จาหํ พฺราหฺมณ อกิริยวาที จาติ. ยถากถํ ปน ภวํ โคตโม กิริยวาที จ อกิริยวาที จาติ.}

\prose{อกิริยํ โข อหํ พฺราหฺมณ วทามิ กาย\-ทุจฺจริตสฺส วจีทุจฺจริตสฺส มโน\-ทุจฺจริตสฺส, อเนกวิหิตานํ ปาปกานํ อกุสลานํ ธมฺมานํ \csromanfolio{64}อกิริยํ วทามิ. กิริยญฺจ โข อหํ พฺราหฺมณ วทามิ กาย\-สุจริตสฺส วจีสุจริตสฺส มโน\-สุจริตสฺส, อนิกวิหิตานํ กุสลานํ ธมฺมานํ กิริยํ วทามิ. เอวํ โข อหํ พฺราหฺมณ กิริยวาที จ อกิริยวาที จาติ.}

\prose{อภิกฺกนฺตํ โภ โคตม ฯเปฯ อุปาสกํ มํ ภวํ โคตโม ธาเรตุ อชฺชตคฺเค ปาณุเปตํ สรณํ คตนฺติ. (3)}

\proseitem{36}{อถ โข อนาถปิณฺฑิโก คหปติ เยน ภควา เตนุปสงฺกมิ, อุปสงฺกมิตฺวา ภควนฺตํ อภิวาเทตฺวา เอกมนฺตํ นิสีทิ, เอกมนฺตํ นิสินฺโน โข อนาถปิณฺฑิโก คหปติ ภควนฺตํ เอตทโวจ “กติ นุ โข ภนฺเต โลเก ทกฺขิเณยฺยา, กตฺถ จ ทานํ ทาตพฺพนฺ”ติ. เทฺว โข คหปติ โลเก ทกฺขิเณยฺยา เสโข จ อเสโข จ. อิเม โข คหปติ เทฺว โลเก ทกฺขิเณยฺยา, เอตฺถ จ ทานํ ทาตพฺพนฺติ.}

\prose{อิทมโวจ ภควา, อิทํ วตฺวาน\csromansharedfootnote{f786312130b697a4}{วตฺวา (สี, อิ) เอวมีทิเสสุ ฐาเนสุ.} สุคโต อถาปรํ เอตทโวจ สตฺถา\csromandash{}}

\setlength{\csromangathapagewidth}{0pt}
\setlength{\csromangathawakwidth}{0pt}
\csromangathameasuregroup{}{“เสโข อเสโข จ อิมสฺมิํ โลเก, \\ อาหุเนยฺยา ยชมานานํ โหนฺติ. \\ เต อุชฺชุภูตา\textsuperscript{1} กาเยน, \\ วาจาย อุท เจตสา.}
\csromangathameasurewak{“เสโข อเสโข จ อิมสฺมิํ โลเก, \\ อาหุเนยฺยา ยชมานานํ โหนฺติ. \\ เต อุชฺชุภูตา\textsuperscript{1} กาเยน, \\ วาจาย อุท เจตสา.}
\setlength{\csromangathaleftcol}{0pt}
\csromangathagroup{\csromangathastanza{\csromangathawak{“เสโข อเสโข จ อิมสฺมิํ โลเก, \\ อาหุเนยฺยา ยชมานานํ โหนฺติ. \\ เต อุชฺชุภูตา\csromansharedfootnote{5573652047f7814f}{อุชุภูตา (สฺยา, กํ, ก)} กาเยน, \\ วาจาย อุท เจตสา.}}}

\prose{เขตฺตํ ตํ ยชมานานํ, เอตฺถ ทินฺนํ มหปฺผลนฺ”ติ. (4)}

\proseitem{37}{เอวํ เม สุตํ\csromandash{}เอกํ สมยํ ภควา สาวตฺถิยํ วิหรติ เชตวเน อนาถ\-ปิณฺฑิกสฺส อาราเม. เตน โข ปน สมเยน อายสฺมา สาริปุตฺโต สาวตฺถิยํ วิหรติ ปุพฺพาราเม มิคาร\-มาตุ\-ปาสาเท. ตตฺร โข อายสฺมา สาริปุตฺโต ภิกฺขู อามนฺเตสิ “อาวุโส ภิกฺขเว”ติ. “อาวุโส”ติ โข เต ภิกฺขู อายสฺมโต สาริปุตฺตสฺส ปจฺจสฺโสสุํ. อายสฺมา สาริปุตฺโต เอตทโวจ “อชฺฌตฺตสํโยชนญฺจ อาวุโส ปุคฺคลํ เทเสสฺสามิ พหิทฺธาสํโยชนญฺจ, ตํ สุณาถ สาธุกํ มนสิ กโรถ, ภาสิสฺสามี”ติ. “เอวมาวุโส”ติ โข เต ภิกฺขู อายสฺมโต สาริปุตฺตสฺส ปจฺจสฺโสสุํ. อายสฺมา สาริปุตฺโต เอตทโวจ “กตโม จาวุโส อชฺฌตฺต\-สํโยชโน ปุคฺคโล? อิธาวุโส ภิกฺขุ สีลวา โหติ ปาติโมกฺข\-สํวร\-สํวุโต วิหรติ \csromanfolio{65}อาจารโคจร\-สมฺปนฺโน อณุมตฺเตสุ วชฺเชสุ ภยทสฺสาวี สมาทาย สิกฺขติ สิกฺขาปเทสุ โส กายสฺส เภทา ปรํ มรณา อญฺญตรํ เทวนิกายํ อุปปชฺชติ, โส ตโต จุโต อาคามี โหติ อาคนฺตา อิตฺถตฺตํ. อยํ วุจฺจติ อาวุโส อชฺฌตฺต\-สํโยชโน ปุคฺคโล อาคามี โหติ อาคนฺตา อิตฺถตฺตํ.}

\prose{กตโม จาวุโส พหิทฺธา\-สํโยชโน ปุคฺคโล? อิธาวุโส ภิกฺขุ สีลวา โหติ ปาติโมกฺข\-สํวร\-สํวุโต วิหรติ อาจารโคจร\-สมฺปนฺโน อณุมตฺเตสุ วชฺเชสุ ภยทสฺสาวี สมาทาย สิกฺขติ สิกฺขาปเทสุ โส อญฺญตรํ สนฺตํ เจโตวิมุตฺติํ อุปสมฺปชฺช วิหรติ, โส กายสฺส เภทา ปรํ มรณา อญฺญตรํ เทวนิกายํ อุปปชฺชติ, โส ตโต จุโต อนาคามี โหติ อนาคนฺตา อิตฺถตฺตํ. อยํ วุจฺจตาวุโส พหิทฺธา\-สํโยชโน ปุคฺคโล อนาคามี โหติ อนาคนฺตา อิตฺถตฺตํ.}

\prose{ปุน จปรํ อาวุโส ภิกฺขุ สีลวา โหติ ฯเปฯ สมาทาย สิกฺขติ สิกฺขาปเทสุ โส กามานํเยว นิพฺพิทาย วิราคาย นิโรธาย ปฏิปนฺโน โหติ. โส ภวานํเยว นิพฺพิทาย วิราคาย นิโรธาย ปฏิปนฺโน โหติ. โส ตณฺหากฺขยาย ปฏิปนฺโน โหติ. โส โลภกฺขยาย ปฏิปนฺโน โหติ, โส กายสฺส เภทา ปรํ มรณา อญฺญตรํ เทวนิกายํ อุปปชฺชติ, โส ตโต จุโต อนาคามี โหติ อนาคนฺตา อิตฺถตฺตํ. อยํ วุจฺจตาวุโส พหิทฺธา\-สํโยชโน ปุคฺคโล อนาคามี โหติ อนาคนฺตา อิตฺถตฺตนฺ”ติ.}

\prose{อถ โข สมฺพหุลา สมจิตฺตา เทวตา เยน ภควา เตนุ\-ปสงฺกมิํสุ, อุปสงฺกมิตฺวา ภควนฺตํ อภิวาเทตฺวา เอกมนฺตํ อฏฺฐํสุ, เอกมนฺตํ ฐิตา โข ตา เทวตา ภควนฺตํ เอตทโวจุํ “เอโส ภนฺเต อายสฺมา สาริปุตฺโต ปุพฺพาราเม มิคาร\-มาตุ\-ปาสาเท ภิกฺขูนํ อชฺฌตฺตสํโยชนญฺจ ปุคฺคลํ เทเสติ พหิทฺธาสํโยชนญฺจ, หฏฺฐา ภนฺเต ปริสา, สาธุ ภนฺเต ภควา เยนายสฺมา สาริปุตฺโต เตนุ\-ปสงฺกมตุ อนุกมฺปํ อุปาทายา”ติ. อธิวาเสสิ ภควา ตุณฺหีภาเวน. อถ โข ภควา เสยฺยถาปิ นาม พลวา ปุริโส สมิญฺชิตํ\csromansharedfootnote{7405ea38caf1a994}{สมฺมิญฺชิตํ (สี, สฺยา, กํ, อิ)} วา พาหํ ปสาเรยฺย, ปสาริตํ วา พาหํ \csromanfolio{66}สมิญฺเชยฺย. เอวเมวํ เชตวเน อนฺตรหิโต ปุพฺพาราเม มิคาร\-มาตุ\-ปาสาเท อายสฺมโต สาริปุตฺตสฺส สมฺมุเข ปาตุรโหสิ. นิสีทิ ภควา ปญฺญตฺเต อาสเน. อายสฺมาปิ โข สาริปุตฺโต ภควนฺตํ อภิวาเทตฺวา เอกมนฺตํ นิสีทิ, เอกมนฺตํ นิสินฺนํ โข อายสฺมนฺตํ สาริปุตฺตํ ภควา เอตทโวจ\csromandash{}}

\proseitem{4}{อิธ สาริปุตฺต สมฺพหุลา สมจิตฺตา เทวตา เยนาหํ เตนุ\-ปสงฺกมิํสุ, อุปสงฺกมิตฺวา มํ อภิวาเทตฺวา เอกมนฺตํ อฏฺฐํสุ, เอกมนฺตํ ฐิตา โข สาริปุตฺต ตา เทวตา มํ เอตทโวจุํ “เอโส ภนฺเต อายสฺมา สาริปุตฺโต ปุพฺพาราเม มิคาร\-มาตุ\-ปาสาเท ภิกฺขูนํ อชฺฌตฺตสํโยชนญฺจ ปุคฺคลํ เทเสติ พหิทฺธาสํโยชนญฺจ, หฏฺฐา ภนฺเต ปริสา, สาธุ ภนฺเต ภควา เยน อายสฺมา สาริปุตฺโต เตนุ\-ปสงฺกมตุ อนุกมฺปํ อุปาทายา”ติ. ตา โข ปน สาริปุตฺต เทวตา ทสปิ หุตฺวา วีสมฺปิ หุตฺวา ติํสมฺปิ หุตฺวา จตฺตาลีสมฺปิ หุตฺวา ปญฺญาสมฺปิ หุตฺวา สฏฺฐิปิ หุตฺวา อารคฺค\-โกฏิ\-นิตุทน\-มตฺเตปิ ติฏฺฐนฺติ, น จ อญฺญมญฺญํ พฺยาพาเธนฺติ\csromansharedfootnote{c1f2d148b52efc58}{พฺยาพาเธนฺตีติ (สพฺพตฺถ)}. สิยา โข ปน\csromansharedfootnote{90166a53e07c2aba}{ปน เต (สี, สฺยา, กํ, อิ)} สาริปุตฺต เอวมสฺส “ตตฺถ นูน ตาสํ เทวตานํ ตถาจิตฺตํ ภาวิตํ, เยน ตา เทวตา ทสปิ หุตฺวา วีสมฺปิ หุตฺวา ติํสมฺปิ หุตฺวา จตฺตาลีสมฺปิ หุตฺวา ปญฺญาสมฺปิ หุตฺวา สฏฺฐิปิ หุตฺวา อารคฺค\-โกฏิ\-นิตุทน\-มตฺเตปิ ติฏฺฐนฺติ, น จ อญฺญมญฺญํ พฺยาพาเธนฺตี”ติ. น โข ปเนตํ สาริปุตฺต เอวํ ทฏฺฐพฺพํ. อิเธว โข สาริปุตฺต ตาสํ เทวตานํ ตถาจิตฺตํ ภาวิตํ, เยน ตา เทวตา ทสปิ หุตฺวา ฯเปฯ น จ อญฺญมญฺญํ พฺยาพาเธนฺติ. ตสฺมาติห สาริปุตฺต เอวํ สิกฺขิตพฺพํ “สนฺตินฺทฺริยา ภวิสฺสาม สนฺตมานสา”ติ. เอวํ หิ โว สาริปุตฺต สิกฺขิตพฺพํ. สนฺตินฺทฺริยานํ หิ โว สาริปุตฺต สนฺตมานสานํ สนฺตํเยว กายกมฺมํ ภวิสฺสติ, สนฺตํ วจีกมฺมํ, สนฺตํ มโนกมฺมํ. สนฺตํเยว อุปหารํ อุปหริสฺสาม สพฺรหฺมจารีสูติ, เอวํ หิ โว สาริปุตฺต สิกฺขิตพฺพํ. อนสฺสุํ โข สาริปุตฺต อญฺญติตฺถิยา ปริพฺพาชกา, เย อิมํ ธมฺม\-ปริยายํ นาสฺโสสุนฺติ. (5)}

\proseitem{38}{เอวํ เม สุตํ\csromandash{}เอกํ สมยํ อายสฺมา มหากจฺจาโน วรณายํ วิหรติ ภทฺทสาริตีเร\csromansharedfootnote{0f8398fdccc2bb27}{กทฺทมทหตีเร (สี, สฺยา, กํ, อิ)}. อถ โข อารามทณฺโฑ พฺราหฺมโณ เยนายสฺมา มหากจฺจาโน เตนุปสงฺกมิ, อุปสงฺกมิตฺวา \csromanfolio{67}อายสฺมตา มหากจฺจาเนน สทฺธิํ สมฺโมทิ, สมฺโมทนียํ กถํ สารณียํ วีติสาเรตฺวา เอกมนฺตํ นิสีทิ, เอกมนฺตํ นิสินฺโน โข อารามทณฺโฑ พฺราหฺมโณ อายสฺมนฺตํ มหากจฺจานํ เอตทโวจ “โก นุ โข โภ กจฺจาน เหตุ โก ปจฺจโย, เยน ขตฺติยาปิ ขตฺติเยหิ วิวทนฺติ, พฺราหฺมณาปิ พฺราหฺมเณหิ วิวทนฺติ, คหปติกาปิ คหปติเกหิ วิวทนฺตี”ติ. กามราคา\-ภินิเวส\-วินิพนฺธ\csromansharedfootnote{43523a0b416050d0}{กามราควินิเวสวินิพทฺธ (สี, สฺยา, กํ, อิ)} ปลิเคธ\-ปริยุฏฺฐาน\-ชฺโฌสาน\-เหตุ โข พฺราหฺมณ ขตฺติยาปิ ขตฺติเยหิ วิวทนฺติ, พฺราหฺมณาปิ พฺราหฺมเณหิ วิวทนฺติ, คหปติกาปิ คหปติเกหิ วิวทนฺตีติ.}

\prose{โก ปน โภ กจฺจาน เหตุ โก ปจฺจโย, เยน สมณาปิ สมเณหิ วิวทนฺตีติ. ทิฏฺฐิราคา\-ภินิเวส\-วินิพนฺธ\-ปลิเคธ\-ปริยุฏฺฐาน\-ชฺโฌสาน\-เหตุ โข พฺราหฺมณ สมณาปิ สมเณหิ วิวทนฺตีติ.}

\prose{อตฺถิ ปน โภ กจฺจาน โกจิ โลกสฺมิํ โย อิมํ เจว กามราคา\-ภินิเวส\-วินิพนฺธ\-ปลิเคธ\-ปริยุฏฺฐาน\-ชฺโฌสานํ สมติกฺกนฺโต, อิมญฺจ ทิฏฺฐิราคา\-ภินิเวส\-วินิพนฺธ\-ปลิเคธ\-ปริยุฏฺฐาน\-ชฺโฌสานํ สมติกฺกนฺโตติ. อตฺถิ พฺราหฺมณ โลกสฺมิํ โย อิมํ เจว กามราคา\-ภินิเวส\-วินิพนฺธ\-ปลิเคธ\-ปริยุฏฺฐาน\-ชฺโฌสานํ สมติกฺกนฺโต, อิมญฺจ ทิฏฺฐิราคา\-ภินิเวส\-วินิพนฺธ\-ปลิเคธ\-ปริยุฏฺฐาน\-ชฺโฌสานํ สมติกฺกนฺโตติ.}

\prose{โก ปน โส โภ กจฺจาน โลกสฺมิํ, โย อิมญฺเจว กามราคา\-ภินิเวส\-วินิพนฺธ\-ปลิเคธ\-ปริยุฏฺฐาน\-ชฺโฌสานํ สมติกฺกนฺโต, อิมญฺจ ทิฏฺฐิราคา\-ภินิเวส\-วินิพนฺธ\-ปลิเคธ\-ปริยุฏฺฐาน\-ชฺโฌสานํ สมติกฺกนฺโตติ. อตฺถิ พฺราหฺมณ ปุรตฺถิเมสุ ชนปเทสุ สาวตฺถี นาม นครํ, ตตฺถ โส ภควา เอตรหิ วิหรติ อรหํ สมฺมาสมฺพุทฺโธ. โส หิ พฺราหฺมณ ภควา อิมญฺเจว กามราคา\-ภินิเวส\-วินิพนฺธ\-ปลิเคธ\-ปริยุฏฺฐาน\-ชฺโฌสานํ สมติกฺกนฺโต, อิมญฺจ ทิฏฺฐิราคา\-ภินิเวส\-วินิพนฺธ\-ปลิเคธ\-ปริยุฏฺฐาน\-ชฺโฌสานํ สมติกฺกนฺโตติ.}

\prose{เอวํ วุตฺเต อารามทณฺโฑ พฺราหฺมโณ อุฏฺฐายาสนา เอกํสํ อุตฺตราสงฺคํ กริตฺวา ทกฺขิณํ ชาณุมณฺฑลํ ปถวิยํ นิหนฺตฺวา เยน ภควา เตนญฺชลิํ ปณาเมตฺวา ติกฺขตฺตุํ อุทานํ อุทาเนสิ\csromandash{}}

\prose{\csromanfolio{68}“นโม ตสฺส ภควโต อรหโต สมฺมา\-สมฺพุทฺธสฺส, นโม ตสฺส ภควโต อรหโต สมฺมา\-สมฺพุทฺธสฺส, นโม ตสฺส ภควโต อรหโต สมฺมา\-สมฺพุทฺธสฺส, โย หิ โส ภควา อิมญฺเจว กามราคา\-ภินิเวส\-วินิพนฺธ\-ปลิเคธ\-ปริยุฏฺฐาน\-ชฺโฌสานํ สมติกฺกนฺโต, อิมญฺจ ทิฏฺฐิราคา\-ภินิเวส\-วินิพนฺธ\-ปลิเคธ\-ปริยุฏฺฐาน\-ชฺโฌสานํ สมติกฺกนฺโต”ติ.}

\prose{อภิกฺกนฺตํ โภ กจฺจาน, อภิกฺกนฺตํ โภ กจฺจาน, เสยฺยถาปิ โภ กจฺจาน นิกฺกุชฺชิตํ วา อุกฺกุชฺเชยฺย, ปฏิจฺฉนฺนํ วา วิวเรยฺย, มูฬฺหสฺส วา มคฺคํ อาจิกฺเขยฺย, อนฺธกาเร วา เตลปชฺโชตํ ธาเรยฺย “จกฺขุมนฺโต รูปานิ ทกฺขนฺตี”ติ. เอวเมวํ โภตา กจฺจาเนน อเนก\-ปริยาเยน ธมฺโม ปกาสิโต. เอสาหํ โภ กจฺจาน ตํ ภวนฺตํ โคตมํ สรณํ คจฺฉามิ ธมฺมญฺจ ภิกฺขุสํฆญฺจ, อุปาสกํ มํ ภวํ กจฺจาโน ธาเรตุ อชฺชตคฺเค ปาณุเปตํ สรณํ คตนฺติ. (6)}

\proseitem{39}{เอกํ สมยํ อายสฺมา มหากจฺจาโน มธุรายํ วิหรติ คุนฺทาวเน. อถ โข กนฺทรายโน\csromansharedfootnote{433a559bce1e1f43}{กณฺฑรายโน (สี, สฺยา, กํ, อิ)} พฺราหฺมโณ เยนายสฺมา มหากจฺจาโน เตนุปสงฺกมิ, อุปสงฺกมิตฺวา อายสฺมตา มหากจฺจาเนน สทฺธิํ ฯเปฯ เอกมนฺตํ นิสินฺโน โข กนฺทรายโน พฺราหฺมโณ อายสฺมนฺตํ มหากจฺจานํ เอตทโวจ “สุตํ เมตํ โภ กจฺจาน น สมโณ กจฺจาโน พฺราหฺมเณ ชิณฺเณ วุทฺเธ มหลฺลเก อทฺธคเต วโยอนุปฺปตฺเต อภิวาเทติ วา ปจฺจุฏฺเฐติ วา อาสเนน วา นิมนฺเตตี”ติ. ตยิทํ โภ กจฺจาน ตเถว, น หิ ภวํ กจฺจาโน พฺราหฺมเณ ชิณฺเณ วุทฺเธ มหลฺลเก อทฺธคเต วโยอนุปฺปตฺเต อภิวาเทติ วา ปจฺจุฏฺเฐติ วา อาสเนน วา นิมนฺเตติ, ตยิทํ โภ กจฺจาน น สมฺปนฺนเมวาติ.}

\prose{อตฺถิ พฺราหฺมณ เตน ภควตา ชานตา ปสฺสตา อรหตา สมฺมา\-สมฺพุทฺเธน วุทฺธภูมิ จ อกฺขาตา ทหรภูมิ จ. วุทฺโธ เจปิ พฺราหฺมณ โหติ อาสีติโก วา นาวุติโก วา วสฺสสติโก วา ชาติยา, โส จ กาเม ปริภุญฺชติ, กามม\-ชฺฌาวสติ, กามปริฬาเหน ปริฑยฺหติ, กามวิตกฺเกหิ ขชฺชติ, กาม\-ปริเยสนาย อุสฺสุโก. อถ โข โส พาโล น เถโรเตฺวว สงฺขฺยํ คจฺฉติ. ทหโร เจปิ พฺราหฺมณ โหติ ยุวา สุสุกาฬเกโส ภเทฺรน โยพฺพเนน สมนฺนาคโต ปฐเมน \csromanfolio{69}วยสา, โส จ น กาเม ปริภุญฺชติ, น กามม\-ชฺฌาวสติ, น กามปริฬาเหน ปริฑยฺหติ, น กามวิตกฺเกหิ ขชฺชติ, น กาม\-ปริเยสนาย อุสฺสุโก. อถ โข โส ปณฺฑิโต เถโรเตฺวว สงฺขฺยํ คจฺฉตีติ.}

\prose{เอวํ วุตฺเต กนฺทรายโน พฺราหฺมโณ อุฏฺฐายาสนา เอกํสํ อุตฺตราสงฺคํ กริตฺวา ทหรานํ สตํ\csromansharedfootnote{be00db3ffef522c6}{สุทํ (สี, สฺยา, กํ, อิ)} ภิกฺขูนํ ปาเท สิรสา วนฺทติ “วุทฺธา ภวนฺโต วุทฺธภูมิยํ ฐิตา, ทหรา มยํ ทหรภูมิยํ ฐิตา”ติ.}

\prose{อภิกฺกนฺตํ โภ กจฺจาน ฯเปฯ อุปาสกํ มํ ภวํ กจฺจาโน ธาเรตุ อชฺชตคฺเค ปาณุเปตํ สรณํ คตนฺติ. (7)}

\proseitem{40}{ยสฺมิํ ภิกฺขเว สมเย โจรา พลวนฺโต โหนฺติ, ราชาโน ตสฺมิํ สมเย ทุพฺพลา โหนฺติ. ตสฺมิํ ภิกฺขเว สมเย รญฺโญ น ผาสุ โหติ อติยาตุํ วา นิยฺยาตุํ วา ปจฺจนฺติเม วา ชนปเท อนุสญฺญาตุํ, พฺราหฺมณ\-คหปติกา\-นมฺปิ ตสฺมิํ สมเย น ผาสุ โหติ อติยาตุํ วา นิยฺยาตุํ วา พาหิรานิ วา กมฺมนฺตานิ ปฏิเวกฺขิตุํ. เอวเมวํ โข ภิกฺขเว ยสฺมิํ สมเย ปาปภิกฺขู พลวนฺโต โหนฺติ, เปสลา ภิกฺขู ตสฺมิํ สมเย ทุพฺพลา โหนฺติ. ตสฺมิํ ภิกฺขเว สมเย เปสลา ภิกฺขู ตุณฺหีภูตา ตุณฺหีภูตาว สํฆมชฺเฌ สงฺกสายนฺติ\csromansharedfootnote{b303e891150ebb72}{สงฺกมฺม ฌายนฺติ (ก), สญฺฌายนฺติ (สีฏฺฐ)}. ปจฺจนฺติเม วา ชนปเท อจฺฉนฺติ\csromansharedfootnote{32775a069e7b7e68}{ภชนฺติ (สี, สฺยา, กํ, อิ)}, ตยิทํ ภิกฺขเว โหติ พหุชนาหิตาย พหุชนา\-สุขาย พหุโน ชนสฺส อนตฺถาย อหิตาย ทุกฺขาย เทวมนุสฺสานํ.}

\prose{ยสฺมิํ ภิกฺขเว สมเย ราชาโน พลวนฺโต โหนฺติ, โจรา ตสฺมิํ สมเย ทุพฺพลา โหนฺติ. ตสฺมิํ ภิกฺขเว สมเย รญฺโญ ผาสุ โหติ อติยาตุํ วา นิยฺยาตุํ วา ปจฺจนฺติเม วา ชนปเท อนุสญฺญาตุํ, พฺราหฺมณ\-คหปติกา\-นมฺปิ ตสฺมิํ สมเย ผาสุ โหติ อติยาตุํ วา นิยฺยาตุํ วา พาหิรานิ วา กมฺมนฺตานิ ปฏิเวกฺขิตุํ. เอวเมวํ โข ภิกฺขเว ยสฺมิํ สมเย เปสลา ภิกฺขู พลวนฺโต โหนฺติ. ปาปภิกฺขู ตสฺมิํ สมเย ทุพฺพลา โหนฺติ. ตสฺมิํ ภิกฺขเว สมเย ปาปภิกฺขู ตุณฺหีภูตา ตุณฺหีภูตาว สํฆมชฺเฌ สงฺกสายนฺติ. เยน วา ปน เตน ปกฺกมนฺติ\csromansharedfootnote{281c64c0aaca9b6b}{ปปตนฺติ (สี, สฺยา, กํ, อิ)}. ตยิทํ ภิกฺขเว โหติ พหุชนหิตาย พหุชน\-สุขาย พหุโน ชนสฺส อตฺถาย หิตาย สุขาย เทว\-มนุสฺสา\-นนฺติ. (8)}

\proseitem{41}{\csromanfolio{70}ทฺวินฺนาหํ ภิกฺขเว มิจฺฉา\-ปฏิปตฺติํ น วณฺเณมิ คิหิสฺส วา ปพฺพชิตสฺส วา. คิหี วา ภิกฺขเว ปพฺพชิโต วา มิจฺฉา\-ปฏิปนฺโน มิจฺฉาปฏิปตฺตา\-ธิกรณ\-เหตุ น อาราธโก โหติ ญายํ ธมฺมํ กุสลํ.}

\prose{ทฺวินฺนาหํ ภิกฺขเว สมฺมา\-ปฏิปตฺติํ วณฺเณมิ คิหิสฺส วา ปพฺพชิตสฺส วา. คิหี วา ภิกฺขเว ปพฺพชิโต วา สมฺมาปฏิปนฺโน สมฺมาปฏิปตฺตา\-ธิกรณ\-เหตุ อาราธโก โหติ ญายํ ธมฺมํ กุสลนฺติ. (9)}

\proseitem{42}{เย เต ภิกฺขเว ภิกฺขู ทุคฺคหิเตหิ สุตฺตนฺเตหิ พฺยญฺชน\-ปฺปติรูปเกหิ อตฺถญฺจ ธมฺมญฺจ ปฏิพาหนฺติ, เต ภิกฺขเว ภิกฺขู พหุชนาหิตาย ปฏิปนฺนา พหุชนา\-สุขาย พหุโน ชนสฺส อนตฺถาย อหิตาย ทุกฺขาย เทวมนุสฺสานํ. พหุญฺจ เต ภิกฺขเว ภิกฺขู อปุญฺญํ ปสวนฺติ, เต จิมํ สทฺธมฺมํ อนฺตรธาเปนฺติ.}

\prosewithclosermidruled{เย เต ภิกฺขเว ภิกฺขู สุคฺคหิเตหิ สุตฺตนฺเตหิ พฺยญฺชน\-ปฺปติรูปเกหิ อตฺถญฺจ ธมฺมญฺจ อนุโลเมนฺติ, เต ภิกฺขเว ภิกฺขู พหุชนหิตาย ปฏิปนฺนา พหุชน\-สุขาย พหุโน ชนสฺส อตฺถาย หิตาย สุขาย เทวมนุสฺสานํ. พหุญฺจ เต ภิกฺขเว ภิกฺขู ปุญฺญํ ปสวนฺติ, เต จิมํ สทฺธมฺมํ ฐเปนฺตีติ. (10)}{สมจิตฺตวคฺโค จตุตฺโถ.}
\chapterhead{5. \textbf{ปริสวคฺค}}
\proseitem{43}{เทฺวมา ภิกฺขเว ปริสา. กตมา เทฺว? อุตฺตานา จ ปริสา, คมฺภีรา จ ปริสา. กตมา จ ภิกฺขเว อุตฺตานา ปริสา? อิธ ภิกฺขเว ยสฺสํ ปริสายํ ภิกฺขู อุทฺธตา โหนฺติ อุนฺนฬา จปลา มุขรา วิกิณฺณวาจา มุฏฺฐสฺสตี อสมฺปชานา อสมาหิตา วิพฺภนฺตจิตฺตา ปากตินฺทฺริยา. อยํ วุจฺจติ ภิกฺขเว อุตฺตานา ปริสา.}

\prose{กตมา จ ภิกฺขเว คมฺภีรา ปริสา? อิธ ภิกฺขเว ยสฺสํ ปริสายํ ภิกฺขู อนุทฺธตา โหนฺติ อนุนฺนฬา อจปลา อมุขรา อวิกิณฺณวาจา อุปฏฺฐิตสฺสตี สมฺปชานา สมาหิตา เอกคฺคจิตฺตา สํวุตินฺทฺริยา. อยํ \csromanfolio{71}วุจฺจติ ภิกฺขเว คมฺภีรา ปริสา. อิมา โข ภิกฺขเว เทฺว ปริสา. เอตทคฺคํ ภิกฺขเว อิมาสํ ทฺวินฺนํ ปริสานํ ยทิทํ คมฺภีรา ปริสาติ. (1)}

\proseitem{44}{เทฺวมา ภิกฺขเว ปริสา. กตมา เทฺว? วคฺคา จ ปริสา, สมคฺคา จ ปริสา. กตมา จ ภิกฺขเว วคฺคา ปริสา? อิธ ภิกฺขเว ยสฺสํ ปริสายํ ภิกฺขู ภณฺฑนชาตา กลหชาตา วิวาทาปนฺนา อญฺญมญฺญํ มุขสตฺตีหิ วิตุทนฺตา วิหรนฺติ. อยํ วุจฺจติ ภิกฺขเว วคฺคา ปริสา.}

\prose{กตมา จ ภิกฺขเว สมคฺคา ปริสา? อิธ ภิกฺขเว ยสฺสํ ปริสายํ ภิกฺขู สมคฺคา สมฺโมทมานา อวิวทมานา ขีโรทกีภูตา อญฺญมญฺญํ ปิยจกฺขูหิ สมฺปสฺสนฺตา วิหรนฺติ. อยํ วุจฺจติ ภิกฺขเว สมคฺคา ปริสา. อิมา โข ภิกฺขเว เทฺว ปริสา. เอตทคฺคํ ภิกฺขเว อิมาสํ ทฺวินฺนํ ปริสานํ ยทิทํ สมคฺคา ปริสาติ. (2)}

\proseitem{45}{เทฺวมา ภิกฺขเว ปริสา. กตมา เทฺว? อนคฺควตี จ ปริสา, อคฺควตี จ ปริสา. กตมา จ ภิกฺขเว อนคฺควตี ปริสา? อิธ ภิกฺขเว ยสฺสํ ปริสายํ เถรา ภิกฺขู พาหุลิกา\csromansharedfootnote{2e2251c74eed60f8}{พาหุลฺลิกา (สฺยา, กํ, ก) ฏีกา โอโลเกตพฺพา.} โหนฺติ สาถลิกา, โอกฺกมเน ปุพฺพงฺคมา, ปวิเวเก นิกฺขิตฺตธุรา, น วีริยํ อารภนฺติ อปฺปตฺตสฺส ปตฺติยา อนธิคตสฺส อธิคมาย อสจฺฉิกตสฺส สจฺฉิกิริยาย, เตสํ ปจฺฉิมา ชนตา ทิฏฺฐินุคติํ อาปชฺชติ. สาปิ โหติ พาหุลิกา สาถลิกา, โอกฺกมเน ปุพฺพงฺคมา, ปวิเวเก นิกฺขิตฺตธุรา, น วีริยํ อารภติ อปฺปตฺตสฺส ปตฺติยา อนธิคตสฺส อธิคมาย อสจฺฉิกตสฺส สจฺฉิกิริยาย. อยํ วุจฺจติ ภิกฺขเว อนคฺควตี ปริสา.}

\prose{กตมา จ ภิกฺขเว อคฺควตี ปริสา? อิธ ภิกฺขเว ยสฺสํ ปริสายํ เถรา ภิกฺขู น พาหุลิกา โหนฺติ น สาถลิกา, โอกฺกมเน นิกฺขิตฺตธุรา, ปวิเวเก ปุพฺพงฺคมา, วีริยํ อารภนฺติ อปฺปตฺตสฺส ปตฺติยา อนธิคตสฺส อธิคมาย อสจฺฉิกตสฺส สจฺฉิกิริยาย. เตสํ ปจฺฉิมา ชนตา ทิฏฺฐานุคติํ อาปชฺชติ. สาปิ โหติ น พาหุลิกา น สาถลิกา, โอกฺกมเน นิกฺขิตฺตธุรา, ปวิเวเก ปุพฺพงฺคมา, วีริยํ อารภติ อปฺปตฺตสฺส ปตฺติยา อนธิคตสฺส อธิคมาย อสจฺฉิกตสฺส สจฺฉิกิริยาย. อยํ วุจฺจติ ภิกฺขเว อคฺควตี ปริสา. อิมา โข ภิกฺขเว \csromanfolio{72}เทฺว ปริสา. เอตทคฺคํ ภิกฺขเว อิมาสํ ทฺวินฺนํ ปริสานํ ยทิทํ อคฺควตี ปริสาติ. (3)}

\proseitem{46}{เทฺวมา ภิกฺขเว ปริสา. กตมา เทฺว? อนริยา จ ปริสา, อริยา จ ปริสา. กตมา จ ภิกฺขเว อนริยา ปริสา? อิธ ภิกฺขเว ยสฺสํ ปริสายํ ภิกฺขู “อิทํ ทุกฺขนฺ”ติ ยถาภูตํ นปฺปชานนฺติ, “อยํ ทุกฺขสมุทโย”ติ ยถาภูตํ นปฺปชานนฺติ, “อยํ ทุกฺขนิโรโธ”ติ ยถาภูตํ นปฺปชานนฺติ, “อยํ ทุกฺขนิโรธ\-คามินี ปฏิปทา”ติ ยถาภูตํ นปฺปชานนฺติ. อยํ วุจฺจติ ภิกฺขเว อนริยา ปริสา.}

\prose{กตมา จ ภิกฺขเว อริยา ปริสา? อิธ ภิกฺขเว ยสฺสํ ปริสายํ ภิกฺขู “อิทํ ทุกฺขนฺ”ติ ยถาภูตํ ปชานนฺติ, อยํ “ทุกฺขสมุทโย”ติ ยถาภูตํ ปชานนฺติ, “อยํ ทุกฺขนิโรโธ”ติ ยถาภูตํ ปชานนฺติ, “อยํ ทุกฺขนิโรธ\-คามินี ปฏิปทา”ติ ยถาภูตํ ปชานนฺติ. อยํ วุจฺจติ ภิกฺขเว อริยา ปริสา. อิมา โข ภิกฺขเว เทฺว ปริสา. เอตทคฺคํ ภิกฺขเว อิมาสํ ทฺวินฺนํ ปริสานํ ยทิทํ อริยา ปริสาติ. (4)}

\proseitem{47}{เทฺวมา ภิกฺขเว ปริสา. กตมา เทฺว? ปริสากสโฏ จ ปริสามณฺโฑ จ. กตโม จ ภิกฺขเว ปริสากสโฏ? อิธ ภิกฺขเว ยสฺสํ ปริสายํ ภิกฺขู ฉนฺทาคติํ คจฺฉนฺติ, โทสาคติํ คจฺฉนฺติ, โมหาคติํ คจฺฉนฺติ, ภยาคติํ คจฺฉนฺติ. อยํ วุจฺจติ ภิกฺขเว ปริสากสโฏ.}

\prose{กตโม จ ภิกฺขเว ปริสามณฺโฑ? อิธ ภิกฺขเว ยสฺสํ ปริสายํ ภิกฺขู น ฉนฺทาคติํ คจฺฉนฺติ, น โทสาคติํ คจฺฉนฺติ, น โมหาคติํ คจฺฉนฺติ, น ภยาคติํ คจฺฉนฺติ. อยํ วุจฺจติ ภิกฺขเว ปริสามณฺโฑ. อิมา โข ภิกฺขเว เทฺว ปริสา. เอตทคฺคํ ภิกฺขเว อิมาสํ ทฺวินฺนํ ปริสานํ ยทิทํ ปริสามณฺโฑติ. (5)}

\proseitem{48}{เทฺวมา ภิกฺขเว ปริสา. กตมา เทฺว? โอกฺกาจิตวินีตา ปริสา โนปฏิปุจฺฉา\-วินีตา, ปฏิปุจฺฉา\-วินีตา ปริสา โนโอกฺกาจิต\-วินีตา. กตมา จ ภิกฺขเว โอกฺกาจิตวินีตา ปริสา โนปฏิปุจฺฉา\-วินีตา? อิธ ภิกฺขเว ยสฺสํ ปริสายํ ภิกฺขู เย เต สุตฺตนฺตา ตถาคต\-ภาสิตา \csromanfolio{73}คมฺภีรา คมฺภีรตฺถา โลกุตฺตรา สุญฺญตา\-ปฏิสํยุตฺตา, เตสุ ภญฺญมาเนสุ น สุสฺสูสนฺติ, น โสตํ โอทหนฺติ, น อญฺญา จิตฺตํ อุปฏฺฐเปนฺติ, น จ เต ธมฺเม อุคฺคเหตพฺพํ ปริยาปุณิตพฺพํ มญฺญนฺติ. เย ปน เต สุตฺตนฺตา กวิตา\csromansharedfootnote{3003b0663c50c9fd}{กวิกตา (สพฺพตฺถ) ฏีกา โอโลเกตพฺพา.} กาเวยฺยา จิตฺตกฺขรา จิตฺตพฺยญฺชนา พาหิรกา สาวกภาสิตา, เตสุ ภญฺญมาเนสุ สุสฺสูสนฺติ, โสตํ โอทหนฺติ, อญฺญา จิตฺตํ อุปฏฺฐเปนฺติ, เต ธมฺเม อุคฺคเหตพฺพํ ปริยาปุณิตพฺพํ มญฺญนฺติ. เต จ ตํ ธมฺมํ ปริยาปุณิตฺวา น เจว อญฺญมญฺญํ ปฏิปุจฺฉนฺติ, น จ ปฏิวิจรนฺติ “อิทํ กถํ, อิมสฺส โก อตฺโถ”ติ. เต อวิวฏญฺเจว น วิวรนฺติ, อนุตฺตานีกตญฺจ น อุตฺตานีกโรนฺติ, อเนกวิหิเตสุ จ กงฺขา\-ฐานิเยสุ ธมฺเมสุ กงฺขํ นปฺปฏิวิโนเทนฺติ. อยํ วุจฺจติ ภิกฺขเว โอกฺกาจิตวินีตา ปริสา โนปฏิปุจฺฉา\-วินีตา.}

\prose{กตมา จ ภิกฺขเว ปฏิปุจฺฉา\-วินีตา ปริสา โนโอกฺกาจิต\-วินีตา? อิธ ภิกฺขเว ยสฺสํ ปริสายํ ภิกฺขู เย เต สุตฺตนฺตา กวิตา กาเวยฺยา จิตฺตกฺขรา จิตฺตพฺยญฺชนา พาหิรกา สาวกภาสิตา, เตสุ ภญฺญมาเนสุ น สุสฺสูสนฺติ, น โสตํ โอทหนฺติ, น อญฺญา จิตฺตํ อุปฏฺฐเปนฺติ, น จ เต ธมฺเม อุคฺคเหตพฺพํ ปริยาปุณิตพฺพํ มญฺญนฺติ. เย ปน เต สุตฺตนฺตา ตถาคต\-ภาสิตา คมฺภีรา คมฺภีรตฺถา โลกุตฺตรา สุญฺญตา\-ปฏิสํยุตฺตา, เตสุ ภญฺญมาเนสุ สุสฺสูสนฺติ, โสตํ โอทหนฺติ, อญฺญา จิตฺตํ อุปฏฺฐเปนฺติ, เต จ ธมฺเม อุคฺคเหตพฺพํ ปริยาปุณิตพฺพํ มญฺญนฺติ. เต ตํ ธมฺมํ ปริยาปุณิตฺวา อญฺญมญฺญํ ปฏิปุจฺฉนฺติ, ปฏิวิจรนฺติ “อิทํ กถํ, อิมสฺส โก อตฺโถ”ติ. เต อวิวฏญฺเจว วิวรนฺติ, อนุตฺตานีกตญฺจ อุตฺตานีกโรนฺติ, อเนกวิหิเตสุ จ กงฺขา\-ฐานิเยสุ ธมฺเมสุ กงฺขํ ปฏิวิโนเทนฺติ. อยํ วุจฺจติ ภิกฺขเว ปฏิปุจฺฉา\-วินีตา ปริสา โนโอกฺกาจิต\-วินีตา. อิมา โข ภิกฺขเว เทฺว ปริสา. เอตทคฺคํ ภิกฺขเว อิมาสํ ทฺวินฺนํ ปริสานํ ยทิทํ ปฏิปุจฺฉา\-วินีตา ปริสา โนโอกฺกาจิต\-วินีตาติ. (6)}

\proseitem{49}{เทฺวมา ภิกฺขเว ปริสา. กตมา เทฺว? อามิสครุ ปริสา โนสทฺธมฺมครุ, สทฺธมฺมครุ ปริสา โนอามิสครุ. กตมา จ ภิกฺขเว อามิสครุ ปริสา โนสทฺธมฺมครุ? อิธ ภิกฺขเว ยสฺสํ ปริสายํ ภิกฺขู คิหีนํ โอทาตวสนานํ สมฺมุขา อญฺญมญฺญสฺส วณฺณํ ภาสนฺติ “อสุโก ภิกฺขุ อุภโตภาค\-วิมุตฺโต, อสุโก ปญฺญาวิมุตฺโต, อสุโก \csromanfolio{74}กายสกฺขี, อสุโก ทุฏฺฐิปฺปตฺโต, อสุโก สทฺธาวิมุตฺโต, อสุโก ธมฺมานุสารี, อสุโก สทฺธานุสารี, อสุโก สีลวา กลฺยาณธมฺโม, อสุโก ทุสฺสีโล ปาปธมฺโม”ติ. เต เตน ลาภํ ลภนฺติ, เต ตํ ลาภํ ลภิตฺวา คถิตา\csromansharedfootnote{7f366a1506c3cf51}{คธิตา (ก)} มุจฺฉิตา อชฺโฌปนฺนา\csromansharedfootnote{8c343db603043724}{อชฺโฌสานา (ก), อนชฺโฌปนฺนา (สี, สฺยา, ก) ติกนิปาเต กุสินารวคฺเค ปฐมสุตฺตฏีกา โอโลเกตพฺพา.} อนาทีนวทสฺสาวิโน อนิสฺสรณปญฺญา ปริภุญฺชนฺติ. อยํ วุจฺจติ ภิกฺขเว อามิสครุ ปริสา โนสทฺธมฺมครุ.}

\prose{กตมา จ ภิกฺขเว สทฺธมฺมครุ ปริสา โนอามิสครุ? อิธ ภิกฺขเว ยสฺสํ ปริสายํ ภิกฺขู คิหีนํ โอทาตวสนานํ สมฺมุขา อญฺญมญฺญสฺส วณฺณํ น ภาสนฺติ “อสุโก ภิกฺขุ อุภโตภาค\-วิมุตฺโต, อสุโก ปญฺญาวิมุตฺโต, อสุโก กายสกฺขี, อสุโก ทิฏฺฐิปฺปตฺโต, อสุโก สทฺธาวิมุตฺโต, อสุโก ธมฺมานุสารี, อสุโก สทฺธานุสารี, อสุโก สีลวา กลฺยาณธมฺโม, อสุโก ทุสฺสีโล ปาปธมฺโม”ติ. เต เตน ลาภํ ลภนฺติ, เต ตํ ลาภํ ลภิตฺวา อคถิตา อมุจฺฉิตา อนชฺโฌสนฺนา อาทีนว\-ทสฺสาวิโน นิสฺสรณปญฺญา ปริภุญฺชนฺติ. อยํ วุจฺจติ ภิกฺขเว สทฺธมฺมครุ ปริสา โนอามิสครุ. อิมา โข ภิกฺขเว เทฺว ปริสา. เอตทคฺคํ ภิกฺขเว อิมาสํ ทฺวินฺนํ ปริสานํ ยทิทํ สทฺธมฺมครุ ปริสา โนอามิสครูติ. (7)}

\proseitem{50}{เทฺวมา ภิกฺขเว ปริสา. กตมา เทฺว? วิสมา จ ปริสา, สมา จ ปริสา. กตมา จ ภิกฺขเว วิสมา ปริสา? อิธ ภิกฺขเว ยสฺสํ ปริสายํ อธมฺมกมฺมานิ ปวตฺตนฺติ, ธมฺมกมฺมานิ นปฺปวตฺตนฺติ, อวินยกมฺมานิ ปวตฺตนฺติ, วินยกมฺมานิ นปฺปวตฺตนฺติ. อธมฺมกมฺมานิ ทิปฺปนฺติ, ธมฺมกมฺมานิ น ทิปฺปนฺติ. อวินยกมฺมานิ ทิปฺปนฺติ, วินยกมฺมานิ น ทิปฺปนฺติ. อยํ วุจฺจติ ภิกฺขเว วิสมา ปริสา. ( )\csromansharedfootnote{e0bebefc10f42748}{(วิสมตฺตา ภิกฺขเว ปริสาย อธมฺมกมฺมานิ ปวตฺตนฺติ, ... วินยกมฺมานิ น ทิปฺปนฺติ.) (สี, อิ)}}

\prose{กตมา จ ภิกฺขเว สมา ปริสา? อิธ ภิกฺขเว ยสฺสํ ปริสายํ ธมฺมกมฺมานิ ปวตฺตนฺติ, อธมฺมกมฺมานิ นปฺปวตฺตนฺติ, วินยกมฺมานิ ปวตฺตนฺติ, อวินยกมฺมานิ นปฺปวตฺตนฺติ. ธมฺมกมฺมานิ ทิปฺปนฺติ, อธมฺมกมฺมานิ น ทิปฺปนฺติ. วินยกมฺมานิ ทิปฺปนฺติ, อวินยกมฺมานิ น ทิปฺปนฺติ. อยํ วุจฺจติ ภิกฺขเว}

\prose{\csromanfolio{75}สมา ปริสา. ( )\csromansharedfootnote{00beb6e15e2e14fb}{(สมตฺตา ภิกฺขเว ปริสาย ธมฺมกมฺมานิ ปวตฺตนฺติ, ... อวินยกมฺมานิ น ทิปฺปนฺติ.) (สี, อิ)} อิมา โข ภิกฺขเว เทฺว ปริสา. เอตทคฺคํ ภิกฺขเว อิมาสํ ทฺวินฺนํ ปริสานํ ยทิทํ สมา ปริสาติ. (8)}

\proseitem{51}{เทฺวมา ภิกฺขเว ปริสา. กตมา เทฺว? อธมฺมิกา จ ปริสา, ธมฺมิกา จ ปริสา ฯเปฯ. อิมา โข ภิกฺขเว เทฺว ปริสา. เอตทคฺคํ ภิกฺขเว อิมาสํ ทฺวินฺนํ ปริสานํ ยทิทํ ธมฺมิกา ปริสาติ. (9)}

\proseitem{52}{เทฺวมา ภิกฺขเว ปริสา. กตมา เทฺว? อธมฺมวาทินี จ ปริสา, ธมฺมวาทินี จ ปริสา. กตมา จ ภิกฺขเว อธมฺมวาทินี ปริสา, อิธ ภิกฺขเว ยสฺสํ ปริสายํ ภิกฺขู อธิกรณํ อาทิยนฺติ ธมฺมิกํ วา อธมฺมิกํ วา, เต ตํ อธิกรณํ อาทิยิตฺวา น เจว อญฺญมญฺญํ สญฺญาเปนฺติ, น จ สญฺญตฺติํ อุปคจฺฉนฺติ, น จ นิชฺฌาเปนฺติ, น จ นิชฺฌตฺติํ อุปคจฺฉนฺติ. เต อสญฺญตฺติพลา อนิชฺฌตฺติพลา อปฺปฏินิสฺสคฺค\-มนฺติโน ตเมว อธิกรณํ ถามสา ปรามาสา\csromansharedfootnote{3646da9b122f47ad}{ปรามสฺส (สี, อิ)} อภินิวิสฺส โวหรนฺติ “อิทเมว สจฺจํ โมฆมญฺญนฺ”ติ. อยํ วุจฺจติ ภิกฺขเว อธมฺมวาทินี ปริสา.}

\prosewithclosermidruled{กตมา จ ภิกฺขเว ธมฺมวาทินี ปริสา? อิธ ภิกฺขเว ยสฺสํ ปริสายํ ภิกฺขู อธิกรณํ อาทิยนฺติ ธมฺมิกํ วา อธมฺมิกํ วา, เต ตํ อธิกรณํ อาทิยิตฺวา อญฺญมญฺญํ สญฺญาเปนฺติ เจว สญฺญตฺติญฺจ อุปคจฺฉนฺติ, นิชฺฌาเปนฺติ เจว นิชฺฌตฺติญฺจ อุปคจฺฉนฺติ. เต สญฺญตฺติพลา นิชฺฌตฺติพลา ปฏินิสฺสคฺค\-มนฺติโน น ตเมว อธิกรณํ ถามสา ปรามาสา อภินิวิสฺส โวหรนฺติ “อิทเมว สจฺจํ โมฆมญฺญนฺ”ติ. อยํ วุจฺจติ ภิกฺขเว ธมฺมวาทินี ปริสา. อิมา โข ภิกฺขเว เทฺว ปริสา. เอตทคฺคํ ภิกฺขเว อิมาสํ ทฺวินฺนํ ปริสานํ ยทิทํ ธมฺมวาทินี ปริสาติ.}{(10) ปริสวคฺโค ปญฺจโม.}
\csromancenter{\textbf{ตสฺสุทฺทานํ}}
\setlength{\csromangathapagewidth}{0pt}
\setlength{\csromangathawakwidth}{0pt}
\csromangathameasuregroup{อุตฺตานา วคฺคา อคฺควตี, \\ โอกฺกาจิตอามิสญฺเจว,}{\csromangathabat{อุตฺตานา วคฺคา อคฺควตี,}{อริยา กสโฏ จ ปญฺจโม.} \\ \csromangathabat{โอกฺกาจิตอามิสญฺเจว,}{วิสมา อธมฺมาธมฺมิเยน จาติ.} \\ \textbf{ปฐโม ปณฺณาสโก สมตฺโต.}}
\csromangathasetleft{อุตฺตานา วคฺคา อคฺควตี, \\ โอกฺกาจิตอามิสญฺเจว,}
\csromangathagroup{\csromangathastanza{\csromangathabat{อุตฺตานา วคฺคา อคฺควตี,}{อริยา กสโฏ จ ปญฺจโม.} \\ \csromangathabat{โอกฺกาจิตอามิสญฺเจว,}{วิสมา อธมฺมาธมฺมิเยน จาติ.} \\ \textbf{ปฐโม ปณฺณาสโก สมตฺโต.}}}

\csromanmatikaanchor{mtk.1}
\csromanmatikaanchor{mtk.103}
\csromantocmarknopage{chapter}{3. ตติยปณฺณาสก}{mtk.1}
\bookmark[dest={mtk.1},level=0]{3. ตติยปณฺณาสก}
\chapterheadpageread{2. ทุติย\-ปณฺณาสก}
\chapterhead{\textbf{(6) 1. ปุคฺคลวคฺค}}
\proseitem{53}{\csromanfolio{76}เทฺวเม ภิกฺขเว ปุคฺคลา โลเก อุปฺปชฺชมานา อุปฺปชฺชนฺติ พหุชนหิตาย พหุชน\-สุขาย พหุโน ชนสฺส อตฺถาย หิตาย สุขาย เทวมนุสฺสานํ. กตเม เทฺว? ตถาคโต จ อรหํ สมฺมาสมฺพุทฺโธ, ราชา จ จกฺกวตฺตี. อิเม โข ภิกฺขเว เทฺว ปุคฺคลา โลเก อุปฺปชฺชมานา อุปฺปชฺชนฺติ พหุชนหิตาย พหุชน\-สุขาย พหุโน ชนสฺส อตฺถาย หิตาย สุขาย เทว\-มนุสฺสา\-นนฺติ. (1)}

\proseitem{54}{เทฺวเม ภิกฺขเว ปุคฺคลา โลเก อุปฺปชฺชมานา อุปฺปชฺชนฺติ อจฺฉริย\-มนุสฺสา. กตเม เทฺว? ตถาคโต จ อรหํ สมฺมาสมฺพุทฺโธ, ราชา จ จกฺกวตฺตี. อิเม โข ภิกฺขเว เทฺว ปุคฺคลา โลเก อุปฺปชฺชมานา อุปฺปชฺชนฺติ อจฺฉริย\-มนุสฺสาติ. (2)}

\proseitem{55}{ทฺวินฺนํ ภิกฺขเว ปุคฺคลานํ กาลกิริยา พหุโน ชนสฺส อนุตปฺปา โหติ. กตเมสํ ทฺวินฺนํ? ตถาคตสฺส จ อรหโต สมฺมา\-สมฺพุทฺธสฺส, รญฺโญ จ จกฺกวตฺติสฺส. อิเมสํ โข ภิกฺขเว ทฺวินฺนํ ปุคฺคลานํ กาลกิริยา พหุโน ชนสฺส อนุตปฺปา โหตีติ. (3)}

\proseitem{56}{เทฺวเม ภิกฺขเว ถูปารหา. กตเม เทฺว? ตถาคโต จ อรหํ สมฺมาสมฺพุทฺโธ, ราชา จ จกฺกวตฺตี. อิเม โข ภิกฺขเว เทฺว ถูปารหาติ. (4)}

\proseitem{57}{เทฺวเม ภิกฺขเว พุทฺธา. กตเม เทฺว? ตถาคโต จ อรหํ สมฺมาสมฺพุทฺโธ, ปจฺเจกพุทฺโธ จ. อิเม โข ภิกฺขเว เทฺว พุทฺธาติ. (5)}

\proseitem{58}{เทฺวเม ภิกฺขเว อสนิยา ผลนฺติยา น สนฺตสนฺติ. กตเม เทฺว? ภิกฺขุ จ ขีณาสโว, หตฺถาชานีโย จ. อิเม โข ภิกฺขเว เทฺว อสนิยา ผลนฺติยา น สนฺตสนฺตีติ. (6)}

\proseitem{59}{เทฺวเม ภิกฺขเว อสนิยา ผลนฺติยา น สนฺตสนฺติ. กตเม เทฺว? ภิกฺขุ จ ขีณาสโว, อสฺสาชานีโย จ. อิเม โข ภิกฺขเว เทฺว อสนิยา ผลนฺติยา น สนฺตสนฺตีติ. (7)}

\chapterheadpageread{(6) 1. ปุคฺคลวคฺค}
\proseitem{60}{\csromanfolio{77}เทฺวเม ภิกฺขเว อสนิยา ผลนฺติยา น สนฺตสนฺติ. กตเม เทฺว? ภิกฺขุ จ ขีณาสโว, สีโห จ มิคราชา. อิเม โข ภิกฺขเว เทฺว อสนิยา ผลนฺติยา น สนฺตสนฺตีติ. (8)}

\proseitem{61}{เทฺวเม ภิกฺขเว อตฺถวเส สมฺปสฺสมานา กิํปุริสา มานุสิํ วาจํ น ภาสนฺติ. กตเม เทฺว? มา จ มุสา ภณิมฺหา, มา จ ปรํ อภูเตน อพฺภาจิกฺขิมฺหาติ. อิเม โข ภิกฺขเว เทฺว อตฺถวเส สมฺปสฺสมานา กิํปุริสา มานุสิํ วาจํ น ภาสนฺตีติ. (9)}

\proseitem{62}{ทฺวินฺนํ ธมฺมานํ ภิกฺขเว อติตฺโต อปฺปฏิวาโน มาตุคาโม กาลํ กโรติ. กตเมสํ ทฺวินฺนํ? เมถุน\-สมาปตฺติยา จ วิชายนสฺส จ. อิเมสํ โข ภิกฺขเว ทฺวินฺนํ ธมฺมานํ อติตฺโต อปฺปฏิวาโน มาตุคาโม กาลํ กโรตีติ. (10)}

\proseitem{63}{อสนฺตสนฺนิวาสญฺจ โว ภิกฺขเว เทเสสฺสามิ สนฺตสนฺนิวาสญฺจ, ตํ สุณาถ, สาธุกํ มนสิ กโรถ ภาสิสฺสามีติ. “เอวํ ภนฺเต”ติ โข เต ภิกฺขู ภควโต ปจฺจสฺโสสุํ. ภควา เอตทโวจ\csromandash{}}

\prose{กถญฺจ ภิกฺขเว อสนฺต\-สนฺนิวาโส โหติ, กถญฺจ อสนฺโต สนฺนิวสนฺติ? อิธ ภิกฺขเว เถรสฺส ภิกฺขุโน เอวํ โหติ “เถโรปิ มํ น วเทยฺย, มชฺฌิโมปิ มํ น วเทยฺย, นโวปิ มํ น วเทยฺย. เถรมฺปาหํ น วเทยฺยํ, มชฺฌิมมฺปาหํ น วเทยฺยํ, นวมฺปาหํ น วเทยฺยํ. เถโร เจปิ มํ วเทยฺย, อหิตานุกมฺปี มํ วเทยฺย โน หิตานุกมฺปี. “โน”ติ นํ วเทยฺยํ, วิเหเฐยฺยํ\csromansharedfootnote{39013b339d670cc3}{วิเหเสยฺยํ (สี, สฺยา, กํ, อิ)}, ปสฺสมฺปิสฺส นปฺปฏิกเรยฺยํ. มชฺฌิโม เจปิ มํ วเทยฺย ฯเปฯ. นโว เจปิ มํ วเทยฺย, อหิตานุกมฺปี มํ วเทยฺย โน หิตานุกมฺปี. “โน”ติ นํ วเทยฺยํ, วิเหเฐยฺยํ, ปสฺสมฺปิสฺส นปฺปฏิกเรยฺยํ”. มชฺฌิมสฺสปิ ภิกฺขุโน เอวํ โหติ ฯเปฯ. นวสฺสปิ ภิกฺขุโน เอวํ โหติ “เถโรปิ มํ น วเทยฺย, มชฺฌิโมปิ มํ น วเทยฺย, นโวปิ มํ น วเทยฺย. เถรมฺปาหํ น วเทยฺยํ, มชฺฌิมมฺปาหํ น วเทยฺยํ, นวมฺปาหํ น วเทยฺยํ. เถโร เจปิ มํ วเทยฺย, อหิตานุกมฺปี มํ วเทยฺย โน หิตานุกมฺมี. “โน”ติ นํ วเทยฺยํ, วิเหเฐยฺยํ, ปสฺสมฺปิสฺส นปฺปฏิกเรยฺยํ. มชฺฌิโม เจปิ มํ วเทยฺย, นโว เจปิ \csromanfolio{78}มํ วเทยฺย, อหิตานุกมฺปี มํ วเทยฺย โน หิตานุกมฺปี. “โน”ติ นํ วเทยฺยํ, วิเหเฐยฺยํ, ปสฺสมฺปิสฺส นปฺปฏิกเรยฺยํ”. เอวํ โข ภิกฺขเว อสนฺต\-สนฺนิวาโส โหติ, เอวญฺจ อสนฺโต สนฺนิวสนฺติ.}

\prose{กถญฺจ ภิกฺขเว สนฺตสนฺนิวาโส โหติ, กถญฺจ สนฺโต สนฺนิวสนฺติ. อิธ ภิกฺขเว เถรสฺส ภิกฺขุโน เอวํ โหติ “เถโรปิ มํ วเทยฺย, มชฺฌิโมปิ มํ วเทยฺย, นโวปิ มํ วเทยฺย. เถรมฺปาหํ วเทยฺยํ, มชฺฌิมมฺปาหํ วเทยฺยํ, นวมฺปาหํ วเทยฺยํ. เถโร เจปิ มํ วเทยฺย, หิตานุกมฺปี มํ วเทยฺย โน อหิตานุกมฺปิ. “สาธู”ติ นํ วเทยฺยํ, น นํ วิเหเฐยฺยํ, ปสฺสมฺปิสฺส ปฏิกเรยฺยํ. มชฺฌิโม เจปิ มํ วเทยฺย ฯเปฯ. นโว เจปิ มํ วเทยฺย, หิตานุกมฺปี มํ วเทยฺย โน อหิตานุกมฺปี. “สาธู”ติ นํ วเทยฺยํ, น นํ วิเหเฐยฺยํ, ปสฺสมฺปิสฺส ปฏิกเรยฺยํ”. มชฺฌิมสฺสปิ ภิกฺขุโน เอวํ โหติ ฯเปฯ. นวสฺสปิ ภิกฺขุโน เอวํ โหติ “เถโรปิ มํ วเทยฺย, มชฺฌิโมปิ มํ วเทยฺย, นโวปิ มํ วเทยฺย. เถรมฺปาหํ วเทยฺยํ, มชฺฌิมมฺปาหํ วเทยฺยํ, นวมฺปาหํ วเทยฺยํ. เถโร เจปิ มํ วเทยฺย, หิตานุกมฺปี มํ วเทยฺย, โน อหิตานุกมฺปี. “สาธู”ติ นํ วเทยฺยํ, น นํ วิเหเฐยฺยํ, ปสฺสมฺปิสฺส ปฏิกเรยฺยํ. มชฺฌิโม เจปิ มํ วเทยฺย ฯเปฯ. นโว เจปิ มํ วเทยฺย, หิตานุกมฺปี มํ วเทยฺย โน อหิตานุกมฺปี. “สาธู”ติ นํ วเทยฺยํ, น นํ วิเหเฐยฺยํ, ปสฺสมฺปิสฺส ปฏิกเรยฺยํ”. เอวํ โข ภิกฺขเว สนฺตสนฺนิวาโส โหติ, เอวญฺจ สนฺโต สนฺนิวสนฺตีติ. (11)}

\proseitemwithclosermidruled{64}{ยสฺมิํ ภิกฺขเว อธิกรเณ อุภโต วจีสํสาโร ทิฏฺฐิปฬาโส เจตโส อาฆาโต อปฺปจฺจโย อนภิรทฺธิ อชฺฌตฺตํ อวูปสนฺตํ โหติ, ตสฺเมตํ ภิกฺขเว อธิกรเณ ปาฏิกงฺขํ\csromandash{}ทีฆตฺตาย ขรตฺตาย วาฬตฺตาย สํวตฺติสฺสติ, ภิกฺขู จ น ผาสุํ\csromansharedfootnote{bef724bc99034e48}{ผาสุ (ก)} วิหริสฺสนฺติ. ยสฺมิญฺจ โข ภิกฺขเว อธิกรเณ อุภโต วจีสํสาโร ทิฏฺฐิปฬาโส เจตโส อาฆาโต อปฺปจฺจโย อนภิรทฺธิ อชฺฌตฺตํ สุวูปสนฺตํ โหติ, ตสฺเมตํ ภิกฺขเว อธิกรเณ ปาฏิกงฺขํ\csromandash{}น ทีฆตฺตาย ขรตฺตาย วาฬตฺตาย สํวตฺติสฺสติ, ภิกฺขู จ ผาสุํ วิหริสฺสนฺตีติ. (12)}{ปุคฺคลวคฺโค ปฐโม.}
\chapterheadpageread{(7) 2. สุขวคฺค \textbf{(7) 2. สุขวคฺค}}
\proseitem{65}{\csromanfolio{79}เทฺวมานิ ภิกฺขเว สุขานิ. กตมานิ เทฺว? คิหิสุขญฺจ ปพฺพชิตสุขญฺจ. อิมานิ โข ภิกฺขเว เทฺว สุขานิ. เอตทคฺคํ ภิกฺขเว อิเมสํ ทฺวินฺนํ สุขานํ ยทิทํ ปพฺพชิต\-สุขนฺติ. (1)}

\proseitem{66}{เทฺวมานิ ภิกฺขเว สุขานิ. กตมานิ เทฺว? กามสุขญฺจ เนกฺขมฺมสุขญฺจ. อิมานิ โข ภิกฺขเว เทฺว สุขานิ. เอตทคฺคํ ภิกฺขเว อิเมสํ ทฺวินฺนํ สุขานํ ยทิทํ เนกฺขมฺม\-สุขนฺติ. (2)}

\proseitem{67}{เทฺวมานิ ภิกฺขเว สุขานิ. กตมานิ เทฺว? อุปธิสุขญฺจ นิรุปธิสุขญฺจ. อิมานิ โข ภิกฺขเว เทฺว สุขานิ. เอตทคฺคํ ภิกฺขเว อิเมสํ ทฺวินฺนํ สุขานํ ยทิทํ นิรุปธิสุ\-ขนฺติ. (3)}

\proseitem{68}{เทฺวมานิ ภิกฺขเว สุขานิ. กตมานิ เทฺว? สาสวสุขญฺจ อนาสวสุขญฺจ. อิมานิ โข ภิกฺขเว เทฺว สุขานิ. เอตทคฺคํ ภิกฺขเว อิเมสํ ทฺวินฺนํ สุขานํ ยทิทํ อนาสว\-สุขนฺติ. (4)}

\proseitem{69}{เทฺวมานิ ภิกฺขเว สุขานิ. กตมานิ เทฺว? สามิสญฺจ สุขํ, นิรามิสญฺจ สุขํ. อิมานิ โข ภิกฺขเว เทฺว สุขานิ. เอตทคฺคํ ภิกฺขเว อิเมสํ ทฺวินฺนํ สุขานํ ยทิทํ นิรามิสํ สุขนฺติ. (5)}

\proseitem{70}{เทฺวมานิ ภิกฺขเว สุขานิ. กตมานิ เทฺว? อริยสุขญฺจ อนริยสุขญฺจ. อิมานิ โข ภิกฺขเว เทฺว สุขานิ. เอตทคฺคํ ภิกฺขเว อิเมสํ ทฺวินฺนํ สุขานํ ยทิทํ อริยสุขนฺติ. (6)}

\proseitem{71}{เทฺวมานิ ภิกฺขเว สุขานิ. กตมานิ เทฺว? กายิกญฺจ สุขํ, เจตสิกญฺจ สุขํ. อิมานิ โข ภิกฺขเว เทฺว สุขานิ. เอตทคฺคํ ภิกฺขเว อิเมสํ ทฺวินฺนํ สุขานํ ยทิทํ เจตสิกํ สุขนฺติ. (7)}

\proseitem{72}{เทฺวมานิ ภิกฺขเว สุขานิ. กตมานิ เทฺว? สปฺปีติกญฺจ สุขํ, นิปฺปีติกญฺจ สุขํ. อิมานิ โข ภิกฺขเว เทฺว สุขานิ. เอตทคฺคํ ภิกฺขเว อิเมสํ ทฺวินฺนํ สุขานํ ยทิทํ นิปฺปีติกํ สุขนฺติ. (8)}

\proseitem{73}{เทฺวมานิ ภิกฺขเว สุขานิ. กตมานิ เทฺว? สาตสุขญฺจ อุเปกฺขาสุขญฺจ. อิมานิ โข ภิกฺขเว เทฺว สุขานิ. เอตทคฺคํ ภิกฺขเว อิเมสํ ทฺวินฺนํ สุขานํ ยทิทํ อุเปกฺขาสุ\-ขนฺติ. (9)}

\proseitem{74}{\csromanfolio{80}เทฺวมานิ ภิกฺขเว สุขานิ. กตมานิ เทฺว? สมาธิสุขญฺจ อสมาธิสุขญฺจ. อิมานิ โข ภิกฺขเว เทฺว สุขานิ. เอตทคฺคํ ภิกฺขเว อิเมสํ ทฺวินฺนํ สุขานํ ยทิทํ สมาธิสุ\-ขนฺติ. (10)}

\proseitem{75}{เทฺวมานิ ภิกฺขเว สุขานิ. กตมานิ เทฺว? สปฺปีติการมฺมณญฺจ สุขํ, นิปฺปีติการมฺมณญฺจ สุขํ. อิมานิ โข ภิกฺขเว เทฺว สุขานิ. เอตทคฺคํ ภิกฺขเว อิเมสํ ทฺวินฺนํ สุขานํ ยทิทํ นิปฺปีติกา\-รมฺมณํ สุขนฺติ. (11)}

\proseitem{76}{เทฺวมานิ ภิกฺขเว สุขานิ. กตมานิ เทฺว? สาตารมฺมณญฺจ สุขํ, อุเปกฺขารมฺมณญฺจ สุขํ. อิมานิ โข ภิกฺขเว เทฺว สุขานิ. เอตทคฺคํ ภิกฺขเว อิเมสํ ทฺวินฺนํ สุขานํ ยทิทํ อุเปกฺขา\-รมฺมณํ สุขนฺติ. (12)}

\proseitemwithclosermidruled{77}{เทฺวมานิ ภิกฺขเว สุขานิ. กตมานิ เทฺว? รูปารมฺมณญฺจ สุขํ, อรูปารมฺมณญฺจ สุขํ. อิมานิ โข ภิกฺขเว เทฺว สุขานิ. เอตทคฺคํ ภิกฺขเว อิเมสํ ทฺวินฺนํ สุขานํ ยทิทํ อรูปารมฺมณํ สุขนฺติ. (13)}{สุขวคฺโค ทุติโย.}
\chapterhead{\textbf{(8) 3. สนิมิตฺตวคฺค}}
\proseitem{78}{สนิมิตฺตา ภิกฺขเว อุปฺปชฺชนฺติ ปาปกา อกุสลา ธมฺมา, โน อนิมิตฺตา, ตสฺเสว นิมิตฺตสฺส ปหานา เอวํ เต ปาปกา อกุสลา ธมฺมา น โหนฺตีติ. (1)}

\proseitem{79}{สนิทานา ภิกฺขเว อุปฺปชฺชนฺติ ปาปกา อกุสลา ธมฺมา, โน อนิทานา, ตสฺเสว นิทานสฺส ปหานา เอวํ เต ปาปกา อกุสลา ธมฺมา น โหนฺตีติ. (2)}

\proseitem{80}{สเหตุกา ภิกฺขเว อุปฺปชฺชนฺติ ปาปกา อกุสลา ธมฺมา, โน อเหตุกา, ตสฺเสว เหตุสฺส ปหานา เอวํ เต ปาปกา อกุสลา ธมฺมา น โหนฺตีติ. (3)}

\proseitem{81}{สสงฺขารา ภิกฺขเว อุปฺปชฺชนฺติ ปาปกา อกุสลา ธมฺมา, โน อสงฺขารา, เตสํเยว สงฺขารานํ ปหานา เอวํ เต ปาปกา อกุสลา ธมฺมา น โหนฺตีติ. (4)}

\chapterheadpageread{(9) 4. ธมฺมวคฺค}
\proseitem{82}{\csromanfolio{81}สปฺปจฺจยา ภิกฺขเว อุปฺปชฺชนฺติ ปาปกา อกุสลา ธมฺมา, โน อปฺปจฺจยา, ตสฺเสว ปจฺจยสฺส ปหานา เอวํ เต ปาปกา อกุสลา ธมฺมา น โหนฺตีติ. (5)}

\proseitem{83}{สรูปา ภิกฺขเว อุปฺปชฺชนฺติ ปาปกา อกุสลา ธมฺมา, โน อรูปา, ตสฺเสว รูปสฺส ปหานา เอวํ เต ปาปกา อกุสลา ธมฺมา น โหนฺตีติ. (6)}

\proseitem{84}{สเวทนา ภิกฺขเว อุปฺปชฺชนฺติ ปาปกา อกุสลา ธมฺมา, โน อเวทนา, ตสฺสาเยว เวทนาย ปหานา เอวํ เต ปาปกา อกุสลา ธมฺมา น โหนฺตีติ. (7)}

\proseitem{85}{สสญฺญา ภิกฺขเว อุปฺปชฺชนฺติ ปาปกา อกุสลา ธมฺมา, โน อสญฺญา, ตสฺสาเยว สญฺญาย ปหานา เอวํ เต ปาปกา อกุสลา ธมฺมา น โหนฺตีติ. (8)}

\proseitem{86}{สวิญฺญาณา ภิกฺขเว อุปฺปชฺชนฺติ ปาปกา อกุสลา ธมฺมา, โน อวิญฺญาณา, ตสฺเสว วิญฺญาณสฺส ปหานา เอวํ เต ปาปกา อกุสลา ธมฺมา น โหนฺตีติ. (9)}

\proseitemwithclosermidruled{87}{สงฺขตา\-รมฺมณา ภิกฺขเว อุปฺปชฺชนฺติ ปาปกา อกุสลา ธมฺมา, โน อสงฺขตา\-รมฺมณา, ตสฺเสว สงฺขตสฺส ปหานา เอวํ เต ปาปกา อกุสลา ธมฺมา น โหนฺตีติ. (10)}{สนิมิตฺตวคฺโค ตติโย.}
\chapterhead{(9) 4. ธมฺมวคฺค}
\proseitem{88}{เทฺวเม ภิกฺขเว ธมฺมา. กตเม เทฺว? เจโตวิมุตฺติ จ ปญฺญาวิมุตฺติ จ. อิเม โข ภิกฺขเว เทฺว ธมฺมาติ. (1)}

\proseitem{89}{เทฺวเม ภิกฺขเว ธมฺมา. กตเม เทฺว? ปคฺคาโห จ อวิกฺเขโป จ. อิเม โข ภิกฺขเว เทฺว ธมฺมาติ. (2)}

\proseitem{90}{\csromanfolio{82}เทฺวเม ภิกฺขเว ธมฺมา. กตเม เทฺว? นามญฺจ รูปญฺจ. อิเม โข ภิกฺขเว เทฺว ธมฺมาติ. (3)}

\proseitem{91}{เทฺวเม ภิกฺขเว ธมฺมา. กตเม เทฺว? วิชฺชา จ วิมุตฺติ จ. อิเม โข ภิกฺขเว เทฺว ธมฺมาติ. (4)}

\proseitem{92}{เทฺวเม ภิกฺขเว ธมฺมา. กตเม เทฺว? ภวทิฏฺฐิ จ วิภวทิฏฺฐิ จ. อิเม โข ภิกฺขเว เทฺว ธมฺมาติ. (5)}

\proseitem{93}{เทฺวเม ภิกฺขเว ธมฺมา. กตเม เทฺว? อหิริกญฺจ อโนตฺตปฺปญฺจ. อิเม โข ภิกฺขเว เทฺว ธมฺมาติ. (6)}

\proseitem{94}{เทฺวเม ภิกฺขเว ธมฺมา. กตเม เทฺว? หิรี จ โอตฺตปฺปญฺจ. อิเม โข ภิกฺขเว เทฺว ธมฺมาติ. (7)}

\proseitem{95}{เทฺวเม ภิกฺขเว ธมฺมา. กตเม เทฺว? โทวจสฺสตา จ ปาปมิตฺตตา จ. อิเม โข ภิกฺขเว เทฺว ธมฺมาติ. (8)}

\proseitem{96}{เทฺวเม ภิกฺขเว ธมฺมา. กตเม เทฺว? โสวจสฺสตา จ กลฺยาณามิตฺตตา จ. อิเม โข ภิกฺขเว เทฺว ธมฺมาติ. (9)}

\proseitem{97}{เทฺวเม ภิกฺขเว ธมฺมา. กตเม เทฺว? ธาตุกุสลตา จ มนสิการ\-กุสลตา จ. อิเม โข ภิกฺขเว เทฺว ธมฺมาติ. (10)}

\proseitemwithclosermidruled{98}{เทฺวเม ภิกฺขเว ธมฺมา. กตเม เทฺว? อาปตฺติกุสลตา จ อาปตฺติวุฏฺฐาน\-กุสลตา จ. อิเม โข ภิกฺขเว เทฺว ธมฺมาติ. (11)}{ธมฺมวคฺโค จตุตฺโถ.}
\chapterhead{\textbf{(10) 5. พาลวคฺค}}
\proseitem{99}{เทฺวเม ภิกฺขเว พาลา. กตเม เทฺว? โย จ อนาคตํ ภารํ วหติ, โย จ อาคตํ ภารํ น วหติ. อิเม โข ภิกฺขเว เทฺว พาลาติ. (1)}

\chapterheadpageread{(10) 5. พาลวคฺค}
\proseitem{100}{\csromanfolio{83}เทฺวเม ภิกฺขเว ปณฺฑิตา. กตเม เทฺว? โย จ อนาคตํ ภารํ น วหติ, โย จ อาคตํ ภารํ วหติ. อิเม โข ภิกฺขเว เทฺว ปณฺฑิตาติ. (2)}

\proseitem{101}{เทฺวเม ภิกฺขเว พาลา. กตเม เทฺว? โย จ อกปฺปิเย กปฺปิยสญฺญี, โย จ กปฺปิเย อกปฺปิยสญฺญี. อิเม โข ภิกฺขเว เทฺว พาลาติ. (3)}

\proseitem{102}{เทฺวเม ภิกฺขเว ปณฺฑิตา. กตเม เทฺว? โย จ อกปฺปิเย อกปฺปิยสญฺญี, โย จ กปฺปิเย กปฺปิยสญฺญี. อิเม โข ภิกฺขเว เทฺว ปณฺฑิตาติ. (4)}

\proseitem{103}{เทฺวเม ภิกฺขเว พาลา. กตเม เทฺว? โย จ อนาปตฺติยา อาปตฺติสญฺญี, โย จ อาปตฺติยา อนาปตฺติสญฺญี. อิเม โข ภิกฺขเว เทฺว พาลาติ. (5)}

\proseitem{104}{เทฺวเม ภิกฺขเว ปณฺฑิตา. กตเม เทฺว? โย จ อนาปตฺติยา อนาปตฺติสญฺญี, โย จ อาปตฺติยา อาปตฺติสญฺญี. อิเม โข ภิกฺขเว เทฺว ปณฺฑิตาติ. (6)}

\proseitem{105}{เทฺวเม ภิกฺขเว พาลา. กตเม เทฺว? โย จ อธมฺเม ธมฺมสญฺญี, โย จ ธมฺเม อธมฺมสญฺญี. อิเม โข ภิกฺขเว เทฺว พาลาติ. (7)}

\proseitem{106}{เทฺวเม ภิกฺขเว ปณฺฑิตา. กตเม เทฺว? โย จ ธมฺเม ธมฺมสญฺญี, โย จ อธมฺเม อธมฺมสญฺญี. อิเม โข ภิกฺขเว เทฺว ปณฺฑิตาติ. (8)}

\proseitem{107}{เทฺวเม ภิกฺขเว พาลา. กตเม เทฺว? โย จ อวินเย วินยสญฺญี, โย จ วินเย อวินยสญฺญี. อิเม โข ภิกฺขเว เทฺว พาลาติ. (9)}

\proseitem{108}{เทฺวเม ภิกฺขเว ปณฺฑิตา. กตเม เทฺว? โย จ อวินเย อวินยสญฺญี, โย จ วินเย วินยสญฺญี. อิเม โข ภิกฺขเว เทฺว ปณฺฑิตาติ. (10)}

\proseitem{109}{\csromanfolio{84}ทฺวินฺนํ ภิกฺขเว อาสวา วฑฺฒนฺติ. กตเมสํ ทฺวินฺนํ? โย จ น กุกฺกุจฺจายิ\-ตพฺพํ กุกฺกุจฺจายติ, โย จ กุกฺกุจฺจายิ\-ตพฺพํ น กุกฺกุจฺจายติ. อิเมสํ โข ภิกฺขเว ทฺวินฺนํ อาสวา วฑฺฒนฺตีติ. (11)}

\proseitem{110}{ทฺวินฺนํ ภิกฺขเว อาสวา น วฑฺฒนฺติ. กตเมสํ ทฺวินฺนํ? โย จ น กุกฺกุจฺจายิ\-ตพฺพํ น กุกฺกุจฺจายติ, โย จ กุกฺกุจฺจายิ\-ตพฺพํ กุกฺกุจฺจายติ. อิเมสํ โข ภิกฺขเว ทฺวินฺนํ อาสวา น วฑฺฒนฺตีติ. (12)}

\proseitem{111}{ทฺวินฺนํ ภิกฺขเว อาสวา วฑฺฒนฺติ. กตเมสํ ทฺวินฺนํ? โย จ อกปฺปิเย กปฺปิยสญฺญี, โย จ กปฺปิเย อกปฺปิยสญฺญี. อิเมสํ โข ภิกฺขเว ทฺวินฺนํ อาสวา วฑฺฒนฺตีติ. (13)}

\proseitem{112}{ทฺวินฺนํ ภิกฺขเว อาสวา น วฑฺฒนฺติ. กตเมสํ ทฺวินฺนํ? โย จ อกปฺปิเย อกปฺปิยสญฺญี, โย จ กปฺปิเย กปฺปิยสญฺญี. อิเมสํ โข ภิกฺขเว ทฺวินฺนํ อาสวา น วฑฺฒนฺตีติ. (14)}

\proseitem{113}{ทฺวินฺนํ ภิกฺขเว อาสวา วฑฺฒนฺติ. กตเมสํ ทฺวินฺนํ? โย จ อาปตฺติยา อนาปตฺติสญฺญี โย จ อนาปตฺติยา อาปตฺติสญฺญี. อิเมสํ โข ภิกฺขเว ทฺวินฺนํ อาสวา วฑฺฒนฺตีติ. (15)}

\proseitem{114}{ทฺวินฺนํ ภิกฺขเว อาสวา น วฑฺฒนฺติ. กตเมสํ ทฺวินฺนํ? โย จ อาปตฺติยา อาปตฺติสญฺญี, โย จ อนาปตฺติยา อนาปตฺติสญฺญี. อิเมสํ โข ภิกฺขเว ทฺวินฺนํ อาสวา น วฑฺฒนฺตีติ. (16)}

\proseitem{115}{ทฺวินฺนํ ภิกฺขเว อาสวา วฑฺฒนฺติ. กตเมสํ ทฺวินฺนํ? โย จ อธมฺเม ธมฺมสญฺญี, โย จ ธมฺเม อธมฺมสญฺญี. อิเมสํ โข ภิกฺขเว ทฺวินฺนํ อาสวา วฑฺฒนฺตีติ. (17)}

\proseitem{116}{ทฺวินฺนํ ภิกฺขเว อาสวา น วฑฺฒนฺติ. กตเมสํ ทฺวินฺนํ? โย จ ธมฺเม ธมฺมสญฺญี, โย จ อธมฺเม อธมฺมสญฺญี. อิเมสํ โข ภิกฺขเว ทฺวินฺนํ อาสวา น วฑฺฒนฺตีติ. (18)}

\chapterheadpageread{(10) 5. พาลวคฺค}
\proseitem{117}{\csromanfolio{85}ทฺวินฺนํ ภิกฺขเว อาสวา วฑฺฒนฺติ. กตเมสํ ทฺวินฺนํ? โย จ อวินเย วินยสญฺญี, โย จ วินเย อวินยสญฺญี. อิเมสํ โข ภิกฺขเว ทฺวินฺนํ อาสวา วฑฺฒนฺตีติ. (19)}

\proseitemwithclosermidruled{118}{ทฺวินฺนํ ภิกฺขเว อาสวา น วฑฺฒนฺติ. กตเมสํ ทฺวินฺนํ? โย จ อวินเย อวินยสญฺญี, โย จ วินเย วินยสญฺญี. อิเมสํ โข ภิกฺขเว ทฺวินฺนํ อาสวา น วฑฺฒนฺตีติ.}{(20) พาลวคฺโค ปญฺจโม.}
\nitthitammajor{\textbf{ทุติโย ปณฺณาสโก สมตฺโต.}}
\chapterheadpageread{3. ตติย\-ปณฺณาสก}
\chapterhead{(11) 1. อาสา\-ทุปฺปชห\-วคฺค}
\proseitem{119}{\csromanfolio{86}เทฺวมา ภิกฺขเว อาสา ทุปฺปชหา. กตมา เทฺว? ลาภาสา จ ชีวิตาสา จ. อิมา โข ภิกฺขเว เทฺว อาสา ทุปฺปชหาติ. (1)}

\proseitem{120}{เทฺวเม ภิกฺขเว ปุคฺคลา ทุลฺลภา โลกสฺมิํ. กตเม เทฺว? โย จ ปุพฺพการี, โย จ กตญฺญู กตเวที. อิเม โข ภิกฺขเว เทฺว ปุคฺคลา ทุลฺลภา โลกสฺมินฺติ. (2)}

\proseitem{121}{เทฺวเม ภิกฺขเว ปุคฺคลา ทุลฺลภา โลกสฺมิํ. กตเม เทฺว? ติตฺโต จ ตปฺเปตา จ. อิเม โข ภิกฺขเว เทฺว ปุคฺคลา ทุลฺลภา โลกสฺมินฺติ. (3)}

\proseitem{122}{เทฺวเม ภิกฺขเว ปุคฺคลา ทุตฺตปฺปยา. กตเม เทฺว? โย จ ลทฺธํ ลทฺธํ นิกฺขิปติ, โย จ ลทฺธํ ลทฺธํ วิสฺสชฺเชติ. อิเม โข ภิกฺขเว เทฺว ปุคฺคลา ทุตฺตปฺปยาติ. (4)}

\proseitem{123}{เทฺวเม ภิกฺขเว ปุคฺคลา สุตปฺปยา. กตเม เทฺว? โย จ ลทฺธํ ลทฺธํ น นิกฺขิปติ, โย จ ลทฺธํ ลทฺธํ น วิสฺสชฺเชติ. อิเม โข ภิกฺขเว เทฺว ปุคฺคลา สุตปฺปยาติ. (5)}

\proseitem{124}{เทฺวเม ภิกฺขเว ปจฺจยา ราคสฺส อุปฺปาทาย. กตเม เทฺว? สุภนิมิตฺตญฺจ, อโยนิโส จ มนสิกาโร. อิเม โข ภิกฺขเว เทฺว ปจฺจยา ราคสฺส อุปฺปาทายาติ. (6)}

\proseitem{125}{เทฺวเม ภิกฺขเว ปจฺจยา โทสสฺส อุปฺปาทาย. กตเม เทฺว? ปฏิฆนิมิตฺตญฺจ, อโยนิโส จ มนสิกาโร. อิเม โข ภิกฺขเว เทฺว ปจฺจยา โทสสฺส อุปฺปาทายาติ. (7)}

\proseitem{126}{เทฺวเม ภิกฺขเว ปจฺจยา มิจฺฉาทิฏฺฐิยา อุปฺปาทาย. กตเม เทฺว? ปรโต จ โฆโส, อโยนิโส จ มนสิกาโร. อิเม โข ภิกฺขเว เทฺว ปจฺจยา มิจฺฉาทิฏฺฐิยา อุปฺปาทายาติ. (8)}

\chapterheadpageread{(12) 2. อายาจนวคฺค}
\proseitem{127}{\csromanfolio{87}เทฺวเม ภิกฺขเว ปจฺจยา สมฺมาทิฏฺฐิยา อุปฺปาทาย. กตเม เทฺว? ปรโต จ โฆโส, โยนิโส จ มนสิกาโร. อิเม โข ภิกฺขเว เทฺว ปจฺจยา สมฺมาทิฏฺฐิยา อุปฺปาทายาติ. (9)}

\proseitem{128}{เทฺวมา ภิกฺขเว อาปตฺติโย. กตมา เทฺว? ลหุกา จ อาปตฺติ, ครุกา จ อาปตฺติ. อิมา โข ภิกฺขเว เทฺว อาปตฺติโยติ. (10)}

\proseitem{129}{เทฺวมา ภิกฺขเว อาปตฺติโย. กตมา เทฺว? ทุฏฺฐุลฺลา จ อาปตฺติ, อทุฏฺฐุลฺลา จ อาปตฺติ. อิมา โข ภิกฺขเว เทฺว อาปตฺติโยติ. (11)}

\proseitemwithclosermidruled{130}{เทฺวมา ภิกฺขเว อาปตฺติโย. กตมา เทฺว? สาวเสสา จ อาปตฺติ, อนวเสสา จ อาปตฺติ. อิมา โข ภิกฺขเว เทฺว อาปตฺติโยติ. (12)}{อาสา\-ทุปฺปชห\-วคฺโค ปฐโม.}
\chapterhead{(12) 2. อายาจนวคฺค}
\proseitem{131}{สทฺโธ ภิกฺขเว ภิกฺขุ เอวํ สมฺมา อายาจมาโน อายาเจยฺย “ตาทิโส โหมิ, ยาทิสา สาริปุตฺต\-โมคฺคลฺลานา”ติ. เอสา ภิกฺขเว ตุลา เอตํ ปมาณํ มม สาวกานํ ภิกฺขูนํ ยทิทํ สาริปุตฺต\-โมคฺคลฺลานาติ. (1)}

\proseitem{132}{สทฺธา ภิกฺขเว ภิกฺขุนี เอวํ สมฺมา อายาจมานา อายาเจยฺย “ตาทิสี โหมิ, ยาทิสี เขมา จ ภิกฺขุนี อุปฺปลวณฺณา จา”ติ. เอสา ภิกฺขเว ตุลา เอตํ ปมาณํ มม สาวิกานํ ภิกฺขุนีนํ ยทิทํ เขมา จ ภิกฺขุนี อุปฺปลวณฺณา จาติ. (2)}

\proseitem{133}{สทฺโธ ภิกฺขเว อุปาสโก เอวํ สมฺมา อายาจมาโน อายาเจยฺย “ตาทิโส โหมิ, ยาทิโส จิตฺโต จ คหปติ, หตฺถโก จ อาฬวโก”ติ. เอสา ภิกฺขเว ตุลา เอตํ ปมาณํ มม สาวกานํ อุปาสกานํ ยทิทํ จิตฺโต จ คหปติ, หตฺถโก จ อาฬวโกติ. (3)}

\proseitem{134}{\csromanfolio{88}สทฺธา ภิกฺขเว อุปาสิกา เอวํ สมฺมา อายาจมานา อายาเจยฺย “ตาทิสี โหมิ, ยาทิสี ขุชฺชุตฺตรา จ อุปาสิกา, เวฬุกณฺฑกิยา\csromansharedfootnote{251f38876738b720}{เวฬุกณฺฑกี (อํ 2. 295 ปิฏฺเฐ; อํ 1. 483 ปิฏฺเฐปิ อาคตํ.)} จ นนฺทมาตา”ติ. เอสา ภิกฺขเว ตุลา เอตํ ปมาณํ มม สาวิกานํ อุปาสิกานํ ยทิทํ ขุชฺชุตฺตรา จ อุปาสิกา, เวฬุกณฺฑกิยา จ นนฺทมาตาติ. (4)}

\proseitem{135}{ทฺวีหิ ภิกฺขเว ธมฺเมหิ สมนฺนาคโต พาโล อพฺยตฺโต อสปฺปุริโส ขตํ อุปหตํ อตฺตานํ ปริหรติ, สาวชฺโช จ โหติ สานุวชฺโช จ วิญฺญูนํ, พหุญฺจ อปุญฺญํ ปสวติ. กตเมหิ ทฺวีหิ? อนนุวิจฺจ อปริโยคาเหตฺวา อวณฺณารหสฺส วณฺณํ ภาสติ, อนนุวิจฺจ อปริโยคาเหตฺวา วณฺณารหสฺส อวณฺณํ ภาสติ. อิเมหิ โข ภิกฺขเว ทฺวีหิ ธมฺเมหิ สมนฺนาคโต พาโล อพฺยตฺโต อสปฺปุริโส ขตํ อุปหตํ อตฺตานํ ปริหรติ, สาวชฺโช จ โหติ สานุวชฺโช จ วิญฺญูนํ, พหุญฺจ อปุญฺญํ ปสวตีติ.}

\prose{ทฺวีหิ ภิกฺขเว ธมฺเมหิ สมนฺนาคโต ปณฺฑิโต วิยตฺโต สปฺปุริโส อกฺขตํ อนุปหตํ อตฺตานํ ปริหรติ, อนวชฺโช จ โหติ อนนุวชฺโช จ วิญฺญูนํ, พหุญฺจ ปุญฺญํ ปสวติ. กตเมหิ ทฺวีหิ? อนุวิจฺจ ปริโยคาเหตฺวา อวณฺณารหสฺส อวณฺณํ ภาสติ, อนุวิจฺจ ปริโยคาเหตฺวา วณฺณารหสฺส วณฺณํ ภาสติ. อิเมหิ โข ภิกฺขเว ทฺวีหิ ธมฺเมหิ สมนฺนาคโต ปณฺฑิโต วิยตฺโต สปฺปุริโส อกฺขตํ อนุปหตํ อตฺตานํ ปริหรติ, อนวชฺโช จ โหติ อนนุวชฺโช จ วิญฺญูนํ, พหุญฺจ ปุญฺญํ ปสวตีติ. (5)}

\proseitem{136}{ทฺวีหิ ภิกฺขเว ธมฺเมหิ สมนฺนาคโต พาโล อพฺยตฺโต อสปฺปุริโส ขตํ อุปหตํ อตฺตานํ ปริหรติ, สาวชฺโช จ โหติ สานุวชฺโช จ วิญฺญูนํ, พหุญฺจ อปุญฺญํ ปสวติ. กตเมหิ ทฺวีหิ? อนนุวิจฺจ อปริโยคาเหตฺวา อปฺปสาทนีเย ฐาเน ปสาทํ อุปทํเสติ, อนนุวิจฺจ อปริโยคาเหตฺวา ปสาทนีเย ฐาเน อปฺปสาทํ อุปทํเสติ. อิเมหิ โข ภิกฺขเว ทฺวีหิ ธมฺเมหิ สมนฺนาคโต พาโล อพฺยตฺโต อสปฺปุริโส ขตํ อุปหตํ อตฺตานํ ปริหรติ, สาวชฺโช จ โหติ สานุวชฺโช จ วิญฺญูนํ, พหุญฺจ อปุญฺญํ ปสวตีติ.}

\chapterheadpageread{(12) 2. อายาจนวคฺค}
\prose{\csromanfolio{89}ทฺวีหิ ภิกฺขเว ธมฺเมหิ สมนฺนาคโต ปณฺฑิโต วิยตฺโต สปฺปุริโส อกฺขตํ อนุปหตํ อตฺตานํ ปริหรติ, อนวชฺโช จ โหติ อนนุวชฺโช จ วิญฺญูนํ, พหุญฺจ ปุญฺญํ ปสวติ. กตเมหิ ทฺวีหิ? อนุวิจฺจ ปริโยคาเหตฺวา อปฺปสาทนีเย ฐาเน อปฺปสาทํ อุปทํเสติ, อนุวิจฺจ ปริโยคาเหตฺวา ปสาทนีเย ฐาเน ปสาทํ อุปทํเสติ. อิเมหิ โข ภิกฺขเว ทฺวีหิ ธมฺเมหิ สมนฺนาคโต ปณฺฑิโต วิยตฺโต สปฺปุริโส อกฺขตํ อนุปหตํ อตฺตานํ ปริหรติ, อนวชฺโช จ โหติ อนนุวชฺโช จ วิญฺญูนํ, พหุญฺจ ปุญฺญํ ปสวตีติ. (6)}

\proseitem{137}{ทฺวีสุ ภิกฺขเว มิจฺฉา\-ปฏิปชฺชมาโน พาโล อพฺยตฺโต อสปฺปุริโส ขตํ อุปหตํ อตฺตานํ ปริหรติ, สาวชฺโช จ โหติ สานุวชฺโช จ วิญฺญูนํ, พหุญฺจ อปุญฺญํ ปสวติ. กตเมสุ ทฺวีสุ? มาตริ จ ปิตริ จ. อิเมสุ โข ภิกฺขเว ทฺวีสุ มิจฺฉา\-ปฏิปชฺชมาโน พาโล อพฺยตฺโต อสปฺปุริโส ขตํ อุปหตํ อตฺตานํ ปริหรติ, สาวชฺโช จ โหติ สานุวชฺโช จ วิญฺญูนํ, พหุญฺจ อปุญฺญํ ปสวตีติ.}

\prose{ทฺวีสุ ภิกฺขเว สมฺมา\-ปฏิปชฺชมาโน ปณฺฑิโต วิยตฺโต สปฺปุริโส อกฺขตํ อนุปหตํ อตฺตานํ ปริหรติ, อนวชฺโช จ โหติ อนนุวชฺโช จ วิญฺญูนํ, พหุญฺจ ปุญฺญํ ปสวติ. กตเมสุ ทฺวีสุ? มาตริ จ ปิตริ จ. อิเมสุ โข ภิกฺขเว ทฺวีสุ สมฺมา\-ปฏิปชฺชมาโน ปณฺฑิโต วิยตฺโต สปฺปุริโส อกฺขตํ อนุปหตํ อตฺตานํ ปริหรติ, อนวชฺโช จ โหติ อนนุวชฺโช จ วิญฺญูนํ, พหุญฺจ ปุญฺญํ ปสวตีติ. (7)}

\proseitem{138}{ทฺวีสุ ภิกฺขเว มิจฺฉา\-ปฏิปชฺชมาโน พาโล อพฺยตฺโต อสปฺปุริโส ขตํ อุปหตํ อตฺตานํ ปริหรติ, สาวชฺโช จ โหติ สานุวชฺโช จ วิญฺญูนํ, พหุญฺจ อปุญฺญํ ปสวติ. กตเมสุ ทฺวีสุ? ตถาคเต จ ตถาคต\-สาวเก จ. อิเมสุ โข ภิกฺขเว มิจฺฉา\-ปฏิปชฺชมาโน พาโล อพฺยตฺโต อสปฺปุริโส ขตํ อุปหตํ อตฺตานํ ปริหรติ, สาวชฺโช จ โหติ สานุวชฺโช จ วิญฺญูนํ, พหุญฺจ อปุญฺญํ ปสวตีติ.}

\prose{ทฺวีสุ ภิกฺขเว สมฺมา\-ปฏิปชฺชมาโน ปณฺฑิโต วิยตฺโต สปฺปุริโส อกฺขตํ อนุปหตํ อตฺตานํ ปริหรติ, อนวชฺโช จ โหติ อนนุวชฺโช จ วิญฺญูนํ, พหุญฺจ ปุญฺญํ ปสวติ. กตเมสุ ทฺวีสุ? ตถาคเต จ \csromanfolio{90}ตถาคต\-สาวเก จ. อิเมสุ โข ภิกฺขเว ทฺวีสุ สมฺมา\-ปฏิปชฺชมาโน ปณฺฑิโต วิยตฺโต สปฺปุริโส อกฺขตํ อนุปหตํ อตฺตานํ ปริหรติ, อนวชฺโช จ โหติ อนนุวชฺโช จ วิญฺญูนํ, พหุญฺจ ปุญฺญํ ปสวตีติ. (8)}

\proseitem{139}{เทฺวเม ภิกฺขเว ธมฺมา. กตเม เทฺว? สจิตฺตโวทานญฺจ น จ กิญฺจิ โลเก อุปาทิยติ. อิเม โข ภิกฺขเว เทฺว ธมฺมาติ. (9)}

\proseitem{140}{เทฺวเม ภิกฺขเว ธมฺมา. กตเม เทฺว? โกโธ จ อุปนาโห จ. อิเม โข ภิกฺขเว เทฺว ธมฺมาติ. (10)}

\proseitemwithclosermidruled{141}{เทฺวเม ภิกฺขเว ธมฺมา. กตเม เทฺว? โกธวินโย จ อุปนาหวินโย จ. อิเม โข ภิกฺขเว เทฺว ธมฺมาติ. (11)}{อายาจนวคฺโค ทุติโย.}
\chapterhead{\textbf{(13) 3. ทานวคฺค}}
\proseitem{142}{เทฺวมานิ ภิกฺขเว ทานานิ. กตมานิ เทฺว? อามิสทานญฺจ ธมฺมทานญฺจ. อิมานิ โข ภิกฺขเว เทฺว ทานานิ. เอตทคฺคํ ภิกฺขเว อิเมสํ ทฺวินฺนํ ทานานํ, ยทิทํ ธมฺมทานนฺติ. (1)}

\proseitem{143}{เทฺวเม ภิกฺขเว ยาคา. กตเม เทฺว? อามิสยาโค จ ธมฺมยาโค จ. อิเม โข ภิกฺขเว เทฺว ยาคา. เอตทคฺคํ ภิกฺขเว อิเมสํ ทฺวินฺนํ ยาคานํ, ยทิทํ ธมฺมยาโคติ. (2)}

\proseitem{144}{เทฺวเม ภิกฺขเว จาคา. กตเม เทฺว? อามิสจาโค จ ธมฺมจาโค จ. อิเม โข ภิกฺขเว เทฺว จาคา. เอตทคฺคํ ภิกฺขเว อิเมสํ ทฺวินฺนํ จาคานํ, ยทิทํ ธมฺมจาโคติ. (3)}

\proseitem{145}{เทฺวเม ภิกฺขเว ปริจฺจาคา. กตเม เทฺว? อามิสปริจฺจาโค จ ธมฺม\-ปริจฺจาโค จ. อิเม โข ภิกฺขเว เทฺว ปริจฺจาคา. เอตทคฺคํ ภิกฺขเว อิเมสํ ทฺวินฺนํ ปริจฺจาคานํ, ยทิทํ ธมฺม\-ปริจฺจาโคติ. (4)}

\chapterheadpageread{(14) 4. สนฺถารวคฺค}
\proseitem{146}{\csromanfolio{91}เทฺวเม ภิกฺขเว โภคา. กตเม เทฺว? อามิสโภโค จ ธมฺมโภโค จ. อิเม โข ภิกฺขเว เทฺว โภคา. เอตทคฺคํ ภิกฺขเว อิเมสํ ทฺวินฺนํ โภคานํ, ยทิทํ ธมฺมโภโคติ. (5)}

\proseitem{147}{เทฺวเม ภิกฺขเว สมฺโภคา. กตเม เทฺว? อามิสสมฺโภโค จ ธมฺมสมฺโภโค จ. อิเม โข ภิกฺขเว เทฺว สมฺโภคา. เอตทคฺคํ ภิกฺขเว อิเมทํ ทฺวินฺนํ สมฺโภคานํ, ยทิทํ ธมฺม\-สมฺโภโคติ. (6)}

\proseitem{148}{เทฺวเม ภิกฺขเว สํวิภาคา. กตเม เทฺว? อามิส\-สํวิภาโค จ ธมฺม\-สํวิภาโค จ. อิเม โข ภิกฺขเว เทฺว สํวิภาคา. เอตทคฺคํ ภิกฺขเว อิเมสํ ทฺวินฺนํ สํวิภาคานํ, ยทิทํ ธมฺม\-สํวิภาโคติ. (7)}

\proseitem{149}{เทฺวเม ภิกฺขเว สงฺคหา. กตเม เทฺว? อามิสสงฺคโห จ ธมฺมสงฺคโห จ. อิเม โข ภิกฺขเว เทฺว สงฺคหา. เอตทคฺคํ ภิกฺขเว อิเมสํ ทฺวินฺนํ สงฺคหานํ, ยทิทํ ธมฺม\-สงฺคโหติ. (8)}

\proseitem{150}{เทฺวเม ภิกฺขเว อนุคฺคหา. กตเม เทฺว? อามิสานุคฺคโห จ ธมฺมานุคฺคโห จ. อิเม โข ภิกฺขเว เทฺว อนุคฺคหา. เอตทคฺคํ ภิกฺขเว อิเมสํ ทฺวินฺนํ อนุคฺคหานํ, ยทิทํ ธมฺมา\-นุคฺคโหติ. (9)}

\proseitemwithclosermidruled{151}{เทฺวมา ภิกฺขเว อนุกมฺปา. กตมา เทฺว? อามิสานุกมฺปา จ ธมฺมานุกมฺปา จ. อิมา โข ภิกฺขเว เทฺว อนุกมฺปา. เอตทคฺคํ ภิกฺขเว อิมาสํ ทฺวินฺนํ อนุกมฺปานํ, ยทิทํ ธมฺมานุกมฺปาติ. (10)}{ทานวคฺโค ตติโย.}
\chapterhead{(14) 4. สนฺถารวคฺค}
\proseitem{152}{เทฺวเม ภิกฺขเว สนฺถารา\csromansharedfootnote{deb333f47cb4cdd6}{สนฺธารา (ก)}. กตเม เทฺว? อามิสสนฺถาโร จ ธมฺมสนฺถาโร จ. อิเม โข ภิกฺขเว เทฺว สนฺถารา. เอตทคฺคํ ภิกฺขเว อิเมสํ ทฺวินฺนํ สนฺถารานํ, ยทิทํ ธมฺม\-สนฺถาโรติ. (1)}

\proseitem{153}{\csromanfolio{92}เทฺวเม ภิกฺขเว ปฏิสนฺถารา\csromansharedfootnote{632d3b37ce9c92a8}{ปฏิสนฺธารา (ก)}. กตเม เทฺว? อามิส\-ปฏิสนฺถาโร จ ธมฺม\-ปฏิสนฺถาโร จ. อิเม โข ภิกฺขเว เทฺว ปฏิสนฺถารา. เอตทคฺคํ ภิกฺขเว อิเมสํ ทฺวินฺนํ ปฏิสนฺถารานํ, ยทิทํ ธมฺม\-ปฏิสนฺถาโรติ. (2)}

\proseitem{154}{เทฺวมา ภิกฺขเว เอสนา. กตมา เทฺว? อามิเสสนา จ ธมฺเมสนา จ. อิมา โข ภิกฺขเว เทฺว เอสนา. เอตทคฺคํ ภิกฺขเว อิมาสํ ทฺวินฺนํ เอสนานํ, ยทิทํ ธมฺเมสนาติ. (3)}

\proseitem{155}{เทฺวมา ภิกฺขเว ปริเยสนา. กตมา เทฺว? อามิส\-ปริเยสนา จ ธมฺม\-ปริเยสนา จ. อิมา โข ภิกฺขเว เทฺว ปริเยสนา. เอตทคฺคํ ภิกฺขเว อิมาสํ ทฺวินฺนํ ปริเยสนานํ, ยทิทํ ธมฺม\-ปริเยสนาติ. (4)}

\proseitem{156}{เทฺวมา ภิกฺขเว ปริเยฏฺฐิโย. กตมา เทฺว? อามิส\-ปริเยฏฺฐิ จ ธมฺม\-ปริเยฏฺฐิ จ. อิมา โข ภิกฺขเว เทฺว ปริเยฏฺฐิโย. เอตทคฺคํ ภิกฺขเว อิมาสํ ทฺวินฺนํ ปริเยฏฺฐีนํ, ยทิทํ ธมฺมปริเยฏฺฐีติ. (5)}

\proseitem{157}{เทฺวมา ภิกฺขเว ปูชา. กตมา เทฺว? อามิสปูชา จ ธมฺมปูชา จ. อิมา โข ภิกฺขเว เทฺว ปูชา. เอตทคฺคํ ภิกฺขเว อิมาสํ ทฺวินฺนํ ปูชานํ, ยทิทํ ธมฺมปูชาติ. (6)}

\proseitem{158}{เทฺวมานิ ภิกฺขเว อาติเถยฺยานิ. กตมานิ เทฺว? อามิสาติเถยฺยญฺจ ธมฺมาติเถยฺยญฺจ. อิมานิ โข ภิกฺขเว เทฺว อาติเถยฺยานิ. เอตทคฺคํ ภิกฺขเว อิเมสํ ทฺวินฺนํ อาติเถยฺยานํ, ยทิทํ ธมฺมา\-ติเถยฺยนฺติ. (7)}

\proseitem{159}{เทฺวมา ภิกฺขเว อิทฺธิโย. กตมา เทฺว? อามิสิทฺธิ จ ธมฺมิทฺธิ จ. อิมา โข ภิกฺขเว เทฺว อิทฺธิโย. เอตทคฺคํ ภิกฺขเว อิมาสํ ทฺวินฺนํ อิทฺธีนํ, ยทิทํ ธมฺมิทฺธีติ. (8)}

\proseitem{160}{เทฺวมา ภิกฺขเว วุทฺธิโย. กตมา เทฺว? อามิสวุทฺธิ จ ธมฺมวุทฺธิ จ. อิมา โข ภิกฺขเว เทฺว วุทฺธิโย. เอตทคฺคํ ภิกฺขเว อิมาสํ ทฺวินฺนํ วุทฺธีนํ, ยทิทํ ธมฺมวุทฺธีติ. (9)}

\chapterheadpageread{(15) 5. สมาปตฺติวคฺค}
\proseitem{161}{\csromanfolio{93}เทฺวมานิ ภิกฺขเว รตนานิ. กตมานิ เทฺว? อามิสรตนญฺจ ธมฺมรตนญฺจ. อิมานิ โข ภิกฺขเว เทฺว รตนานิ. เอตทคฺคํ ภิกฺขเว อิเมสํ ทฺวินฺนํ รตนานํ, ยทิทํ ธมฺม\-รตนนฺติ. (10)}

\proseitem{162}{เทฺวเม ภิกฺขเว สนฺนิจยา. กตเม เทฺว? อามิสสนฺนิจโย จ ธมฺม\-สนฺนิจโย จ. อิเม โข ภิกฺขเว เทฺว สนฺนิจยา. เอตทคฺคํ ภิกฺขเว อิเมสํ ทฺวินฺนํ สนฺนิจยานํ, ยทิทํ ธมฺม\-สนฺนิจโยติ. (11)}

\proseitemwithclosermidruled{163}{เทฺวมานิ ภิกฺขเว เวปุลฺลานิ. กตมานิ เทฺว? อามิสเวปุลฺลญฺจ ธมฺมเวปุลฺลญฺจ. อิมานิ โข ภิกฺขเว เทฺว เวปุลฺลานิ. เอตทคฺคํ ภิกฺขเว อิเมสํ ทฺวินฺนํ เวปุลฺลานํ, ยทิทํ ธมฺม\-เวปุลฺลนฺติ. (12)}{สนฺถารวคฺโค จตุตฺโถ.}
\chapterhead{(15) 5. สมาปตฺติวคฺค}
\proseitem{164}{เทฺวเม ภิกฺขเว ธมฺมา. กตเม เทฺว? สมาปตฺติ\-กุสลตา จ สมาปตฺติ\-วุฏฺฐาน\-กุสลตา จ. อิเม โข ภิกฺขเว เทฺว ธมฺมาติ. (1)}

\proseitem{165}{เทฺวเม ภิกฺขเว ธมฺมา. กตเม เทฺว? อชฺชวญฺจ มทฺทวญฺจ. อิเม โข ภิกฺขเว เทฺว ธมฺมา. (2)}

\proseitem{166}{เทฺวเม ภิกฺขเว ธมฺมา. กตเม เทฺว? ขนฺติ จ โสรจฺจญฺจ. อิเม โข ภิกฺขเว เทฺว ธมฺมา. (3)}

\proseitem{167}{เทฺวเม ภิกฺขเว ธมฺมา. กตเม เทฺว? สาขลฺยญฺจ ปฏิสนฺถาโร จ. อิเม โข ภิกฺขเว เทฺว ธมฺมา. (4)}

\proseitem{168}{เทฺวเม ภิกฺขเว ธมฺมา. กตเม เทฺว? อวิหิํสา จ โสเจยฺยญฺจ. อิเม โข ภิกฺขเว เทฺว ธมฺมา. (5)}

\proseitem{169}{เทฺวเม ภิกฺขเว ธมฺมา. กตเม เทฺว? อินฺทฺริเยสุ อคุตฺตทฺวารตา จ, โภชเน อมตฺตญฺญุตา จ. อิเม โข ภิกฺขเว เทฺว ธมฺมา. (6)}

\proseitem{170}{\csromanfolio{94}เทฺวเม ภิกฺขเว ธมฺมา. กตเม เทฺว? อินฺทฺริเยสุ คุตฺตทฺวารตา จ, โภชเน มตฺตญฺญุตา จ. อิเม โข ภิกฺขเว เทฺว ธมฺมา. (7)}

\proseitem{171}{เทฺวเม ภิกฺขเว ธมฺมา. กตเม เทฺว? ปฏิสงฺขานพลญฺจ ภาวนาพลญฺจ. อิเม โข ภิกฺขเว เทฺว ธมฺมา. (8)}

\proseitem{172}{เทฺวเม ภิกฺขเว ธมฺมา. กตเม เทฺว? สติพลญฺจ สมาธิพลญฺจ. อิเม โข ภิกฺขเว เทฺว ธมฺมา. (9)}

\proseitem{173}{เทฺวเม ภิกฺขเว ธมฺมา. กตเม เทฺว? สมโถ จ วิปสฺสนา จ. อิเม โข ภิกฺขเว เทฺว ธมฺมา. (10)}

\proseitem{174}{เทฺวเม ภิกฺขเว ธมฺมา. กตเม เทฺว? สีลวิปตฺติ จ ทิฏฺฐิวิปตฺติ จ. อิเม โข ภิกฺขเว เทฺว ธมฺมา. (11)}

\proseitem{175}{เทฺวเม ภิกฺขเว ธมฺมา. กตเม เทฺว? สีลสมฺปทา จ ทิฏฺฐิสมฺปทา จ. อิเม โข ภิกฺขเว เทฺว ธมฺมา. (12)}

\proseitem{176}{เทฺวเม ภิกฺขเว ธมฺมา. กตเม เทฺว? สีลวิสุทฺธิ จ ทิฏฺฐิวิสุทฺธิ จ. อิเม โข ภิกฺขเว เทฺว ธมฺมา. (13)}

\proseitem{177}{เทฺวเม ภิกฺขเว ธมฺมา. กตเม เทฺว? ทิฏฺฐิวิสุทฺธิ จ, ยถาทิฏฺฐิสฺส จ ปธานํ. อิเม โข ภิกฺขเว เทฺว ธมฺมา. (14)}

\proseitem{178}{เทฺวเม ภิกฺขเว ธมฺมา. กตเม เทฺว? อสนฺตุฏฺฐิตา จ กุสเลสุ ธมฺเมสุ, อปฺปฏิวานิตา จ ปธานสฺมิํ. อิเม โข ภิกฺขเว เทฺว ธมฺมา. (15)}

\proseitem{179}{เทฺวเม ภิกฺขเว ธมฺมา. กตเม เทฺว? มุฏฺฐสฺสจฺจญฺจ อสมฺปชญฺญญฺจ. อิเม โข ภิกฺขเว เทฺว ธมฺมา. (16)}

\proseitemwithclosermidruled{180}{เทฺวเม ภิกฺขเว ธมฺมา, กตเม เทฺว? สติ จ สมฺปชญฺญญฺจ. อิเม โข ภิกฺขเว เทฺว ธมฺมาติ.}{(17) สมาปตฺติวคฺโค ปญฺจโม.}
\nitthitammajor{\textbf{ตติโย ปณฺณาสโก สมตฺโต.}}
\csromanheader{2}{1. โกธเปยฺยาล \textbf{1. โกธเปยฺยาล}}
\proseitem{181}{\csromanfolio{95}เทฺวเม ภิกฺขเว ธมฺมา. กตเม เทฺว? โกโธ จ อุปนาโห จ. มกฺโข จ ปฬาโส\csromansharedfootnote{88cc1b9f108bfc20}{ปลาโส (ก)} จ. อิสฺสา จ มจฺฉริยญฺจ. มายา จ สาเฐยฺยญฺจ. อหิริกญฺจ อโนตฺตปฺปญฺจ. อิเม โข ภิกฺขเว เทฺว ธมฺมา. (\csromansharedfootnote{88cc1b9f108bfc20}{ปลาโส (ก)}-5)}

\proseitem{182}{เทฺวเม ภิกฺขเว ธมฺมา. กตเม เทฺว? อกฺโกโธ จ อนุปนาโห จ. อมกฺโข จ อปฬาโส จ. อนิสฺสา จ อมจฺฉริยญฺจ. อมายา จ อสาเฐยฺยญฺจ. หิรี จ โอตฺตปฺปญฺจ. อิเม โข ภิกฺขเว เทฺว ธมฺมา. \mbox{(6-10)}}

\proseitem{183}{ทฺวีหิ ภิกฺขเว ธมฺเมหิ สมนฺนาคโต ทุกฺขํ วิหรติ. กตเมหิ ทฺวีหิ? โกเธน จ อุปนาเหน จ. มกฺเขน จ ปฬาเสน จ. อิสฺสาย จ มจฺฉริเยน จ. มายาย จ สาเฐยฺเยน จ. อหิริเกน จ อโนตฺตปฺเปน จ. อิเมหิ โข ภิกฺขเว ทฺวีหิ ธมฺเมหิ สมนฺนาคโต ทุกฺขํ วิหรติ. \mbox{(11-15)}}

\proseitem{184}{ทฺวีหิ ภิกฺขเว ธมฺเมหิ สมนฺนาคโต สุขํ วิหรติ. กตเมหิ ทฺวีหิ? อกฺโกเธน จ อนุปนาเหน จ. อมกฺเขน จ อปฬาเสน จ. อนิสฺสาย จ อมจฺฉริเยน จ. อมายาย จ อสาเฐยฺเยน จ. หิริยา จ โอตฺตปฺเปน จ. อิเมหิ โข ภิกฺขเว ทฺวีหิ ธมฺเมหิ สมนฺนาคโต สุขํ วิหรติ. \mbox{(16-20)}}

\proseitem{185}{เทฺวเม ภิกฺขเว ธมฺมา เสขสฺส ภิกฺขุโน ปริหานาย สํวตฺตนฺติ. กตเม เทฺว? โกโธ จ อุปนาโห จ. มกฺโข จ ปฬาโส จ. อิสฺสา จ มจฺฉริยญฺจ. มายา จ สาเฐยฺยญฺจ. อหิริกญฺจ อโนตฺตปฺปญฺจ. อิเม โข ภิกฺขเว เทฺว ธมฺมา เสขสฺส ภิกฺขุโน ปริหานาย สํวตฺตนฺติ. \mbox{(21-25)}}

\proseitem{186}{เทฺวเม ภิกฺขเว ธมฺมา เสขสฺส ภิกฺขุโน อปริหานาย สํวตฺตนฺติ. กตเม เทฺว? อกฺโกโธ จ อนุปนาโห จ. อมกฺโข จ อปฬาโส จ. อนิสฺสา จ อมจฺฉริยญฺจ. อมายา จ อสาเฐยฺยญฺจ. หิรี จ โอตฺตปฺปญฺจ. อิเม โข ภิกฺขเว เทฺว ธมฺมา เสขสฺส ภิกฺขุโน อปริหานาย สํวตฺตนฺติ. \mbox{(26-30)}}

\proseitem{187}{ทฺวีหิ ภิกฺขเว ธมฺเมหิ สมนฺนาคโต ยถาภตํ นิกฺขิตฺโต, เอวํ นิรเย. กตเมหิ ทฺวีหิ? โกเธน จ อุปนาเหน จ. มกฺเขน จ ปฬาเสน จ. อิสฺสาย จ มจฺฉริเยน จ. มายาย จ สาเฐยฺเยน จ. อหิริเกน จ อโนตฺตปฺเปน จ. อิเมหิ โข ภิกฺขเว ทฺวีหิ ธมฺเมหิ สมนฺนาคโต ยถาภตํ นิกฺขิตฺโต, เอวํ นิรเย. \mbox{(31-35)}}

\proseitem{188}{\csromanfolio{96}ทฺวีหิ ภิกฺขเว ธมฺเมหิ สมนฺนาคโต ยถาภตํ นิกฺขิตฺโต, เอวํ สคฺเค. กตเมหิ ทฺวีหิ? อกฺโกเธน จ อนุปนาเหน จ. อมกฺเขน จ อปฬาเสน จ. อนิสฺสาย จ อมจฺฉริเยน จ. อมายาย จ อสาเฐยฺเยน จ. หิริยา จ โอตฺตปฺเปน จ. อิเมหิ โข ภิกฺขเว ทฺวีหิ ธมฺเมหิ สมนฺนาคโต ยถาภตํ นิกฺขิตฺโต, เอวํ สคฺเค. \mbox{(36-40)}}

\proseitem{189}{ทฺวีหิ ภิกฺขเว ธมฺเมหิ สมนฺนาคโต อิเธกจฺโจ กายสฺส เภทา ปรํ มรณา อปายํ ทุคฺคติํ วินิปาตํ นิรยํ อุปปชฺชติ. กตเมหิ ทฺวีหิ? โกเธน จ อุปนาเหน จ. มกฺเขน จ ปฬาเสน จ. อิสฺสาย จ มจฺฉริเยน จ. มายาย จ สาเฐยฺเยน จ. อหิริเกน จ อโนตฺตปฺเปน จ. อิเมหิ โข ภิกฺขเว ทฺวีหิ ธมฺเมหิ สมนฺนาคโต อิเธกจฺโจ กายสฺส เภทา ปรํ มรณา อปายํ ทุคฺคติํ วินิปาตํ นิรยํ อุปปชฺชติ. \mbox{(41-45)}}

\proseitemwithcloser{190}{ทฺวีหิ ภิกฺขเว ธมฺเมหิ สมนฺนาคโต อิเธกจฺโจ กายสฺส เภทา ปรํ มรณา สุคติํ สคฺคํ โลกํ อุปปชฺชติ. กตเมหิ ทฺวีหิ? อกฺโกเธน จ อนุปนาเหน จ. อมกฺเขน จ อปฬาเสน จ. อนิสฺสาย จ อมจฺฉริเยน จ. อมายาย จ อสาเฐยฺเยน จ. หิริยา จ โอตฺตปฺเปน จ. อิเมหิ โข ภิกฺขเว ทฺวีหิ ธมฺเมหิ สมนฺนาคโต อิเธกจฺโจ กายสฺส เภทา ปรํ มรณา สุคติํ สคฺคํ โลกํ อุปปชฺชติ. \mbox{(46-50)}}{โกธเปยฺยาลํ นิฏฺฐิตํ.}
\csromanheader{2}{2. \textbf{อกุสลเปยฺยาล}}
\proseitemwithcloser{191-200}{เทฺวเม ภิกฺขเว ธมฺมา อกุสลา. เทฺวเม ภิกฺขเว ธมฺมา กุสลา. เทฺวเม ภิกฺขเว ธมฺมา สาวชฺชา. เทฺวเม ภิกฺขเว ธมฺมา อนวชฺชา. เทฺวเม ภิกฺขเว ธมฺมา ทุกฺขุทฺรยา. เทฺวเม ภิกฺขเว ธมฺมา สุขุทฺรยา. เทฺวเม ภิกฺขเว ธมฺมา ทุกฺขวิปากา. เทฺวเม ภิกฺขเว ธมฺมา สุขวิปากา. เทฺวเม ภิกฺขเว ธมฺมา สพฺยาพชฺฌา. เทฺวเม ภิกฺขเว ธมฺมา อพฺยาพชฺฌา. กตเม เทฺว? อกฺโกโธ จ อนุปนาโห จ. อมกฺโข จ อปฬาโส จ. อนิสฺสา จ อมจฺฉริยญฺจ. อมายา จ อสาเฐยฺยญฺจ. หิรี จ โอตฺตปฺปญฺจ. อิเม โข ภิกฺขเว เทฺว ธมฺมา อพฺยาพชฺฌาติ. \mbox{(1-50)}}{อกุสล\-เปยฺยาลํ นิฏฺฐิตํ.}
\csromanheader{2}{3. วินยเปยฺยาล \textbf{3. วินยเปยฺยาล}}
\proseitem{201}{\csromanfolio{97}เทฺวเม ภิกฺขเว อตฺถวเส ปฏิจฺจ ตถาคเตน สาวกานํ สิกฺขาปทํ ปญฺญตฺตํ. กตเม เทฺว? สํฆสุฏฺฐุตาย, สํฆผาสุตาย. ทุมฺมงฺกูนํ ปุคฺคลานํ นิคฺคหาย, เปสลานํ ภิกฺขูนํ ผาสุวิหาราย. ทิฏฺฐ\-ธมฺมิกานํ อาสวานํ สํวราย, สมฺปรายิกานํ อาสวานํ ปฏิฆาตาย. ทิฏฺฐ\-ธมฺมิกานํ เวรานํ สํวราย, สมฺปรายิกานํ เวรานํ ปฏิฆาตาย. ทิฏฺฐ\-ธมฺมิกานํ วชฺชานํ สํวราย, สมฺปรายิกานํ วชฺชานํ ปฏิฆาตาย. ทิฏฺฐ\-ธมฺมิกานํ ภยานํ สํวราย, สมฺปรายิกานํ ภยานํ ปฏิฆาตาย. ทิฏฺฐ\-ธมฺมิกานํ อกุสลานํ ธมฺมานํ สํวราย, สมฺปรายิกานํ อกุสลานํ ธมฺมานํ ปฏิฆาตาย. คิหีนํ อนุกมฺปาย, ปาปิจฺฉานํ ภิกฺขูนํ ปกฺขุ\-ปจฺเฉทาย. อปฺปสนฺนานํ ปสาทาย, ปสนฺนานํ ภิยฺโยภาวาย. สทฺธมฺม\-ฏฺฐิติยา, วินยา\-นุคฺคหาย. อิเม โข ภิกฺขเว เทฺว อตฺถวเส ปฏิจฺจ ตถาคเตน สาวกานํ สิกฺขาปทํ ปญฺญตฺตนฺติ. \mbox{(1-10)}}

\proseitemwithcloser{202-230}{เทฺวเม ภิกฺขเว อตฺถวเส ปฏิจฺจ ตถาคเตน สาวกานํ ปาติโมกฺขํ ปญฺญตฺตํ ฯเปฯ ปาติโมกฺขุ\-ทฺเทโส ปญฺญตฺโต. ปาติโมกฺข\-ฏฺฐปนํ ปญฺญตฺตํ. ปวารณา ปญฺญตฺตา. ปวารณ\-ฏฺฐปนํ ปญฺญตฺตํ. ตชฺชนีย\-กมฺมํ ปญฺญตฺตํ. นิยสฺสกมฺมํ ปญฺญตฺตํ. ปพฺพาชนีย\-กมฺมํ ปญฺญตฺตํ. ปฏิสารณีย\-กมฺมํ ปญฺญตฺตํ. อุกฺเขปนีย\-กมฺมํ ปญฺญตฺตํ. ปริวาสทานํ ปญฺญตฺตํ. มูลาย ปฏิกสฺสนํ ปญฺญตฺตํ. มานตฺตทานํ ปญฺญตฺตํ. อพฺภานํ ปญฺญตฺตํ. โอสารณียํ ปญฺญตฺตํ. นิสฺสารณียํ ปญฺญตฺตํ. อุปสมฺปทา ปญฺญตฺตา. ญตฺติกมฺมํ ปญฺญตฺตํ. ญตฺติ\-ทุติย\-กมฺมํ ปญฺญตฺตํ. ญตฺติ\-จตุตฺถ\-กมฺมํ ปญฺญตฺตํ. อปญฺญตฺเต ปญฺญตฺตํ. ปญฺญตฺเต อนุปญฺญตฺตํ. สมฺมุขาวินโย ปญฺญตฺโต. สติวินโย ปญฺญตฺโต. อมูฬฺหวินโย ปญฺญตฺโต. ปฏิญฺญาต\-กรณํ ปญฺญตฺตํ. เยภุยฺยสิกา ปญฺญตฺตา. ตสฺสปาปิยสิกา ปญฺญตฺตา. ติณวตฺถารโก ปญฺญตฺโต. กตเม เทฺว? สํฆสุฏฺฐุตาย, สํฆผาสุตาย. ทุมฺมงฺกูนํ ปุคฺคลานํ นิคฺคหาย, เปสลานํ ภิกฺขูนํ ผาสุวิหาราย. ทิฏฺฐ\-ธมฺมิกานํ อาสวานํ สํวราย, สมฺปรายิกานํ อาสวานํ ปฏิฆาตาย. ทิฏฺฐ\-ธมฺมิกานํ เวรานํ สํวราย, สมฺปรายิกานํ เวรานํ ปฏิฆาตาย. ทิฏฺฐ\-ธมฺมิกานํ วชฺชานํ สํวราย, สมฺปรายิกานํ วชฺชานํ ปฏิฆาตาย. ทิฏฺฐ\-ธมฺมิกานํ ภยานํ สํวราย, สมฺปรายิกานํ ภยานํ ปฏิฆาตาย. ทิฏฺฐ\-ธมฺมิกานํ อกุสลานํ ธมฺมานํ สํวราย, สมฺปรายิกานํ อกุสลานํ ธมฺมานํ ปฏิฆาตาย. คิหีนํ อนุกมฺปาย, ปาปิจฺฉานํ ภิกฺขูนํ ปกฺขุ\-ปจฺเฉทาย. \csromanfolio{98}อปฺปสนฺนานํ ปสาทาย, ปสนฺนานํ ภิยฺโยภาวาย. สทฺธมฺม\-ฏฺฐิติยา วินยา\-นุคฺคหาย. อิเม โข ภิกฺขเว เทฺว อตฺถวเส ปฏิจฺจ ตถาคเตน สาวกานํ ติณวตฺถารโก ปญฺญตฺโตติ. \mbox{(11-300)}}{วินเปยฺยาลํ นิฏฺฐิตํ.}
\csromanmatikaanchor{mtk.2}
\csromanmatikaanchor{mtk.286}
\csromantocmarknopage{section}{(11) 1. สมฺโพธวคฺค}{mtk.2}
\bookmark[dest={mtk.2},level=1]{(11) 1. สมฺโพธวคฺค}
\csromanheader{1}{4. \textbf{ราคเปยฺยาล}}
\proseitem{231}{ราคสฺส ภิกฺขเว อภิญฺญาย เทฺว ธมฺมา ภาเวตพฺพา. กตเม เทฺว? สมโถ จ วิปสฺสนา จ. ราคสฺส ภิกฺขเว อภิญฺญาย อิเม เทฺว ธมฺมา ภาเวตพฺพาติ. (1)}

\prose{ราคสฺส ภิกฺขเว ปริญฺญาย. ปริกฺขยาย. ปหานาย. ขยาย. วยาย. วิราคาย. นิโรธาย. จาคาย. ปฏินิสฺสคฺคาย เทฺว ธมฺมา ภาเวตพฺพา ฯเปฯ. (\csromansharedfootnote{58456e94fd777360}{( ) เอตฺถนฺตเร ปาโฐ สีสฺยากํอิโปตฺถเกสุ นตฺถิ.}-10)}

\proseitem{232-246}{โทสสฺส ฯเปฯ. โมหสฺส. โกธสฺส. อุปนาหาสฺส. มกฺขสฺส. ปฬาสสฺส. อิสฺสาย. มจฺฉริยสฺส. มายาย. สาเฐยฺยสฺส. ถมฺภสฺส. สารมฺภสฺส. มานสฺส. อติมานสฺส. มทสฺส. ปมาทสฺส. อภิญฺญาย. ปริญฺญาย. ปริกฺขยาย. ปหานาย. ขยาย. วยาย. วิราคาย. นิโรธาย. จาคาย. ปฏินิสฺสคฺคาย เทฺว ธมฺมา ภาเวตพฺพา. กตเม เทฺว, สมโถ จ วิปสฺสนา จ. ปมาทสฺส ภิกฺขเว ปฏินิสฺสคฺคาย อิเม เทฺว ธมฺมา ภาเวตพฺพาติ. \mbox{(11-170)}}

\prosewithcloser{(อิทมโวจ ภควา, อตฺตมนา เต ภิกฺขู ภควโต ภาสิตํ อภินนฺทุนฺติ.)\csromansharedfootnote{58456e94fd777360}{( ) เอตฺถนฺตเร ปาโฐ สีสฺยากํอิโปตฺถเกสุ นตฺถิ.}}{ราคเปยฺยาลํ นิฏฺฐิตํ.}
\nitthitammajor{ทุกนิปาตปาฬิ นิฏฺฐิตา.}
\csromanpalirecto
\gambhira{\textbf{ติกนิปาตปาฬิ}}
\namakkaram{นโม ตสฺส ภควโต อรหโต สมฺมา\-สมฺพุทฺธสฺส.}
\csromanheader{2}{1. ปฐม\-ปณฺณาสก}
\chapterhead{1. \textbf{พาลวคฺค}}
\csromanheader{2}{1. \textbf{ภยสุตฺต}}
\proseitem{1}{\csromanfolio{99}เอวํ เม สุตํ\csromandash{}เอกํ สมยํ ภควา สาวตฺถิยํ วิหรติ เชตวเน อนาถ\-ปิณฺฑิกสฺส อาราเม. ตตฺร โข ภควา ภิกฺขู อามนฺเตสิ “ภิกฺขโว”ติ. “ภทนฺเต\csromansharedfootnote{5fd99653347407b8}{ภทฺทนฺเต (ก)}”ติ เต ภิกฺขู ภควโต ปจฺจสฺโสสุํ. ภควา เอตทโวจ\csromandash{}}

\prose{ยานิ กานิจิ ภิกฺขเว ภยานิ อุปฺปชฺชนฺติ, สพฺพานิ ตานิ พาลโต อุปฺปชฺชนฺติ, โน ปณฺฑิตโต. เย เกจิ อุปทฺทวา อุปฺปชฺชนฺติ, สพฺเพ เต พาลโต อุปฺปชฺชนฺติ, โน ปณฺฑิตโต. เย เกจิ อุปสคฺคา อุปฺปชฺชนฺติ, สพฺเพ เต พาลโต อุปฺปชฺชนฺติ, โน ปณฺฑิตโต. เสยฺยถาปิ ภิกฺขเว นฬาคารา วา ติณาคารา วา\csromansharedfootnote{d838b8fe93ea3309}{นฬาคารํ วา ติณาคารํ วา (สี)} อคฺคิ มุตฺโต\csromansharedfootnote{2c91271c0a14e740}{อคฺคิมุกฺโก (สี), อคฺคิ มุกฺโก (สฺยา, กํ, อิ)} กูฏาคารานิปิ ฑหติ อุลฺลิตฺตา\-วลิตฺตานิ นิวาตานิ ผุสิตคฺคฬานิ ปิหิต\-วาตปานานิ. เอวเมวํ โข ภิกฺขเว ยานิ กานิจิ ภายานิ อุปฺปชฺชนฺติ, สพฺพานิ ตานิ พาลโต อุปฺปชฺชนฺติ, โน ปณฺฑิตโต. เย เกจิ อุปทฺทวา อุปฺปชฺชนฺติ, สพฺเพ เต พาลโต อุปฺปชฺชนฺติ, โน ปณฺฑิตโต. เย เกจิ อุปสคฺคา อุปฺปชฺชนฺติ, สพฺเพ เต พาลโต อุปฺปชฺชนฺติ, โน ปณฺฑิตโต.}

\prose{\csromanfolio{100}อิติ โข ภิกฺขเว สปฺปฏิภโย พาโล, อปฺปฏิภโย ปณฺฑิโต, สอุปทฺทโว พาโล, อนุปทฺทโว ปณฺฑิโต, สอุปสคฺโค พาโล, อนุปสคฺโค ปณฺฑิโต, นตฺถิ ภิกฺขเว ปณฺฑิตโต ภยํ, นตฺถิ ปณฺฑิตโต อุปทฺทโว, นตฺถิ ปณฺฑิตโต อุปสคฺโค.}

\prose{ตสฺมาติห ภิกฺขเว เอวํ สิกฺขิตพฺพํ “เยหิ ตีหิ ธมฺเมหิ สมนฺนาคโต พาโล เวทิตพฺโพ, เต ตโย ธมฺเม อภินิวชฺเชตฺวา, เยหิ ตีหิ ธมฺเมหิ สมนฺนาคโต ปณฺฑิโต เวทิตพฺโพ, เต ตโย ธมฺเม สมาทาย วตฺติสฺสามา”ติ. เอวํ หิ โว ภิกฺขเว สิกฺขิตพฺพนฺติ.  ปฐมํ.}

\csromanheader{2}{2. \textbf{ลกฺขณสุตฺต}}
\proseitem{2}{กมฺมลกฺขโณ ภิกฺขเว พาโล, กมฺมลกฺขโณ ปณฺฑิโต, อปทานโสภนี\csromansharedfootnote{057bc6d1014e2c59}{อปทาเน โสภติ (สฺยา, กํ, อิ)} ปญฺญาติ\csromansharedfootnote{1cb1498b383d0d2a}{ปญฺญตฺติ (?)}. ตีหิ ภิกฺขเว ธมฺเมหิ สมนฺนาคโต พาโล เวทิตพฺโพ. กตเมหิ ตีหิ? กาย\-ทุจฺจริเตน วจีทุจฺจริเตน มโน\-ทุจฺจริเตน. อิเมหิ โข ภิกฺขเว ตีหิ ธมฺเมหิ สมนฺนาคโต พาโล เวทิตพฺโพ.}

\prose{ตีหิ ภิกฺขเว ธมฺเมหิ สมนฺนาคโต ปณฺฑิโต เวทิตพฺโพ. กตเมหิ ตีหิ? กายสุจริเตน วจีสุจริเตน มโนสุจริเตน. อิเมหิ โข ภิกฺขเว ตีหิ ธมฺเมหิ สมนฺนาคโต ปณฺฑิโต เวทิตพฺโพ.}

\prose{ตสฺมาติห ภิกฺขเว เอวํ สิกฺขิตพฺพํ “เยหิ ตีหิ ธมฺเมหิ สมนฺนาคโต พาโล เวทิตพฺโพ, เต ตโย ธมฺเม อภินิวชฺเชตฺวา, เยหิ ตีหิ ธมฺเมหิ สมนฺนาคโต ปณฺฑิโต เวทิตพฺโพ, เต ตโย ธมฺเม สมาทาย วตฺติสฺสามา”ติ. เอวํ หิ โว ภิกฺขเว สิกฺขิตพฺพนฺติ.  ทุติยํ.}

\csromanheader{2}{3. \textbf{จินฺตีสุตฺต}}
\proseitem{3}{ตีณิมานิ ภิกฺขเว พาลสฺส พาลลกฺขณานิ พาลนิมิตฺตานิ พาลาปทานานิ. กตมานิ ตีณิ? อิธ ภิกฺขเว พาโล ทุจฺจินฺติต\-จินฺตี จ โหติ ทุพฺภาสิต\-ภาสี จ ทุกฺกฏ\-กมฺม\-การี จ. โน เจทํ\csromansharedfootnote{ee68b814aedcedaf}{โน เจตํ (สฺยา, กํ, ก)} ภิกฺขเว พาโล ทุจฺจินฺติต\-จินฺตี จ อภวิสฺส ทุพฺภาสิต\-ภาสี จ ทุกฺกฏ\-กมฺม\-การี จ, เกน นํ ปณฺฑิตา ชาเนยฺยุํ\csromansharedfootnote{9600cbeefdb63312}{เตน นํ ปณฺฑิตา น ชาเนยฺยุํ (ก), น นํ ปณฺฑิตา ชาเนยฺยุํ (?)} “พาโล อยํ ภวํ อสปฺปุริโส”ติ.}

\vspace{\tipitakaparskip}
\csromancenter{พาลวคฺค ยสฺมา จ โข ภิกฺขเว พาโล ทุจฺจินฺติต\-จินฺตี จ โหติ ทุพฺภาสิต\-ภาสี จ ทุกฺกฏ\-กมฺม\-การี จ, ตสฺมา นํ ปณฺฑิตา ชานนฺติ “พาโล อยํ ภวํ อสปฺปุริโส”ติ. อิมานิ โข ภิกฺขเว ตีณิ พาลสฺส พาลลกฺขณานิ พาลนิมิตฺตานิ พาลาปทานานิ.}
\prose{\csromanfolio{101}ตีณิมานิ ภิกฺขเว ปณฺฑิตสฺส ปณฺฑิต\-ลกฺขณานิ ปณฺฑิต\-นิมิตฺตานิ ปณฺฑิตา\-ปทานานิ. กตมานิ ตีณิ? อิธ ภิกฺขเว ปณฺฑิโต สุจินฺติตจินฺตี จ โหติ สุภาสิตภาสี จ สุกต\-กมฺม\-การี จ. โน เจทํ ภิกฺขเว ปณฺฑิโต สุจินฺติตจินฺตี จ อภวิสฺส สุภาสิตภาสี จ สุกต\-กมฺม\-การี จ, เกน นํ ปณฺฑิตา ชาเนยฺยุํ “ปณฺฑิโต อยํ ภวํ สปฺปุริโส”ติ. ยสฺมา จ โข ภิกฺขเว ปณฺฑิโต สุจินฺติตจินฺตี จ โหติ สุภาสิตภาสี จ สุกต\-กมฺม\-การี จ, ตสฺมา นํ ปณฺฑิตา ชานนฺติ “ปณฺฑิโต อยํ ภวํ สปฺปุริโส”ติ. อิมานิ โข ภิกฺขเว ตีณิ ปณฺฑิตสฺส ปณฺฑิต\-ลกฺขณานิ ปณฺฑิต\-นิมิตฺตานิ ปณฺฑิตา\-ปทานานิ. ตสฺมาติห ฯเปฯ.  ตติยํ.}

\csromanheader{2}{4. \textbf{อจฺจย สุตฺต}}
\proseitem{4}{ตีหิ ภิกฺขเว ธมฺเมหิ สมนฺนาคโต พาโล เวทิตพฺโพ. กตเมหิ ตีหิ? อจฺจยํ อจฺจยโต น ปสฺสติ, อจฺจยํ อจฺจยโต ทิสฺวา ยถาธมฺมํ นปฺปฏิกโรติ, ปรสฺส โข ปน อจฺจยํ เทเสนฺตสฺส ยถาธมฺมํ นปฺปฏิคฺคณฺหาติ. อิเมหิ โข ภิกฺขเว ตีหิ ธมฺเมหิ สมนฺนาคโต พาโล เวทิตพฺโพ.}

\prose{ตีหิ ภิกฺขเว ธมฺเมหิ สมนฺนาคโต ปณฺฑิโต เวทิตพฺโพ. กตเมหิ ตีหิ? อจฺจยํ อจฺจยโต ปสฺสติ, อจฺจยํ อจฺจยโต ทิสฺวา ยถาธมฺมํ นปฏิกโรติ, ปรสฺส โข ปน อจฺจยํ เทเสนฺตสฺส ยถาธมฺมํ ปฏิคฺคณฺหาติ. อิเมหิ โข ภิกฺขเว ตีหิ ธมฺเมหิ สมนฺนาคโต ปณฺฑิโต เวทิตพฺโพ. ตสฺมาติห ฯเปฯ.  จตุตฺถํ.}

\csromanheader{2}{5. \textbf{อโยนิโสสุตฺต}}
\proseitem{5}{ตีหิ ภกฺขเว ธมฺเมหิ สมนฺนาคโต พาโล เวทิตพฺโพ. กตเมหิ ตีหิ? อโยนิโส ปญฺหํ กตฺตา โหติ, อโยนิโส ปญฺหํ วิสฺสชฺเชตา โหติ, ปรสฺส โข ปน โยนิโส ปญฺหํ วิสฺสชฺชิตํ \csromanfolio{102}ปริมณฺฑเลหิ ปทพฺยญฺชเนหิ สิลิฏฺเฐหิ อุปคเตหิ นาพฺภนุโมทิตา โหติ. อิเมหิ โข ภิกฺขเว ตีหิ ธมฺเมหิ สมนฺนาคโต พาโล เวทิตพฺโพ.}

\prose{ตีหิ ภิกฺขเว ธมฺเมหิ สมนฺนาคโต ปณฺฑิโต เวทิตพฺโพ. กตเมหิ ตีหิ? โยนิโส ปญฺหํ กตฺตา โหติ, โยนิโส ปญฺหํ วิสฺสชฺเชตา โหติ, ปรสฺส โข ปน โยนิโส ปญฺหํ วิสฺสชฺชิตํ ปริมณฺฑเลหิ ปทพฺยญฺชเนหิ สิลิฏฺเฐหิ อุปคเตหิ อพฺภนุโมทิตา โหติ. อิเมหิ โข ภิกฺขเว ตีหิ ธมฺเมหิ สมนฺนาคโต ปณฺฑิโต เวทิตพฺโพ. ตสฺมาติห ฯเปฯ.  ปญฺจมํ.}

\csromanheader{2}{6. \textbf{อกุสลสุตฺต}}
\proseitem{6}{ตีหิ ภิกฺขเว ธมฺเมหิ สมนฺนาคโต พาโล เวทิตพฺโพ. กตเมหิ ตีหิ? อกุสเลน กายกมฺเมน, อกุสเลน วจีกมฺเมน, อกุสเลน มโนกมฺเมน. อิเมหิ โข ภิกฺขเว ตีหิ ธมฺเมหิ สมนฺนาคโต พาโล เวทิตพฺโพ.}

\prose{ตีหิ ภิกฺขเว ธมฺเมหิ สมนฺนาคโต ปณฺฑิโต เวทิตพฺโพ. กตเมหิ ตีหิ? กุสเลน กายกมฺเมน, กุสเลน วจีกมฺเมน, กุสเลน มโนกมฺเมน. อิเมหิ โข ภิกฺขเว ตีหิ ธมฺเมหิ สมนฺนาคโต ปณฺฑิโต เวทิตพฺโพ. ตสฺมาติห ฯเปฯ. ฉฏฺฐํ.}

\csromanheader{2}{7. \textbf{สาวชฺชสุตฺต}}
\proseitem{7}{ตีหิ ภิกฺขเว ธมฺเมหิ สมนฺนาคโต พาโล เวทิตพฺโพ. กตเมหิ ตีหิ? สาวชฺเชน กายกมฺเมน, สาวชฺเชน วจีกมฺเมน, สาวชฺเชน มโนกมฺเมน ฯเปฯ อนวชฺเชน กายกมฺเมน, อนวชฺเชน วจีกมฺเมน, อนวชฺเชน มโนกมฺเมน ฯเปฯ.  สตฺตมํ.}

\csromanheader{2}{8. สพฺยาพชฺฌ\-สุตฺต}
\proseitem{8}{ตีหิ ภิกฺขเว ธมฺเมหิ สมนฺนาคโต พาโล เวทิตพฺโพ. กตเมหิ ตีหิ? สพฺยาพชฺเฌน กายกมฺเมน, สพฺยาพชฺเฌน วจีกมฺเมน, สพฺยาพชฺเฌน มโนกมฺเมน ฯเปฯ อพฺยาพชฺเฌน กายกมฺเมน, อพฺยาพชฺเชน วจี-}

\vspace{\tipitakaparskip}
\csromancenter{พาลวคฺค กมฺเมน, อพฺยาพชฺเชน มโนกมฺเมน. อิเมหิ โข ภิกฺขเว ตีหิ ธมฺเมหิ สมนฺนาคโต ปณฺฑิโต เวทิตพฺโพ.}
\prose{\csromanfolio{103}ตสฺมาติห ภิกฺขเว เอวํ สิกฺขิตพฺพํ “เยหิ ตีหิ ธมฺเมหิ สมนฺนาคโต พาโล เวทิตพฺโพ, เต ตโย ธมฺเม อภินิวชฺเชตฺวา, เยหิ ตีหิ ธมฺเมหิ สมนฺนาคโต ปณฺฑิโต เวทิตพฺโพ, เต ตโย ธมฺเม สมาทาย วตฺติสฺสามา”ติ. เอวํ หิ โว ภิกฺขเว สิกฺขิตพฺพนฺติ.  อฏฺฐมํ.}

\csromanheader{2}{9. \textbf{ขตสุตฺต}}
\proseitem{9}{ตีหิ ภิกฺขเว ธมฺเมหิ สมนฺนาคโต พาโล อพฺยตฺโต อสปฺปุริโส ขตํ อุปหตํ อตฺตานํ ปริหรติ, สาวชฺโช จ โหติ สานุวชฺโช จ วิญฺญูนํ, พหุญฺจ อปุญฺญํ ปสวติ. กตเมหิ ตีหิ? กาย\-ทุจฺจริเตน วจีทุจฺจริเตน มโน\-ทุจฺจริเตน. อิเมหิ โข ภิกฺขเว ตีหิ ธมฺเมหิ สมนฺนาคโต พาโล อพฺยตฺโต อสปฺปุริโส ขตํ อุปหตํ อตฺตานํ ปริหรติ, สาวชฺโช จ โหติ สานุวชฺโช จ วิญฺญูนํ, พหุญฺจ อปุญฺญํ ปสวติ.}

\prose{ตีหิ ภิกฺขเว ธมฺเมหิ สมนฺนาคโต ปณฺฑิโต วิยตฺโต สปฺปุริโส อกฺขตํ อนุปหตํ อตฺตานํ ปริหรติ, อนวชฺโช จ โหติ อนนุวชฺโช จ วิญฺญูนํ, พหุญฺจ ปุญฺญํ ปสวติ. กตเมหิ ตีหิ? กายสุจริเตน วจีสุจริเตน มโนสุจริเตน. อิเมหิ โข ภิกฺขเว ตีหิ ธมฺเมหิ สมนฺนาคโต ปณฺฑิโต วิยตฺโต อปฺปุริโส อกฺขตํ อนุปหตํ อตฺตานํ ปริหรติ, อนวชฺโช จ โหติ อนนุวชฺโช จ วิญฺญูนํ, พหุญฺจ ปุญฺญํ ปสวตีติ.  นวมํ.}

\csromanheader{2}{10. \textbf{มลสุตฺต}}
\proseitemwithclosermidruled{10}{ตีหิ ภิกฺขเว ธมฺเมหิ สมนฺนาคโต ตโย มเล อปฺปหาย ยถาภตํ นิกฺขิตฺโต, เอวํ นิรเย. กตเมหิ ตีหิ? ทุสฺสีโล จ โหติ, ทุสฺสีลฺยมล\-ญฺจสฺส อปฺปหีนํ โหติ. อิสฺสุกี จ โหติ, อิสฺสามลญฺจสฺส อปฺปหีนํ โหติ. มจฺฉรี จ โหติ, มจฺเฉรมล\-ญฺจสฺส อปฺปหีนํ โหติ. อิเมหิ โข ภิกฺขเว ตีหิ ธมฺเมหิ สมนฺนาคโต อิเม ตโย มเล อปฺปหาย ยถาภตํ นิกฺขิตฺโต, เอวํ นิรเย. \csromanfolio{104}ตีหิ ภิกฺขเว ธมฺเมหิ สมนฺนาคโต ตโย มเล ปหาย ยถาภตํ นิกฺขิตฺโต, เอวํ สคฺเค. กตเมหิ ตีหิ? สีลวา จ โหติ, ทุสฺสีลฺยมล\-ญฺจสฺส ปหีนํ โหติ. อนิสฺสุกี จ โหติ, อิสฺสามลญฺจสฺส ปหีนํ โหติ. อมจฺฉรี จ โหติ, มจฺเฉรมล\-ญฺจสฺส ปหีนํ โหติ. อิเมหิ โข ภิกฺขเว ตีหิ ธมฺเมหิ สมนฺนาคโต อิเม ตโย มเล ปหาย ยถาภตํ นิกฺขิตฺโต, เอวํ สคฺเคติ.  ทสมํ.}{พาลวคฺโค ปฐโม.}
\csromancenter{\textbf{ตสฺสุทฺทานํ}}
\setlength{\csromangathapagewidth}{0pt}
\setlength{\csromangathawakwidth}{0pt}
\csromangathameasuregroup{ภยํ ลกฺขณจินฺตี จ, \\ อกุสลญฺจ สาวชฺชํ,}{\csromangathabat{ภยํ ลกฺขณจินฺตี จ,}{อจฺจยญฺจ อโยนิโส.} \\ \csromangathabat{อกุสลญฺจ สาวชฺชํ,}{สพฺยาพชฺฌขตํ มลนฺติ.}}
\csromangathasetleft{ภยํ ลกฺขณจินฺตี จ, \\ อกุสลญฺจ สาวชฺชํ,}
\csromangathagroup{\csromangathastanza{\csromangathabat{ภยํ ลกฺขณจินฺตี จ,}{อจฺจยญฺจ อโยนิโส.} \\ \csromangathabat{อกุสลญฺจ สาวชฺชํ,}{สพฺยาพชฺฌขตํ มลนฺติ.}}}

\chapterhead{2. \textbf{รถการวคฺค}}
\csromanheader{2}{1. \textbf{ญาตสุตฺต}}
\proseitem{11}{ตีหิ ภิกฺขเว ธมฺเมหิ สมนฺนาคโต ญาโต ภิกฺขุ พหุชน\-อหิตาย ปฏิปนฺโน โหติ พหุชน\-ทุกฺขาย พหุโน ชนสฺส อนตฺถาย อหิตาย ทุกฺขาย เทวมนุสฺสานํ. กตเมหิ ตีหิ? อนนุโลมิเก กายกมฺเม สมาทเปติ, อนนุโลมิเก วจีกมฺเม สมาทเปติ, อนนุโลมิเกสุ ธมฺเมสุ สมาทเปติ. อิเมหิ โข ภิกฺขเว ตีหิ ธมฺเมหิ สมนฺนาคโต ญาโต ภิกฺขุ พหุชน\-อหิตาย ปฏิปนฺโน โหติ พหุชน\-ทุกฺขาย พหุโน ชนสฺส อนตฺถาย อหิตาย ทุกฺขาย เทวมนุสฺสานํ.}

\prose{ตีหิ ภิกฺขเว ธมฺเมหิ สมนฺนาคโต ญาโต ภิกฺขุ พหุชนหิตาย ปฏิปนฺโน โหติ พหุชน\-สุขาย พหุโน ชนสฺส อตฺถาย หิตาย สุขาย เทวมนุสฺสานํ. กตเมหิ ตีหิ? อนุโลมิเก กายกมฺเม สมาทเปติ, อนุโลมิเก วจีกมฺเม สมาทเปติ, อนุโลมิเกสุ ธมฺเมสุ สมาทเปติ. อิเมหิ โข ภิกฺขเว ตีหิ ธมฺเมหิ สมนฺนาคโต ญาโต ภิกฺขุ พหุชนหิตาย ปฏิปนฺโน โหติ พหุชน\-สุขาย พหุโน ชนสฺส อตฺถาย หิตาย สุขาย เทว\-มนุสฺสา\-นนฺติ.  ปฐมํ.}

\chapterheadpageread{2. รถการวคฺค \textbf{2. สารณียสุตฺต}}
\proseitem{12}{\csromanfolio{105}ตีณิมานิ ภิกฺขเว รญฺโญ ขตฺติยสฺส มุทฺธา\-วสิตฺตสฺส ยาวชีวํ สารณียานิ\csromansharedfootnote{42665d464db059f9}{สรณียานิ (สี, สฺยา, กํ, อิ)} ภวนฺติ. กตมานิ ตีณิ? ยสฺมิํ ภิกฺขเว ปเทเส ราชา ขตฺติโย มุทฺธาวสิตฺโต ชาโต โหติ, อิทํ ภิกฺขเว ปฐมํ รญฺโญ ขตฺติยสฺส มุทฺธา\-วสิตฺตสฺส ยาวชีวํ สารณียํ โหติ.}

\prose{ปุน จปรํ ภิกฺขเว ยสฺมิํ ปเทเส ราชา ขตฺติโย มุทฺธาวสิตฺโต โหติ, อิทํ ภิกฺขเว ทุติยํ รญฺโญ ขตฺติยสฺส มุทฺธา\-วสิตฺตสฺส ยาวชีวํ สารณียํ โหติ.}

\prose{ปุน จปรํ ภิกฺขเว ยสฺมิํ ปเทเส ราชา ขตฺติโย มุทฺธาวสิตฺโต สงฺคามํ อภิวิชินิตฺวา วิชิตสงฺคาโม ตเมว สงฺคามสีสํ อชฺฌาวสติ, อิทํ ภิกฺขเว ตติยํ รญฺโญ ขตฺติยสฺส มุทฺธวสิตฺตสฺส ยาวชีวํ สารณียํ โหติ. อิมานิ โข ภิกฺขเว ตีณิ รญฺโญ ขตฺติยสฺส มุทฺธา\-วสิตฺตสฺส ยาวชีวํ สารณียานิ ภวนฺติ.}

\prose{เอวเมวํ โข ภิกฺขเว ตีณิมานิ ภิกฺขุสฺส ยาวชีวํ สารณียานิ ภวนฺติ. กตมานิ ตีณิ? ยสฺมิํ ภิกฺขเว ปเทเส ภิกฺขุ เกสมสฺสุํ โอหาเรตฺวา กาสายานิ วตฺถานิ อจฺฉาเทตฺวา อคารสฺมา อนคาริยํ ปพฺพชิโต โหติ, อิทํ ภิกฺขเว ปฐมํ ภิกฺขุสฺส ยาวชีวํ สารณียํ โหติ.}

\prose{ปุน จปรํ ภิกฺขเว ยสฺมิํ ปเทเส ภิกฺขุ “อิทํ ทุกฺขนฺ”ติ ยถาภูตํ ปชานาติ, “อยํ ทุกฺขสมุทโย”ติ ยถาภูตํ ปชานาติ, “อยํ ทุกฺขนิโรโธ”ติ ยถาภูตํ ปชานาติ, “อยํ ทุกฺขนิโรธ\-คามินี ปฏิปทา”ติ ยถาภูตํ ปชานาติ, อิทํ ภิกฺขเว ทุติยํ ภิกฺขุสฺส ยาวชีวํ สารณียํ โหติ.}

\prose{ปุน จปรํ ภิกฺขเว ยสฺมิํ ปเทเส ภิกฺขุ อาสวานํ ขยา อนาสวํ เจโตวิมุตฺติํ ปญฺญาวิมุตฺติํ ทิฏฺเฐว ธมฺเม สยํ อภิญฺญา สจฺฉิกตฺวา อุปสมฺปชฺช วิหรติ, อิทํ ภิกฺขเว ตติยํ ภิกฺขุสฺส ยาวชีวํ สารณียํ โหติ. อิมานิ โข ภิกฺขเว ตีณิ ภิกฺขุสฺส ยาวชีวํ สารณียานิ ภวนฺตีติ.  ทุติยํ.}

\csromanpalirecto
\gambhira{ติกนิปาตปาฬิ \textbf{3. อาสํสสุตฺต}}
\proseitem{13}{\csromanfolio{106}ตโยเม ภิกฺขเว ปุคฺคลา สนฺโต สํวิชฺชมานา โลกสฺมิํ. กตเม ตโย? นิราโส อาสํโส วิคตาโส. กตโม จ ภิกฺขเว ปุคฺคโล นิราโส? อิธ ภิกฺขเว เอกจฺโจ ปุคฺคโล นีเจ กุเล ปจฺจาชาโต โหติ จณฺฑาลกุเล วา เวนกุเล\csromansharedfootnote{19214b6e764cb04d}{เวณกุเล (สฺยา, กํ, อิ)} วา เนสาทกุเล วา รถการกุเล วา ปุกฺกุสกุเล วา ทลิทฺเท อปฺป\-นฺน\-ปาน\-โภชเน กสิรวุตฺติเก, ยตฺถ กสิเรน ฆาสจฺฉาโท ลพฺภติ. โส จ โหติ ทุพฺพณฺโณ ทุทฺทสิโก โอโกฏิมโก พวฺหาพาโธ\csromansharedfootnote{9c99195a1531f803}{พหฺวาพาโธ (สฺยา, กํ, อิ, ก)} กาโณ วา กุณี วา ขญฺโช วา ปกฺขหโต วา, น ลาภี อนฺนสฺส ปานสฺส วตฺถสฺส ยานสฺส มาลา\-คนฺธ\-วิเลปนสฺส เสยฺยา\-วสถ\-ปทีเปยฺยสฺส. โส สุณาติ “อิตฺถนฺนาโม กิร ขตฺติโย ขตฺติเยหิ ขตฺติยา\-ภิเสเกน อภิสิตฺโต”ติ. ตสฺส น เอวํ โหติ “กุทาสฺสุ นาม มมฺปิ ขตฺติยา ขตฺติยา\-ภิเสเกน อภิสิญฺจิสฺสนฺตี”ติ. อยํ วุจฺจติ ภิกฺขเว ปุคฺคโล นิราโส.}

\prose{กตโม จ ภิกฺขเว ปุคฺคโล อาสํโส? อิธ ภิกฺขเว รญฺโญ ขตฺติยสฺส มุทฺธา\-วสิตฺตสฺส เชฏฺโฐ ปุตฺโต โหติ อาภิเสโก อนภิสิตฺโต อจลปฺปตฺโต\csromansharedfootnote{926aeba0d3c6fa22}{มจลปฺปตฺโต (สี, อิ)}. โส สุณาติ “อิตฺถนฺนาโม กิร ขตฺติโย ขตฺติเยหิ ขตฺติยา\-ภิเสเกน อภิสิตฺโต”ติ. ตสฺส เอวํ โหติ “กุทาสฺสุ นาม มมฺปิ ขตฺติยา ขตฺติยา\-ภิเสเกน อภิสิญฺจิสฺสนฺตี”ติ. อยํ วุจฺจติ ภิกฺขเว ปุคฺคโล อาสํโส.}

\prose{กตโม จ ภิกฺขเว ปุคฺคโล วิคตาโส? อิธ ภิกฺขเว ราชา โหติ ขตฺติโย มุทฺธาวสิตฺโต. โส สุณาติ “อิตฺถนฺนาโม กิร ขตฺติโย ขตฺติเยหิ ขตฺติยา\-ภิเสเกน อภิสิตฺโต”ติ. ตสฺส น เอวํ โหติ “กุทาสฺสุ นาม มมฺปิ ขตฺติยา ขตฺติยา\-ภิเสเกน อภิสิญฺจิสฺสนฺตี”ติ. ตํ กิสฺส เหตุ? ยา หิสฺส ภิกฺขเว ปุพฺเพ อนภิสิตฺตสฺส อภิเสกาสา, สา\csromansharedfootnote{e5a1c78e5ada5517}{สาสฺส (สี, สฺยา, กํ, อิ)} ปฏิปฺปสฺสทฺธา. อยํ วุจฺจติ ภิกฺขเว ปุคฺคโล วิคตาโส. อิเม โข ภิกฺขเว ตโย ปุคฺคลา สนฺโต สํวิชฺชมานา โลกสฺมิํ.}

\chapterheadpageread{2. รถการวคฺค}
\prose{\csromanfolio{107}เอวเมวํ โข ภิกฺขเว ตโย ปุคฺคลา สนฺโต สํวิชฺชมานา ภิกฺขูสุ. กตเม ตโย? นิราโส อาสํโส วิคตาโส. กตโม จ ภิกฺขเว ปุคฺคโล นิราโส? อิธ ภิกฺขเว เอกจฺโจ ปุคฺคโล ทุสฺสีโล โหติ ปาปธมฺโม อสุจิ สงฺกสฺสร\-สมาจาโร ปฏิจฺฉนฺน\-กมฺมนฺโต อสฺสมโณ สมณปฏิญฺโญ อพฺรหฺมจารี พฺรหฺมจาริ\-ปฏิญฺโญ อนฺโตปูติ อวสฺสุโต กสมฺพุชาโต. โส สุณาติ “อิตฺถนฺนาโม กิร ภิกฺขุ อาสวานํ ขยา อนาสวํ เจโตวิมุตฺติํ ปญฺญาวิมุตฺติํ ทิฏฺเฐว ธมฺเม สยํ อภิญฺญา สจฺฉิกตฺวา อุปสมฺปชฺช วิหรตี”ติ. ตสฺส น เอวํ โหติ “กุทาสฺสุ นาม อหมฺปิ อาสวานํ ขยา อนาสวํ เจโตวิมุตฺติํ ปญฺญาวิมุตฺติํ ทิฏฺเฐว ธมฺเม สยํ อภิญฺญา สจฺฉิกตฺวา อุปสมฺปชฺช วิหริสฺสามี”ติ. อยํ วุจฺจติ ภิกฺขเว ปุคฺคโล นิราโส.}

\prose{กตโม จ ภิกฺขเว ปุคฺคโล อาสํโส? อิธ ภิกฺขเว ภิกฺขุ สีลวา โหติ กลฺยาณธมฺโม. โส สุณาติ “อิตฺถนฺนาโม กิร ภิกฺขุ อาสวานํ ขยา อนาสวํ เจโตวิมุตฺติํ ปญฺญาวิมุตฺติํ ทิฏฺเฐว ธมฺเม สยํ อภิญฺญา สจฺฉิกตฺวา อุปสมฺปชฺช วิหรตี”ติ. ตสฺส เอวํ โหติ “กุทาสฺสุ นาม อหมฺปิ อาสวานํ ขยา อนาสวํ เจโตวิมุตฺติํ ปญฺญาวิมุตฺติํ ทิฏฺเฐว ธมฺเม สยํ อภิญฺญา สจฺฉิกตฺวา อุปสมฺปชฺช วิหริสฺสามี”ติ. อยํ วุจฺจติ ภิกฺขเว ปุคฺคโล อาสํโส.}

\prose{กตโม จ ภิกฺขเว ปุคฺคโล วิคตาโส? อิธ ภิกฺขเว ภิกฺขุ อรหํ โหติ ขีณาสโว. โส สุณาติ “อิตฺถนฺนาโม กิร ภิกฺขุ อาสวานํ ขยา อนาสวํ เจโตวิมุตฺติํ ปญฺญาวิมุตฺติํ ทิฏฺเฐว ธมฺเม สยํ อภิญฺญา สจฺฉิกตฺวา อุปสมฺปชฺช วิหรตี”ติ. ตสฺส น เอวํ โหติ “กุทาสฺสุ นาม อหมฺปิ อาสวานํ ขยา ฯเปฯ สจฺฉิกตฺวา อุปสมฺปชฺช วิหริสฺสามี”ติ. ตํ กิสฺส เหตุ, ยา หิสฺส ภิกฺขเว ปุพฺเพ อวิมุตฺตสฺส วิมุตฺตาสา, สา ปฏิปฺปสฺสทฺธา. อยํ วุจฺจติ ภิกฺขเว ปุคฺคโล วิคตาโส. อิเม โข ภิกฺขเว ตโย ปุคฺคลา สนฺโต สํวิชฺชมานา ภิกฺขูสูติ.  ตติยํ.}

\csromanheader{2}{4. จกฺกวตฺติ\-สุตฺต}
\proseitem{14}{โยปิ โส ภิกฺขเว ราชา จกฺกวตฺตี ธมฺมิโก ธมฺมราชา, โสปิ น อราชกํ จกฺกํ วตฺเตตีติ. เอวํ วุตฺเต อญฺญตโร ภิกฺขุ ภควนฺตํ \csromanfolio{108}เอตทโวจ “โก ปน ภนฺเต รญฺโญ จกฺกวตฺติสฺส ธมฺมิกสฺส ธมฺมรญฺโญ ราชา”ติ\csromansharedfootnote{37e7b2f6d6087eea}{จกฺกนฺติ (ก)}. “ธมฺโม ภิกฺขู”ติ ภควา อโวจ, อิธ ภิกฺขุ ราชา จกฺกวตฺตี ธมฺมิโก ธมฺมราชา ธมฺมํเยว นิสฺสาย ธมฺมํ สกฺกโรนฺโต ธมฺมํ ครุํ กโรนฺโต\csromansharedfootnote{120bc1e00d2609b2}{ครุกโรนฺโต (สี, สฺยา, กํ, อิ)} ธมฺมํ อปจายมาโน ธมฺมทฺธโช ธมฺมเกตุ ธมฺมา\-ธิปเตยฺโย ธมฺมิกํ รกฺขา\-วรณ\-คุตฺติํ สํวิทหติ อนฺโตชนสฺมิํ.}

\prose{ปุน จปรํ ภิกฺขุ ราชา จกฺกวตฺตี ธมฺมิโก ธมฺมราชา ธมฺมํเยว นิสฺสาย ธมฺมํ สกฺกโรนฺโต ธมฺมํ ครุํ กโรนฺโต ธมฺมํ อปจายมาโน ธมฺมทฺธโช ธมฺมเกตุ ธมฺมา\-ธิปเตยฺโย ธมฺมิกํ รกฺขา\-วรณ\-คุตฺติํ สํวิทหติ ขตฺติเยสุ, อนุยนฺเตสุ\csromansharedfootnote{8e58a7c6bc50a7a4}{อนุยุตฺเตสุ (สี, สฺยา, กํ, อิ)}, พลกายสฺมิํ, พฺราหฺมณคหปติเกสุ, เนคม\-ชานปเทสุ, สมณ\-พฺราหฺมเณสุ, มิคปกฺขีสุ. ส โข โส ภิกฺขุ ราชา จกฺกวตฺตี ธมฺมิโก ธมฺมราชา ธมฺมํเยว นิสฺสาย ธมฺมํ สกฺกโรนฺโต ธมฺมํ ครุํ กโรนฺโต ธมฺมํ อปจายมาโน ธมฺมทฺธโช ธมฺมเกตุ ธมฺมา\-ธิปเตยฺโย ธมฺมิกํ รกฺขา\-วรณ\-คุตฺติํ สํวิทหิตฺวา อนฺโตชนสฺมิํ, ธมฺมิกํ รกฺขา\-วรณ\-คุตฺติํ สํวิทหิตฺวา ขตฺติเยสุ ฯเปฯ อนุยนฺเตสุ. พลกายสฺมิํ. พฺราหฺมณคหปติเกสุ. เนคม\-ชานปเทสุ. สมณ\-พฺราหฺมเณสุ. มิคปกฺขีสุ ธมฺเมเนว จกฺกํ วตฺเตติ, ตํ โหติ จกฺกํ อปฺปฏิวตฺติยํ เกนจิ มนุสฺสภูเตน ปจฺจตฺถิเกน ปาณินา.}

\prose{เอวเมวํ โข ภิกฺขุ\csromansharedfootnote{17b89de435939d95}{ภิกฺขเว (ก)} ตถาคโต อรหํ สมฺมาสมฺพุทฺโธ ธมฺมิโก ธมฺมราชา ธมฺมํเยว นิสฺสาย ธมฺมํ สกฺกโรนฺโต ธมฺมํ ครุํ กโรนฺโต ธมฺมํ อปจายมาโน ธมฺมทฺธโช ธมฺมเกตุ ธมฺมา\-ธิปเตยฺโย ธมฺมิกํ รกฺขา\-วรณ\-คุตฺติํ สํวิทหติ กายกมฺมสฺมิํ “เอวรูปํ กายกมฺมํ เสวิตพฺพํ, เอวรูปํ กายกมฺมํ น เสวิตพฺพนฺ”ติ.}

\prose{ปุน จปรํ ภิกฺขุ ตถาคโต อรหํ สมฺมาสมฺพุทฺโธ ธมฺมิโก ธมฺมราชา ธมฺมํเยว นิสฺสาย ธมฺมํ สกฺกโรนฺโต ธมฺมํ ครุํ กโรนฺโต ธมฺมํ อปจายมาโน ธมฺมทฺธโช ธมฺมเกตุ ธมฺมา\-ธิปเตยฺโย ธมฺมิกํ รกฺขา\-วรณ\-คุตฺติํ สํวิทหติ วจีกมฺมสฺมิํ “เอวรูปํ วจีกมฺมํ เสวิตพฺพํ, เอวรูปํ วจีกมฺมํ น เสวิตพฺพนฺ”ติ ฯเปฯ มโนกมฺมสฺมิํ “เอวรูปํ มโนกมฺมํ เสวิตพฺพํ, เอวรูปํ มโนกมฺมํ น เสวิตพฺพนฺ”ติ.}

\chapterheadpageread{2. รถการวคฺค}
\prose{\csromanfolio{109}ส โข โส ภิกฺขุ ตถาคโต อรหํ สมฺมาสมฺพุทฺโธ ธมฺมิโก ธมฺมราชา ธมฺมํเยว นิสฺสาย ธมฺมํ สกฺกโรนฺโต ธมฺมํ ครุํ กโรนฺโต ธมฺมํ อปจายมาโน ธมฺมทฺธโช ธมฺมเกตุ ธมฺมา\-ธิปเตยฺโย ธมฺมิกํ รกฺขา\-วรณ\-คุตฺติํ สํวิทหิตฺวา กายกมฺมสฺมิํ, ธมฺมิกํ รกฺขา\-วรณ\-คุตฺติํ สํวิทหิตฺวา วจีกมฺมสฺมิํ, ธมฺมิกํ รกฺขา\-วรณ\-คุตฺติํ สํวิทหิตฺวา มโนกมฺมสฺมิํ ธมฺเมเนว อนุตฺตรํ ธมฺมจกฺกํ ปวตฺเตติ, ตํ โหติ จกฺกํ อปฺปฏิวตฺติยํ สมเณน วา พฺราหฺมเณน วา เทเวน วา มาเรน วา พฺรหฺมุนา วา เกนจิ วา โลกสฺมินฺติ.  จตุตฺถํ.}

\csromanheader{2}{5. \textbf{สเจตนสุตฺต}}
\proseitem{5}{เอกํ สมยํ ภควา พาราณสิยํ วิหรติ อิสิปตเน มิคทาเย. ตตฺร โข ภควา ภิกฺขู อามนฺเตสิ “ภิกฺขโว”ติ. “ภทนฺเต”ติ เต ภิกฺขู ภควโต ปจฺจสฺโสสุํ. ภควา เอตทโวจ\csromandash{}}

\prose{ภูตปุพฺพํ ภิกฺขเว ราชา อโหสิ สเจตโน\csromansharedfootnote{43aad547a71ca78d}{ปเจตโน (สี, สฺยา, กํ, อิ)} นาม. อถ โข ภิกฺขเว ราชา สเจตโน รถการํ อามนฺเตสิ “อิโต เม สมฺม รถการ ฉนฺนํ มาสานํ อจฺจเยน สงฺคาโม ภวิสฺสติ, สกฺขิสฺสสิ\csromansharedfootnote{bb1d9b197465b8a5}{สกฺขสิ (สฺยา, กํ, อิ)} เม สมฺม รถการ นวํ จกฺกยุคํ กาตุนฺ”ติ. “สกฺโกมิ เทวา”ติ โข ภิกฺขเว รถกาโร รญฺโญ สเจตนสฺส ปจฺจสฺโสสิ. อถ โข ภิกฺขเว รถกาโร ฉหิ มาเสหิ ฉารตฺตูเนหิ เอกํ จกฺกํ นิฏฺฐาเปสิ. อถ โข ภิกฺขเว ราชา สเจตโน รถการํ อามนฺเตสิ “อิโต เม สมฺม รถการ ฉนฺนํ ทิวสานํ อจฺจเยน สงฺคาโม ภวิสฺสติ, นิฏฺฐิตํ นวํ จกฺกยุคนฺ”ติ. อิเมหิ โข เทว ฉหิ มาเสหิ ฉารตฺตูเนหิ เอกํ จกฺกํ นิฏฺฐิตนฺติ. สกฺขิสฺสสิ ปน เม สมฺม รถการ อิเมหิ ฉหิ ทิวเสหิ ทุติยํ จกฺกํ นิฏฺฐาเปตุนฺติ. “สกฺโกมิ เทวา”ติ โข ภิกฺขเว รถกาโร ฉหิ ทิวเสหิ ทุติยํ จกฺกํ นิฏฺฐาเปตฺวา นวํ จกฺกยุคํ อาทาย เยน ราชา สเจตโน เตนุปสงฺกมิ, อุปสงฺกมิตฺวา ราชานํ สเจตนํ เอตทโวจ “อิทํ เต เทว นวํ จกฺกยุคํ นิฏฺฐิตนฺ”ติ. ยญฺจ เต อิทํ สมฺม รถการ จกฺกํ ฉหิ มาเสหิ นิฏฺฐิตํ ฉารตฺตูเนหิ, ยญฺจ เต อิทํ จกฺกํ ฉหิ ทิวเสหิ นิฏฺฐิตํ, อิเมสํ กิํ นานากรณํ, เนสาหํ กิญฺจิ นานากรณํ ปสฺสามีติ. อตฺเถสํ เทว นานากรณํ, ปสฺสตุ เทโว นานากรณนฺติ.}

\prose{\csromanfolio{110}อถ โข ภิกฺขเว รถกาโร ยํ ตํ จกฺกํ ฉหิ ทิวเสหิ นิฏฺฐิตํ, ตํ ปวตฺเตสิ. ตํ ปวตฺติตํ สมานํ ยาวติกา อภิสงฺขารสฺส คติ, ตาวติกํ คนฺตฺวา จิงฺคุลายิตฺวา ภูมิยํ ปปติ. ยํ ปน ตํ จกฺกํ ฉหิ มาเสหิ นิฏฺฐิตํ ฉารตฺตูเนหิ, ตํ ปวตฺเตสิ. ตํ ปวตฺติตํ สมานํ ยาวติกา อภิสงฺขารสฺส คติ, ตาวติกํ คนฺตฺวา อกฺขาหตํ มญฺเญ อฏฺฐาสิ.}

\prose{โก นุ โข สมฺม รถการ เหตุ โก ปจฺจโย, ยมิทํ\csromansharedfootnote{bcdfd19f11be3c35}{ยทิทํ (ก)} จกฺกํ ฉหิ ทิวเสหิ นิฏฺฐิตํ, ตํ ปวตฺติตํ สมานํ ยาวติกา อภิสงฺขารสฺส คติ, ตาวติกํ คนฺตฺวา จิงฺคุลายิตฺวา ภูมิยํ ปปติ. โก ปน สมฺม รถการ เหตุ โก ปจฺจโย, ยมิทํ จกฺกํ ฉหิ มาเสหิ นิฏฺฐิตํ ฉารตฺตูเนหิ, ตํ ปวตฺติตํ สมานํ ยาวติกา อภิสงฺขารสฺส คติ, ตาวติกํ คนฺตฺวา อกฺขาหตํ มญฺเญ อฏฺฐาสีติ. ยมิทํ เทว จกฺกํ ฉหิ ทิวเสหิ นิฏฺฐิตํ, ตสฺส เนมิปิ สวงฺกา สโทสา สกสาวา อราปิ สวงฺกา สโทสา สกสาวา, นาภิปิ สวงฺกา สโทสา สกสาวา. ตํ เนมิยาปิ สวงฺกตฺตา สโทสตฺตา สกสาวตฺตา, อรานมฺปิ สวงฺกตฺตา สโทสตฺตา สกสาวตฺตา, นาภิยาปิ สวงฺกตฺตา สโทสตฺตา สกสาวตฺตา ปวตฺติตํ สมานํ ยาวติกา อภิสงฺขารสฺส คติ, ตาวติกํ คนฺตฺวา จิงฺคุลายิตฺวา ภูมิยํ ปปติ. ยํ ปน ตํ เทว จกฺกํ ฉหิ มาเสหิ นิฏฺฐิตํ ฉารตฺตูเนหิ, ตสฺส เนมิปิ อวงฺกา อโทสา อกสาวา, อราปิ อวงฺกา อโทสา อกสาวา, นาภิปิ อวงฺกา อโทสา อกสาวา. ตํ เนมิยาปิ อวงฺกตฺตา อโทสตฺตา อกสาวตฺตา, อรานมฺปิ อวงฺกตฺตา อโทสตฺตา อกสาวตฺตา, นาภิยาปิ อวงฺกตฺตา อโทสตฺตา อกสาวตฺตา ปวตฺติตํ สมานํ ยาวติกา อภิสงฺขารสฺส คติ, ตาวติกํ คนฺตฺวา อกฺขาหตํ มญฺเญ อฏฺฐาสีติ.}

\prose{สิยา โข ปน ภิกฺขเว ตุมฺหากํ เอวมสฺส “อญฺโญ นูน เตน สมเยน โส รถกาโร อโหสี”ติ. น โข ปเนตํ ภิกฺขเว เอวํ ทฏฺฐพฺพํ, อหํ เตน สมเยน โส รถกาโร อโหสิํ, ตทาหํ ภิกฺขเว กุสโล ทารุวงฺกานํ ทารุโทสานํ ทารุกสาวานํ. เอตรหิ โข ปนาหํ ภิกฺขเว อรหํ สมฺมาสมฺพุทฺโธ กุสโล กายวงฺกานํ \csromanfolio{111}กายโทสานํ กายกสาวานํ, กุสโล วจีวงฺกานํ วจีโทสานํ วจีกสาวานํ, กุสโล มโนวงฺกานํ มโนโทสานํ มโนกสาวานํ. ยสฺส กสฺสจิ ภิกฺขเว ภิกฺขุสฺส วา ภิกฺขุนิยา วา กายวงฺโก อปฺปหีโน กายโทโส กายกสาโว, วจีวงฺโก อปฺปหีโน วจีโทโส วจีกสาโว, มโนวงฺโก อปฺปหีโน มโนโทโส มโนกสาโว. เอวํ ปปติตา เต ภิกฺขเว อิมสฺมา ธมฺมวินยา, เสยฺยถาปิ ตํ จกฺกํ ฉหิ ทิวเสหิ นิฏฺฐิตํ. ยสฺส กสฺสจิ ภิกฺขเว ภิกฺขุสฺส วา ภิกฺขุนิยา วา กายวงฺโก ปหีโน กายโทโส กายกสาโว, วจีวงฺโก ปหีโน วจีโทโส วจีกสาโว, มโนวงฺโก ปหีโน มโนโทโส มโนกสาโว. เอวํ ปติฏฺฐิตา เต ภิกฺขเว อิมสฺมิํ ธมฺมวินเย, เสยฺยถาปิ ตํ จกฺกํ ฉหิ มาเสหิ นิฏฺฐิตํ ฉารตฺตูเนหิ.}

\proseitem{5}{ตสฺมาติห ภิกฺขเว เอวํ สิกฺขิตพฺพํ “กายวงฺกํ ปชหิสฺสาม กายโทสํ กายกสาวํ, วจีวงฺกํ ปชหิสฺสาม วจีโทสํ วจีกสาวํ, มโนวงฺกํ ปชหิสฺสาม มโนโทสํ มโนกสาวนฺ”ติ. เอวํ หิ โว ภิกฺขเว สิกฺขิตพฺพนฺติ.  ปญฺจมํ.}

\csromanheader{2}{6. \textbf{อปณฺณกสุตฺต}}
\proseitem{5}{ตีหิ ภิกฺขเว ธมฺเมหิ สมนฺนาคโต ภิกฺขุ อปณฺณก\-ปฏิปทํ\csromansharedfootnote{366602ac7b4062e1}{อปณฺณกตํ ปฏิปทํ (สี, อิ) ฏีกาย ปน สเมติ.} ปฏิปนฺโน โหติ, โยนิ จสฺส อารทฺธา โหติ อาสวานํ ขยาย. กตเมหิ ตีหิ? อิธ ภิกฺขเว ภิกฺขุ อินฺทฺริเยสุ คุตฺตทฺวาโร โหติ, โภชเน มตฺตญฺญู โหติ, ชาคริยํ อนุยุตฺโต โหติ.}

\prose{กถญฺจ ภิกฺขเว ภิกฺขุ อินฺทฺริเยสุ คุตฺตทฺวาโร โหติ? อิธ ภิกฺขเว ภิกฺขุ จกฺขุนา รูปํ ทิสฺวา น นิมิตฺตคฺคาหี โหติ นานุพฺยญฺชน\-คฺคาหี, ยตฺวาธิกรณเมนํ\csromansharedfootnote{1858f9cd7a6e970f}{ยตฺวาธิกรณเมตํ (สี)} จกฺขุนฺทฺริยํ อสํวุตํ วิหรนฺตํ อภิชฺฌา\-โทมนสฺสา ปาปกา อกุสลา ธมฺมา อนฺวาสฺสเวยฺยุํ, ตสฺส สํวราย ปฏิปชฺชติ, รกฺขติ จกฺขุนฺทฺริยํ, จกฺขุนฺทฺริเย สํวรํ อาปชฺชติ. โสเตน สทฺทํ สุตฺวา. ฆาเนน คนฺธํ ฆายิตฺวา. ชิวฺหาย รสํ สายิตฺวา. กาเยน โผฏฺฐพฺพํ ผุสิตฺวา. มนสา ธมฺมํ วิญฺญาย น นิมิตฺตคฺคาหี โหติ นานุพฺยญฺชน\-คฺคาหี, \csromanfolio{112}ยตฺวาธิกรณเมนํ มนินฺทฺริยํ อสํวุตํ วิหรนฺตํ อภิชฺฌา\-โทมนสฺสา ปาปกา อกุสลา ธมฺมา อนฺวาสฺสเวยฺยุํ, ตสฺส สํวราย ปฏิปชฺชติ, รกฺขติ มนินฺทฺริยํ, มนินฺทฺริเย สํวรํ อาปชฺชติ. เอวํ โข ภิกฺขเว ภิกฺขุ อินฺทฺริเยสุ คุตฺตทฺวาโร โหติ.}

\proseitem{16}{กถญฺจ ภิกฺขเว ภิกฺขุ โภชเน มตฺตญฺญู โหติ? อิธ ภิกฺขเว ภิกฺขุ. ปฏิสงฺขา โยนิโส อาหารํ อาหาเรติ เนว ทวาย น มทาย น มณฺฑานาย น วิภูสนาย, ยาวเทว อิมสฺส กายสฺส ฐิติยา ยาปนาย วิหิํสู\-ปรติยา พฺรหฺมจริยา\-นุคฺคหาย, อิติ ปุราณญฺจ เวทนํ ปฏิหงฺขามิ, นวญฺจ เวทนํ น อุปฺปาเทสฺสามิ, ยาตฺรา จ เม ภวิสฺสติ, อนวชฺชตา จ ผาสุวิหาโร จาติ. เอวํ โข ภิกฺขเว ภิกฺขุ โภชเน มตฺตญฺญู โหติ.}

\prose{กถญฺจ ภิกฺขเว ภิกฺขุ ชาคริยํ อนุยุตฺโต โหติ? อิธ ภิกฺขเว ภิกฺขุ ทิวสํ จงฺกเมน นิสชฺชาย อาวรณีเยหิ ธมฺเมหิ จิตฺตํ ปริโสเธติ, รตฺติยา ปฐมํ ยามํ จงฺกเมน นิสชฺชาย อาวรณีเยหิ ธมฺเมหิ จิตฺตํ ปริโสเธติ, รตฺติยา มชฺฌิมํ ยามํ ทกฺขิเณน ปสฺเสน สีหเสยฺยํ กปฺเปติ ปาเท ปาทํ อจฺจาธาย สโต สมฺปชาโน อุฏฺฐานสญฺญํ มนสิ กริตฺวา, รตฺติยา ปจฺฉิมํ ยามํ ปจฺจุฏฺฐาย จงฺกเมน นิสชฺชาย อาวรณีเยหิ ธมฺเมหิ จิตฺตํ ปริโสเธติ. เอวํ โข ภิกฺขเว ภิกฺขุ ชาคริยํ อนุยุตฺโต โหติ. อิเมหิ โข ภิกฺขเว ตีหิ ธมฺเมหิ สมนฺนาคโต ภิกฺขุ อปณฺณก\-ปฏิปทํ ปฏิปนฺโน โหติ, โยนิ จสฺส อารทฺธา โหติ อาสวานํ ขยายาติ.  ฉฏฺฐํ.}

\csromanheader{2}{7. อตฺตพฺยาพาธ\-สุตฺต}
\proseitem{17}{ตโยเม ภิกฺขเว ธมฺมา อตฺตพฺยาพาธายปิ สํวตฺตนฺติ, ปรพฺยาพาธายปิ สํวตฺตนฺติ, อุภย\-พฺยาพาธายปิ สํวตฺตนฺติ. กตเม ตโย? กายทุจฺจริตํ วจีทุจฺจริตํ มโนทุจฺจริตํ. อิเม โข ภิกฺขเว ตโย ธมฺมา อตฺตพฺยาพาธายปิ สํวตฺตนฺติ, ปรพฺยาพาธายปิ สํวตฺตนฺติ, อุภย\-พฺยาพาธายปิ สํวตฺตนฺติ.}

\prose{ตโยเม ภิกฺขเว ธมฺมา เนว\-ตฺตพฺยาพาธายปิ สํวตฺตนฺติ, น ปรพฺยาพาธายปิ สํวตฺตนฺติ, น อุภย\-พฺยาพาธายปิ สํวตฺตนฺติ. กตเม ตโย? กายสุจริตํ วจีสุจริตํ มโนสุจริตํ. อิเม โข ภิกฺขเว \csromanfolio{113}ตโย ธมฺมา เนว\-ตฺตพฺยาพาธายปิ สํวตฺตนฺติ, น ปรพฺยาพาธายปิ สํวตฺตนฺติ, น อุภย\-พฺยาพาธายปิ สํวตฺตนฺตีติ.  สตฺตมํ.}

\csromanheader{2}{8. \textbf{เทวโลกสุตฺต}}
\proseitem{18}{สเจ โว ภิกฺขเว อญฺญติตฺถิยา ปริพฺพาชกา เอวํ ปุจฺเฉยฺยุํ “เทวโลกู\-ปปตฺติยา อาวุโส สมเณ โคตเม พฺรหฺมจริยํ วุสฺสถา”ติ. นนุ ตุเมฺห ภิกฺขเว เอวํ ปุฏฺฐา อฏฺฏีเยยฺยาถ หราเยยฺยาถ ชิคุจฺเฉยฺยาถาติ. เอวํ ภนฺเต. อิติ กิร ตุเมฺห ภิกฺขเว ทิพฺเพน อายุนา อฏฺฏียถ หรายถ ชิคุจฺฉถ, ทิพฺเพน วณฺเณน, ทิพฺเพน สุเขน, ทิพฺเพน ยเสน, ทิพฺเพนา\-ธิปเตยฺเยน อฏฺฏียถ หรายถ ชิคุจฺฉถ, ปเคว โข ปน ภิกฺขเว ตุเมฺหหิ กาย\-ทุจฺจริเตน อฏฺฏียิตพฺพํ หรายิตพฺพํ ชิคุจฺฉิตพฺพํ. วจีทุจฺจริเตน. มโน\-ทุจฺจริเตน อฏฺฏียิตพฺพํ หรายิตพฺพํ ชิคุจฺฉิ\-ตพฺพนฺติ.  อฏฺฐมํ.}

\csromanheader{2}{9. ปฐม\-ปาปณิก\-สุตฺต}
\proseitem{19}{ตีหิ ภิกฺขเว องฺเคหิ สมนฺนาคโต ปาปณิโก อภพฺโพ อนธิคตํ วา โภคํ อธิคนฺตุํ, อธิคตํ วา โภคํ ผาติํ กาตุํ. กตเมหิ ตีหิ? อิธ ภิกฺขเว ปาปณิโก ปุพฺพณฺหสมยํ น สกฺกจฺจํ กมฺมนฺตํ อธิฏฺฐาติ, มชฺฌนฺหิกสมยํ\csromansharedfootnote{8410d587878334f4}{มชฺฌนฺติกสมยํ (สี, สฺยา, กํ, อิ)} น สกฺกจฺจํ กมฺมนฺตํ อธิฏฺฐาติ, สายนฺหสมยํ น สกฺกจฺจํ กมฺมนฺตํ อธิฏฺฐาติ. อิเมหิ โข ภิกฺขเว ตีหิ องฺเคหิ สมนฺนาคโต ปาปณิโก อภพฺโพ อนธิคตํ วา โภคํ อธิคนฺตุํ, อธิคตํ วา โภคํ ผาติํ กาตุํ\csromansharedfootnote{640036299927b0ac}{ผาติกตฺตุํ (สี), ผาติกาตุํ (สฺยา, กํ, อิ)}.}

\prose{เอวเมวํ โข ภิกฺขเว ตีหิ ธมฺเมหิ สมนฺนาคโต ภิกฺขุ อภพฺโพ อนธิคตํ วา กุสลํ ธมฺมํ อธิคนฺตุํ, อธิคตํ วา กุสลํ ธมฺมํ ผาติํ กาตุํ. กตเมหิ ตีหิ? อิธ ภิกฺขเว ภิกฺขุ ปุพฺพณฺหสมยํ น สกฺกจฺจํ สมาธิ\-นิมิตฺตํ อธิฏฺฐาติ, มชฺฌนฺหิกสมยํ น สกฺกจฺจํ สมาธิ\-นิมิตฺตํ อธิฏฺฐาติ, สายนฺหสมยํ น สกฺกจฺจํ สมาธิ\-นิมิตฺตํ อธิฏฺฐาติ. อิเมหิ โข ภิกฺขเว ตีหิ ธมฺเมหิ สมนฺนาคโต ภิกฺขุ อภพฺโพ อนธิคตํ วา กุสลํ ธมฺมํ อธิคนฺตุํ, อธิคตํ วา กุสลํ ธมฺมํ ผาติํ กาตุํ.}

\prose{\csromanfolio{114}ตีหิ ภิกฺขเว องฺเคหิ สมนฺนาคโต ปาปณิโก ภพฺโพ อนธิคตํ วา โภคํ อธิคนฺตุํ, อธิคตํ วา โภคํ ผาติํ กาตุํ. กตเมหิ ตีหิ? อิธ ภิกฺขเว ปาปณิโก ปุพฺพณฺหสมยํ สกฺกจฺจํ กมฺมนฺตํ อธิฏฺฐาติ, มชฺฌนฺหิกสมยํ ฯเปฯ สายนฺหสมยํ สกฺกจฺจํ กมฺมนฺตํ อธิฏฺฐาติ. อิเมหิ โข ภิกฺขเว ตีหิ องฺเคหิ สมนฺนาคโต ปาปณิโก ภพฺโพ อนธิคตํ วา โภคํ อธิคนฺตุํ, อธิคตํ วา โภคํ ผาติํ กาตุํ.}

\prose{เอวเมวํ โข ภิกฺขเว ตีหิ ธมฺเมหิ สมนฺนาคโต ภิกฺขุ ภพฺโพ อนธิคตํ วา กุสลํ ธมฺมํ อธิคนฺตุํ, อธิคตํ วา กุสลํ ธมฺมํ ผาติํ กาตุํ. กตเมหิ ตีหิ? อิธ ภิกฺขเว ภิกฺขุ ปุพฺพณฺหสมยํ สกฺกจฺจํ สมาธิ\-นิมิตฺตํ อธิฏฺฐาติ, มชฺฌนฺหิกสมยํ ฯเปฯ สายนฺหสมยํ สกฺกจฺจํ สมาธิ\-นิมิตฺตํ อธิฏฺฐาติ. อิเมหิ โข ภิกฺขเว ตีหิ ธมฺเมหิ สมนฺนาคโต ภิกฺขุ ภพฺโพ อนธิคตํ วา กุสลํ ธมฺมํ อธิคนฺตุํ, อธิคตํ วา กุสลํ ธมฺมํ ผาติํ กาตุนฺติ.  นวมํ.}

\csromanheader{2}{10. ทุติย\-ปาปณิก\-สุตฺต}
\proseitem{20}{ตีหิ ภิกฺขเว องฺเคหิ สมนฺนาคโต ปาปณิโก นจิรสฺเสว มหตฺตํ เวปุลฺลตฺตํ\csromansharedfootnote{1a2f93e2c3e22af0}{มหนฺตตฺตํ วา เวปุลฺลตฺตํ วา (อิ, ก)} ปาปุณาติ โภเคสุ. กตเมหิ ตีหิ? อิธ ภิกฺขเว ปาปณิโก จกฺขุมา จ โหติ วิธุโร จ นิสฺสย\-สมฺปนฺโน จ. กถญฺจ ภิกฺขเว ปาปณิโก จกฺขุมา โหติ? อิธ ภิกฺขเว ปาปณิโก ปณิยํ ชานาติ “อิทํ ปณิยํ เอวํ กีตํ, เอวํ วิกฺกยมานํ\csromansharedfootnote{17380c69a41bef59}{วิกฺกียมานํ (?)} เอตฺตกํ มูลํ ภวิสฺสติ เอตฺตโก อุทโย”ติ\csromansharedfootnote{f8f1f762dd4c2459}{อุทฺทโยติ (สี)}. เอวํ โข ภิกฺขเว ปาปณิโก จกฺขุมา โหติ.}

\prose{กถญฺจ ภิกฺขเว ปาปณิโก วิธุโร โหติ? อิธ ภิกฺขเว ปาปณิโก กุสโล โหติ ปณิยํ เกตุญฺจ วิกฺเกตุญฺจ. เอวํ โข ภิกฺขเว ปาปณิโก วิธุโร โหติ.}

\prose{กถญฺจ ภิกฺขเว ปาปณิโก นิสฺสย\-สมฺปนฺโน โหติ? อิธ ภิกฺขเว ปาปณิกํ เย เต คหปตี วา คหปติปุตฺตา วา อฑฺฒา มหทฺธนา มหาโภคา, เต เอวํ ชานนฺติ “อยํ โข ภวํ ปาปณิโก จกฺขุมา \csromanfolio{115}วิธุโร จ ปฏิพโล ปุตฺตทารญฺจ โปเสตุํ, อมฺหากญฺจ กาเลน กาลํ อนุปฺปทาตุนฺ”ติ. เต นํ โภเคหิ นิปตนฺติ “อิโต สมฺม ปาปณิก โภเค กริตฺวา\csromansharedfootnote{ca6ae6beceafc266}{หริตฺวา (สี, สฺยา, กํ)} ปุตฺตทารญฺจ โปเสหิ, อมฺหากญฺจ กาเลน กาลํ อนุปฺปเทหี”ติ. เอวํ โข ภิกฺขเว ปาปณิโก นิสฺสย\-สมฺปนฺโน โหติ. อิเมหิ โข ภิกฺขเว ตีหิ องฺเคหิ สมนฺนาคโต ปาปณิโก นจิรสฺเสว มหตฺตํ เวปุลฺลตฺตํ ปาปุณาติ โภเคสุ.}

\prose{เอวเมวํ โข ภิกฺขเว ตีหิ ธมฺเมหิ สมนฺนาคโต ภิกฺขุ นจิรสฺเสว มหตฺตํ เวปุลฺลตฺตํ ปาปุณาติ กุสเลสุ ธมฺเมสุ. กตเมหิ ตีหิ? อิธ ภิกฺขเว ภิกฺขุ จกฺขุมา จ โหติ วิธุโร จ นิสฺสย\-สมฺปนฺโน จ. กถญฺจ ภิกฺขเว ภิกฺขุ จกฺขุมา โหติ? อิธ ภิกฺขเว ภิกฺขุ “อิทํ ทุกฺขนฺ”ติ ยถาภูตํ ปชานาติ, “อยํ ทุกฺขสมุทโย”ติ ยถาภูตํ ปชานาติ, “อยํ ทุกฺขนิโรโธ”ติ ยถาภูตํ ปชานาติ, “อยํ ทุกฺขนิโรธ\-คามินี ปฏิปทา”ติ ยถาภูตํ ปชานาติ. เอวํ โข ภิกฺขเว ภิกฺขุ จกฺขุมา โหติ.}

\prose{กถญฺจ ภิกฺขเว ภิกฺขุ วิธุโร โหติ? อิธ ภิกฺขเว ภิกฺขุ อารทฺธวีริโย วิหรติ อกุสลานํ ธมฺมานํ ปหานาย, กุสลานํ ธมฺมานํ อุปสมฺปทาย ถามวา ทฬฺหปรกฺกโม อนิกฺขิตฺตธุโร กุสเลสุ ธมฺเมสุ. เอวํ โข ภิกฺขเว ภิกฺขุ วิธุโร โหติ.}

\prosewithclosermidruled{กถญฺจ ภิกฺขเว ภิกฺขุ นิสฺสย\-สมฺปนฺโน โหติ? อิธ ภิกฺขเว ภิกฺขุ เย เต ภิกฺขู พหุสฺสุตา อาคตาคมา ธมฺมธรา วินยธรา มาติกาธรา, เต กาเลน กาลํ อุปสงฺกมิตฺวา ปริปุจฺฉติ ปริปญฺหติ “อิทํ ภนฺเต กถํ, อิมสฺส โก อตฺโถ”ติ. ตสฺส เต อายสฺมนฺโต อวิวฏํ เจว วิวรนฺติ, อนุตฺตานีกตญฺจ อุตฺตานีกโรนฺติ, อเนกวิหิเตสุ จ กงฺขา\-ฐานิเยสุ ธมฺเมสุ กงฺขํ ปฏิวิโนเทนฺติ. เอวํ โข ภิกฺขเว ภิกฺขุ นิสฺสย\-สมฺปนฺโน โหติ. อิเมหิ โข ภิกฺขเว ตีหิ ธมฺเมหิ สมนฺนาคโต ภิกฺขุ นจิรสฺเสว มหตฺตํ เวปุลฺลตฺตํ ปาปุณาติ กุสเลสุ ธมฺเมสูติ.  ทสมํ.}{รถการวคฺโค ทุติโย.}
\nitthitam{ปฐม\-ภาณวาโร นิฏฺฐิโต.}
\csromanpalirecto
\gambhira{ติกนิปาตปาฬิ \textbf{ตสฺสุทฺทานํ}}
\setlength{\csromangathapagewidth}{0pt}
\setlength{\csromangathawakwidth}{0pt}
\csromangathameasuregroup{ญาโต\textsuperscript{1} สารณีโย ภิกฺขุ, \\ อปณฺณกตฺตา เทโว จ,}{\csromangathabat{ญาโต\textsuperscript{1} สารณีโย ภิกฺขุ,}{จกฺกวตฺตี สเจตโน.} \\ \csromangathabat{อปณฺณกตฺตา เทโว จ,}{ทุเว ปาปณิเกน จาติ.}}
\csromangathasetleft{ญาโต\textsuperscript{1} สารณีโย ภิกฺขุ, \\ อปณฺณกตฺตา เทโว จ,}
\csromangathagroup{\csromangathastanza{\csromanfolio{116}\csromangathabat{ญาโต\csromansharedfootnote{7aedd28a4643db72}{ญาตโก (สฺยา, กํ)} สารณีโย ภิกฺขุ,}{จกฺกวตฺตี สเจตโน.} \\ \csromangathabat{อปณฺณกตฺตา เทโว จ,}{ทุเว ปาปณิเกน จาติ.}}}

\chapterhead{3. \textbf{ปุคฺคลวคฺค}}
\csromanheader{2}{1. \textbf{สมิทฺธสุตฺต}}
\proseitem{21}{เอวํ เม สุตํ\csromandash{}เอกํ สมยํ ภควา สาวตฺถิยํ วิหรติ เชตวเน อนาถ\-ปิณฺฑิกสฺส อาราเม. อถ โข อายสฺมา จ สมิทฺโธ,\csromansharedfootnote{102ce71ce39a4276}{สวิฏฺโฐ (สี, สฺยา, กํ, อิ)} อายสฺมา จ มหาโกฏฺฐิโก\csromansharedfootnote{b8379db2e670fe9e}{มหาโกฏฺฐิโต (สี, สฺยา, กํ, อิ)} เยนายสฺมา สาริปุตฺโต เตนุ\-ปสงฺกมิํสุ, อุปสงฺกมิตฺวา อายสฺมตา สาริปุตฺเตน สทฺธิํ สมฺโมทิํสุ, สมฺโมทนียํ กถํ สารณียํ วีติสาเรตฺวา เอกมนฺตํ นิสีทิํสุ, เอกมนฺตํ นิสินฺนํ โข อายสฺมนฺตํ สมิทฺธํ อายสฺมา สาริปุตฺโต เอตทโวจ\csromandash{}}

\prose{“ตโยเม อาวุโส สมิทฺธ ปุคฺคลา สนฺโต สํวิชฺชมานา โลกสฺมิํ. กตเม ตโย? กายสกฺขี ทิฏฺฐิปฺปตฺโต\csromansharedfootnote{d26e210e51614829}{ทิฏฺฐปฺปตฺโต (ก)} สทฺธาวิมุตฺโต. อิเม โข อาวุโส ตโย ปุคฺคลา สนฺโต สํวิชฺชมานา โลกสฺมิํ. อิเมสํ อาวุโส ติณฺณํ ปุคฺคลานํ กตโม เต ปุคฺคโล ขมติ อภิกฺกนฺตตโร จ ปณีตตโร จา”ติ.}

\prose{ตโยเม อาวุโส สาริปุตฺต ปุคฺคลา สนฺโต สํวิชฺชมานา โลกสฺมิํ. กตเม ตโย? กายสกฺขี ทิฏฺฐิปฺปตฺโต สทฺธาวิมุตฺโต. อิเม โข อาวุโส ตโย ปุคฺคลา สนฺโต สํวิชฺชมานา โลกสฺมิํ. อิเมสํ อาวุโส ติณฺณํ ปุคฺคลานํ ยฺวายํ\csromansharedfootnote{846a1fc9a87f8ba2}{โยยํ (ก)} ปุคฺคโล สทฺธาวิมุตฺโต, อยํ เม ปุคฺคโล ขมติ อิเมสํ ติณฺณํ ปุคฺคลานํ อภิกฺกนฺตตโร จ ปณีตตโร จ. ตํ กิสฺส เหตุ? อิมสฺส อาวุโส ปุคฺคลสฺส สทฺธินฺทฺริยํ อธิมตฺตนฺติ.}

\prose{อถ โข อายสฺมา สาริปุตฺโต อายสฺมนฺตํ มหาโกฏฺฐิกํ เอตทโวจ “ตโยเม อาวุโส โกฏฺฐิก ปุคฺคลา สนฺโต สํวิชฺชมานา \csromanfolio{117}โลกสฺมิํ. กตเม ตโย? กายสกฺขี ทิฏฺฐิปฺปตฺโต สทฺธาวิมุตฺโต. อิเม โข อาวุโส ตโย ปุคฺคลา สนฺโต สํวิชฺชมานา โลกสฺมิํ. อิเมสํ อาวุโส ติณฺณํ ปุคฺคลานํ กตโม เต ปุคฺคโล ขมติ อภิกฺกนฺตตโร จ ปณีตตโร จา”ติ.}

\prose{ตโยเม อาวุโส สาริปุตฺต ปุคฺคลา สนฺโต สํวิชฺชมานา โลกสฺมิํ. กตเม ตโย? กายสกฺขี ทิฏฺฐิปฺปตฺโต สทฺธาวิมุตฺโต. อิเม โข อาวุโส ตโย ปุคฺคลา สนฺโต สํวิชฺชมานา โลกสฺมิํ. อิเมสํ อาวุโส ติณฺณํ ปุคฺคลานํ ยฺวายํ ปุคฺคโล กายสกฺขี, อยํ เม ปุคฺคโล ขมติ อิเมสํ ติณฺณํ ปุคฺคลานํ อภิกฺกนฺตตโร จ ปณีตตโร จ. ตํ กิสฺส เหตุ? อิมสฺส อาวุโส ปุคฺคลสฺส สมาธินฺทฺริยํ อธิมตฺตนฺติ.}

\prose{อถ โข อายสฺมา มหาโกฏฺฐิโก อายสฺมนฺตํ สาริปุตฺตํ เอตทโวจ “ตโยเม อาวุโส สาริปุตฺต ปุคฺคลา สนฺโต สํวิชฺชมานา โลกสฺมิํ. กตเม ตโย? กายสกฺขี ทิฏฺฐิปฺปตฺโต สทฺธาวิมุตฺโต. อิเม โข อาวุโส ตโย ปุคฺคลา สนฺโต สํวิชฺชมานา โลกสฺมิํ. อิเมสํ อาวุโส ติณฺณํ ปุคฺคลานํ กตโม เต ปุคฺคโล ขมติ อภิกฺกนฺตตโร จ ปณีตตโร จา”ติ.}

\prose{ตโยเม อาวุโส โกฏฺฐิก ปุคฺคลา สนฺโต สํวิชฺชมานา โลกสฺมิํ. กตเม ตโย? กายสกฺขี ทิฏฺฐิปฺปตฺโต สทฺธาวิมุตฺโต. อิเม โข อาวุโส ตโย ปุคฺคลา สนฺโต สํวิชฺชมานา โลกสฺมิํ. อิเมสํ อาวุโส ติณฺณํ ปุคฺคลานํ ยฺวายํ ปุคฺคโล ทิฏฺฐิปฺปตฺโต, อยํ เม ปุคฺคโล ขมติ อิเมสํ ติณฺณํ ปุคฺคลานํ อภิกฺกนฺตตโร จ ปณีตตโร จ. ตํ กิสฺส เหตุ? อิมสฺส อาวุโส ปุคฺคลสฺส ปญฺญินฺทฺริยํ อธิมตฺตนฺติ.}

\prose{อถ โข อายสฺมา สาริปุตฺโต อายสฺมนฺตญฺจ สมิทฺธํ, อายสฺมนฺตญฺจ มหาโกฏฺฐิกํ เอตทโวจ “พฺยากตํ โข อาวุโส อเมฺหหิ สพฺเพเหว ยถาสกํ ปฏิภานํ, อายามาวุโส, เยน ภควา เตนุปสงฺกมิสฺสาม, อุปสงฺกมิตฺวา ภควโต เอตมตฺถํ อาโรเจสฺสาม, ยถา โน ภควา พฺยากริสฺสติ, ตถา นํ ธาเรสฺสามา”ติ. “เอวมาวุโส”ติ โข อายสฺมา จ สมิทฺโธ, อายสฺมา จ มหาโกฏฺฐิโก อายสฺมโต \csromanfolio{118}สาริปุตฺตสฺส ปจฺจสฺโสสุํ. อถ โข อายสฺมา จ สาริปุตฺโต, อายสฺมา จ สมิทฺโธ, อายสฺมา จ มหาโกฏฺฐิโก เยน ภควา เตนุ\-ปสงฺกมิํสุ, อุปสงฺกมิตฺวา ภควนฺตํ อภิวาเทตฺวา เอกมนฺตํ นิสีทิํสุ, เอกมนฺตํ นิสินฺโน โข อายสฺมา สาริปุตฺโต ยาวตโก อโหสิ อายสฺมตา จ สมิทฺเธน, อายสฺมตา จ มหาโกฏฺฐิเกน สทฺธิํ กถาสลฺลาโป, ตํ สพฺพํ ภควโต อาโรเจสิ.}

\proseitem{3}{น เขฺวตฺถ สาริปุตฺต สุกรํ เอกํเสน พฺยากาตุํ “อยํ อิเมสํ ติณฺณํ ปุคฺคลานํ อภิกฺกนฺตตโร จ ปณีตตโร จา”ติ. ฐานํ เหตํ สาริปุตฺต วิชฺชติ. ยฺวายํ ปุคฺคโล สทฺธาวิมุตฺโต, สฺวาสฺส\csromansharedfootnote{85e48d80c91276ed}{สฺวายํ (สฺยา, กํ, อิ), โสยํ (ก)} อรหตฺตาย ปฏิปนฺโน, ยฺวายํ ปุคฺคโล กายสกฺขี, สฺวาสฺส สกทาคามี วา อนาคามี วา, โย จายํ ปุคฺคโล ทิฏฺฐิปฺปตฺโต, โสปสฺส\csromansharedfootnote{578cb219a4e452bc}{โสยํ (ก)} สกทาคามี วา อนาคามี วา.}

\prose{น เขฺวตฺถ สาริปุตฺต สุกรํ เอกํเสน พฺยากาตุํ “อยํ อิเมสํ ติณฺณํ ปุคฺคลานํ อภิกฺกนฺตตโร จ ปณีตตโร จา”ติ. ฐานํ เหตํ สาริปุตฺต วิชฺชติ, ยฺวายํ ปุคฺคโล กายสกฺขี, สฺวาสฺส อรหตฺตาย ปฏิปนฺโน, ยฺวายํ ปุคฺคโล สทฺธาวิมุตฺโต, สฺวาสฺส สกทาคามี วา อนาคามี วา, โย จายํ ปุคฺคโล ทิฏฺฐิปฺปตฺโต, โสปสฺส สกทาคามี วา อนาคามี วา.}

\prose{น เขฺวตฺถ สาริปุตฺต สุกรํ เอกํเสน พฺยากาตุํ “อยํ อิเมสํ ติณฺณํ ปุคฺคลานํ อภิกฺกนฺตโร จ ปณีตตโร จา”ติ. ฐานํ เหตํ สาริปุตฺต วิชฺชติ, ยฺวายํ ปุคฺคโล ทิฏฺฐิปฺปตฺโต, สฺวาสฺส อรหตฺตาย ปฏิปนฺโน, ยฺวายํ ปุคฺคโล สทฺธาวิมุตฺโต, สฺวาสฺส สกทาคามี วา อนาคามี วา, โย จายํ ปุคฺคโล กายสกฺขี, โสปสฺส สกทาคามี วา อนาคามี วา.}

\prose{น เขฺวตฺถ สาริปุตฺต สุกรํ เอกํเสน พฺยากาตุํ “อยํ อิเมสํ ติณฺณํ ปุคฺคลานํ อภิกฺกนฺตตโร จ ปณีตตโร จา”ติ.  ปฏฺฐมํ.}

\csromanheader{2}{2. \textbf{คิลานสุตฺต}}
\proseitem{22}{3 ตโยเม ภิกฺขเว คิลานา สนฺโต สํวิชฺชมานา โลกสฺมิํ. กตเม ตโย? อิธ ภิกฺขเว เอกจฺโจ คิลาโน ลภนฺโต วา สปฺปายานิ โภชนานิ อลภนฺโต วา สปฺปายานิ โภชนานิ, ลภนฺโต วา สปฺปายานิ เภสชฺชานิ อลภนฺโต วา สปฺปายานิ เภสชฺชานิ, ลภนฺโต \csromanfolio{119}วา ปติรูปํ อุปฏฺฐากํ อลภนฺโต วา ปติรูปํ อุปฏฺฐากํ เนว วุฏฺฐาติ ตมฺหา อาพาธา.}

\prose{อิธ ปน ภิกฺขเว เอกจฺโจ คิลาโน ลภนฺโต วา สปฺปายานิ โภชนานิ อลภนฺโต วา สปฺปายานิ โภชนานิ, ลภนฺโต วา สปฺปายานิ เภสชฺชานิ อลภนฺโต วา สปฺปายานิ เภสชฺชานิ, ลภนฺโต วา ปติรูปํ อุปฏฺฐากํ อลภนฺโต วา ปติรูปํ อุปฏฺฐากํ วุฏฺฐาติ ตมฺหา อาพาธา.}

\prose{อิธ ปน ภิกฺขเว เอกจฺโจ คิลาโน ลภนฺโตว สปฺปายานิ โภชนานิ โน อลภนฺโต, ลภนฺโตว สปฺปายานิ เภสชฺชานิ โน อลภนฺโต, ลภนฺโตว ปติรูปํ อุปฏฺฐากํ โน อลภนฺโต วุฏฺฐาติ ตมฺหา อาพาธา.}

\prose{ตตฺร ภิกฺขเว ยฺวายํ คิลาโน ลภนฺโตว สปฺปายานิ โภชนานิ โน อลภนฺโต, ลภนฺโตว สปฺปายานิ เภสชฺชานิ โน อลภนฺโต, ลภนฺโตว ปติรูปํ อุปฏฺฐากํ โน อลภนฺโต วุฏฺฐาติ ตมฺหา อาพาธา. อิมํ โข ภิกฺขเว คิลานํ ปฏิจฺจ คิลานภตฺตํ อนุญฺญาตํ, คิลาน\-เภสชฺชํ อนุญฺญาตํ, คิลานุปฏฺฐาโก อนุญฺญาโต. อิมญฺจ ปน ภิกฺขเว คิลานํ ปฏิจฺจ อญฺเญปิ คิลานา อุปฏฺฐาตพฺพา. อิเม โข ภิกฺขเว ตโย คิลานา สนฺโต สํวิชฺชมานา โลกสฺมิํ.}

\prose{เอวเมวํ โข ภิกฺขเว ตโยเม คิลานูปมา ปุคฺคลา สนฺโต สํวิชฺชมานา โลกสฺมิํ. กตเม ตโย? อิธ ภิกฺขเว เอกจฺโจ ปุคฺคโล ลภนฺโต วา ตถาคตํ ทสฺสนาย อลภนฺโต วา ตถาคตํ ทสฺสนาย, ลภนฺโต วา ตถาคต\-ปฺปเวทิตํ ธมฺมวินยํ สวนาย อลภนฺโต วา ตถาคต\-ปฺปเวทิตํ ธมฺมวินยํ สวนาย เนว โอกฺกมติ นิยามํ กุสเลสุ ธมฺเมสุ สมฺมตฺตํ.}

\prose{อิธ ปน ภิกฺขเว เอกจฺโจ ปุคฺคโล ลภนฺโต วา ตถาคตํ ทสฺสนาย อลภนฺโต วา ตถาคตํ ทสฺสนาย, ลภนฺโต วา ตถาคต\-ปฺปเวทิตํ ธมฺมวินยํ สวนาย อลภนฺโต วา ตถาคต\-ปฺปเวทิตํ ธมฺมวินยํ สวนาย โอกฺกมติ นิยามํ กุสเลสุ ธมฺเมสุ สมฺมตํ.}

\prose{อิธ ปน ภิกฺขเว เอกจฺโจ ปุคฺคโล ลภนฺโตว ตถาคตํ ทสฺสานาย โน อลภนฺโต, ลภนฺโตว ตถาคต\-ปฺปเวทิตํ ธมฺมวินยํ สวนาย โน อลภนฺโต โอกฺกมติ นิยามํ กุสเลสุ ธมฺเมสุ สมฺมตฺตํ.}

\prose{\csromanfolio{120}ตตฺร ภิกฺขเว ยฺวายํ ปุคฺคโล ลภนฺโตว ตถาคตํ ทสฺสนาย โน อลภนฺโต, ลภนฺโตว ตถาคต\-ปฺปเวทิตํ ธมฺมวินยํ สวนาย โน อลภนฺโต โอกฺกมติ นิยามํ กุสเลสุ ธมฺเมสุ สมฺมตฺตํ. อิมํ โข ภิกฺขเว ปุคฺคลํ ปฏิจฺจ ธมฺมเทสนา อนุญฺญาตา, อิมญฺจ ปน ภิกฺขเว ปุคฺคลํ ปฏิจฺจ อญฺเญสมฺปิ ธมฺโม เทเสตพฺโพ. อิเม โข ภิกฺขเว ตโย คิลานูปมา ปุคฺคลา สนฺโต สํวิชฺชมานา โลกสฺมินฺติ.  ทุติยํ.}

\csromanheader{2}{3. \textbf{สงฺขารสุตฺต}}
\proseitem{23}{ตโยเม ภิกฺขเว ปุคฺคลา สนฺโต สํวิชฺชมานา โลกสฺมิํ. กตเม ตโย? อิธ ภิกฺขเว เอกจฺโจ ปุคฺคโล สพฺยาพชฺฌํ\csromansharedfootnote{1a720caa051e79b2}{สพฺยาปชฺฌํ (สพฺพตฺถ) เอวมุปริปิ.} กายสงฺขารํ อภิสงฺขโรติ, สพฺยาพชฺฌํ วจีสงฺขารํ อภิสงฺขโรติ, สพฺยาพชฺฌํ มโนสงฺขารํ อภิสงฺขโรติ. โส สพฺยาพชฺฌํ กายสงฺขารํ อภิส\-งฺขริ\-ตฺวา สพฺยาพชฺฌํ วจีสงฺขารํ อภิส\-งฺขริ\-ตฺวา สพฺยาพชฺฌํ มโนสงฺขารํ อภิส\-งฺขริ\-ตฺวา สพฺยาพชฺฌํ โลกํ อุปปชฺชติ. ตเมนํ สพฺยาพชฺฌํ โลกํ อุปปนฺนํ สมานํ สพฺยาพชฺฌา ผสฺสา ผุสนฺติ. โส สพฺยาพชฺเฌหิ ผสฺเสหิ ผุฏฺโฐ สมาโน สพฺยาพชฺฌํ เวทนํ เวทยติ เอกนฺตทุกฺขํ, เสยฺยถาปิ สตฺตา เนรยิกา.}

\prose{อิธ ปน ภิกฺขเว เอกจฺโจ ปุคฺคโล อพฺยาพชฺฌํ กายสงฺขารํ อภิสงฺขโรติ, อพฺยาพชฺฌํ วจีสงฺขารํ อภิสงฺขโรติ, อพฺยาพชฺฌํ มโนสงฺขารํ อภิสงฺขโรติ. โส อพฺยาพชฺฌํ กายสงฺขารํ อภิส\-งฺขริ\-ตฺวา อพฺยาพชฺฌํ วจีสงฺขารํ อภิส\-งฺขริ\-ตฺวา อพฺยาพชฺฌํ มโนสงฺขารํ อภิส\-งฺขริ\-ตฺวา อพฺยาพชฺฌํ โลกํ อุปปชฺชติ. ตเมนํ อพฺยาพชฺฌํ โลกํ อุปปนฺนํ สมานํ อพฺยาพชฺฌา ผสฺสา ผุสนฺติ. โส อพฺยาพชฺเฌหิ ผสฺเสหิ ผุฏฺโฐ สมาโน อพฺยาพชฺฌํ เวทนํ เวทยติ เอกนฺตสุขํ, เสยฺยถาปิ เทวา สุภกิณฺหา.}

\prose{อิธ ปน ภิกฺขเว เอกจฺโจ ปุคฺคโล สพฺยาพชฺฌมฺปิ อพฺยาพชฺฌมฺปิ กายสงฺขารํ อภิสงฺขโรติ, สพฺยาพชฺฌมฺปิ อพฺยาพชฺฌมฺปิ วจีสงฺขารํ อภิสงฺขโรติ, สพฺยาพชฺฌมฺปิ อพฺยาพชฺฌมฺปิ มโนสงฺขารํ อภิสงฺขโรติ. โส สพฺยาพชฺฌมฺปิ อพฺยาพชฺฌมฺปิ กายสงฺขารํ อภิส\-งฺขริ\-ตฺวา สพฺยาพชฺฌมฺปิ อพฺยาพชฺฌมฺปิ วจีสงฺขารํ อภิส\-งฺขริ\-ตฺวา สพฺยาพชฺฌมฺปิ อพฺยาพชฺฌมฺปิ มโนสงฺขารํ อภิส\-งฺขริ\-ตฺวา สพฺยาพชฺฌมฺปิ อพฺยาพชฺฌมฺปิ โลกํ อุปปชฺชติ. ตเมนํ สพฺยาพชฺฌมฺปิ อพฺยาพชฺฌมฺปิ \csromanfolio{121}โลกํ อุปปนฺนํ สมานํ สพฺยาพชฺฌาปิ อพฺยาพชฺฌาปิ ผสฺสา ผุสนฺติ. โส สพฺยาพชฺเฌหิปิ อพฺยาพชฺเฌหิปิ ผสฺเสหิ ผุฏฺโฐ สมาโน สพฺยาพชฺฌมฺปิ อพฺยาพชฺฌมฺปิ เวทนํ เวทยติ โวกิณฺณสุขทุกฺขํ, เสยฺยถาปิ มนุสฺสา เอกจฺเจ จ เทวา เอกจฺเจ จ วินิปาติกา. อิเม โข ภิกฺขเว ตโย ปุคฺคลา สนฺโต สํวิชฺชมานา โลกสฺมินฺติ.  ตติยํ.}

\csromanheader{2}{4. \textbf{พหุการสุตฺต}}
\proseitem{24}{ตโยเม ภิกฺขเว ปุคฺคลา ปุคฺคลสฺส พหุการา. กตเม ตโย? ยํ ภิกฺขเว ปุคฺคลํ อาคมฺม ปุคฺคโล พุทฺธํ สรณํ คโต โหติ, ธมฺมํ สรณํ คโต โหติ, สํฆํ สรณํ คโต โหติ. อยํ ภิกฺขเว ปุคฺคโล อิมสฺส ปุคฺคลสฺส พหุกาโร.}

\prose{ปุน จปรํ ภิกฺขเว ยํ ปุคฺคลํ อาคมฺม ปุคฺคโล “อิทํ ทุกฺขนฺ”ติ ยถาภูตํ ปชานาติ, “อยํ ทุกฺขสมุทโย”ติ ยถาภูตํ ปชานาติ, “อยํ ทุกฺขนิโรโธ”ติ ยถาภูตํ ปชานาติ, “อยํ ทุกฺขนิโรธ\-คามินี ปฏิปทา”ติ ยถาภูตํ ปชานาติ. อยํ ภิกฺขเว ปุคฺคโล อิมสฺส ปุคฺคลสฺส พหุกาโร.}

\prose{ปุน จปรํ ภิกฺขเว ยํ ปุคฺคลํ อาคมฺม ปุคฺคโล อาสวานํ ขยา อนาสวํ เจโตวิมุตฺติํ ปญฺญาวิมุตฺติํ ทิฏฺเฐว ธมฺเม สยํ อภิญฺญา สจฺฉิกตฺวา อุปสมฺปชฺช วิหรติ. อยํ ภิกฺขเว ปุคฺคโล อิมสฺส ปุคฺคลสฺส พหุกาโร. อิเม โข ภิกฺขเว ตโย ปุคฺคลา ปุคฺคลสฺส พหุการา.}

\prose{อิเมหิ จ ปน ภิกฺขเว ตีหิ ปุคฺคเลหิ อิมสฺส ปุคฺคลสฺส นตฺถญฺโญ ปุคฺคโล พหุกาโรติ วทามิ. อิเมสํ ปน ภิกฺขเว ติณฺณํ ปุคฺคลานํ อิมินา ปุคฺคเลน น สุปฺปติการํ วทามิ, ยทิทํ อภิวาทน\-ปจฺจุฏฺฐาน\-อญฺชลิกมฺม\-สามีจิกมฺม\-จีวร\-ปิณฺฑปาต\-เสนาสน- คิลานปจฺจย\-เภสชฺชปริกฺขารา\-นุปฺปทาเนนาติ.  จตุตฺถํ.}

\csromanheader{2}{5. วชิรูปม\-สุตฺต}
\proseitem{25}{ตโยเม ภิกฺขเว ปุคฺคลา สนฺโต สํวิชฺชมานา โลกสฺมิํ. กตเม ตโย? อรุกูปมจิตฺโต ปุคฺคโล, วิชฺชูปมจิตฺโต ปุคฺคโล, วชิรูปม\-จิตฺโต ปุคฺคโล. กตโม จ ภิกฺขเว อรุกูปมจิตฺโต ปุคฺคโล? อิธ ภิกฺขเว เอกจฺโจ ปุคฺคโล โกธโน โหติ อุปายาสพหุโล, \csromanfolio{122}อปฺปมฺปิ วุตฺโต สมาโน อภิสชฺชติ กุปฺปติ พฺยาปชฺชติ ปติตฺถียติ, โกปญฺจ โทสญฺจ อปฺปจฺจยญฺจ ปาตุกโรติ. เสยฺยถาปิ ภิกฺขเว ทุฏฺฐารุโก\csromansharedfootnote{c8c9aaad04d63e22}{ทุฏฺฐารุกา (สี)} กฏฺเฐน วา กฐลาย\csromansharedfootnote{7bb204732127628d}{กถลาย (สฺยา, กํ, ก), กฐเลน - กถเลน (อฏฺฐกถา) 3. ฆฏฺฏิตา (สี)} วา ฆฏฺฏิโต\csromansharedfootnote{240e0831a04ac08b}{อสฺสวโนติ (สี)} ภิยฺโยโส มตฺตาย อาสวํ เทติ\csromansharedfootnote{9c65b1bd40da9dfb}{ฆฏฺฏิตา (สี)}. เอวเมวํ โข ภิกฺขเว อิเธกจฺโจ ปุคฺคโล โกธโน โหติ อุปายาสพหุโล, อปฺปมฺปิ วุตฺโต สมาโน อภิสชฺชติ กุปฺปติ พฺยาปชฺชติ ปติตฺถียติ, โกปญฺจ โทสญฺจ อปฺปจฺจยญฺจ ปาตุกโรติ. อยํ วุจฺจติ ภิกฺขเว อรุกูปมจิตฺโต ปุคฺคโล.}

\prose{กตโม จ ภิกฺขเว วิชฺชูปมจิตฺโต ปุคฺคโล? อิธ ภิกฺขเว เอกจฺโจ ปุคฺคโล “อิทํ ทุกฺขนฺ”ติ ยถาภูตํ ปชานาติ, “อยํ ทุกฺขสมุทโย”ติ ยถาภูตํ ปชานาติ, “อยํ ทุกฺขนิโรโธ”ติ ยถาภูตํ ปชานาติ, “อยํ ทุกฺขนิโรธ\-คามินี ปฏิปทา”ติ ยถาภูตํ ปชานาติ. เสยฺยถาปิ ภิกฺขเว จกฺขุมา ปุริโส รตฺต\-นฺธการ\-ติมิสายํ วิชฺชนฺตริกาย รูปานิ ปสฺเสยฺย. เอวเมวํ โข ภิกฺขเว อิเธกจฺโจ ปุคฺคโล “อิทํ ทุกฺขนฺ”ติ ยถาภูตํ ปชานาติ ฯเปฯ “อยํ ทุกฺขนิโรธ\-คามินี ปฏิปทา”ติ ยถาภูตํ ปชานาติ. อยํ วุจฺจติ ภิกฺขเว วิชฺชูปมจิตฺโต ปุคฺคโล.}

\prose{กตโม จ ภิกฺขเว วชิรูปม\-จิตฺโต ปุคฺคโล? อิธ ภิกฺขเว เอกจฺโจ ปุคฺคโล อาสวานํ ขยา อนาสวํ เจโตวิมุตฺติํ ปญฺญาวิมุตฺติํ ทิฏฺเฐว ธมฺเม สยํ อภิญฺญา สจฺฉิกตฺวา อุปสมฺปชฺช วิหรติ. เสยฺยถาปิ ภิกฺขเว วชิรสฺส นตฺถิ กิญฺจิ อเภชฺชํ มณิ วา ปาสาโณ วา. เอวเมวํ โข ภิกฺขเว อิเธกจฺโจ ปุคฺคโล อาสวานํ ขยา ฯเปฯ อุปสมฺปชฺช วิหรติ. อยํ วุจฺจติ ภิกฺขเว วชิรูปม\-จิตฺโต ปุคฺคโล. อิเม โข ภิกฺขเว ตโย ปุคฺคลา สนฺโต สํวิชฺชมานา โลกสฺมินฺติ\csromansharedfootnote{1501f58e287ec29a}{อภิ 3. 133 ปิฏฺเฐปิ.}.  ปญฺจมํ.}

\csromanheader{2}{6. \textbf{เสวิตพฺพสุตฺต}}
\proseitem{26}{ตโยเม ภิกฺขเว ปุคฺคลา สนฺโต สํวิชฺชมานา โลกสฺมิํ. กตเม ตโย? อตฺถิ ภิกฺขเว ปุคฺคโล น เสวิตพฺโพ น ภชิตพฺโพ น ปยิรุปาสิตพฺโพ, อตฺถิ ภิกฺขเว ปุคฺคโล เสวิตพฺโพ ภชิตพฺโพ ปยิรุปาสิตพฺโพ, อตฺถิ ภิกฺขเว ปุคฺคโล สกฺกตฺวา ครุํ กตฺวา เสวิตพฺโพ ภชิตพฺโพ ปยิรุปาสิตพฺโพ. กตโม จ ภิกฺขเว \csromanfolio{123}ปุคฺคโล น เสวิตพฺโพ น ภชิตพฺโพ น ปยิรุปาสิตพฺโพ? อิธ ภิกฺขเว เอกจฺโจ ปุคฺคโล หีโน โหติ สีเลน สมาธินา ปญฺญาย. เอวรูโป ภิกฺขเว ปุคฺคโล น เสวิตพฺโพ น ภชิตพฺโพ น ปยิรุปาสิตพฺโพ อญฺญตฺร อนุทฺทยา อญฺญตฺร อนุกมฺปา.}

\prose{กตโม จ ภิกฺขเว ปุคฺคโล เสวิตพฺโพ ภชิตพฺโพ ปยิรุปาสิตพฺโพ? อิธ ภิกฺขเว เอกจฺโจ ปุคฺคโล สทิโส โหติ สีเลน สมาธินา ปญฺญาย, เอวรูโป ภิกฺขเว ปุคฺคโล เสวิตพฺโพ ภชิตพฺโพ ปยิรุปาสิตพฺโพ. ตํ กิสฺส เหตุ? สีลสามญฺญคตานํ สตํ สีลกถา จ โน ภวิสฺสติ, สา จ โน ปวตฺตินี\csromansharedfootnote{efefdfe4e5e0c144}{ปวตฺตนี (สี, สฺยา, กํ, อิ) อภิ 3. 140 ปิฏฺเฐปิ ปสฺสิตพฺพํ.} ภวิสฺสติ, สา จ โน ผาสุ ภวิสฺสติ. สมาธิ\-สามญฺญคตานํ สตํ สมาธิกถา จ โน ภวิสฺสติ, สา จ โน ปวตฺตินี ภวิสฺสติ, สา จ โน ผาสุ ภวิสฺสติ. ปญฺญา\-สามญฺญคตานํ สตํ ปญฺญากถา จ โน ภวิสฺสติ, สา จ โน ปวตฺตินี ภวิสฺสติ, สา จ โน ผาสุ ภวิสฺสตีติ. ตสฺมา เอวรูโป ปุคฺคโล เสวิตพฺโพ ภชิตพฺโพ ปยิรุปาสิตพฺโพ.}

\prose{กตโม จ ภิกฺขเว ปุคฺคโล สกฺกตฺวา ครุํ กตฺวา เสวิตพฺโพ ภชิตพฺโพ ปยิรุปาสิตพฺโพ? อิธ ภิกฺขเว เอกจฺโจ ปุคฺคโล อธิโก โหติ สีเลน สมาธินา ปญฺญาย. เอวรูโป ภิกฺขเว ปุคฺคโล สกฺกตฺวา ครุํ กตฺวา เสวิตพฺโพ ภชิตพฺโพ ปยิรุปาสิตพฺโพ. ตํ กิสฺส เหตุ? อิติ อปริปูรํ วา สีลกฺขนฺธํ ปริปูเรสฺสามิ, ปริปูรํ วา สีลกฺขนฺธํ ตตฺถ ตตฺถ ปญฺญาย อนุคฺคเหสฺสามิ. อปริปูรํ วา สมาธิ\-กฺขนฺธํ ปริปูเรสฺสามิ, ปริปูรํ วา สมาธิ\-กฺขนฺธํ ตตฺถ ตตฺถ ปญฺญาย อนุคฺคเหสฺสามิ. อปริปูรํ วา ปญฺญากฺขนฺธํ ปริปูเรสฺสามิ, ปริปูรํ วา ปญฺญากฺขนฺธํ ตตฺถ ตตฺถ ปญฺญาย อนุคฺคเหสฺสามีติ. ตสฺมา เอวรูโป ปุคฺคโล สกฺกตฺวา ครุํ กตฺวา เสวิตพฺโพ ภชิตพฺโพ ปยิรุปาสิตพฺโพ. อิเม โข ภิกฺขเว ตโย ปุคฺคลา สนฺโต สํวิชฺชมานา โลกสฺมินฺติ.}

\setlength{\csromangathapagewidth}{0pt}
\setlength{\csromangathawakwidth}{0pt}
\csromangathameasuregroup{}{นิหียติ ปุริโส นิหีนเสวี, \\ น จ หาเยถ กทาจิ ตุลฺยเสวี. \\ เสฏฺฐมุปนมํ อุเทติ ขิปฺปํ, \\ ตสฺมา อตฺตโน อุตฺตริํ ภเชถาติ.ฉฏฺฐํ.}
\csromangathameasurewak{นิหียติ ปุริโส นิหีนเสวี, \\ น จ หาเยถ กทาจิ ตุลฺยเสวี. \\ เสฏฺฐมุปนมํ อุเทติ ขิปฺปํ, \\ ตสฺมา อตฺตโน อุตฺตริํ ภเชถาติ.ฉฏฺฐํ.}
\setlength{\csromangathaleftcol}{0pt}
\csromangathagroup{\csromangathastanza{\csromangathawak{นิหียติ ปุริโส นิหีนเสวี, \\ น จ หาเยถ กทาจิ ตุลฺยเสวี. \\ เสฏฺฐมุปนมํ อุเทติ ขิปฺปํ, \\ ตสฺมา อตฺตโน อุตฺตริํ ภเชถาติ.  ฉฏฺฐํ.}}}

\csromanpalirecto
\gambhira{ติกนิปาตปาฬิ 7. ชิคุจฺฉิตพฺพ\-สุตฺต}
\proseitem{27}{\csromanfolio{124}ตโยเม ภิกฺขเว ปุคฺคลา สนฺโต สํวิชฺชมานา โลกสฺมิํ. กตเม ตโย? อตฺถิ ภิกฺขเว ปุคฺคโล ชิคุจฺฉิตพฺโพ น เสวิตพฺโพ น ภชิตพฺโพ น ปยิรุปาสิตพฺโพ, อตฺถิ ภิกฺขเว ปุคฺคโล อชฺฌุเปกฺขิตพฺโพ น เสวิตพฺโพ น ภชิตพฺโพ น ปยิรุปาสิตพฺโพ, อตฺถิ ภิกฺขเว ปุคฺคโล เสวิตพฺโพ ภชิตพฺโพ ปยิรุปาสิตพฺโพ. กตโม จ ภิกฺขเว ปุคฺคโล ชิคุจฺฉิตพฺโพ น เสวิตพฺโพ น ภชิตพฺโพ น ปยิรุปาสิตพฺโพ? อิธ ภิกฺขเว เอกจฺโจ ปุคฺคโล ทุสฺสีโล โหติ ปาปธมฺโม อสุจิ สงฺกสฺสร\-สมาจาโร ปฏิจฺฉนฺน\-กมฺมนฺโต อสฺสมโณ สมณปฏิญฺโญ อพฺรหฺมจารี พฺรหฺมจาริ\-ปฏิญฺโญ อนฺโตปูติ อวสฺสุโต กสมฺพุชาโต, เอวรูโป ภิกฺขเว ปุคฺคโล ชิคุจฺฉิตพฺโพ น เสวิตพฺโพ น ภชิตพฺโพ น ปยิรุปาสิตพฺโพ. ตํ กิสฺส เหตุ? กิญฺจาปิ ภิกฺขเว เอวรูปสฺส ปุคฺคลสฺส น ทิฏฺฐานุคติํ อาปชฺชติ, อถ โข นํ ปาปโก กิตฺติสทฺโท อพฺภุคฺคจฺฉติ “ปาปมิตฺโต ปุริสปุคฺคโล ปาปสหาโย ปาปสมฺปวงฺโก”ติ. เสยฺยถาปิ ภิกฺขเว อหิ คูถคโต กิญฺจาปิ น ทํสติ\csromansharedfootnote{6f52f2ee1f4ae874}{qอํสติ (สี, สฺยา), ฑสฺสติ (อิ)}, อถ โข นํ มกฺเขติ. เอวเมวํ โข ภิกฺขเว กิญฺจาปิ เอวรูปสฺส ปุคฺคลสฺส น ทิฏฺฐานุคติํ อาปชฺชติ, อถ โข นํ ปาปโก กิตฺติสทฺโท อพฺภุคฺคจฺฉติ “ปาปมิตฺโต ปุริสปุคฺคโล ปาปสหาโย ปาปสมฺปวงฺโก”ติ. ตสฺมา เอวรูโป ปุคฺคโล ชิคุจฺฉิตพฺโพ น เสวิตพฺโพ น ภชิตพฺโพ น ปยิรูปาสิตพฺโพ.}

\prose{กตโม จ ภิกฺขเว ปุคฺคโล อชฺฌุเปกฺขิตพฺโพ น เสวิตพฺโพ น ภชิตพฺโพ น ปยิรุปาสิตพฺโพ? อิธ ภิกฺขเว เอกจฺโจ ปุคฺคโล โกธโน โหติ อุปายาสพหุโล, อปฺปมฺปิ วุตฺโต สมาโน อภิสชฺชติ กุปฺปติ พฺยาปชฺชติ ปติตฺถียติ, โกปญฺจ โทสญฺจ อปฺปจฺจยญฺจ ปาตุกโรติ. เสยฺยถาปิ ภิกฺขเว ทุฏฺฐารุโก กฏฺเฐน วา กฐลาย วา ฆฏฺฏิโต ภิยฺโยโส มตฺตาย อาสวํ เทติ. เอวเมวํ โข ภิกฺขเว ฯเปฯ. เสยฺยถาปิ ภิกฺขเว ตินฺทุกาลาตํ กฏฺเฐน วา กฐลาย วา ฆฏฺฏิตํ ภิยฺโยโส มตฺตาย จิจฺจิฏายติ จิฏิจิฏายติ. เอวเมวํ โข ภิกฺขเว ฯเปฯ. เสยฺยถาปิ ภิกฺขเว คูถกูโป กฏฺเฐน วา กฐลาย วา ฆฏฺฏิโต ภิยฺโยโส มตฺตาย ทุคฺคนฺโธ โหติ. เอวเมวํ โข ภิกฺขเว อิเธกจฺโจ \csromanfolio{125}ปุคฺคโล โกธโน โหติ อุปายาสพหุโล, อปฺปมฺปิ วุตฺโต สมาโน อภิสชฺชติ กุปฺปติ พฺยาปชฺชติ ปติตฺถียติ, โกปญฺจ โทสญฺจ อปฺปจฺจยญฺจ ปาตุกโรติ. เอวรูโป ภิกฺขเว ปุคฺคโล อชฺฌุเปกฺขิตพฺโพ น เสวิตพฺโพ น ภชิตพฺโพ น ปยิรุปาสิตพฺโพ. ตํ กิสฺส เหตุ? อกฺโกเสยฺยปิ มํ, ปริภาเสยฺยปิ มํ, อนตฺถมฺปิ มํ กเรยฺยาติ. ตสฺมา เอวรูโป ปุคฺคโล อชฺฌุเปกฺขิตพฺโพ น เสวิตพฺโพ น ภชิตพฺโพ น ปยิรุปาสิตพฺโพ.}

\prose{กตโม จ ภิกฺขเว ปุคฺคโล เสวิตพฺโพ ภชิตพฺโพ ปยิรุปาสิตพฺโพ? อิธ ภิกฺขเว เอกจฺโจ ปุคฺคโล สีลวา โหติ กลฺยาณธมฺโม, เอวรูโป ภิกฺขเว ปุคฺคโล เสวิตพฺโพ ภชิตพฺโพ ปยิรุปาสิตพฺโพ. ตํ กิสฺส เหตุ? กิญฺจาปิ ภิกฺขเว เอวรูปสฺส ปุคฺคลสฺส น ทิฏฺฐานุคติํ อาปชฺชติ, อถ โข นํ กลฺยาโณ กิตฺติสทฺโท อพฺภุคฺคจฺฉติ “กลฺยาณมิตฺโต ปุริสปุคฺคโล กลฺยาณสหาโย กลฺยาณ\-สมฺปวงฺโก”ติ. ตสฺมา เอวรูโป ปุคฺคโล เสวิตพฺโพ ภชิตพฺโพ ปยิรุปาสิตพฺโพ. อิเม โข ภิกฺขเว ตโย ปุคฺคลา สนฺโต สํวิชฺชมานา โลกสฺมินฺติ.}

\setlength{\csromangathapagewidth}{0pt}
\setlength{\csromangathawakwidth}{0pt}
\csromangathameasuregroup{}{นิหียติ ปุริโส นิหีนเสวี, \\ น จ หาเยถ กทาจิ ตุลฺยเสวี. \\ เสฏฺฐมุปนมํ อุเทติ ขิปฺปํ, \\ ตสฺมา อตฺตโน อุตฺตริํ ภเชถาติ.สตฺตมํ.}
\csromangathameasurewak{นิหียติ ปุริโส นิหีนเสวี, \\ น จ หาเยถ กทาจิ ตุลฺยเสวี. \\ เสฏฺฐมุปนมํ อุเทติ ขิปฺปํ, \\ ตสฺมา อตฺตโน อุตฺตริํ ภเชถาติ.สตฺตมํ.}
\setlength{\csromangathaleftcol}{0pt}
\csromangathagroup{\csromangathastanza{\csromangathawak{นิหียติ ปุริโส นิหีนเสวี, \\ น จ หาเยถ กทาจิ ตุลฺยเสวี. \\ เสฏฺฐมุปนมํ อุเทติ ขิปฺปํ, \\ ตสฺมา อตฺตโน อุตฺตริํ ภเชถาติ.  สตฺตมํ.}}}

\csromanheader{2}{8. คูถภาณี\-สุตฺต}
\proseitem{28}{ตโยเม ภิกฺขเว ปุคฺคลา สนฺโต สํวิชฺชมานา โลกสฺมิํ. กตเม ตโย? คูถภาณี ปุปฺผภาณี มธุภาณี. กตโม จ ภิกฺขเว ปุคฺคโล คูถภาณี? อิธ ภิกฺขเว เอกจฺโจ ปุคฺคโล สภคฺคโต วา ปริสคฺคโต วา\csromansharedfootnote{a45207ecd9f31d19}{สภาคโต วา ปริสาคโต วา (สฺยา, กํ)} ญาติมชฺฌคโต วา ปูคมชฺฌคโต วา ราชกุล\-มชฺฌคโต วา อภินีโต สกฺขิปุฏฺโฐ “เอหมฺโภ ปุริส, ยํ ชานาสิ, ตํ วเทหี”ติ. โส อชานํ วา อาห “ชานามี”ติ, ชานํ วา อาห “น ชานามี”ติ, อปสฺสํ วา อาห “ปสฺสามี”ติ, ปสฺสํ วา อาห \csromanfolio{126}“น ปสฺสามี”ติ\csromansharedfootnote{e5701b64d10e56b2}{ม 1. 355; อภิ 3. 13 ปิฏฺเฐสุปิ.}, อิติ อตฺตเหตุ วา ปรเหตุ วา อามิส\-กิญฺจิกฺข\-เหตุ วา สมฺปชานมุสา ภาสิตา โหติ. อยํ วุจฺจติ ภิกฺขเว ปุคฺคโล คูถภาณี.}

\prose{กตโม จ ภิกฺขเว ปุคฺคโล ปุปฺผภาณี? อิธ ภิกฺขเว เอกจฺโจ ปุคฺคโล สภคฺคโต วา ปริสคฺคโต วา ญาติมชฺฌคโต วา ปูคมชฺฌคโต วา ราชกุล\-มชฺฌคโต วา อภินีโต สกฺขิปุฏฺโฐ “เอหมฺโภ ปุริส, ยํ ปชานาสิ, ตํ วเทหี”ติ. โส อชานํ วา อาห “น ชานามี”ติ, ชานํ วา อาห “ชานามี”ติ, อปสฺสํ วา อาห “น ปสฺสมี”ติ, ปสฺสํ วา อาห “ปสฺสามี”ติ, อิติ อตฺตเหตุ วา ปรเหตุ วา อามิส\-กิญฺจิกฺข\-เหตุ วา น สมฺปชานมุสา ภาสิตา โหติ. อยํ วุจฺจติ ภิกฺขเว ปุคฺคโล ปุปฺผภาณี.}

\prose{กตโม จ ภิกฺขเว ปุคฺคโล มธุภาณี? อิธ ภิกฺขเว เอกจฺโจ ปุคฺคโล ผรุสํ วาจํ ปหาย ผรุสาย วาจาย ปฏิวิรโต โหติ, ยา สา วาจา เนลา กณฺณสุขา เปมนียา หทยงฺคมา โปรี พหุชนกนฺตา พหุชนมนาปา, ตถารูปิํ วาจํ ภาสิตา โหติ. อยํ วุจฺจติ ภิกฺขเว ปุคฺคโล มธุภาณี. อิเม โข ภิกฺขเว ตโย ปุคฺคลา สนฺโต สํวิชฺชมานา โลกสฺมินฺติ.  อฏฺฐมํ.}

\csromanheader{2}{9. \textbf{อนฺธสุตฺต}}
\proseitem{29}{ตโยเม ภิกฺขเว ปุคฺคลา สนฺโต สํวิชฺชมานา โลกสฺมิํ. กตเม ตโย? อนฺโธ เอกจกฺขุ ทฺวิจกฺขุ. กตโม จ ภิกฺขเว ปุคฺคโล อนฺโธ? อิธ ภิกฺขเว เอกจฺจสฺส ปุคฺคลสฺส ตถารูปํ จกฺขุ น โหติ, ยถารูเปน จกฺขุนา อนธิคตํ วา โภคํ อธิคจฺเฉยฺย, อธิคตํ วา โภคํ ผาติํ กเรยฺย\csromansharedfootnote{cc017033290a67cd}{ผาติกเรยฺย (สี)}. ตถารูปมฺปิสฺส จกฺขุ น โหติ, ยถารูเปน จกฺขุนา กุสลากุสเล ธมฺเม ชาเนยฺย, สาวชฺชานวชฺเช ธมฺเม ชาเนยฺย, หีนปฺปณีเต ธมฺเม ชาเนยฺย, กณฺห\-สุกฺก\-สปฺปฏิภาเค ธมฺเม ชาเนยฺย. อยํ วุจฺจติ ภิกฺขเว ปุคฺคโล อนฺโธ.}

\prose{กตโม จ ภิกฺขเว ปุคฺคโล เอกจกฺขุ? อิธ ภิกฺขเว เอกจฺจสฺส ปุคฺคลสฺส ตถารูปํ จกฺขุ โหติ, ยถารูเปน จกฺขุนา อนธิคตํ วา \csromanfolio{127}โภคํ อธิคจฺเจยฺย, อธิคตํ วา โภคํ ผาติํ กเรยฺย. ตถารูปํ ปนสฺส\csromansharedfootnote{3fefac9516293d64}{ตถารูปมฺปิสฺส (สฺยา, กํ, อิ, ก)} จกฺขุ น โหติ, ยถารูเปน จกฺขุนา กุสลากุสเล ธมฺเม ชาเนยฺย, สาวชฺชานวชฺเช ธมฺเม ชาเนยฺย, หีนปฺปณีเต ธมฺเม ชาเนยฺย, กณฺห\-สุกฺก\-สปฺปฏิภาเค ธมฺเม ชาเนยฺย. อยํ วุจฺจติ ภิกฺขเว ปุคฺคโล เอกจกฺขุ.}

\prose{กตโม จ ภิกฺขเว ปุคฺคโล ทฺวิจกฺขุ? อิธ ภิกฺขเว เอกจฺจสฺส ปุคฺคลสฺส ตถารูปํ จกฺขุ โหติ, ยถารูเปน จกฺขุนา อนธิคตํ วา โภคํ อธิคจฺเฉยฺย, อธิคตํ วา โภคํ ผาติํ กเรยฺย. ตถารูปมฺปิสฺส จกฺขุ โหติ, ยถารูเปน จกฺขุนา กุสลากุสเล ธมฺเม ชาเนยฺย, สาวชฺชานวชฺเช ธมฺเม ชาเนยฺย, หีนปฺปณีเต ธมฺเม ชาเนยฺย, กณฺห\-สุกฺก\-สปฺปฏิภาเค ธมฺเม ชาเนยฺย. อยํ วุจฺจติ ภิกฺขเว ปุคฺคโล ทฺวิจกฺขุ. อิเม โข ภิกฺขเว ตโย ปุคฺคลา สนฺโต สํวิชฺชมานา โลกสฺมินฺติ.}

\setlength{\csromangathapagewidth}{0pt}
\setlength{\csromangathawakwidth}{0pt}
\csromangathameasuregroup{น เจว โภคา ตถารูปา, \\ อุภยตฺถ กลิคฺคาโห, \\ อถาปรายํ อกฺขาโต, \\ ธมฺมาธมฺเมน สโฐ โส\textsuperscript{1}, \\ เถยฺเยน กูฏกมฺเมน, \\ กุสโล โหติ สงฺฆาตุํ\textsuperscript{1}, \\ อิโต โส นิรยํ คนฺตฺวา, \\ ทฺวิจกฺขุ ปน อกฺขาโต, \\ ธมฺมลทฺเธหิ โภเคหิ, \\ ททาติ เสฏฺฐสงฺกปฺโป, \\ อุเปติ ภทฺทกํ ฐานํ, \\ อนฺธญฺจ เอกจกฺขุญฺจ, \\ ทฺวิจกฺขุํ ปน เสเวถ,}{\csromangathabat{น เจว โภคา ตถารูปา,}{น จ ปุญฺญานิ กุพฺพติ.} \\ \csromangathabat{อุภยตฺถ กลิคฺคาโห,}{อนฺธสฺส หตจกฺขุโน.} \\ \csromangathabat{อถาปรายํ อกฺขาโต,}{เอกจกฺขุ จ ปุคฺคโล.} \\ \csromangathabat{ธมฺมาธมฺเมน สโฐ โส\textsuperscript{1},}{โภคานิ ปริเยสติ.} \\ \csromangathabat{เถยฺเยน กูฏกมฺเมน,}{มุสาวาเทน จูภยํ.} \\ \csromangathabat{กุสโล โหติ สงฺฆาตุํ\textsuperscript{1},}{กามโภคี จ มานโว.} \\ \csromangathabat{อิโต โส นิรยํ คนฺตฺวา,}{เอกจกฺขุ วิหญฺญติ.} \\ \csromangathabat{ทฺวิจกฺขุ ปน อกฺขาโต,}{เสฏฺโฐ ปุริสปุคฺคโล.} \\ \csromangathabat{ธมฺมลทฺเธหิ โภเคหิ,}{อุฏฺฐานาธิคตํ ธนํ.} \\ \csromangathabat{ททาติ เสฏฺฐสงฺกปฺโป,}{อพฺยคฺคมานโส นโร.} \\ \csromangathabat{อุเปติ ภทฺทกํ ฐานํ,}{ยตฺถ คนฺตฺวา น โสจติ.} \\ \csromangathabat{อนฺธญฺจ เอกจกฺขุญฺจ,}{อารกา ปริวชฺชเย.} \\ \csromangathabat{ทฺวิจกฺขุํ ปน เสเวถ,}{เสฏฺฐํ ปุริสปุคฺคลนฺติ.นวมํ.}}
\csromangathasetleft{น เจว โภคา ตถารูปา, \\ อุภยตฺถ กลิคฺคาโห, \\ อถาปรายํ อกฺขาโต, \\ ธมฺมาธมฺเมน สโฐ โส\textsuperscript{1}, \\ เถยฺเยน กูฏกมฺเมน, \\ กุสโล โหติ สงฺฆาตุํ\textsuperscript{1}, \\ อิโต โส นิรยํ คนฺตฺวา, \\ ทฺวิจกฺขุ ปน อกฺขาโต, \\ ธมฺมลทฺเธหิ โภเคหิ, \\ ททาติ เสฏฺฐสงฺกปฺโป, \\ อุเปติ ภทฺทกํ ฐานํ, \\ อนฺธญฺจ เอกจกฺขุญฺจ, \\ ทฺวิจกฺขุํ ปน เสเวถ,}
\csromangathagroup{\csromangathastanza{\csromangathabat{น เจว โภคา ตถารูปา,}{น จ ปุญฺญานิ กุพฺพติ.} \\ \csromangathabat{อุภยตฺถ กลิคฺคาโห,}{อนฺธสฺส หตจกฺขุโน.}}\csromangathastanza{\csromangathabat{อถาปรายํ อกฺขาโต,}{เอกจกฺขุ จ ปุคฺคโล.} \\ \csromangathabat{ธมฺมาธมฺเมน สโฐ โส\csromansharedfootnote{cf7c3f079c5224b4}{สํสฏฺโฐ (สี, สฺยา, กํ, อิ), สโฐติ (ก) 3. สํหาตุํ (สฺยา)},}{โภคานิ ปริเยสติ.}}\csromangathastanza{\csromangathabat{เถยฺเยน กูฏกมฺเมน,}{มุสาวาเทน จูภยํ.} \\ \csromangathabat{กุสโล โหติ สงฺฆาตุํ\csromansharedfootnote{c73fc0fbbd421c55}{สํหาตุํ (สฺยา)},}{กามโภคี จ มานโว.}}\csromangathastanza{\csromangathabat{อิโต โส นิรยํ คนฺตฺวา,}{เอกจกฺขุ วิหญฺญติ.} \\ \csromangathabat{ทฺวิจกฺขุ ปน อกฺขาโต,}{เสฏฺโฐ ปุริสปุคฺคโล.}}\csromangathastanza{\csromangathabat{ธมฺมลทฺเธหิ โภเคหิ,}{อุฏฺฐานาธิคตํ ธนํ.} \\ \csromangathabat{ททาติ เสฏฺฐสงฺกปฺโป,}{อพฺยคฺคมานโส นโร.}}\csromangathastanza{\csromangathabat{อุเปติ ภทฺทกํ ฐานํ,}{ยตฺถ คนฺตฺวา น โสจติ.} \\ \csromangathabat{อนฺธญฺจ เอกจกฺขุญฺจ,}{อารกา ปริวชฺชเย.} \\ \csromangathabat{ทฺวิจกฺขุํ ปน เสเวถ,}{เสฏฺฐํ ปุริสปุคฺคลนฺติ.  นวมํ.}}}

\csromanpalirecto
\gambhira{ติกนิปาตปาฬิ \textbf{10. อวกุชฺชสุตฺต}}
\proseitem{30}{\csromanfolio{128}ตโยเม ภิกฺขเว\csromansharedfootnote{8b523ce99daffcec}{อภิ 3. 135 ปิฏฺเฐปิ.} ปุคฺคลา สนฺโต สํวิชฺชมานา โลกสฺมิํ. กตเม ตโย? อวกุชฺชปญฺโญ ปุคฺคโล, อุจฺฉงฺคปญฺโญ ปุคฺคโล, ปุถุปญฺโญ ปุคฺคโล. กตโม จ ภิกฺขเว อวกุชฺชปญฺโญ ปุคฺคโล? อิธ ภิกฺขเว เอกจฺโจ ปุคฺคโล อารามํ คนฺตา โหติ อภิกฺขณํ ภิกฺขูนํ สนฺติเก ธมฺม\-สฺสวนาย, ตสฺส ภิกฺขู ธมฺมํ เทเสนฺติ อาทิกลฺยาณํ มชฺเฌกลฺยาณํ ปริโยสาน\-กลฺยาณํ สาตฺถํ สพฺยญฺชนํ เกวล\-ปริปุณฺณํ ปริสุทฺธํ พฺรหฺมจริยํ ปกาเสนฺติ. โส ตสฺมิํ อาสเน นิสินฺโน ตสฺสา กถาย เนว อาทิํ มนสิ กโรติ, น มชฺฌํ มนสิ กโรติ, น ปริโยสานํ มนสิ กโรติ, วุฏฺฐิโตปิ ตมฺหา อาสนา ตสฺสา กถาย เนว อาทิํ มนสิ กโรติ, น มชฺฌํ มนสิ กโรติ, น ปริโยสานํ มนสิ กโรติ. เสยฺยถาปิ ภิกฺขเว กุมฺโภ นิกุชฺโช\csromansharedfootnote{42ff413519112778}{นิกฺกุชฺโช (สี, อิ)} ตตฺร อุทกํ อาสิตฺตํ วิวฏฺฏติ โน สณฺฐาติ. เอวเมวํ โข ภิกฺขเว อิเธกจฺโจ ปุคฺคโล อารามํ คนฺตา โหติ อภิกฺขณํ ภิกฺขูนํ สนฺติเก ธมฺม\-สฺสวนาย, ตสฺส ภิกฺขู ธมฺมํ เทเสนฺติ อาทิกลฺยาณํ มชฺเฌกลฺยานํ ปริโยสาน\-กลฺยาณํ สาตฺถํ สพฺยญฺชนํ เกวล\-ปริปุณฺณํ ปริสุทฺธํ พฺรหฺมจริยํ ปกาเสนฺติ. โส ตสฺมิํ อาสเน นิสินฺโน ตสฺสา กถาย เนว อาทิํ มนสิ กโรติ, น มชฺฌํ มนสิ กโรติ, น ปริโยสานํ มนสิ กโรติ, วุฏฺฐิโตปิ ตมฺหา อาสนา ตสฺสา กถาย เนว อาทิํ มนสิ กโรติ, น มชฺฌํ มนสิ กโรติ, น ปริโยสานํ มนสิ กโรติ. อยํ วุจฺจติ ภิกฺขเว อวกุชฺชปญฺโญ ปุคฺคโล.}

\prose{กตโม จ ภิกฺขเว อุจฺฉงฺคปญฺโญ ปุคฺคโล? อิธ ภิกฺขเว เอกจฺโจ ปุคฺคโล อารามํ คนฺตา โหติ อภิกฺขณํ ภิกฺขูนํ สนฺติเก ธมฺม\-สฺสวนาย, ตสฺส ภิกฺขู ธมฺมํ เทเสนฺติ อาทิกลฺยาณํ มชฺเฌกลฺยาณํ ปริโยสาน\-กลฺยาณํ สาตฺถํ สพฺยญฺชนํ เกวล\-ปริปุณฺณํ ปริสุทฺธํ พฺรหฺมจริยํ ปกาเสนฺติ. โส ตสฺมิํ อาสเน นิสินฺโน ตสฺสา กถาย อาทิมฺปิ มนสิ กโรติ, มชฺฌมฺปิ มนสิ กโรติ, ปริโยสานมฺปิ มนสิ กโรติ, วุฏฺฐิโต จ โข ตมฺหา อาสนา ตสฺสา กถาย เนว อาทิํ มนสิ กโรติ, น มชฺฌํ มนสิ กโรติ, น ปริโยสานํ มนสิ กโรติ. เสยฺยถาปิ ภิกฺขเว ปุริสสฺส อุจฺฉงฺเค นานาขชฺชกานิ อากิณฺณานิ ติลา ตณฺฑุลา โมทกา พทรา. โส ตมฺหา อาสนา วุฏฺฐหนฺโต สติสมฺโมสา \csromanfolio{129}ปกิเรยฺย. เอวเมวํ โข ภิกฺขเว อิเธกจฺโจ ปุคฺคโล อารามํ คนฺตา โหติ อภิกฺขณํ ภิกฺขูนํ สนฺติเก ธมฺม\-สฺสวนาย, ตสฺส ภิกฺขู ธมฺมํ เทเสนฺติ อาทิกลฺยาณํ มชฺเฌกลฺยาณํ ปริโยสาน\-กลฺยาณํ สาตฺถํ สพฺยญฺชนํ เกวล\-ปริปุณฺณํ ปริสุทฺธํ พฺรหฺมจริยํ ปกาเสนฺติ. โส ตสฺมิํ อาสเน นิสินฺโน ตสฺสา กถาย อาทิมฺปิ มนสิ กโรติ, มชฺฌมฺปิ มนสิ กโรติ, ปริโยสานมฺปิ มนสิ กโรติ, วุฏฺฐิโต จ โข ตมฺหา อาสนา ตสฺสา กถาย เนว อาทิํ มนสิ กโรติ, น มชฺฌํ มนสิ กโรติ, น ปริโยสานํ มนสิ กโรติ. อยํ วุจฺจติ ภิกฺขเว อุจฺฉงฺคปญฺโญ ปุคฺคโล.}

\prose{กตโม จ ภิกฺขเว ปุถุปญฺโญ ปุคฺคโล? อิธ ภิกฺขเว เอกจฺโจ ปุคฺคโล อารามํ คนฺตา โหติ อภิกฺขณํ ภิกฺขูนํ สนฺติเก ธมฺม\-สฺสวนาย, ตสฺส ภิกฺขู ธมฺมํ เทเสนฺติ อาทิกลฺยาณํ มชฺเฌกลฺยาณํ ปริโยสาน\-กลฺยาณํ สาตฺถํ สพฺยญฺชนํ เกวล\-ปริปุณฺณํ ปริสุทฺธํ พฺรหฺมจริยํ ปกาเสนฺติ. โส ตสฺมิํ อาสเน นิสินฺโน ตสฺสา กถาย อาทิมฺปิ มนสิ กโรติ, มชฺฌมฺปิ มนสิ กโรติ, ปริโยสานมฺปิ มนสิ กโรติ, วุฏฺฐิโตปิ ตมฺหา อาสนา ตสฺสา กถาย อาทิมฺปิ มนสิ กโรติ, มชฺฌมฺปิ มนสิ กโรติ, ปริโยสานมฺปิ มนสิ กโรติ. เสยฺยถาปิ ภิกฺขเว กุมฺโภ อุกฺกุชฺโช ตตฺร อุทกํ อาสิตฺตํ สณฺฐาติ โน วิวฏฺฏติ. เอวเมวํ โข ภิกฺขเว อิเธกจฺโจ ปุคฺคโล อารามํ คนฺตา โหติ อภิกฺขณํ ภิกฺขูนํ สนฺติเก ธมฺม\-สฺสวนาย, ตสฺส ภิกฺขู ธมฺมํ เทเสนฺติ อาทิกลฺยาณํ มชฺเฌกลฺยาณํ ปริโยสาน\-กลฺยาณํ สาตฺถํ สพฺยญฺชนํ เกวล\-ปริปุณฺณํ ปริสุทฺธํ พฺรหฺมจริยํ ปกาเสนฺติ. โส ตสฺมิํ อาสเน นิสินฺโน ตสฺสา กถาย อาทิมฺปิ มนสิ กโรติ, มชฺฌมฺปิ มนสิ กาโรติ, ปริโยสานมฺปิ มนสิ กโรติ, วุฏฺฐิโตปิ ตมฺหา อาสนา ตสฺสา กถาย อาทิมฺปิ มนสิ กโรติ, มชฺฌมฺปิ มนสิ กโรติ, ปริโยสานมฺปิ มนสิ กโรติ. อยํ วุจฺจติ ภิกฺขเว ปุถุปญฺโญ ปุคฺคโล. อิเม โข ภิกฺขเว ตโย ปุคฺคลา สนฺโต สํวิชฺชมานา โลกสฺมินฺติ.}

\prose{อวกุชฺชปญฺโญ ปุริโส, ทุมฺเมโธ อวิจกฺขโณ.}

\prose{อภิกฺขณมฺปิ เจ โหติ, คนฺตา ภิกฺขูน สนฺติเก.}

\prose{อาทิํ กถาย มชฺฌญฺจ, ปริโยสานญฺจ ตาทิโส.}

\prose{อุคฺคเหตุํ น สกฺโกติ, ปญฺญา หิสฺส น วิชฺชติ.}

\prose{\csromanfolio{130}อุจฺฉงฺคปญฺโญ ปุริโส, เสยฺโย เอเตน วุจฺจติ.}

\prose{อภิกฺขณมฺปิ เจ โหติ, คนฺตา ภิกฺขูน สนฺติเก.}

\prose{อาทิํ กถาย มชฺฌญฺจ, ปริโยสานญฺจ ตาทิโส.}

\prose{นิสินฺโน อาสเน ตสฺมิํ, อุคฺคเหตฺวาน พฺยญฺชนํ.}

\prose{วุฏฺฐิโต นปฺปชานาติ, คหิตํ หิสฺส\csromansharedfootnote{4095b84d88573f86}{คหิตมฺปิสฺส (ก)} มุสฺสติ.}

\prose{ปุถุปญฺโญ จ ปุริโส, เสยฺโย เอเตหิ\csromansharedfootnote{54afba802e6208cb}{เอเตน (ก)} วุจฺจติ.}

\prose{อภิกฺขณมฺปิ เจ โหติ, คนฺตา ภิกฺขูน สนฺติเก.}

\prose{อาทิํ กถาย มชฺฌญฺจ, ปริโยสานญฺจ ตาทิโส.}

\prose{นิสินฺโน อาสเน ตสฺมิํ, อุคฺคเหตฺวาน พฺยญฺชนํ.}

\prose{ธาเรติ เสฏฺฐสงฺกปฺโป, อพฺยคฺคมานโส นโร.}

\prosewithclosermidruled{ธมฺมา\-นุธมฺม\-ปฺปฏิปนฺโน, ทุกฺขสฺส\-นฺตกโร สิยาติ.  ทสมํ.}{ปุคฺคลวคฺโค ตติโย.}
\csromancenter{\textbf{ตสฺสุทฺทานํ}}
\setlength{\csromangathapagewidth}{0pt}
\setlength{\csromangathawakwidth}{0pt}
\csromangathameasuregroup{สมิทฺธ\textsuperscript{1} คิลาน สงฺขารา, \\ เสวิ ชิคุจฺฉ คูถภาณี,}{\csromangathabat{สมิทฺธ\textsuperscript{1} คิลาน สงฺขารา,}{พหุการา วชิเรน จ.} \\ \csromangathabat{เสวิ ชิคุจฺฉ คูถภาณี,}{อนฺโธ จ อวกุชฺชตาติ.}}
\csromangathasetleft{สมิทฺธ\textsuperscript{1} คิลาน สงฺขารา, \\ เสวิ ชิคุจฺฉ คูถภาณี,}
\csromangathagroup{\csromangathastanza{\csromangathabat{สมิทฺธ\csromansharedfootnote{5165798e98208e8a}{กายสกฺขิ (สี), สวิฏฺฐ (สฺยา, กํ), เสฏฺฐ (ก)} คิลาน สงฺขารา,}{พหุการา วชิเรน จ.} \\ \csromangathabat{เสวิ ชิคุจฺฉ คูถภาณี,}{อนฺโธ จ อวกุชฺชตาติ.}}}

\chapterhead{4. \textbf{เทวทูตวคฺค}}
\csromanheader{2}{1. สพฺรหฺมก\-สุตฺต}
\proseitem{31}{สพฺรหฺมกานิ ภิกฺขเว ตานิ กุลานิ, เยสํ ปุตฺตานํ มาตาปิตโร อชฺฌาคาเร ปูชิตา โหนฺติ. สปุพฺพา\-จริย\-กานิ ภิกฺขเว ตานิ กุลานิ, เยสํ ปุตฺตานํ มาตาปิตโร อชฺฌาคาเร ปูชิตา โหนฺติ. สาหุเนยฺยานิ ภิกฺขเว ตานิ กุลานิ, เยสํ ปุตฺตานํ มาตาปิตโร อชฺฌาคาเร ปูชิตา โหนฺติ. “พฺรหฺมา”ติ ภิกฺขเว มาตาปิตูนํ เอตํ อธิวจนํ, “ปุพฺพาจริยา”ติ ภิกฺขเว มาตาปิตูนํ เอตํ อธิวจนํ, \csromanfolio{131}“อหุเนยฺยา”ติ ภิกฺขเว มาตาปิตูนํ เอตํ อธิวจนํ. ตํ กิสฺส เหตุ? พหุการา ภิกฺขเว มาตาปิตโร ปุตฺตานํ อาปาทกา โปสกา อิมสฺส โลกสฺส ทสฺเสตาโรติ.}

\prose{“พฺรหฺมา”ติ มาตาปิตโร,}

\prose{“ปุพฺพาจริยา”ติ วุจฺจเร.}

\setlength{\csromangathapagewidth}{0pt}
\setlength{\csromangathawakwidth}{0pt}
\csromangathameasuregroup{อาหุเนยฺยา จ ปุตฺตานํ, \\ ตสฺมา หิ เน นมสฺเสยฺย, \\ อนฺเนน อถ ปาเนน, \\ อุจฺฉาทเนน นฺหาปเนน\textsuperscript{1}, \\ ตาย นํ ปาริจริยาย, \\ อิเธว\textsuperscript{1} นํ ปสํสนฺติ,}{\csromangathabat{อาหุเนยฺยา จ ปุตฺตานํ,}{ปชาย อนุกมฺปกา.} \\ \csromangathabat{ตสฺมา หิ เน นมสฺเสยฺย,}{สกฺกเรยฺย จ ปณฺฑิโต.} \\ \csromangathabat{อนฺเนน อถ ปาเนน,}{วตฺเถน สยเนน จ.} \\ \csromangathabat{อุจฺฉาทเนน นฺหาปเนน\textsuperscript{1},}{ปาทานํ โธวเนน จ.} \\ \csromangathabat{ตาย นํ ปาริจริยาย,}{มาตาปิตูสุ ปณฺฑิตา.} \\ \csromangathabat{อิเธว\textsuperscript{1} นํ ปสํสนฺติ,}{เปจฺจ สคฺเค ปโมทตีติ\textsuperscript{1}.ปฐมํ.}}
\csromangathasetleft{อาหุเนยฺยา จ ปุตฺตานํ, \\ ตสฺมา หิ เน นมสฺเสยฺย, \\ อนฺเนน อถ ปาเนน, \\ อุจฺฉาทเนน นฺหาปเนน\textsuperscript{1}, \\ ตาย นํ ปาริจริยาย, \\ อิเธว\textsuperscript{1} นํ ปสํสนฺติ,}
\csromangathagroup{\csromangathastanza{\csromangathabat{อาหุเนยฺยา จ ปุตฺตานํ,}{ปชาย อนุกมฺปกา.}}\csromangathastanza{\csromangathabat{ตสฺมา หิ เน นมสฺเสยฺย,}{สกฺกเรยฺย จ ปณฺฑิโต.} \\ \csromangathabat{อนฺเนน อถ ปาเนน,}{วตฺเถน สยเนน จ.}}\csromangathastanza{\csromangathabat{อุจฺฉาทเนน นฺหาปเนน\csromansharedfootnote{f916334071f658af}{นหาปเนน (สี)},}{ปาทานํ โธวเนน จ.} \\ \csromangathabat{ตาย นํ ปาริจริยาย,}{มาตาปิตูสุ ปณฺฑิตา.} \\ \csromangathabat{อิเธว\csromansharedfootnote{f916334071f658af}{นหาปเนน (สี)} นํ ปสํสนฺติ,}{เปจฺจ สคฺเค ปโมทตีติ\csromansharedfootnote{f916334071f658af}{นหาปเนน (สี)}.  ปฐมํ.}}}

\csromanheader{2}{2. \textbf{อานนฺทสุตฺต}}
\proseitem{32}{อถ โข อายสฺมา อานนฺโท เยน ภควา เตนุปสงฺกมิ, อุปสงฺกมิตฺวา ภควนฺตํ อภิวาเทตฺวา เอกมนฺตํ นิสีทิ, เอกมนฺตํ นิสินฺโน โข อายสฺมา อานนฺโท ภควนฺตํ เอตทโวจ\csromandash{}}

\prose{“สิยา นุ โข ภนฺเต ภิกฺขุโน ตถารูโป สมาธิ\-ปฏิลาโภ, ยถา อิมสฺมิญฺจ สวิญฺญาณเก กาเย อหงฺการ\-มมงฺการ\-มานานุสยา นาสฺสุ, พหิทฺธา จ สพฺพนิมิตฺเตสุ อหงฺการ\-มมงฺการ\-มานานุสยา นาสฺสุ, ยญฺจ เจโตวิมุตฺติํ ปญฺญาวิมุตฺติํ อุปสมฺปชฺช วิหรโต อหงฺการ\-มมงฺการ\-มานานุสยา น โหนฺติ, ตญฺจ เสโตวิมุตฺติํ ปญฺญาวิมุตฺติํ อุปสมฺปชฺช วิหเรยฺยา”ติ. สิย อานนฺท ภิกฺขุโน ตถารูโป สมาธิ\-ปฏิลาโภ, ยถา อิมสฺมิญฺจ สวิญฺญาณเก กาเย อหงฺการ\-มมงฺการ\-มานานุสยา นาสฺสุ, พหิทฺธา จ สพฺพนิมิตฺเตสุ อหงฺการ\-มมงฺการ\-มานานุสยา นาสฺสุ, ยญฺจ เจโตวิมุตฺติํ ปญฺญาวิมุตฺติํ อุปสมฺปชฺช วิหรโต อหงฺการ\-มมงฺการ\-มานานุสยา น โหนฺติ, ตญฺจ เจโตวิมุตฺติํ ปญฺญาวิมุตฺติํ อุปสมฺปชฺช วิหเรยฺยาติ.}

\prose{ยถา กถํ ปน ภนฺเต สิยา ภิกฺขุโน ตถารูโป สมาธิ\-ปฏิลาโภ, ยถา อิมสฺมิญฺจ สวิญฺญาณเก กาเย อหงฺการ\-มมงฺการ\-มานานุสยา นาสฺสุ, พหิทฺธา จ สพฺพนิมิตฺเตสุ \csromanfolio{132}อหงฺการ\-มมงฺการ\-มานานุสยา นาสฺสุ, ยญฺจ เจโตวิมุตฺติํ ปญฺญาวิมุตฺติํ อุปสมฺปชฺช วิหรโต อหงฺการ\-มมงฺการ\-มานานุสยา น โหนฺติ, ตญฺจ เจโตวิมุตฺติํ ปญฺญาวิมุตฺติํ อุปสมฺปชฺช วิหเรยฺยาติ.}

\prose{อิธานนฺท ภิกฺขุโน เอวํ โหติ “เอตํ สนฺตํ เอตํ ปณีตํ, ยทิทํ สพฺพ\-สงฺขาร\-สมโถ สพฺพู\-ปธิ\-ปฏินิสฺสคฺโค ตณฺหากฺขโย วิราโค นิโรโธ นิพฺพานนฺ”ติ. เอวํ โข อานนฺท สิยา ภิกฺขุโน ตถารูโป สมาธิ\-ปฏิลาโภ, ยถา อิมสฺมิญฺจ สวิญฺญาณเก กาเย อหงฺการ\-มมงฺการ\-มานานุสยา นาสฺสุ, พหิทฺธา จ สพฺพนิมิตฺเตสุ อหงฺการ\-มมงฺการ\-มานานุสยา นาสฺสุ, ยญฺจ เจโตวิมุตฺติํ ปญฺญาวิมุตฺติํ อุปสมฺปชฺช วิหรโต อหงฺการ\-มมงฺการ\-มานานุสยา น โหนฺติ, ตญฺจ เจโตวิมุตฺติํ ปญฺญาวิมุตฺติํ อุปสมฺปชฺช วิหเรยฺยาติ. อิทญฺจ ปน เมตํ อานนฺท สนฺธาย ภาสิตํ ปารายเน ปุณฺณกปเญฺห\csromandash{}}

\setlength{\csromangathapagewidth}{0pt}
\setlength{\csromangathawakwidth}{0pt}
\csromangathameasuregroup{}{“สงฺขาย โลกสฺมิํ ปโรปรานิ\textsuperscript{1}, \\ ยสฺสิญฺชิตํ นตฺถิ กุหิญฺจิ โลเก. \\ สนฺโต วิธูโม อนีโฆ\textsuperscript{1} นิราโส, \\ อตาริ โส ชาติชรนฺติ พฺรูมี”ติ. \\ ทุติยํ.}
\csromangathameasurewak{“สงฺขาย โลกสฺมิํ ปโรปรานิ\textsuperscript{1}, \\ ยสฺสิญฺชิตํ นตฺถิ กุหิญฺจิ โลเก. \\ สนฺโต วิธูโม อนีโฆ\textsuperscript{1} นิราโส, \\ อตาริ โส ชาติชรนฺติ พฺรูมี”ติ. \\ ทุติยํ.}
\setlength{\csromangathaleftcol}{0pt}
\csromangathagroup{\csromangathastanza{\csromangathawak{“สงฺขาย โลกสฺมิํ ปโรปรานิ\csromansharedfootnote{3d275e7928bd9c34}{ปโรวรานิ (สี, อิ) ขุ 1. 436; ขุ 8. 8 ปิฏฺเฐสุ.}, \\ ยสฺสิญฺชิตํ นตฺถิ กุหิญฺจิ โลเก. \\ สนฺโต วิธูโม อนีโฆ\csromansharedfootnote{8ffe93f5a76e5a22}{อนิโฆ (สี, สฺยา, กํ, อิ), อนโฆ (?)} นิราโส, \\ อตาริ โส ชาติชรนฺติ พฺรูมี”ติ.}}\csromangathastanza{\csromangathawak{ทุติยํ.}}}

\csromanheader{2}{3. \textbf{สาริปุตฺตสุตฺต}}
\proseitem{33}{อถ โข อายสฺมา สาริปุตฺโต เยน ภควา เตนุปสงฺกมิ, อุปสงฺกมิตฺวา ภควนฺตํ อภิวาเทตฺวา เอกมนฺตํ นิสีทิ, เอกมนฺตํ นิสินฺนํ โข อายสฺมนฺตํ สาริปุตฺตํ ภควา เอตทโวจ “สํขิตฺเตนปิ โข อหํ สาริปุตฺต ธมฺมํ เทเสยฺยํ, วิตฺถาเรนปิ โข อหํ สาริปุตฺต ธมฺมํ เทเสยฺยํ, สํขิตฺต\-วิตฺถาเรนปิ โข อหํ สาริปุตฺต ธมฺมํ เทเสยฺยํ, อญฺญาตาโร จ ทุลฺลภา”ติ. เอตสฺส ภควา กาโล, เอตสฺส สุคต กาโล, ยํ ภควา สํขิตฺเตนปิ ธมฺมํ เทเสยฺย, วิตฺถาเรนปิ ธมฺมํ เทเสยฺย, สํขิตฺต\-วิตฺถาเรนปิ ธมฺมํ เทเสยฺย, ภวิสฺสนฺติ ธมฺมสฺส อญฺญาตาโรติ.}

\prose{ตสฺมาติห สาริปุตฺต เอวํ สิกฺขิตพฺพํ “อิมสฺมิญฺจ สวิญฺญาณเก กาเย อหงฺการ\-มมงฺการ\-มานานุสยา น ภวิสฺสนฺติ, พหิทฺธา จ สพฺพนิมิตฺเตสุ อหงฺการ\-มมงฺการ\-มานานุสยา น ภวิสฺสนฺติ, ยญฺจ เจโตนิมุตฺติํ \csromanfolio{133}ปญฺญาวิมุตฺติํ อุปสมฺปชฺช วิหรโต อหงฺการ\-มมงฺการ\-มานานุสยา น โหนฺติ, ตญฺจ เจโตวิมุตฺติํ ปญฺญาวิมุตฺติํ อุปสมฺปชฺช วิหริสฺสามา”ติ. เอวํ หิ โข สาริปุตฺต สิกฺขิตพฺพํ.}

\prose{ยโต จ โข สาริปุตฺต ภิกฺขุโน อิมสฺมิญฺจ สวิญฺญาณเก กาเย อหงฺการ\-มมงฺการ\-มานานุสยา น โหนฺติ, พหิทฺธา จ สพฺพนิมิตฺเตสุ อหงฺการ\-มมงฺการ\-มานานุสยา น โหนฺติ, ยญฺจ เจโตวิมุตฺติํ ปญฺญาวิมุตฺติํ อุปสมฺปชฺช วิหรโต อหงฺการ\-มมงฺการ\-มานานุสยา น โหนฺติ, ตญฺจ เจโตวิมุตฺติํ ปญฺญาวิมุตฺติํ อุปสมฺปชฺช วิหรติ. อยํ วุจฺจติ สาริปุตฺต ภิกฺขุ อจฺเฉจฺฉิ\csromansharedfootnote{8dd57c38ea38414e}{อจฺเฉชฺชิ (สฺยา, กํ, ก)} ตณฺหํ, วิวตฺตยิ\csromansharedfootnote{ad9184755c407031}{วาวตฺตยิ (สี, อิ)} สํโยชนํ, สมฺมา มานาภิสมยา อนฺตมกาสิ ทุกฺขสฺส. อิทญฺจ ปน เมตํ สาริปุตฺต สนฺทฺธาย ภาสิตํ ปารายเน\csromansharedfootnote{41f9a78387d41028}{ปารายเณ (สี)} อุทยปเญฺห\csromandash{}}

\setlength{\csromangathapagewidth}{0pt}
\setlength{\csromangathawakwidth}{0pt}
\csromangathameasuregroup{“ปหานํ กามสญฺญานํ, \\ ถินสฺส จ ปนูทนํ, \\ อุเปกฺขาสติสํสุทฺธํ,}{\csromangathabat{“ปหานํ กามสญฺญานํ,}{โทมนสฺสาน จูภยํ.} \\ \csromangathabat{ถินสฺส จ ปนูทนํ,}{กุกฺกุจฺจานํ นิวารณํ.} \\ \csromangathabat{อุเปกฺขาสติสํสุทฺธํ,}{ธมฺมตกฺกปุเรชวํ.} \\ ตติยํ.}
\csromangathameasurewak{ตติยํ.}
\csromangathasetleft{“ปหานํ กามสญฺญานํ, \\ ถินสฺส จ ปนูทนํ, \\ อุเปกฺขาสติสํสุทฺธํ,}
\csromangathagroup{\csromangathastanza{\csromangathabat{“ปหานํ กามสญฺญานํ,}{โทมนสฺสาน จูภยํ.} \\ \csromangathabat{ถินสฺส จ ปนูทนํ,}{กุกฺกุจฺจานํ นิวารณํ.} \\ \csromangathabat{อุเปกฺขา\-สติ\-สํสุทฺธํ,}{ธมฺมตกฺก\-ปุเรชวํ.}}\csromangathastanza{\csromangathawak{ตติยํ.}}}

\prose{อญฺญาวิโมกฺขํ ปพฺรูมิ, อวิชฺชาย ปเภทนนฺ”ติ\csromansharedfootnote{7cd32c96c2a3fe58}{ขุ 1. 446; ขุ 8. 17 ปิฏฺเฐสุปิ.}.}

\csromanheader{2}{4. \textbf{นิทานสุตฺต}}
\proseitem{34}{ตีณิมานิ ภิกฺขเว นิทานานิ กมฺมานํ สมุทยาย. กตมานิ ตีณิ? โลโภ นิทานํ กมฺมานํ สมุทยาย, โทโส นิทานํ กมฺมานํ สมุทยาย, โมโห นิทานํ กมฺมานํ สมุทยาย.}

\prose{ยํ ภิกฺขเว โลภปกตํ กมฺมํ โลภชํ โลภนิทานํ โลภสมุทยํ. ยตฺถสฺส อตฺตภาโว นิพฺพตฺตติ, ตตฺถ ตํ กมฺมํ วิปจฺจติ. ยตฺถ ตํ กมฺมํ วิปจฺจติ, ตตฺถ ตสฺส กมฺมสฺส วิปากํ ปฏิสํเวเทติ ทิฏฺเฐ วา ธมฺเม, อุปปชฺช วา\csromansharedfootnote{5319a0b406ab6ec0}{อุปปชฺเช วา (สี, สฺยา, กํ) อุปปชฺชิตฺวาติ ม 3. 257 ปิฏฺเฐ ปาฬิยา สํวณฺณนา.}, อปเร วา\csromansharedfootnote{e842d3181eb4c7e8}{อปราปเร วา (ก)} ปริยาเย.}

\prose{\csromanfolio{134}ยํ ภิกฺขเว โทสปกตํ กมฺมํ โทสชํ โทสนิทานํ โทสสมุทยํ. ยตฺถสฺส อตฺตภาโว นิพฺพตฺตติ, ตตฺถ ตํ กมฺมํ วิปจฺจติ. ยตฺถ ตํ กมฺมํ วิปจฺจติ, ตตฺถ ตสฺส กมฺมสฺส วิปากํ ปฏิสํเวเทติ ทิฏฺเฐ วา ธมฺเม, อุปปชฺช วา, อปเร วา ปริยาเย.}

\prose{ยํ ภิกฺขเว โมหปกตํ กมฺมํ โมหชํ โมหนิทานํ โมหสมุทยํ. ยตฺถสฺส อตฺตภาโว นิพฺพตฺตติ, ตตฺถ ตํ กมฺมํ วิปจฺจติ. ยตฺถ ตํ กมฺมํ วิปจฺจติ, ตตฺถ ตสฺส กมฺมสฺส วิปากํ ปฏิสํเวเทติ ทิฏฺเฐ วา ธมฺเม, อุปปชฺช วา, อปเร วา ปริยาเย.}

\prose{เสยฺยถาปิ ภิกฺขเว พีชานิ อขณฺฑานิ อปูตีนิ อวาตาตปหตานิ สาราทานิ สุขสยิตานิ สุเขตฺเต สุปริกมฺม\-กตาย ภูมิยา นิกฺขิตฺตานิ, เทโว จ สมฺมาธารํ อนุปฺปเวจฺเฉยฺย, เอวสฺสุ ตานิ ภิกฺขเว พีชานิ วุทฺธิํ วิรุฬฺหิํ เวปุลฺลํ อาปชฺเชยฺยุํ. เอวเมวํ โข ภิกฺขเว ยํ โลภปกตํ กมฺมํ โลภชํ โลภนิทานํ โลภสมุทยํ. ยตฺถสฺส อตฺตภาโว นิพฺพตฺตติ, ตตฺถ ตํ กมฺมํ วิปจฺจติ. ยตฺถ ตํ กมฺมํ วิปจฺจติ, ตตฺถ ตสฺส กมฺมสฺส วิปากํ ปฏิสํเวเทติ ทิฏฺเฐ วา ธมฺเม, อุปปชฺช วา, อปเร วา ปริยาเย.}

\prose{ยํ โทสปกตํ กมฺมํ ฯเปฯ. ยํ โมหปกตํ กมฺมํ โมหชํ โมหนิทานํ โมหสมุทยํ. ยตฺถสฺส อตฺตภาโว นิพฺพตฺตติ, ตตฺถ ตํ กมฺมํ วิปจฺจติ. ยตฺถ ตํ กมฺมํ วิปจฺจติ, ตตฺถ ตสฺส กมฺมสฺส วิปากํ ปฏิสํเวเทติ ทิฏฺเฐ วา ธมฺเม, อุปปชฺช วา, อปเร วา ปริยาเย. อิมานิ โข ภิกฺขเว ตีณี นิทานานิ กมฺมานํ สมุทยาย.}

\prose{ตีณีมานิ ภิกฺขเว นิทานานิ กมฺมานํ สมุทยาย. กตมานิ ตีณิ? อโลโภ นิทานํ กมฺมานํ สมุทยาย, อโทโส นิทานํ กมฺมานํ สมุทยาย, อโมโห นิทานํ กมฺมานํ สมุทยาย.}

\prose{ยํ ภิกฺขเว อโลภปกตํ กมฺมํ อโลภชํ อโลภนิทานํ อโลภ\-สมุทยํ, โลเภ วิคเต เอวํ ตํ กมฺมํ ปหีนํ โหติ อุจฺฉินฺนมูลํ ตาลา\-วตฺถุกตํ อนภาวํกตํ อายติํ อนุปฺปาท\-ธมฺมํ.}

\prose{ยํ ภิกฺขเว อโทสปกตํ กมฺมํ อโทสชํ อโทสนิทานํ อโทสสมุทยํ, โทเส วิคเต เอวํ ตํ กมฺมํ ปหีนํ โหติ อุจฺฉินฺนมูลํ ตาลา\-วตฺถุกตํ อนภาวํกตํ อายติํ อนุปฺปาท\-ธมฺมํ.}

\chapterheadpageread{4. เทวทูตวคฺค}
\prose{\csromanfolio{135}ยํ ภิกฺขเว อโมหปกตํ กมฺมํ อโมหชํ อโมหนิทานํ อโมหสมุทยํ, โมเห วิคเต เอวํ ตํ กมฺมํ ปหีนํ โหติ อุจฺฉินฺนมูลํ ตาลา\-วตฺถุกตํ อนภาวํกตํ อายติํ อนุปฺปาท\-ธมฺมํ.}

\prose{เสยฺยถาปิ ภิกฺขเว พีชานิ อขณฺฑานิ อปูตีนิ อวาตาตปหตานิ สาราทานิ สุขสยิตานิ, ตานิ ปุริโส อคฺคินา ฑเหยฺย, อคฺคินา ฑหิตฺวา มสิํ กเรยฺย, มสิํ กริตฺวา มหาวาเต วา โอผุเณยฺย\csromansharedfootnote{d037f3107d5a6dae}{โอปุเนยฺย (สี, อิ)}, นทิยา วา สีฆโสตาย ปวาเหยฺย. เอวสฺสุ ตานิ ภิกฺขเว พีชานิ อุจฺฉินฺนมูลานิ ตาลา\-วตฺถุ\-กตานิ อนภาวํ\-กตานิ\csromansharedfootnote{aca9ac8bad58300a}{อนภาวกตานิ (สี, อิ)} อายติํ อนุปฺปาท\-ธมฺมานิ. เอวเมวํ โข ภิกฺขเว ยํ อโลภปกตํ กมฺมํ อโลภชํ อโลภนิทานํ อโลภ\-สมุทยํ, โลเภ วิคเต เอวํ ตํ กมฺมํ ปหีนํ โหติ อุจฺฉินฺนมูลํ ตาลา\-วตฺถุกตํ อนภาวํกตํ อายติํ อนุปฺปาท\-ธมฺมํ.}

\prose{ยํ อโทสปกตํ กมฺมํ ฯเปฯ. ยํ อโมหปกตํ กมฺมํ อโมหชํ อโมหนิทานํ อโมหสมุทยํ, โมเห วิคเต เอวํ ตํ กมฺมํ ปหีนํ โหติ ฯเปฯ อายติํ อนุปฺปาท\-ธมฺมํ. อิมานิ โข ภิกฺขเว ตีณิ นิทานานิ กมฺมานํ สมุทยายาติ.}

\setlength{\csromangathapagewidth}{0pt}
\setlength{\csromangathawakwidth}{0pt}
\csromangathameasuregroup{โลภชํ โทสชํ เจว\textsuperscript{1}, \\ ยํ เตน ปกตํ กมฺมํ, \\ อิเธว ตํ เวทนิยํ, \\ ตสฺมา โลภญฺจ โทสญฺจ, \\ วิชฺชํ อุปฺปาทยํ ภิกฺขุ,}{\csromangathabat{โลภชํ โทสชํ เจว\textsuperscript{1},}{โมหชญฺจาปวิทฺทสุ.} \\ \csromangathabat{ยํ เตน ปกตํ กมฺมํ,}{อปฺปํ วา ยทิ วา พหุํ.} \\ \csromangathabat{อิเธว ตํ เวทนิยํ,}{วตฺถุ อญฺญํ น วิชฺชติ.} \\ \csromangathabat{ตสฺมา โลภญฺจ โทสญฺจ,}{โมหชญฺจาปิ วิทฺทสุ.} \\ \csromangathabat{วิชฺชํ อุปฺปาทยํ ภิกฺขุ,}{สพฺพา ทุคฺคติโย ชเหติ.จตุตฺถํ.}}
\csromangathasetleft{โลภชํ โทสชํ เจว\textsuperscript{1}, \\ ยํ เตน ปกตํ กมฺมํ, \\ อิเธว ตํ เวทนิยํ, \\ ตสฺมา โลภญฺจ โทสญฺจ, \\ วิชฺชํ อุปฺปาทยํ ภิกฺขุ,}
\csromangathagroup{\csromangathastanza{\csromangathabat{โลภชํ โทสชํ เจว\csromansharedfootnote{cb35578c9df38803}{โทสชํ กมฺมํ (ก)},}{โมหชญฺจาปวิทฺทสุ.} \\ \csromangathabat{ยํ เตน ปกตํ กมฺมํ,}{อปฺปํ วา ยทิ วา พหุํ.}}\csromangathastanza{\csromangathabat{อิเธว ตํ เวทนิยํ,}{วตฺถุ อญฺญํ น วิชฺชติ.} \\ \csromangathabat{ตสฺมา โลภญฺจ โทสญฺจ,}{โมหชญฺจาปิ วิทฺทสุ.} \\ \csromangathabat{วิชฺชํ อุปฺปาทยํ ภิกฺขุ,}{สพฺพา ทุคฺคติโย ชเหติ.  จตุตฺถํ.}}}

\csromanheader{2}{5. \textbf{หตฺถกสุตฺต}}
\proseitem{35}{เอวํ เม สุตํ\csromandash{}เอกํ สมยํ ภควา อาฬวิยํ วิหรติ โคมคฺเค สิํสปาวเน ปณฺณสนฺถเร. อถ โข หตฺถโก อาฬวโก ชงฺฆาวิหารํ อนุจงฺกมมาโน อนุวิจรมาโน อทฺทส ภควนฺตํ โคมคฺเค สิํสปาวเน ปณฺณสนฺถเร นิสินฺนํ, ทิสฺวา เยน ภควา เตนุปสงฺกมิ, อุปสงฺกมิตฺวา ภควนฺตํ อภิวาเทตฺวา เอกมนฺตํ นิสีทิ, เอกมนฺตํ นิสินฺโน โข หตฺถโก อาฬวโก ภควนฺตํ เอตทโวจ “กจฺจิ ภนฺเต ภควา สุขม\-สยิตฺถา”ติ. \csromanfolio{136}เอวํ กุมาร สุขม\-สยิตฺถํ, เย จ ปน โลเก สุขํ เสนฺติ, อหํ เตสํ อญฺญตโรติ.}

\prose{สีตา ภนฺเต เหมนฺติกา รตฺติ, อนฺตรฏฺฐโก หิมปาตสมโย, ขรา โคกณฺฏกหตา ภูมิ, ตนุโก ปณฺณสนฺถโร, วิรฬานิ รุกฺขสฺส ปตฺตานิ, สีตานิ, กาสายานิ วตฺถานิ, สีโต จ เวรมฺโภ วาโต วายติ, อถ จ ปน ภควา เอวมาห “เอวํ กุมาร สุขม\-สยิตฺถํ, เย จ ปน โลเก สุขํ เสนฺติ, อหํ เตสํ อญฺญตโร”ติ.}

\prose{เตน หิ กุมาร ตญฺเญเวตฺถ ปฏิปุจฺฉิสฺสามิ, ยถา เต ขเมยฺย, ตถา นํ พฺยากเรยฺยาสิ. ตํ กิํ มญฺญสิ กุมาร, อิธสฺส คหปติสฺส วา คหปติ\-ปุตฺตสฺส วา กูฏาคารํ อุลฺลิตฺตา\-วลิตฺตํ นิวาตํ ผุสิตคฺคฬํ ปิหิต\-วาตปานํ. ตตฺรสฺส ปลฺลงฺโก โคนกตฺถโต ปฏิกตฺถโต ปฏลิกตฺถโต กทลิ\-มิค\-ปวร\-ปจฺจตฺถรโณ\csromansharedfootnote{036357b51e476b1b}{กาทลิมิคปวรปจฺจตฺถรโณ (สี)} สอุตฺตรจฺฉโท อุภโต\-โลหิตกู\-ปธาโน, เตลปฺปทีโป เจตฺถ ฌาเยยฺย\csromansharedfootnote{f5160afdaab71fbf}{ชาเลยฺย (ก)}, จตสฺโส จ\csromansharedfootnote{91b315c4d76e9240}{ตสฺเสว (ก)} ปชาปติโย มนาปามนาเปน ปจฺจุปฏฺฐิตา อสฺสุ. ตํ กิํ มญฺญสิ กุมาร, สุขํ วา โส สเยยฺย โน วา, กถํ วา เต เอตฺถ โหตีติ. สุขํ โส ภนฺเต สเยยฺย, เย จ ปน โลเก สุขํ เสนฺติ, โส เตสํ อญฺญตโรติ.}

\prose{ตํ กิํ มญฺญสิ กุมาร, อปิ นุ ตสฺส คหปติสฺส วา คหปติ\-ปุตฺตสฺส วา อุปฺปชฺเชยฺยุํ ราคชา ปริฬาหา กายิกา วา เจตสิกา วา, เยหิ โส ราคเชหิ ปริฬาเหหิ ปริฑยฺหมาโน ทุกฺขํ สเยยฺยาติ. เอวํ ภนฺเตติ.}

\prose{เยหิ โข โส กุมาร คหปติ วา คหปติปุตฺโต วา ราคเชหิ ปริฬาเหหิ ปริฑยฺหมาโน ทุกฺขํ สเยยฺย. โส ราโค ตถาคตสฺส ปหีโน อุจฺฉินฺนมูโล ตาลาวตฺถุกโต อนภาวํกโต อายติํ อนุปฺปาทธมฺโม, ตสฺมาหํ สุขม\-สยิตฺถํ.}

\prose{ตํ กิํ มญฺญสิ กุมาร, อปิ นุ ตสฺส คหปติสฺส วา คหปติ\-ปุตฺตสฺส วา อุปฺปชฺเชยฺยุํ โทสชา ปริฬาหา ฯเปฯ โมหชา ปริฬาหา กายิกา วา เจตสิกา วา, เยหิ โส โมหเชหิ ปริฬาเหหิ ปริฑยฺหมาโน ทุกฺขํ สเยยฺยาติ. เอวํ ภนฺเตติ.}

\chapterheadpageread{4. เทวทูตวคฺค}
\prose{\csromanfolio{137}เยหิ โข โส กุมาร คหปติ วา คหปติปุตฺโต วา โมหเชหิ ปริฬาเหหิ ปริฑยฺหมาโน ทุกฺขํ สเยยฺย. โส โมโห ตถาคตสฺส ปหีโน อุจฺฉินฺนมูโล ตาลาวตฺถุกโต อนภาวํกโต อายติํ อนุปฺปาทธมฺโม. ตสฺมาหํ สุขม\-สยิตฺถนฺติ.}

\setlength{\csromangathapagewidth}{0pt}
\setlength{\csromangathawakwidth}{0pt}
\csromangathameasuregroup{สพฺพทา เว สุขํ เสติ, \\ โย น ลิมฺปติ\textsuperscript{1} กาเมสุ, \\ สพฺพา อาสตฺติโย เฉตฺวา, \\ อุปสนฺโต สุขํ เสติ,}{\csromangathabat{สพฺพทา เว สุขํ เสติ,}{พฺรหฺมโณ ปรินิพฺพุโต.} \\ \csromangathabat{โย น ลิมฺปติ\textsuperscript{1} กาเมสุ,}{สีติภูโต นิรูปธิ.} \\ \csromangathabat{สพฺพา อาสตฺติโย เฉตฺวา,}{วิเนยฺย หทเย ทรํ.} \\ \csromangathabat{อุปสนฺโต สุขํ เสติ,}{สนฺติํ ปปฺปุยฺย เจตโสติ.} \\ ปญฺจมํ.}
\csromangathameasurewak{ปญฺจมํ.}
\csromangathasetleft{สพฺพทา เว สุขํ เสติ, \\ โย น ลิมฺปติ\textsuperscript{1} กาเมสุ, \\ สพฺพา อาสตฺติโย เฉตฺวา, \\ อุปสนฺโต สุขํ เสติ,}
\csromangathagroup{\csromangathastanza{\csromangathabat{สพฺพทา เว สุขํ เสติ,}{พฺรหฺมโณ ปรินิพฺพุโต.} \\ \csromangathabat{โย น ลิมฺปติ\csromansharedfootnote{cedcdf48bf26d589}{ลิปฺปติ (สี, สฺยา, กํ, ก)} กาเมสุ,}{สีติภูโต นิรูปธิ.}}\csromangathastanza{\csromangathabat{สพฺพา อาสตฺติโย เฉตฺวา,}{วิเนยฺย หทเย ทรํ.} \\ \csromangathabat{อุปสนฺโต สุขํ เสติ,}{สนฺติํ ปปฺปุยฺย เจตโสติ.}}\csromangathastanza{\csromangathawak{ปญฺจมํ.}}}

\csromanheader{2}{6. \textbf{เทวทุตสุตฺต}}
\proseitem{36}{ตีณิมานิ ภิกฺขเว เทวทูตานิ. กตมานิ ตีณิ? อิธ ภิกฺขเว เอกจฺโจ กาเยน ทุจฺจริตํ จรติ, วาจาย ทุจฺจริตํ จรติ, มนสา ทุจฺจริตํ จรติ. โส กาเยน ทุจฺจริตํ จริตฺวา วาจาย ทุจฺจริตํ จริตฺวา มนสา ทุจฺจริตํ จริตฺวา กายสฺส เภทา ปรํ มรณา อปายํ ทุคฺคติํ วินิปาตํ นิรยํ อุปปชฺชติ, ตเมนํ ภิกฺขเว นิรยปาลา นานาพาหาสุ คเหตฺวา ยมสฺส รญฺโญ ทสฺเสนฺติ “อยํ เทว ปุริโส อมตฺเตยฺโย อเปตฺเตยฺโย อสามญฺโญ อพฺรหฺมญฺโญ, น กุเล เชฏฺฐาปจายี, อิมสฺส เทโว ทณฺฑํ ปเณตู”ติ.}

\prose{ตเมนํ ภิกฺขเว ยโม ราชา ปฐมํ เทวทูตํ สมนุยุญฺชติ สมนุคาหติ สมนุภาสติ “อมฺโภ ปุริส น ตฺวํ อทฺทส มนุสฺเสสุ ปฐมํ เทวทูตํ ปาตุภูตนฺ”ติ? โส เอวมาห “นาทฺทสํ ภนฺเต”ติ.}

\prose{ตเมนํ ภิกฺขเว ยโม ราชา เอวมาห “อมฺโภ ปุริส น ตฺวํ อทฺทส มนุสฺเสสุ อิตฺถิํ วา ปุริสํ วา อาสีติกํ วา นาวุติกํ วา วสฺสสติกํ วา\csromansharedfootnote{fb60e4a153c0ce54}{ปสฺส ม 3. 218 ปิฏฺเฐ.} ชาติยา ชิณฺณํ โคปานสิวงฺกํ โภคฺคํ ทณฺฑปรายณํ\csromansharedfootnote{d5a85ec0404224d8}{ทณฺฑปรายนํ (สฺยา, กํ, อิ)} ปเวธมานํ คจฺฉนฺตํ อาตุรํ คตโยพฺพนํ ขณฺฑทนฺตํ ปลิตเกสํ วิลูนํ ขลฺลิตสิรํ\csromansharedfootnote{d131a182f5d14059}{ขลิตํ สิโร (สี, อิ), ขลิตสิรํ (สฺยา, กํ) ม 3. 218 ปิฏฺเฐปิ.} วลิตํ ติลกาหตคตฺตนฺ”ติ? โส เอวมาห “อทฺทสํ ภนฺเต”ติ.}

\prose{\csromanfolio{138}ตเมนํ ภิกฺขเว ยโม ราชา เอวมาห “อมฺโภ ปุริส ตสฺส เต วิญฺญุสฺส สโต มหลฺลกสฺส น เอตทโหสิ ‘อหมฺปิ โขมฺหิ ชราธมฺโม ชรํ อนตีโต, หนฺทาหํ กลฺยาณํ กโรมิ กาเยน วาจาย มนสา’ติ”? โส เอวมาห “นาสกฺขิสฺสํ ภนฺเต, ปมาทสฺสํ ภนฺเต”ติ.}

\prose{ตเมนํ ภิกฺขเว ยโม ราชา เอวมาห “อมฺโภ ปุริสปมาทตาย\csromansharedfootnote{7b22cde9d31cfa18}{ปมาทวตาย (สี, สฺยา, กํ, อิ) ม 3. 218 ปิฏฺเฐปิ.} น กลฺยาณมกาสิ กาเยน วาจาย มนสา, ตคฺฆ ตฺวํ\csromansharedfootnote{672b3cfb0804b7fc}{ตํ (ก)} อมฺโภ ปุริส ตถา กริสฺสนฺติ, ยถา ตํ\csromansharedfootnote{7829f85bcb36afaa}{เต (ก)} ปมตฺตํ. ตํ โข ปน เต เอตํ\csromansharedfootnote{87e916e7d5926830}{ตํ โข ปเนตํ (สี, สฺยา, กํ, อิ)} ปาปกมฺมํ\csromansharedfootnote{9dee3fc78620f566}{ปาปํ กมฺมํ (สี)} เนว มาตรา กตํ, น ปิตรา กตํ, น ภาตรา กตํ, น ภคินิยา กตํ, น มิตฺตามจฺเจหิ กตํ, น ญาติสาโลหิเตหิ กตํ, น เทวตาหิ กตํ, น สมณ\-พฺราหฺมเณหิ กตํ, อถ โข ตยาเวตํ ปาปกมฺมํ กตํ, ตฺวญฺเญเวตสฺส วิปากํ ปฏิสํเวทิสฺสสี”ติ. (1)}

\prose{ตเมนํ ภิกฺขเว ยโม ราชา ปฐมํ เทวทูตํ สมนุยุญฺชิตฺวา สมนุคาหิตฺวา สมนุภาสิตฺวา ทุติยํ เทวทูตํ สมนุยุญฺชติ สมนุคาหติ สมนุภาสติ “อมฺโภ ปุริส น ตฺวํ อทฺทส มนุสฺเสสุ ทุติยํ เทวทูตํ ปาตุภูตนฺ”ติ? โส เอวมาห “นาทฺทสํ ภนฺเต”ติ. ตเมนํ ภิกฺขเว ยโม ราชา เอวมาห “อมฺโภ ปุริส น ตฺวํ อทฺทส มนุสฺเสสุ อิตฺถิํ วา ปุริสํ วา อาพาธิกํ ทุกฺขิตํ พาฬฺหคิลานํ สเก มุตฺตกรีเส ปลิปนฺนํ เสมานํ อญฺเญหิ วุฏฺฐาปิยมานํ อญฺเญหิ สํเวสิยมานนฺ”ติ? โส เอวมาห “อทฺทสํ ภนฺเต”ติ.}

\prose{ตเมนํ ภิกฺขเว ยโม ราชา เอวมาห “อมฺโภ ปุริส ตสฺส เต วิญฺญุสฺส สโต มหลฺลกสฺส น เอตทโหสิ ‘อหมฺปิ โขมฺหิ พฺยาธิธมฺโม พฺยาธิํ อนตีโต, หนฺทาหํ กลฺยาณํ กโรมิ กาเยน วาจาย มนสา’ติ”? โส เอวมาห “นาสกฺขิสฺสํ ภนฺเต, ปมาทสฺสํ ภนฺเต”ติ.}

\prose{ตเมนํ ภิกฺขเว ยโม ราชา เอวมาห “อมฺโภ ปุริส ปมาทตาย น กลฺยาณมกาสิ กาเยน วาจาย มนสา, ตคฺฆ ตฺวํ อมฺโภ ปุริส ตถา กริสฺสนฺติ, ยถา ตํ ปมตฺตํ. ตํ โข ปน เต เอตํ ปาปกมฺมํ เนว มาตรา กตํ, น ปิตรา กตํ, น ภาตรา กตํ, น ภคินิยา กตํ, น \csromanfolio{139}มิตฺตามจฺเจหิ กตํ, น ญาติสาโลหิเตหิ กตํ, น เทวตาหิ กตํ, น สมณ\-พฺราหฺมเณหิ กตํ. อถ โข ตยาเวตํ ปาปกมฺมํ กตํ, ตฺวญฺเญเวตสฺส วิปากํ ปฏิสํเวทิสฺสสี”ติ. (2)}

\proseitem{36}{ตเมนํ ภิกฺขเว ยโม ราชา ทุติยํ เทวทูตํ สมนุยุญฺชิตฺวา สมนุคาหิตฺวา สมนุภาสิตฺวา ตติยํ เทวทูตํ สมนุยุญฺชติ สมนุคาหติ สมนุภาสติ “อมฺโภ ปุริส น ตฺวํ อทฺทส มนุสฺเสสุ ตติยํ เทวทูตํ ปาตุภูตนฺ”ติ? โส เอวมาห “นาทฺทสํ ภนฺเต”ติ.}

\prose{ตเมนํ ภิกฺขเว ยโม ราชา เอวมาห “อมฺโภ ปุริส น ตฺวํ อทฺทส มนุสฺเสสุ อิตฺถิํ วา ปุริสํ วา เอกาหมตํ วา ทฺวีหมตํ วา ตีหมตํ วา อุทฺธุมาตกํ วินีลกํ วิปุพฺพกชาตนฺ”ติ? โส เอวมาห “อทฺทสํ ภนฺเต”ติ.}

\prose{ตเมนํ ภิกฺขเว ยโม ราชา เอวมาห “อมฺโภ ปุริส ตสฺส เต วิญฺญุสฺส สโต มหลฺลกสฺส น เอตทโหสิ ‘อหมฺปิ โขมฺหิ มรณธมฺโม มรณํ อนตีโต, หนฺทาหํ กลฺยาณํ กโรมิ กาเยน วาจาย มนสา’ติ”? โส เอวมาห “นาสกฺขิสฺสํ ภนฺเต, ปมาทสฺสํ ภนฺเต”ติ.}

\prose{ตเมนํ ภิกฺขเว ยโม ราชา เอวมาห “อมฺโภ ปุริส ปมาทตาย น กลฺยาณมกาสิ กาเยน วาจาย มนสา, ตคฺฆ ตฺวํ อมฺโภ ปุริส ตถา กริสฺสนฺติ, ยถา ตํ ปมตฺตํ. ตํ โข ปน เต เอตํ ปาปกมฺมํ เนว มาตรา กตํ, น ปิตรา กตํ, น ภาตรา กตํ, น ภคินิยา กตํ, น มิตฺตามจฺเจหิ กตํ, น ญาติสาโลหิเตหิ กตํ, น เทวตาหิ กตํ, น สมณ\-พฺราหฺมเณหิ กตํ. อถ โข ตยาเวตํ ปาปกมฺมํ กตํ, ตฺวญฺเญเวตสฺส วิปากํ ปฏิสํเวทิสฺสสี”ติ. (3)}

\prose{ตเมนํ ภิกฺขเว ยโม ราชา ตติยํ เทวทูตํ สมนุยุญฺชิตฺวา สมนุคาหิตฺวา สมนุภาสิตฺวา ตุณฺหี โหติ. ตเมนํ ภิกฺขเว นิรยปาลา ปญฺจ\-วิธ\-พนฺธนํ นาม การณํ กโรนฺติ, ตตฺตํ อโยขิลํ หตฺเถ คเมนฺติ, ตตฺถํ อโยขิลํ ทุติยสฺมิํ หตฺเถ คเมนฺติ, ตตฺตํ อโยขิลํ ปาเท คเมนฺติ, ตตฺตํ อโยขิลํ ทุติยสฺมิํ ปาเท คเมนฺติ, ตตฺตํ อโยขิลํ มชฺเฌ อุรสฺมิํ คเมนฺติ. โส ตตฺถ ทุกฺขา ติพฺพา\csromansharedfootnote{dbad2df778e3784d}{ติปฺปา (สี)} ขรา กฏุกา เวทนา เวทิยติ, น จ ตาว กาลํ กโรติ, ยาว น ตํ ปาปกมฺมํ พฺยนฺตีโหติ.}

\prose{\csromanfolio{140}ตเมนํ ภิกฺขเว นิรยปาลา สํเวเสตฺวา\csromansharedfootnote{908704402fe04e7a}{สํกฑฺฒิภฺวา (ก)} กุธารีหิ ตจฺฉนฺติ. โส ตตฺถ ทุกฺขา ติพฺพา ขรา กฏุกา เวทนา เวทิยติ, น จ ตาว กาลํ กโรติ, ยาว น ตํ ปาปกมฺมํ พฺยนฺตีโหติ.}

\prose{ตเมนํ ภิกฺขเว นิรยปาลา อุทฺธํปาทํ อโธสิรํ คเหตฺวา วาสีหิ ตจฺฉนฺติ ฯเปฯ. ตเมนํ ภิกฺขเว นิรยปาลา รเถ โยเชตฺวา อาทิตฺตาย ภูมิยา สมฺปชฺชลิตาย สโชติภูตาย\csromansharedfootnote{d3f7d3c37f771d61}{สญฺโชติภูตาย (สฺยา, กํ)} สาเรนฺติปิ ปจฺจาสาเรนฺติปิ ฯเปฯ. ตเมนํ ภิกฺขเว นิรยปาลา มหนฺตํ องฺคารปพฺพตํ อาทิตฺตํ สมฺปชฺชลิตํ สโชติภูตํ อาโรเปนฺติปิ โอโรเปนฺติปิ ฯเปฯ. ตเมนํ ภิกฺขเว นิรยปาลา อุทฺธํปาทํ อโธสิรํ คเหตฺวา ตตฺตาย โลหกุมฺภิยา ปกฺขิปนฺติ อาทิตฺตาย สมฺปชฺชลิตาย สโชติภูตาย. โส ตตฺถ เผณุทฺเทหกํ ปจฺจมาโน สกิมฺปิ อุทฺธํ คจฺฉติ, สกิมฺปิ อโธ คจฺฉติ, สกิมฺปิ ติริยํ คจฺฉติ. โส ตตฺถ ทุกฺขา ติพฺพา ขรา กฏุกา เวทนา เวทิยติ, น จ ตาว กาลํ กโรติ, ยาว น ตํ ปาปกมฺมํ พฺยนฺตีโหติ. ตเมนํ ภิกฺขเว นิรยปาลา มหานิรเย ปกฺขิปนฺติ. โส โข ปน ภิกฺขเว มหานิรโย\csromandash{}}

\setlength{\csromangathapagewidth}{0pt}
\setlength{\csromangathawakwidth}{0pt}
\csromangathameasuregroup{จตุกฺกณฺโณ จตุทฺวาโร, \\ อโยปาการปริยนฺโต, \\ ตสฺส อโยมยา ภูมิ,}{\csromangathabat{จตุกฺกณฺโณ จตุทฺวาโร,}{วิภตฺโต ภาคโส มิโต.} \\ \csromangathabat{อโยปาการปริยนฺโต,}{อยสา ปฏิกุชฺชิโต.} \\ \csromangathabat{ตสฺส อโยมยา ภูมิ,}{ชลิตา เตชสา ยุตา.}}
\csromangathasetleft{จตุกฺกณฺโณ จตุทฺวาโร, \\ อโยปาการปริยนฺโต, \\ ตสฺส อโยมยา ภูมิ,}
\csromangathagroup{\csromangathastanza{\csromangathabat{จตุกฺกณฺโณ จตุทฺวาโร,}{วิภตฺโต ภาคโส มิโต.} \\ \csromangathabat{อโยปาการปริยนฺโต,}{อยสา ปฏิกุชฺชิโต.} \\ \csromangathabat{ตสฺส อโยมยา ภูมิ,}{ชลิตา เตชสา ยุตา.}}}

\vspace{\tipitakaparskip}
\csromancenter{สมนฺตา โยชนสตํ, ผริตฺวา ติฏฺฐติ สพฺพทาติ.*}
\prose{ภูตปุพฺพํ ภิกฺขเว ยมสฺส รญฺโญ เอตทโหสิ “เย กิร โภ โลเก ปาปกานิ กมฺมานิ กโรนฺติ, เต เอวรูปา วิวิธา กมฺมการณา กรียนฺติ. อโห วตาหํ มนุสฺสตฺตํ ลเภยฺยํ, ตถาคโต จ โลเก อุปฺปชฺเชยฺย อรหํ สมฺมาสมฺพุทฺโธ, ตญฺจาหํ ภควนฺตํ ปยิรุปาเสยฺยํ, โส จ เม ภควา ธมฺมํ เทเสยฺย, ตสฺส จาหํ ภควโต ธมฺมํ อาชาเนยฺยนฺ”ติ. ตํ โข ปนาหํ ภิกฺขเว น อญฺญสฺส สมณสฺส วา พฺราหฺมณสฺส วา สุตฺวา เอวํ วทามิ, อปิ จ โข ภิกฺขเว ยเทว เม สามํ ญาตํ สามํ ทิฏฺฐํ สามํ วิทิตํ, ตเทวาหํ วทามีติ.}

\chapterheadpageread{4. เทวทูตวคฺค}
\setlength{\csromangathapagewidth}{0pt}
\setlength{\csromangathawakwidth}{0pt}
\csromangathameasuregroup{“โจทิตา เทวทูเตหิ, \\ เต ทีฆรตฺตํ โสจนฺติ, \\ เย จ โข เทวทูเตหิ, \\ โจทิตา นปฺปมชฺชนฺติ, \\ อุปาทาเน ภยํ ทิสฺวา, \\ อนุปาทา วิมุจฺจนฺติ, \\ เต อปฺปมตฺตา\textsuperscript{1} สุขิโน\textsuperscript{1}, \\ สพฺพเวรภยาตีตา,}{\csromangathabat{“โจทิตา เทวทูเตหิ,}{เย ปมชฺชนฺติ มาณวา.} \\ \csromangathabat{เต ทีฆรตฺตํ โสจนฺติ,}{หีนกายูปคา นรา.} \\ \csromangathabat{เย จ โข เทวทูเตหิ,}{สนฺโต สปฺปุริสา อิธ.} \\ \csromangathabat{โจทิตา นปฺปมชฺชนฺติ,}{อริยธมฺเม กุทาจนํ.} \\ \csromangathabat{อุปาทาเน ภยํ ทิสฺวา,}{ชาติมรณสมฺภเว.} \\ \csromangathabat{อนุปาทา วิมุจฺจนฺติ,}{ชาติมรณสงฺขเย.} \\ \csromangathabat{เต อปฺปมตฺตา\textsuperscript{1} สุขิโน\textsuperscript{1},}{ทิฏฺฐธมฺมาภินิพฺพุตา.} \\ \csromangathabat{สพฺพเวรภยาตีตา,}{สพฺพทุกฺขํ อุปจฺจคุ นฺติ.ฉฏฺฐํ.}}
\csromangathasetleft{“โจทิตา เทวทูเตหิ, \\ เต ทีฆรตฺตํ โสจนฺติ, \\ เย จ โข เทวทูเตหิ, \\ โจทิตา นปฺปมชฺชนฺติ, \\ อุปาทาเน ภยํ ทิสฺวา, \\ อนุปาทา วิมุจฺจนฺติ, \\ เต อปฺปมตฺตา\textsuperscript{1} สุขิโน\textsuperscript{1}, \\ สพฺพเวรภยาตีตา,}
\csromangathagroup{\csromangathastanza{\csromanfolio{141}\csromangathabat{“โจทิตา เทวทูเตหิ,}{เย ปมชฺชนฺติ มาณวา.} \\ \csromangathabat{เต ทีฆรตฺตํ โสจนฺติ,}{หีนกายูปคา นรา.}}\csromangathastanza{\csromangathabat{เย จ โข เทวทูเตหิ,}{สนฺโต สปฺปุริสา อิธ.} \\ \csromangathabat{โจทิตา นปฺปมชฺชนฺติ,}{อริยธมฺเม กุทาจนํ.}}\csromangathastanza{\csromangathabat{อุปาทาเน ภยํ ทิสฺวา,}{ชาติมรณสมฺภเว.} \\ \csromangathabat{อนุปาทา วิมุจฺจนฺติ,}{ชาติมรณสงฺขเย.}}\csromangathastanza{\csromangathabat{เต อปฺปมตฺตา\csromansharedfootnote{b1c2cefb7c257d4e}{เต โขปฺปมตฺตา (สี), เต เขมปฺปตฺตา (สฺยา, กํ, อิ) ม 3. 225 ปิฏฺเฐปิ.} สุขิโน\csromansharedfootnote{18268e84946c746d}{สุขิตา (สี, สฺยา)},}{ทิฏฺฐธมฺมา\-ภินิพฺพุตา.} \\ \csromangathabat{สพฺพเวรภยาตีตา,}{สพฺพทุกฺขํ อุปจฺจคุ นฺติ.  ฉฏฺฐํ.}}}

\csromanheader{2}{7. จตุมหาราช\-สุตฺต}
\proseitem{37}{อฏฺฐมิยํ ภิกฺขเว ปกฺขสฺส จตุนฺนํ มหาราชานํ อมจฺจา ปาริสชฺชา อิมํ โลกํ อนุวิจรนฺติ “กจฺจิ พหู มนุสฺสา มนุสฺเสสุ มตฺเตยฺยา เปตฺเตยฺยา สามญฺญา พฺรหฺมญฺญา กุเล เชฏฺฐาปจายิโน, อุโปสถํ อุปวสนฺติ ปฏิชาคโรนฺติ, ปุญฺญานิ กโรนฺตี”ติ. จาตุทฺทสิํ ภิกฺขเว ปกฺขสฺส จตุนฺนํ มหาราชานํ ปุตฺตา อิมํ โลกํ อนุวิจรนฺติ “กจฺจิ พหู มนุสฺสา มนุสฺเสสุ มตฺเตยฺยา เปตฺเตยฺยา สามญฺญา พฺรหฺมญฺญา กุเล เชฏฺฐาปจายิโน อุโปสถํ อุปวสนฺติ, ปฏิชาคโรนฺติ, ปุญฺญานิ กโรนฺตี”ติ. ตทหุ ภิกฺขเว อุโปสเถ ปนฺนรเส จตฺตาโร มหาราชาโน สามญฺเญว อิมํ โลกํ อนุวิจรนฺติ “กจฺจิ พหู มนุสฺสา มนุสฺเสสุ มตฺเตยฺยา เปตฺเตยฺยา สามญฺญา พฺรหฺมญฺญา กุเล เชฏฺฐาปจายิโน, อุโปสถํ อุปวสนฺติ ปฏิชาคโรนฺติ, ปุญฺญานิ กโรนฺตี”ติ.}

\prose{สเจ ภิกฺขเว อปฺปกา โหนฺติ มนุสฺสา มนุสฺเสสุ มตฺเตยฺยา เปตฺเตยฺยา สามญฺญา พฺรหฺมญฺญา กุเล เชฏฺฐาปจายิโน, อุโปสถํ อุปวสนฺติ ปฏิชาคโรนฺติ, ปุญฺญานิ กโรนฺติ. ตเมนํ ภิกฺขเว จตฺตาโร มหาราชาโน เทวานํ ตาวติํสานํ สุธมฺมาย สภาย สนฺนิสินฺนานํ สนฺนิปติตานํ อาโรเจนฺติ “อปฺปกา โข มาริสา มนุสฺสา มนุสฺเสสุ มตฺเตยฺยา เปตฺเตยฺยา สามญฺญา พฺรหฺมญฺญา กุเล เชฏฺฐาปจายิโน, อุโปสถํ อุปวสนฺติ ปฏิชาคโรนฺติ, ปุญฺญานิ กโรนฺตี”ติ. เตน โข \csromanfolio{142}ภิกฺขเว เทวา ตาวติํสา อนตฺตมนา โหนฺติ “ทิพฺพา วต โภ กายา ปริหายิสฺสนฺติ, ปริปูริสฺสนฺติ อสุรกายา”ติ.}

\prose{สเจ ปน ภิกฺขเว พหู โหนฺติ มนุสฺสา มนุสฺเสสุ มตฺเตยฺยา เปตฺเตยฺยา สามญฺญา พฺรหฺมญฺญา กุเล เชฏฺฐาปจายิโน, อุโปสถํ อุปวสนฺติ ปฏิชาคโรนฺติ, ปุญฺญานิ กโรนฺติ. ตเมนํ ภิกฺขเว จตฺตาโร มหาราชาโน เทวานํ ตาวติํสานํ สุธมฺมาย สภาย สนฺนิสินฺนานํ สนฺนิปติตานํ อาโรเจนฺติ “พหู โข มาริสา มนุสฺสา มนุสฺเสสุ มตฺเตยฺยา เปตฺเตยฺยา สามญฺญา พฺรหฺมญฺญา กุเล เชฏฺฐาปจายิโน, อุโปสถํ อุปวสนฺติ ปฏิชาคโรนฺติ, ปุญฺญานิ กโรนฺตี”ติ. เตน ภิกฺขเว เทวา ตาวติํสา อตฺตมนา โหนฺติ “ทิพฺพา วต โภ กายา ปริปูริสฺสนฺติ, ปริหายิสฺสนฺติ อสุรกายา”ติ.}

\prose{ภูตปุพฺพํ ภิกฺขเว สกฺโก เทวานมินฺโท เทเว ตาวติํเส อนุนยมาโน ตายํ เวลายํ อิมํ คาถํ อภาสิ\csromandash{}}

\setlength{\csromangathapagewidth}{0pt}
\setlength{\csromangathawakwidth}{0pt}
\csromangathameasuregroup{“จาตุทฺทสิํ ปญฺจทสิํ, \\ ปาฏิหาริยปกฺขญฺจ, \\ อุโปสถํ อุปวเสยฺย,}{\csromangathabat{“จาตุทฺทสิํ ปญฺจทสิํ,}{ยา จ ปกฺขสฺส อฏฺฐมี.} \\ \csromangathabat{ปาฏิหาริยปกฺขญฺจ,}{อฏฺฐงฺคสุสมาคตํ.} \\ \csromangathabat{อุโปสถํ อุปวเสยฺย,}{โยปิสฺส\textsuperscript{1} มาทิโส นโร”ติ.}}
\csromangathasetleft{“จาตุทฺทสิํ ปญฺจทสิํ, \\ ปาฏิหาริยปกฺขญฺจ, \\ อุโปสถํ อุปวเสยฺย,}
\csromangathagroup{\csromangathastanza{\csromangathabat{“จาตุทฺทสิํ ปญฺจทสิํ,}{ยา จ ปกฺขสฺส อฏฺฐมี.} \\ \csromangathabat{ปาฏิหาริยปกฺขญฺจ,}{อฏฺฐงฺคสุสมาคตํ.} \\ \csromangathabat{อุโปสถํ อุปวเสยฺย,}{โยปิสฺส\csromansharedfootnote{e72905a4cd583100}{โยปสฺส (สี, สฺยา, กํ, อิ)} มาทิโส นโร”ติ.}}}

\prose{สา โข ปเนสา ภิกฺขเว สกฺเกน เทวานมินฺเทน คาถา ทุคฺคีตา น สุคีตา, ทุพฺภาสิตา น สุภาสิตา. ตํ กิสฺส เหตุ? สกฺโก หิ ภิกฺขเว เทวานมินฺโท อวีตราโค อวีตโทโส อวีตโมโห.}

\prose{โย จ โข โส ภิกฺขเว ภิกฺขุ อรหํ ขีณาสโว วุสิตวา พฺรหฺมจริโย กตกรณีโย โอหิตภาโร อนุปฺปตฺต\-สทตฺโถ ปริกฺขีณ\-ภวสํโยชโน สมฺมทญฺญา วิมุตฺโต, ตสฺส โข เอตํ ภิกฺขเว ภิกฺขุโน\csromansharedfootnote{4513346a21d2be9b}{ตสฺส โข เอตํ ภิกฺขุโน (สี, สฺยา), ตสฺส โข เอวํ ภิกฺขเว ภิกฺขุโน (ก)} กลฺลํ วจนาย\csromandash{}}

\setlength{\csromangathapagewidth}{0pt}
\setlength{\csromangathawakwidth}{0pt}
\csromangathameasuregroup{“จาตุทฺทสิํ ปญฺจทสิํ, \\ ปาฏิหาริยปกฺขญฺจ, \\ อุโปสถํ อุปวเสยฺย,}{\csromangathabat{“จาตุทฺทสิํ ปญฺจทสิํ,}{ยา จ ปกฺขสฺส อฏฺฐมี.} \\ \csromangathabat{ปาฏิหาริยปกฺขญฺจ,}{อฏฺฐงฺคสุสมาคตํ.} \\ \csromangathabat{อุโปสถํ อุปวเสยฺย,}{โยปิสฺส มาทิโส นโร”ติ.}}
\csromangathasetleft{“จาตุทฺทสิํ ปญฺจทสิํ, \\ ปาฏิหาริยปกฺขญฺจ, \\ อุโปสถํ อุปวเสยฺย,}
\csromangathagroup{\csromangathastanza{\csromangathabat{“จาตุทฺทสิํ ปญฺจทสิํ,}{ยา จ ปกฺขสฺส อฏฺฐมี.} \\ \csromangathabat{ปาฏิหาริยปกฺขญฺจ,}{อฏฺฐงฺคสุสมาคตํ.} \\ \csromangathabat{อุโปสถํ อุปวเสยฺย,}{โยปิสฺส มาทิโส นโร”ติ.}}}

\chapterheadpageread{4. เทวทูตวคฺค}
\prose{\csromanfolio{143}ตํ กิสฺส เหตุ, โส หิ ภิกฺขเว ภิกฺขุ วีตราโค วีตโทโส วีตโมโหติ.  สตฺตมํ.}

\csromanheader{2}{8. ทุติย\-จตุ\-มหาราช\-สุตฺต}
\proseitem{38}{ภูตปุพฺพํ ภิกฺขเว สกฺโก เทวานมินฺโท เทเว ตาวติํเส อนุนยมาโน ตายํ เวลายํ อิมํ คาถํ อภาสิ\csromandash{}}

\setlength{\csromangathapagewidth}{0pt}
\setlength{\csromangathawakwidth}{0pt}
\csromangathameasuregroup{“จาตุทฺทสิํ ปญฺจทสิํ, \\ ปาฏิหาริยปกฺขญฺจ, \\ อุโปสถํ อุปวเสยฺย,}{\csromangathabat{“จาตุทฺทสิํ ปญฺจทสิํ,}{ยา จ ปกฺขสฺส อฏฺฐมี.} \\ \csromangathabat{ปาฏิหาริยปกฺขญฺจ,}{อฏฺฐงฺคสุสมาคตํ.} \\ \csromangathabat{อุโปสถํ อุปวเสยฺย,}{โยปิสฺส มาทิโส นโร”ติ.}}
\csromangathasetleft{“จาตุทฺทสิํ ปญฺจทสิํ, \\ ปาฏิหาริยปกฺขญฺจ, \\ อุโปสถํ อุปวเสยฺย,}
\csromangathagroup{\csromangathastanza{\csromangathabat{“จาตุทฺทสิํ ปญฺจทสิํ,}{ยา จ ปกฺขสฺส อฏฺฐมี.} \\ \csromangathabat{ปาฏิหาริยปกฺขญฺจ,}{อฏฺฐงฺคสุสมาคตํ.} \\ \csromangathabat{อุโปสถํ อุปวเสยฺย,}{โยปิสฺส มาทิโส นโร”ติ.}}}

\prose{สา โข ปเนสา ภิกฺขเว สกฺเกน เทวานมินฺเทน คาถา ทุคฺคีตา น สุคีตา, ทุพฺภาสิตา น สุภาสิตา. ตํ กิสฺส เหตุ? สกฺโก หิ ภิกฺขเว เทวานมินฺโท อปริมุตฺโต ชาติยา ชราย มรเณน โสเกหิ ปริเทเวหิ ทุกฺเขหิ โทมนสฺเสหิ อุปายาเสหิ, อปริมุตฺโต ทุกฺขสฺมาติ วทามิ.}

\prose{โย จ โข โส ภิกฺขเว ภิกฺขุ อรหํ ขีณาสโว วุสิตวา กตกรณีโย โอหิตภาโร อนุปฺปตฺต\-สทตฺโถ ปริกฺขีณ\-ภวสํโยชโน สมฺมทญฺญา วิมุตฺโต. ตสฺส โข เอตํ ภิกฺขเว ภิกฺขุโน กลฺลํ วจนาย\csromandash{}}

\setlength{\csromangathapagewidth}{0pt}
\setlength{\csromangathawakwidth}{0pt}
\csromangathameasuregroup{“จาตุทฺทสิํ ปญฺจทสิํ, \\ ปาฏิหาริยปกฺขญฺจ, \\ อุโปสถํ อุปวเสยฺย,}{\csromangathabat{“จาตุทฺทสิํ ปญฺจทสิํ,}{ยา จ ปกฺขสฺส อฏฺฐมี.} \\ \csromangathabat{ปาฏิหาริยปกฺขญฺจ,}{อฏฺฐงฺคสุสมาคตํ.} \\ \csromangathabat{อุโปสถํ อุปวเสยฺย,}{โยปิสฺส มาทิโส นโร”ติ.}}
\csromangathasetleft{“จาตุทฺทสิํ ปญฺจทสิํ, \\ ปาฏิหาริยปกฺขญฺจ, \\ อุโปสถํ อุปวเสยฺย,}
\csromangathagroup{\csromangathastanza{\csromangathabat{“จาตุทฺทสิํ ปญฺจทสิํ,}{ยา จ ปกฺขสฺส อฏฺฐมี.} \\ \csromangathabat{ปาฏิหาริยปกฺขญฺจ,}{อฏฺฐงฺคสุสมาคตํ.} \\ \csromangathabat{อุโปสถํ อุปวเสยฺย,}{โยปิสฺส มาทิโส นโร”ติ.}}}

\prose{ตํ กิสฺส เหตุ? โส หิ ภิกฺขเว ภิกฺขุ ปริมุตฺโต ชาติยา ชราย มรเณน โสเกหิ ปริเทเวหิ ทุกฺเขหิ โทมนสฺเสหิ อุปายาเสหิ, ปริมุตฺโต ทุกฺขสฺมาติ วทามีติ.  อฏฺฐมํ.}

\csromanheader{2}{9. \textbf{สุขุมาลสุตฺต}}
\proseitem{39}{สุขุมาโล อหํ ภิกฺขเว ปรม\-สุขุมาโล อจฺจนฺต\-สุขุมาโล, มม สุทํ ภิกฺขเว ปิตุ นิเวสเน โปกฺขรณิโย การิตา โหนฺติ. เอกตฺถ สุทํ ภิกฺขเว อุปฺปลํ วปฺปติ\csromansharedfootnote{f291f6baa165d764}{ปุปฺผติ (สี, อิ)} เอกตฺถ ปทุมํ เอกตฺถ ปุณฺฑรีกํ ยาวเทว มมตฺถาย. น โข ปนสฺสาหํ ภิกฺขเว อกาสิกํ จนฺทนํ ธาเรมิ\csromansharedfootnote{1c6efa4179f78c4c}{กาสิกํ จนฺทนํ ธาเรมิ (สฺยา, กํ, ก), อกาสิกํ ธาเรมิ(?)}, \csromanfolio{144}กาสิกํ ภิกฺขเว สุ เม ตํ เวฐนํ โหติ, กาสิกา กญฺจุกา กาสิกํ นิวาสนํ กาสิโก อุตฺตราสงฺโค. รตฺตินฺทิวํ\csromansharedfootnote{dd3fab0193bd20cb}{รตฺติทิวํ (ก)} โข ปน เม สุ ตํ ภิกฺขเว เสตจฺฉตฺตํ ธารียติ “มา นํ ผุสิ สีตํ วา อุณฺหํ วา ติณํ วา รโช วา อุสฺสาโว วา”ติ.}

\prose{ตสฺส มยฺหํ ภิกฺขเว ตโย ปาสาทา อเหสุํ\csromandash{}เอโก เหมนฺติโก, เอโก คิมฺหิโก, เอโก วสฺสิโก. โส โข อหํ ภิกฺขเว วสฺสิเก ปาสาเท วสฺสิเก จตฺตาโร มาเส นิปฺปุริเสหิ ตูริเยหิ ปริจารยมาโน\csromansharedfootnote{d56e7a7c21e56aff}{ปริจาริยมาโน (สฺยา, กํ, อิ, ก)} น เหฏฺฐาปาสาทํ โอโรหามิ. ยถา โข ปน ภิกฺขเว อญฺเญสํ นิเวสเน ทาส\-กมฺมกร\-โปริสสฺส กณาชกํ โภชนํ ทียติ พิลงฺคทุติยํ. เอวเมวสฺสุ เม ภิกฺขเว ปิตุ นิเวสเน ทาส\-กมฺมกร\-โปริสสฺส สาลิมํโสทโน ทียติ.}

\prose{ตสฺส มยฺหํ ภิกฺขเว เอวรูปาย อิทฺธิยา สมนฺนาคตสฺส, เอวรูเปน จ สุขุมาเลน เอตทโหสิ “อสฺสุตวา โข ปุถุชฺชโน อตฺตนา ชราธมฺโม สมาโน ชรํ อนตีโต ปรํ ชิณฺณํ ทิสฺวา อฏฺฏียติ หรายติ ชิคุจฺฉติ อตฺตานํเยว อติสิตฺวา. อหมฺปิ โขมฺหิ ชราธมฺโม ชรํ อนตีโต, อหญฺเจว\csromansharedfootnote{cf278ffcd2fa33b2}{อหญฺเจ (?)} โข ปน ชราธมฺโม สมาโน ชรํ อนตีโต ปรํ ชิณฺณํ ทิสฺวา อฏฺฏีเยยฺยํ หราเยยฺยํ ชิคุจฺเฉยฺยํ, น เมตํ อสฺส ปติรูปนฺ”ติ. ตสฺส มยฺหํ ภิกฺขเว อิติ ปฏิสญฺจิกฺขโต โย โยพฺพเน โยพฺพนมโท, โส สพฺพโส ปหียิ.}

\prose{“อสฺสุตวา โข ปุถุชฺชโน อตฺตนา พฺยาธิธมฺโม สมาโน พฺยาธิํ อนตีโต ปรํ พฺยาธิตํ ทิสฺวา อฏฺฏียติ หรายติ ชิคุจฺฉติ อตฺตานํเยว อติสิตฺวา. อหมฺปิ โขมฺหิ พฺยาธิธมฺโม พฺยาธิํ อนตีโต, อหญฺเจว โข ปน พฺยาธิธมฺโม สมาโน พฺยาธิํ อนตีโต ปรํ พฺยาธิตํ ทิสฺวา อฏฺฏีเยยฺยํ หราเยยฺยํ ชิคุจฺเฉยฺยํ, น เมตํ อสฺส ปติรูปนฺ”ติ. ตสฺส มยฺหํ ภิกฺขเว อิติ ปฏิสญฺจิกฺขโต โย อาโรเคฺย อาโรคฺยมโท, โส สพฺพโส ปหียิ.}

\prose{“อสฺสุตวา โข ปุถุชฺชโน อตฺตนา มรณธมฺโม สมาโน มรณํ อนตีโต ปรํ มตํ ทิสฺวา อฏฺฏียติ หรายติ ชิคุจฺฉติ อตฺตานํเยว อติสิตฺวา. อหิมฺปิ โขมฺหิ มรณธมฺโม มรณํ อนตีโต, อหญฺเจว โข \csromanfolio{145}ปน มรณธมฺโม สมาโน มรณํ อนตีโต ปรํ มตํ ทิสฺวา อฏฺฏีเยยฺยํ หราเยยฺยํ ชิคุจฺเฉยฺยํ, น เมตํ อสฺส ปติรูปนฺ”ติ. ตสฺส มยฺหํ ภิกฺขเว อิติ ปฏิสญฺจิกฺขโต โย ชีวิเต ชีวิตมโท, โส สพฺพโส ปหียีติ.}

\prose{ตโยเม ภิกฺขเว มทา. กตเม ตโย? โยพฺพนมโท อาโรคฺยมโท ชีวิตมโท. โยพฺพนมท\-มตฺโต วา ภิกฺขเว อสฺสุตวา ปุถุชฺชโน กาเยน ทุจฺจริตํ จรติ, วาจาย ทุจฺจริตํ จรติ, มนสา ทุจฺจริตํ จรติ. โส กาเยน ทุจฺจริตํ จริตฺวา วาจาย ทุจฺจริตํ จริตฺวา มนสา ทุจฺจริตํ จริตฺวา กายสฺส เภทา ปรํ มรณา อปายํ ทุคฺคติํ วินิปาตํ นิรยํ อุปปชฺชติ. อาโรคฺยมท\-มตฺโต วา ภิกฺขเว อสฺสุตวา ปุถุชฺชโน ฯเปฯ. ชีวิตมท\-มตฺโต วา ภิกฺขเว อสฺสุตวา ปุถุชฺชโน กาเยน ทุจฺจริตํ จรติ, วาจาย ทุจฺจริตํ จรติ, มนสา ทุจฺจริตํ จรติ. โส กาเยน ทุจฺจริตํ จริตฺวา วาจาย ทุจฺจริตํ จริตฺวา มนสา ทุจฺจริตํ จริตฺวา กายสฺส เภทา ปรํ มรณา อปายํ ทุคฺคติํ วินิปาตํ นิรยํ อุปปชฺชติ.}

\prose{โยพฺพนมท\-มตฺโต วา ภิกฺขเว ภิกฺขุ สิกฺขํ ปจฺจกฺขาย หีนายาวตฺตติ. อาโรคฺยมท\-มตฺโต วา ภิกฺขเว ภิกฺขุ ฯเปฯ. ชีวิตมท\-มตฺโต วา ภิกฺขเว ภิกฺขุ สิกฺขํ ปจฺจกฺขาย หีนายาวตฺตตีติ.}

\setlength{\csromangathapagewidth}{0pt}
\setlength{\csromangathawakwidth}{0pt}
\csromangathameasuregroup{พฺยาธิธมฺมา ชราธมฺมา, \\ ยถาธมฺมา\textsuperscript{1} ตถาสนฺตา, \\ อหญฺเจ ตํ ชิคุจฺเฉยฺยํ, \\ น เมตํ ปติรูปสฺส, \\ โสหํ เอวํ วิหรนฺโต, \\ อาโรเคฺย โยพฺพนสฺมิญฺจ, \\ สพฺเพ มเท อภิโภสฺมิ\textsuperscript{1}, \\ ตสฺส เม อหุ อุสฺสาโห, \\ นาหํ ภพฺโพ เอตรหิ, \\ อนิวตฺติ ภวิสฺสามิ,}{\csromangathabat{พฺยาธิธมฺมา ชราธมฺมา,}{อโถ มรณธมฺมิโน.} \\ \csromangathabat{ยถาธมฺมา\textsuperscript{1} ตถาสนฺตา,}{ชิคุจฺฉนฺติ ปุถุชฺชนา.} \\ \csromangathabat{อหญฺเจ ตํ ชิคุจฺเฉยฺยํ,}{เอวํธมฺเมสุ ปาณิสุ.} \\ \csromangathabat{น เมตํ ปติรูปสฺส,}{มม เอวํ วิหาริโน.} \\ \csromangathabat{โสหํ เอวํ วิหรนฺโต,}{ญตฺวา ธมฺมํ นิรูปธิํ.} \\ \csromangathabat{อาโรเคฺย โยพฺพนสฺมิญฺจ,}{ชีวิตสฺมิญฺจ เย มทา.} \\ \csromangathabat{สพฺเพ มเท อภิโภสฺมิ\textsuperscript{1},}{เนกฺขมฺเม ทฏฺฐุ เขมตํ.} \\ \csromangathabat{ตสฺส เม อหุ อุสฺสาโห,}{นิพฺพานํ อภิปสฺสโต.} \\ \csromangathabat{นาหํ ภพฺโพ เอตรหิ,}{กามานิ ปฏิเสวิตุํ.} \\ \csromangathabat{อนิวตฺติ ภวิสฺสามิ,}{พฺรหฺมจริยปรายโณติ.นวมํ.}}
\csromangathasetleft{พฺยาธิธมฺมา ชราธมฺมา, \\ ยถาธมฺมา\textsuperscript{1} ตถาสนฺตา, \\ อหญฺเจ ตํ ชิคุจฺเฉยฺยํ, \\ น เมตํ ปติรูปสฺส, \\ โสหํ เอวํ วิหรนฺโต, \\ อาโรเคฺย โยพฺพนสฺมิญฺจ, \\ สพฺเพ มเท อภิโภสฺมิ\textsuperscript{1}, \\ ตสฺส เม อหุ อุสฺสาโห, \\ นาหํ ภพฺโพ เอตรหิ, \\ อนิวตฺติ ภวิสฺสามิ,}
\csromangathagroup{\csromangathastanza{\csromangathabat{พฺยาธิธมฺมา ชราธมฺมา,}{อโถ มรณธมฺมิโน.} \\ \csromangathabat{ยถาธมฺมา\csromansharedfootnote{69463ddc48c63fdb}{พฺยาธิธมฺโม ชราธมฺโม, อโถ มรณธมฺมิโก. ยถา ธมฺโม (ก)} ตถาสนฺตา,}{ชิคุจฺฉนฺติ ปุถุชฺชนา.}}\csromangathastanza{\csromangathabat{อหญฺเจ ตํ ชิคุจฺเฉยฺยํ,}{เอวํธมฺเมสุ ปาณิสุ.} \\ \csromangathabat{น เมตํ ปติรูปสฺส,}{มม เอวํ วิหาริโน.}}\csromangathastanza{\csromangathabat{โสหํ เอวํ วิหรนฺโต,}{ญตฺวา ธมฺมํ นิรูปธิํ.} \\ \csromangathabat{อาโรเคฺย โยพฺพนสฺมิญฺจ,}{ชีวิตสฺมิญฺจ เย มทา.}}\csromangathastanza{\csromangathabat{สพฺเพ มเท อภิโภสฺมิ\csromansharedfootnote{1146c19a51b7a6ca}{อตีโตสฺมิ (ก)},}{เนกฺขมฺเม ทฏฺฐุ เขมตํ.} \\ \csromangathabat{ตสฺส เม อหุ อุสฺสาโห,}{นิพฺพานํ อภิปสฺสโต.}}\csromangathastanza{\csromangathabat{นาหํ ภพฺโพ เอตรหิ,}{กามานิ ปฏิเสวิตุํ.} \\ \csromangathabat{อนิวตฺติ ภวิสฺสามิ,}{พฺรหฺมจริยปรายโณติ.  นวมํ.}}}

\csromanpalirecto
\gambhira{ติกนิปาตปาฬิ 10. อาธิปเตยฺย\-สุตฺต}
\proseitem{40}{\csromanfolio{146}ตีณิมานิ ภิกฺขเว อาธิปเตยฺยานิ. กตมานิ ตีณิ? อตฺตา\-ธิปเตยฺยํ โลกา\-ธิปเตยฺยํ ธมฺมา\-ธิปเตยฺยํ. กตมญฺจ ภิกฺขเว อตฺตา\-ธิปเตยฺยํ? อิธ ภิกฺขเว ภิกฺขุ อรญฺญคโต วา รุกฺขมูลคโต วา สุญฺญาคารคโต วา อิติ ปฏิสญฺจิกฺขติ “น โข ปนาหํ จีวรเหตุ อคารสฺมา อนคาริยํ ปพฺพชิโต น ปิณฺฑปาตเหตุ น เสนาสนเหตุ น อิติภวา\-ภ\-ว\-เหตุ อคารสฺมา อนคาริยํ ปพฺพชิโต. อปิ จ โขมฺหิ โอติณฺโณ ชาติยา ชราย มรเณน โสเกหิ ปริเทเวหิ ทุกฺเขหิ โทมนสฺเสหิ อุปายาเสหิ ทุกฺโขติณฺโณ ทุกฺขปเรโต, อปฺเปว นาม อิมสฺส เกวลสฺส ทุกฺข\-กฺขนฺธสฺส อนฺตกิริยา ปญฺญาเยถาติ. อหญฺเจว โข ปน ยาทิสเก\csromansharedfootnote{5b187f3250c6e408}{ยาทิสเก วา (สี, อิ, ก)} กาเม โอหาย อคารสฺมา อนคาริยํ ปพฺพชิโต, ตาทิสเก วา\csromansharedfootnote{89bc7870b446eb08}{จ (ก)} กาเม ปริเยเสยฺยํ ตโต วา\csromansharedfootnote{89bc7870b446eb08}{จ (ก)} ปาปิฏฺฐตเร, น เมตํ ปติรูปนฺ”ติ. โส อิติ ปฏิสญฺจิกฺขติ “อารทฺธํ โข ปน เม วีริยํ ภวิสฺสติ อสลฺลีนํ, อุปฏฺฐิตา สติ อสมฺมุฏฺฐา, ปสฺสทฺโธ กาโย อสารทฺโธ, สมาหิตํ จิตฺตํ เอกคฺคนฺ”ติ. โส อตฺตานํเยว อธิปติํ กริตฺวา อกุสลํ ปชหหิ, กุสลํ ภาเวติ, สาวชฺชํ ปชหติ, อนวชฺชํ ภาเวติ, สุทฺธํ อตฺตานํ ปริหรติ. อิทํ วุจฺจติ ภิกฺขเว อตฺตา\-ธิปเตยฺยํ.}

\prose{กตมญฺจ ภิกฺขเว โลกา\-ธิปเตยฺยํ? อิธ ภิกฺขเว ภิกฺขุ อรญฺญคโต วา รุกฺขมูลคโต วา สุญฺญาคารคโต วา อิติ ปฏิสญฺจิกฺขติ “น โข ปนาหํ จีวรเหตุ อคารสฺมา อนคาริยํ ปพฺพชิโต น ปิณฺฑปาตเหตุ น เสนาสนเหตุ น อิติภวา\-ภ\-ว\-เหตุ อคารสฺมา อนคาริยํ ปพฺพชิโต. อปิ จ โขมฺหิ โอติณฺโณ ชาติยา ชราย มรเณน โสเกหิ ปริเทเวหิ ทุกฺเขหิ โทมนสฺเสหิ อุปายาเสหิ ทุกฺโขติณฺโณ ทุกฺขปเรโต, อปฺเปว นาม อิมสฺส เกวลสฺส ทุกฺข\-กฺขนฺธสฺส อนฺตกิริยา ปญฺญาเยถาติ. อหญฺเจว โข ปน เอวํ ปพฺพชิโต สมาโน กามวิตกฺกํ วา วิตกฺเกยํ, พฺยาปาท\-วิตกฺกํ วา วิตกฺเกยฺยํ, วิหิํสา\-วิตกฺกํ วา วิตกฺเกยฺยํ, มหา โข ปนายํ โลกสนฺนิวาโส, มหนฺตสฺมิํ โข ปน โลกสนฺนิวาเส สนฺติ สมณพฺราหฺมณา อิทฺธิมนฺโต ทิพฺพจกฺขุกา ปรจิตฺต\-วิทุโน, เต ทูรโตปิ ปสฺสนฺติ, อาสนฺนาปิ น ทิสฺสนฺติ, เจตสาปิ จิตฺตํ ปชานนฺติ\csromansharedfootnote{5fa9b0375c0e45b0}{ชานนฺติ (ก)}, เตปิ มํ เอวํ ชาเนยฺยุํ “ปสฺสถ โภ อิมํ \csromanfolio{147}กุลปุตฺตํ, สทฺธา อคารสฺมา อนคาริยํ ปพฺพชิโต สมาโน โวกิณฺโณ วิหรติ ปาปเกหิ อกุสเลหิ ธมฺเมหี”ติ. เทวตาปิ โข สนฺติ อิทฺธิมนฺตินิโย ทิพฺพจกฺขุกา ปรจิตฺต\-วิทุนิโย, ตา ทูรโตปิ ปสฺสนฺติ, อาสนฺนาปิ น ทิสฺสนฺติ, เจตสาปิ จิตฺตํ ปชานนฺติ, ตาปิ มํ เอวํ ชาเนยฺยุํ “ปสฺสถ โภ อิมํ กุลปุตฺตํ, สทฺธา อคารสฺมา อนคาริยํ ปพฺพชิโต สมาโน โวกิณฺโณ วิหรติ ปาปเกหิ อกุสเลหิ ธมฺเมหี”ติ. โส อิติ ปฏิสญฺจิกฺขติ “อารทฺธํ โข ปน เม วีริยํ ภวิสฺสติ อสลฺลีนํ, อุปฏฺฐิตา สติ อสมฺมุฏฺฐา, ปสฺสทฺโธ กาโย อสารทฺโธ, สมาหิตํ จิตฺตํ เอกคฺคนฺ”ติ. โส โลกํเยว อธิปติํ กริตฺวา อกุสลํ ปชหติ, กุสลํ ภาเวติ, สาวชฺชํ ปชหติ, อนวชฺชํ ภาเวติ, สุทฺธํ อตฺตานํ ปริหรติ. อิทํ วุจฺจติ ภิกฺขเว โลกา\-ธิปเตยฺยํ.}

\prose{กตมญฺจ ภิกฺขเว ธมฺมา\-ธิปเตยฺยํ? อิธ ภิกฺขเว ภิกฺขุ อรญฺญคโต วา รุกฺขมูลคโต วา สุญฺญาคารคโต วา อิติ ปฏิสญฺจิกฺขติ “น โข ปนาหํ จีวรเหตุ อคารสฺมา อนคาริยํ ปพฺพชิโต น ปิณฺฑปาตเหตุ น เสนาสนเหตุ น อิติภวา\-ภ\-ว\-เหตุ อคารสฺมา อนคาริยํ ปพฺพชฺชิโต. อปิ จ โขมฺหิ โอติณฺโณ ชาติยา ชราย มรเณน โสเกหิ ปริเทเวหิ ทุกฺเขหิ โทมนสฺเสหิ อุปายาเสหิ ทุกฺโขติณฺโณ ทุกฺขปเรโต, อปฺเปว นาม อิมสฺส เกวลสฺส ทุกฺข\-กฺขนฺธสฺส อนฺตกิริยา ปญฺญาเยถาติ. สฺวากฺขาโต ภควตา ธมฺโม สนฺทิฏฺฐิโก อกาลิโก เอหิปสฺสิโก โอปเนยฺยิโก ปจฺจตฺตํ เวทิตพฺโพ วิญฺญูหีติ. สนฺติ โข ปน เม สพฺรหฺมจารี ชานํ ปสฺสํ วิหรนฺติ. อหญฺเจว โข ปน เอวํ สฺวากฺขาเต ธมฺมวินเย ปพฺพชิโต สมาโน กุสีโต วิหเรยฺยํ ปมตฺโต, น เมตํ อสฺส ปติรูปนฺ”ติ. โส อิติ ปฏิสญฺจิกฺขติ “อารทฺธํ โข ปน เม วีริยํ ภวิสฺสติ อสลฺลีนํ, อุปฏฺฐิตา สติ อสมฺมุฏฺฐา, ปสฺสทฺโธ กาโย อสารทฺโธ, สมาหิตํ จิตฺตํ เอกคฺคนฺ”ติ. โส ธมฺมํเยว อธิปติํ กริตฺวา อกุสลํ ปชหติ, กุสลํ ภาเวติ, สาวชฺชํ ปชหติ, อนวชฺชํ ภาเวติ, สุทฺธํ อตฺตานํ ปริหรติ. อิทํ วุจฺจติ ภิกฺขเว ธมฺมา\-ธิปเตยฺยํ. อิมานิ โข ภิกฺขเว ตีณิ อาธิปเตยฺยานีติ.}

\setlength{\csromangathapagewidth}{0pt}
\setlength{\csromangathawakwidth}{0pt}
\csromangathameasuregroup{นตฺถิ โลเก รโห นาม, \\ อตฺตา เต ปุริส ชานาติ, \\ กลฺยาณํ วต โภ สกฺขิ, \\ โย สนฺตํ อตฺตนิ ปาปํ,}{\csromangathabat{นตฺถิ โลเก รโห นาม,}{ปาปกมฺมํ ปกุพฺพโต.} \\ \csromangathabat{อตฺตา เต ปุริส ชานาติ,}{สจฺจํ วา ยทิ วา มุสา.} \\ \csromangathabat{กลฺยาณํ วต โภ สกฺขิ,}{อตฺตานํ อติมญฺญสิ.} \\ \csromangathabat{โย สนฺตํ อตฺตนิ ปาปํ,}{อตฺตานํ ปริคูหสิ.} \\ ปสฺสนฺติ เทวา จ ตถาคตา จ, \\ โลกสฺมิํ พาลํ วิสมํ จรนฺตํ. \\ ตสฺมา หิ อตฺตาธิปเตยฺยโก จ\textsuperscript{1}, \\ โลกาธิโป จ นิปโก จ ฌายี. \\ ธมฺมาธิโป จ อนุธมฺมจารี, \\ น หียติ สจฺจปรกฺกโม มุนิ. \\ ปสยฺห มารํ อภิภุยฺย อนฺตกํ, \\ โย จ ผุสี ชาติกฺขยํ ปธานวา. \\ โส ตาทิโส โลกวิทู สุเมโธ, \\ สพฺเพสุ ธมฺเมสุ อตมฺมโย มุนีติ.ทสมํ.}
\csromangathameasurewak{ปสฺสนฺติ เทวา จ ตถาคตา จ, \\ โลกสฺมิํ พาลํ วิสมํ จรนฺตํ. \\ ตสฺมา หิ อตฺตาธิปเตยฺยโก จ\textsuperscript{1}, \\ โลกาธิโป จ นิปโก จ ฌายี. \\ ธมฺมาธิโป จ อนุธมฺมจารี, \\ น หียติ สจฺจปรกฺกโม มุนิ. \\ ปสยฺห มารํ อภิภุยฺย อนฺตกํ, \\ โย จ ผุสี ชาติกฺขยํ ปธานวา. \\ โส ตาทิโส โลกวิทู สุเมโธ, \\ สพฺเพสุ ธมฺเมสุ อตมฺมโย มุนีติ.ทสมํ.}
\csromangathasetleft{นตฺถิ โลเก รโห นาม, \\ อตฺตา เต ปุริส ชานาติ, \\ กลฺยาณํ วต โภ สกฺขิ, \\ โย สนฺตํ อตฺตนิ ปาปํ,}
\csromangathagroupwithclosermidruled{\csromangathastanza{\csromangathabat{นตฺถิ โลเก รโห นาม,}{ปาปกมฺมํ ปกุพฺพโต.} \\ \csromangathabat{อตฺตา เต ปุริส ชานาติ,}{สจฺจํ วา ยทิ วา มุสา.}}\csromangathastanza{\csromanfolio{148}\csromangathabat{กลฺยาณํ วต โภ สกฺขิ,}{อตฺตานํ อติมญฺญสิ.} \\ \csromangathabat{โย สนฺตํ อตฺตนิ ปาปํ,}{อตฺตานํ ปริคูหสิ.}}\csromangathastanza{\csromangathawak{ปสฺสนฺติ เทวา จ ตถาคตา จ, \\ โลกสฺมิํ พาลํ วิสมํ จรนฺตํ. \\ ตสฺมา หิ อตฺตาธิปเตยฺยโก จ\csromansharedfootnote{1d5d2975551bea8f}{อตฺตาธิปโก สโต จเร (สี, สฺยา, กํ, อิ)}, \\ โลกาธิโป จ นิปโก จ ฌายี.}}\csromangathastanza{\csromangathawak{ธมฺมาธิโป จ อนุธมฺมจารี, \\ น หียติ สจฺจปรกฺกโม มุนิ. \\ ปสยฺห มารํ อภิภุยฺย อนฺตกํ, \\ โย จ ผุสี ชาติกฺขยํ ปธานวา. \\ โส ตาทิโส โลกวิทู สุเมโธ, \\ สพฺเพสุ ธมฺเมสุ อตมฺมโย มุนีติ.  ทสมํ.}}}{เทวทูตวคฺโค จตุตฺโถ.}
\csromancenter{\textbf{ตสฺสุทฺทานํ}}
\setlength{\csromangathapagewidth}{0pt}
\setlength{\csromangathawakwidth}{0pt}
\csromangathameasuregroup{พฺรหฺม อานนฺท สาริปุตฺโต, \\ ทูตา ทุเว จ ราชาโน,}{\csromangathabat{พฺรหฺม อานนฺท สาริปุตฺโต,}{นิทานํ หตฺถเกน จ.} \\ \csromangathabat{ทูตา ทุเว จ ราชาโน,}{สุขุมาลาธิปเตยฺเยน จาติ.}}
\csromangathasetleft{พฺรหฺม อานนฺท สาริปุตฺโต, \\ ทูตา ทุเว จ ราชาโน,}
\csromangathagroup{\csromangathastanza{\csromangathabat{พฺรหฺม อานนฺท สาริปุตฺโต,}{นิทานํ หตฺถเกน จ.} \\ \csromangathabat{ทูตา ทุเว จ ราชาโน,}{สุขุมาลาธิปเตยฺเยน จาติ.}}}

\chapterhead{5. \textbf{จูฬวคฺค}}
\csromanheader{2}{1. สมฺมุขีภาว\-สุตฺต}
\proseitem{41}{ติณฺณํ ภิกฺขเว สมฺมุขีภาวา สทฺโธ กุลปุตฺโต พหุํ ปุญฺญํ ปสวติ. กตเมสํ ติณฺณํ? สทฺธาย ภิกฺขเว สมฺมุขีภาวา สทฺโธ กุลปุตฺโต พหุํ ปุญฺญํ ปสวติ, เทยฺยธมฺมสฺส ภิกฺขเว สมฺมุขีภาวา สทฺโธ กุลปุตฺโต พหุํ ปุญฺญํ ปสวติ, ทกฺขิเณยฺยานํ ภิกฺขเว สมฺมุขีภาวา สทฺโธ กุลปุตฺโต พหุํ ปุญฺญํ ปสวติ. อิเมสํ โข ภิกฺขเว ติณฺณํ สมฺมุขีภาวา สทฺโธ กุลปุตฺโต พหุํ ปุญฺญํ ปสวตีติ.  ปฐมํ.}

\chapterheadpageread{5. จูฬวคฺค \textbf{2. ติฐานสุตฺต}}
\proseitem{42}{\csromanfolio{149}ตีหิ ภิกฺขเว ฐาเนหิ สทฺโธ ปสนฺโน เวทิตพฺโพ. กตเมหิ ตีหิ? สีลวนฺตานํ ทสฺสนกาโม โหติ, สทฺธมฺมํ โสตุกาโม โหติ, วิคต\-มล\-มจฺเฉเรน เจตสา อคารํ อชฺฌาวสติ มุตฺตจาโค ปยตปาณิ โวสคฺครโต ยาจโยโค ทาน\-สํวิภาค\-รโต. อิเมหิ โข ภิกฺขเว ตีหิ ฐาเนหิ สทฺโธ ปสนฺโน เวทิตพฺโพ.}

\setlength{\csromangathapagewidth}{0pt}
\setlength{\csromangathawakwidth}{0pt}
\csromangathameasuregroup{ทสฺสนกาโม สีลวตํ, \\ วินเย มจฺเฉรมลํ,}{\csromangathabat{ทสฺสนกาโม สีลวตํ,}{สทฺธมฺมํ โสตุมิจฺฉติ.} \\ \csromangathabat{วินเย มจฺเฉรมลํ,}{ส เว “สทฺโธ”ติ วุจฺจตีติ.ทุติยํ.}}
\csromangathasetleft{ทสฺสนกาโม สีลวตํ, \\ วินเย มจฺเฉรมลํ,}
\csromangathagroup{\csromangathastanza{\csromangathabat{ทสฺสนกาโม สีลวตํ,}{สทฺธมฺมํ โสตุมิจฺฉติ.} \\ \csromangathabat{วินเย มจฺเฉรมลํ,}{ส เว “สทฺโธ”ติ วุจฺจตีติ.  ทุติยํ.}}}

\csromanheader{2}{3. \textbf{อตฺถวสสุตฺต}}
\proseitem{43}{ตโย ภิกฺขเว อตฺถวเส สมฺปสฺสมาเนน อลเมว ปเรสํ ธมฺมํ เทเสตุํ. กตเม ตโย? โย ธมฺมํ เทเสติ, โส อตฺถ\-ปฺปฏิสํเวที จ โหติ ธมฺม\-ปฺปฏิสํเวที จ. โย ธมฺมํ สุณาติ, โส อตฺถ\-ปฺปฏิสํเวที จ โหติ ธมฺม\-ปฺปฏิสํเวที จ. โย เจว ธมฺมํ เทเสติ, โย จ ธมฺมํ สุณาติ, อุโภ อตฺถ\-ปฺปฏิสํเวทิโน จ โหนฺติ ธมฺม\-ปฺปฏิสํเวทิโน จ. อิเม โข ภิกฺขเว ตโย อตฺถวเส สมฺปสฺสมาเนน อลเมว ปเรสํ ธมฺมํ เทเสตุนฺติ.  ตติยํ.}

\csromanheader{2}{4. กถา\-ปวตฺติ\-สุตฺต}
\proseitem{44}{ตีหิ ภิกฺขเว ฐาเนหิ กถา ปวตฺตินี โหติ. กตเมหิ ตีหิ? โย ธมฺมํ เทเสติ, โส อตฺถ\-ปฺปฏิสํเวที จ โหติ ธมฺม\-ปฺปฏิสํเวที จ. โย ธมฺมํ สุณาติ, โส อตฺถ\-ปฺปฏิสํเวที จ โหติ ธมฺม\-ปฺปฏิสํเวที จ. โย เจว ธมฺมํ เทเสติ, โย จ ธมฺมํ สุณาติ, อุโภ อตฺถ\-ปฺปฏิสํเวทิโน จ โหนฺติ ธมฺม\-ปฺปฏิสํเวทิโน จ. อิเมหิ โข ภิกฺขเว ตีหิ ฐาเนหิ กถา ปวตฺตินี โหติ.  จตุตฺถํ.}

\csromanheader{2}{5. \textbf{ปณฺฑิตสุตฺต}}
\proseitem{45}{ตีณิมานิ ภิกฺขเว ปณฺฑิต\-ปญฺญตฺตานิ สปฺปุริส\-ปญฺญตฺตานิ. กตมานิ ตีณิ? ทานํ ภิกฺขเว ปณฺฑิต\-ปญฺญตฺตํ สปฺปุริส\-ปญฺญตฺตํ, ปพฺพชฺชา \csromanfolio{150}ภิกฺขเว ปณฺฑิต\-ปญฺญตฺตา สปฺปุริส\-ปญฺญตฺตา, มาตาปิตูนํ ภิกฺขเว อุปฏฺฐานํ ปณฺฑิต\-ปญฺญตฺตํ สปฺปุริส\-ปญฺญตฺตํ. อิมานิ โข ภิกฺขเว ตีณิ ปณฺฑิต\-ปญฺญตฺตานิ สปฺปุริส\-ปญฺญตฺตานีติ.}

\setlength{\csromangathapagewidth}{0pt}
\setlength{\csromangathawakwidth}{0pt}
\csromangathameasuregroup{สพฺภิ ทานํ อุปญฺญตฺตํ, \\ มาตาปิตุ อุปฏฺฐานํ, \\ สตํ เอตานิ ฐานานิ, \\ อริโย ทสฺสนสมฺปนฺโน,}{\csromangathabat{สพฺภิ ทานํ อุปญฺญตฺตํ,}{อหิํสา สํยโม ทโม.} \\ \csromangathabat{มาตาปิตุ อุปฏฺฐานํ,}{สนฺตานํ พฺรหฺมจารินํ.} \\ \csromangathabat{สตํ เอตานิ ฐานานิ,}{ยานิ เสเวถ ปณฺฑิโต.} \\ \csromangathabat{อริโย ทสฺสนสมฺปนฺโน,}{ส โลกํ ภชเต สิวนฺติ.} \\ ปญฺจมํ.}
\csromangathameasurewak{ปญฺจมํ.}
\csromangathasetleft{สพฺภิ ทานํ อุปญฺญตฺตํ, \\ มาตาปิตุ อุปฏฺฐานํ, \\ สตํ เอตานิ ฐานานิ, \\ อริโย ทสฺสนสมฺปนฺโน,}
\csromangathagroup{\csromangathastanza{\csromangathabat{สพฺภิ ทานํ อุปญฺญตฺตํ,}{อหิํสา สํยโม ทโม.} \\ \csromangathabat{มาตาปิตุ อุปฏฺฐานํ,}{สนฺตานํ พฺรหฺมจารินํ.}}\csromangathastanza{\csromangathabat{สตํ เอตานิ ฐานานิ,}{ยานิ เสเวถ ปณฺฑิโต.} \\ \csromangathabat{อริโย ทสฺสน\-สมฺปนฺโน,}{ส โลกํ ภชเต สิวนฺติ.}}\csromangathastanza{\csromangathawak{ปญฺจมํ.}}}

\csromanheader{2}{6. \textbf{สีลวนฺตสุตฺต}}
\proseitem{46}{ยํ ภิกฺขเว สีลวนฺโต ปพฺพชิตา คามํ วา นิคมํ วา อุปนิสฺสาย วิหรนฺติ, ตตฺถ มนุสฺสา ตีหิ ฐาเนหิ พหุํ ปุญฺญํ ปสวนฺติ. กตเมหิ ตีหิ? กาเยน วาจาย มนสา. ยํ ภิกฺขเว สีลวนฺโต ปพฺพชิตา คามํ วา นิคมํ วา อุปนิสฺสาย วิหรนฺติ, ตตฺถ มนุสฺสา อิเมหิ ตีหิ ฐาเนหิ พหุํ ปุญฺญํ ปสวนฺตีติ.  ฉฏฺฐํ.}

\csromanheader{2}{7. สงฺขตลกฺขณ\-สุตฺต}
\proseitem{47}{ตีณิมานิ ภิกฺขเว สงฺขตสฺส สงฺขต\-ลกฺขณานิ. กตมานิ ตีณิ? อุปฺปาโท ปญฺญายติ, วโย ปญฺญายติ, ฐิตสฺส อญฺญถตฺตํ ปญฺญายติ. อิมานิ โข ภิกฺขเว ตีณิ สงฺขตสฺส สงฺขตลกฺขณานีติ.  สตฺตมํ.}

\csromanheader{2}{8. อสงฺขตลกฺขณ\-สุตฺต}
\proseitem{48}{ตีณิมานิ ภิกฺขเว อสงฺขตสฺส อสงฺขตลกฺขณานิ. กตมานิ ตีณิ? น อุปฺปาโท ปญฺญายติ, น วโย ปญฺญายติ, น ฐิตสฺส อญฺญถตฺตํ ปญฺญายติ. อิมานิ โข ภิกฺขเว ตีณิ อสงฺขตสฺส อสงฺขตลกฺขณานีติ.  อฏฺฐมํ.}

\csromanheader{2}{9. ปพฺพตราช\-สุตฺต}
\proseitem{49}{หิมวนฺตํ ภิกฺขเว ปพฺพตราชํ นิสฺสาย มหาสาลา ตีหิ วฑฺฒีหิ วฑฺฒนฺติ. กตมาหิ ตีหิ? สาขา\-ปตฺต\-ปลาเสน วฑฺฒนฺติ, ตจปปฏิกาย \csromanfolio{151}วฑฺฒนฺติ, เผคฺคุสาเรน วฑฺฒนฺติ. หิมวนฺตํ ภิกฺขเว ปพฺพตราชํ นิสฺสาย มหาสาลา อิมาหิ ตีหิ วฑฺฒีหิ วฑฺฒนฺติ.}

\prose{เอวเมวํ โข ภิกฺขเว สทฺธํ กุลปติํ นิสฺสาย อนฺโต ชโน ตีหิ วฑฺฒีหิ วฑฺฒติ. กตเมหิ ตีหิ? สทฺธาย วฑฺฒติ, สีเลน วฑฺฒติ, ปญฺญาย วฑฺฒติ. สทฺธํ ภิกฺขเว กุลปติํ นิสฺสาย อนฺโต ชโน อิมาหิ ตีหิ วฑฺฒีหิ วฑฺฒตีติ.}

\setlength{\csromangathapagewidth}{0pt}
\setlength{\csromangathawakwidth}{0pt}
\csromangathameasuregroup{ยถาปิ ปพฺพโต เสโล, \\ ตํ รุกฺขา อุปนิสฺสาย, \\ ตเถว สีลสมฺปนฺนํ, \\ อุปนิสฺสาย วฑฺฒนฺติ, \\ อมจฺจา ญาติสํฆา จ, \\ ตฺยาสฺส สีลวโต สีลํ, \\ ปสฺสมานานุกุพฺพนฺติ, \\ อิธ ธมฺมํ จริตฺวาน, \\ นนฺทิโน เทวโลกสฺมิํ,}{\csromangathabat{ยถาปิ ปพฺพโต เสโล,}{อรญฺญสฺมิํ พฺรหาวเน.} \\ \csromangathabat{ตํ รุกฺขา อุปนิสฺสาย,}{วฑฺฒนฺเต เต วนปฺปตี.} \\ \csromangathabat{ตเถว สีลสมฺปนฺนํ,}{สทฺธํ กุลปติํ อิธ.} \\ \csromangathabat{อุปนิสฺสาย วฑฺฒนฺติ,}{ปุตฺตทารา จ พนฺธวา.} \\ \csromangathabat{อมจฺจา ญาติสํฆา จ,}{เย จสฺส อนุชีวิโน.} \\ \csromangathabat{ตฺยาสฺส สีลวโต สีลํ,}{จาคํ สุจริตานิ จ.} \\ \csromangathabat{ปสฺสมานานุกุพฺพนฺติ,}{อตฺตมตฺถํ\textsuperscript{1} วิจกฺขณา.} \\ \csromangathabat{อิธ ธมฺมํ จริตฺวาน,}{มคฺคํ สุคติคามินํ.} \\ \csromangathabat{นนฺทิโน เทวโลกสฺมิํ,}{โมทนฺติ กามกามิโนติ.} \\ นวมํ.}
\csromangathasetleft{ยถาปิ ปพฺพโต เสโล, \\ ตํ รุกฺขา อุปนิสฺสาย, \\ ตเถว สีลสมฺปนฺนํ, \\ อุปนิสฺสาย วฑฺฒนฺติ, \\ อมจฺจา ญาติสํฆา จ, \\ ตฺยาสฺส สีลวโต สีลํ, \\ ปสฺสมานานุกุพฺพนฺติ, \\ อิธ ธมฺมํ จริตฺวาน, \\ นนฺทิโน เทวโลกสฺมิํ,}
\csromangathagroup{\csromangathastanza{\csromangathabat{ยถาปิ ปพฺพโต เสโล,}{อรญฺญสฺมิํ พฺรหาวเน.} \\ \csromangathabat{ตํ รุกฺขา อุปนิสฺสาย,}{วฑฺฒนฺเต เต วนปฺปตี.}}\csromangathastanza{\csromangathabat{ตเถว สีลสมฺปนฺนํ,}{สทฺธํ กุลปติํ อิธ.} \\ \csromangathabat{อุปนิสฺสาย วฑฺฒนฺติ,}{ปุตฺตทารา จ พนฺธวา.}}\csromangathastanza{\csromangathabat{อมจฺจา ญาติสํฆา จ,}{เย จสฺส อนุชีวิโน.} \\ \csromangathabat{ตฺยาสฺส สีลวโต สีลํ,}{จาคํ สุจริตานิ จ.}}\csromangathastanza{\csromangathabat{ปสฺสมานานุกุพฺพนฺติ,}{อตฺตมตฺถํ\csromansharedfootnote{338dd9d64fd7ee1e}{[มิสฺสินฺคฺ โนเต 0]} วิจกฺขณา.} \\ \csromangathabat{อิธ ธมฺมํ จริตฺวาน,}{มคฺคํ สุคติคามินํ.}}\csromangathastanza{\csromangathabat{นนฺทิโน เทวโลกสฺมิํ,}{โมทนฺติ กามกามิโนติ.} \\ นวมํ.}}

\csromanheader{2}{10. อาตปฺป\-กรณีย\-สุตฺต}
\proseitem{50}{ตีหิ ภิกฺขเว ฐาเนหิ อาตปฺปํ กรณียํ. กตเมหิ ตีหิ? อนุปฺปนฺนานํ ปาปกานํ อกุสลานํ ธมฺมานํ อนุปฺปาทาย อาตปฺปํ กรณียํ, อนุปฺปนฺนานํ กุสลานํ ธมฺมานํ อุปฺปาทาย อาตปฺปํ กรณียํ, อุปฺปนฺนานํ สารีริกานํ เวทนานํ ทุกฺขานํ ติพฺพานํ ขรานํ กฏุกานํ อสาตานํ อมนาปานํ ปาณหรานํ อธิวาสนาย อาตปฺปํ กรณียํ. อิเมหิ ตีหิ ภิกฺขเว ฐาเนหิ อาตปฺปํ กรณียํ.}

\prose{ยโต โข ภิกฺขเว ภิกฺขุ อนุปฺปนฺนานํ ปาปกานํ อกุสลานํ ธมฺมานํ อนุปฺปาทาย อาตปฺปํ กโรติ, อนุปฺปนฺนานํ กุสลานํ ธมฺมานํ อุปฺปาทาย อาตปฺปํ กโรติ, อุปฺปนฺนานํ สารีริกานํ เวทนานํ ทุกฺขานํ ติพฺพานํ ขรานํ กฏุกานํ อสาตานํ อมนาปานํ ปาณหรานํ อธิวาสนาย \csromanfolio{152}อาตปฺปํ กโรติ. อยํ วุจฺจติ ภิกฺขเว ภิกฺขุ “อาตาปี นิปโก สโต สมฺมา ทุกฺขสฺส อนฺตกิริยายา”ติ.  ทสมํ.}

\csromanheader{2}{11. \textbf{มหาโจรสุตฺต}}
\proseitem{51}{ตีหิ ภิกฺขเว องฺเคหิ สมนฺนาคโต มหาโจโร สนฺธิมฺปิ ฉินฺทติ, นิลฺโลปมฺปิ หรติ, เอกาคาริกมฺปิ กโรติ, ปริปนฺเถปิ ติฏฺฐติ. กตเมหิ ตีหิ? อิธ ภิกฺขเว มหาโจโร วิสมนิสฺสิโต จ โหติ, คหนนิสฺสิโต จ โหติ, พลวนิสฺสิโต จ โหติ. กถญฺจ ภิกฺขเว มหาโจโร วิสมนิสฺสิโต โหติ? อิธ ภิกฺขเว มหาโจโร นทีวิทุคฺคํ วา นิสฺสิโต โหติ ปพฺพตวิสมํ วา. เอวํ โข ภิกฺขเว มหาโจโร วิสมนิสฺสิโต โหติ.}

\prose{กถญฺจ ภิกฺขเว มหาโจโร คหนนิสฺสิโต โหติ? อิธ ภิกฺขเว มหาโจโร ติณคหนํ วา นิสฺสิโต โหติ รุกฺขคหนํ วา โรธํ\csromansharedfootnote{08fa3979358ca9a1}{เคธํ (สี, อิ)} วา มหาวนสณฺฑํ วา. เอวํ โข ภิกฺขเว มหาโจโร คหนนิสฺสิโต โหติ.}

\prose{กถญฺจ ภิกฺขเว มหาโจโร พลวนิสฺสิโต โหติ? อิธ ภิกฺขเว มหาโจโร ราชานํ วา ราช\-มหามตฺตานํ วา นิสฺสิโต โหติ. ตสฺส เอวํ โหติ “สเจ มํ โกจิ กิญฺจิ วกฺขติ, อิเม เม ราชาโน วา ราชมหามตฺตา วา ปริโยธาย อตฺถํ ภณิสฺสนฺตี”ติ. สเจ นํ โกจิ กิญฺจิ อาห, ตฺยาสฺส ราชาโน วา ราชมหามตฺตา วา ปริโยธาย อตฺถํ ภณนฺติ. เอวํ โข ภิกฺขเว มหาโจโร พลวนิสฺสิโต โหติ. อิเม โข ภิกฺขเว ตีหิ องฺเคหิ สมนฺนาคโต มหาโจโร สนฺธิมฺปิ ฉินฺทติ, นิลฺโลปมฺปิ หรติ, เอกาคาริกมฺปิ กโรติ, ปริปนฺเถปิ ติฏฺฐติ.}

\prose{เอวเมวํ โข ภิกฺขเว ตีหิ องฺเคหิ สมนฺนาคโต ปาปภิกฺขุ ขตํ อุปหตํ อตฺตานํ ปริหรติ, สาวชฺโช จ โหติ สานุวชฺโช จ วิญฺญูนํ, พหุญฺจ อปุญฺญํ ปสวติ. กตเมหิ ตีหิ? อิธ ภิกฺขเว ปาปภิกฺขุ วิสมนิสฺสิโต จ โหติ คหนนิสฺสิโต จ พลวนิสฺสิโต จ.}

\prose{กถญฺจ ภิกฺขเว ปาปภิกฺขุ วิสมนิสฺสิโต โหติ? อิธ ภิกฺขเว ปาปภิกฺขุ วิสเมน กายกมฺเมน สมนฺนาคโต โหติ, วิสเมน \csromanfolio{153}วจีกมฺเมน สมนฺนาคโต โหติ, วิสเมน มโนกมฺเมน สมนฺนาคโต โหติ. เอวํ โข ภิกฺขเว ปาปภิกฺขุ วิสมนิสฺสิโต โหติ.}

\prose{กถญฺจ ภิกฺขเว ปาปภิกฺขุ คหนนิสฺสิโต โหติ? อิธ ภิกฺขเว ปาปภิกฺขุ มิจฺฉาทิฏฺฐิโก โหติ, อนฺตคฺคาหิกาย ทิฏฺฐิยา สมนฺนาคโต โหติ. เอวํ โข ภิกฺขเว ปาปภิกฺขุ คหนนิสฺสิโต โหติ.}

\prosewithclosermidruled{กถญฺจ ภิกฺขเว ปาปภิกฺขุ พลวนิสฺสิโต โหติ? อิธ ภิกฺขเว ปาปภิกฺขุ ราชานํ วา ราช\-มหามตฺตานํ วา นิสฺสิโต โหติ. ตสฺส เอวํ โหติ “สเจ มํ โกจิ กิญฺจิ วกฺขติ, อิเม เม ราชาโน วา ราชมหามตฺตา วา ปริโยธาย อตฺถํ ภณิสฺสนฺตี”ติ. สเจ นํ โกจิ กิญฺจิ อาห, ตฺยาสฺส ราชาโน วา ราชมหามตฺตา วา ปริโยธาย อตฺถํ ภณนฺติ. เอวํ โข ภิกฺขเว ปาปภิกฺขุ พลวนิสฺสิโต โหติ. อิเมหิ โข ภิกฺขเว ตีหิ ธมฺเมหิ สมนฺนาคโต ปาปภิกฺขุ ขตํ อุปหตํ อตฺตานํ ปริหรติ, สาวชฺโช จ โหติ สานุวชฺโช จ วิญฺญูนํ, พหุญฺจ อปุญฺญํ ปสวตีติ.  เอกาทสมํ.}{จูฬวคฺโค ปญฺจโม.}
\csromancenter{\textbf{ตสฺสุทฺทานํ}}
\setlength{\csromangathapagewidth}{0pt}
\setlength{\csromangathawakwidth}{0pt}
\csromangathameasuregroup{สมฺมุขี ฐานตฺถวสํ, \\ สงฺขตํ ปพฺพตาตปฺปํ,}{\csromangathabat{สมฺมุขี ฐานตฺถวสํ,}{ปวตฺติ ปณฺฑิต สีลวํ.} \\ \csromangathabat{สงฺขตํ ปพฺพตาตปฺปํ,}{มหาโจเรเนกาทสาติ\textsuperscript{1}.} \\ \textbf{ปฐโม ปณฺณาสโก สมตฺโต.}}
\csromangathasetleft{สมฺมุขี ฐานตฺถวสํ, \\ สงฺขตํ ปพฺพตาตปฺปํ,}
\csromangathagroup{\csromangathastanza{\csromangathabat{สมฺมุขี ฐานตฺถวสํ,}{ปวตฺติ ปณฺฑิต สีลวํ.} \\ \csromangathabat{สงฺขตํ ปพฺพตาตปฺปํ,}{มหาโจเรเนกาทสาติ\csromansharedfootnote{963d90228621d0d1}{มหาโจเรน เต ทสาติ (ก)}.} \\ \textbf{ปฐโม ปณฺณาสโก สมตฺโต.}}}

\chapterheadpageread{\textbf{(6) 1. พฺราหฺมณวคฺค}}
\csromanheader{2}{1. ปฐม\-เทฺว\-พฺราหฺมณ\-สุตฺต}
\proseitem{52}{\csromanfolio{154}อถ โข เทฺว พฺราหฺมณา ชิณฺณา วุทฺธา มหลฺลกา อทฺธคตา วโยอนุปฺปตฺตา วีสวสฺส\-สติกา ชาติยา เยน ภควา เตนุ\-ปสงฺกมิํสุ, อุปสงฺกมิตฺวา ภควตา สทฺธิํ สมฺโมทิํสุ, สมฺโมทนียํ กถํ สารณียํ วีติสาเรตฺวา เอกมนฺตํ นิสีทิํสุ, เอกมนฺตํ นิสินฺนา โข เต พฺราหฺมณา ภควนฺตํ เอตทโวจุํ “มยมสฺสุ โภ โคตม พฺราหฺมณา ชิณฺณา วุทฺธา มหลฺลกา อทฺธคตา วโยอนุปฺปตฺตา วีสวสฺส\-สติกา ชาติยา, เต จมฺหา อกตกลฺยาณา อกตกุสลา อกต\-ภีรุตฺตาณา, โอวทตุ โน ภวํ โคตโม, อนุสาสตุ โน ภวํ โคตโม, ยํ อมฺหากํ อสฺส ทีฆรตฺตํ หิตาย สุขายา”ติ.}

\prose{ตคฺฆ ตุเมฺห พฺราหฺมณา ชิณฺณา วุทฺธา มหลฺลกา อทฺธคตา วโยอนุปฺปตฺตา วีสวสฺส\-สติกา ชาติยา, เต จตฺถ อกตกลฺยาณา อกตกุสลา อกต\-ภีรุตฺตาณา, อุปนียติ โข อยํ พฺราหฺมณา โลโก ชราย พฺยาธินา มรเณน, เอวํ อุปนียมาเน โข พฺราหฺมณา โลเก ชราย พฺยาธินา มรเณน โย อิธ กาเยน สํยโม วาจาย สํยโม มนสา สํยโม, ตํ ตสฺส เปตสฺส ตาณญฺจ เลณญฺจ ทีปญฺจ สรณญฺจ ปรายณญฺจาติ.}

\setlength{\csromangathapagewidth}{0pt}
\setlength{\csromangathawakwidth}{0pt}
\csromangathameasuregroup{}{อุปนียติ ชีวิตมปฺปมายุ, \\ ชรูปนีตสฺส น สนฺติ ตาณา. \\ เอตํ ภยํ มรเณ เปกฺขมาโน, \\ ปุญฺญานิ กยิราถ สุขาวหานิ. \\ โยธ กาเยน สํยโม, \\ วาจาย อุท เจตสา. \\ ตํ ตสฺส เปตสฺส สุขาย โหติ, \\ ยํ ชีวมาโน ปกโรติ ปุญฺญนฺติ.\textsuperscript{1} . \\ ปฐมํ.}
\csromangathameasurewak{อุปนียติ ชีวิตมปฺปมายุ, \\ ชรูปนีตสฺส น สนฺติ ตาณา. \\ เอตํ ภยํ มรเณ เปกฺขมาโน, \\ ปุญฺญานิ กยิราถ สุขาวหานิ. \\ โยธ กาเยน สํยโม, \\ วาจาย อุท เจตสา. \\ ตํ ตสฺส เปตสฺส สุขาย โหติ, \\ ยํ ชีวมาโน ปกโรติ ปุญฺญนฺติ.\textsuperscript{1} . \\ ปฐมํ.}
\setlength{\csromangathaleftcol}{0pt}
\csromangathagroup{\csromangathastanza{\csromangathawak{อุปนียติ ชีวิตม\-ปฺปมายุ, \\ ชรูปนีตสฺส น สนฺติ ตาณา. \\ เอตํ ภยํ มรเณ เปกฺขมาโน, \\ ปุญฺญานิ กยิราถ สุขาวหานิ.}}\csromangathastanza{\csromangathawak{โยธ กาเยน สํยโม, \\ วาจาย อุท เจตสา. \\ ตํ ตสฺส เปตสฺส สุขาย โหติ, \\ ยํ ชีวมาโน ปกโรติ ปุญฺญนฺติ.\csromansharedfootnote{338dd9d64fd7ee1e}{[มิสฺสินฺคฺ โนเต 0]} .}}\csromangathastanza{\csromangathawak{ปฐมํ.}}}

\chapterheadpageread{(6) 1. พฺราหฺมณวคฺค 2. ทุติย\-เทฺว\-พฺราหฺมณ\-สุตฺต}
\proseitem{53}{\csromanfolio{155}อถ โข เทฺว พฺราหฺมณา ชิณฺณา วุทฺธา มหลฺลกา อทฺธคตา วโยอนุปฺปตฺตา วีสวสฺส\-สติกา ชาติยา เยน ภควา เตนุ\-ปสงฺกมิํสุ, อุปสงฺกมิตฺวา ภควนฺตํ อภิวาเทตฺวา เอกมนฺตํ นิสีทิํสุ, เอกมนฺตํ นิสินฺนา โข เต พฺราหฺมณา ภควนฺตํ เอตทโวจุํ “มยมสฺสุ โภ โคตม พฺราหฺมณา ชิณฺณา วุทฺธา มหลฺลกา อทฺธคตา วโยอนุปฺปตฺตา วีสวสฺส\-สติกา ชาติยา, เต จมฺหา อกตกลฺยาณา อกตกุสลา อกต\-ภีรุตฺตาณา, โอวทตุ โน ภวํ โคตโม, อนุสาสตุ โน ภวํ โคตโม, ยํ อมฺหากํ อสฺส ทีฆรตฺตํ หิตาย สุขายา”ติ.}

\prose{ตคฺฆ ตุเมฺห พฺราหฺมณา ชิณฺณา วุทฺธา มหลฺลกา อทฺธคตา วโยอนุปฺปตฺตา วีสวสฺส\-สติกา ชาติยา. เต จตฺถ อกตกลฺยาณา อกตกุสลา อกต\-ภีรุตฺตาณา, อาทิตฺโต โข อยํ พฺราหฺมณา โลโก ชราย พฺยาธินา มรเณน, เอวํ อาทิตฺเต โข พฺราหฺมณา โลเก ชราย พฺยาธินา มรเณน โย อิธ กาเยน สํยโม วาจาย สํยโม มนสา สํยโม, ตํ ตสฺส เปตสฺส ตาณญฺจ เลณญฺจ ทีปญฺจ สรณญฺจ ปรายณญฺจาติ.}

\setlength{\csromangathapagewidth}{0pt}
\setlength{\csromangathawakwidth}{0pt}
\csromangathameasuregroup{อาทิตฺตสฺมิํ อคารสฺมิํ, \\ ตํ ตสฺส โหติ อตฺถาย, \\ เอวํ อาทิตฺโต โข\textsuperscript{1} โลโก, \\ นีหเรเถว ทาเนน,}{\csromangathabat{อาทิตฺตสฺมิํ อคารสฺมิํ,}{ยํ นีหรติ ภาชนํ.} \\ \csromangathabat{ตํ ตสฺส โหติ อตฺถาย,}{โน จ ยํ ตตฺถ ฑยฺหติ.} \\ \csromangathabat{เอวํ อาทิตฺโต โข\textsuperscript{1} โลโก,}{ชราย มรเณน จ.} \\ \csromangathabat{นีหเรเถว ทาเนน,}{ทินฺนํ โหติ สุนีหตํ\textsuperscript{1}.} \\ โยธ กาเยน สํยโม, \\ วาจาย อุท เจตสา. \\ ตํ ตสฺส เปตสฺส สุขาย โหติ, \\ ยํ ชีวมาโน ปกโรติ ปุญฺญนฺติ.ทุติยํ.}
\csromangathameasurewak{โยธ กาเยน สํยโม, \\ วาจาย อุท เจตสา. \\ ตํ ตสฺส เปตสฺส สุขาย โหติ, \\ ยํ ชีวมาโน ปกโรติ ปุญฺญนฺติ.ทุติยํ.}
\csromangathasetleft{อาทิตฺตสฺมิํ อคารสฺมิํ, \\ ตํ ตสฺส โหติ อตฺถาย, \\ เอวํ อาทิตฺโต โข\textsuperscript{1} โลโก, \\ นีหเรเถว ทาเนน,}
\csromangathagroup{\csromangathastanza{\csromangathabat{อาทิตฺตสฺมิํ อคารสฺมิํ,}{ยํ นีหรติ ภาชนํ.} \\ \csromangathabat{ตํ ตสฺส โหติ อตฺถาย,}{โน จ ยํ ตตฺถ ฑยฺหติ.}}\csromangathastanza{\csromangathabat{เอวํ อาทิตฺโต โข\csromansharedfootnote{144af1a7e1585147}{เอวํ อาทีปิโต (สี, อิ), เอวํ อาทิตฺตโก (สฺยา, กํ) สํ 1. 29 ปิฏฺเฐปิ.} โลโก,}{ชราย มรเณน จ.} \\ \csromangathabat{นีหเรเถว ทาเนน,}{ทินฺนํ โหติ สุนีหตํ\csromansharedfootnote{144af1a7e1585147}{เอวํ อาทีปิโต (สี, อิ), เอวํ อาทิตฺตโก (สฺยา, กํ) สํ 1. 29 ปิฏฺเฐปิ.}.}}\csromangathastanza{\csromangathawak{โยธ กาเยน สํยโม, \\ วาจาย อุท เจตสา. \\ ตํ ตสฺส เปตสฺส สุขาย โหติ, \\ ยํ ชีวมาโน ปกโรติ ปุญฺญนฺติ.  ทุติยํ.}}}

\csromanheader{2}{3. อญฺญตร\-พฺราหฺมณ\-สุตฺต}
\proseitem{54}{อถ โข อญฺญตโร พฺราหฺมโณ เยน ภควา เตนุปสงฺกมิ, อุปสงฺกมิตฺวา ภควตา สทฺธิํ สมฺโมทิ ฯเปฯ เอกมนฺตํ นิสินฺโน โข โส \csromanfolio{156}พฺราหฺมโณ ภควนฺตํ เอตทโวจ “สนฺทิฏฺฐิโก ธมฺโม สนฺทิฏฺฐิโก ธมฺโมติ โภ โคตม วุจฺจติ, กิตฺตาวตา นุ โข โภ โคตม สนฺทิฏฺฐิโก ธมฺโม โหติ อกาลิโก เอหิปสฺสิโก โอปเนยฺยิโก\csromansharedfootnote{551a4ef641e488e2}{โอปนยิโก (สี, สฺยา, กํ, อิ) ปญฺจมสุตฺตสฺส ฏีกา โอโลเกตพฺพา. 2. ( ) เอตฺถนฺตเร ปาฏฺโฐ สีสฺยากํอิโปตฺถเกสุ น ทิสฺสติ, อิธปิ ทุฏฺฐมูฬฺหวาเรสุ. 3. น เจตสิกํ (สี, สฺยา, ก)} ปจฺจตฺตํ เวทิตพฺโพ วิญฺญูหี”ติ.}

\prose{รตฺโต โข พฺราหฺมณ ราเคน อภิภูโต ปริยาทินฺน\-จิตฺโต อตฺตพฺยาพาธายปิ เจเตติ, ปรพฺยาพาธายปิ เจเตติ, อุภย\-พฺยาพาธายปิ เจเตติ, เจตสิกมฺปิ ทุกฺขํ โทมนสฺสํ ปฏิสํเวเทติ. ราเค ปหีเน เนว\-ตฺตพฺยาพาธายปิ เจเตติ, น ปรพฺยาพาธายปิ เจเตติ, น อุภย\-พฺยาพาธายปิ เจเตติ, น เจตสิกํ ทุกฺขํ โทมนสฺสํ ปฏิสํเวเทติ. (รตฺโต โข ฯเปฯ กาเยน ทุจฺจริตํ จรติ, วาจาย ทุจฺจริตํ จรติ, มนสา ทุจฺจริตํ จรติ. ราเค ปหีเน เนว กาเยน ทุจฺจริตํ จรติ, น วาจาย ทุจฺจริตํ จรติ, น มนสา ทุจฺจริตํ จรติ. รตฺโต โข ฯเปฯ อตฺตตฺถมฺปิ ยถาภูตํ นปฺปชานาติ, ปรตฺถมฺปิ ยถาภูตํ นปฺปชานาติ, อุภยตฺถมฺปิ ยถาภูตํ นปฺปชานาติ. ราเค ปหีเน อตฺตตฺถมฺปิ ยถาภูตํ ปชานาติ, ปรตฺถมฺปิ ยถาภูตํ ปชานาติ, อุภยตฺถมฺปิ ยถาภูตํ ปชานาติ.)2 เอวมฺปิ โข พฺราหฺมณ สนฺทิฏฺฐิโก ธมฺโม โหติ ฯเปฯ.}

\prose{ทุฏฺโฐ โข พฺราหฺมณ โทเสน อภิภูโต ปริยาทินฺน\-จิตฺโต อตฺตพฺยาพาธายปิ เจเตติ, ปรพฺยาพาธายปิ เจเตติ, อุภย\-พฺยาพาธายปิ เจเตติ, เจตสิกมฺปิ ทุกฺขํ โทมนสฺสํ ปฏิสํเวเทติ. โทเส ปหีเน เนว\-ตฺตพฺยาพาธายปิ เจเตติ, น ปรพฺยาพาธายปิ เจเตติ, น อุภย\-พฺยาพาธายปิ เจเตติ, น เจตสิกมฺปิ\csromansharedfootnote{351abaa1618f4a18}{น เจตสิกํ (สี, สฺยา, ก)} ทุกฺขํ โทมนสฺสํ ปฏิสํเวเทติ. เอวมฺปิ โข พฺราหฺมณ สนฺทิฏฺฐิโก ธมฺโม โหติ ฯเปฯ.}

\prose{มูโฬฺห โข พฺราหฺมณ โมเหน อภิภูโต ปริยาทินฺน\-จิตฺโต อตฺตพฺยาพาธายปิ เจเตติ, ปรพฺยาพาธายปิ เจเตติ, อุภย\-พฺยาพาธายปิ เจเตติ, เจตสิกมฺปิ ทุกฺขํ โทมนสฺสํ ปฏิสํเวเทติ. โมเห ปหีเน เนว\-ตฺตพฺยาพาธายปิ เจเตติ, น ปรพฺยาพาธายปิ เจเตติ, น อุภย\-พฺยาพาธายปิ เจเตติ, น เจตสิกํ\csromansharedfootnote{351abaa1618f4a18}{น เจตสิกํ (สี, สฺยา, ก)} ทุกฺขํ โทมนสฺสํ}

\vspace{\tipitakaparskip}
\csromancenter{(6) 1. พฺราหฺมณวคฺค ปฏิสํเวเทติ. เอวํ โข พฺราหฺมณ สนฺทิฏฺฐิโก ธมฺโม โหติ อกาลิโก เอหิปสฺสิโก โอปเนยฺยิโก ปจฺจตฺตํ เวทิตพฺโพ วิญฺญูหีติ.}
\prose{\csromanfolio{157}อภิกฺกนฺตํ โภ โคตม, อภิกฺกนฺตํ โภ โคตม. เสยฺยถาปิ โภ โคตม นิกฺกุชฺชิตํ วา อุกฺกุชฺเชยฺย, ปฏิจฺฉนฺนํ วา วิวเรยฺย, มูฬฺหสฺส วา มคฺคํ อาจิกฺเขยฺย, อนฺธกาเร วา เตลปชฺโชตํ ธาเรยฺย “จกฺขุมนฺโต รูปานิ ทกฺขนฺตี”ติ. เอวเมวํ โภตา โคตเมน อเนก\-ปริยาเยน ธมฺโม ปกาสิโต, เอสาหํ ภวนฺตํ โคตมํ สรณํ คจฺฉามิ ธมฺมญฺจ ภิกฺขุสํฆญฺจ, อุปาสกํ มํ ภวํ โคตโม ธาเรตุ อชฺชตคฺเค ปาณุเปตํ สรณํ คตนฺติ.  ตติยํ.}

\csromanheader{2}{4. ปริพฺพาชก\-สุตฺต}
\proseitem{55}{อถ โข อญฺญตโร พฺราหฺมณ\-ปริพฺพาชโก เยน ภควา เตนุปสงฺกมิ, อุปสงฺกมิตฺวา ฯเปฯ เอกมนฺตํ นิสินฺโน โข โส พฺราหฺมณ\-ปริพฺพาชโก ภควนฺตํ เอตทโวจ “สนฺทิฏฺฐิโก ธมฺโม สนฺทิฏฺฐิโก ธมฺโมติ โภ โคตม วุจฺจติ, กิตฺตาวตา นุ โข โภ โคตม สนฺทิฏฺฐิโก ธมฺโม โหติ อกาลิโก เอหิปสฺสิโก โอปเนยฺยิโก ปจฺจตฺตํ เวทิตพฺโพ วิญฺญูหี”ติ.}

\prose{รตฺโต โข พฺราหฺมณ ราเคน อภิภูโต ปริยาทินฺน\-จิตฺโต อตฺตพฺยาพาธายปิ เจเตติ, ปรพฺยาพาธายปิ เจเตติ, อุภย\-พฺยาพาธายปิ เจเตติ, เจตสิกมฺปิ ทุกฺขํ โทมนสฺสํ ปฏิสํเวเทติ. ราเค ปหีเน เนว\-ตฺตพฺยาพาธายปิ เจเตติ, น ปรพฺยาพาธายปิ เจเตติ, น อุภย\-พฺยาพาธายปิ เจเตติ, น เจตสิกมฺปิ ทุกฺขํ โทมนสฺสํ ปฏิสํเวเทติ.}

\prose{รตฺโต โข พฺราหฺมณ ราเคน อภิภูโต ปริยาทินฺน\-จิตฺโต กาเยน ทุจฺจริตํ จรติ, วาจาย ทุจฺจริตํ จรติ, มนสา ทุจฺจริตํ จรติ. ราเค ปหีเน เนว กาเยน ทุจฺจริตํ จรติ, น วาจาย ทุจฺจริตํ จรติ, น มนสา ทุจฺจริตํ จรติ.}

\prose{รตฺโต โข พฺราหฺมณ ราเคน อภิภูโต ปริยาทินฺน\-จิตฺโต อตฺตตฺถมฺปิ ยถาภูตํ นปฺปชานาติ, ปรตฺถมฺปิ ยถาภูตํ นปฺปชานาติ, อุภยตฺถมฺปิ ยถาภูตํ นปฺปชานาติ. ราเค ปหีเน อตฺตตฺถมฺปิ ยถาภูตํ \csromanfolio{158}ปชานาติ, ปรตฺถมฺปิ ยถาภูตํ ปชานาติ, อุภยตฺถมฺปิ ยถาภูตํ ปชานาติ. เอวมฺปิ โข พฺราหฺมณ สนฺทิฏฺฐิโก ธมฺโม โหติ ฯเปฯ.}

\prose{ทุฏฺโฐ โข พฺราหฺมณ โทเสน ฯเปฯ. มูโฬฺห โข พฺราหฺมณ โมเหน อภิภูโต ปริยาทินฺน\-จิตฺโต อตฺตพฺยาพาธายปิ เจเตติ, ปรพฺยาพาธายปิ เจเตติ, อุภย\-พฺยาพาธายปิ เจเตติ, เจตสิกมฺปิ ทุกฺขํ โทมนสฺสํ ปฏิสํเวเทติ. โมเห ปหีเน เนว\-ตฺตพฺยาพาธายปิ เจเตติ, น ปรพฺยาพาธายปิ เจเตติ, น อุภย\-พฺยาพาธายปิ เจเตติ, น เจตสิกํ ทุกฺขํ โทมนสฺสํ ปฏิสํเวเทติ.}

\prose{มูโฬฺห โข พฺราหฺมณ โมเหน อภิภูโต ปริยาทินฺน\-จิตฺโต กาเยน ทุจฺจริตํ จรติ, วาจาย ทุจฺจริตํ จรติ, มนสา ทุจฺจริตํ จรติ. โมเห ปหีเน เนว กาเยน ทุจฺจริตํ จรติ, น วาจาย ทุจฺจริตํ จรติ, น มนสา ทุจฺจริตํ จรติ.}

\prose{มูโฬฺห โข พฺราหฺมณ โมเหน อภิภูโต ปริยาทินฺน\-จิตฺโต อตฺตตฺถมฺปิ ยถาภูตํ นปฺปชานาติ, ปรตฺถมฺปิ ยถาภูตํ นปฺปชานาติ, อุภยตฺถมฺปิ ยถาภูตํ นปฺปชานาติ. โมเห ปหีเน อตฺตตฺถมฺปิ ยถาภูตํ ปชานาติ, ปรตฺถมฺปิ ยถาภูตํ ปชานาติ, อุภยตฺถมฺปิ ยถาภูตํ ปชานาติ. เอวํ โข พฺราหฺมณ สนฺทิฏฺฐิโก ธมฺโม โหติ อกาลิโก เอหิปสฺสิโก โอปเนยฺยิโก ปจฺจตฺตํ เวทิตพฺโพ วิญฺญูหีติ.}

\prose{อภิกฺกนฺตํ โภ โคตม ฯเปฯ อุปาสกํ มํ ภวํ โคตโม ธาเรตุ อชฺชตคฺเค ปาณุเปตํ สรณํ คตนฺติ.  จตุตฺถํ.}

\csromanheader{2}{5. \textbf{นิพฺภุตสุตฺต}}
\proseitem{56}{อถ โข ชาณุสฺโสณิ พฺราหฺมโณ เยน ภควา เตนุปสงฺกมิ, อุปสงฺกมิตฺวา ภควนฺตํ อภิวาเทตฺวา เอกมนฺตํ นิสีทิ, เอกมนฺตํ นิสินฺโน โข ชาณุสฺโสณิ พฺราหฺมโณ ภควนฺตํ เอตทโวจ “สนฺทิฏฺฐิกํ นิพฺพานํ สนฺทิฏฺฐิกํ นิพฺพานนฺติ โภ โคตม วุจฺจติ, กิตฺตาวตา นุ โข โภ โคตม สนฺทิฏฺฐิกํ นิพฺพานํ โหติ อกาลิกํ เอหิปสฺสิกํ โอปเนยฺยิกํ ปจฺจตฺตํ เวทิตพฺพํ วิญฺญูหี”ติ.}

\chapterheadpageread{(6) 1. พฺราหฺมณวคฺค}
\prose{\csromanfolio{159}รตฺโต โข พฺราหฺมณ ราเคน อภิภูโต ปริยาทินฺน\-จิตฺโต อตฺตพฺยาพาธายปิ เจเตติ, ปรพฺยาพาธายปิ เจเตติ, อุภย\-พฺยาพาธายปิ เจเตติ, เจตสิกมฺปิ ทุกฺขํ โทมนสฺสํ ปฏิสํเวเทติ. ราเค ปหีเน เนว\-ตฺตพฺยาพาธายปิ เจเตติ, น ปรพฺยาพาธายปิ เจเตติ, น อุภย\-พฺยาพาธายปิ เจเตติ, น เจตสิกมฺปิ ทุกฺขํ โทมนสฺสํ ปฏิสํเวเทติ. เอวมฺปิ พฺราหฺมณ สนฺทิฏฺฐิกํ นิพฺพานํ โหติ.}

\prose{ทุฏฺโฐ โข พฺราหฺมณ ฯเปฯ. มูโฬฺห โข พฺราหฺมณ โมเหน อภิภูโต ปริยาทินฺน\-จิตฺโต อตฺตพฺยาพาธายปิ เจเตติ, ปรพฺยาธายปิ เจเตติ, อุภย\-พฺยาพาธายปิ เจเตติ, เจตสิกมฺปิ ทุกฺขํ โทมนสฺสํ ปฏิสํเวเทติ. โมเห ปหีเน เนว\-ตฺตพฺยาพาธายปิ เจเตติ, น ปรพฺยาพาธายปิ เจเตติ, น อุภย\-พฺยาพาธายปิ เจเตติ, น เจตสิกํ ทุกฺขํ โทมนสฺสํ ปฏิสํเวเทติ. เอวมฺปิ โข พฺราหฺมณ สนฺทิฏฺฐิกํ นิพฺพานํ โหติ,}

\prose{ยโต โข อยํ พฺราหฺมณ\csromansharedfootnote{413ab1ae1194c22f}{ยโต จ โข อยํ พฺราหฺมณ (สี), ยโต โข พฺราหฺมณ อกาลิกํ เอหิปสฺสิกํ โอปเนยฺยิกํ ปจฺจตฺตํ เวทิตพฺพํ (ก)} อนวเสสํ ราคกฺขยํ ปฏิสํเวเทติ, อนวเสสํ โทสกฺขยํ ปฏิสํเวเทติ, อนวเสสํ โมหกฺขยํ ปฏิสํเวเทติ. เอวํ โข พฺราหฺมณ สนฺทิฏฺฐิกํ นิพฺพานํ โหติ อกาลิกํ เอหิปสฺสิกํ โอปเนยฺยิกํ ปจฺจตฺตํ เวทิตพฺพํ วิญฺญูหีติ.}

\prose{อภิกฺกนฺตํ โภ โคตม ฯเปฯ อุปาสกํ มํ ภวํ โคตโม ธาเรตุ อชฺชตคฺเค ปาณุเปตํ สรณํ คตนฺติ.  ปญฺจมํ.}

\csromanheader{2}{6. \textbf{ปโลกสุตฺต}}
\proseitem{57}{อถ โข อญฺญตโร พฺราหฺมณ\-มหาสาโล เยน ภควา เตนุปสงฺกมิ ฯเปฯ เอกมนฺตํ นิสินฺโน โข โส พฺราหฺมณ\-มหาสาโล ภควนฺตํ เอตทโวจ “สุตํ เมตํ โภ โคตม ปุพฺพกานํ พฺราหฺมณานํ วุทฺธานํ มหลฺลกานํ อาจริย\-ปาจริยานํ ภาสมานานํ ‘ปุพฺเพ สุทํ\csromansharedfootnote{6eff657447e5ab94}{ปุพฺพสฺสุทํ (สี, สฺยา, กํ, อิ)} อยํ โลโก อวีจิ มญฺเญ ผุโฏ อโหสิ มนุสฺเสหิ กุกฺกุฏ\-สํปาติกา คามนิคม\-ราชธานิโย’ติ, โก นุ โภ โคตม เหตุ โก ปจฺจโย, เยเนตรหิ มนุสฺสานํ ขโย โหติ ตนุตฺตํ ปญฺญายติ, \csromanfolio{160}คามาปิ อคามา โหนฺติ, นิคมาปิ อนิคมา โหนฺติ, นคราปิ อนครา โหนฺติ, ชนปทาปิ อชนปทา โหนฺตี”ติ.}

\prose{เอตรหิ พฺราหฺมณ มนุสฺสา อธมฺมราค\-รตฺตา วิสมโลภา\-ภิภูตา มิจฺฉาธมฺม\-ปเรตา. เต อธมฺมราค\-รตฺตา วิสมโลภา\-ภิภูตา มิจฺฉาธมฺม\-ปเรตา ติณฺหานิ สตฺถานิ คเหตฺวา อญฺญมญฺญํ\csromansharedfootnote{cf12ff2181f32cb7}{อญฺญมญฺญสฺส (สพฺพตฺถ)} ชีวิตา โวโรเปนฺติ, เตน พหู มนุสฺสา กาลํ กโรนฺติ. อยมฺปิ โข พฺราหฺมณ เหตุ อยํ ปจฺจโย, เยเนตรหิ มนุสฺสานํ ขโย โหติ ตนุตฺตํ ปญฺญายติ, คามาปิ อคามา โหนฺติ, นิคมาปิ อนิคมา โหนฺติ, นคราปิ อนครา โหนฺติ, ชนปทาปิ อชนปทา โหนฺติ.}

\prose{ปุน จปรํ พฺราหฺมณ เอตรหิ มนุสฺสา อธมฺมราค\-รตฺตา วิสมโลภา\-ภิภูตา มิจฺฉาธมฺม\-ปเรตา. เตสํ อธมฺมราค\-รตฺตานํ วิสมโลภา\-ภิภูตานํ มิจฺฉาธมฺม\-ปเรตานํ เทโว น สมฺมาธารํ อนุปฺปเวจฺฉติ, เตน ทุพฺภิกฺขํ โหติ ทุสฺสสฺสํ เสตฏฺฐิกํ สลากาวุตฺตํ, เตน พหู มนุสฺสา กาลํ กโรนฺติ. อยมฺปิ โข พฺราหฺมณ เหตุ อยํ ปจฺจโย, เยเนตรหิ มนุสฺสานํ ขโย โหติ ตนุตฺตํ ปญฺญายติ, คามาปิ อคามา โหนฺติ, นิคมาปิ อนิคมา โหนฺติ, นคราปิ อนครา โหนฺติ, ชนปทาปิ อชนปทา โหนฺติ.}

\prose{ปุน จปรํ พฺราหฺมณ เอตรหิ มนุสฺสา อธมฺมราค\-รตฺตา วิสมโลภา\-ภิภูตา มิจฺฉาธมฺม\-ปเรตา. เตสํ อธมฺมราค\-รตฺตานํ วิสมโลภา\-ภิภูตานํ มิจฺฉาธมฺม\-ปเรตานํ ยกฺขา วาเฬ อมนุสฺเส โอสฺสชฺชนฺติ\csromansharedfootnote{09ec393868053cca}{โอสฺสชนฺติ (สี)}, เตน พหู มนุสฺสา กาลํ กโรนฺติ. อยมฺปิ โข พฺราหฺมณ เหตุ อยํ ปจฺจโย, เยเนตรหิ มนุสฺสานํ ขโย โหติ ตนุตฺตํ ปญฺญายติ, คามาปิ อคามา โหนฺติ, นิคมาปิ อนิคมา โหนฺติ, นคราปิ อนครา โหนฺติ, ชนปทาปิ อชนปทา โหนฺตีติ.}

\prose{อภิกฺกนฺตํ โภ โคตม ฯเปฯ อุปาสกํ มํ ภวํ โคตโม ธาเรตุ อชฺชตคฺเค ปาณุเปตํ สรณํ คตนฺติ.  ฉฏฺฐํ.}

\chapterheadpageread{(6) 1. พฺราหฺมณวคฺค 7. วจฺฉโคตฺต\-สุตฺต}
\proseitem{58}{\csromanfolio{161}อถ โข วจฺฉโคตฺโต\csromansharedfootnote{60e7e8e2f389d381}{วจฺฉปุตฺโต (ก)} ปริพฺพาชโก เยน ภควา เตนุปสงฺกมิ, อุปสงฺกมิตฺวา ภควตา สทฺธิํ สมฺโมทิ, สมฺโมทนียํ กถํ สารณียํ วีติสาเรตฺวา เอกมนฺตํ นิสีทิ, เอกมนฺตํ นิสินฺโน โข วจฺฉโคตฺโต ปริพฺพาชโก ภควนฺตํ เอตทโวจ “สุตํ เมตํ โภ โคตม สมโณ โคตโม เอวมาห ‘มยฺหเมว ทานํ ทาตพฺพํ, นาญฺเญสํ ทานํ ทาตพฺพํ. มยฺหเมว สาวกานํ ทานํ ทาตพฺพํ, นาญฺเญสํ สาวกานํ ทานํ ทาตพฺพํ. มยฺหเมว ทินฺนํ มหปฺผลํ, นาญฺเญสํ ทินฺนํ มหปฺผลํ. มยฺหเมว สาวกานํ ทินฺนํ มหปฺผลํ, นาญฺเญสํ สาวกานํ ทินฺนํ มหปฺผลนฺ’ติ. เย เต โภ โคตม เอวมาหํสุ ‘สมโณ โคตโม เอวมาห มยฺหเมว ทานํ ทาตพฺพํ, นาญฺเญสํ ทานํ ทาตพฺพํ. มยฺหเมว สาวกานํ ทานํ ทาตพฺพํ, นาญฺเญสํ สาวกานํ ทานํ ทาตพฺพํ. มยฺหเมว ทินฺนํ มหปฺผลํ, นาญฺเญสํ ทินฺนํ มหปฺผลํ. มยฺหเมว สาวกานํ ทินฺนํ มหปฺผลํ, นาญฺเญสํ สาวกานํ ทินฺนํ มหปฺผลนฺ’ติ. กิจฺจิ เต โภโต โคตมสฺส วุตฺตวาทิโน, น จ ภวนฺตํ โคตมํ อภูเตน อพฺภาจิกฺขนฺติ, ธมฺมสฺส จานุธมฺมํ พฺยากโรนฺติ. น จ โกจิ สหธมฺมิโก วาทานุปาโต\csromansharedfootnote{3b5234c5bf819f99}{วาทานุวาโท (ก)} คารยฺหํ ฐานํ อาคจฺฉติ, อนพฺภกฺขาตุ\-กามา หิ มยํ ภวนฺตํ โคตมนฺ”ติ.}

\prose{เย เต วจฺฉ เอวมาหํสุ “สมโณ โคตโม เอวมาห ‘มยฺหเมว ทานํ ทาตพฺพํ ฯเปฯ นาญฺเญสํ สาวกานํ ทินฺนํ มหปฺผลนฺ’ติ”. น เม เต วุตฺตวาทิโน, อพฺภาจิกฺขนฺติ จ ปน มํ\csromansharedfootnote{15411179c6798272}{จ ปน มํ เต (สี, สฺยา, กํ, อิ)} อสตา อภูเตน. โย โข วจฺฉ ปรํ ทานํ ททนฺตํ วาเรติ, โส ติณฺณํ อนฺตรายกโร โหติ ติณฺณํ ปาริปนฺถิโก. กตเมสํ ติณฺณํ? ทายกสฺส ปุญฺญ\-นฺตรายกโร โหติ, ปฏิคฺคาหกานํ ลาภ\-นฺตรายกโร โหติ, ปุพฺเพว โข ปนสฺส อตฺตา ขโต จ โหติ อุปหโต จ. โย โข วจฺฉ ปรํ ทานํ ททนฺตํ วาเรติ, โส อิเมสํ ติณฺณํ อนฺตรายกโร โหติ ติณฺณํ ปาริปนฺถิโก.}

\prose{อหํ โข ปน วจฺฉ เอวํ วทามิ\csromandash{}เยหิ เต จนฺทนิกาย วา โอลิคลฺเล วา ปาณา, ตตฺรปิ โย ถาลิโธวนํ\csromansharedfootnote{2ac611dc7334061f}{ถาลกโธวนํ (ก)} วา สราวโธวนํ วา ฉฑฺเฑติ \csromanfolio{162}“เย ตตฺถ ปาณา เต เตน ยาเปนฺตู”ติ, ตโตนิทานมฺปาหํ วจฺฉ ปุญฺญสฺส อาคมํ วทามิ, โก ปน วาโท มนุสฺสภูเต. อปิ จาหํ วจฺฉ สีลวโต ทินฺนํ มหปฺผลํ วทามิ, โน ตถา ทุสฺสีลสฺส. โส จ โหติ ปญฺจ\-งฺค\-วิปฺปหีโน ปญฺจ\-งฺค\-สมนฺนาคโต.}

\prose{กตมานิ ปญฺจงฺคานิ ปหีนานิ โหนฺติ? กามจฺฉนฺโท ปหีโน โหติ, พฺยาปาโท ปหีโน โหติ, ถินมิทฺธํ ปหีนํ โหติ, อุทฺธจฺจ\-กุกฺกุจฺจํ ปหีนํ โหติ, วิจิกิจฺฉา ปหีนา โหติ. อิมานิ ปญฺจงฺคานิ วิปฺปหีนานิ โหนฺติ.}

\prose{กตเมหิ ปญฺจหิ องฺเคหิ สมนฺนาคโต โหติ? อเสกฺเขน สีลกฺขนฺเธน สมนฺนาคโต โหติ, อเสกฺเขน สมาธิ\-กฺขนฺเธน สมนฺนาคโต โหติ, อเสกฺเขน ปญฺญา\-กฺขนฺเธน สมนฺนาคโต โหติ, อเสกฺเขน วิมุตฺติ\-กฺขนฺเธน สมนฺนาคโต โหติ, อเสกฺเขน วิมุตฺติ\-ญาณ\-ทสฺสนกฺขนฺเธน สมนฺนาคโต โหติ. อิเมหิ ปญฺจหิ องฺเคหิ สมนฺนาคโต โหติ. อิติ ปญฺจ\-งฺค\-วิปฺปหีเน ปญฺจ\-งฺค\-สมนฺนาคเต ทินฺนํ มหปฺผลนฺติ วทามีติ.}

\setlength{\csromangathapagewidth}{0pt}
\setlength{\csromangathawakwidth}{0pt}
\csromangathameasuregroup{อิติ กณฺหาสุ เสตาสุ, \\ กมฺมาสาสุ สรูปาสุ, \\ ยาสุ กาสุจิ เอตาสุ, \\ โธรโยฺห พลสมฺปนฺโน, \\ ตเมว ภาเร ยุญฺชนฺติ,}{\csromangathabat{อิติ กณฺหาสุ เสตาสุ,}{โรหิณีสุ หรีสุ วา.} \\ \csromangathabat{กมฺมาสาสุ สรูปาสุ,}{โคสุ ปาเรวตาสุ วา.} \\ \csromangathabat{ยาสุ กาสุจิ เอตาสุ,}{ทนฺโต ชายติ ปุงฺคโว.} \\ \csromangathabat{โธรโยฺห พลสมฺปนฺโน,}{กลฺยาณชวนิกฺกโม.} \\ \csromangathabat{ตเมว ภาเร ยุญฺชนฺติ,}{นาสฺส วณฺณํ ปริกฺขเร.}}
\csromangathasetleft{อิติ กณฺหาสุ เสตาสุ, \\ กมฺมาสาสุ สรูปาสุ, \\ ยาสุ กาสุจิ เอตาสุ, \\ โธรโยฺห พลสมฺปนฺโน, \\ ตเมว ภาเร ยุญฺชนฺติ,}
\csromangathagroup{\csromangathastanza{\csromangathabat{อิติ กณฺหาสุ เสตาสุ,}{โรหิณีสุ หรีสุ วา.} \\ \csromangathabat{กมฺมาสาสุ สรูปาสุ,}{โคสุ ปาเรวตาสุ วา.}}\csromangathastanza{\csromangathabat{ยาสุ กาสุจิ เอตาสุ,}{ทนฺโต ชายติ ปุงฺคโว.} \\ \csromangathabat{โธรโยฺห พลสมฺปนฺโน,}{กลฺยาณชวนิกฺกโม.} \\ \csromangathabat{ตเมว ภาเร ยุญฺชนฺติ,}{นาสฺส วณฺณํ ปริกฺขเร.}}}

\vspace{\tipitakaparskip}
\csromancenter{เอวเมวํ มนุสฺเสสุ ยสฺมิํ กสฺมิญฺจิ ชาติเย.}
\setlength{\csromangathapagewidth}{0pt}
\setlength{\csromangathawakwidth}{0pt}
\csromangathameasuregroup{ขตฺติเย พฺราหฺมเณ เวสฺเส, \\ ยาสุ กาสุจิ เอตาสุ, \\ ธมฺมฏฺโฐ สีลสมฺปนฺโน, \\ ปหีนชาติมรโณ, \\ ปนฺนภาโร วิสํยุตฺโต, \\ ปารคู สพฺพธมฺมานํ, \\ ตสฺมิํเยว\textsuperscript{1} วิรเช เขตฺเต,}{\csromangathabat{ขตฺติเย พฺราหฺมเณ เวสฺเส,}{สุทฺเท จณฺฑาลปุกฺกุเส.} \\ \csromangathabat{ยาสุ กาสุจิ เอตาสุ,}{ทนฺโต ชายติ สุพฺพโต.} \\ \csromangathabat{ธมฺมฏฺโฐ สีลสมฺปนฺโน,}{สจฺจวาที หิรีมโน.} \\ \csromangathabat{ปหีนชาติมรโณ,}{พฺรหฺมจริยสฺส เกวลี.} \\ \csromangathabat{ปนฺนภาโร วิสํยุตฺโต,}{กตกิจฺโจ อนาสโว.} \\ \csromangathabat{ปารคู สพฺพธมฺมานํ,}{อนุปาทาย นิพฺพุโต.} \\ \csromangathabat{ตสฺมิํเยว\textsuperscript{1} วิรเช เขตฺเต,}{วิปุลา โหติ ทกฺขิณา.}}
\csromangathasetleft{ขตฺติเย พฺราหฺมเณ เวสฺเส, \\ ยาสุ กาสุจิ เอตาสุ, \\ ธมฺมฏฺโฐ สีลสมฺปนฺโน, \\ ปหีนชาติมรโณ, \\ ปนฺนภาโร วิสํยุตฺโต, \\ ปารคู สพฺพธมฺมานํ, \\ ตสฺมิํเยว\textsuperscript{1} วิรเช เขตฺเต,}
\csromangathagroup{\csromangathastanza{\csromangathabat{ขตฺติเย พฺราหฺมเณ เวสฺเส,}{สุทฺเท จณฺฑาลปุกฺกุเส.} \\ \csromangathabat{ยาสุ กาสุจิ เอตาสุ,}{ทนฺโต ชายติ สุพฺพโต.}}\csromangathastanza{\csromangathabat{ธมฺมฏฺโฐ สีลสมฺปนฺโน,}{สจฺจวาที หิรีมโน.} \\ \csromangathabat{ปหีนชาติมรโณ,}{พฺรหฺม\-จริยสฺส เกวลี.}}\csromangathastanza{\csromangathabat{ปนฺนภาโร วิสํยุตฺโต,}{กตกิจฺโจ อนาสโว.} \\ \csromangathabat{ปารคู สพฺพธมฺมานํ,}{อนุปาทาย นิพฺพุโต.} \\ \csromangathabat{ตสฺมิํเยว\csromansharedfootnote{72ac9be6bc8673ab}{ตสฺมิํ เว (สฺยา, กํ)} วิรเช เขตฺเต,}{วิปุลา โหติ ทกฺขิณา.}}}

\chapterheadpageread{(6) 1. พฺราหฺมณวคฺค}
\setlength{\csromangathapagewidth}{0pt}
\setlength{\csromangathawakwidth}{0pt}
\csromangathameasuregroup{พาลา จ อวิชานนฺตา, \\ พหิทฺธา เทนฺติ ทานานิ, \\ เย จ สนฺเต อุปาสนฺติ, \\ สทฺธา จ เนสํ สุคเต, \\ เทวโลกญฺจ เต ยนฺติ, \\ อนุปุพฺเพน นิพฺพานํ,}{\csromangathabat{พาลา จ อวิชานนฺตา,}{ทุมฺเมธา อสฺสุตาวิโน.} \\ \csromangathabat{พหิทฺธา เทนฺติ ทานานิ,}{น หิ สนฺเต อุปาสเร.} \\ \csromangathabat{เย จ สนฺเต อุปาสนฺติ,}{สปฺปญฺเญ ธีรสมฺมเต.} \\ \csromangathabat{สทฺธา จ เนสํ สุคเต,}{มูลชาตา ปติฏฺฐิตา.} \\ \csromangathabat{เทวโลกญฺจ เต ยนฺติ,}{กุเล วา อิธ ชายเร.} \\ \csromangathabat{อนุปุพฺเพน นิพฺพานํ,}{อธิคจฺฉนฺติ ปณฺฑิตาติ.สตฺตมํ.}}
\csromangathasetleft{พาลา จ อวิชานนฺตา, \\ พหิทฺธา เทนฺติ ทานานิ, \\ เย จ สนฺเต อุปาสนฺติ, \\ สทฺธา จ เนสํ สุคเต, \\ เทวโลกญฺจ เต ยนฺติ, \\ อนุปุพฺเพน นิพฺพานํ,}
\csromangathagroup{\csromangathastanza{\csromanfolio{163}\csromangathabat{พาลา จ อวิชานนฺตา,}{ทุมฺเมธา อสฺสุตาวิโน.} \\ \csromangathabat{พหิทฺธา เทนฺติ ทานานิ,}{น หิ สนฺเต อุปาสเร.}}\csromangathastanza{\csromangathabat{เย จ สนฺเต อุปาสนฺติ,}{สปฺปญฺเญ ธีรสมฺมเต.} \\ \csromangathabat{สทฺธา จ เนสํ สุคเต,}{มูลชาตา ปติฏฺฐิตา.}}\csromangathastanza{\csromangathabat{เทวโลกญฺจ เต ยนฺติ,}{กุเล วา อิธ ชายเร.} \\ \csromangathabat{อนุปุพฺเพน นิพฺพานํ,}{อธิคจฺฉนฺติ ปณฺฑิตาติ.  สตฺตมํ.}}}

\csromanheader{2}{8. \textbf{ติกณฺณสุตฺต}}
\proseitem{59}{อถ โข ติกณฺโณ พฺราหฺมโณ เยน ภควา เตนุปสงฺกมิ, อุปสงฺกมิตฺวา ภควตา สทฺธิํ ฯเปฯ เอกมนฺตํ นิสินฺโน โข ติกณฺโณ พฺราหฺมโณ ภควโต สมฺมุขา เตวิชฺชานํ สุทํ พฺราหฺมณานํ วณฺณํ ภาสติ “เอวมฺปิ เตวิชฺชา พฺราหฺมณา, อิติปิ เตวิชฺชา พฺราหฺมณา”ติ.}

\prose{ยถา กถํ ปน พฺราหฺมณ พฺราหฺมณา พฺราหฺมณํ เตวิชฺชํ ปญฺญาเปนฺตีติ? อิธ โภ โคตม พฺราหฺมโณ อุภโต สุชาโต โหติ มาติโต จ ปิติโต จ สํสุทฺธ\-คหณิโก ยาว สตฺตมา ปิตามหยุคา อกฺขิตฺโต อนุปกฺกุฏฺโฐ ชาติวาเทน, อชฺฌายโก มนฺตธโร ติณฺณํ เวทานํ ปารคู สนิฆณฺฑุ\-เกฏุภานํ สากฺขร\-ปฺปเภทานํ อิติหาส\-ปญฺจมานํ ปทโก เวยฺยากรโณ โลกายต\-มหา\-ปุริส\-ลกฺขเณสุ อนวโยติ. เอวํ โข โภ โคตม พฺราหฺมณา พฺราหฺมณํ เตวิชฺชํ ปญฺญาเปนฺตีติ.}

\prose{อญฺญถา โข พฺราหฺมณ พฺราหฺมณา พฺราหฺมณํ เตวิชฺชํ ปญฺญาเปนฺติ, อญฺญถา จ ปน อริยสฺส วินเย เตวิชฺโช โหตีติ. ยถา กถํ ปน โภ โคตม อริยสฺส วินเย เตวิชฺโช โหติ? สาธุ เม ภวํ โคตโม ตถา ธมฺมํ เทเสตุ ยถา อริยสฺส วินเย เตวิชฺโช โหตีติ. เตน หิ พฺราหฺมณ สุณาหิ สาธุกํ มนสิ กโรหิ, ภาสิสฺสามีติ. “เอวํ โภ”ติ โข ติกณฺโณ พฺราหฺมโณ ภควโต ปจฺจสฺโสสิ. ภควา เอตทโวจ\csromandash{}}

\prose{อิธ พฺราหฺมณ ภิกฺขุ วิวิจฺเจว กาเมหิ วิวิจฺจ อกุสเลหิ ธมฺเมหิ สวิตกฺกํ สวิจารํ วิเวกชํ ปีติสุขํ ปฐมํ ฌานํ อุปสมฺปชฺช วิหรติ, วิตกฺก\-วิจารานํ วูปสมา อชฺฌตฺตํ สมฺปสาทนํ เจตโส เอโกทิภาวํ อวิตกฺกํ อวิจารํ สมาธิชํ ปีติสุขํ ทุติยํ ฌานํ อุปสมฺปชฺช วิหรติ, \csromanfolio{164}ปีติยา จ วิราคา อุเปกฺขโก จ วิหรติ สโต จ สมฺปชาโน, สุขญฺจ กาเยน ปฏิสํเวเทติ, ยํ ตํ อริยา อาจิกฺขนฺติ “อุเปกฺขโก สติมา สุขวิหารี”ติ ตติยํ ฌานํ อุปสมฺปชฺช วิหรติ, สุขสฺส จ ปหานา ทุกฺขสฺส จ ปหานา ปุพฺเพว โสมนสฺส\-โทมนสฺสานํ อตฺถงฺคมา อทุกฺขมสุขํ อุเปกฺขา\-สติ\-ปาริสุทฺธิํ จตุตฺถํ ฌานํ อุปสมฺปชฺช วิหรติ.}

\prose{โส เอวํ สมาหิเต จิตฺเต ปริสุทฺเธ ปริโยทาเต อนงฺคเณ วิคตู\-ปกฺกิเลเส มุทุภูเต กมฺมนิเย ฐิเต อาเนญฺชปฺปตฺเต ปุพฺเพ\-นิวาสา\-นุสฺสติ\-ญาณาย จิตฺตํ อภินินฺนาเมติ. โส อเนกวิหิตํ ปุพฺเพนิวาสํ อนุสฺสรติ. เสยฺยถิทํ? เอกมฺปิ ชาติํ เทฺวปิ ชาติโย ติสฺโสปิ ชาติโย จตสฺโสปิ ชาติโย ปญฺจปิ ชาติโย ทสปิ ชาติโย วีสมฺปิ ชาติโย ติํสมฺปิ ชาติโย จตฺตารีสมฺปิ ชาติโย ปญฺญาสมฺปิ ชาติโย ชาติสตมฺปิ ชาติสหสฺสมฺปิ ชาติ\-สต\-สหสฺสมฺปิ อเนเกปิ สํวฏฺฏกปฺเป อเนเกปิ วิวฏฺฏกปฺเป อเนเกปิ สํวฏฺฏ\-วิวฏฺฏ\-กปฺเป “อมุตฺราสิํ เอวํนาโม เอวํโคตฺโต เอวํวณฺโณ เอวมาหาโร เอวํ\-สุขทุกฺข\-ปฺปฏิสํเวที เอวมา\-ยุปริยนฺโต, โส ตโต จุโต อมุตฺร อุทปาทิํ. ตตฺราปาสิํ เอวํนาโม เอวํโคตฺโต เอวํวณฺโณ เอวมาหาโร เอวํ\-สุขทุกฺข\-ปฺปฏิสํเวที เอวมา\-ยุปริยนฺโต, โส ตโต จุโต อิธูปปนฺโน”ติ. อิติ สาการํ สอุทฺเทสํ อเนกวิหิตํ ปุพฺเพนิวาสํ อนุสฺสรติ. อยมสฺส ปฐมา วิชฺชา อธิคตา โหติ, อวิชฺชา วิหตา, วิชฺชา อุปฺปนฺนา, ตโม วิหโต, อาโลโก อุปฺปนฺโน, ยถา ตํ อปฺปมตฺตสฺส อาตาปิโน ปหิตตฺตสฺส วิหรโต.}

\prose{โส เอวํ สมาหิเต จิตฺเต ปริสุทฺเธ ปริโยทาเต อนงฺคเณ วิคตู\-ปกฺกิเลเส มุทุภูเต กมฺมนิเย ฐิเต อาเนญฺชปฺปตฺเต สตฺตานํ จุตูปปาต\-ญาณาย จิตฺตํ อภินินฺนาเมติ. โส ทิพฺเพน จกฺขุนา วิสุทฺเธน อติกฺกนฺต\-มานุสเกน สตฺเต ปสฺสติ จวมาเน อุปปชฺชมาเน หีเน ปณีเต สุวณฺเณ ทุพฺพณฺเณ สุคเต ทุคฺคเต, ยถากมฺมูปเค สตฺเต ปชานาติ “อิเม วต โภนฺโต สตฺตา กาย\-ทุจฺจริเตน สมนฺนาคตา ฯเปฯ มโน\-ทุจฺจริเตน สมนฺนาคตา อริยานํ อุปวาทกา มิจฺฉาทิฏฺฐิกา มิจฺฉาทิฏฺฐิ\-กมฺม\-สมาทานา, เต กายสฺส เภทา ปรํ มรณา อปายํ ทุคฺคติํ วินิปาตํ นิรยํ อุปปนฺนา. อิเม วา ปน โภนฺโต สตฺตา กายสุจริเตน สมนฺนาคตา วจีสุจริเตน สมนฺนาคตา มโนสุจริเตน}

\vspace{\tipitakaparskip}
\csromancenter{(6) 1. พฺราหฺมณวคฺค สมนฺนาคตา อริยานํ อนุปวาทกา สมฺมาทิฏฺฐิกา สมฺมา\-ทิฏฺฐิ\-กมฺม\-สมาทานา, เต กายสฺส เภทา ปรํ มรณา สุคติํ สคฺคํ โลกํ อุปปนฺนา”ติ. อิติ ทิพฺเพน จกฺขุนา วิสุทฺเธน อติกฺกนฺต\-มานุสเกน สตฺเต ปสฺสติ จวมาเน อุปปชฺชมาเน หีเน ปณีเต สุวณฺเณ ทุพฺพณฺเณ สุคเต ทุคฺคเต, ยถากมฺมูปเค สตฺเต ปชานาติ. อยมสฺส ทุติยา วิชฺชา อธิคตา โหติ, อวิชฺชา วิหตา, วิชฺชา อุปฺปนฺนา, ตโม วิหโต, อาโลโก อุปฺปนฺโน, ยถา ตํ อปฺปมตฺตสฺส อาตาปิโน ปหิตตฺตสฺส วิหรโต.}
\prose{\csromanfolio{165}โส เอวํ สมาหิเต จิตฺเต ปริสุทฺเธ ปริโยทาเต อนงฺคเณ วิคตู\-ปกฺกิเลเส มุทุภูเต กมฺมนิเย ฐิเต อาเนญฺชปฺปตฺเต อาสวานํ ขยญาณาย จิตฺตํ อภินินฺนาเมติ, โส “อิทํ ทุกฺขนฺ”ติ ยถาภูตํ ปชานาติ, “อยํ ทุกฺขสมุทโย”ติ ยถาภูตํ ปชานาติ, “อยํ ทุกฺขนิโรโธ”ติ ยถาภูตํ ปชานาติ, “อยํ ทุกฺขนิโรธ\-คามินี ปฏิปทา”ติ ยถาภูตํ ปชานาติ. “อิเม อาสวา”ติ ยถาภูตํ ปชานาติ, “อยํ อาสวสมุทโย”ติ ยถาภูตํ ปชานาติ, “อยํ อาสวนิโรโธ”ติ ยถาภูตํ ปชานาติ, “อยํ อาสว\-นิโรธ\-คามินี ปฏิปทา”ติ ยถาภูตํ ปชานาติ. ตสฺส เอวํ ชานโต เอวํ ปสฺสโต กามาสวาปิ จิตฺตํ วิมุจฺจติ, ภวาสวาปิ จิตฺตํ วิมุจฺจติ, อวิชฺชาสวาปิ จิตฺตํ วิมุจฺจติ, วิมุตฺตสฺมิํ “วิมุตฺตมฺ”อิติ ญาณํ โหติ, “ขีณา ชาติ, วุสิตํ พฺรหฺมจริยํ, กตํ กรณียํ, นาปรํ อิตฺถตฺตายา”ติ ปชานาติ. อยมสฺส ตติยา วิชฺชา อธิคตา โหติ, อวิชฺชา วิหตา, วิชฺชา อุปฺปนฺนา, ตโม วิหโต, อาโลโก อุปฺปนฺโน, ยถา ตํ อปฺปมตฺตสฺส อาตาปิโน ปหิตตฺตสฺส วิหรโตติ.}

\setlength{\csromangathapagewidth}{0pt}
\setlength{\csromangathawakwidth}{0pt}
\csromangathameasuregroup{อนุจฺจาวจสีลสฺส, \\ จิตฺตํ ยสฺส วสีภูตํ, \\ ตํ เว ตโมนุทํ ธีรํ, \\ หิตํ เทวมนุสฺสานํ, \\ ตีหิ วิชฺชาหิ สมฺปนฺนํ, \\ พุทฺธํ อนฺติมเทหินํ\textsuperscript{1}, \\ \textsuperscript{*}ปุพฺเพนิวาสํ โย เวที, \\ อโถ ชาติกฺขยํ ปตฺโต, \\ เอตาหิ ตีหิ วิชฺชาหิ, \\ ตมหํ วทามิ เตวิชฺชํ,}{\csromangathabat{อนุจฺจาวจสีลสฺส,}{นิปกสฺส จ ฌายิโน.} \\ \csromangathabat{จิตฺตํ ยสฺส วสีภูตํ,}{เอกคฺคํ สุสมาหิตํ.} \\ \csromangathabat{ตํ เว ตโมนุทํ ธีรํ,}{เตวิชฺชํ มจฺจุหายินํ.} \\ \csromangathabat{หิตํ เทวมนุสฺสานํ,}{อาหุ สพฺพปฺปหายินํ.} \\ \csromangathabat{ตีหิ วิชฺชาหิ สมฺปนฺนํ,}{อสมฺมูฬฺหวิหารินํ.} \\ \csromangathabat{พุทฺธํ อนฺติมเทหินํ\textsuperscript{1},}{ตํ นมสฺสนฺติ โคตมํ.} \\ \csromangathabat{\textsuperscript{*}ปุพฺเพนิวาสํ โย เวที,}{สคฺคาปายญฺจ ปสฺสติ.} \\ \csromangathabat{อโถ ชาติกฺขยํ ปตฺโต,}{อภิญฺญาโวสิโต มุนิ.} \\ \csromangathabat{เอตาหิ ตีหิ วิชฺชาหิ,}{เตวิชฺโช โหติ พฺราหฺมโณ.} \\ \csromangathabat{ตมหํ วทามิ เตวิชฺชํ,}{นาญฺญํ ลปิตลาปนนฺติ.}}
\csromangathasetleft{อนุจฺจาวจสีลสฺส, \\ จิตฺตํ ยสฺส วสีภูตํ, \\ ตํ เว ตโมนุทํ ธีรํ, \\ หิตํ เทวมนุสฺสานํ, \\ ตีหิ วิชฺชาหิ สมฺปนฺนํ, \\ พุทฺธํ อนฺติมเทหินํ\textsuperscript{1}, \\ \textsuperscript{*}ปุพฺเพนิวาสํ โย เวที, \\ อโถ ชาติกฺขยํ ปตฺโต, \\ เอตาหิ ตีหิ วิชฺชาหิ, \\ ตมหํ วทามิ เตวิชฺชํ,}
\csromangathagroup{\csromangathastanza{\csromangathabat{อนุจฺจาวจสีลสฺส,}{นิปกสฺส จ ฌายิโน.} \\ \csromangathabat{จิตฺตํ ยสฺส วสีภูตํ,}{เอกคฺคํ สุสมาหิตํ.}}\csromangathastanza{\csromangathabat{ตํ เว ตโมนุทํ ธีรํ,}{เตวิชฺชํ มจฺจุหายินํ.} \\ \csromangathabat{หิตํ เทวมนุสฺสานํ,}{อาหุ สพฺพ\-ปฺปหายินํ.}}\csromangathastanza{\csromangathabat{ตีหิ วิชฺชาหิ สมฺปนฺนํ,}{อสมฺมูฬฺหวิหารินํ.} \\ \csromangathabat{พุทฺธํ อนฺติมเทหินํ\csromansharedfootnote{d20f849c4e1e6a93}{อนฺติมสารีรํ (สี, สฺยา, กํ, อิ)},}{ตํ นมสฺสนฺติ โคตมํ.}}\csromangathastanza{\csromanfolio{166}\csromangathabat{\csromansymbolfootnote{*}{ขุ 1. 74, 262 ปิฏฺเฐสุปิ.}ปุพฺเพนิวาสํ โย เวที,}{สคฺคาปายญฺจ ปสฺสติ.} \\ \csromangathabat{อโถ ชาติกฺขยํ ปตฺโต,}{อภิญฺญาโวสิโต มุนิ.}}\csromangathastanza{\csromangathabat{เอตาหิ ตีหิ วิชฺชาหิ,}{เตวิชฺโช โหติ พฺราหฺมโณ.} \\ \csromangathabat{ตมหํ วทามิ เตวิชฺชํ,}{นาญฺญํ ลปิตลาปนนฺติ.}}}

\prose{เอวํ โข พฺราหฺมณ อริยสฺส วินเย เตวิชฺโช โหตีติ. อญฺญาถา โภ โคตม พฺราหฺมณานํ เตวิชฺโช, อญฺญถา จ ปน อริยสฺส วินเย เตวิชฺโช โหติ. อิมสฺส จ ปน โภ โคตม อริยสฺส วินเย เตวิชฺชสฺส พฺราหฺมณานํ เตวิชฺโช กลํ นาคฺฆติ โสฬสิํ.}

\prose{อภิกฺกนฺตํ โภ โคตม ฯเปฯ อุปาสกํ มํ ภวํ โคตโม ธาเรตุ อชฺชตคฺเค ปาณุเปตํ สรณํ คตนฺติ.  อฏฺฐมํ.}

\csromanheader{2}{9. \textbf{ชาณุสฺโสณิสุตฺต}}
\proseitem{60}{อถ โข ชาณุสฺโสณิ พฺราหฺมโณ เยน ภควา เตนุปสงฺกมิ, อุปสงฺกมิตฺวา ภควตา สทฺธิํ ฯเปฯ เอกมนฺตํ นิสินฺโน โข ชาณุสฺโสณิ พฺราหฺมโณ ภควนฺตํ เอตทโวจ “ยสฺสสฺสุ โภ โคตม ยญฺโญ วา สทฺธํ วา ถาลิปาโก วา เทยฺยธมฺมํ วา เตวิชฺเชสุ พฺราหฺมเณสุ ทานํ ทเทยฺยา”ติ. ยถา กถํ ปน พฺราหฺมณ พฺราหฺมณา เตวิชฺชํ ปญฺญเปนฺตีติ. อิธ โข โภ โคตม พฺราหฺมโณ อุภโต สุชาโต โหติ มาติโต จ ปิติโต จ สํสุทฺธ\-คหณิโก ยาว สตฺตมา ปิตามหยุคา อกฺขิตฺโต อนุปกฺกุฏฺโฐ ชาติวาเทน, อชฺฌายโก มนฺตธโร ติณฺณํ เวทานํ ปารคู สนิฆณฺฑุ\-เกฏุภานํ สากฺขร\-ปฺปเภทานํ อิติหาส\-ปญฺจมานํ ปทโก เวยฺยากรโณ โลกายต\-มหา\-ปุริส\-ลกฺขเณสุ อนวโยติ. เอวํ โข โภ โคตม พฺราหฺมณา เตวิชฺชํ ปญฺญเปนฺตีติ.}

\prose{อญฺญถา โข พฺราหฺมณ พฺราหฺมณา พฺราหฺมณํ เตวิชฺชํ ปญฺญเปนฺติ, อญฺญถา จ ปน อริยสฺส วินเย เตวิชฺโช โหตีติ. ยถา กถํ ปน โภ โคตม อริยสฺส วินเย เตวิชฺโช โหติ? สาธุ เม ภวํ โคตโม ตถา ธมฺมํ เทเสตุ, ยถา อริยสฺส วินเย เตวิชฺโช โหตีติ. เตน หิ พฺราหฺมณ สุณาหิ สาธุกํ มนสิ กโรหิ,}

\vspace{\tipitakaparskip}
\csromancenter{(6) 1. พฺราหฺมณวคฺค ภาสิสฺสามีติ. “เอวํ โภ”ติ โข ชาณุสฺโสณิ พฺราหฺมโณ ภควโต ปจฺจสฺโสสิ. ภควา เอตทโวจ\csromandash{}}
\prose{\csromanfolio{167}อิธ ปน พฺราหฺมณ ภิกฺขุ วิวิจฺเจว กาเมหิ ฯเปฯ จตุตฺถํ ฌานํ อุปสมฺปชฺช วิหรติ.}

\prose{โส เอวํ สมาหิเต จิตฺเต ปริสุทฺเธ ปริโยทาเต อนงฺคเณ วิคตู\-ปกฺกิเลเส มุทุภูเต กมฺมนิเย ฐิเต อาเนญฺชปฺปตฺเต ปุพฺเพ\-นิวาสา\-นุสฺสติ\-ญาณาย จิตฺตํ อภินินฺนาเมติ. โส อเนกวิหิตํ ปุพฺเพนิวาสํ อนุสฺสรติ. เสยฺยถิทํ? เอกมฺปิ ชาติํ เทฺวปิ ชาติโย ฯเปฯ. อิติ สาการํ สอุทฺเทสํ อเนกวิหิตํ ปุพฺเพนิวาสํ อนุสฺสรติ. อยมสฺส ปฐมา วิชฺชา อธิคตา โหติ, อวิชฺชา วิหตา, วิชฺชา อุปฺปนฺนา, ตโม วิหโต, อาโลโก อุปฺปนฺโน, ยถา ตํ อปฺปมตฺตสฺส อาตาปิโน ปหิตตฺตสฺส วิหรโต.}

\prose{โส เอวํ สมาหิเต จิตฺเต ปริสุทฺเธ ปริโยทาเต อนงฺคเณ วิคตู\-ปกฺกิเลเส มุทุภูเต กมฺมนิเย ฐิเต อาเนญฺชปฺปตฺเต สตฺตานํ จุตูปปาต\-ญาณาย จิตฺตํ อภินินฺนาเมติ. โส ทิพฺเพน จกฺขุนา วิสุทฺเธน อติกฺกนฺต\-มานุสเกน ฯเปฯ ยถากมฺมูปเค สตฺเต ปชานาติ. อยมสฺส ทุติยา วิชฺชา อธิคตา โหติ, อวิชฺชา วิหตา, วิชฺชา อุปฺปนฺนา, ตโม วิหโต, อาโลโก อุปฺปนฺโน, ยถา ตํ อปฺปมตฺตสฺส อาตาปิโน ปหิตตฺตสฺส วิหรโต.}

\prose{โส เอวํ สมาหิเต จิตฺเต ปริสุทฺเธ ปริโยทาเต อนงฺคเณ วิคตู\-ปกฺกิเลเส มุทุภูเต กมฺมนิเย ฐิเต อาเนญฺชปฺปตฺเต อาสวานํ ขยญาณาย จิตฺตํ อภินินฺนาเมติ, โส “อิทํ ทุกฺขนฺ”ติ ยถาภูตํ ปชานาติ ฯเปฯ “อยํ ทุกฺขนิโรธ\-คามินี ปฏิปทา”ติ ยถาภูตํ ปชานาติ. “อิเม อาสวา”ติ ยถาภูตํ ปชานาติ ฯเปฯ “อยํ อาสว\-นิโรธ\-คามินี ปฏิปทา”ติ ยถาภูตํ ปชานาติ. ตสฺส เอวํ ชานโต เอวํ ปสฺสโต กามาสวาปิ จิตฺตํ วิมุจฺจติ, ภวาสวาปิ จิตฺตํ วิมุจฺจติ, อวิชฺชาสวาปิ จิตฺตํ วิมุจฺจติ, วิมุตฺตสฺมิํ “วิมุตฺตมฺ”อิติ ญาณํ โหติ, “ขีณา ชาติ, วุสิตํ พฺรหฺมจริยํ, กตํ กรณียํ, นาปรํ อิตฺถตฺตายา”ติ ปชานาติ. อยมสฺส ตติยา วิชฺชา อธิคตา โหติ, อวิชฺชา วิหตา, วิชฺชา อุปฺปนฺนา, ตโม \csromanfolio{168}วิหโต, อาโลโก อุปฺปนฺโน, ยถา ตํ อปฺปมตฺตสฺส อาตาปิโน ปหิตตฺตสฺส วิหรโตติ.}

\setlength{\csromangathapagewidth}{0pt}
\setlength{\csromangathawakwidth}{0pt}
\csromangathameasuregroup{โย สีลพฺพตสมฺปนฺโน, \\ จิตฺตํ ยสฺส วสีภูตํ, \\ \textsuperscript{*}ปุพฺเพนิวาสํ โย เวที, \\ อโถ ชาติกฺขยํ ปตฺโต, \\ เอตาหิ ตีหิ วิชฺชาหิ, \\ ตมหํ วทามิ เตวิชฺชํ,}{\csromangathabat{โย สีลพฺพตสมฺปนฺโน,}{ปหิตตฺโต สมาหิโต.} \\ \csromangathabat{จิตฺตํ ยสฺส วสีภูตํ,}{เอกคฺคํ สุสมาหิตํ.} \\ \csromangathabat{\textsuperscript{*}ปุพฺเพนิวาสํ โย เวที,}{สคฺคาปายญฺจ ปสฺสติ.} \\ \csromangathabat{อโถ ชาติกฺขยํ ปตฺโต,}{อภิญฺญาโวสิโต มุนิ.} \\ \csromangathabat{เอตาหิ ตีหิ วิชฺชาหิ,}{เตวิชฺโช โหติ พฺราหฺมโณ.} \\ \csromangathabat{ตมหํ วทามิ เตวิชฺชํ,}{นาญฺญํ ลปิตลาปนนฺติ.}}
\csromangathasetleft{โย สีลพฺพตสมฺปนฺโน, \\ จิตฺตํ ยสฺส วสีภูตํ, \\ \textsuperscript{*}ปุพฺเพนิวาสํ โย เวที, \\ อโถ ชาติกฺขยํ ปตฺโต, \\ เอตาหิ ตีหิ วิชฺชาหิ, \\ ตมหํ วทามิ เตวิชฺชํ,}
\csromangathagroup{\csromangathastanza{\csromangathabat{โย สีลพฺพต\-สมฺปนฺโน,}{ปหิตตฺโต สมาหิโต.} \\ \csromangathabat{จิตฺตํ ยสฺส วสีภูตํ,}{เอกคฺคํ สุสมาหิตํ.}}\csromangathastanza{\csromangathabat{\csromansymbolfootnote{*}{ขุ 1. 74, 262 ปิฏฺเฐสุปิ.}ปุพฺเพนิวาสํ โย เวที,}{สคฺคาปายญฺจ ปสฺสติ.} \\ \csromangathabat{อโถ ชาติกฺขยํ ปตฺโต,}{อภิญฺญาโวสิโต มุนิ.}}\csromangathastanza{\csromangathabat{เอตาหิ ตีหิ วิชฺชาหิ,}{เตวิชฺโช โหติ พฺราหฺมโณ.} \\ \csromangathabat{ตมหํ วทามิ เตวิชฺชํ,}{นาญฺญํ ลปิตลาปนนฺติ.}}}

\prose{เอวํ โข พฺราหฺมณ อริยสฺส วินเย เตวิชฺโช โหตีติ. อญฺญถา โภ โคตม พฺราหฺมณานํ เตวิชฺโช, อญฺญถา จ ปน อริยสฺส วินเย เตวิชฺโช โหติ. อิมสฺส จ โภ โคตม อริยสฺส วินเย เตวิชฺชสฺส พฺราหฺมณานํ เตวิชฺโช กลํ นาคฺฆติ โสฬสิํ.}

\prose{อภิกฺกนฺตํ โภ โคตม ฯเปฯ อุปาสกํ มํ ภวํ โคตโม ธาเรตุ อชฺชตคฺเค ปาณุเปตํ สรณํ คตนฺติ.  นวมํ.}

\csromanheader{2}{10. \textbf{สงฺคารวสุตฺต}}
\proseitem{61}{อถ โข สงฺคารโว พฺราหฺมโณ เยน ภควา เตนุปสงฺกมิ, อุปสงฺกมิตฺวา ภควตา สทฺธิํ สมฺโมทิ, สมฺโมทนียํ กถํ สารณียํ วีติสาเรตฺวา เอกมนฺตํ นิสีทิ, เอกมนฺตํ นิสินฺโน โข สงฺคารโว พฺราหฺมโณ ภควนฺตํ เอตทโวจ “มยมสฺสุ โภ โคตม พฺราหฺมณา นาม ยญฺญํ ยชามปิ ยชาเปมปิ, ตตฺร โภ โคตม โย เจว ยชติ\csromansharedfootnote{820840381701523c}{โย เจว ยญฺญํ ยชติ (สฺยา, กํ)} โย จ ยชาเปติ, สพฺเพ เต อเนกสารีริกํ ปุญฺญ\-ปฺปฏิปทํ ปฏิปนฺนา โหนฺติ, ยทิทํ ยญฺญาธิกรณํ. โย ปนายํ โภ โคตม ยสฺส วา ตสฺส วา กุลา อคารสฺมา อนคาริยํ ปพฺพชิโต เอกมตฺตานํ ทเมติ, เอกมตฺตานํ สเมติ, เอกมตฺตานํ ปรินิพฺพาเปติ. เอวมสฺสายํ เอกสารีริกํ ปุญฺญ\-ปฺปฏิปทํ ปฏิปนฺโน โหติ, ยทิทํ ปพฺพชฺชา\-ธิกรณนฺติ.}

\prose{เตน หิ พฺราหฺมณ ตญฺเญเวตฺถ ปฏิปุจฺฉิสฺสามิ. ยถา เต ขเมยฺย, ตถา นํ พฺยากเรยฺยาสิ. ตํ กิํ มญฺญสิ พฺราหฺมณ, อิธ ตถาคโต}

\vspace{\tipitakaparskip}
\csromancenter{(6) 1. พฺราหฺมณวคฺค โลเก อุปฺปชฺชติ อรหํ สมฺมาสมฺพุทฺโธ วิชฺชา\-จรณ\-สมฺปนฺโน สุคโต โลกวิทู อนุตฺตโร ปุริส\-ทมฺม\-สารถิ สตฺถา เทวมนุสฺสานํ พุทฺโธ ภควา, โส เอวมาห “เอถายํ มคฺโค อยํ ปฏิปทา, ยถาปฏิปนฺโน อหํ อนุตฺตรํ พฺรหฺมจริโยคธํ สยํ อภิญฺญา สจฺฉิกตฺวา ปเวเทมิ, เอถ\csromansharedfootnote{0a95b81431f872be}{เอตํ (ก)} ตุเมฺหปิ ตถา ปฏิปชฺชถ, ยถาปฏิปนฺนา ตุเมฺหปิ อนุตฺตรํ พฺรหฺมจริโยคธํ สยํ อภิญฺญา สจฺฉิกตฺวา อุปสมฺปชฺช วิหริสฺสถา”ติ. อิติ อยํ เจว\csromansharedfootnote{57d37d5d488ee0fb}{สยํ เจว (ก)} สตฺถา ธมฺมํ เทเสติ, ปเร จ ตถตฺถาย ปฏิปชฺชนฺติ. ตานิ โข ปน โหนฺติ อเนกานิปิ สตานิ, อเนกานิปิ สหสฺสานิ, อเนกานิปิ สตสหสฺสานิ.}
\prose{\csromanfolio{169}ตํ กิํ มญฺญสิ พฺราหฺมณ, อิจฺจายํ เอวํ สนฺเต เอกสารีริกา วา ปุญฺญปฺปฏปทา โหติ อเนกสารีริกา วา, ยทิทํ ปพฺพชฺชา\-ธิกรณนฺติ. อิจฺจายมฺปิ\csromansharedfootnote{3581a539adfa71d1}{อิจฺจายนฺเต (ก)} โภ โคตม เอวํ สนฺเต อเนกสารีริกา ปุญฺญปฺปฏิปทา โหติ, ยทิทํ ปพฺพชฺชา\-ธิกรณนฺติ.}

\prose{เอวํ วุตฺเต อายสฺมา อานนฺโท สงฺคารวํ พฺราหฺมณํ เอตทโวจ “อิมาสํ เต พฺราหฺมณ ทฺวินฺนํ ปฏิปทานํ กตมา ปฏิปทา ขมติ อปฺปตฺถตรา จ อปฺป\-สมารมฺภ\-ตรา จ มหปฺผลตรา จ มหา\-นิสํส\-ตรา จา”ติ. เอวํ วุตฺเต สงฺคารโว พฺราหฺมโณ อายสฺมนฺตํ อานนฺทํ เอตทโวจ “เสยฺยถาปิ ภวํ โคตโม ภวํ จานนฺโท, เอเต เม ปุชฺชา, เอเต เม ปาสํสา”ติ.}

\prose{ทุติยมฺปิ โข อายสฺมา อานนฺโท สงฺคารวํ พฺราหฺมณํ เอตทโวจ “น โข ตฺยาหํ พฺราหฺมณ เอวํ ปุจฺฉามิ ‘เก วา เต ปุชฺชา, เก วา เต ปาสํสา’ติ. เอวํ โข ตฺยาหํ พฺราหฺมณ ปุจฺฉามิ ‘อิมาสํ เต พฺราหฺมณ ทฺวินฺนํ ปฏิปทานํ กตมา ปฏิปทา ขมติ อปฺปตฺถตรา จ อปฺป\-สมารมฺภ\-ตรา จ มหปฺผลตรา จ มหา\-นิสํส\-ตรา จา’ติ. ทุติยมฺปิ โข สงฺคารโว พฺราหฺมโณ อายสฺมนฺตํ อานนฺทํ เอตทโวจ “เสยฺยถาปิ ภวํ โคตโม ภวญฺจานนฺโท, เอเต เม ปุชฺชา, เอเต เม ปาสํสา”ติ.}

\prose{ตติยมฺปิ โข อายสฺมา อานนฺโท สงฺคารวํ พฺราหฺมณํ เอตทโวจ “น โข ตฺยาหํ พฺราหฺมณ เอวํ ปุจฺฉามิ ‘เก วา เต ปุชฺชา, เก วา เต ปาสํสา’ติ. เอวํ โข ตฺยาหํ พฺราหฺมณ ปุจฺฉามิ ‘อิมาสํ เต พฺราหฺมณ ทฺวินฺนํ ปฏิปทานํ กตมา ปฏิปทา ขมติ อปฺปตฺถตรา จ อปฺป\-สมารมฺภ\-ตรา จ \csromanfolio{170}มหปฺผลตรา จ มหา\-นิสํส\-ตรา จา’ติ”. ตติยมฺปิ โข สงฺคารโว พฺราหฺมโณ อายสฺมนฺตํ อานนฺทํ เอตทโวจ “เสยฺยถาปิ ภวํ โคตโม ภวญฺจานนฺโท, เอเต เม ปุชฺชา, เอเต เม ปาสํสา”ติ.}

\prose{อถ โข ภควโต เอตทโหสิ “ยาว ตติยมฺปิ โข สงฺคารโว พฺราหฺมโณ อานนฺเทน สหธมฺมิกํ ปญฺหํ ปุฏฺโฐ สํสาเทติ\csromansharedfootnote{1a9b30e400989d98}{ม 1. 275 ปิฏฺเฐปิ.} โน วิสฺสชฺเชติ, ยํนูนาหํ ปริโมเจยฺยนฺ”ติ. อถ โข ภควา สงฺคารวํ พฺราหฺมณํ เอตทโวจ “กา นฺวชฺช พฺราหฺมณ ราชนฺเตปุเร ราชปุริสานํ\csromansharedfootnote{87f069ad5c136faa}{ราชปริสายํ (สี, สฺยา, กํ, อิ) * ขุ 9. 401; ที 1. 206 ปิฏฺเฐสุปิ.} สนฺนิสินฺนานํ สนฺนิปติตานํ อนฺตรากถา อุทปาที”ติ. อยํ ขฺวชฺช โภ โคตม ราชนฺเตปุเร ราชปุริสานํ สนฺนิสินฺนานํ สนฺนิปติตานํ อนฺตรากถา อุทปาทิ “ปุพฺเพ สุทํ อปฺปตรา เจว ภิกฺขู อเหสุํ, พหุตรา จ อุตฺตริ มนุสฺสธมฺมา อิทฺธิ\-ปาฏิหาริยํ ทสฺเสสุํ. เอตรหิ ปน พหุตรา เจว ภิกฺขู, อปฺปตรา จ อุตฺตริ มนุสฺสธมฺมา อิทฺธิ\-ปาฏิหาริยํ ทสฺเสนฺตี”ติ. อยํ ขฺวชฺช โภ โคตม ราชนฺเตปุเร ราชปุริสานํ สนฺนิสินฺนานํ สนฺนิปติตานํ อนฺตรากถา อุทปาทีติ.}

\prose{\csromansymbolmark{*}ตีณิ โข อิมานิ พฺราหฺมณ ปาฏิหาริยานิ. กตมานิ ตีณิ? อิทฺธิ\-ปาฏิหาริยํ อาเทสนา\-ปาฏิหาริยํ อนุสาสนี\-ปาฏิหาริยํ. กตมญฺจ พฺราหฺมณ อิทฺธิ\-ปาฏิหาริยํ? อิธ พฺราหฺมณ เอกจฺโจ อเนกวิหิตํ อิทฺธิวิธํ ปจฺจนุโภติ. เอโกปิ หุตฺวา พหุธา โหติ, พหุธาปิ หุตฺวา เอโก โหติ, อาวิภาวํ ติโรภาวํ ติโรกุฏฺฏํ ติโรปาการํ ติโรปพฺพตํ อสชฺชมาโน คจฺฉติ เสยฺยถาปิ อากาเส. ปถวิยาปิ อุมฺมุชฺช\-นิมุชฺชํ กโรติ เสยฺยถาปิ อุทเก. อุทเกปิ อภิชฺชมาเน คจฺฉติ เสยฺยถาปิ ปถวิยํ. อากาเสปิ ปลฺลงฺเกน กมติ เสยฺยถาปิ ปกฺขี สกุโณ. อิเมปิ จนฺทิมสูริเย เอวํมหิทฺธิเก เอวํ\-มหา\-นุภาเว ปาณินา ปริมสติ\csromansharedfootnote{b1b3fdbab061a4f3}{ปรามสติ (ที 1. 206; ขุ 9. 107 ปิฏฺเฐสุ.)} ปริมชฺชติ. ยาว พฺรหฺมโลกาปิ กาเยน วสํ วตฺเตติ. อิทํ วุจฺจติ พฺราหฺมณ อิทฺธิ\-ปาฏิหาริยํ.}

\prose{กตมญฺจ พฺราหฺมณ อาเทสนา\-ปาฏิหาริยํ? อิธ พฺราหฺมณ เอกจฺโจ นิมิตฺเตน อาทิสติ “เอวมฺปิ เต มโน, อิตฺถมฺปิ เต มโน, อิติปิ เต จิตฺตนฺ”ติ. โส พหุํ เจปิ อาทิสติ, ตเถว ตํ โหติ โน อญฺญถา.}

\chapterheadpageread{(6) 1. พฺราหฺมณวคฺค}
\prose{\csromanfolio{171}อิธ ปน พฺราหฺมณ เอกจฺโจ น เหว โข นิมิตฺเตน อาทิสติ, อปิ จ โข มนุสฺสานํ วา อมนุสฺสานํ วา เทวตานํ วา สทฺทํ สุตฺวา อาทิสติ “เอวมฺปิ เต มโน, อิตฺถมฺปิ เต มโน, อิติปิ เต จิตฺตนฺ”ติ. โส พหุํ เจปิ อาทิสติ, ตเถว ตํ โหติ โน อญฺญถา.}

\prose{อิธ ปน พฺราหฺมณ เอกจฺโจ น เหว โข นิมิตฺเตน อาทิสติ, นปิ มนุสฺสานํ วา อมนุสฺสานํ วา เทวตานํ วา สทฺทํ สุตฺวา อาทิสติ, อปิ จ โข วิตกฺกยโต วิจารยโต วิตกฺก\-วิปฺผาร\-สทฺทํ สุตฺวา อาทิสติ “เอวมฺปิ เต มโน, อิตฺถมฺปิ เต มโน, อิติปิ เต จิตฺตนฺ”ติ. โส พหุํ เจปิ อาทิสติ, ตเถว ตํ โหติ โน อญฺญถา.}

\prose{อิธ ปน พฺราหฺมณ เอกจฺโจ น เหว โข นิมิตฺเตน อาทิสติ, นปิ มนุสฺสานํ วา อมนุสฺสานํ วา เทวตานํ วา สทฺทํ สุตฺวา อาทิสติ, นปิ วิตกฺกยโต วิจารยโต วิตกฺก\-วิปฺผาร\-สทฺทํ สุตฺวา อาทิสติ. อปิ จ โข อวิตกฺกํ อวิจารํ สมาธิํ สมาปนฺนสฺส เจตสา เจโต ปริจฺจ ปชานาติ “ยถา อิมสฺส โภโต มโนสงฺขารา ปณิหิตา, อิมสฺส จิตฺตสฺส อนนฺตรา อมุํ นาม วิตกฺกํ วิตกฺเกสฺสตี”ติ. โส พหุํ เจปิ อาทิสติ, ตเถว ตํ โหติ โน อญฺญถา. อิทํ วุจฺจติ พฺราหฺมณ อาเทสนา\-ปาฏิหาริยํ.}

\prose{กตมญฺจ พฺราหฺมณ อนุสาสนี\-ปาฏิหาริยํ? อิธ พฺราหฺมณ เอกจฺโจ เอวมนุสาสติ “เอวํ วิตกฺเกถ, มา เอวํ วิตกฺกยิตฺถ. เอวํ มนสิ กโรถ, มา เอวํ มนสากตฺถ. อิทํ ปชหถ, อิทํ อุปสมฺปชฺช วิหรถา”ติ. อิทํ วุจฺจติ พฺราหฺมณ อนุสาสนี\-ปาฏิหาริยํ. อิมานิ โข พฺราหฺมณ ตีณิ ปาฏิหาริยานิ. อิเมสํ เต พฺราหฺมณ ติณฺณํ ปาฏิหาริยานํ กตมํ ปาฏิหาริยํ ขมติ อภิกฺกนฺตตรญฺจ ปณีตตร\-ญฺจาติ.}

\prose{ตตฺร โภ โคตม ยทิทํ\csromansharedfootnote{ea09ead0b9ad4456}{ยมิทํ (สฺยา, กํ, อิ)} ปาฏิหาริยํ, อิเธกจฺโจ อเนกวิหิตํ อิทฺธิวิธํ ปจฺจนุโภติ ฯเปฯ. ยาว พฺราหฺมโลกาปิ กาเยน วสํ วตฺเตติ. อิทํ โภ โคตม ปาฏิหาริยํ. โยว\csromansharedfootnote{fbfd8c11c59df0f2}{โย จ (สฺยา, กํ, อิ, ก)} นํ กโรติ, โสว\csromansharedfootnote{ce966cdb1d1c8315}{โส จ (สฺยา, กํ, อิ, ก)} นํ ปฏิสํเวเทติ. โยว\csromansharedfootnote{1cca7ba2e8749ddc}{ตสฺสเมว (สี, ก), ตสฺส เจว (สฺยา, กํ, อิ)} นํ กโรติ, ตสฺเสว\csromansharedfootnote{fbfd8c11c59df0f2}{โย จ (สฺยา, กํ, อิ, ก)} ตํ โหติ. อิทํ เม โภ โคตม ปาฏิหาริยํ มายา\-สหธมฺมรูปํ วิย ขายติ.}

\prose{\csromanfolio{172}ยมฺปิทํ โภ โคตม ปาฏิหาริยํ อิเธกจฺโจ นิมิตฺเตน อาทิสติ “เอวมฺปิ เต มโน, อิตฺถมฺปิ เต มโน, อิติปิ เต จิตฺตนฺ”ติ. โส พหุํ เจปิ อาทิสติ, ตเถว ตํ โหติ โน อญฺญถา. อิธ ปน โภ โคตม เอกจฺโจ น เหว โข นิมิตฺเตน อาทิสติ. อปิ จ โข มนุสฺสานํ วา อมนุสฺสานํ วา เทวตานํ วา สทฺทํ สุตฺวา อาทิสติ ฯเปฯ นปิ มนุสฺสานํ วา อมนุสฺสานํ วา เทวตานํ วา สทฺทํ สุตฺวา อาทิสติ. อปิ จ โข วิตกฺกยโต วิจารยโต วิตกฺก\-วิปฺผาร\-สทฺทํ สุตฺวา อาทิสติ ฯเปฯ นปิ วิตกฺกยโต วิจารยโต วิตกฺก\-วิปฺผาร\-สทฺทํ สุตฺวา อาทิสติ. อปิ จ โข อวิตกฺกํ อวิจารํ สมาธิํ สมาปนฺนสฺส เจตสา เจโต ปริจฺจ ปชานาติ “ยถา อิมสฺส โภโต มโนสงฺขารา ปณิหิตา, อิมสฺส จิตฺตสฺส อนนฺตรา อมุํ นาม วิตกฺกํ วิตกฺเกสฺสตี”ติ. โส พหุํ เจปิ อาทิสติ, ตเถว ตํ โหติ โน อญฺญถา. อิทมฺปิ โภ โคตม ปาฏิหาริยํ. โยว นํ กโรติ, โสว นํ ปฏิสํเวเทติ. โยว นํ กโรติ, ตสฺเสว ตํ โหติ. อิทมฺปิ เม โภ โคตม ปาฏิหาริยํ มายา\-สหธมฺมรูปํ วิย ขายติ.}

\prose{ยญฺจ โข อิทํ โภ โคตม ปาฏิหาริยํ, อิเธกจฺโจ เอวํ อนุสาสติ “เอวํ วิตกฺเกถ, มา เอวํ วิตกฺกยิตฺถ. เอวํ มนสิ กโรถ, มา เอวํ มนสากตฺถ. อิทํ ปชหถ, อิทํ อุปสมฺปชฺช วิหรถา”ติ. อิทเมว โภ โคตม ปาฏิหาริยํ ขมติ อิเมสํ ติณฺณํ ปาฏิหาริยานํ อภิกฺกนฺตตรญฺจ ปณีตตรญฺจ.}

\prose{อจฺฉริยํ โภ โคตม, อพฺภุตํ โภ โคตม. ยาว สุภาสิตมิทํ โภตา โคตเมน อิเมหิ จ มยํ ตีหิ ปาฏิหาริเยหิ สมนฺนาคตํ ภวนฺตํ โคตมํ ธาเรม. ภวํ หิ โคตโม อเนกวิหิตํ อิทฺธิวิธํ ปจฺจนุโภติ ฯเปฯ. ยาว พฺรหฺมโลกาปิ กาเยน วสํ วตฺเตติ. ภวํ หิ โคตโม อวิตกฺกํ อวิจารํ สมาธิํ สมาปนฺนสฺส เจตสา เจโต ปริจฺจ ปชานาติ “ยถา อิมสฺส โภโต มโนสงฺขารา ปณิหิตา, อิมสฺส จิตฺตสฺส อนนฺตรา อมุํ นาม วิตกฺกํ วิตกฺเกสฺสตี”ติ. ภวํ หิ โคตโม เอวมนุสาสติ “เอวํ วิตกฺเกถ, มา เอวํ วิตกฺกยิตฺถ. เอวํ มนสิ กโรถ, มา เอวํ มนสากตฺถ. อิทํ ปชหถ, อิทํ อุปสมฺปชฺช วิหรถา”ติ.}

\chapterheadpageread{(6) 1. พฺราหฺมณวคฺค}
\prose{\csromanfolio{173}อทฺธา โข ตฺยายํ พฺราหฺมณ อาสชฺช อุปนีย วาจา ภาสิตา, อปิ จ ตฺยาหํ พฺยากริสฺสามิ, อหํ หิ พฺราหฺมณ อเนกวิหิตํ อิทฺธิวิธํ ปจฺจนุโภมิ ฯเปฯ. ยาว พฺรหฺมโลกาปิ กาเยน วสํ วตฺเตมิ. อหํ หิ พฺราหฺมณ อวิตกฺกํ อวิจารํ สมาธิํ สมาปนฺนสฺส เจตสา เจโต ปริจฺจ ปชานามิ “ยถา อิมสฺส โภโต มโนสงฺขารา ปณิหิตา, อิมสฺส จิตฺตสฺส อนนฺตรา อมุํ นาม วิตกฺกํ วิตกฺเกสฺสตี”ติ. อหํ หิ พฺราหฺมณ เอวมนุสาสามิ “เอวํ วิตกฺเกถ, มา เอวํ วิตกฺกยิตฺถ. เอวํ มนสิ กโรถ, มา เอวํ มนสากตฺถ. อิทํ ปชหถ, อิทํ อุปสมฺปชฺช วิหรถา”ติ.}

\prose{อตฺถิ ปน โภ โคตม อญฺโญ เอกภิกฺขุปิ, โย อิเมหิ ตีหิ ปาฏิหาริเยหิ สมนฺนาคโต อญฺญตฺร โภตา โคตเมนาติ. น โข พฺราหฺมณ เอกํเยว สตํ, น เทฺว สตานิ, น ตีณิ สตานิ, น จตฺตาริ สตานิ, น ปญฺจ สตานิ, อถ โข ภิยฺโยว เย\csromansharedfootnote{8fffc3d0eeea86fc}{เต (ก) ปสฺส ม 2. 158 ปิฏฺเฐ.} ภิกฺขู อิเมหิ ตีหิ ปาฏิหาริเยหิ สมนฺนาคตาติ. กหํ ปน โภ โคตม เอตรหิ เต ภิกฺขู วิหรนฺตีติ. อิมสฺมิํเยว โข พฺราหฺมณ ภิกฺขุ\-สํเฆติ.}

\prosewithclosermidruled{อภิกฺกนฺตํ โภ โคตม, อภิกฺกนฺตํ โภ โคตม, เสยฺยถาปิ โภ โคตม นิกฺกุชฺชิตํ วา อุกฺกุชฺเชยฺย, ปฏิจฺฉนฺนํ วา วิวเรยฺย, มูฬฺหสฺส วา มคฺคํ อาจิกฺเขยฺย, อนฺธกาเร วา เตลปชฺโชตํ ธาเรยฺย “จกฺขุมนฺโต รูปานิ ทกฺขนฺตี”ติ. เอวเมวํ โภตา โคตเมน อเนก\-ปริยาเยน ธมฺโม ปกาสิโต, เอสาหํ ภวนฺตํ โคตมํ สรณํ คจฺฉามิ ธมฺมญฺจ ภิกฺขุสํฆญฺจ, อุปาสกํ มํ ภวํ โคตโม ธาเรตุ อชฺชตคฺเค ปาณุเปตํ สรณํ คตนฺติ.  ทสมํ.}{พฺราหฺมณวคฺโค ปฐโม.}
\csromancenter{\textbf{ตสฺสุทฺทานํ}}
\setlength{\csromangathapagewidth}{0pt}
\setlength{\csromangathawakwidth}{0pt}
\csromangathameasuregroup{เทฺว พฺราหฺมณา จญฺญตโร, \\ ปโลกวจฺโฉ ติกณฺโณ,}{\csromangathabat{เทฺว พฺราหฺมณา จญฺญตโร,}{ปริพฺพาชเกน นิพฺพุตํ.} \\ \csromangathabat{ปโลกวจฺโฉ ติกณฺโณ,}{โสณิ สงฺคารเวน จาติ.}}
\csromangathasetleft{เทฺว พฺราหฺมณา จญฺญตโร, \\ ปโลกวจฺโฉ ติกณฺโณ,}
\csromangathagroup{\csromangathastanza{\csromangathabat{เทฺว พฺราหฺมณา จญฺญตโร,}{ปริพฺพาชเกน นิพฺพุตํ.} \\ \csromangathabat{ปโลกวจฺโฉ ติกณฺโณ,}{โสณิ สงฺคารเวน จาติ.}}}

\csromanpalirecto
\gambhira{ติกนิปาตปาฬิ \textbf{(7) 2. มหาวคฺค}}
\csromanheader{2}{1. ติตฺถายตน\-สุตฺต}
\proseitem{62}{\csromanfolio{174}ตีณิมานิ ภิกฺขเว ติตฺถายตนานิ, ยานิ ปณฺฑิเตหิ สมนุ\-ยุญฺชิย\-มานานิ สมนุ\-คาหิย\-มานานิ\csromansharedfootnote{4b8704fcd17aa7ef}{สมนุคฺคาหิยมานานิ (สฺยา, กํ, ก)} สมนุ\-ภาสิย\-มานานิ ปรมฺปิ คนฺตฺวา อกิริยาย สณฺฐหนฺติ. กตมานิ ตีณิ? สนฺติ ภิกฺขเว เอเก สมณพฺราหฺมณา เอวํวาทิโน เอวํทิฏฺฐิโน “ยํ กิญฺจายํ ปุริสปุคฺคโล ปฏิสํเวเทติ สุขํ วา ทุกฺขํ วา อทุกฺขมสุขํ วา, สพฺพํ ตํ ปุพฺเพกตเหตู”ติ. สนฺติ ภิกฺขเว เอเก สมณพฺราหฺมณา เอวํวาทิโน เอวํทิฏฺฐิโน “ยํ กิญฺจายํ ปุริสปุคฺคโล ปฏิสํเวเทติ สุขํ วา ทุกฺขํ วา อทุกฺขมสุขํ วา, สพฺพํ ตํ อิสฺสรนิมฺมานเหตู”ติ. สนฺติ ภิกฺขเว เอเก สมณพฺราหฺมณา เอวํวาทิโน เอวํทิฏฺฐิโน “ยํ กิญฺจายํ ปุริสปุคฺคโล ปฏิสํเวเทติ สุขํ วา ทุกฺขํ วา อทุกฺขมสุขํ วา, สพฺพํ ตํ อเหตุอปฺปจฺจยา”ติ.}

\prose{ตตฺร ภิกฺขเว เย เต สมณพฺราหฺมณา เอวํวาทิโน เอวํทิฏฺฐิโน “ยํ กิญฺจายํ ปุริสปุคฺคโล ปฏิสํเวเทติ สุขํ วา ทุกฺขํ วา อทุกฺขมสุขํ วา, สพฺพํ ตํ ปุพฺเพกตเหตู”ติ. ตฺยาหํ อุปสงฺกมิตฺวา เอวํ วทามิ “สจฺจํ กิร ตุเมฺห อายสฺมนฺโต เอวํวาทิโน เอวํทิฏฺฐิโน ‘ยํ กิญฺจายํ ปุริสปุคฺคโล ปฏิสํเวเทติ สุขํ วา ทุกฺขํ วา อทุกฺขมสุขํ วา, สพฺพํ ตํ ปุพฺเพกตเหตู’ติ”. เต จ เม\csromansharedfootnote{d08bf1da8b35bd24}{เต เจ เม (สี, สฺยา, กํ, อิ)} เอวํ ปุฏฺฐา “อามา”ติ\csromansharedfootnote{c538567820e9d75f}{อาโมติ (สี)} ปฏิชานนฺติ. ตฺยาหํ เอวํ วทามิ “เตนหายสฺมนฺโต ปาณาติปาติโน ภวิสฺสนฺติ ปุพฺเพกตเหตุ, อทินฺนาทายิโน ภวิสฺสนฺติ ปุพฺเพกตเหตุ, อพฺรหฺมจาริโน ภวิสฺสนฺติ ปุพฺเพกตเหตุ, มุสาวาทิโน ภวิสฺสนฺติ ปุพฺเพกตเหตุ, ปิสุณวาจา ภวิสฺสนฺติ ปุพฺเพกตเหตุ, ผรุสวาจา ภวิสฺสนฺติ ปุพฺเพกตเหตุ, สมฺผ\-ปฺปลาปิโน ภวิสฺสนฺติ ปุพฺเพกตเหตุ, อภิชฺฌาลุโน ภวิสฺสนฺติ ปุพฺเพกตเหตุ, พฺยาปนฺนจิตฺตา ภวิสฺสนฺติ ปุพฺเพกตเหตุ, มิจฺฉาทิฏฺฐิกา ภวิสฺสนฺติ ปุพฺเพกตเหตุ”.}

\prose{ปุพฺเพกตํ โข ปน ภิกฺขเว สารโต ปจฺจาคจฺฉตํ น โหติ ฉนฺโท วา วายาโม วา “อิทํ วา กรณียํ, อิทํ วา อกรณียนฺ”ติ. อิติ กรณียา\-กรณีเย โข ปน สจฺจโต เถตโต อนุปลพฺภิยมาเน มุฏฺฐสฺสตีนํ อนารกฺขานํ วิหรตํ น โหติ ปจฺจตฺตํ สหธมฺมิโก}

\vspace{\tipitakaparskip}
\csromancenter{(7) 2. มหาวคฺค สมณวาโท. อยํ โข เม ภิกฺขเว เตสุ สมณ\-พฺราหฺมเณสุ เอวํวาทีสุ เอวํทิฏฺฐีสุ ปฐโม สหธมฺมิโก นิคฺคโห โหติ. (1)}
\prose{\csromanfolio{175}ตตฺร ภิกฺขเว เย เต สมณพฺราหฺมณา เอวํวาทิโน เอวํทิฏฺฐิโน “ยํ กิญฺจายํ ปุริสปุคฺคโล ปฏิสํเวเทติ สุขํ วา ทุกฺขํ วา อทุกฺขมสุขํ วา, สพฺพํ ตํ อิสฺสรนิมฺมานเหตู”ติ. ตฺยาหํ อุปสงฺกมิตฺวา เอวํ วทามิ “สจฺจํ กิร ตุเมฺห อายสฺมนฺโต เอวํวาทิโน เอวํทิฏฺฐิโน ‘ยํ กิญฺจายํ ปุริสปุคฺคโล ปฏิสํเวเทติ สุขํ วา ทุกฺขํ วา อทุกฺขมสุขํ วา, สพฺพํ ตํ อิสฺสรนิมฺมานเหตู’ติ”. เต จ เม เอวํ ปุฏฺฐา “อมา”ติ ปฏิชานนฺติ. ตฺยาหํ เอวํ วทามิ “เตนหายสฺมนฺโต ปาณาติปาติโน ภวิสฺสนฺติ อิสฺสร\-นิมฺมาน\-เหตุ, อทินฺนาทายิโน ภวิสฺสนฺติ อิสฺสร\-นิมฺมาน\-เหตุ, อพฺรหฺมจาริโน ภวิสฺสนฺติ อิสฺสร\-นิมฺมาน\-เหตุ, มุสาวาทิโน ภวิสฺสนฺติ อิสฺสร\-นิมฺมาน\-เหตุ, ปิสุณวาจา ภวิสฺสนฺติ อิสฺสร\-นิมฺมาน\-เหตุ, ผรุสวาจา ภวิสฺสนฺติ อิสฺสร\-นิมฺมาน\-เหตุ, สมฺผ\-ปฺปลาปิโน ภวิสฺสนฺติ อิสฺสร\-นิมฺมาน\-เหตุ, อภิชฺฌาลุโน ภวิสฺสนฺติ อิสฺสร\-นิมฺมาน\-เหตุ, พฺยาปนฺนจิตฺตา ภวิสฺสนฺติ อิสฺสร\-นิมฺมาน\-เหตุ, มิจฺฉาทิฏฺฐิกา ภวิสฺสนฺติ อิสฺสร\-นิมฺมาน\-เหตุ”.}

\prose{อิสฺสรนิมฺมานํ โข ปน ภิกฺขเว สารโต ปจฺจาคจฺฉตํ น โหติ ฉนฺโท วา วายาโม วา อิทํ วา กรณียํ, อิทํ วา อกรณียนฺ”ติ. อิติ กรณียา\-กรณีเย โข ปน สจฺจโต เถตโต อนุปลพฺภิยมาเน มุฏฺฐสฺสตีนํ อนารกฺขานํ วิหรตํ น โหติ ปจฺจตฺตํ สหธมฺมิโก สมณวาโท. อยํ โข เม ภิกฺขเว เตสุ สมณ\-พฺราหฺมเณสุ เอวํวาทีสุ เอวํทิฏฺฐีสุ ทุติโย สหธมฺมิโก นิคฺคโห โหติ. (2)}

\prose{ตตฺร ภิกฺขเว เย เต สมณพฺราหฺมณา เอวํวาทิโน เอวํทิฏฺฐิโน “ยํ กิญฺจายํ ปุริสปุคฺคโล ปฏิสํเวเทติ สุขํ วา ทุกฺขํ วา อทุกฺขมสุขํ วา, สพฺพํ ตํ อเหตุอปฺปจฺจยา”ติ. ตฺยาหํ อุปสงฺกมิตฺวา เอวํ วทามิ “สจฺจํ กิร ตุเมฺห อายสฺมนฺโต เอวํวาทิโน เอวํทิฏฺฐิโน ‘ยํ กิญฺจายํ ปุริสปุคฺคโล ปฏิสํเวเทติ สุขํ วา ทุกฺขํ วา อทุกฺขมสุขํ วา, สพฺพํ ตํ อเหตุอปฺปจฺจยา’ติ”. เต จ เม เอวํ ปุฏฺฐา “อามา”ติ ปฏิชานนฺติ. ตฺยาหํ เอวํ วทามิ “เตนหายสฺมนฺโต ปาณาติปาติโน ภวิสฺสนฺติ อเหตุอปฺปจฺจยา ฯเปฯ มิจฺฉาทิฏฺฐิกา ภวิสฺสนฺติ อเหตุอปฺปจฺจยา”.}

\prose{\csromanfolio{176}อเหตุ\-อปฺปจฺจยํ\csromansharedfootnote{d4f33c6cb88f433f}{อเหตุํ (สี), อเหตุ (สฺยา, กํ), อเหตุอปฺปจฺจยา (อิ), อเหตุํ อปฺปจฺจยํ (ก)} โข ปน ภิกฺขเว สารโต ปจฺจาคจฺฉตํ น โหติ ฉนฺโท วา วายาโม วา “อิทํ วา กรณียํ, อิทํ วา อกรณียนฺ”ติ. อิติ กรณียา\-กรณีเย โข ปน สจฺจโต เถตโต อนุปลพฺภิยมาเน มุฏฺฐสฺสาตีนํ อนารกฺขานํ วิหรตํ น โหติ ปจฺจตฺตํ สหธมฺมิโก สมณวาโท. อยํ โข เม ภิกฺขเว เตสุ สมณ\-พฺราหฺมเณสุ เอวํวาทีสุ เอวํทิฏฺฐีสุ ตติโย สหธมฺมิโก นิคฺคโห โหติ. (3)}

\prose{อิมานิ โข ภิกฺขเว ตีณิ ติตฺถายตนานิ, ยานิ ปณฺฑิเตหิ สมนุ\-ยุญฺชิย\-มานานิ สมนุ\-คาหิย\-มานานิ สมนุ\-ภาสิย\-มานานิ ปรมฺปิ คนฺตฺวา อกิริยาย สณฺฐหนฺติ.}

\prose{อยํ โข ปน ภิกฺขเว มยา ธมฺมา เทสิโต อนิคฺคหิโต อสํกิลิฏฺโฐ อนุปวชฺโช อปฺปฏิกุฏฺโฐ สมเณหิ พฺราหฺมเณหิ วิญฺญูหิ. กตโม จ ภิกฺขเว มยา ธมฺโม เทสิโต อนิคฺคหิโต อสํกิลิฏฺโฐ อนุปวชฺโช อปฺปฏิกุฏฺโฐ สมเณหิ พฺราหฺมเณหิ วิญฺญูหิ? “อิมา ฉ ธาตุโย”ติ ภิกฺขเว มยา ธมฺโม เทสิโต อนิคฺคหิโต อสํกิลิฏฺโฐ อนุปวชฺโช อปฺปฏิกุฏฺโฐ สมเณหิ พฺราหฺมเณหิ วิญฺญูหิ. “อิมานิ ฉ ผสฺสายตนานี”ติ ภิกฺขเว มยา ธมฺโม เทสิโต อนิคฺคหิโต อสํกิลิฏฺโฐ อนุปวชฺโช อปฺปฏิกุฏฺโฐ สมเณหิ พฺราหฺมเณหิ วิญฺญูหิ. “อิเม อฏฺฐารส มโนปวิจารา”ติ ภิกฺขเว มยา ธมฺโม เทสิโต อนิคฺคหิโต อสํกิลิฏฺโฐ อนุปวชฺโช อปฺปฏิกุฏฺโฐ สมเณหิ พฺราหฺมเณหิ วิญฺญูหิ. “อิมานิ จตฺตาริ อริยสจฺจานี”ติ ภิกฺขเว มยา ธมฺโม เทสิโต อนิคฺคหิโต อสํกิลิฏฺโฐ อนุปวชฺโช อปฺปฏิกุฏฺโฐ สมเณหิ พฺราหฺมเณหิ วิญฺญูหิ.}

\prose{“อิมา ฉ ธาตุโยติ ภิกฺขเว มยา ธมฺโม เทสิโต อนิคฺคหิโต อสํกิลิฏฺโฐ อนุปวชฺโช อปฺปฏิกุฏฺโฐ สมเณหิ พฺราหฺมเณหิ วิญฺญูหี”ติ อิติ โข ปเนตํ วุตฺตํ, กิญฺเจตํ ปฏิจฺจ วุตฺตํ? ฉยิมา ภิกฺขเว ธาตุโย ปถวีธาตุ อาโปธาตุ เตโชธาตุ วาโยธาตุ อากาสธาตุ วิญฺญาณธาตุ. “อิมา ฉ ธาตุโยติ ภิกฺขเว มยา ธมฺโม เทสิโต อนิคฺคหิโต อสํกิลิฏฺโฐ อนุปวชฺโช อปฺปฏิกุฏฺโฐ สมเณหิ พฺราหฺมเณหิ วิญฺญูหี”ติ อิติ ยํ ตํ วุตฺตํ, อิทเมตํ ปฏิจฺจ วุตฺตํ.}

\chapterheadpageread{(7) 2. มหาวคฺค}
\prose{\csromanfolio{177}“อิมานิ ฉ ผสฺสายตนานีติ ภิกฺขเว มยา ธมฺโม เทสิโต อนิคฺคหิโต อสํกิลิฏฺโฐ อนุปวชฺโช อปฺปฏิกุฏฺโฐ สมเณหิ พฺราหฺมเณหิ วิญฺญูหี”ติ อิติ โข ปเนตํ วุตฺตํ, กิญฺเจตํ ปฏิจฺจ วุตฺตํ? ฉยิมานิ ภิกฺขเว ผสฺสายตนานิ จกฺขุ ผสฺสายตนํ, โสตํ ผสฺสายตนํ, ฆานํ ผสฺสายตนํ, ชิวฺหา ผสฺสายตนํ, กาโย ผสฺสายตนํ, มโน ผสฺสายตนํ. “อิมานิ ฉ ผสฺสายตนานีติ ภิกฺขเว มยา ธมฺโม เทสิโต อนิคฺคหิโต อสํกิลิฏฺโฐ อนุปวชฺโช อปฺปฏิกุฏฺโฐ สมเณหิ พฺราหฺมเณหิ วิญฺญูหี”ติ อิติ ยํ ตํ วุตฺตํ, อิทเมตํ ปฏิจฺจ วุตฺตํ.}

\prose{“อิเม อฏฺฐารส มโนปวิจาราติ ภิกฺขเว มยา ธมฺโม เทสิโต อนิคฺคหิโต อสํกิลิฏฺโฐ อนุปวชฺโช อปฺปฏิกุฏฺโฐ สมเณหิ พฺราหฺมเณหิ วิญฺญูหี”ติ อิติ โข ปเนตํ วุตฺตํ, กิญฺเจตํ ปฏิจฺจ วุตฺตํ. จกฺขุนา รูปํ ทิสฺวา โสมนสฺส\-ฏฺฐานิยํ รูปํ อุปวิจรติ, โทมนสฺส\-ฏฺฐานิยํ รูปํ อุปวิจรติ, อุเปกฺขา\-ฏฺฐานิยํ รูปํ อุปวิจรติ. โสเตน สทฺทํ สุตฺวา. ฆาเนน คนฺธํ ฆายิตฺวา. ชิวฺหาย รสํ สายิตฺวา. กาเยน โผฏฺฐพฺพํ ผุสิตฺวา. มนสา ธมฺมํ วิญฺญาย โสมนสฺส\-ฏฺฐานิยํ ธมฺมํ อุปวิจรติ, โทมนสฺส\-ฏฺฐานิยํ ธมฺมํ อุปวิจรติ, อุเปกฺขา\-ฏฺฐานิยํ ธมฺมํ อุปวิจรติ. “อิเม อฏฺฐารส มโนปวิจาราติ ภิกฺขเว มยา ธมฺโม เทสิโต อนิคฺคหิโต อสํกิลิฏฺโฐ อนุปวชฺโช อปฺปฏิกุฏฺโฐ สมเณหิ พฺราหฺมเณหิ วิญฺญูหี”ติ อิติ ยํ ตํ วุตฺตํ, อิทเมตํ ปฏิจฺจ วุตฺตํ.}

\prose{“อิมานิ จตฺตาริ อริยสจฺจานีติ ภิกฺขเว มยา ธมฺโม เทสิโต อนิคฺคหิโต อสํกิลิฏฺโฐ อนุปวชฺโช อปฺปฏิกุฏฺโฐ สมเณหิ พฺราหฺมเณหิ วิญฺญูหี”ติ อิติ โข ปเนตํ วุตฺตํ, กิญฺเจตํ ปฏิจฺจ วุตฺตํ. ฉนฺนํ ภิกฺขเว ธาตูนํ อุปาทาย คพฺภสฺสา\-วกฺกนฺติ โหติ, โอกฺกนฺติยา สติ นามรูปํ, นามรูป\-ปจฺจยา สฬายตนํ, สฬายตน\-ปจฺจยา ผสฺโส, ผสฺสปจฺจยา เวทนา, เวทิยมานสฺส โข ปนาหํ ภิกฺขเว “อิทํ ทุกฺขนฺ”ติ ปญฺญเปมิ, “อยํ ทุกฺขสมุทโย”ติ ปญฺญเปมิ, “อยํ ทุกฺขนิโรโธ”ติ ปญฺญเปมิ, “อยํ ทุกฺขนิโรธ\-คามินี ปฏิปทา”ติ ปญฺญเปมิ.}

\prose{กตมญฺจ ภิกฺขเว ทุกฺขํ อริยสจฺจํ? ชาติปิ ทุกฺขา, ชราปิ ทุกฺขา, ( )\csromansharedfootnote{d97138d9c2d757fe}{(พฺยาธิปิ ทุกฺโข) (สี, อิ, ก) อฏฺฐกถาย สํสนฺเทตพฺพํ วิสุทฺธิ 2. 128 ปิฏฺเฐ จ ปสฺสิตพฺพํ.} มรณมฺปิ ทุกฺขํ, โสก\-ปริเทว\-ทุกฺข\-โทมนสฺสุ\-ปายาสาปิ ทุกฺขา,}

\prose{\csromanfolio{178}(อปฺปิเยหิ สมฺปโยโค ทุกฺโข, ปิเยหิ วิปฺปโยโค ทุกฺโข)\csromansharedfootnote{52e12187d79b1208}{(นตฺถิ กตฺถจิ)} ยมฺปิจฺฉํ น ลภติ ตมฺปิ ทุกฺขํ, สํขิตฺเตน ปญฺจุ\-ปาทาน\-กฺขนฺธา ทุกฺขา. อิทํ วุจฺจติ ภิกฺขเว ทุกฺขํ อริยสจฺจํ.}

\prose{กตมญฺจ ภิกฺขเว ทุกฺข\-สมุทยํ\csromansharedfootnote{92abb159022a2cc5}{ทุกฺขสมุทโย (สฺยา, กํ)} อริยสจฺจํ? อวิชฺชาปจฺจยา สงฺขารา, สงฺขาร\-ปจฺจยา วิญฺญาณํ, วิญฺญาณปจฺจยา นามรูปํ, นามรูปจฺจยา สฬายตนํ, สฬายตน\-ปจฺจยา ผสฺโส, ผสฺสปจฺจยา เวทนา, เวทนาปจฺจยา ตณฺหา, ตณฺหาปจฺจยา อุปาทานํ, อุปาทานปจฺจยา ภโว, ภวปจฺจยา ชาติ, ชาติปจฺจยา ชรามรณํ โสกปริเทว\-ทุกฺข- โทมนสฺสุ\-ปายาสา สมฺภวนฺติ. เอวเมตสฺส เกวลสฺส ทุกฺข\-กฺขนฺธสฺส สมุทโย โหติ. อิทํ วุจฺจติ ภิกฺขเว ทุกฺข\-สมุทยํ อริยสจฺจํ.}

\prose{กตมญฺจ ภิกฺขเว ทุกฺขนิโรธํ\csromansharedfootnote{72cf75c19b00299e}{ทุกฺขนิโรโธ (สฺยา, กํ)} อริยสจฺจํ? อวิชฺชาย เตฺวว อเสส\-วิราค\-นิโรธา สงฺขาร\-นิโรโธ, สงฺขาร\-นิโรธา วิญฺญาณนิโรโธ, วิญฺญาณนิโรธา นามรูป\-นิโรโธ, นามรูป\-นิโรธา สฬายตน\-นิโรโธ, สฬายตน\-นิโรธา ผสฺสนิโรโธ, ผสฺสนิโรธา เวทนานิโรโธ, เวทนานิโรธา ตณฺหานิโรโธ, ตณฺหานิโรธา อุปาทานนิโรโธ, อุปาทานนิโรธา ภวนิโรโธ, ภวนิโรธา ชาตินิโรโธ, ชาตินิโรธา ชรามรณํ โสก\-ปริเทว\-ทุกฺข\-โทมนสฺสุ\-ปายาสา นิรุชฺฌนฺติ. เอวเมตสฺส เกวลสฺส ทุกฺข\-กฺขนฺธสฺส นิโรโธ โหติ. อิทํ วุจฺจติ ภิกฺขเว ทุกฺขนิโรธํ อริยสจฺจํ.}

\prose{กตมญฺจ ภิกฺขเว ทุกฺขนิโรธ\-คามินี ปฏิปทา อริยสจฺจํ? อยเมว อริโย อฏฺฐงฺคิโก มคฺโค. เสยฺยถิทํ? สมฺมาทิฏฺฐิ สมฺมาสงฺกปฺโป สมฺมาวาจา สมฺมากมฺมนฺโต สมฺมาอาชีโว สมฺมาวายาโม สมฺมาสติ สมฺมาสมาธิ. อิทํ วุจฺจติ ภิกฺขเว ทุกฺขนิโรธ\-คามินี ปฏิปทา อริยสจฺจํ. “อิมานิ จตฺตาริ อริยสจฺจานีติ ภิกฺขเว มยา ธมฺโม เทสิโต อนิคฺคหิโต อสํกิลิฏฺโฐ อนุปวชฺโช อปฺปฏิกุฏฺโฐ สมเณหิ พฺราหฺมเณหิ วิญฺญูหี”ติ อิติ ยํ ตํ วุตฺตํ, อิทเมตํ ปฏิจฺจ วุตฺตนฺติ.  ปฐมํ.}

\csromanheader{2}{2. \textbf{ภยสุตฺต}}
\proseitem{63}{ตีณิมานิ ภิกฺขเว “อมาตาปุตฺติกานิ ภยานี”ติ อสฺสุตวา ปุถุชฺชโน ภาสติ. กตมานิ ตีณิ? โหติ โส ภิกฺขเว สมโย,}

\vspace{\tipitakaparskip}
\csromancenter{(7) 2. มหาวคฺค ยํ มหาอคฺคิฑาโห วุฏฺฐาติ. มหาอคฺคิฑาเห โข ปน ภิกฺขเว วุฏฺฐิเต เตน คามาปิ ฑยฺหนฺติ, นิคมาปิ ฑยฺหนฺติ, นคราปิ ฑยฺหนฺติ. คาเมสุปิ ฑยฺหมาเนสุ นิคเมสุปิ ฑยฺหมาเนสุ นคเรสุปิ ฑยฺหมาเนสุ ตตฺถ มาตาปิ ปุตฺตํ นปฺปฏิลภติ, ปุตฺโตปิ มาตรํ นปฺปฏิลภติ. อิทํ ภิกฺขเว ปฐมํ “อมาตาปุตฺติกํ ภยนฺ”ติ อสฺสุตวา ปุถุชฺชโน ภาสติ.}
\prose{\csromanfolio{179}ปุน จปรํ ภิกฺขเว โหติ โส สมโย, ยํ มหาเมโฆ วุฏฺฐาติ. มหาเมเฆ โข ปน ภิกฺขเว วุฏฺฐิเต มหา\-อุทกวาหโก สญฺชายติ, มหา\-อุทกวาหเก โข ปน ภิกฺขเว สญฺชายนฺเต เตน คามาปิ วุยฺหนฺติ, นิคมาปิ วุยฺหนฺติ, นคราปิ วุยฺหนฺติ. คาเมสุปิ วุยฺหมาเนสุ นิคเมสุปิ วุยฺหมาเนสุ นคเรสุปิ วุยฺหมาเนสุ ตตฺถ มาตาปิ ปุตฺตํ นปฺปฏิลภติ, ปุตฺโตปิ มาตรํ นปฺปฏิลภติ, อิทํ ภิกฺขเว ทุติยํ “อมาตาปุตฺติกํ ภยนฺ”ติ อสฺสุตวา ปุถุชฺชโน ภาสติ.}

\prose{ปุน จปรํ ภิกฺขเว โหติ โส สมโย, ยํ ภยํ โหติ อฏวิสงฺโกโป จกฺกสมารุฬฺหา ชานปทา ปริยายนฺติ. ภเย โข ปน ภิกฺขเว สติ อฏวิสงฺโกเป จกฺก\-สมารุเฬฺหสุ ชานปเทสุ ปริยายนฺเตสุ ตตฺถ มาตาปิ ปุตฺตํ นปฺปฏิลภติ, ปุตฺโตปิ มาตรํ นปฺปฏิลภติ, อิทํ ภิกฺขเว ตติยํ “อมาตาปุตฺติกํ ภยนฺ”ติ อสฺสุตวา ปุถุชฺชโน ภาสติ. อิมานิ โข ภิกฺขเว ตีณิ “อมาตาปุตฺติกานิ ภยานี”ติ อสฺสุตวา ปุถุชฺชโน ภาสติ.}

\prose{ตานิ โข ปนิมานิ\csromansharedfootnote{f8b47e6e68faa4bc}{อิมานิ โข (สี), อิมานิ โข ปน (ก)} ภิกฺขเว ตีณิ สมาตาปุตฺติกานิเยว ภยานิ “อมาตาปุตฺติกานิ ภยานี”ติ อสฺสุตวา ปุถุชฺชโน ภาสติ. กตมานิ ตีณิ? โหติ โส ภิกฺขเว สมโย, ยํ มหาอคฺคิฑาโห วุฏฺฐาติ, มหาอคฺคิฑาเห โข ปน ภิกฺขเว วุฏฺฐิเต เตน คามาปิ ฑยฺหนฺติ, นิคมาปิ ฑยฺหนฺติ, นคราปิ ฑยฺหนฺติ. คาเมสุปิ ฑยฺหมาเนสุ นิคเมสุปิ ฑยฺหมาเนสุ นคเรสุปิ ฑยฺหมาเนสุ โหติ โส สมโย, ยํ กทาจิ กรหจิ มาตาปิ ปุตฺตํ ปฏิลภติ, ปุตฺโตปิ มาตรํ ปฏิลภติ, อิทํ ภิกฺขเว ปฐมํ สมาตาปุตฺติกํเยว ภยํ “อมาตาปุตฺติกํ ภยนฺ”ติ อสฺสุตวา ปุถุชฺชโน ภาสติ.}

\prose{\csromanfolio{180}ปุน จปรํ ภิกฺขเว โหติ โส สมโย, ยํ มหาเมโฆ วุฏฺฐาติ. มหาเมเฆ โข ปน ภิกฺขเว วุฏฺฐิเต มหา\-อุทกวาหโก สญฺชายติ, มหา\-อุทกวาหเก โข ปน ภิกฺขเว สญฺชาเต เตน คามาปิ วุยฺหนฺติ, นิคมาปิ วุยฺหนฺติ, นคราปิ วุยฺหนฺติ. คาเมสุปิ วุยฺหมาเนสุ นิคเมสุปิ วุยฺหมาเนสุ นคเรสุปิ วุยฺหมาเนสุ โหติ โส สมโย, ยํ กทาจิ กรหจิ มาตาปิ ปุตฺตํ ปฏิลภติ, ปุตฺโตปิ มาตรํ ปฏิลภติ, อิทํ ภิกฺขเว ทุติยํ สมาตาปุตฺติกํเยว ภยํ “อมาตาปุตฺติกํ ภยนฺ”ติ อสฺสุตวา ปุถุชฺชโน ภาสติ.}

\prose{ปุน จปรํ ภิกฺขเว โหติ โส สมโย, ยํ ภยํ โหติ อฏวิสงฺโกโป จกฺกสมารุฬฺหา ชานปทา ปริยายนฺติ. ภเย โข ปน ภิกฺขเว สติ อฏวิสงฺโกเป จกฺก\-สมารุเฬฺหสุ ชานปเทสุ ปริยายนฺเตสุ โหติ โส สมโย, ยํ กทาจิ กรหจิ มาตาปิ ปุตฺตํ ปฏิลภติ, ปุตฺโตปิ มาตรํ ปฏิลภติ. อิทํ ภิกฺขเว ตติยํ สมาตาปุตฺติกํเยว ภยํ “อมาตาปุตฺติกํ ภยนฺ”ติ อสฺสุตวา ปุถุชฺชโน ภาสติ. อิมานิ โข ภิกฺขเว ตีณิ สมาตาปุตฺติกานิเยว ภยานิ “อมาตาปุตฺติกานิ ภยานี”ติ อสฺสุตวา ปุถุชฺชโน ภาสติ.}

\prose{ตีณิมานิ ภิกฺขเว อมาตาปุตฺติกานิ ภยานิ. กตมานิ ตีณิ? ชราภยํ พฺยาธิภยํ มรณภยนฺติ. น ภิกฺขเว มาตา ปุตฺตํ ชีรมานํ เอวํ ลภติ “อหํ ชีรามิ, มา เม ปุตฺโต ชีรี”ติ, ปุตฺโต วา ปน มาตรํ ชีรมานํ น เอวํ ลภติ “อหํ ชีรามิ, มา เม มาตา ชีรี”ติ.}

\prose{น ภิกฺขเว มาตา ปุตฺตํ พฺยาธิยมานํ เอวํ ลภติ “อหํ พฺยาธิยามิ, มา เม ปุตฺโต พฺยาธิยี”ติ, ปุตฺโต วา ปน มาตรํ พฺยาธิยมานํ น เอวํ ลภติ “อหํ พฺยาธิยามิ, มา เม มาตา พฺยาธิยี”ติ.}

\prose{น ภิกฺขเว มาตา ปุตฺตํ มียมานํ เอวํ ลภติ “อหํ มียามิ, มา เม ปุตฺโต มียี”ติ, ปุตฺโต วา ปน มาตรํ มียมานํ น เอวํ ลภติ “อหํ มียามิ, มา เม มาตา มียี”ติ. อิมานิ โข ภิกฺขเว ตีณิ อมาตาปุตฺติกานิ ภยานีติ.}

\prose{อตฺถิ ภิกฺขเว มคฺโค อตฺถิ ปฏิปทา, อิเมสญฺจ ติณฺณํ สมาตาปุตฺติกานํ ภยานํ, อิเมสญฺจ ติณฺณํ อมาตาปุตฺติกานํ ภยานํ ปหานาย สมติกฺกมาย สํวตฺตติ. กตโม จ ภิกฺขเว มคฺโค กตมา จ}

\vspace{\tipitakaparskip}
\csromancenter{(7) 2. มหาวคฺค ปฏิปทา, อิเมสญฺจ ติณฺณํ สมาตาปุตฺติกานํ ภยานํ, อิเมสญฺจ ติณฺณํ อมาตาปุตฺติกานํ ภยานํ ปหานาย สมติกฺกมาย สํวตฺตติ? อยเมว อริโย อฏฺฐงฺคิโก มคฺโค. เสยฺยถิทํ? สมฺมาทิฏฺฐิ สมฺมาสงฺกปฺโป สมฺมาวาจา สมฺมากมฺมนฺโต สมฺมาอาชีโว สมฺมาวายาโม สมฺมาสติ สมฺมาสมาธิ. อยํ โข ภิกฺขเว มคฺโค อยํ ปฏิปทา, อิเมสญฺจ ติณฺณํ สมาตาปุตฺติกานํ ภยานํ, อิเมสญฺจ ติณฺณํ อมาตาปุตฺติกานํ ภยานํ ปหานาย สมติกฺกมาย สํวตฺตตีติ.  ทุติยํ.}
\csromanheader{2}{3. เวนาคปุร\-สุตฺต}
\proseitem{64}{\csromanfolio{181}เอกํ สมยํ ภควา โกสเลสุ จาริกํ จรมาโน มหตา ภิกฺขุ\-สํเฆน สทฺธิํ เยน เวนาคปุรํ นาม โกสลานํ พฺราหฺมณคาโม ตทวสริ. อสฺโสสุํ โข เวนาคปุริกา พฺราหฺมณ\-คหปติกา “สมโณ ขลุ โภ โคตโม สกฺยปุตฺโต สกฺยกุลา ปพฺพชิโต เวนาคปุรํ อนุปฺปตฺโต, ตํ โข ปน ภวนฺตํ โคตมํ เอวํ กลฺยาโณ กิตฺติสทฺโท อพฺภุคฺคโต ‘อิติปิ โส ภควา อรหํ สมฺมาสมฺพุทฺโธ วิชฺชา\-จรณ\-สมฺปนฺโน สุคโต โลกวิทู อนุตฺตโร ปุริส\-ทมฺม\-สารถิ สตฺถา เทวมนุสฺสานํ พุทฺโธ ภควา’ติ\csromansharedfootnote{a6bd117a922df51e}{ภควา (สี, สฺยา, กํ, อิ) อิทํ สุตฺตวณฺณนาย อฏฺฐกถาย สํสนฺเทตพฺพํ. วิ 1. 1; ที 1. 82 ปิฏฺเฐสุปิ ปสฺสิตพฺพํ.}, โส อิมํ โลกํ สเทวกํ สมารกํ สพฺรหฺมกํ สสฺสมณ\-พฺราหฺมณิํ ปชํ สเทวมนุสฺสํ สยํ อภิญฺญา สจฺฉิกตฺวา ปเวเทติ, โส ธมฺมํ เทเสติ อาทิกลฺยาณํ มชฺเฌกลฺยาณํ ปริโยสาน\-กลฺยาณํ สาตฺถํ สพฺยญฺชนํ เกวล\-ปริปุณฺณํ ปริสุทฺธํ พฺรหฺมจริยํ ปกาเสติ. สาธุ โข ปน ตถารูปานํ อรหตํ ทสฺสนํ โหตี”ติ.}

\prose{อถ โข เวนาคปุริกา พฺราหฺมณ\-คหปติกา เยน ภควา เตนุ\-ปสงฺกมิํสุ, อุปสงฺกมิตฺวา อปฺเปกจฺเจ ภควนฺตํ อภิวาเทตฺวา เอกมนฺตํ นิสีทิํสุ, อปฺปกจฺเจ ภควตา สทฺธิํ สมฺโมทิํสุ, สมฺโมทนียํ กถํ สารณียํ วีติสาเรตฺวา เอกมนฺตํ นิสีทิํสุ, อปฺเปกจฺเจ เยน ภควา เตนญฺชลิํ ปณาเมตฺวา เอกมนฺตํ นิสีทิํสุ, อปฺเปกจฺเจ นามโคตฺตํ สาเวตฺวา เอกมนฺตํ นิสีทิํสุ, อปฺเปกจฺเจ ตุณฺหีภูตา เอกมนฺตํ นิสีทิํสุ, เอกมนฺตํ นิสินฺโน โข เวนาคปุริโก วจฺฉโคตฺโต พฺราหฺมโณ ภควนฺตํ เอตทโวจ\csromandash{}}

\prose{\csromanfolio{182}อจฺฉริยํ โภ โคตม, อพฺภุตํ โภ โคตม, ยาวญฺจิทํ โภโต โคตมสฺส วิปฺปสนฺนานิ อินฺทฺริยานิ, ปริสุทฺโธ ฉวิวณฺโณ ปริโยทาโต. เสยฺยถาปิ โภ โคตม สารทํ พทรปณฺฑุํ\csromansharedfootnote{698b1364b87f49b2}{มณฺฑํ (ก)} ปริสุทฺธํ โหติ ปริโยทาตํ. เอวเมวํ โภโต โคตมสฺส วิปฺปสนฺนานิ อินฺทฺริยานิ, ปริสุทฺโธ ฉวิวณฺโณ ปริโยทาโต. เสยฺยถาปิ โภ โคตม ตาลปกฺกํ สมฺปติ พนฺธนา ปมุตฺตํ\csromansharedfootnote{c2bcc01686816024}{มุตฺตํ (สี, อิ, ก)} ปริสุทฺธํ โหติ ปริโยทาตํ. เอวเมวํ โภโต โคตมสฺส วิปฺปสนฺนานิ อินฺทฺริยานิ, ปริสุทฺโธ ฉวิวณฺโณ ปริโยทาโต. เสยฺยถาปิ โภ โคตม เนกฺขํ\csromansharedfootnote{6999fa58bd4937e6}{นิกฺขํ อิติปิ (ม 3. 143 ปิฏฺเฐ)} ชมฺโพนทํ ทกฺข\-กมฺมารปุตฺต\-สุปริกมฺมกตํ อุกฺกามุเข สุกุสล\-สมฺปหฏฺฐํ ปณฺฑุกมฺพเล นิกฺขิตฺตํ ภาสเต จ ตปเต จ วิโรจติ จ. เอวเมวํ โภโต โคตมสฺส วิปฺปสนฺนานิ อินฺทฺริยานิ, ปริสุทฺโธ ฉวิวณฺโณ ปริโยทาโต. ยานิ ตานิ โภ โคตม อุจฺจาสยน\-มหาสยนานิ. เสยฺยถิทํ? อาสนฺทิ ปลฺลงฺโก โคนโก จิตฺตโก ปฏิกา ปฏลิกา ตูลิกา วิกติกา อุทฺทโลมี เอกนฺตโลมี กฏฺฏิสฺสํ โกเสยฺยํ กุตฺตกํ หตฺถตฺถรํ อสฺสตฺถรํ รถตฺถรํ อชินปฺปเวณี กทลิ\-มิค\-ปวร\-ปจฺจตฺถรณํ\csromansharedfootnote{de04de3bf90d5847}{กาทลิมิคปวรปจฺจตฺถรณํ (สี)} สอุตฺตรจฺฉทํ อุภโต\-โลหิตกู\-ปธานํ. เอวรูปานํ นูน ภวํ โคตโม อุจฺจาสยน\-มหาสยนานํ นิกามลาภี อกิจฺฉลาภี อกสิรลาภีติ.}

\prose{ยานิ โข ปน ตานิ พฺราหฺมณ อุจฺจาสยน\-มหาสยนานิ. เสยฺยถิทํ? อาสนฺทิ ปลฺลงฺโก โคนโก จิตฺตโก ปฏิกา ปฏลิกา ตูลิกา วิกติกา อุทฺทโลมี เอกนฺตโลมี กฏฺฏิสฺสํ โกเสยฺยํ กุตฺตกํ หตฺถตฺถรํ อสฺสตฺถรํ รถตฺถรํ อชินปฺปเวณี กทลิ\-มิค\-ปวร\-ปจฺจตฺถรณํ สอุตฺตรจฺฉทํ อุภโต\-โลหิตกู\-ปธานํ. ทุลฺลภานิ ตานิ ปพฺพชิตานํ, ลทฺธา จ ปน\csromansharedfootnote{43db1a429c7216a9}{ลทฺธานิ จ (สี, สฺยา, กํ), ลทฺธา จ (อิ)} น กปฺปนฺติ.}

\prose{ตีณิ โข อิมานิ พฺราหฺมณ อุจฺจาสยน\-มหาสยนานิ, เยสาหํ เอตรหิ นิกามลาภี อกิจฺฉลาภี อกสิรลาภี. กตมานิ ตีณิ? ทิพฺพํ อุจฺจาสยน\-มหาสยนํ, พฺรหฺมํ อุจฺจาสยน\-มหาสยนํ, อริยํ อุจฺจาสยน\-มหาสยนํ. อิมานิ โข พฺราหฺมณ ตีณิ อุจฺจาสยน\-มหาสยนานิ, เยสาหํ เอตรหิ นิกามลาภี อกิจฺฉลาภี อกสิรลาภีติ.}

\chapterheadpageread{(7) 2. มหาวคฺค}
\prose{\csromanfolio{183}กตมํ ปน ตํ โภ โคตม ทิพฺพํ อุจฺจาสยน\-มหาสยนํ, ยสฺส ภวํ โคตโม เอตรหิ นิกามลาภี อกิจฺฉลาภี อกสิรลาภีติ? อิธาหํ พฺราหฺมณ ยํ คามํ วา นิคมํ วา อุปนิสฺสาย วิหรามิ, โส ปุพฺพณฺหสมยํ นิวาเสตฺวา ปตฺต\-จีวรมา\-ทาย ตเมว คามํ วา นิคมํ วา ปิณฺฑาย ปวิสามิ. โส ปจฺฉาภตฺตํ ปิณฺฑปาต\-ปฏิกฺกนฺโต วนนฺตญฺเญว ปวิสามิ\csromansharedfootnote{b77d04305357c946}{ปจารยามิ (สี, สฺยา, กํ)}, โส ยเทว ตตฺถ โหนฺติ ติณานิ วา ปณฺณานิ วา, ตานิ เอกชฺฌํ สงฺฆริตฺวา นิสีทามิ ปลฺลงฺกํ อาภุชิตฺวา อุชุํ กายํ ปณิธาย ปริมุขํ สติํ อุปฏฺฐเปตฺวา. โส วิวิจฺเจว กาเมหิ วิวิจฺจ อกุสเลหิ ธมฺเมหิ สวิตกฺกํ สวิจารํ วิเวกชํ ปีติสุขํ ปฐมํ ฌานํ อุปสมฺปชฺช วิหรามิ. วิตกฺก\-วิจารานํ วูปสมา อชฺฌตฺตํ สมฺปสาทนํ เจตโส เอโกทิภาวํ อวิตกฺกํ อวิจารํ สมาธิชํ ปีติสุขํ ทุติยํ ฌานํ อุปสมฺปชฺช วิหรามิ. ปีติยา จ วิราคา อุเปกฺขโก จ วิหรามิ สโต จ สมฺปชาโน, สุขญฺจ กาเยน ปฏิสํเวเทมิ, ยํ ตํ อริยา อาจิกฺขนฺติ “อุเปกฺขโก สติมา สุขวิหารี”ติ ตติยํ ฌานํ อุปสมฺปชฺช วิหรามิ. สุขสฺส จ ปหานา ทุกฺขสฺส จ ปหานา ปุพฺเพว โสมนสฺส\-โทมนสฺสานํ อตฺถงฺคมา อทุกฺขมสุขํ อุเปกฺขา\-สติ\-ปาริสุทฺธิํ จตุตฺถํ ฌานํ อุปสมฺปชฺช วิหรามิ. โส เจ อหํ พฺราหฺมณ เอวํภูโต จงฺกมามิ, ทิพฺโพ เม เอโส ตสฺมิํ สมเย จงฺกโม โหติ. โส เจ อหํ พฺราหฺมณ เอวํภูโต ติฏฺฐามิ, ทิพฺพํ เม เอตํ ตสฺมิํ สมเย ฐานํ โหติ. โส เจ อหํ พฺราหฺมณ เอวํภูโต นิสีทามิ, ทิพฺพํ เม เอตํ ตสฺมิํ สมเย อาสนํ โหติ. โส เจ อหํ พฺราหฺมณ เอวํภูโต เสยฺยํ กปฺเปมิ, ทิพฺพํ เม เอตํ ตสฺมิํ สมเย อุจฺจาสยน\-มหาสยนํ โหติ. อิทํ โข พฺราหฺมณ ทิพฺพํ อุจฺจาสยน\-มหาสยนํ, ยสฺสาหํ เอตรหิ นิกามลาภี อกิจฺฉลาภี อกสิรลาภีติ.}

\prose{อจฺฉริยํ โภ โคตม, อพฺภุตํ โภ โคตม, โก จญฺโญ เอวรูปสฺส ทิพฺพสฺส อุจฺจาสยน\-มหาสยนสฺส นิกามลาภี ภวิสฺสติ อกิจฺฉลาภี อกสิรลาภี อญฺญตฺร โภตา โคตเมน.}

\prose{กตมํ ปน โภ โคตม พฺรหฺมํ อุจฺจาสยน\-มหาสยนํ, ยสฺส ภวํ โคตโม เอตรหิ นิกามลาภี อกิจฺฉลาภี อกสิรลาภีติ. อิธาหํ พฺราหฺมณ ยํ คามํ วา นิคมํ วา อุปนิสฺสาย วิหรามิ, โส}

\prose{\csromanfolio{184}ปุพฺพณฺหสมยํ นิวาเสตฺวา ปตฺต\-จีวรมา\-ทาย ตเมว คามํ วา นิคมํ วา ปิณฺฑาย ปวิสามิ, โส ปจฺฉาภตฺตํ ปิณฺฑปาต\-ปฏิกฺกนฺโต วนนฺตญฺเญว ปวิสามิ. โส ยเทว ตตฺถ โหนฺติ ติณานิ วา ปณฺณานิ วา, ตานิ เอกชฺฌํ สงฺฆริตฺวา นิสีทามิ ปลฺลงฺกํ อาภุชิตฺวา อุชุํ กายํ ปณิธาย ปริมุขํ สติํ อุปฏฺฐเปตฺวา. โส เมตฺตา\-สหคเตน เจตสา เอกํ ทิสํ ผริตฺวา วิหรามิ. ตถา ทุติยํ. ตถา ตติยํ. ตถา จตุตฺถํ. อิติ อุทฺธมโธ ติริยํ สพฺพธิ สพฺพตฺตตาย สพฺพาวนฺตํ โลกํ เมตฺตา\-สหคเตน เจตสา วิปุเลน มหคฺคเตน อปฺปมาเณน อเวเรน อพฺยาปชฺเชน\csromansharedfootnote{37bd4dcb51a235ae}{อพฺยาปชฺเฌน (สพฺพตฺถ)} ผริตฺวา วิหรามิ. กรุณา\-สหคเตน เจตสา ฯเปฯ มฺõทิตาสหคเตน เจตสา. อุเปกฺขา\-สหคเตน เจตสา เอกํ ทิสํ ผริตฺวา วิหรามิ. ตถา ทุติยํ. ตถา ตติยํ. ตถา จตุตฺถํ\csromansharedfootnote{dc62e42f90200188}{จตุตฺถิํ (สี)}. อิติ อุทฺธมโธ ติริยํ สพฺพธิ สพฺพตฺตตาย สพฺพาวนฺตํ โลกํ อุเปกฺขา\-สหคเตน เจตสา วิปุเลน มหคฺคเตน อปฺปฺqมาเณน อเวเรน อพฺยาปชฺเชน ผริตฺวา วิหรามฺé. โส เจ อหํ พฺราหฺมณ เอวํภูโต จงฺกมามิ, พฺราหฺมา \}เอ เอโส ตสฺมิํ สมเย จงฺกโม โหติ. โส เจ อหํ พฺราหฺมณ เอวํภูโต ติฏฺฐามิ ฯเปฯ นิสีทามิ ฯเปฯ เสยฺยํ กปฺเปมิ, พฺรหฺมํ เม เอตํ ตสฺมิํ สมเย อุจฺจาสยน\-มหาสยนํ โหติ. อิทํ โข พฺราหฺมณ พฺรหฺมํ อุจฺจาสยน\-มหาสยนํ, ยสฺสาหํ เอตรหิ นิกามลาภี อกิจฺฉาลาภี อกสิรลาภีติ.}

\prose{อจฺฉริยํ โภ โคตม, อพฺภุตํ โภ โคตม, โก จญฺโญ เอวรูปสฺส พฺรหฺมสฺส อุจฺจาสยน\-มหาสยนสฺส นิกามลาภี ภวิสฺสติ อกิจฺฉลาภี อกสิรลาภี อญฺญตฺร โภตา โคตเมน.}

\prose{กตมํ ปน ตํ โภ โคตม อริยํ อุจฺจาสยน\-มหาสยนํ, ยสฺส ภวํ โคตโม เอตรหิ นิกามลาภี อกิจฺฉลาภี อกสิรลาภีติ? อิธาหํ พฺราหฺมณ ยํ คามํ วา นิคมํ วา อุปนิสฺสาย วิหรามิ, โส ปุพฺพณฺหสมยํ นิวาเสตฺวา ปตฺต\-จีวรมา\-ทาย ตเมว คามํ วา นิคมํ วา ปิณฺฑาย ปวิสามิ, โส ปจฺฉาภตฺตํ ปิณฺฑปาต\-ปฏิกฺกนฺโต วนนฺตญฺเญว ปวิสามิ, โส ยเทว ตตฺถ โหนฺติ ติณานิ วา ปณฺณานิ วา, ตานิ เอกชฺฌํ สงฺฆริตฺวา นิสีทามิ ปลฺลงฺกํ อาภุชิตฺวา อุชุํ กายํ ปณิธาย ปริมุขํ สติํ อุปฏฺฐเปตฺวา. โส เอวํ ชานามิ “ราโค เม ปหีโน}

\vspace{\tipitakaparskip}
\csromancenter{(7) 2. มหาวคฺค อุจฺฉินฺนมูโล ตาลาวตฺถุกโต อนภาวํกโต อายติํ อนุปฺปาทธมฺโม, โทโส เม ปหีโน อุจฺฉินฺนมูโล ตาลาวตฺถุกโต อนภาวํกโต อายติํ อนุปฺปาทธมฺโม. โมโห เม ปหีโน อุจฺฉินฺนมูโล ตาลาวตฺตุกโต อนภาวํกโต อายติํ อนุปฺปาทธมฺโม”. โส เจ อหํ พฺราหฺมณ เอวํภูโต จงฺกมามิ, อริโย เม เอโส ตสฺมิํ สมเย จงฺกโม โหติ. โส เจ อหํ พฺราหฺมณ เอวํ ภูโต ติฏฺฐามิ ฯเปฯ นิสีทามิ ฯเปฯ เสยฺยํ กปฺเปมิ. อริยํ เม เอตํ ตสฺมิํ สมเย อุจฺจาสยน\-มหาสยนํ โหติ. อิทํ โข พฺราหฺมณ อริยํ อุจฺจาสยน\-มหาสยนํ, ยสฺสาหํ เอตรหิ นิกามลาภี อกิจฺฉลาภี อกสิรลาภีติ.}
\prose{\csromanfolio{185}อจฺฉริยํ โภ โคตม, อพฺภุตํ โภ โคตม, โก จญฺโญ เอวรูปสฺส อริยสฺส อุจฺจาสยน\-มหาสยนสฺส นิกามลาภี ภวิสฺสติ อกิจฺฉลาภี อกสิรลาภี อญฺญตฺร โภตา โคตเมน.}

\prose{อภิกฺกนฺตํ โภ โคตม, อภิกฺกนฺตํ โภ โคตม, เสยฺยถาปิ โภ โคตม นิกฺกุชฺชิตํ วา อุกฺกุชฺเชยฺย, ปฏิจฺฉนฺนํ วา วิวเรยฺย, มูฬฺหสฺส วา มคฺคํ อาจิกฺเขยฺย, อนฺธกาเร วา เตลปชฺโชตํ ธาเรยฺย “จกฺขุมนฺโต รูปานิ ทกฺขนฺตี”ติ, เอวเมวํ โข โภตา โคตเมน อเนก\-ปริยาเยน ธมฺโม ปกาสิโต, เอเต มยํ ภวนฺตํ โคตมํ สรณํ คจฺฉาม ธมฺมญฺจ ภิกฺขุสํฆญฺจ, อุปาสเก โน ภวํ โคตโม ธาเรตุ อชฺชตคฺเค ปาณุเปเต สรณํ คเตติ.  ตติยํ.}

\csromanheader{2}{4. \textbf{สรภสุตฺต}}
\proseitem{65}{เอวํ เม สุตํ\csromandash{}เอกํ สมยํ ภควา ราชคเห วิหรติ คิชฺฌกูเฏ ปพฺพเต. เตน โข ปน สมเยน สรโภ นาม ปริพฺพาชโก อจิรปกฺกนฺโต โหติ อิมสฺมา ธมฺมวินยา, โส ราชคเห ปริสติ\csromansharedfootnote{33039af806ccf922}{ปริสติํ (สี, อิ)} เอวํ วาจํ ภาสติ “อญฺญาโต มยา สมณานํ สกฺย\-ปุตฺติ\-กานํ ธมฺโม, อญฺญาย จ ปนาหํ สมณานํ สกฺย\-ปุตฺติ\-กานํ ธมฺมํ, เอวาหํ ตสฺมา ธมฺมวินยา อปกฺกนฺโต”ติ. อถ โข สมฺพหุลา ภิกฺขู ปุพฺพณฺหสมยํ นิวาเสตฺวา ปตฺต\-จีวรมา\-ทาย ราชคหํ ปิณฺฑาย ปวิสิํสุ. อสฺโสสุํ โข เต ภิกฺขู สรภสฺส ปริพฺพาชกสฺส ราชคเห ปริสติ เอวํ วาจํ \csromanfolio{186}ภาสมานสฺส “อญฺญาโต มยา สมณานํ สกฺย\-ปุตฺติ\-กานํ ธมฺโม, อญฺญาย จ ปนาหํ สมณานํ สกฺย\-ปุตฺติ\-กานํ ธมฺมํ, เอวาหํ ตสฺมา ธมฺมวินยา อปกฺกนฺโต”ติ.}

\prose{อถ โข เต ภิกฺขู ราชคเห ปิณฺฑาย จริตฺวา ปจฺฉาภตฺตํ ปิณฺฑปาต\-ปฏิกฺกนฺตา เยน ภควา เตนุ\-ปสงฺกมิํสุ, อุปสงฺกมิตฺวา ภควนฺตํ อภิวาเทตฺวา เอกมนฺตํ นิสีทิํสุ, เอกมนฺตํ นิสินฺนา โข เต ภิกฺขู ภควนฺตํ เอตทโวจุํ “สรโภ นาม ภนฺเต ปริพฺพาชโก อจิรปกฺกนฺโต อิมสฺมา ธมฺมวินยา, โส ราชคเห ปริสติ เอวํ วาจํ ภาสติ ‘อญฺญาโต มยา สมณานํ สกฺย\-ปุตฺติ\-กานํ ธมฺโม, อญฺญาย จ ปนาหํ สมณานํ สกฺย\-ปุตฺติ\-กานํ ธมฺมํ, เอวาหํ ตสฺมา ธมฺมวินยา อปกฺกนฺโต’ติ. สาธุ ภนฺเต ภควา เยน สิปฺปินิกาตีรํ\csromansharedfootnote{701a8d6acb2c57ba}{สปฺปินิกาตีรํ (สี, อิ), สปฺปินิยา ตีรํ (สฺยา, กํ)} ปริพฺพาชกา\-ราโม, เยน สรโภ ปริพฺพาชโก เตนุ\-ปสงฺกมตุ อนุกมฺปํ อุปาทายา”ติ. อธิวาเสสิ ภควา ตุณฺหีภาเวน.}

\prose{อถ โข ภควา สายนฺหสมยํ ปฏิสลฺลานา วุฏฺฐิโต เยน สิปฺปินิกาตีรํ ปริพฺพาชกา\-ราโม, เยน สรโภ ปริพฺพาชโก เตนุปสงฺกมิ, อุปสงฺกมิตฺวา ปญฺญตฺเต อาสเน นิสีทิ, นิสชฺช โข ภควา สรภํ ปริพฺพาชกํ เอตทโวจ “สจฺจํ กิร ตฺวํ สรภ เอวํ วเทสิ ‘อญฺญาโต มยา สมณานํ สกฺย\-ปุตฺติ\-กานํ ธมฺโม, อญฺญาย จ ปนาหํ สมณานํ สกฺย\-ปุตฺติ\-กานํ ธมฺมํ, เอวาหํ ตสฺมา ธมฺมวินยา อปกฺกนฺโต’ติ”. เอวํ วุตฺเต สรโภ ปริพฺพาชโก ตุณฺหี อโหสิ.}

\prose{ทุติยมฺปิ โข ภควา สรภํ ปริพฺพาชกํ เอตทโวจ “วเทหิ สรภ, กินฺติ เต อญฺญาโต สมณานํ สกฺย\-ปุตฺติ\-กานํ ธมฺโม, สเจ เต อปริปูรํ ภวิสฺสติ, อหํ ปริปูเรสฺสามิ. สเจ ปน เต ปริปูรํ ภวิสฺสติ, อหํ อนุโมทิสฺสามี”ติ. ทุติยมฺปิ โข สรโภ ปริพฺพาชโก ตุณฺหี อโหสิ.}

\prose{ตติยมฺปิ โข ภควา สรภํ ปริพฺพาชกํ เอตทโวจ (“โย\csromansharedfootnote{ee40768de16ae275}{มยา (สฺยา, กํ, อิ)} โข สรภ ปญฺญายติ สมณานํ สกฺย\-ปุตฺติ\-กานํ ธมฺโม)\csromansharedfootnote{d6166eef043da89d}{( ) สีโปตฺถเก นตฺถิ.} วเทหิ สรภ, กินฺติ เต อญฺญาโต สมณานํ สกฺย\-ปุตฺติ\-กานํ ธมฺโม, สเจ เต}

\vspace{\tipitakaparskip}
\csromancenter{(7) 2. มหาวคฺค อปริปูรํ ภวิสฺสติ, อหํ ปริปูเรสฺสามิ, สเจ ปน เต ปริปูรํ ภวิสฺสติ, อหํ อนุโมทิสฺสามี”ติ. ตติยมฺปิ โข สรโภ ปริพฺพาชโก ตุณฺหี อโหสิ.}
\prose{\csromanfolio{187}อถ โข เต ปริพฺพาชกา สรภํ ปริพฺพาชกํ เอตทโวจุํ “ยเทว โข ตฺวํ อาวุโส สรภ สมณํ โคตมํ ยาเจยฺยาสิ, ตเทว เต สมโณ โคตโม ปวาเรติ. วเทหาวุโส สรภ, กินฺติ เต อญฺญาโต สมณานํ สกฺย\-ปุตฺติ\-กานํ ธมฺโม, สเจ เต อปริปูรํ ภวิสฺสติ, สมโณ โคตโม ปริปูเรสฺสติ. สเจ ปน เต ปริปูรํ ภวิสฺสติ, สมโณ โคตโม อนุโมทิสฺสตี”ติ. เอวํ วุตฺเต สรโภ ปริพฺพาชโก ตุณฺหีภูโต มงฺกุภูโต ปตฺตกฺขนฺโธ อโธมุโข ปชฺฌายนฺโต อปฺปฏิภาโน นิสีทิ.}

\prose{อถ โข ภควา สรภํ ปริพฺพาชกํ ตุณฺหีภูตํ มงฺกุภูตํ ปตฺตกฺขนฺทํ อโธมุขํ ปชฺฌายนฺตํ อปฺปฏิภานํ วิทิตฺวา เต ปริพฺพาชเก เอตทโวจ\csromandash{}}

\prose{โย โข มํ ปริพฺพาชกา\csromansharedfootnote{687a856cace06ee4}{ปริพฺพาชโก (อิ, ก)} เอวํ วเทยฺย “สมฺมา\-สมฺพุทฺธสฺส เต ปฏิชานโต อิเม ธมฺมา อนภิสมฺพุทฺธา”ติ. ตมหํ ตตฺถ สาธุกํ สมนุยุญฺเชยฺยํ สมนุคาเหยฺยํ สมนุภาเสยฺยํ, โส วต มยา สาธุกํ สมนุ\-ยุญฺชิย\-มาโน สมนุ\-คาหิย\-มาโน สมนุ\-ภาสิย\-มาโน อฏฺฐานเมตํ อนวกาโส ยํ โส ติณฺณํ ฐานานํ นาญฺญตรํ\csromansharedfootnote{7c17f49c10eee0bb}{อญฺญตรํ (ก)} ฐานํ นิคจฺเฉยฺย\csromandash{}อญฺเญน วา อญฺญํ ปฏิจริสฺสติ, พหิทฺธา กถํ อปนาเมสฺสติ, โกปญฺจ โทสญฺจ อปฺปจฺจยญฺจ ปาตุกริสฺสติ. ตุณฺหีภูโต มงฺกุภูโต\csromansharedfootnote{b7afbab80d423073}{ตุณฺหีภูโต วา มงฺกุภูโต (สี, สฺยา, กํ), ตุณฺหีภูโต วา มงฺกุภูโต วา (อิ)} ปตฺตกฺขนฺโธ อโธมุโข ปชฺฌายนฺโต อปฺปฏิภาโน นิสีทิสฺสติ, เสยฺยถาปิ สรโภ ปริพฺพาชโก.}

\prose{โย โข มํ ปริพฺพาชกา เอวํ วเทยฺย “ขีณาสวสฺส เต ปฏิชานโต อิเม อาสวา อปริกฺขีณา”ติ. ตมหํ ตตฺถ สาธุกํ สมนุยุญฺเชยฺยํ สมนุคาเหยฺยํ สมนุภาเสยฺยํ, โส วต มยา สาธุกํ สมนุ\-ยุญฺชิย\-มาโน สมนุ\-คาหิย\-มาโน สมนุ\-ภาสิย\-มาโน อฏฺฐานเมตํ อนวกาโส ยํ โส ติณฺณํ ฐานานํ นาญฺญตรํ ฐานํ นิคจฺเฉยฺย\csromandash{} อญฺเญน วา อญฺญํ ปฏิจริสฺสติ, พหิทฺธา กถํ อปนาเมสฺสติ, โกปญฺจ โทสญฺจ \csromanfolio{188}อปฺปจฺจยญฺจ ปาตุกริสฺสติ. ตุณฺหีภูโต มงฺกุภูโต ปตฺตกฺขนฺโธ อโธมุโข ปชฺฌายนฺโต อปฺปฏิภาโน นิสีทิสฺสติ, เสยฺยถาปิ สรโภ ปริพฺพาชโก.}

\prose{โย โข มํ ปริพฺพาชกา เอวํ วเทยฺย “ยสฺส โข ปน เต อตฺถาย ธมฺโม เทสิโต, โส น นิยฺยาติ ตกฺกรสฺส สมฺมา ทุกฺขกฺขยายา”ติ. ตมหํ ตตฺถ สาธุกํ สมนุยุญฺเชยฺยํ สมนุคาเหยฺยํ สมนุภาเสยฺยํ, โส วต มยา สาธุกํ สมนุ\-ยุญฺชิย\-มาโน สมนุ\-คาหิย\-มาโน สมนุ\-ภาสิย\-มาโน อฏฺฐานเมตํ อนวกาโส ยํ โส ติณฺณํ ฐานานํ นาญฺญตรํ ฐานํ นิคจฺเฉยฺย\csromandash{}อญฺเญน วา อญฺญํ ปฏิจริสฺสติ, พหิทฺธา กถํ อปนาเมสฺสติ, โกปญฺจ โทสญฺจ อปฺปจฺจยญฺจ ปาตุกริสฺสติ. ตุณฺหีภูโต มงฺกุภูโต ปตฺตกฺขนฺโธ อโธมุโข ปชฺฌายนฺโต อปฺปฏิภาโน นิสีทิสฺสติ, เสยฺยถาปิ สรโภ ปริพฺพาชโกติ. อถ โข ภควา สิปฺปินิกาตีเร ปริพฺพาชกา\-ราเม ติกฺขตฺตุํ สีหนาทํ นทิตฺวา เวหาสํ ปกฺกามิ.}

\prose{อถ โข เต ปริพฺพาชกา อจิร\-ปกฺกนฺตสฺส ภควโต สรภํ ปริพฺพาชกํ สมนฺตโต วาจาย\-สนฺนิโตทเกน\csromansharedfootnote{200f6ec4f3058610}{วาจาสตฺติโตทเกน (สี) 2. สญฺฌมฺภริมกํสุ (อิ, ที 1. 175 ปิฏฺเฐปิ) สํ 1. 469 ปิฏฺเฐ ปน อุปริปาโฐ วิย. 3. เสคาลกํเยว (สี, สฺยา, กํ, อิ)} สญฺชมฺภริมกํสุ\csromansharedfootnote{2f485c8925d39cca}{เภทณฺฑกํ (ก)}. “เสยฺยถาปิ อาวุโส สรภ พฺรหารญฺเญ ชรสิงฺคาโล ‘สีหนาทํ นทิสฺสามี’ติ สิงฺคาลกํเยว\csromansharedfootnote{b7deaa7a53d902c0}{อมฺพกมทฺทรี (สี) 6. ผุสฺสกรวิตํ (สี), ปุสฺสกรวิตํ (สฺยา, กํ, อิ)} นทติ, เภรณฺฑ\-กํเยว นทติ\csromansharedfootnote{9d12109420ad765d}{สญฺฌมฺภริมกํสุ (อิ, ที 1. 175 ปิฏฺเฐปิ) สํ 1. 469 ปิฏฺเฐ ปน อุปริปาโฐ วิย.}. เอวเมว โข ตฺวํ อาวุโส สรภ อญฺญเตฺรว สมเณน โคตเมน ‘สีหนาทํ นทิสฺสามี’ติ สิงฺคาลกํเยว นทสิ เภรณฺฑ\-กํเยว นทสิ. เสยฺยถาปิ อาวุโส สรภ อมฺพุกสญฺจรี\csromansharedfootnote{56ba8899d5de8c64}{เสคาลกํเยว (สี, สฺยา, กํ, อิ)} ‘ปุริสก\-รวิตํ\csromansharedfootnote{7efc991b9d9e6ce4}{ผุสฺสกรวิตํ (สี), ปุสฺสกรวิตํ (สฺยา, กํ, อิ)} รวิสฺสามี’ติ อมฺพุกสญฺจริ\-รวิตํเยว รวติ. เอวเมวํ โข ตฺวํ อาวุโส สรภ อญฺญเตฺรว สมเณน โคตเมน ‘ปุริสก\-รวิตํ รวิสฺสามี’ติ อมฺพุกสญฺจริ\-รวิตํเยว รวสิ. เสยฺยถาปิ อาวุโส สรภ อุสโภ สุญฺญาย โคสาลาย คมฺภีรํ นทิตพฺพํ มญฺญติ. เอวเมวํ โข ตฺวํ อาวุโส สรภ อญฺญเตฺรว สมเณน โคตเมน คมฺภีรํ นทิตพฺพํ มญฺญสีติ. อถ โข เต ปริพฺพาชกา สรภํ}

\vspace{\tipitakaparskip}
\csromancenter{(7) 2. มหาวคฺค ปริพฺพาชกํ สมนฺตโต วาจาย\-สนฺนิโตทเกน สญฺชมฺภริมกํสูติ.  จตุตฺถํ.}
\csromanheader{2}{5. \textbf{เกสมุตฺติสุตฺต}}
\proseitem{66}{\csromanfolio{189}เอวํ เม สุตํ\csromandash{}เอกํ สมยํ ภควา โกสเลสุ จาริกํ จรมาโน มหตา ภิกฺขุ\-สํเฆน สทฺธิํ เยน เกสมุตฺตํ\csromansharedfootnote{b9591b4565f30ca6}{เกสปุตฺตํ (สี, สฺยา, กํ, อิ)} นาม กาลามานํ นิคโม ตทวสริ. อสฺโสสุํ โข เกสมุตฺติยา กาลามา “สมโณ ขลุ โภ โคตโม สกฺยปุตฺโต สกฺยกุลา ปพฺพชิโต เกสมุตฺตํ อนุปฺปตฺโต, ตํ โข ปน ภวนฺตํ โคตมํ เอวํ กลฺยาโณ กิตฺติสทฺโท อพฺภุคฺคโต อิติปิ โส ภควา ฯเปฯ. สาธุ โข ปน ตถารูปานํ อรหตํ ทสฺสนํ โหตี”ติ.}

\prose{อถ โข เกสมุตฺติยา กาลามา เยน ภควา เตนุ\-ปสงฺกมิํสุ, อุปสงฺกมิตฺวา อปฺเปกจฺเจ ภควนฺตํ อภิวาเทตฺวา เอกมนฺตํ นิสีทิํสุ, อปฺเปกจฺเจ ภควตา สทฺธิํ สมฺโมทิํสุ, สมฺโมทนียํ กถํ สารณียํ วีติสาเรตฺวา เอกมนฺตํ นิสีทิํสุ, อปฺเปกจฺเจ เยน ภควา เตนญฺชลิํ ปณาเมตฺวา เอกมนฺตํ นิสีทิํสุ, อปฺเปกจฺเจ นามโคตฺตํ สาเวตฺวา เอกมนฺตํ นิสีทิํสุ, อปฺเปกจฺเจ ตุณฺหีภูตา เอกมนฺตํ นิสีทิํสุ, เอกมนฺตํ นิสินฺนา โข เต เกสมุตฺติยา กาลามา ภควนฺตํ เอตทโวจุํ\csromandash{}}

\prose{สนฺติ ภนฺเต เอเก สมณพฺราหฺมณา เกสมุตฺตํ อาคจฺฉนฺติ, เต สกํเยว วาทํ ทีเปนฺติ โชเตนฺติ, ปรปฺปวาทํ ปน ขุํเสนฺติ วมฺเภนฺติ ปริภวนฺติ, โอมกฺขิํ\csromansharedfootnote{b2eddba6284202cc}{โอปปกฺขิํ (สี, สฺยา, กํ, อิ), โอมกฺขิกํ (ก)} กโรนฺติ. อปิเรปิ ภนฺเต เอเก สมณพฺราหฺมณา เกสมุตฺตํ อาคจฺฉนฺติ, เตปิ สกํเยว วาทํ ทีเปนฺติ โชเตนฺติ, ปรปฺปวาทํ ปน ขุํเสนฺติ วมฺเภนฺติ ปริภวนฺติ, โอมกฺขิํ กโรนฺติ. เตสํ โน ภนฺเต อมฺหากํ โหเตว กงฺขา, โหติ วิจิกิจฺฉา “โก สุ นาม อิเมสํ ภวตํ สมณ\-พฺราหฺมณานํ สจฺจํ อาห, โก มุสา”ติ. อลํ หิ โว กาลามา กงฺขิตุํ อลํ วิจิกิจฺฉิตุํ, กงฺขนีเยว ปน\csromansharedfootnote{7bf63a38f5cc9a72}{กงฺขนีเย จ ปน (สํ 2. 527, 564 ปิฏฺเฐสุ)} โว ฐาเน วิจิกิจฺฉา อุปฺปนฺนา.}

\prose{เอถ ตุเมฺห กาลามา, มา อนุสฺสเวน, มา ปรมฺปราย, มา อิติกิราย, มา ปิฏก\-สมฺปทาเนน, มา ตกฺกเหตุ, มา นยเหตุ, มา \csromanfolio{190}อาการ\-ปริวิตกฺเกน, มา ทิฏฺฐิ\-นิชฺฌาน\-กฺขนฺติยา, มา ภพฺพรูปตาย, มา สมโณ โน ครูติ. ยทา ตุเมฺห กาลามา อตฺตนาว ชาเนยฺยาถ “อิเม ธมฺมา อกุสลา, อิเม ธมฺมา สาวชฺชา, อิเม ธมฺมา วิญฺญุครหิตา, อิเม ธมฺมา สมตฺตา สมาทินฺนา\csromansharedfootnote{e2b24e09269cf714}{สมาทิณฺณา (ก)} อหิตาย ทุกฺขาย สํวตฺตนฺตี”ติ. อถ ตุเมฺห กาลามา ปชเหยฺยาถ.}

\prose{ตํ กิํ มญฺญถ กาลามา, โลโภ ปุริสสฺส อชฺฌตฺถํ อุปฺปชฺชมาโน อุปฺปชฺชติ หิตาย วา อหิตาย วาติ. อหิตาย ภนฺเต. ลุทฺโธ ปนายํ กาลามา ปุริสปุคฺคโล โลเภน อภิภูโต ปริยาทินฺน\-จิตฺโต ปาณมฺปิ หนติ, อทินฺนมฺปิ อาทิยติ, ปรทารมฺปิ คจฺฉติ, มุสาปิ ภณติ, ปรมฺปิ ตถตฺตาย\csromansharedfootnote{6fb7c822656a82ef}{ตทตฺถาย (ก)} สมาทเปติ, ยํ’ส\csromansharedfootnote{4ab508d8b57c9bf2}{ยํ ตสฺส (ก) อนนฺตรสุตฺเต ปน “ยํ’ส” อิเตฺวว สพฺพตฺถปิ ทิสฺสติ.} โหติ ทีฆรตฺตํ อหิตาย ทุกฺขายาติ. เอวํ ภนฺเต.}

\prose{ตํ กิํ มญฺญถ กาลามา, โทโส ปุริสสฺส อชฺฌตฺตํ อุปฺปชฺชมาโน อุปฺปชฺชติ หิตาย วา อหิตาย วาติ. อหิตาย ภนฺเต. ทุฏฺโฐ ปนายํ กาลามา ปุริสปุคฺคโล โทเสน อภิภูโต ปริยาทินฺน\-จิตฺโต ปาณมฺปิ หนติ\csromansharedfootnote{082999f5ff776188}{หนฺติ (สี, อิ)}, อทินฺนมฺปิ อาทิยติ, ปรทารมฺปิ คจฺฉติ, มุสาปิ ภณติ, ปรมฺปิ ตถตฺตาย สมาทเปติ, ยํ’ส โหติ ทีฆรตฺตํ อหิตาย ทุกฺขายาติ. เอวํ ภนฺเต.}

\prose{ตํ กิํ มญฺญถ กาลามา, โมโห ปุริสสฺส อชฺฌตฺตํ อุปฺปชฺชมาโน อุปฺปชฺชติ หิตาย วา อหิตาย วาติ. อหิตาย ภนฺเต. มูโฬฺห ปนายํ กาลามา ปุริสปุคฺคโล โมเหน อภิภูโต ปริยาทินฺน\-จิตฺโต ปาณมฺปิ หนติ, อทินฺนมฺปิ อาทิยติ, ปรทารมฺปิ คจฺฉติ, มุสาปิ ภณติ, ปรมฺปิ ตถตฺตาย สมาทเปติ, ยํ’ส โหติ ทีฆรตฺตํ อหิตาย ทุกฺขายาติ. เอวํ ภนฺเต.}

\prose{ตํ กิํ มญฺญถ กาลามา, อิเม ธมฺมา กุสลา วา อกุสลา วาติ. อกุสลา ภนฺเต. สาวชฺชา วา อนวชฺชา วาติ. สาวชฺชา ภนฺเต. วิญฺญุครหิตา วา วิญฺญุปฺปสตฺถา วาติ. วิญฺญุครหิตา ภนฺเต. สมตฺตา สมาทินฺนา อหิตาย ทุกฺขาย สํวตฺตนฺติ โน วา, กถํ วา\csromansharedfootnote{c1d44201fe973c73}{กถํ วา โว (?)} เอตฺถ}

\vspace{\tipitakaparskip}
\csromancenter{(7) 2. มหาวคฺค โหตีติ. สมตฺตา ภนฺเต สมาทินฺนา อหิตาย ทุกฺขาย สํวตฺตนฺตีติ, เอวํ โน เอตฺถ โหตีติ.}
\prose{\csromanfolio{191}อิติ โข กาลามา ยํ ตํ อโวจุมฺหา\csromansharedfootnote{7981be099e9b76ac}{อาโวจุมฺห (สี, สฺยา, กํ, อิ) อุปริ 512 ปิฏฺเฐปิ.} “เอถ ตุเมฺห กาลามา, มา อนุสฺสเวน, มา ปรมฺปราย, มา อิติกิราย, มา ปิฏก\-สมฺปทาเนน, มา ตกฺกเหตุ, มา นยเหตุ, มา อาการ\-ปริวิตกฺเกน, มา ทิฏฺฐิ\-นิชฺฌาน\-กฺขนฺติยา, มา ภพฺพรูปตาย, มา สมโณ โน ครูติ. ยทา ตุเมฺห กาลามา อตฺตนาว ชาเนยฺยาถ ‘อิเม ธมฺมา อกุสลา, อิเม ธมฺมา สาวชฺชา, อิเม ธมฺมา วิญฺญุครหิตา, อิเม ธมฺมา สมตฺตา สมาทินฺนา อหิตาย ทุกฺขาย สํวตฺตนฺตี’ติ. อถ ตุเมฺห กาลามา ปชเหยฺยาถา”ติ อิติ ยํ ตํ วุตฺตํ, อิทเมตํ ปฏิจฺจ วุตฺตํ.}

\prose{เอถ ตุเมฺห กาลามา, มา อนุสฺสเวน, มา ปรมฺปราย, มา อิติกิราย, มา ปิฏก\-สมฺปทาเนน, มา ตกฺกเหตุ, มา นยเหตุ, มา อาการ\-ปริวิตกฺเกน, มา ทิฏฺฐิ\-นิชฺฌาน\-กฺขนฺติยา, มา ภพฺพรูปตาย, มา สมโณ โน ครูติ. ยทา ตุเมฺห กาลามา อตฺตนาว ชาเนยฺยาถ “อิเม ธมฺมา กุสลา, อิเม ธมฺมา อนวชฺชา, อิเม ธมฺมา วิญฺญุปฺปสตฺถา, อิเม ธมฺมา สมตฺตา สมาทินฺนา หิตาย สุขาย สํวตฺตนฺตี”ติ. อถ ตุเมฺห กาลามา อุปสมฺปชฺช วิหเรยฺยาถ.}

\prose{ตํ กิํ มญฺญถ กาลามา, อโลโภ ปุริสสฺส อชฺฌตฺตํ อุปฺปชฺชมาโน อุปฺปชฺชติ หิตาย วา อหิตาย วาติ. หิตาย ภนฺเต. อลุทฺโธ ปนายํ กาลามา ปุริสปุคฺคโล โลเภน อนภิภูโต อปริยาทินฺนจิตฺโต เนว ปาณํ หนติ, น อทินฺนํ อาทิยติ, น ปรทารํ คจฺฉติ, น มุสา ภณติ, น ปรมฺปิ ตถตฺตาย สมาทเปติ, ยํ’ส โหติ ทีฆรตฺตํ หิตาย สุขายาติ. เอวํ ภนฺเต.}

\prose{ตํ กิํ มญฺญถ กาลามา, อโทโส ปุริสสฺส อชฺฌตฺตํ อุปฺปชฺชมาโน อุปฺปชฺชติ ฯเปฯ อโมโห ปุริสสฺส อชฺฌตฺตํ อุปฺปชฺชมาโน อุปฺปชฺชติ ฯเปฯ หิตาย สุขายาติ. เอวํ ภนฺเต.}

\prosewithcloser{ตํ กิํ มญฺญถ กาลามา, อิเม ธมฺมา กุสลา วา อกุสลา วาติ. กุสลา ภนฺเต. สาวชฺชา วา อนวชฺชา วาติ. อนวชฺชา ภนฺเต. วิญฺญุครหิตา วา วิญฺญุปฺปสตฺถา วาติ. วิญฺญุปฺปสตฺถา ภนฺเต.}{สมตฺตา}
\prosecont{\csromanfolio{192}สมาทินฺนา หิตาย สุขาย สํวตฺตนฺติ โน วา, กถํ วา เอตฺถ โหตีติ. สมตฺตา ภนฺเต สมาทินฺนา หิตาย สุขาย สํวตฺตนฺติ, เอวํ โน เอตฺถ โหตีติ.}

\prose{อิติ โข กาลามา ยํ ตํ อโวจุมฺหา “เอถ ตุเมฺห กาลามา, มา อนุสฺสเวน, มา ปรมฺปราย, มา อิติกิราย, มา ปิฏก\-สมฺปทาเนน, มา ตกฺกเหตุ, มา นยเหตุ, มา อาการ\-ปริวิตกฺเกน, มา ทิฏฺฐิ\-นิชฺฌาน\-กฺขนฺติยา, มา ภพฺพรูปตาย, มา สมโณ โน ครูติ. ยทา ตุเมฺห กาลามา อตฺตนาว ชาเนยฺยาถ ‘อิเม ธมฺมา กุสลา, อิเม ธมฺมา อนวชฺชา, อิเม ธมฺมา วิญฺญุปฺปสตฺถา, อิเม ธมฺมา สมตฺตา สมาทินฺนา หิตาย สุขาย สํวตฺตนฺตี’ติ. อถ ตุเมฺห กาลามา อุปสมฺปชฺช วิหเรยฺยาถา”ติ อิติ ยํ ตํ วุตฺตํ, อิทเมตํ ปฏิจฺจ วุตฺตํ.}

\prose{ส โข โส\csromansharedfootnote{9bb916d9c94e8344}{โย โข (ก)} กาลามา อริยสาวโก เอวํ วิคตาภิชฺโช วิคตพฺยาปาโท อสมฺมูโฬฺห สมฺปชาโน ปติสฺสโต\csromansharedfootnote{63871ae55b4e25d9}{สโต (ก)} เมตฺตา\-สหคเตน เจตสา เอกํ ทิสํ ผริตฺวา วิหรติ. ตถา ทุติยํ. ตถา ตติยํ. ตถา จตุตฺถํ. อิติ อุทฺธมโธ ติริยํ สพฺพธิ สพฺพตฺตตาย สพฺพาวนฺตํ โลกํ เมตฺตา\-สหคเตน เจตสา วิปุเลน มหคฺคเตน อปฺปมาเณน อเวเรน อพฺยาปชฺเชน ผริตฺวา วิหรติ. กรุณา\-สหคเตน เจตสา ฯเปฯ มุทิตา\-สหคเตน เจตสา ฯเปฯ อุเปกฺขา\-สหคเตน เจตสา เอกํ ทิสํ ผริตฺวา วิหรติ. ตถา ทุติยํ. ตถา ตติยํ. ตถา จตุตฺถํ. อิติ อุทฺธมโธ ติริยํ สพฺพธิ สพฺพตฺตตาย สพฺพาวนฺตํ โลกํ อุเปกฺขา\-สหคเตน เจตสา วิปุเลน มหคฺคเตน อปฺปมาเณน อเวเรน อพฺยาปชฺชน ผริตฺวา วิหรติ.}

\prose{ส\csromansharedfootnote{3b180d75344e4b4b}{สเจ (ก)} โข โส กาลามา อริยสาวโก เอวํ อเวรจิตฺโต เอวํ อพฺยาปชฺชจิตฺโต เอวํ อสํกิลิฏฺฐ\-จิตฺโต เอวํ วิสุทฺธจิตฺโต, ตสฺส ทิฏฺเฐว ธมฺเม จตฺตาโร อสฺสาสา อธิคตา โหนฺติ. สเจ โข ปน อตฺถิ ปโร โลโก, อตฺถิ สุกต\-ทุกฺกฏานํ\csromansharedfootnote{b38fa52368346108}{สุกฏทุกฺกฏานํ (สี, สฺยา, กํ, อิ)} กมฺมานํ ผลํ วิปาโก. อถาหํ\csromansharedfootnote{97b9b8a47b260e84}{ฐานมหํ (สี, อิ), ฐานเมตํ เยนาหํ (สฺยา, กํ)} กายสฺส เภทา ปรํ มรณา สุคติํ สคฺคํ โลกํ อุปปชฺชิสฺสามีติ. อยมสฺส ปฐโม อสฺสาโส อธิคโต โหติ.}

\chapterheadpageread{(7) 2. มหาวคฺค}
\prose{\csromanfolio{193}สเจ โข ปน นตฺถิ ปโร โลโก, นตฺถิ สุกต\-ทุกฺกฏานํ กมฺมานํ ผลํ วิปาโก. อถาหํ\csromansharedfootnote{6cd8fc51b898ea1d}{อิธาหํ (สี, สฺยา, กํ, อิ)} ทิฏฺเฐว ธมฺเม อเวรํ อพฺยาปชฺชํ อนีฆํ สุขิํ\csromansharedfootnote{0f7d45ec5bcbdc7c}{สุขํ (สี), สุขี (สฺยา, กํ)} อตฺตานํ ปติหรามีติ. อยมสฺส ทุติโย อสฺสาโส อธิคโต โหติ.}

\prose{สเจ โข ปน กโรโต กรียติ ปาปํ, น โข ปนาหํ กสฺสจิ ปาปํ เจเตมิ, อกโรนฺตํ โข ปน มํ ปาปกมฺมํ กุโต ทุกฺขํ ผุสิสฺสตีติ. อยมสฺส ตติโย อสฺสาโส อธิคโต โหติ.}

\prose{สเจ โข ปน กโรโต น กรียติ ปาปํ. อถาหํ อุภเยเนว วิสุทฺธํ อตฺตานํ สมนุปสฺสามีติ. อยมสฺส จตุตฺโถ อสฺสาโส อธิคโต โหติ.}

\prose{ส โข โส กาลามา อริยสาวโก เอวํ อเวรจิตฺโต เอวํ อพฺยาปชฺชจิตฺโต เอวํ อสํกิลิฏฺฐ\-จิตฺโต เอวํ วิสุทฺธจิตฺโต, ตสฺส ทิฏฺเฐว ธมฺเม อิเม จตฺตาโร อสฺสาสา อธิคตา โหนฺตีติ.}

\prose{เอวเมตํ ภควา, เอวเมตํ สุคต, ส โข โส ภนฺเต อริยสาวโก เอวํ อเวรจิตฺโต เอวํ อพฺยาปชฺชจิตฺโต เอวํ อสํกิลิฏฺฐ\-จิตฺโต เอวํ วิสุทฺธจิตฺโต, ตสฺส ทิฏฺเฐว ธมฺเม จตฺตาโร อสฺสาสา อธิคตา โหนฺติ. สเจ โข ปน อตฺถิ ปโร โลโก, อตฺถิ สุกต\-ทุกฺกฏานํ กมฺมานํ ผลํ วิปาโก. อถาหํ กายสฺส เภทา ปรํ มรณา สุคติํ สคฺคํ โลกํ อุปปชฺชิสฺสามีติ. อยมสฺส ปฐโม อสฺสาโส อธิคโต โหติ.}

\prose{สเจ โข ปน นตฺถิ ปโร โลโก, นตฺถิ สุกต\-ทุกฺกฏานํ กมฺมานํ ผลํ วิปาโก. อถาหํ ทิฏฺเฐว ธมฺเม อเวรํ อพฺยาปชฺชํ อนีฆํ สุขิํ อตฺตานํ ปริหรามีติ. อยมสฺส ทุติโย อสฺสาโส อธิคโต โหติ.}

\prose{สเจ โข ปน กโรโต กรียติ ปาปํ, น โข ปนาหํ กสฺสจิ ปาปํ เจเตมิ. อกโรนฺตํ โข ปน มํ ปาปกมฺมํ กุโต ทุกฺขํ ผุสิสฺสตีติ. อยมสฺส ตติโย อสฺสาโส อธิคโต โหติ.}

\prose{\csromanfolio{194}สเจ โข ปน กโรโต น กรียติ ปาปํ. อถาหํ อุภเยเนว วิสุทฺธํ อตฺตานํ สมนุปสฺสามีติ. อยมสฺส จตุตฺโถ อสฺสาโส อธิคโต โหติ.}

\prose{ส โข โส ภนฺเต อริยสาวโก เอวํ อเวรจิตฺโต เอวํ อพฺยาปชฺชจิตฺโต เอวํ อสํกิลิฏฺฐ\-จิตฺโต เอวํ วิสุทฺธจิตฺโต, ตสฺส ทิฏฺเฐว ธมฺเม อิเม จตฺตาโร อสฺสาสา อธิคตา โหนฺติ.}

\prose{อภิกฺกนฺตํ ภนฺเต ฯเปฯ เอเต มยํ ภนฺเต ภควนฺตํ สรณํ คจฺฉาม ธมฺมญฺจ ภิกฺขุสํฆญฺจ, อุปาสเก โน ภนฺเต ภควา ธาเรตุ อชฺชตคฺเค ปาณุเปเต สรณํ คเตติ.  ปญฺจมํ.}

\csromanheader{2}{6. \textbf{สาฬฺหสุตฺต}}
\proseitem{67}{เอวํ เม สุตํ\csromandash{}เอกํ สมยํ อายสฺมา นนฺทโก สาวตฺถิยํ วิหรติ ปุพฺพาราเม มิคาร\-มาตุ\-ปาสาเท. อถ โข สาโฬฺห จ มิคารนตฺตา, สาโณ จ เสขุนิยนตฺตา\csromansharedfootnote{5bfabfe4a57881cf}{โรหโณ จ เปขุณิยนตฺตา (สี, สฺยา, กํ, อิ)} เยนายสฺมา นนฺทโก เตนุ\-ปสงฺกมิํสุ, อุปสงฺกมิตฺวา อายสฺมนฺตํ นนฺทกํ อภิวาเทตฺวา เอกมนฺตํ นิสีทิํสุ, เอกมนฺตํ นิสินฺนํ โข สาฬฺหํ มิคารนตฺตารํ อายสฺมา นนฺทโก เอตทโวจ\csromandash{}}

\prose{เอถ ตุเมฺห สาฬฺหา, มา อนุสฺสเวน, มา ปรมฺปราย, มา อิติกิราย, มา ปิฏก\-สมฺปทาเนน, มา ตกฺกเหตุ, มา นยเหตุ, มา อาการ\-ปริวิตกฺเกน, มา ทิฏฺฐิ\-นิชฺฌาน\-กฺขนฺติยา, มา ภพฺพรูปตาย, มา สมโณ โน ครูติ. ยทา ตุเมฺห สาฬฺหา อตฺตนาว ชาเนยฺยาถ “อิเม ธมฺมา อกุสลา, อิเม ธมฺมา สาวชฺชา, อิเม ธมฺมา วิญฺญุครหิตา, อิเม ธมฺมา สมตฺตา สมาทินฺนา อหิตาย ทุกฺขาย สํวตฺตนฺตี”ติ. อถ ตุเมฺห สาฬฺหา ปชเหยฺยาถ.}

\prose{ตํ กิํ มญฺญถ สาฬฺหา, อตฺถิ โลโภติ. เอวํ ภนฺเต. “อภิชฺฌา”ติ โข อหํ สาฬฺหา เอตมตฺถํ วทามิ, ลุทฺโธ โข อยํ สาฬฺหา อภิชฺฌาลุ ปาณมฺปิ หนติ, อทินฺนมฺปิ อาทิยติ, ปรทารมฺปิ คจฺฉติ, มุสาปิ ภณติ, ปรมฺปิ ตถตฺตาย สมาทเปติ, ยํ’ส โหติ ทีฆรตฺตํ อหิตาย ทุกฺขายาติ. เอวํ ภนฺเต.}

\chapterheadpageread{(7) 2. มหาวคฺค}
\prose{\csromanfolio{195}ตํ กิํ มญฺญถ สาฬฺหา, อตฺถิ โทโสติ. เอวํ ภนฺเต. “พฺยาปาโท”ติ โข อหํ สาฬฺหา เอตมตฺถํ วทามิ, ทุฏฺโฐ โข อยํ สาฬฺหา พฺยาปนฺนจิตฺโต ปาณมฺปิ หนติ, อทินฺนมฺปิ อาทิยติ, ปรทารมฺปิ คจฺฉติ, มุสาปิ ภณติ, ปรมฺปิ ตถตฺตาย สมาทเปติ, ยํ’ส โหติ ทีฆรตฺตํ อหิตาย ทุกฺขายาติ. เอวํ ภนฺเต.}

\prose{ตํ กิํ มญฺญถ สาฬฺหา, อตฺถิ โมโหติ. เอวํ ภนฺเต. “อวิชฺชา”ติ โข อหํ สาฬฺหา เอตมตฺถํ วทามิ, มูโฬฺห โข อยํ สาฬฺหา อวิชฺชาคโต ปาณมฺปิ หนติ, อทินฺนมฺปิ อาทิยติ, ปรทารมฺปิ คจฺฉติ, มุสาปิ ภณติ, ปรมฺปิ ตถตฺตาย สมาทเปติ, ยํ’ส โหติ ทีฆรตฺตํ อหิตาย ทุกฺขายาติ. เอวํ ภนฺเต.}

\prose{ตํ กิํ มญฺญถ สาฬฺหา, อิเม ธมฺมา กุสลา วา อกุสลา วาติ. อกุสลา ภนฺเต. สาวชฺชา วา อนวชฺชา วาติ. สาวชฺชา ภนฺเต. วิญฺญุครหิตา วา วิญฺญุปฺปสตฺถา วาติ. วิญฺญุครหิตา ภนฺเต. สมตฺตา สมาทินฺนา อหิตาย ทุกฺขาย สํวตฺตนฺติ โน วา, กถํ วา เอตฺถ โหตีติ. สมตฺตา ภนฺเต สมาทินฺนา อหิตาย ทุกฺขาย สํวตฺตนฺตีติ เอวํ โน เอตฺถ โหตีติ.}

\prose{อิติ โข สาฬฺหา ยํ ตํ อโวจุมฺหา “เอถ ตุเมฺห สาฬฺหา, มา อนุสฺสเวน, มา ปรมฺปราย, มา อิติกิราย, มา ปิฏก\-สมฺปทาเนน, มา ตกฺกเหตุ, มา นยเหตุ, มา อาการ\-ปริวิตกฺเกน, มา ทิฏฺฐิ\-นิชฺฌาน\-กฺขนฺติยา, มา ภพฺพรูปตาย, มา สมโณ โน ครูติ. ยทา ตุเมฺห สาฬฺหา อตฺตนาว ชาเนยฺยาถ ‘อิเม ธมฺมา อกุสลา, อิเม ธมฺมา สาวชฺชา, อิเม ธมฺมา วิญฺญุครหิตา, อิเม ธมฺมา สมตฺตา สมาทินฺนา อหิตาย ทุกฺขาย สํวตฺตนฺตี’ติ. อถ ตุเมฺห สาฬฺหา ปชเหยฺยาถา”ติ อิติ ยํ ตํ วุตฺตํ, อิทเมตํ ปฏิจฺจ วุตฺตํ.}

\prose{เอถ ตุเมฺห สาฬฺหา, มา อนุสฺสเวน, มา ปรมฺปราย, มา อิติกิราย, มา ปิฏก\-สมฺปทาเนน, มา ตกฺกเหตุ, มา นยเหตุ, มา อาการ\-ปริวิตกฺเกน, มา ทิฏฺฐิ\-นิชฺฌาน\-กฺขนฺติยา, มา ภพฺพรูปตาย, มา สมโณ โน ครูติ. ยทา ตุเมฺห สาฬฺหา อตฺตนาว ชาเนยฺยาถ “อิเม ธมฺมา กุสลา, อิเม ธมฺมา อนวชฺชา, อิเม ธมฺมา วิญฺญุปฺปสตฺถา, อิเม ธมฺมา}

\prose{\csromanfolio{196}สมตฺตา สมาทินฺนา หิตาย สุขาย สํวตฺตนฺตี”ติ. อถ ตุเมฺห สาฬฺหา อุปสมฺปชฺช วิหเรยฺยาถ.}

\prose{ตํ กิํ มญฺญถ สาฬฺหา, อตฺถิ อโลโภติ. เอวํ ภนฺเต. “อนภิชฺฌา”ติ โข อหํ สาฬฺหา เอตมตฺถํ วทามิ, อลุทฺโธ โข อยํ สาฬฺหา อนภิชฺฌาลุ เนว ปาณํ หนติ, น อทินฺนํ อาทิยติ, น ปรทารํ คจฺฉติ, น มุสา ภณติ, ปรมฺปิ น ตถตฺตาย สมาทเปติ, ยํ’ส โหติ ทีฆรตฺตํ หิตาย สุขายาติ. เอวํ ภนฺเต.}

\prose{ตํ กิํ มญฺญถ สาฬฺหา, อตฺถิ อโทโสติ. เอวํ ภนฺเต. “อพฺยาปาโท”ติ โข อหํ สาฬฺหา เอตมตฺถํ วทามิ, อทุฏฺโฐ โข อยํ สาฬฺหา อพฺยาปนฺนจิตฺโต เนว ปาณํ หนติ, น อทินฺนํ อาทิยติ, น ปรทารํ คจฺฉติ, น มุสา ภณติ, ปรมฺปิ น ตถตฺตาย สมาทเปติ, ยํ’ส โหติ ทีฆรตฺตํ หิตาย สุขายาติ. เอวํ ภนฺเต.}

\prose{ตํ กิํ มญฺญถ สาฬฺหา, อตฺถิ อโมโหติ. เอวํ ภนฺเต. “วิชฺชา”ติ โข อหํ สาฬฺหา เอตมตฺถํ วทามิ, อมูโฬฺห โข อยํ สาฬฺหา วิชฺชาคโต เนว ปาณํ หวติ, น อทินฺนํ อาทิยติ, น ปรทารํ คจฺฉติ, น มุสา ภณติ, ปรมฺปิ น ตถตฺตาย สมาทเปติ, ยํ’ส โหติ ทีฆรตฺตํ หิตาย สุขายาติ. เอวํ ภนฺเต.}

\prose{ตํ กิํ มญฺญถ สาฬฺหา, อิเม ธมฺมา กุสลา วา อกุสลา วาติ. กุสลา ภนฺเต. สาวชฺชา วา อนวชฺชา วาติ. อนวชฺชา ภนฺเต. วิญฺญุครหิตา วา วิญฺญุปฺปสตฺถา วาติ. วิญฺญุปฺปสตฺถา ภนฺเต. สมตฺตา สมาทินฺนา หิตาย สุขาย สํวตฺตนฺติ โน วา, กถํ วา เอตฺถ โหตีติ. สมตฺตา ภนฺเต สมาทินฺนา หิตาย สุขาย สํวตฺตนฺตีติ เอวํ โน เอตฺถ โหตีติ.}

\prose{อิติ โข สาฬฺหา ยํ ตํ อโวจุมฺหา “เอถ ตุเมฺห สาฬฺหา, มา อนุสฺสเวน, มา ปรมฺปราย, มา อิติกิราย, มา ปิฏก\-สมฺปทาเนน, มา ตกฺกเหตุ, มา นยเหตุ, มา อาการ\-ปริวิตกฺเกน, มา ทิฏฺฐิ\-นิชฺฌาน\-กฺขนฺติยา, มา ภพฺพรูปตาย, มา สมโณ โน ครูติ. ยทา ตุเมฺห สาฬฺหา อตฺตนาว ชาเนยฺยาถ ‘อิเม ธมฺมา กุสลา, อิเม ธมฺมา อนวชฺชา, อิเม ธมฺมา วิญฺญุปฺปสตฺถา, อิเม ธมฺมา สมตฺตา}

\vspace{\tipitakaparskip}
\csromancenter{(7) 2. มหาวคฺค สมาทินฺนา ทีฆรตฺตํ หิราย สุขาย สํวตฺตนฺตี’ติ. อถ ตุเมฺห สาฬฺหา อุปสมฺปชฺช วิหเรยฺยาถา”ติ อิติ ยํ ตํ วุตฺตํ, อิทเมตํ ปฏิจฺจ วุตฺตํ.}
\prose{\csromanfolio{197}ส โข โส สาฬฺหา อริยสาวโก เอวํ วิคตาภิชฺโฌ วิคตพฺยาปาโท อสมฺมูโฬฺห สมฺปชาโน ปติสฺสโต เมตฺตา\-สหคเตน เจตสา ฯเปฯ กรุณา. มุทิตา. อุเปกฺขา\-สหคเตน เจตสา เอกํ ทิสํ ผริตฺวา วิหรติ. ตถา ทุติยํ. ตถา ตติยํ. ตถา จตุตฺถํ. อิติ อุทฺธมโธ ติริยํ สพฺพธิ สพฺพตฺตตาย สพฺพาวนฺตํ โลกํ อุเปกฺขา\-สหคเตน เจตสา วิปุเลน มหคฺคเตน อปฺปมาเณน อเวเรน อพฺยาปชฺเชน ผริตฺวา วิหรติ. โส เอวํ ปชานาติ “อตฺถิ อิทํ, อตฺถิ หีนํ, อตฺถิ ปณีตํ, อตฺถิ อิมสฺส สญฺญาคตสฺส อุตฺตริ\csromansharedfootnote{3d336730ed7a9e01}{อุตฺตริํ (สี, สฺยา, กํ, อิ)} นิสฺสรณนฺ”ติ. ตสฺส เอวํ ชานโต เอวํ ปสฺสโต กามาสวาปิ จิตฺตํ วิมุจฺจติ, ภวาสวาปิ จิตฺตํ วิมุจฺจติ, อวิชฺชาสวาปิ จิตฺตํ วิมุจฺจติ, วิมุตฺตสฺมิํ “วิมุตฺตมฺ”อิติ ญาณํ โหติ, “ขีณา ชาติ, วุสิตํ พฺรหฺมจริยํ, กตํ กรณียํ, นาปรํ อิตฺถตฺตายา”ติ ปชานาติ.}

\prose{โส เอวํ ปชานาติ “อหุ ปุพฺเพ โลโภ, ตทหุ อกุสลํ, โส เอตรหิ นตฺถิ, อิจฺเจตํ กุสลํ. อหุ ปุพฺเพ โทโส ฯเปฯ. อหุ ปุพฺเพ โมโห, ตทหุ อกุสลํ, โส เอตรหิ นตฺถิ, อิจฺเจตํ กุสลนฺ”ติ. โส ทิฏฺเฐว ธมฺเม นิจฺฉาโต นิพฺพุโต สีติภูโต สุขปฺปฏิสํเวที พฺรหฺมภูเตน อตฺตนา วิหรตีติ.  ฉฏฺฐํ.}

\csromanheader{2}{7. กถาวตฺถุ\-สุตฺต}
\proseitem{68}{ตีณิมานิ ภิกฺขเว กถาวตฺถูนิ. กตมานิ ตีณิ? อตีตํ วา ภิกฺขเว อทฺธานํ อารพฺภ กถํ กเถยฺย “เอวํ อโหสิ อตีตมทฺธานนฺ”ติ, อนาคตํ วา ภิกฺขเว อทฺธานํ อารพฺภ กถํ กเถยฺย “เอวํ ภวิสฺสติ อนาคตมทฺธานนฺ”ติ, เอตรหิ วา ภิกฺขเว ปจฺจุปฺปนฺนํ อทฺธานํ อารพฺภ กถํ กเถยฺย “เอวํ โหติ เอตรหิ ปจฺจุปฺปนฺนมทฺธานนฺ”ติ\csromansharedfootnote{130cd975c9399720}{เอวํ เอตรหิ ปจฺจุปฺปนฺนนฺติ (สี, อิ, ก), เอวํ โหติ เอตรหิ ปจฺจุปฺปนฺนนฺติ (สฺยา, กํ)}.}

\prose{กถา\-สมฺปโยเคน ภิกฺขเว ปุคฺคโล เวทิตพฺโพ “ยทิ วา กจฺโฉ, ยทิ วา อกจฺโฉ”ติ. สจายํ ภิกฺขเว ปุคฺคโล ปญฺหํ ปุฏฺโฐ สมาโน \csromanfolio{198}เอกํส\-พฺยากรณียํ ปญฺหํ น เอกํเสน พฺยากโรติ, วิภชฺช\-พฺยากรณียํ ปญฺหํ น วิภชฺช พฺยากโรติ, ปฏิปุจฺฉา\-พฺยากรณียํ ปญฺหํ น ปฏิปุจฺฉา พฺยากโรติ, ฐปนียํ ปญฺหํ น ฐเปติ\csromansharedfootnote{48eda34f46b0ba4e}{ถปนียํ ปญฺหํ น ถเปติ (ก)}. เอวํ สนฺตายํ ภิกฺขเว ปุคฺคโล อกจฺโฉ โหติ. สเจ ปนายํ ภิกฺขเว ปุคฺคโล ปญฺหํ ปุฏฺโฐ สมาโน เอกํส\-พฺยากรณียํ ปญฺหํ เอกํเสน พฺยากโรติ, วิภชฺช\-พฺยากรณียํ ปญฺหํ วิภชฺช พฺยากโรติ, ปฏิปุจฺฉา\-พฺยากรณียํ ปญฺหํ ปฏิปุจฺฉา พฺยากโรติ, ฐปนียํ ปญฺหํ ฐเปติ. เอวํ สนฺตายํ ภิกฺขเว ปุคฺคโล กจฺโฉ โหติ.}

\prose{กถา\-สมฺปโยเคน ภิกฺขเว ปุคฺคโล เวทิตพฺโพ “ยทิ วา กจฺโฉ, ยทิ วา อกจฺโฉ”ติ. สจายํ ภิกฺขเว ปุคฺคโล ปญฺหํ ปุฏฺโฐ สมาโน ฐานาฐาเน น สณฺฐาติ, ปริกปฺเป น สณฺฐาติ, อญฺญาตวาเท\csromansharedfootnote{599ee710b1a1bfb2}{อญฺญวาเท (สี, สฺยา, กํ, อิ), อญฺญาตวาเร (ก)} น สณฺฐาติ, ปฏิปทาย น สณฺฐาติ. เอวํ สนฺตายํ ภิกฺขเว ปุคฺคโล อกจฺโฉ โหติ. สเจ ปนายํ ภิกฺขเว ปุคฺคโล ปญฺหํ ปุฏฺโฐ สมาโน ฐานาฐาเน สณฺฐาติ, ปริกปฺเป สณฺฐาติ, อญฺญาตวาเท สณฺฐาติ, ปฏิปทาย สณฺฐาติ. เอวํ สนฺตายํ ภิกฺขเว ปุคฺคโล กจฺโฉ โหติ.}

\prose{กถา\-สมฺปโยเคน ภิกฺขเว ปุคฺคโล เวทิตพฺโพ “ยทิ วา กจฺโฉ, ยทิ วา อกจฺโฉ”ติ. สจายํ ภิกฺขเว ปุคฺคโล ปญฺหํ ปุฏฺโฐ สมาโน อญฺเญนญฺญํ ปฏิจรติ, พหิทฺธา กถํ อปนาเมติ, โกปญฺจ โทสญฺจ อปฺปจฺจยญฺจ ปาตุกโรติ. เอวํ สนฺตายํ ภิกฺขเว ปุคฺคโล อกจฺโฉ โหติ. สเจ ปนายํ ภิกฺขเว ปุคฺคโล ปญฺหํ ปุฏฺโฐ สมาโน น อญฺเญนญฺญํ ปฏิจรติ, น พหิทฺธา กถํ อปนาเมติ, น โกปญฺจ โทสญฺจ อปฺปจฺจยญฺจ ปาตุกโรติ. เอวํ สนฺตายํ ภิกฺขเว ปุคฺคโล กจฺโฉ โหติ.}

\prose{กถา\-สมฺปโยเคน ภิกฺขเว ปุคฺคโล เวทิตพฺโพ “ยทิ วา กจฺโฉ, ยทิ วา อกจฺโฉ”ติ. สจายํ ภิกฺขเว ปุคฺคโล ปญฺหํ ปุฏฺโฐ สมาโน อภิหรติ, อภิมทฺทติ, อนุปชคฺฆติ\csromansharedfootnote{60a50250fa1ae322}{อนุสํชคฺฆติ (ก)}, ขลิตํ คณฺหาติ. เอวํ สนฺตายํ ภิกฺขเว ปุคฺคโล อกจฺโฉ โหติ. สเจ ปนายํ ภิกฺขเว ปุคฺคโล ปญฺหํ ปุฏฺโฐ สมาโน นาภิหรติ, นาภิมทฺทติ, น อนุปชคฺฆติ, น ขลิตํ คณฺหาติ. เอวํ สนฺตายํ ภิกฺขเว ปุคฺคโล กจฺโฉ โหติ.}

\chapterheadpageread{(7) 2. มหาวคฺค}
\prose{\csromanfolio{199}กถา\-สมฺปโยเคน ภิกฺขเว ปุคฺคโล เวทิตพฺโพ “ยทิ วา สอุปนิโส, ยทิ วา อนุปนิโส”ติ. อโนหิตโสโต ภิกฺขเว อนุปนิโส โหติ, โอหิตโสโต สอุปนิโส โหติ. โส สอุปนิโส สมาโน อภิชานาติ เอกํ ธมฺมํ, ปริชานาติ เอกํ ธมฺมํ, ปชหติ เอกํ ธมฺมํ, สจฺฉิกโรติ เอกํ ธมฺมํ. โส อภิชานนฺโต เอกํ ธมฺมํ, ปริชานนฺโต เอกํ ธมฺมํ, ปชหนฺโต เอกํ ธมฺมํ, สจฺฉิกโรนฺโต เอกํ ธมฺมํ สมฺมาวิมุตฺติํ ผุสติ. เอตทตฺถา ภิกฺขเว กถา, เอตทตฺถา มนฺตนา, เอตทตฺถา อุปนิสา, เอตทตฺถํ โสตาวธานํ, ยทิทํ อนุปาทา จิตฺตสฺส วิโมกฺโขติ.}

\setlength{\csromangathapagewidth}{0pt}
\setlength{\csromangathawakwidth}{0pt}
\csromangathameasuregroup{เย วิรุทฺธา สลฺลปนฺติ, \\ อนริยคุณมาสชฺช, \\ ทุพฺภาสิตํ วิกฺขลิตํ, \\ อญฺโญญฺญสฺสาภินนฺทนฺติ, \\ สเจ จสฺส กถากาโม, \\ ธมฺมฏฺฐปฏิสํยุตฺตา, \\ ตํ กถํ กถเย ธีโร, \\ อนุนฺนเตน มนสา, \\ อนุสูยายมาโน โส, \\ สุภาสิตํ อนุโมเทยฺย, \\ อุปารมฺภํ น สิกฺเขยฺย, \\ นาภิหเร นาภิมทฺเท, \\ อญฺญาตตฺถํ ปสาทตฺถํ, \\ เอวํ โข อริยา มนฺเตนฺติ, \\ เอตทญฺญาย เมธาวี,}{\csromangathabat{เย วิรุทฺธา สลฺลปนฺติ,}{วินิวิฏฺฐา สมุสฺสิตา.} \\ \csromangathabat{อนริยคุณมาสชฺช,}{อญฺโญญฺญวิวเรสิโน.} \\ \csromangathabat{ทุพฺภาสิตํ วิกฺขลิตํ,}{สมฺปโมหํ\textsuperscript{1} ปราชยํ.} \\ \csromangathabat{อญฺโญญฺญสฺสาภินนฺทนฺติ,}{ตทริโย กถนาจเร\textsuperscript{1}.} \\ \csromangathabat{สเจ จสฺส กถากาโม,}{กาลมญฺญาย ปณฺฑิโต.} \\ \csromangathabat{ธมฺมฏฺฐปฏิสํยุตฺตา,}{ยา อริยจริตา\textsuperscript{1} กถา.} \\ \csromangathabat{ตํ กถํ กถเย ธีโร,}{อวิรุทฺโธ อนุสฺสิโต.} \\ \csromangathabat{อนุนฺนเตน มนสา,}{อปฬาโส อสาหโส.} \\ \csromangathabat{อนุสูยายมาโน โส,}{สมฺมทญฺญาย ภาสติ.} \\ \csromangathabat{สุภาสิตํ อนุโมเทยฺย,}{ทุพฺภฏฺเฐ นาปสาทเย\textsuperscript{1}.} \\ \csromangathabat{อุปารมฺภํ น สิกฺเขยฺย,}{ขลิตญฺจ น คาหเย\textsuperscript{1}.} \\ \csromangathabat{นาภิหเร นาภิมทฺเท,}{น วาจํ ปยุตํ ภเณ.} \\ \csromangathabat{อญฺญาตตฺถํ ปสาทตฺถํ,}{สตํ เว โหติ มนฺตนา.} \\ \csromangathabat{เอวํ โข อริยา มนฺเตนฺติ,}{เอสา อริยาน มนฺตนา.} \\ \csromangathabat{เอตทญฺญาย เมธาวี,}{น สมุสฺเสยฺย มนฺตเยติ.สตฺตมํ.}}
\csromangathasetleft{เย วิรุทฺธา สลฺลปนฺติ, \\ อนริยคุณมาสชฺช, \\ ทุพฺภาสิตํ วิกฺขลิตํ, \\ อญฺโญญฺญสฺสาภินนฺทนฺติ, \\ สเจ จสฺส กถากาโม, \\ ธมฺมฏฺฐปฏิสํยุตฺตา, \\ ตํ กถํ กถเย ธีโร, \\ อนุนฺนเตน มนสา, \\ อนุสูยายมาโน โส, \\ สุภาสิตํ อนุโมเทยฺย, \\ อุปารมฺภํ น สิกฺเขยฺย, \\ นาภิหเร นาภิมทฺเท, \\ อญฺญาตตฺถํ ปสาทตฺถํ, \\ เอวํ โข อริยา มนฺเตนฺติ, \\ เอตทญฺญาย เมธาวี,}
\csromangathagroup{\csromangathastanza{\csromangathabat{เย วิรุทฺธา สลฺลปนฺติ,}{วินิวิฏฺฐา สมุสฺสิตา.} \\ \csromangathabat{อนริยคุณมาสชฺช,}{อญฺโญญฺญวิวเรสิโน.}}\csromangathastanza{\csromangathabat{ทุพฺภาสิตํ วิกฺขลิตํ,}{สมฺปโมหํ\csromansharedfootnote{6c89b3236ab95e9c}{สสมฺโมหํ (ก)} ปราชยํ.} \\ \csromangathabat{อญฺโญญฺญสฺสาภินนฺทนฺติ,}{ตทริโย กถนาจเร\csromansharedfootnote{6c89b3236ab95e9c}{สสมฺโมหํ (ก)}.}}\csromangathastanza{\csromangathabat{สเจ จสฺส กถากาโม,}{กาลมญฺญาย ปณฺฑิโต.} \\ \csromangathabat{ธมฺมฏฺฐปฏิสํยุตฺตา,}{ยา อริยจริตา\csromansharedfootnote{12aff01a67a38f14}{อริยญฺจริตา (สี), อริยาทิกา (ก)} กถา.}}\csromangathastanza{\csromangathabat{ตํ กถํ กถเย ธีโร,}{อวิรุทฺโธ อนุสฺสิโต.} \\ \csromangathabat{อนุนฺนเตน มนสา,}{อปฬาโส อสาหโส.}}\csromangathastanza{\csromangathabat{อนุสูยายมาโน โส,}{สมฺมทญฺญาย ภาสติ.} \\ \csromangathabat{สุภาสิตํ อนุโมเทยฺย,}{ทุพฺภฏฺเฐ นาปสาทเย\csromansharedfootnote{f6dddfc7c7d3fd04}{นาวสาทเย (สี, อิ)}.}}\csromangathastanza{\csromangathabat{อุปารมฺภํ น สิกฺเขยฺย,}{ขลิตญฺจ น คาหเย\csromansharedfootnote{a787cfc0283e58f3}{น ภาสเย (ก)}.} \\ \csromangathabat{นาภิหเร นาภิมทฺเท,}{น วาจํ ปยุตํ ภเณ.}}\csromangathastanza{\csromangathabat{อญฺญาตตฺถํ ปสาทตฺถํ,}{สตํ เว โหติ มนฺตนา.} \\ \csromangathabat{เอวํ โข อริยา มนฺเตนฺติ,}{เอสา อริยาน มนฺตนา.} \\ \csromangathabat{เอตทญฺญาย เมธาวี,}{น สมุสฺเสยฺย มนฺตเยติ.  สตฺตมํ.}}}

\csromanheader{2}{8. อญฺญติตฺถิย\-สุตฺต}
\proseitem{69}{สเจ ภิกฺขเว อญฺญติตฺถิยา ปริพฺพาชกา เอวํ ปุจฺเฉยฺยุํ “ตโยเม อาวุโส ธมฺมา. กตเม ตโย? ราโค โทโส \csromanfolio{200}โมโห. อิเม โข อาวุโส ตโย ธมฺมา. อิเมสํ อาวุโส ติณฺณํ ธมฺมานํ โก วิเสโส, โก อธิปฺปยาโส\csromansharedfootnote{c096b86c8f341691}{อธิปฺปาโย (สี), อธิปฺปายาโส (สฺยา, กํ, อิ) อธิ + ป + ยสุ + ณ = อธิปฺปยาโส.}, กิํ นานากรณนฺ”ติ. เอวํ ปุฏฺฐา ตุเมฺห ภิกฺขเว เตสํ อญฺญ\-ติตฺถิยานํ ปริพฺพาชกานํ กินฺติ พฺยากเรยฺยาถาติ. ภควํมูลกา โน ภนฺเต ธมฺมา, ภควํ\-เนตฺติกา ภควํ\-ปฏิสรณา. สาธุ วต ภนฺเต ภควนฺตํเยว ปฏิภาตุ เอตสฺส ภาสิตสฺส อตฺโถ, ภควโต สุตฺวา ภิกฺขู ธาเรสฺสนฺตีติ. เตน หิ ภิกฺขเว สุณาถ สาธุกํ มนสิ กโรถ, ภาสิสฺสามีติ. “เอวํ ภนฺเต”ติ โข เต ภิกฺขู ภควโต ปจฺจสฺโสสุํ. ภควา เอตทโวจ\csromandash{}}

\prose{สเจ ภิกฺขเว อญฺญติตฺถิยา ปริพฺพาชกา เอวํ ปุจฺเฉยฺยุํ “ตโยเม อาวุโส ธมฺมา. กตเม ตโย? ราโค โทโส โมโห. อิเม โข อาวุโส ตโย ธมฺมา. อิเมสํ อาวุโส ติณฺณํ ธมฺมานํ โก วิเสโส, โก อธิปฺปยาโส, กิํ นานากรณนฺ”ติ. เอวํ ปุฏฺฐา ตุเมฺห ภิกฺขเว เตสํ อญฺญ\-ติตฺถิยานํ ปริพฺพาชกานํ เอวํ พฺยากเรยฺยาถ “ราโค โข อาวุโส อปฺปสาวชฺโช ทนฺธวิราคี, โทโส มหาสาวชฺโช ขิปฺปวิราคี, โมโห มหาสาวชฺโช ทนฺธวิราคี”ติ.}

\prose{โก ปนาวุโส เหตุ โก ปจฺจโย, เยน อนุปฺปนฺโน วา ราโค อุปฺปชฺชติ, อุปฺปนฺโน วา ราโค ภิยฺโยภาวาย เวปุลฺลาย สํวตฺตตีติ, “สุภนิมิตฺตนฺ”ติสฺส วจนียํ. ตสฺส สุภนิมิตฺตํ อโยนิโส มนสิ กโรโต อนุปฺปนฺโน วา ราโค อุปฺปชฺชติ, อุปฺปนฺโน วา ราโค ภิยฺโยภาวาย เวปุลฺลาย สํวตฺตติ. อยํ โข อาวุโส เหตุ อยํ ปจฺจโย, เยน อนุปฺปนฺโน วา ราโค อุปฺปชฺชติ, อุปฺปนฺโน วา ราโค ภิยฺโยภาวาย เวปุลฺลาย สํวตฺตตีติ.}

\prose{โก ปนาวุโส เหตุ โก ปจฺจโย, เยน อนุปฺปนฺโน วา โทโส อุปฺปชฺชติ, อุปฺปนฺโน วา โทโส ภิยฺโยภาวาย เวปุลฺลาย สํวตฺตตีติ, “ปฏิฆนิมิตฺตนฺ”ติสฺส วจนียํ. ตสฺส ปฏิฆ\-นิมิตฺตํ อโยนิโส มนสิ กโรโต อนุปฺปนฺโน วา โทโส อุปฺปชฺชติ, อุปฺปนฺโน วา โทโส}

\vspace{\tipitakaparskip}
\csromancenter{(7) 2. มหาวคฺค ภิยฺโยภาวาย เวปุลฺลาย สํวตฺตติ. อยํ โข อาวุโส เหตุ อยํ ปจฺจโย, เยน อนุปฺปนฺโน วา โทโส อุปฺปชฺชติ, อุปฺปนฺโน วา โทโส ภิยฺโยภาวาย เวปุลฺลาย สํวตฺตตีติ.}
\prose{\csromanfolio{201}โก ปนาวุโส เหตุ โก ปจฺจโย, เยน อนุปฺปนฺโน วา โมโห อุปฺปชฺชติ, อุปฺปนฺโน วา โมโห ภิยฺโยภาวาย เวปุลฺลาย สํวตฺตตีติ, “อโยนิโส\-มนสิกาโร”ติสฺส วจนียํ. ตสฺส อโยนิโส มนสิ กโรโต อนุปฺปนฺโน วา โมโห อุปฺปชฺชติ, อุปฺปนฺโน วา โมโห ภิยฺโยภาวาย เวปุลฺลาย สํวตฺตติ. อยํ โข อาวุโส เหตุ อยํ ปจฺจโย, เยน อนุปฺปนฺโน วา โมโห อุปฺปชฺชติ, อุปฺปนฺโน วา โมโห ภิยฺโยภาวาย เวปุลฺลาย สํวตฺตตีติ.}

\prose{โก ปนาวุโส เหตุ โก ปจฺจโย, เยน อนุปฺปนฺโน เจว ราโค นุปฺปชฺชติ, อุปฺปนฺโน จ ราโค ปหียตีติ, “อสุภนิมิตฺตนฺ”ติสฺส วจนียํ. ตสฺส อสุภนิมิตฺตํ โยนิโส มนสิ กโรโต อนุปฺปนฺโน เจว ราโค นุปฺปชฺชติ, อุปฺปนฺโน จ ราโค ปหียติ. อยํ โข อาวุโส เหตุ อยํ ปจฺจโย, เยน อนุปฺปนฺโน เจว ราโค นุปฺปชฺชติ, อุปฺปนฺโน จ ราโค ปหียตีติ.}

\prose{โก ปนาวุโส เหตุ โก ปจฺจโย, เยน อนุปฺปนฺโน เจว โทโส นุปฺปชฺชติ, อุปฺปนฺโน จ โทโส ปหียตีติ, “เมตฺตา เจโตวิมุตฺตี”ติสฺส วจนียํ. ตสฺส เมตฺตํ เจโตวิมุตฺติํ โยนิโส มนสิ กโรโต อนุปฺปนฺโน เจว โทโส นุปฺปชฺชติ, อุปฺปนฺโน จ โทโส ปหียติ. อยํ โข อาวุโส เหตุ อยํ ปจฺจโย, เยน อนุปฺปนฺโน เจว โทโส นุปฺปชฺชติ, อุปฺปนฺโน จ โทโส ปหียตีติ.}

\prose{โก ปนาวุโส เหตุ โก ปจฺจโย, เยน อนุปฺปนฺโน เจว โมโห นุปฺปชฺชติ, อุปฺปนฺโน จ โมโห ปหียตีติ, “โยนิโส\-มนสิกาโร”ติสฺส วจนียํ. ตสฺส โยนิโส มนสิ กโรโต อนุปฺปนฺโน เจว โมโห นุปฺปชฺชติ, อุปฺปนฺโน จ โมโห ปหียติ. อยํ โข อาวุโส เหตุ อยํ ปจฺจโย, เยน อนุปฺปนฺโน วา โมโห นุปฺปชฺชติ, อุปฺปนฺโน วา โมโห ปหียตีติ.  อฏฺฐมํ.}

\csromanpalirecto
\gambhira{ติกนิปาตปาฬิ 9. อกุสลมูล\-สุตฺต}
\proseitem{70}{\csromanfolio{202}ตีณิมานิ ภิกฺขเว อกุสลมูลานิ. กตมานิ ตีณิ? โลโภ อกุสลมูลํ, โทโส อกุสลมูลํ, โมโห อกุสลมูลํ.}

\prose{ยทปิ ภิกฺขเว โลโภ, ตทปิ อกุสลมูลํ\csromansharedfootnote{22c90b448e844ea8}{อกุสลํ (สี, สฺยา, กํ, อิ)}. ยทปิ ลุทฺโธ อภิสงฺขโรติ กาเยน วาจาย มนสา, ตทปิ อกุสลํ\csromansharedfootnote{a1472f1dec45c599}{อกุสลมูลํ (ก)}. ยทปิ ลุทฺโธ โลเภน อภิภูโต ปริยาทินฺน\-จิตฺโต ปรสฺส อสตา ทุกฺขํ อุปฺปาทยติ\csromansharedfootnote{a6409281931f5acc}{อุปทหติ (สี, สฺยา, กํ, อิ)} วเธน วา พนฺธเนน วา ชานิยา วา ครหาย วา ปพฺพาชนาย วา พลวมฺหิ พลตฺโถ อิติปิ, ตทปิ อกุสลํ\csromansharedfootnote{742ab91c794225c8}{อิทํ ปน สพฺพตฺถปิ เอวเมว ทิสฺสติ.}. อิติสฺส’เม โลภชา โลภนิทานา โลภสมุทยา โลภปจฺจยา อเนเก ปาปกา อกุสลา ธมฺมา สมฺภวนฺติ.}

\prose{ยทปิ ภิกฺขเว โทโส, ตทปิ อกุสลมูลํ. ยทปิ ทุฏฺโฐ อภิสงฺขโรติ กาเยน วาจาย มนสา, ตทปิ อกุสลํ. ยทปิ ทุฏฺโฐ โทเสน อภิภูโต ปริยาทินฺน\-จิตฺโต ปรสฺส อสตา ทุกฺขํ อุปฺปาทยติ วเธน วา พนฺธเนน วา ชานิยา วา ครหาย วา ปพฺพาชนาย วา พลวมฺหิ พลตฺโถ อิติปิ, ตทปิ อกุสลํ. อิติสฺส’เม โทสชา โทสนิทานา โทสสมุทยา โทสปจฺจยา อเนเก ปาปกา อกุสลา ธมฺมา สมฺภวนฺติ.}

\prose{ยทปิ ภิกฺขเว โมโห, ตทปิ อกุสลมูลํ. ยทปิ มูโฬฺห อภิสงฺขโรติ กาเยน วาจาย มนสา, ตทปิ อกุสลํ. ยทปิ มูโฬฺห โมเหน อภิภูโต ปริยาทินฺน\-จิตฺโต ปรสฺส อสตา ทุกฺขํ อุปฺปาทยติ วเธน วา พนฺธเนน วา ชานิยา วา ครหาย วา ปพฺพาชนาย วา พลวมฺหิ พลตฺโถ อิติปิ, ตทปิ อกุสลํ. อิติสฺส’เม โมหชา โมหนิทานา โมหสมุทยา โมหปจฺจยา อเนเก ปาปกา อกุสลา ธมฺมา สมฺภวนฺติ. เอวรูโป จายํ ภิกฺขเว ปุคฺคโล วุจฺจติ “อกาลวาที”ติปิ “อภูตวาที”ติปิ “อนตฺถวาที”ติปิ “อธมฺมวาที”ติปิ “อวินยวาที”ติปิ.}

\chapterheadpageread{(7) 2. มหาวคฺค}
\prose{\csromanfolio{203}กสฺมา จายํ ภิกฺขเว เอวรูโป ปุคฺคโล วุจฺจติ “อกาลวาที”ติปิ “อภูตวาที”ติปิ “อนตฺถวาที”ติปิ “อธมฺมวาที”ติปิ “อวินยวาที”ติปิ? ตถาหายํ ภิกฺขเว ปุคฺคโล ปรสฺส อสตา ทุกฺขํ อุปฺปาทยติ วเธน วา พนฺธเนน วา ชานิยา วา ครหาย วา ปพฺพาชนาย วา พลวมฺหิ พลตฺโถ อิติปิ, ภูเตน โข ปน วุจฺจมาโน อวชานาติ โน ปฏิชานาติ, อภูเตน วุจฺจมาโน น อาตปฺปํ กโรติ ตสฺส นิพฺเพฐนาย “อิติเปตํ อตจฺฉํ อิติเปตํ อภูตนฺ”ติ. ตสฺมา เอวรูโป ปุคฺคโล วุจฺจติ “อกาลวาที”ติปิ “อภูตวาที”ติปิ “อนตฺถวาที”ติปิ “อธมฺมวาที”ติปิ “อวินยวาที”ติปิ.}

\prose{เอวรูโป ภิกฺขเว ปุคฺคโล โลภเชหิ ปาปเกหิ อกุสเลหิ ธมฺเมหิ อภิภูโต ปริยาทินฺน\-จิตฺโต ทิฏฺเฐ เจว ธมฺเม ทุกฺขํ วิหรติ สวิฆาตํ สอุปายาสํ สปริฬาหํ, กายสฺส จ เภทา ปรํ มรณา ทุคฺคติ ปาฏิกงฺขา.}

\prose{โทสเชหิ ฯเปฯ. โมหเชหิ ปาปเกหิ อกุสเลหิ ธมฺเมหิ อภิภูโต ปริยาทินฺน\-จิตฺโต ทิฏฺเฐ เจว ธมฺเม ทุกฺขํ วิหรติ สวิฆาตํ สอุปายาสํ สปริฬาหํ, กายสฺส จ เภทา ปรํ มรณา ทุคฺคติ ปาฏิกงฺขา. เสยฺยถาปิ ภิกฺขเว สาโล วา ธโว วา ผนฺทโน วา ตีหิ มาลุวาลตาหิ อุทฺธสฺโต ปริโยนทฺโธ อนยํ อาปชฺชติ, พฺยสนํ อาปชฺชติ, อนยพฺยสนํ อาปชฺชติ. เอวเมวํ โข ภิกฺขเว เอวรูโป ปุคฺคโล โลภเชหิ ปาปเกหิ อกุสเลหิ ธมฺเมหิ อภิภูโต ปริยาทินฺน\-จิตฺโต ทิฏฺเฐ เจว ธมฺเม ทุกฺขํ วิหรติ สวิฆาตํ สอุปายาสํ สปริฬาหํ, กายสฺส จ เภทา ปรํ มรณา ทุคฺคติ ปาฏิกงฺขา.}

\prose{โทสเชหิ ฯเปฯ. โมหเชหิ ปาปเกหิ อกุสเลหิ ธมฺเมหิ อภิภูโต ปริยาทินฺน\-จิตฺโต ทิฏฺเฐ เจว ธมฺเม ทุกฺขํ วิหรติ สวิฆาตํ สอุปายาสํ สปริฬาหํ, กายสฺส จ เภทา ปรํ มรณา ทุคฺคติ ปาฏิกงฺขา. อิมานิ โข ภิกฺขเว ตีณิ อกุสลมูลานีติ.}

\prose{ตีณิมานิ ภิกฺขเว กุสลมูลานิ. กตมานิ ตีณิ? อโลโภ กุสลมูลํ, อโทโส กุสลมูลํ, อโมโห กุสลมูลํ.}

\prose{\csromanfolio{204}ยทปิ ภิกฺขเว อโลโภ, ตทปิ กุสลมูลํ\csromansharedfootnote{9b7ee7738afec699}{กุสลํ (สี, สฺยา, กํ, อิ)}. ยทปิ อลุทฺโธ อภิสงฺขโรติ กาเยน วาจาย มนสา, ตทปิ กุสลํ\csromansharedfootnote{fb503336b55a1ca7}{กุสลมูลํ (ก)}. ยทปิ อลุทฺโธ โลเภน อนภิภูโต อปริยาทินฺนจิตฺโต น ปรสฺส อสตา ทุกฺขํ อุปฺปาทยติ วเธน วา พนฺธเนน วา ชานิยา วา ครหาย วา ปพฺพาชนาย วา พลวมฺหิ พลตฺโถ อิติปิ, ตทปิ กุสลํ. อิติสฺส’เม อโลภชา อโลภนิทานา อโลภสมุทยา อโลภปจฺจยา อเนเก กุสลา ธมฺมา สมฺภวนฺติ.}

\prose{ยทปิ ภิกฺขเว อโทโส, ตทปิ กุสลมูลํ. ยทปิ อทุฏฺโฐ อภิสงฺขโรติ กาเยน วาจาย มนสา, ตทปิ กุสลํ. ยทปิ อทุฏฺโฐ โทเสน อนภิภูโต อปริยาทินฺนจิตฺโต น ปรสฺส อสตา ทุกฺขํ อุปฺปาทยติ วเธน วา พนฺธเนน วา ชานิยา วา ครหาย วา ปพฺพาชนาย วา พลวมฺหิ พลตฺโถ อิติปิ, ตทปิ กุสลํ. อิติสฺส’เม อโทสชา อโทสนิทานา อโทสสมุทยา อโทสปจฺจยา อเนเก กุสลา ธมฺมา สมฺภวนฺติ.}

\prose{ยทปิ ภิกฺขเว อโมโห, ตทปิ กุสลมูลํ. ยทปิ อมูโฬฺห อภิสงฺขโรติ กาเยน วาจาย มนสา, ตทปิ กุสลํ. ยทปิ อมูโฬฺห โมเหน อนภิภูโต อปริยาทินฺนจิตฺโต น ปรสฺส อสตา ทุกฺขํ อุปฺปาทยติ วเธน วา พนฺธเนน วา ชานิยา วา ครหาย วา ปพฺพาชนาย วา พลวมฺหิ พลตฺโถ อิติปิ, ตทปิ กุสลํ. อิติสฺส’เม อโมหชา อโมหนิทานา อโมหสมุทยา อโมหปจฺจยา อเนเก กุสลา ธมฺมา สมฺภวนฺติ. เอวรูโป จายํ ภิกฺขเว ปุคฺคโล วุจฺจติ “กาลวาที”ติปิ “ภูตวาที”ติปิ “อตฺถวาที”ติปิ “ธมฺมวาที”ติปิ “วินยวาที”ติปิ.}

\prose{กสฺมา จายํ ภิกฺขเว เอวรูโป ปุคฺคโล วุจฺจติ “กาลวาที”ติปิ “ภูตวาที”ติปิ “อตฺถวาที”ติปิ “วินยวาที”ติปิ? ตถาหายํ ภิกฺขเว ปุคฺคโล น ปรสฺส อสตา ทุกฺขํ อุปฺปาทยติ วเธน วา พนฺธเนน วา ชานิยา วา ครหาย วา ปพฺพาชนาย วา พลวมฺหิ พลตฺโถ อิติปิ. ภูเตน โข ปน วุจฺจมาโน ปฏิชานาติ โน อวชานาติ, อภูเตน วุจฺจมาโน อาตปฺปํ กโรติ ตสฺส นิพฺเพฐนาย “อิติเปตํ อตจฺฉํ, อิติเปตํ}

\vspace{\tipitakaparskip}
\csromancenter{(7) 2. มหาวคฺค อภูตนฺ”ติ. ตสฺมา เอวรูโป ปุคฺคโล วุจฺจติ “กาลวาที”ติปิ “อตฺถวาที”ติปิ “ธมฺมวาที”ติปิ “วินยวาที”ติปิ.}
\prose{\csromanfolio{205}เอวรูปสฺส ภิกฺขเว ปุคฺคลสฺส โลภชา ปาปกา อกุสลา ธมฺมา ปหีนา อุจฺฉินฺนมูลา ตาลาวตฺถุกตา อนภาวํกตา อายติํ อนุปฺปาทธมฺมา ทิฏฺเฐว ธมฺเม สุขํ วิหรติ อวิฆาตํ อนุปายาสํ อปริฬาหํ, ทิฏฺเฐว ธมฺเม ปรินิพฺพายติ.}

\prose{โทสชา ฯเปฯ ปรินิพฺพายติ. โมหชา ฯเปฯ ปรินิพฺพายติ. เสยฺยถาปิ ภิกฺขเว สาโล วา ธโว วา ผนฺทโน วา ตีหิ มาลุวาลตาหิ อุทฺธสฺโต ปริโยนทฺโธ. อถ ปุริโส อาคจฺเฉยฺย กุทาลปิฏกํ\csromansharedfootnote{2dea7970ca7edc24}{กุทฺทาลปิฏกํ (สี, สฺยา, กํ, อิ)} อาทาย, โส ตํ มาลุวาลตํ มูเล ฉินฺเทยฺย, มูเล เฉตฺวา ปลิขเณยฺย, ปลิขณิตฺวา มูลานิ อุทฺธเรยฺยอนฺตมโส อุสีร\-นาฬิ\-มตฺตานิปิ\csromansharedfootnote{f8b332ef56b2ebbb}{อุสีรนาฬมตฺตานิปิ (สี, สฺยา, กํ, อิ)}. โส ตํ มาลุวาลตํ ขณฺฑาขณฺฑิกํ ฉินฺเทยฺย, ขณฺฑาขณฺฑิกํ เฉตฺวา ผาเลยฺย, ผาเลตฺวา สกลิกํ สกลิกํ กเรยฺย, สกลิกํ สกลิกํ กริตฺวา วาตาตเป วิโสเสยฺย, วาตาตเป วิโสเสตฺวา อคฺคินา ฑเหยฺย, อคฺคินา ฑหิตฺวา มสิํ กเรยฺย, มสิํ กริตฺวา มหาวาเต วา โอผุเณยฺย, นทิยา วา สีฆโสตาย ปวาเหยฺย. เอวมสฺส\csromansharedfootnote{6c49524ed036c7df}{เอวมสฺสุ (สี), เอวสฺสุ (ก)} ตา ภิกฺขเว มาลุวาลตา อุจฺฉินฺนมูลา ตาลาวตฺถุกตา อนภาวํกตา อายติํ อนุปฺปาทธมฺมา. เอวเมวํ โข ภิกฺขเว เอวรูปสฺส ปุคฺคลสฺส โลภชา ปาปกา อกุสลา ธมฺมา ปหีนา อุจฺฉินฺนมูลา ตาลาวตฺถุกตา อนภาวํกตา อายติํ อนุปฺปาทธมฺมา ทิฏฺเฐว ธมฺเม สุขํ วิหรติ อวิฆาตํ อนุปายาสํ อปริฬาหํ, ทิฏฺเฐว ธมฺเม ปรินิพฺพายติ.}

\prose{โทสชา ฯเปฯ. โมหชา ปาปกา อกุสลา ธมฺมา ปหีนา อุจฺฉินฺนมูลา ตาลาวตฺถุกตา อนภาวํกตา อายติํ อนุปฺปาทธมฺมา ทิฏฺเฐว ธมฺเม สุขํ วิหรติ อวิฆาตํ อนุปายาสํ อปริฬาหํ, ทิฏฺเฐว ธมฺเม ปรินิพฺพายติ. อิมานิ โข ภิกฺขเว ตีณิ กุสลมูลานีติ.  นวมํ.}

\csromanpalirecto
\gambhira{ติกนิปาตปาฬิ \textbf{10. อุโปสถสุตฺต}}
\proseitem{71}{\csromanfolio{206}เอวํ เม สุตํ\csromandash{}เอกํ สมยํ ภควา สาวตฺถิยํ วิหรติ ปุพฺพาราเม มิคาร\-มาตุ\-ปาสาเท. อถ โข วิสาขา มิคารมาตา ตทหุโปสเถ เยน ภควา เตนุปสงฺกมิ, อุปสงฺกมิตฺวา ภควนฺตํ อภิวาเทตฺวา เอกมนฺตํ นิสีทิ, เอกมนฺตํ นิสินฺนํ โข วิสาขํ มิคารมาตรํ ภควา เอตทโวจ “หนฺท กุโต นุ ตฺวํ วิสาเข อาคจฺฉสิ ทิวา ทิวสฺสา”ติ. อุโปสถาหํ ภนฺเต อชฺช อุปวสามีติ.}

\prose{ตโย โข เม วิสาเข อุโปสถา. กตเม ตโย? โคปาลกุโปสโถ นิคณฺฐุ\-โปสโถ อริยุโปสโถ. กถญฺจ วิสาเข โคปาลกุโปสโถ โหติ? เสยฺยถาปิ วิสาเข โคปาลโก สายนฺหสมเย สามิกานํ คาโว นิยฺยาเตตฺวา อิติ ปฏิสญฺจิกฺขติ “อชฺช โข คาโว อมุกสฺมิญฺจ อมุกสฺมิญฺจ ปเทเส จริํสุ, อมุกสฺมิญฺจ อมุกสฺมิญฺจ ปเทเส ปานียานิ ปิวิํสุ. เสฺว ทานิ คาโว อมุกสฺมิญฺจ อมุกสฺมิญฺจ ปเทเส จริสฺสนฺติ, อมุกสฺมิญฺจ อมุกสฺมิญฺจ ปเทเส ปานียานิ ปิวิสฺสนฺตี”ติ. เอวเมวํ โข วิสาเข อิเธกจฺโจ อุโปสถิโก อิติ ปฏิสญฺจิกฺขติ “อหํ ขฺวชฺช อิทญฺจิทญฺจ ขาทนียํ ขาทิํ, อิทญฺจิทญฺจ โภชนียํ ภุญฺชิํ. เสฺว ทานาหํ อิทญฺจิทญฺจ ขาทนียํ ขาทิสฺสามิ, อิทญฺจิทญฺจ โภชนียํ ภุญฺชิสฺสามี”ติ. โส เตน อภิชฺฌา\-สหคเตน เจตสา ทิวสํ อตินาเมติ. เอวํ โข วิสาเข โคปาลกุโปสโถ โหติ. เอวํ อุปวุตฺโถ โข วิสาเข โคปาลกุโปสโถ น มหปฺผโล โหติ น มหานิสํโส น มหาชุติโก น มหาวิปฺผาโร. (1)}

\prose{กถญฺจ วิสาเข นิคณฺฐุ\-โปสโถ โหติ? อตฺถิ วิสาเข นิคณฺฐา นาม สมณชาติกา, เต สาวกํ เอวํ สมาทเปนฺติ “เอหิ ตฺวํ อมฺโภ ปุริส เย ปุรตฺถิมาย ทิสาย ปาณา ปรํ โยชนสตํ, เตสุ ทณฺฑํ นิกฺขิปาหิ. เย ปจฺฉิมาย ทิสาย ปาณา ปรํ โยชนสตํ, เตสุ ทณฺฑํ นิกฺขิปาหิ. เย อุตฺตราย ทิสาย ปาณา ปรํ โยชนสตํ, เตสุ ทณฺฑํ นิกฺขิปาหิ. เย ทกฺขิณาย ทิสาย ปาณา ปรํ โยชนสตํ, เตสุ ทณฺฑํ นิกฺขิปาหี”ติ. อิติ เอกจฺจานํ ปาณานํ อนุทฺทยาย อนุกมฺปาย สมาทเปนฺติ, เอกจฺจานํ ปาณานํ นานุทฺทยาย นานุกมฺปาย สมาทเปนฺติ. เต ตทหุโปสเถ สาวกํ เอวํ สมาทเปนฺติ “เอหิ ตฺวํ อมฺโภ ปุริส}

\vspace{\tipitakaparskip}
\csromancenter{(7) 2. มหาวคฺค สพฺพเจลานิ\csromansharedfootnote{d4c358b7b485108b}{สพฺพเวรานิ (ก)} นิกฺขิปิตฺวา เอวํ วเทหิ ‘นาหํ กฺวจนิ กสฺสจิ กิญฺจนตสฺมิํ\csromansharedfootnote{581d8518f6826195}{กิญฺจนตสฺมิ (?) กิริยาปทเมตํ ยถา กิญฺจนตตฺถีติ.}, น จ มม กฺวจนิ กตฺถจิ กิญฺจนตตฺถี’ติ”. ชานนฺติ โข ปนสฺส มาตาปิตโร “อยํ อมฺหากํ ปุตฺโต”ติ, โสปิ ชานาติ “อิเม มยฺหํ มาตาปิตโร”ติ. ชานาติ โข ปนสฺส ปุตฺตทาโร “อยํ มยฺหํ ภตฺตา”ติ, โสปิ ชานาติ “อยํ มยฺหํ ปุตฺตทาโร”ติ. ชานนฺติ โข ปนสฺส ทาส\-กมฺมกร\-โปริสา “อยํ อมฺหากํ อยฺโย”ติ, โสปิ ชานาติ “อิเม มยฺหํ ทาส\-กมฺมกร\-โปริสา”ติ. อิติ ยสฺมิํ สมเย สจฺเจ สมาทเปตพฺพา, มุสาวาเท ตสฺมิํ สมเย สมาทเปนฺติ, อิทํ ตสฺส มุสาวาทสฺมิํ วทามิ. โส ตสฺสา รตฺติยา อจฺจเยน โภเค อทินฺนํเยว ปริภุญฺชติ, อิทํ ตสฺส อทินฺนาทานสฺมิํ วทามิ. เอวํ โข วิสาเข นิคณฺฐุ\-โปสโถ โหติ, เอวํ อุปวุตฺโถ โข วิสาเข นิคณฺฐุ\-โปสโถ น มหปฺผโล โหติ น มหานิสํโส น มหาชุติโก น มหาวิปฺผาโร. (2)}
\prose{\csromanfolio{207}กถญฺจ วิสาเข อริยุโปสโถ โหติ? อุปกฺกิลิฏฺฐสฺส วิสาเข จิตฺตสฺส อุปกฺกเมน ปริโยทปนา\csromansharedfootnote{a78d6cf9994ba8fd}{ปริโยทาปนา (?)} โหติ. กถญฺจ วิสาเข อุปกฺกิลิฏฺฐสฺส จิตฺตสฺส อุปกฺกเมน ปริโยทปนา โหติ? อิธ วิสาเข อริยสาวโก ตถาคตํ อนุสฺสรติ “อิติปิ โส ภควา อรหํ สมฺมาสมฺพุทฺโธ วิชฺชา\-จรณ\-สมฺปนฺโน สุคโต โลกวิทู อนุตฺตโร ปุริส\-ทมฺม\-สารถิ สตฺถา เทวมนุสฺสานํ พุทฺโธ ภควา”ติ. ตสฺส ตถาคตํ อนุสฺสรโต จิตฺตํ ปสีทติ, ปาโมชฺชํ อุปฺปชฺชติ, เย จิตฺตสฺส อุปกฺกิเลสา, เต ปหียนฺติ. เสยฺยถาปิ วิสาเข อุปกฺกิลิฏฺฐสฺส สีสสฺส อุปกฺกเมน ปริโยทปนา โหติ.}

\prose{กถญฺจ วิสาเข อุปกฺกิลิฏฺฐสฺส สีสสฺส อุปกฺกเมน ปริโยทปนา โหติ? กกฺกญฺจ ปฏิจฺจ, มตฺติกญฺจ ปฏิจฺจ, อุทกญฺจ ปฏิจฺจ, ปุริสสฺส จ ตชฺชํ วายามํ ปฏิจฺจ. เอวํ โข วิสาเข อุปกฺกิลิฏฺฐสฺส สีสสฺส อุปกฺกเมน ปริโยทปนา โหติ. เอวเมวํ โข วิสาเข อุปกฺกิลิฏฺฐสฺส จิตฺตสฺส อุปกฺกเมน ปริโยทปนา โหติ.}

\prose{กถญฺจ วิสาเข อุปกฺกิลิฏฺฐสฺส จิตฺตสฺส อุปกฺกเมน ปริโยทปนา โหติ? อิธ วิสาเข อริยสาวโก ตถาคตํ อนุสฺสรติ “อิติปิ \csromanfolio{208}โส ภควา อรหํ สมฺมาสมฺพุทฺโธ วิชฺชา\-จรณ\-สมฺปนฺโน สุคโต โลกวิทู อนุตฺตโร ปุริส\-ทมฺม\-สารถิ สตฺถา เทวมนุสฺสานํ พุทฺโธ ภควา”ติ. ตสฺส ตถาคตํ อนุสฺสรโต จิตฺตํ ปสีทติ, ปาโมชฺชํ อุปฺปชฺชติ, เย จิตฺตสฺส อุปกฺกิเลสา, เต ปหียนฺติ. อยํ วุจฺจติ วิสาเข อริยสาวโก พฺรหฺมุโปสถํ อุปวสติ, พฺรหฺมุนา สทฺธิํ สํวสติ, พฺรหฺมญฺจสฺส\csromansharedfootnote{85a51dee17a65175}{พฺรหฺมญฺจ (ก)} อารพฺภ จิตฺตํ ปสีทติ, ปาโมชฺชํ อุปฺปชฺชติ, เย จิตฺตสฺส อุปกฺกิเลสา, เต ปหียนฺติ. เอวํ โข วิสาเข อุปกฺกิลิฏฺฐสฺส จิตฺตสฺส อุปกฺกเมน ปริโยทปนา โหติ. (3-ก)}

\prose{อุปกฺกิลิฏฺฐสฺส วิสาเข จิตฺตสฺส อุปกฺกเมน ปริโยทปนา โหติ. กถญฺจ วิสาเข อุปกฺกิลิฏฺฐสฺส จิตฺตสฺส อุปกฺกเมน ปริโยทปนา โหติ? อิธ วิสาเข อริยสาวโก ธมฺมํ อนุสฺสรติ “สฺวากฺขาโต ภควตา ธมฺโม สนฺทิฏฺฐิโก อกาลิโก เอหิปสฺสิโก โอปเนยฺยิโก ปจฺจตฺตํ เวทิตพฺโพ วิญฺญูหี”ติ. ตสฺส ธมฺมํ อนุสฺสรโต จิตฺตํ ปสีทติ, ปาโมชฺชํ อุปฺปชฺชติ, เย จิตฺตสฺส อุปกฺกิเลสา, เต ปหียนฺติ. เสยฺยถาปิ วิสาเข อุปกฺกิลิฏฺฐสฺส กายสฺส อุปกฺกเมน ปริโยทปนา โหติ.}

\prose{กถญฺจ วิสาเข อุปกฺกิลิฏฺฐสฺส กายสฺส อุปกฺกเมน ปริโยทปนา โหติ, โสตฺติญฺจ ปฏิจฺจ, จุณฺณญฺจ ปฏิจฺจ, อุทกญฺจ ปฏิจฺจ, ปุริสสฺส จ ตชฺชํ วายามํ ปฏิจฺจ. เอวํ โข วิสาเข อุปกฺกิลิฏฺฐสฺส กายสฺส อุปกฺกเมน ปริโยทปนา โหติ. เอวเมวํ โข วิสาเข อุปกฺกิลิฏฺฐสฺส จิตฺตสฺส อุปกฺกเมน ปริโยทปนา โหติ. กถญฺจ วิสาเข อุปกฺกิลิฏฺฐสฺส จิตฺตสฺส อุปกฺกเมน ปริโยทปนา โหติ, อิธ วิสาเข อริยสาวโก ธมฺมํ อนุสฺสรติ “สฺวากฺขาโต ภควตา ธมฺโม สนฺทิฏฺฐิโก อกาลิโก เอหิปสฺสิโก โอปเนยฺยิโก ปจฺจตฺตํ เวทิตพฺโพ วิญฺญูหี”ติ. ตสฺส ธมฺมํ อนุสฺสรโต จิตฺตํ ปสีทติ, ปาโมชฺชํ อุปฺปชฺชติ, เย จิตฺตสฺส อุปกฺกิเลสา, เต ปหียนฺติ. อยํ วุจฺจติ วิสาเข อริยสาวโก ธมฺมุโปสถํ อุปวสติ, ธมฺเมน สทฺธิํ สํวสติ, ธมฺมญฺจสฺส อารพฺภ จิตฺตํ ปสีทติ, ปาโมชฺชํ อุปฺปชฺชติ, เย จิตฺตสฺส อุปกฺกิเลสา, เต ปหียนฺติ. เอวํ โข วิสาเข อุปกฺกิลิฏฺฐสฺส จิตฺตสฺส อุปกฺกเมน ปริโยทปนา โหติ. (3-ข)}

\chapterheadpageread{(7) 2. มหาวคฺค}
\prose{\csromanfolio{209}อุปกฺกิลิฏฺฐสฺส วิสาเข จิตฺตสฺส อุปกฺกเมน ปริโยทปนา โหติ. กถญฺจ วิสาเข อุปกฺกิลิฏฺฐสฺส จิตฺตสฺส อุปกฺกเมน ปริโยทปนา โหติ? อิธ วิสาเข อริยสาวโก สํฆํ อนุสฺสรติ “สุปฺปฏิปนฺโน ภควโต สาวกสํโฆ, อุชุปฺปฏิปนฺโน ภควโต สาวกสํโฆ, ญายปฺปฏิปนฺโน ภควโต สาวกสํโฆ, สามีจิ\-ปฺปฏิปนฺโน ภควโต สาวกสํโฆ, ยทิทํ จตฺตาริ ปุริสยุคานิ อฏฺฐ ปุริสปุคฺคลา, เอส ภควโต สาวกสํโฆ อาหุเนยฺโย ปาหุเนยฺโย ทกฺขิเณยฺโย อญฺชลิกรณีโย อนุตฺตรํ ปุญฺญกฺเขตฺตํ โลกสฺสา”ติ. ตสฺส สํฆํ อนุสฺสรโต จิตฺตํ ปสีทติ, ปาโมชฺชํ อุปฺปชฺชติ, เย จิตฺตสฺส อุปกฺกิเลสา, เต ปหียนฺติ. เสยฺยถาปิ วิสาเข อุปกฺกิลิฏฺฐสฺส วตฺถสฺส อุปกฺกเมน ปริโยทปนา โหติ.}

\prose{กถญฺจ วิสาเข อุปกฺกิลิฏฺฐสฺส วตฺถสฺส อุปกฺกเมน ปริโยทปนา โหติ, อุสฺมญฺจ\csromansharedfootnote{b78b9fb1d2850b90}{อูสญฺจ (สฺยา, กํ, อฏฺฐกถายมฺปิ ปาฐนฺตรํ, สํ 2. 107 ปิฏฺเฐ เขมกสุตฺตปาฬิยาปิ สเมติ.) อุสุมญฺจ (สี)} ปฏิจฺจ, ขารญฺจ ปฏิจฺจ, โคมยญฺจ ปฏิจฺจ, อุทกญฺจ ปฏิจฺจ, ปุริสสฺส จ ตชฺชํ วายามํ ปฏิจฺจ. เอวํ โข วิสาเข อุปกฺกิลิฏฺฐสฺส วตฺถสฺส อุปกฺกเมน ปริโยทปนา โหติ. เอวเมวํ โข วิสาเข อุปกฺกิลิฏฺฐสฺส จิตฺตสฺส อุปกฺกเมน ปริโยทปนา โหติ.}

\prose{กถญฺจ วิสาเข อุปกฺกิลิฏฺฐสฺส จิตฺตสฺส อุปกฺกเมน ปริโยทปนา โหติ, อิธ วิสาเข อริยสาวโก สํฆํ อนุสฺสรติ “สุปฺปฏิปนฺโน ภควโต สาวกสํโฆ ฯเปฯ อนุตฺตรํ ปุญฺญกฺเขตฺตํ โลกสฺสา”ติ. ตสฺส สํฆํ อนุสฺสรโต จิตฺตํ ปสีทติ, ปาโมชฺชํ อุปฺปชฺชติ, เย จิตฺตสฺส อุปกฺกิเลสา, เต ปหียนฺติ. อยํ วุจฺจติ วิสาเข อริยสาวโก สํฆุโปสถํ อุปวสติ, สํเฆน สทฺธิํ สํวสติ, สํฆญฺจสฺส อารพฺภ จิตฺตํ ปสีทติ, ปาโมชฺชํ อุปฺปชฺชติ, เย จิตฺตสฺส อุปกฺกิเลสา, เต ปหียนฺติ. เอวํ โข วิสาเข อุปกฺกิลิฏฺฐสฺส จิตฺตสฺส อุปกฺกเมน ปริโยทปนา โหติ. (3-ค)}

\prose{อุปกฺกิลิฏฺฐสฺส วิสาเข จิตฺตสฺส อุปกฺกเมน ปริโยทปนา โหติ. กถญฺจ วิสาเข อุปกฺกิลิฏฺฐสฺส จิตฺตสฺส อุปกฺกเมน ปริโยทปนา โหติ? อิธ วิสาเข อริยสาวโก อตฺตโน สีลานิ อนุสฺสรติ}

\prose{\csromanfolio{210}อขณฺฑานิ อจฺฉิทฺทานิ อสพลานิ อกมฺมาสานิ ภุชิสฺสานิ วิญฺญุ\-ปฺปสตฺถานิ อปรามฏฺฐานิ สมาธิ\-สํวตฺตนิกานิ. ตสฺส สีลํ อนุสฺสรโต จิตฺตํ ปสีทติ, ปาโมชฺชํ อุปฺปชฺชติ, เย จิตฺตสฺส อุปกฺกิเลสา, เต ปหียนฺติ. เสยฺยถาปิ วิสาเข อุปกฺกิลิฏฺฐสฺส อาทาสสฺส อุปกฺกเมน ปริโยทปนา โหติ.}

\prose{กถญฺจ วิสาเข อุปกฺกิลิฏฺฐสฺส อาทาสสฺส อุปกฺกเมน ปริโยทปนา โหติ? เตลญฺจ ปฏิจฺจ, ฉาริกญฺจ ปฏิจฺจ, วาลณฺฑุปกญฺจ ปฏิจฺจ, ปุริสสฺส จ ตชฺชํ วายามํ ปฏิจฺจ. เอวํ โข วิสาเข อุปกฺกิลิฏฺฐสฺส อาทาสสฺส อุปกฺกเมน ปริโยทปนา โหติ. เอวเมวํ โข วิสาเข อุปกฺกิลิฏฺฐสฺส จิตฺตสฺส อุปกฺกเมน ปริโยทปนา โหติ.}

\prose{กถญฺจ วิสาเข อุปกฺกิลิฏฺฐสฺส จิตฺตสฺส อุปกฺกเมน ปริโยทปนา โหติ? อิธ วิสาเข อริยสาวโก อตฺตโน สีลานิ อนุสฺสรติ อขณฺฑานิ ฯเปฯ สมาธิ\-สํวตฺตนิกานิ. ตสฺส สีลํ อนุสฺสรโต จิตฺตํ ปสีทติ, ปาโมชฺชํ อุปฺปชฺชติ, เย จิตฺตสฺส อุปกฺกิเลสา, เต ปหียนฺติ. อยํ วุจฺจติ วิสาเข อริยสาวโก สีลุโปสถํ อุปวสติ, สีเลน สทฺธิํ สํวสติ, สีลญฺจสฺส อารพฺภ จิตฺตํ ปสีทติ, ปาโมชฺชํ อุปฺปชฺชติ, เย จิตฺตสฺส อุปกฺกิเลสา, เต ปหียนฺติ. เอวํ โข วิสาเข อุปกฺกิลิฏฺฐสฺส จิตฺตสฺส อุปกฺกเมน ปริโยทปนา โหติ. (\csromansharedfootnote{611572329e7d70b9}{ตตฺถุปฺปนฺนา (สี, อิ)}-ฆ)}

\prose{อุปกฺกิลิฏฺฐสฺส วิสาเข จิตฺตสฺส อุปกฺกเมน ปริโยทปนา โหติ. กถญฺจ วิสาเข อุปกฺกิลิฏฺฐสฺส จิตฺตสฺส อุปกฺกเมน ปริโยทปนา โหติ? อิธ วิสาเข อริยสาวโก เทวตา อนุสฺสรติ “สนฺติ เทวา จาตุมหาราชิกา\csromansharedfootnote{1560c226e04cb43c}{จาตุมฺมหาราชิกา (สี, สฺยา, กํ, อิ)}, สนฺติ เทวา ตาวติํสา, สนฺติ เทวา ยามา, สนฺติ เทวา ตุสิตา, สนฺติ เทวา นิมฺมานรติโน, สนฺติ เทวา ปรนิมฺมิต\-วส\-วตฺติโน, สนฺติ เทวา พฺรหฺมกายิกา, สนฺติ เทวา ตตุตฺตริ\csromansharedfootnote{9b03c61a8460ff8f}{ตตุตฺตริํ (สี, อิ)}, ยถารูปาย สทฺธาย สมนฺนาคตา ตา เทวตา อิโต จุตา ตตฺถุปปนฺนา\csromansharedfootnote{611572329e7d70b9}{ตตฺถุปฺปนฺนา (สี, อิ)}, มยฺหมฺปิ ตถารูปา สทฺธา สํวิชฺชติ. ยถารูเปน สีเลน สมนฺนาคตา ตา เทวตา อิโต จุตา ตตฺถุปปนฺนา, มยฺหมฺปิ ตถารูปํ สีลํ สํวิชฺชติ. ยถารูเปน สุเตน สมนฺนาคตา ตา เทวตา อิโต จุตา ตตฺถุปปนฺนา, มยฺหมฺปิ ตถารูปํ สุตํ สํวิชฺชติ. ยถารูเปน จาเคน สมนฺนาคตา ตา เทวตา อิโต จุตา ตตฺถุปปนฺนา, มยฺหมฺปิ ตถารูโป จาโค สํวิชฺชติ. ยถารูปาย}

\vspace{\tipitakaparskip}
\csromancenter{(7) 2. มหาวคฺค ปญฺญาย สมนฺนาคตา ตา เทวตา อิโต จุตา ตตฺถุปปนฺนา, มยฺหมฺปิ ตถารูปา ปญฺญา สํวิชฺชตี”ติ. ตสฺส อตฺตโน จ ตาสญฺจ เทวตานํ สทฺธญฺจ สีลญฺจ สุตญฺจ จาคญฺจ ปญฺญญฺจ อนุสฺสรโต จิตฺตํ ปสีทติ, ปาโมชฺชํ อุปฺปชฺชติ, เย จิตฺตสฺส อุปกฺกิเลสา, เต ปหียนฺติ. เสยฺยถาปิ วิสาเข อุปกฺกิลิฏฺฐสฺส ชาตรูปสฺส อุปกฺกเมน ปริโยทปนา โหติ.}
\prose{\csromanfolio{211}กถญฺจ วิสาเข อุปกฺกิลิฏฺฐสฺส ชาตรูปสฺส อุปกฺกเมน ปริโยทปนา โหติ? อุกฺกญฺจ ปฏิจฺจ, โลณญฺจ ปฏิจฺจ, เครุกญฺจ ปฏิจฺจ, นาฬิกสณฺฑาสญฺจ\csromansharedfootnote{f15a668fa9523253}{นาฬิกญฺจ ปฏิจฺจ สณฺฑาสญฺจ (อิ, ก)} ปฏิจฺจ, ปุริสสฺส จ ตชฺชํ วายามํ ปฏิจฺจ. เอวํ โข วิสาเข อุปกฺกิลิฏฺฐสฺส ชาตรูปสฺส อุปกฺกเมน ปริโยทปนา โหติ. เอวเมวํ โข วิสาเข อุปกฺกิลิฏฺฐสฺส จิตฺตสฺส อุปกฺกเมน ปริโยทปนา โหติ.}

\prose{กถญฺจ วิสาเข อุปกฺกิลิฏฺฐสฺส จิตฺตสฺส อุปกฺกเมน ปริโยทปนา โหติ? อิธ วิสาเข อริยสาวโก เทวตา อนุสฺสรติ “สนฺติ เทวา จาตุมหาราชิกา. สนฺติ เทวา ตาวติํสา ฯเปฯ สนฺติ เทวา ตตุตฺตริ, ยถารูปาย สทฺธาย สมนฺนาคตา ตา เทวตา อิโต จุตา ตตฺถุปปนฺนา, มยฺหมฺปิ ตถารูปา สทฺธา สํวิชฺชติ, ยถารูเปน สีเลน. สุเตน. จาเคน. ปญฺญาย สมนฺนาคตา ตา เทวตา อิโต จุตา ตตฺถุปปนฺนา, มยฺหมฺปิ ตถารูปา ปญฺญา สํวิชฺชตี”ติ. ตสฺส อตฺตโน จ ตาสญฺจ เทวตานํ สทฺธญฺจ สีลญฺจ สุตญฺจ จาคญฺจ ปญฺญญฺจ อนุสฺสรโต จิตฺตํ ปสีทติ, ปาโมชฺชํ อุปฺปชฺชติ, เย จิตฺตสฺส อุปกฺกิเลสา, เต ปหียนฺติ. อยํ วุจฺจติ วิสาเข อริยสาวโก เทวตุโปสถํ อุปวสติ, เทวตาหิ สทฺธิํ สํวสติ, เทวตา อารพฺภ จิตฺตํ ปสีทติ, ปาโมชฺชํ อุปฺปชฺชติ, เย จิตฺตสฺส อุปกฺกิเลสา, เต ปหียนฺติ. เอวํ โข วิสาเข อุปกฺกิลิฏฺฐสฺส จิตฺตสฺส อุปกฺกเมน ปริโยทปนา โหติ. (3-\^{}อ)}

\prose{ส โข โส วิสาเข อริยสาวโก อิติ ปฏิสญฺจิกฺขติ\csromandash{}ยาวชีวํ อรหนฺโต ปาณาติปาตํ ปหาย ปาณาติปาตา ปฏิวิรตา นิหิตทณฺฑา นิหิตสตฺถา ลชฺชี ทยาปนฺนา สพฺพปาณภูต\-หิตา\-นุกมฺปี วิหรนฺติ, อหมฺปชฺช อิมญฺจ รตฺติํ อิมญฺจ ทิวสํ ปาณาติปาตํ ปหาย ปาณาติปาตา ปฏิวิรโต นิหิตทณฺโฑ นิหิตสตฺโถ ลชฺชี ทยาปนฺโน}

\prose{\csromanfolio{212}สพฺพปาณภูต\-หิตา\-นุกมฺปี วิหรามิ, อิมินาปิ\csromansharedfootnote{f969edf4dcdc5729}{อิมินาปหํ (สี) อํ 3. 79 ปิฏฺเฐปิ.} องฺเคน อรหตํ อนุกโรมิ, อุโปสโถ จ เม อุปวุตฺโถ ภวิสฺสติ.}

\prose{ยาวชีวํ อรหนฺโต อทินฺนาทานํ ปหาย อทินฺนาทานา ปฏิวิรตา ทินฺนาทายี ทินฺน\-ปาฏิกงฺขี อเถเนน สุจิภูเตน อตฺตนา วิหรนฺติ, อหมฺปชฺช อิมญฺจ รตฺติํ อิมญฺจ ทิวสํ อทินฺนาทานํ ปหาย อทินฺนาทานา ปฏิวิรโต ทินฺนาทายี ทินฺน\-ปาฏิกงฺขี อเถเนน สุจิภูเตน อตฺตนา วิหรามิ, อิมินาปิ องฺเคน อรหตํ อนุกโรมิ, อุโปสโถ จ เม อุปวุตฺโถ ภวิสฺสติ.}

\prose{ยาวชีวํ อรหนฺโต อพฺรหฺมจริยํ ปหาย พฺรหฺมจารี อาราจารี\csromansharedfootnote{c170a12139b6eb31}{อนาจารี (อิ)} วิรตา เมถุนา คามธมฺมา, อหมฺปชฺช อิมญฺจ รตฺติํ อิมญฺจ ทิวสํ อพฺรหฺมจริยํ ปหาย พฺรหฺมจารี อาราจารี วิรโต เมถุนา คามธมฺมา, อิมินาปิ องฺเคน อรหตํ อนุกโรมิ, อุโปสโถ จ เม อุปวุตฺโถ ภวิสฺสติ.}

\prose{ยาวชีวํ อรหนฺโต มุสาวาทํ ปหาย มุสาวาทา ปฏิวิรตา สจฺจวาที สจฺจสนฺธา เถตา ปจฺจยิกา อวิสํวาทกา โลกสฺส, อหมฺปชฺช อิมญฺจ รตฺติํ อิมญฺจ ทิวสํ มุสาวาทํ ปหาย มุสาวาทา ปฏิวิรโต สจฺจวาที สจฺจสนฺโธ เถโต ปจฺจยิโก อวิสํวาทโก โลกสฺส, อิมินาปิ องฺเคน อรหตํ อนุกโรมิ, อุโปสโถ จ เม อุปวุตฺโถ ภวิสฺสติ.}

\prose{ยาวชีวํ อรหนฺโต สุรา\-เมรย\-มชฺช\-ปมาทฏฺฐานํ ปหาย สุรา\-เมรย\-มชฺช\-ปมาทฏฺฐานา ปฏิวิรตา, อหมฺปชฺช อิมญฺจ รตฺติํ อิมญฺจ ทิวสํ สุรา\-เมรย\-มชฺช\-ปมาทฏฺฐานํ ปหาย สุรา\-เมรย\-มชฺช\-ปมาทฏฺฐานา ปฏิวิรโต, อิมินาปิ องฺเคน อรหตํ อนุกโรมิ, อุโปสโถ จ เม อุปวุตฺโถ ภวิสฺสติ.}

\prose{ยาวชีวํ อรหนฺโต เอกภตฺติกา รตฺตูปรตา วิรตา วิกาลโภชนา, อหมฺปชฺช อิมญฺจ รตฺติํ อิมญฺจ ทิวสํ เอกภตฺติโก รตฺตูปรโต วิรโต วิกาลโภชนา, อิมินาปิ องฺเคน อรหตํ อนุกโรมิ, อุโปสโถ จ เม อุปวุตฺโถ ภวิสฺสติ.}

\prose{ยาวชีวํ อรหนฺโต นจฺจ คีต วาทิต วิสูก ทสฺสน มาลาคนฺธ วิเลปน\-ธารณ มณฺฑน วิภูสน\-ฏฺฐานา ปฏิวิรตา, อหมฺปชฺช อิมญฺจ รตฺติํ อิมญฺจ ทิวสํ นจฺจ คีต วาทิต วิสูก ทสฺสน มาลาคนฺธ วิเลปน ธารณ\-มณฺฑน\-วิภูสน\-ฏฺฐานา}

\vspace{\tipitakaparskip}
\csromancenter{(7) 2. มหาวคฺค ปฏิวิรโต, อิมินาปิ องฺเคน อรหตํ อนุกโรมิ, อุโปสโถ จ เม อุปวุตฺโถ ภวิสฺสติ.}
\prose{\csromanfolio{213}ยาวชีวํ อรหนฺโต อุจฺจาสยน\-มหาสยนํ ปหาย อุจฺจาสยน\-มหาสยนา ปฏิวิรตา นีจเสยฺยํ กปฺเปนฺติ มญฺจเก วา ติณสนฺถารเก วา, อหมฺปชฺช อิมญฺจ รตฺติํ อิมญฺจ ทิวสํ อุจฺจาสยน\-มหาสยนํ ปหาย อุจฺจาสยน\-มหาสยนา ปฏิวิรโต นีจเสยฺยํ กปฺเปมิ มญฺจเก วา ติณสนฺถารเก วา, อิมินาปิ องฺเคน อรหตํ อนุกโรมิ, อุโปสโถ จ เม อุปวุตฺโถ ภวิสฺสตีติ.}

\prose{เอวํ โข วิสาเข อริยุโปสโถ โหติ. เอวํ อุปวุตฺโถ โข วิสาเข อริยุโปสโถ มหปฺผโล โหติ มหานิสํโส มหาชุติโก มหาวิปฺผาโร.}

\prose{กีวมหปฺผโล โหติ, กีวมหานิสํโส, กีวมหาชุติโก, กีวมหาวิปฺผาโร? เสยฺยถาปิ วิสาเข โย อิเมสํ โสฬสนฺนํ มหา\-ชนปทานํ ปหูตรตฺตรตนานํ\csromansharedfootnote{6d6d848054ab45f4}{ปหูตสตฺตรตนานํ (กสี, สฺยา, กํ, อิ) ฏีกายํ ทสฺสิตปาฬิเยว. อํ 3. 81 ปิฏฺเฐปิ. 2. มจฺจานํ (ก)} อิสฺสริยาธิปจฺจํ รชฺชํ กาเรยฺย. เสยฺยถิทํ? องฺคานํ มคธานํ กาสีนํ โกสลานํ วชฺชีนํ มลฺลานํ เจตีนํ วงฺคานํ กุรูนํ ปญฺจาลานํ มจฺฉานํ\csromansharedfootnote{56a05de1fef1fe93}{เอกํ (ก)} สูรเสนานํ อสฺสกานํ อวนฺตีนํ คนฺธารานํ กมฺโพชานํ. อฏฺฐงฺค\-สมนฺนาคตสฺส อุโปสถสฺส เอตํ\csromansharedfootnote{1ec1fcd3a43233f9}{มจฺจานํ (ก)} กลํ นาคฺฆติ โสฬสิํ. ตํ กิสฺส เหตุ? กปณํ วิสาเข มานุสกํ รชฺชํ ทิพฺพํ สุขํ อุปนิธาย.}

\prose{ยานิ วิสาเข มานุสกานิ ปญฺญาส วสฺสานิ, จาตุมหาราชิกานํ เทวานํ เอโส เอโก รตฺตินฺทิโว\csromansharedfootnote{a50ebba343cb6852}{รตฺติทิโว (ก)}, ตาย รตฺติยา ติํสรตฺติโย มาโส, เตน มาเสน ทฺวาทสมาสิโย สํวจฺฉโร, เตน สํวจฺฉเรน ทิพฺพานิ ปญฺจ วสฺสสตานิ จาตุมหาราชิกานํ เทวานํ อายุปฺปมาณํ. ฐานํ โข ปเนตํ วิสาเข วิชฺชติ, ยํ อิเธกจฺโจ อิตฺถี วา ปุริโส วา อฏฺฐงฺค\-สมนฺนาคตํ อุโปสถํ อุปวสิตฺวา กายสฺส เภทา ปรํ มรณา จาตุมหาราชิกานํ เทวานํ สหพฺยตํ อุปปชฺเชยฺย. อิทํ โข ปเนตํ วิสาเข สนฺธาย ภาสิตํ กปณํ มานุสกํ รชฺชํ ทิพฺพํ สุขํ อุปนิธาย.}

\prose{\csromanfolio{214}ยํ วิสาเข มานุสกํ วสฺสสตํ, ตาวติํสานํ เทวานํ เอโส เอโก รตฺตินฺทิโว, ตาย รตฺติยา ติํสรตฺติโย มาโส, เตน มาเสน ทฺวาทสมาสิโย สํวจฺฉโร, เตน สํวจฺฉเรน ทิพฺพํ วสฺสสหสฺสํ ตาวติํสานํ เทวานํ อายุปฺปมาณํ. ฐานํ โข ปเนตํ วิสาเข วิชฺชติ, ยํ อิเธกจฺโจ อิตฺถี วา ปุริโส วา อฏฺฐงฺค\-สมนฺนาคตํ อุโปสถํ อุปวสิตฺวา กายสฺส เภทา ปรํ มรณา ตาวติํสานํ เทวานํ สหพฺยตํ อุปปชฺเชยฺย. อิทํ โข ปเนตํ วิสาเข สนฺธาย ภาสิตํ กปณํ มานุสกํ รชฺชํ ทิพฺพํ สุขํ อุปนิธาย.}

\prose{ยานิ วิสาเข มานุสกานิ เทฺว วสฺสสตานิ, ยามานํ เทวานํ เอโส เอโก รตฺตินฺทิโว, ตาย รตฺติยา ติํสรตฺติโย มาโส, เตน มาเสน ทฺวาทสมาสิโย สํวจฺฉโร, เตน สํวจฺฉเรน ทิพฺพานิ เทฺว วสฺสสหสฺสานิ ยามานํ เทวานํ อายุปฺปมาณํ. ฐานํ โข ปเนตํ วิสาเข วิชฺชติ, ยํ อิเธกจฺโจ อิตฺถี วา ปุริโส วา อฏฺฐงฺค\-สมนฺนาคตํ อุโปสถํ อุปวสิตฺวา กายสฺส เภทา ปรํ มรณา ยามานํ เทวานํ สหพฺยตํ อุปปชฺเชยฺย. อิทํ โข ปเนตํ วิสาเข สนฺธาย ภาสิตํ กปณํ มานุสกํ รชฺชํ ทิพฺพํ สุขํ อุปนิธาย.}

\prose{ยานิ วิสาเข มานุสกานิ จตฺตาริ วสฺสสตานิ, ตุสิตานํ เทวานํ เอโส เอโก รตฺตินฺทิโว, ตาย รตฺติยา ติํสรตฺติโย มาโส, เตน มาเสน ทฺวาทสมาสิโย สํวจฺฉโร, เตน สํวจฺฉเรน ทิพฺพานิ จตฺตาริ วสฺสสหสฺสานิ ตุสิตานํ เทวานํ อายุปฺปมาณํ. ฐานํ โข ปเนตํ วิสาเข วิชฺชติ, ยํ อิเธกจฺโจ อิตฺถี วา ปุริโส วา อฏฺฐงฺค\-สมนฺนาคตํ อุโปสถํ อุปวสิตฺวา กายสฺส เภทา ปรํ มรณา ตุสิตานํ เทวานํ สหพฺยตํ อุปปชฺเชยฺย. อิทํ โข ปเนตํ วิสาเข สนฺธาย ภาสิตํ กปณํ มานุสกํ รชฺชํ ทิพฺพํ สุขํ อุปนิธาย.}

\prose{ยานิ วิสาเข มานุสกานิ อฏฺฐ วสฺสสตานิ, นิมฺมานรตีนํ เทวานํ เอโส เอโก รตฺตินฺทิโว, ตาย รตฺติยา ติํสรตฺติโย มาโส, เตน มาเสน ทฺวาทสมาสิโย สํวจฺฉโร, เตน สํวจฺฉเรน ทิพฺพานิ อฏฺฐ วสฺสสหสฺสานิ นิมฺมานรตีนํ เทวานํ อายุปฺปมาณํ. ฐานํ โข ปเนตํ วิสาเข วิชฺชติ, ยํ อิเธกจฺโจ อิตฺถี วา ปุริโส วา อฏฺฐงฺค\-สมนฺนาคตํ อุโปสถํ อุปวสิตฺวา กายสฺส เภทา ปรํ มรณา นิมฺมานรตีนํ เทวานํ}

\vspace{\tipitakaparskip}
\csromancenter{(7) 2. มหาวคฺค สหพฺยตํ อุปปชฺเชยฺย. อิทํ โข ปเนตํ วิสาเข สนฺธาย ภาสิตํ กปณํ มานุสกํ รชฺชํ ทิพฺพํ สุขํ อุปนิธาย.}
\prose{\csromanfolio{215}ยานิ วิสาเข มานุสกานิ โสฬส วสฺสสตานิ, ปรนิมฺมิต\-วส\-วตฺตีนํ เทวานํ เอโส เอโก รตฺตินฺทิโว, ตาย รตฺติยา ติํสรตฺติโย มาโส, เตน มาเสน ทฺวาทสมาสิโย สํวจฺฉโร, เตน สํวจฺฉเรน ทิพฺพานิ โสฬส วสฺสสหสฺสานิ ปรนิมฺมิต\-วส\-วตฺตีนํ เทวานํ อายุปฺปมาณํ. ฐานํ โข ปเนตํ วิสาเข วิชฺชติ, ยํ อิเธกจฺโจ อิตฺถี วา ปุริโส วา อฏฺฐงฺค\-สมนฺนาคตํ อุโปสถํ อุปวสิตฺวา กายสฺส เภทา ปรํ มรณา ปรนิมฺมิต\-วส\-วตฺตีนํ เทวานํ สหพฺยตํ อุปปชฺเชยฺย. อิทํ โข ปเนตํ วิสาเข สนฺธาย ภาสิตํ กปณํ มานุสกํ รชฺชํ ทิพฺพํ สุขํ อุปนิธายาติ.}

\setlength{\csromangathapagewidth}{0pt}
\setlength{\csromangathawakwidth}{0pt}
\csromangathameasuregroup{อฏฺฐงฺคุเปตสฺส อุโปสถสฺส,}{ปาณํ น หญฺเญ\textsuperscript{1} น จทินฺนมาทิเย, \\ มุสา น ภาเส น จ มชฺชโป สิยา. \\ อพฺรหฺมจริยา วิรเมยฺย เมถุนา, \\ รตฺติํ น ภุญฺเชยฺย วิกาลโภชนํ. \\ มาลํ น ธาเร น จ คนฺธมาจเร, \\ มญฺเจ ฉมายํ ว สเยถ สนฺถเต. \\ เอตํ หิ อฏฺฐงฺคิกมาหุโปสถํ, \\ พุทฺเธน ทุกฺขนฺตคุนา ปกาสิตํ. \\ จนฺโท จ สุริโย จ อุโภ สุทสฺสนา, \\ โอภาสยํ อนุปริยนฺติ ยาวตา. \\ ตโมนุทา เต ปน อนฺตลิกฺขคา, \\ นเภ ปภาสนฺติ ทิสาวิโรจนา. \\ เอตสฺมิํ ยํ วิชฺชติ อนฺตเร ธนํ, \\ มุตฺตา มณิ เวฬุริยญฺจ ภทฺทกํ. \\ สิงฺคีสุวณฺณํ อถ วาปิ กญฺจนํ, \\ ยํ ชาตรูปํ “หฏกนฺ”ติ วุจฺจติ. \\ \csromangathabat{อฏฺฐงฺคุเปตสฺส อุโปสถสฺส,}{กลมฺปิ เต นานุภวนฺติ โสฬสิํ.} \\ จนฺทปฺปภา ตารคณา จ สพฺเพ. \\ ตสฺมา หิ นารี จ นโร จ สีลวา, \\ อฏฺฐงฺคุเปตํ อุปวสฺสุโปสถํ. \\ ปุญฺญานิ กตฺวาน สุขุทฺรยานิ, \\ อนินฺทิตา สคฺคมุเปนฺติ ฐานนฺติ.ทสมํ.}
\csromangathameasurewak{ปาณํ น หญฺเญ\textsuperscript{1} น จทินฺนมาทิเย, \\ มุสา น ภาเส น จ มชฺชโป สิยา. \\ อพฺรหฺมจริยา วิรเมยฺย เมถุนา, \\ รตฺติํ น ภุญฺเชยฺย วิกาลโภชนํ. \\ มาลํ น ธาเร น จ คนฺธมาจเร, \\ มญฺเจ ฉมายํ ว สเยถ สนฺถเต. \\ เอตํ หิ อฏฺฐงฺคิกมาหุโปสถํ, \\ พุทฺเธน ทุกฺขนฺตคุนา ปกาสิตํ. \\ จนฺโท จ สุริโย จ อุโภ สุทสฺสนา, \\ โอภาสยํ อนุปริยนฺติ ยาวตา. \\ ตโมนุทา เต ปน อนฺตลิกฺขคา, \\ นเภ ปภาสนฺติ ทิสาวิโรจนา. \\ เอตสฺมิํ ยํ วิชฺชติ อนฺตเร ธนํ, \\ มุตฺตา มณิ เวฬุริยญฺจ ภทฺทกํ. \\ สิงฺคีสุวณฺณํ อถ วาปิ กญฺจนํ, \\ ยํ ชาตรูปํ “หฏกนฺ”ติ วุจฺจติ. \\ อฏฺฐงฺคุเปตํ อุปวสฺสุโปสถํ. \\ ปุญฺญานิ กตฺวาน สุขุทฺรยานิ, \\ อนินฺทิตา สคฺคมุเปนฺติ ฐานนฺติ.ทสมํ.}
\csromangathasetleft{อฏฺฐงฺคุเปตสฺส อุโปสถสฺส,}
\csromangathagroupwithclosermidruled{\csromangathastanza{\csromangathawak{ปาณํ น หญฺเญ\csromansharedfootnote{b2b256138e82e2a2}{น หาเน (สี, อิ), น หเน (ก)} น จทินฺนมาทิเย, \\ มุสา น ภาเส น จ มชฺชโป สิยา. \\ อพฺรหฺมจริยา วิรเมยฺย เมถุนา, \\ รตฺติํ น ภุญฺเชยฺย วิกาลโภชนํ.}}\csromangathastanza{\csromangathawak{มาลํ น ธาเร น จ คนฺธมาจเร, \\ มญฺเจ ฉมายํ ว สเยถ สนฺถเต. \\ เอตํ หิ อฏฺฐงฺคิกมา\-หุ\-โปสถํ, \\ พุทฺเธน ทุกฺขนฺตคุนา ปกาสิตํ.}}\csromangathastanza{\csromangathawak{จนฺโท จ สุริโย จ อุโภ สุทสฺสนา, \\ โอภาสยํ อนุปริยนฺติ ยาวตา. \\ ตโมนุทา เต ปน อนฺตลิกฺขคา, \\ นเภ ปภาสนฺติ ทิสาวิโรจนา.}}\csromangathastanza{\csromangathawak{เอตสฺมิํ ยํ วิชฺชติ อนฺตเร ธนํ, \\ มุตฺตา มณิ เวฬุริยญฺจ ภทฺทกํ. \\ สิงฺคีสุวณฺณํ อถ วาปิ กญฺจนํ, \\ ยํ ชาตรูปํ “หฏกนฺ”ติ วุจฺจติ.}}\csromangathastanza{\csromanfolio{216}\csromangathabat{อฏฺฐงฺคุเปตสฺส อุโปสถสฺส,}{กลมฺปิ เต นานุภวนฺติ โสฬสิํ.} \\ จนฺทปฺปภา ตารคณา จ สพฺเพ. \\ ตสฺมา หิ นารี จ นโร จ สีลวา,}\csromangathastanza{\csromangathawak{อฏฺฐงฺคุเปตํ อุปวสฺสุโปสถํ. \\ ปุญฺญานิ กตฺวาน สุขุทฺรยานิ, \\ อนินฺทิตา สคฺคมุเปนฺติ ฐานนฺติ.  ทสมํ.}}}{มหาวคฺโค ทุติโย.}
\csromancenter{\textbf{ตสฺสุทฺทานํ}}
\setlength{\csromangathapagewidth}{0pt}
\setlength{\csromangathawakwidth}{0pt}
\csromangathameasuregroup{ติตฺถภยญฺจ เวนาโค, \\ สาโฬฺห จาปิ กถาวตฺถุ,}{\csromangathabat{ติตฺถภยญฺจ เวนาโค,}{สรโภ เกสมุตฺติยา.} \\ \csromangathabat{สาโฬฺห จาปิ กถาวตฺถุ,}{ติตฺถิยมูลุโปสโถติ.}}
\csromangathasetleft{ติตฺถภยญฺจ เวนาโค, \\ สาโฬฺห จาปิ กถาวตฺถุ,}
\csromangathagroup{\csromangathastanza{\csromangathabat{ติตฺถภยญฺจ เวนาโค,}{สรโภ เกสมุตฺติยา.} \\ \csromangathabat{สาโฬฺห จาปิ กถาวตฺถุ,}{ติตฺถิยมูลุโปสโถติ.}}}

\chapterhead{\textbf{(8) 3. อานนฺทวคฺค}}
\csromanheader{2}{1. \textbf{ฉนฺนสุตฺต}}
\proseitem{72}{เอกํ สมยํ ภควา สาวตฺถิยํ วิหรติ เชตวเน อนาถ\-ปิณฺฑิกสฺส อาราเม. อถ โข ฉนฺโน ปริพฺพาชโก เยนายสฺมา อานนฺโท เตนุปสงฺกมิ, อุปสงฺกมิตฺวา อายสฺมตา อานนฺเทน สทฺธิํ สมฺโมทิ, สมฺโมทนียํ กถํ สารณียํ วีติสาเรตฺวา เอกมนฺตํ นิสีทิ, เอกมนฺตํ นิสินฺโน โข ฉนฺโน ปริพฺพาชโก อายสฺมนฺตํ อานนฺทํ เอตทโวจ “ตุเมฺหปิ อาวุโส อานนฺท ราคสฺส ปหานํ ปญฺญาเปถ, โทสสฺส ปหานํ ปญฺญาเปถ, โมหสฺส ปหานํ ปญฺญาเปถา”ติ. มยํ โข อาวุโส ราคสฺส ปหานํ ปญฺญาเปม, โทสสฺส ปหานํ ปญฺญาเปม, โมหสฺส ปหานํ ปญฺญาเปมาติ.}

\prose{กิํ ปน ตุเมฺห อาวุโส ราเค อาทีนวํ ทิสฺวา ราคสฺส ปหานํ ปญฺญาเปถ, กิํ โทเส อาทีนวํ ทิสฺวา โทสสฺส ปหานํ ปญฺญาเปถ, กิํ โมเห อาทีนวํ ทิสฺวา โมหสฺส ปหานํ ปญฺญาเปถาติ.}

\chapterheadpageread{(8) 3. อานนฺทวคฺค}
\prose{\csromanfolio{217}รตฺโต โข อาวุโส ราเคน อภิภูโต ปริยาทินฺน\-จิตฺโต อตฺตพฺยาพาธายปิ เจเตติ, ปรพฺยาพาธายปิ เจเตติ, อุภย\-พฺยาพาธายปิ เจเตติ, เจตสิกมฺปิ ทุกฺขํ โทมนสฺสํ ปฏิสํเวเทติ. ราเค ปหีเน เนว\-ตฺตพฺยาพาธายปิ เจเตติ, น ปรพฺยาพาธายปิ เจเตติ, น อุภย\-พฺยาพาธายปิ เจเตติ, น เจตสิกํ ทุกฺขํ โทมนสฺสํ ปฏิสํเวเทติ. รตฺโต โข อาวุโส ราเคน อภิภูโต ปริยาทินฺน\-จิตฺโต กาเยน ทุจฺจริตํ จรติ, วาจาย ทุจฺจริตํ จรติ, มนสา ทุจฺจริตํ จรติ. ราเค ปหีเน เนว กาเยน ทุจฺจริตํ จรติ, น วาจาย ทุจฺจริตํ จรติ, น มนสา ทุจฺจริตํ จรติ. รตฺโต โข อาวุโส ราเคน อภิภูโต ปริยาทินฺน\-จิตฺโต อตฺตตฺถมฺปิ ยถาภูตํ นปฺปชานาติ, ปรตฺถมฺปิ ยถาภูตํ นปฺปชานาติ, อุภยตฺถมฺปิ ยถาภูตํ นปฺปชานาติ. ราเค ปหีเน อตฺตตฺถมฺปิ ยถาภูตํ ปชานาติ, ปรตฺถมฺปิ ยถาภูตํ ปชานาติ, อุภยตฺถมฺปิ ยถาภูตํ ปชานาติ. ราโค โข อาวุโส อนฺธกรโณ อจกฺขุกรโณ อญฺญาณกรโณ ปญฺญานิโรธิโก วิฆาต\-ปกฺขิโก อนิพฺพานสํวตฺตนิโก.}

\prose{ทุฏฺโฐ โข อาวุโส โทเสน ฯเปฯ. มูโฬฺห โข อาวุโส โมเหน อภิภูโต ปริยาทินฺน\-จิตฺโต อตฺตพฺยาพาธายปิ เจเตติ, ปรพฺยาพาธายปิ เจเตติ, อุภย\-พฺยาพาธายปิ เจเตติ, เจตสิกมฺปิ ทุกฺขํ โทมนสฺสํ ปฏิสํเวเทติ. โมเห ปหีเน เนว\-ตฺตพฺยาพาธายปิ เจเตติ, น ปรพฺยาพาธายปิ เจเตติ, น อุภย\-พฺยาพาธายปิ เจเตติ, น เจตสิกํ ทุกฺขํ โทมนสฺสํ ปฏิสํเวเทติ. มูโฬฺห โข อาวุโส โมเหน อภิภูโต ปริยาทินฺน\-จิตฺโต กาเยน ทุจฺจริตํ จรติ, วาจาย ทุจฺจริตํ จรติ, มนสา ทุจฺจริตํ จรติ. โมเห ปหีเน เนว กาเยน ทุจฺจริตํ จรติ, น วาจาย ทุจฺจริตํ จรติ, น มนสา ทุจฺจริตํ จรติ. มูโฬฺห โข อาวุโส โมเหน อภิภูโต ปริยาทินฺน\-จิตฺโต อตฺตตฺถมฺปิ ยถาภูตํ นปฺปชานาติ, ปรตฺถมฺปิ ยถาภูตํ นปฺปชานาติ อุภยตฺถมฺปิ ยถาภูตํ นปฺปชานาติ. โมเห ปหีเน อตฺตตฺถมฺปิ ยถาภูตํ ปชานาติ, ปรตฺถมฺปิ ยถาภูตํ ปชานาติ, อุภยตฺถมฺปิ ยถาภูตํ ปชานาติ. โมโห โข อาวุโส อนฺธกรโณ อจกฺขุกรโณ อญฺญาณกรโณ ปญฺญานิโรธิโก วิฆาต\-ปกฺขิโก อนิพฺพานสํวตฺตนิโก. อิทํ โข มยํ อาวุโส ราเค อาทีนวํ ทิสฺวา ราคสฺส}

\prose{\csromanfolio{218}ปหานํ ปญฺญาเปม. อิทํ โทเส อาทีนวํ ทิสฺวา โทสสฺส ปหานํ ปญฺญาเปม. อิทํ โมเห อาทีนวํ ทิสฺวา โมหสฺส ปหานํ ปญฺญาเปมาติ.}

\prose{อตฺถิ ปนาวุโส มคฺโค, อตฺถิ ปฏิปทา เอตสฺส ราคสฺส โทสสฺส โมหสฺส ปหานายาติ. อตฺถาวุโส มคฺโค, อตฺถิ ปฏิปทา เอตสฺส ราคสฺส โทสสฺส โมหสฺส ปหานายาติ. กตโม ปนาวุโส มคฺโค, กตมา ปฏิปทา เอตสฺส ราคสฺส โทสสฺส โมหสฺส ปหานายาติ. อยเมว อริโย อฏฺฐงฺคิโก มคฺโค. เสยฺยถิทํ? สมฺมาทิฏฺฐิ ฯเปฯ สมฺมาสมาธิ. อยํ โข อาวุโส มคฺโค, อยํ ปฏิปทา เอตสฺส ราคสฺส โทสสฺส โมหสฺส ปหานายาติ. ภทฺทโก โข อาวุโส มคฺโค, ภทฺทิกา ปฏิปทา เอตสฺส ราคสฺส โทสสฺส โมหสฺส ปหานาย, อลํ จ ปนาวุโส อานนฺท อปฺปมาทายาติ.  ปฐมํ.}

\csromanheader{2}{2. \textbf{อาชีวกสุตฺต}}
\proseitem{73}{เอกํ สมยํ อายสฺมา อานนฺโท โกสมฺพิยํ วิหรติ โฆสิตาราเม. อถ โข อญฺญตโร อาชีวกสาวโก คหปติ เยนายสฺมา อานนฺโท เตนุปสงฺกมิ, อุปสงฺกมิตฺวา อายสฺมนฺตํ อานนฺทํ อภิวาเทตฺวา เอกมนฺตํ นิสีทิ, เอกมนฺตํ นิสินฺโน โข โส อาชีวกสาวโก คหปติ อายสฺมนฺตํ อานนฺทํ เอตทโวจ\csromandash{}}

\prose{“เกสํ โน ภนฺเต อานนฺท ธมฺโม สฺวากฺขาโต, เก โลเก สุปฺปฏิปนฺนา, เก โลเก สุกตา”ติ\csromansharedfootnote{b20866a9ee71e6ba}{สุคตาติ (สี, สฺยา, กํ, อิ)}. เตน หิ คหปติ ตญฺเญเวตฺถ ปฏิปุจฺฉิสฺสามิ, ยถา เต ขเมยฺย, ตถา นํ พฺยากเรยฺยาสิ. ตํ กิํ มญฺญสิ คหปติ, เย ราคสฺส ปหานาย ธมฺมํ เทเสนฺติ, โทสสฺส ปหานาย ธมฺมํ เทเสนฺติ, โมหสฺส ปหานาย ธมฺมํ เทเสนฺติ, เตสํ ธมฺโม สฺวากฺขาโต โน วา, กถํ วา เต เอตฺถ โหตีติ. เย ภนฺเต ราคสฺส ปหานาย ธมฺมํ เทเสนฺติ, โทสสฺส ปหานาย ธมฺมํ เทเสนฺติ, โมหสฺส ปหานาย ธมฺมํ เทเสนฺติ, เตสํ ธมฺโม สฺวากฺขาโต, เอวํ เม เอตฺถ โหตีติ.}

\prose{ตํ กิํ มญฺญสิ คหปติ, เย ราคสฺส ปหานาย ปฏิปนฺนา, โทสสฺส ปหานาย ปฏิปนฺนา, โมหสฺส ปหานาย ปฏิปนฺนา, เต โลเก}

\vspace{\tipitakaparskip}
\csromancenter{(8) 3. อานนฺทวคฺค สุปฺปฏิปนฺนา โน วา, กถํ วา เต เอตฺถ โหตีติ. เย ภนฺเต ราคสฺส ปหานาย ปฏิปนฺนาย, โทสสฺส ปหานาย ปฏิปนฺนา, โมหสฺส ปหานาย ปฏิปนฺนา, เต โลเก สุปฺปฏิปนฺนา, เอวํ เม เอตฺถ โหตีติ.}
\prose{\csromanfolio{219}ตํ กิํ มญฺญสิ คหปติ, เยสํ ราโค ปหีโน อุจฺฉินฺนมูโล ตาลาวตฺถุกโต อนภาวํกโต อายติํ อนุปฺปาทธมฺโม, เยสํ โทโส ปหีโน อุจฺฉินฺนมูโล ตาลาวตฺถุกโต อนภาวํกโต อายติํ อนุปฺปาทธมฺโม, เยสํ โมโห ปหีโน อุจฺฉินฺนมูโล ตาลาวตฺถุกโต อนภาวํกโต อายติํ อนุปฺปาทธมฺโม, เต โลเก สุกตา โน วา, กถํ วา เต เอตฺถ โหตีติ. เยสํ ภนฺเต ราโค ปหีโน อุจฺฉินฺนมูโล ตาลาวตฺถุกโต อนภาวํกโต อายติํ อนุปฺปาทธมฺโม, เยสํ โทโส ปหีโน ฯเปฯ เยสํ โมโห ปหีโน อุจฺฉินฺนมูโล ตาลาวตฺถุกโต อนภาวํกโต อายติํ อนุปฺปาทธมฺโม, เต โลเก สุกตา, เอวํ เม เอตฺถ โหตีติ.}

\prose{อิติ โข คหปติ ตยาเวตํ\csromansharedfootnote{b433b4ec6145c0d7}{ตยา เจตํ (สี, อิ, ก)} พฺยากตํ “เย ภนฺเต ราคสฺส ปหานาย ธมฺมํ เทเสนฺติ, โทสสฺส ปหานาย ธมฺมํ เทเสนฺติ, โมหสฺส ปหานาย ธมฺมํ เทเสนฺติ, เตสํ ธมฺโม สฺวากฺขาโต”ติ. ตยาเวตํ พฺยากตํ “เย ภนฺเต ราคสฺส ปหานาย ปฏิปนฺนา โทสสฺส ปหานาย ปฏิปนฺนา, โมหสฺส ปหานาย ปฏิปนฺนา, เต โลเก สุปฺปฏิปนฺนา”ติ. ตยาเวตํ พฺยากตํ “เยสํ ภนฺเต ราโค ปหีโน อุจฺฉินฺนมูโล ตาลาวตฺถุกโต อนภาวํกโต อายติํ อนุปฺปาทธมฺโม, เยสํ โทโส ปหีโน ฯเปฯ เยสํ โมโห ปหีโน อุจฺฉินฺนมูโล ตาลาวตฺถุกโต อนภาวํกโต อายติํ อนุปฺปาทธมฺโม, เต โลเก สุกตา”ติ.}

\prose{อจฺฉริยํ ภนฺเต, อพฺภุตํ ภนฺเต, น เจว นาม สธมฺมุ\-กฺกํสนา ภวิสฺสติ, น จ ปรธมฺมา\-ปสาทนา\csromansharedfootnote{ba6112bb50a941c1}{น ปรธมฺมาปสาทนา (สี, อิ), น ปรธมฺมวมฺภนา (ม 2. 193 ปิฏฺเฐ)}, อายตเนว\csromansharedfootnote{93d6f17732c0561b}{อายตเน จ (ม 2. 193 ปิฏฺเฐ)} ธมฺมเทสนา, อตฺโถ จ วุตฺโต, อตฺถา จ อนุปนีโต. ตุเมฺห ภนฺเต อานนฺท ราคสฺส ปหานาย ธมฺมํ เทเสถ, โทสสฺส ฯเปฯ โมหสฺส ปหานาย ธมฺมํ เทเสถ, ตุมฺหากํ ภนฺเต อานนฺท ธมฺโม สฺวากฺขาโต. ตุเมฺห ภนฺเต อานนฺท ราคสฺส \csromanfolio{220}ปหานาย ปฏิปนฺนา, โทสสฺส ฯเปฯ โมหสฺส ปหานาย ปฏิปนฺนา, ตุเมฺห ภนฺเต โลเก สุปฺปฏิปนฺนา. ตุมฺหากํ ภนฺเต อานนฺท ราโค ปหีโน อุจฺฉินฺนมูโล ตาลาวตฺถุกโต อนภาวํกโต อายติํ อนุปฺปาทธมฺโม, ตุมฺหากํ โทโส ปหีโน ฯเปฯ ตุมฺหากํ โมโห ปหีโน อุจฺฉินฺนมูโล ตาลาวตฺถุกโต อนภาวํกโต อายติํ อนุปฺปาทธมฺโม, ตุเมฺห โลเก สุกตา.}

\prose{อภิกฺกนฺตํ ภนฺเต, อภิกฺกนฺตํ ภนฺเต, เสยฺยถาปิ ภนฺเต นิกฺกุชฺชิตํ วา อุกฺกุชฺเชยฺย, ปฏิจฺฉนฺนํ วา วิวเรยฺย, มูฬฺหสฺส วา มคฺคํ อาจิกฺเขยฺย, อนฺธกาเร วา เตลปชฺโชตํ ธาเรยฺย “จกฺขุมนฺโต รูปานิ ทกฺขนฺตี”ติ. เอวเมวํ อยฺเยน อานนฺเทน อเนก\-ปริยาเยน ธมฺโม ปกาสิโต, เอสาหํ ภนฺเต อานนฺท ตํ ภควนฺตํ สรณํ คจฺฉามิ ธมฺมญฺจ ภิกฺขุสํฆญฺจ, อุปาสกํ มํ อยฺโย อานนฺโท ธาเรตุ อชฺชตคฺเค ปาณุเปตํ สรณํ คตนฺติ.  ทุติยํ.}

\csromanheader{2}{3. มหานาม\-สกฺก\-สุตฺต}
\proseitem{74}{เอวํ เม สุตํ\csromandash{}เอกํ สมยํ ภควา สกฺเกสุ วิหรติ กปิล\-วตฺถุสฺมิํ นิโคฺรธาราเม. เตน โข ปน สมเยน ภควา คิลานา วุฏฺฐิโต\csromansharedfootnote{7ad0dbd929e80867}{คิลานวุฏฺฐิโต (สทฺทนีติ)} โหติ อจิรวุฏฺฐิโต เคลญฺญา. อถ โข มหานาโม สกฺโก เยน ภควา เตนุปสงฺกมิ, อุปสงฺกมิตฺวา ภควนฺตํ อภิวาเทตฺวา เอกมนฺตํ นิสีทิ, เอกมนฺตํ นิสินฺโน โข มหานาโม สกฺโก ภควนฺตํ เอตทโวจ “ทีฆรตฺตาหํ ภนฺเต ภควตา เอวํ ธมฺมํ เทสิตํ อาชานามิ ‘สมาหิตสฺส ญาณํ, โน อสมาหิตสฺสา’ติ. สมาธิ นุ โข ภนฺเต ปุพฺเพ, ปจฺฉา ญาณํ, อุทาหุ ญาณํ ปุพฺเพ, ปจฺฉา สมาธี”ติ. อถ โข อายสฺมโต อานนฺทสฺส เอตทโหสิ “ภควา โข คิลานวุฏฺฐิโต อจิรวุฏฺฐิโต เคลญฺญา, อยญฺจ มหานาโม สกฺโก ภควนฺตํ อติคมฺภีรํ ปญฺหํ ปุจฺฉติ, ยํนูนาหํ มหานามํ สกฺกํ เอกมนฺตํ อปเนตฺวา ธมฺมํ เทเสยฺยนฺ”ติ.}

\prose{อถ โข อายสฺมา อานนฺโท มหานามํ สกฺกํ พาหายํ คเหตฺวา เอกมนฺตํ อปเนตฺวา มหานามํ สกฺกํ เอตทโวจ “เสขมฺปิ โข มหานาม สีลํ วุตฺตํ ภควตา, อเสขมฺปิ สีลํ วุตฺตํ ภควตา. เสโขปิ สมาธิ วุตฺโต ภควตา, อเสโขปิ สมาธิ วุตฺโต ภควตา. เสขาปิ ปญฺญา วุตฺตา ภควตา, อเสขาปิ ปญฺญา วุตฺตา ภควตา. กตมญฺจ}

\vspace{\tipitakaparskip}
\csromancenter{(8) 3. อานนฺทวคฺค มหานาม เสขํ สีลํ? อิธ มหานาม ภิกฺขุ สีลวา โหติ, ปาติโมกฺข\-สํวร\-สํวุตฺโต วิหรติ ฯเปฯ สมาทาย สิกฺขติ สิกฺขาปเทสุ. อิทํ วุจฺจติ มหานาม เสขํ สีลํ.}
\prose{\csromanfolio{221}กตโม จ มหานาม เสโข สมาธิ? อิธ มหานาม ภิกฺขุ วิวิจฺเจว กาเมหิ ฯเปฯ จตุตฺถํ ฌานํ อุปสมฺปชฺช วิหรติ. อยํ วุจฺจติ มหานาม เสโข สมาธิ.}

\prose{กตมา จ มหานาม เสขา ปญฺญา? อิธ มหานาม ภิกฺขุ “อิทํ ทุกฺขนฺ”ติ ยถาภูตํ ปชานาติ ฯเปฯ “อยํ ทุกฺขนิโรธ\-คามินี ปฏิปทา”ติ ยถาภูตํ ปชานาติ, อยํ วุจฺจติ มหานาม เสขา ปญฺญา.}

\prose{ส โข โส มหานาม อริยสาวโก เอวํ สีลสมฺปนฺโน เอวํ สมาธิ\-สมฺปนฺโน เอวํ ปญฺญาสมฺปนฺโน อาสวานํ ขยา อนาสวํ เจโตวิมุตฺติํ ปญฺญาวิมุตฺติํ ทิฏฺเฐว ธมฺเม สยํ อภิญฺญา สจฺฉิกตฺวา อุปสมฺปชฺช วิหรติ. เอวํ โข มหานาม เสขมฺปิ สีลํ วุตฺตํ ภควตา, อเสขมฺปิ สีลํ วุตฺตํ ภควตา. เสโขปิ สมาธิ วุตฺโต ภควตา, อเสโขปิ สมาธิ วุตฺโต ภควตา. เสขาปิ ปญฺญา วุตฺตา ภควตา, อเสขาปิ ปญฺญา วุตฺตา ภควตาติ.  ตติยํ.}

\csromanheader{2}{4. \textbf{นิคณฺฐสุตฺต}}
\proseitem{75}{เอกํ สมยํ อายสฺมา อานนฺโท เวสาลิยํ วิหรติ มหาวเน กูฏาคาร\-สาลายํ. อถ โข อภโย จ ลิจฺฉวิ, ปณฺฑิต\-กุมารโก จ ลิจฺฉวิ เยนายสฺมา อานนฺโท เตนุ\-ปสงฺกมิํสุ, อุปสงฺกมิตฺวา อายสฺมนฺตํ อานนฺทํ อภิวาเทตฺวา เอกมนฺตํ นิสีทิํสุ, เอกมนฺตํ นิสินฺโน โข อภโย ลิจฺฉวิ อายสฺมนฺตํ อานนฺทํ เอตทโวจ “นิคณฺโฐ ภนฺเต นาฏปุตฺโต\csromansharedfootnote{be303b1eabc46dd2}{นาถปุตฺโต (สี, อิ)} สพฺพญฺญู สพฺพทสฺสาวี อปริเสสํ ญาณทสฺสนํ ปฏิชานาติ ‘จรโต จ เม ติฏฺฐโต จ สุตฺตสฺส จ ชาครสฺส จ สตตํ สมิตํ ญาณทสฺสนํ ปจฺจุปฏฺฐิตนฺ’ติ. โส ปุราณานํ กมฺมานํ ตปสา พฺยนฺตีภาวํ ปญฺญาเปติ นวานํ กมฺมานํ อกรณา เสตุฆาตํ. อิติ กมฺมกฺขยา ทุกฺขกฺขโย, ทุกฺขกฺขยา เวทนากฺขโย, เวทนากฺขยา สพฺพํ ทุกฺขํ นิชฺชิณฺณํ ภวิสฺสติ, เอวเมติสฺสา สนฺทิฏฺฐิกาย \csromanfolio{222}นิชฺชราย วิสุทฺธิยา สมติกฺกโม โหติ. อิธ ภนฺเต ภควา กิมาหา”ติ.}

\prose{ติสฺโส โข อิมา อภย นิชฺชรา วิสุทฺธิโย เตน ภควตา ชานตา ปสฺสตา อรหตา สมฺมา\-สมฺพุทฺเธน สมฺมทกฺขาตา สตฺตานํ วิสุทฺธิยา โสก\-ปริเทวานํ สมติกฺกมาย ทุกฺข\-โทมนสฺสานํ อตฺถงฺคมาย ญายสฺส อธิคมาย นิพฺพานสฺส สจฺฉิกิริยาย. กตมา ติสฺโส? อิธ อภย ภิกฺขุ สีลวา โหติ ฯเปฯ สมาทาย สิกฺขติ สิกฺขาปเทสุ. โส นวญฺจ กมฺมํ น กโรติ, ปุราณญฺจ กมฺมํ ผุสฺส ผุสฺส พฺยนฺตีกโรติ. สนฺทิฏฺฐิกา นิชฺชรา อกาลิกา เอหิปสฺสิกา โอปเนยฺยิกา ปจฺจตฺตํ เวทิตพฺพา วิญฺญูหีติ.}

\prose{ส โข โส อภย ภิกฺขุ เอวํ สีลสมฺปนฺโน วิวิจฺเจว กาเมหิ ฯเปฯ จตุตฺถํ ฌานํ อุปสมฺปชฺช วิหรติ. โส นวญฺจ กมฺมํ น กโรติ, ปุราณญฺจ กมฺมํ ผุสฺส ผุสฺส พฺยนฺตีกโรติ, สนฺทิฏฺฐิกา นิชฺชรา อกาลิกา เอหิปสฺสิกา โอปเนยฺยิกา ปจฺจตฺตํ เวทิตพฺพา วิญฺญูหีติ.}

\prose{ส โข โส อภย ภิกฺขุ เอวํ สมาธิ\-สมฺปนฺโน\csromansharedfootnote{e71e76d674bdf6c4}{เอวํ สีลสมฺปนฺโน เอวํ สมาธิสมฺปนฺโน (สี, สฺยา, กํ)} อาสวานํ ขยา อนาสวํ เจโตวิมุตฺติํ ปญฺญาวิมุตฺติํ ทิฏฺเฐว ธมฺเม สยํ อภิญฺญา สจฺฉิกตฺวา อุปสมฺปชฺช วิหรติ. โส นวญฺจ กมฺมํ น กโรติ, ปุราณญฺจ กมฺมํ ผุสฺส ผุสฺส พฺยนฺตีกโรติ, สนฺทิฏฺฐิกา นิชฺชรา อกาลิกา เอหิปสฺสิกา โอปเนยฺยิกา ปจฺจตฺตํ เวทิตพฺพา วิญฺญูหีติ. อิมา โข อภย ติสฺโส นิชฺชรา วิสุทฺธิโย, เตน ภควตา ชานตา ปสฺสตา อรหตา สมฺมา\-สมฺพุทฺเธน สมฺมทกฺขาตา สตฺตานํ วิสุทฺธิยา โสก\-ปริเทวานํ สมติกฺกมาย ทุกฺข\-โทมนสฺสานํ อตฺถงฺคมาย ญายสฺส อธิคมาย นิพฺพานสฺส สจฺฉิกิริยายาติ.}

\prose{เอวํ วุตฺเต ปณฺฑิต\-กุมารโก ลิจฺฉวิ อภยํ ลิจฺฉวิํ เอตทโวจ “กิํ ปน ตฺวํ สมฺม อภย อายสฺมโต อานนฺทสฺส สุภาสิตํ สุภาสิตโต นาพฺภนุโมทสี”ติ. กฺยาหํ สมฺม ปณฺฑิต\-กุมารก อายสฺมโต อานนฺทสฺส สุภาสิตํ สุภาสิตโต นาพฺภนุโมทิสฺสามิ, มุทฺธาปิ ตสฺส วิปเตยฺย, โย อายสฺมโต อานนฺทสฺส สุภาสิตํ สุภาสิตโต นาพฺภนุโมเทยฺยาติ.  จตุตฺถํ.}

\chapterheadpageread{(8) 3. อานนฺทวคฺค \textbf{5. นิเวสกสุตฺต}}
\proseitem{76}{\csromanfolio{223}อถ โข อายสฺมา อานนฺโท เยน ภควา เตนุปสงฺกมิ, อุปสงฺกมิตฺวา ภควนฺตํ อภิวาเทตฺวา เอกมนฺตํ นิสีทิ, เอกมนฺตํ นิสินฺนํ โข อายสฺมนฺตํ อานนฺทํ ภควา เอตทโวจ\csromandash{}}

\prose{เย อานนฺท อนุกมฺเปยฺยาถ, เย จ โสตพฺพํ มญฺเญยฺยุํ มิตฺตา วา อมจฺจา วา ญาตี วา สาโลหิตา วา, เต โว อานนฺท ตีสุ ฐาเนสุ สมาทเปตพฺพา\csromansharedfootnote{b00da63212b9ba77}{สมาทาเปตพฺพา (?)} นิเวเสตพฺพา ปติฏฺฐาเป\-ตพฺพา. กตเมสุ ตีสุ? พุทฺเธ อเวจฺจปฺปสาเท สมาทเปตพฺพา นิเวเสตพฺพา ปติฏฺฐาเป\-ตพฺพา “อิติปิ โส ภควา อรหํ สมฺมาสมฺพุทฺโธ วิชฺชา\-จรณ\-สมฺปนฺโน สุคโต โลกวิทู อนุตฺตโร ปุริส\-ทมฺม\-สารถิ สตฺถา เทวมนุสฺสานํ พุทฺโธ ภควา”ติ. ธมฺเม อเวจฺจปฺปสาเท สมาทเปตพฺพา นิเวเสตพฺพา ปติฏฺฐาเป\-ตพฺพา “สฺวากฺขาโต ภควตา ธมฺโม สนฺทิฏฺฐิโก อกาลิโก เอหิปสฺสิโก โอปเนยฺยิโก ปจฺจตฺตํ เวทิตพฺโพ วิญฺญูหี”ติ. สํเฆ อเวจฺจปฺปสาเท สมาทเปตพฺพา นิเวเสตพฺพา ปติฏฺฐาเป\-ตพฺพา “สุปฺปฏิปนฺโน ภควโต สาวกสํโฆ, อุชุปฺปฏิปนฺโน ภควโต สาวกสํโฆ, ญายปฺปฏิปนฺโน ภควโต สาวกสํโฆ, สามีจิ\-ปฺปฏิปนฺโน ภควโต สาวกสํโฆ, ยทิทํ จตฺตาริ ปุริสยุคานิ อฏฺฐ ปุริสปุคฺคลา, เอส ภควโต สฺวกสํโฆ อาหุเนยฺโย ปาหุเนยฺโย ทกฺขิเณยฺโย อญฺชลิกรณีโย อนุตฺตรํ ปุญฺญกฺเขตฺตํ โลกสฺสา”ติ.}

\prose{สิยา อานนฺท จตุนฺนํ มหาภูตานํ อญฺญถตฺตํ ปถวีธาตุยา อาโปธาตุยา เตโชธาตุยา วาโยธาตุยา, น เตฺวว พุทฺเธ อเวจฺจ\-ปฺปสาเทน สมนฺนาคตสฺส อริยสาวกสฺส สิยา อญฺญถตฺตํ. ตตฺริทํ อญฺญถตฺตํ, โส วตานนฺท พุทฺเธ อเวจฺจ\-ปฺปสาเทน สมนฺนาคโต อริยสาวโก นิรยํ วา ติรจฺฉาน\-โยนิํ วา เปตฺติวิสยํ วา อุปปชฺชิสฺสตีติ เนตํ ฐานํ วิชฺชติ.}

\prose{สิยา อานนฺท จตุนฺนํ มหาภูตานํ อญฺญถตฺตํ ปถวีธาตุยา อาโปธาตุยา เตโชธาตุยา วาโยธาตุยา, น เตฺวว ธมฺเม ฯเปฯ น เตฺวว สํเฆ อเวจฺจ\-ปฺปสาเทน สมนฺนาคตสฺส อริยสาวกสฺส สิยา อญฺญถตฺตํ. ตตฺริทํ อญฺญถตฺตํ, โส วตานนฺท สํเฆ อเวจฺจ\-ปฺปสาเทน \csromanfolio{224}สมนฺนาคโต อริยสาวโก นิรยํ วา ติรจฺฉาน\-โยนิํ วา เปตฺติวิสยํ วา อุปปชฺชิสฺสตีติ เนตํ ฐานํ วิชฺชติ.}

\prose{เย อานนฺท อนุกมฺเปยฺยาถ, เย จ โสตพฺพํ มญฺเญยฺยุํ มิตฺตา วา อมจฺจา วา ญาตี วา สาโลหิตา วา, เต โว อานนฺท อิเมสุ ตีสุ ฐาเนสุ สมาทเปตพฺพา นิเวเสตพฺพา ปติฏฺฐาเป\-ตพฺพาติ.  ปญฺจมํ.}

\csromanheader{2}{6. ปฐม\-ภว\-สุตฺต}
\proseitem{77}{อถ โข อายสฺมา อานนฺโท เยน ภควา เตนุปสงฺกมิ, อุปสงฺกมิตฺวา ภควนฺตํ อภิวาเทตฺวา เอกมนฺตํ นิสีทิ, เอกมนฺตํ นิสินฺโน โข อายสฺมา อานนฺโท ภควนฺตํ เอตทโวจ “ภโว ภโวติ ภนฺเต วุจฺจติ, กิตฺตาวตา นุ โข ภนฺเต ภโว โหตี”ติ.}

\prose{กามธาตุเวปกฺกญฺจ อานนฺท กมฺมํ นาภวิสฺส, อปิ นุ โข กามภโว ปญฺญาเยถาติ. โน เหตํ ภนฺเต. อิติ โข อานนฺท กมฺมํ เขตฺตํ, วิญฺญาณํ พีชํ, ตณฺหา สฺเนโห\csromansharedfootnote{c83e6a18bd2dab01}{สิเนโห (สี, สฺยา, กํ, อิ)}. อวิชฺชา\-นีวรณานํ สตฺตานํ ตณฺหา\-สํโยชนานํ หีนาย ธาตุยา วิญฺญาณํ ปติฏฺฐิตํ, เอวํ อายติํ\csromansharedfootnote{df890a47de56e9f4}{อายติ (สี)} ปุนพฺภวา\-ภินิพฺพตฺติ โหติ. ( )\csromansharedfootnote{e60cfbd578e1065b}{(เอวํ โข อานนฺท ภโว โหตีติ) (ก) ทุติยสุตฺเต ปน อิทํ ปาฐนานตฺตํ นตฺติ.}}

\prose{รูปธาตุเวปกฺกญฺจ อานนฺท กมฺมํ นาภวิสฺส, อปิ นุ โข รูปภโว ปญฺญาเยถาติ. โน เหตํ ภนฺเต. อิติ โข อานนฺท กมฺมํ เขตฺตํ, วิญฺญาณํ พีชํ, ตณฺหา สฺเนโห. อวิชฺชา\-นีวรณานํ สตฺตานํ ตณฺหา\-สํโยชนานํ มชฺฌิมาย ธาตุยา วิญฺญาณํ ปติฏฺฐิตํ, เอวํ อายติํ ปุนพฺภวา\-ภินิพฺพตฺติ โหติ. ( )\csromansharedfootnote{e60cfbd578e1065b}{(เอวํ โข อานนฺท ภโว โหตีติ) (ก) ทุติยสุตฺเต ปน อิทํ ปาฐนานตฺตํ นตฺติ.}}

\prose{อรูปธาตุเวปกฺกญฺจ อานนฺท กมฺมํ นาภวิสฺส, อปิ นุ โข อรูปภโว ปญฺญาเยถาติ. โน เหตํ ภนฺเต. อิติ โข อานนฺท กมฺมํ เขตฺตํ, วิญฺญาณํ พีชํ, ตณฺหา สฺเนโห. อวิชฺชา\-นีวรณานํ สตฺตานํ ตณฺหา\-สํโยชนานํ ปณีตาย ธาตุยา วิญฺญาณํ ปติฏฺฐิตํ, เอวํ อายติํ ปุนพฺภวา\-ภินิพฺพตฺติ โหติ. เอวํ โข อานนฺท ภโว โหตีติ.  ฉฏฺฐํ.}

\chapterheadpageread{(8) 3. อานนฺทวคฺค 7. ทุติย\-ภว\-สุตฺต}
\proseitem{78}{\csromanfolio{225}อถ โข อายสฺมา อานนฺโท เยน ภควา เตนุปสงฺกมิ ฯเปฯ อายสฺมา อานนฺโท ภควนฺตํ เอตทโวจ “ภโว ภโวติ ภนฺเต วุจฺจติ, กิตฺตาวตา นุ โข ภนฺเต ภโว โหตี”ติ.}

\prose{กามธาตุเวปกฺกญฺจ อานนฺท กมฺมํ นาภวิสฺส, อปิ นุ โข กามภโว ปญฺญาเยถาติ. โน เหตํ ภนฺเต. อิติ โข อานนฺท กมฺมํ เขตฺตํ, วิญฺญาณํ พีชํ, ตณฺหา สฺเนโห. อวิชฺชา\-นีวรณานํ สตฺตานํ ตณฺหา\-สํโยชนานํ หีนาย ธาตุยา เจตนา ปติฏฺฐิตา ปตฺถนา ปติฏฺฐิตา, เอวํ อายติํ ปุนพฺภวา\-ภินิพฺพตฺติ โหติ.}

\prose{รูปธาตุเวปกฺกญฺจ อานนฺท กมฺมํ นาภวิสฺส, อปิ นุ โข รูปภโว ปญฺญาเยถาติ. โน เหตํ ภนฺเต. อิติ โข อานนฺท กมฺมํ เขตฺตํ, วิญฺญาณํ พีชํ, ตณฺหา สฺเนโห. อวิชฺชา\-นีวรณานํ สตฺตานํ ตณฺหา\-สํโยชนานํ มชฺฌิมาย ธาตุยา เจตนา ปติฏฺฐิตา ปตฺถนา ปติฏฺฐิตา, เอวํ อายติํ ปุนพฺภวา\-ภินิพฺพตฺติ โหติ.}

\prose{อรูปธาตุเวปกฺกญฺจ อานนฺท กมฺมํ นาภวิสฺส, อปิ นุ โข อรูปภโว ปญฺญาเยถาติ. โน เหตํ ภนฺเต. อิติ โข อานนฺท กมฺมํ เขตฺตํ, วิญฺญาณํ พีชํ, ตณฺหา สฺเนโห. อวิชฺชา\-นีวรณานํ สตฺตานํ ตณฺหา\-สํโยชนานํ ปณีตาย ธาตุยา เจตนา ปติฏฺฐิตา ปตฺถนา ปติฏฺฐิตา, เอวํ อายติํ ปุนพฺภวา\-ภินิพฺพตฺติ โหติ. เอวํ โข อานนฺท ภโว โหตีติ.  }

\csromanheader{2}{8. \textbf{สีลพฺพตสุตฺต}}
\proseitem{79}{อถ โข อายสฺมา อานนฺโท เยน ภควา เตนุปสงฺกมิ, อุปสงฺกมิตฺวา ภควนฺตํ อภิวาเทตฺวา เอกมนฺตํ นิสีทิ, เอกมนฺตํ นิสินฺนํ โข อายสฺมนฺตํ อานนฺทํ ภควา เอตทโวจ “สพฺพํ นุ โข อานนฺท สีลพฺพตํ ชีวิตํ พฺรหฺมจริยํ อุปฏฺฐานสารํ สผลนฺ”ติ. น เขฺวตฺถ ภนฺเต เอกํเสนาติ. เตนหานนฺท วิภชสฺสูติ.}

\prose{ยญฺหิสฺส\csromansharedfootnote{28251c6fddc8f03b}{ยถารูปํ หิสฺส (?) ม 3. 94 ปิฏฺเฐ เสวิตพฺพาเสวิตพฺพสุตฺตานุรูปํ.} ภนฺเต สีลพฺพตํ ชีวิตํ พฺรหฺมจริยํ อุปฏฺฐานสารํ เสวโต อกุสลา ธมฺมา อภิวฑฺฒนฺติ, กุสลา ธมฺมา ปริหายนฺติ. เอวรูปํ สีลพฺพตํ \csromanfolio{226}ชีวิตํ พฺรหฺมจริยํ อุปฏฺฐานสารํ อผลํ. ยญฺจ ขฺวาสฺส\csromansharedfootnote{c1a081dd439c98a1}{ยญฺหิสฺส (ก), ยถารูปญฺจ ขฺวาสฺส (?)} ภนฺเต สีลพฺพตํ ชีวิตํ พฺรหฺมจริยํ อุปฏฺฐานสารํ เสวโต อกุสลา ธมฺมา ปริหายนฺติ, กุสลา ธมฺมา อภิวฑฺฒนฺติ. เอวรูปํ สีลพฺพตํ ชีวิตํ พฺรหฺมจริยํ อุปฏฺฐานสารํ สผลนฺติ. อิทมโวจ อายสฺมา อานนฺโท, สมนุญฺโญ สตฺถา อโหสิ.}

\prose{อถ โข อายสฺมา อานนฺโท “สมนุญฺโญ เม สตฺถา”ติ อุฏฺฐายาสนา ภควนฺตํ อภิวาเทตฺวา ปทกฺขิณํ กตฺวา ปกฺกามิ. อถ โข ภควา อจิรปกฺกนฺเต อายสฺมนฺเต อานนฺเท ภิกฺขู อามนฺเตสิ “เสโข ภิกฺขเว อานนฺโท, น จ ปนสฺส สุลภรูโป สมสโม ปญฺญายา”ติ.  อฏฺฐมํ.}

\csromanheader{2}{9. คนฺธชาต\-สุตฺต}
\proseitem{80}{อถ โข อายสฺมา อานนฺโท เยน ภควา เตนุปสงฺกมิ, อุปสงฺกมิตฺวา ภควนฺตํ อภิวาเทตฺวา เอกมนฺตํ นิสีทิ, เอกมนฺตํ นิสินฺโน โข อายสฺมา อานนฺโท ภควนฺตํ เอตทโวจ\csromandash{}}

\prose{ตีณิมานิ ภนฺเต คนฺธชาตานิ, เยสํ อนุวาตํเยว คนฺโธ คจฺฉติ โน ปฏิวาตํ. กตมานิ ตีณิ? มูลคนฺโธ สารคนฺโธ ปุปฺผคนฺโธ, อิมานิ โข ภนฺเต ตีณิ คนฺธชาตานิ, เยสํ อนุวาตํเยว คนฺโธ คจฺฉติ โน ปฏิวาตํ. อตฺถิ นุ โข ภนฺเต กิญฺจิ คนฺธชาตํ, ยสฺส อนุวาตมฺปิ คนฺโธ คจฺฉติ, ปฏิวาตมฺปิ คนฺโธ คจฺฉติ, อนุวาต\-ปฏิวาตมฺปิ คนฺโธ คจฺฉตีติ.}

\prose{อตฺถานนฺท กิญฺจิ คนฺธชาตํ\csromansharedfootnote{7f2225f9d5bb59ca}{อตฺถานนฺท คนฺธชาตํ (สี, สฺยา, กํ, อิ)}, ยสฺส อนุวาตมฺปิ คนฺโธ คจฺฉติ, ปฏิวาตมฺปิ คนฺโธ คจฺฉติ, อนุวาต\-ปฏิวาตมฺปิ คนฺโธ คจฺฉตีติ. กตมญฺจ ปน ภนฺเต คนฺธชาตํ, ยสฺส อนุวาตมฺปิ คนฺโธ คจฺฉติ, ปฏิวาตมฺปิ คนฺโธ คจฺฉติ อนุวาต\-ปฏิวาตมฺปิ คนฺโธ คจฺฉตีติ.}

\prose{อิธานนฺท ยสฺมิํ คาเม วา นิคเม วา อิตฺถี วา ปุริโส วา พุทฺธํ สรณํ คโต โหติ, ธมฺมํ สรณํ คโต โหติ, สํฆํ สรณํ คโต โหติ, ปาณาติปาตา ปฏิวิรโต โหติ, อทินฺนาทานา ปฏิวิรโต โหติ, กาเมสุ\-มิจฺฉาจารา ปฏิวิรโต โหติ, มุสาวาทา ปฏิวิรโต โหติ, สุรา\-เมรย\-มชฺช\-ปมาทฏฺฐานา ปฏิวิรโต}

\vspace{\tipitakaparskip}
\csromancenter{(8) 3. อานนฺทวคฺค โหติ, สีลวา โหติ กลฺยาณธมฺโม, วิคต\-มล\-มจฺเฉเรน เจตสา อคารํ อชฺฌาวสติ, มุตฺตจาโค ปยตปาณิ โวสคฺครโต ยาจโยโค ทาน\-สํวิภาค\-รโต.}
\prose{\csromanfolio{227}ตสฺส ทิสาสุ สมณพฺราหฺมณา วณฺณํ ภาสนฺติ “อมุกสฺมิํ\csromansharedfootnote{fab7e2ec3552ceb6}{อสุกสฺมิํ (สี, สฺยา, กํ, อิ)} นาม คาเม วา นิคเม วา อิตฺถี วา ปุริโส วา พุทฺธํ สรณํ คโต โหติ, ธมฺมํ สรณํ คโต โหติ, สํฆํ สรณํ คโต โหติ, ปาณาติปาตา ปฏิวิรโต โหติ, อทินฺนาทานา ปฏิวิรโต โหติ, กาเมสุ\-มิจฺฉาจารา ปฏิวิรโต โหติ, มุสาวาทา ปฏิวิรโต โหติ, สุรา\-เมรย\-มชฺช\-ปมาทฏฺฐานา ปฏิวิรโต โหติ, สีลวา โหติ กลฺยาณธมฺโม, วิคต\-มล\-มจฺเฉเรน เจตสา อคารํ อชฺฌาวสติ, มุตฺตจาโค ปยตปาณิ โวสคฺครโต ยาจโยโค ทาน\-สํวิภาค\-รโต”ติ.}

\prose{เทวตาปิสฺส\csromansharedfootnote{ad179d49d8140d9a}{เทวตาปิสฺส อมนุสฺสา (สี, อิ), เทวตาปิสฺส อมนุสฺสาปิ (ก), เทวตาปิสฺส ฯเปฯ มนุสฺสาปิสฺส (?)} วณฺณํ ภาสนฺติ “อมุกสฺมิํ นาม คาเม วา นิคเม วา อิตฺถี วา ปุริโส วา พุทฺธํ สรณํ คโต โหติ, ธมฺมํ สรณํ คโต โหติ, สํฆํ สรณํ คโต โหติ, ปาณาติปาตา ปฏิวิรโต โหติ ฯเปฯ สุรา\-เมรย\-มชฺช\-ปมาทฏฺฐานา ปฏิวิรโต โหติ, สีลวา โหติ กลฺยาณธมฺโม, วิคต\-มล\-มจฺเฉเรน เจตสา อคารํ อชฺฌาวสติ, มุตฺตจาโค ปยตปาณิ โวสคฺครโต ยาจโยโค ทาน\-สํวิภาค\-รโต”ติ. อิทํ โข ตํ อานนฺท คนฺธชาตํ, ยสฺส อนุวาตมฺปิ คนฺโธ คจฺฉติ, ปฏิวาตมฺปิ คนฺโธ คจฺฉติ, อนุวาต\-ปฏิวาตมฺปิ คนฺโธ คจฺฉตีติ.}

\setlength{\csromangathapagewidth}{0pt}
\setlength{\csromangathawakwidth}{0pt}
\csromangathameasuregroup{}{น ปุปฺผคนฺโธ ปฏิวาตเมติ, \\ น จนฺทนํ ตครมลฺลิกา\textsuperscript{1} วา. \\ สตญฺจ คนฺโธ ปฏิวาตเมติ, \\ สพฺพา ทิสา สปฺปุริโส ปวายตีติ.นวมํ.}
\csromangathameasurewak{น ปุปฺผคนฺโธ ปฏิวาตเมติ, \\ น จนฺทนํ ตครมลฺลิกา\textsuperscript{1} วา. \\ สตญฺจ คนฺโธ ปฏิวาตเมติ, \\ สพฺพา ทิสา สปฺปุริโส ปวายตีติ.นวมํ.}
\setlength{\csromangathaleftcol}{0pt}
\csromangathagroup{\csromangathastanza{\csromangathawak{น ปุปฺผคนฺโธ ปฏิวาตเมติ, \\ น จนฺทนํ ตครมลฺลิกา\csromansharedfootnote{3cfc4fc3c19bf837}{ตคฺครมลฺลิกา (อิ)} วา. \\ สตญฺจ คนฺโธ ปฏิวาตเมติ, \\ สพฺพา ทิสา สปฺปุริโส ปวายตีติ.  นวมํ.}}}

\csromanheader{2}{10. \textbf{จูฬนิกาสุตฺต}}
\proseitem{81}{อถ โข อายสฺมา อานนฺโท เยน ภควา เตนุปสงฺกมิ, อุปสงฺกมิตฺวา ภควนฺตํ อภิวาเทตฺวา เอกมนฺตํ นิสีทิ, เอกมนฺตํ นิสินฺโน โข อายสฺมา อานนฺโท ภควนฺตํ เอตทโวจ “สมฺมุขา เมตํ ภนฺเต \csromanfolio{228}ภควโต สุตํ, สมฺมุขา ปฏิคฺคหิตํ ‘ภควโต อานนฺท สิขิสฺส อภิภู นาม สาวโก พฺรหฺมโลเก ฐิโต สหสฺสิ\-โลกธาตุํ\csromansharedfootnote{b43c315b6d800858}{สหสฺสีโลกธาตุํ (อิ) สํ 1. 157 ปิฏฺเฐ วิตฺถาโร.} สเรน วิญฺญาเปสี’ติ, ภควา ปน ภนฺเต อรหํ สมฺมาสมฺพุทฺโธ กีวตกํ ปโหติ สเรน วิญฺญาเปตุนฺ”ติ. สาวโก โส อานนฺท, อปฺปเมยฺยา ตถาคตาติ.}

\prose{ทุติยมฺปิ โข อายสฺมา อานนฺโท ภควนฺตํ เอตทโวจ “สมฺมุขา เมตํ ภนฺเต ภควโต สุตํ, สมฺมุขา ปฏิคฺคหิตํ ‘ภควโต อานนฺท สิขิสฺส อภิภู นาม สาวโก พฺรหฺมโลเก ฐิโต สหสฺสิ\-โลกธาตุํ สเรน วิญฺญาเปสี’ติ, ภควา ปน ภนฺเต อรหํ สมฺมาสมฺพุทฺโธ กีวตกํ ปโหติ สเรน วิญฺญาเปตุนฺ”ติ. สาวโก โส อานนฺท, อปฺปเมยฺยา ตถาคตาติ.}

\prose{ตติยมฺปิ โข อายสฺมา อานนฺโท ภควนฺตํ เอตทโวจ “สมฺมุขา เมตํ ภนฺเต ภควโต สุตํ, สมฺมุขา ปฏิคฺคหิตํ ‘ภควโต อานนฺท สิขิสฺส อภิภู นาม สาวโก พฺรหฺมโลเก ฐิโต สหสฺสิ\-โลกธาตุํ สเรน วิญฺญาเปสี’ติ, ภควา ปน ภนฺเต อรหํ สมฺมาสมฺพุทฺโธ กีวตกํ ปโหติ สเรน วิญฺญาเปตุนฺ”ติ. สุตา เต อานนฺท สหสฺสี จูฬนิกา โลกธาตูติ. เอตสฺส ภควา กาโล เอตสฺส สุคต กาโล, ยํ ภควา ภาเสยฺย, ภควโต สุตฺวา ภิกฺขู ธาเรสฺสนฺตีติ. เตนหานนฺท สุณาหิ สาธุกํ มนสิ กโรหิ, ภาสิสฺสามีติ. “เอวํ ภนฺเต”ติ โข อายสฺมา อานนฺโท ภควโต ปจฺจสฺโสสิ. ภควา เอตทโวจ\csromandash{}}

\prose{ยาวตา อานนฺท จนฺทิมสูริยา\csromansharedfootnote{444420b70b60ca28}{จนฺทิมสุริยา (สี, สฺยา, กํ, อิ)} ปริหรนฺติ ทิสา ภนฺติ วิโรจนา, ตาว สหสฺสธา โลโก, ตสฺมิํ สหสฺสธา โลเก สหสฺสํ\csromansharedfootnote{e5401418499089ea}{ตสฺมิํ สหสฺสํ (สฺยา, กํ, อิ)} จนฺทานํ, สหสฺสํ สูริยานํ, สหสฺสํ สิเนรุ\-ปพฺพตราชานํ, สหสฺสํ ชมฺพุทีปานํ, สหสฺสํ อปรโคยานานํ, สหสฺสํ อุตฺตรกุรูนํ, สหสฺสํ ปุพฺพวิเทหานํ, จตฺตาริ มหาสมุทฺท\-สหสฺสานิ, จตฺตาริ มหาราช\-สหสฺสานิ, สหสฺสํ จาตุมหาราชิกานํ, สหสฺสํ ตาวติํสานํ, สหสฺสํ ยามานํ, สหสฺสํ ตุสิตานํ, สหสฺสํ นิมฺมานรตีนํ, สหสฺสํ ปรนิมฺมิต\-วส\-วตฺตีนํ, สหสฺสํ พฺรหฺมโลกานํ. อยํ วุจฺจตานนฺท สหสฺสี จูฬนิกา โลกธาตุ.}

\chapterheadpageread{(8) 3. อานนฺทวคฺค}
\prose{\csromanfolio{229}ยาวตานนฺท สหสฺสี จูฬนิกา โลกธาตุ, ตาว สหสฺสธา โลโก. อยํ วุจฺจตานนฺท ทฺวิสหสฺสี มชฺฌิมิกา โลกธาตุ.}

\prose{ยาวตานนฺท ทฺวิสหสฺสี มชฺฌิมิกา โลกธาตุ, ตาว สหสฺสธา โลโก. อยํ วุจฺจตานนฺท ติสหสฺสี มหาสหสฺสี โลกธาตุ.}

\prose{อากงฺขมาโน อานนฺท ตถาคโต ติสหสฺสิ\-มหาสหสฺสิ\-โลกธาตุํ\csromansharedfootnote{60682c4233c0584a}{ติสหสฺสิํ มหาสหสฺสิํ โลกธาตุํ (สฺยา, กํ), ติสหสฺสีมหาสหสฺสีโลกธาตุํ (อิ)} สเรน วิญฺญาเปยฺย, ยาวตา ปน อากงฺเขยฺยาติ.}

\prose{ยถา กถํ ปน ภนฺเต ภควา ติสหสฺสิ\-มหาสหสฺสิ\-โลกธาตุํ สเรน วิญฺญาเปยฺย, ยาวตา ปน อากงฺเขยฺยาติ. อิธานนฺท ตถาคโต ติสหสฺสิ\-มหาสหสฺสิ\-โลกธาตุํ โอภาเสน ผเรยฺย, ยทา เต สตฺตา ตํ อาโลกํ สญฺชาเนยฺยุํ. อถ ตถาคโต โฆสํ กเรยฺย สทฺทมนุสฺสาเวยฺย. เอวํ โข อานนฺท ตถาคโต ติสหสฺสิ\-มหาสหสฺสิ\-โลกธาตุํ สเรน วิญฺญาเปยฺย, ยาวตา ปน อากงฺเขยฺยาติ.}

\prosewithclosermidruled{เอวํ วุตฺเต อายสฺมา อานนฺโท (อายสฺมนฺตํ อุทายิํ)\csromansharedfootnote{afe5c1d18f329173}{(ภควนฺตํ) (สี), ( ) นตฺถิ สฺยากํโปตฺถเกสุ. อฏฺฐกถาย สเมติ.} เอตทโวจ “ลาภา วต เม, สุลทฺธํ วต เม, ยสฺส เม สตฺถา เอวํมหิทฺธิโก เอวํ\-มหา\-นุภาโว”ติ. เอวํ วุตฺเต อายสฺมา อุทายี อายสฺมนฺตํ อานนฺทํ เอตทโวจ “กิํ ตุเยฺหตฺถ อาวุโส อานนฺท ยทิ เต สตฺถา เอวํมหิทฺธิโก เอวํ\-มหา\-นุภาโว”ติ. เอวํ วุตฺเต ภาควา อายสฺมนฺตํ อุทายิํ เอตทโวจ “มา เหวํ อุทายิ, มา เหวํ อุทายิ, สเจ อุทายิ อานนฺโท อวีตราโค กาลํ กเรยฺย, เตน จิตฺตปฺปสาเทน สตฺตกฺขตฺตุํ เทเวสุ เทวรชฺชํ กาเรยฺย, สตฺตกฺขตฺตุํ อิมสฺมิํ เยว ชมฺพุทีเป มหารชฺชํ กาเรยฺย, อปิ จ อุทายิ อานนฺโท ทิฏฺเฐว ธมฺเม ปรินิพฺพายิสฺสตี”ติ.  ทสมํ.}{อานนฺทวคฺโค ตติโย.}
\csromancenter{\textbf{ตสฺสุทฺทานํ}}
\setlength{\csromangathapagewidth}{0pt}
\setlength{\csromangathawakwidth}{0pt}
\csromangathameasuregroup{ฉนฺโน อาชีวโก สกฺโก, \\ ทุเว ภวา สีลพฺพตํ,}{\csromangathabat{ฉนฺโน อาชีวโก สกฺโก,}{นิคณฺโฐ จ นิเวสโก.} \\ \csromangathabat{ทุเว ภวา สีลพฺพตํ,}{คนฺธชาตญฺจ จูฬนีติ.}}
\csromangathasetleft{ฉนฺโน อาชีวโก สกฺโก, \\ ทุเว ภวา สีลพฺพตํ,}
\csromangathagroup{\csromangathastanza{\csromangathabat{ฉนฺโน อาชีวโก สกฺโก,}{นิคณฺโฐ จ นิเวสโก.} \\ \csromangathabat{ทุเว ภวา สีลพฺพตํ,}{คนฺธชาตญฺจ จูฬนีติ.}}}

\csromanpalirecto
\gambhira{ติกนิปาตปาฬิ \textbf{(9) 4. สมณวคฺค}}
\csromanheader{2}{1. \textbf{สมณสุตฺต}}
\proseitem{82}{\csromanfolio{230}ตีณิมานิ ภิกฺขเว สมณสฺส สมณิยานิ สมณ\-กรณียานิ. กตมานิ ตีณิ? อธิสีล\-สิกฺขา\-สมาทานํ อธิจิตฺต\-สิกฺขา\-สมาทานํ อธิปญฺญา\-สิกฺขา\-สมาทานํ. อิมานิ โข ภิกฺขเว ตีณิ สมณสฺส สมณิยานิ สมณ\-กรณียานิ.}

\prose{ตสฺมาติห ภิกฺขเว เอวํ สิกฺขิตพฺพํ “ติพฺโพ โน ฉนฺโท ภวิสฺสติ อธิสีล\-สิกฺขาสมาทาเน, ติพฺโพ โน ฉนฺโท ภวิสฺสติ อธิจิตฺต\-สิกฺขาสมาทาเน, ติพฺโพ โน ฉนฺโท ภวิสฺสติ อธิปญฺญา\-สิกฺขาสมาทาเน”ติ. เอวํ หิ โว ภิกฺขเว สิกฺขิตพฺพนฺติ.  ปฐมํ. 2. คทฺรภสุตฺต}

\proseitem{83}{เสยฺยถาปิ ภิกฺขเว คทฺรโภ โคคณํ ปิฏฺฐิโต ปิฏฺฐิโต อนุพนฺโธ โหติ “อหมฺปิ ทมฺโม อหมฺปิ ทมฺโม”ติ\csromansharedfootnote{38b643afdf0fcd7b}{อหมฺปิ โค อมฺหา อหมฺปิ โค อมฺหาติ (สี), อหมฺปิ อมฺหา อหมฺปิ อมฺหาติ (สฺยา, กํ, อิ), อหมฺปิ โค อหมฺปิ โคติ (?)}. ตสฺส น ตาทิโส วณฺโณ โหติ เสยฺยถาปิ คุนฺนํ. น ตาทิโส สโร โหติ เสยฺยถาปิ คุนฺนํ. น ตาทิสํ ปทํ โหติ เสยฺยถาปิ คุนฺนํ. โส โคคณํเยว ปิฏฺฐิโต ปิฏฺฐิโต อนุพนฺโธ โหติ “อหมฺปิ ทมฺโม อหมฺปิ ทมฺโม”ติ.}

\prose{เอวเมวํ โข ภิกฺขเว อิเธกจฺโจ ภิกฺขุ ภิกฺขุสํฆํ ปิฏฺฐิโต ปิฏฺฐิโต อนุพนฺโธ โหติ “อหมฺปิ ภิกฺขุ อหมฺปิ ภิกฺขู”ติ. ตสฺส น ตาทิโส ฉนฺโท โหติ อธิสีล\-สิกฺขาสมาทาเน เสยฺยถาปิ อญฺเญสํ ภิกฺขูนํ. น ตาทิโส ฉนฺโท โหติ อธิจิตฺต\-สิกฺขาสมาทาเน เสยฺยถาปิ อญฺเญสํ ภิกฺขูนํ. น ตาทิโส ฉนฺโท โหติ อธิปญฺญา\-สิกฺขาสมาทาเน เสยฺยถาปิ อญฺเญสํ ภิกฺขูนํ. โส ภิกฺขุสํฆํเยว ปิฏฺฐิโต ปิฏฺฐิโต อนุพนฺโธ โหติ “อหมฺปิ ภิกฺขุ อหมฺปิ ภิกฺขู”ติ.}

\chapterheadpageread{(9) 4. สมณวคฺค}
\prose{\csromanfolio{231}ตสฺมาติห ภิกฺขเว เอวํ สิกฺขิตพฺพํ “ติพฺโพ โน ฉนฺโท ภวิสฺสติ อธิสีล\-สิกฺขาสมาทาเน, ติพฺโพ โน ฉนฺโท ภวิสฺสติ อธิจิตฺต\-สิกฺขาสมาทาเน, ติพฺโพ โน ฉนฺโท ภวิสฺสติ อธิปญฺญา\-สิกฺขาสมาทาเน”ติ. เอวํ หิ โว ภิกฺขเว สิกฺขิตพฺพนฺติ.  ทุติยํ.}

\csromanheader{2}{3. \textbf{เขตฺตสุตฺต}}
\proseitem{84}{ตีณิมานิ ภิกฺขเว กสฺสกสฺส คหปติสฺส ปุพฺเพ กรณียานิ. กตมานิ ตีณิ? อิธ ภิกฺขเว กสฺสโก คหปติ ปฏิกจฺเจว\csromansharedfootnote{aa443bc1a45a8a11}{ปฏิคจฺเจว (สี, อิ)} เขตฺตํ สุกฏฺฐํ กโรติ สุมติกตํ\csromansharedfootnote{ea2581a4eb68d14c}{สุมตฺติกตํ (ก), เอตฺถ มติสทฺโท กฏฺฐเขตฺตสฺส สมีกรณสาธเน ทารุภณฺเฑ วตฺตตีติ สกฺกตอภิธาเนสุ อาคตํ. ตํ “มติยา สุฏฺฐุ สมีกตนฺ”ติ อฏฺฐกถาย สเมติ.}, ปฏิกจฺเจว เขตฺตํ สุกฏฺฐํ กริตฺวา สุมติกตํ กาเลน พีชานิ ปติฏฺฐาเปติ, กาเลน พีชานิ ปติฏฺฐาเปตฺวา สมเยน อุทกํ อภิเนติปิ อปเนติปิ. อิมานิ โข ภิกฺขเว ตีณิ กสฺสกสฺส คหปติสฺส ปุพฺเพ กรณียานิ.}

\prose{เอวเมวํ โข ภิกฺขเว ตีณิมานิ ภิกฺขุสฺส ปุพฺเพ กรณียานิ. กตมานิ ตีณิ? อธิสีล\-สิกฺขา\-สมาทานํ อธิจิตฺต\-สิกฺขา\-สมาทานํ อธิปญฺญา\-สิกฺขา\-สมาทานํ. อิมานิ โข ภิกฺขเว ตีณิ ภิกฺขุสฺส ปุพฺเพ กรณียานิ.}

\prose{ตสฺมาติห ภิกฺขเว เอวํ สิกฺขิตพฺพํ “ติพฺโพ โน ฉนฺโท ภวิสฺสติ อธิสีล\-สิกฺขาสมาทาเน, ติพฺโพ โน ฉนฺโท ภวิสฺสติ อธิจิตฺต\-สิกฺขาสมาทาเน, ติพฺโพ โน ฉนฺโท ภวิสฺสติ อธิปญฺญา\-สิกฺขาสมาทาเน”ติ. เอวํ หิ โว ภิกฺขเว สิกฺขิตพฺพนฺติ.  ตติยํ.}

\csromanheader{2}{4. วชฺชิปุตฺต\-สุตฺต}
\proseitem{85}{เอกํ สมยํ ภควา เวสาลิยํ วิหรติ มหาวเน กูฏาคาร\-สาลายํ. อถ โข อญฺญตโร วชฺชิปุตฺตโก ภิกฺขุ เยน ภควา เตนุปสงฺกมิ, อุปสงฺกมิตฺวา ภควนฺตํ อภิวาเทตฺวา เอกมนฺตํ นิสีทิ, เอกมนฺตํ นิสินฺโน โข โส วชฺชิปุตฺตโก ภิกฺขุ ภควนฺตํ เอตทโวจ\csromandash{} “สาธิกมิทํ ภนฺเต ทิยฑฺฒ\-สิกฺขาปทสตํ\csromansharedfootnote{56b8684112704f03}{ทิยฑฺฒํ สิกฺขาปทสตํ (สี)} อนฺวทฺธมาสํ อุทฺเทสํ อาคจฺฉติ. \csromanfolio{232}นาหํ ภนฺเต เอตฺถ สกฺโกมิ สิกฺขิตุนฺ”ติ. สกฺขิสฺสสิ ปน ตฺวํ ภิกฺขุ ตีสุ สิกฺขาสุ สิกฺขิตุํ อธิสีล\-สิกฺขาย อธิจิตฺต\-สิกฺขาย อธิปญฺญา\-สิกฺขายาติ. สกฺโกมหํ ภนฺเต ตีสุ สิกฺขาสุ สิกฺขิตุํ อธิสีล\-สิกฺขาย อธิจิตฺต\-สิกฺขาย อธิปญฺญา\-สิกฺขายาติ. ตสฺมาติห ตฺวํ ภิกฺขุ ตีสุ สิกฺขาสุ สิกฺขสฺสุ อธิสีล\-สิกฺขาย อธิจิตฺต\-สิกฺขาย อธิปญฺญา\-สิกฺขาย.}

\prose{ยโต โข ตฺวํ ภิกฺขุ อธิสีลมฺปิ สิกฺขิสฺสสิ, อธิจิตฺตมฺปิ สิกฺขิสฺสสิ, อธิปญฺญมฺปิ สิกฺขิสฺสสิ. ตสฺส ตุยฺหํ ภิกฺขุ อธิสีลมฺปิ สิกฺขโต อธิจิตฺตมฺปิ สิกฺขโต อธิปญฺญมฺปิ สิกฺขโต ราโค ปหียิสฺสติ, โทโส ปหียิสฺสติ, โมโห ปหียิสฺสติ. โส ตฺวํ ราคสฺส ปหานา โทสสฺส ปหานา โมหสฺส ปหานา ยํ อกุสลํ, น ตํ กริสฺสสิ. ยํ ปาปํ, น ตํ เสวิสฺสสีติ.}

\prose{อถ โข โส ภิกฺขุ อปเรน สมเยน อธิสีลมฺปิ สิกฺขิ, อธิจิตฺตมฺปิ สิกฺขิ, อธิปญฺญมฺปิ สิกฺขิ. ตสฺส อธิสีลมฺปิ สิกฺขโต อธิจิตฺตมฺปิ สิกฺขโต อธิปญฺญมฺปิ สิกฺขโต ราโค ปหียิ, โทโส ปหียิ, โมโห ปหียิ. โส ราคสฺส ปหานา โทสสฺส ปหานา โมหสฺส ปหานา ยํ อกุสลํ, ตํ นากาสิ. ยํ ปาปํ, ตํ น เสวีติ.  จตุตฺถํ.}

\csromanheader{2}{5. \textbf{เสกฺขสุตฺต}}
\proseitem{86}{อถ โข อญฺญตโร ภิกฺขุ เยน ภควา เตนุปสงฺกมิ, อุปสงฺกมิตฺวา ภควนฺตํ อภิวาเทตฺวา เอกมนฺตํ นิสีทิ, เอกมนฺตํ นิสินฺโน โข โส ภิกฺขุ ภควนฺตํ เอตทโวจ\csromandash{}}

\prose{“เสโข เสโขติ ภนฺเต วุจฺจติ. กิตฺตาวตา นุ โข ภนฺเต เสโข โหตี”ติ. สิกฺขตีติ โข ภิกฺขุ ตสฺมา “เสโข”ติ วุจฺจติ. กิญฺจ สิกฺขติ, อธิสีลมฺปิ สิกฺขติ, อธิจิตฺตมฺปิ สิกฺขติ, อธิปญฺญมฺปิ สิกฺขติ, สิกฺขตีติ โข ภิกฺขุ ตสฺมา “เสโข”ติ วุจฺจตีติ.}

\setlength{\csromangathapagewidth}{0pt}
\setlength{\csromangathawakwidth}{0pt}
\csromangathameasuregroup{เสขสฺส สิกฺขมานสฺส, \\ ขยสฺมิํ ปฐมํ ญาณํ, \\ ตโต อญฺญาวิมุตฺตสฺส\textsuperscript{1},}{\csromangathabat{เสขสฺส สิกฺขมานสฺส,}{อุชุมคฺคานุสาริโน.} \\ \csromangathabat{ขยสฺมิํ ปฐมํ ญาณํ,}{ตโต อญฺญา อนนฺตรา.} \\ \csromangathabat{ตโต อญฺญาวิมุตฺตสฺส\textsuperscript{1},}{ญาณํ เว\textsuperscript{1} โหติ ตาทิโน.}}
\csromangathasetleft{เสขสฺส สิกฺขมานสฺส, \\ ขยสฺมิํ ปฐมํ ญาณํ, \\ ตโต อญฺญาวิมุตฺตสฺส\textsuperscript{1},}
\csromangathagroup{\csromangathastanza{\csromangathabat{เสขสฺส สิกฺขมานสฺส,}{อุชุมคฺคานุสาริโน.} \\ \csromangathabat{ขยสฺมิํ ปฐมํ ญาณํ,}{ตโต อญฺญา อนนฺตรา.} \\ \csromangathabat{ตโต อญฺญาวิมุตฺตสฺส\csromansharedfootnote{d2ac3e8b83a800c0}{อญฺญาวิมุตฺติยา (ก)},}{ญาณํ เว\csromansharedfootnote{cece35fcde56a9c2}{ญาณญฺจ (ก)} โหติ ตาทิโน.}}}

\prose{อกุปฺปา เม วิมุตฺตีติ, ภวสํโยชน\-กฺขเยติ.  ปญฺจมํ. ( )\csromansharedfootnote{1b39c4457cd2c0d2}{(อฏฺฐมํ ภาณวารํ นิฏฺฐิตํ) (ก)}}

\chapterheadpageread{(9) 4. สมณวคฺค 6. ปฐม\-สิกฺขา\-สุตฺต}
\proseitem{87}{\csromanfolio{233}สาธิกมิทํ ภิกฺขเว ทิยฑฺฒ\-สิกฺขาปทสตํ อนฺวทฺธมาสํ อุทฺเทสํ อาคจฺฉติ, ยตฺถ อตฺตกามา กุลปุตฺตา สิกฺขนฺติ. ติสฺโส อิมา ภิกฺขเว สิกฺขา, ยตฺเถตํ สพฺพํ สโมธานํ คจฺฉติ. กตมา ติสฺโส? อธิสีลสิกฺขา อธิจิตฺต\-สิกฺขา อธิปญฺญา\-สิกฺขา. อิมา โข ภิกฺขเว ติสฺโส สิกฺขา, ยตฺเถตํ สพฺพํ สโมธานํ คจฺฉติ.}

\prose{อิธ ภิกฺขเว ภิกฺขุ สีเลสุ ปริปูรการี โหติ, สมาธิสฺมิํ มตฺตโส การี, ปญฺญาย มตฺตโส การี. โส ยานิ ตานิ ขุทฺทา\-นุขุทฺท\-กานิ สิกฺขาปทานิ, ตานิ อาปชฺชติปิ วุฏฺฐาติปิ. ตํ กิสฺส เหตุ? น หิ เมตฺถ ภิกฺขเว อภพฺพตา วุตฺตา. ยานิ จ โข ตานิ สิกฺขาปทานิ อาทิ\-พฺรหฺม\-จริย\-กานิ พฺรหฺม\-จริย\-สารุปฺปานิ, ตตฺถ ธุวสีโล\csromansharedfootnote{38082ac1ce717eb9}{ธุวสีลี (สี) อภิ 3. 142 ปิฏฺเฐปิ (โถกํ วิสทิสํ)} จ โหติ ฐิตสีโล\csromansharedfootnote{dadafecd4ae6ca8b}{ฐิตสีลี (สี)} จ, สมาทาย สิกฺขติ สิกฺขาปเทสุ. โส ติณฺณํ สํโยชนานํ ปริกฺขยา โสตาปนฺโน โหติ อวินิปาต\-ธมฺโม นิยโต สมฺโพธิ\-ปรายโณ.}

\prose{อิธ ปน ภิกฺขเว ภิกฺขุ สีเลสุ ปริปูรการี โหติ, สมาธิสฺมิํ มตฺตโส การี, ปญฺญาย มตฺตโส การี. โส ยานิ ตานิ ขุทฺทา\-นุขุทฺท\-กานิ สิกฺขาปทานิ, ตานิ อาปชฺชติปิ วุฏฺฐาติปิ. ตํ กิสฺส เหตุ? น หิ เมตฺถ ภิกฺขเว อภพฺพตา วุตฺตา. ยานิ จ โข ตานิ สิกฺขาปทานิ อาทิ\-พฺรหฺม\-จริย\-กานิ พฺรหฺม\-จริย\-สารุปฺปานิ, ตตฺถ ธุวสีโล จ โหติ ฐิตสีโล จ, สมาทาย สิกฺขติ สิกฺขาปเทสุ. โส ติณฺณํ สํโยชนานํ ปริกฺขยา ราค\-โทส\-โมหานํ ตนุตฺตา สกทาคามี โหติ, สกิเทว อิมํ โลกํ อาคนฺตฺวา ทุกฺขสฺสนฺตํ กโรติ.}

\prose{อิธ ปน ภิกฺขเว ภิกฺขุ สีเลสุ ปริปูรการี โหติ, สมาธิสฺมิํ ปริปูรการี, ปญฺญาย มตฺตโส การี. โส ยานิ ตานิ ขุทฺทา\-นุขุทฺท\-กานิ สิกฺขาปทานิ, ตานิ อาปชฺชติปิ วุฏฺฐาติปิ. ตํ กิสฺส เหตุ? น หิ เมตฺถ ภิกฺขเว อภพฺพตา วุตฺตา. ยานิ จ โข ตานิ สิกฺขาปทานิ อาทิ\-พฺรหฺม\-จริย\-กานิ พฺรหฺม\-จริย\-สารุปฺปานิ, ตตฺถ ธุวสีโล จ โหติ ฐิตสีโล จ, สมาทาย สิกฺขติ สิกฺขาปเทสุ. โส ปญฺจนฺนํ \csromanfolio{234}โอรมฺภาคิยานํ สํโยชนานํ ปริกฺขยา โอปปาติโก โหติ, ตตฺถ ปรินิพฺพายี อนาวตฺติธมฺโม ตสฺมา โลกา.}

\prose{อิธ ปน ภิกฺขเว ภิกฺขุ สีเลสุ ปริปูรการี โหติ, สมาธิสฺมิํ ปริปูรการี, ปญฺญาย ปริปูรการี. โส ยานิ ตานิ ขุทฺทา\-นุขุทฺท\-กานิ สิกฺขาปทานิ, ตานิ อาปชฺชติปิ วุฏฺฐาติปิ. ตํ กิสฺส เหตุ? น หิ เมตฺถ ภิกฺขเว อภพฺพตา วุตฺตา. ยานิ จ โข ตานิ สิกฺขาปทานิ อาทิ\-พฺรหฺม\-จริย\-กานิ พฺรหฺม\-จริย\-สารุปฺปานิ, ตตฺถ ธุวสีโล จ โหติ ฐิตสีโล จ, สมาทาย สิกฺขติ สิกฺขาปเทสุ. โส อาสวานํ ขยา อนาสวํ เจโตวิมุตฺติํ ปญฺญาวิมุตฺติํ ทิฏฺเฐว ธมฺเม สยํ อภิญฺญา สจฺฉิกตฺวา อุปสมฺปชฺช วิหรติ.}

\prose{อิติ โข ภิกฺขเว ปเทสํ ปเทสการี อาราเธติ, ปริปูรํ ปริปูรการี, อวญฺฌานิเตฺววาหํ\csromansharedfootnote{d5cab7e661dc80f0}{อวญฺฌาเนวาหํ (ก)} ภิกฺขเว สิกฺขาปทานิ วทามีติ.  ฉฏฺฐํ.}

\csromanheader{2}{7. ทุติย\-สิกฺขา\-สุตฺต}
\proseitem{88}{สาธิกมิทํ ภิกฺขเว ทิยฑฺฒ\-สิกฺขาปทสตํ อนฺวทฺธมาสํ อุทฺเทสํ อาคจฺฉติ, ยตฺถ อตฺตกามา กุลปุตฺตา สิกฺขนฺติ. ติสฺโส อิมา ภิกฺขเว สิกฺขา, ยตฺเถตํ สพฺพํ สโมธานํ คจฺฉติ. กตมา ติสฺโส? อธิสีลสิกฺขา อธิจิตฺต\-สิกฺขา อธิปญฺญา\-สิกฺขา. อิมา โข ภิกฺขเว ติสฺโส สิกฺขา, ยตฺเถตํ สพฺพํ สโมธานํ คจฺฉติ.}

\prose{อิธ ภิกฺขเว ภิกฺขุ สีเลสุ ปริปูรการี โหติ, สมาธิสฺมิํ มตฺตโส การี, ปญฺญาย มตฺตโส การี. โส ยานิ ตานิ ขุทฺทา\-นุขุทฺท\-กานิ สิกฺขาปทานิ, ตานิ อาปชฺชติปิ วุฏฺฐาติปิ. ตํ กิสฺส เหตุ? น หิ เมตฺถ ภิกฺขเว อภพฺพตา วุตฺตา. ยานิ จ โข ตานิ สิกฺขาปทานิ อาทิ\-พฺรหฺม\-จริย\-กานิ พฺรหฺม\-จริย\-สารุปฺปานิ, ตตฺถ ธุวสีโล จ โหติ ฐิตสีโล จ, สมาทาย สิกฺขติ สิกฺขาปเทสุ. โส ติณฺณํ สํโยชนานํ ปริกฺขยา สตฺตกฺขตฺตุ\-ปรโม โหติ, สตฺตกฺขตฺตุปรมํ เทเว จ มนุสฺเส จ สนฺธาวิตฺวา สํสริตฺวา ทุกฺขสฺสนฺตํ กโรติ. โส ติณฺณํ สํโยชนานํ ปริกฺขยา โกลํโกโล โหติ, เทฺว วา ตีณิ วา กุลานิ สนฺธาวิตฺวา สํสริตฺวา ทุกฺขสฺสนฺตํ กโรติ. โส}

\vspace{\tipitakaparskip}
\csromancenter{(9) 4. สมณวคฺค ติณฺณํ สํโยชนานํ ปริกฺขยา เอกพีชี โหติ, เอกํเยว มานุสกํ ภวํ นิพฺพตฺเตตฺวา ทุกฺขสฺสนฺตํ กโรติ. โส ติณฺณํ สํโยชนานํ ปริกฺขยา ราค\-โทส\-โมหานํ ตนุตฺตา สกทาคามี โหติ, สกิเทว อิมํ โลกํ อาคนฺตฺวา ทุกฺขสฺสนฺตํ กโรติ.}
\prose{\csromanfolio{235}อิธ ปน ภิกฺขเว ภิกฺขุ สีเลสุ ปริปูรการี โหติ, สมาธิสฺมิํ ปริปูรการี, ปญฺญาย มตฺตโส การี. โส ยานิ ตานิ ขุทฺทา\-นุขุทฺท\-กานิ สิกฺขาปทานิ, ตานิ อาปชฺชติปิ วุฏฺฐาติปิ. ตํ กิสฺส เหตุ? น หิ เมตฺถ ภิกฺขเว อภพฺพตา วุตฺตา. ยานิ จ โข ตานิ สิกฺขาปทานิ อาทิ\-พฺรหฺม\-จริย\-กานิ พฺรหฺม\-จริย\-สารุปฺปานิ, ตตฺถ ธุวสีโล จ โหติ ฐิตสีโล จ, สมาทาย สิกฺขติ สิกฺขาปเทสุ. โส ปญฺจนฺนํ โอรมฺภาคิยานํ สํโยชนานํ ปริกฺขยา อุทฺธํโสโต อกนิฏฺฐคามี. โส ปญฺจนฺนํ โอรมฺภาคิยานํ สํโยชนานํ ปริกฺขยา สสงฺขาร\-ปรินิพฺพายี โหติ. โส ปญฺจนฺนํ โอรมฺภาคิยานํ สํโยชนานํ ปริกฺขยา อสงฺขารปรินิพฺพายี โหติ. โส ปญฺจนฺนํ โอรมฺภาคิยานํ สํโยชนานํ ปริกฺขยา อุปหจฺจ\-ปรินิพฺพายี โหติ. โส ปญฺจนฺนํ โอรมฺภาคิยานํ สํโยชนานํ ปริกฺขยา อนฺตรา\-ปรินิพฺพายี โหติ.}

\prose{อิธ ปน ภิกฺขเว ภิกฺขุ สีเลสุ ปริปูรการี โหติ, สมาธิสฺมิํ ปริปูรการี, ปญฺญาย ปริปูรการี. โส ยานิ ตานิ ขุทฺทา\-นุขุทฺท\-กานิ สิกฺขาปทานิ, ตานิ อาปชฺชติปิ วุฏฺฐาติปิ. ตํ กิสฺส เหตุ? น หิ เมตฺถ ภิกฺขเว อภพฺพตา วุตฺตา. ยานิ จ โข ตานิ สิกฺขาปทานิ อาทิ\-พฺรหฺม\-จริย\-กานิ พฺรหฺม\-จริย\-สารุปฺปานิ, ตตฺถ ธุวสีโล จ โหติ ฐิตสีโล จ, สมาทาย สิกฺขติ สิกฺขาปเทสุ. โส อาสวานํ ขยา อนาสวํ เจโตวิมุตฺติํ ปญฺญาวิมุตฺติํ ทิฏฺเฐว ธมฺเม สยํ อภิญฺญา สจฺฉิกตฺวา อุปสมฺปชฺช วิหรติ.}

\prose{อิธ โข ภิกฺขเว ปเทสํ ปเทสการี อาราเธติ, ปริปูรํ ปริปูรการี, อวญฺฌานิเตฺววาหํ ภิกฺขเว สิกฺขาปทานิ วทามีติ.  }

\csromanheader{2}{8. ตติย\-สิกฺขา\-สุตฺต}
\proseitem{89}{สาธิกมิทํ ภิกฺขเว ทิยฑฺฒ\-สิกฺขาปทสตํ อนฺวทฺธมาสํ อุทฺเทสํ อาคจฺฉติ, ยตฺถ อตฺตกามา กุลปุตฺตา สิกฺขนฺติ. ติสฺโส อิมา ภิกฺขเว \csromanfolio{236}สิกฺขา, ยตฺเถตํ สพฺพํ สโมธานํ คจฺฉติ. กตมา ติสฺโส? อธิสีลสิกฺขา อธิจิตฺต\-สิกฺขา อธิปญฺญา\-สิกฺขา. อิมา โข ภิกฺขเว ติสฺโส สิกฺขา, ยตฺเถตํ สพฺพํ สโมธานํ คจฺฉติ.}

\prose{อิธ ภิกฺขเว ภิกฺขุ สีเลสุ ปริปูรการี โหติ, สมาธิสฺมิํ ปริปูรการี, ปญฺญาย ปริปูรการี. โส ยานิ ตานิ ขุทฺทา\-นุขุทฺท\-กานิ สิกฺขาปทานิ, ตานิ อาปชฺชติปิ วุฏฺฐาติปิ. ตํ กิสฺส เหตุ? น หิ เมตฺถ ภิกฺขเว อภพฺพตา วุตฺตา. ยานิ จ โข ตานิ สิกฺขาปทานิ อาทิ\-พฺรหฺม\-จริย\-กานิ พฺรหฺม\-จริย\-สารุปฺปานิ, ตตฺถ ธุวสีโล จ โหติ ฐิตสีโล จ, สมาทาย สิกฺขติ สิกฺขาปเทสุ. โส อาสวานํ ขยา อนาสวํ เจโตวิมุตฺติํ ปญฺญาวิมุตฺติํ ทิฏฺเฐว ธมฺเม สยํ อภิญฺญา สจฺฉิกตฺวา อุปสมฺปชฺช วิหรติ. ตํ วา ปน อนภิสมฺภวํ อปฺปฏิวิชฺฌํ ปญฺจนฺนํ โอรมฺภาคิยานํ สํโยชนานํ ปริกฺขยา อนฺตรา\-ปรินิพฺพายี โหติ. ตํ วา ปน อนภิสมฺภวํ อปฺปฏิวิชฺฌํ ปญฺจนฺนํ โอรมฺภาคิยานํ สํโยชนานํ ปริกฺขยา อุปหจฺจ\-ปรินิพฺพายี โหติ. ตํ วา ปน อนภิสนฺภวํ อปฺปฏิวิชฺฌํ ปญฺจนฺนํ โอรมฺภาคิยานํ สํโยชนานํ ปริกฺขยา อสงฺขารปรินิพฺพายี โหติ. ตํ วา ปน อนภิสมฺภวํ อปฺปฏิวิชฺฌํ ปญฺจนฺนํ โอรมฺภาคิยานํ สํโยชนานํ ปริกฺขยา สสงฺขาร\-ปรินิพฺพายี โหติ. ตํ วา ปน อนภิสมฺภวํ อปฺปฏิวิชฺฌํ ปญฺจนฺนํ โอรมฺภาคิยานํ สํโยชนานํ ปริกฺขยา อุทฺธํโสโต โหติ อกนิฏฺฐคามี. ตํ วา ปน อนภิสมฺภวํ อปฺปฏิวิชฺฌํ ติณฺณํ สํโยชนานํ ปริกฺขยา ราค\-โทส\-โมหานํ ตนุตฺตา สกทาคามี โหติ, สกิเทว อิมํ โลกํ อาคนฺตฺวา ทุกฺขสฺสนฺตํ กโรติ. ตํ วา ปน อนภิสมฺภวํ อปฺปฏิวิชฺฌํ ติณฺณํ สํโยชนานํ ปริกฺขยา เอกพีชี โหติ, เอกํเยว มานุสกํ ภวํ นิพฺพตฺเตตฺวา ทุกฺขสฺสนฺตํ กโรติ. ตํ วา ปน อนภิสมฺภวํ อปฺปฏิวิชฺฌํ ติณฺณํ สํโยชนานํ ปริกฺขยา โกลํโกโล โหติ, เทฺว วา ตีณิ วา กุลานิ สนฺธาวิตฺวา สํสริตฺวา ทุกฺขสฺสนฺตํ กโรติ. ตํ วา ปน อนภิสมฺภวํ อปฺปฏิวิชฺฌํ ติณฺณํ สํโยชนานํ ปริกฺขยา สตฺตกฺขตฺตุ\-ปรโม โหติ, สตฺตกฺขตฺตุปรมํ เทเว จ มนุสฺเส จ สนฺธาวิตฺวา สํสริตฺวา ทุกฺขสฺสนฺตํ กโรติ.}

\prose{อิติ โข ภิกฺขเว ปริปูรํ ปริปูรการี อาราเธติ, ปเทสํ ปเทสการี, อวญฺฌานิเตฺววาหํ ภิกฺขเว สิกฺขาปทานิ วทามีติ.  อฏฺฐมํ.}

\chapterheadpageread{(9) 4. สมณวคฺค 9. ปฐม\-สิกฺขตฺตย\-สุตฺต}
\proseitem{90}{\csromanfolio{237}ติสฺโส อิมา ภิกฺขเว สิกฺขา. กตมา ติสฺโส? อธิสีลสิกฺขา อธิจิตฺต\-สิกฺขา อธิปญฺญา\-สิกฺขา.}

\prose{กตมา จ ภิกฺขเว อธิสีลสิกฺขา? อิธ ภิกฺขเว ภิกฺขุ สีลวา โหติ ฯเปฯ สมาทาย สิกฺขติ สิกฺขาปเทสุ. อยํ วุจฺจติ ภิกฺขเว อธิสีลสิกฺขา.}

\prose{กตมา จ ภิกฺขเว อธิจิตฺต\-สิกฺขา? อิธ ภิกฺขเว ภิกฺขุ วิวิจฺเจว กาเมหิ ฯเปฯ จตุตฺถํ ฌานํ อุปสมฺปชฺช วิหรติ. อยํ วุจฺจติ ภิกฺขเว อธิจิตฺต\-สิกฺขา.}

\prose{กตมา จ ภิกฺขเว อธิปญฺญา\-สิกฺขา? อิธ ภิกฺขเว ภิกฺขุ “อิทํ ทุกฺขนฺ”ติ ยถาภูตํ ปชานาติ ฯเปฯ “อยํ ทุกฺขนิโรธ\-คามินี ปฏิปทา”ติ ยถาภูตํ ปชานาติ. อยํ วุจฺจติ ภิกฺขเว อธิปญฺญา\-สิกฺขา. อิมา โข ภิกฺขเว ติสฺโส สิกฺขาติ.  นวมํ.}

\csromanheader{2}{10. ทุติย\-สิกฺขตฺตย\-สุตฺต}
\proseitem{91}{ติสฺโส อิมา ภิกฺขเว สิกฺขา. กตมา ติสฺโส? อธิสีลสิกฺขา อธิจิตฺต\-สิกฺขา อธิปญฺญา\-สิกฺขา.}

\prose{กตมา จ ภิกฺขเว อธิสีลสิกฺขา? อิธ ภิกฺขเว ภิกฺขุ สีลวา โหติ ฯเปฯ สมาทาย สิกฺขติ สิกฺขาปเทสุ. อยํ วุจฺจติ ภิกฺขเว อธิสีลสิกฺขา.}

\prose{กตมา จ ภิกฺขเว อธิจิตฺต\-สิกฺขา? อิธ ภิกฺขเว ภิกฺขุ วิวิจฺเจว กาเมหิ ฯเปฯ จตุตฺถํ ฌานํ อุปสมฺปชฺช วิหรติ. อยํ วุจฺจติ ภิกฺขเว อธิจิตฺต\-สิกฺขา.}

\prose{กตมา จ ภิกฺขเว อธิปญฺญา\-สิกฺขา? อิธ ภิกฺขเว ภิกฺขุ อาสวานํ ขยา อนาสวํ เจโตวิมุตฺติํ ปญฺญาวิมุตฺติํ ทิฏฺเฐว ธมฺเม สยํ อภิญฺญา สจฺฉิกตฺวา อุปสมฺปชฺช วิหรติ, อยํ วุจฺจติ ภิกฺขเว อธิปญฺญา\-สิกฺขา, อิมา โข ภิกฺขเว ติสฺโส สิกฺขาติ.}

\setlength{\csromangathapagewidth}{0pt}
\setlength{\csromangathawakwidth}{0pt}
\csromangathameasuregroup{อธิสีลํ อธิจิตฺตํ, \\ ถามวา ธิติมา ฌายี, \\ ยถา ปุเร ตถา ปจฺฉา, \\ ยถา อโธ ตถา อุทฺธํ, \\ ยถา ทิวา ตถา รตฺติํ, \\ อภิภุยฺย ทิสา สพฺพา, \\ ตมาหุ เสขํ ปฏิปทํ\textsuperscript{1}, \\ ตมาหุ โลเก สมฺพุทฺทํ, \\ วิญฺญาณสฺส นิโรเธน, \\ ปชฺโชตสฺเสว นิพฺพานํ,}{\csromangathabat{อธิสีลํ อธิจิตฺตํ,}{อธิปญฺญญฺจ วีริยวา.} \\ \csromangathabat{ถามวา ธิติมา ฌายี,}{สโต คุตฺตินฺทฺริโย\textsuperscript{1} จเร.} \\ \csromangathabat{ยถา ปุเร ตถา ปจฺฉา,}{ยถา ปจฺฉา ตถา ปุเร.} \\ \csromangathabat{ยถา อโธ ตถา อุทฺธํ,}{ยถา อุทฺธํ ตถา อโธ.} \\ \csromangathabat{ยถา ทิวา ตถา รตฺติํ,}{ยถา รตฺติํ ตถา ทิวา.} \\ \csromangathabat{อภิภุยฺย ทิสา สพฺพา,}{อปฺปมาณสมาธินา.} \\ \csromangathabat{ตมาหุ เสขํ ปฏิปทํ\textsuperscript{1},}{อโถ สํสุทฺธจาริยํ\textsuperscript{1}.} \\ \csromangathabat{ตมาหุ โลเก สมฺพุทฺทํ,}{ธีรํ ปฏิปทนฺตคุํ.} \\ \csromangathabat{วิญฺญาณสฺส นิโรเธน,}{ตณฺหากฺขยวิมุตฺติโน.} \\ \csromangathabat{ปชฺโชตสฺเสว นิพฺพานํ,}{วิโมกฺโข โหติ เจตโสติ.} \\ ทสมํ.}
\csromangathameasurewak{ทสมํ.}
\csromangathasetleft{อธิสีลํ อธิจิตฺตํ, \\ ถามวา ธิติมา ฌายี, \\ ยถา ปุเร ตถา ปจฺฉา, \\ ยถา อโธ ตถา อุทฺธํ, \\ ยถา ทิวา ตถา รตฺติํ, \\ อภิภุยฺย ทิสา สพฺพา, \\ ตมาหุ เสขํ ปฏิปทํ\textsuperscript{1}, \\ ตมาหุ โลเก สมฺพุทฺทํ, \\ วิญฺญาณสฺส นิโรเธน, \\ ปชฺโชตสฺเสว นิพฺพานํ,}
\csromangathagroup{\csromangathastanza{\csromanfolio{238}\csromangathabat{อธิสีลํ อธิจิตฺตํ,}{อธิปญฺญญฺจ วีริยวา.} \\ \csromangathabat{ถามวา ธิติมา ฌายี,}{สโต คุตฺตินฺทฺริโย\csromansharedfootnote{24c3853c4929f3ee}{อุปฺปตฺตินฺทฺริโย (ก)} จเร.}}\csromangathastanza{\csromangathabat{ยถา ปุเร ตถา ปจฺฉา,}{ยถา ปจฺฉา ตถา ปุเร.} \\ \csromangathabat{ยถา อโธ ตถา อุทฺธํ,}{ยถา อุทฺธํ ตถา อโธ.}}\csromangathastanza{\csromangathabat{ยถา ทิวา ตถา รตฺติํ,}{ยถา รตฺติํ ตถา ทิวา.} \\ \csromangathabat{อภิภุยฺย ทิสา สพฺพา,}{อปฺปมาณสมาธินา.}}\csromangathastanza{\csromangathabat{ตมาหุ เสขํ ปฏิปทํ\csromansharedfootnote{4411ace901396724}{ปาฏิปทํ (?) ม 2. 19 ปิฏฺเฐ ปสฺสิตพฺพํ.},}{อโถ สํสุทฺธจาริยํ\csromansharedfootnote{0fe4b6df019efc2a}{สํสุทฺธจารณํ (สี, อิ), สํสุทฺธจารินํ (สฺยา, กํ)}.} \\ \csromangathabat{ตมาหุ โลเก สมฺพุทฺทํ,}{ธีรํ ปฏิปทนฺตคุํ.}}\csromangathastanza{\csromangathabat{วิญฺญาณสฺส นิโรเธน,}{ตณฺหากฺขยวิมุตฺติโน.} \\ \csromangathabat{ปชฺโชตสฺเสว นิพฺพานํ,}{วิโมกฺโข โหติ เจตโสติ.}}\csromangathastanza{\csromangathawak{ทสมํ.}}}

\csromanheader{2}{11. \textbf{สงฺกวาสุตฺต}}
\proseitem{92}{เอกํ สมยํ ภควา โกสเลสุ จาริกํ จรมาโน มหตา ภิกฺขุ\-สํเฆน สทฺธิํ เยน สงฺกวา\csromansharedfootnote{f9e0b1ffb6af52ac}{ปงฺกธา (สี, สฺยา, กํ, อิ)} นาม โกสลานํ นิคโม ตทวสริ. ตตฺร สุทํ ภควา สงฺกวายํ วิหรติ. เตน โข ปน สมเยน กสฺสปโคตฺโต นาม ภิกฺขุ สงฺกวายํ อาวาสิโก โหติ. ตตฺร สุทํ ภควา สิกฺขาปท\-ปฏิสํยุตฺตาย ธมฺมิยา กถาย ภิกฺขู สนฺทสฺเสติ สมาทเปติ สมุตฺเตเชติ สมฺปหํเสติ. อถ โข กสฺสป\-โคตฺตสฺส ภิกฺขุโน ภควติ สิกฺขาปท\-ปฏิสํยุตฺตาย ธมฺมิยา กถาย ภิกฺขู สนฺทสฺเสนฺเต สมาทเปนฺเต สมุตฺเตเชนฺเต สมฺปหํเสนฺเต อหุเทว อกฺขนฺติ, อหุ อปฺปจฺจโย “อธิสลฺลิขเต\-วายํ\csromansharedfootnote{c2cea35f1ec85971}{อธิสลฺเลขเตวายํ (สฺยา, กํ, ก)} สมโณ”ติ. อถ โข ภควา สงฺกวายํ ยถาภิรนฺตํ วิหริตฺวา เยน ราชคหํ เตน จาริกํ ปกฺกามิ, อนุปุพฺเพน จาริกํ จรมาโน เยน ราชคหํ ตทวสริ. ตตฺร สุทํ ภควา ราชคเห วิหรติ.}

\prose{อถ โข กสฺสป\-โคตฺตสฺส ภิกฺขุโน อจิร\-ปกฺกนฺตสฺส ภควโต อหุเทว กุกฺกุจฺจํ, อหุ วิปฺปฏิสาโร “อลาภา วต เม, น วต เม}

\vspace{\tipitakaparskip}
\csromancenter{(9) 4. สมณวคฺค ลาภา. ทุลฺลทฺธํ วต เม, น วต เม สุลทฺธํ. ยสฺส เม ภควติ สิกฺขาปท\-ปฏิสํยุตฺตาย ธมฺมิยา กถาย ภิกฺขู สนฺทสฺเสนฺเต สมาทเปนฺเต สมุตฺเตเชนฺเต สมฺปหํเสนฺเต อหุเทว อกฺขนฺติ, อหุ อปฺปจฺจโย ‘อธิสลฺลิขเต\-วายํ สมโณ’ติ. ยํนูนาหํ เยน ภควา เตนุปสงฺกเมยฺยํ, อุปสงฺกมิตฺวา ภควโต สนฺติเก อจฺจยํ อจฺจยโต เทเสยฺยนฺ”ติ. อถ โข กสฺสปโคตฺโต ภิกฺขุ เสนาสนํ สํสาเมตฺวา ปตฺต\-จีวรมา\-ทาย เยน ราชคหํ เตน ปกฺกามิ อนุปุพฺเพน เยน ราชคหํ เยน คิชฺฌกูโฏ ปพฺพโต เยน ภควา เตนุปสงฺกมิ, อุปสงฺกมิตฺวา ภควนฺตํ อภิวาเทตฺวา เอกมนฺตํ นิสีทิ, เอกมนฺตํ นิสินฺโน โข กสฺสปโคตฺโต ภิกฺขุ ภควนฺตํ เอตทโวจ\csromandash{}}
\prose{\csromanfolio{239}เอกมิทํ ภนฺเต สมยํ ภควา สงฺกวายํ วิหรติ สงฺกวา นาม โกสลานํ นิคโม. ตตฺร ภนฺเต ภควา สิกฺขาปท\-ปฏิสํยุตฺตาย ธมฺมิยา กถาย ภิกฺขู สนฺทสฺเสสิ สมาทเปสิ สมุตฺเตเชสิ สมฺปหํเสสิ. ตสฺส มยฺหํ ภควติ สิกฺขาปท\-ปฏิสํยุตฺตาย ธมฺมิยา กถาย ภิกฺขู สนฺทสฺเสนฺเต สมาทเปนฺเต สมุตฺเตเชนฺเต สมฺปหํเสนฺเต อหุเทว อกฺขนฺติ, อหุ อปฺปจฺจโย “อธิสลฺลิขเต\-วายํ สมโณ”ติ. อถ โข ภควา สงฺกวายํ ยถาภิรนฺตํ วิหริตฺวา เยน ราชคหํ เตน จาริกํ ปกฺกามิ, ( )\csromansharedfootnote{e2d9fb1bfa6aba96}{(อนุปุพฺเพน จาริกํ จรมาโน เยน ราชคหํ ตทวสริ. ตตฺร สุทํ ภควา ราชคเห วิหรติ. อถ โข) (ก)} ตสฺส มยฺหํ ภนฺเต อจิร\-ปกฺกนฺตสฺส ภควโต อหุเทว กุกฺกุจฺจํ, อหุ วิปฺปฏิสาโร “อลาภา วต เม, น วต เม ลาภา. ทุลฺลทฺธํ วต เม, น วต เม สุลทฺธํ. ยสฺส เม ภควติ สิกฺขาปท\-ปฏิสํยุตฺตาย ธมฺมิยา กถาย ภิกฺขู สนฺทสฺเสนฺเต สมาทเปนฺเต สมุตฺเตเชนฺเต สมฺปหํเสนฺเต อหุเทว อกฺขนฺติ, อหุ อปฺปจฺจโย ‘อธิสลฺลิขเต\-วายํ สมโณ’ติ. ยํนูนาหํ เยน ภควา เตนุปสงฺกเมยฺยํ, อุปสงฺกมิตฺวา ภควโต สนฺติเก อจฺจยํ อจฺจยโต เทเสยฺยนฺ”ติ. อจฺจโย มํ ภนฺเต อจฺจคมา ยถาพาลํ ยถามูฬฺหํ ยถาอกุสลํ. ยสฺส เม ภควติ สิกฺขาปท\-ปฏิสํยุตฺตาย ธมฺมิยา กถาย ภิกฺขู สนฺทสฺเสนฺเต สมาทเปนฺเต สมุตฺเตเชนฺเต สมฺปหํเสนฺเต อหุเทว อกฺขนฺติ, อหุ อปฺปจฺจโย “อธิสลฺลิขเต\-วายํ สมโณ”ติ. \csromanfolio{240}ตสฺส เม ภนฺเต ภควา อจฺจยํ อจฺจยโต ปฏิคฺคณฺหาตุ อายติํ สํวรายาติ.}

\prose{ตคฺฆ ตํ\csromansharedfootnote{0741458b35bd512d}{ตคฺฆ ตฺวํ (สี, อิ)} กสฺสป อจฺจโย อจฺจคมา ยถาพาลํ ยถามูฬฺหํ ยถาอกุสลํ. ยสฺส เต มยิ สิกฺขาปท\-ปฏิสํยุตฺตาย ธมฺมิยา กถาย ภิกฺขู สนฺทสฺเสนฺเต สมาทเปนฺเต สมุตฺเตเชนฺเต สมฺปหํเสนฺเต อหุเทว อกฺขนฺติ, อหุ อปฺปจฺจโย “อธิสลฺลิขเต\-วายํ สมโณ”ติ. ยโต จ โข ตฺวํ กสฺสป อจฺจยํ อจฺจยโต ทิสฺวา ยถาธมฺมํ ปฏิกโรสิ. ตํ เต มยํ ปฏิคฺคณฺหาม. วุทฺธิเหสา กสฺสป อริยสฺส วินเย, โย อจฺจยํ อจฺจยโต ทิสฺวา ยถาธมฺมํ ปฏิกโรติ, อายติํ สํวรํ อาปชฺชติ.}

\prose{เถโร เจปิ กสฺสป ภิกฺขุ โหติ น สิกฺขากาโม น สิกฺขา\-สมาทานสฺส วณฺณวาที. เย จญฺเญ ภิกฺขู น สิกฺขากามา, เต จ น สิกฺขาย สมาทเปติ. เย จญฺเญ ภิกฺขู สิกฺขากามา, เตสญฺจ น วณฺณํ ภณติ ภูตํ ตจฺฉํ กาเลน. เอวรูปสฺสาหํ กสฺสป เถรสฺส ภิกฺขุโน น วณฺณํ ภณามิ. ตํ กิสฺส เหตุ? “สตฺถา หิสฺส วณฺณํ ภณาตี”ติ อญฺเญ นํ\csromansharedfootnote{c5707113bbe8a186}{ตํ (สี, อิ)} ภิกฺขู ภเชยฺยุํ. เย นํ ภเชยฺยุํ, ตฺยาสฺส ทิฏฺฐานุคติํ อาปชฺเชยฺยุํ. ยฺยาสฺส ทิฏฺฐานุคติํ อาปชฺเชยฺยุํ, เตสํ ตํ อสฺส ทีฆรตฺตํ อหิตาย ทุกฺขายาติ. ตสฺมาหํ กสฺสป เอวรูปสฺส เถรสฺส ภิกฺขุโน น วณฺณํ ภณามิ.}

\prose{มชฺฌิโม เจปิ กสฺสป ภิกฺขุ โหติ ฯเปฯ. นโว เจปิ กสฺสป ภิกฺขุ โหติ น สิกฺขากาโม น สิกฺขา\-สมาทานสฺส วณฺณวาที. เย จญฺเญ ภิกฺขู น สิกฺขากามา, เต จ น สิกฺขาย สมาทเปติ. เย จญฺเญ ภิกฺขู สิกฺขากามา, เตสญฺจ น วณฺณํ ภณติ ภูตํ ตจฺฉํ กาเลน. เอวรูปสฺสาหํ กสฺสป นวสฺส ภิกฺขุโน น วณฺณํ ภณามิ. ตํ กิสฺส เหตุ? “สตฺถา หิสฺส วณฺณํ ภณตี”ติ อญฺเญ นํ ภิกฺขู ภเชยฺยุํ. เย นํ ภเชยฺยุํ, ตฺยาสฺส ทิฏฺฐานุคติํ อาปชฺเชยฺยุํ. ยฺยาสฺส ทิฏฺฐานุคติํ อาปชฺเชยฺยุํ, เตสํ ตํ อสฺส ทีฆรตฺตํ อหิตาย ทุกฺขายาติ. ตสฺมาหํ กสฺสป เอวรูปสฺส นวสฺส ภิกฺขุโน น วณฺณํ ภณามิ.}

\chapterheadpageread{(10) 5. โลณกปลฺล\-วคฺค}
\prose{\csromanfolio{241}เถโร เจปิ กสฺสป ภิกฺขุ โหติ สิกฺขากาโม สิกฺขา\-สมาทานสฺส วณฺณวาที. เย จญฺเญ ภิกฺขู น สิกฺขากามา, เต จ สิกฺขาย สมาทเปติ. เย จญฺเญ ภิกฺขู สิกฺขากามา, เตสญฺจ วณฺณํ ภณติ ภูตํ ตจฺฉํ กาเลน. เอวรูปสฺสาหํ กสฺสป เถรสฺส ภิกฺขุโน วณฺณํ ภณามิ. ตํ กิสฺส เหตุ? “สตฺถา หิสฺส วณฺณํ ภณตี”ติ อญฺเญ นํ ภิกฺขู ภเชยฺยุํ. เย นํ ภเชยฺยุํ, ตฺยาสฺส ทิฏฺฐานุคติํ อาปชฺเชยฺยุํ. ยฺยาสฺส ทิฏฺฐานุคติํ อาปชฺเชยฺยุํ, เตสํ ตํ อสฺส ทีฆรตฺตํ หิตาย สุขายาติ. ตสฺมาหํ กสฺสป เอวรูปสฺส เถรสฺส ภิกฺขุโน วณฺณํ ภณามิ.}

\prosewithclosermidruled{มชฺฌิโม เจปิ กสฺสป ภิกฺขุ โหติ สิกฺขากาโม ฯเปฯ. นโว เจปิ กสฺสป ภิกฺขุ โหติ สิกฺขากาโม สิกฺขา\-สมาทานสฺส วณฺณวาที. เย จญฺเญ ภิกฺขู น สิกฺขากามา, เต จ สิกฺขาย สมาทเปติ. เย จญฺเญ ภิกฺขู สิกฺขากามา, เตสญฺจ วณฺณํ ภณติ ภูตํ ตจฺฉํ กาเลน. เอวรูปสฺสาหํ กสฺสป นวสฺส ภิกฺขุโน วณฺณํ ภณามิ, ตํ กิสฺส เหตุ? “สตฺถา หิสฺส วณฺณํ ภณตี”ติ อญฺเญ นํ ภิกฺขู ภเชยฺยุํ. เย นํ ภเชยฺยุํ, ตฺยาสฺส ทิฏฺฐานุคติํ อาปชฺเชยฺยุํ. ยฺยาสฺส ทิฏฺฐานุคติํ อาปชฺเชยฺยุํ, เตสํ ตํ อสฺส ทีฆรตฺตํ หิตาย สุขายาติ. ตสฺมาหํ กสฺสป เอวรูปสฺส นวสฺส ภิกฺขุโน วณฺณํ ภณามีติ.  เอกาทสมํ.}{สมณวคฺโค จตุตฺโถ.}
\csromancenter{\textbf{ตสฺสุทฺทานํ}}
\setlength{\csromangathapagewidth}{0pt}
\setlength{\csromangathawakwidth}{0pt}
\csromangathameasuregroup{สมโณ คทฺรโภ เขตฺตํ, \\ ตโย จ สิกฺขนา วุตฺตา,}{\csromangathabat{สมโณ คทฺรโภ เขตฺตํ,}{วชฺชิปุตฺโต จ เสกฺขกํ.} \\ \csromangathabat{ตโย จ สิกฺขนา วุตฺตา,}{เทฺว สิกฺขา สงฺกวาย จาติ.}}
\csromangathasetleft{สมโณ คทฺรโภ เขตฺตํ, \\ ตโย จ สิกฺขนา วุตฺตา,}
\csromangathagroup{\csromangathastanza{\csromangathabat{สมโณ คทฺรโภ เขตฺตํ,}{วชฺชิปุตฺโต จ เสกฺขกํ.} \\ \csromangathabat{ตโย จ สิกฺขนา วุตฺตา,}{เทฺว สิกฺขา สงฺกวาย จาติ.}}}

\chapterhead{(10) 5. โลณกปลฺล\-วคฺค}
\csromanheader{2}{1. \textbf{อจฺจายิกสุตฺต}}
\proseitem{93}{ตีณิมานิ ภิกฺขเว กสฺสกสฺส คหปติสฺส อจฺจายิกานิ กรณียานิ. กตมานิ ตีณิ? อิธ ภิกฺขเว กสฺสโก คหปติ สีฆํ สีฆํ เขตฺตํ สุกฏฺฐํ กโรติ สุมติกตํ. สีฆํ สีฆํ เขตฺตํ สุกฏฺฐํ \csromanfolio{242}กริตฺวา สุมติกตํ สีฆํ สีฆํ พีชานิ ปติฏฺฐาเปติ. สีฆํ สีฆํ พีชานิ ปติฏฺฐาเปตฺวา สีฆํ สีฆํ อุทกํ อภิเนติปิ อปเนติปิ. อิมนิ โข ภิกฺขเว ตีณิ กสฺสกสฺส คหปติสฺส อจฺจายิกานิ กรณียานิ. ตสฺส โข ตํ ภิกฺขเว กสฺสกสฺส คหปติสฺส นตฺถิ สา อิทฺธิ วา อานุภาโว วา “อชฺเชว เม ธญฺญานิ ชายนฺตุ, เสฺวว คพฺภีนิ โหนฺตุ, อุกฺกรเสฺวว ปจฺจนฺตู”ติ. อถ โข ภิกฺขเว โหติ โส สมโย, ยํ ตสฺส กสฺสกสฺส คหปติสฺส ตานิ ธญฺญานิ อุตุปริณามีนิ ชายนฺติปิ, คพฺภีนิปิ โหนฺติ ปจฺจนฺติปิ.}

\prose{เอวเมวํ โข ภิกฺขเว ตีณิมานิ ภิกฺขุสฺส อจฺจายิกานิ กรณียานิ. กตมานิ ตีณิ? อธิสีล\-สิกฺขา\-สมาทานํ อธิจิตฺต\-สิกฺขา\-สมาทานํ อธิปญฺญา\-สิกฺขา\-สมาทานํ. อิมานิ โข ภิกฺขเว ตีณิ ภิกฺขุสฺส อจฺจายิกานิ กรณียานิ. ตสฺส โข ตํ ภิกฺขเว ภิกฺขุโน นตฺถิ สา อิทฺธิ วา อานุภาโว วา “อชฺเชว เม อนุปาทาย อาสเวหิ จิตฺตํ วิมุจฺจตุ เสฺว วา อุตฺตรเสฺว วา”ติ. อถ โข ภิกฺขเว โหติ โส สมโย, ยํ ตสฺส ภิกฺขุโน อธิสีลมฺปิ สิกฺขโต อธิจิตฺตมฺปิ สิกฺขโต อธิปญฺญมฺปิ สิกฺขโต อนุปาทาย อาสเวหิ จิตฺตํ วิมุจฺจติ.}

\prose{ตสฺมาติห ภิกฺขเว เอวํ สิกฺขิตพฺพํ “ติพฺโพ โน ฉนฺโท ภวิสฺสติ อธิสีล\-สิกฺขาสมาทาเน, ติพฺโพ ฉนฺโท ภวิสฺสติ อธิจิตฺต\-สิกฺขาสมาทาเน, ติพฺโพ ฉนฺโท ภวิสฺสติ อธิปญฺญา\-สิกฺขาสมาทาเน”ติ. เอวํ หิ โว ภิกฺขเว สิกฺขิตพฺพนฺติ.  ปฐมํ.}

\csromanheader{2}{2. \textbf{ปวิเวกสุตฺต}}
\proseitem{94}{ตีณิมานิ ภิกฺขเว อญฺญติตฺถิยา ปริพฺพาชกา ปวิเวกานิ ปญฺญาเปนฺติ. กตมานิ ตีณิ? จีวร\-ปวิเวกํ ปิณฺฑปาต\-ปวิเวกํ เสนาสน\-ปวิเวกํ.}

\prose{ตตฺริทํ ภิกฺขเว อญฺญติตฺถิยา ปริพฺพาชกา จีวร\-ปวิเวกสฺมิํ ปญฺญาเปนฺติ, สาณานิปิ ธาเรนฺติ, มสาณานิปิ ธาเรนฺติ, ฉวทุสฺสานิปิ ธาเรนฺติ, ปํสุกูลานิปิ ธาเรนฺติ, ติรีฏานิปิ ธาเรนฺติ, อชินมฺปิ ธาเรนฺติ, อชินกฺขิปมฺปิ ธาเรนฺติ, กุสจีรมฺปิ ธาเรนฺติ, วากจีรมฺปิ ธาเรนฺติ, ผลกจีรมฺปิ ธาเรนฺติ, เกสกมฺพลมฺปิ ธาเรนฺติ, วาลกมฺพลมฺปิ ธาเรนฺติ,}

\vspace{\tipitakaparskip}
\csromancenter{(10) 5. โลณกปลฺล\-วคฺค อุลูก\-ปกฺขิกมฺปิ ธาเรนฺติ. อิทํ โข ภิกฺขเว อญฺญติตฺถิยา ปริพฺพาชกา จีวร\-ปวิเวกสฺมิํ ปญฺญาเปนฺติ.}
\prose{\csromanfolio{243}ตตฺริทํ ภิกฺขเว อญฺญติตฺถิยา ปริพฺพาชกา ปิณฺฑปาต\-ปวิเวกสฺมิํ ปญฺญาเปนฺติ, สากภกฺขาปิ โหนฺติ, สามากภกฺขาปิ โหนฺติ, นีวารภกฺขาปิ โหนฺติ, ททฺทุลภกฺขาปิ โหนฺติ, หฏภกฺขาปิ โหนฺติ, กณภกฺขาปิ โหนฺติ, อาจามภกฺขาปิ โหนฺติ, ปิญฺญากภกฺขาปิ โหนฺติ, ติณภกฺขาปิ โหนฺติ, โคมยภกฺขาปิ โหนฺติ, วนมูล\-ผลา\-หารา ยาเปนฺติ ปวตฺต\-ผล\-โภชี. อิทํ โข ภิกฺขเว อญฺญติตฺถิยา ปริพฺพาชกา ปิณฺฑปาต\-ปวิเวกสฺมิํ ปญฺญาเปนฺติ.}

\prose{ตตฺริทํ ภิกฺขเว อญฺญติตฺถิยา ปริพฺพาชกา เสนาสน\-ปวิเวกสฺมิํ ปญฺญาเปนฺติ, อรญฺญํ รุกฺขมูลํ สุสานํ\csromansharedfootnote{cfc0a44499390a1e}{รุกฺขมูลํ ภุสาคารํ สุสานํ (ก)} วนปตฺถํ อพฺโภกาสํ ปลาลปุญฺชํ ภุสาคารํ\csromansharedfootnote{552557f9bc38b8e3}{สุญฺญาคารํ (ก)}. อิทํ โข ภิกฺขเว อญฺญติตฺถิยา ปริพฺพาชกา เสนาสน\-ปวิเวกสฺมิํ ปญฺญาเปนฺติ. อิมานิ โข ภิกฺขเว ตีณิ อญฺญติตฺถิยา ปริพฺพาชกา ปวิเวกานิ ปญฺญาเปนฺติ.}

\prose{ตีณิ โข ปนิมานิ ภิกฺขเว อิมสฺมิํ ธมฺมวินเย ภิกฺขุโน ปวิเวกานิ. กตมานิ ตีณิ? อิธ ภิกฺขเว ภิกฺขุ สีลวา จ โหติ, ทุสฺสีลฺยญฺจสฺส ปหีนํ โหติ, เตน จ วิวิตฺโต โหติ. สมฺมาทิฏฺฐิโก จ โหติ, มิจฺฉาทิฏฺฐิ จสฺส ปหีนา โหติ, ตาย จ วิวิตฺโต โหติ. ขีณาสโว จ โหติ, อาสวา จสฺส ปหีนา โหนฺติ, เตหิ จ วิวิตฺโต โหติ. ยโต โข ภิกฺขเว ภิกฺขุ สีลวา โหติ, ทุสฺสีลฺยญฺจสฺส ปหีนํ โหติ, เตน จ วิวิตฺโต โหติ. สมฺมาทิฏฺฐิโก จ โหติ, มิจฺฉาทิฏฺฐิ จสฺส ปหีนา โหติ, ตาย จ วิวิตฺโต โหติ. ขีณาสโว จ โหติ, อาสวา จสฺส ปหีนา โหนฺติ, เตหิ จ วิวิตฺโต โหติ. อยํ วุจฺจติ ภิกฺขเว ภิกฺขุ อคฺคปฺปตฺโต สารปฺปตฺโต สุทฺโธ สาเร ปติฏฺฐิโต.}

\prose{เสยฺยถาปิ ภิกฺขเว กสฺสกสฺส คหปติสฺส สมฺปนฺนํ สาลิกฺเขตฺตํ, ตเมนํ กสฺสโก คหปติ สีฆํ สีฆํ\csromansharedfootnote{d9d5d77fda431127}{สีฆสีฆํ (สี, สฺยา, กํ, อิ)} ลวาเปยฺย, สีฆํ สีฆํ ลวาเปตฺวา สีฆํ สีฆํ สงฺฆาราเปยฺย, สีฆํ สีฆํ สงฺฆราเปตฺวา สีฆํ \csromanfolio{244}สีฆํ อุพฺพหาเปยฺย\csromansharedfootnote{1e39247c1462b137}{อุพฺพาหาเปยฺย (สฺยา, กํ)}, สีฆํ สีฆํ อุพฺพหาเปตฺวา สีฆํ สีฆํ ปุญฺชํ การาเปยฺย, สีฆํ สีฆํ ปุญฺชํ การาเปตฺวา สีฆํ สีฆํ มทฺทาเปยฺย, สีฆํ สีฆํ มทฺทาเปตฺวา สีฆํ สีฆํ ปลาลานิ อุทฺธราเปยฺย, สีฆํ สีฆํ ปลาลานิ อุทฺธราเปตฺวา สีฆํ สีฆํ ภุสิกํ อุทฺธราเปยฺย, สีฆํ สีฆํ ภุสิกํ อุทฺธราเปตฺวา สีฆํ สีฆํ โอปุนาเปยฺย, สีฆํ สีฆํ โอปุนาเปตฺวา สีฆํ สีฆํ อติหราเปยฺย, สีฆํ สีฆํ อติหราเปตฺวา สีฆํ สีฆํ โกฏฺฏาเปยฺย, สีฆํ สีฆํ โกฏฺฏาเปตฺวา สีฆํ สีฆํ ถุสานิ อุทฺธราเปยฺย. เอวมสฺสุ\csromansharedfootnote{70bcadee0f04c6ff}{เอวสฺสุ (ก)} ตานิ ภิกฺขเว กสฺสกสฺส คหปติสฺส ธญฺญานิ อคฺคปฺปตฺตานิ สารปฺปตฺตานิ สุทฺธานิ สาเร ปติฏฺฐิตานิ.}

\prose{เอวเมวํ โข ภิกฺขเว ยโต ภิกฺขุ สีลวา จ โหติ, ทุสฺสีลฺยญฺจสฺส ปหีนํ โหติ, เตน จ วิวิตฺโต โหติ. สมฺมาทิฏฺฐิโก จ โหติ, มิจฺฉาทิฏฺฐิ จสฺส ปหีนา โหติ, ตาย จ วิวิตฺโต โหติ. ขีณาสโว จ โหติ, อาสวา จสฺส ปหีนา โหนฺติ, เตหิ จ วิวิตฺโต โหติ. อยํ วุจฺจติ ภิกฺขเว ภิกฺขุ อคฺคปฺปตฺโต สารปฺปตฺโต สุทฺโธ สาเร ปติฏฺฐิโตติ.  ทุติยํ.}

\csromanheader{2}{3. \textbf{สรทสุตฺต}}
\proseitem{95}{เสยฺยถาปิ ภิกฺขเว สรทสมเย วิทฺเธ วิคตวลาหเก เทเว อาทิจฺโจ นภํ อพฺภุสฺสกฺกมาโน\csromansharedfootnote{cd143aa481dd025a}{อพฺภุสฺสุกฺกมาโน (สี, อิ)} สพฺพํ อากาสคตํ ตมคตํ อภิวิหจฺจ ภาสเต จ ตปเต จ วิโรจติ จ.}

\prose{เอวเมวํ โข ภิกฺขเว ยโต อริยสาวกสฺส วิรชํ วีตมลํ ธมฺมจกฺขุํ อุปฺปชฺชติ\csromansharedfootnote{73dffa30ba116534}{อุทปาทิ (สพฺพตฺถ)}, สห ทสฺสนุปฺปาทา ภิกฺขเว อริยสาวกสฺส ตีณิ สํโยชนานิ ปหียนฺติ สกฺกายทิฏฺฐิ วิจิกิจฺฉา สีลพฺพต\-ปรามาโส.}

\prose{อถาปรํ ทฺวีหิ ธมฺเมหิ นิยฺยาติ อภิชฺฌาย จ พฺยาปาเทน จ, โส วิวิจฺเจว กาเมหิ วิวิจฺจ อกุสเลหิ ธมฺเมหิ สวิตกฺกํ สวิจารํ วิเวกชํ ปีติสุกํ ปฐมํ ฌานํ อุปสมฺปชฺช วิหรติ. ตสฺมิํ เจ ภิกฺขเว สมเย อริยสาวโก กาลํ กเรยฺย. นตฺถิ ตํ\csromansharedfootnote{f4c06bebf72a231c}{ตสฺส (ก)} สํโยชนํ}

\vspace{\tipitakaparskip}
\csromancenter{(10) 5. โลณกปลฺล\-วคฺค เยน สํโยชเนน สํยุตฺโต อริยสาวโก ปุน อิมํ\csromansharedfootnote{0d6736e7805ea331}{ปุนยิมํ (สฺยา, กํ, ก)} โลกํ อาคจฺเฉยฺยาติ.  ตติยํ.}
\csromanheader{2}{4. \textbf{ปริสาสุตฺต}}
\proseitem{96}{\csromanfolio{245}ติสฺโส อิมา ภิกฺขเว ปริสา. กตมา ติสฺโส? อคฺควตี ปริสา, วคฺคา ปริสา, สมคฺคา ปริสา.}

\prose{กตมา จ ภิกฺขเว อคฺควตี ปริสา? อิธ ภิกฺขเว ยสฺสํ ปริสายํ เถรา ภิกฺขู น พาหุลิกา โหนฺติ น สาถลิกา โอกฺกมเน นิกฺขิตฺตธุรา ปวิเวเก ปุพฺพงฺคมา, วีริยํ อารภนฺติ อปฺปตฺตสฺส ปตฺติยา อนธิคตสฺส อธิคมาย อสจฺฉิกตสฺส สจฺฉิกิริยาย, เตสํ ปจฺฉิมา ชนตา ทิฏฺฐานุคติํ อาปชฺชติ. สาปิ โหติ น พาหุลิกา น สาถลิกา โอกฺกมเน นิกฺขิตฺตธุรา ปวิเวเก ปุพฺพงฺคมา, วีริยํ อารภติ อปฺปตฺตสฺส ปตฺติยา อนธิคตสฺส อธิคมาย อสจฺฉิกตสฺส สจฺฉิกิริยาย. อยํ วุจฺจติ ภิกฺขเว อคฺควตี ปริสา.}

\prose{กตมา จ ภิกฺขเว วคฺคา ปริสา? อิธ ภิกฺขเว ยสฺสํ ปริสายํ ภิกฺขู ภณฺฑนชาตา กลหชาตา วิวาทาปนฺนา อญฺญมญฺญํ มุขสตฺตีหิ วิตุทนฺตา วิหรนฺติ. อยํ วุจฺจติ ภิกฺขเว วคฺคา ปริสา.}

\prose{กตมา จ ภิกฺขเว สมคฺคา ปริสา? อิธ ภิกฺขเว ยสฺสํ ปริสายํ ภิกฺขู สมคฺคา สมฺโมทมานา อวิวทมานา ขีโรทกีภูตา อญฺญมญฺญํ ปิยจกฺขูหิ สมฺปสฺสนฺตา วิหรนฺติ. อยํ วุจฺจติ ภิกฺขเว สมคฺคา ปริสา.}

\prose{ยสฺมิํ ภิกฺขเว สมเย ภิกฺขู สมคฺคา สมฺโมทมานา อวิวทมานา ขีโรทกีภูตา อญฺญมญฺญํ ปิยจกฺขูหิ สมฺปสฺสนฺตา วิหรนฺติ. พหุํ ภิกฺขเว ภิกฺขู ตสฺมิํ สมเย ปุญฺญํ ปสวนฺติ, พฺรหฺมํ ภิกฺขเว วิหารํ ตสฺมิํ สมเย ภิกฺขู วิหรนฺติ, ยทิทํ มุทิตาย เจโตวิมุตฺติยา. ปมุทิตสฺส ปีติ ชายติ, ปีติมนสฺส กาโย ปสฺสมฺภติ, ปสฺสทฺธกาโย สุขํ เวทิยติ, สุขิโน จิตฺตํ สมาธิยติ.}

\prose{เสยฺยถาปิ ภิกฺขเว อุปริปพฺพเต ถุลฺล\-ผุสิตเก เทเว วสฺสนฺเต ตํ อุทกํ ยถานินฺนํ ปวตฺตมานํ ปพฺพต\-กนฺทร\-ปทร\-สาขา ปริปูเรติ, ปพฺพต\-กนฺทร\-ปทร\-สาขา ปริปูรา กุโสพฺเภ\csromansharedfootnote{46a9b7fb2e3ae2a3}{กุสฺสุพฺเภ (สี, อิ), กุสุมฺเภ (สฺยา, กํ, ก)} ปริปูเรนฺติ, กุโสพฺภา ปริปูรา \csromanfolio{246}มหาโสพฺเภ ปริปูเรนฺติ, มหาโสพฺภา ปริปูรา กุนฺนทิโย ปริปูเรนฺติ, กุนฺนทิโย ปริปูรา มหานทิโย ปริปูเรนฺติ, มหานทิโย ปริปูรา สมุทฺทํ\csromansharedfootnote{b6eab0c2da86d09c}{สมุทฺทสาคเร (ก)} ปริปูเรนฺติ.}

\prose{เอวเมวํ โข ภิกฺขเว ยสฺมิํ สมเย ภิกฺขู สมคฺคา สมฺโมทมานา อวิวทมานา ขีโรทกีภูตา อญฺญมญฺญํ ปิยจกฺขูหิ สมฺปสฺสนฺตา วิหรนฺติ, พหุํ ภิกฺขเว ภิกฺขู ตสฺมิํ สมเย ปุญฺญํ ปสวนฺติ, พฺรหฺมํ ภิกฺขเว วิหารํ ตสฺมิํ สมเย ภิกฺขู วิหรนฺติ, ยทิทํ มุทิตาย เจโตวิมุตฺติยา. ปมุทิตสฺส ปีติ ชายติ, ปีติมนสฺส กาโย ปสฺสมฺภติ, ปสฺสทฺธกาโย สุขํ เวทิยติ, สุขิโน จิตฺตํ สมาธิยติ. อิมา โข ภิกฺขเว ติสฺโส ปริสาติ.  จตุตฺถํ.}

\prose{5. ปฐม\-อาชานีย\-สุตฺต}

\proseitem{97}{ตีหิ ภิกฺขเว องฺเคหิ สมนฺนาคโต รญฺโญ ภโทฺร\csromansharedfootnote{8fe84201a79db3ce}{ภทฺโท (ก)} อสฺสาชานีโย ราชารโห โหติ ราชโภคฺโค, รญฺโญ องฺคนฺเตว สงฺขฺยํ\csromansharedfootnote{f3b84ba968ed1c5b}{สงฺขํ (สี, สฺยา, กํ, อิ)} คจฺฉติ. กตเมหิ ตีหิ? อิธ ภิกฺขเว รญฺโญ ภโทฺร อสฺสาชานีโย วณฺณสมฺปนฺโน จ โหติ พลสมฺปนฺโน จ ชวสมฺปนฺโน จ. อิเมหิ โข ภิกฺขเว ตีหิ องฺเคหิ สมนฺนาคโต รญฺโญ ภโทฺร อสฺสาชานีโย ราชารโห โหติ ราชโภคฺโค, รญฺโญ องฺคนฺเตว สงฺขฺยํ คจฺฉติ. เอวเมวํ โข ภิกฺขเว ตีหิ ธมฺเมหิ สมนฺนาคโต ภิกฺขุ อาหุเนยฺโย โหติ ปาหุเนยฺโย ทกฺขิเณยฺโย อญฺชลิกรณีโย อนุตฺตรํ ปุญฺญกฺเขตฺตํ โลกสฺส. กตเมหิ ตีหิ? อิธ ภิกฺขเว ภิกฺขุ วณฺณสมฺปนฺโน จ โหติ พลสมฺปนฺโน จ ชวสมฺปนฺโน จ.}

\prose{กถญฺจ ภิกฺขเว ภิกฺขุ วณฺณสมฺปนฺโน โหติ? อิธ ภิกฺขเว ภิกฺขุ สีลวา โหติ, ปาติโมกฺข\-สํวร\-สํวุโต วิหรติ อาจารโคจร\-สมฺปนฺโน, อณุมตฺเตสุ วชฺเชสุ ภยทสฺสาวี สมาทาย สิกฺขติ สิกฺขาปเทสุ. เอวํ โข ภิกฺขเว ภิกฺขุ วณฺณสมฺปนฺโน โหติ.}

\prose{กถญฺจ ภิกฺขเว ภิกฺขุ พลสมฺปนฺโน โหติ? อิธ ภิกฺขเว ภิกฺขุ อารทฺธวีริโย วิหรติ อกุสลานํ ธมฺมานํ ปหานาย กุสลานํ}

\vspace{\tipitakaparskip}
\csromancenter{(10) 5. โลณกปลฺล\-วคฺค ธมฺมานํ อุปสมฺปทาย ถามวา ทฬฺหปรกฺกโม อนิกฺขิตฺตธุโร กุสเลสุ ธมฺเมสุ. เอวํ โข ภิกฺขเว ภิกฺขุ พลสมฺปนฺโน โหติ.}
\prose{\csromanfolio{247}กถญฺจ ภิกฺขเว ภิกฺขุ ชวสมฺปนฺโน โหติ? อิธ ภิกฺขเว ภิกฺขุ “อิทํ ทุกฺขนฺ”ติ ยถาภูตํ ปชานาติ, “อยํ ทุกฺขสมุทโย”ติ ยถาภูตํ ปชานาติ, “อยํ ทุกฺขนิโรโธ”ติ ยถาภูตํ ปชานาติ, “อยํ ทุกฺขนิโรธ\-คามินี ปฏิปทา”ติ ยถาภูตํ ปชานาติ. เอวํ โข ภิกฺขเว ภิกฺขุ ชวสมฺปนฺโน โหติ. อิเมหิ โข ภิกฺขเว ตีหิ ธมฺเมหิ สมนฺนาคโต ภิกฺขุ อาหุเนยฺโย โหติ ปาหุเนยฺโย ทกฺขิเณยฺโย อญฺชลิกรณีโย อนุตฺตรํ ปุญฺญกฺเขตฺตํ โลกสฺสาติ.  ปญฺจมํ.}

\prose{6. ทุติย\-อาชานีย\-สุตฺต}

\proseitem{98}{ตีหิ ภิกฺขเว องฺเคหิ สมนฺนาคโต รญฺโญ ภโทฺร อสฺสาชานีโย ราชารโห โหติ ราชโภคฺโค, รญฺโญ องฺคนฺเตว สงฺขฺยํ คจฺฉติ. กตเมหิ ตีหิ? อิธ ภิกฺขเว รญฺโญ ภโทฺร อสฺสาชานีโย วณฺณสมฺปนฺโน จ โหติ พลสมฺปนฺโน จ ชวสมฺปนฺโน จ. อิเมหิ โข ภิกฺขเว ตีหิ องฺเคหิ สมนฺนาคโต รญฺโญ ภโทฺร อสฺสาชานีโย ราชารโห โหติ ราชโภคฺโค, รญฺโญ องฺคนฺเตว สงฺขฺยํ คจฺฉติ. เอวเมวํ โข ภิกฺขเว ตีหิ ธมฺเมหิ สมนฺนาคโต ภิกฺขุ อาหุเนยฺโย โหติ ฯเปฯ อนุตฺตรํ ปุญฺญกฺเขตฺตํ โลกสฺส. กตเมหิ ตีหิ? อิธ ภิกฺขเว ภิกฺขุ วณฺณสมฺปนฺโน จ โหติ พลสมฺปนฺโน จ ชวสมฺปนฺโน จ.}

\prose{กถญฺจ ภิกฺขเว ภิกฺขุ วณฺณสมฺปนฺโน โหติ? อิธ ภิกฺขเว ภิกฺขุ สีลวา โหติ ฯเปฯ สมาทาย สิกฺขติ สิกฺขาปเทสุ. เอวํ โข ภิกฺขเว ภิกฺขุ วณฺณสมฺปนฺโน โหติ.}

\prose{กถญฺจ ภิกฺขเว ภิกฺขุ พลสมฺปนฺโน โหติ? อิธ ภิกฺขเว ภิกฺขุ อารทฺธวีริโย วิหรติ อกุสลานํ ธมฺมานํ ปหานาย กุสลานํ ธมฺมานํ อุปสมฺปทาย ถามวา ทฬฺหปรกฺกโม อนิกฺขิตฺตธุโร กุสเลสุ ธมฺเมสุ. เอวํ โข ภิกฺขเว ภิกฺขุ พลสมฺปนฺโน โหติ.}

\prose{กถญฺจ ภิกฺขเว ภิกฺขุ ชวสมฺปนฺโน โหติ? อิธ ภิกฺขเว ภิกฺขุ ปญฺจนฺนํ โอรมฺภาคิยานํ สํโยชนานํ ปริกฺขยา โอปปาติโก โหติ ตตฺถ \csromanfolio{248}ปรินิพฺพายี อนาวตฺติธมฺโม ตสฺมา โลกา. เอวํ โข ภิกฺขเว ภิกฺขุ ชวสมฺปนฺโน โหติ. อิเมหิ โข ภิกฺขเว ตีหิ ธมฺเมหิ สมนฺนาคโต ภิกฺขุ อาหุเนยฺโย โหติ ฯเปฯ อนุตฺตรํ ปุญฺญกฺเขตฺตํ โลกสฺสาติ.  ฉฏฺฐํ.}

\prose{7. ตติย\-อาชานีย\-สุตฺต}

\proseitem{99}{ตีหิ ภิกฺขเว องฺเคหิ สมนฺนาคโต รญฺโญ ภโทฺร อสฺสาชานีโย ราชารโห โหติ ราชโภคฺโค, รญฺโญ องฺคนฺเตว สงฺขฺยํ คจฺฉติ. กตเมหิ ตีหิ? อิธ ภิกฺขเว รญฺโญ ภโทฺร อสฺสาชานีโย วณฺณสมฺปนฺโน จ โหติ พลสมฺปนฺโน จ ชวสมฺปนฺโน จ. อิเมหิ โข ภิกฺขเว ตีหิ องฺเคหิ สมนฺนาคโต รญฺโญ ภโทฺร อสฺสาชานีโย ราชารโห โหติ ราชโภคฺโค, รญฺโญ องฺคนฺเตว สงฺขฺยํ คจฺฉติ. เอวเมวํ โข ภิกฺขเว ตีหิ ธมฺเมหิ สมนฺนาคโต ภิกฺขุ อาหุเนยฺโย โหติ ปาหุเนยฺโย ทกฺขิเณยฺโย อญฺชลิกรณีโย อนุตฺตรํ ปุญฺญกฺเขตฺตํ โลกสฺส. กตเมหิ ตีหิ? อิธ ภิกฺขเว ภิกฺขุ วณฺณสมฺปนฺโน จ โหติ พลสมฺปนฺโน จ ชวสมฺปนฺโน จ.}

\prose{กถญฺจ ภิกฺขเว ภิกฺขุ วณฺณสมฺปนฺโน โหติ? อิธ ภิกฺขเว ภิกฺขุ สีลวา โหติ, ปาติโมกฺข\-สํวร\-สํวุโต วิหรติ, อาจารโคจร\-สมฺปนฺโน อณุมตฺเตสุ วชฺเชสุ ภยทสฺสาวี สมาทาย สิกฺขติ สิกฺขาปเทสุ. เอวํ โข ภิกฺขเว ภิกฺขุ วณฺณสมฺปนฺโน โหติ.}

\prose{กถญฺจ ภิกฺขเว ภิกฺขุ พลสมฺปนฺโน โหติ? อิธ ภิกฺขเว ภิกฺขุ อารทฺธวีริโย วิหรติ อกุสลานํ ธมฺมานํ ปหานาย กุสลานํ ธมฺมานํ อุปสมฺปทาย ถามวา ทฬฺหปรกฺกโม อนิกฺขิตฺตธุโร กุสเลสุ ธมฺเมสุ. เอวํ โข ภิกฺขเว ภิกฺขุ พลสมฺปนฺโน โหติ.}

\prose{กถญฺจ ภิกฺขเว ภิกฺขุ ชวสมฺปนฺโน โหติ? อิธ ภิกฺขเว ภิกฺขุ อาสวานํ ขยา อนาสวํ เจโตวิมุตฺติํ ปญฺญาวิมุตฺติํ ทิฏฺเฐว ธมฺเม สยํ อภิญฺญา สจฺฉิกตฺวา อุปสมฺปชฺช วิหรติ. เอวํ โข ภิกฺขเว ภิกฺขุ ชวสมฺปนฺโน โหติ. อิเมหิ โข ภิกฺขเว ตีหิ ธมฺเมหิ สมนฺนาคโต ภิกฺขุ อาหุเนยฺยา โหติ ฯเปฯ อนุตฺตรํ ปุญฺญกฺเขตฺตํ โลกสฺสาติ.  สตฺตมํ.}

\chapterheadpageread{(10) 5. โลณกปลฺล\-วคฺค 8. โปตฺถกสุตฺต}
\proseitem{100}{\csromanfolio{249}นโวปิ ภิกฺขเว โปตฺถโก ทุพฺพณฺโณ จ โหติ ทุกฺข\-สมฺผสฺโส จ อปฺปคฺโฆ จ, มชฺฌิโมปิ ภิกฺขเว โปตฺถโก ทุพฺพณฺโณ จ โหติ ทุกฺข\-สมฺผสฺโส จ อปฺปคฺโฆ จ, ชิณฺโณปิ ภิกฺขเว โปตฺถโก ทุพฺพณฺโณ จ โหติ ทุกฺข\-สมฺผสฺโส จ อปฺปคฺโฆ จ. ชิณฺณมฺปิ ภิกฺขเว โปตฺถกํ อุกฺขลิ\-ปริมชฺชนํ วา กโรนฺติ, สงฺการกูเฏ วา นํ\csromansharedfootnote{ee133feb3b41562a}{ตํ (สี), ฐาเน (ก)} ฉฑฺเฑนฺติ.}

\prose{เอวเมวํ โข ภิกฺขเว นโว เจปิ ภิกฺขุ โหติ ทุสฺสีโล ปาปธมฺโม, อิทมสฺส ทุพฺพณฺณตาย วทามิ. เสยฺยถาปิ โส ภิกฺขเว โปตฺถโก ทุพฺพณฺโณ, ตถูปมาหํ ภิกฺขเว อิมํ ปุคฺคลํ วทามิ. เย โข ปนสฺส เสวนฺติ ภชนฺติ ปยิรุปาสนฺติ ทิฏฺฐานุคติํ อาปชฺชนฺติ, เตสํ ตํ โหติ ทีฆรตฺตํ อหิตาย ทุกฺขาย, อิทมสฺส ทุกฺข\-สมฺผสฺสตาย วทามิ. เสยฺยถาปิ โส ภิกฺขเว โปตฺถโก ทุกฺข\-สมฺผสฺโส, ตถูปมาหํ ภิกฺขเว อิมํ ปุคฺคลํ วทามิ. เยสํ โข ปน โส\csromansharedfootnote{08daae3f82072ea7}{เยสํ โข ปน (สี, สฺยา, กํ, อิ), เยสํ โส (ก) อภิ 3. 138 ปิฏฺเฐปิ ปสฺสิตพฺพํ.} ปฏิคฺคณฺหาติ จีวร\-ปิณฺฑปาต\-เสนาสน\-คิลานปฺปจฺจย\-เภสชฺช\-ปริกฺขารํ, เตสํ ตํ น มหปฺผลํ โหติ น มหานิสํสํ, อิทมสฺส อปฺปคฺฆตาย วทามิ. เสยฺยถาปิ โส ภิกฺขเว โปตฺถโก อปฺปคฺโฆ, ตถูปมาหํ ภิกฺขเว อิมํ ปุคฺคลํ วทามิ. มชฺฌิโม เจปิ ภิกฺขเว ภิกฺขุ โหติ ฯเปฯ. เถโร เจปิ ภิกฺขเว ภิกฺขุ โหติ ทุสฺสีโล ปาปธมฺโม, อิทมสฺส ทุพฺพณฺณตาย วทามิ. เสยฺยถาปิ โส ภิกฺขเว โปตฺถโก ทุพฺพณฺโณ, ตถูปมาหํ ภิกฺขเว อิมํ ปุคฺคลํ วทามิ. เย โข ปนสฺส เสวนฺติ ภชนฺติ ปยิรุปาสนฺติ ทิฏฺฐานุคติํ อาปชฺชนฺติ, เตสํ ตํ โหติ ทีฆรตฺตํ อหิตาย ทุกฺขาย. อิทมสฺส ทุกฺข\-สมฺผสฺสตาย วทามิ. เสยฺยถาปิ โส ภิกฺขเว โปตฺถโก ทุกฺข\-สมฺผสฺโส, ตถูปมาหํ ภิกฺขเว อิมํ ปุคฺคลํ วทามิ. เยสํ โข ปน โส ปฏิคฺคณฺหาติ จีวร\-ปิณฺฑปาต\-เสนาสน\-คิลานปฺปจฺจย\-เภสชฺช\-ปริกฺขารํ, เตสํ ตํ น มหปฺผลํ โหติ น มหานิสํสํ, อิทมสฺส อปฺปคฺฆตาย วทามิ. เสยฺยถาปิ โส ภิกฺขเว โปตฺถโก อปฺปคฺโฆ, ตถูปมาหํ ภิกฺขเว อิมํ ปุคฺคลํ วทามิ.}

\prose{เอวรูโป จายํ ภิกฺขเว เถโร ภิกฺขุ สํฆมชฺเฌ ภณติ, ตเมนํ ภิกฺขู เอวมาหํสุ “กิํ นุ โข ตุยฺหํ พาลสฺส อพฺยตฺตสฺส ภณิเตน, ตฺวมฺปิ นาม ภณิตพฺพํ มญฺญสี”ติ. โส กุปิโต อนตฺตมโน ตถารูปิํ วาจํ \csromanfolio{250}นิจฺฉาเรติ, ยถารูปาย วาจาย สํโฆ ตํ อุกฺขิปติ สงฺการกูเฏว นํ โปตฺถกํ.}

\prose{นวมฺปิ ภิกฺขเว กาสิกํ วตฺถํ วณฺณวนฺต\-ญฺเจว โหติ สุขสมฺผสฺสญฺจ มหคฺฆญฺจ, มชฺฌิมมฺปิ ภิกฺขเว กาสิกํ วตฺถํ วณฺณวนฺต\-ญฺเจว โหติ สุขสมฺผสฺสญฺจ มหคฺฆญฺจ, ชิณฺณมฺปิ ภิกฺขเว กาสิกํ วตฺถํ วณฺณวนฺต\-ญฺเจว โหติ สุขสมฺผสฺสญฺจ มหคฺฆญฺจ, ชิณฺณมฺปิ ภิกฺขเว กาสิกํ วตฺถํ รตน\-ปลิเวฐนํ วา กโรติ คนฺธ\-กรณฺฑเก วา นํ ปกฺขิปนฺติ.}

\prose{เอวเมวํ โข ภิกฺขเว นโว เจปิ ภิกฺขุ โหติ สีลวา กลฺยาณธมฺโม, อิทมสฺส สุวณฺณตาย วทามิ. เสยฺยถาปิ ตํ ภิกฺขเว กาสิกํ วตฺถํ วณฺณวนฺตํ, ตถูปมาหํ ภิกฺขเว อิมํ ปุคฺคลํ วทามิ. เย โข ปนสฺส เสวนฺติ ภชนฺติ ปยิรุปาสนฺติ ทิฏฺฐานุคติํ อาปชฺชนฺติ, เตสํ ตํ โหติ ทีฆรตฺตํ หิตาย สุขาย, อิทมสฺส สุขสมฺผสฺสตาย วทามิ. เสยฺยถาปิ ตํ ภิกฺขเว กาสิกํ วตฺถํ สุขสมฺผสฺสํ, ตถูปมาหํ ภิกฺขเว อิมํ ปุคฺคลํ วทามิ. เยสํ โข ปน โส ปฏิคฺคณฺหาติ จีวร\-ปิณฺฑปาต\-เสนาสน\-คิลานปฺปจฺจย\-เภสชฺช\-ปริกฺขารํ, เตสํ ตํ มหปฺผลํ โหติ มหานิสํสํ, อิทมสฺส มหคฺฆตาย วทามิ. เสยฺยถาปิ ตํ ภิกฺขเว กาสิกํ วตฺถํ มหคฺฆํ, ตถูปมาหํ ภิกฺขเว อิมํ ปุคฺคลํ วทามิ. มชฺฌิโม เจปิ ภิกฺขเว ภิกฺขุ โหติ ฯเปฯ. เถโร เจปิ ภิกฺขเว ภิกฺขุ โหติ ฯเปฯ ปุคฺคลํ วทามิ.}

\prose{เอวรูโป จายํ ภิกฺขเว เถโร ภิกฺขุ สํฆมชฺเฌ ภณติ, ตเมนํ ภิกฺขู เอวมาหํสุ “อปฺปสทฺทา อายสฺมนฺโต โหถ, เถโร ภิกฺขุ ธมฺมญฺจ วินยญฺจ ภณตี”ติ. ตสฺมาติห ภิกฺขเว เอวํ สิกฺขิตพฺพํ “กาสิก\-วตฺถู\-ปมา ภวิสฺสาม, น โปตฺถกูปมา”ติ\csromansharedfootnote{5d5222dbf438ae22}{กาสิกํ วตฺถํ ตถูปมาหํ ภวิสฺสามิ, น โปตฺถกูปมาหนฺติ (ก)}. เอวํ หิ โว ภิกฺขเว สิกฺขิตพฺพนฺติ.  อฏฺฐมํ.}

\csromanheader{2}{9. โลณกปลฺล\-สุตฺต}
\proseitem{101}{โย\csromansharedfootnote{038a2deeed89c980}{โย โข (สฺยา, กํ), โย จ โข (ก)} ภิกฺขเว เอวํ วเทยฺย “ยถา ยถายํ ปุริโส กมฺมํ กโรติ, ตถา ตถา ตํ ปฏิสํเวทิยตี”ติ. เอวํ สนฺตํ ภิกฺขเว พฺรหฺมจริย\-วาโส น โหติ, โอกาโส น ปญฺญายติ สมฺมา ทุกฺขสฺส}

\vspace{\tipitakaparskip}
\csromancenter{(10) 5. โลณกปลฺล\-วคฺค อนฺตกิริยาย. โย จ โข ภิกฺขเว เอวํ วเทยฺย “ยถา ยถา เวทนียํ อยํ ปุริโส กมฺมํ กโรติ, ตถา ตถาสฺส วิปากํ ปฏิสํเวทิยตี”ติ. เอวํ สนฺตํ ภิกฺขเว พฺรหฺมจริย\-วาโส โหติ, โอกาโส ปญฺญายติ สมฺมา ทุกฺขสฺส อนฺตกิริยาย. อิธ ภิกฺขเว เอกจฺจสฺส ปุคฺคลสฺส อปฺปมตฺตกมฺปิ ปาปกมฺมํ\csromansharedfootnote{65d3d041fe6c9a76}{ปาปํ กมฺมํ (สี, อิ)} กตํ, ตเมนํ นิรยํ อุปเนติ. อิธ ปน ภิกฺขเว เอกจฺจสฺส ปุคฺคลสฺส ตาทิสํเยว อปฺปมตฺตกํ ปาปกมฺมํ กตํ ทิฏฺฐธมฺม\-เวทนียํ โหติ, นา’ณุปิ ขายติ กิํ พหุเทว.}
\prose{\csromanfolio{251}กถํรูปสฺส ภิกฺขเว ปุคฺคลสฺส อปฺปมตฺตกมฺปิ ปาปกมฺมํ กตํ, ตเมนํ นิรยํ อุปเนติ? อิธ ปน ภิกฺขเว เอกจฺโจ ปุคฺคโล อภาวิตกาโย โหติ อภาวิตสีโล อภาวิตจิตฺโต อภาวิตปญฺโญ ปริตฺโต อปฺปาตุโม อปฺปทุกฺข\-วิหารี, เอวรูปสฺส ภิกฺขเว ปุคฺคลสฺส อปฺปมตฺตกมฺปิ ปาปกมฺมํ กตํ, ตเมนํ นิรยํ อุปเนติ.}

\prose{กถํรูปสฺส ภิกฺขเว ปุคฺคลสฺส ตาทิสํเยว อปฺปมตฺตกํ ปาปกมฺมํ กตํ ทิฏฺฐธมฺม\-เวทนียํ โหติ, นา’ณุปิ ขายติ กิํ พหุเทว? อิธ ภิกฺขเว เอกจฺโจ ปุคฺคโล ภาวิตกาโย โหติ ภาวิตสีโล ภาวิตจิตฺโต ภาวิตปญฺโญ อปริตฺโต มหตฺโต\csromansharedfootnote{5b37a458f0e83f64}{มหตฺตา (สี, สฺยา, กํ, อิ)} อปฺปมาณวิหารี. เอวรูปสฺส ภิกฺขเว ปุคฺคลสฺส ตาทิสํเยว อปฺปมตฺตกํ ปาปกมฺมํ กตํ ทิฏฺฐธมฺม\-เวทนียํ โหติ, นา’ณุปิ ขายติ กิํ พหุเทว.}

\prose{เสยฺยถาปิ ภิกฺขเว ปุริโส โลณกปลฺลํ\csromansharedfootnote{3bf62b129adeb2f7}{โลณผลํ (สี, สฺยา, กํ, อิ)} ปริตฺเต อุทกมลฺลเก\csromansharedfootnote{c1164152f1f9818c}{อุทกกปลฺลเก (ก)} ปกฺขิเปยฺย. ตํ กิํ มญฺญถ ภิกฺขเว, อปิ นุ ตํ ปริตฺตํ อุทกํ\csromansharedfootnote{3e711e57cdaa5d0b}{อุทกมลฺลเก อุทกํ (สี, สฺยา, กํ, อิ)} อมุนา โลณกปลฺเลน โลณํ อสฺส อเปยฺยนฺติ. เอวํ ภนฺเต. ตํ กิสฺส เหตุ? อทุํ หิ ภนฺเต ปริตฺตํ อุทกกปลฺลเก อุทกํ, ตํ อมุนา โลณกปลฺเลน โลณํ อสฺส อเปยฺยนฺติ. เสยฺยถาปิ ภิกฺขเว ปุริโส โลณกปลฺลกํ คงฺคาย นทิยา ปกฺขิเปยฺย. ตํ กิํ มญฺญถ ภิกฺขเว, อปิ นุ สา คงฺคา นที อมุนา โลณกปลฺเลน โลณํ อสฺส อเปยฺยนฺติ. โน เหตํ ภนฺเต. ตํ กิสฺส เหตุ? อสุ หิ ภนฺเต คงฺคาย นทิยา มหา อุทกกฺขนฺโธ, โส อมุนา โลณกปลฺเลน โลโณ น อสฺส อเปยฺโยติ\csromansharedfootnote{eb072ff10ae9d3ea}{โลณํ เนวสฺส อเปยฺยนฺติ (สี), น โลโณ อสฺส อเปยฺโยติ (อิ)}.}

\prose{\csromanfolio{252}เอวเมวํ โข ภิกฺขเว อิเธกจฺจสฺส ปุคฺคลสฺส อปฺปมตฺตกมฺปิ ปาปกมฺมํ กตํ, ตเมนํ นิรยํ อุปเนติ. อิธ ภิกฺขเว เอกจฺจสฺส ปุคฺคลสฺส ตาทิสํเยว อปฺปมตฺตกํ ปาปกมฺมํ กตํ ทิฏฺฐธมฺม\-เวทนียํ โหติ, นา’ณุปิ ขายติ กิํ พหุเทว.}

\prose{กถํรูปสฺส ภิกฺขเว ปุคฺคลสฺส อปฺปมตฺตกมฺปิ ปาปกมฺมํ กตํ, ตเมนํ นิรยํ อุปเนติ? อิธ ภิกฺขเว เอกจฺโจ ปุคฺคโล อภาวิตกาโย โหติ อภาวิตสีโล อภาวิตจิตฺโต อภาวิตปญฺโญ ปริตฺโต อปฺปาตุโม อปฺปทุกฺข\-วิหารี, เอวรูปสฺส ภิกฺขเว ปุคฺคลสฺส อปฺปมตฺตกมฺปิ ปาปกมฺมํ กตํ, ตเมนํ นิรยํ อุปเนติ.}

\prose{กถํรูปสฺส ภิกฺขเว ปุคฺคลสฺส ตาทิสํเยว อปฺปมตฺตกํ ปาปกมฺมํ กตํ ทิฏฺฐธมฺม\-เวทนียํ โหติ, นา’ณุปิ ขายติ กิํ พหุเทว? อิธ ภิกฺขเว เอกจฺโจ ปุคฺคโล ภาวิตกาโย โหติ ภาวิตสีโล ภาวิตจิตฺโต ภาวิตปญฺโญ อปริตฺโต มหตฺโต อปฺปมาณวิหารี, เอวรูปสฺส ภิกฺขเว ปุคฺคลสฺส ตาทิสํเยว อปฺปมตฺตกํ ปาปกมฺมํ กตํ ทิฏฺฐธมฺม\-เวทนียํ โหติ, นา’ณุปิ ขายติ กิํ พหุเทว. (1)}

\prose{อิธ ภิกฺขเว เอกจฺโจ อฑฺฒกหาปเณนปิ พนฺธนํ นิคจฺฉติ, กหาปเณนปิ พนฺธนํ นิคจฺฉติ, กหาปณสเตนปิ พนฺธนํ นิคจฺฉติ. อิธ ภิกฺขเว เอกจฺโจ อฑฺฒกหาปเณนปิ น พนฺธนํ นิคจฺฉติ, กหาปเณนปิ น พนฺธนํ นิคจฺฉติ, กหาปณสเตนปิ น พนฺธนํ นิคจฺฉติ.}

\prose{กถํรูโป ภิกฺขเว อฑฺฒกหาปเณนปิ พนฺธนํ นิคจฺฉติ, กหาปเณนปิ พนฺธนํ นิคจฺฉติ, กหาปณสเตนปิ พนฺธนํ นิคจฺฉติ? อิธ ภิกฺขเว เอกจฺโจ ทลิทฺโท โหติ อปฺปสฺสโก อปฺปโภโค, เอวรูโป ภิกฺขเว อฑฺฒกหาปเณนปิ พนฺธนํ นิคจฺฉติ, กหาปเณนปิ พนฺธนํ นิคจฺฉติ, กหาปณสเตนปิ พนฺธนํ นิคจฺฉติ.}

\prose{กถํรูโป ภิกฺขเว อฑฺฒกหาปเณนปิ น พนฺธนํ นิคจฺฉติ, กหาปเณนปิ น พนฺธนํ นิคจฺฉติ, กหาปณสเตนปิ น พนฺธนํ นิคจฺฉติ? อิธ ภิกฺขเว เอกจฺโจ อฑฺโฒ โหติ มหทฺธโน มหาโภโค, เอวรูโป ภิกฺขเว อฑฺฒกหาปเณนปิ น พนฺธนํ นิคจฺฉติ, กหาปเณนปิ น พนฺธนํ นิคจฺฉติ, กหาปณสเตนปิ น พนฺธนํ นิคจฺฉติ. เอวเมวํ โข ภิกฺขเว อิเธกจฺจสฺส ปุคฺคลสฺส อปฺปมตฺตกํ ปาปกมฺมํ กตํ, ตเมนํ นิรยํ อุปเนติ.}

\vspace{\tipitakaparskip}
\csromancenter{(10) 5. โลณกปลฺล\-วคฺค อิธ ภิกฺขเว เอกจฺจสฺส ปุคฺคลสฺส ตาทิสํเยว อปฺปมตฺตกํ ปาปกมฺมํ กตํ ทิฏฺฐธมฺม\-เวทนียํ โหติ, นา’ณุปิ ขายติ กิํ พหุเทว.}
\prose{\csromanfolio{253}กถํรูปสฺส ภิกฺขเว ปุคฺคลสฺส อปฺปมตฺตกํ ปาปกมฺมํ กตํ, ตเมนํ นิรยํ อุปเนติ? อิธ ภิกฺขเว เอกจฺโจ ปุคฺคโล อภาวิตกาโย โหติ อภาวิตสีโล อภาวิตจิตฺโต อภาวิตปญฺโญ ปริตฺโต อปฺปาตุโม อปฺปทุกฺข\-วิหารี, เอวรูปสฺส ภิกฺขเว ปุคฺคลสฺส ตาทิสํเยว อปฺปมตฺตกํ ปาปกมฺมํ กตํ, ตเมนํ นิรยํ อุปเนติ.}

\prose{กถํรูปสฺส ภิกฺขเว ปุคฺคลสฺส ตาทิสํเยว อปฺปมตฺตกํ ปาปกมฺมํ กตํ ทิฏฺฐธมฺม\-เวทนียํ โหติ, นา’ณุปิ ขายติ กิํ พหุเทว? อิธ ภิกฺขเว เอกจฺโจ ปุคฺคโล ภาวิตกาโย โหติ ภาวิตสีโล ภาวิตจิตฺโต ภาวิตปญฺโญ อปริตฺโต มหตฺโต อปฺปมาณวิหารี, เอวรูปสฺส ภิกฺขเว ปุคฺคลสฺส ตาทิสํเยว อปฺปมตฺตกํ ปาปกมฺมํ กตํ ทิฏฺฐธมฺม\-เวทนียํ โหติ, นา’ณุปิ ขายติ กิํ พหุเทว. (2)}

\prose{อิธ ภิกฺขเว เอกจฺโจ ปุคฺคโล ภาวิตกาโย โหติ ภาวิตสีโล ภาวิตจิตฺโต ภาวิตปญฺโญ อปริตฺโต มหตฺโต อปฺปมาณวิหารี, เอวรูปสฺส ภิกฺขเว ปุคฺคลสฺส ตาทิสํเยว อปฺปมตฺตกํ ปาปกมฺมํ กตํ ทิฏฺฐธมฺม\-เวทนียํ โหติ, นา’ณุปิ ขายติ กิํ พหุเทว. เสยฺยถาปิ ภิกฺขเว โอรพฺภิโก วา อุรพฺภฆาตโก วา อปฺเปกจฺจํ อุรพฺภํ อทินฺนํ อาทิยมานํ ปโหติ หนฺตุํ วา พนฺธิตุํ วา ชาเปตุํ วา ยถาปจฺจยํ วา กาตุํ. อปฺเปกจฺจํ อุรพฺภํ อทินฺนํ อาทิยมานํ นปฺปโหติ หนฺตุํ วา พนฺธิตุํ วา ชาเปตุํ วา ยถาปจฺจยํ วา กาตุํ.}

\prose{กถํรูปํ ภิกฺขเว โอรพฺภิโก วา อุรพฺภฆาตโก วา อุรพฺภํ อทินฺนํ อาทิยมานํ ปโหติ หนฺตุํ วา พนฺธิตุํ วา ชาเปตุํ วา ยถาปจฺจยํ วา กาตุํ? อิธ ภิกฺขเว เอกจฺโจ ทลิทฺโท โหติ อปฺปสฺสโก อปฺปโภโค, เอวรูปํ ภิกฺขเว โอรพฺภิโก วา อุรพฺภฆาตโก วา อุรพฺภํ อทินฺนํ อาทิยมานํ ปโหติ หนฺตุํ วา พนฺธิตุํ วา ชาเปตุํ วา ยถาปจฺจยํ วา กาตุํ.}

\prose{กถํรูปํ ภิกฺขเว โอรพฺภิโก วา อุรพฺภฆาตโก วา อุรพฺภํ อทินฺนํ อาทิยมานํ นปฺปโหติ หนฺตุํ วา พนฺธิตุํ วา ชาเปตุํ วา ยถาปจฺจยํ วา}

\prose{\csromanfolio{254}กาตุํ? อิธ ภิกฺขเว เอกจฺโจ อฑฺโฒ โหติ มหทฺธโน มหาโภโค ราชา วา ราชมหามตฺโต วา. เอวรูปํ ภิกฺขเว โอรพฺภิโก วา อุรพฺภฆาตโก วา อุรพฺภํ อทินฺนํ อาทิยมานํ นปฺปโหติ หนฺตุํ วา พนฺธิตุํ วา ชาเปตุํ วา ยถาปจฺจยํ วา กาตุํ, อญฺญทตฺถุ ปญฺชลิโกว\csromansharedfootnote{910c3f291596c447}{ปญฺชลิโก (ก)} นํ\csromansharedfootnote{a689abe11df2483f}{ปรํ (ก)} ยาจติ “เทหิ เม มาริส อุรพฺภํ วา อุรพฺภธนํ วา”ติ. เอวเมวํ โข ภิกฺขเว อิเธกจฺจสฺส ปุคฺคลสฺส ตาทิสํเยว อปฺปมตฺตกมฺปิ ปาปกมฺมํ กตํ, ตเมนํ นิรยํ อุปเนติ. อิธ ปน ภิกฺขเว เอกจฺจสฺส ปุคฺคลสฺส ตาทิสํเยว อปฺปมตฺตกํ ปาปกมฺมํ กตํ ทิฏฺฐธมฺม\-เวทนียํ โหติ, นา’ณุปิ ขายติ กิํ พหุเทว.}

\prose{กถํรูปสฺส ภิกฺขเว ปุคฺคลสฺส อปฺปมตฺตกมฺปิ ปาปกมฺมํ กตํ, ตเมนํ นิรยํ อุปเนติ? อิธ ภิกฺขเว เอกจฺโจ ปุคฺคโล อภาวิตกาโย โหติ อภาวิตสีโล อภาวิตจิตฺโต อภาวิตปญฺโญ ปริตฺโต อปฺปาตุโม อปฺปทุกฺข\-วิหารี, เอวรูปสฺส ภิกฺขเว ปุคฺคลสฺส อปฺปมตฺตกมฺปิ ปาปกมฺมํ กตํ, ตเมนํ นิรยํ อุปเนติ.}

\prose{กถํรูปสฺส ภิกฺขเว ปุคฺคลสฺส ตาทิสํเยว อปฺปมตฺตกํ ปาปกมฺมํ กตํ ทิฏฺฐธมฺม\-เวทนียํ โหติ, นาณุปิ ขายติ กิํ พหุเทว? อิธ ภิกฺขเว เอกจฺโจ ปุคฺคโล ภาวิตกาโย โหติ ภาวิตสีโล ภาวิตจิตฺโต ภาวิตปญฺโญ อปริตฺโต มหตฺโต อปฺปมาณวิหารี, เอวรูปสฺส ภิกฺขเว ปุคฺคลสฺส ตาทิสํเยว อปฺปมตฺตกํ ปาปกมฺมํ กตํ ทิฏฺฐธมฺม\-เวทนียํ โหติ, นา’ณุปิ ขายติ กิํ พหุเทว. (3)}

\prose{โย ภิกฺขเว เอวํ วเทยฺย “ยถา ยถายํ ปุริโส กมฺมํ กโรติ, ตถา ตถา ตํ ปฏิสํเวเทตี”ติ. เอวํ สนฺตํ ภิกฺขเว พฺรหฺมจริย\-วาโส น โหติ, โอกาโส น ปญฺญายติ สมฺมา ทุกฺขสฺส อนฺตกิริยาย. โย จ โข ภิกฺขเว เอวํ วเทยฺย “ยถา ยถา เวทนียํ อยํ ปุริโส กมฺมํ กโรติ, ตถา ตถา ตสฺส วิปากํ ปฏิสํเวเทตี”ติ. เอวํ สนฺตํ ภิกฺขเว พฺรหฺมจริย\-วาโส โหติ, โอกาโส ปญฺญายติ สมฺมา ทุกฺขสฺส อนฺตกิริยายาติ.  นวมํ.}

\chapterheadpageread{(10) 5. โลณกปลฺล\-วคฺค 10. ปํสุโธวก\-สุตฺต}
\proseitem{102}{\csromanfolio{255}สนฺติ ภิกฺขเว ชาตรูปสฺส โอฬาริกา อุปกฺกิเลสา ปํสุวาลุกา\csromansharedfootnote{a1924d3a673152e6}{ปํสุวาลิกา (สี, สฺยา, กํ, อิ)} สกฺขร\-กฐลา, ตเมนํ ปํสุโธวโก วา ปํสุโธวก\-นฺเตวาสี วา โทณิยํ อากิริตฺวา โธวติ สนฺโธวติ นิทฺโธวติ. ตสฺมิํ ปหีเน ตสฺมิํ พฺยนฺตีกเต สนฺติ ชาตรูปสฺส มชฺฌิม\-สหคตา อุปกฺกิเลสา สุขุม\-สกฺขรา ถูลวาสุกา\csromansharedfootnote{14b1cd4318ad3a66}{ถูลวาลิกา (สี, อิ), ถุลฺลวาลิกา (สฺยา, กํ) 3. สุวณฺณชาตรูปกาวสิสฺสนฺติ (ก)}, ตเมนํ ปํสุโธวโก วา ปํสุโธวก\-นฺเตวาสี วา โธวติ สนฺโธคติ นิทฺโธวติ. ตสฺมิํ ปหีเน ตสฺมิํ พฺยนฺตีกเต สนฺติ ชาตรูปสฺส สุขุม\-สหคตา อุปกฺกิเลสา สุขุมวาลุกา กาฬชลฺลิกา, ตเมนํ ปํสุโธวโก วา ปํสุโธวก\-นฺเตวาสี วา โธวติ สนฺโธวติ นิทฺโธวติ. ตสฺมิํ ปหีเน ตสฺมิํ พฺยนฺตีกเต อถาปรํ สุวณฺณสิกตาวสิสฺสนฺติ\csromansharedfootnote{b3bf15001c0e92bb}{อธนฺตํ อสนฺธนฺตํ (สฺยา, กํ) 5. อนิทฺธนฺตํ อนิหิตํ อนินฺนีตกสาวํ (สี, สฺยา, กํ, อิ) 6. นิหิตํ นินฺนีตกสาวํ (สี, สฺยา, กํ, อิ)}, ตเมนํ สุวณฺณกาโร วา สุวณฺณการ\-นฺเตวาสี วา ชาตรูปํ มูสายํ ปกฺขิปิตฺวา ธมติ สนฺธมติ นิทฺธมติ, ตํ โหติ ชาตรูปํ ธนฺตํ สนฺธนฺตํ\csromansharedfootnote{e2f44f7f8044249d}{มุทฺทิกาย (อํ 2. 14 ปิฏฺเฐ)} นิทฺธนฺตํ อนิทฺธนฺต\-กสาวํ\csromansharedfootnote{f4bc605552f11589}{คีเวยฺยเกน (ก), คีเวยฺยกาย (?)} น เจว มุทุ โหติ น จ กมฺมนิยํ น จ ปภสฺสรํ, ปภงฺคุ จ, น จ สมฺมา อุเปติ กมฺมาย. โหติ โส ภิกฺขเว สมโย, ยํ สุวณฺณกาโร วา สุวณฺณการ\-นฺเตวาสี วา ตํ ชาตรูปํ ธมติ สนฺธมติ นิทฺธมติ, ตํ โหติ ชาตรูปํ ธนฺตํ สนฺธนฺตํ นิทฺธนฺตํ นิทฺธนฺต\-กสาวํ\csromansharedfootnote{871367d3a1842408}{สุวณฺณชาตรูปกาวสิสฺสนฺติ (ก)} มุทุ จ โหติ กมฺมนิยญฺจ ปภสฺสรญฺจ, น จ ปภงฺคุ, สมฺมา อุเปติ กมฺมาย. ยสฺสา ยสฺสา จ ปิลนฺธน\-วิกติยา อากงฺขติ ยทิ ปฏฺฏิกาย\csromansharedfootnote{83ff22fed8020089}{อนิทฺธนฺตํ อนิหิตํ อนินฺนีตกสาวํ (สี, สฺยา, กํ, อิ)} ยทิ กุณฺฑลาย ยทิ คีเวยฺยเก\csromansharedfootnote{fe02ef3725fb2724}{นิหิตํ นินฺนีตกสาวํ (สี, สฺยา, กํ, อิ)} ยทิ สุวณฺณมาลาย, ตญฺจสฺส อตฺถํ อนุโภติ.}

\prose{เอวเมวํ โข ภิกฺขเว สนฺติ อธิจิตฺตม\-นุยุตฺตสฺส ภิกฺขุโน โอฬาริกา อุปกฺกิเลสา กายทุจฺจริตํ วจีทุจฺจริตํ มโนทุจฺจริตํ, ตเมนํ สเจตโส ภิกฺขุ ทพฺพชาติโก ปชหติ วิโนเทติ พฺยนฺตีกโรติ อนภาวํ คเมติ. ตสฺมิํ ปหีเน ตสฺมิํ พฺยนฺตีกเต สนฺติ อธิจิตฺตม\-นุยุตฺตสฺส ภิกฺขุโน มชฺฌิม\-สหคตา อุปกฺกิเลสา กามวิตกฺโก พฺยาปาทวิตกฺโก วิหิํสาวิตกฺโก, ตเมนํ สเจตโส ภิกฺขุ ทพฺพชาติโก ปชหติ วิโนเทติ พฺยนฺตีกโรติ อนภาวํ คเมติ. \csromanfolio{256}ตสฺมิํ ปหีเน ตสฺมิํ พฺยนฺตีกเต สนฺติ อธิจิตฺตม\-นุยุตฺตสฺส ภิกฺขุโน สุขุม\-สหคตา อุปกฺกิเลสา ญาติวิตกฺโก ชนปท\-วิตกฺโก อนวญฺญตฺติ\-ปฏิสํยุตฺโต วิตกฺโก, ตเมนํ สเจตโส ภิกฺขุ ทพฺพชาติโก ปชหติ วิโนเทติ พฺยนฺตีกโรติ อนภาวํ คเมติ. ตสฺมิํ ปหีเน ตสฺมิํ พฺยนฺตีกเต อถาปรํ ธมฺมวิตกฺกาวสิสฺสนฺติ\csromansharedfootnote{51e1c9ad6bba555d}{ธมฺมวิตกฺโกวสิสฺสติ (ก)}, โส โหติ สมาธิ น เจว สนฺโต, น จ ปณีโต, นปฺปฏิปฺปสฺสทฺธ\-ลทฺโธ, น เอโกทิภาวา\-ธิคโต, สสงฺขาร\-นิคฺคยฺห\-วาริต\-คโต\csromansharedfootnote{4a75c4ec95357bee}{สสงฺขารนิคฺคยฺหวาริตวโต (สี, สฺยา, กํ, อิ), สสงฺขารนิคฺคยฺหวาริวาวโต (ก), สสงฺขารนิคฺคยฺหวาริยาธิคโต (?) อํ 3. 225; ที 3. 234 ปิฏฺเฐสุปิ.}. โหติ โส ภิกฺขเว สมโย, ยํ ตํ จิตฺตํ อชฺฌตฺตํเยว สนฺติฏฺฐติ สนฺนิสีทติ, เอโกทิ โหติ\csromansharedfootnote{0675ec216101e8ab}{เอโกทิภาวํ คจฺฉติ (สี), เอโกทิภาโว โหติ (สฺยา, กํ, ก), เอโกทิโหติ (อิ)} สมาธิยติ, โส โหติ สมาธิ สนฺโต ปณีโต ปฏิปฺปสฺสทฺธิ\-ลทฺโธ เอโกทิภาวา\-ธิคโต, น สสงฺขาร\-นิคฺคยฺห\-วาริต\-คโต. ยสฺส ยสฺส จ อภิญฺญา สจฺฉิ\-กรณียสฺส ธมฺมสฺส จิตฺตํ อภินินฺนาเมติ อภิญฺญา สจฺฉิกิริยาย, ตตฺร ตเตฺรว สกฺขิ\-ภพฺพตํ ปาปุณาติ สติ สติอายตเน.}

\prose{โส สเจ อากงฺขติ “อเนกวิหิตํ อิทฺธิวิธํ ปจฺจนุภเวยฺยํ เอโกปิ หุตฺวา พหุธา อสฺสํ, พหุธาปิ หุตฺวา เอโก อสฺสํ, อาวิภาวํ ติโรภาวํ ติโรกุฏฺฏํ ติโรปาการํ ติโรปพฺพตํ อสชฺชมาโน คจฺเฉยฺยํ เสยฺยาถาปิ อากาเส. ปถวิยาปิ อุมฺมุชฺช\-นิมุชฺชํ กเรยฺยํ เสยฺยถาปิ อุทเก. อุทเกปิ อภิชฺชมาเน\csromansharedfootnote{351f9c41fcc0a63b}{อภิชฺชมาโน (สี, อิ, ก)} คจฺเฉยฺยํ เสยฺยถาปิ ปถวิยํ. อากาเสปิ ปลฺลงฺเกน กเมยฺยํ เสยฺยถาปิ ปกฺขี สกุโณ. อิเมปิ จนฺทิมสูริเย เอวํมหิทฺธิเก เอวํ\-มหา\-นุภาเว ปาณินา ปริมเสยฺยํ ปริมชฺเชยฺยํ, ยาว พฺรหฺมโลกาปิ กาเยน วสํ วตฺเตยฺยนฺ”ติ. ตตฺร ตเตฺรว สกฺขิ\-ภพฺพตํ ปาปุณาติ สติ สติอายตเน.}

\prose{โส สเจ อากงฺขติ “ทิพฺพาย โสตธาตุยา วิสุทฺธาย อติกฺกนฺตมานุสิกาย อุโภ สทฺเท สุเณยฺยํ ทิพฺเพ จ มานุเส จ เย ทูเร สนฺติเก จา”ติ. ตตฺร ตเตฺรว สกฺขิ\-ภพฺพตํ ปาปุณาติ สติ สติอายตเน.}

\chapterheadpageread{(10) 5. โลณกปลฺล\-วคฺค}
\prose{\csromanfolio{257}โส สเจ อากงฺขติ “ปรสตฺตานํ ปรปุคฺคลานํ เจตสา เจโต ปริจฺจ ปชาเนยฺยํ, สราคํ วา จิตฺตํ ‘สราคํ จิตฺตนฺ’ติ ปชาเนยฺยํ, วีตราคํ วา จิตฺตํ ‘วีตราคํ จิตฺตนฺ’ติ ปชาเนยฺยํ, สโทสํ วา จิตฺตํ ‘สโทสํ จิตฺตนฺ’ติ ปชาเนยฺยํ, วีตโทสํ วา จิตฺตํ ‘วีตโทสํ จิตฺตนฺ’ติ ปชาเนยฺยํ, สโมหํ วา จิตฺตํ ‘สโมหํ จิตฺตนฺ’ติ ปชาเนยฺยํ, วีตโมหํ วา จิตฺตํ ‘วีตโมหํ จิตฺตนฺ’ติ ปชาเนยฺยํ, สํขิตฺตํ วา จิตฺตํ ‘สํขิตฺตํ จิตฺตนฺ’ติ ปชาเนยฺยํ, วิกฺขิตฺตํ วา จิตฺตํ ‘วิกฺขิตฺตํ จิตฺตนฺ’ติ ปชาเนยฺยํ, มหคฺคตํ วา จิตฺตํ ‘มหคฺคตํ จิตฺตนฺ’ติ ปชาเนยฺยํ, อมหคฺคตํ วา จิตฺตํ ‘อมหคฺคตํ จิตฺตนฺ’ติ ปชาเนยฺยํ, สอุตฺตรํ วา จิตฺตํ ‘สอุตฺตรํ จิตฺตนฺ’ติ ปชาเนยฺยํ, อนุตฺตรํ วา จิตฺตํ ‘อนุตฺตรํ จิตฺตนฺ’ติ ปชาเนยฺยํ, สมาหิตํ วา จิตฺตํ ‘สมาหิต จิตฺตนฺ’ติ ปชาเนยฺยํ, อสมาหิตํ วา จิตฺตํ ‘อสมาหิตํ จิตฺตนฺ’ติ ปชาเนยฺยํ, วิมุตฺตํ วา จิตฺตํ ‘วิมุตฺตํ จิตฺตนฺ’ติ ปชาเนยฺยํ, อวิมุตฺตํ วา จิตฺตํ ‘อวิมุตฺตํ จิตฺตนฺ’ติ ปชาเนยฺยนฺ”ติ. ตตฺร ตเตฺรว สกฺขิ\-ภพฺพตํ ปาปุณาติ สติ สติอายตเน.}

\prose{โส สเจ อากงฺขติ “อเนกวิหิตํ ปุพฺเพนิวาสํ อนุสฺสเรยฺยํ. เสยฺยถิทํ, เอกมฺปิ ชาติํ เทฺวปิ ชาติโย ติสฺโสปิ ชาติโย จตสฺโสปิ ชาติโย ปญฺจปิ ชาติโย ทสปิ ชาติโย วีสมฺปิ ชาติโย ติํสมฺปิ ชาติโย จตฺตาลีสมฺปิ ชาติโย ปญฺญาสมฺปิ ชาติโย ชาติสตมฺปิ ชาติสหสฺสมฺปิ ชาติ\-สต\-สหสฺสมฺปิ อเนเกปิ สํวฏฺฏกปฺเป อเนเกปิ วิวฏฺฏกปฺเป อเนเกปิ สํวฏฺฏ\-วิวฏฺฏ\-กปฺเป ‘อมุตฺราสิํ เอวํนาโม เอวํโคตฺโต เอวํวณฺโณ เอวมาหาโร เอวํ\-สุขทุกฺข\-ปฺปฏิสํเวที เอวมา\-ยุปริยนฺโต, โส ตโต จุโต อมุตฺร อุทปาทิํ. ตตฺราปาสิํ เอวํนาโม เอวํโคตฺโต เอวํวณฺโณ เอวมาหาโร เอวํ\-สุขทุกฺข\-ปฺปฏิสํเวที เอวมา\-ยุปริยนฺโต, โส ตโต จุโต อิธูปปนฺโน’ติ. อิติ สาการํ สอุทฺเทสํ อเนกวิหิตํ ปุพฺเพนิวาสํ อนุสฺสเรยฺยนฺ”ติ. ตตฺร ตเตฺรว สกฺขิ\-ภพฺพตํ ปาปุณาติ สติ สติอายตเน.}

\prose{โส สเจ อากงฺขติ “ทิพฺเพน จกฺขุนา วิสุทฺเธน อติกฺกนฺต\-มานุสเกน สตฺเต ปสฺเสยฺยํ จวมาเน อุปปชฺชมาเน หีเน ปณีเต สุวณฺเณ ทุพฺพณฺเณ, สุคเต ทุคฺคเต ยถากมฺมูปเค สตฺเต ปชาเนยฺยํ ‘อิเม วต โภนฺโต สตฺตา กายทุจฺจริเยน สมนฺนาคตา วจีทุจฺจริเตน สมนฺนาคตา มโน\-ทุจฺจริเตน สมนฺนาคตา อริยานํ อุปวาทกา มิจฺฉาทิฏฺฐิกา มิจฺฉาทิฏฺฐิติกนิปาตปาฬิ \csromanfolio{258}กมฺมสมาทานา, เต กายสฺส เภทา ปรํ มรณา อปายํ ทุคฺคติํ วินิปาตํ นิรยํ อุปปนฺนา. อิเม วา ปน โภนฺโต สตฺตา กายสุจริเตน สมนฺนาคตา วจีสุจริเตน สมนฺนาคตา มโนสุจริเตน สมนฺนาคตา อริยานํ อนุปวาทกา สมฺมาทิฏฺฐิกา สมฺมา\-ทิฏฺฐิ\-กมฺม\-สมาทานา, เต กายสฺส เภทา ปรํ มรณา สุคติํ สคฺคํ โลกํ อุปปนฺนา’ติ. อิติ ทิพฺเพน จกฺขุนา วิสุทฺเธน อติกฺกนฺต\-มานุสเกน สตฺเต ปสฺเสยฺยํ จวมาเน อุปปชฺชมาเน หีเน ปณีเต สุวณฺเณ ทุพฺพณฺเณ, สุคเต ทุคฺคเต ยถากมฺมูปเค สตฺเต ปชาเนยฺยนฺ”ติ. ตตฺร ตเตฺรว สกฺขิ\-ภพฺพตํ ปาปุณาติ สติ สติอายตเน.}

\prose{โส สเจ อากงฺขติ “อาสวานํ ขยา อนาสวํ เจโตวิมุตฺติํ ปญฺญาวิมุตฺติํ ทิฏฺเฐว ธมฺเม สยํ อภิญฺญา สจฺฉิกตฺวา อุปสมฺปชฺช วิหเรยฺยนฺ”ติ. ตตฺร ตเตฺรว สกฺขิ\-ภพฺพตํ ปาปุณาติ สติ สติอายตเนติ.  ทสมํ.}

\csromanheader{2}{11. \textbf{นิมิตฺตสุตฺต}}
\proseitem{103}{อธิจิตฺตม\-นุยุตฺเตน ภิกฺขเว ภิกฺขุนา ตีณิ นิมิตฺตานิ กาเลน กาลํ มนสิ กาตพฺพานิ. กาเลน กาลํ สมาธิ\-นิมิตฺตํ มนสิ กาตพฺพํ, กาเลน กาลํ ปคฺคห\-นิมิตฺตํ มนสิ กาตพฺพํ, กาเลน กาลํ อุเปกฺขา\-นิมิตฺตํ มนสิ กาตพฺพํ. สเจ ภิกฺขเว อธิจิตฺตม\-นุยุตฺโต ภิกฺขุ เอกนฺตํ สมาธิ\-นิมิตฺตํเยว มนสิ กเรยฺย, ฐานํ ตํ จิตฺตํ โกสชฺชาย สํวตฺเตยฺย. สเจ ภิกฺขเว อธิจิตฺตม\-นุยุตฺโต ภิกฺขุ เอกนฺตํ ปคฺคห\-นิมิตฺตํเยว มนสิ กเรยฺย, ฐานํ ตํ จิตฺตํ อุทฺธจฺจาย สํวตฺเตยฺย. สเจ ภิกฺขเว อธิจิตฺตม\-นุยุตฺโต ภิกฺขุ เอกนฺตํ อุเปกฺขา\-นิมิตฺตํเยว มนสิ กเรยฺย, ฐานํ ตํ จิตฺตํ น สมฺมา สมาธิเยยฺย อาสวานํ ขยาย. ยโต จ โข ภิกฺขเว อธิจิตฺตม\-นุยุตฺโต ภิกฺขุ กาเลน กาลํ สมาธิ\-นิมิตฺตํ มนสิ กโรติ, กาเลน กาลํ ปคฺคห\-นิมิตฺตํ มนสิ กโรติ, กาเลน กาลํ อุเปกฺขา\-นิมิตฺตํ มนสิ กโรติ. ตํ โหติ จิตฺตํ มุทุญฺจ กมฺมนิยญฺจ ปภสฺสรญฺจ, น จ ปภงฺคุ, สมฺมา สมาธิยติ อาสวานํ ขยาย.}

\prose{เสยฺยถาปิ ภิกฺขเว สุวณฺณกาโร วา สุวณฺณการ\-นฺเตวาสี วา อุกฺกํ พนฺเธยฺย\csromansharedfootnote{6d9f727562eef049}{พนฺธติ ... อาลิมฺปติ (วิสุทฺธิ 1. 239 ปิฏฺเฐ ตํฏีกายํ จ) มฏฺฐ 1. 179; ม 3. 286 ปิฏฺเฐสุ ตํ อฏฺฐกถาฏีกาสุ จ ปสฺสิตพฺพํ.}, อุกฺกํ พนฺธิตฺวา อุกฺกามุขํ อาลิมฺเปยฺย\csromansharedfootnote{6d9f727562eef049}{พนฺธติ ... อาลิมฺปติ (วิสุทฺธิ 1. 239 ปิฏฺเฐ ตํฏีกายํ จ) มฏฺฐ 1. 179; ม 3. 286 ปิฏฺเฐสุ ตํ อฏฺฐกถาฏีกาสุ จ ปสฺสิตพฺพํ.}, อุกฺกามุขํ อาลิมฺเปตฺวา}

\vspace{\tipitakaparskip}
\csromancenter{(10) 5. โลณกปลฺล\-วคฺค สณฺฑาเสน ชาตรูปํ คเหตฺวา อุกฺกามุเข ปกฺขิเปยฺย\csromansharedfootnote{3046e1bcd841fb89}{ปกฺขิปติ (วิสุทฺธิ 1. 239 ปิฏฺเฐ.)}, อุกฺกามุเข ปกฺขิปิตฺวา กาเลน กาลํ อภิธมติ, กาเลน กาลํ อุทเกน ปริปฺโผเสติ, กาเลน กาลํ อชฺฌุเปกฺขติ. สเจ ภิกฺขเว สุวณฺณกาโร วา สุวณฺณการ\-นฺเตวาสี วา ตํ ชาตรูปํ เอกนฺตํ อภิธเมยฺย, ฐานํ ตํ ชาตรูปํ ฑเหยฺย. สเจ ภิกฺขเว สุวณฺณกาโร วา สุวณฺณการ\-นฺเตวาสี วา ตํ ชาตรูปํ เอกนฺตํ อุทเกน ปริปฺโผเสยฺย, ฐานํ ตํ ชาตรูปํ นิพฺพาเปยฺย\csromansharedfootnote{8649889d50947f97}{นิพฺพาเยยฺย (สี)}. สเจ ภิกฺขเว สุวณฺณกาโร วา สุวณฺณการ\-นฺเตวาสี วา ตํ ชาตรูปํ เอกนฺตํ อชฺฌุเปกฺเขยฺย, ฐานํ ตํ ชาตรูปํ น สมฺมา ปริปากํ คจฺเฉยฺย. ยโต จ โข ภิกฺขเว สุวณฺณกาโร วา สุวณฺณการ\-นฺเตวาสี วา ตํ ชาตรูปํ กาเลน กาลํ อภิธมติ, กาเลน กาลํ อุทเกน ปริปฺโผเสติ, กาเลน กาลํ อชฺฌุเปกฺขติ, ตํ โหติ ชาตรูปํ มุทุญฺจ กมฺมนิยญฺจ ปภสฺสรญฺจ, น จ ปภงฺคุ, สมฺมา อุเปติ กมฺมาย. ยสฺสา ยสฺสา จ ปิลนฺธน\-วิกติยา อากงฺขติ ยทิ ปฏฺฏิกาย ยทิ กุณฺฑลาย ยทิ คีเวยฺยเก ยทิ สุวณฺณมาลาย, ตํ จสฺส อตฺถํ อนุโภติ.}
\prose{\csromanfolio{259}เอวเมวํ โข ภิกฺขเว อธิจิตฺตม\-นุยุตฺเตน ภิกฺขุนา ตีณิ นิมิตฺตานิ กาเลน กาลํ มนสิ กาตพฺพานิ, กาเลน กาลํ สมาธิ\-นิมิตฺตํ มนสิ กาตพฺพํ, กาเลน กาลํ ปคฺคห\-นิมิตฺตํ มนสิ กาตพฺพํ, กาเลน กาลํ อุเปกฺขา\-นิมิตฺตํ มนสิ กาตพฺพํ. สเจ ภิกฺขเว อธิจิตฺตม\-นุยุตฺโต ภิกฺขุ เอกนฺตํ สมาธิ\-นิมิตฺตํเยว มนสิ กเรยฺย, ฐานํ ตํ จิตฺตํ โกสชฺชาย สํวตฺเตยฺย. สเจ ภิกฺขเว อธิจิตฺตม\-นุยุตฺโต ภิกฺขุ เอกนฺตํ ปคฺคห\-นิมิตฺตํ เยว มนสิ กเรยฺย, ฐานํ ตํ จิตฺตํ อุทฺธจฺจาย สํวตฺเตยฺย. สเจ ภิกฺขเว อธิจิตฺตม\-นุยุตฺโต ภิกฺขุ เอกนฺตํ อุเปกฺขา\-นิมิตฺตํเยว มนสิ กเรยฺย, ฐานํ ตํ จิตฺตํ น สมฺมา สมาธิเยยฺย อาสวานํ ขยาย. ยโต จ โข ภิกฺขเว อธิจิตฺตม\-นุยุตฺโต ภิกฺขุ กาเลน กาลํ สมาธิ\-นิมิตฺตํ มนสิ กโรติ, กาเลน กาลํ ปคฺคห\-นิมิตฺตํ มนสิ กโรติ, กาเลน กาลํ อุเปกฺขา\-นิมิตฺตํ มนสิ กโรติ, ตํ โหติ จิตฺตํ มุทุญฺจ กมฺมนิยญฺจ ปภสฺสรญฺจ, น จ ปภงฺคุ, สมฺมา สมาธิยติ อาสวานํ ขยาย. ยสฺส ยสฺส จ อภิญฺญา สจฺฉิ\-กรณียสฺส ธมฺมสฺส จิตฺตํ อภินินฺนาเมติ อภิญฺญา สจฺฉิกิริยาย, ตตฺร ตเตฺรว สกฺขิ\-ภพฺพตํ ปาปุณาติ สติ สติอายตเน.}

\prosewithclosermidruled{\csromanfolio{260}โส สเจ อากงฺขติ “อเนกวิหิตํ อิทฺธิวิธํ ปจฺจนุภเวยฺยํ ฯเปฯ (ฉ อภิญฺญา วิตฺถาเรตพฺพา) อาสวานํ ขยา ฯเปฯ สจฺฉิกตฺวา อุปสมฺปชฺช วิหเรยฺยนฺ”ติ. ตตฺร ตเตฺรว สกฺขิ\-ภพฺพตํ ปาปุณาติ สติ สติอายตเนติ.  เอกาทสมํ.}{โลณกปลฺล\-วคฺโค\csromansharedfootnote{338dd9d64fd7ee1e}{[มิสฺสินฺคฺ โนเต 0]} ปญฺจโม.}
\csromancenter{\textbf{ตสฺสุทฺทานํ}}
\setlength{\csromangathapagewidth}{0pt}
\setlength{\csromangathawakwidth}{0pt}
\csromangathameasuregroup{อจฺจายิกํ ปวิเวกํ, \\ อาชานียา โปตฺถโก จ,}{\csromangathabat{อจฺจายิกํ ปวิเวกํ,}{สรโท ปริสา ตโย.} \\ \csromangathabat{อาชานียา โปตฺถโก จ,}{โลณํ โธวติ นิมิตฺตานีติ.} \\ \textbf{ทุติโย ปณฺณาสโก สมตฺโต.}}
\csromangathasetleft{อจฺจายิกํ ปวิเวกํ, \\ อาชานียา โปตฺถโก จ,}
\csromangathagroup{\csromangathastanza{\csromangathabat{อจฺจายิกํ ปวิเวกํ,}{สรโท ปริสา ตโย.} \\ \csromangathabat{อาชานียา โปตฺถโก จ,}{โลณํ โธวติ นิมิตฺตานีติ.} \\ \textbf{ทุติโย ปณฺณาสโก สมตฺโต.}}}

\csromanmarkh{(11)}\csromanmarkhii{1. สมฺโพธวคฺค}\csromanheadernomark{1}{\textbf{(11) 1. สมฺโพธวคฺค}}
\csromanmatikaanchor{mtk.3}
\csromantocmarkref{subsection}{1. ปุพฺเพวสมฺโพธสุตฺต}{mtk.3}
\bookmark[dest={mtk.3},level=2]{1. ปุพฺเพวสมฺโพธสุตฺต}
\csromanheader{2}{1. ปุพฺเพวสมฺโพธ\-สุตฺต}
\csromanmatikaanchor{mtk.4}
\csromantocmarkref{subsection}{2-3. อสฺสาทสุตฺต}{mtk.4}
\bookmark[dest={mtk.4},level=2]{2-3. อสฺสาทสุตฺต}
\proseitem{104}{\csromanfolio{261}ปุพฺเพว เม ภิกฺขเว สมฺโพธา อนภิสมฺพุทฺธสฺส โพธิสตฺตสฺเสว สโต เอตทโหสิ “โก นุ โข โลเก อสฺสาโท, โก อาทีนโว, กิํ นิสฺสรณนฺ”ติ. ตสฺส มยฺหํ ภิกฺขเว เอตทโหสิ “ยํ โข โลกํ\csromansharedfootnote{20756ecb477c5c82}{โลเก (สี, สฺยา, กํ, อิ)} ปฏิจฺจ อุปฺปชฺชติ สุขํ โสมนสฺสํ, อยํ โลเก อสฺสาโท. ยํ โลโก\csromansharedfootnote{f150147308c05474}{โลเก (อิ, ก)} อนิจฺโจ ทุกฺโข วิปริณาม\-ธมฺโม, อยํ โลเก อาทีนโว. โย โลเก ฉนฺทราค\-วินโย ฉนฺท\-ราค\-ปฺปหานํ, อิทํ โลเก นิสฺสรณนฺ”ติ\csromansharedfootnote{2cf7278b2861cf73}{โลกนิสฺสรณํ (ฏฺฐ). 262 ปิฏฺเฐ “โลเก นิสฺสรณนฺ”ติ ปเทน สํสนฺทิตพฺพํ.}. ยาวกีวญฺจาหํ ภิกฺขเว เอวํ โลกสฺส อสฺสาทญฺจ อสฺสาทโต, อาทีนวญฺจ อาทีนวโต, นิสฺสรณญฺจ นิสฺสรณโต ยถาภูตํ นาพฺภญฺญาสิํ, เนว ตาวาหํ ภิกฺขเว สเทวเก โลเก สมารเก สพฺรหฺมเก สสฺสมณ\-พฺราหฺมณิยา ปชาย สเทว\-มนุสฺสาย “อนุตฺตรํ สมฺมาสมฺโพธิํ อภิสมฺพุทฺโธ”ติ\csromansharedfootnote{cd8f4e72603be4e7}{อภิสมฺพุทฺโธ (สี, สฺยา, กํ, ก)} ปจฺจญฺญาสิํ. ยโต จ ขฺวาหํ\csromansharedfootnote{1e97f87ae701e165}{โข อหํ (สี, อิ), โขหํ (สฺยา, กํ, ก)} ภิกฺขเว เอวํ โลกสฺส อสฺสาทญฺจ อสฺสาทโต, อาทีนวญฺจ อาทีนวโต, นิสฺสรณญฺจ นิสฺสรณโต ยถาภูตํ อพฺภญฺญาสิํ, อถาหํ ภิกฺขเว สเทวเก โลเก สมารเก สพฺรหฺมเก สสฺสมณ\-พฺราหฺมณิยา ปชาย สเทว\-มนุสฺสาย “อนุตฺตรํ สมฺมาสมฺโพธิํ อภิสมฺพุทฺโธ”ติ ปจฺจญฺญาสิํ. ญาณญฺจ ปน เม ทสฺสนํ อุทปาทิ “อกุปฺปา เม วิมุตฺติ\csromansharedfootnote{4a253074c6d5b4b9}{เจโตวิมุตฺติ (สี, อิ, ก)}, อยมนฺติมา ชาติ, นตฺถิ ทานิ ปุนพฺภโว”ติ.  ปฐมํ.}

\prose{2. ปฐม\-อสฺสาท\-สุตฺต}

\proseitem{105}{โลกสฺสาหํ ภิกฺขเว อสฺสาท\-ปริเยสนํ อจริํ, โย โลเก อสฺสาโท ตทชฺฌคมํ. ยาวตโก โลเก อสฺสาโท, ปญฺญาย เม โส สุทิฏฺโฐ. โลกสฺสาหํ ภิกฺขเว อาทีนว\-ปริเยสนํ \csromanfolio{262}อจริํ, โย โลเก อาทีนโว ตทชฺฌคมํ. ยาวตโก โลเก อาทีนโว, ปญฺญาย เม โส สุทิฏฺโฐ. โลกสฺสาหํ ภิกฺขเว นิสฺสรณ\-ปริเยสนํ อจริํ, ยํ โลเก นิสฺสรณํ ตทชฺฌคมํ. ยาวตกํ โลเก นิสฺสรณํ, ปญฺญาย เม ตํ สุทิฏฺฐํ. ยาวกีวญฺจาหํ ภิกฺขเว โลกสฺส อสฺสาทญฺจ อสฺสาทโต, อาทีนวญฺจ อาทีนวโต, นิสฺสรณญฺจ นิสฺสรณโต ยถาภูตํ นาพฺภญฺญาสิํ, เนว ตาวาหํ ภิกฺขเว สเทวเก โลเก สมารเก สพฺรหฺมเก สสฺสมณพฺรหฺมณิยา ปชาย สเทว\-มนุสฺสาย “อนุตฺตรํ สมฺมาสมฺโพธิํ อภิสมฺพุทฺโธ”ติ ปจฺจญฺญาสิํ. ยโต จ ขฺวาหํ ภิกฺขเว โลกสฺส อสฺสาทญฺจ อสฺสาทโต, อาทีนวญฺจ อาทีนวโต, นิสฺสรณญฺจ นิสฺสรณโก ยถาภูตํ อพฺภญฺญาสิํ, อถาหํ ภิกฺขเว สเทวเก โลเก สมารเก สพฺรหฺมเก สสฺสมณ\-พฺราหฺมณิยา ปชาย สเทว\-มนุสฺสาย “อนุตฺตรํ สมฺมาสมฺโพธิํ อภิสมฺพุทฺโธ”ติ ปจฺจญฺญาสิํ. ญาณญฺจ ปน เม ทสฺสนํ อุทปาทิ “อกุปฺปา เม วิมุตฺติ, อยมนฺติมา ชาติ, นตฺถิ ทานิ ปุนพฺภโว”ติ. ทุติยํ.}

\prose{3. ทุติย\-อสฺสาท\-สุตฺต}

\proseitem{106}{โน เจทํ\csromansharedfootnote{b266845d257b25b5}{โน เจตํ (สฺยา, กํ, อิ, ก) สํ 2. 25 ปิฏฺเฐ ปสฺสิตพฺพํ.} ภิกฺขเว โลเก อสฺสาโท อภวิสฺส, นยิทํ สตฺตา โลเก สารชฺเชยฺยุํ. ยสฺมา จ โข ภิกฺขเว อตฺถิ โลเก อสฺสาโท, ตสฺมา สตฺตา โลเก สารชฺชนฺติ. โน เจทํ ภิกฺขเว โลเก อาทีนโว อภวิสฺส, นยิทํ สตฺตา โลเก นิพฺพินฺเทยฺยุํ. ยสฺมา จ โข ภิกฺขเว อตฺถิ โลเก อาทีนโว, ตสฺมา สตฺตา โลเก นิพฺพินฺทนฺติ. โน เจทํ ภิกฺขเว โลเก นิสฺสรณํ อภวิสฺส นยิทํ สตฺตา โลกมฺหา\csromansharedfootnote{5eb1e64df7115202}{โลเก (ก)} นิสฺสเรยฺยุํ. ยสฺมา จ โข ภิกฺขเว อตฺถิ โลเก นิสฺสรณํ, ตสฺมา สตฺตา โลกมฺหา นิสฺสรนฺติ. ยาวกีวญฺจ ภิกฺขเว สตฺตา โลกสฺส อสฺสาทญฺจ อสฺสาทโต อาทีนวญฺจ อาทีนวโต นิสฺสรณญฺจ นิสฺสรณโต ยถาภูตํ นาพฺภญฺญาสุํ\csromansharedfootnote{10f3d9bc8c6923ac}{นาพฺภญฺญํสุ (สํ 2. 26 ปิฏฺเฐ)}. เนว ตาว ภิกฺขเว สตฺตา สเทวก โลกา สมารกา สพฺรหฺมกา สสฺสมณ\-พฺราหฺมณิยา ปชาย สเทว\-มนุสฺสาย นิสฺสฏา วิสํยุตฺตา วิปฺปมุตฺตา\csromansharedfootnote{c0b209dc654c962c}{วิปฺปยุตฺตา (ก)} วิมริยาทีกเตน\csromansharedfootnote{9635861967e0dede}{วิมริยาทิกเตน (สี, อิ, ก)} เจตสา วิหริํสุ. ยโต จ โข ภิกฺขเว สตฺตา}

\chapterheadpageread{(11) 1. สมฺโพธวคฺค}
\prose{\csromanfolio{263}โลกสฺส อสฺสาทญฺจ อสฺสาทโต อาทีนวญฺจ อาทีนวโต นิสฺสรณญฺจ นิสฺสรณโต ยถาภูตํ อพฺภาญฺญาสุํ, อถ ภิกฺขเว สตฺตา สเทวกา โลกา สมารกา สพฺรหฺมกา สสฺสมณ\-พฺราหฺมณิยา ปชาย สเทว\-มนุสฺสาย นิสฺสฏา วิสํยุตฺตา วิปฺปมุตฺตา วิมริยาทีกเตน เจตสา วิหรนฺตีติ.  ตติยํ.}

\csromanmatikaanchor{mtk.5}
\csromantocmarkref{subsection}{4. สมณพฺราหฺมณสุตฺต}{mtk.5}
\bookmark[dest={mtk.5},level=2]{4. สมณพฺราหฺมณสุตฺต}
\csromanheader{2}{4. สมณพฺราหฺมณ\-สุตฺต}
\proseitem{107}{เย เกจิ ภิกฺขเว สมณา วา พฺราหฺมณา วา โลกสฺส อสฺสาทญฺจ อสฺสาทโต, อาทีนวญฺจ อาทีนวโต, นิสฺสรณญฺจ นิสฺสรณโต ยถาภูตํ นปฺปชานนฺติ, น เม เต\csromansharedfootnote{2f5238ac63aeb022}{น เต (ก)} ภิกฺขเว สมณา วา พฺราหฺมณา วา สมเณสุ วา สมณสมฺมตา พฺราหฺมเณสุ วา พฺราหฺมณ\-สมฺมตา, น จ ปน เต อายสฺมนฺโต สามญฺญตฺถํ วา พฺรหฺมญฺญตฺถํ วา ทิฏฺเฐว ธมฺเม สยํ อภิญฺญา สจฺฉิกตฺวา อุปสมฺปชฺช วิหรนฺติ. เย จ โข เกจิ ภิกฺขเว สมณา วา พฺราหฺมณา วา โลกสฺส อสฺสาทญฺจ อสฺสาทโต, อาทีนวญฺจ อาทีนวโต, นิสฺสรณญฺจ นิสฺสรณโต ยถาภูตํ ปชานนฺติ, เต โข ภิกฺขเว สมณา วา พฺราหฺมณา วา สมเณสุ วา สมณสมฺมตา พฺราหฺมเณสุ วา พฺราหฺมณ\-สมฺมตา, เต จ ปนายสฺมนฺโต สามญฺญตฺถญฺจ พฺรหฺมญฺญตฺถญฺจ ทิฏฺเฐว ธมฺเม สยํ อภิญฺญา สจฺฉิกตฺวา อุปสมฺปชฺช วิหรนฺตีติ\csromansharedfootnote{c455fefcf0067874}{วิหริสฺสนฺติ (สี, อิ)}.  จตุตฺถํ.}

\csromanmatikaanchor{mtk.6}
\csromantocmarkref{subsection}{5. รุณฺณสุตฺต}{mtk.6}
\bookmark[dest={mtk.6},level=2]{5. รุณฺณสุตฺต}
\csromanheader{2}{5. \textbf{รุณฺณสุตฺต}}
\proseitem{108}{รุณฺณมิทํ ภิกฺขเว อริยสฺส วินเย ยทิทํ คีตํ, อุมฺมตฺตกมิทํ ภิกฺขเว อริยสฺส วินเย ยทิทํ นจฺจํ, โกมารกมิทํ ภิกฺขเว อริยสฺส วินเย ยทิทํ อติเวลํ ทนฺต\-วิทํสก\-หสิตํ\csromansharedfootnote{3b64b1282def982b}{ทนฺตวิทํสกํ หสิตํ (สี, อิ)}. ตสฺมาติห ภิกฺขเว เสตุฆาโต คีเต เสตุฆาโต นจฺเจ, อลํ โว ธมฺม\-ปฺปโมทิตานํ สตํ สิตํ สิตมตฺตายาติ.  ปญฺจมํ.}

\csromanmatikaanchor{mtk.7}
\csromantocmarkref{subsection}{6. อติตฺติสุตฺต}{mtk.7}
\bookmark[dest={mtk.7},level=2]{6. อติตฺติสุตฺต}
\csromanheader{2}{6. \textbf{อติตฺติสุตฺต}}
\proseitem{109}{ติณฺณํ ภิกฺขเว ปฏิเสวนาย นตฺถิ ติตฺติ. กตเมสํ ติณฺณํ? โสปฺปสฺส ภิกฺขเว ปฏิเสวนาย นตฺถิ ติตฺติ, สุรา\-เมรย\-ปานสฺส ภิกฺขเว \csromanfolio{264}ปฏิเสวนาย นตฺถิ ติตฺติ, เมถุนธมฺม\-สมาปตฺติยา ภิกฺขเว ปฏิเสวนาย นตฺถิ ติตฺติ. อิเมสํ ภิกฺขเว ติณฺณํ ปฏิเสวนาย นตฺถิ ติตฺตีติ.  ฉฏฺฐํ.}

\csromanmatikaanchor{mtk.8}
\csromantocmarkref{subsection}{7. อรกฺขิตสุตฺต}{mtk.8}
\bookmark[dest={mtk.8},level=2]{7. อรกฺขิตสุตฺต}
\csromanheader{2}{7. \textbf{อรกฺขิตสุตฺต}}
\proseitem{110}{อถ โข อนาถปิณฺฑิโก คหปติ เยน ภควา เตนุปสงฺกมิ, อุปสงฺกมิตฺวา ภควนฺตํ อภิวาเทตฺวา เอกมนฺตํ นิสีทิ, เอกมนฺตํ นิสินฺนํ โข อนาถปิณฺฑิกํ คหปติํ ภควา เอตทโวจ “จิตฺเต คหปติ อรกฺขิเต กายกมฺมมฺปิ อรกฺขิตํ โหติ, วจีกมฺมมฺปิ อรกฺขิตํ โหติ, มโนกมฺมมฺปิ อรกฺขิตํ โหติ. ตสฺส อรกฺขิต\-กายกมฺมนฺตสฺส อรกฺขิต\-วจีกมฺมนฺตสฺส อรกฺขิต\-มโนกมฺมนฺตสฺส กายกมฺมมฺปิ อวสฺสุตํ โหติ, วจีกมฺมมฺปิ อวสฺสุตํ โหติ, มโนกมฺมมฺปิ อวสฺสุตํ โหติ. ตสฺส อวสฺสุต\-กายกมฺมนฺตสฺส อวสฺสุต\-วจีกมฺมนฺตสฺส อวสฺสุต\-มโนกมฺมนฺตสฺส กายกมฺมมฺปิ ปูติกํ โหติ, วจีกมฺมมฺปิ ปูติกํ โหติ, มโนกมฺมมฺปิ ปูติกํ โหติ. ตสฺส ปูติกายกมฺมนฺตสฺส ปูติวจีกมฺมนฺตสฺส ปูติมโนกมฺมนฺตสฺส น ภทฺทกํ มรณํ โหติ น ภทฺทิกา กาลกิริยา.}

\prose{เสยฺยถาปิ คหปติ กูฏาคาเร ทุจฺฉนฺเน กูฏมฺปิ อรกฺขิตํ โหติ, โคปานสิโยปิ อรกฺขิตา โหนฺติ, ภิตฺติปิ อรกฺขิตา โหติ, กูฏมฺปิ อวสฺสุตํ โหติ, โคปานสิโยปิ อวสฺสุตา โหนฺติ, ภิตฺติปิ อวสฺสุตา โหติ. กูฏมฺปิ ปูติกํ โหติ, โคปานสิโยปิ ปูติกา โหนฺติ, ภิตฺติปิ ปูติกา โหติ.}

\prose{เอวเมวํ โข คหปติ จิตฺเต อรกฺขิเต กายกมฺมมฺปิ อรกฺขิตํ โหติ, วจีกมฺมมฺปิ อรกฺขิตํ โหติ, มโนกมฺมมฺปิ อรกฺขิตํ โหติ. ตสฺส อรกฺขิต\-กายกมฺมนฺตสฺส อรกฺขิต\-วจีกมฺมนฺตสฺส อรกฺขิต\-มโนกมฺมนฺตสฺส กายกมฺมมฺปิ อวสฺสุตํ โหติ, วจีกมฺมมฺปิ อวสฺสุตํ โหติ, มโนกมฺมมฺปิ อวสฺสุตํ โหติ. ตสฺส อวสฺสุต\-กายกมฺมนฺตสฺส อวสฺสุต\-วจีกมฺมนฺตสฺส อวสฺสุต\-มโนกมฺมนฺตสฺส กายกมฺมมฺปิ ปูติกํ โหติ, วจีกมฺมมฺปิ ปูติกํ โหติ, มโนกมฺมมฺปิ ปูติกํ โหติ. ตสฺส ปูติกายกมฺมนฺตสฺส ปูติวจีกมฺมนฺตสฺส ปูติมโนกมฺมนฺตสฺส น ภทฺทกํ มรณํ โหติ น ภทฺทิกา กาลงฺกิริยา.}

\chapterheadpageread{(11) 1. สมฺโพธวคฺค}
\prose{\csromanfolio{265}จิตฺเต คหปติ รกฺขิเต กายกมฺมมฺปิ รกฺขิตํ โหติ, วจีกมฺมมฺปิ รกฺขิตํ โหติ, มโนกมฺมมฺปิ รกฺขิตํ โหติ. ตสฺส รกฺขิต\-กายกมฺมนฺตสฺส รกฺขิต\-วจีกมฺมนฺตสฺส รกฺขิต\-มโนกมฺมนฺตสฺส กายกมฺมมฺปิ อนวสฺสุตํ โหติ, วจีกมฺมมฺปิ อนวสฺสุตํ โหติ, มโนกมฺมมฺปิ อนวสฺสุตํ โหติ. ตสฺส อนวสฺสุต\-กายกมฺมนฺตสฺส อนวสฺสุต\-วจีกมฺมนฺตสฺส อนวสฺสุต\-มโนกมฺมนฺตสฺส กายกมฺมมฺปิ อปูติกํ โหติ, วจีกมฺมมฺปิ อปูติกํ โหติ, มโนกมฺมมฺปิ อปูติกํ โหติ. ตสฺส อปูติ\-กายกมฺมนฺตสฺส อปูติ\-วจีกมฺมนฺตสฺส อปูติ\-มโนกมฺมนฺตสฺส ภทฺทกํ มรณํ โหติ ภทฺทิกา กาลกิริยา.}

\prose{เสยฺยถาปิ คหปติ กูฏาคาเร สุจฺฉนฺเน กูฏมฺปิ รกฺขิตํ โหติ, โคปานสิโยปิ รกฺขิตา โหนฺติ, ภิตฺติปิ รกฺขิตา โหติ. กูฏมฺปิ อนวสฺสุตํ โหติ, โคปานสิโยปิ อนวสฺสุตา โหนฺติ, ภิตฺติปิ อนวสฺสุตา โหติ. กูฏมฺปิ อปูติกํ โหติ, โคปานสิโยปิ อปูติกา โหนฺติ, ภิตฺติปิ อปูติกา โหติ.}

\prose{เอวเมวํ โข คหปติ จิตฺเต รกฺขิเต กายกมฺมมฺปิ รกฺขิตํ โหติ, วจีกมฺมมฺปิ รกฺขิตํ โหติ, มโนกมฺมมฺปิ รกฺขิตํ โหติ. ตสฺส รกฺขิต\-กายกมฺมนฺตสฺส รกฺขิต\-วจีกมฺมนฺตสฺส รกฺขิต\-มโนกมฺมนฺตสฺส กายกมฺมมฺปิ อนวสฺสุตํ โหติ, วจีกมฺมมฺปิ อนวสฺสุตํ โหติ, มโนกมฺมมฺปิ อนวสฺสุตํ โหติ. ตสฺส อนวสฺสุต\-กายกมฺมนฺตสฺส อนวสฺสุต\-วจีกมฺมนฺตสฺส อนวสฺสุต\-มโนกมฺมนฺตสฺส กายกมฺมมฺปิ อปูติกํ โหติ, วจีกมฺมมฺปิ อปูติกํ โหติ, มโนกมฺมมฺปิ อปูติกํ โหติ. ตสฺส อปูติ\-กายกมฺมนฺตสฺส อปูติ\-วจีกมฺมนฺตสฺส อปูติ\-มโนกมฺมนฺตสฺส ภทฺทิกํ มรณํ โหติ ภทฺทิกา กาลกิริยาติ.  สตฺตมํ.}

\csromanmatikaanchor{mtk.9}
\csromantocmarkref{subsection}{8. พฺยาปนฺนสุตฺต}{mtk.9}
\bookmark[dest={mtk.9},level=2]{8. พฺยาปนฺนสุตฺต}
\csromanheader{2}{8. \textbf{พฺยาปนฺนสุตฺต}}
\csromanmatikaanchor{mtk.10}
\csromantocmarkref{subsection}{9-10. นิทานสุตฺตานิ}{mtk.10}
\bookmark[dest={mtk.10},level=2]{9-10. นิทานสุตฺตานิ}
\proseitem{111}{เอกมนฺตํ นิสินฺนํ โข อนาถปิณฺฑิกํ คหปติํ ภควา เอตทโวจ “จิตฺเต คหปติ พฺยาปนฺเน กายกมฺมมฺปิ พฺยาปนฺนํ โหติ, วจีกมฺมมฺปิ พฺยาปนฺนํ โหติ, มโนกมฺมมฺปิ พฺยาปนฺนํ โหติ. ตสฺส พฺยาปนฺน\-กายกมฺมนฺตสฺส พฺยาปนฺน\-วจีกมฺมนฺตสฺส พฺยาปนฺน\-มโนกมฺมนฺตสฺส น ภทฺทกํ มรณํ โหติ น ภทฺทิกา กาลกิริยา. เสยฺยถาปิ คหปติ กูฏาคาเร ทุจฺฉนฺเน กูฏมฺปิ พฺยาปนฺนํ โหติ, โคปานสิโยปิ พฺยาปนฺนา โหนฺติ, ภิตฺติปิ \csromanfolio{266}พฺยาปนฺนา โหติ. เอวเมวํ โข คหปติ จิตฺเต พฺยาปนฺเน กายกมฺมมฺปิ พฺยาปนฺนํ โหติ, วจีกมฺมมฺปิ พฺยาปนฺนํ โหติ, มโนกมฺมมฺปิ พฺยาปนฺนํ โหติ. ตสฺส พฺยาปนฺน\-กายกมฺมนฺตสฺส พฺยาปนฺน\-วจีกมฺมนฺตสฺส พฺยาปนฺน\-มโนกมฺมนฺตสฺส น ภทฺทกํ มรณํ โหติ น ภทฺทิกา กาลกิริยา.}

\prose{จิตฺเต คหปติ อพฺยาปนฺเน กายกมฺมมฺปิ อพฺยาปนฺนํ โหติ, วจีกมฺมมฺปิ อพฺยาปนฺนํ โหติ, มโนกมฺมมฺปิ อพฺยาปนฺนํ โหติ. ตสฺส อพฺยาปนฺน\-กายกมฺมนฺตสฺส อพฺยาปนฺน\-วจีกมฺมนฺตสฺส อพฺยาปนฺน\-มโนกมฺมนฺตสฺส ภทฺทกํ มรณํ โหติ ภทฺทิกา กาลกิริยา. เสยฺยาถาปิ คหปติ กูฏาคาเร สุจฺฉนฺเน กูฏมฺปิ อพฺยาปนฺนํ โหติ, โคปานสิโยปิ อพฺยาปนฺนา โหนฺติ, ภิตฺติปิ อพฺยาปนฺนา โหติ. เอวเมวํ โข คหปติ จิตฺเต อพฺยาปนฺเน กายกมฺมมฺปิ อพฺยาปนฺนํ โหติ, วจีกมฺมมฺปิ อพฺยาปนฺนํ โหติ, มโนกมฺมมฺปิ อพฺยาปนฺนํ โหติ. ตสฺส อพฺยาปนฺน\-กายกมฺมนฺตสฺส ฯเปฯ อพฺยาปนฺน\-มโนกมฺมนฺตสฺส ภทฺทกํ มรณํ โหติ ภทฺทิกา กาลงฺกิริยาติ.  อฏฺฐมํ.}

\csromanheader{2}{9. ปฐม\-นิทาน\-สุตฺต}
\proseitem{112}{ตีณิมานิ ภิกฺขเว นิทานานิ กมฺมานํ สมุทยาย. กตมานิ ตีณิ? โลโภ นิทานํ กมฺมานํ สมุทยาย, โทโส นิทานํ กมฺมานํ สมุทยาย, โมโห นิทานํ กมฺมานํ สมุทยาย. ยํ ภิกฺขเว โลภปกตํ กมฺมํ โลภชํ โลภนิทานํ โลภสมุทยํ, ตํ กมฺมํ อกุสลํ, ตํ กมฺมํ สวชฺชํ, ตํ กมฺมํ ทุกฺขวิปากํ, ตํ กมฺมํ กมฺม\-สมุทยาย สํวตฺตติ, น ตํ กมฺมํ กมฺมนิโรธาย สํวตฺตติ. ยํ ภิกฺขเว โทสปกตํ กมฺมํ โทสชํ โทสนิทานํ โทสสมุทยํ, ตํ กมฺมํ อกุสลํ, ตํ กมฺมํ สาวชฺชํ, ตํ กมฺมํ ทุกฺขวิปากํ, ตํ กมฺมํ กมฺม\-สมุทยาย สํวตฺตติ, น ตํ กมฺมํ กมฺมนิโรธาย สํวตฺตติ. ยํ ภิกฺขเว โมหปกตํ กมฺมํ โมหชํ โมหนิทานํ โมหสมุทยํ, ตํ กมฺมํ อกุสลํ, ตํ กมฺมํ สาวชฺชํ, ตํ กมฺมํ ทุกฺขวิปากํ, ตํ กมฺมํ กมฺม\-สมุทยาย สํวตฺตติ, น ตํ กมฺมํ กมฺมนิโรธาย สํวตฺตติ. อิมานิ โข ภิกฺขเว ตีณิ นิทานานิ กมฺมานํ สมุทยาย.}

\prose{ตีณิมานิ ภิกฺขเว นิทานานิ กมฺมานํ สมุทยาย. กตมานิ ตีณิ? อโลโภ นิทานํ กมฺมานํ สมุทยาย, อโทโส นิทานํ กมฺมานํ สมุทยาย, อโมโห นิทานํ กมฺมานํ สมุทยาย. ยํ ภิกฺขเว อโลภปกตํ กมฺมํ อโลภชํ อโลภนิทานํ อโลภ\-สมุทยํ,}

\chapterheadpageread{(11) 1. สมฺโพธวคฺค}
\prose{\csromanfolio{267}ตํ กมฺมํ กุสลํ, ตํ กมฺมํ อนวชฺชํ, ตํ กมฺมํ สุขวิปากํ, ตํ กมฺมํ กมฺมนิโรธาย สํวตฺตติ, น ตํ กมฺมํ กมฺม\-สมุทยาย สํวตฺตติ. ยํ ภิกฺขเว อโทสปกตํ กมฺมํ อโทสชํ อโทสนิทานํ อโทสสมุทยํ, ตํ กมฺมํ กุสลํ, ตํ กมฺมํ อนวชฺชํ, ตํ กมฺมํ สุขวิปากํ, ตํ กมฺมํ กมฺมนิโรธาย สํวตฺตติ, น ตํ กมฺมํ กมฺม\-สมุทยาย สํวตฺตติ. ยํ ภิกฺขเว อโมหปกตํ กมฺมํ อโมหชํ อโมหนิทานํ อโมหสมุทยํ, ตํ กมฺมํ กุสลํ, ตํ กมฺมํ อนวชฺชํ, ตํ กมฺมํ สุขวิปากํ, ตํ กมฺมํ กมฺมนิโรธาย สํวตฺตติ, น ตํ กมฺมํ กมฺม\-สมุทยาย สํวตฺตติ. อิมานิ โข ภิกฺขเว ตีณิ นิทานานิ กมฺมานํ สมุทยายาติ.  นวมํ.}

\csromanheader{2}{10. ทุติย\-นิทาน\-สุตฺต}
\proseitem{113}{ตีณิมานิ ภิกฺขเว นิทานานิ กมฺมานํ สมุทยาย. กตมานิ ตีณิ? อตีเต ภิกฺขเว ฉนฺท\-ราค\-ฏฺฐานิเย ธมฺเม อารพฺภ ฉนฺโท ชายติ, อนาคเต ภิกฺขเว ฉนฺท\-ราค\-ฏฺฐานิเย ธมฺเม อารพฺภ ฉนฺโท ชายติ, ปจฺจุปฺปนฺเน ภิกฺขเว ฉนฺท\-ราค\-ฏฺฐานิเย ธมฺเม อารพฺภ ฉนฺโท ชายติ. กถญฺจ ภิกฺขเว อตีเต ฉนฺท\-ราค\-ฏฺฐานิเย ธมฺเม อารพฺภ ฉนฺโท ชายติ? อตีเต ภิกฺขเว ฉนฺท\-ราค\-ฏฺฐานิเย ธมฺเม อารพฺภ เจตสา อนุวิตกฺเกติ อนุวิจาเรติ, ตสฺส อตีเต ฉนฺท\-ราค\-ฏฺฐานิเย ธมฺเม อารพฺภ เจตสา อนุวิตกฺกยโต อนุวิจารยโต ฉนฺโท ชายติ, ฉนฺทชาโต เตหิ ธมฺเมหิ สํยุตฺโต โหติ. เอตมหํ ภิกฺขเว สํโยชนํ วทามิ, โย เจตโส สาราโค. เอวํ โข ภิกฺขเว อตีเต ฉนฺท\-ราค\-ฏฺฐานิเย ธมฺเม อารพฺภ ฉนฺโท ชายติ.}

\prose{กถญฺจ ภิกฺขเว อนาคเต ฉนฺท\-ราค\-ฏฺฐานิเย ธมฺเม อารพฺภ ฉนฺโท ชายติ? อนาคเต ภิกฺขเว ฉนฺท\-ราค\-ฏฺฐานิเย ธมฺเม อารพฺภ เจตสา อนุวิตกฺเกติ อนุวิจาเรติ, ตสฺส อนาคเต ฉนฺท\-ราค\-ฏฺฐานิเย ธมฺเม อารพฺภ เจตสา อนุวิตกฺกยโต อนุวิจารยโต ฉนฺโท ชายติ, ฉนฺทชาโต เตหิ ธมฺเมหิ สํยุตฺโต โหติ. เอตมหํ ภิกฺขเว สํโยชนํ วทามิ, โย เจตโส สาราโค. เอวํ โข ภิกฺขเว อนาคเต ฉนฺท\-ราค\-ฏฺฐานิเย ธมฺเม อารพฺภ ฉนฺโท ชายติ.}

\prose{กถญฺจ ภิกฺขเว ปจฺจุปฺปนฺเน ฉนฺท\-ราค\-ฏฺฐานิเย ธมฺเม อารพฺภ ฉนฺโท ชายติ? ปจฺจุปฺปนฺเน ภิกฺขเว ฉนฺท\-ราค\-ฏฺฐานิเย ธมฺเม อารพฺภ เจตสา \csromanfolio{268}อนุวิตกฺเกติ อนุวิจาเรติ, ตสฺส ปจฺจุปฺปนฺเน ฉนฺท\-ราค\-ฏฺฐานิเย ธมฺเม อารพฺภ เจตสา อนุวิตกฺกยโต อนุวิจารยโต ฉนฺโท ชายติ, ฉนฺทชาโต เตหิ ธมฺเมหิ สํยุตฺโต โหติ. เอตมหํ ภิกฺขเว สํโยชนํ วทามิ, โย เจตโส สาราโค. เอวํ โข ภิกฺขเว ปจฺจุปฺปนฺเน ฉนฺท\-ราค\-ฏฺฐานิเย ธมฺเม อารพฺภ ฉนฺโท ชายติ. อิมานิ โข ภิกฺขเว ตีณิ นิทานานิ กมฺมานํ สมุทยาย.}

\prose{ตีณิมานิ ภิกฺขเว นิทานานิ กมฺมานํ สมุทยาย. กตมานิ ตีณิ? อตีเต ภิกฺขเว ฉนฺท\-ราค\-ฏฺฐานิเย ธมฺเม อารพฺภ ฉนฺโท น ชายติ, อนาคเต ภิกฺขเว ฉนฺท\-ราค\-ฏฺฐานิเย ธมฺเม อารพฺภ ฉนฺโท น ชายติ, ปจฺจุปฺปนฺเน ภิกฺขเว ฉนฺท\-ราค\-ฏฺฐานิเย ธมฺเม อารพฺภ ฉนฺโท น ชายติ. กถญฺจ ภิกฺขเว อตีเต ฉนฺท\-ราค\-ฏฺฐานิเย ธมฺเม อารพฺภ ฉนฺโท น ชายติ? อตีตานํ ภิกฺขเว ฉนฺท\-ราค\-ฏฺฐานิยานํ ธมฺมานํ อายติํ วิปากํ ปชานาติ, อายติํ วิปากํ วิทิตฺวา ตทภินิวตฺเตติ, ตทภินิวตฺเตตฺวา\csromansharedfootnote{ae03b073a7c5f58c}{ตทภินิวชฺเชติ, ตทภินิวชฺเชตฺวา (สี, สฺยา, กํ)} เจตสา อภินิ\-วิชฺฌิ\-ตฺวา\csromansharedfootnote{e980f961535125ed}{อภิวิราเชตฺวา (สี, สฺยา, กํ, อิ)} ปญฺญาย อติวิชฺฌ\csromansharedfootnote{63773739c41900fd}{อภินิวิชฺฌ (ก)} ปสฺสติ. เอวํ โข ภิกฺขเว อตีเต ฉนฺท\-ราค\-ฏฺฐานิเย ธมฺเม อารพฺภ ฉนฺโท น ชายติ.}

\prose{กถญฺจ ภิกฺขเว อนาคเต ฉนฺท\-ราค\-ฏฺฐานิเย ธมฺเม อารพฺภ ฉนฺโท น ชายติ? อนาคตานํ ภิกฺขเว ฉนฺท\-ราค\-ฏฺฐานิยานํ ธมฺมานํ อายติํ วิปากํ ปชานาติ, อายติํ วิปากํ วิทิตฺวา ตทภินิวตฺเตติ, ตทภินิวตฺเตตฺวา เจตสา อภินิ\-วิชฺฌิ\-ตฺวา ปญฺญาย อติวิชฺฌ ปสฺสติ. เอวํ โข ภิกฺขเว อนาคเต ฉนฺท\-ราค\-ฏฺฐานิเย ธมฺเม อารพฺภ ฉนฺโท น ชายติ.}

\prosewithclosermidruled{กถญฺจ ภิกฺขเว ปจฺจุปฺปนฺเน ฉนฺท\-ราค\-ฏฺฐานิเย ธมฺเม อารพฺภ ฉนฺโท น ชายติ? ปจฺจุปฺปนฺนานํ ภิกฺขเว ฉนฺท\-ราค\-ฏฺฐานิยานํ ธมฺมานํ อายติํ วิปากํ ปชานาติ, อายติํ วิปากํ วิทิตฺวา ตทภินิวตฺเตติ, ตทภินิวตฺเตตฺวา เจตสา อภินิ\-วิชฺฌิ\-ตฺวา ปญฺญาย อติวิชฺฌ ปสฺสติ. เอวํ โข ภิกฺขเว ปจฺจุปฺปนฺเน ฉนฺท\-ราค\-ฏฺฐานิเย ธมฺเม อารพฺภ ฉนฺโท น ชายติ. อิมานิ โข ภิกฺขเว ตีณิ นิทานานิ กมฺมานํ สมุทยายาติ.  ทสมํ.}{สมฺโพธวคฺโค ปฐโม.}
\chapterheadpageread{(12) 2. อาปายิกวคฺค \textbf{ตสฺสุทฺทานํ}}
\setlength{\csromangathapagewidth}{0pt}
\setlength{\csromangathawakwidth}{0pt}
\csromangathameasuregroup{ปุพฺเพว ทุเว อสฺสาทา, \\ อติตฺติ เทฺว จ วุตฺตานิ,}{\csromangathabat{ปุพฺเพว ทุเว อสฺสาทา,}{สมโณ รุณฺณปญฺจมํ.} \\ \csromangathabat{อติตฺติ เทฺว จ วุตฺตานิ,}{นิทานานิ อปเร ทุเวติ.}}
\csromangathasetleft{ปุพฺเพว ทุเว อสฺสาทา, \\ อติตฺติ เทฺว จ วุตฺตานิ,}
\csromangathagroup{\csromangathastanza{\csromanfolio{269}\csromangathabat{ปุพฺเพว ทุเว อสฺสาทา,}{สมโณ รุณฺณปญฺจมํ.} \\ \csromangathabat{อติตฺติ เทฺว จ วุตฺตานิ,}{นิทานานิ อปเร ทุเวติ.}}}

\csromanmatikaanchor{mtk.11}
\csromantocmarknopage{section}{(12) 2. อาปายิกวคฺค}{mtk.11}
\bookmark[dest={mtk.11},level=1]{(12) 2. อาปายิกวคฺค}
\csromanmarkh{(12)}\csromanmarkhii{2. อาปายิกวคฺค}\csromanheadernomark{1}{(12) 2. อาปายิกวคฺค}
\csromanmatikaanchor{mtk.12}
\csromantocmarkref{subsection}{1. อาปายิกสุตฺต}{mtk.12}
\bookmark[dest={mtk.12},level=2]{1. อาปายิกสุตฺต}
\csromanheader{2}{1. \textbf{อาปายิกสุตฺต}}
\proseitem{114}{ตโยเม ภิกฺขเว อาปายิกา เนรยิกา อิทมปฺปหาย. กตเม ตโย? โย จ อพฺรหฺมจารี พฺรหฺมจาริ\-ปฏิญฺโญ, โย จ สุทฺธํ พฺรหฺมจริยํ จรนฺตํ อมูลเกน\csromansharedfootnote{6ec5a81be7b30518}{อภูเตน (ก)} อพฺรหฺมจริเยน อนุทฺธํเสติ, โย จายํ เอวํวาที เอวํทิฏฺฐิ “นตฺถิ กาเมสุ โทโส”ติ, โส ตาย กาเมสุ ปาตพฺยตํ อาปชฺชติ. อิเม โข ภิกฺขเว ตโย อาปายิกา เนรยิกา อิทมปฺปหายาติ.  ปฐมํ.}

\csromanmatikaanchor{mtk.13}
\csromantocmarkref{subsection}{2. ทุลฺลภสุตฺต}{mtk.13}
\bookmark[dest={mtk.13},level=2]{2. ทุลฺลภสุตฺต}
\csromanheader{2}{2. \textbf{ทุลฺลภสุตฺต}}
\proseitem{115}{ติณฺณํ ภิกฺขเว ปาตุภาโว ทุลฺลโภ โลกสฺมิํ. กตเมสํ ติณฺณํ? ตถาคตสฺส ภิกฺขเว อรหโต สมฺมา\-สมฺพุทฺธสฺส ปาตุภาโว ทุลฺลโภ โลกสฺมิํ, ตถาคต\-ปฺปเวทิตสฺส ธมฺม\-วินยสฺส เทเสตา ปุคฺคโล ทุลฺลโภ โลกสฺมิํ, กตญฺญู กตเวที ปุคฺคโล ทุลฺลโภ โลกสฺมิํ. อิเมสํ โข ภิกฺขเว ติณฺณํ ปาตุภาโว ทุลฺลโภ โลกสฺมินฺติ.  ทุติยํ.}

\csromanmatikaanchor{mtk.14}
\csromantocmarkref{subsection}{3. อปฺปเมยฺยสุตฺต}{mtk.14}
\bookmark[dest={mtk.14},level=2]{3. อปฺปเมยฺยสุตฺต}
\csromanheader{2}{3. \textbf{อปฺปเมยฺยสุตฺต}}
\proseitem{116}{ตโยเม ภิกฺขเว ปุคฺคลา สนฺโต สํวิชฺชมานา โลกสฺมิํ. กตเม ตโย? สุปฺปเมยฺโย ทุปฺปเมยฺโย อปฺปเมยฺโย. กตโม จ ภิกฺขเว ปุคฺคโล สุปฺปเมยฺโย? อิธ ภิกฺขเว เอกจฺโจ ปุคฺคโล อุทฺธโต โหติ อุนฺนโฬ จปโล มุขโร วิกิณฺณวาโจ มุฏฺฐสฺสติ อสมฺปชาโน อสมาหิโต วิพฺภนฺตจิตฺโต ปากตินฺทฺริโย. อยํ วุจฺจติ ภิกฺขเว ปุคฺคโล สุปฺปเมยฺโย.}

\prose{\csromanfolio{270}กตโม จ ภิกฺขเว ปุคฺคโล ทุปฺปเมยฺโย? อิธ ภิกฺขเว เอกจฺโจ ปุคฺคโล อนุทฺธโต โหติ อนุนฺนโฬ อจปโล อมุขโร อวิกิณฺณวาโจ อุปฏฺฐิตสฺสติ สมฺปชาโน สมาหิโต เอกคฺคจิตฺโต สํวุตินฺทฺริโย. อยํ วุจฺจติ ภิกฺขเว ปุคฺคโล ทุปฺปเมยฺโย.}

\prose{กตโม จ ภิกฺขเว ปุคฺคโล อปฺปเมยฺโย? อิธ ภิกฺขเว ภิกฺขุ อรหํ โหติ ขีณาสโว. อยํ วุจฺจติ ภิกฺขเว ปุคฺคโล อปฺปเมยฺโย. อิเม โข ภิกฺขเว ตโย ปุคฺคลา สนฺโต สํวิชฺชมานา โลกสฺมินฺติ.  ตติยํ.}

\csromanmatikaanchor{mtk.15}
\csromantocmarkref{subsection}{4. อาเนญฺชสุตฺต}{mtk.15}
\bookmark[dest={mtk.15},level=2]{4. อาเนญฺชสุตฺต}
\csromanheader{2}{4. \textbf{อาเนญฺชสุตฺต}}
\proseitem{117}{ตโยเม ภิกฺขเว ปุคฺคลา สนฺโต สํวิชฺชมานา โลกสฺมิํ. กตเม ตโย? อิธ ภิกฺขเว เอกจฺโจ ปุคฺคโล สพฺพโส รูปสญฺญานํ สมติกฺกมา ปฏิฆ\-สญฺญานํ อตฺถงฺคมา นานตฺต\-สญฺญานํ อมนสิการา “อนนฺโต อากาโส”ติ อากาสา\-นญฺจา\-ยตนํ อุปสมฺปชฺช วิหรติ, โส ตทสฺสาเทติ, ตํ นิกาเมติ, เตน จ วิตฺติํ อาปชฺชติ. ตตฺร ฐิโต ตทธิมุตฺโต ตพฺพหุล\-วิหารี อปริหีโน กาลํ กุรุมาโน อากาสานญฺจายตนู\-ปคานํ เทวานํ สหพฺยตํ อุปปชฺชติ. อากาสานญฺจายตนู\-ปคานํ ภิกฺขเว เทวานํ วีสติ กปฺปสหสฺสานิ อายุปฺปมาณํ, ตตฺถ ปุถุชฺชโน ยาวตายุกํ ฐตฺวา ยาวตกํ เตสํ เทวานํ อายุปฺปมาณํ, ตํ สพฺพํ เขเปตฺวา นิรยมฺปิ คจฺฉติ, ติรจฺฉาน\-โยนิมฺปิ คจฺฉติ, เปตฺติวิสยมฺปิ คจฺฉติ. ภควโต ปน สาวโก ตตฺถ ยาวตายุกํ ฐตฺวา ยาวตกํ เตสํ เทวานํ อายุปฺปมาณํ, ตํ สพฺพํ เขเปตฺวา ตสฺมิํเยว ภเว ปรินิพฺพายติ. อยํ โข ภิกฺขเว วิเสโส อยํ อธิปฺปยาโส อิทํ นานากรณํ สุตวโต อริยสาวกสฺส อสฺสุตวตา ปุถุชฺชเนน ยทิทํ คติยา อุปปตฺติยา.}

\prose{ปุน จปรํ ภิกฺขเว อิเธกจฺโจ ปุคฺคโล สพฺพโส อากาสา\-นญฺจา\-ยตนํ สมติกฺกมฺม “อนนฺตํ วิญฺญาณนฺ”ติ วิญฺญาณญฺจา\-ยตนํ อุปสมฺปชฺช วิหรติ. โส ตทสฺสาเทติ, ตํ นิกาเมติ, เตน จ วิตฺติํ อาปชฺชติ. ตตฺร ฐิโต ตทธิมุตฺโต ตพฺพหุล\-วิหารี อปริหีโน กาลํ กุรุมาโน วิญฺญาณญฺจายตนู\-ปคานํ เทวานํ สหพฺยตํ อุปปชฺชติ. วิญฺญาณญฺจายตนู\-ปคานํ ภิกฺขเว เทวานํ จตฺตารีสํ กปฺปสหสฺสานิ อายุปฺปมาณํ,}

\vspace{\tipitakaparskip}
\csromancenter{(12) 2. อาปายิกวคฺค ตตฺถ ปุถุชฺชโน ยาวตายุกํ ฐตฺวา ยาวตกํ เตสํ เทวานํ อายุปฺปมาณํ, ตํ สพฺพํ เขเปตฺวา นิรยมฺปิ คจฺฉติ, ติรจฺฉาน\-โยนิมฺปิ คจฺฉติ, เปตฺติวิสยมฺปิ คจฺฉติ. ภควโต ปน สาวโก ตตฺถ ยาวตายุกํ ฐตฺวา ยาวตกํ เตสํ เทวานํ อายุปฺปมาณํ, ตํ สพฺพํ เขเปตฺวา ตสฺมิํ เยว ภเว ปรินิพฺพายติ. อยํ โข ภิกฺขเว วิเสโส อยํ อธิปฺปยาโส อิทํ นานากรณํ สุตวโต อริยสาวกสฺส อสฺสุตวตา ปุถุชฺชเนน ยทิทํ คติยา อุปปตฺติยา.}
\prose{\csromanfolio{271}ปุน จปรํ ภิกฺขเว อิเธกจฺโจ ปุคฺคโล สพฺพโส วิญฺญาณญฺจา\-ยตนํ สมติกฺกมฺม “นตฺถิ กิญฺจี”ติ อากิญฺจญฺญา\-ยตนํ อุปสมฺปชฺช วิหรติ. โส ตทสฺสาเทติ, ตํ นิกาเมติ, เตน จ วิตฺติํ อาปชฺชติ. ตตฺร ฐิโต ตทธิมุตฺโต ตพฺพหุล\-วิหารี อปริหีโน กาลํ กุรุมาโน อากิญฺจญฺญายตนู\-ปคานํ เทวานํ สหพฺยตํ อุปปชฺชติ. อากิญฺจญฺญายตนู\-ปคานํ ภิกฺขเว เทวานํ สฏฺฐิ กปฺปสหสฺสานิ อายุปฺปมาณํ, ตตฺถ ปุถุชฺชโน ยาวตายุกํ ฐตฺวา ยาวตกํ เตสํ เทวานํ อายุปฺปมาณํ, ตํ สพฺพํ เขเปตฺวา นิรยมฺปิ คจฺฉติ, ติรจฺฉาน\-โยนิมฺปิ คจฺฉติ, เปตฺติวิสยมฺปิ คจฺฉติ. ภควโต ปน สาวโก ตตฺถ ยาวตายุกํ ฐตฺวา ยาวตกํ เตสํ เทวานํ อายุปฺปมาณํ, ตํ สพฺพํ เขเปตฺวา ตสฺมิํเยว ภเว ปรินิพฺพายติ. อยํ โข ภิกฺขเว วิเสโส อยํ อธิปฺปยาโส อิทํ นานากรณํ สุตวโต อริยสาวกสฺส อสฺสุตวตา ปุถุชฺชเนน ยทิทํ คติยา อุปปตฺติยา. อิเม โข ภิกฺขเว ตโย ปุคฺคลา สนฺโต สํวิชฺชมานา โลกสฺมินฺติ.  จตุตฺถํ.}

\csromanmatikaanchor{mtk.16}
\csromantocmarkref{subsection}{5. วิปตฺติสมฺปทาสุตฺต}{mtk.16}
\bookmark[dest={mtk.16},level=2]{5. วิปตฺติสมฺปทาสุตฺต}
\csromanheader{2}{5. วิปตฺติ\-สมฺปทา\-สุตฺต}
\proseitem{118}{ติสฺโส อิมา ภิกฺขเว วิปตฺติโย. กตมา ติสฺโส? สีลวิปตฺติ จิตฺตวิปตฺติ ทิฏฺฐิวิปตฺติ. กตมา จ ภิกฺขเว สีลวิปตฺติ? อิธ ภิกฺขเว เอกจฺโจ ปาณาติปาตี โหติ, อทินฺนาทายี โหติ, กาเมสุ\-มิจฺฉาจารี โหติ, มุสาวาที โหติ, ปิสุณวาโจ โหติ, ผรุสวาโจ โหติ, สมฺผปฺปลาปี โหติ. อยํ วุจฺจติ ภิกฺขเว สีลวิปตฺติ.}

\prose{กตมา จ ภิกฺขเว จิตฺตวิปตฺติ? อิธ ภิกฺขเว เอกจฺโจ อภิชฺฌาลุ โหติ พฺยาปนฺนจิตฺโต. อยํ วุจฺจติ ภิกฺขเว จิตฺตวิปตฺติ.}

\prose{\csromanfolio{272}กตมา จ ภิกฺขเว ทิฏฺฐิวิปตฺติ? อิธ ภิกฺขเว เอกจฺโจ มิจฺฉาทิฏฺฐิโก โหติ วิปรีต\-ทสฺสโน “นตฺถิ ทินฺนํ, นตฺถิ ยิฏฺฐํ, นตฺถิ หุตํ, นตฺถิ สุกต\-ทุกฺกฏานํ กมฺมานํ ผลํ วิปาโก, นตฺถิ อยํ โลโก, นตฺถิ ปโร โลโก, นตฺถิ มาตา, นตฺถิ ปิตา, นตฺถิ สตฺตา โอปปาติกา, นตฺถิ โลเก สมณพฺราหฺมณา สมฺมคฺคตา สมฺมาปฏิปนฺนา, เย อิมญฺจ โลกํ ปรญฺจ โลกํ สยํ อภิญฺญา สจฺฉิกตฺวา ปเวเทนฺตี”ติ. อยํ วุจฺจติ ภิกฺขเว ทิฏฺฐิวิปตฺติ. สีลวิปตฺติ\-เหตุ วา ภิกฺขเว สตฺตา กายสฺส เภทา ปรํ มรณา อปายํ ทุคฺคติํ วินิปาตํ นิรยํ อุปปชฺชนฺติ. จิตฺตวิปตฺติ\-เหตุ วา ภิกฺขเว สตฺตา กายสฺส เภทา ปรํ มรณา อปายํ ทุคฺคติํ วินิปาตํ นิรยํ อุปปชฺชนฺติ. ทิฏฺฐิวิปตฺติ\-เหตุ วา ภิกฺขเว สตฺตา กายสฺส เภทา ปรํ มรณา อปายํ ทุคฺคติํ วินิปาตํ นิรยํ อุปปชฺชนฺติ. อิมา โข ภิกฺขเว ติสฺโส วิปตฺติโยติ.}

\prose{ติสฺโส อิมา ภิกฺขเว สมฺปทา. กตมา ติสฺโส? สีลสมฺปทา จิตฺตสมฺปทา ทิฏฺฐิสมฺปทา. กตมา จ ภิกฺขเว สีลสมฺปทา? อิธ ภิกฺขเว เอกจฺโจ ปาณาติปาตา ปฏิวิรโต โหติ, อทินฺนาทานา ปฏิวิรโต โหติ, กาเมสุ\-มิจฺฉาจารา ปฏิวิรโต โหติ, มุสาวาทา ปฏิวิรโต โหติ, ปิสุณาย วาจาย ปฏิวิรโต โหติ, ผรุสาย วาจาย ปฏิวิรโต โหติ, สมฺผปฺปลาปา ปฏิวิรโต โหติ. อยํ วุจฺจติ ภิกฺขเว สีลสมฺปทา.}

\prose{กตมา จ ภิกฺขเว จิตฺตสมฺปทา? อิธ ภิกฺขเว เอกจฺโจ อนภิชฺฌาลุ โหติ อพฺยาปนฺนจิตฺโต. อยํ วุจฺจติ ภิกฺขเว จิตฺตสมฺปทา.}

\prose{กตมา จ ภิกฺขเว ทิฏฺฐิสมฺปทา? อิธ ภิกฺขเว เอกจฺโจ สมฺมาทิฏฺฐิโก โหติ อวิปรีต\-ทสฺสโน “อตฺถิ ทินฺนํ, อตฺถิ ยิฏฺฐํ, อตฺถิ หุตํ, อตฺถิ สุกต\-ทุกฺกฏานํ กมฺมานํ ผลํ วิปาโก, อตฺถิ อยํ โลโก, อตฺถิ ปโร โลโก, อตฺถิ มาตา, อตฺถิ ปิตา, อตฺถิ สตฺตา โอปปาติกา, อตฺถิ โลเก สมณพฺรหฺมณา สมฺมคฺคตา สมฺมาปฏิปนฺนา, เย อิมญฺจ โลกํ ปรญฺจ โลกํ สยํ อภิญฺญา สจฺฉิกตฺวา ปเวเทนฺตี”ติ. อยํ วุจฺจติ ภิกฺขเว ทิฏฺฐิสมฺปทา. สีลสมฺปทา\-เหตุ วา ภิกฺขเว สตฺตา กายสฺส เภทา ปรํ มรณา สุคติํ สคฺคํ โลกํ อุปปชฺชนฺติ. จิตฺตสมฺปทา\-เหตุ วา ภิกฺขเว สตฺตา กายสฺส เภทา ปรํ มรณา สุคติํ สคฺคํ โลกํ}

\vspace{\tipitakaparskip}
\csromancenter{(12) 2. อาปายิกวคฺค อุปปชฺชนฺติ. ทิฏฺฐิสมฺปทา\-เหตุ วา ภิกฺขเว สตฺตา กายสฺส เภทา ปรํ มรณา สุคติํ สคฺคํ โลกํ อุปปชฺชนฺติ. อิมา โข ภิกฺขเว ติสฺโส สมฺปทาติ.  ปญฺจมํ.}
\csromanmatikaanchor{mtk.17}
\csromantocmarkref{subsection}{6. อปณฺณกสุตฺต}{mtk.17}
\bookmark[dest={mtk.17},level=2]{6. อปณฺณกสุตฺต}
\csromanheader{2}{6. \textbf{อปณฺณกสุตฺต}}
\proseitem{119}{\csromanfolio{273}ติสฺโส อิมา ภิกฺขเว วิปตฺติโย. กตมา ติสฺโส? สีลวิปตฺติ จิตฺตวิปตฺติ ทิฏฺฐิวิปตฺติ. กตมา จ ภิกฺขเว สีลวิปตฺติ? อิธ ภิกฺขเว เอกจฺโจ ปาณาติปาตี โหติ ฯเปฯ สมฺผปฺปลาปี โหติ. อยํ วุจฺจติ ภิกฺขเว สีลวิปตฺติ.}

\prose{กตมา จ ภิกฺขเว จิตฺตวิปตฺติ? อิธ ภิกฺขเว เอกจฺโจ อภิชฺฌาลุ โหติ พฺยาปนฺนจิตฺโต. อยํ วุจฺจติ ภิกฺขเว จิตฺตวิปตฺติ.}

\prose{กตมา จ ภิกฺขเว ทิฏฺฐิวิปตฺติ? อิธ ภิกฺขเว เอกจฺโจ มิจฺฉาทิฏฺฐิโก โหติ วิปรีต\-ทสฺสโน “นตฺถิ ทินฺนํ, นตฺถิ ยิฏฺฐํ ฯเปฯ เย อิมญฺจ โลกํ ปรญฺจ โลกํ สยํ อภิญฺญา สจฺฉิกตฺวา ปเวเทนฺตี”ติ. อยํ วุจฺจติ ภิกฺขเว ทิฏฺฐิวิปตฺติ. สีลวิปตฺติ\-เหตุ วา ภิกฺขเว ฯเปฯ. ทิฏฺฐิวิปตฺติ\-เหตุ วา ภิกฺขเว สตฺตา กายสฺส เภทา ปรํ มรณา อปายํ ทุคฺคติํ วินิปาตํ นิรยํ อุปปชฺชนฺติ. เสยฺยถาปิ ภิกฺขเว อปณฺณโก มณิ อุทฺธํ ขิตฺโต เยน เยเนว ปติฏฺฐาติ, สุปฺปติฏฺฐิตํเยว ปติฏฺฐาติ. เอวเมวํ โข ภิกฺขเว สีลวิปตฺติ\-เหตุ วา สตฺตา ฯเปฯ อุปปชฺชนฺติ. อิมา โข ภิกฺขเว ติสฺโส วิปตฺติโยติ.}

\prose{ติสฺโส อิมา ภิกฺขเว สมฺปทา. กตมา ติสฺโส? สีลสมฺปทา จิตฺตสมฺปทา ทิฏฺฐิสมฺปทา. กตมา จ ภิกฺขเว สีลสมฺปทา, อิธ ภิกฺขเว เอกจฺโจ ปาณาติปาตา ปฏิวิรโต โหติ ฯเปฯ. อยํ วุจฺจติ ภิกฺขเว สีลสมฺปทา.}

\prose{กตมา จ ภิกฺขเว จิตฺตสมฺปทา? อิธ ภิกฺขเว เอกจฺโจ อนภิชฺฌาลุ โหติ อพฺยาปนฺนจิตฺโต. อยํ วุจฺจติ ภิกฺขเว จิตฺตสมฺปทา.}

\prose{กตมา จ ภิกฺขเว ทิฏฺฐิสมฺปทา? อิธ ภิกฺขเว เอกจฺโจ สมฺมาทิฏฺฐิโก โหติ อวิปรีต\-ทสฺสโน “อตฺถิ ทินฺนํ, อตฺถิ ยิฏฺฐํ ฯเปฯ เย อิมญฺจ โลกํ ปรญฺจ โลกํ สยํ อภิญฺญา สจฺฉิกตฺวา ปเวเทนฺตี”ติ. อยํ วุจฺจติ ภิกฺขเว ทิฏฺฐิสมฺปทา. สีลสมฺปทา\-เหตุ วา ภิกฺขเว สตฺตา กายสฺส เภทา ปรํ มรณา สุคติํ สคฺคํ โลกํ อุปปชฺชนฺติ. จิตฺตสมฺปทา\-เหตุ วา ฯเปฯ. \csromanfolio{274}ทิฏฺฐิสมฺปทา\-เหตุ วา ภิกฺขเว สตฺตา กายสฺส เภทา ปรํ มรณา สุคหิํ สคฺคํ โลกํ อุปปชฺชนฺติ. เสยฺยถาปิ ภิกฺขเว อปณฺณโก มณิ อุทฺธํ ขิตฺโต เยน เยเนว ปติฏฺฐาติ, สุปฺปติฏฺฐิตํเยว ปติฏฺฐาติ. เอวเมวํ โข ภิกฺขเว สีลสมฺปทา\-เหตุ วา สตฺตา กายสฺส เภทา ปรํ มรณา สุคติํ สคฺคํ โลกํ อุปปชฺชนฺติ. จิตฺตสมฺปทา\-เหตุ วา สตฺตา ฯเปฯ. ทิฏฺฐิสมฺปทา\-เหตุ วา สตฺตา กายสฺส เภทา ปรํ มรณา สุคติํ สคฺคํ โลกํ อุปปชฺชนฺติ. อิมา โข ภิกฺขเว ติสฺโส สมฺปทาติ.  ฉฏฺฐํ.}

\csromanmatikaanchor{mtk.18}
\csromantocmarkref{subsection}{7. กมฺมนฺตสุตฺต}{mtk.18}
\bookmark[dest={mtk.18},level=2]{7. กมฺมนฺตสุตฺต}
\csromanheader{2}{7. \textbf{กมฺมนฺตสุตฺต}}
\proseitem{120}{ติสฺโส อิมา ภิกฺขเว วิปตฺติโย. กตมา ติสฺโส? กมฺมนฺต\-วิปตฺติ อาชีววิปตฺติ ทิฏฺฐิวิปตฺติ. กตมา จ ภิกฺขเว กมฺมนฺต\-วิปตฺติ? อิธ ภิกฺขเว เอกจฺโจ ปาณาติปาตี โหติ ฯเปฯ สมฺผปฺปลาปี โหติ. อยํ วุจฺจติ ภิกฺขเว กมฺมนฺต\-วิปตฺติ.}

\prose{กตมา จ ภิกฺขเว อาชีววิปตฺติ? อิธ ภิกฺขเว เอกจฺโจ มิจฺฉาชีโว โหติ, มิจฺฉาอาชีเวน ชีวิกํ\csromansharedfootnote{5c432ab1fb9f51bd}{ชีวิตํ (สฺยา, กํ, ก)} กปฺเปติ. อยํ วุจฺจติ ภิกฺขเว อาชีววิปตฺติ.}

\prose{กตมา จ ภิกฺขเว ทิฏฺฐิวิปตฺติ? อิธ ภิกฺขเว เอกจฺโจ มิจฺฉาทิฏฺฐิโก โหติ วิปรีต\-ทสฺสโน “นตฺถิ ทินฺนํ, นตฺถิ ยิฏฺฐํ ฯเปฯ เย อิมญฺจ โลกํ ปรญฺจ โลกํ สยํ อภิญฺญา สจฺฉิกตฺวา ปเวเทนฺตี”ติ. อยํ วุจฺจติ ภิกฺขเว ทิฏฺฐิวิปตฺติ. อิมา โข ภิกฺขเว ติสฺโส วิปตฺติโยติ.}

\prose{ติสฺโส อิมา ภิกฺขเว สมฺปทา. กตมา ติสฺโส? กมฺมนฺต\-สมฺปทา อาชีวสมฺปทา ทิฏฺฐิสมฺปทา. กตมา จ ภิกฺขเว กมฺมนฺต\-สมฺปทา? อิธ ภิกฺขเว เอกจฺโจ ปาณาติปาตา ปฏิวิรโต โหติ ฯเปฯ สมฺผปฺปลาปา ปฏิวิรโต โหติ. อยํ วุจฺจติ ภิกฺขเว กมฺมนฺต\-สมฺปทา.}

\prose{กตมา จ ภิกฺขเว อาชีวสมฺปทา? อิธ ภิกฺขเว เอกจฺโจ สมฺมาอาชีโว โหติ, สมฺมาอาชีเวน ชีวิกํ กปฺเปติ. อยํ วุจฺจติ ภิกฺขเว อาชีวสมฺปทา.}

\prose{กตมา จ ภิกฺขเว ทิฏฺฐิสมฺปทา? อิธ ภิกฺขเว เอกจฺโจ สมฺมาทิฏฺฐิโก โหติ อวิปรีต\-ทสฺสโน “อตฺถิ ทินฺนํ, อตฺถิ ยิฏฺฐํ ฯเปฯ เย อิมญฺจ โลกํ}

\vspace{\tipitakaparskip}
\csromancenter{(12) 2. อาปายิกวคฺค ปรญฺจ โลกํ สยํ อภิญฺญา สจฺฉิกตฺวา ปเวเทนฺตี”ติ. อยํ วุจฺจติ ภิกฺขเว ทิฏฺฐิสมฺปทา. อิมา โข ภิกฺขเว ติสฺโส สมฺปทาติ.  สตฺตมํ.}
\csromanheader{2}{8. ปฐม\-โสเจยฺย\-สุตฺต}
\csromanmatikaanchor{mtk.19}
\csromantocmarkref{subsection}{8-9. โสเจยฺยสุตฺตานิ}{mtk.19}
\bookmark[dest={mtk.19},level=2]{8-9. โสเจยฺยสุตฺตานิ}
\proseitem{121}{\csromanfolio{275}ตีณิมานิ ภิกฺขเว โสเจยฺยานิ. กตมานิ ตีณิ? กายโสเจยฺยํ วจีโสเจยฺยํ มโนโสเจยฺยํ. กตมญฺจ ภิกฺขเว กายโสเจยฺยํ? อิธ ภิกฺขเว เอกจฺโจ ปาณาติปาตา ปฏิวิรโต โหติ, อทินฺนาทานา ปฏิวิรโต โหติ, กาเมสุ\-มิจฺฉาจารา ปฏิวิรโต โหติ. อิทํ วุจฺจติ ภิกฺขเว กายโสเจยฺยํ.}

\prose{กตมญฺจ ภิกฺขเว วจีโสเจยฺยํ? อิธ ภิกฺขเว เอกจฺโจ มุสาวาทา ปฏิวิรโต โหติ, ปิสุณาย วาจาย ปฏิวิรโต โหติ, ผรุสาย วาจาย ปฏิวิรโต โหติ, สมฺผปฺปลาปา ปฏิวิรโต โหติ. อิทํ วุจฺจติ ภิกฺขเว วจีโสเจยฺยํ.}

\prose{กตมญฺจ ภิกฺขเว มโนโสเจยฺยํ? อิธ ภิกฺขเว เอกจฺโจ อนภิชฺฌาลุ โหติ อพฺยาปนฺนจิตฺโต สมฺมาทิฏฺฐิโก. อิทํ วุจฺจติ ภิกฺขเว มโนโสเจยฺยํ. อิมานิ โข ภิกฺขเว ตีณิ โสเจยฺยานีติ.  อฏฺฐมํ.}

\csromanheader{2}{9. ทุติย\-โสเจยฺย\-สุตฺต}
\proseitem{122}{ตีณิมานิ ภิกฺขเว โสเจยฺยานิ. กตมานิ ตีณิ? กายโสเจยฺยํ วจีโสเจยฺยํ มโนโสเจยฺยํ. กตมญฺจ ภิกฺขเว กายโสเจยฺยํ? อิธ ภิกฺขเว ภิกฺขุ ปาณาติปาตา ปฏิวิรโต โหติ, อทินฺนาทานา ปฏิวิรโต โหติ, อพฺรหฺมจริยา ปฏิวิรโต โหติ. อิทํ วุจฺจติ ภิกฺขเว กายโสเจยฺยํ.}

\prose{กตมญฺจ ภิกฺขเว วจีโสเจยฺยํ? อิธ ภิกฺขเว ภิกฺขุ มุสาวาทา ปฏิวิรโต โหติ. ปิสุณาย วาจาย ปฏิวิรโต โหติ, ผรุสาย วาจาย ปฏิวิรโต โหติ, สมฺผปฺปลาปา ปฏิวิรโต โหติ. อิทํ วุจฺจติ ภิกฺขเว วจีโสเจยฺยํ.}

\prose{กตมญฺจ ภิกฺขเว มโนโสเจยฺยํ, อิธ ภิกฺขเว ภิกฺขุ สนฺตํ วา อชฺฌตฺตํ กามจฺฉนฺทํ “อตฺถิ เม อชฺฌตฺตํ กามจฺฉนฺโท”ติ ปชานาติ, อสนฺตํ วา อชฺฌตฺตํ กามจฺฉนฺทํ “นตฺถิ เม อชฺฌตฺตํ กามจฺฉนฺโท”ติ ปชานาติ, ยถา จ อนุปฺปนฺนสฺส \csromanfolio{276}กามจฺฉนฺทสฺส อุปฺปาโท โหติ ตญฺจ ปชานาติ, ยถา จ อุปฺปนฺนสฺส กามจฺฉนฺทสฺส ปหานํ โหติ ตญฺจ ปชานาติ. ยถา จ ปหีนสฺส กามจฺฉนฺทสฺส อายติํ อนุปฺปาโท โหติ ตญฺจ ปชานาติ. สนฺตํ วา อชฺฌตฺตํ พฺยาปาทํ “อตฺถิ เม อชฺฌตฺตํ พฺยาปาโท”ติ ปชานาติ, อสนฺตํ วา อชฺฌตฺตํ พฺยาปาทํ “นตฺถิ เม อชฺฌตฺตํ พฺยาปาโท”ติ ปชานาติ, ยถา จ อนุปฺปนฺนสฺส พฺยาปาทสฺส อุปฺปาโท โหติ ตญฺจ ปชานาติ, ยถา จ อนุปฺปนฺนสฺส พฺยาปาทสฺส ปหานํ โหติ ตญฺจ ปชานาติ, ยถา จ ปหีนสฺส พฺยาปาทสฺส อายติํ อนุปฺปาโท โหติ ตญฺจ ปชานาติ. สนฺตํ วา อชฺฌตฺตํ ถินมิทฺธํ “อตฺถิ เม อชฺฌตฺตํ ถินมิทฺธนฺ”ติ ปชานาติ, อสนฺตํ วา อชฺฌตฺตํ ถินมิทฺธํ “นตฺถิ เม อชฺฌตฺตํ ถินมิทฺธนฺ”ติ ปชานาติ, ยถา จ อนุปฺปนฺนสฺส ถินมิทฺธสฺส อุปฺปาโท โหติ ตญฺจ ปชานาติ, ยถา จ อุปฺปนฺนสฺส ถินมิทฺธสฺส ปหานํ โหติ ตญฺจ ปชานาติ, ยถา จ ปหีนสฺส ถินมิทฺธสฺส อายติํ อนุปฺปาโท โหติ ตญฺจ ปชานาติ. สนฺตํ วา อชฺฌตฺตํ อุทฺธจฺจ\-กุกฺกุจฺจํ “อตฺถิ เม อชฺฌตฺตํ อุทฺธจฺจกุกฺกุจฺจนฺ”ติ ปชานาติ, อสนฺตํ วา อชฺฌตฺตํ อุทฺธจฺจ\-กุกฺกุจฺจํ “นตฺถิ เม อชฺฌตฺตํ อุทฺธจฺจกุกฺกุจฺจนฺ”ติ ปชานาติ, ยถา จ อนุปฺปนฺนสฺส อุทฺธจฺจ\-กุกฺกุจฺจสฺส อุปฺปาโท โหติ ตญฺจ ปชานาติ, ยถา จ อุปฺปนฺนสฺส อุทฺธจฺจ\-กุกฺกุจฺจสฺส ปหานํ โหติ ตญฺจ ปชานาติ, ยถา จ ปหีนสฺส อุทฺธจฺจ\-กุกฺกุจฺจสฺส อายติํ อนุปฺปาโท โหติ ตญฺจ ปชานาติ. สนฺตํ วา อชฺฌตฺตํ วิจิกิจฺฉํ “อตฺถิ เม อชฺฌตฺตํ วิจิกิจฺฉา”ติ ปชานาติ, อสนฺตํ วา อชฺฌตฺตํ วิจิกิจฺฉํ “นตฺถิ เม อชฺฌตฺตํ วิจิกิจฺฉา”ติ ปชานาติ. ยถา จ อนุปฺปนฺนาย วิจิกิจฺฉาย อุปฺปาโท โหติ ตญฺจ ปชานาติ, ยถา จ อุปฺปนฺนาย วิจิกิจฺฉาย ปหานํ โหติ ตญฺจ ปชานาติ, ยถา จ ปหีนาย วิจิกิจฺฉาย อายติํ อนุปฺปาโท โหติ ตญฺจ ปชานาติ. อิทํ วุจฺจติ ภิกฺขเว มโนโสเจยฺยํ. อิมานิ โข ภิกฺขเว ตีณิ โสเจยฺยานีติ.}

\setlength{\csromangathapagewidth}{0pt}
\setlength{\csromangathawakwidth}{0pt}
\csromangathameasuregroup{\textsuperscript{*}กายสุจิํ วจีสุจิํ, \\ สุจิํ โสเจยฺยสมฺปนฺนํ,}{\csromangathabat{\textsuperscript{*}กายสุจิํ วจีสุจิํ,}{เจโตสุจิํ อนาสวํ.} \\ \csromangathabat{สุจิํ โสเจยฺยสมฺปนฺนํ,}{อาหุ นินฺหาตปาปกนฺติ.} \\ นวมํ.}
\csromangathameasurewak{นวมํ.}
\csromangathasetleft{\textsuperscript{*}กายสุจิํ วจีสุจิํ, \\ สุจิํ โสเจยฺยสมฺปนฺนํ,}
\csromangathagroup{\csromangathastanza{\csromangathabat{\csromansymbolmark{*}กายสุจิํ วจีสุจิํ,}{เจโตสุจิํ อนาสวํ.} \\ \csromangathabat{สุจิํ โสเจยฺยสมฺปนฺนํ,}{อาหุ นินฺหาต\-ปาปกนฺติ.}}\csromangathastanza{\csromangathawak{นวมํ.}}}

\csromanmatikaanchor{mtk.20}
\csromantocmarkref{subsection}{10. โมเนยฺยสุตฺต}{mtk.20}
\bookmark[dest={mtk.20},level=2]{10. โมเนยฺยสุตฺต}
\csromanheader{2}{10. \textbf{โมเนยฺยสุตฺต}}
\proseitem{123}{ตีณิมานิ ภิกฺขเว โมเนยฺยานิ. กตมานิ ตีณิ? กายโมเนยฺยํ วจีโมเนยฺยํ มโนโมเนยฺยํ. กตมญฺจ ภิกฺขเว}

\vspace{\tipitakaparskip}
\csromancenter{(13) 3. กุสินารวคฺค กายโมเนยฺยํ? อิธ ภิกฺขเว ภิกฺขุ ปาณาติปาตา ปฏิวิรโต โหติ, อทินฺนาทานา ปฏิวิรโต โหติ, อพฺรหฺมจริยา ปฏิวิรโต โหติ. อิทํ วุจฺจติ ภิกฺขเว กายโมเนยฺยํ.}
\prose{\csromanfolio{277}กตมญฺจ ภิกฺขเว วจีโมเนยฺยํ? อิธ ภิกฺขเว ภิกฺขุ มุสาวาทา ปฏิวิรโต โหติ, ปิสุณาย วาจาย ปฏิวิรโต โหติ, ผรุสาย วาจาย ปฏิวิรโต โหติ, สมฺผปฺปลาปา ปฏิวิรโต โหติ. อิทํ วุจฺจติ ภิกฺขเว วจีโมเนยฺยํ.}

\prose{กตมญฺจ ภิกฺขเว มโนโมเนยฺยํ? อิธ ภิกฺขเว ภิกฺขุ อาสวานํ ขยา อนาสวํ เจโตวิมุตฺติํ ปญฺญาวิมุตฺติํ ทิฏฺเฐว ธมฺเม สยํ อภิญฺญา สจฺฉิกตฺวา อุปสมฺปชฺช วิหรติ. อิทํ วุจฺจติ ภิกฺขเว มโนโมเนยฺยํ. อิมานิ โข ภิกฺขเว ตีณิ โมเนยฺยานีติ.}

\setlength{\csromangathapagewidth}{0pt}
\setlength{\csromangathawakwidth}{0pt}
\csromangathameasuregroup{กายมุนิํ วจีมุนิํ, \\ มุนิํ โมเนยฺยสมฺปนฺนํ,}{\csromangathabat{กายมุนิํ วจีมุนิํ,}{เจโตมุนิํ อนาสวํ.} \\ \csromangathabat{มุนิํ โมเนยฺยสมฺปนฺนํ,}{อาหุ สพฺพปฺปหายินนฺติ.} \\ ทสมํ.}
\csromangathameasurewak{ทสมํ.}
\csromangathasetleft{กายมุนิํ วจีมุนิํ, \\ มุนิํ โมเนยฺยสมฺปนฺนํ,}
\csromangathagroupwithclosermidruled{\csromangathastanza{\csromangathabat{กายมุนิํ วจีมุนิํ,}{เจโตมุนิํ อนาสวํ.} \\ \csromangathabat{มุนิํ โมเนยฺย\-สมฺปนฺนํ,}{อาหุ สพฺพ\-ปฺปหายินนฺติ.}}\csromangathastanza{\csromangathawak{ทสมํ.}}}{อาปายิกวคฺโค ทุติโย.}
\csromancenter{\textbf{ตสฺสุทฺทานํ}}
\setlength{\csromangathapagewidth}{0pt}
\setlength{\csromangathawakwidth}{0pt}
\csromangathameasuregroup{อาปายิโก ทุลฺลโภ อปฺปเมยฺยํ, \\ อปณฺณโก จ กมฺมนฺโต,}{\csromangathabat{อาปายิโก ทุลฺลโภ อปฺปเมยฺยํ,}{อาเนญฺชวิปตฺติสมฺปทา.} \\ \csromangathabat{อปณฺณโก จ กมฺมนฺโต,}{เทฺว โสเจยฺยานิ โมเนยฺยนฺติ.}}
\csromangathasetleft{อาปายิโก ทุลฺลโภ อปฺปเมยฺยํ, \\ อปณฺณโก จ กมฺมนฺโต,}
\csromangathagroup{\csromangathastanza{\csromangathabat{อาปายิโก ทุลฺลโภ อปฺปเมยฺยํ,}{อาเนญฺชวิปตฺติสมฺปทา.} \\ \csromangathabat{อปณฺณโก จ กมฺมนฺโต,}{เทฺว โสเจยฺยานิ โมเนยฺยนฺติ.}}}

\csromanmatikaanchor{mtk.21}
\csromantocmarknopage{section}{(13) 3. กุสินารวคฺค}{mtk.21}
\bookmark[dest={mtk.21},level=1]{(13) 3. กุสินารวคฺค}
\csromanmarkh{(13)}\csromanmarkhii{3. กุสินารวคฺค}\csromanheadernomark{1}{(13) 3. กุสินารวคฺค}
\csromanmatikaanchor{mtk.22}
\csromantocmarkref{subsection}{1. กุสินารสุตฺต}{mtk.22}
\bookmark[dest={mtk.22},level=2]{1. กุสินารสุตฺต}
\csromanheader{2}{1. \textbf{กุสินารสุตฺต}}
\proseitem{124}{เอกํ สมยํ ภควา กุสินารายํ วิหรติ พลิหรเณ วนสณฺเฑ. ตตฺร โข ภควา ภิกฺขู อามนฺเตสิ “ภิกฺขโว”ติ. “ภทนฺเต”ติ เต ภิกฺขู ภควโต ปจฺจสฺโสสุํ. ภควา เอตทโวจ\csromandash{}}

\prose{อิธ ภิกฺขเว ภิกฺขุ อญฺญตรํ คามํ วา นิคมํวา อุปนิสฺสาย วิหรติ, ตเมนํ คหปติ วา คหปติปุตฺโต วา อุปสงฺกมิตฺวา สฺวาตนาย ภตฺเตน \csromanfolio{278}นิมนฺเตติ, อากงฺขมาโน ภิกฺขเว ภิกฺขุ อธิวาเสติ. โส ตสฺสา รตฺติยา อจฺจเยน ปุพฺพณฺหสมยํ นิวาเสตฺวา ปตฺต\-จีวรมา\-ทาย เยน ตสฺส คหปติสฺส วา คหปติ\-ปุตฺตสฺส วา นิเวสนํ เตนุ\-ปสงฺกมติ, อุปสงฺกมิตฺวา ปญฺญตฺเต อาสเน นิสีทติ, ตเมนํ โส คหปติ วา คหปติปุตฺโต วา ปณีเตน ขาทนีเยน โภชนีเยน สหตฺถา สนฺตปฺเปติ สมฺปวาเรติ.}

\prose{ตสฺส เอวํ โหติ “สาธุ วต มฺยายํ คหปติ วา คหปติปุตฺโต วา ปณีเตน ขาทนีเยน โภชนีเยน สหตฺถา สนฺตปฺเปติ สมฺปวาเรตี”ติ. เอวมฺปิสฺส โหติ “อโห วต มายํ คหปติ วา คหปติปุตฺโต วา อายติมฺปิ เอวรูเปน ปณีเตน ขาทนีเยน โภชนีเยน สหตฺถา สนฺตปฺเปยฺย สมฺปวาเรยฺยา”ติ. โส ตํ ปิณฺฑปาตํ คถิโต\csromansharedfootnote{2e43a61a3e241a6d}{คธิโต (สฺยา, กํ, ก)} มุจฺฉิโต อชฺโฌสนฺโน\csromansharedfootnote{58ac98d7d755009d}{อชฺฌาปนฺโน (สี, ก) อชฺโฌปนฺโน (ฏีกา)} อนาทีนวทสฺสาวี อนิสฺสรณปญฺโญ ปริภุญฺชติ, โส ตตฺถ กามวิตกฺกมฺปิ วิตกฺเกติ, พฺยาปาท\-วิตกฺกมฺปิ วิตกฺเกติ, วิหิํสา\-วิตกฺกมฺปิ วิตกฺเกติ. เอวรูปสฺสาหํ ภิกฺขเว ภิกฺขุโน ทินฺนํ น มหปฺผลนฺติ วทามิ. ตํ กิสฺส เหตุ? ปมตฺโต หิ ภิกฺขเว ภิกฺขุ วิหรติ.}

\prose{อิธ ปน ภิกฺขเว ภิกฺขุ อญฺญตรํ คามํ วา นิคมํ วา อุปนิสฺสาย วิหรติ, ตเมนํ คหปติ วา คหปติปุตฺโต วา อุปสงฺกมิตฺวา สฺวาตนาย ภตฺเตน นิมนฺเตติ, อากงฺขมาโน ภิกฺขเว ภิกฺขุ อธิวาเสติ. โส ตสฺสา รตฺติยา อจฺจเยน ปุพฺพณฺหสมยํ นิวาเสตฺวา ปตฺต\-จีวรมา\-ทาย เยน ตสฺส คหปติสฺส วา คหปติ\-ปุตฺตสฺส วา นิเวสนํ เตนุ\-ปสงฺกมติ, อุปสณฺกมิตฺวา ปญฺญตฺเต อาสเน นิสีทติ, ตเมนํ โส คหปติ วา คหปติปุตฺโต วา ปณีเตน ขาทนีเยน โภชนีเยน สหตฺถา สนฺตปฺเปติ สมฺปวาเรติ.}

\prose{ตสฺส น เอวํ โหติ “สาธุ วต มฺยายํ คหปติ วา คหปติปุตฺโต วา ปณีเตน ขาทนีเยน โภชนีเยน สหตฺถา สนฺตปฺเปติ สมฺปวาเรตี”ติ. เอวมฺปิสฺส น โหติ “อโห วต มายํ คหปติ วา คหปติปุตฺโต วา อายติมฺปิ เอวรูเปน ปณีเตน ขาทนีเยน โภชนีเยน สหตฺถา สนฺตปฺเปยฺย สมฺปวาเรยฺยา”ติ. โส ตํ ปิณฺฑปาตํ อคถิโต อมุจฺฉิโต อนชฺโฌสนฺโน อาทีนวทสฺสาวี นิสฺสรณปญฺโญ ปริภุญฺชติ, โส ตตฺถ}

\vspace{\tipitakaparskip}
\csromancenter{(13) 3. กุสินารวคฺค เนกฺขมฺม\-วิตกฺกมฺปิ วิตกฺเกติ, อพฺยาปาทวิตกฺกมฺปิ วิตกฺเกติ, อวิหิํสา\-วิตกฺกมฺปิ วิตกฺเกติ. เอวรูปสฺสาหํ ภิกฺขเว ภิกฺขุโน ทินฺนํ มหปฺผลนฺติ วทามิ. ตํ กิสฺส เหตุ? อปฺปมตฺโต หิ ภิกฺขเว ภิกฺขุ วิหรตีติ.  ปฐมํ.}
\csromanmatikaanchor{mtk.23}
\csromantocmarkref{subsection}{2. ภณฺฑนสุตฺต}{mtk.23}
\bookmark[dest={mtk.23},level=2]{2. ภณฺฑนสุตฺต}
\csromanheader{2}{2. \textbf{ภณฺฑนสุตฺต}}
\proseitem{125}{\csromanfolio{279}ยสฺสํ ภิกฺขเว ทิสายํ ภิกฺขู ภณฺฑนชาตา กลหชาตา วิวาทาปนฺนา อญฺญมญฺญํ มุขสตฺตีหิ วิตุทนฺตา วิหรนฺติ. มนสิ กาตุมฺปิ เม เอสา ภิกฺขเว ทิสา น ผาสุ โหติ, ปเคว คนฺตุํ. นิฏฺฐเมตฺถ คจฺฉามิ\csromandash{}อทฺธา เต อายสฺมนฺโต ตโย ธมฺเม ปชหิํสุ, ตโย ธมฺเม พหุลมกํสุ\csromansharedfootnote{40fadb4bcaef5917}{พหุลีมกํสุ (สฺยา, กํ, อิ)}. กตเม ตโย ธมฺเม ปชหิํสุ? เนกฺขมฺม\-วิตกฺกํ อพฺยาปาท\-วิตกฺกํ อวิหิํสา\-วิตกฺกํ. อิเม ตโย ธมฺเม ปชหิํสุ. กตเม ตโย ธมฺเม พหุลมกํสุ? กามวิตกฺกํ พฺยาปาท\-วิตกฺกํ วิหิํสา\-วิตกฺกํ. อิเม ตโย ธมฺเม พหุลมกํสุ. ยสฺสํ ภิกฺขเว ทิสายํ ภิกฺขู ภณฺฑนชาตา กลหชาตา วิวาทาปนฺนา อญฺญมญฺญํ มุขสตฺตีหิ วิตุทนฺตา วิหรนฺติ. มนสิ กาตุมฺปิ เม เอสา ภิกฺขเว ทิสา น ผาสุ โหติ, ปเคว คนฺตุํ. นิฏฺฐเมตฺถ คจฺฉามิ “อทฺธา เต อายสฺมนฺโต อิเม ตโย ธมฺเม ปชหิํสุ, อิเม ตโย ธมฺเม พหุลมกํสุ”.}

\prose{ยสฺสํ ปน ภิกฺขเว ทิสายํ ภิกฺขู สมคฺคา สมฺโมทมานา อวิวทมานา ขีโรทกีภูตา อญฺญมญฺญํ ปิยจกฺขูหิ สมฺปสฺสานฺตา วิหรนฺติ. คนฺตุมฺปิ เม เอสา ภิกฺขเว ทิสา ผาสุ โหติ, ปเคว มนสิ กาตุํ. นิฏฺฐเมตฺถ คจฺฉามิ\csromandash{}อทฺธา เต อายสฺมนฺโต ตโย ธมฺเม ปชหิํสุ, ตโย ธมฺเม พหุลมกํสุ. กตเม ตโย ธมฺเม ปชหิํสุ? กามวิตกฺกํ พฺยาปาท\-วิตกฺกํ วิหิํสา\-วิตกฺกํ. อิเม ตโย ธมฺเม ปชหิํสุ. กตเม ตโย ธมฺเม พหุลมกํสุ? เนกฺขมฺม\-วิตกฺกํ อพฺยาปาท\-วิตกฺกํ อวิหิํสา\-วิตกฺกํ. อิเม ตโย ธมฺเม พหุลมกํสุ. ยสฺสํ ภิกฺขเว ทิสายํ ภิกฺขู สมคฺคา สมฺโมทมานา อวิวทมานา ขีโรทกีภูตา อญฺญมญฺญํ ปิยจกฺขูหิ สมฺปสฺสนฺตา วิหรนฺติ. คนฺตุมฺปิ เม เอสา ภิกฺขเว ทิสา ผาสุ โหติ, ปเคว มนสิ กาตุํ. นิฏฺฐเมตฺถ คจฺฉามิ “อทฺธา เต อายสฺมนฺโต อิเม ตโย ธมฺเม ปชหิํสุ, อิเม ตโย ธมฺเม พหุลมกํสู”ติ.  ทุติยํ.}

\csromanpalirecto
\csromanmatikaanchor{mtk.24}
\csromantocmarkref{subsection}{3. โคตมกเจติยสุตฺต}{mtk.24}
\bookmark[dest={mtk.24},level=2]{3. โคตมกเจติยสุตฺต}
\csromanmarkh{ติกนิปาตปาฬิ}\csromanmarkhii{3. โคตมก\-เจติย\-สุตฺต}\csromanheadernomark{2}{ติกนิปาตปาฬิ 3. โคตมก\-เจติย\-สุตฺต}
\proseitem{126}{\csromanfolio{280}เอกํ สมยํ ภควา เวสาลิยํ วิหรติ โคตมเก เจติเย. ตตฺร โข ภควา ภิกฺขู อามนฺเตสิ “ภิกฺขโว”ติ. “ภทนฺเต”ติ เต ภิกฺขู ภควโต ปจฺจสฺโสสุํ. ภควา เอตทโวจ\csromandash{}}

\prose{“อภิญฺญายาหํ ภิกฺขเว ธมฺมํ เทเสมิ โน อนภิญฺญาย, สนิทานาหํ ภิกฺขเว ธมฺมํ เทเสมิ โน อนิทานํ, สปฺปาฏิหาริยาหํ ภิกฺขเว ธมฺมํ เทเสมิ โน อปฺปาฏิหาริยํ. ตสฺส มยฺหํ ภิกฺขเว อภิญฺญาย ธมฺมํ เทสยโต โน อนภิญฺญาย, สนิทานํ ธมฺมํ เทสยโต โน อนิทานํ, สปฺปาฏิหาริยํ ธมฺมํ เทสยโต โน อปฺปาฏิหาริยํ กรณีโย โอวาโท กรณียา อนุสาสนี, อลญฺจ ปน โว ภิกฺขเว ตุฏฺฐิยา อลํ อตฺตมนตาย อลํ โสมนสฺสาย “สมฺมาสมฺพุทฺโธ ภควา, สฺวากฺขาโต ธมฺโม, สุปฺปฏิปนฺโน สํโฆ”ติ.}

\prose{อิทมโวจ ภควา, อตฺตมนา เต ภิกฺขู ภควโต ภาสิตํ อภินนฺทุนฺติ. อิมสฺมิํ จ ปน เวยฺยากรณสฺมิํ ภญฺญมาเน สหสฺสี โลกธาตุ อกมฺปิตฺถาติ.  ตติยํ.}

\csromanmatikaanchor{mtk.25}
\csromantocmarkref{subsection}{4. ภรณฺฑุกาลามสุตฺต}{mtk.25}
\bookmark[dest={mtk.25},level=2]{4. ภรณฺฑุกาลามสุตฺต}
\csromanheader{2}{4. ภรณฺฑุ\-กาลาม\-สุตฺต}
\proseitem{127}{เอกํ สมยํ ภควา โกสเลสุ จาริกํ จรมาโน เยน กปิลวตฺถุ ตทวสริ. อสฺโสสิ โข มหานาโม สกฺโก “ภควา กิร กปิลวตฺถุํ อนุปฺปตฺโต”ติ. อถ โข มหานาโม สกฺโก เยน ภควา เตนุปสงฺกมิ, อุปสงฺกมิตฺวา ภควนฺตํ อภิวาเทตฺวา เอกมนฺตํ อฏฺฐาสิ. เอกมนฺตํ ฐิตํ โข มหานามํ สกฺกํ ภควา เอตทโวจ\csromandash{}}

\prose{“คจฺฉ มหานาม, กปิล\-วตฺถุสฺมิํ ตถารูปํ อาวสถํ ชาน, ยตฺถชฺช มยํ เอกรตฺติํ วิหเรยฺยามา”ติ. “เอวํ ภนฺเต”ติ โข มหานาโม สกฺโก ภควโต ปฏิสฺสุตฺวา กปิลวตฺถุํ ปวิสิตฺวา เกวลกปฺปํ กปิลวตฺถุํ อนฺวาหิณฺฑนฺโต\csromansharedfootnote{041769e68ba32abe}{อาหิณฺฑนฺโต (สฺยา, กํ)} นาทฺทส กปิล\-วตฺถุสฺมิํ ตถารูปํ อาวสถํ, ยตฺถชฺช ภควา เอกรตฺติํ วิหเรยฺย.}

\prose{อถ โข มหานาโม สกฺโก เยน ภควา เตนุปสงฺกมิ, อุปสงฺกมิตฺวา ภควนฺตํ เอตทโวจ “นตฺถิ ภนฺเต กปิล\-วตฺถุสฺมิํ ตถารูโป}

\vspace{\tipitakaparskip}
\csromancenter{(13) 3. กุสินารวคฺค อาวสโถ, ยตฺถชฺช ภควา เอกรตฺติํ วิหเรยฺย. อยํ ภนฺเต ภรณฺฑุ กาลาโม ภควโต ปุราณ\-สพฺรหฺมจารี, ตสฺสชฺช ภควา อสฺสเม เอกรตฺติํ วิหรตู”ติ, คจฺฉ มหานาม สนฺถรํ ปญฺญาเปหีติ. “เอวํ ภนฺเต”ติ โข มหานาโม สกฺโก ภควโต ปฏิสฺสุตฺวา เยน ภรณฺฑุสฺส กาลามสฺส อสฺสโม เตนุปสงฺกมิ, อุปสงฺกมิตฺวา สนฺถรํ ปญฺญาเปตฺวา อุทกํ ฐเปตฺวา ปาทานํ โธวนาย เยน ภควา เตนุปสงฺกมิ, อุปสงฺกมิตฺวา ภควนฺตํ เอตทโวจ “สนฺถโต ภนฺเต สนฺถาโร อุทกํ ฐปิตํ ปาทานํ โธวนาย, ยสฺสทานิ ภนฺเต ภควา กาลํ มญฺญตี”ติ.}
\prose{\csromanfolio{281}อถ โข ภควา เยน ภรณฺฑุสฺส กาลามสฺส อสฺสโม เตนุปสงฺกมิ, อุปสงฺกมิตฺวา ปญฺญตฺเต อาสเน นิสีทิ, นิสชฺช โข ภควา ปาเท ปกฺขาเลสิ. อถ โข มหานามสฺส สกฺกสฺส เอตทโหสิ “อกาโล โข อชฺช ภควนฺตํ ปยิรุปาสิตุํ, กิลนฺโต ภควา, เสฺว ทานาหํ ภควนฺตํ ปยิรุปาสิสฺสามี”ติ ภควนฺตํ อภิวาเทตฺวา ปทกฺขิณํ กตฺวา ปกฺกามิ.}

\prose{อถ โข มหานาโม สกฺโก ตสฺสา รตฺติยา อจฺจเยน เยน ภควา เตนุปสงฺกมิ, อุปสงฺกมิตฺวา เอกมนฺตํ นิสีทิ, เอกมนฺตํ นิสินฺนํ โข มหานามํ สกฺกํ ภควา เอตทโวจ “ตโย โขเม มหานาม สตฺถาโร สนฺโต สํวิชฺชมานา โลกสฺมิํ. กตเม ตโย? อิธ มหานาม เอกจฺโจ สตฺถา กามานํ ปริญฺญํ ปญฺญาเปติ, น รูปานํ ปริญฺญํ ปญฺญาเปติ, น เวทนานํ ปริญฺญํ ปญฺญาเปติ. อิธ ปน มหานาม เอกจฺโจ สตฺถา กามานํ ปริญฺญํ ปญฺญาเปติ, รูปานํ ปริญฺญํ ปญฺญาเปติ, น เวทนานํ ปริญฺญํ ปญฺญาเปติ. อิธ ปน มหานาม เอกจฺโจ สตฺถา กามานํ ปริญฺญํ ปญฺญาเปติ, รูปานํ ปริญฺญํ ปญฺญาเปติ, เวทนานํ ปริญฺญํ ปญฺญาเปติ. อิเม โข มหานาม ตโย สตฺถาโร สนฺโต สํวิชฺชามานา โลกสฺมิํ. อิเมสํ มหานาม ติณฺณํ สตฺถารานํ เอกา นิฏฺฐา อุทาหุ ปุถุ นิฏฺฐาติ.}

\prose{เอวํ วุตฺเต ภรณฺฑุ กาลาโม มหานามํ สกฺกํ เอตทโวจ “เอกาติ มหานาม วเทหี”ติ. เอวํ วุตฺเต ภควา มหานามํ สกฺกํ เอตทโวจ “นานาติ มหานาม วเทหี”ติ. ทุติยมฺปิ โข ภรณฺฑุ กาลาโม มหานามํ สกฺกํ เอตทโวจ “เอกาติ มหานาม วเทหี”ติ. ทุติยมฺปิ โข ภควา มหานามํ สกฺกํ เอตทโวจ “นานาติ มหานาม วเทหี”ติ. ตติยมฺปิ โข ภรณฺฑุ กาลาโม มหานามํ สกฺกํ เอตทโวจ “เอกาติ \csromanfolio{282}มหานาม วเทหี”ติ. ตติยมฺปิ โข ภควา มหานามํ สกฺกํ เอตทโวจ “นานาติ มหานาม วเทหี”ติ.}

\prose{อถ โข ภรณฺฑุ กาลามสฺส เอตทโหสิ “มเหสกฺขสฺส วตมฺหิ มหานามสฺส สกฺกสฺส สมฺมุขา สมเณน โคตเมน ยาวตติยํ อปสาทิโต, ยํนูนาหํ กปิล\-วตฺถุมฺหา ปกฺกเมยฺยนฺ”ติ. อถ โข ภรณฺฑุ กาลาโม กปิล\-วตฺถุมฺหา ปกฺกามิ. ยํ กปิล\-วตฺถุมฺหา ปกฺกามิ, ตถา ปกฺกนฺโตว อโหสิ, น ปุน ปจฺจาคจฺฉีติ.  จตุตฺถํ.}

\csromanmatikaanchor{mtk.26}
\csromantocmarkref{subsection}{5. หตฺถกสุตฺต}{mtk.26}
\bookmark[dest={mtk.26},level=2]{5. หตฺถกสุตฺต}
\csromanheader{2}{5. \textbf{หตฺถกสุตฺต}}
\proseitem{128}{เอกํ สมยํ ภควา สาวตฺถิยํ วิหรติ เชตวเน อนาถ\-ปิณฺฑิกสฺส อาราเม. อถ โข หตฺถโก เทวปุตฺโต อภิกฺกนฺตาย รตฺติยา อภิกฺกนฺต\-วณฺโณ เกวลกปฺปํ เชตวนํ โอภาเสตฺวา เยน ภควา เตนุปสงฺกมิ, อุปสงฺกมิตฺวา ภควโต ปุรโต “ฐสฺสามี”ติ โอสีทติเมว สํสีทติเมว\csromansharedfootnote{770bb9239f2c9cbc}{โอสีทติ เจว สํสีทติ จ (สี, อิ), โอสีทติ สํสีทติ (สฺยา, กํ)}, น สกฺโกติ สณฺฐาตุํ. เสยฺยถาปิ นาม สปฺปิ วา เตลํ วา วาลุกาย อาสิตฺตํ โอสีทติเมว สํสีทติเมว น สณฺฐาติ. เอวเมวํ หตฺถโก เทวปุตฺโต ภควโต ปุรโต “ฐสฺสามี”ติ โอสีทติเมว สํสีทติเมว, น สกฺโกติ สณฺฐาตุํ.}

\prose{อถ โข ภควา หตฺถกํ เทวปุตฺตํ เอตทโวจ “โอฬาริกํ หตฺถก อตฺตภาวํ อภินิมฺมินาหี”ติ. “เอวํ ภนฺเต”ติ โข หตฺถโก เทวปุตฺโต ภควโต ปฏิสฺสุตฺวา โอฬาริกํ อตฺตภาวํ อภินิมฺมินิตฺวา ภควนฺตํ อภิวาเทตฺวา เอกมนฺตํ อฏฺฐาสิ. เอกมนฺตํ ฐิตํ โข หตฺถกํ เทวปุตฺตํ ภควา เอตทโวจ\csromandash{}}

\prose{“เย เต หตฺถก ธมฺมา ปุพฺเพ มนุสฺส\-ภูตสฺส ปวตฺติโน อเหสุํ, อปิ นุ เต เต ธมฺมา เอตรหิ ปวตฺติโน”ติ. เย จ เม ภนฺเต ธมฺมา ปุพฺเพ มนุสฺส\-ภูตสฺส ปวตฺติโน อเหสุํ, เต จ เม ธมฺมา เอตรหิ ปวตฺติโน. เย จ เม ภนฺเต ธมฺมา ปุพฺเพ มนุสฺส\-ภูตสฺส นปฺปวตฺติโน อเหสุํ, เต จ เม ธมฺมา เอตรหิ ปวตฺติโน. เสยฺยถาปิ ภนฺเต ภควา เอตรหิ อากิณฺโณ วิหรติ ภิกฺขูหิ ภิกฺขุนีหิ อุปาสเกหิ อุปาสิกาหิ}

\vspace{\tipitakaparskip}
\csromancenter{(13) 3. กุสินารวคฺค ราชูหิ ราช\-มหามตฺเตหิ ติตฺถิเยหิ ติตฺถิย\-สาวเกหิ. เอวเมวํ โข อหํ ภนฺเต อากิณฺโณ วิหรามิ เทวปุตฺเตหิ. ทูรโตปิ ภนฺเต เทวปุตฺตา อาคจฺฉนฺติ “หตฺถกสฺส เทวปุตฺตสฺส สนฺติเก ธมฺมํ โสสฺสามา”ติ. ติณฺณาหํ ภนฺเต ธมฺมานํ อติตฺโต อปฺปฏิวาโน กาลงฺกโต. กตเมสํ ติณฺณํ? ภควโต อหํ ภนฺเต ทสฺสนสฺส อติตฺโต อปฺปฏิวาโน กาลงฺกโต, สทฺธมฺมสวนสฺสาหํ ภนฺเต อติตฺโต อปฺปฏิวาโน กาลงฺกโต, สํฆสฺสาหํ ภนฺเต อุปฏฺฐานสฺส อติตฺโต อปฺปฏิวาโน กาลงฺกโต. อิเมสํ โข อหํ ภนฺเต ติณฺณํ ธมฺมานํ อติตฺโต อปฺปฏิวาโน กาลงฺกโตติ.}
\setlength{\csromangathapagewidth}{0pt}
\setlength{\csromangathawakwidth}{0pt}
\csromangathameasuregroup{นาหํ ภควโต ทสฺสนสฺส, \\ สํฆสฺส อุปฏฺฐานสฺส, \\ อธิสีลํ สิกฺขมาโน, \\ ติณฺณํ ธมฺมานํ อติตฺโต,}{\csromangathabat{นาหํ ภควโต ทสฺสนสฺส,}{ติตฺติมชฺฌคา\textsuperscript{1} กุทาจนํ.} \\ \csromangathabat{สํฆสฺส อุปฏฺฐานสฺส,}{สทฺธมฺมสวนสฺส จ.} \\ \csromangathabat{อธิสีลํ สิกฺขมาโน,}{สทฺธมฺมสวเน รโต.} \\ \csromangathabat{ติณฺณํ ธมฺมานํ อติตฺโต,}{หตฺถโก อวิหํ คโตติ.} \\ ปญฺจมํ.}
\csromangathameasurewak{ปญฺจมํ.}
\csromangathasetleft{นาหํ ภควโต ทสฺสนสฺส, \\ สํฆสฺส อุปฏฺฐานสฺส, \\ อธิสีลํ สิกฺขมาโน, \\ ติณฺณํ ธมฺมานํ อติตฺโต,}
\csromangathagroup{\csromangathastanza{\csromanfolio{283}\csromangathabat{นาหํ ภควโต ทสฺสนสฺส,}{ติตฺติมชฺฌคา\csromansharedfootnote{4309f0ece85625f1}{ติตฺติ ติตฺติสมฺภวํ (ก)} กุทาจนํ.} \\ \csromangathabat{สํฆสฺส อุปฏฺฐานสฺส,}{สทฺธมฺมสวนสฺส จ.}}\csromangathastanza{\csromangathabat{อธิสีลํ สิกฺขมาโน,}{สทฺธมฺม\-สวเน รโต.} \\ \csromangathabat{ติณฺณํ ธมฺมานํ อติตฺโต,}{หตฺถโก อวิหํ คโตติ.}}\csromangathastanza{\csromangathawak{ปญฺจมํ.}}}

\csromanmatikaanchor{mtk.27}
\csromantocmarkref{subsection}{6. กฏุวิยสุตฺต}{mtk.27}
\bookmark[dest={mtk.27},level=2]{6. กฏุวิยสุตฺต}
\csromanheader{2}{6. \textbf{กฏุวิยสุตฺต}}
\proseitem{129}{เอกํ สมยํ ภควา พาราณสิยํ วิหรติ อิสิปตเน มิคทาเย. อถ โข ภควา ปุพฺพณฺหสมยํ นิวาเสตฺวา ปตฺต\-จีวรมา\-ทาย พาราณสิํ ปิณฺฑาย ปาวิสิ. อทฺทสา โข ภควา โคโยค\-ปิลกฺขสฺมิํ\csromansharedfootnote{d9847a8d6ddd6b5d}{โคโยคมิลกฺขสฺมิํ (สฺยา, กํ, ก)} ปิณฺฑาย จรมาโน\csromansharedfootnote{5ebb8c4d8c2e63bf}{จรมานํ(ก)} อญฺญตรํ ภิกฺขุํ ริตฺตสฺสาทํ พาหิรสฺสาทํ มุฏฺฐสฺสติํ อสมฺปชานํ อสมาหิตํ วิพฺภนฺต\-จิตฺตํ ปากตินฺทฺริยํ, ทิสฺวา ตํ ภิกฺขุํ เอตทโวจ\csromandash{}}

\csromanmatikaanchor{mtk.28}
\csromantocmarkref{subsection}{7-8. อนุรุทฺธสุตฺตานิ}{mtk.28}
\bookmark[dest={mtk.28},level=2]{7-8. อนุรุทฺธสุตฺตานิ}
\prose{“มา โข ตฺวํ ภิกฺขุ อตฺตานํ กฏุวิยมกาสิ, ตํ วต ภิกฺขุ กฏุวิยกตํ อตฺตานํ อามคนฺเธน\csromansharedfootnote{9f10dfe75c623c01}{อามคนฺเธ (สี, สฺยา, กํ, อิ)} อวสฺสุตํ มกฺขิกา นานุปติสฺสนฺติ นานฺวาสฺสวิสฺสนฺตีติ\csromansharedfootnote{0bacba1a8757c2dc}{นานุพนฺธิสฺสนฺติ (ก)} เนตํ ฐานํ วิชฺชตี”ติ. อถ โข โส ภิกฺขุ ภควตา อิมินา โอวาเทน โอวทิโต สํเวคมาปาทิ. อถ โข ภควา พาราณสิยํ ปิณฺฑาย จริตฺวา ปจฺฉาภตฺตํ ปิณฺฑปาต\-ปฏิกฺกนฺโต ภิกฺขู อามนฺเตสิ\csromandash{} “อิธาหํ ภิกฺขเว ปุพฺพณฺหสมยํ นิวาเสตฺวา ปตฺต\-จีวรมา\-ทาย \csromanfolio{284}พาราณสิํ ปณฺฑาย ปาวิสิํ, อทฺทสํ โข อหํ ภิกฺขเว โคโยค\-ปิลกฺขสฺมิํ ปิณฺฑาย จรมาโน อญฺญตรํ ภิกฺขุ ริตฺตสฺสาทํ พาหิรสฺสาทํ มุฏฺฐสฺสติํ อสมฺปชานํ อสมาหิตํ วิพฺภนฺต\-จิตฺตํ ปากตินฺทฺริยํ, ทิสฺวา ตํ ภิกฺขุ เอตทโวจํ\csromandash{}}

\prose{‘มา โข ตฺวํ ภิกฺขุ อตฺตานํ กฏุวิยมกาสิ, ตํ วต ภิกฺขุ กฏุวิยกตํ อตฺตานํ อามคนฺเธน อวสฺสุตํ มกฺขิกา นานุปติสฺสนฺติ นานฺวาสฺสวิสฺสนฺตีติ เนตํ ฐานํ วิชฺชตี’ติ. อถ โข ภิกฺขเว โส ภิกฺขุ มยา อิมินา โอวาเทน โอวทิโต สํเวคมาปาที”ติ. เอวํ วุตฺเต อญฺญตโร ภิกฺขุ ภควนฺตํ เอตทโวจ “กิํ นุ โข ภนฺเต กฏุวิยํ, โก อามคนฺโธ, กา มกฺขิกา”ติ.}

\prose{อภิชฺฌา โข ภิกฺขุ กฏุวิยํ, พฺยาปาโท อามคนฺโธ, ปาปกา อกุสลา วิตกฺกา มกฺขิกา, ตํ วต ภิกฺขุ กฏุวิยกตํ อตฺตานํ อามคนฺเธน อวสฺสุตํ มกฺขิกา นานุปติสฺสนฺติ นานฺวาสฺสวิสฺสนฺตีติ เนตํ ฐานํ วิชฺชตีติ.}

\setlength{\csromangathapagewidth}{0pt}
\setlength{\csromangathawakwidth}{0pt}
\csromangathameasuregroup{อคุตฺตํ จกฺขุโสตสฺมิํ, \\ มกฺขิกานุปติสฺสนฺติ, \\ กฏุวิยกโต ภิกฺขุ, \\ อารกา โหติ นิพฺพานา, \\ คาเม วา ยทิ วารญฺเญ, \\ ปเรติ\textsuperscript{1} พาโล ทุมฺเมโธ, \\ เย จ สีเลน สมฺปนฺนา, \\ อุปสนฺตา สุขํ เสนฺติ,}{\csromangathabat{อคุตฺตํ จกฺขุโสตสฺมิํ,}{อินฺทฺริเยสุ อสํวุตํ.} \\ \csromangathabat{มกฺขิกานุปติสฺสนฺติ,}{สงฺกปฺปา ราคนิสฺสิตา.} \\ \csromangathabat{กฏุวิยกโต ภิกฺขุ,}{อามคนฺเธ อวสฺสุโต.} \\ \csromangathabat{อารกา โหติ นิพฺพานา,}{วิฆาตสฺเสว ภาควา.} \\ \csromangathabat{คาเม วา ยทิ วารญฺเญ,}{อลทฺธา สมถมตฺตโน\textsuperscript{1}.} \\ \csromangathabat{ปเรติ\textsuperscript{1} พาโล ทุมฺเมโธ,}{มกฺขิกาหิ ปุรกฺขโต.} \\ \csromangathabat{เย จ สีเลน สมฺปนฺนา,}{ปญฺญายูปสเมรตา.} \\ \csromangathabat{อุปสนฺตา สุขํ เสนฺติ,}{นาสยิตฺวาน มกฺขิกาติ.} \\ ฉฏฺฐํ.}
\csromangathameasurewak{ฉฏฺฐํ.}
\csromangathasetleft{อคุตฺตํ จกฺขุโสตสฺมิํ, \\ มกฺขิกานุปติสฺสนฺติ, \\ กฏุวิยกโต ภิกฺขุ, \\ อารกา โหติ นิพฺพานา, \\ คาเม วา ยทิ วารญฺเญ, \\ ปเรติ\textsuperscript{1} พาโล ทุมฺเมโธ, \\ เย จ สีเลน สมฺปนฺนา, \\ อุปสนฺตา สุขํ เสนฺติ,}
\csromangathagroup{\csromangathastanza{\csromangathabat{อคุตฺตํ จกฺขุโสตสฺมิํ,}{อินฺทฺริเยสุ อสํวุตํ.} \\ \csromangathabat{มกฺขิกานุปติสฺสนฺติ,}{สงฺกปฺปา ราคนิสฺสิตา.}}\csromangathastanza{\csromangathabat{กฏุวิยกโต ภิกฺขุ,}{อามคนฺเธ อวสฺสุโต.} \\ \csromangathabat{อารกา โหติ นิพฺพานา,}{วิฆาตสฺเสว ภาควา.}}\csromangathastanza{\csromangathabat{คาเม วา ยทิ วารญฺเญ,}{อลทฺธา สมถมตฺตโน\csromansharedfootnote{77388f08591897eb}{สมมตฺตโน (สี, สฺยา, กํ), สมฺมมตฺตโน (อิ)}.} \\ \csromangathabat{ปเรติ\csromansharedfootnote{77388f08591897eb}{สมมตฺตโน (สี, สฺยา, กํ), สมฺมมตฺตโน (อิ)} พาโล ทุมฺเมโธ,}{มกฺขิกาหิ ปุรกฺขโต.}}\csromangathastanza{\csromangathabat{เย จ สีเลน สมฺปนฺนา,}{ปญฺญายูปสเมรตา.} \\ \csromangathabat{อุปสนฺตา สุขํ เสนฺติ,}{นาสยิตฺวาน มกฺขิกาติ.}}\csromangathastanza{\csromangathawak{ฉฏฺฐํ.}}}

\prose{7. ปฐม\-อนุรุทฺธ\-สุตฺต}

\proseitem{130}{อถ โข อายสฺมา อนุรุทฺโธ เยน ภควา เตนุปสงฺกมิ, อุปสงฺกมิตฺวา ภควนฺตํ อภิวาเทตฺวา เอกมนฺตํ นิสีทิ, เอกมนฺตํ นิสินฺโน โข อายสฺมา อนุรุทฺโธ ภควนฺตํ เอตทโวจ “อิธาหํ ภนฺเต ทิพฺเพน}

\vspace{\tipitakaparskip}
\csromancenter{(13) 3. กุสินารวคฺค จกฺขุนา วิสุทฺเธน อติกฺกนฺต\-มานุสเกน เยภุยฺเยน ปสฺสามิ มาตุคามํ กายสฺส เภทา ปรํ มรณา อปายํ ทุคฺคติํ วินิปาตํ นิรยํ อุปปชฺชมานํ. กติหิ นุ โข ภนฺเต ธมฺเมหิ สมนฺนาคโต มาตุคาโม กายสฺส เภทา ปรํ มรณา อปายํ ทุคฺคติํ วินิปาตํ นิรยํ อุปปชฺชตี”ติ.}
\prose{\csromanfolio{285}ตีหิ โข อนุรุทฺธ ธมฺเมหิ สมนฺนาคโต มาตุคาโม กายสฺส เภทา ปรํ มรณา อปายํ ทุคฺคติํ วินิปาตํ นิรยํ อุปปชฺชติ. กตเมหิ ตีหิ? อิธ อนุรุทฺธ มาตุคาโม ปุพฺพณฺหสมยํ มจฺเฉรมล\-ปริยุฏฺฐิเตน เจตสา อคารํ อชฺฌาวสติ, มชฺฌนฺหิกสมยํ อิสฺสา\-ปริยุฏฺฐิเตน เจตสา อคารํ อชฺฌาวสติ, สายนฺหสมยํ กาม\-ราค\-ปริยุฏฺฐิเตน เจตสา อคารํ อชฺฌาวสติ. อิเมหิ โข อนุรุทฺธ ตีหิ ธมฺเมหิ สมนฺนาคโต มาตุคาโม กยสฺส เภทา ปรํ มรณา อปายํ ทุคฺคติํ วินิปาตํ นิรยํ อุปปชฺชตีติ.  }

\prose{8. ทุติย\-อนุรุทฺธ\-สุตฺต}

\proseitem{131}{อถ โข อายสฺมา อนุรุทฺโธ เยนายสฺมา สาริปุตฺโต เตนุปสงฺกมิ, อุปสงฺกมิตฺวา อายสฺมตา สาริปุตฺเตน สทฺธิํ สมฺโมทิ, สมฺโมทนียํ กถํ สารณียํ วีติสาเรตฺวา เอกมนฺตํ นิสีทิ, เอกมนฺตํ นิสินฺโน โข อายสฺมา อนุรุทฺโธ อายสฺมนฺตํ สาริปุตฺตํ เอตทโวจ “อิธาหํ อาวุโส สาริปุตฺต ทิพฺเพน จกฺขุนา วิสุทฺเธน อติกฺกนฺต\-มานุสเกน สหสฺสํ โลกํ โอโลเกมิ, อารทฺธํ โข ปน เม วีริยํ อสลฺลีนํ, อุปฏฺฐิตา สติ อสมฺมุฏฺฐา\csromansharedfootnote{0777691e8e31d9c6}{อปมฺมุฏฺฐา (สี), อปมุฏฺฐา (สฺยา, กํ)}, ปสฺสทฺโธ กาโย อสารทฺโธ, สมาหิตํ จิตฺตํ เอกคฺคํ, อถ จ ปน เม นานุปาทาย\csromansharedfootnote{2b812e690c007f71}{น อนุปาทาย (สี, สฺยา, กํ, อิ)} อาสเวหิ จิตฺตํ วิมุจฺจตี”ติ.}

\prose{ยํ โข เต อาวุโส อนุรุทฺธ เอวํ โหติ “อหํ ทิพฺเพน จกฺขุนา วิสุทฺเธน อติกฺกนฺต\-มานุสเกน สหสฺสํ โลกํ โวโลเกมี”ติ, อิทํ เต มานสฺมิํ. ยมฺปิ เต อาวุโส อนุรุทฺธ เอวํ โหติ “อารทฺธํ โข ปน เม วีริยํ อสลฺลีนํ, อุปฏฺฐิตา สติ อสมฺมุฏฺฐา, ปสฺสทฺโธ กาโย อสารทฺโธ, สมาหิตํ จิตฺตํ เอกคฺคนฺ”ติ, อิทํ เต อุทฺธจฺจสฺมิํ. ยมฺปิ เต อาวุโส อนุรุทฺธ เอวํ โหติ “อถ จ ปน เม นานุปาทาย อาสเวหิ จิตฺตํ วิมุจฺจตี”ติ, \csromanfolio{286}อิทํ เต กุกฺกุจฺจสฺมิํ. สาธุ วตายสฺมา อนุรุทฺโธ อิเม ตโย ธมฺเม ปหาย อิเม ตโย ธมฺเม อมนสิกริตฺวา อมตาย ธาตุยา จิตฺตํ อุปสํหรตูติ.}

\prose{อถ โข อายสฺมา อนุรุทฺโธ อปเรน สมเยน อิเม ตโย ธมฺเม ปหาย อิเม ตโย ธมฺเม อมนสิกริตฺวา อมตาย ธาตุยา จิตฺตํ อุปสํหริ\csromansharedfootnote{4d1bb7a09b32ac1c}{อุปสํหาสิ (สฺยา, กํ, อิ), อุปสํหรติ (ก)}. อถ โข อายสฺมา อนุรุทฺโธ เอโก วูปกฏฺโฐ อปฺปมตฺโต อาตาปี ปหิตตฺโต วิหรนฺโต นจิรสฺเสว ยสฺสตฺถาย กุลปุตฺตา สมฺมเทว อคารสฺมา อนคาริยํ ปพฺพชนฺติ, ตทนุตฺตรํ พฺรหฺมจริย\-ปริโยสานํ ทิฏฺเฐว ธมฺเม สยํ อภิญฺญา สจฺฉิกตฺวา อุปสมฺปชฺช วิหาสิ, “ขีณา ชาติ, วุสิตํ พฺรหฺมจริยํ, กตํ กรณียํ, นาปรํ อิตฺถตฺตายา”ติ อพฺภญฺญาสิ. อญฺญตโร จ ปนายสฺมา อนุรุทฺโธ อรหตํ อโหสีติ.  อฏฺฐมํ.}

\csromanmatikaanchor{mtk.29}
\csromantocmarkref{subsection}{9. ปฏิจฺฉนฺนสุตฺต}{mtk.29}
\bookmark[dest={mtk.29},level=2]{9. ปฏิจฺฉนฺนสุตฺต}
\csromanheader{2}{9. ปฏิจฺฉนฺน\-สุตฺต}
\proseitem{132}{ตีณิมานิ ภิกฺขเว ปฏิจฺฉนฺนานิ อาวหนฺติ\csromansharedfootnote{774c41d2d120a91a}{วหนฺติ (สี, สฺยา, กํ, อิ)} โน วิวฏานิ. กตมานิ ตีณิ? มาตุคาโม ภิกฺขเว ปฏิจฺฉนฺโน อาวหติ โน วิวโฏ, พฺราหฺมณานํ ภิกฺขเว มนฺตา ปฏิจฺฉนฺนา อาวหนฺติ โน วิวฏา, มิจฺฉาทิฏฺฐิ ภิกฺขเว ปฏิจฺฉนฺนา อาวหติ โน วิวฏา. อิมานิ โข ภิกฺขเว ตีณิ ปฏิจฺฉนฺนานิ อาวหนฺติ โน วิวฏานิ.}

\prose{ตีณิมานิ ภิกฺขเว วิวฏานิ วิโรจนฺติ โน ปฏิจฺฉนฺนานิ. กตมานิ ตีณิ? จนฺทมณฺฑลํ ภิกฺขเว วิวฏํ วิโรจติ โน ปฏิจฺฉนฺนํ, สูริยมณฺฑลํ ภิกฺขเว วิวฏํ วิโรจติ โน ปฏิจฺฉนฺนํ, ตถาคต\-ปฺปเวทิโต ธมฺมวินโย ภิกฺขเว วิวโฏ วิโรจติ โน ปฏิจฺฉนฺโน. อิมานิ โข ภิกฺขเว ตีณิ วิวฏานิ วิโรจนฺติ โน ปฏิจฺฉนฺนานีติ.  นวมํ.}

\csromanmatikaanchor{mtk.30}
\csromantocmarkref{subsection}{10. เลขสุตฺต}{mtk.30}
\bookmark[dest={mtk.30},level=2]{10. เลขสุตฺต}
\csromanheader{2}{10. \textbf{เลขสุตฺต}}
\proseitem{133}{ตโยเม ภิกฺขเว ปุคฺคลา สนฺโต สํวิชฺชมานา โลกสฺมิํ. กตเม ตโย? ปาสาณ\-เลขู\-ปโม ปุคฺคโล, ปถวิ\-เลขู\-ปโม ปุคฺคโล, อุทกเลขูปโม ปุคฺคโล. กตโม จ ภิกฺขเว ปาสาณ\-เลขู\-ปโม ปุคฺคโล? อิธ ภิกฺขเว เอกจฺโจ ปุคฺคโล อภิณฺหํ กุชฺฌติ, โส จ ขฺวสฺส}

\vspace{\tipitakaparskip}
\csromancenter{(14) 4. โยธาชีววคฺค โกโธ ทีฆรตฺตํ อนุเสติ. เสยฺยถาปิ ภิกฺขเว ปาสาเณ เลขา น ขิปฺปํ ลุชฺชติ วาเตน วา อุทเกน วา, จิรฏฺฐิติกา โหติ. เอวเมวํ โข ภิกฺขเว อิเธกจฺโจ ปุคฺคโล อภิณฺหํ กุชฺชติ, โส จ ขฺวสฺส โกโธ ทีฆรตฺตํ อนุเสติ. อยํ วุจฺจติ ภิกฺขเว ปาสาณ\-เลขู\-ปโม ปุคฺคโล.}
\prose{\csromanfolio{287}กตโม จ ภิกฺขเว ปถวิ\-เลขู\-ปโม ปุคฺคโล? อิธ ภิกฺขเว เอกจฺโจ ปุคฺคโล อภิณฺหํ กุชฺฌติ, โส จ ขฺวสฺส โกโธ น ทีฆรตฺตํ อนุเสติ. เสยฺยถาปิ ภิกฺขเว ปถวิยา เลขา ขิปฺปํ ลุชฺชติ วาเตน วา อุทเกน วา, น จิรฏฺฐิติกา โหติ, เอวเมวํ โข ภิกฺขเว อิเธกจฺโจ ปุคฺคโล อภิณฺหํ กุชฺฌติ, โส จ ขฺวสฺส โกโธ น ทีฆรตฺตํ อนุเสติ. อยํ วุจฺจติ ภิกฺขเว ปถวิ\-เลขู\-ปโม ปุคฺคโล.}

\prosewithclosermidruled{กตโม จ ภิกฺขเว อุทกเลขูปโม ปุคฺคโล? อิธ ภิกฺขเว เอกจฺโจ ปุคฺคโล อาคาเฬฺหนปิ วุจฺจมาโน ผรุเสนปิ วุจฺจมาโน อมนาเปนปิ วุจฺจมาโน สนฺธิยติเมว\csromansharedfootnote{6bd63ff728725a30}{...เยว (สฺยา, กํ), ...เจว (อิ)} สํสนฺทติเมว\csromansharedfootnote{5c9b32c0a4e725ac}{อภิ 3. 137 ปิฏฺเฐปิ.} สมฺโมทติเมว\csromansharedfootnote{6bd63ff728725a30}{...เยว (สฺยา, กํ), ...เจว (อิ)}. เสยฺยถาปิ ภิกฺขเว อุทเก เลขา ขิปฺปํเยว ปฏิวิคจฺฉติ, น จิรฏฺฐิติกา โหติ. เอวเมวํ โข ภิกฺขเว อิเธกจฺโจ ปุคฺคโล อาคาเฬฺหนปิ วุจฺจมาโน ผรุเสนปิ วุจฺจมาโน อมนาเปนปิ วุจฺจมาโน สนฺธิยติเมว สํสนฺทติเมว สมฺโมทติเมว. อยํ วุจฺจติ ภิกฺขเว อุทกเลขูปโม ปุคฺคโล. อิเม โข ภิกฺขเว ตโย ปุคฺคลา สนฺโต สํวิชฺชมานา โลกสฺมินฺติ\csromansharedfootnote{6bd63ff728725a30}{...เยว (สฺยา, กํ), ...เจว (อิ)}. ทสมํ.}{กุสินารวคฺโค ตติโย.}
\csromancenter{\textbf{ตสฺสุทฺทานํ}}
\setlength{\csromangathapagewidth}{0pt}
\setlength{\csromangathawakwidth}{0pt}
\csromangathameasuregroup{กุสินารภณฺฑนา เจว, \\ กฏุวิยํ เทฺว อนุรุทฺธา,}{\csromangathabat{กุสินารภณฺฑนา เจว,}{โคตมภรณฺฑุหตฺถโก.} \\ \csromangathabat{กฏุวิยํ เทฺว อนุรุทฺธา,}{ปฏิจฺฉนฺนํ เลเขน เต ทสาติ.}}
\csromangathasetleft{กุสินารภณฺฑนา เจว, \\ กฏุวิยํ เทฺว อนุรุทฺธา,}
\csromangathagroup{\csromangathastanza{\csromangathabat{กุสินารภณฺฑนา เจว,}{โคตมภรณฺฑุหตฺถโก.} \\ \csromangathabat{กฏุวิยํ เทฺว อนุรุทฺธา,}{ปฏิจฺฉนฺนํ เลเขน เต ทสาติ.}}}

\csromanmatikaanchor{mtk.31}
\csromantocmarknopage{section}{(14) 4. โยธาชีววคฺค}{mtk.31}
\bookmark[dest={mtk.31},level=1]{(14) 4. โยธาชีววคฺค}
\csromanmarkh{(14)}\csromanmarkhii{4. โยธาชีววคฺค}\csromanheadernomark{1}{(14) 4. โยธาชีววคฺค}
\csromanmatikaanchor{mtk.32}
\csromantocmarkref{subsection}{1. โยธาชีวสุตฺต}{mtk.32}
\bookmark[dest={mtk.32},level=2]{1. โยธาชีวสุตฺต}
\csromanheader{2}{1. \textbf{โยธาชีวสุตฺต}}
\proseitem{134}{ตีติ ภิกฺขเว องฺเคหิ สมนฺนาคโต โยธาชีโว ราชารโต โหติ ราชโภคฺโค, รญฺโญ องฺคนฺเตว สงฺขฺยํ คจฺฉติ. กตเมหิ \csromanfolio{288}ตีหิ? อิธ ภิกฺขเว โยธาชีโว ทูเร ปาตี จ โหติ, อกฺขณเวธี จ, มหโต จ กายสฺส ปทาเลตา. อิเมหิ โข ภิกฺขเว ตีหิ องฺเคหิ สมนฺนาคโต โยธาชีโว ราชารโห โหติ ราชโภคฺโค, รญฺโญ องฺคนฺเตว สงฺขฺยํ คจฺฉติ. เอวเมวํ โข ภิกฺขเว ตีหิ องฺเคหิ สมนฺนาคโต ภิกฺขุ อาหุเนยฺโย โหติ ฯเปฯ อนุตฺตรํ ปุญฺญกฺเขตฺตํ โลกสฺส. กตเมหิ ตีหิ? อิธ ภิกฺขเว ภิกฺขุ ทูเร ปาตี จ โหติ, อกฺขณเวธี จ, มหโต จ กายสฺส ปทาเลตา.}

\prose{กถญฺจ ภิกฺขเว ภิกฺขุ ทูเร ปาตี โหติ? อิธ ภิกฺขเว ภิกฺขุ ยํ กิญฺจิ รูปํ อตีตา\-นาคต\-ปจฺจุปฺปนฺนํ อชฺชตฺตํ วา พหิทฺธา วา โอฬาริกํ วา สุขุมํ วา หีนํ วา ปณีตํ วา ยํ ทูเร สนฺติเก วา, สพฺพํ รูปํ “เนตํ มม, เนโสหมสฺมิ, น เมโส อตฺตา”ติ เอวเมตํ ยถาภูตํ สมฺมปฺปญฺญาย ปสฺสติ. ยา กาจิ เวทนา อตีตา\-นาคต\-ปจฺจุปฺปนฺนา อชฺฌาตฺตา วา พหิทฺธา วา โอฬาริกา วา สุขุมา วา หีนา วา ปณีตา วา ยา ทูเร สนฺติเก วา, สพฺพํ เวทนํ\csromansharedfootnote{7b826905ea46bc64}{สพฺพา เวทนา (สฺยา, กํ, อิ, ก); อุปริ 490 ปิฏฺเฐปิ.} “เนตํ มม, เนโสหมสฺมิ, น เมโส อตฺตา”ติ เอวเมตํ ยถาภูตํ สมฺมปฺปญฺญาย ปสฺสติ. ยา กาจิ สญฺญา อตีตา\-นาคต\-ปจฺจุปฺปนฺนา อชฺฌตฺตา วา พหิทฺธา วา โอฬาริกา วา สุขุมา วา หีนา วา ปณีตา วา ยา ทูเร สนฺติเก วา, สพฺพํ สญฺญํ\csromansharedfootnote{08004e5ad6b3f94a}{สพฺพา สญฺญา (สฺยา, กํ, อิ, ก)} “เนตํ มม, เนโสหมสฺมิ, น เมโส อตฺตา”ติ เอวเมตํ ยถาภูตํ สมฺมปฺปญฺญาย ปสฺสติ. เย เกจิ สงฺขารา อตีตา\-นาคต\-ปจฺจุปฺปนฺนา อชฺฌตฺตา วา พหิทฺธา วา โอฬาริกา วา สุขุมา วา หีนา วา ปณีตา วา เย ทูเร สนฺติเก วา, สพฺเพ สงฺขาเร “เนตํ มม, เนโสหมสฺมิ, น เมโส อตฺตา”ติ เอวเมตํ ยถาภูตํ สมฺมปฺปญฺญาย ปสฺสติ. ยํ กิญฺจิ วิญฺญาณํ อตีตา\-นาคต\-ปจฺจุปฺปนฺนํ อชฺฌตฺตํ วา พหิทฺธา วา โอฬาริกํ วา สุขุมํ วา หีนํ วา ปณีตํ วา ยํ ทูเร สนฺติเก วา, สพฺพํ วิญฺญาณํ “เนตํ มม, เนโสหมสฺมิ, น เมโส อตฺตา”ติ เอวเมตํ ยถาภูตํ สมฺมปฺปญฺญาย ปสฺสติ. เอวํ โข ภิกฺขเว ภิกฺขุ ทูเร ปาตี โหติ.}

\prose{กถญฺจ ภิกฺขเว ภิกฺขุ อกฺขณเวธี โหติ? อิธ ภิกฺขเว ภิกฺขุ “อิทํ ทุกฺขนฺ”ติ ยถาภูตํ ปชานาติ “อยํ ทุกฺขสมุทโย”ติ ยถาภูตํ ปชานาติ, “อยํ ทุกฺขนิโรโธ”ติ ยถาภูตํ ปชานาติ, “อยํ}

\vspace{\tipitakaparskip}
\csromancenter{(14) 4. โยธาชีววคฺค ทุกฺขนิโรธ\-คามินี ปฏิปทา”ติ ยถาภูตํ ปชานาติ. เอวํ โข ภิกฺขเว ภิกฺขุ อกฺขณเวธี โหติ.}
\prose{\csromanfolio{289}กถญฺจ ภิกฺขเว ภิกฺขุ มหโต กายสฺส ปทาเลตา โหติ? อิธ ภิกฺขเว ภิกฺขุ มหนฺตํ อวิชฺชา\-กฺขนฺธํ ปทาเลติ. เอวํ โข ภิกฺขเว ภิกฺขุ มหโต กายสฺส ปทาเลตา โหติ. อิเมหิ โข ภิกฺขเว ตีหิ ธมฺเมหิ สมนฺนาคโต ภิกฺขุ อาหุเนยฺโย โหติ ฯเปฯ อนุตฺตรํ ปุญฺญกฺเขตฺตํ โลกสฺสาติ.  ปฐมํ.}

\csromanmatikaanchor{mtk.33}
\csromantocmarkref{subsection}{2. ปริสาสุตฺต}{mtk.33}
\bookmark[dest={mtk.33},level=2]{2. ปริสาสุตฺต}
\csromanheader{2}{2. \textbf{ปริสาสุตฺต}}
\proseitem{135}{ติสฺโส อิมา ภิกฺขเว ปริสา. กตมา ติสฺโส? อุกฺกาจิตวินีตา ปริสา, ปฏิปุจฺฉา\-วินีตา ปริสา, ยาวตาวินีตา\csromansharedfootnote{67aca38444845529}{ยาวตชฺฌาวินีตา (อฏฺฐกถายํ ปาฐนฺตรํ)} ปริสา. อิมา โข ภิกฺขเว ติสฺโส ปริสาติ.  ทุติยํ.}

\csromanmatikaanchor{mtk.34}
\csromantocmarkref{subsection}{3. มิตฺตสุตฺต}{mtk.34}
\bookmark[dest={mtk.34},level=2]{3. มิตฺตสุตฺต}
\csromanheader{2}{3. \textbf{มิตฺตสุตฺต}}
\proseitem{136}{ตีหิ ภิกฺขเว องฺเคหิ สมนฺนาคโต มิตฺโต เสวิตพฺโพ. กตเมหิ ตีหิ? ( )\csromansharedfootnote{3a586c852dac52b1}{(อิธ ภิกฺขเว ภิกฺขุ) (อิ, ก)} ทุทฺททํ ททาติ, ทุกฺกรํ กโรติ, ทุกฺขมํ ขมติ. อิเมหิ โข ภิกฺขเว ตีหิ องฺเคหิ สมนฺนาคโต มิตฺโต เสวิตพฺโพติ.  ตติยํ.}

\csromanmatikaanchor{mtk.35}
\csromantocmarkref{subsection}{4. อุปฺปาทาสุตฺต}{mtk.35}
\bookmark[dest={mtk.35},level=2]{4. อุปฺปาทาสุตฺต}
\csromanheader{2}{4. \textbf{อุปฺปาทาสุตฺต}}
\csromanmatikaanchor{mtk.36}
\csromantocmarkref{subsection}{5-7. เกสกมฺพลสุตฺตาทีนิ}{mtk.36}
\bookmark[dest={mtk.36},level=2]{5-7. เกสกมฺพลสุตฺตาทีนิ}
\proseitem{137}{อุปฺปาทา วา ภิกฺขเว ตถาคตานํ อนุปฺปาทา วา ตถาคตานํ ฐิตาว สา ธาตุ ธมฺมฏฺฐิตตา ธมฺมนิยามตา. สพฺเพ สงฺขารา อนิจฺจา ตํ ตถาคโต อภิสมฺ\-พุชฺฌติ อภิสเมติ, อภิสมฺ\-พุชฺฌิ\-ตฺวา อภิสเมตฺวา อาจิกฺขติ เทเสติ ปญฺญาเปติ ปฏฺฐเปติ วิวรติ วิภชติ อุตฺตานีกโรติ “สพฺเพ สงฺขารา อนิจฺจา”ติ. อุปฺปาทา วา ภิกฺขเว ตถาคตานํ อนุปฺปาทา วา ตถาคตานํ ฐิตาว สา ธาตุ ธมฺมฏฺฐิตตา ธมฺมนิยามตา สพฺเพ สงฺขารา ทุกฺขา, ตํ ตถาคโต อภิสมฺ\-พุชฺฌติ อภิสเมติ, \csromanfolio{290}อภิสมฺ\-พุชฺฌิ\-ตฺวา อภิสเมตฺวา อาจิกฺขติ เทเสติ ปญฺญาเปติ ปฏฺฐเปติ วิวรติ วิภชติ อุตฺตานีกโรติ “สพฺเพ สงฺขารา ทุกฺขา”ติ. อุปฺปาทา วา ภิกฺขเว ตถาคตานํ อนุปฺปาทา วา ตถาคตานํ ฐิตาว สา ธาตุ ธมฺมฏฺฐิตตา ธมฺมนิยามตา สพฺเพ ธมฺมา อนตฺตา, ตํ ตถาคโต อภิสมฺ\-พุชฺฌติ อภิสเมติ, อภิสมฺ\-พุชฺฌิ\-ตฺวา อภิสเมตฺวา อาจิกฺขติ เทเสติ ปญฺญาเปติ ปฏฺฐเปติ วิวรติ วิภชติ อุตฺตานีกโรติ “สพฺเพ ธมฺมา อนตฺตา”ติ.  จตุตฺถํ.}

\csromanheader{2}{5. เกสกมฺพล\-สุตฺต}
\proseitem{138}{เสยฺยถาปิ ภิกฺขเว ยานิ กานิจิ ตนฺตาวุตานํ วตฺถานํ, เกสกมฺพโล เตสํ ปฏิกิฏฺโฐ อกฺขายติ. เกสกมฺพโล ภิกฺขเว สีเต สีโต, อุเณฺห อุโณฺห ทุพฺพณฺโณ ทุคฺคนฺโธ ทุกฺข\-สมฺผสฺโส. เอวเมวํ โข ภิกฺขเว ยานิ กานิจิ ปุถุสมณพฺราหฺมณ\-วาทานํ\csromansharedfootnote{f72a5fa0ade1c1a9}{สมณปฺปวาทานํ (สี, สฺยา, กํ, อิ)}, มกฺขลิวาโท เตสํ ปฏิกิฏฺโฐ อกฺขายติ.}

\prose{มกฺขลิ ภิกฺขเว โมฆปุริโส เอวํวาที เอวํทิฏฺฐิ “นตฺถิ กมฺมํ, นตฺถิ กิริยํ, นตฺถิ วีริยนฺ”ติ. เยปิ เต ภิกฺขเว อเหสุํ อตีตมทฺธานํ อรหนฺโต สมฺมาสมฺพุทฺธา, เตปิ ภควนฺโต กมฺมวาทา เจว อเหสุํ กิริยวาทา จ วีริยวาทา จ, เตปิ ภิกฺขเว มกฺขลิ โมฆปุริโส ปฏิพาหติ “นตฺถิ กมฺมํ, นตฺถิ กิริยํ, นตฺถิ วีริยนฺ”ติ. เยปิ เต ภิกฺขเว ภวิสฺสนฺติ อนาคตมทฺธานํ อรหนฺโต สมฺมาสมฺพุทฺธา, เตปิ ภควนฺโต กมฺมวาทา เจว ภวิสฺสนฺติ กิริยวาทา จ วีริยวาทา จ, เตปิ ภิกฺขเว มกฺขลิ โมฆปุริโส ปฏิพาหติ “นตฺถิ กมฺมํ, นตฺถิ กิริยํ, นตฺถิ วีริยนฺ”ติ. อหมฺปิ ภิกฺขเว เอตรหิ อรหํ สมฺมาสมฺพุทฺโธ กมฺมวาโท เจว กิริยวาโท จ วีริยวาโท จ, มมฺปิ ภิกฺขเว มกฺขลิ โมฆปุริโส ปฏิพาหติ “นตฺถิ กมฺมํ, นตฺถิ กิริยํ, นตฺถิ วีริยนฺ”ติ.}

\prose{เสยฺยถาปิ ภิกฺขเว นทีมุเธ ขิปฺปํ อุฑฺเฑยฺย\csromansharedfootnote{4f5f951fece1af2e}{โอฑฺเฑยฺย (สี)} พหูนํ\csromansharedfootnote{4644053058c5f146}{พหุนฺนํ (สี, สฺยา, กํ, อิ)} มจฺฉานํ อหิตาย ทุกฺขาย อนยาย พฺยสนาย. เอวเมวํ โข ภิกฺขเว มกฺขลิ โมฆปุริโส มนุสฺสขิปฺปํ มญฺเญ โลเก อุปฺปนฺโน พหูนํ สตฺตานํ อหิตาย ทุกฺขาย อนยาย พฺยสนายาติ.  ปญฺจมํ.}

\chapterheadpageread{(14) 4. โยธาชีววคฺค \textbf{6. สมฺปทาสุตฺต}}
\proseitem{139}{\csromanfolio{291}ติสฺโส อิมา ภิกฺขเว สมฺปทา. กตมา ติสฺโส? สทฺธาสมฺปทา สีลสมฺปทา ปญฺญาสมฺปทา. อิมา โข ภิกฺขเว ติสฺโส สมฺปทาติ.  ฉฏฺฐํ.}

\csromanheader{2}{7. \textbf{วุทฺธิสุตฺต}}
\proseitem{140}{ติสฺโส อิมา ภิกฺขเว วุทฺธิโย. กตมา ติสฺโส? สทฺธาวุทฺธิ สีลวุทฺธิ ปญฺญาวุทฺธิ. อิมา โข ภิกฺขเว ติสฺโส วุทฺธิโยติ.  สตฺตมํ.}

\csromanmatikaanchor{mtk.37}
\csromantocmarkref{subsection}{8. อสฺสขฬุงฺกสุตฺต}{mtk.37}
\bookmark[dest={mtk.37},level=2]{8. อสฺสขฬุงฺกสุตฺต}
\csromanheader{2}{8. อสฺสขฬุงฺก\-สุตฺต}
\proseitem{141}{ตโย จ ภิกฺขเว อสฺสขฬุงฺเก เทเสสฺสามิ, ตโย จ ปุริสขฬุงฺเก, ตํ สุณาถ สาธุกํ มนสิ กโรถ, ภาสิสฺสามีติ. “เอวํ ภนฺเต”ติ โข เต ภิกฺขู ภควโต ปจฺจสฺโสสุํ. ภควา เอตทโวจ\csromandash{}}

\prose{กตเม จ ภิกฺขเว ตโย อสฺสขฬุงฺกา? อิธ ภิกฺขเว เอกจฺโจ อสฺสขฬุงฺโก ชวสมฺปนฺโน โหติ น วณฺณสมฺปนฺโน น อาโรหปริณาห\-สมฺปนฺโน, อิธ ปน ภิกฺขเว เอกจฺโจ อสฺสขฬุงฺโก ชวสมฺปนฺโน จ โหติ วณฺณสมฺปนฺโน จ น อาโรหปริณาห\-สมฺปนฺโน, อิธ ปน ภิกฺขเว เอกจฺโจ อสฺสขฬุงฺโก ชวสมฺปนฺโน จ โหติ วณฺณสมฺปนฺโน จ อาโรหปริณาห\-สมฺปนฺโน จ. อิเม โข ภิกฺขเว ตโย อสฺสขฬุงฺกา.}

\prose{กตเม จ ภิกฺขเว ตโย ปุริสขฬุงฺกา? อิธ ภิกฺขเว เอกจฺโจ ปุริสขฬุงฺโก ชวสมฺปนฺโน โหติ น วณฺณสมฺปนฺโน น อาโรหปริณาห\-สมฺปนฺโน, อิธ ปน ภิกฺขเว เอกจฺโจ ปุริสขฬุงฺโก ชวสมฺปนฺโน จ โหติ วณฺณสมฺปนฺโน จ น อาโรหปริณาห\-สมฺปนฺโน, อิธ ปน ภิกฺขเว เอกจฺโจ ปุริสขฬุงฺโก ชวสมฺปนฺโน จ โหติ วณฺณสมฺปนฺโน จ อาโรหปริณาห\-สมฺปนฺโน จ.}

\prose{กถญฺจ ภิกฺขเว ปุริสขฬุงฺโก ชวสมฺปนฺโน โหติ น วณฺณสมฺปนฺโน น อาโรหปริณาห\-สมฺปนฺโน? อิธ ภิกฺขเว ภิกฺขุ “อิทํ ทุกฺขนฺ”ติ ยถาภูตํ ปชานาติ ฯเปฯ “อยํ ทุกฺขนิโรธ\-คามินี ปฏิปทา”ติ ยถาภูตํ ปชานาติ, อิทมสฺส ชวสฺมิํ วทามิ. อภิธมฺเม โข ปน อภิวินเย ปญฺหํ ปุฏฺโฐ \csromanfolio{292}สํสาเทติ\csromansharedfootnote{4a43d6bf1811d6ef}{สํหีเรติ (ก)} โน วิสฺสชฺเชติ, อิทมสฺส น วณฺณสฺมิํ วทามิ. น โข ปน ลาภี โหติ จีวร\-ปิณฺฑปาต\-เสนาสน\-คิลานปฺปจฺจย\-เภสชฺช\-ปริกฺขารานํ, อิทมสฺส น อาโรห\-ปริณาหสฺมิํ วทามิ. เอวํ โข ภิกฺขเว ปุริสขฬุงฺโก ชวสมฺปนฺโน โหติ น วณฺณสมฺปนฺโน น อาโรหปริณาห\-สมฺปนฺโน.}

\prose{กถญฺจ ภิกฺขเว ปุริสขฬุงฺโก ชวสมฺปนฺโน จ โหติ วณฺณสมฺปนฺโน จ น อาโรหปริณาห\-สมฺปนฺโน? อิธ ภิกฺขเว ภิกฺขุ “อิทํ ทุกฺขนฺ”ติ ยถาภูตํ ปชานาติ ฯเปฯ “อยํ ทุกฺขนิโรธ\-คามินี ปฏิปทา”ติ ยถาภูตํ ปชานาติ, อิทมสฺส ชวสฺมิํ วทามิ. อภิธมฺเม โข ปน อภิวินเย ปญฺหํ ปุฏฺโฐ วิสฺสชฺเชติ โน สํสาเทติ, อิทมสฺส วณฺณสฺมิํ วทามิ. น ปน ลาภี โหติ จีวร\-ปิณฺฑปาต\-เสนาสน\-คิลานปฺปจฺจย\-เภสชฺช\-ปริกฺขารานํ, อิทมสฺส น อาโรห\-ปริณาหสฺมิํ วทามิ. เอวํ โข ภิกฺขเว ปุริสขฬุงฺโก ชวสมฺปนฺโน จ โหติ วณฺณสมฺปนฺโน จ น อาโรหปริณาห\-สมฺปนฺโน.}

\prose{กถญฺจ ภิกฺขเว ปุริสขฬุงฺโก ชวสมฺปนฺโน จ โหติ วณฺณสมฺปนฺโน จ อาโรหปริณาห\-สมฺปนฺโน จ? อิธ ภิกฺขเว ภิกฺขุ “อิทํ ทุกฺขนฺ”ติ ยถาภูตํ ปชานาติ ฯเปฯ “อยํ ทุกฺขนิโรธ\-คามินี ปฏิปทา”ติ ยถาภูตํ ปชานาติ, อิทมสฺส ชวสฺมิํ วทามิ. อภิธมฺเม โข ปน อภิวินเย ปญฺหํ ปุฏฺโฐ วิสฺสชฺเชติ โน สํสาเทติ, อิทมสฺส วณฺณสฺมิํ วทามิ. ลาภี โข ปน โหติ จีวรวิณฺฑปาตเสนาสนคิลานปฺปจฺจยเภสชฺชปริกฺขารานํ, อิทมสฺส อาโรห\-ปริณาหสฺมิํ วทามิ. เอวํ โข ภิกฺขเว ปุริสขฬุงฺโก ชวสมฺปนฺโน จ โหติ วณฺณสมฺปนฺโน จ อาโรหปริณาห\-สมฺปนฺโน จ. อิเม โข ภิกฺขเว ตโย ปุริสขฬุงฺกาติ.  อฏฺฐมํ.}

\csromanmatikaanchor{mtk.38}
\csromantocmarkref{subsection}{9. อสฺสปรสฺสสุตฺต}{mtk.38}
\bookmark[dest={mtk.38},level=2]{9. อสฺสปรสฺสสุตฺต}
\csromanheader{2}{9. อสฺสปรสฺส\-สุตฺต}
\proseitem{142}{ตโย จ ภิกฺขเว อสฺสปรสฺเส\csromansharedfootnote{416ac4583dbe3350}{อสฺสสทสฺเส (สี, สฺยา, กํ, อิ)} เทเสสฺสามิ ตโย จ ปุริสปรสฺเส\csromansharedfootnote{fae29402cebca3ee}{ปุริสสทสฺเส (สี, สฺยา, กํ, อิ)}, ตํ สุณาถ สาธุกํ มนสิ กโรถ, ภาสิสฺสามีติ. “เอวํ ภนฺเต”ติ โข เต ภิกฺขู ภควโต ปจฺจสฺโสสุํ. ภควา เอตทโวจ\csromandash{}}

\chapterheadpageread{(14) 4. โยธาชีววคฺค}
\prose{\csromanfolio{293}กตเม จ ภิกฺขเว ตโย อสฺสปรสฺสา? อิธ ภิกฺขเว เอกจฺโจ อสฺสปรสฺโส ชวสมฺปนฺโน โหติ น วณฺณสมฺปนฺโน น อาโรหปริณาห\-สมฺปนฺโน, อิธ ปน ภิกฺขเว เอกจฺโจ อสฺสปรสฺโส ชวสมฺปนฺโน โหติ วณฺณสมฺปนฺโน จ น อาโรหปริณาห\-สมฺปนฺโน, อิธ ปน ภิกฺขเว เอกจฺโจ อสฺสปรสฺโส ชวสมฺปนฺโน จ โหติ วณฺณสมฺปนฺโน จ อาโรหปริณาห\-สมฺปนฺโน จ. อิเม โข ภิกฺขเว ตโย อสฺสปรสฺสา.}

\prose{กตเม จ ภิกฺขเว ตโย ปุริสปรสฺสา? อิธ ภิกฺขเว เอกจฺโจ ปุริสปรสฺโส ชวสมฺปนฺโน โหติ น วณฺณสมฺปนฺโน น อาโรหปริณาห\-สมฺปนฺโน, อิธ ปน ภิกฺขเว เอกจฺโจ ปุริสปรสฺโส ชวสมฺปนฺโน จ โหติ วณฺณสมฺปนฺโน จ น อาโรหปริณาห\-สมฺปนฺโน, อิธ ปน ภิกฺขเว เอกจฺโจ ปุริสปรสฺโส ชวสมฺปนฺโน จ โหติ วณฺณสมฺปนฺโน จ อาโรหปริณาห\-สมฺปนฺโน จ.}

\prose{กถญฺจ ภิกฺขเว ปุริสปรสฺโส ชวสมฺปนฺโน โหติ น วณฺณสมฺปนฺโน น อาโรหปริณาห\-สมฺปนฺโน? อิธ ภิกฺขเว ภิกฺขุ ปญฺจนฺนํ โอรมฺภาคิยานํ สํโยชนานํ ปริกฺขยา โอปปาติโก โหติ ตตฺถ ปรินิพฺพายี อนาวตฺติธมฺโม ตสฺมา โลกา, อิทมสฺส ชวสฺมิํ วทามิ. อภิธมฺเม โข ปน อภิวินเย ปญฺหํ ปุฏฺโฐ สํสาเทติ โน วิสฺสชฺเชติ, อิทมสฺส น วณฺณสฺมิํ วทามิ. น โข ปน ลาภี โหติ จีวร\-ปิณฺฑปาต\-เสนาสน\-คิลานปฺปจฺจย\-เภสชฺช\-ปริกฺขารานํ, อิทมสฺส น อาโรห\-ปริณาหสฺมิํ วทามิ. เอวํ โข ภิกฺขเว ปุริสปรสฺโส ชวสมฺปนฺโน โหติ น วณฺณสมฺปนฺโน น อาโรหปริณาห\-สมฺปนฺโน.}

\prose{กถญฺจ ภิกฺขเว ปุริสปรสฺโส ชวสมฺปนฺโน จ โหติ วณฺณสมฺปนฺโน จ น อาโรหปริณาห\-สมฺปนฺโน? อิธ ภิกฺขเว ภิกฺขุ ปญฺจนฺนํ โอรมฺภาคิยานํ สํโยชนานํ ปริกฺขยา โอปปาติโก โหติ ตตฺถ ปรินิพฺพายี อนาวตฺติธมฺโม ตสฺมา โลกา, อิทมสฺส ชวสฺมิํ วทามิ. อภิธมฺเม โข ปน อภิวินเย ปญฺหํ ปุฏฺโฐ วิสฺสชฺเชติ โน สํสาเทติ, อิทมสฺส วณฺณสฺมิํ วทามิ. น โข ปน ลาภี โหติ จีวร\-ปิณฺฑปาต\-เสนาสน\-คิลานปฺปจฺจย\-เภสชฺช\-ปริกฺขารานํ, อิทมสฺส น อาโรห\-ปริณาหสฺมิํ วทามิ. เอวํ โข ภิกฺขเว ปุริสปรสฺโส ชวสมฺปนฺโน จ โหติ วณฺณสมฺปนฺโน จ น อาโรหปริณาห\-สมฺปนฺโน.}

\prose{\csromanfolio{294}กถญฺจ ภิกฺขเว ปุริสปรสฺโส ชวสมฺปนฺโน จ โหติ วณฺณสมฺปนฺโน จ อาโรหปริณาห\-สมฺปนฺโน จ? อิธ ภิกฺขเว ภิกฺขุ ปญฺจนฺนํ โอรมฺภาคิยานํ สํโยชนานํ ปริกฺขยา โอปปาติโก โหติ ตตฺถ ปรินิพฺพายี อนาวตฺติธมฺโม ตสฺมา โลกา, อิทมสฺส ชวสฺมิํ วทามิ. อภิธมฺเม โข ปน อภิวินเย ปญฺหํ ปุฏฺโฐ วิสฺสชฺเชติ โน สํสาเทติ, อิทมสฺส วณฺณสฺมิํ วทามิ. ลาภี โข ปน โหติ จีวร\-ปิณฺฑปาต\-เสนาสน\-คิลานปฺปจฺจย\-เภสชฺช\-ปริกฺขารานํ, อิทมสฺส อาโรห\-ปริณาหสฺมิํ วทามิ. เอวํ โข ภิกฺขเว ปุริสปรสฺโส ชวสมฺปนฺโน จ โหติ วณฺณสมฺปนฺโน จ อาโรหปริณาห\-สมฺปนฺโน จ. อิเม โข ภิกฺขเว ตโย ปุริสปรสฺสาติ.  นวมํ.}

\csromanmatikaanchor{mtk.39}
\csromantocmarkref{subsection}{10. อสฺสาชานียสุตฺต}{mtk.39}
\bookmark[dest={mtk.39},level=2]{10. อสฺสาชานียสุตฺต}
\csromanheader{2}{10. อสฺสาชานีย\-สุตฺต}
\proseitem{143}{ตโย จ ภิกฺขเว ภเทฺร อสฺสาชานีเย เทเสสฺสามิ ตโย จ ภเทฺร ปุริสาชานีเย, ตํ สุณาถ สาธุกํ มนสิ กโรถ, ภาสิสฺสามีติ. “เอวํ ภนฺเต”ติ โข เต ภิกฺขู ภควโต ปจฺจสฺโสสุํ. ภควา เอตทโวจ\csromandash{}}

\prose{กตเม จ ภิกฺขเว ตโย ภทฺรา อสฺสาชานียา? อิธ ภิกฺขเว เอกจฺโจ ภโทฺร อสฺสาชานีโย ฯเปฯ ชวสมฺปนฺโน จ โหติ วณฺณสมฺปนฺโน จ อาโรหปริณาห\-สมฺปนฺโน จ. อิเม โข ภิกฺขเว ตโย ภทฺรา อสฺสาชานียา.}

\prose{กตเม จ ภิกฺขเว ตโย ภทฺรา ปุริสาชานียา? อิธ ภิกฺขเว เอกจฺโจ ภโทฺร ปุริสาชานีโย ฯเปฯ ชวสมฺปนฺโน จ โหติ วณฺณสมฺปนฺโน จ อาโรหปริณาห\-สมฺปนฺโน จ.}

\prose{กถญฺจ ภิกฺขเว ภโทฺร ปุริสาชานีโย ฯเปฯ ชวสมฺปนฺโน จ โหติ วณฺณสมฺปนฺโน จ อาโรหปริณาห\-สมฺปนฺโน จ? อิธ ภิกฺขเว ภิกฺขุ อาสวานํ ขยา อนาสวํ เจโตวิมุตฺติํ ปญฺญาวิมุตฺติํ ทิฏฺเฐว ธมฺเม สยํ อภิญฺญา สจฺฉิกตฺวา อุปสมฺปชฺช วิหรติ, อิทมสฺส ชวสฺมิํ วทามิ. อภิธมฺเม โข ปน อภิวินเย ปญฺหํ ปุฏฺโฐ วิสฺสชฺเชติ โน สํสาเทติ, อิทมสฺส วณฺณสฺมิํ วทามิ. ลาภี โข ปน โหติ จีวร\-ปิณฺฑปาต\-เสนาสน\-คิลานปฺปจฺจย\-เภสชฺช\-ปริกฺขารานํ, อิทมสฺส อาโรห\-ปริณาหสฺมิํ วทามิ. เอวํ โข ภิกฺขเว ภโทฺร ปุริสาชานีโย ชวสมฺปนฺโน จ โหติ}

\vspace{\tipitakaparskip}
\csromancenter{(14) 4. โยธาชีววคฺค วณฺณสมฺปนฺโน จ อาโรหปริณาห\-สมฺปนฺโน จ. อิเม โข ภิกฺขเว ตโย ภทฺรา ปุริสาชานียาติ.  ทสมํ.}
\csromanheader{2}{11. ปฐม\-โมรนิวาป\-สุตฺต}
\csromanmatikaanchor{mtk.40}
\csromantocmarkref{subsection}{13. โมรนิวาปสุตฺตานิ}{mtk.40}
\bookmark[dest={mtk.40},level=2]{13. โมรนิวาปสุตฺตานิ}
\proseitem{144}{\csromanfolio{295}เอกํ สมยํ ภควา ราชคเห วิหรติ โมรนิวาเป ปริพฺพาชกา\-ราเม. ตตฺร โข ภควา ภิกฺขู อามนฺเตสิ “ภิกฺขโว”ติ. “ภทนฺเต”ติ เต ภิกฺขู ภควโต ปจฺจสฺโสสุํ. ภควา เอตทโวจ\csromandash{}}

\prose{ตีหิ ภิกฺขเว ธมฺเมหิ สมนฺนาคโต ภิกฺขุ อจฺจนฺตนิฏฺโฐ โหติ อจฺจนฺต\-โยคกฺเขมี อจฺจนฺต\-พฺรหฺมจารี อจฺจนฺต\-ปริโยสาโน เสฏฺโฐ เทวมนุสฺสานํ. กตเมหิ ตีหิ? อเสกฺเขน สีลกฺขนฺเธน, อเสกฺเขน สมาธิ\-กฺขนฺเธน, อเสกฺเขน ปญฺญา\-กฺขนฺเธน. อิเมหิ โข ภิกฺขเว ตีหิ ธมฺเมหิ สมนฺนาคโต ภิกฺขุ อจฺจนฺตนิฏฺโฐ โหติ อจฺจนฺต\-โยคกฺเขมี อจฺจนฺต\-พฺรหฺมจารี อจฺจนฺต\-ปริโยสาโน เสฏฺโฐ เทว\-มนุสฺสา\-นนฺติ.  เอกาทสมํ.}

\csromanheader{2}{12. ทุติย\-โมรนิวาป\-สุตฺต}
\proseitem{145}{ตีหิ ภิกฺขเว ธมฺเมหิ สมนฺนาคโต ภิกฺขุ อจฺจนฺตนิฏฺโฐ โหติ อจฺจนฺต\-โยคกฺเขมี อจฺจนฺต\-พฺรหฺมจารี อจฺจนฺต\-ปริโยสาโน เสฏฺโฐ เทวมนุสฺสานํ. กตเมหิ ตีหิ? อิทฺธิ\-ปาฏิหาริเยน อาเทสนา\-ปาฏิหาริเยน อนุสาสนี\-ปาฏิหาริเยน. อิเมหิ โข ภิกฺขเว ตีหิ ธมฺเมหิ สมนฺนาคโต ภิกฺขุ อจฺจนฺตนิฏฺโฐ โหติ อจฺจนฺต\-โยคกฺเขมี อจฺจนฺต\-พฺรหฺมจารี อจฺจนฺต\-ปริโยสาโน เสฏฺโฐ เทว\-มนุสฺสา\-นนฺติ.  ทฺวาทสมํ.}

\csromanheader{2}{13. ตติย\-โมรนิวาป\-สุตฺต}
\csromanmatikaanchor{mtk.41}
\csromanmatikaanchor{mtk.42}
\csromantocmarknopage{section}{(15) 5. มงฺคลวคฺค}{mtk.41}
\bookmark[dest={mtk.41},level=1]{(15) 5. มงฺคลวคฺค}
\csromantocmarkref{subsection}{1-4. อกุสลสุตฺตาทีนิ}{mtk.42}
\bookmark[dest={mtk.42},level=2]{1-4. อกุสลสุตฺตาทีนิ}
\proseitemwithclosermidruled{146}{ตีหิ ภิกฺขเว ธมฺเมหิ สมนฺนาคโต ภิกฺขุ อจฺจนฺตนิฏฺโฐ โหติ อจฺจนฺต\-โยคกฺเขมี อจฺจนฺต\-พฺรหฺมจารี อจฺจนฺต\-ปริโยสาโน เสฏฺโฐ เทวมนุสฺสานํ. กตเมหิ ตีหิ? สมฺมาทิฏฺฐิยา สมฺมาญาเณน สมฺมาวิมุตฺติยา. อิเมหิ โข ภิกฺขเว ตีหิ ธมฺเมหิ สมนฺนาคโต ภิกฺขุ \csromanfolio{296}อจฺจนฺตนิฏฺโฐ โหติ อจฺจนฺต\-โยคกฺเขมี อจฺจนฺต\-พฺรหฺมจารี อจฺจนฺต\-ปริโยสาโน เสฏฺโฐ เทว\-มนุสฺสา\-นนฺติ.  เตรสมํ.}{โยธาชีววคฺโค จตุตฺโถ.}
\csromancenter{\textbf{ตสฺสุทฺทานํ}}
\setlength{\csromangathapagewidth}{0pt}
\setlength{\csromangathawakwidth}{0pt}
\csromangathameasuregroup{โยโธ ปริสมิตฺตญฺจ, \\ สมฺปทา วุทฺธิ ตโย อสฺสา,}{\csromangathabat{โยโธ ปริสมิตฺตญฺจ,}{อุปฺปาทา เกสกมฺพโล.} \\ \csromangathabat{สมฺปทา วุทฺธิ ตโย อสฺสา,}{ตโย โมรนิวาปิโนติ.}}
\csromangathasetleft{โยโธ ปริสมิตฺตญฺจ, \\ สมฺปทา วุทฺธิ ตโย อสฺสา,}
\csromangathagroup{\csromangathastanza{\csromangathabat{โยโธ ปริสมิตฺตญฺจ,}{อุปฺปาทา เกสกมฺพโล.} \\ \csromangathabat{สมฺปทา วุทฺธิ ตโย อสฺสา,}{ตโย โมรนิวาปิโนติ.}}}

\csromanmarkh{(15)}\csromanmarkhii{5. มงฺคลวคฺค}\csromanheadernomark{1}{\textbf{(15) 5. มงฺคลวคฺค}}
\csromanheader{2}{1. \textbf{อกุสลสุตฺต}}
\proseitem{147}{ตีหิ ภิกฺขเว ธมฺเมหิ สมนฺนาคโต ยถาภตํ นิกฺขิตฺโต เอวํ นิรเย. กตเมหิ ตีหิ? อกุสเลน กายกมฺเมน, อกุสเลน วจีกมฺเมน, อกุสเลน มโนกมฺเมน. อิเมหิ โข ภิกฺขเว ตีหิ ธมฺเมหิ สมนฺนาคโต ยถาภตํ นิกฺขิตฺโต เอวํ นิรเย.}

\prose{ตีหิ ภิกฺขเว ธมฺเมหิ สมนฺนาคโต ยถาภตํ นิกฺขิตฺโต เอวํ สคฺเค. กตเมหิ ตีหิ? กุสเลน กายกมฺเมน, กุสเลน วจีกมฺเมน, กุสเลน มโนกมฺเมน. อิเมหิ โข ภิกฺขเว ตีหิ ธมฺเมหิ สมนฺนาคโต ยถาภตํ นิกฺขิตฺโต เอวํ สคฺเคติ.  ปฐมํ.}

\csromanheader{2}{2. \textbf{สาวชฺชสุตฺต}}
\proseitem{148}{ตีหิ ภิกฺขเว ธมฺเมหิ สมนฺนาคโต ยถาภตํ นิกฺขิตฺโต เอวํ นิรเย. กตเมหิ ตีหิ? สาวชฺเชน กายกมฺเมน, สาวชฺเชน วจีกมฺเมน, สาวชฺเชน มโนกมฺเมน. อิเมหิ โข ภิกฺขเว ตีหิ ธมฺเมหิ สมนฺนาคโต ยถาภตํ นิกฺขิตฺโต เอวํ นิรเย.}

\prose{ตีหิ ภิกฺขเว ธมฺเมหิ สมนฺนาคโต ยถาภตํ นิกฺขิตฺโต เอวํ สคฺเค. กตเมหิ ตีหิ? อนวชฺเชน กายกมฺเมน, อนวชฺเชน วจีกมฺเมน, อนวชฺเชน มโนกมฺเมน. อิเมหิ โข ฯเปฯ เอวํ สคฺเคติ.  ทุติยํ.}

\chapterheadpageread{(15) 5. มงฺคลวคฺค \textbf{3. วิสมสุตฺต}}
\csromanmatikaanchor{mtk.43}
\csromantocmarkref{subsection}{5-8. ขตสุตฺตานิ}{mtk.43}
\bookmark[dest={mtk.43},level=2]{5-8. ขตสุตฺตานิ}
\proseitem{149}{\csromanfolio{297}ตีหิ ภิกฺขเว ฯเปฯ วิสเมน กายกมฺเมน, วิสเมน วจีกมฺเมน, วิสเมน มโนกมฺเมน. อิเมหิ โข ฯเปฯ เอวํ นิรเย.}

\prose{ตีหิ ภิกฺขเว ธมฺเมหิ ฯเปฯ สเมน กายกมฺเมน, สเมน วจีกมฺเมน, สเมน มโนกมฺเมน. อิเมหิ โข ฯเปฯ เอวํ สคฺเคติ.  ตติยํ.}

\csromanheader{2}{4. \textbf{อสุจิสุตฺต}}
\proseitem{150}{ตีหิ ภิกฺขเว ฯเปฯ อสุจินา กายกมฺเมน, อสุจินา วจีกมฺเมน, อสุจินา มโนกมฺเมน. อิเมหิ โข ฯเปฯ เอวํ นิรเย.}

\prose{ตีหิ ภิกฺขเว ฯเปฯ สุจินา กายกมฺเมน, สุจินา วจีกมฺเมน, สุจินา มโนกมฺเมน. อิเมหิ โข ภิกฺขเว ตีหิ ธมฺเมหิ สมนฺนาคโต ยถาภตํ นิกฺขิตฺโต เอวํ สคฺเคติ.  จตุตฺถํ.}

\csromanheader{2}{5. ปฐม\-ขต\-สุตฺต}
\proseitem{151}{ตีหิ ภิกฺขเว ธมฺเมหิ สมนฺนาคโต พาโล อพฺยตฺโต อสปฺปุริโส ขตํ อุปหตํ อตฺตานํ ปริหรติ, สาวชฺโช จ โหติ สานุวชฺโช จ วิญฺญูนํ, พหุญฺจ อปุญฺญํ ปสวติ. กตเมหิ ตีหิ? อกุสเลน กายกมฺเมน, อกุสเลน วจีกมฺเมน, อกุสเลน มโนกมฺเมน. อิเมหิ โข ภิกฺขเว ตีหิ ธมฺเมหิ สมนฺนาคโต พาโล อพฺยตฺโต อสปฺปุริโส ขตํ อุปหตํ อตฺตานํ ปริหรติ, สาวชฺโช จ โหติ สานุวชฺโช จ วิญฺญูนํ, พหุญฺจ อปุญฺญํ ปสวติ.}

\prose{ตีหิ ภิกฺขเว ธมฺเมหิ สมนฺนาคโต ปณฺฑิโต วิยตฺโต สปฺปุริโส อกฺขตํ อนุปหตํ อตฺตานํ ปริหรติ, อนวชฺโช จ โหติ อนนุวชฺโช จ วิญฺญูนํ, พหุญฺจ ปุญฺญํ ปสวติ. กตเมหิ ตีหิ? กุสเลน กายกมฺเมน, กุสเลน วจีกมฺเมน, กุสเลน มโนกมฺเมน ฯเปฯ.  ปญฺจมํ.}

\csromanheader{2}{6. ทุติยขต\-สุตฺต}
\proseitem{152}{ตีหิ ภิกฺขเว ฯเปฯ สาวชฺเชน กายกมฺเมน, สาวชฺเชน วจีกมฺเมน, สาวชฺเชน มโนกมฺเมน ฯเปฯ.}

\prose{\csromanfolio{298}ตีหิ ภิกฺขเว ฯเปฯ อนวชฺเชน กายกมฺเมน, อนวชฺเชน วจีกมฺเมน, อนวชฺเชน มโนกมฺเมน ฯเปฯ.  ฉฏฺฐํ.}

\csromanheader{2}{7. ตติย\-ขต\-สุตฺต}
\proseitem{153}{ตีหิ ภิกฺขเว ฯเปฯ สิสเมน กายกมฺเมน, วิสเมน วจีกมฺเมน, วิสเมน มโนกมฺเมน ฯเปฯ.}

\prose{ตีหิ ภิกฺขเว ฯเปฯ สเมน กายกมฺเมน, สเมน วจีกมฺเมน, สเมน มโนกมฺเมน ฯเปฯ.  สตฺตมํ.}

\csromanheader{2}{8. จตุตฺถ\-ขต\-สุตฺต}
\proseitem{154}{ตีหิ ภิกฺขเว ฯเปฯ อสุจินา กายกมฺเมน, อสุจินา วจีกมฺเมน, อสุจินา มโนกมฺเมน ฯเปฯ.}

\prose{ตีหิ ภิกฺขเว ฯเปฯ สุจินา กายกมฺเมน, สุจินา วจีกมฺเมน, สุจินา มโนกมฺเมน. อิเมหิ โข ภิกฺขเว ตีหิ ธมฺเมหิ สมนฺนาคโต ปณฺฑิโต วิยตฺโต สปฺปุริโส อกฺขตํ อนุปหตํ อตฺตานํ ปริหรติ, อนวชฺโช จ โหติ อนนุวชฺโช จ วิญฺญูนํ, พหุญฺจ ปุญฺญํ ปสวตีติ.  อฏฺฐมํ.}

\csromanmatikaanchor{mtk.44}
\csromantocmarkref{subsection}{9. วนฺทนาสุตฺต}{mtk.44}
\bookmark[dest={mtk.44},level=2]{9. วนฺทนาสุตฺต}
\csromanheader{2}{9. \textbf{วนฺทนาสุตฺต}}
\proseitem{155}{ติสฺโส อิมา ภิกฺขเว วนฺทนา. กตมา ติสฺโส? กาเยน วาจาย มนสา. อิมา โข ภิกฺขเว ติสฺโส วนฺทนาติ.  นวมํ.}

\csromanmatikaanchor{mtk.45}
\csromantocmarkref{subsection}{10. ปุพฺพณฺหสุตฺต}{mtk.45}
\bookmark[dest={mtk.45},level=2]{10. ปุพฺพณฺหสุตฺต}
\csromanheader{2}{10. \textbf{ปุพฺพณฺหสุตฺต}}
\proseitem{156}{เย ภิกฺขเว สตฺตา ปุพฺพณฺหสมยํ กาเยน สุจริตํ จรนฺติ, วาจาย สุจริตํ จรนฺติ, มนสา สุจริตํ จรนฺติ. สุปุพฺพโณฺห ภิกฺขเว เตสํ สตฺตานํ.}

\prose{เย ภิกฺขเว สตฺตา มชฺฌนฺหิกสมยํ กาเยน สุจริตํ จรนฺติ, วาจาย สุจริตํ จรนฺติ, มนสา สุจริตํ จรนฺติ. สุมชฺฌนฺหิโก ภิกฺขเว เตสํ สตฺตนํ.}

\chapterheadpageread{(15) 5. มงฺคลวคฺค}
\prose{\csromanfolio{299}เย ภิกฺขเว สตฺตา สายนฺหสมยํ กาเยน สุจริตํ จรนฺติ, วาจาย สุจริตํ จรนฺติ, มนสา สุจริตํ จรนฺติ. สุสายโนฺห ภิกฺขเว เตสํ สตฺตานนฺติ.  ทสมํ.}

\setlength{\csromangathapagewidth}{0pt}
\setlength{\csromangathawakwidth}{0pt}
\csromangathameasuregroup{สุนกฺขตฺตํ สุมงฺคลํ, \\ สุขโณ สุมุหุตฺโต จ, \\ ปทกฺขิณํ กายกมฺมํ, \\ ปทกฺขิณํ มโนกมฺมํ, \\ ปทกฺขิณานิ กตฺวาน, \\ เต อตฺถลทฺธา สุขิตา, \\ อโรคา สุขิตา โหถ,}{\csromangathabat{สุนกฺขตฺตํ สุมงฺคลํ,}{สุปฺปภาตํ สุหุฏฺฐิตํ\textsuperscript{1}.} \\ \csromangathabat{สุขโณ สุมุหุตฺโต จ,}{สุยิฏฺฐํ พฺรหฺมจาริสุ.} \\ \csromangathabat{ปทกฺขิณํ กายกมฺมํ,}{วาจากมฺมํ ปทกฺขิณํ.} \\ \csromangathabat{ปทกฺขิณํ มโนกมฺมํ,}{ปณีธิ เต ปทกฺขิเณ\textsuperscript{1}.} \\ \csromangathabat{ปทกฺขิณานิ กตฺวาน,}{ลภนฺตตฺเถ\textsuperscript{1} ปทกฺขิเณ.} \\ \csromangathabat{เต อตฺถลทฺธา สุขิตา,}{วิรุฬฺหา พุทฺธสาสเน.} \\ \csromangathabat{อโรคา สุขิตา โหถ,}{สห สพฺเพหิ ญาติภีติ.}}
\csromangathasetleft{สุนกฺขตฺตํ สุมงฺคลํ, \\ สุขโณ สุมุหุตฺโต จ, \\ ปทกฺขิณํ กายกมฺมํ, \\ ปทกฺขิณํ มโนกมฺมํ, \\ ปทกฺขิณานิ กตฺวาน, \\ เต อตฺถลทฺธา สุขิตา, \\ อโรคา สุขิตา โหถ,}
\csromangathagroupwithclosermidruled{\csromangathastanza{\csromangathabat{สุนกฺขตฺตํ สุมงฺคลํ,}{สุปฺปภาตํ สุหุฏฺฐิตํ\csromansharedfootnote{d82aa94e99660298}{สุวุฏฺฐิตํ (สี, อิ)}.} \\ \csromangathabat{สุขโณ สุมุหุตฺโต จ,}{สุยิฏฺฐํ พฺรหฺมจาริสุ.}}\csromangathastanza{\csromangathabat{ปทกฺขิณํ กายกมฺมํ,}{วาจากมฺมํ ปทกฺขิณํ.} \\ \csromangathabat{ปทกฺขิณํ มโนกมฺมํ,}{ปณีธิ เต ปทกฺขิเณ\csromansharedfootnote{549089005ab18db5}{ปณิธิโย ปทกฺขิณา (สี, อิ), ปณิธิ เต ปทกฺขิณา (สฺยา, กํ)}.}}\csromangathastanza{\csromangathabat{ปทกฺขิณานิ กตฺวาน,}{ลภนฺตตฺเถ\csromansharedfootnote{b72b5b627c4d5323}{ลภตตฺเถ (สี, อิ)} ปทกฺขิเณ.} \\ \csromangathabat{เต อตฺถลทฺธา สุขิตา,}{วิรุฬฺหา พุทฺธสาสเน.} \\ \csromangathabat{อโรคา สุขิตา โหถ,}{สห สพฺเพหิ ญาติภีติ.}}}{มงฺคลวคฺโค ปญฺจโม.}
\csromancenter{ตสฺสุทฺทานํ}
\setlength{\csromangathapagewidth}{0pt}
\setlength{\csromangathawakwidth}{0pt}
\csromangathameasuregroup{อกุสลญฺจ สาวชฺชํ, \\ จตุโร ขตา วนฺทนา,}{\csromangathabat{อกุสลญฺจ สาวชฺชํ,}{วิสมาสุจินา สห.} \\ \csromangathabat{จตุโร ขตา วนฺทนา,}{ปุพฺพเณฺหน จ เต ทสาติ.} \\ \textbf{ตติโย ปณฺณาสโก สมตฺโต.}}
\csromangathasetleft{อกุสลญฺจ สาวชฺชํ, \\ จตุโร ขตา วนฺทนา,}
\csromangathagroup{\csromangathastanza{\csromangathabat{อกุสลญฺจ สาวชฺชํ,}{วิสมาสุจินา สห.} \\ \csromangathabat{จตุโร ขตา วนฺทนา,}{ปุพฺพเณฺหน จ เต ทสาติ.} \\ \textbf{ตติโย ปณฺณาสโก สมตฺโต.}}}

\csromanpalirecto
\csromanmatikaanchor{mtk.46}
\csromantocmarkref{section}{(16) 6. อเจลกวคฺค}{mtk.46}
\bookmark[dest={mtk.46},level=1]{(16) 6. อเจลกวคฺค}
\csromanmarkh{ติกนิปาตปาฬิ (16)}\csromanmarkhii{6. อเจลกวคฺค}\csromanheadernomark{1}{ติกนิปาตปาฬิ \textbf{(16) 6. อเจลกวคฺค}}
\proseitem{157-163}{\csromanfolio{300}ติสฺโส อิมา ภิกฺขเว ปฏิปทา. กตมา ติสฺโส? อาคาฬฺหา ปฏิปทา, นิชฺฌามา ปฏิปทา, มชฺฌิมา ปฏิปทา. กตมา จ ภิกฺขเว อาคาฬฺหา ปฏิปทา? อิธ ภิกฺขเว เอกจฺโจ เอวํวาที โหติ เอวํทิฏฺฐิ “นตฺถิ กาเมสุ โทโส”ติ, โส กาเมสุ ปาตพฺยตํ อาปชฺชติ. อยํ วุจฺจติ ภิกฺขเว อาคาฬฺหา ปฏิปทา.}

\prose{กตมา จ ภิกฺขเว นิชฺฌามา ปฏิปทา? อิธ ภิกฺขเว เอกจฺโจ อเจลโก โหติ มุตฺตาจาโร หตฺถา\-ปเลขโน\csromansharedfootnote{2cf226ba39ab7742}{หตฺถาวเลขโน (สฺยา, กํ) ที 1. 157; ม 1. 111; อุปริ 527; ม 2. 5 ปิฏฺเฐสุปิ ปสฺสิตพฺพา.} น เอหิภทนฺติโก, น ติฏฺฐ\-ภทนฺติโก, นาภิหฏํ, น อุทฺทิสฺสกตํ, น นิมนฺตนํ สาทิยติ, โส น กุมฺภิมุขา ปฏิคฺคณฺหาติ, น กโฬปิมุขา\csromansharedfootnote{b16d0afafdaec71b}{ขโฬปิมุขา (สี, สฺยา, กํ)} ปฏิคฺคณฺหาติ, น เอฬกมนฺตรํ, น ทณฺฑมนฺตรํ, น มุสลมนฺตรํ, น ทฺวินฺนํ ภุญฺชมานานํ, น คพฺภินิยา, น ปายมานาย, น ปุริส\-นฺตร\-คตาย, น สงฺกิตฺตีสุ, น ยตฺถ สา อุปฏฺฐิโต โหติ, น ยตฺถ มกฺขิกา สณฺฑ\-สณฺฑ\-จารินี, น มจฺฉํ, น มํสํ, น สุรํ, น เมรยํ, น ถุโสทกํ ปิวติ. โส เอกาคาริโก วา โหติ เอกาโลปิโก, ทฺวาคาริโก วา โหติ ทฺวาโลปิโก, สตฺตาคาริโก วา โหติ สตฺตาโลปิโก, เอกิสฺสาปิ ทตฺติยา ยาเปติ, ทฺวีหิปิ ทตฺตีหิ ยาเปติ, สตฺตหิปิ ทตฺตีหิ ยาเปติ, เอกาหิกมฺปิ อาหารํ อาหาเรติ, ทฺวาหิกมฺปิ อาหารํ อาหาเรติ, สตฺตาหิกมฺปิ อาหารํ อาหาเรติ. อิติ เอวรูปํ อทฺธมาสิกมฺปิ ปริยาย\-ภตฺตโภชนานุโยคมนุยุตฺโต วิหรติ.}

\prose{โส สากภกฺโขปิ โหติ, สามาก\-ภกฺโขปิ โหติ, นีวาร\-ภกฺโขปิ โหติ, ททฺทุล\-ภกฺโขปิ โหติ, หฏภกฺโขปิ โหติ, กณภกฺโขปิ โหติ, อาจามภกฺโขปิ โหติ, ปิญฺญาก\-ภกฺโขปิ โหติ, ติณภกฺโขปิ โหติ, โคมย\-ภกฺโขปิ โหติ, วนมูล\-ผลา\-หาโร ยาเปติ, ปวตฺต\-ผล\-โภชี.}

\prose{โส สาณานิปิ ธาเรติ, มสาณานิปิ ธาเรติ, ฉวทุสฺสานิปิ ธาเรติ, ปํสุกูลานิปิ ธาเรติ, ติรีฏานิปิ ธาเรติ, อชินมฺปิ ธาเรติ, อชินกฺขิปมฺปิ ธาเรติ, กุสจีรมฺปิ ธาเรติ, วากจีรมฺปิ ธาเรติ, ผลกจีรมฺปิ ธาเรติ, เกสกมฺพลมฺปิ ธาเรติ, วาฬกมฺพลมฺปิ ธาเรติ,}

\vspace{\tipitakaparskip}
\csromancenter{(16) 6. อเจลกวคฺค อุลูก\-ปกฺขิกมฺปิ ธาเรติ, เกสมสฺสุ\-โลจโกปิ โหติ เกสมสฺสุ\-โลจนา\-นุโยคม\-นุยุตฺโต, อุพฺภฏฺฐโกปิ โหติ อาสน\-ปฏิกฺขิตฺโต, อุกฺกุฏิโกปิ โหติ อุกฺกุฏิก\-ปฺปธานมนุยุตฺโต, กณฺฏกา\-ปสฺสยิ\-โกปิ โหติ กณฺฏกาปสฺสเย เสยฺยํ กปฺเปติ, สาย\-ตติยกมฺปิ อุทโก\-โรหนา\-นุโยคม\-นุยุตฺโต วิหรติ. อิติ เอวรูปํ อเนกวิหิตํ กายสฺส อาตาปน\-ปริตาปนา\-นุโยคม\-นุยุตฺโต วิหรติ. อยํ วุจฺจติ ภิกฺขเว นิชฺฌามา ปฏิปทา.}
\prose{\csromanfolio{301}กตมา จ ภิกฺขเว มชฺฌิมา ปฏิปทา? อิธ ภิกฺขเว ภิกฺขุ กาเย กายานุปสฺสี วิหรติ อาตาปี สมฺปชาโน สติมา วิเนยฺย โลเก อภิชฺฌา- โทมนสฺสํ. เวทนาสุ ฯเปฯ. จิตฺเต ฯเปฯ. ธมฺเมสุ ธมฺมานุปสฺสี วิหรติ อาตาปี สมฺปชาโน สติมา วิเนยฺย โลเก อภิชฺฌา\-โทมนสฺสํ. อยํ วุจฺจติ ภิกฺขเว มชฺฌิมา ปฏิปทา. อิมา โข ภิกฺขเว ติสฺโส ปฏิปทาติ. (1)}

\prose{ติสฺโส อิมา ภิกฺขเว ปฏิปทา. กตมา ติสฺโส? อาคาฬฺหา ปฏิปทา, นิชฺฌามา ปฏิปทา, มชฺฌิมา ปฏิปทา. กตมา จ ภิกฺขเว อาคาฬฺหา ปฏิปทา ฯเปฯ. อยํ วุจฺจติ ภิกฺขเว อาคาฬฺหา ปฏิปทา.}

\prose{กตมา จ ภิกฺขเว นิชฺฌามา ปฏิปทา ฯเปฯ. อยํ วุจฺจติ ภิกฺขเว นิชฺฌามา ปฏิปทา.}

\prose{กตมา จ ภิกฺขเว มชฺฌิมา ปฏิปทา? อิธ ภิกฺขเว ภิกฺขุ อนุปฺปนฺนานํ ปาปกานํ อกุสลานํ ธมฺมานํ อนุปฺปาทาย ฉนฺทํ ชเนติ วายมติ วีริยํ อารภติ จิตฺตํ ปคฺคณฺหาติ ปทหติ. อุปฺปนฺนานํ ปาปกานํ อกุสลานํ ธมฺมานํ ปหานาย ฉนฺทํ ชเนติ วายมติ วีริยํ อารภติ จิตฺตํ ปคฺคณฺหาติ ปทหติ. อนุปฺปนฺนานํ กุสลานํ ธมฺมานํ อุปฺปาทาย ฉนฺทํ ชเนติ วายมติ วีริยํ อารภติ จิตฺตํ ปคฺคณฺหาติ ปทหติ. อุปฺปนฺนานํ กุสลานํ ธมฺมานํ ฐิติยา อสมฺโมสาย ภิยฺโยภาวาย เวปุลฺลาย ภาวนาย ปาริปูริยา ฉนฺทํ ชเนติ วายมติ วีริยํ อารภติ จิตฺตํ ปคฺคณฺหาติ ปทหติ. (2)}

\prose{ฉนฺท\-สมาธิ\-ปธาน\-สงฺขาร\-สมนฺนาคตํ อิทฺธิปาทํ ภาเวติ. วีริยสมาธิ ฯเปฯ. จิตฺตสมาธิ ฯเปฯ. วีมํสาสมาธิ\-ปธานสงฺขาร- สมนฺนาคตํ อิทฺธิปาทํ ภาเวติ ฯเปฯ. (3)}

\prose{สทฺธินฺทฺริยํ ภาเวติ. วีริยินฺทฺริยํ ภาเวติ. สตินฺทฺริยํ ภาเวติ. สมาธินฺทฺริยํ ภาเวติ. ปญฺญินฺทฺริยํ ภาเวติ ฯเปฯ. (4)}

\csromanmatikaanchor{mtk.47}
\csromanmatikaanchor{mtk.48}
\csromantocmarkref{section}{อุทฺทานคาถา}{mtk.47}
\bookmark[dest={mtk.47},level=1]{อุทฺทานคาถา}
\csromantocmarkref{section}{(17) 7. กมฺมปถเปยฺยาล}{mtk.48}
\bookmark[dest={mtk.48},level=1]{(17) 7. กมฺมปถเปยฺยาล}
\prose{\csromanfolio{302}สทฺธาพลํ ภาเวติ. วีริยพลํ ภาเวติ. สติพลํ ภาเวติ. สมาธิพลํ ภาเวติ. ปญฺญาพลํ ภาเวติ ฯเปฯ. (5)}

\prose{สติสมฺโพชฺฌงฺคํ ภาเวติ. ธมฺมวิจย\-สมฺโพชฺฌงฺคํ ภาเวติ. วีริย\-สมฺโพชฺฌงฺคํ ภาเวติ. ปีติสมฺโพชฺฌงฺคํ ภาเวติ. ปสฺสทฺธิ\-สมฺโพชฺฌงฺคํ ภาเวติ. สมาธิ\-สมฺโพชฺฌงฺคํ ภาเวติ. อุเปกฺขา\-สมฺโพชฺฌงฺคํ ภาเวติ ฯเปฯ. (6)}

\prose{สมฺมาทิฏฺฐิํ ภาเวติ. สมฺมาสงฺกปฺปํ ภาเวติ. สมฺมาวาจํ ภาเวติ. สมฺมากมฺมนฺตํ ภาเวติ. สมฺมาอาชีวํ ภาเวติ. สมฺมาวายามํ ภาเวติ. สมฺมาสติํ ภาเวติ. สมฺมาสมาธิํ ภาเวติ. อยํ วุจฺจติ ภิกฺขเว มชฺฌิมา ปฏิปทา. อิมา โข ภิกฺขเว ติสฺโส ปฏิปทาติ. (7) อเจลกวคฺโค ฉฏฺโฐ.}

\vspace{\tipitakaparskip}
\csromancenter{\textbf{ตสฺสุทฺทานํ}}
\setlength{\csromangathapagewidth}{0pt}
\setlength{\csromangathawakwidth}{0pt}
\csromangathameasuregroup{สติปฏฺฐานํ สมฺมปฺปธานํ, \\ พลํ โพชฺฌงฺโค มคฺโค จ,}{\csromangathabat{สติปฏฺฐานํ สมฺมปฺปธานํ,}{อิทฺธิปาทินฺทฺริเยน จ.} \\ \csromangathabat{พลํ โพชฺฌงฺโค มคฺโค จ,}{ปฏิปทาย โยชเยติ.} \\ (17) 7. \\ กมฺมปถเปยฺยาล}
\csromangathameasurewak{(17) 7. \\ กมฺมปถเปยฺยาล}
\csromangathasetleft{สติปฏฺฐานํ สมฺมปฺปธานํ, \\ พลํ โพชฺฌงฺโค มคฺโค จ,}
\csromangathagroup{\csromangathastanza{\csromangathabat{สติปฏฺฐานํ สมฺมปฺปธานํ,}{อิทฺธิปาทินฺทฺริเยน จ.} \\ \csromangathabat{พลํ โพชฺฌงฺโค มคฺโค จ,}{ปฏิปทาย โยชเยติ.}}\csromangathastanza{\csromangathawak{(17) 7. \\ กมฺมปถ\-เปยฺยาล}}}

\proseitem{164-183}{ตีหิ ภิกฺขเว ธมฺเมหิ สมนฺนาคโต ยถาภตํ นิกฺขิตฺโต เอวํ นิรเย. กตเมหิ ตีหิ? อตฺตนา จ ปาณาติปาตี โหติ, ปรญฺจ ปาณาติปาเต สมาทเปติ, ปาณาติปาเต จ สมนุญฺโญ โหติ. อิเมหิ โข ภิกฺขเว ตีหิ ธมฺเมหิ สมนฺนาคโต ยถาภตํ นิกฺขิตฺโต เอวํ นิรเย.}

\prose{ตีหิ ภิกฺขเว ธมฺเมหิ สมนฺนาคโต ยถาภตํ นิกฺขิตฺโต เอวํ สคฺเค. กตเมหิ ตีหิ? อตฺตนา จ ปาณาติปาตา ปฏิวิรโต โหติ, ปรญฺจ ปาณาติปาตา เวรมณิยา สมาทเปติ, ปาณาติปาตา เวรมณิยา จ สมนุญฺโญ โหติ ฯเปฯ. (2)}

\prose{อตฺตนา จ อทินฺนาทายี โหติ, ปรญฺจ อทินฺนาทาเน สมาทเปติ, อทินฺนาทาเน จ สมนุญฺโญ โหติฯเปฯ.}

\prose{อตฺตนา จ อทินฺนาทานา ปฏิวิรโต โหติ, ปรญฺจ อทินฺนาทานา เวรมณิยา สมาทเปติ, อทินฺนาทานา เวรมณิยา จ สมนุญฺโญ โหติ ฯเปฯ. (4)}

\vspace{\tipitakaparskip}
\csromancenter{(17) 7. กมฺมปถ\-เปยฺยาล}
\prose{\csromanfolio{303}อตฺตนา จ กาเมสุ\-มิจฺฉาจารี โหติ, ปรญฺจ กาเมสุ\-มิจฺฉาจาเร สมาทเปติ, กาเมสุ\-มิจฺฉาจาเร จ สมนุญฺโญ โหติ ฯเปฯ.}

\prose{อตฺตนา จ กาเมสุ\-มิจฺฉาจารา ปฏิวิรโต โหติ, ปรญฺจ กาเมสุ\-มิจฺฉาจารา เวรมณิยา สมาทเปติ, กาเมสุ\-มิจฺฉาจารา เวรมณิยา จ สมนุญฺโญ โหติ ฯเปฯ. (6)}

\prose{อตฺตนา จ มุสาวาที โหติ, ปรญฺจ มุสาวาเท สมาทเปติ, มุสาวาเท จ สมนุญฺโญ โหติ ฯเปฯ.}

\prose{อตฺตนา จ มุสาวาทา ปฏิวิรโต โหติ, ปรญฺจ มุสาวาทา เวรมณิยา สมาทเปติ, มุสาวาทา เวรมณิยา จ สมนุญฺโญ โหติ ฯเปฯ. (8)}

\prose{อตฺตนา จ ปิสุณวาโจ โหติ, ปรญฺจ ปิสุณาย วาจาย สมาทเปติ, ปิสุณาย วาจาย จ สมนุญฺโญ โหติ ฯเปฯ.}

\prose{อตฺตนา จ ปิสุณาย วาจาย ปฏิวิรโต โหติ, ปรญฺจ ปิสุณาย วาจาย เวรมณิยา สมาทเปติ, ปิสุณาย วาจาย เวรมณิยา จ สมนุญฺโญ โหติ ฯเปฯ. (10)}

\prose{อตฺตนา จ ผรุสวาโจ โหติ, ปรญฺจ ผรุสาย วาจาย สมาทเปติ, ผรุสาย วาจาย จ สมนุญฺโญ โหติ ฯเปฯ.}

\prose{อตฺตนา จ ผรุสาย วาจาย ปฏิวิรโต โหติ, ปรญฺจ ผรุสาย วาจาย เวรมณิยา สมาทเปติ, ผรุสาย วาจาย เวรมณิยา จ สมนุญฺโญ โหติ ฯเปฯ. (12)}

\prose{อตฺตนา จ สมฺผปฺปลาปี โหติ, ปรญฺจ สมฺผปฺปลาเป สมาทเปติ, สมฺผปฺปลาเป จ สมนุญฺโญ โหติ ฯเปฯ.}

\prose{อตฺตนา จ สมฺผปฺปลาปา ปฏิวิรโต โหติ, ปรญฺจ สมฺผปฺปลาปา เวรมณิยา สมาทเปติ, สมฺผปฺปลาปา เวรมณิยา จ สมนุญฺโญ โหติ ฯเปฯ. (14)}

\prose{อตฺตนา จ อภิชฺฌาลุ โหติ, ปรญฺจ อภิชฺฌาย สมาทเปติ, อภิชฺฌาย จ สมนุญฺโญ โหติ ฯเปฯ.}

\prose{อตฺตนา จ อนภิชฺฌาลุ โหติ, ปรญฺจ อนภิชฺฌาย สมาทเปติ, อนภิชฺฌาย จ สมนุญฺโญ โหติ ฯเปฯ. (16)}

\csromanmatikaanchor{mtk.49}
\csromanmatikaanchor{mtk.50}
\csromantocmarkref{section}{อุทฺทานคาถา}{mtk.49}
\bookmark[dest={mtk.49},level=1]{อุทฺทานคาถา}
\csromantocmarkref{section}{(18) 8. ราคเปยฺยาล}{mtk.50}
\bookmark[dest={mtk.50},level=1]{(18) 8. ราคเปยฺยาล}
\prose{\csromanfolio{304}อตฺตนา จ พฺยาปนฺนจิตฺโต โหติ, ปรญฺจ พฺยาปาเท สมาทเปติ, พฺยาปาเท จ สมนุญฺโญ โหติ ฯเปฯ.}

\prose{อตฺตนา จ อพฺยาปนฺนจิตฺโต โหติ, ปรญฺจ อพฺยาปาเท สมาทเปติ, อพฺยาปาเท จ สมนุญฺโญ โหติ ฯเปฯ. (18)}

\prose{อตฺตนา จ มิจฺฉาทิฏฺฐิโก โหติ, ปรญฺจ มิจฺฉาทิฏฺฐิยา สมาทเปติ, มิจฺฉาทิฏฺฐิยา จ สมนุญฺโญ โหติ ฯเปฯ.}

\prosewithcloser{อตฺตนา จ สมฺมาทิฏฺฐิโก โหติ, ปรญฺจ สมฺมาทิฏฺฐิยา สมาทเปติ, สมฺมาทิฏฺฐิยา จ สมนุญฺโญ โหติ. อิเมหิ โข ภิกฺขเว ตีหิ ธมฺเมหิ สมนฺนาคโต ยถาภตํ นิกฺขิตฺโต เอวํ สคฺเคติ.}{(20) กมฺมปถ\-เปยฺยาลํ นิฏฺฐิตํ.}
\csromancenter{\textbf{ตสฺสุทฺทานํ}}
\setlength{\csromangathapagewidth}{0pt}
\setlength{\csromangathawakwidth}{0pt}
\csromangathameasuregroup{ปาณํ อทินฺนมิจฺฉา จ, \\ ผรุสา สมฺผปฺปลาโป จ, \\ กมฺมปเถสุ เปยฺยาลํ,}{\csromangathabat{ปาณํ อทินฺนมิจฺฉา จ,}{มุสาวาที จ ปิสุณา.} \\ \csromangathabat{ผรุสา สมฺผปฺปลาโป จ,}{อภิชฺฌา พฺยาปาททิฏฺฐิ จ.} \\ \csromangathabat{กมฺมปเถสุ เปยฺยาลํ,}{ติกเกน นิโยชเยติ.} \\ (18) 8. \\ ราคเปยฺยาล}
\csromangathasetleft{ปาณํ อทินฺนมิจฺฉา จ, \\ ผรุสา สมฺผปฺปลาโป จ, \\ กมฺมปเถสุ เปยฺยาลํ,}
\csromangathagroup{\csromangathastanza{\csromangathabat{ปาณํ อทินฺนมิจฺฉา จ,}{มุสาวาที จ ปิสุณา.} \\ \csromangathabat{ผรุสา สมฺผปฺปลาโป จ,}{อภิชฺฌา พฺยาปาททิฏฺฐิ จ.}}\csromangathastanza{\csromangathabat{กมฺมปเถสุ เปยฺยาลํ,}{ติกเกน นิโยชเยติ.} \\ (18) 8. \\ ราคเปยฺยาล}}

\proseitem{184}{ราคสฺส ภิกฺขเว อภิญฺญาย ตโย ธมฺมา ภาเวตพฺพา. กตเม ตโย? สุญฺญโต สมาธิ, อนมิตฺโต สมาธิ, อปฺปณิหิโต สมาธิ. ราคสฺส ภิกฺขเว อภิญฺญาย อิเม ตโย ธมฺมา ภาเวตพฺพา. (1) ( )\csromansharedfootnote{65dd646d0cfdb1c7}{(ราคสฺส ภิกฺขเว อภิญฺญาย ตโย ธมฺมา ภาเวตพฺพา. กตเม ตโย? สวิตกฺกสวิจาโร สมาธิ, อวิตกฺกวิจารมตฺโต สมาธิ, อวิตกฺกอวิจาโร สมาธิ. ราคสฺส ภิกฺขเว อภิญฺญาย อิเม ตโย ธมฺมา ภาเวตพฺพา.) เอตฺถนฺตเร ปาโฐ กตฺถจิ ทิสฺสติ, อฏฺฐกถายํ ปสฺสิตพฺโพ.}}

\vspace{\tipitakaparskip}
\csromancenter{(18) 8. ราคเปยฺยาล}
\csromanmatikaanchor{mtk.51}
\csromantocmarkref{section}{อุทฺทานคาถา}{mtk.51}
\bookmark[dest={mtk.51},level=1]{อุทฺทานคาถา}
\prose{\csromanfolio{305}ราคสฺส ภิกฺขเว ปริญฺญาย ฯเปฯ ปริกฺขยาย. ปหานาย. ขยาย. วยาย. วิราคาย. นิโรธาย. จาคาย. ปฏินิสฺสคฺคาย. อิเม ตโย ธมฺมา ภาเวตพฺพา.}

\prose{โทสสฺส. โมหสฺส. โกธสฺส. อุปนาหสฺส. มกฺขสฺส. ปฬาสสฺส. อิสฺสาย. มจฺฉริยสฺส. มายาย. สาเฐยฺยสฺส. ถมฺภสฺส. สารมฺภสฺส. มานสฺส. อติมานสฺส. มทสฺส. ปมาทสฺส. อภิญฺญาย. ปริญฺญาย. ปริกฺขยาย. ปหานาย. ขยาย. วยาย. วิราคาย. นิโรธาย. จาคาย. ปฏินิสฺสคฺคาย. อิเม ตโย ธมฺมา ภาเวตพฺพาติ.}

\prosewithcloser{(อิทมโวจ ภควา.}{อตฺตมนา เต ภิกฺขู ภควโต ภาสิตํ อภินนฺทุนฺติ.)\csromansharedfootnote{338dd9d64fd7ee1e}{[มิสฺสินฺคฺ โนเต 0]} ราคเปยฺยาลํ นิฏฺฐิตํ.}
\csromancenter{\textbf{ตสฺสุทฺทานํ}}
\setlength{\csromangathapagewidth}{0pt}
\setlength{\csromangathawakwidth}{0pt}
\csromangathameasuregroup{2 ราคํ โทสญฺจ โมหญฺจ, \\ มกฺขปฬาสอิสฺสา จ, \\ ถมฺภสารมฺภมานญฺจ, \\ ปมาทา สตฺตรส วุตฺตา, \\ เอเต โอปมฺมยุตฺเตน, \\ ปริญฺญาย ปริกฺขยา, \\ วิราคนิโรธจาคํ, \\ สุญฺญโต อนิมิตฺโต จ,}{\csromangathabat{2 ราคํ โทสญฺจ โมหญฺจ,}{โกธูปนาหปญฺจมํ.} \\ \csromangathabat{มกฺขปฬาสอิสฺสา จ,}{มจฺฉริมายาสาเฐยฺยา.} \\ \csromangathabat{ถมฺภสารมฺภมานญฺจ,}{อติมานมทสฺส จ.} \\ \csromangathabat{ปมาทา สตฺตรส วุตฺตา,}{ราคเปยฺยาลนิสฺสิตา.} \\ \csromangathabat{เอเต โอปมฺมยุตฺเตน,}{อาปาเทน อภิญฺญาย.} \\ \csromangathabat{ปริญฺญาย ปริกฺขยา,}{ปหานกฺขยพฺพเยน.} \\ \csromangathabat{วิราคนิโรธจาคํ,}{ปฏินิสฺสคฺเค อิเม ทส.} \\ \csromangathabat{สุญฺญโต อนิมิตฺโต จ,}{อปฺปณิหิโต จ ตโย.} \\ สมาธิมูลกา เปยฺยาเลสุปิ ววตฺถิตา จาติ.}
\csromangathasetleft{2 ราคํ โทสญฺจ โมหญฺจ, \\ มกฺขปฬาสอิสฺสา จ, \\ ถมฺภสารมฺภมานญฺจ, \\ ปมาทา สตฺตรส วุตฺตา, \\ เอเต โอปมฺมยุตฺเตน, \\ ปริญฺญาย ปริกฺขยา, \\ วิราคนิโรธจาคํ, \\ สุญฺญโต อนิมิตฺโต จ,}
\csromangathagroupwithclosermajor{\csromangathastanza{\csromangathabat{2 ราคํ โทสญฺจ โมหญฺจ,}{โกธูปนาหปญฺจมํ.} \\ \csromangathabat{มกฺขปฬาสอิสฺสา จ,}{มจฺฉริมายาสาเฐยฺยา.}}\csromangathastanza{\csromangathabat{ถมฺภสารมฺภมานญฺจ,}{อติมานมทสฺส จ.} \\ \csromangathabat{ปมาทา สตฺตรส วุตฺตา,}{ราคเปยฺยาลนิสฺสิตา.}}\csromangathastanza{\csromangathabat{เอเต โอปมฺมยุตฺเตน,}{อาปาเทน อภิญฺญาย.} \\ \csromangathabat{ปริญฺญาย ปริกฺขยา,}{ปหานกฺขยพฺพเยน.}}\csromangathastanza{\csromangathabat{วิราคนิโรธจาคํ,}{ปฏินิสฺสคฺเค อิเม ทส.} \\ \csromangathabat{สุญฺญโต อนิมิตฺโต จ,}{อปฺปณิหิโต จ ตโย.} \\ สมาธิมูลกา เปยฺยาเลสุปิ ววตฺถิตา จาติ.}}{\textbf{ติกนิปาตปาฬิ นิฏฺฐิตา.}}
\csromanpalirecto
\csromanmatikaanchor{mtk.52}
\csromantocmarknopage{chapter}{จตุกฺกนิปาตปาฬิ}{mtk.52}
\bookmark[dest={mtk.52},level=0]{จตุกฺกนิปาตปาฬิ}
\csromantocmarknopage{chapter}{1. ปฐมปณฺณาสก}{mtk.53}
\bookmark[dest={mtk.53},level=0]{1. ปฐมปณฺณาสก}
\gambhira{จตุกฺกนิปาต\-ปาฬิ}
\namakkaram{นโม ตสฺส ภควโต อรหโต สมฺมา\-สมฺพุทฺธสฺส.}
\csromanheader{2}{1. ปฐม\-ปณฺณาสก}
\chapterhead{1. ภณฺฑคาม\-วคฺค}
\csromanmatikaanchor{mtk.54}
\csromanmatikaanchor{mtk.55}
\csromantocmarknopage{section}{1. ภณฺฑคามวคฺค}{mtk.54}
\bookmark[dest={mtk.54},level=1]{1. ภณฺฑคามวคฺค}
\csromantocmarkref{subsection}{1. อนุพุทฺธสุตฺต}{mtk.55}
\bookmark[dest={mtk.55},level=2]{1. อนุพุทฺธสุตฺต}
\csromanheader{2}{1. \textbf{อนุพุทฺธสุตฺต}}
\proseitem{1}{\csromanfolio{307}เอวํ เม สุตํ\csromandash{}เอกํ สมยํ ภควา วชฺชีสุ วิหรติ ภณฺฑคาเม. ตตฺร โข ภควา ภิกฺขู อามนฺเตสิ “ภิกฺขโว”ติ. “ภทนฺเต”ติ เต ภิกฺขู ภควโต ปจฺจสฺโสสุํ. ภควา เอตทโวจ\csromandash{}}

\prose{จตุนฺนํ ภิกฺขเว ธมฺมานํ อนนุโพธา อปฺปฏิเวธา เอวมิทํ ทีฆมทฺธานํ สนฺธาวิตํ สํสริตํ มมญฺเจว ตุมฺหากญฺจ. กตเมสํ จตุนฺนํ? อริยสฺส ภิกฺขเว สีลสฺส อนนุโพธา อปฺปฏิเวธา เอวมิทํ ทีฆมทฺธานํ สนฺธาวิตํ สํสริตํ มมญฺเจว ตุมฺหากญฺจ. อริยสฺส ภิกฺขเว สมาธิสฺส อนนุโพธา อปฺปฏิเวธา เอวมิทํ ทีฆมทฺธานํ สนฺธาวิตํ สํสริตํ มมญฺจว ตุมฺหากญฺจ. อริยาย ภิกฺขเว ปญฺญาย อนนุโพธา อปฺปฏิเวธา เอวมิทํ ทีฆมทฺธานํ สนฺธาวิตํ สํสริตํ มมญฺเจว ตุมฺหากญฺจ. อริยาย ภิกฺขเว วิมุตฺติยา อนนุโพธา อปฺปฏิเวธา เอวมิทํ ทีฆมทฺธานํ สนฺธาวิตํ สํสริตํ มมญฺเจว ตุมฺหากญฺจ. ตยิทํ ภิกฺขเว อริยํ สีลํ อนุพุทฺธํ ปฏิวิทฺธํ, อริโย สมาธิ อนุพุทฺธา ปฏิวิทฺโธ, อริยา ปญฺญา อนุพุทฺธา ปฏิวิทฺธา, อริยา วิมุตฺติ อนุพุทฺธา ปฏิวิทฺธา, อุจฺฉินฺนา ภวตณฺหา, ขีณฺหา ภวเนตฺติ, นตฺถิ ทานิ ปุนพฺภโวติ.}

\prose{\csromanfolio{308}อิทมโวจ ภควา. อิทํ วตฺวาน สุคโต อถาปรํ เอตทโวจ สตฺถา\csromandash{}}

\prose{“สีลํ สมาธิ ปญฺญา จ, วิมุตฺติ จ อนุตฺตรา.}

\prose{อนุพุทฺธา อิเม ธมฺมา, โคตเมน ยสสฺสินา.}

\prose{อิติ พุทฺโธ อภิญฺญาย, ธมฺมมกฺขาสิ ภิกฺขุนํ.}

\prose{ทุกฺขสฺส\-นฺตกโร สตฺถา, จกฺขุมา ปรินิพฺพุโต”ติ.  ปฐมํ.}

\csromanmatikaanchor{mtk.56}
\csromantocmarkref{subsection}{2. ปปติตสุตฺต}{mtk.56}
\bookmark[dest={mtk.56},level=2]{2. ปปติตสุตฺต}
\csromanheader{2}{2. \textbf{ปปติตสุตฺต}}
\proseitem{2}{จตูหิ ภิกฺขเว ธมฺเมหิ อสมนฺนาคโต “อิมสฺมา ธมฺมวินยา ปปติโต”ติ วุจฺจติ. กตเมหิ จตูหิ? อริเยน ภิกฺขเว สีเลน อสมนฺนาคโต “อิมสฺมา ธมฺมวินยา ปปติโต”ติ วุจฺจติ, อริเยน ภิกฺขเว สมาธินา อสมนฺนาคโต “อิมสฺมา ธมฺมวินยา ปปติโต”ติ วุจฺจติ, อริยาย ภิกฺขเว ปญฺญาย อสมนฺนาคโต “อิมสฺมา ธมฺมวินยา ปปติโต”ติ วุจฺจติ, อริยาย ภิกฺขเว วิมุตฺติยา อสมนฺนาคโต “อิมสฺมา ธมฺมวินยา ปปติโต”ติ วุจฺจติ. อิเมหิ โข ภิกฺขเว จตูหิ ธมฺเมหิ อสมนฺนาคโต “อิมสฺมา ธมฺมวินยา ปปติโต”ติ วุจฺจติ.}

\prose{จตูหิ ภิกฺขเว ธมฺเมหิ สมนฺนาคโต อิมสฺมา ธมฺมวินยา “อปปติโต”ติ\csromansharedfootnote{8c1c84ead0038c5e}{อปฺปปติโตติ (ก)} วุจฺจติ. กตเมหิ จตูหิ? อริเยน ภิกฺขเว สีเลน สมนฺนาคโต “อิมสฺมา ธมฺมวินยา อปปติโต”ติ วุจฺจติ, อริเยน ภิกฺขเว สมาธินา สมนฺนาคโต “อิมสฺมา ธมฺมวินยา อปปติโต”ติ วุจฺจติ, อริยาย ภิกฺขเว ปญฺญาย สมนฺนาคโต “อิมสฺมา ธมฺมวินยา อปปติโต”ติ วุจฺจติ, อริยาย ภิกฺขเว วิมุตฺติยา สมนฺนาคโต “อิมสฺมา ธมฺมวินยา อปปติโต”ติ วุจฺจติ. อิเมหิ โข ภิกฺขเว จตูหิ ธมฺเมหิ สมนฺนาคโต “อิมสฺมา ธมฺมวินยา อปฺปปติโต”ติ วุจฺจตีติ.}

\prose{จุตา ปตนฺติ ปติตา, คิทฺธา จ ปุนราคตา.}

\prose{กตํ กิจฺจํ รตํ รมฺมํ, สุเขนา\-นฺวาคตํ สุขนฺติ.  ทุติยํ.}

\csromanmarkh{1. ภณฺฑคาม\-วคฺค}\csromanmarkhii{3. ปฐม\-ขต\-สุตฺต}\csromanheadernomark{1}{1. ภณฺฑคาม\-วคฺค 3. ปฐม\-ขต\-สุตฺต}
\csromanmatikaanchor{mtk.57}
\csromantocmarkref{subsection}{3-4. ขตสุตฺตานิ}{mtk.57}
\bookmark[dest={mtk.57},level=2]{3-4. ขตสุตฺตานิ}
\proseitem{3}{\csromanfolio{309}จตูหิ ภิกฺขเว ธมฺเมหิ สมนฺนาคโต พาโล อพฺยตฺโต\csromansharedfootnote{2c2c65f0a772cf10}{อวฺยตฺโต (สี, อิ)} อสปฺปุริโส ขตํ อุปหตํ อตฺตานํ ปริหรติ, สาวชฺโช จ โหติ สานุวชฺโช จ วิญฺญูนํ, พหุญฺจ อปุญฺญํ ปสวติ. กตเมหิ จตูหิ? อนนุวิจฺจ อปริโยคาเหตฺวา อวณฺณารหสฺส วณฺณํ ภาสติ, อนนุวิจฺจ อปริโยคาเหตฺวา วณฺณารหสฺส อวณฺณํ ภาสติ, อนนุวิจฺจ อปริโยคาเหตฺวา อปฺปสาทนีเย ฐาเน ปสาทํ อุปทํเสติ, อนนุวิจฺจ อปริโยคาเหตฺวา ปสาทนีเย ฐาเน อปฺปสาทํ อุปทํเสติ. อิเมหิ โข ภิกฺขเว จตูหิ ธมฺเมหิ สมนฺนาคโต พาโล อพฺยตฺโต อสปฺปุริโส ขตํ อุปหตํ อตฺตานํ ปริหรติ, สาวชฺโช จ โหติ สานุวชฺโช จ วิญฺญูนํ, พหุญฺจ อปุญฺญํ ปสวติ.}

\prose{จตูหิ ภิกฺขเว ธมฺเมหิ สมนฺนาคโต ปณฺฑิโต วิยตฺโต\csromansharedfootnote{9c29039aa05ec2d1}{วฺยตฺโต (สี, อิ), พฺยตฺโต (สฺยา, กํ)} สปฺปุริโส อกฺขตํ อนุปหตํ อตฺตานํ ปริหรติ, อนวชฺโช จ โหติ อนนุวชฺโช จ วิญฺญูนํ, พหุญฺจ ปุญฺญํ ปสวติ. กตเมหิ จตูหิ? อนุวิจฺจ ปริโยคาเหตฺวา อวณฺณารหสฺส อวณฺณํ ภาสติ, อนุวิจฺจ ปริโยคาเหตฺวา วณฺณารหสฺส วณฺณํ ภาสติ, อนุวิจฺจ ปริโยคาเหตฺวา อปฺปสาทนีเย ฐาเน อปฺปสาทํ อุปทํเสติ, อนุวิจฺจ ปริโยคาเหตฺวา ปสาทนีเย ฐาเน ปสาทํ อุปทํเสติ. อิเมหิ โข ภิกฺขเว จตูหิ ธมฺเมหิ สมนฺนาคโต ปณฺฑิโต วิยตฺโต สปฺปุริโส อกฺขตํ อนุปหตํ อตฺตานํ ปริหรติ, อนวชฺโช จ โหติ อนนุวชฺโช จ วิญฺญูนํ, พหุญฺจ ปุญฺญํ ปสวตีติ.}

\setlength{\csromangathapagewidth}{0pt}
\setlength{\csromangathawakwidth}{0pt}
\csromangathameasuregroup{}{\textsuperscript{*}โย นินฺทิยํ ปสํสติ, \\ ตํ วา นินฺทติ โย ปสํสิโย. \\ วิจินาติ มุเขน โส กลิํ, \\ กลินา เตน สุขํ น วินฺทติ. \\ \textsuperscript{*}อปฺปมตฺโต อยํ กลิ, \\ โย อกฺเขสุ ธนปราชโย. \\ สพฺพสฺสาปิ สหาปิ อตฺตนา, \\ อยเมว มหนฺตตโร กลิ. \\ โย สุคเตสุ มนํ ปโทสเย. \\ สตํ สหสฺสานํ นิรพฺพุทานํ, \\ ฉตฺติํสตี ปญฺจ จ อพฺพุทานิ. \\ ยมริยครหี\textsuperscript{1} นิรยํ อุเปติ, \\ วาจํ มนญฺจ ปณิธาย ปาปกนฺติ.ตติยํ.}
\csromangathameasurewak{\textsuperscript{*}โย นินฺทิยํ ปสํสติ, \\ ตํ วา นินฺทติ โย ปสํสิโย. \\ วิจินาติ มุเขน โส กลิํ, \\ กลินา เตน สุขํ น วินฺทติ. \\ \textsuperscript{*}อปฺปมตฺโต อยํ กลิ, \\ โย อกฺเขสุ ธนปราชโย. \\ สพฺพสฺสาปิ สหาปิ อตฺตนา, \\ อยเมว มหนฺตตโร กลิ. \\ สตํ สหสฺสานํ นิรพฺพุทานํ, \\ ฉตฺติํสตี ปญฺจ จ อพฺพุทานิ. \\ ยมริยครหี\textsuperscript{1} นิรยํ อุเปติ, \\ วาจํ มนญฺจ ปณิธาย ปาปกนฺติ.ตติยํ.}
\setlength{\csromangathaleftcol}{0pt}
\csromangathagroup{\csromangathastanza{\csromangathawak{\csromansymbolfootnote{*}{ขุ 1. 381; สํ 1. 151 ปิฏฺเฐสุปิ.}โย นินฺทิยํ ปสํสติ, \\ ตํ วา นินฺทติ โย ปสํสิโย. \\ วิจินาติ มุเขน โส กลิํ, \\ กลินา เตน สุขํ น วินฺทติ.}}\csromangathastanza{\csromangathawak{\csromansymbolmark{*}อปฺปมตฺโต อยํ กลิ, \\ โย อกฺเขสุ ธนปราชโย. \\ สพฺพสฺสาปิ สหาปิ อตฺตนา, \\ อยเมว มหนฺตตโร กลิ.}}\csromangathastanza{โย สุคเตสุ มนํ ปโทสเย.}\csromangathastanza{\csromangathawak{\csromanfolio{310}สตํ สหสฺสานํ นิรพฺพุทานํ, \\ ฉตฺติํสตี ปญฺจ จ อพฺพุทานิ. \\ ยมริยครหี\csromansharedfootnote{8cb20bfa2a75417d}{ยมริยํ ครหีย (สฺยา, กํ)} นิรยํ อุเปติ, \\ วาจํ มนญฺจ ปณิธาย ปาปกนฺติ.  ตติยํ.}}}

\csromanheader{2}{4. ทุติยขต\-สุตฺต}
\proseitem{4}{จตูสุ ภิกฺขเว มิจฺฉา ปฏิปชฺชมาโน พาโล อพฺยตฺโต อสปฺปุริโส ขตํ อุปหตํ อตฺตานํ ปริหรติ, สาวชฺโช จ โหติ สานุวชฺโช จ วิญฺญูนํ, พหุญฺจ อปุญฺญํ ปสวติ. กตเมสุ จตูสุ? มาตริ ภิกฺขเว มิจฺฉา ปฏิปชฺชมาโน พาโล อพฺยตฺโต อสปฺปุริโส ขตํ อุปหตํ อตฺตานํ ปริหรติ, สาวชฺโช จ โหติ สานุวชฺโช จ วิญฺญูนํ, พหุญฺจ อปุญฺญํ ปสวติ. ปิตริ ภิกฺขเว มิจฺฉา ปฏิปชฺชมาโน ฯเปฯ. ตถาคเต ภิกฺขเว มิจฺฉา ปฏิปชฺชมาโน ฯเปฯ. ตถาคต\-สาวเก ภิกฺขเว มิจฺฉา ปฏิปชฺชมาโน พาโล อพฺยตฺโต อสปฺปุริโส ขตํ อุปหตํ อตฺตานํ ปริหรติ, สาวชฺโช จ โหติ สานุวชฺโช จ วิญฺญูนํ, พหุญฺจ อปุญฺญํ ปสวติ. อิเมสุ โข ภิกฺขเว จตูสุ มิจฺฉา ปฏิปชฺชมาโน พาโล อพฺยตฺโต อสปฺปุริโส ขตํ อุปหตํ อตฺตานํ ปริหรติ, สาวชฺโช จ โหติ สานุวชฺโช จ วิญฺญูนํ, พหุญฺจ อปุญฺญํ ปสวติ.}

\prose{จตูสุ ภิกฺขเว สมฺมา ปฏิปชฺชมาโน ปณฺฑิโต วิยตฺโต สปฺปุริโส อกฺขตํ อนุปหตํ อตฺตานํ ปริหรติ, อนวชฺโช จ โหติ อนนุวชฺโช จ วิญฺญูนํ, พหุญฺจ ปุญฺญํ ปสวติ. กตเมสุ จตูสุ? มาตริ ภิกฺขเว สมฺมา ปฏิปชฺชมาโน ปณฺฑิโต วิยตฺโต สปฺปุริโส อกฺขตํ อนุปหตํ อตฺตานํ ปริหรติ, อนวชฺโช จ โหติ อนนุวชฺโช จ วิญฺญูนํ, พหุญฺจ ปุญฺญํ ปสวติ. ปิตริ ภิกฺขเว สมฺมา ปฏิปชฺชมาโน ฯเปฯ. ตถาคเต ภิกฺขเว สมฺมา ปฏิปชฺชมาโน ฯเปฯ. ตถาคต\-สาวเก ภิกฺขเว สมฺมา ปฏิปชฺชมาโน ปณฺฑิโต วิยตฺโต สปฺปุริโส อกฺขตํ อนุปหตํ อตฺตานํ ปริหรติ, อนวชฺโช จ โหติ อนนุวชฺโช จ วิญฺญูนํ, พหุญฺจ ปุญฺญํ ปสวติ. อิเมสุ โข ภิกฺขเว จตูสุ สมฺมา ปฏิปชฺชมาโน ปณฺฑิโต วิยตฺโต สปฺปุริโส อกฺขตํ อนุปหตํ อตฺตานํ ปริหรติ, อนวชฺโช จ โหติ อนนุวชฺโช จ วิญฺญูนํ, พหุญฺจ ปุญฺญํ ปสวตีติ.}

\chapterheadpageread{1. ภณฺฑคาม\-วคฺค}
\setlength{\csromangathapagewidth}{0pt}
\setlength{\csromangathawakwidth}{0pt}
\csromangathameasuregroup{มาตริ ปิตริ จาปิ, \\ ตถาคเต วา สมฺพุทฺเธ, \\ พหุญฺจ โส ปสวติ, \\ ตาย นํ อธมฺมจริยาย\textsuperscript{1}, \\ อิเธว นํ ครหนฺติ, \\ มาตริ ปิตริ จาปิ, \\ ตถาคเต วา สมฺพุทฺเธ, \\ พหุญฺจ โส ปสวติ, \\ ตาย นํ ธมฺมจริยาย, \\ อิเธว\textsuperscript{1} นํ ปสํสนฺติ,}{\csromangathabat{มาตริ ปิตริ จาปิ,}{โย มิจฺฉา ปฏิปชฺชติ.} \\ \csromangathabat{ตถาคเต วา สมฺพุทฺเธ,}{อถ วา ตสฺส สาวเก.} \\ \csromangathabat{พหุญฺจ โส ปสวติ,}{อปุญฺญํ ตาทิโส นโร.} \\ \csromangathabat{ตาย นํ อธมฺมจริยาย\textsuperscript{1},}{มาตาปิตูสุ ปณฺฑิตา.} \\ \csromangathabat{อิเธว นํ ครหนฺติ,}{เปจฺจาปายญฺจ คจฺฉติ.} \\ \csromangathabat{มาตริ ปิตริ จาปิ,}{โย สมฺมา ปฏิปชฺชติ.} \\ \csromangathabat{ตถาคเต วา สมฺพุทฺเธ,}{อถ วา ตสฺส สาวเก.} \\ \csromangathabat{พหุญฺจ โส ปสวติ,}{ปุญฺญํ เอตาทิโส\textsuperscript{1} นโร.} \\ \csromangathabat{ตาย นํ ธมฺมจริยาย,}{มาตาปิตูสุ ปณฺฑิตา.} \\ \csromangathabat{อิเธว\textsuperscript{1} นํ ปสํสนฺติ,}{เปจฺจ สคฺเค ปโมทตีติ\textsuperscript{1}.จตุตฺถํ.}}
\csromangathasetleft{มาตริ ปิตริ จาปิ, \\ ตถาคเต วา สมฺพุทฺเธ, \\ พหุญฺจ โส ปสวติ, \\ ตาย นํ อธมฺมจริยาย\textsuperscript{1}, \\ อิเธว นํ ครหนฺติ, \\ มาตริ ปิตริ จาปิ, \\ ตถาคเต วา สมฺพุทฺเธ, \\ พหุญฺจ โส ปสวติ, \\ ตาย นํ ธมฺมจริยาย, \\ อิเธว\textsuperscript{1} นํ ปสํสนฺติ,}
\csromangathagroup{\csromangathastanza{\csromanfolio{311}\csromangathabat{มาตริ ปิตริ จาปิ,}{โย มิจฺฉา ปฏิปชฺชติ.} \\ \csromangathabat{ตถาคเต วา สมฺพุทฺเธ,}{อถ วา ตสฺส สาวเก.}}\csromangathastanza{\csromangathabat{พหุญฺจ โส ปสวติ,}{อปุญฺญํ ตาทิโส นโร.} \\ \csromangathabat{ตาย นํ อธมฺมจริยาย\csromansharedfootnote{83099e10cb0834ab}{ตาย อธมฺมจริยาย (สี, สฺยา, กํ, อิ)},}{มาตาปิตูสุ ปณฺฑิตา.}}\csromangathastanza{\csromangathabat{อิเธว นํ ครหนฺติ,}{เปจฺจาปายญฺจ คจฺฉติ.} \\ \csromangathabat{มาตริ ปิตริ จาปิ,}{โย สมฺมา ปฏิปชฺชติ.}}\csromangathastanza{\csromangathabat{ตถาคเต วา สมฺพุทฺเธ,}{อถ วา ตสฺส สาวเก.} \\ \csromangathabat{พหุญฺจ โส ปสวติ,}{ปุญฺญํ เอตาทิโส\csromansharedfootnote{b3333786422001a0}{ปุญฺญมฺปิ ตาทิโส (สี, สฺยา, กํ)} นโร.}}\csromangathastanza{\csromangathabat{ตาย นํ ธมฺมจริยาย,}{มาตาปิตูสุ ปณฺฑิตา.} \\ \csromangathabat{อิเธว\csromansharedfootnote{495f432540fc5213}{อิธ เจว (สี)} นํ ปสํสนฺติ,}{เปจฺจ สคฺเค ปโมทตีติ\csromansharedfootnote{6efdb498a5a247f5}{สคฺเค จ โมทตีติ (สี)}.  จตุตฺถํ.}}}

\csromanmatikaanchor{mtk.58}
\csromantocmarkref{subsection}{5. อนุโสตสุตฺต}{mtk.58}
\bookmark[dest={mtk.58},level=2]{5. อนุโสตสุตฺต}
\csromanheader{2}{5. \textbf{อนุโสตสุตฺต}}
\proseitem{5}{จตฺตาโรเม ภิกฺขเว ปุคฺคลา สนฺโต สํวิชฺชมานา โลกสฺมิํ. กตเม จตฺตาโร? อนุโสตคามี ปุคฺคโล, ปฏิโสตคามี ปุคฺคโล, ฐิตตฺโต ปุคฺคโล, ติณฺโณ ปารงฺคโต\csromansharedfootnote{355e7893e60b0adc}{ปารคโต (สี, สฺยา, กํ)} ถเล ติฏฺฐติ พฺราหฺมโณ. กตโม จ ภิกฺขเว อนุโสตคามี ปุคฺคโล? อิธ ภิกฺขเว เอกจฺโจ ปุคฺคโล กาเม จ ปฏิเสวติ, ปาปญฺจ กมฺมํ กโรติ. อยํ วุจฺจติ ภิกฺขเว อนุโสตคามี ปุคฺคโล.}

\prose{กตโม จ ภิกฺขเว ปฏิโสตคามี ปุคฺคโล? อิธ ภิกฺขเว เอกจฺโจ ปุคฺคโล กาเม จ นปฺปฏิเสวติ, ปาปญฺจ กมฺมํ น กโรติ, สหาปิ ทุกฺเขน สหาปิ โทมนสฺเสน อสฺสุมุโขปิ รุทมาโน ปริปุณฺณํ ปริสุทฺธํ พฺรหฺมจริยํ จรติ. อยํ วุจฺจติ ภิกฺขเว ปฏิโสตคามี ปุคฺคโล.}

\prose{กตโม จ ภิกฺขเว ฐิตตฺโต ปุคฺคโล? อิธ ภิกฺขเว เอกจฺโจ ปุคฺคโล ปญฺจนฺนํ โอรมฺภาคิยานํ สํโยชนานํ ปริกฺขยา โอปปาติโก โหติ ตตฺถ ปรินิพฺพายี อนาวตฺติธมฺโม ตสฺมา โลกา. อยํ วุจฺจติ ภิกฺขเว ฐิตตฺโต ปุคฺคโล.}

\prose{\csromanfolio{312}กตโม จ ภิกฺขเว ปุคฺคโล ติณฺโณ ปารงฺคโต ถเล ติฏฺฐติ พฺราหฺมโณ? อิธ ภิกฺขเว เอกจฺโจ ปุคฺคโล อาสวานํ ขยา อนาสวํ เจโตวิมุตฺติํ ปญฺญาวิมุตฺติํ ทิฏฺเฐว ธมฺเม สยํ อภิญฺญา สจฺฉิกตฺวา อุปสมฺปชฺช วิหรติ. อยํ วุจฺจติ ภิกฺขเว ปุคฺคโล ติณฺโณ ปารงฺคโต ถเล ติฏฺฐติ พฺราหฺมโณ. อิเม โข ภิกฺขเว จตฺตาโร ปุคฺคลา สนฺโต สํวิชฺชมานา โลกสฺมินฺติ.}

\setlength{\csromangathapagewidth}{0pt}
\setlength{\csromangathawakwidth}{0pt}
\csromangathameasuregroup{}{เย เกจิ กาเมสุ อสญฺญตา ชนา, \\ อวีตราคา อิธ กามโภคิโน. \\ ปุนปฺปุนํ ชาติชรูปคามิ เต\textsuperscript{1}, \\ ตณฺหาธิปนฺนา อนุโสตคามิโน. \\ ตสฺมา หิ ธีโร อิธุปฏฺฐิตสฺสตี, \\ กาเม จ ปาเป จ อเสวมาโน. \\ สหาปิ ทุกฺเขน ชเหยฺย กาเม, \\ “ปฏิโสตคามี”ติ ตมาหุ ปุคฺคลํ. \\ โย เว กิเลสานิ ปหาย ปญฺจ, \\ ปริปุณฺณเสโข อปริหานธมฺโม. \\ เจโตวสิปฺปตฺโต สมาหิตินฺทฺริโย, \\ ส เว “ฐิตตฺโต”ติ นโร ปวุจฺจติ. \\ ปโรปรา ยสฺส สเมจฺจ ธมฺมา, \\ วิธูปิตา อตฺถคตา น สนฺติ. \\ ส เว มุนิ\textsuperscript{1} วุสิตพฺรหฺมจริโย, \\ โลกนฺตคู “ปารคโต”ติ วุจฺจตีติ.ปญฺจมํ.}
\csromangathameasurewak{เย เกจิ กาเมสุ อสญฺญตา ชนา, \\ อวีตราคา อิธ กามโภคิโน. \\ ปุนปฺปุนํ ชาติชรูปคามิ เต\textsuperscript{1}, \\ ตณฺหาธิปนฺนา อนุโสตคามิโน. \\ ตสฺมา หิ ธีโร อิธุปฏฺฐิตสฺสตี, \\ กาเม จ ปาเป จ อเสวมาโน. \\ สหาปิ ทุกฺเขน ชเหยฺย กาเม, \\ “ปฏิโสตคามี”ติ ตมาหุ ปุคฺคลํ. \\ โย เว กิเลสานิ ปหาย ปญฺจ, \\ ปริปุณฺณเสโข อปริหานธมฺโม. \\ เจโตวสิปฺปตฺโต สมาหิตินฺทฺริโย, \\ ส เว “ฐิตตฺโต”ติ นโร ปวุจฺจติ. \\ ปโรปรา ยสฺส สเมจฺจ ธมฺมา, \\ วิธูปิตา อตฺถคตา น สนฺติ. \\ ส เว มุนิ\textsuperscript{1} วุสิตพฺรหฺมจริโย, \\ โลกนฺตคู “ปารคโต”ติ วุจฺจตีติ.ปญฺจมํ.}
\setlength{\csromangathaleftcol}{0pt}
\csromangathagroup{\csromangathastanza{\csromangathawak{เย เกจิ กาเมสุ อสญฺญตา ชนา, \\ อวีตราคา อิธ กามโภคิโน. \\ ปุนปฺปุนํ ชาติชรูปคามิ เต\csromansharedfootnote{8a704089920464af}{ชาติชรูปคาหิโน (สี), ชาติชรูปคา หิ เต (สฺยา, กํ)}, \\ ตณฺหาธิปนฺนา อนุโสตคามิโน.}}\csromangathastanza{\csromangathawak{ตสฺมา หิ ธีโร อิธุปฏฺฐิตสฺสตี, \\ กาเม จ ปาเป จ อเสวมาโน. \\ สหาปิ ทุกฺเขน ชเหยฺย กาเม, \\ “ปฏิโสตคามี”ติ ตมาหุ ปุคฺคลํ.}}\csromangathastanza{\csromangathawak{โย เว กิเลสานิ ปหาย ปญฺจ, \\ ปริปุณฺณเสโข อปริหาน\-ธมฺโม. \\ เจโตวสิปฺปตฺโต สมาหิตินฺทฺริโย, \\ ส เว “ฐิตตฺโต”ติ นโร ปวุจฺจติ.}}\csromangathastanza{\csromangathawak{ปโรปรา ยสฺส สเมจฺจ ธมฺมา, \\ วิธูปิตา อตฺถคตา น สนฺติ. \\ ส เว มุนิ\csromansharedfootnote{1bfdb24ca1eb4d44}{ส เวทคู (สี, สฺยา, กํ, อิ)} วุสิต\-พฺรหฺมจริโย, \\ โลกนฺตคู “ปารคโต”ติ วุจฺจตีติ.  ปญฺจมํ.}}}

\csromanmatikaanchor{mtk.59}
\csromantocmarkref{subsection}{6. อปฺปสฺสุตสุตฺต}{mtk.59}
\bookmark[dest={mtk.59},level=2]{6. อปฺปสฺสุตสุตฺต}
\csromanheader{2}{6. \textbf{อปฺปสฺสุตสุตฺต}}
\proseitem{6}{จตฺตาโรเม ภิกฺขเว ปุคฺคลา สนฺโต สํวิชฺชมานา โลกสฺมิํ. กตเม จตฺตาโร? อปฺปสฺสุโต สุเตน อนุปปนฺโน, อปฺปสฺสุโต สุเตน อุปปนฺโน, พหุสฺสุโต สุเตน อนุปปนฺโน, พหุสฺสุโต สุเตน อุปปนฺโน. กถญฺจ ภิกฺขเว ปุคฺคโล อปฺปสฺสุโต โหติ สุเตน \csromanfolio{313}อนุปปนฺโน? อิธ ภิกฺขเว เอกจฺจสฺส ปุคฺคลสฺส อปฺปกํ สุตํ โหติ สุตฺตํ เคยฺยํ เวยฺยากรณํ คาถา อุทานํ อิติวุตฺตกํ ชาตกํ อพฺภุตธมฺมํ เวทลฺลํ. โส ตสฺส อปฺปกสฺส สุตสฺส น อตฺถมญฺญาย ธมฺมมญฺญาย\csromansharedfootnote{86dbdf50b9c7ec1d}{น ธมฺมมญฺญาย (อิ, ก)} ธมฺมา\-นุธมฺม\-ปฺปฏิปนฺโน โหติ. เอวํ โข ภิกฺขเว ปุคฺคโล อปฺปสฺสุโต โหติ สุเตน อนุปปนฺโน.}

\prose{กถญฺจ ภิกฺขเว ปุคฺคโล อปฺปสฺสุโต โหติ สุเตน อุปปนฺโน? อิธ ภิกฺขเว เอกจฺจสฺส ปุคฺคลสฺส อปฺปกํ สุตํ โหติ สุตฺตํ เคยฺยํ เวยฺยากรณํ คาถา อุทานํ อิติวุตฺตกํ ชาตกํ อพฺภุตธมฺมํ เวทลฺลํ. โส ตสฺส อปฺปกสฺส สุตสฺส อตฺถมญฺญาย ธมฺมมญฺญาย ธมฺมา\-นุธมฺม\-ปฺปฏิปนฺโน โหติ. เอวํ โข ภิกฺขเว ปุคฺคโล อปฺปสฺสุโต โหติ สุเตน อุปปนฺโน.}

\prose{กถญฺจ ภิกฺขเว ปุคฺคโล พหุสฺสุโต โหติ สุเตน อนุปปนฺโน? อิธ ภิกฺขเว เอกจฺจสฺส ปุคฺคลสฺส พหุกํ สุตํ โหติ สุตฺตํ เคยฺยํ เวยฺยากรณํ คาถา อุทานํ อิติวุตฺตกํ ชาตกํ อพฺภุตธมฺมํ เวทลฺลํ. โส ตสฺส พหุกสฺส สุตสฺส น อตฺถมญฺญาย ธมฺมมญฺญาย\csromansharedfootnote{94467f14d202b935}{น ธมฺมมญฺญาย (อิ)} ธมฺมา\-นุธมฺม\-ปฺปฏิปนฺโน โหติ. เอวํ โข ภิกฺขเว ปุคฺคโล พหุสฺสุโต โหติ สุเตน อนุปปนฺโน.}

\prose{กถญฺจ ภิกฺขเว ปุคฺคโล พหุสฺสุโต โหติ สุเตน อุปปนฺโน? อิธ ภิกฺขเว เอกจฺจสฺส ปุคฺคลสฺส พหุกํ สุตํ โหติ สุตฺตํ เคยฺยํ เวยฺยากรณํ คาถา อุทานํ อิติวุตฺตกํ ชาตกํ อพฺภุตธมฺมํ เวทลฺลํ. โส ตสฺส พหุกสฺส สุตสฺส อตฺถมญฺญาย ธมฺมมญฺญาย ธมฺมา\-นุธมฺม\-ปฺปฏิปนฺโน โหติ. เอวํ โข ภิกฺขเว ปุคฺคโล พหุสฺสุโต โหติ สุเตน อุปปนฺโน. อิเม โข ภิกฺขเว จตฺตาโร ปุคฺคลา สนฺโต สํวิชฺชมานา โลกสฺมินฺติ.}

\setlength{\csromangathapagewidth}{0pt}
\setlength{\csromangathawakwidth}{0pt}
\csromangathameasuregroup{อปฺปสฺสุโตปิ เจ โหติ, \\ อุภเยน นํ ครหนฺติ, \\ อปฺปสฺสุโตปิ เจ โหติ, \\ สีลโต นํ ปสํสนฺติ, \\ พหุสฺสุโตปิ เจ โหติ, \\ สีลโต นํ ครหนฺติ, \\ พหุสฺสุโตปิ เจ โหติ, \\ อุภเยน นํ ปสํสนฺติ, \\ พหุสฺสุตํ ธมฺมธรํ, \\ เนกฺขํ ชมฺโพนทสฺเสว, \\ เทวาปิ นํ ปสํสนฺติ,}{\csromangathabat{อปฺปสฺสุโตปิ เจ โหติ,}{สีเลสุ อสมาหิโต.} \\ \csromangathabat{อุภเยน นํ ครหนฺติ,}{สีลโต จ สุเตน จ.} \\ \csromangathabat{อปฺปสฺสุโตปิ เจ โหติ,}{สีเลสุ สุสมาหิโต.} \\ \csromangathabat{สีลโต นํ ปสํสนฺติ,}{ตสฺส สมฺปชฺชเต สุตํ.} \\ \csromangathabat{พหุสฺสุโตปิ เจ โหติ,}{สีเลสุ อสมาหิโต.} \\ \csromangathabat{สีลโต นํ ครหนฺติ,}{นาสฺส สมฺปชฺชเต สุตํ.} \\ \csromangathabat{พหุสฺสุโตปิ เจ โหติ,}{สีเลสุ สุสมาหิโต.} \\ \csromangathabat{อุภเยน นํ ปสํสนฺติ,}{สีลโต จ สุเตน จ.} \\ \csromangathabat{พหุสฺสุตํ ธมฺมธรํ,}{สปฺปญฺญํ พุทฺธสาวกํ.} \\ \csromangathabat{เนกฺขํ ชมฺโพนทสฺเสว,}{โก ตํ นินฺทิตุมรหติ.} \\ \csromangathabat{เทวาปิ นํ ปสํสนฺติ,}{พฺรหฺมุนาปิ ปสํสิโตติ.ฉฏฺฐํ.}}
\csromangathasetleft{อปฺปสฺสุโตปิ เจ โหติ, \\ อุภเยน นํ ครหนฺติ, \\ อปฺปสฺสุโตปิ เจ โหติ, \\ สีลโต นํ ปสํสนฺติ, \\ พหุสฺสุโตปิ เจ โหติ, \\ สีลโต นํ ครหนฺติ, \\ พหุสฺสุโตปิ เจ โหติ, \\ อุภเยน นํ ปสํสนฺติ, \\ พหุสฺสุตํ ธมฺมธรํ, \\ เนกฺขํ ชมฺโพนทสฺเสว, \\ เทวาปิ นํ ปสํสนฺติ,}
\csromangathagroup{\csromangathastanza{\csromangathabat{อปฺปสฺสุโตปิ เจ โหติ,}{สีเลสุ อสมาหิโต.} \\ \csromangathabat{อุภเยน นํ ครหนฺติ,}{สีลโต จ สุเตน จ.}}\csromangathastanza{\csromangathabat{อปฺปสฺสุโตปิ เจ โหติ,}{สีเลสุ สุสมาหิโต.} \\ \csromangathabat{สีลโต นํ ปสํสนฺติ,}{ตสฺส สมฺปชฺชเต สุตํ.}}\csromangathastanza{\csromanfolio{314}\csromangathabat{พหุสฺสุโตปิ เจ โหติ,}{สีเลสุ อสมาหิโต.} \\ \csromangathabat{สีลโต นํ ครหนฺติ,}{นาสฺส สมฺปชฺชเต สุตํ.}}\csromangathastanza{\csromangathabat{พหุสฺสุโตปิ เจ โหติ,}{สีเลสุ สุสมาหิโต.} \\ \csromangathabat{อุภเยน นํ ปสํสนฺติ,}{สีลโต จ สุเตน จ.}}\csromangathastanza{\csromangathabat{พหุสฺสุตํ ธมฺมธรํ,}{สปฺปญฺญํ พุทฺธสาวกํ.} \\ \csromangathabat{เนกฺขํ ชมฺโพนทสฺเสว,}{โก ตํ นินฺทิตุม\-รหติ.} \\ \csromangathabat{เทวาปิ นํ ปสํสนฺติ,}{พฺรหฺมุนาปิ ปสํสิโตติ.  ฉฏฺฐํ.}}}

\csromanmatikaanchor{mtk.60}
\csromantocmarkref{subsection}{7. โสภนสุตฺต}{mtk.60}
\bookmark[dest={mtk.60},level=2]{7. โสภนสุตฺต}
\csromanheader{2}{7. \textbf{โสภนสุตฺต}}
\proseitem{7}{จตฺตาโรเม ภิกฺขเว วิยตฺตา วินีตา วิสารทา พหุสฺสุตา ธมฺมธรา ธมฺมา\-นุธมฺม\-ปฺปฏิปนฺนา สํฆํ โสเภนฺติ. กตเม จตฺตาโร? ภิกฺขุ ภิกฺขเว วิยตฺโต วินีโต วิสารโท พหุสฺสุโต ธมฺมธโร ธมฺมา\-นุธมฺม\-ปฺปฏิปนฺโน สํฆํ โสเภติ, ภิกฺขุนี ภิกฺขเว วิยตฺตา วินีตา วิสารทา พหุสฺสุตา ธมฺมธรา ธมฺมา\-นุธมฺม\-ปฺปฏิปนฺนา สํฆํ โสเภติ, อุปาสโก ภิกฺขเว วิยตฺโต วินีโต วิสารโท พหุสฺสุโต ธมฺมธโร ธมฺมา\-นุธมฺม\-ปฺปฏิปนฺโน สํฆํ โสเภติ, อุปาสิกา ภิกฺขเว วิยตฺตา วินีตา วิสารทา พหุสฺสุตา ธมฺมธรา ธมฺมา\-นุธมฺม\-ปฺปฏิปนฺนา สํฆํ โสเภติ. อิเม โข ภิกฺขเว จตฺตาโร วิยตฺตา วินีตา วิสารทา พหุสฺสุตา ธมฺมธรา ธมฺมา\-นุธมฺม\-ปฺปฏิปนฺนา สํฆํ โสเภนฺตีติ.}

\setlength{\csromangathapagewidth}{0pt}
\setlength{\csromangathawakwidth}{0pt}
\csromangathameasuregroup{ภิกฺขุ จ สีลสมฺปนฺโน, \\ อุปาสโก จ โย สทฺโธ, \\ เอเต โข สํฆํ โสเภนฺติ,}{โย โหติ วิยตฺโต\textsuperscript{1} จ วิสารโท จ, \\ พหุสฺสุโต ธมฺมธโร จ โหติ. \\ ธมฺมสฺส โหติ อนุธมฺมจารี, \\ ส ตาทิโส วุจฺจติ สํฆโสภโน\textsuperscript{1}. \\ \csromangathabat{ภิกฺขุ จ สีลสมฺปนฺโน,}{ภิกฺขุนี จ พหุสฺสุตา.} \\ \csromangathabat{อุปาสโก จ โย สทฺโธ,}{ยา จ สทฺธา อุปาสิกา.} \\ \csromangathabat{เอเต โข สํฆํ โสเภนฺติ,}{เอเต หิ สํฆโสภนาติ.}}
\csromangathameasurewak{โย โหติ วิยตฺโต\textsuperscript{1} จ วิสารโท จ, \\ พหุสฺสุโต ธมฺมธโร จ โหติ. \\ ธมฺมสฺส โหติ อนุธมฺมจารี, \\ ส ตาทิโส วุจฺจติ สํฆโสภโน\textsuperscript{1}.}
\csromangathasetleft{ภิกฺขุ จ สีลสมฺปนฺโน, \\ อุปาสโก จ โย สทฺโธ, \\ เอเต โข สํฆํ โสเภนฺติ,}
\csromangathagroup{\csromangathastanza{\csromangathawak{โย โหติ วิยตฺโต\csromansharedfootnote{9c29039aa05ec2d1}{วฺยตฺโต (สี, อิ), พฺยตฺโต (สฺยา, กํ)} จ วิสารโท จ, \\ พหุสฺสุโต ธมฺมธโร จ โหติ. \\ ธมฺมสฺส โหติ อนุธมฺมจารี, \\ ส ตาทิโส วุจฺจติ สํฆโสภโน\csromansharedfootnote{d2d552410e0df710}{สํฆโสภโณ (ก)}.}}\csromangathastanza{\csromangathabat{ภิกฺขุ จ สีลสมฺปนฺโน,}{ภิกฺขุนี จ พหุสฺสุตา.} \\ \csromangathabat{อุปาสโก จ โย สทฺโธ,}{ยา จ สทฺธา อุปาสิกา.} \\ \csromangathabat{เอเต โข สํฆํ โสเภนฺติ,}{เอเต หิ สํฆโสภนาติ.}}}

\csromanmatikaanchor{mtk.61}
\csromantocmarkref{subsection}{8. เวสารชฺชสุตฺต}{mtk.61}
\bookmark[dest={mtk.61},level=2]{8. เวสารชฺชสุตฺต}
\csromanmarkh{1. ภณฺฑคาม\-วคฺค}\csromanmarkhii{8. เวสารชฺชสุตฺต}\csromanheadernomark{2}{1. ภณฺฑคาม\-วคฺค 8. เวสารชฺชสุตฺต}
\proseitem{8}{\csromanfolio{315}จตฺตาริมานิ ภิกฺขเว ตถาคตสฺส เวสารชฺชานิ, เยหิ เวสารชฺเชหิ สมนฺนาคโต ตถาคโต อาสภํ ฐานํ ปฏิชานาติ, ปริสาสุ สีหนาทํ นทติ, พฺรหฺมจกฺกํ ปวตฺเตติ. กตมานิ จตฺตาริ? สมฺมา\-สมฺพุทฺธสฺส เต ปฏิชานโต “อิเม ธมฺมา อนภิสมฺพุทฺธา”ติ, ตตฺร วต มํ สมโณ วา พฺราหฺมโณ วา เทโว วา มาโร วา พฺรหฺมา วา โกจิ วา โลกสฺมิํ สหธมฺเมน ปฏิโจเทสฺสตีติ นิมิตฺตเมตํ ภิกฺขเว น สมนุปสฺสามิ, เอตมหํ\csromansharedfootnote{6897c9db1bd7c0de}{เอตมฺปหํ (สี, สฺยา, กํ, อิ)} ภิกฺขเว นิมิตฺตํ อสมนุปสฺสนฺโต เขมปฺปตฺโต อภยปฺปตฺโต เวสารชฺช\-ปฺปตฺโต วิหรามิ.}

\prose{ขีณาสวสฺส เต ปฏิชานโต “อิเม อาสวา อปริกฺขีณา”ติ, ตตฺร วต มํ สมโณ วา พฺราหฺมโณ วา เทโว วา มาโร วา พฺรหฺมา วา โกจิ วา โลกสฺมิํ สหธมฺเมน ปฏิโจเทสฺสตีติ นิมิตฺตเมตํ ภิกฺขเว น สมนุปสฺสามิ, เอตมหํ ภิกฺขเว นิมิตฺตํ อสมนุปสฺสนฺโต เขมปฺปตฺโต อภยปฺปตฺโต เวสารชฺช\-ปฺปตฺโต วิหรามิ.}

\prose{“เย โข ปน เต อนฺตรายิกา ธมฺมา วุตฺตา, เต ปฏิเสวโต นาลํ อนฺตรายายา”ติ, ตตฺร วต มํ สมโณ วา พฺราหฺมโณ วา เทโว วา มาโร วา พฺรหฺมา วา โกจิ วา โลกสฺมิํ สหธมฺเมน ปฏิโจเทสฺสตีติ นิมิตฺตเมตํ ภิกฺขเว น สมนุปสฺสามิ, เอตมหํ ภิกฺขเว นิมิตฺตํ อสมนุปสฺสนฺโต เขมปฺปตฺโต อภยปฺปตฺโต เวสารชฺช\-ปฺปตฺโต วิหรามิ.}

\prose{“ยสฺส โข ปน เต อตฺถาย ธมฺโม เทสิโต, โส น นิยฺยาติ ตกฺกรสฺส สมฺมา ทุกฺขกฺขยายา”ติ, ตตฺร วต มํ สมโณ วา พฺราหฺมโณ วา เทโว วา มาโร วา พฺรหฺมา วา โกจิ วา โลกสฺมิํ สหธมฺเมน ปฏิโจเทสฺสตีติ นิมิตฺตเมตํ ภิกฺขเว น สมนุปสฺสามิ, เอตมหํ ภิกฺขเว นิมิตฺตํ อสมนุปสฺสนฺโต เขมปฺปตฺโต อภยปฺปตฺโต เวสารชฺช\-ปฺปตฺโต วิหรามิ. อิมานิ โข ภิกฺขเว จตฺตาริ ตถาคตสฺส เวสารชฺชานิ, เยหิ เวสารชฺเชหิ สมนฺนาคโต ตถาคโต อาสภํ ฐานํ ปฏิชานาติ, ปริสาสุ สีหนาทํ นทติ, พฺรหฺมจกฺกํ ปวตฺเตตีติ.}

\setlength{\csromangathapagewidth}{0pt}
\setlength{\csromangathawakwidth}{0pt}
\csromangathameasuregroup{}{เย เกจิเม วาทปถา ปุถุสฺสิตา, \\ ยํนิสฺสิตา สมณพฺราหฺมณา จ. \\ ตถาคตํ ปตฺวา น เต ภวนฺติ, \\ วิสารทํ วาทปถาติวตฺตํ\textsuperscript{1}. \\ โย ธมฺมจกฺกํ อภิภุยฺย เกวลี\textsuperscript{1}, \\ ปวตฺตยี สพฺพภูตานุกมฺปี. \\ ตํ ตาทิสํ เทวมนุสฺสเสฏฺฐํ, \\ สตฺตา นมสฺสนฺติ ภวสฺส ปารคุนฺติ.อฏฺฐมํ.}
\csromangathameasurewak{เย เกจิเม วาทปถา ปุถุสฺสิตา, \\ ยํนิสฺสิตา สมณพฺราหฺมณา จ. \\ ตถาคตํ ปตฺวา น เต ภวนฺติ, \\ วิสารทํ วาทปถาติวตฺตํ\textsuperscript{1}. \\ โย ธมฺมจกฺกํ อภิภุยฺย เกวลี\textsuperscript{1}, \\ ปวตฺตยี สพฺพภูตานุกมฺปี. \\ ตํ ตาทิสํ เทวมนุสฺสเสฏฺฐํ, \\ สตฺตา นมสฺสนฺติ ภวสฺส ปารคุนฺติ.อฏฺฐมํ.}
\setlength{\csromangathaleftcol}{0pt}
\csromangathagroup{\csromangathastanza{\csromangathawak{\csromanfolio{316}เย เกจิเม วาทปถา ปุถุสฺสิตา, \\ ยํนิสฺสิตา สมณพฺราหฺมณา จ. \\ ตถาคตํ ปตฺวา น เต ภวนฺติ, \\ วิสารทํ วาทปถาติวตฺตํ\csromansharedfootnote{f7e18d395e031ffa}{วาทปถาภิวตฺตินํ (สี), วาทปถาติ วุตฺตํ (อิ, ก)}.}}\csromangathastanza{\csromangathawak{โย ธมฺมจกฺกํ อภิภุยฺย เกวลี\csromansharedfootnote{2f907f7e7c613b3f}{เกวลํ (สฺยา), เกวโล (ก)}, \\ ปวตฺตยี สพฺพภูตานุกมฺปี. \\ ตํ ตาทิสํ เทวมนุสฺสเสฏฺฐํ, \\ สตฺตา นมสฺสนฺติ ภวสฺส ปารคุนฺติ.  อฏฺฐมํ.}}}

\csromanmatikaanchor{mtk.62}
\csromantocmarkref{subsection}{9. ตณฺหุปฺปาทสุตฺต}{mtk.62}
\bookmark[dest={mtk.62},level=2]{9. ตณฺหุปฺปาทสุตฺต}
\csromanheader{2}{9. ตณฺหุปฺปาท\-สุตฺต}
\proseitem{9}{จตฺตาโรเม ภิกฺขเว ตณฺหุปฺปาทา, ยตฺถ ภิกฺขุโน ตณฺหา อุปฺปชฺชมานา อุปฺปชฺชติ. กตเม จตฺตาโร? จีวรเหตุ วา ภิกฺขเว ภิกฺขุโน ตณฺหา อุปฺปชฺชมานา อุปฺปชฺชติ, ปิณฺฑปาตเหตุ วา ภิกฺขเว ภิกฺขุโน ตณฺหา อุปฺปชฺชมานา อุปฺปชฺชติ, เสนาสนเหตุ วา ภิกฺขเว ภิกฺขุโน ตณฺหา อุปฺปชฺชมานา อุปฺปชฺชติ, อิติภวา\-ภ\-ว\-เหตุ วา ภิกฺขเว ภิกฺขุโน ตณฺหา อุปฺปชฺชมานา อุปฺปชฺชติ. อิเม โข ภิกฺขเว จตฺตาโร ตณฺหุปฺปาทา, ยตฺถ ภิกฺขุโน ตณฺหา อุปฺปชฺชมานา อุปฺปชฺชตีติ.}

\setlength{\csromangathapagewidth}{0pt}
\setlength{\csromangathawakwidth}{0pt}
\csromangathameasuregroup{ตณฺหาทุติโย ปุริโส, \\ อิตฺถภาวญฺญถาภาวํ, \\ เอวมาทีนวํ ญตฺวา, \\ วีตตโณฺห อนาทาโน,}{\csromangathabat{ตณฺหาทุติโย ปุริโส,}{ทีฆมทฺธาน สํสรํ.} \\ \csromangathabat{อิตฺถภาวญฺญถาภาวํ,}{สํสารํ นาติวตฺตติ.} \\ \csromangathabat{เอวมาทีนวํ ญตฺวา,}{ตณฺหํ ทุกฺขสฺส สมฺภวํ.} \\ \csromangathabat{วีตตโณฺห อนาทาโน,}{สโต ภิกฺขุ ปริพฺพเชติ\textsuperscript{1}.} \\ นวมํ.}
\csromangathameasurewak{นวมํ.}
\csromangathasetleft{ตณฺหาทุติโย ปุริโส, \\ อิตฺถภาวญฺญถาภาวํ, \\ เอวมาทีนวํ ญตฺวา, \\ วีตตโณฺห อนาทาโน,}
\csromangathagroup{\csromangathastanza{\csromangathabat{ตณฺหาทุติโย ปุริโส,}{ทีฆมทฺธาน สํสรํ.} \\ \csromangathabat{อิตฺถ\-ภาว\-ญฺญถา\-ภาวํ,}{สํสารํ นาติวตฺตติ.}}\csromangathastanza{\csromangathabat{เอวมาทีนวํ ญตฺวา,}{ตณฺหํ ทุกฺขสฺส สมฺภวํ.} \\ \csromangathabat{วีตตโณฺห อนาทาโน,}{สโต ภิกฺขุ ปริพฺพเชติ\csromansharedfootnote{338dd9d64fd7ee1e}{[มิสฺสินฺคฺ โนเต 0]}.}}\csromangathastanza{\csromangathawak{นวมํ.}}}

\csromanmatikaanchor{mtk.63}
\csromantocmarkref{subsection}{10. โยคสุตฺต}{mtk.63}
\bookmark[dest={mtk.63},level=2]{10. โยคสุตฺต}
\csromanheader{2}{10. \textbf{โยคสุตฺต}}
\proseitem{10}{จตฺตาโรเม ภิกฺขเว โยคา. กตเม จตฺตาโร? กามโยโค ภวโยโค ทิฏฺฐิโยโค อวิชฺชาโยโค. กตโม จ ภิกฺขเว กามโยโค, อิธ ภิกฺขเว เอกจฺโจ กามานํ สมุทยญฺจ อตฺถงฺคมญฺจ\csromansharedfootnote{0f0a54b607be70ec}{อตฺถคมญฺจ (สี, อิ)} อสฺสาทญฺจ อาทีนวญฺจ นิสฺสรณญฺจ ยถาภูตํ นปฺปชานาติ, \csromanfolio{317}ตสฺส กามานํ สมุทยญฺจ อตฺถงฺคมญฺจ อสฺสาทญฺจ อาทีนวญฺจ นิสฺสรณญฺจ ยถาภูตํ อปฺปชานโต\csromansharedfootnote{eb7b61b1122c3471}{นปฺปชานโต (สฺยา, กํ, ก)} โย กาเมสุ กามราโค กามนนฺที\csromansharedfootnote{1c570eb7101d362c}{กามนนฺทิ (สี, สฺยา, กํ)} กามสฺเนโห กามมุจฺฉา กามปิปาสา กามปริฬาโห กามชฺโฌสานํ กามตณฺหา, สานุเสติ. อยํ วุจฺจติ ภิกฺขเว กามโยโค. อิติ กามโยโค.}

\prose{ภวโยโค จ กถํ โหติ? อิธ ภิกฺขเว เอกจฺโจ ภวานํ สมุทยญฺจ อตฺถงฺคมญฺจ อสฺสาทญฺจ อาทีนวญฺจ นิสฺสรณญฺจ ยถาภูตํ นปฺปชานาติ, ตสฺส ภวานํ สมุทยญฺจ อตฺถงฺคมญฺจ อสฺสาทญฺจ อาทีนวญฺจ นิสฺสรณญฺจ ยถาภูตํ อปฺปชานโต โย ภเวสุ ภวราโค ภวนนฺที ภวสฺเนโห ภวมุจฺฉา ภวปิปาสา ภวปริฬาโห ภวชฺโฌสานํ ภวตณฺหา, สานุเสติ. อยํ วุจฺจติ ภิกฺขเว ภวโยโค. อิติ กามโยโค ภวโยโค.}

\prose{ทิฏฺฐิโยโค จ กถํ โหติ? อิธ ภิกฺขเว เอกจฺโจ ทิฏฺฐีนํ สมุทยญฺจ อตฺถงฺคมญฺจ อสฺสาทญฺจ อาทีนวญฺจ นิสฺสรณญฺจ ยถาภูตํ นปฺปชานาติ, ตสฺส ทิฏฺฐีนํ สมุทยญฺจ อตฺถงฺคมญฺจ อสฺสาทญฺจ อาทีนวญฺจ นิสฺสรณญฺจ ยถาภูตํ อปฺปชานโต โย ทิฏฺฐีสุ ทิฏฺฐิราโค ทิฏฺฐินนฺที ทิฏฺฐิสฺเนโห ทิฏฺฐิมุจฺฉา ทิฏฺฐิปิปาสา ทิฏฺฐิปริฬาโห ทิฏฺฐิชฺโฌสานํ\csromansharedfootnote{7dd5bad0f2e1e2b2}{ทิฏฺฐิอชฺโฌสานํ (สี, อิ)} ทิฏฺฐิตณฺหา, สานุเสติ. อยํ วุจฺจติ ภิกฺขเว ทิฏฺฐิโยโค. อิติ กามโยโค ภวโยโค ทิฏฺฐิโยโค.}

\prose{อวิชฺชาโยโค จ กถํ โหติ? อิธ ภิกฺขเว เอกจฺโจ ฉนฺนํ ผสฺสา\-ยตนานํ สมุทยญฺจ อตฺถงฺคมญฺจ อสฺสาทญฺจ อาทีนวญฺจ นิสฺสรณญฺจ ยถาภูตํ นปฺปชานาติ, ตสฺส ฉนฺนํ ผสฺสา\-ยตนานํ สมุทยญฺจ อตฺถงฺคมญฺจ อสฺสาทญฺจ อาทีนวญฺจ นิสฺสรณญฺจ ยถาภูตํ อปฺปชานโต ยา ฉสุ ผสฺสายตเนสุ อวิชฺชา อญฺญาณํ, สานุเสติ. อยํ วุจฺจติ ภิกฺขเว อวิชฺชาโยโค. อิติ กามโยโค ภวโยโค ทิฏฺฐิโยโค อวิชฺชาโยโค. สํยุตฺโต ปาปเกหิ อกุสเลหิ ธมฺเมหิ สํกิเลสิเกหิ โปโนภวิเกหิ\csromansharedfootnote{0ac664b9cbe50e78}{โปโนพฺภวิเกหิ (สี, อิ)} สทเรหิ ทุกฺขวิปาเกหิ อายติํ ชาติ\-ชรา\-มรณิ\-เกหิ, ตสฺมา “อโยคกฺเขมี”ติ วุจฺจติ. อิเม โข ภิกฺขเว จตฺตาโร โยคา.}

\prose{\csromanfolio{318}จตฺตาโรเม ภิกฺขเว วิสํโยคา. กตเม จตฺตาโร? กามโยค\-วิสํโยโค ภวโยค\-วิสํโยโค ทิฏฺฐิโยค\-วิสํโยโค อวิชฺชาโยค\-วิสํโยโค. กตโม จ ภิกฺขเว กามโยค\-วิสํโยโค? อิธ ภิกฺขเว เอกจฺโจ กามานํ สมุทยญฺจ อตฺถงฺคมญฺจ อสฺสาทญฺจ อาทีนวญฺจ นิสฺสรณญฺจ ยถาภูตํ ปชานาติ, ตสฺส กามานํ สมุทยญฺจ อตฺถงฺคมญฺจ อสฺสาทญฺจ อาทีนวญฺจ นิสฺสรณญฺจ ยถาภูตํ ปชานโต โย กาเมสุ กามราโค กามนนฺที กามสฺเนโห กามมุจฺฉา กามปิปาสา กามปริฬาโห กามชฺโฌสานํ กามตณฺหา, สา นานุเสติ. อยํ วุจฺจติ ภิกฺขเว กามโยค\-วิสํโยโค. อิติ กามโยค\-วิสํโยโค.}

\prose{ภวโยค\-วิสํโยโค จ กถํ โหติ? อิธ ภิกฺขเว เอกจฺโจ ภวานํ สมุทยญฺจ อตฺถงฺคมญฺจ อสฺสาทญฺจ อาทีนวญฺจ นิสฺสรณญฺจ ยถาภูตํ ปชานาติ, ตสฺส ภวานํ สมุทยญฺจ อตฺถงฺคมญฺจ อสฺสาทญฺจ อาทีนวญฺจ นิสฺสรณญฺจ ยถาภูตํ ปชานโต โย ภเวสุ ภวราโค ภวนนฺที ภวสฺเนโห ภวมุจฺฉา ภวปิปาสา ภวปริฬาโห ภวชฺโฌสานํ ภวตณฺหา, สา นานุเสติ. อยํ วุจฺจติ ภิกฺขเว ภวโยค\-วิสํโยโค. อิติ กามโยค\-วิสํโยโค ภวโยค\-วิสํโยโค.}

\prose{ทิฏฺฐิโยค\-วิสํโยโค จ กถํ โหติ? อิธ ภิกฺขเว เอกจฺโจ ทิฏฺฐีนํ สมุทยญฺจ อตฺถงฺคมญฺจ อสฺสาทญฺจ อาทีนวญฺจ นิสฺสรณญฺจ ยถาภูตํ ปชานาติ, ตสฺส ทิฏฺฐีนํ สมุทยญฺจ อตฺถงฺคมญฺจ อสฺสาทญฺจ อาทีนวญฺจ นิสฺสรณญฺจ ยถาภูตํ ปชานโต โย ทิฏฺฐีสุ ทิฏฺฐิราโค ทิฏฺฐินนฺที ทิฏฺฐิสฺเนโห ทิฏฺฐิมุจฺฉา ทิฏฺฐิปิปาสา ทิฏฺฐิปริฬาโห ทิฏฺฐิชฺโฌสานํ ทิฏฺฐิตณฺหา, สา นานุเสติ. อยํ วุจฺจติ ภิกฺขเว ทิฏฺฐิโยค\-วิสํโยโค. อิติ กามโยค\-วิสํโยโค ภวโยค\-วิสํโยโค ทิฏฺฐิโยค\-วิสํโยโค.}

\prose{อวิชฺชาโยค\-วิสํโยโค จ กถํ โหติ? อิธ ภิกฺขเว เอกจฺโจ ฉนฺนํ ผสฺสา\-ยตนานํ สมุทยญฺจ อตฺถงฺคมญฺจ อสฺสาทญฺจ อาทีนวญฺจ นิสฺสรณญฺจ ยถาภูตํ ปชานาติ, ตสฺส ฉนฺนํ ผสฺสา\-ยตนานํ สมุทยญฺจ อตฺถงฺคมญฺจ อสฺสาทญฺจ อาทีนวญฺจ นิสฺสรณญฺจ ยถาภูตํ ปชานโต ยา ฉสุ ผสฺสายตเนสุ อวิชฺชา อญฺญาณํ, สา นานุเสติ. อยํ วุจฺจติ ภิกฺขเว \csromanfolio{319}อวิชฺชาโยค\-วิสํโยโค. อิติ กามโยค\-วิสํโยโค ภวโยค\-วิสํโยโค ทิฏฺฐิโยค\-วิสํโยโค อวิชฺชาโยค\-วิสํโยโค. วิสํยุตฺโต ปาปเกหิ อกุสเลหิ ธมฺเมหิ สํกิเลสิเกหิ โปโนภวิเกหิ สทเรหิ ทุกฺขวิปาเกหิ อายติํ ชาติ\-ชรา\-มรณิ\-เกหิ. ตสฺมา “โยคกฺเขมี”ติ วุจฺจติ. อิเม โข ภิกฺขเว จตฺตาโร วิสํโยคาติ.}

\setlength{\csromangathapagewidth}{0pt}
\setlength{\csromangathawakwidth}{0pt}
\csromangathameasuregroup{กามโยเคน สํยุตฺตา, \\ ทิฏฺฐิโยเคน สํยุตฺตา, \\ สตฺตา คจฺฉนฺติ สํสารํ, \\ เย จ กาเม ปริญฺญาย, \\ ทิฏฺฐิโยคํ สมูหจฺจ, \\ สพฺพโยควิสํยุตฺตา,}{\csromangathabat{กามโยเคน สํยุตฺตา,}{ภวโยเคน จูภยํ.} \\ \csromangathabat{ทิฏฺฐิโยเคน สํยุตฺตา,}{อวิชฺชาย ปุรกฺขตา.} \\ \csromangathabat{สตฺตา คจฺฉนฺติ สํสารํ,}{ชาติมรณคามิโน.} \\ \csromangathabat{เย จ กาเม ปริญฺญาย,}{ภวโยคญฺจ สพฺพโส.} \\ \csromangathabat{ทิฏฺฐิโยคํ สมูหจฺจ,}{อวิชฺชญฺจ วิราชยํ.} \\ \csromangathabat{สพฺพโยควิสํยุตฺตา,}{เต เว โยคาติคา มุนีติ.} \\ ทสมํ.}
\csromangathameasurewak{ทสมํ.}
\csromangathasetleft{กามโยเคน สํยุตฺตา, \\ ทิฏฺฐิโยเคน สํยุตฺตา, \\ สตฺตา คจฺฉนฺติ สํสารํ, \\ เย จ กาเม ปริญฺญาย, \\ ทิฏฺฐิโยคํ สมูหจฺจ, \\ สพฺพโยควิสํยุตฺตา,}
\csromangathagroupwithclosermidruled{\csromangathastanza{\csromangathabat{กามโยเคน สํยุตฺตา,}{ภวโยเคน จูภยํ.} \\ \csromangathabat{ทิฏฺฐิโยเคน สํยุตฺตา,}{อวิชฺชาย ปุรกฺขตา.}}\csromangathastanza{\csromangathabat{สตฺตา คจฺฉนฺติ สํสารํ,}{ชาติมรณคามิโน.} \\ \csromangathabat{เย จ กาเม ปริญฺญาย,}{ภวโยคญฺจ สพฺพโส.}}\csromangathastanza{\csromangathabat{ทิฏฺฐิโยคํ สมูหจฺจ,}{อวิชฺชญฺจ วิราชยํ.} \\ \csromangathabat{สพฺพ\-โยค\-วิสํยุตฺตา,}{เต เว โยคาติคา มุนีติ.}}\csromangathastanza{\csromangathawak{ทสมํ.}}}{ภณฺฑคาม\-วคฺโค ปฐโม.}
\csromancenter{\textbf{ตสฺสุทฺทานํ}}
\setlength{\csromangathapagewidth}{0pt}
\setlength{\csromangathawakwidth}{0pt}
\csromangathameasuregroup{อนุพุทฺธํ ปปติตํ เทฺว, \\ อปฺปสฺสุโต จ โสภนํ,}{\csromangathabat{อนุพุทฺธํ ปปติตํ เทฺว,}{ขตา อนุโสตปญฺจมํ.} \\ \csromangathabat{อปฺปสฺสุโต จ โสภนํ,}{เวสารชฺชํ ตณฺหาโยเคน เต ทสาติ.}}
\csromangathasetleft{อนุพุทฺธํ ปปติตํ เทฺว, \\ อปฺปสฺสุโต จ โสภนํ,}
\csromangathagroup{\csromangathastanza{\csromangathabat{อนุพุทฺธํ ปปติตํ เทฺว,}{ขตา อนุโสตปญฺจมํ.} \\ \csromangathabat{อปฺปสฺสุโต จ โสภนํ,}{เวสารชฺชํ ตณฺหาโยเคน เต ทสาติ.}}}

\csromanmatikaanchor{mtk.64}
\csromantocmarknopage{section}{2. จรวคฺค}{mtk.64}
\bookmark[dest={mtk.64},level=1]{2. จรวคฺค}
\csromanheader{1}{2. จรวคฺค}
\csromanmatikaanchor{mtk.65}
\csromantocmarkref{subsection}{1. จรสุตฺต}{mtk.65}
\bookmark[dest={mtk.65},level=2]{1. จรสุตฺต}
\csromanheader{2}{1. \textbf{จรสุตฺต}}
\proseitem{11}{\csromansymbolfootnote{*}{ขุ 1. 272 ปิฏฺเฐปิ.}จรโต เจปิ ภิกฺขเว ภิกฺขุโน อุปฺปชฺชติ กามวิตกฺโก วา พฺยาปาทวิตกฺโก วา วิหิํสาวิตกฺโก วา, ตํ เจ ภิกฺขุ อธิวาเสติ นปฺปชหติ น วิโนเทติ น พฺยนฺตีกโรติ\csromansharedfootnote{6f6613369a05fceb}{พฺยนฺติกโรติ (อิ), พฺยนฺติํ กโรติ (ก)} น อนภาวํ คเมติ. จรมฺปิ ภิกฺขเว ภิกฺขุ เอวํภูโต อนาตาปี อโนตฺตาปี “สตตํ สมิตํ กุสีโต หีนวีริโย”ติ วุจฺจติ.}

\prose{\csromanfolio{320}ฐิตสฺส เจปิ ภิกฺขเว ภิกฺขุโน อุปฺปชฺชติ กามวิตกฺโก วา พฺยาปาทวิตกฺโก วา วิหิํสาวิตกฺโก วา, ตํ เจ ภิกฺขุ อธิวาเสติ นปฺปชหติ น วิโนเทติ น พฺยนฺตีกโรติ น อนภาวํ คเมติ. ฐิโตปิ ภิกฺขเว ภิกฺขุ เอวํภูโต อนาตาปี อโนตฺตาปี “สตตํ สมิตํ กุสีโต หีนวีริโย”ติ วุจฺจติ.}

\prose{นิสินฺนสฺส เจปิ ภิกฺขเว ภิกฺขุโน อุปฺปชฺชติ กามวิตกฺโก วา พฺยาปาทวิตกฺโก วา วิหิํสาวิตกฺโก วา, ตํ เจ ภิกฺขุ อธิวาเสติ นปฺปชหติ น วิโนเทติ น พฺยนฺตีกโรติ น อนภาวํ คเมติ. นิสินฺโนปิ ภิกฺขเว ภิกฺขุ เอวํภูโต อนาตาปี อโนตฺตาปี “สตตํ สมิตํ กุสีโต หีนวีริโย”ติ วุจฺจติ.}

\prose{สยานสฺส เจปิ ภิกฺขเว ภิกฺขุโน ชาครสฺส อุปฺปชฺชติ กามวิตกฺโก วา พฺยาปาทวิตกฺโก วา วิหิํสาวิตกฺโก วา, ตํ เจ ภิกฺขุ อธิวาเสติ นปฺปชหติ น วิโนเทติ น พฺยนฺตีกโรติ น อนภาวํ คเมติ. สยาโนปิ ภิกฺขเว ภิกฺขุ ชาคโร เอวํภูโต อนาตาปี อโนตฺตาปี “สตตํ สมิตํ กุสีโต หีนวีริโย”ติ วุจฺจติ.}

\prose{จรโต เจปิ ภิกฺขเว ภิกฺขุโน อุปฺปชฺชติ กามวิตกฺโก วา พฺยาปาทวิตกฺโก วา วิหิํสาวิตกฺโก วา, ตํ เจ ภิกฺขุ นาธิวาเสติ ปชหติ วิโนเทติ พฺยนฺตีกโรติ อนภาวํ คเมติ. จรมฺปิ ภิกฺขเว ภิกฺขุ เอวํภูโต อาตาปี โอตฺตาปี “สตตํ สมิตํ อารทฺธวีริโย ปหิตตฺโต”ติ วุจฺจติ.}

\prose{ฐิตสฺส เจปิ ภิกฺขเว ภิกฺขุโน อุปฺปชฺชติ กามวิตกฺโก วา พฺยาปาทวิตกฺโก วา วิหิํสาวิตกฺโก วา, ตํ เจ ภิกฺขุ นาธิวาเสติ ปชหติ วิโนเทติ พฺยนฺตีกโรติ อนภาวํ คเมติ. ฐิโตปิ ภิกฺขเว ภิกฺขุ เอวํภูโต อาตาปี โอตฺตาปี “สตตํ สมิตํ อารทฺธวีริโย ปหิตตฺโต”ติ วุจฺจติ.}

\prose{นิสินฺนสฺส เจปิ ภิกฺขเว ภิกฺขุโน อุปฺปชฺชติ กามวิตกฺโก วา พฺยาปาทวิตกฺโก วา วิหิํสาวิตกฺโก วา, ตํ เจ ภิกฺขุ นาธิวาเสติ ปชหติ วิโนเทติ พฺยนฺตีกโรติ อนภาวํ คเมติ. นิสินฺโนปิ ภิกฺขเว ภิกฺขุ เอวํภูโต อาตาปี โอตฺตาปี “สตตํ สมิตํ อารทฺธวีริโย ปหิตตฺโต”ติ วุจฺจติ.}

\chapterheadpageread{2. จรวคฺค}
\prose{\csromanfolio{321}สยานสฺส เจปิ ภิกฺขเว ภิกฺขุโน ชาครสฺส อุปฺปชฺชติ กามวิตกฺโก วา พฺยาปาทวิตกฺโก วา วิหิํสาวิตกฺโก วา, ตํ เจ ภิกฺขุ นาธิวาเสติ ปชหติ วิโนเทติ พฺยนฺตีกโรติ อนภาวํ คเมติ. สยาโนปิ ภิกฺขเว ภิกฺขุ ชาคโร เอวํภูโต อาตาปี โอตฺตาปี “สตตํ สมิตํ อารทฺธวีริโย ปหิตตฺโต”ติ วุจฺจตีติ.}

\setlength{\csromangathapagewidth}{0pt}
\setlength{\csromangathawakwidth}{0pt}
\csromangathameasuregroup{จรํ วา ยทิ วา ติฏฺฐํ, \\ โย วิตกฺกํ วิตกฺเกติ, \\ กุมฺมคฺคปฺปฏิปนฺโน โส, \\ อภพฺโพ ตาทิโส ภิกฺขุ, \\ โย จ จรํ วา ติฏฺฐํ วา, \\ วิตกฺกํ สมยิตฺวาน, \\ ภพฺโพ โส ตาทิโส ภิกฺขุ,}{\csromangathabat{จรํ วา ยทิ วา ติฏฺฐํ,}{นิสินฺโน อุท วา สยํ.} \\ \csromangathabat{โย วิตกฺกํ วิตกฺเกติ,}{ปาปกํ เคหนิสฺสิตํ.} \\ \csromangathabat{กุมฺมคฺคปฺปฏิปนฺโน โส,}{โมหเนยฺเยสุ มุจฺฉิโต.} \\ \csromangathabat{อภพฺโพ ตาทิโส ภิกฺขุ,}{ผุฏฺฐุํ สมฺโพธิมุตฺตมํ.} \\ \csromangathabat{โย จ จรํ วา ติฏฺฐํ วา,}{นิสินฺโน อุท วา สยํ.} \\ \csromangathabat{วิตกฺกํ สมยิตฺวาน,}{วิตกฺกูปสเม รโต.} \\ \csromangathabat{ภพฺโพ โส ตาทิโส ภิกฺขุ,}{ผุฏฺฐุํ สมฺโพธิมุตฺตมนฺติ.} \\ ปฐมํ.}
\csromangathasetleft{จรํ วา ยทิ วา ติฏฺฐํ, \\ โย วิตกฺกํ วิตกฺเกติ, \\ กุมฺมคฺคปฺปฏิปนฺโน โส, \\ อภพฺโพ ตาทิโส ภิกฺขุ, \\ โย จ จรํ วา ติฏฺฐํ วา, \\ วิตกฺกํ สมยิตฺวาน, \\ ภพฺโพ โส ตาทิโส ภิกฺขุ,}
\csromangathagroup{\csromangathastanza{\csromangathabat{จรํ วา ยทิ วา ติฏฺฐํ,}{นิสินฺโน อุท วา สยํ.} \\ \csromangathabat{โย วิตกฺกํ วิตกฺเกติ,}{ปาปกํ เคหนิสฺสิตํ.}}\csromangathastanza{\csromangathabat{กุมฺมคฺค\-ปฺปฏิปนฺโน โส,}{โมหเนยฺเยสุ มุจฺฉิโต.} \\ \csromangathabat{อภพฺโพ ตาทิโส ภิกฺขุ,}{ผุฏฺฐุํ สมฺโพธิมุ\-ตฺตมํ.}}\csromangathastanza{\csromangathabat{โย จ จรํ วา ติฏฺฐํ วา,}{นิสินฺโน อุท วา สยํ.} \\ \csromangathabat{วิตกฺกํ สมยิตฺวาน,}{วิตกฺกูปสเม รโต.}}\csromangathastanza{\csromangathabat{ภพฺโพ โส ตาทิโส ภิกฺขุ,}{ผุฏฺฐุํ สมฺโพธิมุตฺตมนฺติ.} \\ ปฐมํ.}}

\csromanmatikaanchor{mtk.66}
\csromantocmarkref{subsection}{2. สีลสุตฺต}{mtk.66}
\bookmark[dest={mtk.66},level=2]{2. สีลสุตฺต}
\csromanheader{2}{2. \textbf{สีลสุตฺต}}
\proseitem{12}{สมฺปนฺนสีลา ภิกฺขเว วิหรถ สมฺปนฺน\-ปาติโมกฺขา, ปาติโมกฺข\-สํวร\-สํวุตา วิหรถ อาจารโคจร\-สมฺปนฺนา, อณุมตฺเตสุ วชฺเชสุ ภยทสฺสาวิโน สมาทาย สิกฺขถ สิกฺขาปเทสุ, สมฺปนฺน\-สีลานํ โว ภิกฺขเว วิหรตํ สมฺปนฺน\-ปาติโมกฺขานํ ปาติโมกฺข\-สํวร\-สํวุตานํ วิหรตํ อาจารโคจร\-สมฺปนฺนานํ อณุมตฺเตสุ วชฺเชสุ ภยทสฺสาวีนํ สมาทาย สิกฺขตํ สิกฺขาปเทสุ. กิมสฺส อุตฺตริ กรณียํ.}

\prose{จรโต เจปิ ภิกฺขเว ภิกฺขุโน อภิชฺฌา\-พฺยาปาโท วิคโต โหติ, ถินมิทฺธํ อุทฺธจฺจ\-กุกฺกุจฺจํ วิจิกิจฺฉา ปหีนา โหติ, อารทฺธํ โหติ วีริยํ อสลฺลีนํ, อุปฏฺฐิตา สติ อสมฺมุฏฺฐา, ปสฺสทฺโธ กาโย อสารทฺโธ, สมาหิตํ จิตฺตํ เอกคฺคํ. จรมฺปิ ภิกฺขเว ภิกฺขุ เอวํภูโต อาตาปี โอตฺตาปี “สตตํ สมิตํ อารทฺธวีริโย ปหิตตฺโต”ติ วุจฺจติ.}

\prose{ฐิตสฺส เจปิ ภิกฺขเว ภิกฺขุโน อภิชฺฌา\-พฺยาปาโท วิคโต โหติ, ถินมิทฺธํ อุทฺธจฺจ\-กุกฺกุจฺจํ วิจิกิจฺฉา ปหีนา โหติ, อารทฺธํ โหติ วีริยํ อสลฺลีนํ, อุปฏฺฐิตา สติ อสมฺมุฏฺฐา, ปสฺสทฺโธ กาโย อสารทฺโธ, \csromanfolio{322}สมาหิตํ จิตฺตํ เอกคฺคํ. ฐิโตปิ ภิกฺขเว ภิกฺขุ เอวํภูโต อาตาปี โอตฺตาปี “สตตํ สมิตํ อารทฺธวีริโย ปหิตตฺโต”ติ วุจฺจติ.}

\prose{นิสินฺนสฺส เจปิ ภิกฺขเว ภิกฺขุโน อภิชฺฌา พฺยาปาโท วิคโต โหติ, ถินมิทฺธํ อุทฺธจฺจ\-กุกฺกุจฺจํ วิจิกิจฺฉา ปหีนา โหติ, อารทฺธํ โหติ วีริยํ อสลฺลีนํ, อุปฏฺฐิตา สติ อสมฺมุฏฺฐา, ปสฺสทฺโธ กาโย อสารทฺโธ, สมาหิตํ จิตฺตํ เอกคฺคํ. นิสินฺโนปิ ภิกฺขเว ภิกฺขุ เอวํภูโต อาตาปี โอตฺตาปี “สตตํ สมิตํ อารทฺธวีริโย ปหิตตฺโต”ติ วุจฺจติ.}

\prose{สยานสฺส เจปิ ภิกฺขเว ภิกฺขุโน ชาครสฺส อภิชฺฌา พฺยาปาโท วิคโต โหติ, ถินมิทฺธํ อุทฺธจฺจ\-กุกฺกุจฺจํ วิจิกิจฺฉา ปหีนา โหติ, อารทฺธํ โหติ วีริยํ อสลฺลีนํ, อุปฏฺฐิตา สติ อสมฺมุฏฺฐา, ปสฺสทฺโธ กาโย อสารทฺโธ, สมาหิตํ จิตฺตํ เอกคฺคํ. สยาโนปิ ภิกฺขเว ภิกฺขุ ชาคโร เอวํภูโต อาตาปี โอตฺตาปี “สตตํ สมิตํ อารทฺธวีริโย ปหิตตฺโต”ติ วุจฺจตีติ.}

\setlength{\csromangathapagewidth}{0pt}
\setlength{\csromangathawakwidth}{0pt}
\csromangathameasuregroup{ยตํ\textsuperscript{1} จเร ยตํ\textsuperscript{1} ติฏฺเฐ, \\ ยตํ\textsuperscript{1} สมิญฺชเย\textsuperscript{1} ภิกฺขุ, \\ อุทฺธํ ติริยํ อปาจีนํ, \\ สมเวกฺขิตา จ ธมฺมานํ, \\ เจโตสมถสามีจิํ, \\ สตตํ ปหิตตฺโต”ติ,}{\csromangathabat{ยตํ\textsuperscript{1} จเร ยตํ\textsuperscript{1} ติฏฺเฐ,}{ยตํ\textsuperscript{1} อจฺเฉ ยตํ\textsuperscript{1} สเย.} \\ \csromangathabat{ยตํ\textsuperscript{1} สมิญฺชเย\textsuperscript{1} ภิกฺขุ,}{ยตเมนํ\textsuperscript{1} ปสารเย.} \\ \csromangathabat{อุทฺธํ ติริยํ อปาจีนํ,}{ยาวตา ชคโต คติ.} \\ \csromangathabat{สมเวกฺขิตา จ ธมฺมานํ,}{ขนฺธานํ อุทยพฺพยํ.} \\ \csromangathabat{เจโตสมถสามีจิํ,}{สิกฺขมานํ สทา “สตํ.} \\ \csromangathabat{สตตํ ปหิตตฺโต”ติ,}{อาหุ ภิกฺขุํ ตถาวิธนฺติ. ทุติยํ.}}
\csromangathasetleft{ยตํ\textsuperscript{1} จเร ยตํ\textsuperscript{1} ติฏฺเฐ, \\ ยตํ\textsuperscript{1} สมิญฺชเย\textsuperscript{1} ภิกฺขุ, \\ อุทฺธํ ติริยํ อปาจีนํ, \\ สมเวกฺขิตา จ ธมฺมานํ, \\ เจโตสมถสามีจิํ, \\ สตตํ ปหิตตฺโต”ติ,}
\csromangathagroup{\csromangathastanza{\csromangathabat{ยตํ\csromansharedfootnote{88cabebce49c6320}{ยถา (ก) ขุ 1. 275 ปิฏฺเฐปิ.} จเร ยตํ\csromansharedfootnote{88cabebce49c6320}{ยถา (ก) ขุ 1. 275 ปิฏฺเฐปิ.} ติฏฺเฐ,}{ยตํ\csromansharedfootnote{88cabebce49c6320}{ยถา (ก) ขุ 1. 275 ปิฏฺเฐปิ.} อจฺเฉ ยตํ\csromansharedfootnote{88cabebce49c6320}{ยถา (ก) ขุ 1. 275 ปิฏฺเฐปิ.} สเย.} \\ \csromangathabat{ยตํ\csromansharedfootnote{88cabebce49c6320}{ยถา (ก) ขุ 1. 275 ปิฏฺเฐปิ.} สมิญฺชเย\csromansharedfootnote{88cabebce49c6320}{ยถา (ก) ขุ 1. 275 ปิฏฺเฐปิ.} ภิกฺขุ,}{ยตเมนํ\csromansharedfootnote{88cabebce49c6320}{ยถา (ก) ขุ 1. 275 ปิฏฺเฐปิ.} ปสารเย.}}\csromangathastanza{\csromangathabat{อุทฺธํ ติริยํ อปาจีนํ,}{ยาวตา ชคโต คติ.} \\ \csromangathabat{สมเวกฺขิตา จ ธมฺมานํ,}{ขนฺธานํ อุทยพฺพยํ.}}\csromangathastanza{\csromangathabat{เจโตสมถสามีจิํ,}{สิกฺขมานํ สทา “สตํ.} \\ \csromangathabat{สตตํ ปหิตตฺโต”ติ,}{อาหุ ภิกฺขุํ ตถาวิธนฺติ. ทุติยํ.}}}

\csromanmatikaanchor{mtk.67}
\csromantocmarkref{subsection}{3. ปธานสุตฺต}{mtk.67}
\bookmark[dest={mtk.67},level=2]{3. ปธานสุตฺต}
\csromanheader{2}{3. \textbf{ปธานสุตฺต}}
\proseitem{13}{จตฺตาริมานิ ภิกฺขเว สมฺม\-ปฺปธานานิ. กตมานิ จตฺตาริ? อิธ ภิกฺขเว ภิกฺขุ อนุปฺปนฺนานํ ปาปกานํ อกุสลานํ ธมฺมานํ อนุปฺปาทาย ฉนฺทํ ชเนติ วายมติ, วีริยํ อารภติ, จิตฺตํ ปคฺคณฺหาติ ปทหติ. อุปฺปนฺนานํ ปาปกานํ อกุสลานํ ธมฺมานํ ปหานาย ฉนฺทํ ชเนติ วายมติ, วีริยํ อารภติ, จิตฺตํ ปคฺคณฺหาติ ปทหติ. อนุปฺปนฺนานํ กุสลานํ ธมฺมานํ อุปฺปาทาย \csromanfolio{323}ฉนฺทํ ชเนติ วายมติ, วีริยํ อารภติ, จิตฺตํ ปคฺคณฺหาติ ปทหติ. อุปฺปนฺนานํ กุสลานํ ธมฺมานํ ฐิติยา อสมฺโมสาย ภิยฺโยภาวาย เวปุลฺลาย ภาวนาย ปาริปูริยา ฉนฺทํ ชเนติ วายมติ, วีริยํ อารภติ, จิตฺตํ ปคฺคณฺหาติ ปทหติ. อิมานิ โข ภิกฺขเว จตฺตาริ สมฺมปฺปธานานีติ.}

\setlength{\csromangathapagewidth}{0pt}
\setlength{\csromangathawakwidth}{0pt}
\csromangathameasuregroup{}{สมฺมปฺปธานา มารเธยฺยาภิภูตา, \\ เต อสิตา ชาติมรณภยสฺส ปารคู. \\ เต ตุสิตา เชตฺวา มารํ สวาหินิํ\textsuperscript{1} เต อเนชา, \\ สพฺพํ นมุจิพลํ อุปาติวตฺตา เต สุขิตาติ.ตติยํ.}
\csromangathameasurewak{สมฺมปฺปธานา มารเธยฺยาภิภูตา, \\ เต อสิตา ชาติมรณภยสฺส ปารคู. \\ เต ตุสิตา เชตฺวา มารํ สวาหินิํ\textsuperscript{1} เต อเนชา, \\ สพฺพํ นมุจิพลํ อุปาติวตฺตา เต สุขิตาติ.ตติยํ.}
\setlength{\csromangathaleftcol}{0pt}
\csromangathagroup{\csromangathastanza{\csromangathawak{สมฺมปฺปธานา มารเธยฺยาภิภูตา, \\ เต อสิตา ชาติมรณภยสฺส ปารคู. \\ เต ตุสิตา เชตฺวา มารํ สวาหินิํ\csromansharedfootnote{0193d6d033032735}{สวาหนํ (สฺยา, กํ, อิ, ก)} เต อเนชา, \\ สพฺพํ นมุจิพลํ อุปาติวตฺตา เต สุขิตาติ.  ตติยํ.}}}

\csromanmatikaanchor{mtk.68}
\csromantocmarkref{subsection}{4. สํวรสุตฺต}{mtk.68}
\bookmark[dest={mtk.68},level=2]{4. สํวรสุตฺต}
\csromanheader{2}{4. \textbf{สํวรสุตฺต}}
\proseitem{13}{จตฺตาริมานิ ภิกฺขเว ปธานานิ. กตมานิ จตฺตาริ? สํวร\-ปฺปธานํ ปหาน\-ปฺปธานํ ภาวนา\-ปฺปธานํ อนุรกฺขณา\-ปฺปธานํ. กตมญฺจ ภิกฺขเว สํวร\-ปฺปธานํ? อิธ ภิกฺขเว ภิกฺขุ จกฺขุนา รูปํ ทิสฺวา น นิมิตฺตคฺคาหี โหติ นานุพฺยญฺชน\-คฺคาหี, ยตฺวาธิกรณเมนํ จกฺขุนฺทฺริยํ อสํวุตํ วิหรนฺตํ อภิชฺฌา\-โทมนสฺสา ปาปกา อกุสลา ธมฺมา อนฺวาสฺสเวยฺยุํ. ตสฺส สํวราย ปฏิปชฺชติ, รกฺขติ จกฺขุนฺทฺริยํ, จกฺขุนฺทฺริเย สํวรํ อาปชฺชติ. โสเตน สทฺทํ สุตฺวา. ฆาเนน คนฺธํ ฆายิตฺวา. ชิวฺหาย รสํ สายิตฺวา. กาเยน โผฏฺฐพฺพํ ผุสิตฺวา. มนสา ธมฺมํ วิญฺญาย น นิมิตฺตคฺคาหี โหติ นานุพฺยญฺชน\-คฺคาหี, ยตฺวาธิกรณเมนํ มนินฺทฺริยํ อสํวุตํ วิหรนฺตํ อภิชฺฌา\-โทมนสฺสา ปาปกา อกุสลา ธมฺมา อนฺวาสฺสเวยฺยุํ. ตสฺส สํวราย ปฏิปชฺชติ, รกฺขติ มนินฺทฺริยํ, มนินฺทฺริเย สํวรํ อาปชฺชติ. อิทํ วุจฺจติ ภิกฺขเว สํวร\-ปฺปธานํ.}

\prose{กตมญฺจ ภิกฺขเว ปหาน\-ปฺปธานํ? อิธ ภิกฺขเว ภิกฺขุ อุปฺปนฺนํ กามวิตกฺกํ นาธิวาเสติ ปชหติ วิโนเทติ พฺยนฺตีกโรติ อนภาวํ คเมติ, อุปฺปนฺนํ พฺยาปาท\-วิตกฺกํ ฯเปฯ อุปฺปนฺนํ วิหิํสา\-วิตกฺกํ ฯเปฯ อุปฺปนฺนุปฺปนฺเน ปาปเก อกุสเล ธมฺเม นาธิวาเสติ ปชหติ วิโนเทติ พฺยนฺตีกโรติ อนภาวํ คเมติ. อิทํ วุจฺจติ ภิกฺขเว ปหาน\-ปฺปธานํ.}

\prose{กตมญฺจ ภิกฺขเว ภาวนา\-ปฺปธานํ? อิธ ภิกฺขเว ภิกฺขุ สติสมฺโพชฺฌงฺคํ ภาเวติ วิเวกนิสฺสิตํ วิราคนิสฺสิตํ นิโรธ\-นิสฺสิตํ โวสคฺค\-ปริณามิํ, \csromanfolio{324}ธมฺมวิจย\-สมฺโพชฺฌงฺคํ ภาเวติ. วิริย\-สมฺโพชฺฌงฺคํ ภาเวติ. ปีติสมฺโพชฺฌงฺคํ ภาเวติ. ปสฺสทฺธิ\-สมฺโพชฺฌงฺคํ ภาเวติ. สมาธิ\-สมฺโพชฺฌงฺคํ ภาเวติ. อุเปกฺขา\-สมฺโพชฺฌงฺคํ ภาเวติ วิเวกนิสฺสิตํ วิราคนิสฺสิตํ นิโรธ\-นิสฺสิตํ โวสคฺค\-ปริณามิํ. อิทํ วุจฺจติ ภิกฺขเว ภาวนา\-ปฺปธานํ.}

\proseitem{14}{กตมญฺจ ภิกฺขเว อนุรกฺขณา\-ปฺปธานํ? อิธ ภิกฺขเว ภิกฺขุ อุปฺปนฺนํ ภทฺทกํ สมาธิ\-นิมิตฺตํ อนุรกฺขติ อฏฺฐิกสญฺญํ ปุฬวกสญฺญํ วินีลกสญฺญํ วิจฺฉิทฺทก\-สญฺญํ อุทฺธุมาตก\-สญฺญํ. อิทํ วุจฺจติ ภิกฺขเว อนุรกฺขณา\-ปฺปธานํ. อิมานิ โข ภิกฺขเว จตฺตาริ ปธานานีติ.}

\setlength{\csromangathapagewidth}{0pt}
\setlength{\csromangathawakwidth}{0pt}
\csromangathameasuregroup{สํวโร จ ปหานญฺจ, \\ เอเต ปธานา จตฺตาโร, \\ เยหิ ภิกฺขุ อิธาตาปี,}{\csromangathabat{สํวโร จ ปหานญฺจ,}{ภาวนา อนุรกฺขณา.} \\ \csromangathabat{เอเต ปธานา จตฺตาโร,}{เทสิตาทิจฺจพนฺธุนา.} \\ \csromangathabat{เยหิ ภิกฺขุ อิธาตาปี,}{ขยํ ทุกฺขสฺส ปาปุเณติ.} \\ จตุตฺถํ.}
\csromangathasetleft{สํวโร จ ปหานญฺจ, \\ เอเต ปธานา จตฺตาโร, \\ เยหิ ภิกฺขุ อิธาตาปี,}
\csromangathagroup{\csromangathastanza{\csromangathabat{สํวโร จ ปหานญฺจ,}{ภาวนา อนุรกฺขณา.} \\ \csromangathabat{เอเต ปธานา จตฺตาโร,}{เทสิตาทิจฺจพนฺธุนา.}}\csromangathastanza{\csromangathabat{เยหิ ภิกฺขุ อิธาตาปี,}{ขยํ ทุกฺขสฺส ปาปุเณติ.} \\ จตุตฺถํ.}}

\csromanmatikaanchor{mtk.69}
\csromantocmarkref{subsection}{5. ปญฺญตฺติสุตฺต}{mtk.69}
\bookmark[dest={mtk.69},level=2]{5. ปญฺญตฺติสุตฺต}
\csromanheader{2}{5. \textbf{ปญฺญตฺติสุตฺต}}
\proseitem{15}{จตสฺโส อิมา ภิกฺขเว อคฺคปญฺญตฺติโย. กตมา จตสฺโส? เอตทคฺคํ ภิกฺขเว อตฺตภาวีนํ ยทิทํ ราหุ อสุรินฺโท. เอตทคฺคํ ภิกฺขเว กามโภคีนํ ยทิทํ ราชา มนฺธาตา. เอตทคฺคํ ภิกฺขเว อาธิปเตยฺยานํ ยทิทํ มาโร ปาปิมา. สเทวเก ภิกฺขเว โลเก สมารเก สพฺรหฺมเก สสฺสมณพฺรหฺมณิยา ปชาย สเทว\-มนุสฺสาย ตถาคโต อคฺคมกฺขายติ อรหํ สมฺมาสมฺพุทฺโธ. อิมา โข ภิกฺขเว จตสฺโส อคฺค\-ปญฺญตฺติโยติ.}

\setlength{\csromangathapagewidth}{0pt}
\setlength{\csromangathawakwidth}{0pt}
\csromangathameasuregroup{ราหุคฺคํ อตฺตภาวีนํ, \\ มาโร อาธิปเตยฺยานํ, \\ อุทฺธํ ติริยํ อปาจีนํ, \\ สเทวกสฺส โลกสฺส,}{\csromangathabat{ราหุคฺคํ อตฺตภาวีนํ,}{มนฺธาตา กามโภคินํ.} \\ \csromangathabat{มาโร อาธิปเตยฺยานํ,}{อิทฺธิยา ยสสา ชลํ.} \\ \csromangathabat{อุทฺธํ ติริยํ อปาจีนํ,}{ยาวตา ชคโต คติ.} \\ \csromangathabat{สเทวกสฺส โลกสฺส,}{พุทฺโธ อคฺโค ปวุจฺจตีติ.} \\ ปญฺจมํ.}
\csromangathameasurewak{ปญฺจมํ.}
\csromangathasetleft{ราหุคฺคํ อตฺตภาวีนํ, \\ มาโร อาธิปเตยฺยานํ, \\ อุทฺธํ ติริยํ อปาจีนํ, \\ สเทวกสฺส โลกสฺส,}
\csromangathagroup{\csromangathastanza{\csromangathabat{ราหุคฺคํ อตฺตภาวีนํ,}{มนฺธาตา กามโภคินํ.} \\ \csromangathabat{มาโร อาธิปเตยฺยานํ,}{อิทฺธิยา ยสสา ชลํ.}}\csromangathastanza{\csromangathabat{อุทฺธํ ติริยํ อปาจีนํ,}{ยาวตา ชคโต คติ.} \\ \csromangathabat{สเทวกสฺส โลกสฺส,}{พุทฺโธ อคฺโค ปวุจฺจตีติ.}}\csromangathastanza{\csromangathawak{ปญฺจมํ.}}}

\csromanmatikaanchor{mtk.70}
\csromantocmarkref{subsection}{6. โสขุมฺมสุตฺต}{mtk.70}
\bookmark[dest={mtk.70},level=2]{6. โสขุมฺมสุตฺต}
\csromanheader{2}{6. \textbf{โสขุมฺมสุตฺต}}
\csromanmatikaanchor{mtk.71}
\csromantocmarkref{subsection}{7-9. อคติสุตฺตานิ}{mtk.71}
\bookmark[dest={mtk.71},level=2]{7-9. อคติสุตฺตานิ}
\proseitem{14}{จตฺตาริมานิ ภิกฺขเว โสขุมฺมานิ. กตมานิ จตฺตาริ? อิธ ภิกฺขเว ภิกฺขุ รูปโสขุมฺเมน สมนฺนาคโต โหติ ปรเมน, เตน จ \csromanfolio{325}รูปโสขุมฺเมน อญฺญํ รูปโสขุมฺมํ อุตฺตริตรํ วา ปณีตตรํ วา น สมนุปสฺสติ, เตน จ รูปโสขุมฺเมน อญฺญํ รูปโสขุมฺมํ อุตฺตริตรํ วา ปณีตตรํ วา น ปตฺเถติ. เวทนา\-โสขุมฺเมน สมนฺนาคโต โหติ ปรเมน, เตน จ เวทนา\-โสขุมฺเมน อญฺญํ เวทนา\-โสขุมฺมํ อุตฺตริตรํ วา ปณีตตรํ วา น สมนุปสฺสติ, เตน จ เวทนา\-โสขุมฺเมน อญฺญํ เวทนา\-โสขุมฺมํ อุตฺตริตรํ วา ปณีตตรํ วา น ปตฺเถติ. สญฺญา\-โสขุมฺเมน สมนฺนาคโต โหติ ปรเมน, เตน จ สญฺญา\-โสขุมฺเมน อญฺญํ สญฺญาโสขุมฺมํ อุตฺตริตรํ วา ปณีตตรํ วา น สมนุปสฺสติ, เตน จ สญฺญา\-โสขุมฺเมน อญฺญํ สญฺญาโสขุมฺมํ อุตฺตริตรํ วา ปณีตตรํ วา น ปตฺเถติ. สงฺขาร\-โสขุมฺเมน สมนฺนาคโต โหติ ปรเมน, เตน จ สงฺขาร\-โสขุมฺเมน อญฺญํ สงฺขาร\-โสขุมฺมํ อุตฺตริตรํ วา ปณีตตรํ วา น สมนุปสฺสติ, เตน จ สงฺขาร\-โสขุมฺเมน อญฺญํ สงฺขาร\-โสขุมฺมํ อุตฺตริตรํ วา ปณีตตรํ วา น ปตฺเถติ. อิมานิ โข ภิกฺขเว จตฺตาริ โสขุมฺมานีติ.}

\setlength{\csromangathapagewidth}{0pt}
\setlength{\csromangathawakwidth}{0pt}
\csromangathameasuregroup{รูปโสขุมฺมตํ ญตฺวา, \\ สญฺญา ยโต สมุเทติ, \\ สงฺขาเร ปรโต ญตฺวา, \\ ส เว สมฺมทฺทโส ภิกฺขุ, \\ ธาเรติ อนฺติมํ เทหํ,}{\csromangathabat{รูปโสขุมฺมตํ ญตฺวา,}{เวทนานญฺจ สมฺภวํ.} \\ \csromangathabat{สญฺญา ยโต สมุเทติ,}{อตฺถํ คจฺฉติ ยตฺถ จ.} \\ \csromangathabat{สงฺขาเร ปรโต ญตฺวา,}{ทุกฺขโต โน จ อตฺตโต.} \\ \csromangathabat{ส เว สมฺมทฺทโส ภิกฺขุ,}{สนฺโต สนฺติปเท รโต.} \\ \csromangathabat{ธาเรติ อนฺติมํ เทหํ,}{เชตฺวา มารํ สวาหินินฺติ.} \\ ฉฏฺฐํ.}
\csromangathasetleft{รูปโสขุมฺมตํ ญตฺวา, \\ สญฺญา ยโต สมุเทติ, \\ สงฺขาเร ปรโต ญตฺวา, \\ ส เว สมฺมทฺทโส ภิกฺขุ, \\ ธาเรติ อนฺติมํ เทหํ,}
\csromangathagroup{\csromangathastanza{\csromangathabat{รูปโสขุมฺมตํ ญตฺวา,}{เวทนานญฺจ สมฺภวํ.} \\ \csromangathabat{สญฺญา ยโต สมุเทติ,}{อตฺถํ คจฺฉติ ยตฺถ จ.}}\csromangathastanza{\csromangathabat{สงฺขาเร ปรโต ญตฺวา,}{ทุกฺขโต โน จ อตฺตโต.} \\ \csromangathabat{ส เว สมฺมทฺทโส ภิกฺขุ,}{สนฺโต สนฺติปเท รโต.}}\csromangathastanza{\csromangathabat{ธาเรติ อนฺติมํ เทหํ,}{เชตฺวา มารํ สวาหินินฺติ.} \\ ฉฏฺฐํ.}}

\prose{7. ปฐม\-อคติ\-สุตฺต}

\proseitem{17}{จตฺตาริมานิ ภิกฺขเว อคติคมนานิ. กตมานิ จตฺตาริ? ฉนฺทาคติํ คจฺฉติ, โทสาคติํ คจฺฉติ, โมหาคติํ คจฺฉติ, ภยาคติํ คจฺฉติ. อิมานิ โข ภิกฺขเว จตฺตาริ อคติคมนานีติ.}

\setlength{\csromangathapagewidth}{0pt}
\setlength{\csromangathawakwidth}{0pt}
\csromangathameasuregroup{ฉนฺทา โทสา ภยา โมหา, \\ นิหียติ ตสฺส ยโส,}{\csromangathabat{ฉนฺทา โทสา ภยา โมหา,}{โย ธมฺมํ อติวตฺตติ.} \\ \csromangathabat{นิหียติ ตสฺส ยโส,}{กาฬปกฺเขว จนฺติมาติ.}}
\csromangathasetleft{ฉนฺทา โทสา ภยา โมหา, \\ นิหียติ ตสฺส ยโส,}
\csromangathagroup{\csromangathastanza{\csromangathabat{ฉนฺทา โทสา ภยา โมหา,}{โย ธมฺมํ อติวตฺตติ.} \\ \csromangathabat{นิหียติ ตสฺส ยโส,}{กาฬปกฺเขว จนฺติมาติ.}}}

\csromanpalirecto
\gambhira{จตุกฺกนิปาต\-ปาฬิ 8. ทุติย\-อคติ\-สุตฺต}
\proseitem{18}{\csromanfolio{326}จตฺตาริมานิ ภิกฺขเว นาคติคมนานิ. กตมานิ จตฺตาริ? น ฉนฺทาคติํ คจฺฉติ, น โทสาคติํ คจฺฉติ, น โมหาคติํ คจฺฉติ, น ภยาคติํ คจฺฉติ. อิมานิ โข ภิกฺขเว จตฺตาริ นาคติ\-คมนานีติ.}

\setlength{\csromangathapagewidth}{0pt}
\setlength{\csromangathawakwidth}{0pt}
\csromangathameasuregroup{ฉนฺทา โทสา ภยา โมหา, \\ อาปูรติ ตสฺส ยโส,}{\csromangathabat{ฉนฺทา โทสา ภยา โมหา,}{โย ธมฺมํ นาติวตฺตติ.} \\ \csromangathabat{อาปูรติ ตสฺส ยโส,}{สุกฺกปกฺเขว จนฺทิมาติ.} \\ อฏฺฐมํ.}
\csromangathameasurewak{อฏฺฐมํ.}
\csromangathasetleft{ฉนฺทา โทสา ภยา โมหา, \\ อาปูรติ ตสฺส ยโส,}
\csromangathagroup{\csromangathastanza{\csromangathabat{ฉนฺทา โทสา ภยา โมหา,}{โย ธมฺมํ นาติวตฺตติ.} \\ \csromangathabat{อาปูรติ ตสฺส ยโส,}{สุกฺกปกฺเขว จนฺทิมาติ.}}\csromangathastanza{\csromangathawak{อฏฺฐมํ.}}}

\prose{9. ตติย\-อคติ\-สุตฺต}

\proseitem{19}{จตฺตริมานิ ภิกฺขเว อคติคมนานิ. กตมานิ จตฺตาริ? ฉนฺทาคติํ คจฺฉติ, โทสาคติํ คจฺฉติ, โมหาคติํ คจฺฉติ, ภยาคติํ คจฺฉติ. อิมานิ โข ภิกฺขเว จตฺตาริ อคติคมนานิ.}

\prose{จตฺตาริมานิ ภิกฺขเว นาคติคมนานิ. กตมานิ จตฺตาริ? น ฉนฺทาคติํ คจฺฉติ, น โทสาคติํ คจฺฉติ, น โมหาคติํ คจฺฉติ, น ภยาคติํ คจฺฉติ. อิมานิ โข ภิกฺขเว จตฺตาริ นาคติ\-คมนานีติ.}

\setlength{\csromangathapagewidth}{0pt}
\setlength{\csromangathawakwidth}{0pt}
\csromangathameasuregroup{ฉนฺทา โทสา ภยา โมหา, \\ นิหียติ ตสฺส ยโส, \\ ฉนฺทา โทสา ภยา โมหา, \\ อาปูรติ ตสฺส ยโส,}{\csromangathabat{ฉนฺทา โทสา ภยา โมหา,}{โย ธมฺมํ อติวตฺตติ.} \\ \csromangathabat{นิหียติ ตสฺส ยโส,}{กาฬปกฺเขว จนฺทิมา.} \\ \csromangathabat{ฉนฺทา โทสา ภยา โมหา,}{โย ธมฺมํ นาติวตฺตติ.} \\ \csromangathabat{อาปูรติ ตสฺส ยโส,}{สุกฺกปกฺเขว จนฺทิมาติ.} \\ นวมํ.}
\csromangathameasurewak{นวมํ.}
\csromangathasetleft{ฉนฺทา โทสา ภยา โมหา, \\ นิหียติ ตสฺส ยโส, \\ ฉนฺทา โทสา ภยา โมหา, \\ อาปูรติ ตสฺส ยโส,}
\csromangathagroup{\csromangathastanza{\csromangathabat{ฉนฺทา โทสา ภยา โมหา,}{โย ธมฺมํ อติวตฺตติ.} \\ \csromangathabat{นิหียติ ตสฺส ยโส,}{กาฬปกฺเขว จนฺทิมา.}}\csromangathastanza{\csromangathabat{ฉนฺทา โทสา ภยา โมหา,}{โย ธมฺมํ นาติวตฺตติ.} \\ \csromangathabat{อาปูรติ ตสฺส ยโส,}{สุกฺกปกฺเขว จนฺทิมาติ.}}\csromangathastanza{\csromangathawak{นวมํ.}}}

\csromanmatikaanchor{mtk.72}
\csromantocmarkref{subsection}{10. ภตฺตุทฺเทสกสุตฺต}{mtk.72}
\bookmark[dest={mtk.72},level=2]{10. ภตฺตุทฺเทสกสุตฺต}
\csromanheader{2}{10. ภตฺตุทฺเทสก\-สุตฺต}
\proseitem{20}{จตูหิ ภิกฺขเว ธมฺเมหิ สมนฺนาคโต ภตฺตุทฺเทสโก ยถาภตํ นิกฺขิตฺโต เอวํ นิรเย. กตเมหิ จตูหิ? ฉนฺทาคติํ คจฺฉติ, โทสาคติํ คจฺฉติ, โมหาคติํ คจฺฉติ, ภยาคติํ คจฺฉติ. อิเมหิ โข ภิกฺขเว จตูหิ ธมฺเมหิ สมนฺนาคโต ภตฺตุทฺเทสโก ยถาภตํ นิกฺขิตฺโต เอวํ นิรเย.}

\csromanmatikaanchor{mtk.73}
\csromanmatikaanchor{mtk.74}
\csromantocmarknopage{section}{3. อุรุเวลวคฺค}{mtk.73}
\bookmark[dest={mtk.73},level=1]{3. อุรุเวลวคฺค}
\csromantocmarkref{subsection}{1-2. อุรุเวลสุตฺตานิ}{mtk.74}
\bookmark[dest={mtk.74},level=2]{1-2. อุรุเวลสุตฺตานิ}
\prose{จตูหิ ภิกฺขเว ธมฺเมหิ สมนฺนาคโต ภตฺตุทฺเทสโก ยถาภตํ นิกฺขิตฺโต เอวํ สคฺเค. กตเมหิ จตูหิ? น ฉนฺทาคติํ คจฺฉติ, น \csromanfolio{327}โทสาคติํ คจฺฉติ, น โมหาคติํ คจฺฉติ, น ภยาคติํ คจฺฉติ. อิเมหิ โข ภิกฺขเว จตูหิ ธมฺเมหิ สมนฺนาคโต ภตฺตุทฺเทสโก ยถาภตํ นิกฺขิตฺโต เอวํ สคฺเคติ.}

\setlength{\csromangathapagewidth}{0pt}
\setlength{\csromangathawakwidth}{0pt}
\csromangathameasuregroup{}{เย เกจิ กาเมสุ อสญฺญตา ชนา, \\ อธมฺมิกา โหนฺติ อธมฺมคารวา. \\ ฉนฺทา โทสา โมหา จ ภยา คามิโน\textsuperscript{1}, \\ ปริสากสโฏ\textsuperscript{1} จ ปเนส วุจฺจติ. \\ เอวํ หิ วุตฺตํ สมเณน ชานตา, \\ ตสฺมา หิ เต สปฺปุริสา ปสํสิยา. \\ ธมฺเม ฐิตา เย น กโรนฺติ ปาปกํ, \\ น ฉนฺทา น โทสา น โมหา น ภยา จ คามิโน\textsuperscript{1}. \\ ปริสาย มณฺโฑ จ ปเนส วุจฺจติ, \\ เอวํ หิ วุตฺตํ สมเณน ชานตาติ.ทสมํ.}
\csromangathameasurewak{เย เกจิ กาเมสุ อสญฺญตา ชนา, \\ อธมฺมิกา โหนฺติ อธมฺมคารวา. \\ ฉนฺทา โทสา โมหา จ ภยา คามิโน\textsuperscript{1}, \\ ปริสากสโฏ\textsuperscript{1} จ ปเนส วุจฺจติ. \\ เอวํ หิ วุตฺตํ สมเณน ชานตา, \\ ตสฺมา หิ เต สปฺปุริสา ปสํสิยา. \\ ธมฺเม ฐิตา เย น กโรนฺติ ปาปกํ, \\ น ฉนฺทา น โทสา น โมหา น ภยา จ คามิโน\textsuperscript{1}. \\ ปริสาย มณฺโฑ จ ปเนส วุจฺจติ, \\ เอวํ หิ วุตฺตํ สมเณน ชานตาติ.ทสมํ.}
\setlength{\csromangathaleftcol}{0pt}
\csromangathagroupwithclosermidruled{\csromangathastanza{\csromangathawak{เย เกจิ กาเมสุ อสญฺญตา ชนา, \\ อธมฺมิกา โหนฺติ อธมฺมคารวา. \\ ฉนฺทา โทสา โมหา จ ภยา คามิโน\csromansharedfootnote{8ea7f7c64f9192d7}{ฉนฺทา จ โทสา จ ภยา จ คามิโน (สี, สฺยา, กํ, อิ)}, \\ ปริสากสโฏ\csromansharedfootnote{97272936c4fec9c7}{ปริสกฺกสาโว (สี, สฺยา, กํ, อิ)} จ ปเนส วุจฺจติ.}}\csromangathastanza{\csromangathawak{เอวํ หิ วุตฺตํ สมเณน ชานตา, \\ ตสฺมา หิ เต สปฺปุริสา ปสํสิยา. \\ ธมฺเม ฐิตา เย น กโรนฺติ ปาปกํ, \\ น ฉนฺทา น โทสา น โมหา น ภยา จ คามิโน\csromansharedfootnote{12ca2c7975e13c5f}{น ฉนฺทโทสา น ภยา จ คามิโน (สี, สฺยา, กํ, อิ)}. \\ ปริสาย มณฺโฑ จ ปเนส วุจฺจติ, \\ เอวํ หิ วุตฺตํ สมเณน ชานตาติ.  ทสมํ.}}}{จรวคฺโค ทุติโย.}
\csromancenter{\textbf{ตสฺสุทฺทานํ}}
\setlength{\csromangathapagewidth}{0pt}
\setlength{\csromangathawakwidth}{0pt}
\csromangathameasuregroup{จรํ สีลํ ปธานานิ, \\ โสขุมฺมํ ตโย อคตี,}{\csromangathabat{จรํ สีลํ ปธานานิ,}{สํวรํ ปญฺญตฺติปญฺจมํ.} \\ \csromangathabat{โสขุมฺมํ ตโย อคตี,}{ภตฺตุทฺเทเสน เต ทสาติ.}}
\csromangathasetleft{จรํ สีลํ ปธานานิ, \\ โสขุมฺมํ ตโย อคตี,}
\csromangathagroup{\csromangathastanza{\csromangathabat{จรํ สีลํ ปธานานิ,}{สํวรํ ปญฺญตฺติปญฺจมํ.} \\ \csromangathabat{โสขุมฺมํ ตโย อคตี,}{ภตฺตุทฺเทเสน เต ทสาติ.}}}

\chapterhead{3. อุรุเวลวคฺค}
\prose{1. ปฐมอุรุเวลสุตฺต}

\proseitem{21}{เอวํ เม สุตํ\csromansharedfootnote{b0213bef2182ef76}{สํ 1. 140 ปิฏฺเฐปิ อาคตํ.}\csromandash{}เอกํ สมยํ ภควา สาวตฺถิยํ วิหรติ เชตวเน อนาถ\-ปิณฺฑิกสฺส อาราเม. ตตฺร โข ภควา ภิกฺขู อามนฺเตสิ “ภิกฺขโว”ติ. “ภทนฺเต”ติ เต ภิกฺขู ภควโต ปจฺจสฺโสสุํ. ภควา เอตทโวจ\csromandash{}}

\prose{\csromanfolio{328}เอกมิทาหํ ภิกฺขเว สมยํ อุรุเวลายํ วิหรามิ นชฺชา เนรญฺชราย ตีเร อชปาล\-นิโคฺรเธ ปฐมา\-ภิสมฺพุทฺโธ, ตสฺส มยฺหํ ภิกฺขเว รโหคตสฺส ปฏิสลฺลีนสฺส เอวํ เจตโส ปริวิตกฺโก อุทปาทิ “ทุกฺขํ โข อคารโว วิหรติ อปฺปติสฺโส, กิํ นุ โข อหํ สมณํ วา พฺราหฺมณํ วา สกฺกตฺวา ครุํ กตฺวา\csromansharedfootnote{7122a79891601e79}{ครุกตฺวา (สี, อิ)} อุปนิสฺสาย วิหเรยฺยนฺ”ติ.}

\prose{ตสฺส มยฺหํ ภิกฺขเว เอตทโหสิ\csromandash{}อปริปูรสฺส โข อหํ สีลกฺขนฺธสฺส ปาริปูริยา อญฺญํ สมณํ วา พฺราหฺมณํ วา สกฺกตฺวา ครุํ กตฺวา อุปนิสฺสาย วิหเรยฺยํ. น โข ปนาหํ ปสฺสามิ สเทวเก โลเก สมารเก สพฺรหฺมเก สสฺสมณ\-พฺราหฺมณิยา ปชาย สเทว\-มนุสฺสาย อญฺญํ สมณํ วา พฺราหฺมณํ วา อตฺตนา สีลสมฺปนฺนตรํ, ยมหํ สกฺกตฺวา ครุํ กตฺวา อุปนิสฺสาย วิหเรยฺยํ.}

\prose{อปริปูรสฺส โข อหํ สมาธิ\-กฺขนฺธสฺส ปาริปูริยา อญฺญํ สมณํ วา พฺราหฺมณํ วา สกฺกตฺวา ครุํ กตฺวา อุปนิสฺสาย วิหเรยฺยํ. น โข ปนาหํ ปสฺสามิ สเทวเก โลเก สมารเก สพฺรหฺมเก สสฺสมณ\-พฺราหฺมณิยา ปชาย สเทว\-มนุสฺสาย อญฺญํ สมณํ วา พฺราหฺมณํ วา อตฺตนา สมาธิ\-สมฺปนฺนตรํ, ยมหํ สกฺกตฺวา ครุํ กตฺวา อุปนิสฺสาย วิหเรยฺยํ.}

\prose{อปริปูรสฺส โข อหํ ปญฺญา\-กฺขนฺธสฺส ปาริปูริยา อญฺญํ สมณํ วา พฺราหฺมณํ วา สกฺกตฺวา ครุํ กตฺวา อุปนิสฺสาย วิหเรยฺยํ. น โข ปนาหํ ปสฺสามิ สเทวเก โลเก สมารเก สพฺรหฺมเก สสฺสมณพฺราหฺมณิย ปชาย สเทว\-มนุสฺสาย อญฺญํ สมณํ วา พฺราหฺมณํ วา อตฺตนา ปญฺญา\-สมฺปนฺนตรํ, ยมหํ สกฺกตฺวา ครุํ กตฺวา อุปนิสฺสาย วิหเรยฺยํ.}

\prose{อปริปูรสฺส โข อหํ วิมุตฺติ\-กฺขนฺธสฺส ปาริปูริยา อญฺญํ สมณํ วา พฺราหฺมณํ วา สกฺกตฺวา ครุํ กตฺวา อุปนิสฺสาย วิหเรยฺยํ. น โข ปนาหํ ปสฺสามิ สเทวเก โลเก สมารเก สพฺรหฺมเก สสฺสมณ\-พฺราหฺมณิยา ปชาย สเทว\-มนุสฺสาย อญฺญํ สมณํ วา พฺราหฺมณํ วา อตฺตนา วิมุตฺติ\-สมฺปนฺนตรํ, ยมหํ สกฺกตฺวา ครุํ กตฺวา อุปนิสฺสาย วิหเรยฺยนฺติ.}

\chapterheadpageread{3. อุรุเวลวคฺค}
\prose{\csromanfolio{329}ตสฺส มยฺหํ ภิกฺขเว เอตทโหสิ “ยํนูนาหํ ยฺวายํ\csromansharedfootnote{cfbf7854997929ac}{โยปายํ (สี, สฺยา, กํ, อิ)} ธมฺโม มยา อภิสมฺพุทฺโธ, ตเมว ธมฺมํ สกฺกตฺวา ครุํ กตฺวา อุปนิสฺสาย วิหเรยฺยนฺ”ติ.}

\prose{อถ โข ภิกฺขเว พฺรหฺมา สหมฺปติ มม เจตสา เจโต\-ปริวิตกฺกม\-ญฺญาย เสยฺยถาปิ นาม พลวา ปุริโส สมิญฺชิตํ\csromansharedfootnote{6d2aa19b3b2e4dad}{สมฺมิญฺชิตํ (สี, สฺยา, กํ, อิ) 3. เย จพฺภตีตา (สี, อิ, ก)} วา พาหํ ปสาเรยฺย, ปสาริตํ วา พาหํ สมิญฺเชยฺย. เอวเมวํ พฺรหฺมโลเก อนฺตรหิโต มม ปุรโต ปาตุรโหสิ. อถ โข ภิกฺขเว พฺรหฺมา สหมฺปติ เอกํสํ อุตฺตราสงฺคํ กริตฺวา ทกฺขิณํ ชาณุมณฺฑลํ ปถวิยํ นิหนฺตฺวา เยนาหํ เตนญฺชลิํ ปณาเมตฺวา มํ เอตทโวจ “เอวเมตํ ภควา, เอวเมตํ สุคต, เยปิ เต ภนฺเต อเหสุํ อตีตมทฺธานํ อรหนฺโต สมฺมาสมฺพุทฺธา, เตปิ ภควนฺโต ธมฺมํเยว สกฺกตฺวา ครุํ กตฺวา อุปนิสฺสาย วิหริํสุ. เยปิ เต ภนฺเต ภวิสฺสนฺติ อนาคตมทฺธานํ อรหนฺโต สมฺมาสมฺพุทฺธา, เตปิ ภควนฺโต ธมฺมํเยว สกฺกตฺวา ครุํ กตฺวา อุปนิสฺสาย วิหริสฺสนฺติ. ภควาปิ ภนฺเต เอตรหิ อรหํ สมฺมาสมฺพุทฺโธ ธมฺมํเยว สกฺกตฺวา ครุํ กตฺวา อุปนิสฺสาย วิหรตู”ติ. อิทเมโวจ พฺรหฺมา สหมฺปติ, อิทํ วตฺวา อถาปรํ เอตทโวจ\csromandash{}}

\setlength{\csromangathapagewidth}{0pt}
\setlength{\csromangathawakwidth}{0pt}
\csromangathameasuregroup{“เย จ อตีตา\textsuperscript{1} สมฺพุทฺธา, \\ โย เจตรหิ สมฺพุทฺโธ, \\ สพฺเพ สทฺธมฺมครุโน, \\ อโถปิ วิหริสฺสนฺติ, \\ ตสฺมา หิ อตฺตกาเมน\textsuperscript{1}, \\ สทฺธมฺโม ครุกาตพฺโพ,}{\csromangathabat{“เย จ อตีตา\textsuperscript{1} สมฺพุทฺธา,}{เย จ พุทฺธา อนาคตา.} \\ \csromangathabat{โย เจตรหิ สมฺพุทฺโธ,}{พหูนํ\textsuperscript{1} โสกนาสโน.} \\ \csromangathabat{สพฺเพ สทฺธมฺมครุโน,}{วิหํสุ\textsuperscript{1} วิหรนฺติ จ.} \\ \csromangathabat{อโถปิ วิหริสฺสนฺติ,}{เอสา พุทฺธาน ธมฺมตา.} \\ \csromangathabat{ตสฺมา หิ อตฺตกาเมน\textsuperscript{1},}{มหตฺตมภิกงฺขตา.} \\ \csromangathabat{สทฺธมฺโม ครุกาตพฺโพ,}{สรํ พุทฺธาน สาสนนฺ”ติ.}}
\csromangathasetleft{“เย จ อตีตา\textsuperscript{1} สมฺพุทฺธา, \\ โย เจตรหิ สมฺพุทฺโธ, \\ สพฺเพ สทฺธมฺมครุโน, \\ อโถปิ วิหริสฺสนฺติ, \\ ตสฺมา หิ อตฺตกาเมน\textsuperscript{1}, \\ สทฺธมฺโม ครุกาตพฺโพ,}
\csromangathagroup{\csromangathastanza{\csromangathabat{“เย จ อตีตา\csromansharedfootnote{2474c54ceb083c0a}{เย จพฺภตีตา (สี, อิ, ก)} สมฺพุทฺธา,}{เย จ พุทฺธา อนาคตา.} \\ \csromangathabat{โย เจตรหิ สมฺพุทฺโธ,}{พหูนํ\csromansharedfootnote{b970d4c076431ca8}{พหุนฺนํ (สี, สฺยา, กํ, อิ) สํ 1. 142 ปิฏฺเฐปิ. 5. วิหริํสุ (สฺยา, กํ)} โสกนาสโน.}}\csromangathastanza{\csromangathabat{สพฺเพ สทฺธมฺมครุโน,}{วิหํสุ\csromansharedfootnote{2aab8b8a5305b7b2}{วิหริํสุ (สฺยา, กํ)} วิหรนฺติ จ.} \\ \csromangathabat{อโถปิ วิหริสฺสนฺติ,}{เอสา พุทฺธาน ธมฺมตา.}}\csromangathastanza{\csromangathabat{ตสฺมา หิ อตฺตกาเมน\csromansharedfootnote{a575e64f334db9ef}{อตฺถกาเมน (สี, ก) 7. โยปายํ (สพฺพตฺถ)},}{มหตฺตม\-ภิกงฺขตา.} \\ \csromangathabat{สทฺธมฺโม ครุกาตพฺโพ,}{สรํ พุทฺธาน สาสนนฺ”ติ.}}}

\prose{อิทมโวจ ภิกฺขเว พฺรหฺมา สหมฺปติ, อิทํ วตฺวา มํ อภิวาเทตฺวา ปทกฺขิณํ กตฺวา ตตฺเถ\-ว\-นฺตรธายิ. อถ ขฺวาหํ ภิกฺขเว พฺรหฺมุโน จ อชฺเฌสนํ วิทิตฺวา อตฺตโน จ ปติรูปํ ยฺวายํ\csromansharedfootnote{7061fd40f6b07a3c}{โยปายํ (สพฺพตฺถ)} ธมฺโม มยา อภิสมฺพุทฺโธ, ตเมว ธมฺมํ สกฺกตฺวา ครุํ กตฺวา อุปนิสฺสาย วิหาสิํ.}

\prose{\csromanfolio{330}ยโต จ โข ภิกฺขเว สํโฆปิ มหตฺเตน สมนฺนาคโต, อถ เม สํเฆปิ คารโวติ.  ปฐมํ.}

\prose{2. ทุติยอุรุเวลสุตฺต}

\proseitem{22}{เอกมิทาหํ ภิกฺขเว สมยํ อุรุเวลายํ วิหรามิ นชฺชา เนรญฺชราย ตีเร อชปาล\-นิโคฺรเธ ปฐมา\-ภิสมฺพุทฺโธ. อถ โข ภิกฺขเว สมฺพหุลา พฺราหฺมณา ชิณฺณา วุทฺธา มหลฺลกา อทฺธคตา วโยอนุปฺปตฺตา เยนาหํ เตนุ\-ปสงฺกมิํสุ, อุปสงฺกมิตฺวา มยา สทฺธิํ สมฺโมทิํสุ, สมฺโมทนียํ กถํ สารณียํ วีติสาเรตฺวา เอกมนฺตํ นิสีทิํสุ, เอกมนฺตํ นิสินฺนา โข ภิกฺขเว เต พฺราหฺมณา มํ เอตทโวจุํ “สุตํ เนตํ\csromansharedfootnote{dbd09aef21ac5fae}{เมตํ (สี, สฺยา, กํ, ก)} โภ โคตม ‘น สมโณ โคตโม พฺราหฺมเณ ชิณฺเณ วุทฺเธ มหลฺลเก อทฺธคเต วโยอนุปฺปตฺเต อภิวาเทติ วา ปจฺจุฏฺเฐติ วา อาสเนน วา นิมนฺเตตี’ติ, ตยิทํ โภ โคตม ตตฺเถว. น หิ ภวํ โคตโม พฺราหฺมเณ ชิณฺเณ วุทฺเธ มหลฺลเก อทฺธคเต วโยอนุปฺปตฺเต อภิวาเทติ วา ปจฺจุฏฺเฐติ วา อาสเนน วา นิมนฺเตติ, ตยิทํ โภ โคตม น สมฺปนฺนเมวา”ติ.}

\prose{ตสฺส มยฺหํ ภิกฺขเว เอตทโหสิ “นยิเม\csromansharedfootnote{4eb6c5194fa7bef0}{น วต เม (สี, อิ), น จยิเม (สฺยา, กํ), น วติเม (?)} อายสฺมนฺโต ชานนฺติ เถรํ วา เถรกรเณ วา ธมฺเม”ติ. วุทฺโธ เจปิ ภิกฺขเว โหติ อาสีติโก วา นาวุติโก วา วสฺสสติโก วา ชาติยา, โส จ โหติ อกาลวาที อภูตวาที อนตฺถวาที อธมฺมวาที อวินยวาที อนิธานวติํ วาจํ ภาสิตา, อกาเลน อนปเทสํ อปริยนฺตวติํ อนตฺถ\-สํหิตํ. อถ โข โส “พาโล เถโร”เตฺวว\csromansharedfootnote{708bdfd17c3e93d2}{เตว (สี, อิ)} สงฺขํ คจฺฉติ.}

\prose{ทหโร เจปิ ภิกฺขเว โหติ ยุวา สุสุกาฬเกโส ภเทฺรน โยพฺพเนน สมนฺนาคโต ปฐเมน วยสา, โส จ โหติ กาลวาที ภูตวาที อตฺถวาที ธมฺมวาที วินยวาที นิธานวติํ วาจํ ภาสิตา, กาเลน สาปเทสํ ปริยนฺตวติํ อตฺถสํหิตํ. อถ โข โส “ปณฺฑิโต เถโร”เตฺวว สงฺขํ คจฺฉติ.}

\chapterheadpageread{3. อุรุเวลวคฺค}
\prose{\csromanfolio{331}จตฺตาโรเม ภิกฺขเว เถรกรณา ธมฺมา. กตเม จตฺตาโร? อิธ ภิกฺขเว ภิกฺขุ สีลวา โหติ, ปาติโมกฺข\-สํวร\-สํวุโต วิหรติ อาจารโคจร\-สมฺปนฺโน อณุมตฺเตสุ วชฺเชสุ ภยทสฺสาวี สมาทาย สิกฺขติ สิกฺขาปเทสุ. พหุสฺสุโต โหติ สุตธโร สุตสนฺนิจโย, เย เต ธมฺมา อาทิกลฺยาณา มชฺเฌกลฺยาณา ปริโยสาน\-กลฺยาณา สาตฺถํ สพฺยญฺชนํ\csromansharedfootnote{ae962218a7776460}{สาตฺถา สพฺยญฺชนา (สี, อิ)} เกวล\-ปริปุณฺณํ\csromansharedfootnote{535459c3dcd528da}{เกวลปริปุณฺณา (สี)} ปริสุทฺธํ พฺรหฺมจริยํ อภิวทนฺติ, ตถารูปาสฺส ธมฺมา พหุสฺสุตา โหนฺติ ธาตา\csromansharedfootnote{acf84a73c1bbfd91}{ธตา (สี, สฺยา, กํ, อิ)} วจสา ปริจิตา มนสา\-นุเปกฺขิตา ทิฏฺฐิยา สุปฺปฏิวิทฺธา. จตุนฺนํ ฌานานํ อาภิเจตสิกานํ ทิฏฺฐ\-ธมฺม\-สุข\-วิหารานํ นิกามลาภี โหติ อกิจฺฉลาภี อกสิรลาภี. อาสวานํ ขยา อนาสวํ เจโตวิมุตฺติํ ปญฺญาวิมุตฺติํ ทิฏฺเฐว ธมฺเม สยํ อภิญฺญา สจฺฉิกตฺวา อุปสมฺปชฺช วิหรติ. อิเม โข ภิกฺขเว จตฺตาโร เถรกรณา ธมฺมาติ.}

\setlength{\csromangathapagewidth}{0pt}
\setlength{\csromangathawakwidth}{0pt}
\csromangathameasuregroup{โย อุทฺธเตน จิตฺเตน, \\ อสมาหิตสงฺกปฺโป, \\ อารา โส ถาวเรยฺยมฺหา, \\ โย จ สีเลน สมฺปนฺโน, \\ สญฺญโต ธีโร ธมฺเมสุ\textsuperscript{1}, \\ ปารคู สพฺพธมฺมานํ, \\ ปหีนชาติมรโณ,}{\csromangathabat{โย อุทฺธเตน จิตฺเตน,}{สมฺผญฺจ พหุ ภาสติ.} \\ \csromangathabat{อสมาหิตสงฺกปฺโป,}{อสทฺธมฺมรโต มโค.} \\ \csromangathabat{อารา โส ถาวเรยฺยมฺหา,}{ปาปทิฏฺฐิ อนาทโร.} \\ \csromangathabat{โย จ สีเลน สมฺปนฺโน,}{สุตวา ปฏิภานวา.} \\ \csromangathabat{สญฺญโต ธีโร ธมฺเมสุ\textsuperscript{1},}{ปญฺญายตฺถํ วิปสฺสติ.} \\ \csromangathabat{ปารคู สพฺพธมฺมานํ,}{อขิโล ปฏิภานวา.} \\ \csromangathabat{ปหีนชาติมรโณ,}{พฺรหฺมจริยสฺส เกวลี.} \\ ทุติยํ.}
\csromangathameasurewak{ทุติยํ.}
\csromangathasetleft{โย อุทฺธเตน จิตฺเตน, \\ อสมาหิตสงฺกปฺโป, \\ อารา โส ถาวเรยฺยมฺหา, \\ โย จ สีเลน สมฺปนฺโน, \\ สญฺญโต ธีโร ธมฺเมสุ\textsuperscript{1}, \\ ปารคู สพฺพธมฺมานํ, \\ ปหีนชาติมรโณ,}
\csromangathagroup{\csromangathastanza{\csromangathabat{โย อุทฺธเตน จิตฺเตน,}{สมฺผญฺจ พหุ ภาสติ.} \\ \csromangathabat{อสมาหิตสงฺกปฺโป,}{อสทฺธมฺมรโต มโค.}}\csromangathastanza{\csromangathabat{อารา โส ถาวเรยฺยมฺหา,}{ปาปทิฏฺฐิ อนาทโร.} \\ \csromangathabat{โย จ สีเลน สมฺปนฺโน,}{สุตวา ปฏิภานวา.}}\csromangathastanza{\csromangathabat{สญฺญโต ธีโร ธมฺเมสุ\csromansharedfootnote{d7f4e055bba3c3a7}{สญฺญโต ธีรธมฺเมสุ (สี), สํยุตฺโต ถิรธมฺเมสุ (สฺยา, กํ)},}{ปญฺญายตฺถํ วิปสฺสติ.} \\ \csromangathabat{ปารคู สพฺพธมฺมานํ,}{อขิโล ปฏิภานวา.} \\ \csromangathabat{ปหีนชาติมรโณ,}{พฺรหฺม\-จริยสฺส เกวลี.}}\csromangathastanza{\csromangathawak{ทุติยํ.}}}

\prose{ตมหํ วทามิ “เถโร”ติ, ยสฺส โน สนฺติ อาสวา.}

\prose{อาสวานํ ขยา ภิกฺขุ, โส “เถโร”ติ ปวุจฺจตีติ.}

\csromanmatikaanchor{mtk.75}
\csromantocmarkref{subsection}{3. โลกสุตฺต}{mtk.75}
\bookmark[dest={mtk.75},level=2]{3. โลกสุตฺต}
\csromanheader{2}{3. \textbf{โลกสุตฺต}}
\proseitem{23}{โลโก ภิกฺขเว ตถาคเตน อภิสมฺพุทฺโธ, โลกสฺมา ตถาคโต วิสํยุตฺโต. โลกสมุทโย ภิกฺขเว ตถาคเตน อภิสมฺพุทฺโธ, โลกสมุทโย ตถาคตสฺส ปหีโน. โลกนิโรโธ ภิกฺขเว ตถาคเตน อภิสมฺพุทฺโธ, โลกนิโรโธ \csromanfolio{332}ตถาคตสฺส สจฺฉิกโต. โลก\-นิโรธ\-คามินี ปฏิปทา ภิกฺขเว ตถาคเตน อภิสมฺพุทฺธา, โลก\-นิโรธ\-คามินี ปฏิปทา ตถาคตสฺส ภาวิตา.}

\prose{ยํ ภิกฺขเว\csromansharedfootnote{d874bab663b4e1dd}{ที 3. 111 ปิฏฺเฐปิ.} สเทวกสฺส โลกสฺส สมารกสฺส สพฺรหฺมกสฺส สสฺสมณพฺรหฺมณิยา ปชาย สเทว\-มนุสฺสาย ทิฏฺฐํ สุตํ มุตํ วิญฺญาตํ ปตฺตํ ปริเยสิตํ อนุวิจริตํ มนสา, สพฺพํ ตํ ตถาคเตน อภิสมฺพุทฺธํ. ตสฺมา “ตถาคโต”ติ วุจฺจติ.}

\prose{ยญฺจ ภิกฺขเว รตฺติํ ตถาคโต อนุตฺตรํ สมฺมาสมฺโพธิํ อภิสมฺ\-พุชฺฌติ, ยญฺจ รตฺติํ อนุปาทิเสสาย นิพฺพานธาตุยา ปรินิพฺพายติ, ยํ เอตสฺมิํ อนฺตเร ภาสติ ลปติ นิทฺทิสติ, สพฺพํ ตํ ตเถว โหติ โน อญฺญถา. ตสฺมา “ตถาคโต”ติ วุจฺจติ.}

\prose{ยถาวาที ภิกฺขเว ตถาคโต ตถาการี, ยถาการี ตถาวาที. อิติ ยถาวาที ตถาการี, ยถาการี ตถาวาที. ตสฺมา “ตถาคโต”ติ วุจฺจติ.}

\prose{สเทวเก ภิกฺขเว โลเก สมารเก สพฺรหฺมเก สสฺสมณ\-พฺราหฺมณิยา ปชาย สเทว\-มนุสฺสาย ตถาคโต อภิภู อนภิภูโต อญฺญทตฺถุ ทโส วสวตฺตี. ตสฺมา “ตถาคโต”ติ วุจฺจติ.}

\setlength{\csromangathapagewidth}{0pt}
\setlength{\csromangathawakwidth}{0pt}
\csromangathameasuregroup{สพฺพํ โลกํ อภิญฺญาย, \\ สพฺพํ โลกํ\textsuperscript{1} วิสํยุตฺโต, \\ ส เว สพฺพาภิภู ธีโร, \\ ผุฏฺฐ’สฺส ปรมา สนฺติ, \\ เอส ขีณาสโว พุทฺโธ, \\ สพฺพกมฺมกฺขยํ ปตฺโต, \\ เอส โส ภควา พุทฺโธ, \\ สเทวกสฺส โลกสฺส, \\ อิติ เทวา มนุสฺสา จ, \\ สงฺคมฺม ตํ นมสฺสนฺติ,}{\csromangathabat{สพฺพํ โลกํ อภิญฺญาย,}{สพฺพํ โลเก ยถาตถํ.} \\ \csromangathabat{สพฺพํ โลกํ\textsuperscript{1} วิสํยุตฺโต,}{สพฺพโลเก อนูปโย.} \\ \csromangathabat{ส เว สพฺพาภิภู ธีโร,}{สพฺพคนฺถปฺปโมจโน.} \\ \csromangathabat{ผุฏฺฐ’สฺส ปรมา สนฺติ,}{นิพฺพานํ อกุโตภยํ.} \\ \csromangathabat{เอส ขีณาสโว พุทฺโธ,}{อนีโฆ ฉินฺนสํสโย.} \\ \csromangathabat{สพฺพกมฺมกฺขยํ ปตฺโต,}{วิมุตฺโต อุปธิสงฺขเย.} \\ \csromangathabat{เอส โส ภควา พุทฺโธ,}{เอส สีโห อนุตฺตโร.} \\ \csromangathabat{สเทวกสฺส โลกสฺส,}{พฺรหฺมจกฺกํ ปวตฺตยี.} \\ \csromangathabat{อิติ เทวา มนุสฺสา จ,}{เย พุทฺธํ สรณํ คตา.} \\ \csromangathabat{สงฺคมฺม ตํ นมสฺสนฺติ,}{มหนฺตํ วีตสารทํ.}}
\csromangathasetleft{สพฺพํ โลกํ อภิญฺญาย, \\ สพฺพํ โลกํ\textsuperscript{1} วิสํยุตฺโต, \\ ส เว สพฺพาภิภู ธีโร, \\ ผุฏฺฐ’สฺส ปรมา สนฺติ, \\ เอส ขีณาสโว พุทฺโธ, \\ สพฺพกมฺมกฺขยํ ปตฺโต, \\ เอส โส ภควา พุทฺโธ, \\ สเทวกสฺส โลกสฺส, \\ อิติ เทวา มนุสฺสา จ, \\ สงฺคมฺม ตํ นมสฺสนฺติ,}
\csromangathagroup{\csromangathastanza{\csromangathabat{สพฺพํ โลกํ อภิญฺญาย,}{สพฺพํ โลเก ยถาตถํ.} \\ \csromangathabat{สพฺพํ โลกํ\csromansharedfootnote{eafe977344dd10d6}{สพฺพโลก (สี, สฺยา, กํ, อิ)} วิสํยุตฺโต,}{สพฺพโลเก อนูปโย.}}\csromangathastanza{\csromangathabat{ส เว สพฺพาภิภู ธีโร,}{สพฺพคนฺถปฺปโมจโน.} \\ \csromangathabat{ผุฏฺฐ’สฺส ปรมา สนฺติ,}{นิพฺพานํ อกุโตภยํ.}}\csromangathastanza{\csromangathabat{เอส ขีณาสโว พุทฺโธ,}{อนีโฆ ฉินฺนสํสโย.} \\ \csromangathabat{สพฺพกมฺม\-กฺขยํ ปตฺโต,}{วิมุตฺโต อุปธิสงฺขเย.}}\csromangathastanza{\csromangathabat{เอส โส ภควา พุทฺโธ,}{เอส สีโห อนุตฺตโร.} \\ \csromangathabat{สเทวกสฺส โลกสฺส,}{พฺรหฺมจกฺกํ ปวตฺตยี.}}\csromangathastanza{\csromangathabat{อิติ เทวา มนุสฺสา จ,}{เย พุทฺธํ สรณํ คตา.} \\ \csromangathabat{สงฺคมฺม ตํ นมสฺสนฺติ,}{มหนฺตํ วีตสารทํ.}}}

\chapterheadpageread{3. อุรุเวลวคฺค}
\setlength{\csromangathapagewidth}{0pt}
\setlength{\csromangathawakwidth}{0pt}
\csromangathameasuregroup{ทนฺโต ทมยตํ เสฏฺโฐ, \\ มุตฺโต โมจยตํ อคฺโค, \\ อิติ เหตํ นมสฺสนฺติ, \\ สเทวกสฺมิํ โลกสฺมิํ,}{\csromangathabat{ทนฺโต ทมยตํ เสฏฺโฐ,}{สนฺโต สมยตํ อิสิ.} \\ \csromangathabat{มุตฺโต โมจยตํ อคฺโค,}{ติณฺโณ ตารยตํ วโร.} \\ \csromangathabat{อิติ เหตํ นมสฺสนฺติ,}{มหนฺตํ วีตสารทํ.} \\ \csromangathabat{สเทวกสฺมิํ โลกสฺมิํ,}{นตฺถิ เม\textsuperscript{1} ปฏิปุคฺคโลติ.} \\ ตติยํ.}
\csromangathameasurewak{ตติยํ.}
\csromangathasetleft{ทนฺโต ทมยตํ เสฏฺโฐ, \\ มุตฺโต โมจยตํ อคฺโค, \\ อิติ เหตํ นมสฺสนฺติ, \\ สเทวกสฺมิํ โลกสฺมิํ,}
\csromangathagroup{\csromangathastanza{\csromanfolio{333}\csromangathabat{ทนฺโต ทมยตํ เสฏฺโฐ,}{สนฺโต สมยตํ อิสิ.} \\ \csromangathabat{มุตฺโต โมจยตํ อคฺโค,}{ติณฺโณ ตารยตํ วโร.}}\csromangathastanza{\csromangathabat{อิติ เหตํ นมสฺสนฺติ,}{มหนฺตํ วีตสารทํ.} \\ \csromangathabat{สเทวกสฺมิํ โลกสฺมิํ,}{นตฺถิ เม\csromansharedfootnote{338dd9d64fd7ee1e}{[มิสฺสินฺคฺ โนเต 0]} ปฏิปุคฺคโลติ.}}\csromangathastanza{\csromangathawak{ตติยํ.}}}

\csromanmatikaanchor{mtk.76}
\csromantocmarkref{subsection}{4. กาฬการามสุตฺต}{mtk.76}
\bookmark[dest={mtk.76},level=2]{4. กาฬการามสุตฺต}
\csromanheader{2}{4. กาฬการาม\-สุตฺต}
\proseitem{24}{เอกํ สมยํ ภควา สาเกเต วิหรติ กาฬการาเม\csromansharedfootnote{e594a854984ab6bb}{โกฬิการาเม (ก)}, ตตฺร โข ภควา ภิกฺขู อามนฺเตสิ “ภิกฺขโว”ติ. “ภทนฺเต”ติ เต ภิกฺขู ภควโต ปจฺจสฺโสสุํ. ภควา เอตทโวจ\csromandash{}}

\prose{ยํ ภิกฺขเว สเทวกสฺส โลกสฺส สมารกสฺส สพฺรหฺมกสฺส สสฺสมณ\-พฺราหฺมณิยา ปชาย สเทว\-มนุสฺสาย ทิฏฺฐํ สุตํ มุตํ วิญฺญาตํ ปตฺตํ ปริเยสิตํ อนุวิจริตํ มนสา, ตมหํ ชานามิ.}

\prose{ยํ ภิกฺขเว สเทวกสฺส โลกสฺส สมารกสฺส สพฺรหฺมกสฺส สสฺสมณ\-พฺราหฺมณิยา ปชาย สเทว\-มนุสฺสาย ทิฏฺฐํ สุตํ มุตํ วิญฺญาตํ ปตฺตํ ปริเยสิตํ อนุวิจริตํ มนสา, ตมหํ อพฺภญฺญาสิํ, ตํ ตถาคตสฺส วิทิตํ, ตํ ตถาคโต น อุปฏฺฐาสิ.}

\prose{ยํ ภิกฺขเว สเทวกสฺส โลกสฺส สมารกสฺส สพฺรหฺมกสฺส สสฺสมณ\-พฺราหฺมณิยา ปชาย สเทว\-มนุสฺสาย ทิฏฺฐํ สุตํ มุตํ วิญฺญาตํ ปตฺตํ ปริเยสิตํ อนุวิจริตํ มนสา, ตมหํ น ชานามีติ วเทยฺยํ, ตํ มมสฺส มุสา.}

\prose{ยํ ภิกฺขเว ฯเปฯ ตมหํ ชานามิ จ น จ ชานามีติ วเทยฺยํ, ตํปสฺส\csromansharedfootnote{7c7f8bc954fb8ec8}{ตํ ปิสฺส (สฺยา, กํ), ตํ มมสฺส (ก)} ตาทิสเมว.}

\prose{ยํ ภิกฺขเว ฯเปฯ ตมหํ เนว ชานามิ น น ชานามีติ วเทยฺยํ, ตํ มมสฺส กลิ.}

\prose{อิติ โข ภิกฺขเว ตถาคโต ทฏฺฐา ทฏฺฐพฺพํ ทิฏฺฐํ น มญฺญติ, อทิฏฺฐํ น มญฺญติ, ทฏฺฐพฺพํ น มญฺญติ, ทฏฺฐารํ น มญฺญติ. สุตฺวา โสตพฺพํ สุตํ \csromanfolio{334}น มญฺญติ, อสุตํ น มญฺญติ, โสตพฺพํ น มญฺญติ, โสตารํ น มญฺญติ. มุตฺวา โมตพฺพํ มุตํ น มญฺญติ, อมุตํ น มญฺญติ, โมตพฺพํ น มญฺญติ, โมตารํ น มญฺญติ. วิญฺญตฺวา วิญฺญาตพฺพํ วิญฺญาตํ น มญฺญติ, อวิญฺญาตํ น มญฺญติ, วิญฺญาตพฺพํ น มญฺญติ, วิญฺญาตารํ น มญฺญติ. อิติ โข ภิกฺขเว ตถาคโต ทิฏฺฐ\-สุต\-มุต\-วิญฺญาตพฺเพสุ ธมฺเมสุ ตาทีเยว ตาที\csromansharedfootnote{db591486fbf3b7d9}{ตาทิโสว ตาที (สฺยา, กํ), ตาทิเส เยว ตาที (อิ), ตาทีเยว ตาทีเยเวกา (ก)}, ตมฺหา จ ปน ตาทิมฺหา\csromansharedfootnote{6b783b2578ec3512}{ตาทิตมฺหา (สี, อิ)} อญฺโญ ตาที อุตฺตริตโร วา ปณีตตโร วา นตฺถีติ วทามีติ.}

\setlength{\csromangathapagewidth}{0pt}
\setlength{\csromangathawakwidth}{0pt}
\csromangathameasuregroup{}{ยํ กิญฺจิ ทิฏฺฐํว สุตํ มุตํ วา, \\ อชฺโฌสิตํ สจฺจมุตํ ปเรสํ. \\ น เตสุ ตาที สยสํวุเตสุ, \\ สจฺจํ มุสา วาปิ ปรํ ทเหยฺย. \\ เอตญฺจ สลฺลํ ปฏิกจฺจ\textsuperscript{1} ทิสฺวา, \\ อชฺโฌสิตา ยตฺถ ปชา วิสตฺตา. \\ ชานามิ ปสฺสามิ ตเถว เอตํ, \\ อชฺโฌสิตํ นตฺถิ ตถาคตานนฺติ.จตุตฺถํ.}
\csromangathameasurewak{ยํ กิญฺจิ ทิฏฺฐํว สุตํ มุตํ วา, \\ อชฺโฌสิตํ สจฺจมุตํ ปเรสํ. \\ น เตสุ ตาที สยสํวุเตสุ, \\ สจฺจํ มุสา วาปิ ปรํ ทเหยฺย. \\ เอตญฺจ สลฺลํ ปฏิกจฺจ\textsuperscript{1} ทิสฺวา, \\ อชฺโฌสิตา ยตฺถ ปชา วิสตฺตา. \\ ชานามิ ปสฺสามิ ตเถว เอตํ, \\ อชฺโฌสิตํ นตฺถิ ตถาคตานนฺติ.จตุตฺถํ.}
\setlength{\csromangathaleftcol}{0pt}
\csromangathagroup{\csromangathastanza{\csromangathawak{ยํ กิญฺจิ ทิฏฺฐํว สุตํ มุตํ วา, \\ อชฺโฌสิตํ สจฺจมุตํ ปเรสํ. \\ น เตสุ ตาที สยสํวุเตสุ, \\ สจฺจํ มุสา วาปิ ปรํ ทเหยฺย.}}\csromangathastanza{\csromangathawak{เอตญฺจ สลฺลํ ปฏิกจฺจ\csromansharedfootnote{e3bbee4f138f5df0}{ปฏิคจฺจ (สี, อิ)} ทิสฺวา, \\ อชฺโฌสิตา ยตฺถ ปชา วิสตฺตา. \\ ชานามิ ปสฺสามิ ตเถว เอตํ, \\ อชฺโฌสิตํ นตฺถิ ตถาคตานนฺติ.  จตุตฺถํ.}}}

\csromanmatikaanchor{mtk.77}
\csromantocmarkref{subsection}{5. พฺรหฺมจริยสุตฺต}{mtk.77}
\bookmark[dest={mtk.77},level=2]{5. พฺรหฺมจริยสุตฺต}
\csromanheader{2}{5. พฺรหฺมจริย\-สุตฺต}
\proseitem{25}{นยิทํ ภิกฺขเว พฺรหฺมจริยํ วุสฺสติ ชนกุหนตฺถํ, น ชนลปนตฺถํ, น ลาภสกฺการสิโลกานิสํสตฺถํ, น อิติวาทปฺปโมกฺขานิสํสตฺถํ น “อิติ มํ ชโน ชานาตู”ติ. อถ โข อิทํ ภิกฺขเว พฺรหฺมจริยํ วุสฺสติ สํวรตฺถํ ปหานตฺถํ วิราคตฺถํ นิโรธตฺถนฺติ.}

\setlength{\csromangathapagewidth}{0pt}
\setlength{\csromangathawakwidth}{0pt}
\csromangathameasuregroup{สํวรตฺถํ ปหานตฺถํ, \\ อเทสยิ โส ภควา, \\ เอส มคฺโค มหนฺเตหิ\textsuperscript{1}, \\ เย จ ตํ ปฏิปชฺชนฺติ, \\ ทุกฺขสฺสนฺตํ กริสฺสนฺติ,}{\csromangathabat{สํวรตฺถํ ปหานตฺถํ,}{พฺรหฺมจริยํ อนีติหํ.} \\ \csromangathabat{อเทสยิ โส ภควา,}{นิพฺพาโนคธคามินํ.} \\ \csromangathabat{เอส มคฺโค มหนฺเตหิ\textsuperscript{1},}{อนุยาโต มเหสิภิ.} \\ \csromangathabat{เย จ ตํ ปฏิปชฺชนฺติ,}{ยถา พุทฺเธน เทสิตํ.} \\ \csromangathabat{ทุกฺขสฺสนฺตํ กริสฺสนฺติ,}{สตฺถุสาสนการิโนติ.ปญฺจมํ.}}
\csromangathasetleft{สํวรตฺถํ ปหานตฺถํ, \\ อเทสยิ โส ภควา, \\ เอส มคฺโค มหนฺเตหิ\textsuperscript{1}, \\ เย จ ตํ ปฏิปชฺชนฺติ, \\ ทุกฺขสฺสนฺตํ กริสฺสนฺติ,}
\csromangathagroup{\csromangathastanza{\csromangathabat{สํวรตฺถํ ปหานตฺถํ,}{พฺรหฺมจริยํ อนีติหํ.} \\ \csromangathabat{อเทสยิ โส ภควา,}{นิพฺพาโนคธคามินํ.}}\csromangathastanza{\csromangathabat{เอส มคฺโค มหนฺเตหิ\csromansharedfootnote{39c5c5ea0f71351e}{มหตฺเตภิ (ก) ขุ 1. 215 ปิฏฺเฐปิ.},}{อนุยาโต มเหสิภิ.} \\ \csromangathabat{เย จ ตํ ปฏิปชฺชนฺติ,}{ยถา พุทฺเธน เทสิตํ.} \\ \csromangathabat{ทุกฺขสฺสนฺตํ กริสฺสนฺติ,}{สตฺถุสาสนการิโนติ.  ปญฺจมํ.}}}

\csromanmatikaanchor{mtk.78}
\csromantocmarkref{subsection}{6. กุหสุตฺต}{mtk.78}
\bookmark[dest={mtk.78},level=2]{6. กุหสุตฺต}
\csromanmarkh{3. อุรุเวลวคฺค}\csromanmarkhii{6. กุหสุตฺต}\csromanheadernomark{2}{3. อุรุเวลวคฺค \textbf{6. กุหสุตฺต}}
\proseitem{26}{\csromanfolio{335}\csromansymbolfootnote{*}{ขุ 1. 270 อิติวุตฺตเกปิ.}เย เต ภิกฺขเว ภิกฺขู กุหา ถทฺธา ลปา สิงฺคี อุนฺนฬา อสมาหิตา, น เม เต ภิกฺขเว ภิกฺขู มามกา, อปคตา จ เต ภิกฺขเว ภิกฺขู อิมสฺมา ธมฺมวินยา, น จ เต อิมสฺมิํ ธมฺมวินเย วุทฺธิํ วิรุฬฺหิํ เวปุลฺลํ อาปชฺชนฺติ. เย จ โข เต ภิกฺขเว ภิกฺขู นิกฺกุหา นิลฺลปา ธีรา อตฺถทฺธา สุสมาหิตา, เต โข เม ภิกฺขเว ภิกฺขู มามกา, อนปคตา จ เต ภิกฺขเว ภิกฺขู อิมสฺมา ธมฺมวินยา, เต จ อิมสฺมิํ ธมฺมวินเย วุทฺธิํ วิรุฬฺหิํ เวปุลฺลํ อาปชฺชนฺตีติ.}

\prose{\csromansymbolmark{*}กุหา ถทฺธา ลปา สิงฺคี, อุนฺนฬา อสมาหิตา.}

\setlength{\csromangathapagewidth}{0pt}
\setlength{\csromangathawakwidth}{0pt}
\csromangathameasuregroup{น เต ธมฺเม วิรูหนฺติ, \\ นิกฺกุหา นิลฺลปา ธีรา, \\ เต เว ธมฺเม วิรูหนฺติ,}{\csromangathabat{น เต ธมฺเม วิรูหนฺติ,}{สมฺมาสมฺพุทฺธเทสิเต.} \\ \csromangathabat{นิกฺกุหา นิลฺลปา ธีรา,}{อตฺถทฺธา สุสมาหิตา.} \\ \csromangathabat{เต เว ธมฺเม วิรูหนฺติ,}{สมฺมาสมฺพุทฺธเทสิเตติ.ฉฏฺฐํ.}}
\csromangathasetleft{น เต ธมฺเม วิรูหนฺติ, \\ นิกฺกุหา นิลฺลปา ธีรา, \\ เต เว ธมฺเม วิรูหนฺติ,}
\csromangathagroup{\csromangathastanza{\csromangathabat{น เต ธมฺเม วิรูหนฺติ,}{สมฺมาสมฺพุทฺธเทสิเต.} \\ \csromangathabat{นิกฺกุหา นิลฺลปา ธีรา,}{อตฺถทฺธา สุสมาหิตา.} \\ \csromangathabat{เต เว ธมฺเม วิรูหนฺติ,}{สมฺมาสมฺพุทฺธเทสิเตติ.  ฉฏฺฐํ.}}}

\csromanmatikaanchor{mtk.79}
\csromantocmarkref{subsection}{7. สนฺตุฏฺฐิสุตฺต}{mtk.79}
\bookmark[dest={mtk.79},level=2]{7. สนฺตุฏฺฐิสุตฺต}
\csromanheader{2}{7. \textbf{สนฺตุฏฺฐิสุตฺต}}
\proseitem{27}{จตฺตาริมานิ ภิกฺขเว อปฺปานิ จ\csromansharedfootnote{9924f48f668beae7}{เจว (ก)} สุลภานิ จ, ตานิ จ อนวชฺชานิ. กตมานิ จตฺตาริ? ปํสุกูลํ ภิกฺขเว จีวรานํ อปฺปญฺจ สุลภญฺจ, ตญฺจ อนวชฺชํ. ปิณฺฑิยาโลโป ภิกฺขเว โภชนานํ อปฺปญฺจ สุลภญฺจ, ตญฺจ อนวชฺชํ. รุกฺขมูลํ ภิกฺขเว เสนาสนานํ อปฺปญฺจ สุลภญฺจ, ตญฺจ อนวชฺชํ. ปูติมุตฺตํ ภิกฺขเว เภสชฺชานํ อปฺปญฺจ สุลภญฺจ, ตญฺจ อนวชฺชํ. อิมานิ โข ภิกฺขเว จตฺตาริ อปฺปานิ จ สุลภานิ จ, ตานิ จ อนวชฺชานิ. ยโต โข ภิกฺขเว ภิกฺขุ อปฺเปน จ ตุฏฺโฐ โหติ สุลเภน จ, อิทมสฺสาหํ “อญฺญตรํ สามญฺญงฺคนฺ”ติ\csromansharedfootnote{564f4302e60a46a1}{สามญฺญนฺติ (ก)} วทามีติ.}

\setlength{\csromangathapagewidth}{0pt}
\setlength{\csromangathawakwidth}{0pt}
\csromangathameasuregroup{อนวชฺเชน ตุฏฺฐสฺส, \\ น เสนาสนมารพฺภ, \\ วิฆาโต โหติ จิตฺตสฺส, \\ เย จสฺส ธมฺมา อกฺขาตา, \\ อธิคฺคหิตา ตุฏฺฐสฺส,}{\csromangathabat{อนวชฺเชน ตุฏฺฐสฺส,}{อปฺเปน สุลเภน จ.} \\ \csromangathabat{น เสนาสนมารพฺภ,}{จีวรํ ปานโภชนํ.} \\ \csromangathabat{วิฆาโต โหติ จิตฺตสฺส,}{ทิสา นปฺปฏิหญฺญติ.} \\ \csromangathabat{เย จสฺส ธมฺมา อกฺขาตา,}{สามญฺญสฺสานุโลมิกา.} \\ \csromangathabat{อธิคฺคหิตา ตุฏฺฐสฺส,}{อปฺปมตฺตสฺส สิกฺขโตติ.สตฺตมํ.}}
\csromangathasetleft{อนวชฺเชน ตุฏฺฐสฺส, \\ น เสนาสนมารพฺภ, \\ วิฆาโต โหติ จิตฺตสฺส, \\ เย จสฺส ธมฺมา อกฺขาตา, \\ อธิคฺคหิตา ตุฏฺฐสฺส,}
\csromangathagroup{\csromangathastanza{\csromangathabat{อนวชฺเชน ตุฏฺฐสฺส,}{อปฺเปน สุลเภน จ.} \\ \csromangathabat{น เสนาสนมารพฺภ,}{จีวรํ ปานโภชนํ.}}\csromangathastanza{\csromangathabat{วิฆาโต โหติ จิตฺตสฺส,}{ทิสา นปฺปฏิหญฺญติ.} \\ \csromangathabat{เย จสฺส ธมฺมา อกฺขาตา,}{สามญฺญสฺสานุโลมิกา.} \\ \csromangathabat{อธิคฺคหิตา ตุฏฺฐสฺส,}{อปฺปมตฺตสฺส สิกฺขโตติ.  สตฺตมํ.}}}

\csromanpalirecto
\csromanmatikaanchor{mtk.80}
\csromantocmarkref{subsection}{8. อริยวํสสุตฺต}{mtk.80}
\bookmark[dest={mtk.80},level=2]{8. อริยวํสสุตฺต}
\csromanmarkh{จตุกฺกนิปาต\-ปาฬิ}\csromanmarkhii{8. อริยวํส\-สุตฺต}\csromanheadernomark{2}{จตุกฺกนิปาต\-ปาฬิ 8. อริยวํส\-สุตฺต}
\proseitem{28}{\csromanfolio{336}จตฺตาโรเม ภิกฺขเว อริยวํสา อคฺคญฺญา รตฺตญฺญา วํสญฺญา โปราณา, อสํกิณฺณา อสํกิณฺณปุพฺพา น สํกียนฺติ น สํกียิสฺสนฺติ, อปฺปฏิกุฏฺฐา สมเณหิ พฺราหฺมเณหิ วิญฺญูหิ. กตเม จตฺตาโร? อิธ ภิกฺขเว ภิกฺขุ สนฺตุฏฺโฐ โหติ อิตรีตเรน จีวเรน, อิตรีตร\-จีวร\-สนฺตุฏฺฐิยา จ วณฺณวาที, น จ จีวรเหตุ อเนสนํ อปฺปติรูปํ อาปชฺชติ, อลทฺธา จ จีวรํ น ปริตสฺสติ, ลทฺธา จ จีวรํ อคธิโต\csromansharedfootnote{239a5e27e89e47ff}{อคถิโต (สี, อิ)} อมุจฺฉิโต อนชฺโฌสนฺโน อาทีนวทสฺสาวี นิสฺสรณปญฺโญ ปริภุญฺชติ, ตาย จ ปน อิตรีตร\-จีวร\-สนฺตุฏฺฐิยา เนว\-ตฺตานุกฺกํเสติ โน\csromansharedfootnote{f595c2b0487e9391}{น (ที 3. 188 ปิฏฺเฐ)} ปรํ วมฺเภติ. โย หิ ตตฺถ ทกฺโข อนลโส สมฺปชาโน ปติสฺสโต, อยํ วุจฺจติ ภิกฺขเว ภิกฺขุ โปราเณ อคฺคญฺเญ อริยวํเส ฐิโต.}

\prose{ปุน จปรํ ภิกฺขเว ภิกฺขุ สนฺตุฏฺโฐ โหติ อิตรีตเรน ปิณฺฑปาเตน, อิตรีตร\-ปิณฺฑปาต\-สนฺตุฏฺฐิยา จ วณฺณวาที, น จ ปิณฺฑปาตเหตุ อเนสนํ อปฺปติรูปํ อาปชฺชติ, อลทฺธา จ ปิณฺฑปาตํ น ปริตสฺสติ, ลทฺธา จ ปิณฺฑปาตํ อคธิโต อมุจฺฉิโต อนชฺโฌสนฺโน อาทีนวทสฺสาวี นิสฺสรณปญฺโญ ปริภุญฺชติ, ตาย จ ปน อิตรีตร\-ปิณฺฑปาต\-สนฺตุฏฺฐิยา เนว\-ตฺตานุกฺกํเสติ โน ปรํ วมฺเภติ. โย หิ ตตฺถ ทกฺโข อนลโส สมฺปชาโน ปติสฺสโต, อยํ วุจฺจติ ภิกฺขเว ภิกฺขุ โปราเณ อคฺคญฺเญ อริยวํเส ฐิโต.}

\prose{ปุน จปรํ ภิกฺขเว ภิกฺขุ สนฺตุฏฺโฐ โหติ อิตรีตเรน เสนาสเนน, อิตรีตร\-เสนาสน\-สนฺตุฏฺฐิยา จ วณฺณวาที, น จ เสนาสนเหตุ อเนสนํ อปฺปติรูปํ อาปชฺชติ, อลทฺธา จ เสนาสนํ น ปริตสฺสติ, ลทฺธา จ เสนาสนํ อคธิโต อมุจฺฉิโต อนชฺโฌสนฺโน อาทีนวทสฺสาวี นิสฺสรณปญฺโญ ปริภุญฺชติ, ตาย จ ปน อิตรีตร\-เสนาสน\-สนฺตุฏฺฐิยา เนว\-ตฺตานุกฺกํเสติ โน ปรํ วมฺเภติ. โย หิ ตตฺถ ทกฺโข อนลโส สมฺปชาโน ปฏิสฺสโต, อยํ วุจฺจติ ภิกฺขเว ภิกฺขุ โปราเณ อคฺคญฺเญ อริยวํเส ฐิโต.}

\prose{ปุน จปรํ ภิกฺขเว ภิกฺขุ ภาวนาราโม โหติ ภาวนารโต, ปหานาราโม โหติ ปหานรโต, ตาย จ ปน ภาวนา\-รามตาย \csromanfolio{337}ภาวนารติยา ปหานารามตาย ปหานรติยา เนว\-ตฺตานุกฺกํเสติ โน ปรํ วมฺเภติ. โย หิ ตตฺถ ทกฺโข อนลโส สมฺปชาโน ปติสฺสโต, อยํ วุจฺจติ ภิกฺขเว ภิกฺขุ โปราเณ อคฺคญฺเญ อริยวํเส ฐิโต. อิเม โข ภิกฺขเว จตฺตาโร อริยวํสา อคฺคญฺญา รตฺตญฺญา วํสญฺญา โปราณา, อสํกิณฺณา อสํกิณฺณปุพฺพา น สํกียนฺติ น สํกียิสฺสนฺติ, อปฺปฏิกุฏฺฐา สมเณหิ พฺราหฺมเณหิ วิญฺญูหิ.}

\prose{อิเมหิ จ ปน ภิกฺขเว จตูหิ อริยวํเสหิ สมนฺนาคโต ภิกฺขุ ปุรตฺถิมาย เจปิ ทิสาย วิหรติ, เสฺวว อรติํ สหติ, น ตํ อรติ สหติ. ปจฺฉิมาย เจปิ ทิสาย วิหรติ, เสฺวว อรติํ สหติ, น ตํ อรติ สหติ. อุตฺตราย เจปิ ทิสาย วิหรติ, เสฺวว อรติํ สหติ, น ตํ อรติ สหติ. ทกฺขิณาย เจปิ ทิสาย วิหรติ, เสฺวว อรติํ สหติ, น ตํ อรติ สหติ. ตํ กิสฺส เหตุ? อรติรติสโห หิ ภิกฺขเว ธีโรติ.}

\setlength{\csromangathapagewidth}{0pt}
\setlength{\csromangathawakwidth}{0pt}
\csromangathameasuregroup{นารติ สหติ ธีรํ\textsuperscript{1}, \\ ธีโรว อรติํ สหติ, \\ สพฺพกมฺมวิหายีนํ, \\ เนกฺขํ ชมฺโพนทสฺเสว, \\ เทวาปิ นํ ปสํสนฺติ,}{\csromangathabat{นารติ สหติ ธีรํ\textsuperscript{1},}{นารติ ธีรํ สหติ.} \\ \csromangathabat{ธีโรว อรติํ สหติ,}{ธีโร หิ อรติสฺสโห.} \\ \csromangathabat{สพฺพกมฺมวิหายีนํ,}{ปนุณฺณํ\textsuperscript{1} โก นิวารเย.} \\ \csromangathabat{เนกฺขํ ชมฺโพนทสฺเสว,}{โก ตํ นินฺทิตุมรหติ.} \\ \csromangathabat{เทวาปิ นํ ปสํสนฺติ,}{พฺรหฺมุนาปิ ปสํสิโตติ.อฏฺฐมํ.}}
\csromangathasetleft{นารติ สหติ ธีรํ\textsuperscript{1}, \\ ธีโรว อรติํ สหติ, \\ สพฺพกมฺมวิหายีนํ, \\ เนกฺขํ ชมฺโพนทสฺเสว, \\ เทวาปิ นํ ปสํสนฺติ,}
\csromangathagroup{\csromangathastanza{\csromangathabat{นารติ สหติ ธีรํ\csromansharedfootnote{9746a876d8f71e05}{วีรํ (สี)},}{นารติ ธีรํ สหติ.} \\ \csromangathabat{ธีโรว อรติํ สหติ,}{ธีโร หิ อรติสฺสโห.}}\csromangathastanza{\csromangathabat{สพฺพกมฺมวิหายีนํ,}{ปนุณฺณํ\csromansharedfootnote{352fee4c7d839820}{ปณุนฺนํ (?)} โก นิวารเย.} \\ \csromangathabat{เนกฺขํ ชมฺโพนทสฺเสว,}{โก ตํ นินฺทิตุม\-รหติ.} \\ \csromangathabat{เทวาปิ นํ ปสํสนฺติ,}{พฺรหฺมุนาปิ ปสํสิโตติ.  อฏฺฐมํ.}}}

\csromanmatikaanchor{mtk.81}
\csromantocmarkref{subsection}{9. ธมฺมปทสุตฺต}{mtk.81}
\bookmark[dest={mtk.81},level=2]{9. ธมฺมปทสุตฺต}
\csromanheader{2}{9. ธมฺมปท\-สุตฺต}
\proseitem{29}{จตฺตาริมานิ ภิกฺขเว ธมฺมปทานิ อคฺคญฺญานิ รตฺตญฺญานิ วํสญฺญานิ โปราณานิ, อสํกิณฺณานิ อสํกิณฺณ\-ปุพฺพานิ น สํกียนฺติ น สํกียิสฺสนฺติ, อปฺปฏิกุฏฺฐานิ สมเณหิ พฺราหฺมเณหิ วิญฺญูหิ. กตมานิ จตฺตาริ? อนภิชฺฌา ภิกฺขเว ธมฺมปทํ อคฺคญฺญํ รตฺตญฺญํ วํสญฺญํ โปราณํ, อสํกิณฺณํ อสํกิณฺณ\-ปุพฺพํ น สํกียติ น สํกียิสฺสติ, อปฺปฏิกุฏฺฐํ สมเณหิ พฺราหฺมเณหิ วิญฺญูหิ.}

\prose{อพฺยาปาโท ภิกฺขเว ธมฺมปทํ อคฺคญฺญํ รตฺตญฺญํ วํสญฺญํ โปราณํ, อสํกิณฺณํ อสํกิณฺณ\-ปุพฺพํ น สํกียติ น สํกียิสฺสติ, อปฺปฏิกุฏฺฐํ สมเณหิ พฺราหฺมเณหิ วิญฺญูหิ.}

\prose{\csromanfolio{338}สมฺมาสติ ภิกฺขเว ธมฺมปทํ อคฺคญฺญํ รตฺตญฺญํ วํสญฺญํ โปราณํ, อสํกิณฺณํ อสํกิณฺณ\-ปุพฺพํ น สํกียติ น สํกียิสฺสติ, อปฺปฏิกุฏฺฐํ สมเณหิ พฺราหฺมเณหิ วิญฺญูหิ.}

\prose{สมฺมาสมาธิ ภิกฺขเว ธมฺมปทํ อคฺคญฺญํ รตฺตญฺญํ วํสญฺญํ โปราณํ, อสํกิณฺณํ อสํกิณฺณ\-ปุพฺพํ น สํกียติ น สํกียิสฺสติ, อปฺปฏิกุฏฺฐํ สมเณหิ พฺราหฺมเณหิ วิญฺญูหิ. อิมานิ โข ภิกฺขเว จตฺตาริ ธมฺมปทานิ อคฺคญฺญานิ รตฺตญฺญานิ วํสญฺญานิ โปราณานิ, อสํกิณฺณานิ อสํกิณฺณ\-ปุพฺพานิ น สํกียนฺติ น สํกียิสฺสนฺติ, อปฺปฏิกุฏฺฐานิ สมเณหิ พฺราหฺมเณหิ วิญฺญูหีติ.}

\setlength{\csromangathapagewidth}{0pt}
\setlength{\csromangathawakwidth}{0pt}
\csromangathameasuregroup{อนภิชฺฌาลุ วิหเรยฺย, \\ สโต เอกคฺคจิตฺตสฺส\textsuperscript{1},}{\csromangathabat{อนภิชฺฌาลุ วิหเรยฺย,}{อพฺยาปนฺเนน เจตสา.} \\ \csromangathabat{สโต เอกคฺคจิตฺตสฺส\textsuperscript{1},}{อชฺฌตฺตํ สุสมาหิโตติ.นวมํ.}}
\csromangathasetleft{อนภิชฺฌาลุ วิหเรยฺย, \\ สโต เอกคฺคจิตฺตสฺส\textsuperscript{1},}
\csromangathagroup{\csromangathastanza{\csromangathabat{อนภิชฺฌาลุ วิหเรยฺย,}{อพฺยาปนฺเนน เจตสา.} \\ \csromangathabat{สโต เอกคฺคจิตฺตสฺส\csromansharedfootnote{3311e61edee5a9f3}{เอกคฺคจิตฺตายํ (ก)},}{อชฺฌตฺตํ สุสมาหิโตติ.นวมํ.}}}

\csromanmatikaanchor{mtk.82}
\csromantocmarkref{subsection}{10. ปริพฺพาชกสุตฺต}{mtk.82}
\bookmark[dest={mtk.82},level=2]{10. ปริพฺพาชกสุตฺต}
\csromanheader{2}{10. ปริพฺพาชก\-สุตฺต}
\proseitem{30}{เอกํ สมยํ ภควา ราชคเห วิหรติ คิชฺฌกูเฏ ปพฺพเต. เตน โข ปน สมเยน สมฺพหุลา อภิญฺญาตา อภิญฺญาตา ปริพฺพาชกา สิปฺปินิกาตีเร\csromansharedfootnote{8805a6db4ea2597c}{สปฺปินิยา ตีเร (สี, อิ), สิปฺปินิยา ตีเร (สฺยา, กํ), สิปฺปินิยา นทิยา ตีเร (ก)} ปริพฺพาชกา\-ราเม ปฏิวสนฺติ. เสยฺยถิทํ? อนฺนภาโร วรธโร สกุลุทายี จ ปริพฺพาชโก อญฺเญ จ อภิญฺญาตา อภิญฺญาตา ปริพฺพาชกา. อถ โข ภควา สายนฺหสมยํ ปฏิสลฺลานา วุฏฺฐิโต เยน สิปฺปินิกาตีรํ ปริพฺพาชกา\-ราโม เตนุปสงฺกมิ, อุปสงฺกมิตฺวา ปญฺญตฺเต อาสเน นิสีทิ, นิสชฺช โข ภควา เต ปริพฺพาชเก เอตทโวจ\csromandash{}}

\prose{จตฺตาริมานิ ปริพฺพาชกา ธมฺมปทานิ อคฺคญฺญานิ รตฺตญฺญานิ วํสญฺญานิ โปราณานิ, อสํกิณฺณานิ อสํกิณฺณ\-ปุพฺพานิ น สํกียนฺติ น สํกียิสฺสนฺติ, อปฺปฏิกุฏฺฐานิ สมเณหิ พฺราหฺมเณหิ วิญฺญูหิ. กตมานิ จตฺตาริ? อนภิชฺฌา ปริพฺพาชกา ธมฺมปทํ อคฺคญฺญํ รตฺตญฺญํ วํสญฺญํ โปราณํ, อสํกิณฺณํ อสํกิณฺณ\-ปุพฺพํ น สํกียติ น สํกียิสฺสติ, อปฺปฏิกุฏฺฐํ สมเณหิ พฺราหฺมเณหิ วิญฺญูหิ. อพฺยาปาโท ปริพฺพาชกา ธมฺมปทํ ฯเปฯ. สมฺมาสติ ปริพฺพาชกา ธมฺมปทํ ฯเปฯ. สมฺมาสมาธิ ปริพฺพาชกา ธมฺมปทํ อคฺคญฺญํ รตฺตญฺญํ \csromanfolio{339}วํสญฺญํ โปราณํ, อสํกิณฺณํ อสํกิณฺณ\-ปุพฺพํ น สํกียติ น สํกียิสฺสติ, อปฺปฏิกุฏฺฐํ สมเณหิ พฺราหฺมเณหิ วิญฺญูหิ. อิมานิ โข ปริพฺพาชกา จตฺตาริ ธมฺมปทานิ อคฺคญฺญานิ รตฺตญฺญานิ วํสญฺญานิ โปราณานิ, อสํกิณฺณานิ อสํกิณฺณ\-ปุพฺพานิ น สํกียนฺติ น สํกียิสฺสนฺติ, อปฺปฏิกุฏฺฐานิ สมเณหิ พฺราหฺมเณหิ วิญฺญูหิ.}

\prose{โย โข ปริพฺพาชกา เอวํ วเทยฺย “อหเมตํ อนภิชฺฌํ ธมฺมปทํ ปจฺจกฺขาย อภิชฺฌาลุํ กาเมสุ ติพฺพสาราคํ สมณํ วา พฺรหฺมณํ วา ปญฺญาเปสฺสามี”ติ. ตมหํ ตตฺถ เอวํ วเทยฺยํ “เอตุ วทตุ พฺยาหรตุ ปสฺสามิสฺสานุภาวนฺ”ติ. โส วต ปริพฺพาชกา อนภิชฺฌํ ธมฺมปทํ ปจฺจกฺขาย อภิชฺฌาลุํ กาเมสุ ติพฺพสาราคํ สมณํ วา พฺราหฺมณํ วา ปญฺญาเปสฺสตีติ เนตํ ฐานํ วิชฺชติ.}

\prose{โย โข ปริพฺพาชกา เอวํ วเทยฺย “อหเมตํ อพฺยาปาทํ ธมฺมปทํ ปจฺจกฺขาย พฺยาปนฺนจิตฺตํ ปทุฏฺฐ\-มนสงฺกปฺปํ สมณํ วา พฺราหฺมณํ วา ปญฺญาเปสฺสามี”ติ. ตมหํ ตตฺถ เอวํ วเทยฺยํ “เอตุ วทตุ พฺยาหรตุ ปสฺสามิสฺสานุภาวนฺ”ติ. โส วต ปริพฺพาชกา อพฺยาปาทํ ธมฺมปทํ ปจฺจกฺขาย พฺยาปนฺนจิตฺตํ ปทุฏฺฐ\-มนสงฺกปฺปํ สมณํ วา พฺราหฺมณํ วา ปญฺญาเปสฺสตีติ เนตํ ฐานํ วิชฺชติ.}

\prose{โย โข ปริพฺพาชกา เอวํ วเทยฺย “อหเมตํ สมฺมาสติํ ธมฺมปทํ ปจฺจกฺขาย มุฏฺฐสฺสติํ อสมฺปชานํ สมณํ วา พฺราหฺมณํ วา ปญฺญาเปสฺสามี”ติ. ตมหํ ตตฺถ เอวํ วเทยฺยํ “เอตุ วทตุ พฺยาหรตุ ปสฺสามิสฺสานุภาวนฺ”ติ. โส วต ปริพฺพาชกา สมฺมาสติํ ธมฺมปทํ ปจฺจกฺขาย มุฏฺฐสฺสติํ อสมฺปชานํ สมณํ วา พฺราหฺมณํ วา ปญฺญาเปสฺสตีติ เนตํ ฐานํ วิชฺชติ.}

\prose{โย โข ปริพฺพาชกา เอวํ วเทยฺย “อหเมตํ สมฺมาสมาธิํ ธมฺมปทํ ปจฺจกฺขาย อสมาหิตํ วิพฺภนฺต\-จิตฺตํ สมณํ วา พฺราหฺมณํ วา ปญฺญาเปสฺสามี”ติ. ตมหํ ตตฺถ เอวํ วเทยฺยํ “เอตุ วทตุ พฺยาหรตุ ปสฺสามิสฺสานุภาวนฺ”ติ. โส วต ปริพฺพาชกา สมฺมาสมาธิํ ธมฺมปทํ ปจฺจกฺขาย อสมาหิตํ วิพฺภนฺต\-จิตฺตํ สมณํ วา พฺราหฺมณํ วา ปญฺญาเปสฺสตีติ เนตํ ฐานํ วิชฺชติ.}

\prose{โย โข ปริพฺพาชกา อิมานิ จตฺตาริ ธมฺมปทานิ ครหิตพฺพํ ปฏิกฺโกสิตพฺพํ มญฺเญยฺย, ตสฺส ทิฏฺเฐว ธมฺเม จตฺตาโร สหธมฺมิกา วาทานุปาตา \csromanfolio{340}คารยฺหา ฐานา\csromansharedfootnote{a17f7730556037cd}{วาทานุวาทา คารยฺหํ ฐานํ (ม 3. 8 ปิฏฺเฐ)} อาคจฺฉนฺติ. กตเม จตฺตาโร? อนภิชฺฌํ เจ ภวํ ธมฺมปทํ ครหติ ปฏิกฺโกสติ. เย จ หิ\csromansharedfootnote{2732495c79b50a92}{เย จ (ม 3. 121 ปิฏฺเฐ)} อภิชฺฌาลู กาเมสุ ติพฺพสาราคา สมณพฺราหฺมณา, เต โภโต ปุชฺชา, เต โภโต ปาสํสา. อพฺยาปาทํ เจ ภวํ ธมฺมปทํ ครหติ ปฏิกฺโกสติ. เย จ หิ พฺยาปนฺนจิตฺตา ปทุฏฺฐ\-มนสงฺกปฺปา สมณพฺราหฺมณา, เต โภโต ปุชฺชา, เต โภโต ปาสํสา. สมฺมาสติํ เจ ภวํ ธมฺมปทํ ครหติ ปฏิกฺโกสติ. เย จ หิ มุฏฺฐสฺสตี อสมฺปชานา สมณพฺราหฺมณา, เต โภโต ปุชฺชา, เต โภโต ปาสํสา. สมฺมาสมาธิํ เจ ภวํ ธมฺมปทํ ครหติ ปฏิกฺโกสติ. เย จ หิ อสมาหิตา วิพฺภนฺตจิตฺตา สมณพฺราหฺมณา, เต โภโต ปุชฺชา, เต โภโต ปาสํสา.}

\proseitem{3}{โย โข ปริพฺพาชกา อิมานิ จตฺตาริ ธมฺมปทานิ ครหิตพฺพํ ปฏิกฺโกสิตพฺพํ มญฺเญยฺย, ตสฺส ทิฏฺเฐว ธมฺเม อิเม จตฺตาโร สหธมฺมิกา วาทานุปาตา คารยฺหา ฐานา อาคจฺฉนฺติ. เยปิ เต ปริพฺพาชกา อเหสุํ อุกฺกลา วสฺสภญฺญา\csromansharedfootnote{54e1a8e15fa14e05}{วสฺสภิญฺญา (ก) สํ 2. 60 ปิฏฺเฐ ปสฺสิตพฺพํ.} อเหตุกวาทา อกิริยวาทา นตฺถิกวาทา, เตปิ อิมานิ จตฺตาริ ธมฺมปทานิ น ครหิตพฺพํ น ปฏิกฺโกสิตพฺพํ อมญฺญิํสุ. ตํ กิสฺส เหตุ? นินฺทาพฺยาโรสนอุปารมฺภภยาติ\csromansharedfootnote{c446a464861298c7}{อุปวาทภยาติ (ก) ม 3. 121 ปิฏฺเฐ; สํ 2. 60 ปิฏฺเฐ จ ปสฺสิตพฺพํ.}.}

\setlength{\csromangathapagewidth}{0pt}
\setlength{\csromangathawakwidth}{0pt}
\csromangathameasuregroup{อพฺยาปนฺโน สทา สโต, \\ อภิชฺฌาวินเย สิกฺขํ,}{\csromangathabat{อพฺยาปนฺโน สทา สโต,}{อชฺฌตฺตํ สุสมาหิโต.} \\ \csromangathabat{อภิชฺฌาวินเย สิกฺขํ,}{“อปฺปมตฺโต”ติ วุจฺจตีติ.ทสมํ.}}
\csromangathasetleft{อพฺยาปนฺโน สทา สโต, \\ อภิชฺฌาวินเย สิกฺขํ,}
\csromangathagroupwithclosermidruled{\csromangathastanza{\csromangathabat{อพฺยาปนฺโน สทา สโต,}{อชฺฌตฺตํ สุสมาหิโต.} \\ \csromangathabat{อภิชฺฌาวินเย สิกฺขํ,}{“อปฺปมตฺโต”ติ วุจฺจตีติ.  ทสมํ.}}}{อุรุเวลวคฺโค ตติโย.}
\csromancenter{\textbf{ตสฺสุทฺทานํ}}
\setlength{\csromangathapagewidth}{0pt}
\setlength{\csromangathawakwidth}{0pt}
\csromangathameasuregroup{เทฺว อุรุเวลา โลโก กาฬโก\textsuperscript{1}, \\ กุหํ สนฺตุฏฺฐิ วํโส จ,}{\csromangathabat{เทฺว อุรุเวลา โลโก กาฬโก\textsuperscript{1},}{พฺรหฺมจริเยน ปญฺจมํ.} \\ \csromangathabat{กุหํ สนฺตุฏฺฐิ วํโส จ,}{ธมฺมปทํ ปริพฺพาชเกน จาติ.}}
\csromangathasetleft{เทฺว อุรุเวลา โลโก กาฬโก\textsuperscript{1}, \\ กุหํ สนฺตุฏฺฐิ วํโส จ,}
\csromangathagroup{\csromangathastanza{\csromangathabat{เทฺว อุรุเวลา โลโก กาฬโก\csromansharedfootnote{d6da0a3cc164135c}{โกฬิโก (ก)},}{พฺรหฺมจริเยน ปญฺจมํ.} \\ \csromangathabat{กุหํ สนฺตุฏฺฐิ วํโส จ,}{ธมฺมปทํ ปริพฺพาชเกน จาติ.}}}

\csromanmatikaanchor{mtk.83}
\csromantocmarknopage{section}{4. จกฺกวคฺค}{mtk.83}
\bookmark[dest={mtk.83},level=1]{4. จกฺกวคฺค}
\csromanheader{1}{4. จกฺกวคฺค \textbf{4. จกฺกวคฺค}}
\csromanmatikaanchor{mtk.84}
\csromantocmarkref{subsection}{1. จกฺกสุตฺต}{mtk.84}
\bookmark[dest={mtk.84},level=2]{1. จกฺกสุตฺต}
\csromanheader{2}{1. \textbf{จกฺกสุตฺต}}
\proseitem{31}{\csromanfolio{341}จตฺตาริมานิ ภิกฺขเว จกฺกานิ, เยหิ สมนฺนาคตานํ เทวมนุสฺสานํ จตุจกฺกํ วตฺตติ, เยหิ สมนฺนาคตา เทวมนุสฺสา นจิรสฺเสว มหนฺตตฺตํ เวปุลฺลตฺตํ ปาปุณนฺติ โภเคสุ. กตมานิ จตฺตาริ? ปติรูป\-เทส\-วาโส สปฺปุริสา\-วสฺสโย\csromansharedfootnote{f698c75d4878d34f}{สปฺปุริสูปสฺสโย (สี, สฺยา, กํ, อิ)} อตฺต\-สมฺมา\-ปณิธิ ปุพฺเพ จ กตปุญฺญตา. อิมานิ โข ภิกฺขเว จตฺตาริ จกฺกานิ. เยหิ สมนฺนาคตานํ เทวมนุสฺสานํ จตุจกฺกํ วตฺตติ, เยหิ สมนฺนาคตา เทวมนุสฺสา นจิรสฺเสว มหนฺตตฺตํ เวปุลฺลตฺตํ ปาปุณนฺติ โภเคสูติ.}

\setlength{\csromangathapagewidth}{0pt}
\setlength{\csromangathawakwidth}{0pt}
\csromangathameasuregroup{ปติรูเป วเส เทเส, \\ สมฺมาปณิธิสมฺปนฺโน, \\ ธญฺญํ ธนํ ยโส กิตฺติ,}{\csromangathabat{ปติรูเป วเส เทเส,}{อริยมิตฺตกโร สิยา.} \\ \csromangathabat{สมฺมาปณิธิสมฺปนฺโน,}{ปุพฺเพ ปุญฺญกโต นโร.} \\ \csromangathabat{ธญฺญํ ธนํ ยโส กิตฺติ,}{สุขญฺเจตํธิวตฺตตีติ.ปฐมํ.}}
\csromangathasetleft{ปติรูเป วเส เทเส, \\ สมฺมาปณิธิสมฺปนฺโน, \\ ธญฺญํ ธนํ ยโส กิตฺติ,}
\csromangathagroup{\csromangathastanza{\csromangathabat{ปติรูเป วเส เทเส,}{อริยมิตฺตกโร สิยา.} \\ \csromangathabat{สมฺมาปณิธิสมฺปนฺโน,}{ปุพฺเพ ปุญฺญกโต นโร.} \\ \csromangathabat{ธญฺญํ ธนํ ยโส กิตฺติ,}{สุขญฺเจตํธิวตฺตตีติ.  ปฐมํ.}}}

\csromanmatikaanchor{mtk.85}
\csromantocmarkref{subsection}{2. สงฺคหสุตฺต}{mtk.85}
\bookmark[dest={mtk.85},level=2]{2. สงฺคหสุตฺต}
\csromanheader{2}{2. \textbf{สงฺคหสุตฺต}}
\proseitem{32}{จตฺตาริมานิ ภิกฺขเว สงฺคห\-วตฺถูนิ. กตมานิ จตฺตาริ? ทานํ เปยฺยวชฺชํ อตฺถจริยา สมานตฺตตา, อิมานิ โข ภิกฺขเว จตฺตาริ สงฺคหวตฺถูนีติ.}

\setlength{\csromangathapagewidth}{0pt}
\setlength{\csromangathawakwidth}{0pt}
\csromangathameasuregroup{ทานญฺจ เปยฺยวชฺชญฺจ\textsuperscript{1}, \\ สมานตฺตตา จ ธมฺเมสุ, \\ เอเต โข สงฺคหา โลเก, \\ เอเต จ สงฺคหา นาสฺสุ, \\ ลเภถ มานํ ปูชํ วา, \\ ยสฺมา จ สงฺคหา\textsuperscript{1} เอเต, \\ ตสฺมา มหตฺตํ ปปฺโปนฺติ,}{\csromangathabat{ทานญฺจ เปยฺยวชฺชญฺจ\textsuperscript{1},}{อตฺถจริยา จ ยา อิธ.} \\ \csromangathabat{สมานตฺตตา จ ธมฺเมสุ,}{ตตฺถ ตตฺถ ยถารหํ.} \\ \csromangathabat{เอเต โข สงฺคหา โลเก,}{รถสฺสาณีว ยายโต.} \\ \csromangathabat{เอเต จ สงฺคหา นาสฺสุ,}{น มาตา ปุตฺตการณา.} \\ \csromangathabat{ลเภถ มานํ ปูชํ วา,}{ปิตา วา ปุตฺตการณา.} \\ \csromangathabat{ยสฺมา จ สงฺคหา\textsuperscript{1} เอเต,}{สมเวกฺขนฺติ ปณฺฑิตา.} \\ \csromangathabat{ตสฺมา มหตฺตํ ปปฺโปนฺติ,}{ปาสํสา จ ภวนฺติ เตติ.ทุติยํ.}}
\csromangathasetleft{ทานญฺจ เปยฺยวชฺชญฺจ\textsuperscript{1}, \\ สมานตฺตตา จ ธมฺเมสุ, \\ เอเต โข สงฺคหา โลเก, \\ เอเต จ สงฺคหา นาสฺสุ, \\ ลเภถ มานํ ปูชํ วา, \\ ยสฺมา จ สงฺคหา\textsuperscript{1} เอเต, \\ ตสฺมา มหตฺตํ ปปฺโปนฺติ,}
\csromangathagroup{\csromangathastanza{\csromangathabat{ทานญฺจ เปยฺยวชฺชญฺจ\csromansharedfootnote{d7c0a11c4fa8fff7}{สงฺคเห (อฏฺฐกถายํ ปาฐนฺตรํ) อุปริ 571 ปิฏฺเฐ; ที 3. 156 ปิฏฺเฐปิ ปสฺสิตพฺพํ.},}{อตฺถจริยา จ ยา อิธ.} \\ \csromangathabat{สมานตฺตตา จ ธมฺเมสุ,}{ตตฺถ ตตฺถ ยถารหํ.}}\csromangathastanza{\csromangathabat{เอเต โข สงฺคหา โลเก,}{รถสฺสาณีว ยายโต.} \\ \csromangathabat{เอเต จ สงฺคหา นาสฺสุ,}{น มาตา ปุตฺตการณา.}}\csromangathastanza{\csromangathabat{ลเภถ มานํ ปูชํ วา,}{ปิตา วา ปุตฺตการณา.} \\ \csromangathabat{ยสฺมา จ สงฺคหา\csromansharedfootnote{d7c0a11c4fa8fff7}{สงฺคเห (อฏฺฐกถายํ ปาฐนฺตรํ) อุปริ 571 ปิฏฺเฐ; ที 3. 156 ปิฏฺเฐปิ ปสฺสิตพฺพํ.} เอเต,}{สมเวกฺขนฺติ ปณฺฑิตา.} \\ \csromangathabat{ตสฺมา มหตฺตํ ปปฺโปนฺติ,}{ปาสํสา จ ภวนฺติ เตติ.  ทุติยํ.}}}

\csromanpalirecto
\csromanmatikaanchor{mtk.86}
\csromantocmarkref{subsection}{3. สีหสุตฺต}{mtk.86}
\bookmark[dest={mtk.86},level=2]{3. สีหสุตฺต}
\csromanmarkh{จตุกฺกนิปาต\-ปาฬิ}\csromanmarkhii{3. สีหสุตฺต}\csromanheadernomark{2}{จตุกฺกนิปาต\-ปาฬิ 3. สีหสุตฺต}
\proseitem{33}{\csromanfolio{342}สีโห ภิกฺขเว มิคราชา สายนฺหสมยํ อาสยา นิกฺขมติ. อาสยา นิกฺขมิตฺวา วิชมฺภติ, วิชมฺภิตฺวา สมนฺตา จตุทฺทิสา อนุวิโลเกติ, สมนฺตา จตุทฺทิสา อนุวิโลเกตฺวา ติกฺขตฺตุํ สีหนาทํ นทติ, ติกฺขตฺตุํ สีหนาทํ นทิตฺวา โคจราย ปกฺกมติ. เย โข ปน เต ภิกฺขเว ติรจฺฉานคตา ปาณา สีหสฺส มิครญฺโญ นทโต สทฺทํ สุณนฺติ, เต เยภุยฺเยน ภยํ สํเวคํ สนฺตาสํ อาปชฺชนฺติ. พิลํ พิลาสยา ปวิสนฺติ, ทกํ ทกาสยา\csromansharedfootnote{92cc380b51f75213}{อุทกํ อุทกาสยา (ก) สํ 2. 70 ปิฏฺเฐ; ปฏิสํฏฺฐ 1. 187 ปิฏฺเฐปิ ปสฺสิตพฺพํ. 2. สกฺกายสฺส อตฺถงฺคโม (ฏฺฐ, สํ 2. 70 ปิฏฺเฐ.) “สกฺกายนิโรธคามินี”ติ ปจฺฉิมปทํ ปน ปสฺสิตพฺพํ.} ปวิสนฺติ, วนํ วนาสยา ปวิสนฺติ, อากาสํ ปกฺขิโน ภชนฺติ. เยปิ เต ภิกฺขเว รญฺโญ นาคา คามนิคม\-ราชธานีสุ ทเฬฺหหิ วรตฺเตหิ พนฺธเนหิ พทฺธา, เตปิ ตานิ พนฺธนานิ สญฺฉินฺทิตฺวา สมฺปทาเลตฺวา ภีตา มุตฺตกรีสํ จชมานา เยน วา เตน วา ปลายนฺติ. เอวํ มหิทฺธิโก โข ภิกฺขเว สีโห มิคราชา ติรจฺฉาน\-คตานํ ปาณานํ เอวํ มเหสกฺโข เอวํ มหานุภาโว.}

\prose{เอวเมวํ โข ภิกฺขเว ยทา ตถาคโต โลเก อุปฺปชฺชติ อรหํ สมฺมาสมฺพุทฺโธ วิชฺชา\-จรณ\-สมฺปนฺโน สุคโต โลกวิทู อนุตฺตโร ปุริส\-ทมฺม\-สารถิ สตฺถา เทวมนุสฺสานํ พุทฺโธ ภควา, โส ธมฺมํ เทเสติ “อิติ สกฺกาโย, อิติ สกฺกาย\-สมุทโย, อิติ สกฺกายนิโรโธ\csromansharedfootnote{4c0329110f7ae2ac}{สกฺกายสฺส อตฺถงฺคโม (ฏฺฐ, สํ 2. 70 ปิฏฺเฐ.) “สกฺกายนิโรธคามินี”ติ ปจฺฉิมปทํ ปน ปสฺสิตพฺพํ.}, อิติ สกฺกาย\-นิโรธ\-คามินี ปฏิปทา”ติ. เยปิ เต ภิกฺขเว เทวา ทีฆายุกา วณฺณวนฺโต สุขพหุลา อุจฺเจสุ วิมาเนสุ จิรฏฺฐิติกา, เตปิ ตถาคตสฺส ธมฺมเทสนํ สุตฺวา เยภุยฺเยน ภยํ สํเวคํ สนฺตาสํ อาปชฺชนฺติ, อนิจฺจา วต กิร โภ มยํ สมานา “นิจฺจมฺหา”ติ อมญฺญิมฺห, อทฺธุวา วต กิร โภ มยํ สมานา “ธุวมฺหา”ติ อมญฺญิมฺห, อสสฺสตา วต กิร โภ มยํ สมานา “สสฺสตมฺหา”ติ อมญฺญิมฺห, มยํ กิร โภ อนิจฺจา อทฺธุวา อสสฺสตา สกฺกาย\-ปริยาปนฺนาติ. เอวํ มหิทฺธิโก โข ภิกฺขเว ตถาคโต สเทวกสฺส โลกสฺส เอวํ มเหสกฺโข เอวํ มหานุภาโวติ.}

\setlength{\csromangathapagewidth}{0pt}
\setlength{\csromangathawakwidth}{0pt}
\csromangathameasuregroup{ยทา พุทฺโธ อภิญฺญาย, \\ สเทวกสฺส โลกสฺส,}{\csromangathabat{ยทา พุทฺโธ อภิญฺญาย,}{ธมฺมจกฺกํ ปวตฺตยี.} \\ \csromangathabat{สเทวกสฺส โลกสฺส,}{สตฺถา อปฺปฏิปุคฺคโล.}}
\csromangathasetleft{ยทา พุทฺโธ อภิญฺญาย, \\ สเทวกสฺส โลกสฺส,}
\csromangathagroup{\csromangathastanza{\csromangathabat{ยทา พุทฺโธ อภิญฺญาย,}{ธมฺมจกฺกํ ปวตฺตยี.} \\ \csromangathabat{สเทวกสฺส โลกสฺส,}{สตฺถา อปฺปฏิปุคฺคโล.}}}

\chapterheadpageread{4. จกฺกวคฺค}
\setlength{\csromangathapagewidth}{0pt}
\setlength{\csromangathawakwidth}{0pt}
\csromangathameasuregroup{สกฺกายญฺจ นิโรธญฺจ, \\ อริยญฺจฏฺฐงฺคิกํ มคฺคํ, \\ เยปิ ทีฆายุกา เทวา, \\ ภีตา สนฺตาส มาปาทุํ, \\ อวีติวตฺตา สกฺกายํ, \\ สุตฺวา อรหโต วากฺยํ,}{\csromangathabat{สกฺกายญฺจ นิโรธญฺจ,}{สกฺกายสฺส จ สมฺภวํ.} \\ \csromangathabat{อริยญฺจฏฺฐงฺคิกํ มคฺคํ,}{ทุกฺขูปสมคามินํ.} \\ \csromangathabat{เยปิ ทีฆายุกา เทวา,}{วณฺณวนฺโต ยสสฺสิโน.} \\ \csromangathabat{ภีตา สนฺตาส มาปาทุํ,}{สีหสฺเสวิ’ตเรมิคา.} \\ \csromangathabat{อวีติวตฺตา สกฺกายํ,}{อนิจฺจา กิร โภ มยํ.} \\ \csromangathabat{สุตฺวา อรหโต วากฺยํ,}{วิปฺปมุตฺตสฺส ตาทิโนติ.* ตติยํ.}}
\csromangathasetleft{สกฺกายญฺจ นิโรธญฺจ, \\ อริยญฺจฏฺฐงฺคิกํ มคฺคํ, \\ เยปิ ทีฆายุกา เทวา, \\ ภีตา สนฺตาส มาปาทุํ, \\ อวีติวตฺตา สกฺกายํ, \\ สุตฺวา อรหโต วากฺยํ,}
\csromangathagroup{\csromangathastanza{\csromanfolio{343}\csromangathabat{สกฺกายญฺจ นิโรธญฺจ,}{สกฺกายสฺส จ สมฺภวํ.} \\ \csromangathabat{อริยญฺจฏฺฐงฺคิกํ มคฺคํ,}{ทุกฺขู\-ปสม\-คามินํ.}}\csromangathastanza{\csromangathabat{เยปิ ทีฆายุกา เทวา,}{วณฺณวนฺโต ยสสฺสิโน.} \\ \csromangathabat{ภีตา สนฺตาส มาปาทุํ,}{สีหสฺเสวิ’ตเรมิคา.}}\csromangathastanza{\csromangathabat{อวีติวตฺตา สกฺกายํ,}{อนิจฺจา กิร โภ มยํ.} \\ \csromangathabat{สุตฺวา อรหโต วากฺยํ,}{วิปฺปมุตฺตสฺส ตาทิโนติ.* ตติยํ.}}}

\csromanmatikaanchor{mtk.87}
\csromantocmarkref{subsection}{4. ปสาทสุตฺต}{mtk.87}
\bookmark[dest={mtk.87},level=2]{4. ปสาทสุตฺต}
\csromanheader{2}{4. อคฺคปฺปสาท\-สุตฺต}
\proseitem{34}{จตฺตาโรเม ภิกฺขเว อคฺคปฺปสาทา. กตเม จตฺตาโร? ยาวตา ภิกฺขเว สตฺตา อปทา วา ทฺวิปทา\csromansharedfootnote{66e48bc1f7e08436}{ทิปทา (สี, อิ) อํ 2. 29; ขุ 1. 255 ปิฏฺเฐสุปิ.} วา จตุปฺปทา วา พหุปฺปทา วา รูปิโน วา อรูปิโน วา สญฺญิโน วา อสญฺญิโน วา เนว\-สญฺญิ\-นา\-สญฺญิโน วา, ตถาคโต เตสํ อคฺคมกฺขายติ อรหํ สมฺมาสมฺพุทฺโธ. เย ภิกฺขเว พุทฺเธ ปสนฺนา, อคฺเค เต ปสนฺนา, อคฺเค โข ปน ปสนฺนานํ อคฺโค วิปาโก โหติ.}

\prose{ยาวตา ภิกฺขเว ธมฺมา สงฺขตา, อริโย อฏฺฐงฺคิโก มคฺโค เตสํ อคฺคมกฺขายติ. เย ภิกฺขเว อริเย อฏฺฐงฺคิเก มคฺเค ปสนฺนา, อคฺเค เต ปสนฺนา, อคฺเค โข ปน ปสนฺนานํ อคฺโค วิปาโก โหติ.}

\prose{ยาวตา ภิกฺขเว ธมฺมา สงฺขตา วา อสงฺขตา วา, วิราโค เตสํ อคฺคมกฺขายติ, ยทิทํ มทนิมฺมทโน ปิปาสวินโย อาลย\-สมุคฺฆาโต วฏฺฏุปจฺเฉโท ตณฺหากฺขโย วิราโค นิโรโธ นิพฺพานํ. เย ภิกฺขเว วิราเค ธมฺเม ปสนฺนา, อคฺเค เต ปสนฺนา, อคฺเค โข ปน ปสนฺนานํ อคฺโค วิปาโก โหติ.}

\prose{ยาวตา ภิกฺขเว สํฆา วา คณา วา, ตถาคต\-สาวกสํโฆ เตสํ อคฺคมกฺขายติ, ยทิทํ จตฺตาริ ปุริสยุคานิ อฏฺฐ ปุริสปุคฺคลา, เอส ภควโต สาวกสํโฆ อาหุเนยฺโย ปาหุเนยฺโย ทกฺขิเณยฺโย อญฺชลิกรณีโย อนุตฺตรํ ปุญฺญกฺเขตฺตํ โลกสฺส. เย ภิกฺขเว สํเฆ \csromanfolio{344}ปสนฺนา, อคฺเค เต ปสนฺนา, อคฺเค โข ปน ปสนฺนานํ อคฺโค วิปาโก โหติ. อิเม โข ภิกฺขเว จตฺตาโร อคฺคปฺปสาทาติ.}

\setlength{\csromangathapagewidth}{0pt}
\setlength{\csromangathawakwidth}{0pt}
\csromangathameasuregroup{อคฺคโต เว ปสนฺนานํ, \\ อคฺเค พุทฺเธ ปสนฺนานํ, \\ อคฺเค ธมฺเม ปสนฺนานํ, \\ อคฺเค สํเฆ ปสนฺนานํ, \\ อคฺคสฺมิํ ทานํ ททตํ, \\ อคฺคํ อายุ จ วณฺโณ จ, \\ อคฺคสฺส ทาตา เมธาวี, \\ เทวภูโต มนุสฺโส วา,}{\csromangathabat{อคฺคโต เว ปสนฺนานํ,}{อคฺคํ ธมฺมํ วิชานตํ.} \\ \csromangathabat{อคฺเค พุทฺเธ ปสนฺนานํ,}{ทกฺขิเณยฺเย อนุตฺตเร.} \\ \csromangathabat{อคฺเค ธมฺเม ปสนฺนานํ,}{วิราคูปสเม สุเข.} \\ \csromangathabat{อคฺเค สํเฆ ปสนฺนานํ,}{ปุญฺญกฺเขตฺเต อนุตฺตเร.} \\ \csromangathabat{อคฺคสฺมิํ ทานํ ททตํ,}{อคฺคํ ปุญฺญํ ปวฑฺฒติ.} \\ \csromangathabat{อคฺคํ อายุ จ วณฺโณ จ,}{ยโส กิตฺติ สุขํ พลํ.} \\ \csromangathabat{อคฺคสฺส ทาตา เมธาวี,}{อคฺคธมฺมสมาหิโต.} \\ \csromangathabat{เทวภูโต มนุสฺโส วา,}{อคฺคปฺปตฺโต ปโมทตีติ\textsuperscript{1}.} \\ จตุตฺถํ.}
\csromangathameasurewak{จตุตฺถํ.}
\csromangathasetleft{อคฺคโต เว ปสนฺนานํ, \\ อคฺเค พุทฺเธ ปสนฺนานํ, \\ อคฺเค ธมฺเม ปสนฺนานํ, \\ อคฺเค สํเฆ ปสนฺนานํ, \\ อคฺคสฺมิํ ทานํ ททตํ, \\ อคฺคํ อายุ จ วณฺโณ จ, \\ อคฺคสฺส ทาตา เมธาวี, \\ เทวภูโต มนุสฺโส วา,}
\csromangathagroup{\csromangathastanza{\csromangathabat{อคฺคโต เว ปสนฺนานํ,}{อคฺคํ ธมฺมํ วิชานตํ.} \\ \csromangathabat{อคฺเค พุทฺเธ ปสนฺนานํ,}{ทกฺขิเณยฺเย อนุตฺตเร.}}\csromangathastanza{\csromangathabat{อคฺเค ธมฺเม ปสนฺนานํ,}{วิราคูปสเม สุเข.} \\ \csromangathabat{อคฺเค สํเฆ ปสนฺนานํ,}{ปุญฺญกฺเขตฺเต อนุตฺตเร.}}\csromangathastanza{\csromangathabat{อคฺคสฺมิํ ทานํ ททตํ,}{อคฺคํ ปุญฺญํ ปวฑฺฒติ.} \\ \csromangathabat{อคฺคํ อายุ จ วณฺโณ จ,}{ยโส กิตฺติ สุขํ พลํ.}}\csromangathastanza{\csromangathabat{อคฺคสฺส ทาตา เมธาวี,}{อคฺคธมฺมสมาหิโต.} \\ \csromangathabat{เทวภูโต มนุสฺโส วา,}{อคฺคปฺปตฺโต ปโมทตีติ\csromansharedfootnote{338dd9d64fd7ee1e}{[มิสฺสินฺคฺ โนเต 0]}.}}\csromangathastanza{\csromangathawak{จตุตฺถํ.}}}

\csromanmatikaanchor{mtk.88}
\csromantocmarkref{subsection}{5. วสฺสการสุตฺต}{mtk.88}
\bookmark[dest={mtk.88},level=2]{5. วสฺสการสุตฺต}
\csromanheader{2}{5. \textbf{วสฺสการสุตฺต}}
\proseitem{35}{เอกํ สมยํ ภควา ราชคเห วิหรติ เวฬุวเน กลนฺทก\-นิวาเป. อถ โข วสฺสกาโร พฺราหฺมโณ มคธ\-มหามตฺโต เยน ภควา เตนุปสงฺกมิ, อุปสงฺกมิตฺวา ภควตา สทฺธิํ สมฺโมทิ, สมฺโมทนียํ กถํ สารณียํ วีติสาเรตฺวา เอกมนฺตํ นิสีทิ, เอกมนฺตํ นิสินฺโน โข วสฺสกาโร พฺราหฺมโณ มคธ\-มหามตฺโต ภควนฺตํ เอตทโวจ\csromandash{}}

\prose{จตูหิ โข มยํ โภ โคตม ธมฺเมหิ สมนฺนาคตํ มหาปญฺญํ มหาปุริสํ ปญฺญาเปม. กตเมหิ จตูหิ? อิธ โภ โคตม พหุสฺสุโต โหติ ตสฺส ตสฺเสว สุตชาตสฺส. ตสฺส ตสฺเสว โข ปน ภาสิตสฺส อตฺถํ ชานาติ “อยํ อิมสฺส ภาสิตสฺส อตฺโถ, อยํ อิมสฺส ภาสิตสฺส อตฺโถ”ติ. สติมา โข ปน โหติ จิรกตมฺปิ จิรภาสิตมฺปิ สริตา อนุสฺสริตา. ยานิ โข ปน ตานิ คหฏฺฐกานิ กิํกรณียานิ, ตตฺถ ทกฺโข โหติ อนลโส, ตตฺรุปายาย วีมํสาย สมนฺนาคโต อลํ กาตุํ อลํ สํวิธาตุํ. อิเมหิ โข มยํ โภ โคตม จตูหิ ธมฺเมหิ สมนฺนาคตํ มหาปญฺญํ มหาปุริสํ ปญฺญาเปม. สเจ เม\csromansharedfootnote{475636355d3ce276}{สเจ ปน เม (สี, ก), สเจ เม ปน (อิ)} โภ โคตม อนุโมทิตพฺพํ อนุโมทตุ เม ภวํ โคตโม. \csromanfolio{345}สเจ ปน เม โภ โคตม ปฏิกฺโกสิตพฺพํ ปฏิกฺโกสตุ เม ภวํ โคตโมติ.}

\prose{เนว โข ตฺยาหํ พฺราหฺมณ อนุโมทามิ น ปฏิกฺโกสามิ. จตูหิ โข อหํ พฺราหฺมณ ธมฺเมหิ สมนฺนาคตํ มหาปญฺญํ มหาปุริสํ ปญฺญาเปมิ. กตเมหิ จตูหิ? อิธ พฺราหฺมณ พหุชนหิตาย ปฏิปนฺโน โหติ พหุชน\-สุขาย, พหุ’สฺส ชนตา อริเย ญาเย ปติฏฺฐาปิตา, ยทิทํ กลฺยาณ\-ธมฺมตา กุสลธมฺมตา. โส ยํ วิตกฺกํ อากงฺขติ วิตกฺเกตุํ, ตํ วิตกฺกํ วิตกฺเกติ. ยํ วิตกฺกํ นากงฺขติ วิตกฺเกตุํ, น ตํ วิตกฺกํ วิตกฺเกติ. ยํ สงฺกปฺปํ อากงฺขติ สงฺกปฺเปตุํ, ตํ สงฺกปฺปํ สงฺกปฺเปติ. ยํ สงฺกปฺปํ นากงฺขติ สงฺกปฺเปตุํ, น ตํ สงฺกปฺปํ สงฺกปฺเปติ. อิติ เจโตวสิปฺปตฺโต โหติ วิตกฺกปเถ. จตุนฺนํ ฌานานํ อาภิเจตสิกานํ ทิฏฺฐ\-ธมฺม\-สุข\-วิหารานํ นิกามลาภี โหติ อกิจฺฉลาภี อกสิรลาภี. อาสวานํ ขยา อนาสวํ เจโตวิมุตฺติํ ปญฺญาวิมุตฺติํ ทิฏฺเฐว ธมฺเม สยํ อภิญฺญา สจฺฉิกตฺวา อุปสมฺปชฺช วิหรติ. เนว โข ตฺยาหํ พฺราหฺมณ อนุโมทามิ น ปน ปฏิกฺโกสามิ. อิเมหิ โข อหํ พฺราหฺมณ จตูหิ ธมฺเมหิ สมนฺนาคตํ มหาปญฺญํ มหาปุริสํ ปญฺญาเปมีติ.}

\prose{อจฺฉริยํ โภ โคตม, อพฺภุตํ โภ โคตม, ยาว สุภาสิตํ จิทํ โภตา โคตเมน. อิเมหิ จ มยํ โภ โคตม จตูหิ ธมฺเมหิ สมนฺนาคตํ ภวนฺตํ โคตมํ ธาเรม “ภวํ หิ โคตโม พหุชนหิตาย ปฏิปนฺโน พหุชน\-สุขาย, พหุ เต\csromansharedfootnote{1386e36d0a59648a}{พหุสฺส (สฺยา, กํ, ก)} ชนตา อริเย ญาเย ปติฏฺฐาปิตา, ยทิทํ กลฺยาณ\-ธมฺมตา กุสลธมฺมตา. ภวํ หิ โคตโม ยํ วิตกฺกํ อากงฺขติ วิตกฺเกตุํ, ตํ วิตกฺกํ วิตกฺเกติ. ยํ วิตกฺกํ นากงฺขติ วิตกฺเกตุํ, น ตํ วิตกฺกํ วิตกฺเกติ. ยํ สงฺกปฺปํ อากงฺขติ สงฺกปฺเปตุํ, ตํ สงฺกปฺปํ สงฺกปฺเปติ. ยํ สงฺกปฺปํ นากงฺขติ สงฺกปฺเปตุํ, น ตํ สงฺกปฺปํ สงฺกปฺเปติ. ภวํ หิ โคตโม เจโตวสิปฺปตฺโต วิตกฺกปเถ. ภวํ หิ โคตโม จตุนฺนํ ฌานานํ อาภิเจตสิกานํ ทิฏฺฐ\-ธมฺม\-สุข\-วิหารานํ นิกามลาภี อกิจฺฉลาภี อกสิรลาภี. ภวํ หิ โคตโม อาสวานํ ขยา อนาสวํ เจโตวิมุตฺติํ ปญฺญาวิมุตฺติํ ทิฏฺเฐว ธมฺเม สยํ อภิญฺญา สจฺฉิกตฺวา อุปสมฺปชฺช วิหรตี”ติ.}

\prose{\csromanfolio{346}อทฺธา โข ตฺยายํ พฺราหฺมณ อาสชฺช อุปนีย วาจา ภาสิตา, อปิ จ ตฺยาหํ พฺยากริสฺสามิ “อหํ หิ พฺราหฺมณ พหุชนหิตาย ปฏิปนฺโน พหุชน\-สุขาย, พหุ เม\csromansharedfootnote{1386e36d0a59648a}{พหุสฺส (สฺยา, กํ, ก)} ชนตา อริเย ญาเย ปติฏฺฐาปิตา, ยทิทํ กลฺยาณ\-ธมฺมตา กุสลธมฺมตา. อหํ หิ พฺราหฺมณ ยํ วิตกฺกํ อากงฺขามิ วิตกฺเกตุํ, ตํ วิตกฺกํ วิตกฺเกมิ. ยํ วิตกฺกํ นากงฺขามิ วิตกฺเกตุํ, น ตํ วิตกฺกํ วิตกฺเกมิ. ยํ สงฺกปฺปํ อากงฺขามิ สงฺกปฺเปตุํ, ตํ สงฺกปฺปํ สงฺกปฺเปมิ. ยํ สงฺกปฺปํ นากงฺขามิ สงฺกปฺเปตุํ, น ตํ สงฺกปฺปํ สงฺกปฺเปมิ. อหํ หิ พฺราหฺมณ เจโตวสิปฺปตฺโต วิตกฺกปเถ. อหํ หิ พฺราหฺมณ จตุนฺนํ ฌานานํ อาภิเจตสิกานํ ทิฏฺฐ\-ธมฺม\-สุข\-วิหารานํ นิกามลาภี อกิจฺฉลาภี อกสิรลาภี. อหํ หิ พฺราหฺมณ อาสวานํ ขยา อนาสวํ เจโตวิมุตฺติํ ปญฺญาวิมุตฺติํ ทิฏฺเฐว ธมฺเม สยํ อภิญฺญา สจฺฉิกตฺวา อุปสมฺปชฺช วิหรามี”ติ.}

\setlength{\csromangathapagewidth}{0pt}
\setlength{\csromangathawakwidth}{0pt}
\csromangathameasuregroup{โย เวทิ สพฺพสตฺตานํ, \\ หิตํ เทวมนุสฺสานํ, \\ ยํ เว ทิสฺวา จ สุตฺวา จ, \\ มคฺคามคฺคสฺส กุสโล,}{\csromangathabat{โย เวทิ สพฺพสตฺตานํ,}{มจฺจุปาสปฺปโมจนํ.} \\ \csromangathabat{หิตํ เทวมนุสฺสานํ,}{ญายํ ธมฺมํ ปกาสยิ.} \\ \csromangathabat{ยํ เว ทิสฺวา จ สุตฺวา จ,}{ปสีทนฺติ พหู ชนา\textsuperscript{1}.} \\ \csromangathabat{มคฺคามคฺคสฺส กุสโล,}{กตกิจฺโจ อนาสโว.} \\ พุทฺโธ อนฺติมสารีโร\textsuperscript{1},}
\csromangathameasurewak{พุทฺโธ อนฺติมสารีโร\textsuperscript{1},}
\csromangathasetleft{โย เวทิ สพฺพสตฺตานํ, \\ หิตํ เทวมนุสฺสานํ, \\ ยํ เว ทิสฺวา จ สุตฺวา จ, \\ มคฺคามคฺคสฺส กุสโล,}
\csromangathagroup{\csromangathastanza{\csromangathabat{โย เวทิ สพฺพสตฺตานํ,}{มจฺจุปาส\-ปฺปโมจนํ.} \\ \csromangathabat{หิตํ เทวมนุสฺสานํ,}{ญายํ ธมฺมํ ปกาสยิ.}}\csromangathastanza{\csromangathabat{ยํ เว ทิสฺวา จ สุตฺวา จ,}{ปสีทนฺติ พหู ชนา\csromansharedfootnote{9ba1bf3249016bbe}{ปสีทติ พหุชฺชโน (สี, สฺยา, กํ, อิ)}.} \\ \csromangathabat{มคฺคามคฺคสฺส กุสโล,}{กตกิจฺโจ อนาสโว.}}\csromangathastanza{\csromangathawak{พุทฺโธ อนฺติมสารีโร\csromansharedfootnote{4d3698b6b3a80822}{อนฺติมธาริโต (ก)},}}}

\prose{“(มหาปญฺโญ)\csromansharedfootnote{08762a6eccd768ed}{( ) สฺยาโปตฺถเก นตฺถิ.} มหาปุริโส”ติ วุจฺจตีติ. ปญฺจมํ.}

\csromanmatikaanchor{mtk.89}
\csromantocmarkref{subsection}{6. โทณสุตฺต}{mtk.89}
\bookmark[dest={mtk.89},level=2]{6. โทณสุตฺต}
\csromanheader{2}{6. \textbf{โทณสุตฺต}}
\proseitem{36}{เอกํ สมยํ ภควา อนฺตรา จ อุกฺกฏฺฐํ, อนฺตรา จ เสตพฺยํ อทฺธาน\-มคฺค\-ปฺปฏิปนฺโน โหติ. โทโณปิ สุทํ พฺราหฺมโณ อนฺตรา จ อุกฺกฏฺฐํ, อนฺตรา จ เสตพฺยํ อทฺธาน\-มคฺค\-ปฺปฏิปนฺโน โหติ. อทฺทสา โข โทโณ พฺราหฺมโณ ภควโต ปาเทสุ จกฺกานิ สหสฺสรานิ สเนมิกานิ สนาภิกานิ สพฺพา\-การ\-ปริปูรานิ, ทิสฺวานสฺส เอตทโหสิ “อจฺฉริยํ วต โภ, อพฺภุตํ วต โภ, น วติมานิ มนุสฺส\-ภูตสฺส ปทานิ ภวิสฺสนฺตี”ติ. อถ โข ภควา มคฺคา โอกฺกมฺม อญฺญตรสฺมิํ รุกฺขมูเล \csromanfolio{347}นิสีทิ ปลฺลงฺกิํ อาภุชิตฺวา อุชุํ กายํ ปณิธาย ปริมุขํ สติํ อุปฏฺฐเปตฺวา. อถ โข โทโณ พฺราหฺมโณ ภควโต ปทานิ อนุคจฺฉนฺโต อทฺทส ภควนฺตํ อญฺญตรสฺมิํ รุกฺขมูเล นิสินฺนํ ปาสาทิกํ ปสาทนียํ สนฺตินฺทฺริยํ สนฺตมานสํ อุตฺตม\-ทมถ\-สมถม\-นุปฺปตฺตํ ทนฺตํ คุตฺตํ สํยตินฺทฺริยํ\csromansharedfootnote{691652dbd4b9ec12}{ยตินฺทฺริยํ (วิ 3. 285 ปิฏฺเฐ)} นาคํ, ทิสฺวาน เยน ภควา เตนุปสงฺกมิ, อุปสงฺกมิตฺวา ภควนฺตํ เอตทโวจ\csromandash{}}

\prose{เทโว โน ภวํ ภวิสฺสตีติ. น โข อหํ พฺราหฺมณ เทโว ภวิสฺสามีติ. คนฺธพฺโพ โน ภวํ ภวิสฺสตีติ. น โข อหํ พฺราหฺมณ คนฺธพฺโพ ภวิสฺสามีติ. ยกฺโข โน ภวํ ภวิสฺสตีติ. น โข อหํ พฺราหฺมณ ยกฺโข ภวิสฺสามีติ. มนุสฺโส โน ภวํ ภวิสฺสตีติ. น โข อหํ พฺราหฺมณ มนุสฺโส ภวิสฺสามีติ.}

\prose{“เทโว โน ภวํ ภวิสฺสตี”ติ อิติ ปุฏฺโฐ สมาโน “น โข อหํ พฺราหฺมณ เทโว ภวิสฺสามี”ติ วเทสิ. “คนฺธพฺโพ โน ภวํ ภวิสฺสตี”ติ อิติ ปุฏฺโฐ สมาโน “น โข อหํ พฺราหฺมณ คนฺธพฺโพ ภวิสฺสามี”ติ วเทสิ. “ยกฺโข โน ภวํ ภวิสฺสตี”ติ อิติ ปุฏฺโฐ สมาโน “น โข อหํ พฺราหฺมณ ยกฺโข ภวิสฺสามี”ติ วเทสิ. “มนุสฺโส โน ภวํ ภวิสฺสตี”ติ อิติ ปุฏฺโฐ สมาโน “น โข อหํ พฺราหฺมณ มนุสฺโส ภวิสฺสามี”ติ วเทสิ. อถ โก จรหิ ภวํ ภวิสฺสตีติ.}

\prose{เยสํ โข อหํ พฺราหฺมณ อาสวานํ อปฺปหีนตฺตา เทโว ภเวยฺยํ, เต เม อาสวา ปหีนา อุจฺฉินฺนมูลา ตาลาวตฺถุกตา อนภาวํกตา อายติํ อนุปฺปาทธมฺมา. เยสํ โข อหํ พฺราหฺมณ อาสวานํ อปฺปหีนตฺตา คนฺธพฺโพ ภเวยฺยํ, ยกฺโข ภเวยฺยํ, มนุสฺโส ภเวยฺยํ, เต เม อาสวา ปหีนา อุจฺฉินฺนมูลา ตาลาวตฺถุกตา อนภาวํกตา อายติํ อนุปฺปาทธมฺมา. เสยฺยถาปิ พฺราหฺมณ อุปฺปลํ วา ปทุมํ วา ปุณฺฑรีกํ วา อุทเก ชาตํ อุทเก สํวฑฺฒํ อุทกา อจฺจุคฺคมฺม ติฏฺฐติ อนุปลิตฺตํ อุทเกน. เอวเมวํ โข อหํ พฺราหฺมณ โลเก ชาโต โลเก สํวฑฺโฒ โลกํ อภิภุยฺย วิหรามิ อนุปลิตฺโต โลเกน. “พุทฺโธ”ติ มํ พฺราหฺมณ ธาเรหีติ.}

\setlength{\csromangathapagewidth}{0pt}
\setlength{\csromangathawakwidth}{0pt}
\csromangathameasuregroup{เยน เทวูปปตฺยสฺส, \\ ยกฺขตฺตํ เยน คจฺเฉยฺยํ, \\ เต มยฺหํ อาสวา ขีณา,}{\csromangathabat{เยน เทวูปปตฺยสฺส,}{คนฺธพฺโพ วา วิหงฺคโม.} \\ \csromangathabat{ยกฺขตฺตํ เยน คจฺเฉยฺยํ,}{มนุสฺสตฺตญฺจ อพฺพเช.} \\ \csromangathabat{เต มยฺหํ อาสวา ขีณา,}{วิทฺธสฺตา วินฬีกตา.}}
\csromangathasetleft{เยน เทวูปปตฺยสฺส, \\ ยกฺขตฺตํ เยน คจฺเฉยฺยํ, \\ เต มยฺหํ อาสวา ขีณา,}
\csromangathagroup{\csromangathastanza{\csromanfolio{348}\csromangathabat{เยน เทวูปปตฺยสฺส,}{คนฺธพฺโพ วา วิหงฺคโม.} \\ \csromangathabat{ยกฺขตฺตํ เยน คจฺเฉยฺยํ,}{มนุสฺสตฺตญฺจ อพฺพเช.} \\ \csromangathabat{เต มยฺหํ อาสวา ขีณา,}{วิทฺธสฺตา วินฬีกตา.}}}

\prose{ปุณฺฑรีกํ ยถา วคฺคุ โตเยน นุปลิปฺปติ\csromansharedfootnote{c2c3d2aa861136c6}{นุปลิมฺปติ (ก)}.}

\setlength{\csromangathapagewidth}{0pt}
\setlength{\csromangathawakwidth}{0pt}
\csromangathameasuregroup{นุปลิปฺปามิ\textsuperscript{1} โลเกน,}{\csromangathabat{นุปลิปฺปามิ\textsuperscript{1} โลเกน,}{ตสฺมา พุทฺโธสฺมิ พฺราหฺมณาติ.} \\ ฉฏฺฐํ.}
\csromangathasetleft{นุปลิปฺปามิ\textsuperscript{1} โลเกน,}
\csromangathagroup{\csromangathastanza{\csromangathabat{นุปลิปฺปามิ\csromansharedfootnote{338dd9d64fd7ee1e}{[มิสฺสินฺคฺ โนเต 0]} โลเกน,}{ตสฺมา พุทฺโธสฺมิ พฺราหฺมณาติ.} \\ ฉฏฺฐํ.}}

\csromanmatikaanchor{mtk.90}
\csromantocmarkref{subsection}{7. อปริหานิยสุตฺต}{mtk.90}
\bookmark[dest={mtk.90},level=2]{7. อปริหานิยสุตฺต}
\csromanheader{2}{7. อปริหานิย\-สุตฺต}
\proseitem{37}{จตูหิ ภิกฺขเว ธมฺเมหิ สมนฺนาคโต ภิกฺขุ อภพฺโพ ปริหานาย นิพฺพานสฺเสว สนฺติเก. กตเมหิ จตูหิ? อิธ ภิกฺขเว ภิกฺขุ สีลสมฺปนฺโน โหติ, อินฺทฺริเยสุ คุตฺตทฺวาโร โหติ, โภชเน มตฺตญฺญู โหติ, ชาคริยํ อนุยุตฺโต โหติ.}

\prose{กถญฺจ ภิกฺขเว ภิกฺขุ สีลสมฺปนฺโน โหติ? อิธ ภิกฺขเว ภิกฺขุ สีลวา โหติ, ปาติโมกฺข\-สํวร\-สํวุโต วิหรติ, อาจารโคจร\-สมฺปนฺโน อณุมตฺเตสุ วชฺเชสุ ภยทสฺสาวี สมาทาย สิกฺขติ สิกฺขาปเทสุ. เอวํ โข ภิกฺขเว ภิกฺขุ สีลสมฺปนฺโน โหติ.}

\prose{กถญฺจ ภิกฺขเว ภิกฺขุ อินฺทฺริเยสุ คุตฺตทฺวาโร โหติ? อิธ ภิกฺขเว ภิกฺขุ จกฺขุนา รูปํ ทิสฺวา น นิมิตฺตคฺคาหี โหติ นานุพฺยญฺชน\-คฺคาหี, ยตฺวาธิกรณเมนํ จกฺขุนฺทฺริยํ อสํวุตํ วิหรนฺตํ อภิชฺฌา\-โทมนสฺสา ปาปกา อกุสลา ธมฺมา อนฺวาสฺสเวยฺยุํ, ตสฺส สํวราย ปฏิปชฺชติ, รกฺขติ จกฺขุนฺทฺริยํ, จกฺขุนฺทฺริเย สํวรํ อาปชฺชติ. โสเตน สทฺทํ สุตฺวา. ฆาเนน คนฺธํ ฆายิตฺวา. ชิวฺหาย รสํ สายิตฺวา. กาเยน โผฏฺฐพฺพํ ผุสิตฺวา. มนสา ธมฺมํ วิญฺญาย น นิมิตฺตคฺคาหี โหติ นานุพฺยญฺชน\-คฺคาหี, ยตฺวาธิกรณเมนํ มนินฺทฺริยํ อสํวุตํ วิหรนฺตํ อภิชฺฌา\-โทมนสฺสา ปาปกา อกุสลา ธมฺมา อนฺวาสฺสเวยฺยุํ, ตสฺส สํวราย ปฏิปชฺชติ, รกฺขติ มนินฺทฺริยํ, มนินฺทฺริเย สํวรํ อาปชฺชติ. เอวํ โข ภิกฺขเว ภิกฺขุ อินฺทฺริเยสุ คุตฺตทฺวาโร โหติ.}

\chapterheadpageread{4. จกฺกวคฺค}
\prose{\csromanfolio{349}กถญฺจ ภิกฺขเว ภิกฺขุ โภชเน มตฺตญฺญู โหติ? อิธ ภิกฺขเว ภิกฺขุ ปฏิสงฺขา โยนิโส อาหารํ อาหาเรติ, เนว ทวาย น มทาย น มณฺฑนาย น วิภูสนาย, ยาวเทว อิมสฺส กายสฺส ฐิติยา ยาปนาย วิหิํสู\-ปรติยา พฺรหฺมจริยา\-นุคฺคหาย, อิติ ปุราณญฺจ เวทนํ ปฏิหงฺขามิ, นวญฺจ เวทนํ น อุปฺปาเทสฺสามิ, ยาตฺรา จ เม ภวิสฺสติ, อนวชฺชตา จ ผาสุวิหาโร จาติ. เอวํ โข ภิกฺขเว ภิกฺขุ โภชเน มตฺตญฺญู โหติ. กถญฺจ ภิกฺขเว ภิกฺขุ ชาคริยํ อนุยุตฺโต โหติ? อิธ ภิกฺขเว ภิกฺขุ ทิวสํ จงฺกเมน นิสชฺชาย อาวรณีเยหิ ธมฺเมหิ จิตฺตํ ปริโสเธติ, รตฺติยา ปฐมํ ยามํ จงฺกเมน นิสชฺชาย อาวรณีเยหิ ธมฺเมหิ จิตฺตํ ปริโสเธติ, รตฺติยา มชฺฌิมํ ยามํ ทกฺขิเณน ปสฺเสน สีหเสยฺยํ กปฺเปติ ปาเท ปาทํ อจฺจาธาย สโต สมฺปชาโน อุฏฺฐานสญฺญํ มนสิ กริตฺวา, รตฺติยา ปจฺฉิมํ ยามํ ปจฺจุฏฺฐาย จงฺกเมน นิสชฺชาย อาวรณีเยหิ ธมฺเมหิ จิตฺตํ ปริโสเธติ. เอวํ โข ภิกฺขเว ภิกฺขุ ชาคริยํ อนุยุตฺโต โหติ. อิเมหิ โข ภิกฺขเว จตูหิ ธมฺเมหิ สมนฺนาคโต ภิกฺขุ อภพฺโพ ปริหานาย นิพฺพานสฺเสว สนฺติเกติ.}

\setlength{\csromangathapagewidth}{0pt}
\setlength{\csromangathawakwidth}{0pt}
\csromangathameasuregroup{สีเล ปติฏฺฐิโต ภิกฺขุ, \\ โภชนมฺหิ จ มตฺตญฺญู, \\ เอวํ วิหารี อาตาปี, \\ ภาวยํ กุสลํ ธมฺมํ, \\ อปฺปมาทรโต ภิกฺขุ, \\ อภพฺโพ ปริหานาย,}{\csromangathabat{สีเล ปติฏฺฐิโต ภิกฺขุ,}{อินฺทฺริเยสุ จ สํวุโต.} \\ \csromangathabat{โภชนมฺหิ จ มตฺตญฺญู,}{ชาคริยํ อนุยุญฺชติ.} \\ \csromangathabat{เอวํ วิหารี อาตาปี,}{อโหรตฺตมตนฺทิโต.} \\ \csromangathabat{ภาวยํ กุสลํ ธมฺมํ,}{โยคกฺเขมสฺส ปตฺติยา.} \\ \csromangathabat{อปฺปมาทรโต ภิกฺขุ,}{ปมาเท ภยทสฺสิ วา\textsuperscript{1}.} \\ \csromangathabat{อภพฺโพ ปริหานาย,}{นิพฺพานสฺเสว สนฺติเกติ.}}
\csromangathasetleft{สีเล ปติฏฺฐิโต ภิกฺขุ, \\ โภชนมฺหิ จ มตฺตญฺญู, \\ เอวํ วิหารี อาตาปี, \\ ภาวยํ กุสลํ ธมฺมํ, \\ อปฺปมาทรโต ภิกฺขุ, \\ อภพฺโพ ปริหานาย,}
\csromangathagroup{\csromangathastanza{\csromangathabat{สีเล ปติฏฺฐิโต ภิกฺขุ,}{อินฺทฺริเยสุ จ สํวุโต.} \\ \csromangathabat{โภชนมฺหิ จ มตฺตญฺญู,}{ชาคริยํ อนุยุญฺชติ.}}\csromangathastanza{\csromangathabat{เอวํ วิหารี อาตาปี,}{อโหรตฺตมตนฺทิโต.} \\ \csromangathabat{ภาวยํ กุสลํ ธมฺมํ,}{โยคกฺเขมสฺส ปตฺติยา.}}\csromangathastanza{\csromangathabat{อปฺปมาทรโต ภิกฺขุ,}{ปมาเท ภยทสฺสิ วา\csromansharedfootnote{977d7843544a10d1}{ภยทสฺสาวี (ก) ขุ 1. 17 ธมฺมปเทปิ.}.} \\ \csromangathabat{อภพฺโพ ปริหานาย,}{นิพฺพานสฺเสว สนฺติเกติ.}}}

\csromanmatikaanchor{mtk.91}
\csromantocmarkref{subsection}{8. ปติลีนสุตฺต}{mtk.91}
\bookmark[dest={mtk.91},level=2]{8. ปติลีนสุตฺต}
\csromanheader{2}{8. \textbf{ปติลีนสุตฺต}}
\proseitem{38}{ปนุณฺณ\-ปจฺเจกสจฺโจ\csromansharedfootnote{277b528de6af1648}{ปณุนฺนปจฺเจกสจฺโจ (?)} ภิกฺขเว ภิกฺขุ สมวย\-สฏฺเฐ\-สโน ปสฺสทฺธ\-กายสงฺขาโร “ปติลีโน”ติ วุจฺจติ. กถญฺจ ภิกฺขเว ภิกฺขุ ปนุณฺณ\-ปจฺเจกสจฺโจ โหติ? อิธ ภิกฺขเว ภิกฺขุโน ยานิ ตานิ \csromanfolio{350}ปุถุสมณ\-พฺราหฺมณานํ ปุถุปจฺเจกสจฺจานิ. เสยฺยถิทํ? “สสฺสโต โลโก”ติ วา, “อสสฺสโต โลโก”ติ วา, “อนฺตวา โลโก”ติ วา, “อนนฺตวา โลโก”ติ วา, “ตํ ชีวํ ตํ สรีรนฺ”ติ วา, “อญฺญํ ชีวํ อญฺญํ สรีรนฺ”ติ วา, “โหติ ตถาคโต ปรํ มรณา”ติ วา, “น โหติ ตถาคโต ปรํ มรณา”ติ วา, “โหติ จ น จ โหติ ตถาคโต ปรํ มรณา”ติ วา, “เนว โหติ น น โหติ ตถาคโต ปรํ มรณา”ติ วา, สพฺพานิ ตานิ นุณฺณานิ โหนฺติ ปนุณฺณานิ จตฺตานิ วนฺตานิ มุตฺตานิ ปหีนานิ ปฏินิสฺสฏฺฐานิ. เอวํ โข ภิกฺขเว ภิกฺขุ ปนุณฺณ\-ปจฺเจกสจฺโจ โหติ.}

\prose{กถญฺจ ภิกฺขเว ภิกฺขุ สมวย\-สฏฺเฐ\-สโน โหติ? อิธ ภิกฺขเว ภิกฺขุโน กาเมสนา ปหีนา โหติ, ภเวสนา ปหีนา โหติ, พฺรหฺมจริเย\-สนา ปฏิปฺปสฺสทฺธา. เอวํ โข ภิกฺขเว ภิกฺขุ สมวย\-สฏฺเฐ\-สโน โหติ.}

\prose{กถญฺจ ภิกฺขเว ภิกฺขุ ปสฺสทฺธ\-กายสงฺขาโร โหติ? อิธ ภิกฺขเว ภิกฺขุ สุขสฺส จ ปหานา ทุกฺขสฺส จ ปหานา ปุพฺเพว โสมนสฺส\-โทมนสฺสานํ อตฺถงฺคมา อทุกฺขมสุขํ อุเปกฺขา\-สติ\-ปาริสุทฺธิํ จตุตฺถํ ฌานํ อุปสมฺปชฺช วิหรติ. เอวํ โข ภิกฺขเว ภิกฺขุ ปสฺสทฺธ\-กายสงฺขาโร โหติ.}

\prose{กถญฺจ ภิกฺขเว ภิกฺขุ ปติลีโน โหติ? อิธ ภิกฺขเว ภิกฺขุโน อสฺมิมาโน ปหีโน โหติ อุจฺฉินฺนมูโล ตาลาวตฺถุกโต อนภาวํกโต อายติํ อนุปฺปาทธมฺโม. เอวํ โข ภิกฺขเว ภิกฺขุ ปติลีโน โหติ. ปนุณฺณ\-ปจฺเจกสจฺโจ ภิกฺขเว ภิกฺขุ สมวย\-สฏฺเฐ\-สโน ปสฺสทฺธ\-กายสงฺขาโร “ปติลีโน”ติ วุจฺจตีติ.}

\setlength{\csromangathapagewidth}{0pt}
\setlength{\csromangathawakwidth}{0pt}
\csromangathameasuregroup{\textsuperscript{*}กาเมสนา ภเวสนา, \\ อิติ สจฺจปรามาโส, \\ \textsuperscript{*}สพฺพราควิรตฺตสฺส, \\ เอสนา ปฏินิสฺสฏฺฐา, \\ ส เว สนฺโต สโต ภิกฺขุ,}{\csromangathabat{\textsuperscript{*}กาเมสนา ภเวสนา,}{พฺรหฺมจริเยสนา สห.} \\ \csromangathabat{อิติ สจฺจปรามาโส,}{ทิฏฺฐิฏฺฐานา สมุสฺสยา.} \\ \csromangathabat{\textsuperscript{*}สพฺพราควิรตฺตสฺส,}{ตณฺหกฺขยวิมุตฺติโน.} \\ \csromangathabat{เอสนา ปฏินิสฺสฏฺฐา,}{ทิฏฺฐิฏฺฐานา สมูหตา.} \\ \csromangathabat{ส เว สนฺโต สโต ภิกฺขุ,}{ปสฺสทฺโธ อปราชิโต.} \\ อฏฺฐมํ.}
\csromangathameasurewak{อฏฺฐมํ.}
\csromangathasetleft{\textsuperscript{*}กาเมสนา ภเวสนา, \\ อิติ สจฺจปรามาโส, \\ \textsuperscript{*}สพฺพราควิรตฺตสฺส, \\ เอสนา ปฏินิสฺสฏฺฐา, \\ ส เว สนฺโต สโต ภิกฺขุ,}
\csromangathagroup{\csromangathastanza{\csromangathabat{\csromansymbolfootnote{*}{ขุ 1. 228 ปิฏฺเฐ อิติวุตฺตเกปิ.}กาเมสนา ภเวสนา,}{พฺรหฺมจริเย\-สนา สห.} \\ \csromangathabat{อิติ สจฺจปรามาโส,}{ทิฏฺฐิฏฺฐานา สมุสฺสยา.}}\csromangathastanza{\csromangathabat{\csromansymbolmark{*}สพฺพราควิรตฺตสฺส,}{ตณฺหกฺขยวิมุตฺติโน.} \\ \csromangathabat{เอสนา ปฏินิสฺสฏฺฐา,}{ทิฏฺฐิฏฺฐานา สมูหตา.} \\ \csromangathabat{ส เว สนฺโต สโต ภิกฺขุ,}{ปสฺสทฺโธ อปราชิโต.}}\csromangathastanza{\csromangathawak{อฏฺฐมํ.}}}

\prose{มานาภิสมยา พุทฺโธ, “ปติลีโน”ติ วุจฺจตีติ.}

\csromanmatikaanchor{mtk.92}
\csromantocmarkref{subsection}{9. อุชฺชยสุตฺต}{mtk.92}
\bookmark[dest={mtk.92},level=2]{9. อุชฺชยสุตฺต}
\csromanmarkh{4. จกฺกวคฺค}\csromanmarkhii{9. อุชฺชยสุตฺต}\csromanheadernomark{2}{4. จกฺกวคฺค \textbf{9. อุชฺชยสุตฺต}}
\proseitem{39}{\csromanfolio{351}อถ โข อุชฺชโย พฺราหฺมโณ เยน ภควา เตนุปสงฺกมิ, อุปสงฺกมิตฺวา ภควตา สทฺธิํ สมฺโมทิ, สมฺโมทนียํ กถํ สารณียํ วีติสาเรตฺวา เอกมนฺตํ นิสีทิ, เอกมนฺตํ นิสินฺโน โข อุชฺชโย พฺราหฺมโณ ภควนฺตํ เอตทโวจ “ภวมฺปิ โน โคตโม ยญฺญํ วณฺเณตี”ติ. น โข อหํ พฺราหฺมณ สพฺพํ ยญฺญํ วณฺเณมิ, น ปนาหํ พฺราหฺมณ สพฺพํ ยญฺญํ น วณฺเณมิ. ยถารูเป โข พฺราหฺมณ ยญฺเญ คาโว หญฺญนฺติ, อเชฬกา หญฺญนฺติ, กุกฺกุฏสูกรา หญฺญนฺติ, วิวิธา ปาณา สํฆาตํ อาปชฺชนฺติ. เอวรูปํ โข อหํ พฺราหฺมณ สารมฺภํ ยญฺญํ น วณฺเณมิ. ตํ กิสฺส เหตุ? เอวรูปํ หิ พฺราหฺมณ สารมฺภํ ยญฺญํ น อุปสงฺกมนฺติ อรหนฺโต วา อรหตฺตมคฺคํ วา สมาปนฺนา.}

\prose{ยถารูเป จ โข พฺราหฺมณ ยญฺเญ เนว คาโว หญฺญนฺติ, น อเชฬกา หญฺญนฺติ, น กุกฺกุฏสูกรา หญฺญนฺติ, น วิวิธา ปาณา สํฆาตํ อาปชฺชนฺติ. เอวรูปํ โข อหํ พฺราหฺมณ นิรารมฺภํ ยญฺญํ วณฺเณมิ, ยทิทํ นิจฺจทานํ อนุกุลยญฺญํ. ตํ กิสฺส เหตุ? เอวรูปํ หิ พฺราหฺมณ นิรารมฺภํ ยญฺญํ อุปสงฺกมนฺติ อรหนฺโต วา อรหตฺตมคฺคํ วา สมาปนฺนาติ.}

\prose{อสฺสเมธํ ปุริสเมธํ, สมฺมาปาสํ วาชเปยฺยํ นิรคฺคฬํ.}

\prose{มหายญฺญา มหารมฺภา\csromansharedfootnote{5f73c858c7aa46f2}{สมฺมาปาสํ วาชเปยฺยํ. นิรคฺคฬํ มหารมฺภา (อิ) สํ 1. 76 ปิฏฺเฐปิ.}, น เต โหนฺติ มหปฺผลา.}

\prose{อเชฬกา จ คาโว จ, วิวิธา ยตฺถ หญฺญเร.}

\prose{น ตํ สมฺมคฺคตา ยญฺญํ, อุปยนฺติ มเหสิโน.}

\prose{เย จ ยญฺญา นิรารมฺภา, ยชนฺติ อนุกุลํ สทา.}

\prose{อเชฬกา จ คาโว จ, วิวิธา เนตฺถ หญฺญเร\csromansharedfootnote{554c23b9e301a736}{นาเชฬกา จ คาโว จ, วิวิธา ยตฺถ หญฺญเร (สฺยา, กํ)}.}

\prose{ตญฺจ สมฺมคฺคตา ยญฺญํ, อุปยนฺติ มเหสิโน.}

\setlength{\csromangathapagewidth}{0pt}
\setlength{\csromangathawakwidth}{0pt}
\csromangathameasuregroup{เอตํ\textsuperscript{1} ยเชถ เมธาวี, \\ เอตํ\textsuperscript{1} หิ ยชมานสฺส, \\ ยญฺโญ จ วิปุโล โหติ,}{\csromangathabat{เอตํ\textsuperscript{1} ยเชถ เมธาวี,}{เอโส ยญฺโญ มหปฺผโล.} \\ \csromangathabat{เอตํ\textsuperscript{1} หิ ยชมานสฺส,}{เสยฺโย โหติ น ปาปิโย.} \\ \csromangathabat{ยญฺโญ จ วิปุโล โหติ,}{ปสีทนฺติ จ เทวตาติ.} \\ นวมํ.}
\csromangathasetleft{เอตํ\textsuperscript{1} ยเชถ เมธาวี, \\ เอตํ\textsuperscript{1} หิ ยชมานสฺส, \\ ยญฺโญ จ วิปุโล โหติ,}
\csromangathagroup{\csromangathastanza{\csromanfolio{352}\csromangathabat{เอตํ\csromansharedfootnote{338dd9d64fd7ee1e}{[มิสฺสินฺคฺ โนเต 0]} ยเชถ เมธาวี,}{เอโส ยญฺโญ มหปฺผโล.} \\ \csromangathabat{เอตํ\csromansharedfootnote{338dd9d64fd7ee1e}{[มิสฺสินฺคฺ โนเต 0]} หิ ยชมานสฺส,}{เสยฺโย โหติ น ปาปิโย.}}\csromangathastanza{\csromangathabat{ยญฺโญ จ วิปุโล โหติ,}{ปสีทนฺติ จ เทวตาติ.} \\ นวมํ.}}

\csromanmatikaanchor{mtk.93}
\csromantocmarkref{subsection}{10. อุทายีสุตฺต}{mtk.93}
\bookmark[dest={mtk.93},level=2]{10. อุทายีสุตฺต}
\csromanheader{2}{10. \textbf{อุทายีสุตฺต}}
\proseitem{40}{อถ โข อุทายี\csromansharedfootnote{252cb2284ffbf5aa}{อุทายิ (สพฺพตฺถ) 4. วิวตฺตจฺฉทา (สี, อิ), วิวฏฺฏจฺฉทา (ก)} พฺราหฺมโณ เยน ภควา เตนุปสงฺกมิ, อุปสงฺกมิตฺวา ภควตา ฯเปฯ เอกมนฺตํ นิสินฺโน โข อุทายี พฺราหฺมโณ ภควนฺตํ เอตทโวจ “ภวมฺปิ โน โคตโม ยญฺญํ วณฺเณตี”ติ. น โข อหํ พฺราหฺมณ สพฺพํ ยญฺญํ วณฺเณมิ, น ปนาหํ พฺราหฺมณ สพฺพํ ยญฺญํ น วณฺเณมิ. ยถารูเป โข พฺราหฺมณ ยญฺเญ คาโว หญฺญนฺติ, อเชฬกา หญฺญนฺติ, กุกฺกุฏสูกรา หญฺญนฺติ, วิวิธา ปาณา สํฆาตํ อาปชฺชนฺติ. เอวรูปํ โข อหํ พฺราหฺมณ สารมฺภํ ยญฺญํ น วณฺเณมิ. ตํ กิสฺส เหตุ? เอวรูปํ หิ พฺราหฺมณ สารมฺภํ ยญฺญํ น อุปสงฺกมนฺติ อรหนฺโต วา อรหตฺตมคฺคํ วา สมาปนฺนา.}

\prose{ยถารูเป จ โข พฺราหฺมณ ยญฺเญ เนว คาโว หญฺญนฺติ, น อเชฬกา หญฺญนฺติ, น กุกฺกุฏสูกรา หญฺญนฺติ, น วิวิธา ปาณา สํฆาตํ อาปชฺชนฺติ. เอวรูปํ โข อหํ พฺราหฺมณ นิรารมฺภํ ยญฺญํ วณฺเณมิ, ยทิทํ นิจฺจทานํ อนุกุลยญฺญํ. ตํ กิสฺส เหตุ? เอวรูปํ หิ พฺราหฺมณ นิรารมฺภํ ยญฺญํ อุปสงฺกมนฺติ อรหนฺโต วา อรหตฺตมคฺคํ วา สมาปนฺนาติ.}

\setlength{\csromangathapagewidth}{0pt}
\setlength{\csromangathawakwidth}{0pt}
\csromangathameasuregroup{อภิสงฺขตํ นิรารมฺภํ, \\ ตาทิสํ อุปสํยนฺติ, \\ วิวฏจฺฉทา\textsuperscript{1} เย โลเก, \\ ยญฺญเมตํ ปสํสนฺติ, \\ ยญฺเญ วา ยทิ วา สทฺเธ, \\ ปสนฺนจิตฺโต ยชติ\textsuperscript{1},}{\csromangathabat{อภิสงฺขตํ นิรารมฺภํ,}{ยญฺญํ กาเลน กปฺปิยํ.} \\ \csromangathabat{ตาทิสํ อุปสํยนฺติ,}{สญฺญตา พฺรหฺมจารโย.} \\ \csromangathabat{วิวฏจฺฉทา\textsuperscript{1} เย โลเก,}{วีติวตฺตา กุลํ คติํ.} \\ \csromangathabat{ยญฺญเมตํ ปสํสนฺติ,}{พุทฺธา ยญฺญสฺส\textsuperscript{1} โกวิทา.} \\ \csromangathabat{ยญฺเญ วา ยทิ วา สทฺเธ,}{หพฺยํ\textsuperscript{1} กตฺวา ยถารหํ.} \\ \csromangathabat{ปสนฺนจิตฺโต ยชติ\textsuperscript{1},}{สุเขตฺเต พฺรหฺมจาริสุ.}}
\csromangathasetleft{อภิสงฺขตํ นิรารมฺภํ, \\ ตาทิสํ อุปสํยนฺติ, \\ วิวฏจฺฉทา\textsuperscript{1} เย โลเก, \\ ยญฺญเมตํ ปสํสนฺติ, \\ ยญฺเญ วา ยทิ วา สทฺเธ, \\ ปสนฺนจิตฺโต ยชติ\textsuperscript{1},}
\csromangathagroup{\csromangathastanza{\csromangathabat{อภิสงฺขตํ นิรารมฺภํ,}{ยญฺญํ กาเลน กปฺปิยํ.} \\ \csromangathabat{ตาทิสํ อุปสํยนฺติ,}{สญฺญตา พฺรหฺมจารโย.}}\csromangathastanza{\csromangathabat{วิวฏจฺฉทา\csromansharedfootnote{5006a14ec66c339b}{วิวตฺตจฺฉทา (สี, อิ), วิวฏฺฏจฺฉทา (ก)} เย โลเก,}{วีติวตฺตา กุลํ คติํ.} \\ \csromangathabat{ยญฺญเมตํ ปสํสนฺติ,}{พุทฺธา ยญฺญสฺส\csromansharedfootnote{435a1221f14f9545}{ปุญฺญสฺส (สฺยา, กํ, อิ) 6. หวฺยํ (สี, อิ), หุญฺญํ (สฺยา, กํ)} โกวิทา.}}\csromangathastanza{\csromangathabat{ยญฺเญ วา ยทิ วา สทฺเธ,}{หพฺยํ\csromansharedfootnote{15e55b6f7f3d7d0c}{หวฺยํ (สี, อิ), หุญฺญํ (สฺยา, กํ)} กตฺวา ยถารหํ.} \\ \csromangathabat{ปสนฺนจิตฺโต ยชติ\csromansharedfootnote{4878753720e1734b}{ปสนฺนจิตฺตา ยชนฺติ (ก)},}{สุเขตฺเต พฺรหฺมจาริสุ.}}}

\chapterheadpageread{5. โรหิตสฺสวคฺค}
\setlength{\csromangathapagewidth}{0pt}
\setlength{\csromangathawakwidth}{0pt}
\csromangathameasuregroup{สุหุตํ สุยิฏฺฐํ สุปฺปตฺตํ\textsuperscript{1}, \\ ยญฺโญ จ วิปุโล โหติ, \\ เอวํ\textsuperscript{1} ยชิตฺวา เมธาวี, \\ อพฺยาพชฺฌํ สุขํ โลกํ,}{\csromangathabat{สุหุตํ สุยิฏฺฐํ สุปฺปตฺตํ\textsuperscript{1},}{ทกฺขิเณยฺเยสุ ยํ กตํ.} \\ \csromangathabat{ยญฺโญ จ วิปุโล โหติ,}{ปสีทนฺติ จ เทวตา.} \\ \csromangathabat{เอวํ\textsuperscript{1} ยชิตฺวา เมธาวี,}{สทฺโธ มุตฺเตน เจตสา.} \\ \csromangathabat{อพฺยาพชฺฌํ สุขํ โลกํ,}{ปณฺฑิโต อุปปชฺชตีติ.} \\ ทสมํ.}
\csromangathameasurewak{ทสมํ.}
\csromangathasetleft{สุหุตํ สุยิฏฺฐํ สุปฺปตฺตํ\textsuperscript{1}, \\ ยญฺโญ จ วิปุโล โหติ, \\ เอวํ\textsuperscript{1} ยชิตฺวา เมธาวี, \\ อพฺยาพชฺฌํ สุขํ โลกํ,}
\csromangathagroupwithclosermidruled{\csromangathastanza{\csromanfolio{353}\csromangathabat{สุหุตํ สุยิฏฺฐํ สุปฺปตฺตํ\csromansharedfootnote{b1e6c3fc2a6cc4aa}{สมฺปตฺตํ (สฺยา, กํ, ก)},}{ทกฺขิเณยฺเยสุ ยํ กตํ.} \\ \csromangathabat{ยญฺโญ จ วิปุโล โหติ,}{ปสีทนฺติ จ เทวตา.}}\csromangathastanza{\csromangathabat{เอวํ\csromansharedfootnote{338dd9d64fd7ee1e}{[มิสฺสินฺคฺ โนเต 0]} ยชิตฺวา เมธาวี,}{สทฺโธ มุตฺเตน เจตสา.} \\ \csromangathabat{อพฺยาพชฺฌํ สุขํ โลกํ,}{ปณฺฑิโต อุปปชฺชตีติ.}}\csromangathastanza{\csromangathawak{ทสมํ.}}}{จกฺกวคฺโค จตุตฺโถ.}
\csromancenter{\textbf{ตสฺสุทฺทานํ}}
\setlength{\csromangathapagewidth}{0pt}
\setlength{\csromangathawakwidth}{0pt}
\csromangathameasuregroup{จกฺโก สงฺคโห สีโห, \\ โทโณ อปริหานิโย ปติลีโน,}{\csromangathabat{จกฺโก สงฺคโห สีโห,}{ปสาโท วสฺสกาเรน ปญฺจมํ.} \\ \csromangathabat{โทโณ อปริหานิโย ปติลีโน,}{อุชฺชโย อุทายินา เต ทสาติ.}}
\csromangathasetleft{จกฺโก สงฺคโห สีโห, \\ โทโณ อปริหานิโย ปติลีโน,}
\csromangathagroup{\csromangathastanza{\csromangathabat{จกฺโก สงฺคโห สีโห,}{ปสาโท วสฺสกาเรน ปญฺจมํ.} \\ \csromangathabat{โทโณ อปริหานิโย ปติลีโน,}{อุชฺชโย อุทายินา เต ทสาติ.}}}

\chapterhead{5. โรหิตสฺสวคฺค}
\csromanmatikaanchor{mtk.94}
\csromanmatikaanchor{mtk.95}
\csromantocmarknopage{section}{5. โรหิตสฺสวคฺค}{mtk.94}
\bookmark[dest={mtk.94},level=1]{5. โรหิตสฺสวคฺค}
\csromantocmarkref{subsection}{1. สมาธิภาวนาสุตฺต}{mtk.95}
\bookmark[dest={mtk.95},level=2]{1. สมาธิภาวนาสุตฺต}
\csromanheader{2}{1. สมาธิภาวนา\-สุตฺต}
\proseitem{41}{จตสฺโส อิมา ภิกฺขเว สมาธิภาวนา. กตมา จตสฺโส? อตฺถิ ภิกฺขเว สมาธิภาวนา ภาวิตา พหุลีกตา ทิฏฺฐ\-ธมฺม\-สุข\-วิหาราย สํวตฺตติ, อตฺถิ ภิกฺขเว สมาธิภาวนา ภาวิตา พหุลีกตา ญาณทสฺสน\-ปฺปฏิลาภาย สํวตฺตติ, อตฺถิ ภิกฺขเว สมาธิภาวนา ภาวิตา พหุลีกตา สติสมฺปชญฺญาย สํวตฺตติ, อตฺถิ ภิกฺขเว สมาธิภาวนา ภาวิตา พหุลีกตา อาสวานํ ขยาย สํวตฺตติ.}

\prose{กตมา จ ภิกฺขเว สมาธิภาวนา ภาวิตา พหุลีกตา ทิฏฺฐ\-ธมฺม\-สุข\-วิหาราย สํวตฺตติ? อิธ ภิกฺขเว ภิกฺขุ วิวิจฺเจว กาเมหิ ฯเปฯ จตุตฺถํ ฌานํ อุปสมฺปชฺช วิหรติ. อยํ ภิกฺขเว สมาธิภาวนา ภาวิตา พหุลีกตา ทิฏฺฐ\-ธมฺม\-สุข\-วิหาราย สํวตฺตติ.}

\prose{\csromanfolio{354}กตมา จ ภิกฺขเว สมาธิภาวนา ภาวิตา พหุลีกตา ญาณทสฺสน\-ปฺปฏิลาภาย สํวตฺตติ? อิธ ภิกฺขเว ภิกฺขุ อาโลกสญฺญํ มนสิ กโรติ, ทิวาสญฺญํ อธิฏฺฐาติ ยถา ทิวา ตถา รตฺติํ, ยถา รตฺติํ ตถา ทิวา, อิติ วิวเฏน เจตสา อปริโยนทฺเธน สปฺปภาสํ จิตฺตํ ภาเวติ. อยํ ภิกฺขเว สมาธิภาวนา ภาวิตา พหุลีกตา ญาณทสฺสน\-ปฺปฏิลาภาย สํวตฺตติ.}

\prose{กตมา จ ภิกฺขเว สมาธิภาวนา ภาวิตา พหุลีกตา สติสมฺปชญฺญาย สํวตฺตติ? อิธ ภิกฺขเว ภิกฺขุโน วิทิตา เวทนา อุปฺปชฺชนฺติ, วิทิตา อุปฏฺฐหนฺติ, วิทิตา อพฺภตฺถํ คจฺฉนฺติ. วิทิตา สญฺญา ฯเปฯ. วิทิตา วิตกฺกา อุปฺปชฺชนฺติ, วิทิตา อุปฏฺฐหนฺติ, วิทิตา อพฺภตฺถํ คจฺฉนฺติ. อยํ ภิกฺขเว สมาธิภาวนา ภาวิตา พหุลีกตา สติสมฺปชญฺญาย สํวตฺตติ.}

\prose{กตมา จ ภิกฺขเว สมาธิภาวนา ภาวิตา พหุลีกตา อาสวานํ ขยาย สํวตฺตติ, อิธ ภิกฺขเว ภิกฺขุ ปญฺจสุ อุปาทาน\-กฺขนฺเธสุ อุทยพฺพยา\-นุปสฺสี วิหรติ “อิติ รูปํ, อิติ รูปสฺส สมุทโย, อิติ รูปสฺส อตฺถงฺคโม\csromansharedfootnote{273837cbc7e4d212}{อตฺถคโม (สี, อิ)}. อิติ เวทนา, อิติ เวทนาย สมุทโย, อิติ เวทนาย อตฺถงฺคโม. อิติ สญฺญา, อิติ สญฺญาย สมุทโย, อิติ สญฺญาย อตฺถงฺคโม. อิติ สงฺขารา, อิติ สงฺขารานํ สมุทโย, อิติ สงฺขารานํ อตฺถงฺคโม. อิติ วิญฺญาณํ, อิติ วิญฺญาณสฺส สมุทโย, อิติ วิญฺญาณสฺส อตฺถงฺคโม”ติ. อยํ ภิกฺขเว สมาธิภาวนา ภาวิตา พหุลีกตา อาสวานํ ขยาย สํวตฺตติ. อิมา โข ภิกฺขเว จตสฺโส สมาธิภาวนา. อิทญฺจ ปน เมตํ ภิกฺขเว สนฺธาย ภาสิตํ ปารายเน ปุณฺณกปเญฺห\csromandash{}}

\setlength{\csromangathapagewidth}{0pt}
\setlength{\csromangathawakwidth}{0pt}
\csromangathameasuregroup{}{“สงฺขาย โลกสฺมิํ ปโรปรานิ, \\ ยสฺสิญฺชิตํ นตฺถิ กุหิญฺจิ โลเก. \\ สนฺโต วิธูโม อนีโฆ นิราโส, \\ ปฐมํ.}
\csromangathameasurewak{“สงฺขาย โลกสฺมิํ ปโรปรานิ, \\ ยสฺสิญฺชิตํ นตฺถิ กุหิญฺจิ โลเก. \\ สนฺโต วิธูโม อนีโฆ นิราโส, \\ ปฐมํ.}
\setlength{\csromangathaleftcol}{0pt}
\csromangathagroup{\csromangathastanza{\csromangathawak{“สงฺขาย โลกสฺมิํ ปโรปรานิ, \\ ยสฺสิญฺชิตํ นตฺถิ กุหิญฺจิ โลเก. \\ สนฺโต วิธูโม อนีโฆ นิราโส,}}\csromangathastanza{\csromangathawak{ปฐมํ.}}}

\prose{อตาริ โส ชาติชรนฺติ พฺรูมี”ติ\csromansharedfootnote{7fdb0071f75c4537}{ขุ 1. 436; ขุ 8. 8, 58 ปิฏฺเฐสุปิ.}.}

\csromanmatikaanchor{mtk.96}
\csromantocmarkref{subsection}{2. ปญฺหพฺยากรณสุตฺต}{mtk.96}
\bookmark[dest={mtk.96},level=2]{2. ปญฺหพฺยากรณสุตฺต}
\csromanmarkh{5. โรหิตสฺสวคฺค}\csromanmarkhii{2. ปญฺหพฺยากรณ\-สุตฺต}\csromanheadernomark{2}{5. โรหิตสฺสวคฺค 2. ปญฺหพฺยากรณ\-สุตฺต}
\csromanmatikaanchor{mtk.97}
\csromantocmarkref{subsection}{3-4. โกธครุสุตฺตานิ}{mtk.97}
\bookmark[dest={mtk.97},level=2]{3-4. โกธครุสุตฺตานิ}
\proseitem{42}{\csromanfolio{355}จตฺตาริมานิ ภิกฺขเว ปญฺห\-พฺยากรณานิ\csromansharedfootnote{d16998e33102d09e}{ปญฺหาพฺยากรณานิ (ก)}. กตมานิ จตฺตาริ? อตฺถิ ภิกฺขเว ปโญฺห เอกํส\-พฺยากรณีโย, อตฺถิ ภิกฺขเว ปโญฺห วิภชฺช\-พฺยากรณีโย, อตฺถิ ภิกฺขเว ปโญฺห ปฏิปุจฺฉา\-พฺยากรณีโย, อตฺถิ ภิกฺขเว ปโญฺห ฐปนีโย. อิมานิ โข ภิกฺขเว จตฺตาริ ปญฺหพฺยากรณานีติ.}

\setlength{\csromangathapagewidth}{0pt}
\setlength{\csromangathawakwidth}{0pt}
\csromangathameasuregroup{เอกํสวจนํ เอกํ, \\ ตติยํ ปฏิปุจฺเฉยฺย, \\ โย จ เตสํ\textsuperscript{1} ตตฺถ ตตฺถ, \\ จตุปญฺหสฺส กุสโล, \\ ทุราสโท ทุปฺปสโห, \\ อโถ อตฺเถ อนตฺเถ จ, \\ อนตฺถํ ปริวชฺเชติ,}{\csromangathabat{เอกํสวจนํ เอกํ,}{วิภชฺชวจนาปรํ.} \\ \csromangathabat{ตติยํ ปฏิปุจฺเฉยฺย,}{จตุตฺถํ ปน ฐาปเย.} \\ \csromangathabat{โย จ เตสํ\textsuperscript{1} ตตฺถ ตตฺถ,}{ชานาติ อนุธมฺมตํ.} \\ \csromangathabat{จตุปญฺหสฺส กุสโล,}{อาหุ ภิกฺขุํ ตถาวิธํ.} \\ \csromangathabat{ทุราสโท ทุปฺปสโห,}{คมฺภีโร ทุปฺปธํสิโย.} \\ \csromangathabat{อโถ อตฺเถ อนตฺเถ จ,}{อุภยสฺส โหติ โกวิโท\textsuperscript{1}.} \\ \csromangathabat{อนตฺถํ ปริวชฺเชติ,}{อตฺถํ คณฺหาติ ปณฺฑิโต.} \\ ทุติยํ.}
\csromangathameasurewak{ทุติยํ.}
\csromangathasetleft{เอกํสวจนํ เอกํ, \\ ตติยํ ปฏิปุจฺเฉยฺย, \\ โย จ เตสํ\textsuperscript{1} ตตฺถ ตตฺถ, \\ จตุปญฺหสฺส กุสโล, \\ ทุราสโท ทุปฺปสโห, \\ อโถ อตฺเถ อนตฺเถ จ, \\ อนตฺถํ ปริวชฺเชติ,}
\csromangathagroup{\csromangathastanza{\csromangathabat{เอกํสวจนํ เอกํ,}{วิภชฺชวจนาปรํ.} \\ \csromangathabat{ตติยํ ปฏิปุจฺเฉยฺย,}{จตุตฺถํ ปน ฐาปเย.}}\csromangathastanza{\csromangathabat{โย จ เตสํ\csromansharedfootnote{fb5a3c33d7d312a8}{เนสํ (สี, สฺยา, กํ)} ตตฺถ ตตฺถ,}{ชานาติ อนุธมฺมตํ.} \\ \csromangathabat{จตุปญฺหสฺส กุสโล,}{อาหุ ภิกฺขุํ ตถาวิธํ.}}\csromangathastanza{\csromangathabat{ทุราสโท ทุปฺปสโห,}{คมฺภีโร ทุปฺปธํสิโย.} \\ \csromangathabat{อโถ อตฺเถ อนตฺเถ จ,}{อุภยสฺส โหติ โกวิโท\csromansharedfootnote{196ef65e7e75ca11}{อุภยตฺถสฺส โกวิโท (สฺยา, กํ)}.} \\ \csromangathabat{อนตฺถํ ปริวชฺเชติ,}{อตฺถํ คณฺหาติ ปณฺฑิโต.}}\csromangathastanza{\csromangathawak{ทุติยํ.}}}

\prose{อตฺถาภิสมยา ธีโร, “ปณฺฑิโต”ติ ปวุจฺจตีติ.}

\csromanheader{2}{3. ปฐม\-โกธครุ\-สุตฺต}
\proseitem{43}{จตฺตาโรเม ภิกฺขเว ปุคฺคลา สนฺโต สํวิชฺชมานา โลกสฺมิํ. กตเม จตฺตาโร? โกธครุ น สทฺธมฺมครุ, มกฺขครุ น สทฺธมฺมครุ, ลาภครุ น สทฺธมฺมครุ, สกฺการครุ น สทฺธมฺมครุ. อิเม โข ภิกฺขเว จตฺตาโร ปุคฺคลา สนฺโต สํวิชฺชมานา โลกสฺมิํ.}

\prose{จตฺตาโรเม ภิกฺขเว ปุคฺคลา สนฺโต สํวิชฺชมานา โลกสฺมิํ. กตเม จตฺตาโร? สทฺธมฺมครุ น โกธครุ, สทฺธมฺมครุ น มกฺขครุ, สทฺธมฺมครุ น ลาภครุ, สทฺธมฺมครุ น สกฺการครุ. อิเม โข ภิกฺขเว จตฺตาโร ปุคฺคลา สนฺโต สํวิชฺชมานา โลกสฺมินฺติ.}

\csromanmatikaanchor{mtk.98}
\csromantocmarkref{subsection}{5-6. โรหิตสฺสสุตฺตานิ}{mtk.98}
\bookmark[dest={mtk.98},level=2]{5-6. โรหิตสฺสสุตฺตานิ}
\setlength{\csromangathapagewidth}{0pt}
\setlength{\csromangathawakwidth}{0pt}
\csromangathameasuregroup{โกธมกฺขครู ภิกฺขู, \\ น เต ธมฺเม วิรูหนฺติ, \\ เย จ สทฺธมฺมครุโน, \\ เต เว ธมฺเม วิรูหนฺติ,}{\csromangathabat{โกธมกฺขครู ภิกฺขู,}{ลาภสกฺการคารวา.} \\ \csromangathabat{น เต ธมฺเม วิรูหนฺติ,}{สมฺมาสมฺพุทฺธเทสิเต.} \\ \csromangathabat{เย จ สทฺธมฺมครุโน,}{วิหํสุ วิหรนฺติ จ.} \\ \csromangathabat{เต เว ธมฺเม วิรูหนฺติ,}{สมฺมาสมฺพุทฺธเทสิเตติ.} \\ ตติยํ.}
\csromangathameasurewak{ตติยํ.}
\csromangathasetleft{โกธมกฺขครู ภิกฺขู, \\ น เต ธมฺเม วิรูหนฺติ, \\ เย จ สทฺธมฺมครุโน, \\ เต เว ธมฺเม วิรูหนฺติ,}
\csromangathagroup{\csromangathastanza{\csromangathabat{โกธมกฺขครู ภิกฺขู,}{ลาภสกฺการคารวา.} \\ \csromangathabat{น เต ธมฺเม วิรูหนฺติ,}{สมฺมาสมฺพุทฺธเทสิเต.}}\csromangathastanza{\csromanfolio{356}\csromangathabat{เย จ สทฺธมฺมครุโน,}{วิหํสุ วิหรนฺติ จ.} \\ \csromangathabat{เต เว ธมฺเม วิรูหนฺติ,}{สมฺมาสมฺพุทฺธเทสิเตติ.}}\csromangathastanza{\csromangathawak{ตติยํ.}}}

\csromanheader{2}{4. ทุติย\-โกธครุ\-สุตฺต}
\proseitem{44}{จตฺตาโรเม ภิกฺขเว อสทฺธมฺมา. กตเม จตฺตาโร? โกธครุตา น สทฺธมฺม\-ครุตา, มกฺขครุตา น สทฺธมฺม\-ครุตา, ลาภครุตา น สทฺธมฺม\-ครุตา, สกฺการครุตา น สทฺธมฺม\-ครุตา. อิเม โข ภิกฺขเว จตฺตาโร อสทฺธมฺมา.}

\prose{จตฺตาโรเม ภิกฺขเว สทฺธมฺมา. กตเม จตฺตาโร? สทฺธมฺม\-ครุตา น โกธครุตา, สทฺธมฺม\-ครุตา น มกฺขครุตา, สทฺธมฺม\-ครุตา น ลาภครุตา, สทฺธมฺม\-ครุตา น สกฺการครุตา. อิเม โข ภิกฺขเว จตฺตาโร สทฺธมฺมาติ.}

\setlength{\csromangathapagewidth}{0pt}
\setlength{\csromangathawakwidth}{0pt}
\csromangathameasuregroup{โกธมกฺขครุ ภิกฺขุ, \\ สุเขตฺเต ปูติพีชํว, \\ เย จ สทฺธมฺมครุโน, \\ เต เว ธมฺเม วิรูหนฺติ,}{\csromangathabat{โกธมกฺขครุ ภิกฺขุ,}{ลาภสกฺการคารโว.} \\ \csromangathabat{สุเขตฺเต ปูติพีชํว,}{สทฺธมฺเม น วิรูหติ.} \\ \csromangathabat{เย จ สทฺธมฺมครุโน,}{วิหํสุ วิหรนฺติ จ.} \\ \csromangathabat{เต เว ธมฺเม วิรูหนฺติ,}{สฺเนหานฺวยมิโวสธาติ\textsuperscript{1}.} \\ จตุตฺถํ.}
\csromangathameasurewak{จตุตฺถํ.}
\csromangathasetleft{โกธมกฺขครุ ภิกฺขุ, \\ สุเขตฺเต ปูติพีชํว, \\ เย จ สทฺธมฺมครุโน, \\ เต เว ธมฺเม วิรูหนฺติ,}
\csromangathagroup{\csromangathastanza{\csromangathabat{โกธมกฺขครุ ภิกฺขุ,}{ลาภสกฺการคารโว.} \\ \csromangathabat{สุเขตฺเต ปูติพีชํว,}{สทฺธมฺเม น วิรูหติ.}}\csromangathastanza{\csromangathabat{เย จ สทฺธมฺมครุโน,}{วิหํสุ วิหรนฺติ จ.} \\ \csromangathabat{เต เว ธมฺเม วิรูหนฺติ,}{สฺเนหานฺวยมิโวสธาติ\csromansharedfootnote{338dd9d64fd7ee1e}{[มิสฺสินฺคฺ โนเต 0]}.}}\csromangathastanza{\csromangathawak{จตุตฺถํ.}}}

\csromanheader{2}{5. \textbf{โรหิตสฺสสุตฺต}}
\proseitem{45}{เอกํ สมยํ ภควา\csromansharedfootnote{b8d79c18d5825d54}{สํ 1. 60 ปิฏฺเฐปิ.} สาวตฺถิยํ วิหรติ เชตวเน อนาถ\-ปิณฺฑิกสฺส อาราเม. อถ โข โรหิตสฺโส เทวปุตฺโต อภิกฺกนฺตาย รตฺติยา อภิกฺกนฺต\-วณฺโณ เกวลกปฺปํ เชตวนํ โอภาเสตฺวา เยน ภควา เตนุปสงฺกมิ, อุปสงฺกมิตฺวา ภควนฺตํ อภิวาเทตฺวา เอกมนฺตํ อฏฺฐาสิ, เอกมนฺตํ ฐิโต โข โรหิตสฺโส เทวปุตฺโต ภควนฺตํ เอตทโวจ\csromandash{}}

\prose{“ยตฺถ นุ โข ภนฺเต น ชายติ น ชียติ น มียติ น จวติ น อุปปชฺชติ, สกฺกา นุ โข โส ภนฺเต คมเนน โลกสฺส อนฺโต ญาตุํ วา ทฏฺฐุํ วา ปาปุณิตุํ วา”ติ. ยตฺถ โข อาวุโส น ชายติ \csromanfolio{357}น ชียติ น มียติ น จวติ น อุปปชฺชติ, นาหํ “ตํ คมเนน โลกสฺส อนฺตํ ญาเตยฺยํ ทฏฺเฐยฺยํ ปตฺเตยฺยนฺ”ติ วทามีติ.}

\prose{อจฺฉริยํ ภนฺเต, อพฺภุตํ ภนฺเต, ยาว สุภาสิตมิทํ ภนฺเต ภควตา “ยตฺถ โข อาวุโส น ชายติ น ชียติ น มียติ น จวติ น อุปปชฺชติ, นาหํ ‘ตํ คมเนน โลกสฺส อนฺตํ ญาเตยฺยํ ทฏฺเฐยฺยํ ปตฺเตยฺยนฺ’ติ วทามี”ติ.}

\prose{ภูตปุพฺพาหํ ภนฺเต โรหิตสฺโส นาม อิสิ อโหสิํ โภชปุตฺโต อิทฺธิมา เวหาสงฺคโม. ตสฺส มยฺหํ ภนฺเต เอวรูโป ชโว อโหสิ\csromandash{}เสยฺยถาปิ นาม ทฬฺหธมฺมา\csromansharedfootnote{405334b3de868cc3}{ทฬฺหธมฺโม, อํ 3. 226; ม 1. 117 ปิฏฺเฐ. (สพฺพตฺถ) ฏีกา จ โมคฺคลฺลานพฺยากรณํ จ โอโลเกตพฺพํ.} ธนุคฺคโห สิกฺขิโต กตหตฺโถ กตูปาสโน ลหุเกน อสเนน อปฺปกสิเรน ติริยํ ตาลจฺฉายํ อติปาเตยฺย. ตสฺส มยฺหํ ภนฺเต เอวรูโป ปทวีติหาโร อโหสิ\csromandash{}เสยฺยถาปิ นาม ปุรตฺถิมา สมุทฺทา ปจฺฉิโม สมุทฺโท. ตสฺส มยฺหํ ภนฺเต เอวรูเปน ชเวน สมนฺนาคตสฺส, เอวรูเปน จ ปทวีติหาเรน เอวรูปํ อิจฺฉาคตํ อุปฺปชฺชิ “อหํ คมเนน โลกสฺส อนฺตํ ปาปุณิสฺสามี”ติ. โส โข อหํ ภนฺเต อญฺญเตฺรว อสิต\-ปีต\-ขายิต\-สายิตา, อญฺญตฺร อุจฺจาร\-ปสฺสาว\-กมฺมา, อญฺญตฺร นิทฺทา\-กิลมถ\-ปฏิวิโนทนา วสฺสสตายุโก วสฺสสตชีวี วสฺสสตํ คนฺตฺวา อปฺปตฺวาว โลกสฺส อนฺตํ อนฺตราเยว กาลงฺกโต.}

\prose{อจฺฉริยํ ภนฺเต, อพฺภุตํ ภนฺเต, ยาว สุภาสิตมิทํ ภนฺเต ภควตา “ยตฺถ โข อาวุโส น ชายติ น ชียติ น มียติ น จวติ น อุปปชฺชติ, นาหํ ‘ตํ คมเนน โลกสฺส อนฺตํ ญาเตยฺยํ ทฏฺเฐยฺยํ ปตฺเตยฺยนฺ’ติ วทามี”ติ.}

\prose{ยตฺถ โข อาวุโส น ชายติ น ชียติ น มียติ น จวติ น อุปปชฺชติ, นาหํ “ตํ คมเนน โลกสฺส อนฺตํ ญาเตยฺยํ ทฏฺเฐยฺยํ ปตฺเตยฺยนฺ”ติ วทามิ. น จาหํ อาวุโส อปฺปตฺวาว โลกสฺส อนฺตํ ทุกฺขสฺส อนฺตกิริยํ วทามิ. อปิ จาหํ อาวุโส อิมสฺมิํเยว พฺยามมตฺเต กเฬวเร\csromansharedfootnote{9879a487cfb37284}{กเฬพเร (สี, อิ)} สสญฺญิมฺหิ สมนเก โลกญฺจ ปญฺญาเปมิ โลกสมุทยญฺจ โลกนิโรธญฺจ โลก\-นิโรธ\-คามินิญฺจ ปฏิปทนฺติ.}

\setlength{\csromangathapagewidth}{0pt}
\setlength{\csromangathawakwidth}{0pt}
\csromangathameasuregroup{คมเนน น ปตฺตพฺโพ, \\ น จ อปฺปตฺวา โลกนฺตํ,}{\csromangathabat{คมเนน น ปตฺตพฺโพ,}{โลกสฺสนฺโต กุทาจนํ.} \\ \csromangathabat{น จ อปฺปตฺวา โลกนฺตํ,}{ทุกฺขา อตฺถิ ปโมจนํ.} \\ ตสฺมา หเว โลกวิทู สุเมโธ, \\ โลกนฺตคู วุสิตพฺรหฺมจริโย. \\ โลกสฺส อนฺตํ สมิตาวิ ญตฺวา, \\ นาสีสตี\textsuperscript{1} โลกมิมํ ปรญฺจาติ. \\ ปญฺจมํ.}
\csromangathameasurewak{ตสฺมา หเว โลกวิทู สุเมโธ, \\ โลกนฺตคู วุสิตพฺรหฺมจริโย. \\ โลกสฺส อนฺตํ สมิตาวิ ญตฺวา, \\ นาสีสตี\textsuperscript{1} โลกมิมํ ปรญฺจาติ. \\ ปญฺจมํ.}
\csromangathasetleft{คมเนน น ปตฺตพฺโพ, \\ น จ อปฺปตฺวา โลกนฺตํ,}
\csromangathagroup{\csromangathastanza{\csromanfolio{358}\csromangathabat{คมเนน น ปตฺตพฺโพ,}{โลกสฺสนฺโต กุทาจนํ.} \\ \csromangathabat{น จ อปฺปตฺวา โลกนฺตํ,}{ทุกฺขา อตฺถิ ปโมจนํ.}}\csromangathastanza{\csromangathawak{ตสฺมา หเว โลกวิทู สุเมโธ, \\ โลกนฺตคู วุสิต\-พฺรหฺมจริโย. \\ โลกสฺส อนฺตํ สมิตาวิ ญตฺวา, \\ นาสีสตี\csromansharedfootnote{338dd9d64fd7ee1e}{[มิสฺสินฺคฺ โนเต 0]} โลกมิมํ ปรญฺจาติ.}}\csromangathastanza{\csromangathawak{ปญฺจมํ.}}}

\csromanheader{2}{6. ทุติย\-โรหิตสฺส\-สุตฺต}
\proseitem{46}{อถ โข ภควา ตสฺสา รตฺติยา อจฺจเยน ภิกฺขู อามนฺเตสิ “อิมํ ภิกฺขเว รตฺติํ โรหิตสฺโส เทวปุตฺโต อภิกฺกนฺตาย รตฺติยา อภิกฺกนฺต\-วณฺโณ เกวลกปฺปํ เชตวนํ โอภาเสตฺวา เยนาหํ เตนุปสงฺกมิ, อุปสงฺกมิตฺวา มํ อภิวาเทตฺวา เอกมนฺตํ อฏฺฐาสิ, เอกมนฺตํ ฐิโต โข ภิกฺขเว โรหิตสฺโส เทวปุตฺโต มํ เอตทโวจ\csromandash{}}

\prose{ยตฺถ นุ โข ภนฺเต\csromansharedfootnote{b8d79c18d5825d54}{สํ 1. 60 ปิฏฺเฐปิ.} น ชายติ น ชียติ น มียติ น จวติ น อุปปชฺชติ, สกฺกา นุ โข โส ภนฺเต คมเนน โลกสฺส อนฺโต ญาตุํ วา ทฏฺฐุํ วา ปาปุณิตุํ วาติ. เอวํ วุตฺเต อหํ ภิกฺขเว โรหิตสฺสํ เทวปุตฺตํ เอตทโวจํ\csromandash{}}

\prose{ยตฺถ โข อาวุโส น ชายติ น ชียติ น มียติ น จวติ น อุปปชฺชติ, นาหํ “ตํ คมเนน โลกสฺส อนฺตํ ญาเตยฺยํ ทฏฺเฐยฺยํ ปตฺเตยฺยนฺ”ติ วทามีติ. เอวํ วุตฺเต ภิกฺขเว โรหิตสฺโส เทวปุตฺโต มํ เอตทโวจ\csromandash{}}

\prose{อจฺฉริยํ ภนฺเต, อพฺภุตํ ภนฺเต, ยาว สุภาสิตมิทํ ภนฺเต ภควตา, ยตฺถ โข อาวุโส น ชายติ น ชียติ น มียติ น จวติ น อุปปชฺชติ, นาหํ “ตํ คมเนน โลกสฺส อนฺตํ ญาเตยฺยํ ทฏฺเฐยฺยํ ปตฺเตยฺยนฺ”ติ วทามิ.}

\prose{ภูตปุพฺพาหํ ภนฺเต โรหิตสฺโส นาม อิสิ อโหสิํ โภชปุตฺโต อิทฺธิมา เวหาสงฺคโม. ตสฺส มยฺหํ ภนฺเต เอวรูโป ชโว อโหสิ\csromandash{}เสยฺยถาปิ นาม ทฬฺหธมฺมา ธนุคฺคโห สิกฺขิโต กตหตฺโถ กตูปาสโน ลหุเกน อสเนน อปฺปกสิเรน ติริยํ ตาลจฺฉายํ \csromanfolio{359}อติปาเตยฺย. ตสฺส มยฺหํ ภนฺเต เอวรูโป ปทวีติหาโร อโหสิ\csromandash{}เสยฺยถาปิ นาม ปุรตฺถิมา สมุทฺทา ปจฺฉิโม สมุทฺโท. ตสฺส มยฺหํ ภนฺเต เอวรูเปน ชเวน สมนฺนาคตสฺส, เอวรูเปน จ ปทวีติหาเรน เอวรูปํ อิจฺฉาคตํ อุปฺปชฺชิ “อหํ คมเนน โลกสฺส อนฺตํ ปาปุณิสฺสามี”ติ. โส โข อหํ ภนฺเต อญฺญเตฺรว อสิต\-ปีต\-ขายิต\-สายิตา, อญฺญตฺร อุจฺจาร\-ปสฺสาว\-กมฺมา, อญฺญตฺร นิทฺทา\-กิลมถ\-ปฏิวิโนทนา วสฺสสตายุโก วสฺสสตชีวี วสฺสสตํ คนฺตฺวา อปฺปตฺวาว โลกสฺส อนฺตํ อนฺตราเยว กาลงฺกโต.}

\prose{อจฺฉริยํ ภนฺเต, อพฺภุตํ ภนฺเต, ยาว สุภาสิตมิทํ ภนฺเต ภควตา “ยตฺถ โข อาวุโส น ชายติ น ชียติ น มียติ น จวติ น อุปปชฺชติ, นาหํ ‘ตํ คมเนน โลกสฺส อนฺตํ ญาเตยฺยํ ทฏฺเฐยฺยํ ปตฺเตยฺยนฺ’ติ วทามี”ติ. เอวํ วุตฺเต อหํ ภิกฺขเว โรหิตสฺสํ เทวปุตฺตํ เอตทโวจํ\csromandash{}}

\prose{ยตฺถ โข อาวุโส น ชายติ น ชียติ น มียติ น จวติ น อุปปชฺชติ, นาหํ “ตํ คมเนน โลกสฺส อนฺตํ ญาเตยฺยํ ทฏฺเฐยฺยํ ปตฺเตยฺยนฺ”ติ วทามีติ. น จาหํ อาวุโส อปฺปตฺวาว โลกสฺส อนฺตํ ทุกฺขสฺส\-นฺต\-กิริยํ วทามิ, อปิ จาหํ อาวุโส อิมสฺมิํเยว พฺยามมตฺเต กเฬวเร สสญฺญิมฺหิ สมนเก โลกญฺจ ปญฺญาเปมิ โลกสมุทยญฺจ โลกนิโรธญฺจ โลก\-นิโรธ\-คามินิญฺจ ปฏิปทนฺติ.}

\setlength{\csromangathapagewidth}{0pt}
\setlength{\csromangathawakwidth}{0pt}
\csromangathameasuregroup{คมเนน น ปตฺตพฺโพ, \\ น จ อปฺปตฺวา โลกนฺตํ,}{\csromangathabat{คมเนน น ปตฺตพฺโพ,}{โลกสฺสนฺโต กุทาจนํ.} \\ \csromangathabat{น จ อปฺปตฺวา โลกนฺตํ,}{ทุกฺขา อตฺถิ ปโมจนํ.} \\ ตสฺมา หเว โลกวิทู สุเมโธ, \\ โลกนฺตคู วุสิตพฺรหฺมจริโย. \\ โลกสฺส อนฺตํ สมิตาวิ ญตฺวา, \\ นาสีสตี โลกมิมํ ปรญฺจาติ. \\ ฉฏฺฐํ.}
\csromangathameasurewak{ตสฺมา หเว โลกวิทู สุเมโธ, \\ โลกนฺตคู วุสิตพฺรหฺมจริโย. \\ โลกสฺส อนฺตํ สมิตาวิ ญตฺวา, \\ นาสีสตี โลกมิมํ ปรญฺจาติ. \\ ฉฏฺฐํ.}
\csromangathasetleft{คมเนน น ปตฺตพฺโพ, \\ น จ อปฺปตฺวา โลกนฺตํ,}
\csromangathagroup{\csromangathastanza{\csromangathabat{คมเนน น ปตฺตพฺโพ,}{โลกสฺสนฺโต กุทาจนํ.} \\ \csromangathabat{น จ อปฺปตฺวา โลกนฺตํ,}{ทุกฺขา อตฺถิ ปโมจนํ.}}\csromangathastanza{\csromangathawak{ตสฺมา หเว โลกวิทู สุเมโธ, \\ โลกนฺตคู วุสิต\-พฺรหฺมจริโย. \\ โลกสฺส อนฺตํ สมิตาวิ ญตฺวา, \\ นาสีสตี โลกมิมํ ปรญฺจาติ.}}\csromangathastanza{\csromangathawak{ฉฏฺฐํ.}}}

\csromanmatikaanchor{mtk.99}
\csromantocmarkref{subsection}{7. สุวิทูรสุตฺต}{mtk.99}
\bookmark[dest={mtk.99},level=2]{7. สุวิทูรสุตฺต}
\csromanheader{2}{7. \textbf{สุวิทูรสุตฺต}}
\proseitem{47}{จตฺตาริมานิ ภิกฺขเว สุวิทูร\-วิทูรานิ. กตมานิ จตฺตาริ? นภญฺจ ภิกฺขเว ปถวี จ, อิทํ ปฐมํ สุวิทูรวิทูเร. โอริมญฺจ ภิกฺขเว ตีรํ สมุทฺทสฺส ปาริมญฺจ, อิทํ ทุติยํ สุวิทูรวิทูเร. ยโต จ ภิกฺขเว เวโรจโน \csromanfolio{360}อพฺภุเทติ ยตฺถ จ อตฺถเมติ\csromansharedfootnote{499cb556a57da32a}{อตฺถงฺคเมติ (สฺยา), เวติ (ก)}, อิทํ ตติยํ สุวิทูรวิทูเร. สตญฺจ ภิกฺขเว ธมฺโม อสตญฺจ ธมฺโม, อิทํ จตุตฺถํ สุวิทูรวิทูเร. อิมานิ โข ภิกฺขเว จตฺตาริ สุวิทูรวิทูรานีติ.}

\setlength{\csromangathapagewidth}{0pt}
\setlength{\csromangathawakwidth}{0pt}
\csromangathameasuregroup{}{\textsuperscript{*}นภญฺจ ทูเร ปถวี จ ทูเร, \\ ปารํ สมุทฺทสฺส ตทาหุ ทูเร. \\ ยโต จ เวโรจโน อพฺภุเทติ, \\ ปภงฺกโร ยตฺถ จ อตฺถเมติ. \\ ตโต หเว ทูรตรํ วทนฺติ, \\ สตญฺจ ธมฺมํ อสตญฺจ ธมฺมํ. \\ อพฺยายิโก โหติ สตํ สมาคโม, \\ ยาวาปิ\textsuperscript{1} ติฏฺเฐยฺย ตเถว โหติ. \\ ขิปฺปํ หิ เวติ อสตํ สมาคโม, \\ ตสฺมา สตํ ธมฺโม อสพฺภิ อารกาติ.}
\csromangathameasurewak{\textsuperscript{*}นภญฺจ ทูเร ปถวี จ ทูเร, \\ ปารํ สมุทฺทสฺส ตทาหุ ทูเร. \\ ยโต จ เวโรจโน อพฺภุเทติ, \\ ปภงฺกโร ยตฺถ จ อตฺถเมติ. \\ ตโต หเว ทูรตรํ วทนฺติ, \\ สตญฺจ ธมฺมํ อสตญฺจ ธมฺมํ. \\ อพฺยายิโก โหติ สตํ สมาคโม, \\ ยาวาปิ\textsuperscript{1} ติฏฺเฐยฺย ตเถว โหติ. \\ ขิปฺปํ หิ เวติ อสตํ สมาคโม, \\ ตสฺมา สตํ ธมฺโม อสพฺภิ อารกาติ.}
\setlength{\csromangathaleftcol}{0pt}
\csromangathagroup{\csromangathastanza{\csromangathawak{\csromansymbolfootnote{*}{ขุ 6. 137, 142 ปิฏฺเฐสุ.}นภญฺจ ทูเร ปถวี จ ทูเร, \\ ปารํ สมุทฺทสฺส ตทาหุ ทูเร. \\ ยโต จ เวโรจโน อพฺภุเทติ, \\ ปภงฺกโร ยตฺถ จ อตฺถเมติ.}}\csromangathastanza{\csromangathawak{ตโต หเว ทูรตรํ วทนฺติ, \\ สตญฺจ ธมฺมํ อสตญฺจ ธมฺมํ. \\ อพฺยายิโก โหติ สตํ สมาคโม, \\ ยาวาปิ\csromansharedfootnote{ebbf0a0135f6aa86}{ยาวมฺปิ (สี, สฺยา, กํ, อิ)} ติฏฺเฐยฺย ตเถว โหติ. \\ ขิปฺปํ หิ เวติ อสตํ สมาคโม, \\ ตสฺมา สตํ ธมฺโม อสพฺภิ อารกาติ.}}}

\csromanmatikaanchor{mtk.100}
\csromantocmarkref{subsection}{8. วิสาขสุตฺต}{mtk.100}
\bookmark[dest={mtk.100},level=2]{8. วิสาขสุตฺต}
\csromanheader{2}{8. \textbf{วิสาขสุตฺต}}
\proseitem{48}{เอกํ สมยํ ภควา สาวตฺถิยํ วิหรติ เชตวเน อนาถ\-ปิณฺฑิกสฺส อาราเม. เตน โข ปน สมเยน อายสฺมา วิสาโข ปญฺจาลปุตฺโต\csromansharedfootnote{5a397837efb73c63}{ปญฺจาลิปุตฺโต (สี, สฺยา, กํ, อิ)} อุปฏฺฐาน\-สาลายํ ภิกฺขู ธมฺมิยา กถาย สนฺทสฺเสติ สมาทเปติ สมุตฺเตเชติ สมฺปหํเสติ โปริยา วาจาย วิสฺสฏฺฐาย อเนลคลาย อตฺถสฺส วิญฺญาปนิยา ปริยาปนฺนาย อนิสฺสิตาย. อถ โข ภควา สายนฺหสมยํ ปฏิสลฺลานา วุฏฺฐิโต เยน อุปฏฺฐานสาลา เตนุปสงฺกมิ, อุปสงฺกมิตฺวา ปญฺญตฺเต อาสเน นิสีทิ, นิสชฺช โข ภควา ภิกฺขู อามนฺเตสิ\csromandash{}}

\prose{“โก นุ โข ภิกฺขเว อุปฏฺฐาน\-สาลายํ ภิกฺขู ธมฺมิยา กถาย สนฺทสฺเสติ สมาทเปติ สมุตฺเตเชติ สมฺปหํเสติ โปริยา วาจาย วิสฺสฏฺฐาย อเนลคลาย อตฺถสฺส วิญฺญาปนิยา ปริยาปนฺนาย \csromanfolio{361}อนิสฺสิตายา”ติ. อายสฺมา ภนฺเต วิสาโข ปญฺจาลปุตฺโต อุปฏฺฐาน\-สาลายํ ภิกฺขู ธมฺมิยา กถาย สนฺทสฺเสติ สมาทเปติ สมุตฺเตเชติ สมฺปหํเสติ โปริยา วาจาย วิสฺสฏฺฐาย อเนลคลาย อตฺถสฺส วิญฺญาปนิยา ปริยาปนฺนาย อนิสฺสิตายาติ.}

\prose{อถ โข ภควา อายสฺมนฺตํ วิสาขํ ปญฺจาลปุตฺตํ เอตทโวจ “สาธุ สาธุ วิสาข, สาธุ โข ตฺวํ วิสาข ภิกฺขู ธมฺมิยา กถาย สนฺทสฺเสสิ สมาทเปสิ สมุตฺเตเชสิ สมฺปหํเสสิ โปริยา วาจาย วิสฺสฏฺฐาย อเนลคลาย อตฺถสฺส วิญฺญาปนิยา ปริยาปนฺนาย อนิสฺสิตายาติ.}

\setlength{\csromangathapagewidth}{0pt}
\setlength{\csromangathawakwidth}{0pt}
\csromangathameasuregroup{นาภาสมานํ ชานนฺติ, \\ ภาสมานญฺจ ชานนฺติ, \\ ภาสเย โชตเย ธมฺมํ, \\ สุภาสิตธชา อิสโย,}{\csromangathabat{นาภาสมานํ ชานนฺติ,}{มิสฺสํ พาเลหิ ปณฺฑิตํ.} \\ \csromangathabat{ภาสมานญฺจ ชานนฺติ,}{เทเสนฺตํ อมตํ ปทํ.} \\ \csromangathabat{ภาสเย โชตเย ธมฺมํ,}{ปคฺคเณฺห อิสินํ ธชํ.} \\ \csromangathabat{สุภาสิตธชา อิสโย,}{ธมฺโม หิ อิสินํ ธโชติ.} \\ อฏฺฐมํ.}
\csromangathameasurewak{อฏฺฐมํ.}
\csromangathasetleft{นาภาสมานํ ชานนฺติ, \\ ภาสมานญฺจ ชานนฺติ, \\ ภาสเย โชตเย ธมฺมํ, \\ สุภาสิตธชา อิสโย,}
\csromangathagroup{\csromangathastanza{\csromangathabat{นาภาสมานํ ชานนฺติ,}{มิสฺสํ พาเลหิ ปณฺฑิตํ.} \\ \csromangathabat{ภาสมานญฺจ ชานนฺติ,}{เทเสนฺตํ อมตํ ปทํ.}}\csromangathastanza{\csromangathabat{ภาสเย โชตเย ธมฺมํ,}{ปคฺคเณฺห อิสินํ ธชํ.} \\ \csromangathabat{สุภาสิตธชา อิสโย,}{ธมฺโม หิ อิสินํ ธโชติ.}}\csromangathastanza{\csromangathawak{อฏฺฐมํ.}}}

\csromanmatikaanchor{mtk.101}
\csromantocmarkref{subsection}{9. วิปลฺลาสสุตฺต}{mtk.101}
\bookmark[dest={mtk.101},level=2]{9. วิปลฺลาสสุตฺต}
\csromanheader{2}{9. \textbf{วิปลฺลาสสุตฺต}}
\proseitem{49}{จตฺตาโรเม ภิกฺขเว สญฺญาวิปลฺลาสา จิตฺตวิปลฺลาสา ทิฏฺฐิ\-วิปลฺลาสา. กตเม จตฺตาโร? อนิจฺเจ ภิกฺขเว “นิจฺจนฺ”ติ สญฺญาวิปลฺลาโส จิตฺตวิปลฺลาโส ทิฏฺฐิ\-วิปลฺลาโส, ทุกฺเข ภิกฺขเว “สุขนฺ”ติ สญฺญาวิปลฺลาโส จิตฺตวิปลฺลาโส ทิฏฺฐิ\-วิปลฺลาโส, อนตฺตนิ ภิกฺขเว “อตฺตา”ติ สญฺญาวิปลฺลาโส จิตฺตวิปลฺลาโส ทิฏฺฐิ\-วิปลฺลาโส, อสุเภ ภิกฺขเว “สุภนฺ”ติ สญฺญาวิปลฺลาโส จิตฺตวิปลฺลาโส ทิฏฺฐิ\-วิปลฺลาโส. อิเม โข ภิกฺขเว จตฺตาโร สญฺญาวิปลฺลาสา จิตฺตวิปลฺลสา ทิฏฺฐิ\-วิปลฺลาสา.}

\prose{จตฺตาโรเม ภิกฺขเว นสญฺญา\-วิปลฺลาสา นจิตฺต\-วิปลฺลาสา นทิฏฺฐิ\-วิปลฺลาสา. กตเม จตฺตาโร? อนิจฺเจ ภิกฺขเว “อนิจฺจนฺ”ติ นสญฺญา\-วิปลฺลาโส นจิตฺต\-วิปลฺลาโส นทิฏฺฐิ\-วิปลฺลาโส, ทุกฺเข ภิกฺขเว “ทุกฺขนฺ”ติ นสญฺญา\-วิปลฺลาโส นจิตฺต\-วิปลฺลาโส นทิฏฺฐิ\-วิปลฺลาโส, อนตฺตนิ ภิกฺขเว “อนตฺตา”ติ นสญฺญา\-วิปลฺลาโส นจิตฺต\-วิปลฺลาโส นทิฏฺฐิ\-วิปลฺลาโส, อสุเภ ภิกฺขเว “อสุภนฺ”ติ นสญฺญา\-วิปลฺลาโส นจิตฺต\-วิปลฺลาโส \csromanfolio{362}นทิฏฺฐิ\-วิปลฺลาโส. อิเม โข ภิกฺขเว จตฺตาโร นสญฺญา\-วิปลฺลาสา นจิตฺต\-วิปลฺลาสา นทิฏฺฐิ\-วิปลฺลาสาติ.}

\setlength{\csromangathapagewidth}{0pt}
\setlength{\csromangathawakwidth}{0pt}
\csromangathameasuregroup{อนิจฺเจ นิจฺจสญฺญิโน, \\ อนตฺตนิ จ อตฺตาติ, \\ มิจฺฉาทิฏฺฐิหตา สตฺตา, \\ เต โยคยุตฺตา มารสฺส, \\ สตฺตา คจฺฉนฺติ สํสารํ, \\ ยทา จ พุทฺธา โลกสฺมิํ, \\ เต อิมํ ธมฺมํ\textsuperscript{1} ปกาเสนฺติ, \\ เตสํ สุตฺวาน สปฺปญฺญา, \\ อนิจฺจํ อนิจฺจโต ทกฺขุํ, \\ อนตฺตนิ อนตฺตาติ, \\ สมฺมาทิฏฺฐิสมาทานา,}{\csromangathabat{อนิจฺเจ นิจฺจสญฺญิโน,}{ทุกฺเข จ สุขสญฺญิโน.} \\ \csromangathabat{อนตฺตนิ จ อตฺตาติ,}{อสุเภ สุภสญฺญิโน.} \\ \csromangathabat{มิจฺฉาทิฏฺฐิหตา สตฺตา,}{ขิตฺตจิตฺตา วิสญฺญิโน.} \\ \csromangathabat{เต โยคยุตฺตา มารสฺส,}{อโยคกฺเขมิโน ชนา.} \\ \csromangathabat{สตฺตา คจฺฉนฺติ สํสารํ,}{ชาติมรณคามิโน.} \\ \csromangathabat{ยทา จ พุทฺธา โลกสฺมิํ,}{อุปฺปชฺชนฺติ ปภงฺกรา.} \\ \csromangathabat{เต อิมํ ธมฺมํ\textsuperscript{1} ปกาเสนฺติ,}{ทุกฺขูปสมคามินํ.} \\ \csromangathabat{เตสํ สุตฺวาน สปฺปญฺญา,}{สจิตฺตํ ปจฺจลทฺธา เต.} \\ \csromangathabat{อนิจฺจํ อนิจฺจโต ทกฺขุํ,}{ทุกฺขมทฺทกฺขุ’ทุกฺขโต.} \\ \csromangathabat{อนตฺตนิ อนตฺตาติ,}{อสุภํ อสุภต’ทฺทสุํ.} \\ \csromangathabat{สมฺมาทิฏฺฐิสมาทานา,}{สพฺพํ ทุกฺขํ อุปจฺจคุนฺติ\textsuperscript{1}.} \\ นวมํ.}
\csromangathasetleft{อนิจฺเจ นิจฺจสญฺญิโน, \\ อนตฺตนิ จ อตฺตาติ, \\ มิจฺฉาทิฏฺฐิหตา สตฺตา, \\ เต โยคยุตฺตา มารสฺส, \\ สตฺตา คจฺฉนฺติ สํสารํ, \\ ยทา จ พุทฺธา โลกสฺมิํ, \\ เต อิมํ ธมฺมํ\textsuperscript{1} ปกาเสนฺติ, \\ เตสํ สุตฺวาน สปฺปญฺญา, \\ อนิจฺจํ อนิจฺจโต ทกฺขุํ, \\ อนตฺตนิ อนตฺตาติ, \\ สมฺมาทิฏฺฐิสมาทานา,}
\csromangathagroup{\csromangathastanza{\csromangathabat{อนิจฺเจ นิจฺจสญฺญิโน,}{ทุกฺเข จ สุขสญฺญิโน.} \\ \csromangathabat{อนตฺตนิ จ อตฺตาติ,}{อสุเภ สุภสญฺญิโน.}}\csromangathastanza{\csromangathabat{มิจฺฉาทิฏฺฐิหตา สตฺตา,}{ขิตฺตจิตฺตา วิสญฺญิโน.} \\ \csromangathabat{เต โยคยุตฺตา มารสฺส,}{อโยคกฺเขมิโน ชนา.}}\csromangathastanza{\csromangathabat{สตฺตา คจฺฉนฺติ สํสารํ,}{ชาติมรณคามิโน.} \\ \csromangathabat{ยทา จ พุทฺธา โลกสฺมิํ,}{อุปฺปชฺชนฺติ ปภงฺกรา.}}\csromangathastanza{\csromangathabat{เต อิมํ ธมฺมํ\csromansharedfootnote{b3c8d479f89fe7f8}{เตมํ ธมฺมํ (สี, สฺยา, กํ)} ปกาเสนฺติ,}{ทุกฺขู\-ปสม\-คามินํ.} \\ \csromangathabat{เตสํ สุตฺวาน สปฺปญฺญา,}{สจิตฺตํ ปจฺจลทฺธา เต.}}\csromangathastanza{\csromangathabat{อนิจฺจํ อนิจฺจโต ทกฺขุํ,}{ทุกฺขมทฺทกฺขุ’ทุกฺขโต.} \\ \csromangathabat{อนตฺตนิ อนตฺตาติ,}{อสุภํ อสุภต’ทฺทสุํ.}}\csromangathastanza{\csromangathabat{สมฺมาทิฏฺฐิสมาทานา,}{สพฺพํ ทุกฺขํ อุปจฺจคุนฺติ\csromansharedfootnote{338dd9d64fd7ee1e}{[มิสฺสินฺคฺ โนเต 0]}.} \\ นวมํ.}}

\csromanmatikaanchor{mtk.102}
\csromantocmarkref{subsection}{10. อุปกฺกิเลสสุตฺต}{mtk.102}
\bookmark[dest={mtk.102},level=2]{10. อุปกฺกิเลสสุตฺต}
\csromantocmarknopage{chapter}{2. ทุติยปณฺณาสก}{mtk.103}
\bookmark[dest={mtk.103},level=0]{2. ทุติยปณฺณาสก}
\csromanheader{2}{10. อุปกฺกิเลส\-สุตฺต}
\proseitem{50}{จตฺตาโรเม ภิกฺขเว\csromansharedfootnote{5ce57ebba7f19ce1}{วิ 4. 492 ปิฏฺเฐปิ.} จนฺทิม\-สูริยานํ อุปกฺกิเลสา, เยหิ อุปกฺกิเลเสหิ อุปกฺกิลิฏฺฐา จนฺทิมสูริยา น ตปนฺติ น ภาสนฺติ น วิโรจนฺติ. กตเม จตฺตาโร? อพฺภา ภิกฺขเว จนฺทิม\-สูริยานํ อุปกฺกิเลสา, เยน อุปกฺกิเลเสน อุปกฺกิลิฏฺฐา จนฺทิมสูริยา น ตปนฺติ น ภาสนฺติ น วิโรจนฺติ.}

\prose{มหิกา ภิกฺขเว จนฺทิม\-สูริยานํ อุปกฺกิเลสา, เยน อุปกฺกิเลเสน อุปกฺกิลิฏฺฐา จนฺทิมสูริยา น ตปนฺติ น ภาสนฺติ น วิโรจนฺติ.}

\prose{ธูโม รโช ภิกฺขเว จนฺทิม\-สูริยานํ อุปกฺกิเลโส, เยน อุปกฺกิเลเสน อุปกฺกิลิฏฺฐา จนฺทิมสูริยา น ตปนฺติ น ภาสนฺติ น วิโรจนฺติ.}

\prose{ราหุ ภิกฺขเว อสุรินฺโท จนฺทิม\-สูริยานํ อุปกฺกิเลโส, เยน อุปกฺกิเลเสน อุปกฺกิลิฏฺฐา จนฺทิมสูริยา น ตปนฺติ น ภาสนฺติ น วิโรจนฺติ. อิเม \csromanfolio{363}โข ภิกฺขเว จตฺตาโร จนฺทิม\-สูริยานํ อุปกฺกิเลสา, เยหิ อุปกฺกิเลเสหิ อุปกฺกิลิฏฺฐา จนฺทิมสูริยา น ตปนฺติ น ภาสนฺติ น วิโรจนฺติ.}

\prose{เอวํเมวํ โข ภิกฺขเว จตฺตาโรเม สมณ\-พฺราหฺมณานํ อุปกฺกิเลสา, เยหิ อุปกฺกิเลเสหิ อุปกฺกิลิฏฺฐา เอเก สมณพฺราหฺมณา น ตปนฺติ น ภาสนฺติ น วิโรจนฺติ. กตเม จตฺตาโร? สนฺติ ภิกฺขเว เอเก สมณพฺราหฺมณา สุรํ ปิวนฺติ เมรยํ, สุราเมรยปานา อปฺปฏิวิรตา. อยํ ภิกฺขเว ปฐโม สมณ\-พฺราหฺมณานํ อุปกฺกิเลโส, เยน อุปกฺกิเลเสน อุปกฺกิลิฏฺฐา เอเก สมณพฺราหฺมณา น ตปนฺติ น ภาสนฺติ น วิโรจนฺติ.}

\prose{สนฺติ ภิกฺขเว เอเก สมณพฺราหฺมณา เมถุนํ ธมฺมํ ปฏิเสวนฺติ, เมถุนสฺมา ธมฺมา อปฺปฏิวิรตา. อยํ ภิกฺขเว ทุติโย สมณ\-พฺราหฺมณานํ อุปกฺกิเลโส, เยน อุปกฺกิเลเสน อุปกฺกิลิฏฺฐา เอเก สมณพฺราหฺมณา น ตปนฺติ น ภาสนฺติ น วิโรจนฺติ.}

\prose{สนฺติ ภิกฺขเว เอเก สมณพฺราหฺมณา ชาตรูป\-รชตํ สาทิยนฺติ, ชาตรูป\-รชต\-ปฏิคฺคหณา อปฺปฏิวิรตา. อยํ ภิกฺขเว ตติโย สมณ\-พฺราหฺมณานํ อุปกฺกิเลโส, เยน อุปกฺกิเลเสน อุปกฺกิลิฏฺฐา เอเก สมณพฺราหฺมณา น ตปนฺติ น ภาสนฺติ น วิโรจนฺติ.}

\prose{สนฺติ ภิกฺขเว เอเก สมณพฺราหฺมณา มิจฺฉาชีเวน ชีวนฺติ, มิจฺฉาชีวา อปฺปฏิวิรตา. อยํ ภิกฺขเว จตุตฺโถ สมณ\-พฺราหฺมณานํ อุปกฺกิเลโส, เยน อุปกฺกิเลเสน อุปกฺกิลิฏฺฐา เอเก สมณพฺราหฺมณา น ตปนฺติ น ภาสนฺติ น วิโรจนฺติ. อิเม โข ภิกฺขเว จตฺตาโร สมณ\-พฺราหฺมณานํ อุปกฺกิเลสา, เยหิ อุปกฺกิเลเสหิ อุปกฺกิลิฏฺฐา เอเก สมณพฺราหฺมณา น ตปนฺติ น ภาสนฺติ น วิโรจนฺตีติ.}

\setlength{\csromangathapagewidth}{0pt}
\setlength{\csromangathawakwidth}{0pt}
\csromangathameasuregroup{ราคโทสปริกฺกิฏฺฐา, \\ อวิชฺชานิวุตา โปสา, \\ สุรํ ปิวนฺติ เมรยํ, \\ รชตํ ชาตรูปญฺจ, \\ มิจฺฉาชีเวน ชีวนฺติ, \\ เอเต อุปกฺกิเลสา วุตฺตา, \\ เยหิ อุปกฺกิเลเสหิ1, \\ น ตปนฺติ น ภาสนฺติ, \\ อนฺธกาเรน โอนทฺธา, \\ วฑฺเฒนฺติ กฏสิํ โฆรํ,}{\csromangathabat{ราคโทสปริกฺกิฏฺฐา,}{เอเก สมณพฺราหฺมณา.} \\ \csromangathabat{อวิชฺชานิวุตา โปสา,}{ปิยรูปาภินนฺทิโน.} \\ \csromangathabat{สุรํ ปิวนฺติ เมรยํ,}{ปฏิเสวนฺติ เมถุนํ.} \\ \csromangathabat{รชตํ ชาตรูปญฺจ,}{สาทิยนฺติ อวิทฺทสู.} \\ \csromangathabat{มิจฺฉาชีเวน ชีวนฺติ,}{เอเก สมณพฺราหฺมณา.} \\ \csromangathabat{เอเต อุปกฺกิเลสา วุตฺตา,}{พุทฺเธนาทิจฺจพนฺธุนา.} \\ \csromangathabat{เยหิ อุปกฺกิเลเสหิ1,}{เอเก สมณพฺราหฺมณา.} \\ \csromangathabat{น ตปนฺติ น ภาสนฺติ,}{อสุทฺธา สรชา มคา.} \\ \csromangathabat{อนฺธกาเรน โอนทฺธา,}{ตณฺหาทาสา สเนตฺติกา.} \\ \csromangathabat{วฑฺเฒนฺติ กฏสิํ โฆรํ,}{อาทิยนฺติ ปุนพฺภวนฺติ.} \\ ทสมํ.}
\csromangathasetleft{ราคโทสปริกฺกิฏฺฐา, \\ อวิชฺชานิวุตา โปสา, \\ สุรํ ปิวนฺติ เมรยํ, \\ รชตํ ชาตรูปญฺจ, \\ มิจฺฉาชีเวน ชีวนฺติ, \\ เอเต อุปกฺกิเลสา วุตฺตา, \\ เยหิ อุปกฺกิเลเสหิ1, \\ น ตปนฺติ น ภาสนฺติ, \\ อนฺธกาเรน โอนทฺธา, \\ วฑฺเฒนฺติ กฏสิํ โฆรํ,}
\csromangathagroupwithclosermidruled{\csromangathastanza{\csromangathabat{ราคโทสปริกฺกิฏฺฐา,}{เอเก สมณพฺราหฺมณา.} \\ \csromangathabat{อวิชฺชานิวุตา โปสา,}{ปิยรูปาภินนฺทิโน.}}\csromangathastanza{\csromangathabat{สุรํ ปิวนฺติ เมรยํ,}{ปฏิเสวนฺติ เมถุนํ.} \\ \csromangathabat{รชตํ ชาตรูปญฺจ,}{สาทิยนฺติ อวิทฺทสู.} \\ \csromangathabat{มิจฺฉาชีเวน ชีวนฺติ,}{เอเก สมณพฺราหฺมณา.}}\csromangathastanza{\csromanfolio{364}\csromangathabat{เอเต อุปกฺกิเลสา วุตฺตา,}{พุทฺเธนา\-ทิจฺจพนฺธุนา.} \\ \csromangathabat{เยหิ อุปกฺกิเลเสหิ1,}{เอเก สมณพฺราหฺมณา.}}\csromangathastanza{\csromangathabat{น ตปนฺติ น ภาสนฺติ,}{อสุทฺธา สรชา มคา.} \\ \csromangathabat{อนฺธกาเรน โอนทฺธา,}{ตณฺหาทาสา สเนตฺติกา.}}\csromangathastanza{\csromangathabat{วฑฺเฒนฺติ กฏสิํ โฆรํ,}{อาทิยนฺติ ปุนพฺภวนฺติ.} \\ ทสมํ.}}{โรหิตสฺสวคฺโค ปญฺจโม.}
\csromancenter{\textbf{ตสฺสุทฺทานํ}}
\setlength{\csromangathapagewidth}{0pt}
\setlength{\csromangathawakwidth}{0pt}
\csromangathameasuregroup{สมาธิปญฺหา เทฺว โกธา, \\ สุวิทูรวิสาขวิปลฺลาสา,}{\csromangathabat{สมาธิปญฺหา เทฺว โกธา,}{โรหิตสฺสาปเร ทุเว.} \\ \csromangathabat{สุวิทูรวิสาขวิปลฺลาสา,}{อุปกฺกิเลเสน เต ทสาติ.} \\ \textbf{ปฐโม ปณฺณาสโก สมตฺโต.}}
\csromangathasetleft{สมาธิปญฺหา เทฺว โกธา, \\ สุวิทูรวิสาขวิปลฺลาสา,}
\csromangathagroup{\csromangathastanza{\csromangathabat{สมาธิปญฺหา เทฺว โกธา,}{โรหิตสฺสาปเร ทุเว.} \\ \csromangathabat{สุวิทูรวิสาขวิปลฺลาสา,}{อุปกฺกิเลเสน เต ทสาติ.} \\ \textbf{ปฐโม ปณฺณาสโก สมตฺโต.}}}

\csromanmatikaanchor{mtk.104}
\csromantocmarknopage{section}{(6) 1. ปุญฺญาภิสนฺทวคฺค}{mtk.104}
\bookmark[dest={mtk.104},level=1]{(6) 1. ปุญฺญาภิสนฺทวคฺค}
\csromanmarkh{(6)}\csromanmarkhii{1. ปุญฺญาภิสนฺท\-วคฺค}\csromanheadernomark{1}{(6) 1. ปุญฺญาภิสนฺท\-วคฺค}
\csromanheader{2}{1. ปฐม\-ปุญฺญาภิสนฺท\-สุตฺต}
\csromanmatikaanchor{mtk.105}
\csromantocmarkref{subsection}{1-2. ปุญฺญาภิสนฺทสุตฺตานิ}{mtk.105}
\bookmark[dest={mtk.105},level=2]{1-2. ปุญฺญาภิสนฺทสุตฺตานิ}
\proseitem{51}{\csromanfolio{365}สาวตฺตินิทานํ. จตฺตาโรเม ภิกฺขเว ปุญฺญาภิสนฺทา กุสลาภิสนฺทา สุขสฺสาหารา โสวคฺคิกา สุขวิปากา สคฺค\-สํวตฺตนิกา อิฏฺฐาย กนฺตาย มนาปาย หิตาย สุขาย สํวตฺตนฺติ. กตเม จตฺตาโร? ยสฺส ภิกฺขเว ภิกฺขุ จีวรํ ปริภุญฺชมาโน อปฺปมาณํ เจโตสมาธิํ อุปสมฺปชฺช วิหรติ, อปฺปมาโณ ตสฺส ปุญฺญาภิสนฺโท กุสลาภิสนฺโท สุขสฺสาหาโร โสวคฺคิโก สุขวิปาโก สคฺค\-สํวตฺตนิโก อิฏฺฐาย กนฺตาย มนาปาย หิตาย สุขาย สํวตฺตติ.}

\prose{ยสฺส ภิกฺขเว ภิกฺขุ ปิณฺฑปาตํ ปริภุญฺชมาโน อปฺปมาณํ เจโตสมาธิํ อุปสมฺปชฺช วิหรติ, อปฺปมาโณ ตสฺส ปุญฺญาภิสนฺโท กุสลาภิสนฺโท สุขสฺสาหาโร โสวคฺคิโก สุขวิปาโก สคฺค\-สํวตฺตนิโก อิฏฺฐาย กนฺตาย มนาปาย หิตาย สุขาย สํวตฺตติ.}

\prose{ยสฺส ภิกฺขเว เสนาสนํ ปริภุญฺชมาโน อปฺปมาณํ เจโตสมาธิํ อุปสมฺปชฺช วิหรติ, อปฺปมาโณ ตสฺส ปุญฺญาภิสนฺโท กุสลาภิสนฺโท สุขสฺสาหาโร โสวคฺคิโก สุขวิปาโก สคฺค\-สํวตฺตนิโก อิฏฺฐาย กนฺตาย มนาปาย หิตาย สุขาย สํวตฺตติ.}

\prose{ยสฺส ภิกฺขเว ภิกฺขุ คิลานปฺปจฺจย\-เภสชฺช\-ปริกฺขารํ ปริภุญฺชมาโน อปฺปมาณํ เจโตสมาธิํ อุปสมฺปชฺช วิหรติ, อปฺปมาโณ ตสฺส ปุญฺญาภิสนฺโท กุสลาภิสนฺโท สุขสฺสาหาโร โสวคฺคิโก สุขวิปาโก สคฺค\-สํวตฺตนิโก อิฏฺฐาย กนฺตาย มนาปาย หิตาย สุขาย สํวตฺตติ. อิเม โข ภิกฺขเว จตฺตาโร ปุญฺญาภิสนฺทา กุสลาภิสนฺทา สุขสฺสาหารา โสวคฺคิกา สุขวิปากา สคฺค\-สํวตฺตนิกา อิฏฺฐาย กนฺตาย มนาปาย หิตาย สุขาย สํวตฺตนฺติ.}

\prose{อิเมหิ จ ปน ภิกฺขเว จตูหิ ปุญฺญา\-ภิสนฺเทหิ กุสลา\-ภิสนฺเทหิ สมนฺนาคตสฺส อริยสาวกสฺส น สุกรํ ปุญฺญสฺส ปมาณํ คเหตุํ\csromansharedfootnote{d7c02471a3f667a5}{คเณตุํ (ก)} \csromanfolio{366}“เอตฺตโก ปุญฺญาภิสนฺโท กุสลาภิสนฺโท สุขสฺสาหาโร โสวคฺคิโก สุขวิปาโก สคฺค\-สํวตฺตนิโก อิฏฺฐาย กนฺตาย มนาปาย หิตาย สุขาย สํวตฺตตี”ติ. อถ โข อสงฺเขฺยยฺโย\csromansharedfootnote{633e6ae118dd3c93}{อสงฺเขยฺโย (สี, สฺยา, กํ, อิ)} อปฺปเมยฺโย มหาปุญฺญกฺขนฺโธ\-เตฺวว สงฺขฺยํ\csromansharedfootnote{f3b84ba968ed1c5b}{สงฺขํ (สี, สฺยา, กํ, อิ)} คจฺฉติ.}

\prose{เสยฺยถาปิ ภิกฺขเว มหาสมุทฺเท น สุกรํ อุทกสฺส ปมาณํ คเหตุํ “เอตฺตกานิ อุทกาฬฺหกานี”ติ วา “เอตฺตกานิ อุทกาฬฺหกสตานี”ติ วา “เอตฺตกานิ อุทกาฬฺหกสหสฺสานี”ติ วา “เอตฺตกานิ อุทกาฬฺหกสตสหสฺสานี”ติ วา. อถ โข อสงฺเขฺยยฺโย อปฺปเมยฺโย มหาอุทกกฺขนฺโธเตฺวว สงฺขฺยํ คจฺฉติ. เอวเมวํ โข ภิกฺขเว อิเมหิ จตูหิ ปุญฺญา\-ภิสนฺเทหิ กุสลา\-ภิสนฺเทหิ สมนฺนาคตสฺส อริยสาวกสฺส น สุกรํ ปุญฺญสฺส ปมาณํ คเหตุํ “เอตฺตโก ปุญฺญาภิสนฺโท กุสลาภิสนฺโท สุขสฺสาหาโร โสวคฺคิโก สุขวิปาโก สคฺค\-สํวตฺตนิโก อิฏฺฐาย กนฺตาย มนาปาย หิตาย สุขาย สํวตฺตตี”ติ. อถ โข อสงฺเขฺยยฺโย อปฺปเมยฺโย มหาปุญฺญกฺขนฺโธ\-เตฺวว สงฺขฺยํ คจฺฉตีติ.}

\setlength{\csromangathapagewidth}{0pt}
\setlength{\csromangathawakwidth}{0pt}
\csromangathameasuregroup{}{มโหทธิํ อปริมิตํ มหาสรํ, \\ พหุเภรวํ รตนวรานมาลยํ\textsuperscript{1}. \\ นชฺโช ยถา นรคณสํฆเสวิตา\textsuperscript{1}, \\ ปุถู สวนฺตี อุปยนฺติ สาครํ. \\ เอวํ นรํ อนฺนทปานวตฺถทํ\textsuperscript{1}, \\ เสยฺยานิสชฺชตฺถรณสฺส ทายกํ. \\ ปุญฺญสฺส ธารา อุปยนฺติ ปณฺฑิตํ, \\ นชฺโช ยถา วาริวหาว สาครนฺติ.ปฐมํ.}
\csromangathameasurewak{มโหทธิํ อปริมิตํ มหาสรํ, \\ พหุเภรวํ รตนวรานมาลยํ\textsuperscript{1}. \\ นชฺโช ยถา นรคณสํฆเสวิตา\textsuperscript{1}, \\ ปุถู สวนฺตี อุปยนฺติ สาครํ. \\ เอวํ นรํ อนฺนทปานวตฺถทํ\textsuperscript{1}, \\ เสยฺยานิสชฺชตฺถรณสฺส ทายกํ. \\ ปุญฺญสฺส ธารา อุปยนฺติ ปณฺฑิตํ, \\ นชฺโช ยถา วาริวหาว สาครนฺติ.ปฐมํ.}
\setlength{\csromangathaleftcol}{0pt}
\csromangathagroup{\csromangathastanza{\csromangathawak{มโหทธิํ อปริมิตํ มหาสรํ, \\ พหุเภรวํ รตนวรานมาลยํ\csromansharedfootnote{b7f36ca6ec7d6205}{รตนคณานมาลยํ (สี, สฺยา, กํ, อิ)}. \\ นชฺโช ยถา นรคณสํฆเสวิตา\csromansharedfootnote{cf7489296a1ba68f}{มจฺฉคณสํฆเสวิตา (สฺยา, กํ)}, \\ ปุถู สวนฺตี อุปยนฺติ สาครํ.}}\csromangathastanza{\csromangathawak{เอวํ นรํ อนฺนทปานวตฺถทํ\csromansharedfootnote{d404ed2dc571335f}{อนฺนปานวตฺถํ (ก)}, \\ เสยฺยานิสชฺชตฺถรณสฺส ทายกํ. \\ ปุญฺญสฺส ธารา อุปยนฺติ ปณฺฑิตํ, \\ นชฺโช ยถา วาริวหาว สาครนฺติ.  ปฐมํ.}}}

\csromanheader{2}{2. ทุติย\-ปุญฺญาภิสนฺท\-สุตฺต}
\proseitem{52}{จตฺตาโรเม ภิกฺขเว ปุญฺญาภิสนฺทา กุสลาภิสนฺทา สุขสฺสาหารา โสวคฺคิกา สุขวิปากา สคฺค\-สํวตฺตนิกา อิฏฺฐาย กนฺตาย มนาปาย หิตาย สุขาย สํวตฺตนฺติ. กตเม จตฺตาโร? อิธ ภิกฺขเว \csromanfolio{367}อริยสาวโก พุทฺเธ อเวจฺจ\-ปฺปสาเทน สมนฺนาคโต โหติ “อิติปิ โส ภควา อรหํ สมฺมาสมฺพุทฺโธ วิชฺชา\-จรณ\-สมฺปนฺโน สุคโต โลกวิทู อนุตฺตโร ปุริส\-ทมฺม\-สารถิ สตฺถา เทวมนุสฺสานํ พุทฺโธ ภควา”ติ. อยํ ภิกฺขเว ปฐโม ปุญฺญาภิสนฺโท กุสลาภิสนฺโท สุขสฺสาหาโร โสวคฺคิโก สุขวิปาโก สคฺค\-สํวตฺตนิโก อิฏฺฐาย กนฺตาย มนาปาย หิตาย สุขาย สํวตฺตติ.}

\prose{ปุน จปรํ ภิกฺขเว อริยสาวโก ธมฺเม อเวจฺจ\-ปฺปสาเทน สมนฺนาคโต โหติ “สฺวากฺขาโต ภควตา ธมฺโม สนฺทิฏฺฐิโก อกาลิโก เอหิปสฺสิโก โอปเนยฺยิโก ปจฺจตฺตํ เวทิตพฺโพ วิญฺญูหี”ติ. อยํ ภิกฺขเว ทุติโย ปุญฺญาภิสนฺโท กุสลาภิสนฺโท สุขสฺสาหาโร โสวคฺคิโก สุขวิปาโก สคฺค\-สํวตฺตนิโก อิฏฺฐาย กนฺตาย มนาปาย หิตาย สุขาย สํวตฺตติ.}

\prose{ปุน จปรํ ภิกฺขเว อริยสาวโก สํเฆ อเวจฺจ\-ปฺปสาเทน สมนฺนาคโต โหติ “สุปฺปฏิปนฺโน ภควโต สาวกสํโฆ อุชุปฺปฏิปนฺโน ภควโต สาวกสํโฆ ญายปฺปฏิปนฺโน ภควโต สาวกสํโฆ สามีจิ\-ปฺปฏิปนฺโน ภควโต สาวกสํโฆ, ยทิทํ จตฺตาริ ปุริสยุคานิ อฏฺฐ ปุริสปุคฺคลา, เอส ภควโต สาวกสํโฆ อาหุเนยฺโย ปาหุเนยฺโย ทกฺขิเณยฺโย อญฺชลิกรณีโย อนุตฺตรํ ปุญฺญกฺเขตฺตํ โลกสฺสา”ติ. อยํ ภิกฺขเว ตติโย ปุญฺญาภิสนฺโท กุสลาภิสนฺโท สุขสฺสาหาโร โสวคฺคิโก สุขวิปาโก สคฺค\-สํวตฺตนิโก อิฏฺฐาย กนฺตาย มนาปาย หิตาย สุขาย สํวตฺตติ.}

\prose{ปุน จปรํ ภิกฺขเว อริยสาวโก อริยกนฺเตหิ สีเลหิ สมนฺนาคโต โหติ อขณฺเฑหิ อจฺฉิทฺเทหิ อสพเลหิ อกมฺมาเสหิ ภุชิสฺเสหิ วิญฺญุ\-ปฺปสตฺเถหิ อปรามฏฺเฐหิ สมาธิ\-สํวตฺตนิเกหิ. อยํ ภิกฺขเว จตุตฺโถ ปุญฺญาภิสนฺโท กุสลาภิสนฺโท สุขสฺสาหาโร โสวคฺคิโก สุขวิปาโก สคฺค\-สํวตฺตนิโก อิฏฺฐาย กนฺตาย มนาปาย หิตาย สุขาย สํวตตฺตติ. อิเม โข ภิกฺขเว จตฺตาโร ปุญฺญาภิสนฺทา กุสลาภิสนฺทา สุขสฺสาหารา โสวคฺคิกา สุขวิปากา สคฺค\-สํวตฺตนิกา อิฏฺฐาย กนฺตาย มนาปาย หิตาย สุขาย สํวตฺตนฺตีติ.}

\csromanmatikaanchor{mtk.106}
\csromantocmarkref{subsection}{3-4. สํวาสสุตฺตานิ}{mtk.106}
\bookmark[dest={mtk.106},level=2]{3-4. สํวาสสุตฺตานิ}
\setlength{\csromangathapagewidth}{0pt}
\setlength{\csromangathawakwidth}{0pt}
\csromangathameasuregroup{\textsuperscript{*}ยสฺส สทฺธา ตถาคเต, \\ สีลญฺจ ยสฺส กลฺยาณํ, \\ สํเฆ ปสาโท ยสฺสตฺถิ, \\ “อทลิทฺโท”ติ ตํ อาหุ, \\ ตสฺมา สทฺธญฺจ สีลญฺจ, \\ อนุยุญฺเชถ เมธาวี,}{\csromangathabat{\textsuperscript{*}ยสฺส สทฺธา ตถาคเต,}{อจลา สุปฺปติฏฺฐิตา.} \\ \csromangathabat{สีลญฺจ ยสฺส กลฺยาณํ,}{อริยกนฺตํ ปสํสิตํ.} \\ \csromangathabat{สํเฆ ปสาโท ยสฺสตฺถิ,}{อุชุภูตญฺจ ทสฺสนํ.} \\ \csromangathabat{“อทลิทฺโท”ติ ตํ อาหุ,}{อโมฆํ ตสฺส ชีวิตํ.} \\ \csromangathabat{ตสฺมา สทฺธญฺจ สีลญฺจ,}{ปสาทํ ธมฺมทสฺสนํ.} \\ \csromangathabat{อนุยุญฺเชถ เมธาวี,}{สรํ พุทฺธาน สาสนนฺติ.ทุติยํ.}}
\csromangathasetleft{\textsuperscript{*}ยสฺส สทฺธา ตถาคเต, \\ สีลญฺจ ยสฺส กลฺยาณํ, \\ สํเฆ ปสาโท ยสฺสตฺถิ, \\ “อทลิทฺโท”ติ ตํ อาหุ, \\ ตสฺมา สทฺธญฺจ สีลญฺจ, \\ อนุยุญฺเชถ เมธาวี,}
\csromangathagroup{\csromangathastanza{\csromanfolio{368}\csromangathabat{\csromansymbolfootnote{*}{อํ 2. 47 ปิฏฺเฐปิ.}ยสฺส สทฺธา ตถาคเต,}{อจลา สุปฺปติฏฺฐิตา.} \\ \csromangathabat{สีลญฺจ ยสฺส กลฺยาณํ,}{อริยกนฺตํ ปสํสิตํ.}}\csromangathastanza{\csromangathabat{สํเฆ ปสาโท ยสฺสตฺถิ,}{อุชุภูตญฺจ ทสฺสนํ.} \\ \csromangathabat{“อทลิทฺโท”ติ ตํ อาหุ,}{อโมฆํ ตสฺส ชีวิตํ.}}\csromangathastanza{\csromangathabat{ตสฺมา สทฺธญฺจ สีลญฺจ,}{ปสาทํ ธมฺมทสฺสนํ.} \\ \csromangathabat{อนุยุญฺเชถ เมธาวี,}{สรํ พุทฺธาน สาสนนฺติ.  ทุติยํ.}}}

\csromanheader{2}{3. ปฐม\-สํวาส\-สุตฺต}
\proseitem{53}{เอกํ สมยํ ภควา อนฺตรา จ มธุรํ, อนฺตรา จ เวรญฺชํ อทฺธาน\-มคฺค\-ปฺปฏิปนฺโน โหติ, สมฺพหุลาปิ โข คหปตี จ คหปตานิโย จ อนฺตรา จ มธุรํ, อนฺตรา จ เวรญฺชํ อทฺธาน\-มคฺค\-ปฺปฏิปนฺนา โหนฺติ. อถ โข ภควา มคฺคา โอกฺกมฺม อญฺญตรสฺมิํ รุกฺขมูเล ( )\csromansharedfootnote{3cd65102829efe49}{(ปญฺญตฺเต อาสเน) (อิ, ก)} นิสีทิ. อทฺทสํสุ โข คหปตี จ คหปตานิโย จ ภควนฺตํ อญฺญตรสฺมิํ รุกฺขมูเล นิสินฺนํ, ทิสฺวา เยน ภควา เตนุ\-ปสงฺกมิํสุ, อุปสงฺกมิตฺวา ภควนฺตํ อภิวาเทตฺวา เอกมนฺตํ นิสีทิํสุ, เอกมนฺตํ นิสินฺเน โข เต คหปตี จ คหปตานิโย จ ภควา เอตทโวจ\csromandash{}}

\prose{จตฺตาโรเม คหปตโย สํวาสา. กตเม จตฺตาโร? ฉโว ฉวาย สทฺธิํ สํวสติ, ฉโว เทวิยา สทฺธิํ สํวสติ, เทโว ฉวาย สทฺธิํ สํวสติ, เทโว เทวิยา สทฺธิํ สํวสติ.}

\prose{กถญฺจ คหปตโย ฉโว ฉวาย สทฺธิํ สํวสติ? อิธ คหปตโย สามิโก โหติ ปาณาติปาตี อทินฺนาทายี กาเมสุ\-มิจฺฉาจารี มุสาวาที สุรา\-เมรย\-มชฺช\-ปมาท\-ฏฺฐายี ทุสฺสีโล ปาปธมฺโม มจฺเฉรมล\-ปริยุฏฺฐิเตน เจตสา อคารํ อชฺฌาวสติ อกฺโกสก\-ปริภาสโก สมณ\-พฺราหฺมณานํ, ภริยาปิสฺส โหติ ปาณาติปาตินี อทินฺนาทายินี กาเมสุ\-มิจฺฉาจารินี มุสาวาทินี สุรา\-เมรย\-มชฺช\-ปมาท\-ฏฺฐายินี ทุสฺสีลา ปาปธมฺมา มจฺเฉรมล\-ปริยุฏฺฐิเตน เจตสา อคารํ อชฺฌาวสติ อกฺโกสิก\-ปริภาสิกา สมณ\-พฺราหฺมณานํ. เอวํ โข คหปตโย ฉโว ฉวาย สทฺธิํ สํวสติ.}

\chapterheadpageread{(6) ปุญฺญาภิสนฺท\-วคฺค}
\prose{\csromanfolio{369}กถญฺจ คหปตโย ฉโว เทวิยา สทฺธิํ สํวสติ? อิธ คหปตโย สามิโก โหติ ปาณาติปาตี อทินฺนาทายี กาเมสุ\-มิจฺฉาจารี มุสาวาที สุรา\-เมรย\-มชฺช\-ปมาท\-ฏฺฐายี ทุสฺสีโล ปาปธมฺโม มจฺเฉรมล\-ปริยุฏฺฐิเตน เจตสา อคารํ อชฺฌาวสติ อกฺโกสก\-ปริภาสโก สมณ\-พฺราหฺมณานํ, ภริยา ขฺวสฺส โหติ ปาณาติปาตา ปฏิวิรตา อทินฺนาทานา ปฏิวิรตา กาเมสุ\-มิจฺฉาจารา ปฏิวิรตา มุสาวาทา ปฏิวิรตา สุรา\-เมรย\-มชฺช\-ปมาทฏฺฐานา ปฏิวิรตา สีลวตี กลฺยาณธมฺมา วิคต\-มล\-มจฺเฉเรน เจตสา อคารํ อชฺฌาวสติ อนกฺโกสิกปริภาสิกา สมณ\-พฺราหฺมณานํ. เอวํ โข คหปตโย ฉโว เทวิยา สทฺธิํ สํวสติ.}

\prose{กถญฺจ คหปตโย เทโว ฉวาย สทฺธิํ สํวสติ? อิธ คหปตโย สามิโก โหติ ปาณาติปาตา ปฏิวิรโต อทินฺนาทานา ปฏิวิรโต กาเมสุ\-มิจฺฉาจารา ปฏิวิรโต มุสาวาทา ปฏิวิรโต สุรา\-เมรย\-มชฺช\-ปมาทฏฺฐานา ปฏิวิรโต สีลวา กลฺยาณธมฺโม วิคต\-มล\-มจฺเฉเรน เจตสา อคารํ อชฺฌาวสติ อนกฺโกสกปริภาสโก สมณ\-พฺราหฺมณานํ, ภริยา ขฺวสฺส โหติ ปาณาติปาตินี ฯเปฯ สุรา\-เมรย\-มชฺช\-ปมาท\-ฏฺฐายินี ทุสฺสีลา ปาปธมฺมา มจฺเฉรมล\-ปริยุฏฺฐิเตน เจตสา อคารํ อชฺฌาวสติ อกฺโกสิก\-ปริภาสิกา สมณ\-พฺราหฺมณานํ. เอวํ โข คหปตโย เทโว ฉวาย สทฺธิํ สํวสติ.}

\prose{กถญฺจ คหปตโย เทโว เทวิยา สทฺธิํ สํวสติ? อิธ คหปตโย สามิโก โหติ ปาณาติปาตา ปฏิวิรโต ฯเปฯ สีลวา กลฺยาณธมฺโม วิคต\-มล\-มจฺเฉเรน เจตสฺส อคารํ อชฺฌาวสติ อนกฺโกสกปริภาสโก สมณ\-พฺราหฺมณานํ, ภริยาปิสฺส โหติ ปาณาติปาตา ปฏิวิรตา ฯเปฯ สุรา\-เมรย\-มชฺช\-ปมาทฏฺฐานา ปฏิวิรตา สีลวตี กลฺยาณธมฺมา วิคต\-มล\-มจฺเฉเรน เจตสา อคารํ อชฺฌาวสติ อนกฺโกสิกปริภาสิกา สมณ\-พฺราหฺมณานํ. เอวํ โข คหปตโย เทโว เทวิยา สทฺธิํ สํวสติ. อิเม โข คหปตโย จตฺตาโร สํวาสาติ.}

\setlength{\csromangathapagewidth}{0pt}
\setlength{\csromangathawakwidth}{0pt}
\csromangathameasuregroup{อุโภ จ โหนฺติ ทุสฺสีลา, \\ เต โหนฺติ ชานิปตโย, \\ สามิโก โหติ ทุสฺสีโล, \\ ภริยา สีลวตี โหติ, \\ สาปิ เทวี สํวสติ, \\ สามิโก สีลวา โหติ, \\ ภริยา โหติ ทุสฺสีลา, \\ สาปิ ฉวา สํวสติ, \\ อุโภ สทฺธา วทญฺญู จ, \\ เต โหนฺติ ชานิปตโย, \\ อตฺถา’สํ ปจุรา โหนฺติ, \\ อมิตฺตา ทุมฺมนา โหนฺติ, \\ อิธ ธมฺมํ จริตฺวาน, \\ นนฺทิโน เทวโลกสฺมิํ,}{\csromangathabat{อุโภ จ โหนฺติ ทุสฺสีลา,}{กทริยา ปริภาสกา.} \\ \csromangathabat{เต โหนฺติ ชานิปตโย,}{ฉวา สํวาสมาคตา.} \\ \csromangathabat{สามิโก โหติ ทุสฺสีโล,}{กทริโย ปริภาสโก.} \\ \csromangathabat{ภริยา สีลวตี โหติ,}{วทญฺญู วีตมจฺฉรา.} \\ \csromangathabat{สาปิ เทวี สํวสติ,}{ฉเวน ปตินา สห.} \\ \csromangathabat{สามิโก สีลวา โหติ,}{วทญฺญู วีตมจฺฉโร.} \\ \csromangathabat{ภริยา โหติ ทุสฺสีลา,}{กทริยา ปริภาสิกา.} \\ \csromangathabat{สาปิ ฉวา สํวสติ,}{เทเวน ปตินา สห.} \\ \csromangathabat{อุโภ สทฺธา วทญฺญู จ,}{สญฺญตา ธมฺมชีวิโน.} \\ \csromangathabat{เต โหนฺติ ชานิปตโย,}{อญฺญมญฺญํ ปิยํวทา.} \\ \csromangathabat{อตฺถา’สํ ปจุรา โหนฺติ,}{ผาสุกํ\textsuperscript{1} อุปชายติ.} \\ \csromangathabat{อมิตฺตา ทุมฺมนา โหนฺติ,}{อุภินฺนํ สมสีลินํ.} \\ \csromangathabat{อิธ ธมฺมํ จริตฺวาน,}{สมสีลพฺพตา อุโภ.} \\ \csromangathabat{นนฺทิโน เทวโลกสฺมิํ,}{โมทนฺติ กามกามิโนติ.ตติยํ.}}
\csromangathasetleft{อุโภ จ โหนฺติ ทุสฺสีลา, \\ เต โหนฺติ ชานิปตโย, \\ สามิโก โหติ ทุสฺสีโล, \\ ภริยา สีลวตี โหติ, \\ สาปิ เทวี สํวสติ, \\ สามิโก สีลวา โหติ, \\ ภริยา โหติ ทุสฺสีลา, \\ สาปิ ฉวา สํวสติ, \\ อุโภ สทฺธา วทญฺญู จ, \\ เต โหนฺติ ชานิปตโย, \\ อตฺถา’สํ ปจุรา โหนฺติ, \\ อมิตฺตา ทุมฺมนา โหนฺติ, \\ อิธ ธมฺมํ จริตฺวาน, \\ นนฺทิโน เทวโลกสฺมิํ,}
\csromangathagroup{\csromangathastanza{\csromangathabat{อุโภ จ โหนฺติ ทุสฺสีลา,}{กทริยา ปริภาสกา.} \\ \csromangathabat{เต โหนฺติ ชานิปตโย,}{ฉวา สํวาสมาคตา.}}\csromangathastanza{\csromanfolio{370}\csromangathabat{สามิโก โหติ ทุสฺสีโล,}{กทริโย ปริภาสโก.} \\ \csromangathabat{ภริยา สีลวตี โหติ,}{วทญฺญู วีตมจฺฉรา.}}\csromangathastanza{\csromangathabat{สาปิ เทวี สํวสติ,}{ฉเวน ปตินา สห.} \\ \csromangathabat{สามิโก สีลวา โหติ,}{วทญฺญู วีตมจฺฉโร.}}\csromangathastanza{\csromangathabat{ภริยา โหติ ทุสฺสีลา,}{กทริยา ปริภาสิกา.} \\ \csromangathabat{สาปิ ฉวา สํวสติ,}{เทเวน ปตินา สห.}}\csromangathastanza{\csromangathabat{อุโภ สทฺธา วทญฺญู จ,}{สญฺญตา ธมฺมชีวิโน.} \\ \csromangathabat{เต โหนฺติ ชานิปตโย,}{อญฺญมญฺญํ ปิยํวทา.}}\csromangathastanza{\csromangathabat{อตฺถา’สํ ปจุรา โหนฺติ,}{ผาสุกํ\csromansharedfootnote{5bd73c4fece8828c}{ผาสตฺตํ (สี), วาสตฺถํ (อิ)} อุปชายติ.} \\ \csromangathabat{อมิตฺตา ทุมฺมนา โหนฺติ,}{อุภินฺนํ สมสีลินํ.}}\csromangathastanza{\csromangathabat{อิธ ธมฺมํ จริตฺวาน,}{สมสีลพฺพตา อุโภ.} \\ \csromangathabat{นนฺทิโน เทวโลกสฺมิํ,}{โมทนฺติ กามกามิโนติ.  ตติยํ.}}}

\csromanheader{2}{4. ทุติย\-สํวาส\-สุตฺต}
\proseitem{54}{จตฺตาโรเม ภิกฺขเว สํวาสา. กตเม จตฺตาโร? ฉโว ฉวาย สทฺธิํ สํวสติ, ฉโว เทวิยา สทฺธิํ สํวสติ, เทโว ฉวาย สทฺธิํ สํวสติ, เทโว เทวิยา สทฺธิํ สํวสติ.}

\prose{กถญฺจ ภิกฺขเว ฉโว ฉวาย สทฺธิํ สํวสติ? อิธ ภิกฺขเว สามิโก โหติ ปาณาติปาตี อทินฺนาทายี กาเมสุ\-มิจฺฉาจารี มุสาวาที ปิสุณวาโจ ผรุสวาโจ สมฺผปฺปลาปี อภิชฺฌาลุ พฺยาปนฺนจิตฺโต มิจฺฉาทิฏฺฐิโก ทุสฺสีโล ปาปธมฺโม มจฺเฉรมล\-ปริยุฏฺฐิเตน เจตสา อคารํ อชฺฌาวสติ อกฺโกสก\-ปริภาสโก สมณ\-พฺราหฺมณานํ, ภริยาปิสฺส โหติ ปาณาติปาตินี อทินฺนาทายินี กาเมสุ\-มิจฺฉาจารินี มุสาวาทินี ปิสุณวาจา ผรุสวาจา สมฺผ\-ปฺปลาปินี อภิชฺฌาลุนี พฺยาปนฺนจิตฺตา มิจฺฉาทิฏฺฐิกา ทุสฺสีลา ปาปธมฺมา มจฺเฉรมล\-ปริยุฏฺฐิเตน เจตสา อคารํ อชฺฌาวสติ อกฺโกสิก\-ปริภาสิกา สมณ\-พฺราหฺมณานํ. เอวํ โข ภิกฺขเว ฉโว ฉวาย สทฺธิํ สํวสติ.}

\chapterheadpageread{(6) ปุญฺญาภิสนฺท\-วคฺค}
\prose{\csromanfolio{371}กถญฺจ ภิกฺขเว ฉโว เทวิยา สทฺธิํ สํวสติ? อิธ ภิกฺขเว สามิโก โหติ ปาณาติปาตี ฯเปฯ มิจฺฉาทิฏฺฐิโก ทุสฺสีโล ปาปธมฺโม มจฺเฉรมล\-ปริยุฏฺฐิเตน เจตสา อคารํ อชฺฌาวสติ อกฺโกสก\-ปริภาสโก สมณ\-พฺราหฺมณานํ, ภริยา ขฺวสฺส โหติ ปาณาติปาตา ปฏิวิรตา อทินฺนาทานา ปฏิวิรตา กาเมสุ\-มิจฺฉาจารา ปฏิวิรตา มุสาวาทา ปฏิวิรตา ปิสุณาย วาจาย ปฏิวิรตา ผรุสาย วาจาย ปฏิวิรตา สมฺผปฺปลาปา ปฏิวิรตา อนภิชฺฌาลุนี อพฺยาปนฺนจิตฺตา สมฺมาทิฏฺฐิกา สีลวตี กลฺยาณธมฺมา วิคต\-มล\-มจฺเฉเรน เจตสา อคารํ อชฺฌาวสติ อนกฺโกสิกปริภาสิกา สมณ\-พฺราหฺมณานํ. เอวํ โข ภิกฺขเว ฉโว เทวิยา สทฺธิํ สํวสติ.}

\prose{กถญฺจ ภิกฺขเว เทโว ฉวาย สทฺธิํ สํวสติ? อิธ ภิกฺขเว สามิโก โหติ ปาณาติปาตา ปฏิวิรโต อทินฺนาทานา ปฏิวิรโต กาเมสุ\-มิจฺฉาจารา ปฏิวิรโต มุสาวาทา ปฏิวิรโต ปิสุณาย วาจาย ปฏิวิรโต ผรุสาย วาจาย ปฏิวิรโต สมฺผปฺปลาปา ปฏิวิรโต อนภิชฺฌาลุ อพฺยาปนฺนจิตฺโต สมฺมาทิฏฺฐิโก สีลวา กลฺยาณธมฺโม วิคต\-มล\-มจฺเฉเรน เจตสา อคารํ อชฺฌาวสติ อนกฺโกสกปริภาสโก สมณ\-พฺราหฺมณานํ, ภริยา ขฺวสฺส โหติ ปาณาติปาตินี ฯเปฯ มิจฺฉาทิฏฺฐิกา ทุสฺสีลา ปาปธมฺมา มจฺเฉรมล\-ปริยุฏฺฐิเตน เจตสา อคารํ อชฺฌาวสติ อกฺโกสิก\-ปริภาสิกา สมณ\-พฺราหฺมณานํ. เอวํ โข ภิกฺขเว เทโว ฉวาย สทฺธิํ สํวสติ.}

\prose{กถญฺจ ภิกฺขเว เทโว เทวิยา สทฺธิํ สํวสติ, อิธ ภิกฺขเว สามิโก โหติ ปาณาติปาตา ปฏิวิรโต ฯเปฯ สมฺมาทิฏฺฐิโก สีลวา กลฺยาณธมฺโม วิคต\-มล\-มจฺเฉเรน เจตสา อคารํ อชฺฌาวสติ อนกฺโกสกปริภาสโก สมณ\-พฺราหฺมณานํ, ภริยาปิสฺส โหติ ปาณาติปาตา ปฏิวิรตา ฯเปฯ สมฺมาทิฏฺฐิกา สีลวตี กลฺยาณธมฺมา วิคต\-มล\-มจฺเฉเรน เจตสา อคารํ อชฺฌาวสติ อนกฺโกสิกปริภาสิกา สมณ\-พฺราหฺมณานํ. เอวํ โข ภิกฺขเว เทโว เทวิยา สทฺธิํ สํวสติ. อิเม โข ภิกฺขเว จตฺตาโร สํวาสาติ.}

\csromanmatikaanchor{mtk.107}
\csromantocmarkref{subsection}{5-6. สมชีวีสุตฺตานิ}{mtk.107}
\bookmark[dest={mtk.107},level=2]{5-6. สมชีวีสุตฺตานิ}
\setlength{\csromangathapagewidth}{0pt}
\setlength{\csromangathawakwidth}{0pt}
\csromangathameasuregroup{อุโภ จ โหนฺติ ทุสฺสีลา, \\ เต โหนฺติ ชานิปตโย, \\ สามิโก โหติ ทุสฺสีโล, \\ ภริยา สีลวตี โหติ, \\ สาปิ เทวี สํวสติ, \\ สามิโก สีลวา โหติ, \\ ภริยา โหติ ทุสฺสีลา, \\ สาปิ ฉวา สํวสติ, \\ อุโภ สทฺธา วทญฺญู จ, \\ เต โหนฺติ ชานิปตโย, \\ อตฺถา’สํ ปจุรา โหนฺติ, \\ อมิตฺตา ทุมฺมนา โหนฺติ, \\ อิธ ธมฺมํ จริตฺวาน, \\ นนฺทิโน เทวโลกสฺมิํ,}{\csromangathabat{อุโภ จ โหนฺติ ทุสฺสีลา,}{กทริยา ปริภาสกา.} \\ \csromangathabat{เต โหนฺติ ชานิปตโย,}{ฉวา สํวาสมาคตา.} \\ \csromangathabat{สามิโก โหติ ทุสฺสีโล,}{กทริโย ปริภาสโก.} \\ \csromangathabat{ภริยา สีลวตี โหติ,}{วทญฺญู วีตมจฺฉรา.} \\ \csromangathabat{สาปิ เทวี สํวสติ,}{ฉเวน ปตินา สห.} \\ \csromangathabat{สามิโก สีลวา โหติ,}{วทญฺญู วีตมจฺฉโร.} \\ \csromangathabat{ภริยา โหติ ทุสฺสีลา,}{กทริยา ปริภาสิกา.} \\ \csromangathabat{สาปิ ฉวา สํวสติ,}{เทเวน ปตินา สห.} \\ \csromangathabat{อุโภ สทฺธา วทญฺญู จ,}{สญฺญตา ธมฺมชีวิโน.} \\ \csromangathabat{เต โหนฺติ ชานิปตโย,}{อญฺญมญฺญํ ปิยํวทา.} \\ \csromangathabat{อตฺถา’สํ ปจุรา โหนฺติ,}{ผาสุกํ อุปชายติ.} \\ \csromangathabat{อมิตฺตา ทุมฺมนา โหนฺติ,}{อุภินฺนํ สมสีลินํ.} \\ \csromangathabat{อิธ ธมฺมํ จริตฺวาน,}{สมสีลพฺพตา อุโภ.} \\ \csromangathabat{นนฺทิโน เทวโลกสฺมิํ,}{โมทนฺติ กามกามิโนติ.จตุตฺถํ.}}
\csromangathasetleft{อุโภ จ โหนฺติ ทุสฺสีลา, \\ เต โหนฺติ ชานิปตโย, \\ สามิโก โหติ ทุสฺสีโล, \\ ภริยา สีลวตี โหติ, \\ สาปิ เทวี สํวสติ, \\ สามิโก สีลวา โหติ, \\ ภริยา โหติ ทุสฺสีลา, \\ สาปิ ฉวา สํวสติ, \\ อุโภ สทฺธา วทญฺญู จ, \\ เต โหนฺติ ชานิปตโย, \\ อตฺถา’สํ ปจุรา โหนฺติ, \\ อมิตฺตา ทุมฺมนา โหนฺติ, \\ อิธ ธมฺมํ จริตฺวาน, \\ นนฺทิโน เทวโลกสฺมิํ,}
\csromangathagroup{\csromangathastanza{\csromangathabat{อุโภ จ โหนฺติ ทุสฺสีลา,}{กทริยา ปริภาสกา.} \\ \csromangathabat{เต โหนฺติ ชานิปตโย,}{ฉวา สํวาสมาคตา.}}\csromangathastanza{\csromanfolio{372}\csromangathabat{สามิโก โหติ ทุสฺสีโล,}{กทริโย ปริภาสโก.} \\ \csromangathabat{ภริยา สีลวตี โหติ,}{วทญฺญู วีตมจฺฉรา.}}\csromangathastanza{\csromangathabat{สาปิ เทวี สํวสติ,}{ฉเวน ปตินา สห.} \\ \csromangathabat{สามิโก สีลวา โหติ,}{วทญฺญู วีตมจฺฉโร.}}\csromangathastanza{\csromangathabat{ภริยา โหติ ทุสฺสีลา,}{กทริยา ปริภาสิกา.} \\ \csromangathabat{สาปิ ฉวา สํวสติ,}{เทเวน ปตินา สห.}}\csromangathastanza{\csromangathabat{อุโภ สทฺธา วทญฺญู จ,}{สญฺญตา ธมฺมชีวิโน.} \\ \csromangathabat{เต โหนฺติ ชานิปตโย,}{อญฺญมญฺญํ ปิยํวทา.}}\csromangathastanza{\csromangathabat{อตฺถา’สํ ปจุรา โหนฺติ,}{ผาสุกํ อุปชายติ.} \\ \csromangathabat{อมิตฺตา ทุมฺมนา โหนฺติ,}{อุภินฺนํ สมสีลินํ.}}\csromangathastanza{\csromangathabat{อิธ ธมฺมํ จริตฺวาน,}{สมสีลพฺพตา อุโภ.} \\ \csromangathabat{นนฺทิโน เทวโลกสฺมิํ,}{โมทนฺติ กามกามิโนติ.  จตุตฺถํ.}}}

\csromanheader{2}{5. ปฐม\-สมชีวี\-สุตฺต}
\proseitem{55}{เอวํ เม สุตํ\csromandash{}เอกํ สมยํ ภควา ภคฺเคสุ วิหรติ สุสุมารคิเร\csromansharedfootnote{906a675938ec4ab1}{สุํสุมารคิเร (สี, สฺยา, กํ, อิ)} เภสกฬาวเน\csromansharedfootnote{8b1ec7303360269d}{เภสกลาวเน (สี, อิ, ก)} มิคทาเย. อถ โข ภควา ปุพฺพณฺหสมยํ นิวาเสตฺวา ปตฺต\-จีวรมา\-ทาย เยน นกุลปิตุโน คหปติสฺส นิเวสนํ เตนุปสงฺกมิ, อุปสงฺกมิตฺวา ปญฺญตฺเต อาสเน นิสีทิ. อถ โข นกุลปิตา จ คหปติ นกุลมาตา จ คหปตานี เยน ภควา เตนุ\-ปสงฺกมิํสุ อุปสงฺกมิตฺวา ภควนฺตํ อภิวาเทตฺวา เอกมนฺตํ นิสีทิํสุ, เอกมนฺตํ นิสินฺโน โข นกุลปิตา คหปติ ภควนฺตํ เอตทโวจ\csromandash{}}

\prose{ยโต เม ภนฺเต นกุลมาตา คหปตานี ทหรสฺเสว ทหรา อานีตา, นาภิชานามิ นกุลมาตรํ คหปตานิํ มนสาปิ อติจริตา, กุโต ปน กาเยน. อิจฺเฉยฺยาม มยํ ภนฺเต ทิฏฺเฐ เจว ธมฺเม อญฺญมญฺญํ ปสฺสิตุํ, อภิสมฺปรายญฺจ อญฺญมญฺญํ ปสฺสิตุนฺติ. นกุลมาตาปิ โข คหปตานี ภควนฺตํ เอตทโวจ “ยโตหํ ภนฺเต นกุลปิตุโน คหปติสฺส ทหรสฺเสว ทหรา อานีตา, นาภิชานามิ นกุลปิตรํ คหปติํ มนสาปิ อติจริตา, กุโต ปน กาเยน. อิจฺเฉยฺยาม มยํ \csromanfolio{373}ภนฺเต ทิฏฺเฐ เจว ธมฺเม อญฺญมญฺญํ ปสฺสิตุํ, อภิสมฺปรายญฺจ อญฺญมญฺญํ ปสฺสิตุนฺ”ติ.}

\prose{อากงฺเขยฺยุํ เจ คหปตโย อุโภ ชานิปตโย ทิฏฺเฐ เจว ธมฺเม อญฺญมญฺญํ ปสฺสิตุํ, อภิสมฺปรายญฺจ อญฺญมญฺญํ ปสฺสิตุํ, อุโภว\csromansharedfootnote{b1872fe6d8b2a673}{อุโภ จ (สี, อิ)} อสฺสุ สมสทฺธา สมสีลา สมจาคา สมปญฺญา, เต ทิฏฺเฐ เจว ธมฺเม อญฺญมญฺญํ ปสฺสนฺติ, อภิสมฺปรายญฺจ อญฺญมญฺญํ ปสฺสนฺตีติ\csromansharedfootnote{69328cb87fbf9f0e}{ปสฺสิสฺสนฺตีติ (ก)}.}

\setlength{\csromangathapagewidth}{0pt}
\setlength{\csromangathawakwidth}{0pt}
\csromangathameasuregroup{อุโภ สทฺธา วทญฺญู จ, \\ เต โหนฺติ ชานิปตโย, \\ อตฺถา’สํ ปจุรา โหนฺติ, \\ อมิตฺตา ทุมฺมนา โหนฺติ, \\ อิธ ธมฺมํ จริตฺวาน, \\ นนฺทิโน เทวโลกสฺมิํ,}{\csromangathabat{อุโภ สทฺธา วทญฺญู จ,}{สญฺญตา ธมฺมชีวิโน.} \\ \csromangathabat{เต โหนฺติ ชานิปตโย,}{อญฺญมญฺญํ ปิยํวทา.} \\ \csromangathabat{อตฺถา’สํ ปจุรา โหนฺติ,}{ผาสุกํ อุปชายติ.} \\ \csromangathabat{อมิตฺตา ทุมฺมนา โหนฺติ,}{อุภินฺนํ สมสีลินํ.} \\ \csromangathabat{อิธ ธมฺมํ จริตฺวาน,}{สมสีลพฺพตา อุโภ.} \\ \csromangathabat{นนฺทิโน เทวโลกสฺมิํ,}{โมทนฺติ กามกามิโนติ.ปญฺจมํ.}}
\csromangathasetleft{อุโภ สทฺธา วทญฺญู จ, \\ เต โหนฺติ ชานิปตโย, \\ อตฺถา’สํ ปจุรา โหนฺติ, \\ อมิตฺตา ทุมฺมนา โหนฺติ, \\ อิธ ธมฺมํ จริตฺวาน, \\ นนฺทิโน เทวโลกสฺมิํ,}
\csromangathagroup{\csromangathastanza{\csromangathabat{อุโภ สทฺธา วทญฺญู จ,}{สญฺญตา ธมฺมชีวิโน.} \\ \csromangathabat{เต โหนฺติ ชานิปตโย,}{อญฺญมญฺญํ ปิยํวทา.}}\csromangathastanza{\csromangathabat{อตฺถา’สํ ปจุรา โหนฺติ,}{ผาสุกํ อุปชายติ.} \\ \csromangathabat{อมิตฺตา ทุมฺมนา โหนฺติ,}{อุภินฺนํ สมสีลินํ.}}\csromangathastanza{\csromangathabat{อิธ ธมฺมํ จริตฺวาน,}{สมสีลพฺพตา อุโภ.} \\ \csromangathabat{นนฺทิโน เทวโลกสฺมิํ,}{โมทนฺติ กามกามิโนติ.  ปญฺจมํ.}}}

\csromanheader{2}{6. ทุติย\-สมชีวี\-สุตฺต}
\proseitem{56}{อากงฺเขยฺยุํ เจ ภิกฺขเว อุโภ ชานิปตโย ทิฏฺเฐ เจว ธมฺเม อญฺญมญฺญํ ปสฺสิตุํ, อภิสมฺปรายญฺจ อญฺญมญฺญํ ปสฺสิตุํ, อุโภว อสฺสุ สมสทฺธา สมสีลา สมจาคา สมปญฺญา, เต ทิฏฺเฐ เจว ธมฺเม อญฺญมญฺญํ ปสฺสนฺติ, อภิสมฺปรายญฺจ อญฺญมญฺญํ ปสฺสนฺตีติ.}

\setlength{\csromangathapagewidth}{0pt}
\setlength{\csromangathawakwidth}{0pt}
\csromangathameasuregroup{อุโภ สทฺธา วทญฺญู จ, \\ เต โหนฺติ ชานิปตโย, \\ อตฺถา’สํ ปจุรา โหนฺติ, \\ อมิตฺตา ทุมฺมนา โหนฺติ, \\ อิธ ธมฺมํ จริตฺวาน, \\ นนฺทิโน เทวโลกสฺมิํ,}{\csromangathabat{อุโภ สทฺธา วทญฺญู จ,}{สญฺญตา ธมฺมชีวิโน.} \\ \csromangathabat{เต โหนฺติ ชานิปตโย,}{อญฺญมญฺญํ ปิยํวทา.} \\ \csromangathabat{อตฺถา’สํ ปจุรา โหนฺติ,}{ผาสุกํ อุปชายติ.} \\ \csromangathabat{อมิตฺตา ทุมฺมนา โหนฺติ,}{อุภินฺนํ สมสีลินํ.} \\ \csromangathabat{อิธ ธมฺมํ จริตฺวาน,}{สมสีลพฺพตา อุโภ.} \\ \csromangathabat{นนฺทิโน เทวโลกสฺมิํ,}{โมทนฺติ กามกามิโนติ.ฉฏฺฐํ.}}
\csromangathasetleft{อุโภ สทฺธา วทญฺญู จ, \\ เต โหนฺติ ชานิปตโย, \\ อตฺถา’สํ ปจุรา โหนฺติ, \\ อมิตฺตา ทุมฺมนา โหนฺติ, \\ อิธ ธมฺมํ จริตฺวาน, \\ นนฺทิโน เทวโลกสฺมิํ,}
\csromangathagroup{\csromangathastanza{\csromangathabat{อุโภ สทฺธา วทญฺญู จ,}{สญฺญตา ธมฺมชีวิโน.} \\ \csromangathabat{เต โหนฺติ ชานิปตโย,}{อญฺญมญฺญํ ปิยํวทา.}}\csromangathastanza{\csromangathabat{อตฺถา’สํ ปจุรา โหนฺติ,}{ผาสุกํ อุปชายติ.} \\ \csromangathabat{อมิตฺตา ทุมฺมนา โหนฺติ,}{อุภินฺนํ สมสีลินํ.}}\csromangathastanza{\csromangathabat{อิธ ธมฺมํ จริตฺวาน,}{สมสีลพฺพตา อุโภ.} \\ \csromangathabat{นนฺทิโน เทวโลกสฺมิํ,}{โมทนฺติ กามกามิโนติ.  ฉฏฺฐํ.}}}

\csromanmatikaanchor{mtk.108}
\csromantocmarkref{subsection}{7. สุปฺปวาสาสุตฺต}{mtk.108}
\bookmark[dest={mtk.108},level=2]{7. สุปฺปวาสาสุตฺต}
\csromanheader{2}{7. \textbf{สุปฺปวาสาสุตฺต}}
\proseitem{57}{เอกํ สมยํ ภควา โกลิเยสุ วิหรติ ปชฺชนิกํ\csromansharedfootnote{adb769ad27aaf055}{สชฺชเนลํ (สี, อิ), ปชฺชเนลํ (สฺยา, กํ)} นาม โกลิยานํ นิคโม. อถ โข ภควา ปุพฺพณฺหสมยํ นิวาเสตฺวา \csromanfolio{374}ปตฺต\-จีวรมา\-ทาย เยน สุปฺปวาสาย โกลิยธีตุยา นิเวสนํ เตนุปสงฺกมิ, อุปสงฺกมิตฺวา ปญฺญตฺเต อาสเน นิสีทิ. อถ โข สุปฺปวาสา โกลิยธีตา ภควนฺตํ ปณีเตน ขาทนีเยน โภชนีเยน สหตฺถา สนฺตปฺเปสิ สมฺปวาเรสิ. อถ โข สุปฺปวาสา โกลิยธีตา ภควนฺตํ ภุตฺตาวิํ โอนีต\-ปตฺต\-ปาณิํ เอกมนฺตํ นิสีทิ. เอกมนฺตํ นิสินฺนํ โข สุปฺปวาสํ โกลิยธีตรํ ภควา เอตทโวจ\csromandash{}}

\prose{โภชนํ สุปฺปวาเส เทนฺตี อริยสาวิกา ปฏิคฺคาหกานํ จตฺตาริ ฐานานิ เทติ. กตมานิ จตฺตาริ? อายุํ เทติ, วณฺณํ เทติ, สุขํ เทติ, พลํ เทติ. อายุํ โข ปน ทตฺวา อายุสฺส ภาคินี โหติ ทิพฺพสฺส วา มานุสสฺส วา, วณฺณํ ทตฺวา วณฺณสฺส ภาคินี โหติ ทิพฺพสฺส วา มานุสสฺส วา, สุขํ ทตฺวา สุขสฺส ภาคินี โหติ ทิพฺพสฺส วา มานุสสฺส วา, พลํ ทตฺวา พลสฺส ภาคินี โหติ ทิพฺพสฺส วา มนุสสฺส วา, โภชนํ สุปฺปวาเส เทนฺตี อริยสาวิกา ปฏิคฺคาหกานํ อิมานิ จตฺตาริ ฐานานิ เทตีติ.}

\setlength{\csromangathapagewidth}{0pt}
\setlength{\csromangathawakwidth}{0pt}
\csromangathameasuregroup{}{สุสงฺขตํ โภชนํ ยา ททาติ, \\ สุจิํ ปณีตํ\textsuperscript{1} รสสา อุเปตํ. \\ สา ทกฺขิณา อุชฺชุคเตสุ ทินฺนา, \\ จรณูปปนฺเนสุ มหคฺคเตสุ. \\ ปุญฺเญน ปุญฺญํ สํสนฺทมานา, \\ มหปฺผลา โลกวิทูน วณฺณิตา. \\ เอตาทิสํ ยญฺญมนุสฺสรนฺตา, \\ เย เวทชาตา วิจรนฺติ โลเก. \\ วิเนยฺย มจฺเฉรมลํ สมูลํ, \\ อนินฺทิตา สคฺคมุเปนฺติ ฐานนฺติ.สตฺตมํ.}
\csromangathameasurewak{สุสงฺขตํ โภชนํ ยา ททาติ, \\ สุจิํ ปณีตํ\textsuperscript{1} รสสา อุเปตํ. \\ สา ทกฺขิณา อุชฺชุคเตสุ ทินฺนา, \\ จรณูปปนฺเนสุ มหคฺคเตสุ. \\ ปุญฺเญน ปุญฺญํ สํสนฺทมานา, \\ มหปฺผลา โลกวิทูน วณฺณิตา. \\ เอตาทิสํ ยญฺญมนุสฺสรนฺตา, \\ เย เวทชาตา วิจรนฺติ โลเก. \\ วิเนยฺย มจฺเฉรมลํ สมูลํ, \\ อนินฺทิตา สคฺคมุเปนฺติ ฐานนฺติ.สตฺตมํ.}
\setlength{\csromangathaleftcol}{0pt}
\csromangathagroup{\csromangathastanza{\csromangathawak{สุสงฺขตํ โภชนํ ยา ททาติ, \\ สุจิํ ปณีตํ\csromansharedfootnote{6ea6ef81223f6a27}{สุปณีตํ (ก)} รสสา อุเปตํ. \\ สา ทกฺขิณา อุชฺชุคเตสุ ทินฺนา, \\ จรณูปปนฺเนสุ มหคฺคเตสุ.}}\csromangathastanza{\csromangathawak{ปุญฺเญน ปุญฺญํ สํสนฺทมานา, \\ มหปฺผลา โลกวิทูน วณฺณิตา. \\ เอตาทิสํ ยญฺญม\-นุสฺสรนฺตา, \\ เย เวทชาตา วิจรนฺติ โลเก. \\ วิเนยฺย มจฺเฉรมลํ สมูลํ, \\ อนินฺทิตา สคฺคมุเปนฺติ ฐานนฺติ.  สตฺตมํ.}}}

\csromanmatikaanchor{mtk.109}
\csromantocmarkref{subsection}{8. สุทตฺตสุตฺต}{mtk.109}
\bookmark[dest={mtk.109},level=2]{8. สุทตฺตสุตฺต}
\csromanheader{2}{8. \textbf{สุทตฺตสุตฺต}}
\proseitem{58}{อถ โข อนาถปิณฺฑิโก คหปติ เยน ภควา เตนุปสงฺกมิ, อุปสงฺกมิตฺวา ภควนฺตํ อภิวาเทตฺวา เอกมนฺตํ นิสีทิ, เอกมนฺตํ นิสินฺนํ โข อนาถปิณฺฑิกํ คหปติํ ภควา เอตทโวจ\csromandash{}}

\chapterheadpageread{(6) ปุญฺญาภิสนฺท\-วคฺค}
\prose{\csromanfolio{375}โภชนํ คหปติ ททมาโน อริยสาวโก ปฏิคฺคาหกานํ จตฺตาริ ฐานานิ เทติ. กตมานิ จตฺตาริ? อายุํ เทติ, วณฺณํ เทติ, สุขํ เทติ, พลํ เทติ. อายุํ โข ปน ทตฺวา อายุสฺส ภาคี โหติ ทิพฺพสฺส วา มานุสสฺส วา. วณฺณํ ทตฺวา. สุขํ ทตฺวา. พลํ ทตฺวา พลสฺส ภาคี โหติ ทิพฺพสฺส วา มานุสสฺส วา. โภชนํ คหปติ ททมาโน อริยสาวโก ปฏิคฺคาหกานํ อิมานิ จตฺตาริ ฐานานิ เทตีติ.}

\setlength{\csromangathapagewidth}{0pt}
\setlength{\csromangathawakwidth}{0pt}
\csromangathameasuregroup{โส อายุทายี วณฺณทายี\textsuperscript{1}, \\ ทีฆายุ ยสวา โหติ,}{\textsuperscript{*}โย สญฺญตานํ ปรทตฺตโภชินํ, \\ กาเลน สกฺกจฺจ ททาติ โภชนํ. \\ จตฺตาริ ฐานานิ อนุปฺปเวจฺฉติ, \\ อายุญฺจ วณฺณญฺจ สุขํ พลญฺจ. \\ \csromangathabat{โส อายุทายี วณฺณทายี\textsuperscript{1},}{สุขํ พลํ ทโท\textsuperscript{1} นโร.} \\ \csromangathabat{ทีฆายุ ยสวา โหติ,}{ยตฺถ ยตฺถูปปชฺชตีติ.อฏฺฐมํ.}}
\csromangathameasurewak{\textsuperscript{*}โย สญฺญตานํ ปรทตฺตโภชินํ, \\ กาเลน สกฺกจฺจ ททาติ โภชนํ. \\ จตฺตาริ ฐานานิ อนุปฺปเวจฺฉติ, \\ อายุญฺจ วณฺณญฺจ สุขํ พลญฺจ.}
\csromangathasetleft{โส อายุทายี วณฺณทายี\textsuperscript{1}, \\ ทีฆายุ ยสวา โหติ,}
\csromangathagroup{\csromangathastanza{\csromangathawak{\csromansymbolfootnote{*}{วิ 3. 315 ปิฏฺเฐปิ.}โย สญฺญตานํ ปรทตฺตโภชินํ, \\ กาเลน สกฺกจฺจ ททาติ โภชนํ. \\ จตฺตาริ ฐานานิ อนุปฺปเวจฺฉติ, \\ อายุญฺจ วณฺณญฺจ สุขํ พลญฺจ.}}\csromangathastanza{\csromangathabat{โส อายุทายี วณฺณทายี\csromansharedfootnote{dae2d24557866341}{โส อายุทายี พลทายี (สี, อิ), อายุทายี พลทายี (สฺยา, กํ)},}{สุขํ พลํ ทโท\csromansharedfootnote{af0747fb056126e7}{สุขํ วณฺณํ ทโท (สี, สฺยา, กํ, อิ), สุขพลทโท (ก)} นโร.} \\ \csromangathabat{ทีฆายุ ยสวา โหติ,}{ยตฺถ ยตฺถูปปชฺชตีติ.  อฏฺฐมํ.}}}

\csromanmatikaanchor{mtk.110}
\csromantocmarkref{subsection}{9. โภชนสุตฺต}{mtk.110}
\bookmark[dest={mtk.110},level=2]{9. โภชนสุตฺต}
\csromanheader{2}{9. \textbf{โภชนสุตฺต}}
\proseitem{59}{โภชนํ ภิกฺขเว ททมาโน ทายโก ปฏิคฺคาหกานํ จตฺตาริ ฐานานิ เทติ. กตมานิ จตฺตาริ? อายุํ เทติ, วณฺณํ เทติ, สุขํ เทติ, พลํ เทติ. อายุํ โข ปน ทตฺวา อายุสฺส ภาคี โหติ ทิพฺพสฺส วา มานุสสฺส วา. วณฺณํ ทตฺวา. สุขํ ทตฺวา. พลํ ทตฺวา พลสฺส ภาคี โหติ ทิพฺพสฺส วา มานุสสฺส วา. โภชนํ ภิกฺขเว ททมาโน ทายโก ปฏิคฺคาหกานํ อิมานิ จตฺตาริ ฐานานิ เทตีติ.}

\setlength{\csromangathapagewidth}{0pt}
\setlength{\csromangathawakwidth}{0pt}
\csromangathameasuregroup{โส อายุทายี วณฺณทายี, \\ ทีฆายุ ยสวา โหติ,}{\textsuperscript{*}โย สญฺญตานํ ปรทตฺตโภชินํ, \\ กาเลน สกฺกจฺจ ททาติ โภชนํ. \\ จตฺตาริ ฐานานิ อนุปฺปเวจฺฉติ, \\ อายุญฺจ วณฺณญฺจ สุขํ พลญฺจ. \\ \csromangathabat{โส อายุทายี วณฺณทายี,}{สุขํ พลํ ทโท นโร.} \\ \csromangathabat{ทีฆายุ ยสวา โหติ,}{ยตฺถ ยตฺถูปปชฺชตีติ.นวมํ.}}
\csromangathameasurewak{\textsuperscript{*}โย สญฺญตานํ ปรทตฺตโภชินํ, \\ กาเลน สกฺกจฺจ ททาติ โภชนํ. \\ จตฺตาริ ฐานานิ อนุปฺปเวจฺฉติ, \\ อายุญฺจ วณฺณญฺจ สุขํ พลญฺจ.}
\csromangathasetleft{โส อายุทายี วณฺณทายี, \\ ทีฆายุ ยสวา โหติ,}
\csromangathagroup{\csromangathastanza{\csromangathawak{\csromansymbolmark{*}โย สญฺญตานํ ปรทตฺตโภชินํ, \\ กาเลน สกฺกจฺจ ททาติ โภชนํ. \\ จตฺตาริ ฐานานิ อนุปฺปเวจฺฉติ, \\ อายุญฺจ วณฺณญฺจ สุขํ พลญฺจ.}}\csromangathastanza{\csromangathabat{โส อายุทายี วณฺณทายี,}{สุขํ พลํ ทโท นโร.} \\ \csromangathabat{ทีฆายุ ยสวา โหติ,}{ยตฺถ ยตฺถูปปชฺชตีติ.  นวมํ.}}}

\csromanpalirecto
\csromanmatikaanchor{mtk.111}
\csromantocmarkref{subsection}{10. คิหิสามีจิสุตฺต}{mtk.111}
\bookmark[dest={mtk.111},level=2]{10. คิหิสามีจิสุตฺต}
\csromanmarkh{จตุกฺกนิปาต\-ปาฬิ}\csromanmarkhii{10. คิหิสามีจิ\-สุตฺต}\csromanheadernomark{2}{จตุกฺกนิปาต\-ปาฬิ 10. คิหิสามีจิ\-สุตฺต}
\proseitem{60}{\csromanfolio{376}อถ โข อนาถปิณฺฑิโก คหปติ เยน ภควา เตนุปสงฺกมิ, อุปสงฺกมิตฺวา ภควนฺตํ อภิวาเทตฺวา เอกมนฺตํ นิสีทิ, เอกมนฺตํ นิสินฺนํ โข อนาถปิณฺฑิกํ คหปติํ ภควา เอตทโวจ\csromandash{}}

\prose{จตูหิ โข คหปติ ธมฺเมหิ สมนฺนาคโต อริยสาวโก คิหิสามีจิ\-ปฏิปทํ ปฏิปนฺโน โหติ ยโส\-ปฏิลาภินิํ สคฺค\-สํวตฺตนิกํ. กตเมหิ จตูหิ? อิธ คหปติ อริยสาวโก ภิกฺขุสํฆํ ปจฺจุปฏฺฐิโต โหติ จีวเรน, ภิกฺขุสํฆํ ปจฺจุปฏฺฐิโต โหติ ปิณฺฑปาเตน, ภิกฺขุสํฆํ ปจฺจุปฏฺฐิโต โหติ เสนาสเนน, ภิกฺขุสํฆํ ปจฺจุปฏฺฐิโต โหติ คิลานปฺปจฺจย\-เภสชฺช\-ปริกฺขาเรน. อิเมหิ โข คหปติ จตูหิ ธมฺเมหิ สมนฺนาคโต อริยสาวโก คิหิสามีจิ\-ปฏิปทํ ปฏิปนฺโน โหติ ยโส\-ปฏิลาภินิํ สคฺค\-สํวตฺตนิกนฺติ.}

\prose{คิหิสามีจิ\-ปฏิปทํ, ปฏิปชฺชนฺติ ปณฺฑิตา.}

\prose{สมฺมคฺคเต สีลวนฺเต, จีวเรน อุปฏฺฐิตา.}

\prose{ปิณฺฑปาต\-สยเนน, คิลาน\-ปฺปจฺจเยน จ.}

\prose{เตสํ ทิวา จ รตฺโต จ, สทา ปุญฺญํ ปวฑฺฒติ.}

\prosewithclosermidruled{สคฺคญฺจ กมติฏฺฐานํ\csromansharedfootnote{61e18dd2a7999881}{สคฺคญฺจ สปฺปติฏฺฐานํ (ก)}, กมฺมํ กตฺวาน ภทฺทกนฺติ.  ทสมํ.}{ปุญฺญาภิสนฺท\-วคฺโค ปฐโม.}
\csromancenter{\textbf{ตสฺสุทฺทานํ}}
\setlength{\csromangathapagewidth}{0pt}
\setlength{\csromangathawakwidth}{0pt}
\csromangathameasuregroup{เทฺว ปุญฺญาภิสนฺทา เทฺว จ, \\ สุปฺปวาสา สุทตฺโต จ,}{\csromangathabat{เทฺว ปุญฺญาภิสนฺทา เทฺว จ,}{สํวาสา สมชีวิโน.} \\ \csromangathabat{สุปฺปวาสา สุทตฺโต จ,}{โภชนํ คิหิสามิจีติ.}}
\csromangathasetleft{เทฺว ปุญฺญาภิสนฺทา เทฺว จ, \\ สุปฺปวาสา สุทตฺโต จ,}
\csromangathagroup{\csromangathastanza{\csromangathabat{เทฺว ปุญฺญาภิสนฺทา เทฺว จ,}{สํวาสา สมชีวิโน.} \\ \csromangathabat{สุปฺปวาสา สุทตฺโต จ,}{โภชนํ คิหิสามิจีติ.}}}

\csromanmatikaanchor{mtk.112}
\csromantocmarknopage{section}{(7) 2. ปตฺตกมฺมวคฺค}{mtk.112}
\bookmark[dest={mtk.112},level=1]{(7) 2. ปตฺตกมฺมวคฺค}
\csromanmarkh{(7)}\csromanmarkhii{2. ปตฺตกมฺม\-วคฺค}\csromanheadernomark{1}{(7) 2. ปตฺตกมฺม\-วคฺค}
\csromanmatikaanchor{mtk.113}
\csromantocmarkref{subsection}{1. ปตฺตกมฺมสุตฺต}{mtk.113}
\bookmark[dest={mtk.113},level=2]{1. ปตฺตกมฺมสุตฺต}
\csromanheader{2}{1. ปตฺตกมฺม\-สุตฺต}
\proseitem{61}{อถ โข อนาถปิณฺฑิโก คหปติ เยน ภควา เตนุปสงฺกมิ, อุปสงฺกมิตฺวา ภควนฺตํ อภิวาเทตฺวา เอกมนฺตํ นิสีทิ, เอกมนฺตํ นิสินฺนํ โข อนาถปิณฺฑิกํ คหปติํ ภควา เอตทโวจ\csromandash{}}

\chapterheadpageread{(7) 2. ปตฺตกมฺม\-วคฺค}
\prose{\csromanfolio{377}จตฺตาโรเม คหปติ ธมฺมา อิฏฺฐา กนฺตา มนาปา ทุลฺลภา โลกสฺมิํ. กตเม จตฺตาโร? โภคา เม อุปฺปชฺชนฺตุ สหธมฺเมนาติ. อยํ ปฐโม ธมฺโม อิฏฺโฐ กนฺโต มนาโป ทุลฺลโภ โลกสฺมิํ.}

\prose{โภเค ลทฺธา สหธมฺเมน ยโส เม อาคจฺฉตุ สห ญาตีหิ สห อุปชฺฌาเยหีติ. อยํ ทุติโย ธมฺโม อิฏฺโฐ กนฺโต มนาโป ทุลฺลโภ โลกสฺมิํ.}

\prose{โภเค ลทฺธา สหธมฺเมน, ยสํ ลทฺธา สห ญาตีหิ สห อุปชฺฌาเยหิ จิรํ ชีวามิ ทีฆมายุํ ปาเลมีติ. อยํ ตติโย ธมฺโม อิฏฺโฐ กนฺโต มนาโป ทุลฺลโภ โลกสฺมิํ.}

\prose{โภเค ลทฺธา สหธมฺเมน, ยสํ ลทฺธา สห ญาตีหิ สห อุปชฺฌาเยหิ จิรํ ชีวิตฺวา, ทีฆมายุํ ปาเลตฺวา กายสฺส เภทา ปรํ มรณา สุคติํ สคฺคํ โลกํ อุปปชฺชามีติ. อยํ จตุตฺโถ ธมฺโม อิฏฺโฐ กนฺโต มนาโป ทุลฺลโภ โลกสฺมิํ. อิเม โข คหปติ จตฺตาโร ธมฺมา อิฏฺฐา กนฺตา มนาปา ทุลฺลภา โลกสฺมิํ.}

\prose{อิเมสํ โข คหปติ จตุนฺนํ ธมฺมานํ อิฏฺฐานํ กนฺตานํ มนาปานํ ทุลฺลภานํ โลกสฺมิํ จตฺตาโร ธมฺมา ปฏิลาภาย สํวตฺตนฺติ. กตเม จตฺตาโร? สทฺธาสมฺปทา สีลสมฺปทา จาคสมฺปทา ปญฺญาสมฺปทา.}

\prose{กตมา จ คหปติ สทฺธาสมฺปทา? อิธ คหปติ อริยสาวโก สทฺโธ โหติ, สทฺทหติ ตถาคตสฺส โพธิํ “อิติปิ โส ภควา อรหํ สมฺมาสมฺพุทฺโธ วิชฺชา\-จรณ\-สมฺปนฺโน สุคโต โลกวิทู อนุตฺตโร ปุริส\-ทมฺม\-สารถิ สตฺถา เทวมนุสฺสานํ พุทฺโธ ภควา”ติ. อยํ วุจฺจติ คหปติ สทฺธาสมฺปทา.}

\prose{กตมา จ คหปติ สีลสมฺปทา? อิธ คหปติ อริยสาวโก ปาณาติปาตา ปฏิวิรโต โหติ ฯเปฯ สุรา\-เมรย\-มชฺช\-ปมาทฏฺฐานา ปฏิวิรโต โหติ. อยํ วุจฺจติ คหปติ สีลสมฺปทา.}

\prose{กตมา จ คหปติ จาคสมฺปทา? อิธ คหปติ อริยสาวโก วิคต\-มล\-มจฺเฉเรน เจตสา อคารํ อชฺฌาวสติ, มุตฺตจาโค ปยตปาณิ โวสคฺครโต ยาจโยโค ทาน\-สํวิภาค\-รโต. อยํ วุจฺจติ คหปติ จาคสมฺปทา.}

\prose{\csromanfolio{378}กตมา จ คหปติ ปญฺญาสมฺปทา? อภิชฺฌา\-วิสมโลภา\-ภิภูเตน คหปติ เจตสา วิหรนฺโต อกิจฺจํ กโรติ, กิจฺจํ อปราเธติ, อกิจฺจํ กโรนฺโต กิจฺจํ อปราเธนฺโต ยสา จ สุขา จ ธํสติ. พฺยาปาทา\-ภิภูเตน คหปติ เจตสา วิหรนฺโต อกิจฺจํ กโรติ, กิจฺจํ อปราเธติ, อกิจฺจํ กโรนฺโต กิจฺจํ อปราเธนฺโต ยสา จ สุขา จ ธํสติ. ถินมิทฺธา\-ภิภูเตน คหปติ เจตสา วิหรนฺโต อกิจฺจํ กโรติ, กิจฺจํ อปราเธติ, อกิจฺจํ กโรนฺโต กิจฺจํ อปราเธนฺโต ยสา จ สุขา จ ธํสติ. อุทฺธจฺจกุกฺกุจฺจา\-ภิภูเตน คหปติ เจตสา วิหรนฺโต อกิจฺจํ กโรติ, กิจฺจํ อปราเธติ, อกิจฺจํ กโรนฺโต กิจฺจํ อปราเธนฺโต ยสา จ สุขา จ ธํสติ. วิจิกิจฺฉา\-ภิภูเตน คหปติ เจตสา วิหรนฺโต อกิจฺจํ กโรติ, กิจฺจํ อปราเธติ, อกิจฺจํ กโรนฺโต กิจฺจํ อปราเธนฺโต ยสา จ สุขา จ ธํสติ.}

\prose{ส โข โส คหปติ อริยสาวโก “อภิชฺฌา\-วิสมโลโภ จิตฺตสฺส อุปกฺกิเลโส”ติ อิติ วิทิตฺวา อภิชฺฌา\-วิสมโลภํ จิตฺตสฺส อุปกฺกิเลสํ ปชหติ. “พฺยาปาโท จิตฺตสฺส อุปกฺกิเลโส”ติ อิติ วิทิตฺวา พฺยาปาทํ จิตฺตสฺส อุปกฺกิเลสํ ปชหติ. “ถินมิทฺธํ จิตฺตสฺส อุปกฺกิเลโส”ติ อิติ วิทิตฺวา ถินมิทฺธํ จิตฺตสฺส อุปกฺกิเลสํ ปชหติ. “อุทฺธจฺจ\-กุกฺกุจฺจํ จิตฺตสฺส อุปกฺกิเลโส”ติ อิติ วิทิตฺวา อุทฺธจฺจ\-กุกฺกุจฺจํ จิตฺตสฺส อุปกฺกิเลสํ ปชหติ. “วิจิกิจฺฉา จิตฺตสฺส อุปกฺกิเลโส”ติ อิติ วิทิตฺวา วิจิกิจฺฉํ จิตฺตสฺส อุปกฺกิเลสํ ปชหติ.}

\prose{ยโต จ โข คหปติ อริยสาวกสฺส “อภิชฺฌา\-วิสมโลโภ จิตฺตสฺส อุปกฺกิเลโส”ติ อิติ วิทิตฺวา อภิชฺฌา\-วิสมโลโภ จิตฺตสฺส อุปกฺกิเลโส ปหีโน โหติ. “พฺยาปาโท จิตฺตสฺส อุปกฺกิเลโส”ติ อิติ วิทิตฺวา พฺยาปาโท จิตฺตสฺส อุปกฺกิเลโส ปหีโน โหติ. “ถินมิทฺธํ จิตฺตสฺส อุปกฺกิเลโส”ติ อิติ วิทิตฺวา ถินมิทฺธํ จิตฺตสฺส อุปกฺกิเลโส ปหีโน โหติ. “อุทฺธจฺจ\-กุกฺกุจฺจํ จิตฺตสฺส อุปกฺกิเลโส”ติ อิติ วิทิตฺวา อุทฺธจฺจ\-กุกฺกุจฺจํ จิตฺตสฺส อุปกฺกิเลโส ปหีโน โหติ. “วิจิกิจฺฉา จิตฺตสฺส อุปกฺกิเลโส”ติ อิติ วิทิตฺวา วิจิกิจฺฉา จิตฺตสฺส อุปกฺกิเลโส ปหีโน โหติ. อยํ วุจฺจติ คหปติ อริยสาวโก มหาปญฺโญ ปุถุปญฺโญ อาปาตทโส\csromansharedfootnote{b460704486092f3c}{อาปาถทโส (สี, สฺยา, กํ, อิ)} ปญฺญาสมฺปนฺโน\csromansharedfootnote{d6ceeb21157488ca}{หาสปญฺโญ (ก)}. อยํ วุจฺจติ}

\vspace{\tipitakaparskip}
\csromancenter{(7) 2. ปตฺตกมฺม\-วคฺค คหปติ ปญฺญาสมฺปทา. อิเมสํ โข คหปติ จตุนฺนํ ธมฺมานํ อิฏฺฐานํ กนฺตานํ มนาปานํ ทุลฺลภานํ โลกสฺมิํ. อิเม จตฺตาโร ธมฺมา ปฏิลาภาย สํวตฺตนฺติ.}
\prose{\csromanfolio{379}ส โข โส คหปติ อริยสาวโก อุฏฺฐาน\-วีริยา\-ธิคเตหิ โภเคหิ พาหาพล\-ปริจิเตหิ เสทา\-วกฺขิตฺเตหิ ธมฺมิเกหิ ธมฺมลทฺเธหิ จตฺตาริ ปตฺตกมฺมานิ กตฺตา โหติ. กตมานิ จตฺตาริ? อิธ คหปติ อริยสาวโก อุฏฺฐาน\-วีริยา\-ธิคเตหิ โภเคหิ พาหาพล\-ปริจิเตหิ เสทา\-วกฺขิตฺเตหิ ธมฺมิเกหิ ธมฺมลทฺเธหิ อตฺตานํ สุเขติ ปีเณติ, สมฺมา สุขํ ปริหรติ, มาตาปิตโร สุเขติ ปีเณติ, สมฺมา สุขํ ปริหรติ, ปุตฺตทารทาสกมฺมกรโปริเส สุเขติ ปีเณติ, สมฺมา สุขํ ปริหรติ, มิตฺตามจฺเจ สุเขติ ปีเณติ, สมฺมา สุขํ ปริหรติ. อิทมสฺส ปฐมํ ฐานคตํ โหติ ปตฺตคตํ อายตนโส ปริภุตฺตํ.}

\prose{ปุน จปรํ คหปติ อริยสาวโก อุฏฺฐาน\-วีริยา\-ธิคเตหิ โภเคหิ พาหาพล\-ปริจิเตหิ เสทา\-วกฺขิตฺเตหิ ธมฺมิเกหิ ธมฺมลทฺเธหิ, ยา ตา โหนฺติ อาปทา อคฺคิโต วา อุทกโต วา ราชโต วา โจรโต วา อปฺปิยโต วา ทายาทโต\csromansharedfootnote{08d28e0eb9d16f43}{อํ 2. 39, 152 ปิฏฺเฐสุปิ.}, ตถารูปาสุ อาปทาสุ ปริโยธาย สํวตฺตติ, โสตฺถิํ อตฺตานํ กโรติ. อิทมสฺส ทุติยํ ฐานคตํ โหติ ปตฺตคตํ อายตนโส ปริภุตฺตํ.}

\prose{ปุน จปรํ คหปติ อริยสาวโก อุฏฺฐาน\-วีริยา\-ธิคเตหิ โภเคหิ พาหาพล\-ปริจิเตหิ เสทา\-วกฺขิตฺเตหิ ธมฺมิเกหิ ธมฺมลทฺเธหิ ปญฺจพลิํ กตฺตา โหติ ญาติพลิํ อติถิพลิํ ปุพฺพเปตพลิํ ราชพลิํ เทวตาพลิํ. อิทมสฺส ตติยํ ฐานคตํ โหติ ปตฺตคตํ อายตนโส ปริภุตฺตํ.}

\prose{ปุน จปรํ คหปติ อริยสาวโก อุฏฺฐาน\-วีริยา\-ธิคเตหิ โภเคหิ พาหาพล\-ปริจิเตหิ เสทา\-วกฺขิตฺเตหิ ธมฺมิเกหิ ธมฺมลทฺเธหิ, เย เต สมณพฺราหฺมณา มทปฺปมาทา ปฏิวิรตา ขนฺติโสรจฺเจ นิวิฏฺฐา, เอกมตฺตานํ ทเมนฺติ, เอกมตฺตานํ สเมนฺติ, เอกมตฺตานํ ปรินิพฺพาเปนฺติ, ตถารูเปสุ สมณ\-พฺราหฺมเณสุ อุทฺธคฺคิกํ ทกฺขิณํ ปติฏฺฐาเปติ โสวคฺคิกํ สุขวิปากํ}

\prose{\csromanfolio{380}สคฺค\-สํวตฺตนิกํ. อิทมสฺส จตุตฺถํ ฐานคตํ โหติ ปตฺตคตํ อายตนโส ปริภุตฺตํ.}

\prose{ส โข โส คหปติ อริยสาวโก อุฏฺฐาน\-วีริยา\-ธิคเตหิ โภเคหิ พาหาพล\-ปริจิเตหิ เสทา\-วกฺขิตฺเตหิ ธมฺมิเกหิ ธมฺมลทฺเธหิ อิมานิ จตฺตาริ ปตฺตกมฺมานิ กตฺตา โหติ. ยสฺส กสฺสจิ คหปติ อญฺญตฺร อิเมหิ จตูหิ ปตฺตกมฺเมหิ โภคา ปริกฺขยํ คจฺฉนฺติ, อิเม วุจฺจนฺติ คหปติ โภคา อฏฺฐานคตา อปตฺตคตา อนายตนโส ปริภุตฺตา. ยสฺส กสฺสจิ คหปติ อิเมหิ จตูหิ ปตฺตกมฺเมหิ โภคา ปริกฺขยํ คจฺฉนฺติ, อิเม วุจฺจนฺติ คหปติ โภคา ฐานคตา ปตฺตคตา อายตนโส ปริภุตฺตาติ.}

\setlength{\csromangathapagewidth}{0pt}
\setlength{\csromangathawakwidth}{0pt}
\csromangathameasuregroup{ภุตฺตา โภคา ภตา ภจฺจา\textsuperscript{1}, \\ อุทฺธคฺคา ทกฺขิณา ทินฺนา, \\ อุปฏฺฐิตา สีลวนฺโต, \\ ยทตฺถํ โภคํ อิจฺเฉยฺย, \\ โส เม อตฺโถ อนุปฺปตฺโต, \\ เอตํ\textsuperscript{1} อนุสฺสรํ มจฺโจ, \\ อิเธว นํ ปสํสนฺติ,}{\csromangathabat{ภุตฺตา โภคา ภตา ภจฺจา\textsuperscript{1},}{วิติณฺณา อาปทาสุ เม.} \\ \csromangathabat{อุทฺธคฺคา ทกฺขิณา ทินฺนา,}{อโถ ปญฺจพลี กตา.} \\ \csromangathabat{อุปฏฺฐิตา สีลวนฺโต,}{สญฺญตา พฺรหฺมจารโย.} \\ \csromangathabat{ยทตฺถํ โภคํ อิจฺเฉยฺย,}{ปณฺฑิโต ฆรมาวสํ.} \\ \csromangathabat{โส เม อตฺโถ อนุปฺปตฺโต,}{กตํ อนนุตาปิยํ.} \\ \csromangathabat{เอตํ\textsuperscript{1} อนุสฺสรํ มจฺโจ,}{อริยธมฺเม ฐิโต นโร.} \\ \csromangathabat{อิเธว นํ ปสํสนฺติ,}{เปจฺจ สคฺเค ปโมทตีติ.ปฐมํ.}}
\csromangathasetleft{ภุตฺตา โภคา ภตา ภจฺจา\textsuperscript{1}, \\ อุทฺธคฺคา ทกฺขิณา ทินฺนา, \\ อุปฏฺฐิตา สีลวนฺโต, \\ ยทตฺถํ โภคํ อิจฺเฉยฺย, \\ โส เม อตฺโถ อนุปฺปตฺโต, \\ เอตํ\textsuperscript{1} อนุสฺสรํ มจฺโจ, \\ อิเธว นํ ปสํสนฺติ,}
\csromangathagroup{\csromangathastanza{\csromangathabat{ภุตฺตา โภคา ภตา ภจฺจา\csromansharedfootnote{c34290b807e3a718}{คตา ภูตา (ก) ภฏา ภจฺจา (สฺยา, กํ)},}{วิติณฺณา อาปทาสุ เม.} \\ \csromangathabat{อุทฺธคฺคา ทกฺขิณา ทินฺนา,}{อโถ ปญฺจพลี กตา.}}\csromangathastanza{\csromangathabat{อุปฏฺฐิตา สีลวนฺโต,}{สญฺญตา พฺรหฺมจารโย.} \\ \csromangathabat{ยทตฺถํ โภคํ อิจฺเฉยฺย,}{ปณฺฑิโต ฆรมาวสํ.}}\csromangathastanza{\csromangathabat{โส เม อตฺโถ อนุปฺปตฺโต,}{กตํ อนนุตาปิยํ.} \\ \csromangathabat{เอตํ\csromansharedfootnote{865fdde227c93b9a}{เอวํ (ก)} อนุสฺสรํ มจฺโจ,}{อริยธมฺเม ฐิโต นโร.} \\ \csromangathabat{อิเธว นํ ปสํสนฺติ,}{เปจฺจ สคฺเค ปโมทตีติ.  ปฐมํ.}}}

\csromanmatikaanchor{mtk.114}
\csromantocmarkref{subsection}{2. อานณฺยสุตฺต}{mtk.114}
\bookmark[dest={mtk.114},level=2]{2. อานณฺยสุตฺต}
\csromanheader{2}{2. \textbf{อานณฺยสุตฺต}}
\proseitem{62}{อถ โข อนาถปิณฺฑิโก คหปติ เยน ภควา เตนุปสงฺกมิ, อุปสงฺกมิตฺวา ภควนฺตํ อภิวาเทตฺวา เอกมนฺตํ นิสีทิ, เอกมนฺตํ นิสินฺนํ โข อนาถปิณฺฑิกํ คหปติํ ภควา เอตทโวจ\csromandash{}}

\prose{จตฺตาริมานิ คหปติ สุขานิ อธิคมนียานิ คิหินา กามโภคินา กาเลน กาลํ สมเยน สมยํ อุปาทาย. กตมานิ จตฺตาริ? อตฺถิสุขํ โภคสุขํ อานณฺยสุขํ\csromansharedfootnote{357c5daa97150a74}{อณณสุขํ (สี, สฺยา, กํ, อิ)} อนวชฺชสุขํ.}

\prose{กตมญฺจ คหปติ อตฺถิสุขํ? อิธ คหปติ กุลปุตฺตสฺส โภคา โหนฺติ อุฏฺฐาน\-วีริยา\-ธิคตา พาหาพล\-ปริจิตา เสทาวกฺขิตฺตา ธมฺมิกา}

\vspace{\tipitakaparskip}
\csromancenter{(7) 2. ปตฺตกมฺม\-วคฺค ธมฺมลทฺธา, โส “โภคา เม อตฺถิ อุฏฺฐาน\-วีริยา\-ธิคตา พาหาพล\-ปริจิตา เสทาวกฺขิตฺตา ธมฺมิกา ธมฺมลทฺธา”ติ อธิคจฺฉติ สุขํ, อธิคจฺฉติ โสมนสฺสํ. อิทํ วุจฺจติ คหปติ อตฺถิสุขํ.}
\prose{\csromanfolio{381}กตมญฺจ คหปติ โภคสุขํ? อิธ คหปติ กุลปุตฺโต อุฏฺฐาน\-วีริยา\-ธิคเตหิ โภเคหิ พาหาพล\-ปริจิเตหิ เสทา\-วกฺขิตฺเตหิ ธมฺมิเกหิ ธมฺมลทฺเธหิ ปริภุญฺชติ, ปุญฺญานิ จ กโรติ, โส “อุฏฺฐาน\-วีริยา\-ธิคเตหิ โภเคหิ พาหาพล\-ปริจิเตหิ เสทา\-วกฺขิตฺเตหิ ธมฺมิเกหิ ธมฺมลทฺเธหิ ปริภุญฺชามิ, ปุญฺญานิ จ กโรมี”ติ อธิคจฺฉติ สุขํ, อธิคจฺฉติ โสมนสฺสํ. อิทํ วุจฺจติ คหปติ โภคสุขํ.}

\prose{กตมญฺจ คหปติ อานณฺยสุขํ? อิธ คหปติ กุลปุตฺโต น กสฺสจิ กิญฺจิ ธาเรติ\csromansharedfootnote{d89b22c59f3dfeea}{กิญฺจิ วา เทติ (ก)} อปฺปํ วา พหุํ วา. โส “น กสฺสจิ กิญฺจิ ธาเรมิ อปฺปํ วา พหุํ วา”ติ อธิคจฺฉติ สุขํ, อธิคจฺฉติ โสมนสฺสํ. อิทํ วุจฺจติ คหปติ อานณฺยสุขํ.}

\prose{กตมญฺจ คหปติ อนวชฺชสุขํ, อิธ คหปติ อริยสาวโก อนวชฺเชน กายกมฺเมน สมนฺนาคโต โหติ, อนวชฺเชน วจีกมฺเมน สมนฺนาคโต โหติ, อนวชฺเชน มโนกมฺเมน สมนฺนาคโต โหติ. โส “อนวชฺเชนมฺหิ กายกมฺเมน สมนฺนาคโต, อนวชฺเชน วจีกมฺเมน สมนฺนาคโต, อนวชฺเชน มโนกมฺเมน สมนฺนาคโต”ติ อธิคจฺฉติ สุขํ, อธิคจฺฉติ โสมนสฺสํ. อิทํ วุจฺจติ คหปติ อนวชฺชสุขํ. อิมานิ โข คหปติ จตฺตาริ สุขานิ อธิคมนียานิ คิหินา กามโภคินา กาเลน กาลํ สมเยน สมยํ อุปาทายาติ.}

\setlength{\csromangathapagewidth}{0pt}
\setlength{\csromangathawakwidth}{0pt}
\csromangathameasuregroup{อานณฺยสุขํ ญตฺวาน, \\ ภุญฺชํ โภคสุขํ มจฺโจ, \\ วิปสฺสมาโน ชานาติ, \\ อนวชฺชสุขสฺเสตํ,}{\csromangathabat{อานณฺยสุขํ ญตฺวาน,}{อโถ อตฺถิสุขํ ปรํ.} \\ \csromangathabat{ภุญฺชํ โภคสุขํ มจฺโจ,}{ตโต ปญฺญา วิปสฺสติ.} \\ \csromangathabat{วิปสฺสมาโน ชานาติ,}{อุโภ ภาเค สุเมธโส.} \\ \csromangathabat{อนวชฺชสุขสฺเสตํ,}{กลํ นาคฺฆติ โสฬสินฺติ.ทุติยํ.}}
\csromangathasetleft{อานณฺยสุขํ ญตฺวาน, \\ ภุญฺชํ โภคสุขํ มจฺโจ, \\ วิปสฺสมาโน ชานาติ, \\ อนวชฺชสุขสฺเสตํ,}
\csromangathagroup{\csromangathastanza{\csromangathabat{อานณฺยสุขํ ญตฺวาน,}{อโถ อตฺถิสุขํ ปรํ.} \\ \csromangathabat{ภุญฺชํ โภคสุขํ มจฺโจ,}{ตโต ปญฺญา วิปสฺสติ.}}\csromangathastanza{\csromangathabat{วิปสฺสมาโน ชานาติ,}{อุโภ ภาเค สุเมธโส.} \\ \csromangathabat{อนวชฺชสุขสฺเสตํ,}{กลํ นาคฺฆติ โสฬสินฺติ.  ทุติยํ.}}}

\csromanmatikaanchor{mtk.115}
\csromantocmarkref{subsection}{3. พฺรหฺมสุตฺต}{mtk.115}
\bookmark[dest={mtk.115},level=2]{3. พฺรหฺมสุตฺต}
\csromanheader{2}{3. \textbf{พฺรหฺมสุตฺต}}
\proseitem{63}{สพฺรหฺมกานิ ภิกฺขเว\csromansharedfootnote{5cf5b2b23821630a}{ขุ 1. 268 ปิฏฺเฐปิ.} ตานิ กุลานิ, เยสํ ปุตฺตานํ มาตาปิตโร อชฺฌาคาเร ปูชิตา โหนฺติ. สปุพฺพา\-จริย\-กานิ ภิกฺขเว ตานิ \csromanfolio{382}กุลานิ, เยสํ ปุตฺตานํ มาตาปิตโร อชฺฌาคาเร ปูชิตา โหนฺติ. สปุพฺพ\-เทวตานิ\csromansharedfootnote{b6fcdc46526c7c38}{สปุพฺพเทวานิ (สฺยา, กํ)} ภิกฺขเว ตานิ กุลานิ, เยสํ ปุตฺตานํ มาตปิตโร อชฺฌาคาเร ปูชิตา โหนฺติ. สาหุเนยฺยกานิ ภิกฺขเว ตานิ กุลานิ, เยสํ ปุตฺตานํ มาตาปิตโร อชฺฌาคาเร ปูชิตา โหนฺติ.}

\prose{“พฺรหฺมา”ติ ภิกฺขเว มาตาปิตูนํ\csromansharedfootnote{b8923dbdf6b9b55d}{มาตาปิตุนฺนํ (สี, อิ)} เอตํ อธิวจนํ, “ปุพฺพาจริยา”ติ ภิกฺขเว มาตาปิตูนํ เอตํ อธิวจนํ, “ปุพฺพเทวตา”ติ\csromansharedfootnote{0cbc60dd9b35f26f}{ปุพฺพเทวาติ (สี, สฺยา, กํ)} ภิกฺขเว มาตาปิตูนํ เอตํ อธิวจนํ, “อาหุเนยฺยา”ติ ภิกฺขเว มาตาปิตูนํ เอตํ อธิวจนํ. ตํ กิสฺส เหตุ? พหุการา ภิกฺขเว มาตาปิตโร ปุตฺตานํ อาปาทกา โปสกา อิมสฺส โลกสฺส ทสฺเสตาโรติ.}

\prose{“พฺรหฺมา”ติ มาตาปิตโร,}

\prose{“ปุพฺพาจริยา”ติ วุจฺจเร.}

\setlength{\csromangathapagewidth}{0pt}
\setlength{\csromangathawakwidth}{0pt}
\csromangathameasuregroup{อาหุเนยฺยา จ ปุตฺตานํ, \\ ตสฺมา หิ เน นมสฺเสยฺย, \\ อนฺเนน อถ ปาเนน, \\ อุจฺฉาทเนน นฺหาปเนน, \\ ตาย นํ ปาริจริยาย, \\ อิเธว นํ ปสํสนฺติ,}{\csromangathabat{อาหุเนยฺยา จ ปุตฺตานํ,}{ปชาย อนุกมฺปกา.} \\ \csromangathabat{ตสฺมา หิ เน นมสฺเสยฺย,}{สกฺกเรยฺย จ ปณฺฑิโต.} \\ \csromangathabat{อนฺเนน อถ ปาเนน,}{วตฺเถน สยเนน จ.} \\ \csromangathabat{อุจฺฉาทเนน นฺหาปเนน,}{ปาทานํ โธวเนน จ.} \\ \csromangathabat{ตาย นํ ปาริจริยาย,}{มาตาปิตูสุ ปณฺฑิตา.} \\ \csromangathabat{อิเธว นํ ปสํสนฺติ,}{เปจฺจ สคฺเค ปโมทตีติ.ตติยํ.}}
\csromangathasetleft{อาหุเนยฺยา จ ปุตฺตานํ, \\ ตสฺมา หิ เน นมสฺเสยฺย, \\ อนฺเนน อถ ปาเนน, \\ อุจฺฉาทเนน นฺหาปเนน, \\ ตาย นํ ปาริจริยาย, \\ อิเธว นํ ปสํสนฺติ,}
\csromangathagroup{\csromangathastanza{\csromangathabat{อาหุเนยฺยา จ ปุตฺตานํ,}{ปชาย อนุกมฺปกา.}}\csromangathastanza{\csromangathabat{ตสฺมา หิ เน นมสฺเสยฺย,}{สกฺกเรยฺย จ ปณฺฑิโต.} \\ \csromangathabat{อนฺเนน อถ ปาเนน,}{วตฺเถน สยเนน จ.}}\csromangathastanza{\csromangathabat{อุจฺฉาทเนน นฺหาปเนน,}{ปาทานํ โธวเนน จ.} \\ \csromangathabat{ตาย นํ ปาริจริยาย,}{มาตาปิตูสุ ปณฺฑิตา.} \\ \csromangathabat{อิเธว นํ ปสํสนฺติ,}{เปจฺจ สคฺเค ปโมทตีติ.  ตติยํ.}}}

\csromanmatikaanchor{mtk.116}
\csromantocmarkref{subsection}{4. นิรยสุตฺต}{mtk.116}
\bookmark[dest={mtk.116},level=2]{4. นิรยสุตฺต}
\csromanheader{2}{4. \textbf{นิรยสุตฺต}}
\proseitem{64}{จตูหิ ภิกฺขเว ธมฺเมหิ สมนฺนาคโต ยถาภตํ นิกฺขิตฺโต เอวํ นิรเย. กตเมหิ จตูหิ? ปาณาติปาตี โหติ, อทินฺนาทายี โหติ, กาเมสุ\-มิจฺฉาจารี โหติ, มุสาวาที โหติ. อิเมหิ โข ภิกฺขเว จตูหิ ธมฺเมหิ สมนฺนาคโต ยถาภตํ นิกฺขิตฺโต เอวํ นิรเยติ.}

\setlength{\csromangathapagewidth}{0pt}
\setlength{\csromangathawakwidth}{0pt}
\csromangathameasuregroup{ปาณาติปาโต อทินฺนาทานํ, \\ ปรทารคมนญฺจาปิ,}{\csromangathabat{ปาณาติปาโต อทินฺนาทานํ,}{มุสาวาโท จ วุจฺจติ.} \\ \csromangathabat{ปรทารคมนญฺจาปิ,}{นปฺปสํสนฺติ ปณฺฑิตาติ.จตุตฺถํ.}}
\csromangathasetleft{ปาณาติปาโต อทินฺนาทานํ, \\ ปรทารคมนญฺจาปิ,}
\csromangathagroup{\csromangathastanza{\csromangathabat{ปาณาติปาโต อทินฺนาทานํ,}{มุสาวาโท จ วุจฺจติ.} \\ \csromangathabat{ปรทารคมนญฺจาปิ,}{นปฺปสํสนฺติ ปณฺฑิตาติ.  จตุตฺถํ.}}}

\csromanmatikaanchor{mtk.117}
\csromantocmarkref{subsection}{5. รูปสุตฺต (จตุปฺปมาณสุตฺต)}{mtk.117}
\bookmark[dest={mtk.117},level=2]{5. รูปสุตฺต (จตุปฺปมาณสุตฺต)}
\csromanheader{2}{5. \textbf{รูปสุตฺต}}
\proseitem{65}{จตฺตาโรเม ภิกฺขเว ปุคฺคลา สนฺโต สํวิชฺชมานา โลกสฺมิํ. กตเม จตฺตาโร? รูปปฺปมาโณ รูปปฺปสนฺโน, โฆสปฺปมาโณ}

\vspace{\tipitakaparskip}
\csromancenter{(7) 2. ปตฺตกมฺม\-วคฺค โฆสปฺปสนฺโน, ลูขปฺปมาโณ ลูขปฺปสนฺโน, ธมฺมปฺปมาโณ ธมฺมปฺปสนฺโน. อิเม โข ภิกฺขเว จตฺตาโร ปุคฺคลา สนฺโต สํวิชฺชมานา โลกสฺมินฺติ.}
\setlength{\csromangathapagewidth}{0pt}
\setlength{\csromangathawakwidth}{0pt}
\csromangathameasuregroup{เย จ รูเป ปมาณิํสุ\textsuperscript{1}, \\ ฉนฺทราควสูเปตา, \\ อชฺฌตฺตญฺจ น ชานาติ, \\ สมนฺตาวรโณ พาโล, \\ อชฺฌตฺตญฺจ น ชานาติ, \\ พหิทฺธา ผลทสฺสาวี, \\ อชฺฌตฺตญฺจ ปชานาติ, \\ วินีวรณทสฺสาวี,}{\csromangathabat{เย จ รูเป ปมาณิํสุ\textsuperscript{1},}{เย จ โฆเสน อนฺวคู.} \\ \csromangathabat{ฉนฺทราควสูเปตา,}{นาภิชานนฺติ เต ชนา\textsuperscript{1}.} \\ \csromangathabat{อชฺฌตฺตญฺจ น ชานาติ,}{พหิทฺธา จ น ปสฺสติ.} \\ \csromangathabat{สมนฺตาวรโณ พาโล,}{ส เว โฆเสน วุยฺหติ.} \\ \csromangathabat{อชฺฌตฺตญฺจ น ชานาติ,}{พหิทฺธา จ วิปสฺสติ.} \\ \csromangathabat{พหิทฺธา ผลทสฺสาวี,}{โสปิ โฆเสน วุยฺหติ.} \\ \csromangathabat{อชฺฌตฺตญฺจ ปชานาติ,}{พหิทฺธา จ วิปสฺสติ.} \\ \csromangathabat{วินีวรณทสฺสาวี,}{น โส โฆเสน วุยฺหตีติ.ปญฺจมํ.}}
\csromangathasetleft{เย จ รูเป ปมาณิํสุ\textsuperscript{1}, \\ ฉนฺทราควสูเปตา, \\ อชฺฌตฺตญฺจ น ชานาติ, \\ สมนฺตาวรโณ พาโล, \\ อชฺฌตฺตญฺจ น ชานาติ, \\ พหิทฺธา ผลทสฺสาวี, \\ อชฺฌตฺตญฺจ ปชานาติ, \\ วินีวรณทสฺสาวี,}
\csromangathagroup{\csromangathastanza{\csromanfolio{383}\csromangathabat{เย จ รูเป ปมาณิํสุ\csromansharedfootnote{e75f1b7790ceceaf}{เย จ รูเปน ปามิํสุ (สี, สฺยา, กํ, อิ)},}{เย จ โฆเสน อนฺวคู.} \\ \csromangathabat{ฉนฺทราควสูเปตา,}{นาภิชานนฺติ เต ชนา\csromansharedfootnote{e75f1b7790ceceaf}{เย จ รูเปน ปามิํสุ (สี, สฺยา, กํ, อิ)}.}}\csromangathastanza{\csromangathabat{อชฺฌตฺตญฺจ น ชานาติ,}{พหิทฺธา จ น ปสฺสติ.} \\ \csromangathabat{สมนฺตาวรโณ พาโล,}{ส เว โฆเสน วุยฺหติ.}}\csromangathastanza{\csromangathabat{อชฺฌตฺตญฺจ น ชานาติ,}{พหิทฺธา จ วิปสฺสติ.} \\ \csromangathabat{พหิทฺธา ผลทสฺสาวี,}{โสปิ โฆเสน วุยฺหติ.}}\csromangathastanza{\csromangathabat{อชฺฌตฺตญฺจ ปชานาติ,}{พหิทฺธา จ วิปสฺสติ.} \\ \csromangathabat{วินีวรณทสฺสาวี,}{น โส โฆเสน วุยฺหตีติ.  ปญฺจมํ.}}}

\csromanmatikaanchor{mtk.118}
\csromantocmarkref{subsection}{6. สราคสุตฺต}{mtk.118}
\bookmark[dest={mtk.118},level=2]{6. สราคสุตฺต}
\csromanheader{2}{6. \textbf{สราคสุตฺต}}
\proseitem{66}{จตฺตาโรเม ภิกฺขเว ปุคฺคลา สนฺโต สํวิชฺชมานา โลกสฺมิํ. กตเม จตฺตาโร? สราโค สโทโส สโมโห สมาโน, อิเม โข ภิกฺขเว จตฺตาโร ปุคฺคลา สนฺโต สํวิชฺชมานา โลกสฺมินฺติ.}

\setlength{\csromangathapagewidth}{0pt}
\setlength{\csromangathawakwidth}{0pt}
\csromangathameasuregroup{สารตฺตา รชนีเยสุ, \\ โมเหน อาวุตา\textsuperscript{1} สตฺตา, \\ ราคชํ โทสชญฺจาปิ, \\ กโรนฺตากุสลํ กมฺมํ\textsuperscript{1}, \\ อวิชฺชานิวุตา โปสา, \\ ยถา ธมฺมา ตถา สนฺตา,}{\csromangathabat{สารตฺตา รชนีเยสุ,}{ปิยรูปาภินนฺทิโน.} \\ \csromangathabat{โมเหน อาวุตา\textsuperscript{1} สตฺตา,}{พทฺธา\textsuperscript{1} วฑฺเฒนฺติ พนฺธนํ.} \\ \csromangathabat{ราคชํ โทสชญฺจาปิ,}{โมหชํ จาป’วิทฺทสู.} \\ \csromangathabat{กโรนฺตากุสลํ กมฺมํ\textsuperscript{1},}{สวิฆาตํ ทุขุทฺรยํ.} \\ \csromangathabat{อวิชฺชานิวุตา โปสา,}{อนฺธภูตา อจกฺขุกา.} \\ \csromangathabat{ยถา ธมฺมา ตถา สนฺตา,}{น ตสฺเสวนฺติ\textsuperscript{1} มญฺญเรติ.ฉฏฺฐํ.}}
\csromangathasetleft{สารตฺตา รชนีเยสุ, \\ โมเหน อาวุตา\textsuperscript{1} สตฺตา, \\ ราคชํ โทสชญฺจาปิ, \\ กโรนฺตากุสลํ กมฺมํ\textsuperscript{1}, \\ อวิชฺชานิวุตา โปสา, \\ ยถา ธมฺมา ตถา สนฺตา,}
\csromangathagroup{\csromangathastanza{\csromangathabat{สารตฺตา รชนีเยสุ,}{ปิยรูปาภินนฺทิโน.} \\ \csromangathabat{โมเหน อาวุตา\csromansharedfootnote{cb866b4d528c16a7}{อธมา (สี, สฺยา, กํ, อิ)} สตฺตา,}{พทฺธา\csromansharedfootnote{d70ddbbe9aae4984}{พนฺธา (ก)} วฑฺเฒนฺติ พนฺธนํ.}}\csromangathastanza{\csromangathabat{ราคชํ โทสชญฺจาปิ,}{โมหชํ จาป’วิทฺทสู.} \\ \csromangathabat{กโรนฺตากุสลํ กมฺมํ\csromansharedfootnote{e827cde3acc9e5a5}{ธมฺมํ (ก)},}{สวิฆาตํ ทุขุทฺรยํ.}}\csromangathastanza{\csromangathabat{อวิชฺชานิวุตา โปสา,}{อนฺธภูตา อจกฺขุกา.} \\ \csromangathabat{ยถา ธมฺมา ตถา สนฺตา,}{น ตสฺเสวนฺติ\csromansharedfootnote{901c0a3a8d8e17de}{นสฺเสวนฺติ (สี)} มญฺญเรติ.  ฉฏฺฐํ.}}}

\csromanmatikaanchor{mtk.119}
\csromantocmarkref{subsection}{7. อหิราชสุตฺต}{mtk.119}
\bookmark[dest={mtk.119},level=2]{7. อหิราชสุตฺต}
\csromanheader{2}{7. \textbf{อหิราชสุตฺต}}
\proseitem{67}{เอกํ สมยํ ภควา สาวตฺถิยํ วิหรติ เชตวเน อนาถ\-ปิณฺฑิกสฺส อาราเม. เตน โข ปน สมเยน สาวตฺถิยํ อญฺญตโร ภิกฺขุ อหินา ทฏฺโฐ กาลงฺกโต โหติ. อถ โข สมฺพหุลา ภิกฺขู เยน ภควา เตนุ\-ปสงฺกมิํสุ, อุปสงฺกมิตฺวา ภควนฺตํ อภิวาเทตฺวา เอกมนฺตํ นิสีทิํสุ, เอกมนฺตํ นิสินฺนา โข เต ภิกฺขู ภควนฺตํ เอตทโวจุํ \csromanfolio{384}“อิธ ภนฺเต สาวตฺถิยํ อญฺญตโร ภิกฺขุ อหินา ทฏฺโฐ กาลงฺกโต”ติ.}

\prose{น หิ นูน\csromansharedfootnote{ecf2e104dd9aa9af}{น ห นูน (สี, สฺยา, กํ, อิ)} โส ภิกฺขเว ภิกฺขุ จตฺตาริ อหิราชกุลานิ เมตฺเตน จิตฺเตน ผริ. สเจ หิ โส ภิกฺขเว ภิกฺขุ จตฺตาริ อหิราชกุลานิ เมตฺเตน จิตฺเตน ผเรยฺย, น หิ โส ภิกฺขเว ภิกฺขุ อหินา ทฏฺโฐ กาลํ กเรยฺย.}

\prose{กตมานิ จตฺตาริ? วิรูปกฺขํ อหิราชกุลํ, เอราปถํ อหิราชกุลํ, ฉพฺยาปุตฺตํ อหิราชกุลํ, กณฺหาโคตมกํ อหิราชกุลํ. น หิ นูน โส ภิกฺขเว ภิกฺขุ อิมานิ จตฺตาริ อหิราชกุลานิ เมตฺเตน จิตฺเตน ผริ. สเจ หิ โส ภิกฺขเว ภิกฺขุ อิมานิ จตฺตาริ อหิราชกุลานิ เมตฺเตน จิตฺเตน ผเรยฺย, น หิ โส ภิกฺขเว ภิกฺขุ อหินา ทฏฺโฐ กาลํ กเรยฺย.}

\prose{อนุชานามิ ภิกฺขเว อิมานิ จตฺตาริ อหิราชกุลานิ เมตฺเตน จิตฺเตน ผริตุํ อตฺตคุตฺติยา อตฺตรกฺขาย อตฺตปริตฺตายาติ.}

\prose{\csromansymbolfootnote{*}{วิ 4. 245; ขุ 5. 53 ปิฏฺเฐสุปิ ปสฺสิตพฺพํ. 2. ทิปาทเกหิ (สี, สฺยา, กํ, อิ)}วิรูปกฺเขหิ เม เมตฺตํ, เมตฺตํ เอราปเถหิ เม.}

\setlength{\csromangathapagewidth}{0pt}
\setlength{\csromangathawakwidth}{0pt}
\csromangathameasuregroup{ฉพฺยาปุตฺเตหิ เม เมตฺตํ, \\ อปาทเกหิ เม เมตฺตํ, \\ จตุปฺปเทหิ เม เมตฺตํ, \\ มา มํ อปาทโก หิํสิ, \\ มา มํ จตุปฺปโท หิํสิ, \\ สพฺเพ สตฺตา สพฺเพ ปาณา, \\ สพฺเพ ภทฺรานิ ปสฺสนฺตุ, \\ อปฺปมาโณ พุทฺโธ, \\ อปฺปมาโณ สํโฆ, \\ อหิวิจฺฉิกา สตปที, \\ กตา เม รกฺขา กตา เม ปริตฺตา\textsuperscript{1},}{\csromangathabat{ฉพฺยาปุตฺเตหิ เม เมตฺตํ,}{เมตฺตํ กณฺหาโคตมเกหิ จ.} \\ \csromangathabat{อปาทเกหิ เม เมตฺตํ,}{เมตฺตํ ทฺวิปาทเกหิ2 เม.} \\ \csromangathabat{จตุปฺปเทหิ เม เมตฺตํ,}{เมตฺตํ พหุปฺปเทหิ เม.} \\ \csromangathabat{มา มํ อปาทโก หิํสิ,}{มา มํ หิํสิ ทฺวิปาทโก\textsuperscript{1}.} \\ \csromangathabat{มา มํ จตุปฺปโท หิํสิ,}{มา มํ หิํสิ พหุปฺปโท.} \\ \csromangathabat{สพฺเพ สตฺตา สพฺเพ ปาณา,}{สพฺเพ ภูตา จ เกวลา.} \\ \csromangathabat{สพฺเพ ภทฺรานิ ปสฺสนฺตุ,}{มา กญฺจิ\textsuperscript{1} ปาปมาคมา.} \\ \csromangathabat{อปฺปมาโณ พุทฺโธ,}{อปฺปมาโณ ธมฺโม.} \\ \csromangathabat{อปฺปมาโณ สํโฆ,}{ปมาณวนฺตานิ สรีสปานิ\textsuperscript{1}.} \\ \csromangathabat{อหิวิจฺฉิกา สตปที,}{อุณฺณนาภี สรพู มูสิกา.} \\ \csromangathabat{กตา เม รกฺขา กตา เม ปริตฺตา\textsuperscript{1},}{ปฏิกฺกมนฺตุภูตานิ.}}
\csromangathasetleft{ฉพฺยาปุตฺเตหิ เม เมตฺตํ, \\ อปาทเกหิ เม เมตฺตํ, \\ จตุปฺปเทหิ เม เมตฺตํ, \\ มา มํ อปาทโก หิํสิ, \\ มา มํ จตุปฺปโท หิํสิ, \\ สพฺเพ สตฺตา สพฺเพ ปาณา, \\ สพฺเพ ภทฺรานิ ปสฺสนฺตุ, \\ อปฺปมาโณ พุทฺโธ, \\ อปฺปมาโณ สํโฆ, \\ อหิวิจฺฉิกา สตปที, \\ กตา เม รกฺขา กตา เม ปริตฺตา\textsuperscript{1},}
\csromangathagroup{\csromangathastanza{\csromangathabat{ฉพฺยาปุตฺเตหิ เม เมตฺตํ,}{เมตฺตํ กณฺหา\-โคตม\-เกหิ จ.} \\ \csromangathabat{อปาทเกหิ เม เมตฺตํ,}{เมตฺตํ ทฺวิปาทเกหิ2 เม.}}\csromangathastanza{\csromangathabat{จตุปฺปเทหิ เม เมตฺตํ,}{เมตฺตํ พหุปฺปเทหิ เม.} \\ \csromangathabat{มา มํ อปาทโก หิํสิ,}{มา มํ หิํสิ ทฺวิปาทโก\csromansharedfootnote{4516fbe782b97f60}{ทิปาทโก (สี, สฺยา, กํ, อิ)}.}}\csromangathastanza{\csromangathabat{มา มํ จตุปฺปโท หิํสิ,}{มา มํ หิํสิ พหุปฺปโท.} \\ \csromangathabat{สพฺเพ สตฺตา สพฺเพ ปาณา,}{สพฺเพ ภูตา จ เกวลา.}}\csromangathastanza{\csromangathabat{สพฺเพ ภทฺรานิ ปสฺสนฺตุ,}{มา กญฺจิ\csromansharedfootnote{45e3891a7a032d1d}{กิญฺจิ (สฺยา, กํ, ก)} ปาปมาคมา.} \\ \csromangathabat{อปฺปมาโณ พุทฺโธ,}{อปฺปมาโณ ธมฺโม.}}\csromangathastanza{\csromangathabat{อปฺปมาโณ สํโฆ,}{ปมาณวนฺตานิ สรีสปานิ\csromansharedfootnote{38f407814ca3e15a}{สิริํสปานิ (สี, สฺยา, กํ, อิ) 6. กตํ เม ปริตฺตํ (?)}.} \\ \csromangathabat{อหิวิจฺฉิกา สตปที,}{อุณฺณนาภี สรพู มูสิกา.} \\ \csromangathabat{กตา เม รกฺขา กตา เม ปริตฺตา\csromansharedfootnote{3fc0b43edcc06d3d}{กตํ เม ปริตฺตํ (?)},}{ปฏิกฺกมนฺตุภูตานิ.}}}

\prose{โสหํ นโม ภควโต, นโม สตฺตนฺนํ สมฺมา\-สมฺพุทฺธานนฺติ.  }

\csromanmatikaanchor{mtk.120}
\csromantocmarkref{subsection}{8. เทวทตฺตสุตฺต}{mtk.120}
\bookmark[dest={mtk.120},level=2]{8. เทวทตฺตสุตฺต}
\csromanmarkh{(7)}\csromanmarkhii{2. ปตฺตกมฺม\-วคฺค 8. เทวทตฺตสุตฺต}\csromanheadernomark{2}{(7) 2. ปตฺตกมฺม\-วคฺค 8. เทวทตฺตสุตฺต}
\proseitem{68}{\csromanfolio{385}เอกํ สมยํ ภควา ราชคเห วิหรติ คิชฺฌกูเฏ ปพฺพเต อจิรปกฺกนฺเต เทวทตฺเต. ตตฺร โข ภควา เทวทตฺตํ อารพฺภ ภิกฺขู อามนฺเตสิ “อตฺตวธาย ภิกฺขเว\csromansharedfootnote{81653a4de1a48581}{วิ 4. 346; สํ 1. 436 ปิฏฺเฐสุปิ.} เทวทตฺตสฺส ลาภ\-สกฺการ\-สิโลโก อุทปาทิ, ปราภวาย ภิกฺขเว เทวทตฺตสฺส ลาภ\-สกฺการ\-สิโลโก อุทปาทิ.}

\prose{เสยฺยถาปิ ภิกฺขเว กทลี อตฺตวธาย ผลํ เทติ, ปราภวาย ผลํ เทติ. เอวเมวํ โข ภิกฺขเว อตฺตวธาย เทวทตฺตสฺส ลาภ\-สกฺการ\-สิโลโก อุทปาทิ, ปราภวาย เทวทตฺตสฺส ลาภ\-สกฺการ\-สิโลโก อุทปาทิ.}

\prose{เสยฺยถาปิ ภิกฺขเว เวฬุ อตฺตวธาย ผลํ เทติ, ปราภวาย ผลํ เทติ. เอวเมวํ โข ภิกฺขเว อตฺตวธาย เทวทตฺตสฺส ลาภ\-สกฺการ\-สิโลโก อุทปาทิ, ปราภวาย เทวทตฺตสฺส ลาภ\-สกฺการ\-สิโลโก อุทปาทิ.}

\prose{เสยฺยถาปิ ภิกฺขเว นโฬ อตฺตวธาย ผลํ เทติ, ปราภวาย ผลํ เทติ. เอวเมวํ โข ภิกฺขเว อตฺตวธาย เทวทตฺตสฺส ลาภ\-สกฺการ\-สิโลโก อุทปาทิ, ปราภวาย เทวทตฺตสฺส ลาภ\-สกฺการ\-สิโลโก อุทปาทิ.}

\prose{เสยฺยถาปิ ภิกฺขเว อสฺสตรี อตฺตวธาย คพฺภํ คณฺหาติ, ปราภวาย คพฺภํ คณฺหาติ. เอวเมวํ โข ภิกฺขเว อตฺตวธาย เทวทตฺตสฺส ลาภ\-สกฺการ\-สิโลโก อุทปาทิ, ปราภวาย เทวทตฺตสฺส ลาภ\-สกฺการ\-สิโลโก อุทปาที”ติ.}

\setlength{\csromangathapagewidth}{0pt}
\setlength{\csromangathawakwidth}{0pt}
\csromangathameasuregroup{ผลํ เว กทลิํ หนฺติ, \\ สกฺกาโร กาปุริสํ หนฺติ,}{\csromangathabat{ผลํ เว กทลิํ หนฺติ,}{ผลํ เวฬุํ, ผลํ นฬํ.} \\ \csromangathabat{สกฺกาโร กาปุริสํ หนฺติ,}{คพฺโภ อสฺสตริํ ยถาติ\textsuperscript{1}.อฏฺฐมํ.}}
\csromangathasetleft{ผลํ เว กทลิํ หนฺติ, \\ สกฺกาโร กาปุริสํ หนฺติ,}
\csromangathagroup{\csromangathastanza{\csromangathabat{ผลํ เว กทลิํ หนฺติ,}{ผลํ เวฬุํ, ผลํ นฬํ.} \\ \csromangathabat{สกฺกาโร กาปุริสํ หนฺติ,}{คพฺโภ อสฺสตริํ ยถาติ\csromansharedfootnote{2ba45b939145d8c6}{วิ 4. 346; สํ 1. 156, 436; ขุ 10. 109 ปิฏฺเฐสุปิ.}.  อฏฺฐมํ.}}}

\csromanmatikaanchor{mtk.121}
\csromantocmarkref{subsection}{9. ปธานสุตฺต}{mtk.121}
\bookmark[dest={mtk.121},level=2]{9. ปธานสุตฺต}
\csromanheader{2}{9. \textbf{ปธานสุตฺต}}
\proseitem{69}{จตฺตาริมานิ ภิกฺขเว ปธานานิ. กตมานิ จตฺตาริ? สํวร\-ปฺปธานํ ปหาน\-ปฺปธานํ ภาวนา\-ปฺปธานํ อนุรกฺขณา\-ปฺปธานํ. กตมญฺจ ภิกฺขเว สํวร\-ปฺปธานํ? อิธ ภิกฺขเว ภิกฺขุ อนุปฺปนฺนานํ ปาปกานํ อกุสลานํ ธมฺมานํ \csromanfolio{386}อนุปฺปาทาย ฉนฺทํ ชเนติ วายมติ, วีริยํ อารภติ, จิตฺตํ ปคฺคณฺหาติ ปทหติ. อิทํ วุจฺจติ ภิกฺขเว สํวร\-ปฺปธานํ.}

\prose{กตมญฺจ ภิกฺขเว ปหาน\-ปฺปธานํ? อิธ ภิกฺขเว ภิกฺขุ อุปฺปนฺนานํ ปาปกานํ อกุสลานํ ธมฺมานํ ปหานาย ฉนฺทํ ชเนติ วายมติ, วีริยํ อารภติ, จิตฺตํ ปคฺคณฺหาติ ปทหติ. อิทํ วุจฺจติ ภิกฺขเว ปหาน\-ปฺปธานํ.}

\prose{กตมญฺจ ภิกฺขเว ภาวนา\-ปฺปธานํ? อิธ ภิกฺขเว ภิกฺขุ อนุปฺปนฺนานํ กุสลานํ ธมฺมานํ อุปฺปาทาย ฉนฺทํ ชเนติ วายมติ, วีริยํ อารภติ, จิตฺตํ ปคฺคณฺหาติ ปทหติ. อิทํ วุจฺจติ ภิกฺขเว ภาวนา\-ปฺปธานํ.}

\prose{กตมญฺจ ภิกฺขเว อนุรกฺขณา\-ปฺปธานํ? อิธ ภิกฺขเว ภิกฺขุ อุปฺปนฺนานํ กุสลานํ ธมฺมานํ ฐิติยา อสมฺโมสาย ภิยฺโยภาวาย เวปุลฺลาย ภาวนาย ปาริปูริยา ฉนฺทํ ชเนติ วายมติ, วีริยํ อารภติ, จิตฺตํ ปคฺคณฺหาติ ปทหติ. อิทํ วุจฺจติ ภิกฺขเว อนุรกฺขณา\-ปฺปธานํ. อิมานิ โข ภิกฺขเว จตฺตาริ ปธานานีติ.}

\setlength{\csromangathapagewidth}{0pt}
\setlength{\csromangathawakwidth}{0pt}
\csromangathameasuregroup{สํวโร จ ปหานญฺจ, \\ เอเต ปธานา จตฺตาโร, \\ โย หิ\textsuperscript{1} ภิกฺขุ อิธาตาปี,}{\csromangathabat{สํวโร จ ปหานญฺจ,}{ภาวนา อนุรกฺขณา.} \\ \csromangathabat{เอเต ปธานา จตฺตาโร,}{เทสิตาทิจฺจพนฺธุนา.} \\ \csromangathabat{โย หิ\textsuperscript{1} ภิกฺขุ อิธาตาปี,}{ขยํ ทุกฺขสฺส ปาปุเณติ.นวมํ.}}
\csromangathasetleft{สํวโร จ ปหานญฺจ, \\ เอเต ปธานา จตฺตาโร, \\ โย หิ\textsuperscript{1} ภิกฺขุ อิธาตาปี,}
\csromangathagroup{\csromangathastanza{\csromangathabat{สํวโร จ ปหานญฺจ,}{ภาวนา อนุรกฺขณา.} \\ \csromangathabat{เอเต ปธานา จตฺตาโร,}{เทสิตาทิจฺจพนฺธุนา.} \\ \csromangathabat{โย หิ\csromansharedfootnote{5e6e3a99def97b0f}{เยหิ (?)} ภิกฺขุ อิธาตาปี,}{ขยํ ทุกฺขสฺส ปาปุเณติ.  นวมํ.}}}

\csromanmatikaanchor{mtk.122}
\csromantocmarkref{subsection}{10. อธมฺมิกสุตฺต}{mtk.122}
\bookmark[dest={mtk.122},level=2]{10. อธมฺมิกสุตฺต}
\csromanheader{2}{10. \textbf{อธมฺมิกสุตฺต}}
\proseitem{70}{ยสฺมิํ ภิกฺขเว สมเย ราชาโน อธมฺมิกา โหนฺติ, ราชายุตฺตาปิ ตสฺมิํ สมเย อธมฺมิกา โหนฺติ, ราชายุตฺเตสุ อธมฺมิเกส พฺราหฺมณ\-คหปติกาปิ ตสฺมิํ สมเย อธมฺมิกา โหนฺติ, พฺราหฺมณคหปติเกสุ อธมฺมิเกสุ เนคมชานปทาปิ ตสฺมิํ สมเย อธมฺมิกา โหนฺติ, เนคม\-ชานปเทสุ อธมฺมิเกสุ วิสมํ จนฺทิมสูริยา ปริวตฺตนฺติ, วิสมํ จนฺทิม\-สูริเยสุ ปริวตฺตนฺเตสุ วิสมํ นกฺขตฺตานิ ตารกรูปานิ ปริวตฺตนฺติ, วิสมํ นกฺขตฺเตสุ ตารกรูเปสุ ปริวตฺตนฺเตสุ วิสมํ รตฺตินฺทิวา\csromansharedfootnote{b5825402cbce0a9f}{รตฺติทิวา (ก)} ปริวตฺตนฺติ, วิสมํ รตฺตินฺทิเวสุ ปริวตฺตนฺเตสุ วิสมํ มาสทฺธมาสา ปริวตฺตนฺติ, วิสมํ มาสทฺธมาเสสุ ปริวตฺตนฺเตสุ วิสมํ อุตุสํวจฺฉรา ปริวตฺตนฺติ, วิสมํ อุตุสํวจฺฉเรสุ ปริวตฺตนฺเตสุ วิสมํ วาตา วายนฺติ วิสมา}

\vspace{\tipitakaparskip}
\csromancenter{(7) 2. ปตฺตกมฺม\-วคฺค อปญฺชสา, วิสมํ วาเตสุ วายนฺเตสุ วิสเมสุ อปญฺชเสสุ เทวตา ปริกุปิตา ภวนฺติ, เทวตาสุ ปริกุปิตาสุ เทโว น สมฺมา ธารํ อนุปฺปเวจฺฉติ, เทเว น สมฺมา ธารํ อนุปฺป\-เวจฺฉ\-นฺเต วิสมปากานิ\csromansharedfootnote{30d49229dbd394e5}{วิสมปากีนิ (สี, สฺยา, กํ), วิสมํ ปากานิ (ก)} สสฺสานิ ภวนฺติ, วิสมปากานิ ภิกฺขเว สสฺสานิ มนุสฺสา ปริภุญฺชนฺตา อปฺปายุกา โหนฺติ ทุพฺพณฺณา จ พวฺหาพาธา\csromansharedfootnote{32e6e1e55ef6aeec}{พหฺวาพาธา (ก)} จ.}
\prose{\csromanfolio{387}ยสฺมิํ ภิกฺขเว สมเย ราชาโน ธมฺมิกา โหนฺติ, ราชายุตฺตาปิ ตสฺมิํ สมเย ธมฺมิกา โหนฺติ, ราชายุตฺเตสุ ธมฺมิเกสุ พฺราหฺมณ\-คหปติกาปิ ตสฺมิํ สมเย ธมฺมิกา โหนฺติ, พฺราหฺมณคหปติเกสุ ธมฺมิเกสุ เนคมชานปทาปิ ตสฺมิํ สมเย ธมฺมิกา โหนฺติ, เนคม\-ชานปเทสุ ธมฺมิเกสุ สมํ จนฺทิมสูริยา ปริวตฺตนฺติ, สมํ จนฺทิม\-สูริเยสุ ปริวตฺตนฺเตสุ สมํ นกฺขตฺตานิ ตารกรูปาติ ปริวตฺตนฺติ, สมํ นกฺขตฺเตสุ ตารกรูเปสุ ปริวตฺตนฺเตสุ สมํ รตฺตินฺทิวา ปริวตฺตนฺติ, สมํ รตฺตินฺทิเวสุ ปริวตฺตนฺเตสุ สมํ มาสทฺธมาสา ปริวตฺตนฺติ, สมํ มาสทฺธมาเสสุ ปริวตฺตนฺเตสุ สมํ อุตุสํวจฺฉรา ปริวตฺตนฺติ, สมํ อุตุสํวจฺฉเรสุ ปริวตฺตนฺเตสุ สมํ วาตา วายนฺติ สมา ปญฺชสา, สมํ วาเตสุ วายนฺเตสุ สเมสุ ปญฺชเสสุ เทวตา อปริกุปิตา ภวนฺติ, เทวตาสุ อปริกุปิตาสุ เทโว สมฺมา ธารํ อนุปฺปเวจฺฉติ, เทเว สมฺมา ธารํ อนุปฺป\-เวจฺฉ\-นฺเต สมปากานิ สสฺสานิ ภวนฺติ, สมปากานิ ภิกฺขเว สสฺสานิ มนุสฺสา ปริภุญฺชนฺตา ทีฆายุกา จ โหนฺติ วณฺณวนฺโต จ พลวนฺโต จ อปฺปาพาธา จาติ.}

\setlength{\csromangathapagewidth}{0pt}
\setlength{\csromangathawakwidth}{0pt}
\csromangathameasuregroup{คุนฺนํ เจ ตรมานานํ, \\ สพฺพา ตา ชิมฺหํ คจฺฉนฺติ, \\ เอวเมวํ มนุสฺเสสุ, \\ โส เจ อธมฺมํ จรติ, \\ สพฺพํ รฏฺฐํ ทุกฺขํ เสติ, \\ คุนฺนํ เจ ตรมานานํ, \\ สพฺพา ตา อุชุํ คจฺฉนฺติ, \\ เอวเมวํ มนุสฺเสสุ, \\ โส สเจ\textsuperscript{1} ธมฺมํ จรติ, \\ สพฺพํ รฏฺฐํ สุขํ เสติ,}{\csromangathabat{คุนฺนํ เจ ตรมานานํ,}{ชิมฺหํ คจฺฉติ ปุงฺคโว.} \\ \csromangathabat{สพฺพา ตา ชิมฺหํ คจฺฉนฺติ,}{เนตฺเต ชิมฺหํ คเต สติ.} \\ \csromangathabat{เอวเมวํ มนุสฺเสสุ,}{โย โหติ เสฏฺฐสมฺมโต.} \\ \csromangathabat{โส เจ อธมฺมํ จรติ,}{ปเคว อิตรา ปชา.} \\ \csromangathabat{สพฺพํ รฏฺฐํ ทุกฺขํ เสติ,}{ราชา เจ โหติ อธมฺมิโก.} \\ \csromangathabat{คุนฺนํ เจ ตรมานานํ,}{อุชุํ คจฺฉติ ปุงฺคโว.} \\ \csromangathabat{สพฺพา ตา อุชุํ คจฺฉนฺติ,}{เนตฺเต อุชุํ คเต สติ.} \\ \csromangathabat{เอวเมวํ มนุสฺเสสุ,}{โย โหติ เสฏฺฐสมฺมโต.} \\ \csromangathabat{โส สเจ\textsuperscript{1} ธมฺมํ จรติ,}{ปเคว อิตรา ปชา.} \\ \csromangathabat{สพฺพํ รฏฺฐํ สุขํ เสติ,}{ราชา เจ โหติ ธมฺมิโกติ.ทสมํ.}}
\csromangathasetleft{คุนฺนํ เจ ตรมานานํ, \\ สพฺพา ตา ชิมฺหํ คจฺฉนฺติ, \\ เอวเมวํ มนุสฺเสสุ, \\ โส เจ อธมฺมํ จรติ, \\ สพฺพํ รฏฺฐํ ทุกฺขํ เสติ, \\ คุนฺนํ เจ ตรมานานํ, \\ สพฺพา ตา อุชุํ คจฺฉนฺติ, \\ เอวเมวํ มนุสฺเสสุ, \\ โส สเจ\textsuperscript{1} ธมฺมํ จรติ, \\ สพฺพํ รฏฺฐํ สุขํ เสติ,}
\csromangathagroupwithclosermidruled{\csromangathastanza{\csromangathabat{คุนฺนํ เจ ตรมานานํ,}{ชิมฺหํ คจฺฉติ ปุงฺคโว.} \\ \csromangathabat{สพฺพา ตา ชิมฺหํ คจฺฉนฺติ,}{เนตฺเต ชิมฺหํ คเต สติ.}}\csromangathastanza{\csromangathabat{เอวเมวํ มนุสฺเสสุ,}{โย โหติ เสฏฺฐสมฺมโต.} \\ \csromangathabat{โส เจ อธมฺมํ จรติ,}{ปเคว อิตรา ปชา.}}\csromangathastanza{\csromangathabat{สพฺพํ รฏฺฐํ ทุกฺขํ เสติ,}{ราชา เจ โหติ อธมฺมิโก.} \\ \csromangathabat{คุนฺนํ เจ ตรมานานํ,}{อุชุํ คจฺฉติ ปุงฺคโว.} \\ \csromangathabat{สพฺพา ตา อุชุํ คจฺฉนฺติ,}{เนตฺเต อุชุํ คเต สติ.}}\csromangathastanza{\csromanfolio{388}\csromangathabat{เอวเมวํ มนุสฺเสสุ,}{โย โหติ เสฏฺฐสมฺมโต.} \\ \csromangathabat{โส สเจ\csromansharedfootnote{22193afc8e4961c8}{โส เจว (สี, อิ), โส เจ (สฺยา)} ธมฺมํ จรติ,}{ปเคว อิตรา ปชา.} \\ \csromangathabat{สพฺพํ รฏฺฐํ สุขํ เสติ,}{ราชา เจ โหติ ธมฺมิโกติ.  ทสมํ.}}}{ปตฺตกมฺม\-วคฺโค ทุติโย.}
\csromancenter{\textbf{ตสฺสุทฺทานํ}}
\setlength{\csromangathapagewidth}{0pt}
\setlength{\csromangathawakwidth}{0pt}
\csromangathameasuregroup{ปตฺตกมฺมํ อานณฺยโก\textsuperscript{1}, \\ สราคอหิราชา เทวทตฺโต,}{\csromangathabat{ปตฺตกมฺมํ อานณฺยโก\textsuperscript{1},}{สพฺรหฺมนิรยา รูเปน ปญฺจมํ.} \\ \csromangathabat{สราคอหิราชา เทวทตฺโต,}{ปธานํ อธมฺมิเกน จาติ.}}
\csromangathasetleft{ปตฺตกมฺมํ อานณฺยโก\textsuperscript{1}, \\ สราคอหิราชา เทวทตฺโต,}
\csromangathagroup{\csromangathastanza{\csromangathabat{ปตฺตกมฺมํ อานณฺยโก\csromansharedfootnote{4a3479621a7b24e7}{อนณโก (สี, อิ), อนณฺยโก (ก)},}{สพฺรหฺมนิรยา รูเปน ปญฺจมํ.} \\ \csromangathabat{สราคอหิราชา เทวทตฺโต,}{ปธานํ อธมฺมิเกน จาติ.}}}

\csromanmatikaanchor{mtk.123}
\csromantocmarknopage{section}{(8) 3. อปณฺณกวคฺค}{mtk.123}
\bookmark[dest={mtk.123},level=1]{(8) 3. อปณฺณกวคฺค}
\csromanmarkh{(8)}\csromanmarkhii{3. อปณฺณกวคฺค}\csromanheadernomark{1}{\textbf{(8) 3. อปณฺณกวคฺค}}
\csromanmatikaanchor{mtk.124}
\csromantocmarkref{subsection}{1. ปธานสุตฺต}{mtk.124}
\bookmark[dest={mtk.124},level=2]{1. ปธานสุตฺต}
\csromanheader{2}{1. \textbf{ปธานสุตฺต}}
\proseitem{71}{จตูหิ ภิกฺขเว ธมฺเมหิ สมนฺนาคโต ภิกฺขุ อปณฺณก\-ปฺปฏิปทํ ปฏิปนฺโน โหติ, โยนิ จสฺส อารทฺธา โหติ อาสวานํ ขยาย. กตเมหิ จตูหิ? อิธ ภิกฺขเว ภิกฺขุ สีลวา โหติ, พหุสฺสุโต โหติ, อารทฺธวีริโย โหติ, ปญฺญวา โหติ. อิเมหิ โข ภิกฺขเว จตูหิ ธมฺเมหิ สมนฺนาคโต ภิกฺขุ อปณฺณก\-ปฺปฏิปทํ ปฏิปนฺโน โหติ, โยนิ จสฺส อารทฺธา โหติ อาสวานํ ขยายาติ.  ปฐมํ.}

\csromanmatikaanchor{mtk.125}
\csromantocmarkref{subsection}{2. สมฺมาทิฏฺฐิสุตฺต}{mtk.125}
\bookmark[dest={mtk.125},level=2]{2. สมฺมาทิฏฺฐิสุตฺต}
\csromanheader{2}{2. สมฺมาทิฏฺฐิ\-สุตฺต}
\proseitem{72}{จตูหิ ภิกฺขเว ธมฺเมหิ สมนฺนาคโต ภิกฺขุ อปณฺณก\-ปฺปฏิปทํ ปฏิปนฺโน โหติ, โยนิ จสฺส อารทฺธา โหติ อาสวานํ ขยาย. กตเมหิ จตูหิ? เนกฺขมฺม\-วิตกฺเกน อพฺยาปาทวิตกฺเกน อวิหิํสา\-วิตกฺเกน สมฺมาทิฏฺฐิยา. อิเมหิ โข ภิกฺขเว จตูหิ ธมฺเมหิ สมนฺนาคโต ภิกฺขุ อปณฺณก\-ปฺปฏิปทํ ปฏิปนฺโน โหติ, โยนิ จสฺส อารทฺธา โหติ อาสวานํ ขยายาติ.  ทุติยํ.}

\csromanmatikaanchor{mtk.126}
\csromantocmarkref{subsection}{3. สปฺปุริสสุตฺต}{mtk.126}
\bookmark[dest={mtk.126},level=2]{3. สปฺปุริสสุตฺต}
\csromanmarkh{(8)}\csromanmarkhii{3. อปณฺณกวคฺค 3. สปฺปุริสสุตฺต}\csromanheadernomark{2}{(8) 3. อปณฺณกวคฺค \textbf{3. สปฺปุริสสุตฺต}}
\proseitem{73}{\csromanfolio{389}จตูหิ ภิกฺขเว ธมฺเมหิ สมนฺนาคโต อสปฺปุริโส เวทิตพฺโพ. กตเมหิ จตูหิ? อิธ ภิกฺขเว อสปฺปุริโส โย โหติ ปรสฺส อวณฺโณ ตํ อปุฏฺโฐปิ ปาตุ กโรติ, โก ปน วาโท ปุฏฺฐสฺส. ปุฏฺโฐ โข ปน ปญฺหาภินีโต อหาเปตฺวา อลมฺพิตฺวา ปริปูรํ วิตฺถาเรน ปรสฺส อวณฺณํ ภาสิตา โหติ, เวทิตพฺพเมตํ ภิกฺขเว “อสปฺปุริโส อยํ ภวนฺ”ติ.}

\prose{ปุน จปรํ ภิกฺขเว อสปฺปุริโส โย โหติ ปรสฺส วณฺโณ ตํ ปุฏฺโฐปิ น ปาตุ กโรติ, โก ปน วาโท อปุฏฺฐสฺส. ปุฏฺโฐ โข ปน ปญฺหาภินีโต หาเปตฺวา ลมฺพิตฺวา อปริปูรํ อวิตฺถาเรน ปรสฺส วณฺณํ ภาสิตา โหติ, เวทิตพฺพเมตํ ภิกฺขเว “อสปฺปุริโส อยํ ภวนฺ”ติ.}

\prose{ปุน จปรํ ภิกฺขเว อสปฺปุริโส โย โหติ อตฺตโน อวณฺโณ ตํ ปุฏฺโฐปิ น ปาตุ กโรติ, โก ปน วาโท อปุฏฺฐสฺส. ปุฏฺโฐ โข ปน ปญฺหาภินีโต หาเปตฺวา ลมฺพิตฺวา อปริปูรํ อวิตฺถาเรน อตฺตโน อวณฺณํ ภาสิตา โหติ, เวทิตพฺพเมตํ ภิกฺขเว “อสปฺปุริโส อยํ ภวนฺ”ติ.}

\prose{ปุน จปรํ ภิกฺขเว อสปฺปุริโส โย โหติ อตฺตโน วณฺโณ ตํ อปุฏฺโฐปิ ปาตุ กโรติ, โก ปน วาโท ปุฏฺฐสฺส. ปุฏฺโฐ โข ปน ปญฺหาภินีโต อหาเปตฺวา อลมฺพิตฺวา ปริปูรํ วิตฺถาเรน อตฺตโน วณฺณํ ภาสิตา โหติ. เวทิตพฺพเมตํ ภิกฺขเว “อสปฺปุริโส อยํ ภวนฺ”ติ. อิเมหิ โข ภิกฺขเว จตูหิ ธมฺเมหิ สมนฺนาคโต อสปฺปุริโส เวทิตพฺโพ.}

\prose{จตูหิ ภิกฺขเว ธมฺเมหิ สมนฺนาคโต สปฺปุริโส เวทิตพฺโพ. กตเมหิ จตูหิ? อิธ ภิกฺขเว สปฺปุริโส โย โหติ ปรสฺส อวณฺโณ ตํ ปุฏฺโฐปิ น ปาตุ กโรติ, โก ปน วาโท อปุฏฺฐสฺส. ปุฏฺโฐ โข ปน ปญฺหาภินีโต หาเปตฺวา ลมฺพิตฺวา อปริปูรํ อวิตฺถาเรน ปรสฺส อวณฺณํ ภาสิตา โหติ. เวทิตพฺพเมตํ ภิกฺขเว “สปฺปุริโส อยํ ภวนฺ”ติ.}

\csromanmatikaanchor{mtk.127}
\csromantocmarkref{subsection}{4-5. อคฺคสุตฺต}{mtk.127}
\bookmark[dest={mtk.127},level=2]{4-5. อคฺคสุตฺต}
\prose{\csromanfolio{390}ปุน จปรํ ภิกฺขเว สปฺปุริโส โย โหติ ปรสฺส วณฺโณ ตํ อปุฏฺโฐปิ ปาตุ กโรติ, โก ปน วาโท ปุฏฺฐสฺส. ปุฏฺโฐ โข ปน ปญฺหาภินีโต อหาเปตฺวา อลมฺพิตฺวา ปริปูรํ วิตฺถาเรน ปรสฺส วณฺณํ ภาสิตา โหติ. เวทิตพฺพเมตํ ภิกฺขเว “สปฺปุริโส อยํ ภวนฺ”ติ.}

\prose{ปุน จปรํ ภิกฺขเว สปฺปุริโส โย โหติ อตฺตโน อวณฺโณ ตํ อปุฏฺโฐปิ ปาตุ กโรติ, โก ปน วาโท ปุฏฺฐสฺส. ปุฏฺโฐ โข ปน ปญฺหาภินีโต อหาเปตฺวา อลมฺพิตฺวา ปริปูรํ วิตฺถาเรน อตฺตโน อวณฺณํ ภาสิตา โหติ. เวทิตพฺพเมตํ ภิกฺขเว “สปฺปุริโส อยํ ภวนฺ”ติ.}

\prose{ปุน จปรํ ภิกฺขเว สปฺปุริโส โย โหติ อตฺตโน วณฺโณ ตํ ปุฏฺโฐปิ น ปาตุ กโรติ, โก ปน วาโท อปุฏฺฐสฺส. ปุฏฺโฐ โข ปน ปญฺหาภินีโต หาเปตฺวา ลมฺพิตฺวา อปริปูรํ อวิตฺถาเรน อตฺตโน วณฺณํ ภาสิตา โหติ. เวทิตพฺพเมตํ ภิกฺขเว “สปฺปุริโส อยํ ภวนฺ”ติ. อิเมหิ โข ภิกฺขเว จตูหิ ธมฺเมหิ สมนฺนาคโต สปฺปุริโส เวทิตพฺโพ.}

\prose{เสยฺยถาปิ ภิกฺขเว วธุกา ยญฺญเทว รตฺติํ วา ทิวํ วา อานีตา โหติ, ตาวเทวสฺสา ติพฺพํ หิโรตฺตปฺปํ ปจฺจุปฏฺฐิตํ โหติ สสฺสุยาปิ สสุเรปิ สามิเกปิ อนฺตมโส ทาส\-กมฺมกร\-โปริเสสุ. สา อปเรน สมเยน สํวาสมนฺวาย วิสฺสาสมนฺวาย สสฺสุมฺปิ สสุรมฺปิ สามิกมฺปิ เอวมาห “อเปถ กิํ ปน ตุเมฺห ชานาถา”ติ. เอวเมวํ โข ภิกฺขเว อิเธกจฺโจ ภิกฺขุ ยญฺญเทว รตฺติํ วา ทิวํ วา อคารสฺมา อนคาริยํ ปพฺพชิโต โหติ, ตาวเทวสฺส ติพฺพํ หิโรตฺตปฺปํ ปจฺจุปฏฺฐิตํ โหติ ภิกฺขูสุ ภิกฺขุนีสุ อุปาสเกสุ อุปาสิกาสุ อนฺตมโส อารามิกสมณุทฺเทเสสุ. โส อปเรน สมเยน สํวาสมนฺวาย วิสฺสาสมนฺวาย อาจริยมฺปิ อุปชฺฌายมฺปิ เอวมาห “อเปถ กิํ ปน ตุเมฺห ชานาถา”ติ. ตสฺมาติห ภิกฺขเว เอวํ สิกฺขิตพฺพํ “อธุนาคต\-วธุกา\-สเมน เจตสา วิหริสฺสามา”ติ. เอวํ หิ โว ภิกฺขเว สิกฺขิตพฺพนฺติ.  ตติยํ.}

\prose{4. ปฐม\-อคฺค\-สุตฺต}

\proseitem{74}{จตฺตริมานิ ภิกฺขเว อคฺคานิ. กตมานิ จตฺตาริ? สีลคฺคํ สมาธิคฺคํ\csromansharedfootnote{029a71b87390ee5a}{สมาธคฺคํ (สี, สฺยา, กํ)} ปญฺญาคฺคํ วิมุตฺตคฺคํ. อิมานิ โข ภิกฺขเว จตฺตาริ อคฺคานีติ.  จตุตฺถํ.}

\chapterheadpageread{(8) 3. อปณฺณกวคฺค 5. ทุติย\-อคฺค\-สุตฺต}
\proseitem{75}{\csromanfolio{391}จตฺตาริมานิ ภิกฺขเว อคฺคานิ. กตมานิ จตฺตาริ? รูปคฺคํ เวทนาคฺคํ สญฺญาคฺคํ ภวคฺคํ. อิมานิ โข ภิกฺขเว จตฺตาริ อคฺคานีติ.  ปญฺจมํ.}

\csromanmatikaanchor{mtk.128}
\csromantocmarkref{subsection}{6. กุสินารสุตฺต}{mtk.128}
\bookmark[dest={mtk.128},level=2]{6. กุสินารสุตฺต}
\csromanheader{2}{6. \textbf{กุสินารสุตฺต}}
\proseitem{76}{เอกํ สมยํ ภควา กุสินารายํ วิหรติ อุปวตฺตเน มลฺลานํ สาลวเน อนฺตเรน ยมกสาลานํ ปรินิพฺพาน\-สมเย. ตตฺร โข ภควา ภิกฺขู อามนฺเตสิ “ภิกฺขโว”ติ. “ภทนฺเต”ติ เต ภิกฺขู ภควโต ปจฺจสฺโสสุํ. ภควา เอตทโวจ\csromandash{}}

\prose{“สิยา โข ปน ภิกฺขเว\csromansharedfootnote{027037f3ab485ebb}{ที 2. 127 ปิฏฺเฐปิ.} เอกภิกฺขุสฺสปิ กงฺขา วา วิมติ วา พุทฺเธ วา ธมฺเม วา สํเฆ วา มคฺเค วา ปฏิปทาย วา, ปุจฺฉถ ภิกฺขเว, มา ปจฺฉา วิปฺปฏิสาริโน อหุวตฺถ ‘สมฺมุขีภูโต โน สตฺถา อโหสิ, นาสกฺขิมฺห ภควนฺตํ สมฺมุขา ปฏิปุจฺฉิตุนฺ’ติ”. เอวํ วุตฺเต เต ภิกฺขู ตุณฺหี อเหสุํ. ทุติยมฺปิ โข ภควา ภิกฺขู อามนฺเตสิ “สิยา โข ปน ภิกฺขเว เอกภิกฺขุสฺสปิ กงฺขา วา วิมติ วา พุทฺเธ วา ธมฺเม วา สํเฆ วา มคฺเค วา ปฏิปทาย วา, ปุจฺฉถ ภิกฺขเว, มา ปจฺฉา วิปฺปฏิสาริโน อหุวตฺถ ‘สมฺมุขีภูโต โน สตฺถา อโหสิ, นาสกฺขิมฺห ภควนฺตํ สมฺมุขา ปฏิปุจฺฉิตุนฺ’ติ”. ทุติยมฺปิ โข เต ภิกฺขู ตุณฺหี อเหสุํ. ตติยมฺปิ โข ภควา ภิกฺขู อามนฺเตสิ “สิยา โข ปน ภิกฺขเว เอกภิกฺขุสฺสปิ กงฺขา วา วิมติ วา พุทฺเธ วา ธมฺเม วา สํเฆ วา มคฺเค วา ปฏิปทาย วา, ปุจฺฉถ ภิกฺขเว, มา ปจฺฉา วิปฺปฏิสาริโน อหุวตฺถ ‘สมฺมุขีภูโต โน สตฺถา อโหสิ, นาสกฺขิมฺห ภควนฺตํ สมฺมุขา ปฏิปุจฺฉิตุนฺ’ติ”. ตติยมฺปิ โข เต ภิกฺขู ตุณฺหี อเหสุํ.}

\prose{อถ โข ภควา ภิกฺขู อามนฺเตสิ “สิยา โข ปน ภิกฺขเว สตฺถุคารเวนปิ น ปุจฺเฉยฺยาถ, สหายโกปิ ภิกฺขเว สหายกสฺส อาโรเจตู”ติ. เอวํ วุตฺเต เต ภิกฺขู ตุณฺหี อเหสุํ. อถ โข อายสฺมา อานนฺโท ภควนฺตํ เอตทโวจ “อจฺฉริยํ ภนฺเต, อพฺภุตํ ภนฺเต, เอวํ ปสนฺโน อหํ ภนฺเต ‘นตฺถิ อิมสฺมิํ ภิกฺขุสํเฆ เอกภิกฺขุสฺสปิ กงฺขา วา วิมติ วา พุทฺเธ วา ธมฺเม วา สํเฆ วา มคฺเค วา ปฏิปทาย วา’ติ”.}

\prose{\csromanfolio{392}ปสาทา โข ตฺวํ อานนฺท วเทสิ, ญาณเมว เหตฺถ อานนฺท ตถาคตสฺส. นตฺถิ อิมสฺมิํ ภิกฺขุสํเฆ เอกภิกฺขุสฺสปิ กงฺขา วา วิมติ วา พุทฺเธ วา ธมฺเม วา สํเฆ วา มคฺเค วา ปฏิปทาย วา. อิเมสํ หิ อานนฺท ปญฺจนฺนํ ภิกฺขุสตานํ โย ปจฺฉิมโก ภิกฺขุ, โส โสตาปนฺโน อวินิปาต\-ธมฺโม นิยโต สมฺโพธิ\-ปรายโณติ.  ฉฏฺฐํ.}

\csromanmatikaanchor{mtk.129}
\csromantocmarkref{subsection}{7. อจินฺเตยฺยสุตฺต}{mtk.129}
\bookmark[dest={mtk.129},level=2]{7. อจินฺเตยฺยสุตฺต}
\csromanheader{2}{7. \textbf{อจินฺเตยฺยสุตฺต}}
\proseitem{77}{จตฺตาริมานิ ภิกฺขเว อจินฺเตยฺยานิ น จินฺเตตพฺพานิ, ยานิ จินฺเตนฺโต อุมฺมาทสฺส วิฆาตสฺส ภาคี อสฺส. กตมานิ จตฺตาริ? พุทฺธานํ ภิกฺขเว พุทฺธวิสโย อจินฺเตยฺโย น จินฺเตตพฺโพ, ยํ จินฺเตนฺโต อุมฺมาทสฺส วิฆาตสฺส ภาคี อสฺส. ฌายิสฺส ภิกฺขเว ฌานวิสโย อจินฺเตยฺโย น จินฺเตตพฺโพ, ยํ จินฺเตนฺโต อุมฺมาทสฺส วิฆาตสฺส ภาคี อสฺส. กมฺมวิปาโก ภิกฺขเว อจินฺเตยฺโย น จินฺเตตพฺโพ, ยํ จินฺเตนฺโต อุมฺมาทสฺส วิฆาตสฺส ภาคี อสฺส. โลกจินฺตา ภิกฺขเว อจินฺเตยฺยา น จินฺเตตพฺพา, ยํ จินฺเตนฺโต อุมฺมาทสฺส วิฆาตสฺส ภาคี อสฺส. อิมานิ โข ภิกฺขเว จตฺตาริ อจินฺเตยฺยานิ น จินฺเตตพฺพานิ, ยานิ จินฺเตนฺโต อุมฺมาทสฺส วิฆาตสฺส ภาคี อสฺสาติ.  สตฺตมํ.}

\csromanmatikaanchor{mtk.130}
\csromantocmarkref{subsection}{8. ทกฺขิณสุตฺต}{mtk.130}
\bookmark[dest={mtk.130},level=2]{8. ทกฺขิณสุตฺต}
\csromanheader{2}{8. \textbf{ทกฺขิณสุตฺต}}
\proseitem{78}{จตสฺโส อิมา ภิกฺขเว ทกฺขิณาวิสุทฺธิโย. กตมา จตสฺโส? อตฺถิ ภิกฺขเว ทกฺขิณา ทายกโต วิสุชฺฌติ โน ปฏิคฺคาหกโต, อตฺถิ ภิกฺขเว ทกฺขิณา ปฏิคฺคาหกโต วิสุชฺฌติ โน ทายกโต, อตฺถิ ภิกฺขเว ทกฺขิณา เนว ทายกโต วิสุชฺฌติ โน ปฏิคฺคาหกโต, อตฺถิ ภิกฺขเว ทกฺขิณา ทายกโต เจว วิสุชฺฌติ ปฏิคฺคาหกโต จ.}

\prose{กถญฺจ ภิกฺขเว ทกฺขิณา ทายกโต วิสุชฺฌติ โน ปฏิคฺคาหกโต? อิธ ภิกฺขเว ทายโก โหติ สีลวา กลฺยาณธมฺโม, ปฏิคฺคาหกา โหนฺติ ทุสฺสีลา ปาปธมฺมา\csromansharedfootnote{148d347f75e1b68e}{ปฏิคฺคาหโก โหติ ทุสฺสีโล ปาปธมฺโม (สฺยา, กํ, ก) ม 3. 299 ปิฏฺเฐ โอโลเกตพฺพํ.}. เอวํ โข ภิกฺขเว ทกฺขิณา ทายกโต วิสุชฺฌติ โน ปฏิคฺคาหกโต.}

\chapterheadpageread{(8) 3. อปณฺณกวคฺค}
\prose{\csromanfolio{393}กถญฺจ ภิกฺขเว ทกฺขิณา ปฏิคฺคาหกโต วิสุชฺฌติ โน ทายกโต? อิธ ภิกฺขเว ทายโก โหติ ทุสฺสีโล ปาปธมฺโม, ปฏิคฺคาหกา โหนฺติ สีลวนฺโต กลฺยาณธมฺมา. เอวํ โข ภิกฺขเว ทกฺขิณา ปฏิคฺคาหกโต วิสุชฺฌติ โน ทายกโต.}

\prose{กถญฺจ ภิกฺขเว ทกฺขิณา เนว ทายกโต วิสุชฺฌติ โน ปฏิคฺคาหกโต? อิธ ภิกฺขเว ทายโก โหติ ทุสฺสีโล ปาปธมฺโม, ปฏิคฺคาหกาปิ โหนฺติ ทุสฺสฺสีลา ปาปธมฺมา. เอวํ โข ภิกฺขเว ทกฺขิณา เนว ทายกโต วิสุชฺฌติ โน ปฏิคฺคาหกโต.}

\prose{กถญฺจ ภิกฺขเว ทกฺขิณา ทายกโต เจว วิสุชฺฌติ ปฏิคฺคาหกโต จ? อิธ ภิกฺขเว ทายโก โหติ สีลวา กลฺยาณธมฺโม, ปฏิคฺคาหกาปิ โหนฺติ สีลวนฺโต กลฺยาณธมฺมา. เอวํ โข ภิกฺขเว ทกฺขิณา ทายกโต เจว วิสุชฺฌติ ปฏิคฺคาหกโต จ. อิมา โข ภิกฺขเว จตสฺโส ทกฺขิณา วิสุทฺธิโยติ.  อฏฺฐมํ.}

\csromanmatikaanchor{mtk.131}
\csromantocmarkref{subsection}{9. วณิชฺชสุตฺต}{mtk.131}
\bookmark[dest={mtk.131},level=2]{9. วณิชฺชสุตฺต}
\csromanheader{2}{9. \textbf{วณิชฺชสุตฺต}}
\proseitem{79}{อถ โข อายสฺมา สาริปุตฺโต เยน ภควา เตนุปสงฺกมิ, อุปสงฺกมิตฺวา ภควนฺตํ อภิวาเทตฺวา เอกมนฺตํ นิสีทิ, เอกมนฺตํ นิสินฺโน โข อายสฺมา สาริปุตฺโต ภควนฺตํ เอตทโวจ “โก นุ โข ภนฺเต เหตุ โก ปจฺจโย, เยน มิเธกจฺจสฺส ตาทิสาว วณิชฺชา ปยุตฺตา เฉทคามินี โหติ. โก ปน ภนฺเต เหตุ โก ปจฺจโย, เยน มิเธกจฺจสฺส ตาทิสาว วณิชฺชา ปยุตฺตา น ยถาธิปฺปายา\csromansharedfootnote{6daef59121fb33be}{ยถาธิปฺปายํ (สี)} โหติ. โก นุ โข ภนฺเต เหตุ โก ปจฺจโย, เยน มิเธกจฺจสฺส ตาทิสาว วณิชฺชา ปยุตฺตา ยถาธิปฺปายา\csromansharedfootnote{6daef59121fb33be}{ยถาธิปฺปายํ (สี)} โหติ. โก ปน ภนฺเต เหตุ โก ปจฺจโย, เยน มิเธกจฺจสฺส ตาทิสาว วณิชฺชา ปยุตฺตา ปราธิปฺปายา โหตี”ติ.}

\prose{อิธ สาริปุตฺต เอกจฺโจ สมณํ วา พฺราหฺมณํ วา อุปสงฺกมิตฺวา ปวาเรติ “วทตุ ภนฺเต ปจฺจเยนา”ติ. โส เยน ปวาเรติ, ตํ น เทติ, โส เจ ตโต จุโต อิตฺถตฺตํ อาคจฺฉติ, โส ยญฺญเทว วณิชฺชํ ปโยเชติ, สาสฺส โหติ เฉทคามินี.}

\prose{\csromanfolio{394}อิธ ปน สาริปุตฺต เอกจฺโจ สมณํ วา พฺราหฺมณํ วา อุปสงฺกมิตฺวา ปวาเรติ “วทตุ ภนฺเต ปจฺจเยนา”ติ. โส เยน ปวาเรติ, ตํ น ยถาธิปฺปายํ เทติ, โส เจ ตโต จุโต อิตฺถตฺตํ อาคจฺฉติ, โส ยญฺญเทว วณิชฺชํ ปโยเชติ, สาสฺส น โหติ ยถาธิปฺปายา\csromansharedfootnote{013f9a1a8554d489}{ยถาธิปฺปายํ (สี, ก)}.}

\prose{อิธ ปน สาริปุตฺต เอกจฺโจ สมณํ วา พฺราหฺมณํ วา อุปสงฺกมิตฺวา ปวาเรติ “วทตุ ภนฺเต ปจฺจเยนา”ติ, โส เยน ปวาเรติ, ตํ ยถาธิปฺปายํ เทติ, โส เจ ตโต จุโต อิตฺถตฺตํ อาคจฺฉติ, โส ยญฺญเทว วณิชฺชํ ปโยเชติ, สาสฺส โหติ ยถาธิปฺปายา\csromansharedfootnote{013f9a1a8554d489}{ยถาธิปฺปายํ (สี, ก)}.}

\prose{อิธ สาริปุตฺต เอกจฺโจ สมณํ วา พฺราหฺมณํ วา อุปสงฺกมิตฺวา ปวาเรติ “วทตุ ภนฺเต ปจฺจเยนา”ติ. โส เยน ปวาเรติ, ตํ ปราธิปฺปายํ เทติ, โส เจ ตโต จุโต อิตฺถตฺตํ อาคจฺฉติ, โส ยญฺญเทว วณิชฺชํ ปโยเชติ, สาสฺส โหติ ปราธิปฺปายา\csromansharedfootnote{183bfc9b0a6281cd}{ปราธิปฺปายํ (ก)}.}

\prose{อยํ โข สาริปุตฺต เหตุ อยํ ปจฺจโย, เยน มิเธกจฺจสฺส ตาทิสาว วณิชฺชา ปยุตฺตา เฉทคามินี โหติ. อยํ ปน สาริปุตฺต เหตุ อยํ ปจฺจโย, เยน มิเธกจฺจสฺส ตาทิสาว วณิชฺชา ปยุตฺตา น ยถาธิปฺปายา โหติ. อยํ โข ปน สาริปุตฺต เหตุ อยํ ปจฺจโย, เยน มิเธกจฺจสฺส ตาทิสาว วณิชฺชา ปยุตฺตา ยถาธิปฺปายา โหติ. อยํ ปน สาริปุตฺต เหตุ อยํ ปจฺจโย, เยน มิเธกจฺจสฺส ตาทิสาว วณิชฺชา ปยุตฺตา ปราธิปฺปายา โหตีติ.  นวมํ.}

\csromanmatikaanchor{mtk.132}
\csromantocmarkref{subsection}{10. กมฺโพชสุตฺต}{mtk.132}
\bookmark[dest={mtk.132},level=2]{10. กมฺโพชสุตฺต}
\csromanheader{2}{10. \textbf{กมฺโพชสุตฺต}}
\proseitem{80}{เอกํ สมยํ ภควา โกสมฺพิยํ วิหรติ โฆสิตาราเม. อถ โข อายสฺมา อานนฺโท เยน ภควา เตนุปสงฺกมิ, อุปสงฺกมิตฺวา ภควนฺตํ อภิวาเทตฺวา เอกมนฺตํ นิสีทิ, เอกมนฺตํ นิสินฺโน โข อายสฺมา อานนฺโท ภควนฺตํ เอตทโวจ\csromandash{}}

\prose{“โก นุ โข ภนฺเต เหตุ โก ปจฺจโย, เยน มาตุคาโม เนว สภายํ นิสีทติ, น กมฺมนฺตํ ปโยเชติ, น กมฺโพชํ คจฺฉตี”ติ. โกธโน อานนฺท มาตุคาโม, อิสฺสุกี อานนฺท มาตุคาโม, มจฺฉรี อานนฺท มาตุคาโม, ทุปฺปญฺโญ อานนฺท มาตุคาโม. อยํ โข อานนฺท}

\vspace{\tipitakaparskip}
\csromancenter{(9) 4. มจลวคฺค เหตุ อยํ ปจฺจโย, เยน มาตุคาโม เนว สภายํ นิสีทติ, น กมฺมนฺตํ ปโยเชติ, น กมฺโพชํ คจฺฉตีติ.  ทสมํ.}
\nitthitammidruled{อปณฺณกวคฺโค ตติโย.}
\csromancenter{\textbf{ตสฺสุทฺทานํ}}
\csromanmatikaanchor{mtk.133}
\csromanmatikaanchor{mtk.134}
\csromantocmarknopage{section}{(9) 4. มจลวคฺค}{mtk.133}
\bookmark[dest={mtk.133},level=1]{(9) 4. มจลวคฺค}
\csromantocmarkref{subsection}{1-2. ปาณาติปาตสุตฺตาทิ}{mtk.134}
\bookmark[dest={mtk.134},level=2]{1-2. ปาณาติปาตสุตฺตาทิ}
\setlength{\csromangathapagewidth}{0pt}
\setlength{\csromangathawakwidth}{0pt}
\csromangathameasuregroup{ปธานํ ทิฏฺฐิสปฺปุริส, \\ กุสินารอจินฺเตยฺยา,}{\csromangathabat{ปธานํ ทิฏฺฐิสปฺปุริส,}{วธุกา เทฺว จ โหนฺติ อคฺคานิ.} \\ \csromangathabat{กุสินารอจินฺเตยฺยา,}{ทกฺขิณา จ วณิชฺชา กมฺโพชนฺติ.}}
\csromangathasetleft{ปธานํ ทิฏฺฐิสปฺปุริส, \\ กุสินารอจินฺเตยฺยา,}
\csromangathagroup{\csromangathastanza{\csromanfolio{395}\csromangathabat{ปธานํ ทิฏฺฐิสปฺปุริส,}{วธุกา เทฺว จ โหนฺติ อคฺคานิ.} \\ \csromangathabat{กุสินารอจินฺเตยฺยา,}{ทกฺขิณา จ วณิชฺชา กมฺโพชนฺติ.}}}

\csromanmarkh{(9)}\csromanmarkhii{4. มจลวคฺค}\csromanheadernomark{1}{(9) 4. มจลวคฺค}
\csromanheader{2}{1. ปาณาติปาต\-สุตฺต}
\proseitem{81}{จตูหิ ภิกฺขเว ธมฺเมหิ สมนฺนาคโต ยถาภตํ นิกฺขิตฺโต เอวํ นิรเย. กตเมหิ จตูหิ? ปาณาติปาตี โหติ, อทินฺนาทายี โหติ, กาเมสุ\-มิจฺฉาจารี โหติ, มุสาวาที โหติ. อิเมหิ โข ภิกฺขเว จตูหิ ธมฺเมหิ สมนฺนาคโต ยถาภตํ นิกฺขิตฺโต เอวํ นิรเย.}

\prose{จตูหิ ภิกฺขเว ธมฺเมหิ สมนฺนาคโต ยถาภตํ นิกฺขิตฺโต เอวํ สคฺเค. กตเมหิ จตูหิ? ปาณาติปาตา ปฏิวิรโต โหติ, อทินฺนาทานา ปฏิวิรโต โหติ, กาเมสุ\-มิจฺฉาจารา ปฏิวิรโต โหติ, มุสาวาทา ปฏิวิรโต โหติ. อิเมหิ โข ภิกฺขเว จตูหิ ธมฺเมหิ สมนฺนาคโต ยถาภตํ นิกฺขิตฺโต เอวํ สคฺเคติ.  ปฐมํ.}

\csromanheader{2}{2. \textbf{มุสาวาทสุตฺต}}
\proseitem{82}{จตูหิ ภิกฺขเว ธมฺเมหิ สมนฺนาคโต ยถาภตํ นิกฺขิตฺโต เอวํ นิรเย. กตเมหิ จตูหิ? มุสาวาที โหติ, ปิสุณวาโจ โหติ, ผรุสวาโจ โหติ, สมฺผปฺปลาปี โหติ. อิเมหิ โข ภิกฺขเว จตูหิ ธมฺเมหิ สมนฺนาคโต ยถาภตํ นิกฺขิตฺโต เอวํ นิรเย.}

\csromanmatikaanchor{mtk.135}
\csromantocmarkref{subsection}{3-4. อวณฺณารหสุตฺตาทิ}{mtk.135}
\bookmark[dest={mtk.135},level=2]{3-4. อวณฺณารหสุตฺตาทิ}
\prose{จตูหิ ภิกฺขเว ธมฺเมหิ สมนฺนาคโต ยถาภตํ นิกฺขิตฺโต เอวํ สคฺเค. กตเมหิ จตูหิ? มุสาวาทา ปฏิวิรโต โหติ, ปิสุณาย วาจาย ปฏิวิรโต โหติ, ผรุสาย วาจาย ปฏิวิรโต โหติ, \csromanfolio{396}สมฺผปฺปลาปา ปฏิวิรโต โหติ. อิเมหิ โข ภิกฺขเว จตูหิ ธมฺเมหิ สมนฺนาคโต ยถาภตํ นิกฺขิตฺโต เอวํ สคฺเคติ.  ทุติยํ.}

\csromanheader{2}{3. อวณฺณารห\-สุตฺต}
\proseitem{83}{จตูหิ ภิกฺขเว ธมฺเมหิ สมนฺนาคโต ยถาภตํ นิกฺขิตฺโต เอวํ นิรเย. กตเมหิ จตูหิ? อนนุวิจฺจ อปริโยคาเหตฺวา อวณฺณารหสฺส วณฺณํ ภาสติ, อนนุวิจฺจ อปริโยคาเหตฺวา วณฺณารหสฺส อวณฺณํ ภาสติ, อนนุวิจฺจ อปริโยคาเหตฺวา อปฺปสาทนีเย ฐาเน ปสาทํ อุปทํเสติ, อนนุวิจฺจ อปริโยคาเหตฺวา ปสาทนีเย ฐาเน อปฺปสาทํ อุปทํเสติ. อิเมหิ โข ภิกฺขเว จตูหิ ธมฺเมหิ สมนฺนาคโต ยถาภตํ นิกฺขิตฺโต เอวํ นิรเย.}

\prose{จตูหิ ภิกฺขเว ธมฺเมหิ สมนฺนาคโต ยถาภตํ นิกฺขิตฺโต เอวํ สคฺเค. กตเมหิ จตูหิ? อนุวิจฺจ ปริโยคาเหตฺวา อวณฺณารหสฺส อวณฺณํ ภาสติ, อนุวิจฺจ ปริโยคาเหตฺวา วณฺณารหสฺส วณฺณํ ภาสติ, อนุวิจฺจ ปริโยคาเหตฺวา อปฺปสาทนีเย ฐาเน อปฺปสาทํ อุปทํเสติ, อนุวิจฺจ ปริโยคาเหตฺวา ปสาทนีเย ฐาเน ปสาทํ อุปทํเสติ. อิเมหิ โข ภิกฺขเว จตูหิ ธมฺเมหิ สมนฺนาคโต ยถาภตํ นิกฺขิตฺโต เอวํ สคฺเคติ.  ตติยํ.}

\csromanheader{2}{4. \textbf{โกธครุสุตฺต}}
\proseitem{84}{จตูหิ ภิกฺขเว ธมฺเมหิ สมนฺนาคโต ยถาภตํ นิกฺขิตฺโต เอวํ นิรเย. กตเมหิ จตูหิ? โกธครุ โหติ น สทฺธมฺมครุ, มกฺขครุ โหติ น สทฺธมฺมครุ, ลาภครุ โหติ น สทฺธมฺมครุ, สกฺการครุ โหติ น สทฺธมฺมครุ. อิเมหิ โข ภิกฺขเว จตูหิ ธมฺเมหิ สมนฺนาคโต ยถาภตํ นิกฺขิตฺโต เอวํ นิรเย.}

\prose{จตูหิ ภิกฺขเว ธมฺเมหิ สมนฺนาคโต ยถาภตํ นิกฺขิตฺโต เอวํ สคฺเค. กตเมหิ จตูหิ? สทฺธมฺมครุ โหติ น โกธครุ, สทฺธมฺมครุ โหติ น มกฺขครุ, สทฺธมฺมครุ โหติ น ลาภครุ, สทฺธมฺมครุ โหติ น สกฺการครุ. อิเมหิ โข ภิกฺขเว จตูหิ ธมฺเมหิ สมนฺนาคโต ยถาภตํ นิกฺขิตฺโต เอวํ สคฺเคติ.  จตุตฺถํ.}

\csromanmatikaanchor{mtk.136}
\csromantocmarkref{subsection}{5. ตโมตมสุตฺต}{mtk.136}
\bookmark[dest={mtk.136},level=2]{5. ตโมตมสุตฺต}
\csromanmarkh{(9)}\csromanmarkhii{4. มจลวคฺค 5. ตโมตมสุตฺต}\csromanheadernomark{2}{(9) 4. มจลวคฺค \textbf{5. ตโมตมสุตฺต}}
\proseitem{85}{\csromanfolio{397}จตฺตาโรเม ภิกฺขเว ปุคฺคลา สนฺโต สํวิชฺชมานา โลกสฺมิํ. กตเม จตฺตาโร? ตโม ตมปรายโณ\csromansharedfootnote{7f754b853c01771b}{...ปรายโน (สฺยา, กํ, อิ) อภิ 3. 159; สํ 1. 94 ปิฏฺเฐสุปิ.}, ตโม โชติปรายโณ, โชติ ตมปรายโณ, โชติ โชติปรายโณ.}

\prose{กถญฺจ ภิกฺขเว ปุคฺคโล ตโม โหติ ตมปรายโณ? อิธ ภิกฺขเว เอกจฺโจ ปุคฺคโล นีเจ กุเล ปจฺจาชาโต โหติ จณฺฑาลกุเล วา เวนกุเล วา เนสาทกุเล วา รถการกุเล วา ปุกฺกุสกุเล วา ทลิทฺเท อปฺป\-นฺน\-ปาน\-โภชเน กสิรวุตฺติเก, ยตฺถ กสิเรน ฆาสจฺฉาโท ลพฺภติ, โส จ โหติ ทุพฺพณฺโณ ทุทฺทสิโก โอโกฏิมโก พวฺหาพาโธ กาโณ วา กุณี วา ขญฺโช วา ปกฺขหโต วา, น ลาภี อนฺนสฺส ปานสฺส วตฺถสฺส ยานสฺส มาลา\-คนฺธ\-วิเลปนสฺส เสยฺยา\-วสถ\-ปทีเปยฺยสฺส. โส กาเยน ทุจฺจริตํ จรติ, วาจาย ทุจฺจริตํ จรติ, มนสา ทุจฺจริตํ จรติ. โส กาเยน ทุจฺจริตํ จริตฺวา วาจาย ทุจฺจริตํ จริตฺวา มนสา ทุจฺจริตํ จริตฺวา กายสฺส เภทา ปรํ มรณา อปายํ ทุคฺคติํ วินิปาตํ นิรยํ อุปปชฺชติ. เอวํ โข ภิกฺขเว ปุคฺคโล ตโม โหติ ตมปรายโณ.}

\prose{กถญฺจ ภิกฺขเว ปุคฺคโล ตโม โหติ โชติปรายโณ, อิธ ภิกฺขเว เอกจฺโจ ปุคฺคโล นีเจ กุเล ปจฺจาชาโต โหติ จณฺฑาลกุเล วา เวนกุเล วา เนสาทกุเล วา รถการกุเล วา ปุกฺกุสกุเล วา ทลิทฺเท อปฺป\-นฺน\-ปาน\-โภชเน กสิรวุตฺติเก, ยตฺถ กสิเรน ฆาสจฺฉาโท ลพฺภติ, โส จ โหติ ทุพฺพณฺโณ ทุทฺทสิโก โอโกฏิมโก พวฺหาพาโธ กาโณ วา กุณี วา ขญฺโช วา ปกฺขหโต วา, น ลาภี อนฺนสฺส ปานสฺส วตฺถสฺส ยานาสฺส มาลา\-คนฺธ\-วิเลปนสฺส เสยฺยา\-วสถ\-ปทีเปยฺยสฺส. โส กาเยน สุจริตํ จรติ, วาจาย สุจริตํ จรติ, มนสา สุจริตํ จรติ. โส กาเยน สุจริตํ จริตฺวา วาจาย สุจริตํ จริตฺวา มนสา สุจริตํ จริตฺวา กายสฺส เภทา ปรํ มรณา สุคติํ สคฺคํ โลกํ อุปปชฺชติ. เอวํ โข ภิกฺขเว ปุคฺคโล ตโม โหติ โชติปรายโณ.}

\prose{\csromanfolio{398}กถญฺจ ภิกฺขเว ปุคฺคโล โชติ โหติ ตมปรายโณ? อิธ ภิกฺขเว เอกจฺโจ ปุคฺคโล อุจฺเจ กุเล ปจฺจาชาโต โหติ ขตฺติย\-มหาสาล\-กุเล วา พฺราหฺมณ\-มหาสาล\-กุเล วา คหปติ\-มหาสาล\-กุเล วา อฑฺเฒ มหทฺธเน มหาโภเค ปหูต\-ชาตรูป\-รชเต ปหูต\-วิตฺตู\-ปกรเณ ปหูต\-ธน\-ธญฺเญ, โส จ โหติ อภิรูโป ทสฺสนีโย ปาสาทิโก ปรมาย วณฺณ\-โปกฺขร\-ตาย สมนฺนาคโต, ลาภี อนฺนสฺส ปานสฺส วตฺถสฺส ยานสฺส มาลา\-คนฺธ\-วิเลปนสฺส เสยฺยา\-วสถ\-ปทีเปยฺยสฺส. โส กาเยน ทุจฺจริตํ จรติ, วาจาย ทุจฺจริตํ จรติ, มนสา ทุจฺจริตํ จรติ. โส กาเยน ทุจฺจริตํ จริตฺวา วาจาย ทุจฺจริตํ จริตฺวา มนสา ทุจฺจริตํ จริตฺวา กายสฺส เภทา ปรํ มรณา อปายํ ทุคฺคติํ วินิปาตํ นิรยํ อุปปชฺชติ. เอวํ โข ภิกฺขเว ปุคฺคโล โชติ โหติ ตมปรายโณ.}

\prose{กถญฺจ ภิกฺขเว ปุคฺคโล โชติ โหติ โชติปรายโณ? อิธ ภิกฺขเว เอกจฺโจ ปุคฺคโล อุจฺเจ กุเล ปจฺจาชาโต โหติ ขตฺติย\-มหาสาล\-กุเล วา พฺราหฺมณ\-มหาสาล\-กุเล วา คหปติ\-มหาสาล\-กุเล วา อฑฺเฒ มหทฺธเน มหาโภเค ปหูต\-ชาตรูป\-รชเต ปหูต\-วิตฺตู\-ปกรเณ ปหูต\-ธน\-ธญฺเญ. โส จ โหติ อภิรูโป ทสฺสนีโย ปาสาทิโก ปรมาย วณฺณ\-โปกฺขร\-ตาย สมนฺนาคโต, ลาภี อนฺนสฺส ปานสฺส วตฺถสฺส ยานสฺส มาลา\-คนฺธ\-วิเลปนสฺส เสยฺยา\-วสถ\-ปทีเปยฺยสฺส. โส กาเยน สุจริตํ จรติ, วาจาย สุจริตํ จรติ, มนสา สุจริตํ จรติ. โส กาเยน สุจริตํ จริตฺวา วาจาย สุจริตํ จริตฺวา มนสา สุจริตํ จริตฺวา กายสฺส เภทา ปรํ มรณา สุคติํ สคฺคํ โลกํ อุปปชฺชติ. เอวํ โข ภิกฺขเว ปุคฺคโล โชติ โหติ โชติปรายโณ. อิเม โข ภิกฺขเว จตฺตาโร ปุคฺคลา สนฺโต สํวิชฺชมานา โลกสฺมินฺติ.  ปญฺจมํ.}

\csromanmatikaanchor{mtk.137}
\csromantocmarkref{subsection}{6. โอณโตณตสุตฺต}{mtk.137}
\bookmark[dest={mtk.137},level=2]{6. โอณโตณตสุตฺต}
\csromanheader{2}{6. \textbf{โอณโตณตสุตฺต}}
\proseitem{86}{จตฺตาโรเม ภิกฺขเว ปุคฺคลา สนฺโต สํวิชฺชมานา โลกสฺมิํ, กตเม จตฺตาโร? โอณโตณโต โอณตุณฺณโต อุณฺณโตณโต อุณฺณตุณฺณโต. อิเม โข ภิกฺขเว จตฺตาโร ปุคฺคลา สนฺโต สํวิชฺชมานา โลกสฺมินฺติ\csromansharedfootnote{84b79398747e3c29}{อภิ 3. 160 ปิฏฺเฐปิ.}.  ฉฏฺฐํ.}

\csromanmatikaanchor{mtk.138}
\csromantocmarkref{subsection}{7. ปุตฺตสุตฺต}{mtk.138}
\bookmark[dest={mtk.138},level=2]{7. ปุตฺตสุตฺต}
\csromanmarkh{(9)}\csromanmarkhii{4. มจลวคฺค 7. ปุตฺตสุตฺต}\csromanheadernomark{2}{(9) 4. มจลวคฺค \textbf{7. ปุตฺตสุตฺต}}
\proseitem{87}{\csromanfolio{399}จตฺตาโรเม ภิกฺขเว ปุคฺคลา สนฺโต สํวิชฺชมานา โลกสฺมิํ. กตเม จตฺตาโร? สมณมจโล สมณ\-ปุณฺฑรีโก สมณปทุโม สมเณสุ สมณ\-สุขุมาโล.}

\prose{กถญฺจ ภิกฺขเว ปุคฺคโล สมณมจโล โหติ? อิธ ภิกฺขเว ภิกฺขุ เสโข โหติ ปาฏิปโท\csromansharedfootnote{66c5ff0cf4fe9e09}{ปฏิปโท (สฺยา, กํ, อิ) ม 2 เสขสุตฺตวณฺณนา โอโลเกตพฺพา.}, อนุตฺตรํ โยคกฺเขมํ ปตฺถยมาโน วิหรติ. เสยฺยถาปิ ภิกฺขเว รญฺโญ ขตฺติยสฺส มุทฺธา\-วสิตฺตสฺส เชฏฺโฐ ปุตฺโต อาภิเสโก อนภิสิตฺโต มจลปฺปตฺโต. เอวเมวํ โข ภิกฺขเว ภิกฺขุ เสโข โหติ ปาฏิปโท, อนุตฺตรํ โยคกฺเขมํ ปตฺถยมาโน วิหรติ. เอวํ โข ภิกฺขเว ปุคฺคโล สมณมจโล โหติ.}

\prose{กถญฺจ ภิกฺขเว ปุคฺคโล สมณ\-ปุณฺฑรีโก โหติ? อิธ ภิกฺขเว ภิกฺขุ อาสวานํ ขยา อนาสวํ เจโตวิมุตฺติํ ปญฺญาวิมุตฺติํ ทิฏฺเฐว ธมฺเม สยํ อภิญฺญา สจฺฉิกตฺวา อุปสมฺปชฺช วิหรติ, โน จ โข อฏฺฐ วิโมกฺเข กาเยน ผุสิตฺวา\csromansharedfootnote{e87c423abe96dce4}{ผสฺสิตฺวา (สี, อิ)} วิหรติ. เอวํ โข ภิกฺขเว ปุคฺคโล สมณ\-ปุณฺฑรีโก โหติ.}

\prose{กถญฺจ ภิกฺขเว ปุคฺคโล สมณปทุโม โหติ? อิธ ภิกฺขเว ภิกฺขุ อาสวานํ ขยา อนาสวํ เจโตวิมุตฺติํ ปญฺญาวิมุตฺติํ ทิฏฺเฐว ธมฺเม สยํ อภิญฺญา สจฺฉิกตฺวา อุปสมฺปชฺช วิหรติ, อฏฺฐ จ วิโมกฺเข กาเยน ผุสิตฺวา วิหรติ. เอวํ โข ภิกฺขเว ปุคฺคโล สมณปทุโม โหติ.}

\prose{กถญฺจ ภิกฺขเว ปุคฺคโล สมเณสุ สมณ\-สุขุมาโล โหติ? อิธ ภิกฺขเว ภิกฺขุ ยาจิโตว พหุลํ จีวรํ ปริภุญฺชติ อปฺปํ อยาจิโต, ยาจิโตว พหุลํ ปิณฺฑปาตํ ปริภุญฺชติ อปฺปํ อยาจิโต, ยาจิโตว พหุลํ เสนาสนํ ปริภุญฺชติ อปฺปํ อยาจิโต, ยาจิโตว พหุลํ คิลานปฺปจฺจย\-เภสชฺช\-ปริกฺขารํ ปริภุญฺชติ อปฺปํ อยาจิโต. เยหิ โข ปน สพฺรหฺมจารีหิ สทฺธิํ วิหรติ, ตฺยสฺส\csromansharedfootnote{9374ce86d6c41bbe}{ตฺยาสฺส (อิ) อํ 2. 115 ปิฏฺเฐปิ.} มนาเปเนว พหุลํ กายกมฺเมน สมุทาจรนฺติ อปฺปํ อมนาเปน, มนาเปเนว พหุลํ วจีกมฺเมน สมุทาจรนฺติ \csromanfolio{400}อปฺปํ อมนาเปน, มนาเปเนว พหุลํ มโนกมฺเมน สมุทาจรนฺติ อปฺปํ อมนาเปน, มนาปํเยว พหุลํ อุปหารํ อุปหรนฺติ อปฺปํ อมนาปํ. ยานิ โข ปน ตานิ เวทยิตานิ ปิตฺต\-สมุฏฺฐานานิ วา เสมฺห\-สมุฏฺฐานานิ วา วาต\-สมุฏฺฐานานิ วา สนฺนิปาติกานิ วา อุตุปริณาม\-ชานิ วา วิสม\-ปริหาร\-ชานิ วา โอปกฺกมิกานิ วา กมฺม\-วิปาก\-ชานิ วา, ตานิ ปนสฺส น พหุเทว อุปฺปชฺชนฺติ, อปฺปาพาโธ โหติ. จตุนฺนํ ฌานานํ อาภิเจตสิกานํ ทิฏฺฐ\-ธมฺม\-สุข\-วิหารานํ นิกามลาภี โหติ อกิจฺฉลาภี อกสิรลาภี, อาสวานํ ขยา อนาสวํ เจโตวิมุตฺติํ ปญฺญาวิมุตฺติํ ทิฏฺเฐว ธมฺเม สยํ อภิญฺญา สจฺฉิกตฺวา อุปสมฺปชฺช วิหรติ. เอวํ โข ภิกฺขเว ปุคฺคโล สมเณสุ สมณ\-สุขุมาโล โหติ.}

\prose{ยํ หิ ตํ ภิกฺขเว สมฺมา วทมาโน วเทยฺย “สมเณสุ สมณ\-สุขุมาโล”ติ, มเมว ตํ ภิกฺขเว สมฺมา วทมาโน วเทยฺย “สมเณสุ สมณ\-สุขุมาโล”ติ. อหํ หิ ภิกฺขเว ยาจิโตว พหุลํ จีวรํ ปริภุญฺชามิ อปฺปํ อยาจิโต, ยาจิโตว พหุลํ ปิณฺฑปาตํ ปริภุญฺชามิ อปฺปํ อยาจิโต, ยาจิโตว พหุลํ เสนาสนํ ปริภุญฺชามิ อปฺปํ อยาจิโต, ยาจิโตว พหุลํ คิลานปฺปจฺจย\-เภสชฺช\-ปริกฺขารํ ปริภุญฺชามิ อปฺปํ อยาจิโต. เยหิ โข ปน ภิกฺขูหิ สทฺธิํ วิหรามิ, เต เม มนาเปเนว พหุลํ กายกมฺเมน สมุทาจรนฺติ อปฺปํ อมนาเปน, มนาเปเนว พหุลํ วจีกมฺเมน สมุทาจรนฺติ อปฺปํ อมนาเปน, มนาเปเนว พหุลํ มโนกมฺเมน สมุทาจรนฺติ อปฺปํ อมนาเปน, มนาปํเยว พหุลํ อุปหารํ อุปหรนฺติ อปฺปํ อมนาปํ. ยานิ โข ปน ตานิ เวทยิตานิ ปิตฺต\-สมุฏฺฐานานิ วา เสมฺห\-สมุฏฺฐานานิ วา วาต\-สมุฏฺฐานานิ วา สนฺนิปาติกานิ วา อุตุปริณาม\-ชานิ วา วิสม\-ปริหาร\-ชานิ วา โอปกฺกมิกานิ วา กมฺม\-วิปาก\-ชานิ วา, ตานิ เม น พหุเทว อุปฺปชฺชนฺติ, อปฺปาพาโธหมสฺมิ. จตุนฺนํ โข ปนสฺมิ ฌานานํ อาภิเจตสิกานํ ทิฏฺฐ\-ธมฺม\-สุข\-วิหารานํ นิกามลาภี อกิจฺฉลาภี อกสิรลาภี, อาสวานํ ขยา อนาสวํ เจโตวิมุตฺติํ ปญฺญาวิมุตฺติํ ทิฏฺเฐว ธมฺเม สยํ อภิญฺญา สจฺฉิกตฺวา อุปสมฺปชฺช วิหรามิ.}

\prose{ยํ หิ ตํ ภิกฺขเว สมฺมา วทมาโน วเทยฺย “สมเณสุ สมณ\-สุขุมาโล”ติ. มเมว ตํ ภิกฺขเว สมฺมา วทมาโน วเทยฺย}

\vspace{\tipitakaparskip}
\csromancenter{(9) 4. มจลวคฺค “สมเณสุ สมณ\-สุขุมาโล”ติ. อิเม โข ภิกฺขเว จตฺตาโร ปุคฺคลา สนฺโต สํวิชฺชมานา โลกสฺมินฺติ.  สตฺตมํ.}
\csromanmatikaanchor{mtk.139}
\csromantocmarkref{subsection}{8. สํโยชนสุตฺต}{mtk.139}
\bookmark[dest={mtk.139},level=2]{8. สํโยชนสุตฺต}
\csromanheader{2}{8. \textbf{สํโยชนสุตฺต}}
\proseitem{88}{\csromanfolio{401}จตฺตาโรเม ภิกฺขเว ปุคฺคลา สนฺโต สํวิชฺชมานา โลกสฺมิํ. กตเม จตฺตาโร, สมณมจโล สมณ\-ปุณฺฑรีโก สมณปทุโม สมเณสุ สมณ\-สุขุมาโล.}

\prose{กถญฺจ ภิกฺขเว ปุคฺคโล\csromansharedfootnote{11d44e10644ff8fa}{อภิ 3. 173 ปิฏฺเฐปิ (โถกํ วิสทิสํ)} สมณมจโล โหติ? อิธ ภิกฺขเว ภิกฺขุ ติณฺณํ สํโยชนานํ ปริกฺขยา โสตาปนฺโน โหติ อวินิปาต\-ธมฺโม นิยโต สมฺโพธิ\-ปรายโณ. เอวํ โข ภิกฺขเว ปุคฺคโล สมณมจโล โหติ.}

\prose{กถญฺจ ภิกฺขเว ปุคฺคโล สมณ\-ปุณฺฑรีโก โหติ? อิธ ภิกฺขเว ภิกฺขุ ติณฺณํ สํโยชนานํ ปริกฺขยา ราค\-โทส\-โมหานํ ตนุตฺตา สกทาคามี โหติ, สกิเทว อิมํ โลกํ อาคนฺตฺวา ทุกฺขสฺสนฺตํ กโรติ. เอวํ โข ภิกฺขเว ปุคฺคโล สมณ\-ปุณฺฑรีโก โหติ.}

\prose{กถญฺจ ภิกฺขเว ปุคฺคโล สมณปทุโม โหติ? อิธ ภิกฺขเว ภิกฺขุ ปญฺจนฺนํ โอรมฺภาคิยานํ สํโยชนานํ ปริกฺขยา โอปปาติโก โหติ, ตตฺถ ปรินิพฺพายี อนาวตฺติธมฺโม ตสฺมา โลกา. เอวํ โข ภิกฺขเว ปุคฺคโล สมณปทุโม โหติ.}

\prose{กถญฺจ ภิกฺขเว ปุคฺคโล สมเณสุ สมณ\-สุขุมาโล โหติ? อิธ ภิกฺขเว ภิกฺขุ อาสวานํ ขยา อนาสวํ เจโตวิมุตฺติํ ปญฺญาวิมุตฺติํ ทิฏฺเฐว ธมฺเม สยํ อภิญฺญา สจฺฉิกตฺวา อุปสมฺปชฺช วิหรติ. เอวํ โข ภิกฺขเว ปุคฺคโล สมเณสุ สมณ\-สุขุมาโล โหติ. อิเม โข ภิกฺขเว จตฺตาโร ปุคฺคลา สนฺโต สํวิชฺชมานา โลกสฺมินฺติ.  อฏฺฐมํ.}

\csromanmatikaanchor{mtk.140}
\csromantocmarkref{subsection}{9. สมฺมาทิฏฺฐิสุตฺต}{mtk.140}
\bookmark[dest={mtk.140},level=2]{9. สมฺมาทิฏฺฐิสุตฺต}
\csromanheader{2}{9. สมฺมาทิฏฺฐิ\-สุตฺต}
\proseitem{89}{จตฺตาโรเม ภิกฺขเว ปุคฺคลา สนฺโต สํวิชฺชมานา โลกสฺมิํ. กตเม จตฺตาโร? สมณมจโล สมณ\-ปุณฺฑรีโก สมณปทุโม สมเณสุ สมณ\-สุขุมาโล.}

\prose{กถญฺจ ภิกฺขเว ปุคฺคโล สมณมจโล โหติ? อิธ ภิกฺขเว ภิกฺขุ สมฺมาทิฏฺฐิโก โหติ, สมฺมาสงฺกปฺโป โหติ, สมฺมาวาโจ โหติ,}

\prose{\csromanfolio{402}สมฺมากมฺมนฺโต โหติ, สมฺมาอาชีโว โหติ, สมฺมาวายาโม โหติ, สมฺมาสติ\csromansharedfootnote{badce4138ed96e35}{สมฺมาสตี (สี, ก)} โหติ, สมฺมาสมาธิ\csromansharedfootnote{113b27e5e627936a}{สมฺมาสมาธี (สี, ก)} โหติ. เอวํ โข ภิกฺขเว ปุคฺคโล สมณมจโล โหติ.}

\prose{กถญฺจ ภิกฺขเว ปุคฺคโล สมณ\-ปุณฺฑรีโก โหติ? อิธ ภิกฺขเว ภิกฺขุ สมฺมาทิฏฺฐิโก โหติ, สมฺมาสงฺกปฺโป โหติ, สมฺมาวาโจ โหติ, สมฺมากมฺมนฺโต โหติ, สมฺมาอาชีโว โหติ, สมฺมาวายาโม โหติ, สมฺมาสติ โหติ, สมฺมาสมาธิ โหติ, สมฺมาญาณี โหติ, สมฺมาวิมุตฺติ\csromansharedfootnote{1bbddfd4861c4a98}{สมฺมาวิมุตฺตี (สี, ก)} โหติ, โน จ โข อฏฺฐ วิโมกฺเข กาเยน ผุสิตฺวา วิหรติ. เอวํ โข ภิกฺขเว ปุคฺคโล สมณ\-ปุณฺฑรีโก โหติ.}

\prose{กถญฺจ ภิกฺขเว ปุคฺคโล สมณปทุโม โหติ? อิธ ภิกฺขเว ภิกฺขุ สมฺมาทิฏฺฐิโก โหติ ฯเปฯ สมฺมาวิมุตฺติ โหติ, อฏฺฐ จ วิโมกฺเข กาเยน ผุสิตฺวา วิหรติ. เอวํ โข ภิกฺขเว ปุคฺคโล สมณปทุโม โหติ.}

\prose{กถญฺจ ภิกฺขเว ปุคฺคโล สมเณสุ สมณ\-สุขุมาโล โหติ? อิธ ภิกฺขเว ภิกฺขุ ยาจิโตว พหุลํ จีวรํ ปริภุญฺชติ อปฺปํ อยาจิโต ฯเปฯ. ยํ หิ ตํ ภิกฺขเว สมฺมา วทมาโน วเทยฺย “สมเณสุ สมณ\-สุขุมาโล”ติ. มเมว ตํ ภิกฺขเว สมฺมา วทมาโน วเทยฺย “สมเณสุ สมณ\-สุขุมาโล”ติ. อิเม โข ภิกฺขเว จตฺตาโร ปุคฺคลา สนฺโต สํวิชฺชมานา โลกสฺมินฺติ.  นวมํ.}

\csromanmatikaanchor{mtk.141}
\csromantocmarkref{subsection}{10. ขนฺธสุตฺต}{mtk.141}
\bookmark[dest={mtk.141},level=2]{10. ขนฺธสุตฺต}
\csromanheader{2}{10. \textbf{ขนฺธสุตฺต}}
\proseitem{90}{จตฺตาโรเม ภิกฺขเว ปุคฺคลา สนฺโต สํวิชฺชมานา โลกสฺมิํ. กตเม จตฺตาโร? สมณมจโล สมณ\-ปุณฺฑรีโก สมณปทุโม สมเณสุ สมณ\-สุขุมาโล.}

\prose{กถญฺจ ภิกฺขเว ปุคฺคโล สมณมจโล โหติ? อิธ ภิกฺขเว ภิกฺขุ เสโข โหติ อปฺปตฺตมานโส, อนุตฺตรํ โยคกฺเขมํ ปตฺถยมาโน วิหรติ. เอวํ โข ภิกฺขเว ปุคฺคโล สมณมจโล โหติ.}

\csromanmatikaanchor{mtk.142}
\csromantocmarknopage{section}{(10) 5. อสุรวคฺค}{mtk.142}
\bookmark[dest={mtk.142},level=1]{(10) 5. อสุรวคฺค}
\csromanmarkh{(10)}\csromanmarkhii{5. อสุรวคฺค}\csromanheadernomark{1}{(10) 5. อสุรวคฺค}
\prose{\csromanfolio{403}กถญฺจ ภิกฺขเว ปุคฺคโล สมณ\-ปุณฺฑรีโก โหติ? อิธ ภิกฺขเว ภิกฺขุ ปญฺจสุ อุปาทาน\-กฺขนฺเธสุ อุทยพฺพยา\-นุปสฺสี วิหรติ “อิติ รูปํ, อิติ รูปสฺส สมุทโย, อิติ รูปสฺส อตฺถงฺคโม. อิติ เวทนา ฯเปฯ. อิติ สญฺญา ฯเปฯ. อิติ สงฺขารา ฯเปฯ. อิติ วิญฺญาณํ, อิติ วิญฺญาณสฺส สมุทโย, อิติ วิญฺญาณสฺส อตฺถงฺคโม”ติ, โน จ โข อฏฺฐ วิโมกฺเข กาเยน ผุสิตฺวา วิหรติ. เอวํ โข ภิกฺขเว ปุคฺคโล สมณ\-ปุณฺฑรีโก โหติ.}

\prose{กถญฺจ ภิกฺขเว ปุคฺคโล สมณปทุโม โหติ? อิธ ภิกฺขเว ภิกฺขุ ปญฺจสุ อุปาทาน\-กฺขนฺเธสุ อุทยพฺพยา\-นุปสฺสี วิหรติ “อิติ รูปํ, อิติ รูปสฺส สมุทโย, อิติ รูปสฺส อตฺถงฺคโม. อิติ เวทนา ฯเปฯ. อิติ สญฺญา ฯเปฯ. อิติ สงฺขารา ฯเปฯ. อิติ วิญฺญาณํ, อิติ วิญฺญาณสฺส สมุทโย, อิติ วิญฺญาณสฺส อตฺถงฺคโม”ติ, อฏฺฐ จ วิโมกฺเข กาเยน ผุสิตฺวา วิหรติ. เอวํ โข ภิกฺขเว ปุคฺคโล สมณปทุโม โหติ.}

\prosewithclosermidruled{กถญฺจ ภิกฺขเว ปุคฺคโล สมเณสุ สมณ\-สุขุมาโล โหติ? อิธ ภิกฺขเว ภิกฺขุ ยาจิโตว พหุลํ จีวรํ ปริภุญฺชติ อปฺปํ อยาจิโต ฯเปฯ. มเมว ตํ ภิกฺขเว สมฺมา วทมาโน วเทยฺย “สมเณสุ สมณ\-สุขุมาโล”ติ. อิเม โข ภิกฺขเว จตฺตาโร ปุคฺคลา สนฺโต สํวิชฺชมานา โลกสฺมินฺติ.  ทสมํ.}{มจลวคฺโค จตุตฺโถ.}
\csromancenter{\textbf{ตสฺสุทฺทานํ}}
\setlength{\csromangathapagewidth}{0pt}
\setlength{\csromangathawakwidth}{0pt}
\csromangathameasuregroup{ปาณาติปาโต จ มุสา, \\ ปุตฺโต สํโยชนญฺเจว,}{\csromangathabat{ปาณาติปาโต จ มุสา,}{อวณฺณโกธตโมณตา.} \\ \csromangathabat{ปุตฺโต สํโยชนญฺเจว,}{ทิฏฺฐิ ขนฺเธน เต ทสาติ.}}
\csromangathasetleft{ปาณาติปาโต จ มุสา, \\ ปุตฺโต สํโยชนญฺเจว,}
\csromangathagroup{\csromangathastanza{\csromangathabat{ปาณาติปาโต จ มุสา,}{อวณฺณโกธตโมณตา.} \\ \csromangathabat{ปุตฺโต สํโยชนญฺเจว,}{ทิฏฺฐิ ขนฺเธน เต ทสาติ.}}}

\chapterhead{(10) 5. อสุรวคฺค}
\csromanmatikaanchor{mtk.143}
\csromantocmarkref{subsection}{1. อสุรสุตฺต}{mtk.143}
\bookmark[dest={mtk.143},level=2]{1. อสุรสุตฺต}
\csromanheader{2}{1. \textbf{อสุรสุตฺต}}
\proseitem{91}{จตฺตาโรเม ภิกฺขเว ปุคฺคลา สนฺโต สํวิชฺชมานา โลกสฺมิํ. กตเม จตฺตาโร? อสุโร อสุรปริวาโร, อสุโร เทวปริวาโร, เทโว อสุรปริวาโร, เทโว เทวปริวาโร.}

\csromanmatikaanchor{mtk.144}
\csromantocmarkref{subsection}{2-4. สมาธิสุตฺตานิ}{mtk.144}
\bookmark[dest={mtk.144},level=2]{2-4. สมาธิสุตฺตานิ}
\prose{\csromanfolio{404}กถญฺจ ภิกฺขเว ปุคฺคโล อสุโร โหติ อสุรปริวาโร? อิธ ภิกฺขเว เอกจฺโจ ปุคฺคโล ทุสฺสีโล โหติ ปาปธมฺโม, ปริสาปิสฺส โหติ ทุสฺสีลา ปาปธมฺมา. เอวํ โข ภิกฺขเว ปุคฺคโล อสุโร โหติ อสุรปริวาโร.}

\prose{กถญฺจ ภิกฺขเว ปุคฺคโล อสุโร โหติ เทวปริวาโร? อิธ ภิกฺขเว เอกจฺโจ ปุคฺคโล ทุสฺสีโล โหติ ปาปธมฺโม, ปริสา จ ขฺวสฺส โหติ สีลวตี กลฺยาณธมฺมา. เอวํ โข ภิกฺขเว ปุคฺคโล อสุโร โหติ เทวปริวาโร.}

\prose{กถญฺจ ภิกฺขเว ปุคฺคโล เทโว โหติ อสุรปริวาโร? อิธ ภิกฺขเว เอกจฺโจ ปุคฺคโล สีลวา โหติ กลฺยาณธมฺโม, ปริสา จ ขฺวสฺส โหติ ทุสฺสีลา ปาปธมฺมา. เอวํ โข ภิกฺขเว ปุคฺคโล เทโว โหติ อสุรปริวาโร.}

\prose{กถญฺจ ภิกฺขเว ปุคฺคโล เทโว โหติ เทวปริวาโร? อิธ ภิกฺขเว เอกจฺโจ ปุคฺคโล สีลวา โหติ กลฺยาณธมฺโม, ปริสาปิสฺส โหติ สีลวตี กลฺยาณธมฺมา. เอวํ โข ภิกฺขเว ปุคฺคโล เทโว โหติ เทวปริวาโร. อิเม โข ภิกฺขเว จตฺตาโร ปุคฺคลา สนฺโต สํวิชฺชมานา โลกสฺมินฺติ.  ปฐมํ.}

\csromanheader{2}{2. ปฐมสมาธิ\-สุตฺต}
\proseitem{92}{จตฺตาโรเม ภิกฺขเว ปุคฺคลา สนฺโต สํวิชฺชมานา โลกสฺมิํ. กตเม จตฺตาโร? อิธ ภิกฺขเว เอกจฺโจ ปุคฺคโล ลาภี โหติ อชฺฌตฺตํ เจโตสมถสฺส, น ลาภี อธิปญฺญา\-ธมฺม\-วิปสฺสนาย. อิธ ปน ภิกฺขเว เอกจฺโจ ปุคฺคโล ลาภี โหติ อธิปญฺญา\-ธมฺม\-วิปสฺสนาย, น ลาภี อชฺฌตฺตํ เจโตสมถสฺส. อิธ ปน ภิกฺขเว เอกจฺโจ ปุคฺคโล น เจว ลาภี โหติ อชฺฌตฺตํ เจโตสมถสฺส, น จ ลาภี อธิปญฺญา\-ธมฺม\-วิปสฺสนาย. อิธ ปน ภิกฺขเว เอกจฺโจ ปุคฺคโล ลาภี เจว โหติ อชฺฌตฺตํ เจโตสมถสฺส, ลาภี จ อธิปญฺญา\-ธมฺม\-วิปสฺสนาย. อิเม โข ภิกฺขเว จตฺตาโร ปุคฺคลา สนฺโต สํวิชฺชมานา โลกสฺมินฺติ.  ทุติยํ.}

\chapterheadpageread{(10) 5. อสุรวคฺค 3. ทุติย\-สมาธิ\-สุตฺต}
\proseitem{93}{\csromanfolio{405}จตฺตาโรเม ภิกฺขเว ปุคฺคลา สนฺโต สํวิชฺชมานา โลกสฺมิํ. กตเม จตฺตาโร? อิธ ภิกฺขเว เอกจฺโจ ปุคฺคโล ลาภี โหติ อชฺฌตฺตํ เจโตสมถสฺส, น ลาภี อธิปญฺญา\-ธมฺม\-วิปสฺสนาย. อิธ ปน ภิกฺขเว เอกจฺโจ ปุคฺคโล ลาภี โหติ อธิปญฺญา\-ธมฺม\-วิปสฺสนาย, น ลาภี อชฺฌตฺตํ เจโตสมถสฺส. อิธ ปน ภิกฺขเว เอกจฺโจ ปุคฺคโล น เจว ลาภี โหติ อชฺฌตฺตํ เจโตสมถสฺส, น จ ลาภี อธิปญฺญา\-ธมฺม\-วิปสฺสนาย. อิธ ปน ภิกฺขเว เอกจฺโจ ปุคฺคโล ลาภี เจว โหติ อชฺฌตฺตํ เจโตสมถสฺส, ลาภี จ อธิปญฺญา\-ธมฺม\-วิปสฺสนาย.}

\prose{ตตฺร ภิกฺขเว ยฺวายํ ปุคฺคโล ลาภี โหติ อชฺฌตฺตํ เจโตสมถสฺส, น ลาภี อธิปญฺญา\-ธมฺม\-วิปสฺสนาย. เตน ภิกฺขเว ปุคฺคเลน อชฺฌตฺตํ เจโตสมเถ ปติฏฺฐาย อธิปญฺญา\-ธมฺม\-วิปสฺสนาย โยโค กรณีโย, โส อปเรน สมเยน ลาภี เจว โหติ อชฺฌตฺตํ เจโตสมถสฺส, ลาภี จ อธิปญฺญา\-ธมฺม\-วิปสฺสนาย.}

\prose{ตตฺร ภิกฺขเว ยฺวายํ ปุคฺคโล ลาภี อธิปญฺญา\-ธมฺม\-วิปสฺสนาย, น ลาภี อชฺฌตฺตํ เจโตสมถสฺส. เตน ภิกฺขเว ปุคฺคเลน อธิปญฺญา\-ธมฺม\-วิปสฺสนาย ปติฏฺฐาย อชฺฌตฺตํ เจโตสมเถ โยโค กรณีโย, โส อปเรน สมเยน ลาภี เจว โหติ อธิปญฺญา\-ธมฺม\-วิปสฺสนาย, ลาภี จ อชฺฌตฺตํ เจโตสมถสฺส.}

\prose{ตตฺร ภิกฺขเว ยฺวายํ ปุคฺคโล น เจว ลาภี อชฺฌตฺตํ เจโตสมถสฺส, น จ ลาภี อธิปญฺญา\-ธมฺม\-วิปสฺสนาย. เตน ภิกฺขเว ปุคฺคเลน เตสํเยว กุสลานํ ธมฺมานํ ปฏิลาภาย อธิมตฺโต ฉนฺโท จ วายาโม จ อุสฺสาโห จ อุสฺโสฬฺหี จ อปฺปฏิวานี จ สติ จ สมฺปชญฺญญฺจ กรณียํ. เสยฺยถาปิ ภิกฺขเว อาทิตฺตเจโล วา อาทิตฺตสีโส วา ตสฺเสว\csromansharedfootnote{091e98f0e71ae6a2}{ตสฺส ตสฺเสว (สี, สฺยา, กํ)} เจลสฺส วา สีสสฺส วา นิพฺพาปนาย อธิมตฺตํ ฉนฺทญฺจ วายามญฺจ อุสฺสาหญฺจ อุสฺโสฬฺหิญฺจ อปฺปฏิวานิญฺจ สติญฺจ สมฺปชญฺญญฺจ กเรยฺย. เอวเมวํ โข ภิกฺขเว เตน ปุคฺคเลน เตสํเยว กุสลานํ ธมฺมานํ ปฏิลาภาย อธิมตฺโต ฉนฺโท จ วายาโม จ อุสฺสาโห จ อุสฺโสฬฺหี จ อปฺปฏิวานี \csromanfolio{406}จ สติ จ สมฺปชญฺญญฺจ กรณียํ. โส อปเรน สมเยน ลาภี เจว โหติ อชฺฌตฺตํ เจโตสมถสฺส, ลาภี จ อธิปญฺญา\-ธมฺม\-วิปสฺสนาย.}

\prose{ตตฺร ภิกฺขเว ยฺวายํ ปุคฺคโล ลาภี เจว โหติ อชฺฌตฺตํ เจโตสมถสฺส, ลาภี จ อธิปญฺญา\-ธมฺม\-วิปสฺสนาย. เตน ภิกฺขเว ปุคฺคเลน เตสุเยว กุสเลสุ ธมฺเมสุ ปติฏฺฐาย อุตฺตริ\csromansharedfootnote{3d336730ed7a9e01}{อุตฺตริํ (สี, สฺยา, กํ, อิ)} อาสวานํ ขยาย โยโค กรณีโย. อิเม โข ภิกฺขเว จตฺตาโร ปุคฺคลา สนฺโต สํวิชฺชมานา โลกสฺมินฺติ.  ตติยํ.}

\csromanheader{2}{4. ตติย\-สมาธิ\-สุตฺต}
\prose{94 จตฺตาโรเม ภิกฺขเว ปุคฺคลา สนฺโต สํวิชฺชมานา โลกสฺมิํ. กตเม จตฺตาโร? อิธ ภิกฺขเว เอกจฺโจ ปุคฺคโล ลาภี โหติ อชฺฌตฺตํ เจโตสมถสฺส, น ลาภี อธิปญฺญา\-ธมฺม\-วิปสฺสนาย. อิธ ปน ภิกฺขเว เอกจฺโจ ปุคฺคโล ลาภี โหติ อธิปญฺญา\-ธมฺม\-วิปสฺสนาย, น ลาภี อชฺฌตฺตํ เจโตสมถสฺส. อิธ ปน ภิกฺขเว เอกจฺโจ ปุคฺคโล น เจว ลาภี โหติ อชฺฌตฺตํ เจโตสมถสฺส, น จ ลาภี อธิปญฺญา\-ธมฺม\-วิปสฺสนาย. อิธ ปน ภิกฺขเว เอกจฺโจ ปุคฺคโล ลาภี เจว โหติ อชฺฌตฺตํ เจโตสมถสฺส, ลาภี จ อธิปญฺญา\-ธมฺม\-วิปสฺสนาย.}

\prose{ตตฺร ภิกฺขเว ยฺวายํ ปุคฺคโล ลาภี อชฺฌตฺตํ เจโตสมถสฺส, น ลาภี อธิปญฺญา\-ธมฺม\-วิปสฺสนาย. เตน ภิกฺขเว ปุคฺคเลน ยฺวายํ ปุคฺคโล ลาภี อธิปญฺญา\-ธมฺม\-วิปสฺสนาย, โส อุปสงฺกมิตฺวา เอวมสฺส วจนีโย “กถํ นุ โข อาวุโส สงฺขารา ทฏฺฐพฺพา, กถํ สงฺขารา สมฺมสิตพฺพา, กถํ สงฺขารา วิปสฺสิตพฺพา”\csromansharedfootnote{f9237889c86076b0}{ปสฺสิตพฺพา (ก)}ติ. ตสฺส โส ยถาทิฏฺฐํ ยถาวิทิตํ พฺยากโรติ “เอวํ โข อาวุโส สงฺขารา ทฏฺฐพฺพา, เอวํ สงฺขารา สมฺมสิตพฺพา, เอวํ สงฺขารา วิปสฺสิตพฺพา”ติ. โส อปเรน สมเยน ลาภี เจว โหติ อชฺฌตฺตํ เจโตสมถสฺส, ลาภี จ อธิปญฺญา\-ธมฺม\-วิปสฺสนาย.}

\prose{ตตฺร ภิกฺขเว ยฺวายํ ปุคฺคโล ลาภี อธิปญฺญา\-ธมฺม\-วิปสฺสนาย, น ลาภี อชฺฌตฺตํ เจโตสมถสฺส. เตน ภิกฺขเว ปุคฺคเลน ยฺวายํ ปุคฺคโล ลาภี อชฺฌตฺตํ เจโตสมถสฺส, โส อุปสงฺกมิตฺวา เอวมสฺส วจนีโย “กถํ นุ โข อาวุโส จิตฺตํ สณฺฐเปตพฺพํ, กถํ จิตฺตํ}

\vspace{\tipitakaparskip}
\csromancenter{(10) 5. อสุรวคฺค สนฺนิ\-สาเท\-ตพฺพํ, กถํ จิตฺตํ เอโกทิ กาตพฺพํ, กถํ จิตฺตํ สมาทหาตพฺพนฺ”ติ. ตสฺส โส ยถาทิฏฺฐํ ยถาวิทิตํ พฺยากโรติ “เอวํ โข อาวุโส จิตฺตํ สณฺฐเปตพฺพํ, เอวํ จิตฺตํ สนฺนิ\-สาเท\-ตพฺพํ, เอวํ จิตฺตํ เอโกทิ กาตพฺพํ\csromansharedfootnote{dbae4432334ee938}{เอโกทิ กตฺตพฺพํ (อิ)}, เอวํ จิตฺตํ สมาทหาตพฺพนฺ”ติ, โส อปเรน สมเยน ลาภี เจว โหติ อธิปญฺญา\-ธมฺม\-วิปสฺสนาย, ลาภี จ อชฺฌตฺตํ เจโตสมถสฺส.}
\prose{\csromanfolio{407}ตตฺร ภิกฺขเว ยฺวายํ ปุคฺคโล น เจว ลาภี อชฺฌตฺตํ เจโตสมถสฺส, น จ ลาภี อธิปญฺญา\-ธมฺม\-วิปสฺสนาย. เตน ภิกฺขเว ปุคฺคเลน ยฺวายํ ปุคฺคโล ลาภี เจว อชฺฌตฺตํ เจโตสมถสฺส, ลาภี จ อธิปญฺญา\-ธมฺม\-วิปสฺสนาย, โส อุปสงฺกมิตฺวา เอวมสฺส วจนีโย “กถํ นุ โข อาวุโส จิตฺตํ สณฺฐเปตพฺพํ, กถํ จิตฺตํ สนฺนิ\-สาเท\-ตพฺพํ, กถํ จิตฺตํ เอโกทิ กาตพฺพํ, กถํ จิตฺตํ สมาทหาตพฺพํ, กถํ สงฺขารา ทฏฺฐพฺพา, กถํ สงฺขารา สมฺมสิตพฺพา, กถํ สงฺขารา วิปสฺสิตพฺพา”ติ. ตสฺส โส ยถาทิฏฺฐํ ยถาวิทิตํ พฺยากโรติ “เอวํ โข อาวุโส จิตฺตํ สณฺฐเปตพฺพํ, เอวํ จิตฺตํ สนฺนิ\-สาเท\-ตพฺพํ, เอวํ จิตฺตํ เอโกทิ กาตพฺพํ, เอวํ จิตฺตํ สมาทหาตพฺพํ, เอวํ สงฺขารา ทฏฺฐพฺพา, เอวํ สงฺขารา สมฺมสิตพฺพา, เอวํ สงฺขารา วิปสฺสิตพฺพา”ติ. โส อปเรน สมเยน ลาภี เจว โหติ อชฺฌตฺตํ เจโตสมถสฺส, ลาภี จ อธิปญฺญา\-ธมฺม\-วิปสฺสนาย.}

\prose{ตตฺร ภิกฺขเว ยฺวายํ ปุคฺคโล ลาภี เจว โหติ อชฺฌตฺตํ เจโตสมถสฺส, ลาภี จ อธิปญฺญา\-ธมฺม\-วิปสฺสนาย. เตน ภิกฺขเว ปุคฺคเลน เตสุ เจว กุสเลสุ ธมฺเมสุ ปติฏฺฐาย อุตฺตริ อาสวานํ ขยาย โยโค กรณีโย. อิเม โข ภิกฺขเว จตฺตาโร ปุคฺคลา สนฺโต สํวิชฺชมานา โลกสฺมินฺติ.  จตุตฺถํ.}

\csromanmatikaanchor{mtk.145}
\csromantocmarkref{subsection}{5. ฉวาลาตสุตฺต}{mtk.145}
\bookmark[dest={mtk.145},level=2]{5. ฉวาลาตสุตฺต}
\csromanheader{2}{5. \textbf{ฉวาลาตสุตฺต}}
\proseitem{95}{จตฺตาโรเม ภิกฺขเว ปุคฺคลา สนฺโต สํวิชฺชมานา โลกสฺมิํ. กตเม จตฺตาโร? เนวตฺตหิตาย ปฏิปนฺโน โน ปรหิตาย, ปรหิตาย ปฏิปนฺโน โน อตฺตหิตาย, อตฺตหิตาย ปฏิปนฺโน โน ปรหิตาย, อตฺตหิตาย เจว\csromansharedfootnote{4f5cb74853f900f2}{อตฺตหิตาย จ (สี, สฺยา, กํ, อิ)} ปฏิปนฺโน ปรหิตาย จ.}

\prose{\csromanfolio{408}เสยฺยถาปิ ภิกฺขเว ฉวาลาตํ อุภโต ปทิตฺตํ\csromansharedfootnote{6efb2359fd53a138}{อาทิตฺตํ (ก)} มชฺเฌ คูถคตํ เนว คาเม กฏฺฐตฺถํ ผรติ, น อรญฺเญ ( )\csromansharedfootnote{44b8da626eecfdeb}{(กฏฺฐตฺถํ ผรติ) กตฺถจิ.}. ตถูปมาหํ ภิกฺขเว อิมํ ปุคฺคลํ วทามิ, ยฺวายํ ปุคฺคโล เนวตฺตหิตาย ปฏิปนฺโน โน ปรหิตาย.}

\prose{ตตฺร ภิกฺขเว ยฺวายํ ปุคฺคโล ปรหิตาย ปฏิปนฺโน โน อตฺตหิตาย, อยํ อิเมสํ ทฺวินฺนํ ปุคฺคลานํ อภิกฺกนฺตตโร จ ปณีตตโร จ. ตตฺร ภิกฺขเว ยฺวายํ ปุคฺคโล อตฺตหิตาย ปฏิปนฺโน โน ปรหิตาย, อยํ อิเมสํ ติณฺณํ ปุคฺคลานํ อภิกฺกนฺตตโร จ ปณีตตโร จ. ตตฺร ภิกฺขเว ยฺวายํ ปุคฺคโล อตฺตหิตาย เจว ปฏิปนฺโน ปรหิตาย จ, อยํ อิเมสํ จตุนฺนํ ปุคฺคลานํ อคฺโค จ เสฏฺโฐ จ ปาโมกฺโข\csromansharedfootnote{c3dd9b40705d495c}{โมกฺโข (อิ) สํ 2. 224; อํ 2. 193 ปิฏฺเฐสุปิ.} จ อุตฺตโม จ ปวโร จ.}

\prose{เสยฺยถาปิ ภิกฺขเว ควา ขีรํ, ขีรมฺหา ทธิ, ทธิมฺหา นวนีตํ, นวนีตมฺหา สปฺปิ, สปฺปิมฺหา สปฺปิมณฺโฑ, สปฺปิมณฺโฑ\csromansharedfootnote{1377eddc129c8fed}{สปฺปิมฺหา สปฺปิมณฺโฑ (ก)} ตตฺถ\csromansharedfootnote{bca22a05b973a22a}{ตตฺร (สํ 2. 224 ปิฏฺเฐ)} อคฺคมกฺขายติ. เอวเมวํ โข ภิกฺขเว ยฺวายํ ปุคฺคโล อตฺตหิตาย เจว ปฏิปนฺโน ปรหิตาย จ, อยํ อิเมสํ จตุนฺนํ ปุคฺคลานํ อคฺโค จ เสฏฺโฐ จ ปาโมกฺโข จ อุตฺตโม จ ปวโร จ. อิเม โข ภิกฺขเว จตฺตาโร ปุคฺคลา สนฺโต สํวิชฺชมานา โลกสฺมินฺติ.  ปญฺจมํ.}

\csromanmatikaanchor{mtk.146}
\csromantocmarkref{subsection}{6. ราควินยสุตฺต}{mtk.146}
\bookmark[dest={mtk.146},level=2]{6. ราควินยสุตฺต}
\csromanheader{2}{6. ราควินย\-สุตฺต}
\proseitem{96}{จตฺตาโรเม ภิกฺขเว ปุคฺคลา สนฺโต สํวิชฺชมานา โลกสฺมิํ. กตเม จตฺตาโร? อตฺตหิตาย ปฏิปนฺโน โน ปรหิตาย, ปรหิตาย ปฏิปนฺโน โน อตฺตหิตาย, เนวตฺตหิตาย ปฏิปนฺโน โน ปรหิตาย, อตฺตหิตาย เจว ปฏิปนฺโน ปรหิตาย จ.}

\prose{กถญฺจ ภิกฺขเว ปุคฺคโล อตฺตหิตาย ปฏิปนฺโน โหติ โน ปรหิตาย? อิธ ภิกฺขเว เอกจฺโจ ปุคฺคโล อตฺตนา ราควินยาย ปฏิปนฺโน โหติ, โน ปรํ ราควินยาย สมาทเปติ. อตฺตนา โทสวินยาย ปฏิปนฺโน โหติ, โน ปรํ โทสวินยาย สมาทเปติ.}

\vspace{\tipitakaparskip}
\csromancenter{(10) 5. อสุรวคฺค อตฺตนา โมหวินยาย ปฏิปนฺโน โหติ, โน ปรํ โมหวินยาย สมาทเปติ. เอวํ โข ภิกฺขเว ปุคฺคโล อตฺตหิตาย ปฏิปนฺโน โหติ โน ปรหิตาย.}
\prose{\csromanfolio{409}กถญฺจ ภิกฺขเว ปุคฺคโล ปรหิตาย ปฏิปนฺโน โหติ โน อตฺตหิตาย? อิธ ภิกฺขเว เอกจฺโจ ปุคฺคโล อตฺตนา น ราควินยาย ปฏิปนฺโน โหติ, ปรํ ราควินยาย สมาทเปติ. อตฺตนา น โทสวินยาย ปฏิปนฺโน โหติ, ปรํ โทสวินยาย สมาทเปติ. อตฺตนา น โมหวินยาย ปฏิปนฺโน โหติ, ปรํ โมหวินยาย สมาทเปติ. เอวํ โข ภิกฺขเว ปุคฺคโล ปรหิตาย ปฏิปนฺโน โหติ โน อตฺตหิตาย.}

\prose{กถญฺจ ภิกฺขเว ปุคฺคโล เนวตฺตหิตาย ปฏิปนฺโน โหติ โน ปรหิตาย? อิธ ภิกฺขเว เอกจฺโจ ปุคฺคโล อตฺตนา น ราควินยาย ปฏิปนฺโน โหติ, โน ปรํ ราควินยาย สมาทเปติ. อตฺตนา น โทสวินยาย ปฏิปนฺโน โหติ, โน ปรํ โทสวินยาย สมาทเปติ. อตฺตนา น โมหวินยาย ปฏิปนฺโน โหติ, โน ปรํ โมหวินยาย สมาทเปติ. เอวํ โข ภิกฺขเว ปุคฺคโล เนวตฺตหิตาย ปฏิปนฺโน โหติ โน ปรหิตาย.}

\prose{กถญฺจ ภิกฺขเว ปุคฺคโล อตฺตหิตาย เจว ปฏิปนฺโน โหติ ปรหิตาย จ? อิธ ภิกฺขเว เอกจฺโจ ปุคฺคโล อตฺตนา จ ราควินยาย ปฏิปนฺโน โหติ, ปรญฺจ ราควินยาย สมาทเปติ. อตฺตนา จ โทสวินยาย ปฏิปนฺโน โหติ, ปรญฺจ โทสวินยาย สมาทเปติ. อตฺตนา จ โมหวินยาย ปฏิปนฺโน โหติ, ปรญฺจ โมหวินยาย สมาทเปติ. เอวํ โข ภิกฺขเว ปุคฺคโล อตฺตหิตาย เจว ปฏิปนฺโน โหติ ปรหิตาย จ. อิเม โข ภิกฺขเว จตฺตาโร ปุคฺคลา สนฺโต สํวิชฺชมานา โลกสฺมินฺติ.  ฉฏฺฐํ.}

\csromanmatikaanchor{mtk.147}
\csromantocmarkref{subsection}{7. ขิปฺปนิสนฺติสุตฺต}{mtk.147}
\bookmark[dest={mtk.147},level=2]{7. ขิปฺปนิสนฺติสุตฺต}
\csromanheader{2}{7. ขิปฺปนิสนฺติ\-สุตฺต}
\proseitem{97}{จตฺตาโรเม ภิกฺขเว ปุคฺคลา สนฺโต สํวิชฺชมานา โลกสฺมิํ. กตเม จตฺตาโร? อตฺตหิตาย ปฏิปนฺโน โน ปรหิตาย, ปรหิตาย ปฏิปนฺโน โน อตฺตหิตาย, เนวตฺตหิตาย ปฏิปนฺโน โน ปรหิตาย, อตฺตหิตาย เจว ปฏิปนฺโน ปรหิตาย จ.}

\prose{\csromanfolio{410}กถญฺจ ภิกฺขเว ปุคฺคโล อตฺตหิตาย ปฏิปนฺโน โหติ โน ปรหิตาย? อิธ ภิกฺขเว เอกจฺโจ ปุคฺคโล ขิปฺปนิสนฺตี จ โหติ กุสเลสุ ธมฺเมสุ, สุตานญฺจ ธมฺมานํ ธารกชาติโก\csromansharedfootnote{eacfa4ad9dd160a0}{ธารณชาติโก (ก)} โหติ. ธาตานญฺจ\csromansharedfootnote{87d4fc68e029070a}{ธตานญฺจ (สี, สฺยา, กํ, อิ)} ธมฺมานํ อตฺถู\-ปปริกฺขี โหติ, อตฺถมญฺญาย ธมฺมมญฺญาย ธมฺมา\-นุธมฺม\-ปฺปฏิปนฺโน โหติ. โน จ กลฺยาณวาโจ โหติ กลฺยาณ\-วากฺกรโณ โปริยา วาจาย สมนฺนาคโต วิสฺสฏฺฐาย อเนลคลาย อตฺถสฺส วิญฺญาปนิยา, โน จ สนฺทสฺสโก โหติ สมาทปโก\csromansharedfootnote{276e18aaa1c77713}{สมาทาปโก (?)} สมุตฺเตชโก สมฺปหํสโก สพฺรหฺมจารีนํ. เอวํ โข ภิกฺขเว ปุคฺคโล อตฺตหิตาย ปฏิปนฺโน โหติ โน ปรหิตาย.}

\prose{กถญฺจ ภิกฺขเว ปุคฺคโล ปรหิตาย ปฏิปนฺโน โหติ โน อตฺตหิตาย? อิธ ภิกฺขเว เอกจฺโจ ปุคฺคโล น เหว โข ขิปฺปนิสนฺตี โหติ กุสเลสุ ธมฺเมสุ, โน จ สุตานํ ธมฺมานํ ธารกชาติโก โหติ, โน จ ธาตานํ ธมฺมานํ อตฺถู\-ปปริกฺขี โหติ, โน จ อตฺถมญฺญาย ธมฺมมญฺญาย ธมฺมา\-นุธมฺม\-ปฺปฏิปนฺโน โหติ. กลฺยาณวาโจ จ โหติ กลฺยาณ\-วากฺกรโณ โปริยา วาจาย สมนฺนาคโต วิสฺสฏฺฐาย อเนลคลาย อตฺถสฺส วิญฺญาปนิยา, สนฺทสฺสโก จ โหติ สมาทปโก สมุตฺเตชโก สมฺปหํสโก สพฺรหฺมจารีนํ. เอวํ โข ภิกฺขเว ปุคฺคโล ปรหิตาย ปฏิปนฺโน โหติ โน อตฺตหิตาย.}

\prose{กถญฺจ ภิกฺขเว ปุคฺคโล เนวตฺตหิตาย ปฏิปนฺโน โหติ โน ปรหิตาย? อิธ ภิกฺขเว เอกจฺโจ ปุคฺคโล น เหว โข ขิปฺปนิสนฺตี โหติ กุสเลสุ ธมฺเมสุ, โน จ สุตานํ ธมฺมานํ ธารกชาติโก โหติ, โน จ ธาตานํ ธมฺมานํ อตฺถู\-ปปริกฺขี โหติ, โน จ อตฺถมญฺญาย ธมฺมมญฺญาย ธมฺมา\-นุธมฺม\-ปฺปฏิปนฺโน โหติ. โน จ กลฺยาณวาโจ โหติ กลฺยาณ\-วากฺกรโณ โปริยา วาจาย สมนฺนาคโต วิสฺสฏฺฐาย อเนลคลาย อตฺถสฺส วิญฺญาปนิยา, โน จ สนฺทสฺสโก โหติ สมาทปโก สมุตฺเตชโก สมฺปหํสโก สพฺรหฺมจารีนํ. เอวํ โข ภิกฺขเว ปุคฺคโล เนวตฺตหิตาย ปฏิปนฺโน โหติ โน ปรหิตาย.}

\chapterheadpageread{(10) 5. อสุรวคฺค}
\prose{\csromanfolio{411}กถญฺจ ภิกฺขเว ปุคฺคโล อตฺตหิตาย เจว ปฏิปนฺโน โหติ ปรหิตาย จ? อิธ ภิกฺขเว เอกจฺโจ ปุคฺคโล ขิปฺปนิสนฺตี จ โหติ กุสเลสุ ธมฺเมสุ, สุตานญฺจ ธมฺมานํ ธารกชาติโก โหติ, ธาตานญฺจ ธมฺมานํ อตฺถู\-ปปริกฺขี โหติ, อตฺถมญฺญาย ธมฺมมญฺญาย ธมฺมา\-นุธมฺม\-ปฺปฏิปนฺโน โหติ. กลฺยาณวาโจ จ โหติ กลฺยาณ\-วากฺกรโณ โปริยา วาจาย สมนฺนาคโต วิสฺสฏฺฐาย อเนลคลาย อตฺถสฺส วิญฺญาปนิยา, สนฺทสฺสโก จ โหติ สมาทปโก สมุตฺเตชโก สมฺปหํสโก สพฺรหฺมจารีนํ. เอวํ โข ภิกฺขเว ปุคฺคโล อตฺตหิตาย เจว ปฏิปนฺโน โหติ ปรหิตาย จ. อิเม โข ภิกฺขเว จตฺตาโร ปุคฺคลา สนฺโต สํวิชฺชมานา โลกสฺมินฺติ.  สตฺตมํ.}

\csromanmatikaanchor{mtk.148}
\csromantocmarkref{subsection}{8-9. อตฺตหิตสุตฺตทฺวยํ}{mtk.148}
\bookmark[dest={mtk.148},level=2]{8-9. อตฺตหิตสุตฺตทฺวยํ}
\csromanheader{2}{8. \textbf{อตฺตหิตสุตฺต}}
\proseitem{98}{จตฺตาโรเม ภิกฺขเว ปุคฺคลา สนฺโต สํวิชฺชมานา โลกสฺมิํ. กตเม จตฺตาโร? อตฺตหิตาย ปฏิปนฺโน โน ปรหิตาย, ปรหิตาย ปฏิปนฺโน โน อตฺตหิตาย, เนวตฺตหิตาย ปฏิปนฺโน โน ปรหิตาย, อตฺตหิตาย เจว ปฏิปนฺโน ปรหิตาย จ. อิเม โข ภิกฺขเว จตฺตาโร ปุคฺคลา สนฺโต สํวิชฺชมานา โลกสฺมินฺติ.  อฏฺฐมํ.}

\csromanheader{2}{9. สิกฺขาปท\-สุตฺต}
\proseitem{99}{จตฺตาโรเม ภิกฺขเว ปุคฺคลา สนฺโต สํวิชฺชมานา โลกสฺมิํ. กตเม จตฺตาโร? อตฺตหิตาย ปฏิปนฺโน โน ปรหิตาย, ปรหิตาย ปฏิปนฺโน โน อตฺตหิตาย, เนวตฺตหิตาย ปฏิปนฺโน โน ปรหิตาย, อตฺตหิตาย เจว ปฏิปนฺโน ปรหิตาย จ.}

\prose{กถญฺจ ภิกฺขเว ปุคฺคโล อตฺตหิตาย ปฏิปนฺโน โหติ โน ปรหิตาย? อิธ ภิกฺขเว เอกจฺโจ ปุคฺคโล อตฺตนา ปาณาติปาตา ปฏิวิรโต โหติ, โน ปรํ ปาณาติปาตา เวรมณิยา สมาทเปติ. อตฺตนา อทินฺนาทานา ปฏิวิรโต โหติ, โน ปรํ อทินฺนาทานา เวรมณิยา สมาทเปติ. อตฺตนา กาเมสุ\-มิจฺฉาจารา ปฏิวิรโต โหติ, โน ปรํ กาเมสุ\-มิจฺฉาจารา เวรมณิยา สมาทเปติ. อตฺตนา มุสาวาทา ปฏิวิรโต โหติ, โน ปรํ มุสาวาทา เวรมณิยา \csromanfolio{412}สมาทเปติ. อตฺตนา สุรา\-เมรย\-มชฺช\-ปมาทฏฺฐานา ปฏิวิรโต โหติ, โน ปรํ สุรา\-เมรย\-มชฺช\-ปมาทฏฺฐานา เวรมณิยา สมาทเปติ. เอวํ โข ภิกฺขเว ปุคฺคโล อตฺตหิตาย ปฏิปนฺโน โหติ โน ปรหิตาย.}

\prose{กถญฺจ ภิกฺขเว ปุคฺคโล ปรหิตาย ปฏิปนฺโน โหติ โน อตฺตหิตาย? อิธ ภิกฺขเว เอกจฺโจ ปุคฺคโล อตฺตนา ปาณาติปาตา อปฺปฏิวิรโต โหติ, ปรํ ปาณาติปาตา เวรมณิยา สมาทเปติ. อตฺตนา อทินฺนาทานา อปฺปฏิวิรโต โหติ, ปรํ อทินฺนาทานา เวรมณิยา สมาทเปติ. อตฺตนา กาเมสุ\-มิจฺฉาจารา อปฺปฏิวิรโต โหติ, ปรํ กาเมสุ\-มิจฺฉาจารา เวรมณิยา สมาทเปติ. อตฺตนา มุสาวาทา อปฺปฏิวิรโต โหติ, ปรํ มุสาวาทา เวรมณิยา สมาทเปติ. อตฺตนา สุรา\-เมรย\-มชฺช\-ปมาทฏฺฐานา อปฺปฏิวิรโต โหติ, ปรํ สุรา\-เมรย\-มชฺช\-ปมาทฏฺฐานา เวรมณิยา สมาทเปติ. เอวํ โข ภิกฺขเว ปุคฺคโล ปรหิตาย ปฏิปนฺโน โหติ โน อตฺตหิตาย.}

\prose{กถญฺจ ภิกฺขเว ปุคฺคโล เนวตฺตหิตาย ปฏิปนฺโน โหติ โน ปรหิตาย? อิธ ภิกฺขเว เอกจฺโจ ปุคฺคโล อตฺตนา ปาณาติปาตา อปฺปฏิวิรโต โหติ, โน ปรํ ปาณาติปาตา เวรมณิยา สมาทเปติ ฯเปฯ. อตฺตนา สุรา\-เมรย\-มชฺช\-ปมาทฏฺฐานา อปฺปฏิวิรโต โหติ, โน ปรํ สุรา\-เมรย\-มชฺช\-ปมาทฏฺฐานา เวรมณิยา สมาทเปติ. เอวํ โข ภิกฺขเว ปุคฺคโล เนวตฺตหิตาย ปฏิปนฺโน โหติ โน ปรหิตาย.}

\prose{กถญฺจ ภิกฺขเว ปุคฺคโล อตฺตหิตาย เจว ปฏิปนฺโน โหติ ปรหิตาย จ? อิธ ภิกฺขเว เอกจฺโจ ปุคฺคโล อตฺตนา จ ปาณาติปาตา ปฏิวิรโต โหติ, ปรญฺจ ปาณาติปาตา เวรมณิยา สมาทเปติ ฯเปฯ. อตฺตนา จ สุรา\-เมรย\-มชฺช\-ปมาทฏฺฐานา ปฏิวิรโต โหติ, ปรญฺจ สุรา\-เมรย\-มชฺช\-ปมาทฏฺฐานา เวรมณิยา สมาทเปติ. เอวํ โข ภิกฺขเว ปุคฺคโล อตฺตหิตาย เจว ปฏิปนฺโน โหติ ปรหิตาย จ. อิเม โข ภิกฺขเว จตฺตาโร ปุคฺคลา สนฺโต สํวิชฺชมานา โลกสฺมินฺติ.  นวมํ.}

\csromanmatikaanchor{mtk.149}
\csromantocmarkref{subsection}{10. โปตลิยสุตฺต}{mtk.149}
\bookmark[dest={mtk.149},level=2]{10. โปตลิยสุตฺต}
\csromanmarkh{(10)}\csromanmarkhii{5. อสุรวคฺค 10. โปตลิยสุตฺต}\csromanheadernomark{2}{(10) 5. อสุรวคฺค \textbf{10. โปตลิยสุตฺต}}
\proseitem{100}{\csromanfolio{413}อถ โข โปตลิโย ปริพฺพาชโก เยน ภควา เตนุปสงฺกมิ, อุปสงฺกมิตฺวา ภควตา สทฺธิํ สมฺโมทิ, สมฺโมทนียํ กถํ สารณียํ วีติสาเรตฺวา เอกมนฺตํ นิสีทิ, เอกมนฺตํ นิสินฺนํ โข โปตลิยํ ปริพฺพาชกํ ภควา เอตทโวจ\csromandash{}}

\prose{จตฺตาโรเม โปตลิย ปุคฺคลา สนฺโต สํวิชฺชมานา โลกสฺมิํ. กตเม จตฺตาโร? อิธ โปตลิย เอกจฺโจ ปุคฺคโล\csromansharedfootnote{2fa1b01b18ef2646}{อภิ 3. 157 ปิฏฺเฐปิ.} อวณฺณารหสฺส อวณฺณํ ภาสิตา โหติ ภูตํ ตจฺฉํ กาเลน, โน จ โข วณฺณารหสฺส วณฺณํ ภาสิตา โหติ ภูตํ ตจฺฉํ กาเลน. อิธ ปน โปตลิย เอกจฺโจ ปุคฺคโล วณฺณารหสฺส วณฺณํ ภาสิตา โหติ ภูตํ ตจฺฉํ กาเลน, โน จ โข อวณฺณารหสฺส อวณฺณํ ภาสิตา โหติ ภูตํ ตจฺฉํ กาเลน. อิธ ปน โปตลิย เอกจฺโจ ปุคฺคโล เนว อวณฺณารหสฺส อวณฺณํ ภาสิตา โหติ ภูตํ ตจฺฉํ กาเลน, โน จ วณฺณารหสฺส วณฺณํ ภาสิตา โหติ ภูตํ ตจฺฉํ กาเลน. อิธ ปน โปตลิย เอกจฺโจ ปุคฺคโล อวณฺณารหสฺส จ อวณฺณํ ภาสิตา โหติ ภูตํ ตจฺฉํ กาเลน, วณฺณารหสฺส จ วณฺณํ ภาสิตา โหติ ภูตํ ตจฺฉํ กาเลน. อิเม โข โปตลิย จตฺตาโร ปุคฺคลา สนฺโต สํวิชฺชมานา โลกสฺมิํ. อิเมสํ โข โปตลิย จตุนฺนํ ปุคฺคลานํ กตโม เต ปุคฺคโล ขมติ อภิกฺกนฺตตโร จ ปณีตตโร จาติ.}

\prose{จตฺตาโรเม โภ โคตม ปุคฺคลา สนฺโต สํวิชฺชมานา โลกสฺมิํ. กตเม จตฺตาโร? อิธ โภ โคตม เอกจฺโจ ปุคฺคโล อวณฺณารหสฺส อวณฺณํ ภาสิตา โหติ ภูตํ ตจฺฉํ กาเลน, โน จ โข วณฺณารหสฺส วณฺณํ ภาสิตา โหติ ภูตํ ตจฺฉํ กาเลน. อิธ ปน โภ โคตม เอกจฺโจ ปุคฺคโล วณฺณารหสฺส วณฺณํ ภาสิตา โหติ ภูตํ ตจฺฉํ กาเลน, โน จ โข อวณฺณารหสฺส อวณฺณํ ภาสิตา โหติ ภูตํ ตจฺฉํ กาเลน. อิธ ปน โภ โคตม เอกจฺโจ ปุคฺคโล เนว อวณฺณารหสฺส อวณฺณํ ภาสิตา โหติ ภูตํ ตจฺฉํ กาเลน, โน จ วณฺณารหสฺส วณฺณํ ภาสิตา โหติ ภูตํ ตจฺฉํ กาเลน. \csromanfolio{414}อิธ ปน โภ โคตม เอกจฺโจ ปุคฺคโล อวณฺณารหสฺส จ อวณฺณํ ภาสิตา โหติ ภูตํ ตจฺฉํ กาเลน, วณฺณารหสฺส จ วณฺณํ ภาสิตา โหติ ภูตํ ตจฺฉํ กาเลน. อิเม โข โภ โคตม จตฺตาโร ปุคฺคลา สนฺโต สํวิชฺชมานา โลกสฺมิํ. อิเมสํ โภ โคตม จตุนฺนํ ปุคฺคลานํ ยฺวายํ ปุคฺคโล เนว อวณฺณารหสฺส อวณฺณํ ภาสิตา โหติ ภูตํ ตจฺฉํ กาเลน, โน จ วณฺณารหสฺส วณฺณํ ภาสิตา โหติ ภูตํ ตจฺฉํ กาเลน. อยํ เม ปุคฺคโล ขมติ อิเมสํ จตุนฺนํ ปุคฺคลานํ อภิกฺกนฺตตโร จ ปณีตตโร จ. ตํ กิสฺส เหตุ? อภิกฺกนฺตา\csromansharedfootnote{dad4c43acb010824}{อภิกฺกนฺตตรา (ก)} เหสา โภ โคตม ยทิทํ อุเปกฺขาติ.}

\prose{จตฺตาโรเม โปตลิย ปุคฺคลา สนฺโต สํวิชฺชมานา โลกสฺมิํ. กตเม จตฺตาโร ฯเปฯ. อิเม โข โปตลิย จตฺตาโร ปุคฺคลา สนฺโต สํวิชฺชมานา โลกสฺมิํ. อิเมสํ โข โปตลิย จตุนฺนํ ปุคฺคลานํ ยฺวายํ ปุคฺคโล อวณฺณารหสฺส จ อวณฺณํ ภาสิตา โหติ ภูตํ ตจฺฉํ กาเลน, วณฺณารหสฺส จ วณฺณํ ภาสิตา โหติ ภูตํ ตจฺฉํ กาเลน. อยํ อิเมสํ จตุนฺนํ ปุคฺคลานํ อภิกฺกนฺตตโร จ ปณีตตโร จ. ตํ กิสฺส เหตุ? อภิกฺกนฺตา เหสา โปตลิย ยทิทํ ตตฺถ ตตฺถ กาลญฺญุตาติ.}

\prose{จตฺตาโรเม โภ โคตม ปุคฺคลา สนฺโต สํวิชฺชมานา โลกสฺมิํ. กตเม จตฺตาโร ฯเปฯ. อิเม โข โภ โคตม จตฺตาโร ปุคฺคลา สนฺโต สํวิชฺชมานา โลกสฺมิํ. อิเมสํ โภ โคตม จตุนฺนํ ปุคฺคลานํ ยฺวายํ ปุคฺคโล อวณฺณารหสฺส จ อวณฺณํ ภาสิตา ภูตํ ตจฺฉํ กาเลน, วณฺณารหสฺส จ วณฺณํ ภาสิตา ภูตํ ตจฺฉํ กาเลน. อยํ เม ปุคฺคโล ขมติ อิเมสํ จตุนฺนํ ปุคฺคลานํ อภิกฺกนฺตตโร จ ปณีตตโร จ. ตํ กิสฺส เหตุ? อภิกฺกนฺตา เหสา โภ โคตม ยทิทํ ตตฺถ ตตฺถ กาลญฺญุตา.}

\prose{อภิกฺกนฺตํ โภ โคตม อภิกฺกนฺตํ โภ โคตม, เสยฺยถาปิ โภ โคตม นิกฺกุชฺชิตํ วา อุกฺกุชฺเชยฺย, ปฏิจฺฉนฺนํ วา วิวเรยฺย, มูฬฺหสฺส วา มคฺคํ อาจิกฺเขยฺย, อนฺธกาเร วา เตลปชฺโชตํ ธาเรยฺย “จกฺขุมนฺโต รูปานิ ทกฺขนฺตี”ติ. เอวเมวํ โภตา โคตเมน อเนก\-ปริยาเยน ธมฺโม}

\vspace{\tipitakaparskip}
\csromancenter{(10) 5. อสุรวคฺค ปกาสิโต. เอสาหํ ภวนฺตํ โคตมํ สรณํ คจฺฉามิ ธมฺมญฺจ ภิกฺขุสํฆญฺจ, อุปาสกํ มํ ภวํ โคตโม ธาเรตุ อชฺชตคฺเค ปาณุเปตํ สรณํ คตนฺติ.  ทสมํ.}
\nitthitammidruled{อสุรวคฺโค ปญฺจโม.}
\csromancenter{\textbf{ตสฺสุทฺทานํ}}
\setlength{\csromangathapagewidth}{0pt}
\setlength{\csromangathawakwidth}{0pt}
\csromangathameasuregroup{อสุโร ตโย สมาธี, \\ ราโค นิสนฺติ อตฺตหิตํ,}{\csromangathabat{อสุโร ตโย สมาธี,}{ฉวาลาเตน ปญฺจมํ.} \\ \csromangathabat{ราโค นิสนฺติ อตฺตหิตํ,}{สิกฺขา โปตลิเยน จาติ.} \\ \textbf{ทุติโย ปณฺณาสโก สมตฺโต.}}
\csromangathasetleft{อสุโร ตโย สมาธี, \\ ราโค นิสนฺติ อตฺตหิตํ,}
\csromangathagroup{\csromangathastanza{\csromanfolio{415}\csromangathabat{อสุโร ตโย สมาธี,}{ฉวาลาเตน ปญฺจมํ.} \\ \csromangathabat{ราโค นิสนฺติ อตฺตหิตํ,}{สิกฺขา โปตลิเยน จาติ.} \\ \textbf{ทุติโย ปณฺณาสโก สมตฺโต.}}}

\csromanmatikaanchor{mtk.150}
\csromanmatikaanchor{mtk.151}
\csromantocmarknopage{chapter}{3. ตติยปณฺณาสก}{mtk.150}
\bookmark[dest={mtk.150},level=0]{3. ตติยปณฺณาสก}
\csromantocmarknopage{section}{(11) 1. วลาหกวคฺค}{mtk.151}
\bookmark[dest={mtk.151},level=1]{(11) 1. วลาหกวคฺค}
\csromanmarkh{(11)}\csromanmarkhii{1. วลาหกวคฺค}\csromanheadernomark{1}{\textbf{(11) 1. วลาหกวคฺค}}
\csromanheader{2}{1. ปฐมวลาหก\-สุตฺต}
\csromanmatikaanchor{mtk.152}
\csromantocmarkref{subsection}{1-2. วลาหกสุตฺตานิ}{mtk.152}
\bookmark[dest={mtk.152},level=2]{1-2. วลาหกสุตฺตานิ}
\proseitem{101}{\csromanfolio{416}เอวํ เม สุตํ\csromandash{}เอกํ สมยํ ภควา สาวตฺถิยํ วิหรติ เชตวเน อนาถ\-ปิณฺฑิกสฺส อาราเม. ตตฺร โข ภควา ภิกฺขู อามนฺเตสิ “ภิกฺขโว”ติ. “ภทนฺเต”ติ เต ภิกฺขู ภควโต ปจฺจสฺโสสุํ. ภควา เอตทโวจ\csromandash{}}

\prose{จตฺตาโรเม ภิกฺขเว วลาหกา. กตเม จตฺตาโร? คชฺชิตา โน วสฺสิตา, วสฺสิตา โน คชฺชิตา, เนว คชฺชิตา โน วสฺสิตา, คชฺชิตา จ วสฺสิตา จ. อิเม โข ภิกฺขเว จตฺตาโร วลาหกา. เอวเมวํ โข ภิกฺขเว จตฺตาโร วลาหกูปมา\csromansharedfootnote{bc739008f418965b}{อภิ 3. 148 ปิฏฺเฐปิ.} ปุคฺคลา สนฺโต สํวิชฺชมานา โลกสฺมิํ. กตเม จตฺตาโร? คชฺชิตา โน วสฺสิตา, วสฺสิตา โน คชฺชิตา, เนว คชฺชิตา โน วสฺสิตา, คชฺชิตา จ วสฺสิตา จ.}

\prose{กถญฺจ ภิกฺขเว ปุคฺคโล คชฺชิตา โหติ โน วสฺสิตา? อิธ ภิกฺขเว เอกจฺโจ ปุคฺคโล ภาสิตา โหติ โน กตฺตา. เอวํ โข ภิกฺขเว ปุคฺคโล คชฺชิตา โหติ โน วสฺสิตา. เสยฺยถาปิ โส ภิกฺขเว วลาหโก คชฺชิตา โน วสฺสิตา, ตถูปมาหํ ภิกฺขเว อิมํ ปุคฺคลํ วทามิ.}

\prose{กถญฺจ ภิกฺขเว ปุคฺคโล วสฺสิตา โหติ โน คชฺชิตา? อิธ ภิกฺขเว เอกจฺโจ ปุคฺคโล กตฺตา โหติ โน ภาสิตา. เอวํ โข ภิกฺขเว ปุคฺคโล วสฺสิตา โหติ โน คชฺชิตา. เสยฺยถาปิ โส ภิกฺขเว วลาหโก วสฺสิตา โน คชฺชิตา, ตถูปมาหํ ภิกฺขเว อิมํ ปุคฺคลํ วทามิ.}

\prose{กถญฺจ ภิกฺขเว ปุคฺคโล เนว คชฺชิตา โหติ โน วสฺสิตา? อิธ ภิกฺขเว เอกจฺโจ ปุคฺคโล เนว ภาสิตา โหติ โน กตฺตา. เอวํ โข ภิกฺขเว ปุคฺคโล เนว คชฺชิตา โหติ โน วสฺสิตา.}

\vspace{\tipitakaparskip}
\csromancenter{(11) 1. วลาหกวคฺค เสยฺยถาปิ โส ภิกฺขเว วลาหโก เนว คชฺชิตา\csromansharedfootnote{a3b411b30191c307}{เนว คชฺชิตา โหติ (ก)} โน วสฺสิตา, ตถูปมาหํ ภิกฺขเว อิมํ ปุคฺคลํ วทามิ.}
\prose{\csromanfolio{417}กถญฺจ ภิกฺขเว ปุคฺคโล คชฺชิตา จ โหติ วสฺสิตา จ? อิธ ภิกฺขเว เอกจฺโจ ปุคฺคโล ภาสิตา จ โหติ กตฺตา จ. เอวํ โข ภิกฺขเว ปุคฺคโล คชฺชิตา จ โหติ วสฺสิตา จ. เสยฺยถาปิ โส ภิกฺขเว วลาหโก คชฺชิตา จ\csromansharedfootnote{58269fd4aa1ea940}{คชฺชิตา จ โหติ (ก)} วสฺสิตา จ, ตถูปมาหํ ภิกฺขเว อิมํ ปุคฺคลํ วทามิ. อิเม โข ภิกฺขเว จตฺตาโร วลาหกูปมา ปุคฺคลา สนฺโต สํวิชฺชมานา โลกสฺมินฺติ.  ปฐมํ.}

\csromanheader{2}{2. ทุติย\-วลาหก\-สุตฺต}
\proseitem{102}{จตฺตาโรเม ภิกฺขเว วลาหกา. กตเม จตฺตาโร? คชฺชิตา โน วสฺสิตา, วสฺสิตา โน คชฺชิตา, เนว คชฺชิตา โน วสฺสิตา, คชฺชิตา จ วสฺสิตา จ. อิเม โข ภิกฺขเว จตฺตาโร วลาหกา. เอวเมวํ โข ภิกฺขเว จตฺตาโร วลาหกูปมา ปุคฺคลา สนฺโต สํวิชฺชมานา โลกสฺมิํ. กตเม จตฺตาโร? คชฺชิตา โน วสฺสิตา, วสฺสิตา โน คชฺชิตา, เนว คชฺชิตา โน วสฺสิตา, คชฺชิตา จ วสฺสิตา จ.}

\prose{กถญฺจ ภิกฺขเว ปุคฺคโล คชฺชิตา โหติ โน วสฺสิตา? อิธ ภิกฺขเว เอกจฺโจ ปุคฺคโล ธมฺมํ ปริยาปุณาติ สุตฺตํ เคยฺยํ เวยฺยากรณํ คาถํ อุทานํ อิติวุตฺตกํ ชาตกํ อพฺภุตธมฺมํ เวทลฺลํ. โส “อิทํ ทุกฺขนฺ”ติ ยถาภูตํ นปฺปชานาติ, “อยํ ทุกฺขสมุทโย”ติ ยถาภูตํ นปฺปชานาติ, “อยํ ทุกฺขนิโรโธ”ติ ยถาภูตํ นปฺปชานาติ, “อยํ ทุกฺขนิโรธ\-คามินี ปฏิปทา”ติ ยถาภูตํ นปฺปชานาติ. เอวํ โข ภิกฺขเว ปุคฺคโล คชฺชิตา โหติ โน วสฺสิตา. เสยฺยถาปิ โส ภิกฺขเว วลาหโก คชฺชิตา โน วสฺสิตา, ตถูปมาหํ ภิกฺขเว อิมํ ปุคฺคลํ วทามิ.}

\prose{กถญฺจ ภิกฺขเว ปุคฺคโล วสฺสิตา โหติ โน คชฺชิตา? อิธ ภิกฺขเว เอกจฺโจ ปุคฺคโล ธมฺมํ น ปริยาปุณาติ สุตฺตํ เคยฺยํ เวยฺยากรณํ คาถํ อุทานํ อิติวุตฺตกํ ชาตกํ อพฺภุตธมฺมํ เวทลฺลํ. โส “อิทํ ทุกฺขนฺ”ติ ยถาภูตํ ปชานาติ ฯเปฯ “อยํ ทุกฺขนิโรธ\-คามินี \csromanfolio{418}ปฏิปทา”ติ ยถาภูตํ ปชานาติ. เอวํ โข ภิกฺขเว ปุคฺคโล วสฺสิตา โหติ โน คชฺชิตา. เสยฺยถาปิ โส ภิกฺขเว วลาหโก วสฺสิตา โน คชฺชิตา, ตถูปมาหํ ภิกฺขเว อิมํ ปุคฺคลํ วทามิ.}

\prose{กถญฺจ ภิกฺขเว ปุคฺคโล เนว คชฺชิตา โหติ โน วสฺสิตา? อิธ ภิกฺขเว เอกจฺโจ ปุคฺคโล เนว ธมฺมํ ปริยาปุณาติ สุตฺตํ เคยฺยํ เวยฺยากรณํ คาถํ อุทานํ อิติวุตฺตกํ ชาตกํ อพฺภุตธมฺมํ เวทลฺลํ. โส “อิทํ ทุกฺขนฺ”ติ ยถาภูตํ นปฺปชานาติ ฯเปฯ “อยํ ทุกฺขนิโรธ\-คามินี ปฏิปทา”ติ ยถาภูตํ นปฺปชานาติ. เอวํ โข ภิกฺขเว ปุคฺคโล เนว คชฺชิตา โหติ โน วสฺสิตา. เสยฺยถาปิ โส ภิกฺขเว วลาหโก เนว คชฺชิตา โน วสฺสิตา, ตถูปมาหํ ภิกฺขเว อิมํ ปุคฺคลํ วทามิ.}

\prose{กถญฺจ ภิกฺขเว ปุคฺคโล คชฺชิตา จ โหติ วสฺสิตา จ? อิธ ภิกฺขเว เอกจฺโจ ปุคฺคโล ธมฺมํ ปริยาปุณาติ สุตฺตํ เคยฺยํ เวยฺยากรณํ คาถํ อุทานํ อิติวุตฺตกํ ชาตกํ อพฺภุตธมฺมํ เวทลฺลํ. โส “อิทํ ทุกฺขนฺ”ติ ยถาภูตํ ปชานาติ ฯเปฯ “อยํ ทุกฺขนิโรธ\-คามินี ปฏิปทา”ติ ยถาภูตํ ปชานาติ. เอวํ โข ภิกฺขเว ปุคฺคโล คชฺชิตา จ โหติ วสฺสิตา จ. เสยฺยถาปิ โส ภิกฺขเว วลาหโก คชฺชิตา จ วสฺสิตา จ, ตถูปมาหํ ภิกฺขเว อิมํ ปุคฺคลํ วทามิ. อิเม โข ภิกฺขเว จตฺตาโร วลาหกูปมา ปุคฺคลา สนฺโต สํวิชฺชมานา โลกสฺมินฺติ.  ทุติยํ.}

\csromanmatikaanchor{mtk.153}
\csromantocmarkref{subsection}{3. กุมฺภสุตฺต}{mtk.153}
\bookmark[dest={mtk.153},level=2]{3. กุมฺภสุตฺต}
\csromanheader{2}{3. \textbf{กุมฺภสุตฺต}}
\proseitem{103}{จตฺตาโรเม ภิกฺขเว กุมฺภา. กตเม จตฺตาโร? ตุจฺโฉ ปิหิโต, ปูโร วิวโฏ, ตุจฺโฉ วิวโฏ, ปูโร ปิหิโต. อิเม โข ภิกฺขเว จตฺตาโร กุมฺภา. เอวเมวํ โข ภิกฺขเว จตฺตาโร กุมฺภูปมา\csromansharedfootnote{dc0b258e2bc8e788}{อภิ 3. 152 ปิฏฺเฐปิ.} ปุคฺคลา สนฺโต สํวิชฺชมานา โลกสฺมิํ. กตเม จตฺตาโร? ตุจฺโฉ ปิหิโต, ปูโร วิวโฏ, ตุจฺโฉ วิวโฏ, ปูโร ปิหิโต.}

\prose{กถญฺจ ภิกฺขเว ปุคฺคโล ตุจฺโฉ โหติ ปิหิโต? อิธ ภิกฺขเว เอกจฺจสฺส ปุคฺคลสฺส ปาสาทิกํ โหติ อภิกฺกนฺตํ ปฏิกฺกนฺตํ อาโลกิตํ วิโลกิตํ สมิญฺชิตํ ปสาริตํ สงฺฆาฏิ\-ปตฺต\-จีวร\-ธารณํ. โส “อิทํ ทุกฺขนฺ”ติ ยถาภูตํ นปฺปชานาติ “อยํ ทุกฺขนิโรธ\-คามินี ปฏิปทา”ติ ยถาภูตํ นปฺปชานาติ. เอวํ โข ภิกฺขเว ปุคฺคโล ตุจฺโฉ โหติ}

\vspace{\tipitakaparskip}
\csromancenter{(11) 1. วลาหกวคฺค ปิหิโต. เสยฺยถาปิ โส ภิกฺขเว กุมฺโภ ตุจฺโฉ ปิหิโต, ตถูปมาหํ ภิกฺขเว อิมํ ปุคฺคลํ วทามิ.}
\prose{\csromanfolio{419}กถญฺจ ภิกฺขเว ปุคฺคโล ปูโร โหติ วิวโฏ? อิธ ภิกฺขเว เอกจฺจสฺส ปุคฺคลสฺส น ปาสาทิกํ โหติ อภิกฺกนฺตํ ปฏิกฺกนฺตํ อาโลกิตํ วิโลกิตํ สมิญฺชิตํ ปสาริตํ สงฺฆาฏิ\-ปตฺต\-จีวร\-ธารณํ. โส “อิทํ ทุกฺขนฺ”ติ ยถาภูตํ ปชานาติ ฯเปฯ “อยํ ทุกฺขนิโรธ\-คามินี ปฏิปทา”ติ ยถาภูตํ ปชานาติ. เอวํ โข ภิกฺขเว ปุคฺคโล ปูโร โหติ วิวโฏ. เสยฺยถาปิ โส ภิกฺขเว กุมฺโภ ปูโร วิวโฏ, ตถูปมาหํ ภิกฺขเว อิมํ ปุคฺคลํ วทามิ.}

\prose{กถญฺจ ภิกฺขเว ปุคฺคโล ตุจฺโฉ โหติ วิวโฏ? อิธ ภิกฺขเว เอกจฺจสฺส ปุคฺคลสฺส น ปาสาทิกํ โหติ อภิกฺกนฺตํ ปฏิกฺกนฺตํ อาโลกิตํ วิโลกิตํ สมิญฺชิตํ ปสาริตํ สงฺฆาฏิ\-ปตฺต\-จีวร\-ธารณํ. โส “อิทํ ทุกฺขนฺ”ติ ยถาภูตํ นปฺปชานาติ ฯเปฯ “อยํ ทุกฺขนิโรธ\-คามินี ปฏิปทา”ติ ยถาภูตํ นปฺปชานาติ. เอวํ โข ภิกฺขเว ปุคฺคโล ตุจฺโฉ โหติ วิวโฏ. เสยฺยถาปิ โส ภิกฺขเว กุมฺโภ ตุจฺโฉ วิวโฏ, ตถูปมาหํ ภิกฺขเว อิมํ ปุคฺคลํ วทามิ.}

\prose{กถญฺจ ภิกฺขเว ปุคฺคโล ปูโร โหติ ปิหิโต? อิธ ภิกฺขเว เอกจฺจสฺส ปุคฺคลสฺส ปาสาทิกํ โหติ อภิกฺกนฺตํ ปฏิกฺกนฺตํ อาโลกิตํ วิโลกิตํ สมิญฺชิตํ ปสาริตํ สงฺฆาฏิ\-ปตฺต\-จีวร\-ธารณํ. โส “อิทํ ทุกฺขนฺ”ติ ยถาภูตํ ปชานาติ ฯเปฯ “อยํ ทุกฺขนิโรธ\-คามินี ปฏิปทา”ติ ยถาภูตํ ปชานาติ. เอวํ โข ภิกฺขเว ปุคฺคโล ปูโร โหติ ปิหิโต. เสยฺยถาปิ โส ภิกฺขเว กุมฺโภ ปูโร ปิหิโต, ตถูปมาหํ ภิกฺขเว อิมํ ปุคฺคลํ วทามิ. อิเม โข ภิกฺขเว จตฺตาโร กุมฺภูปมา ปุคฺคลา สนฺโต สํวิชฺชมานา โลกสฺมินฺติ.  ตติยํ.}

\csromanmatikaanchor{mtk.154}
\csromantocmarkref{subsection}{4. อุทกรหทสุตฺต}{mtk.154}
\bookmark[dest={mtk.154},level=2]{4. อุทกรหทสุตฺต}
\csromanheader{2}{4. อุทกรหท\-สุตฺต}
\proseitem{104}{จตฺตาโรเม ภิกฺขเว อุทกรหทา. กตเม จตฺตาโร? อุตฺตาโน คมฺภีโรภาโส, คมฺภีโร อุตฺตาโนภาโส, อุตฺตาโน อุตฺตาโนภาโส, คมฺภีโร คมฺภีโรภาโส. อิเม โข ภิกฺขเว จตฺตาโร อุทกรหทา. เอวเมวํ โข ภิกฺขเว จตฺตาโร อุทกรหทู\-ปมา\csromansharedfootnote{a200a4716cf60213}{อภิ 3. 153 ปิฏฺเฐปิ.} ปุคฺคลา สนฺโต สํวิชฺชมานา โลกสฺมิํ. กตเม จตฺตาโร? อุตฺตาโน \csromanfolio{420}คมฺภีโรภาโส, คมฺภีโร อุตฺตาโนภาโส, อุตฺตาโน อุตฺตาโนภาโส, คมฺภีโร คมฺภีโรภาโส.}

\prose{กถญฺจ ภิกฺขเว ปุคฺคโล อุตฺตาโน โหติ คมฺภีโรภาโส? อิธ ภิกฺขเว เอกจฺจสฺส ปุคฺคลสฺส ปาสาทิกํ โหติ อภิกฺกนฺตํ ปฏิกฺกนฺตํ อาโลกิตํ วิโลกิตํ สมิญฺชิตํ ปสาริตํ สงฺฆาฏิ\-ปตฺต\-จีวร\-ธารณํ. โส “อิทํ ทุกฺขนฺ”ติ ยถาภูตํ นปฺปชานาติ ฯเปฯ “อยํ ทุกฺขนิโรธ\-คามินี ปฏิปทา”ติ ยถาภูตํ นปฺปชานาติ. เอวํ โข ภิกฺขเว ปุคฺคโล อุตฺตาโน โหติ คมฺภีโรภาโส. เสยฺยถาปิ โส ภิกฺขเว อุทกรหโท อุตฺตาโน คมฺภีโรภาโส, ตถูปมาหํ ภิกฺขเว อิมํ ปุคฺคลํ วทามิ.}

\prose{กถญฺจ ภิกฺขเว ปุคฺคโล คมฺภีโร โหติ อุตฺตาโนภาโส? อิธ ภิกฺขเว เอกจฺจสฺส ปุคฺคลสฺส น ปาสาทิกํ โหติ อภิกฺกนฺตํ ปฏิกฺกนฺตํ อาโลกิตํ วิโลกิตํ สมิญฺชิตํ ปสาริตํ สงฺฆาฏิ\-ปตฺต\-จีวร\-ธารณํ. โส “อิทํ ทุกฺขนฺ”ติ ยถาภูตํ ปชานาติ ฯเปฯ “อยํ ทุกฺขนิโรธ\-คามินี ปฏิปทา”ติ ยถาภูตํ ปชานาติ. เอวํ โข ภิกฺขเว ปุคฺคโล คมฺภีโร โหติ อุตฺตาโนภาโส. เสยฺยถาปิ โส ภิกฺขเว อุทกรหโท คมฺภีโร อุตฺตาโนภาโส, ตถูปมาหํ ภิกฺขเว อิมํ ปุคฺคลํ วทามิ.}

\prose{กถญฺจ ภิกฺขเว ปุคฺคโล อุตฺตาโน โหติ อุตฺตาโนภาโส? อิธ ภิกฺขเว เอกจฺจสฺส ปุคฺคลสฺส น ปาสาทิกํ โหติ อภิกฺกนฺตํ ปฏิกฺกนฺตํ อาโลกิตํ วิโลกิตํ สมิญฺชิตํ ปสาริตํ สงฺฆาฏิ\-ปตฺต\-จีวร\-ธารณํ. โส “อิทํ ทุกฺขนฺ”ติ ยถาภูตํ นปฺปชานาติ ฯเปฯ “อยํ ทุกฺขนิโรธ\-คามินี ปฏิปทา”ติ ยถาภูตํ นปฺปชานาติ. เอวํ โข ภิกฺขเว ปุคฺคโล อุตฺตาโน โหติ อุตฺตาโนภาโส. เสยฺยถาปิ โส ภิกฺขเว อุทกรหโท อุตฺตาโน อุตฺตาโนภาโส, ตถูปมาหํ ภิกฺขเว อิมํ ปุคฺคลํ วทามิ.}

\prose{กถญฺจ ภิกฺขเว ปุคฺคโล คมฺภีโร โหติ คมฺภีโรภาโส? อิธ ภิกฺขเว เอกจฺจสฺส ปุคฺคลสฺส ปาสาทิกํ โหติ อภิกฺกนฺตํ ปฏิกฺกนฺตํ อาโลกิตํ วิโลกิตํ สมิญฺชิตํ ปสาริตํ สงฺฆาฏิ\-ปตฺต\-จีวร\-ธารณํ. โส “อิทํ ทุกฺขนฺ”ติ ยถาภูตํ ปชานาติ ฯเปฯ “อยํ ทุกฺขนิโรธ\-คามินี ปฏิปทา”ติ ยถาภูตํ ปชานาติ. เอวํ โข ภิกฺขเว ปุคฺคโล คมฺภีโร โหติ คมฺภีโรภาโส. เสยฺยถาปิ โส ภิกฺขเว อุทกรหโท คมฺภีโร คมฺภีโรภาโส, ตถูปมาหํ ภิกฺขเว อิมํ ปุคฺคลํ วทามิ. อิเม}

\vspace{\tipitakaparskip}
\csromancenter{(11) 1. วลาหกวคฺค โข ภิกฺขเว จตฺตาโร อุทกรหทู\-ปมา ปุคฺคลา สนฺโต สํวิชฺชมานา โลกสฺมินฺติ.  จตุตฺถํ.}
\csromanmatikaanchor{mtk.155}
\csromantocmarkref{subsection}{5. อมฺพสุตฺต}{mtk.155}
\bookmark[dest={mtk.155},level=2]{5. อมฺพสุตฺต}
\csromanheader{2}{5. \textbf{อมฺพสุตฺต}}
\proseitem{105}{\csromanfolio{421}จตฺตาริมานิ ภิกฺขเว อมฺพานิ. กตมานิ จตฺตาริ? อามํ ปกฺกวณฺณิ\csromansharedfootnote{81754a7bc198c24e}{ปกฺกวณฺณี (อภิ 3. 151 ปิฏฺเฐปิ.)}, ปกฺกํ อามวณฺณิ\csromansharedfootnote{9a4ec21ba47d3a78}{อามวณฺณี (สี, สฺยา, กํ, อิ)}, อามํ อามวณฺณิ, ปกฺกํ ปกฺกวณฺณิ. อิมานิ โข ภิกฺขเว จตฺตาริ อมฺพานิ. เอวเมวํ โข ภิกฺขเว จตฺตาโร อมฺพูปมา\csromansharedfootnote{81754a7bc198c24e}{ปกฺกวณฺณี (อภิ 3. 151 ปิฏฺเฐปิ.)} ปุคฺคลา สนฺโต สํวิชฺชมานา โลกสฺมิํ. กตเม จตฺตาโร, อาโม ปกฺกวณฺณี, ปกฺโก อามวณฺณี, อาโม อามวณฺณี, ปกฺโก ปกฺกวณฺณี.}

\prose{กถญฺจ ภิกฺขเว ปุคฺคโล อาโม โหติ ปกฺกวณฺณี? อิธ ภิกฺขเว เอกจฺจสฺส ปุคฺคลสฺส ปาสาทิกํ โหติ อภิกฺกนฺตํ ปฏิกฺกนฺตํ อาโลกิตํ วิโลกิตํ สมิญฺชิตํ ปสาริตํ สงฺฆาฏิ\-ปตฺต\-จีวร\-ธารณํ. โส “อิทํ ทุกฺขนฺ”ติ ยถาภูตํ นปฺปชานาติ ฯเปฯ “อยํ ทุกฺขนิโรธ\-คามินี ปฏิปทา”ติ ยถาภูตํ นปฺปชานาติ. เอวํ โข ภิกฺขเว ปุคฺคโล อาโม โหติ ปกฺกวณฺณี. เสยฺยถาปิ ตํ ภิกฺขเว อมฺพํ อามํ ปกฺกวณฺณิ, ตถูปมาหํ ภิกฺขเว อิมํ ปุคฺคลํ วทามิ.}

\prose{กถญฺจ ภิกฺขเว ปุคฺคโล ปกฺโก โหติ อามวณฺณี? อิธ ภิกฺขเว เอกจฺจสฺส ปุคฺคลสฺส น ปาสาทิกํ โหติ อภิกฺกนฺตํ ปฏิกฺกนฺตํ อาโลกิตํ วิโลกิตํ สมิญฺชิตํ ปสาริตํ สงฺฆาฏิ\-ปตฺต\-จีวร\-ธารณํ. โส “อิทํ ทุกฺขนฺ”ติ ยถาภูตํ ปชานาติ ฯเปฯ “อยํ ทุกฺขนิโรธ\-คามินี ปฏิปทา”ติ ยถาภูตํ ปชานาติ. เอวํ โข ภิกฺขเว ปุคฺคโล ปกฺโก โหติ อามวณฺณี. เสยฺยถาปิ ตํ ภิกฺขเว อมฺพํ ปกฺกํ อามวณฺณิ, ตถูปมาหํ ภิกฺขเว อิมํ ปุคฺคลํ วทามิ.}

\prose{กถญฺจ ภิกฺขเว ปุคฺคโล อาโม โหติ อามวณฺณี? อิธ ภิกฺขเว เอกจฺจสฺส ปุคฺคลสฺส น ปาสาทิกํ โหติ อภิกฺกนฺตํ ปฏิกฺกนฺตํ อาโลกิตํ วิโลกิตํ สมิญฺชิตํ ปสาริตํ สงฺฆาฏิ\-ปตฺต\-จีวร\-ธารณํ. โส “อิทํ ทุกฺขนฺ”ติ ยถาภูตํ นปฺปชานาติ ฯเปฯ “อยํ ทุกฺขนิโรธ\-คามินี ปฏิปทา”ติ ยถาภูตํ นปฺปชานาติ. เอวํ โข ภิกฺขเว ปุคฺคโล อาโม โหติ \csromanfolio{422}อามวณฺณี. เสยฺยถาปิ ตํ ภิกฺขเว อมฺพํ อามํ อามวณฺณิ, ตถูปมาหํ ภิกฺขเว อิมํ ปุคฺคลํ วทามิ.}

\prose{กถญฺจ ภิกฺขเว ปุคฺคโล ปกฺโก โหติ ปกฺกวณฺณี? อิธ ภิกฺขเว เอกจฺจสฺส ปุคฺคลสฺส ปาสาทิกํ โหติ อภิกฺกนฺตํ ปฏิกฺกนฺตํ อาโลกิตํ วิโลกิตํ สมิญฺชิตํ ปสาริตํ สงฺฆาฏิ\-ปตฺต\-จีวร\-ธารณํ. โส “อิทํ ทุกฺขนฺ”ติ ยถาภูตํ ปชานาติ ฯเปฯ “อยํ ทุกฺขนิโรธ\-คามินี ปฏิปทา”ติ ยถาภูตํ ปชานาติ. เอวํ โข ภิกฺขเว ปุคฺคโล ปกฺโก โหติ ปกฺกวณฺณี. เสยฺยถาปิ ตํ ภิกฺขเว อมฺพํ ปกฺกํ ปกฺกวณฺณิ, ตถูปมาหํ ภิกฺขเว อิมํ ปุคฺคลํ วทามิ. อิเม โข ภิกฺขเว จตฺตาโร อมฺพูปมา ปุคฺคลา สนฺโต สํวิชฺชมานา โลกสฺมินฺติ.  ปญฺจมํ.}

\vspace{\tipitakaparskip}
\csromancenter{( )\csromansharedfootnote{12fa6ae2a5f557b2}{( ) “ฉฏฺฐํ อุตฺตานตฺถเมวา”ติ อฏฺฐกถายํ ทสฺสิตํ, ปาฬิโปตฺถเกสุ ปน กตฺถจิปิ น ทิสฺสติ. “...อามํ ปกฺโกภาสํ, ปกฺกํ อาโมภาสนฺ”ติอาทินา ปาโฐ ภเวยฺย.}}
\csromanmatikaanchor{mtk.156}
\csromantocmarkref{subsection}{7. มูสิกสุตฺต}{mtk.156}
\bookmark[dest={mtk.156},level=2]{7. มูสิกสุตฺต}
\csromanheader{2}{7. \textbf{มูสิกสุตฺต}}
\proseitem{107}{จตสฺโส อิมา ภิกฺขเว มูสิกา. กตมา จตสฺโส? คาธํ กตฺตา โน วสิตา, วสิตา โน คาธํ กตฺตา, เนว คาธํ กตฺตา โน วสิตา, คาธํ กตฺตา จ วสิตา จ. อิมา โข ภิกฺขเว จตสฺโส มูสิกา. เอวเมวํ โข ภิกฺขเว จตฺตาโร มูสิกูปมา ปุคฺคลา สนฺโต สํวิชฺชมานา โลกสฺมิํ. กตเม จตฺตาโร? คาธํ กตฺตา โน วสิตา, วสิตา โน คาธํ กตฺตา, เนว คาธํ กตฺตา โน วสิตา, คาธํ กตฺตา จ วสิตา จ.}

\prose{กถญฺจ ภิกฺขเว ปุคฺคโล คาธํ กตฺตา โหติ โน วสิตา? อิธ ภิกฺขเว เอกจฺโจ ปุคฺคโล ธมฺมํ ปริยาปุณาติ สุตฺตํ เคยฺยํ เวยฺยากรณํ คาถํ อุทานํ อิติวุตฺตกํ ชาตกํ อพฺภุตธมฺมํ เวทลฺลํ. โส “อิทํ ทุกฺขนฺ”ติ ยถาภูตํ นปฺปชานาติ ฯเปฯ “อยํ ทุกฺขนิโรธ\-คามินี ปฏิปทา”ติ ยถาภูตํ นปฺปชานาติ. เอวํ โข ภิกฺขเว ปุคฺคโล คาธํ กตฺตา โหติ โน วสิตา. เสยฺยถาปิ สา ภิกฺขเว มูสิกา คาธํ กตฺตา โน วสิตา, ตถูปมาหํ ภิกฺขเว อิมํ ปุคฺคลํ วทามิ.}

\prose{กถญฺจ ภิกฺขเว ปุคฺคโล วสิตา โหติ โน คาธํ กตฺตา? อิธ ภิกฺขเว เอกจฺโจ ปุคฺคโล ธมฺมํ น ปริยาปุณาติ สุตฺตํ เคยฺยํ เวยฺยากรณํ}

\vspace{\tipitakaparskip}
\csromancenter{(11) 1. วลาหกวคฺค คาถํ อุทานํ อิติวุตฺตกํ ชาตกํ อพฺภุตธมฺมํ เวทลฺลํ. โส “อิทํ ทุกฺขนฺ”ติ ยถาภูตํ ปชานาติ ฯเปฯ “อยํ ทุกฺขนิโรธ\-คามินี ปฏิปทา”ติ ยถาภูตํ ปชานาติ. เอวํ โข ภิกฺขเว ปุคฺคโล วสิตา โหติ โน คาธํ กตฺตา. เสยฺยถาปิ สา ภิกฺขเว มูสิกา วสิตา โหติ โน คาธํ กตฺตา, ตถูปมาหํ ภิกฺขเว อิมํ ปุคฺคลํ วทามิ.}
\prose{\csromanfolio{423}กถญฺจ ภิกฺขเว ปุคฺคโล เนว คาธํ กตฺตา โหติ โน วสิตา? อิธ ภิกฺขเว เอกจฺโจ ปุคฺคโล ธมฺมํ น ปริยาปุณาติ สุตฺตํ เคยฺยํ เวยฺยากรณํ คาถํ อุทานํ อิติวุตฺตกํ ชาตกํ อพฺภุตธมฺมํ เวทลฺลํ. โส “อิทํ ทุกฺขนฺ”ติ ยถาภูตํ นปฺปชานาติ ฯเปฯ “อยํ ทุกฺขนิโรธ\-คามินี ปฏิปทา”ติ ยถาภูตํ นปฺปชานาติ. เอวํ โข ภิกฺขเว ปุคฺคโล เนว คาธํ กตฺตา โหติ โน วสิตา. เสยฺยถาปิ สา ภิกฺขเว มูสิกา เนว คาธํ กตฺตา โหติ โน วสิตา, ตถูปมาหํ ภิกฺขเว อิมํ ปุคฺคลํ วทามิ.}

\prose{กถญฺจ ภิกฺขเว ปุคฺคโล คาธํ กตฺตา จ โหติ วสิตา จ? อิธ ภิกฺขเว เอกจฺโจ ปุคฺคโล ธมฺมํ ปริยาปุณาติ สุตฺตํ เคยฺยํ เวยฺยากรณํ คาถํ อุทานํ อิติวุตฺตกํ ชาตกํ อพฺภุตธมฺมํ เวทลฺลํ. โส “อิทํ ทุกฺขนฺ”ติ ยถาภูตํ ปชานาติ ฯเปฯ “อยํ ทุกฺขนิโรธ\-คามินี ปฏิปทา”ติ ยถาภูตํ ปชานาติ. เอวํ โข ภิกฺขเว ปุคฺคโล คาธํ กตฺตา จ โหติ วสิตา จ. เสยฺยถาปิ สา ภิกฺขเว มูสิกา คาธํ กตฺตา จ โหติ วสิตา จ, ตถูปมาหํ ภิกฺขเว อิมํ ปุคฺคลํ วทามิ. อิเม โข ภิกฺขเว จตฺตาโร มูสิกูปมา ปุคฺคลา สนฺโต สํวิชฺชมานา โลกสฺมินฺติ.  สตฺตมํ.}

\csromanmatikaanchor{mtk.157}
\csromantocmarkref{subsection}{8. พลีพทฺทสุตฺต}{mtk.157}
\bookmark[dest={mtk.157},level=2]{8. พลีพทฺทสุตฺต}
\csromanheader{2}{8. \textbf{พลีพทฺทสุตฺต}}
\proseitem{108}{จตฺตาโรเม ภิกฺขเว พลีพทฺทา\csromansharedfootnote{114ac3dfcf8b9cde}{พลิวทฺทา (สี, สฺยา, กํ, อิ), พลิพทฺธา (ก) อภิ 3. 154 ปิฏฺเฐปิ.}. กตเม จตฺตาโร? สควจณฺโฑ โน ปรควจณฺโฑ, ปรควจณฺโฑ โน สควจณฺโฑ, สควจณฺโฑ จ ปรควจณฺโฑ จ, เนว สควจณฺโฑ โน ปรควจณฺโฑ. อิเม โข ภิกฺขเว จตฺตาโร พลีพทฺทา. เอวเมวํ โข ภิกฺขเว จตฺตาโร พลีพทฺทูปมา ปุคฺคลา สนฺโต สํวิชฺชมานา โลกสฺมิํ. กตเม จตฺตาโร? สควจณฺโฑ โน ปรควจณฺโฑ, ปรควจณฺโฑ โน สควจณฺโฑ, สควจณฺโฑ จ ปรควจณฺโฑ จ, เนว สควจณฺโฑ โน ปรควจณฺโฑ.}

\prose{\csromanfolio{424}กถญฺจ ภิกฺขเว ปุคฺคโล สควจณฺโฑ โหติ โน ปรควจณฺโฑ? อิธ ภิกฺขเว เอกจฺโจ ปุคฺคโล สกปริสํ อุพฺเพเชตา โหติ โน ปรปริสํ. เอวํ โข ภิกฺขเว ปุคฺคโล สควจณฺโฑ โหติ โน ปรควจณฺโฑ. เสยฺยถาปิ โส ภิกฺขเว พลีพทฺโท สควจณฺโฑ โน ปรควจณฺโฑ, ถตูปมาหํ ภิกฺขเว อิมํ ปุคฺคลํ วทามิ.}

\prose{กถญฺจ ภิกฺขเว ปุคฺคโล ปรควจณฺโฑ โหติ โน สควจณฺโฑ? อิธ ภิกฺขเว เอกจฺโจ ปุคฺคโล ปรปริสํ อุพฺเพเชตา โหติ โน สกปริสํ. เอวํ โข ภิกฺขเว ปุคฺคโล ปรควจณฺโฑ โหติ โน สควจณฺโฑ. เสยฺยถาปิ โส ภิกฺขเว พลีพทฺโท ปรควจณฺโฑ โน สควจณฺโฑ, ตถูปมาหํ ภิกฺขเว อิมํ ปุคฺคลํ วทามิ.}

\prose{กถญฺจ ภิกฺขเว ปุคฺคโล สควจณฺโฑ จ โหติ ปรควจณฺโฑ จ? อิธ ภิกฺขเว เอกจฺโจ ปุคฺคโล สกปริสํ อุพฺเพเชตา โหติ ปรปริสญฺจ. เอวํ โข ภิกฺขเว ปุคฺคโล สควจณฺโฑ จ โหติ ปรควจณฺโฑ จ. เสยฺยถาปิ โส ภิกฺขเว พลีพทฺโท สควจณฺโฑ จ ปรควจณฺโฑ จ, ตถูปมาหํ ภิกฺขเว อิมํ ปุคฺคลํ วทามิ.}

\prose{กถญฺจ ภิกฺขเว ปุคฺคโล เนว สควจณฺโฑ โหติ โน ปรควจณฺโฑ? อิธ ภิกฺขเว เอกจฺโจ ปุคฺคโล เนว สกปริสํ อุพฺเพเชตา โหติ โน ปรปริสญฺจ. เอวํ โข ภิกฺขเว ปุคฺคโล เนว สควจณฺโฑ โหติ โน ปรควจณฺโฑ. เสยฺยถาปิ โส ภิกฺขเว พลีพทฺโท เนว สควจณฺโฑ โน ปรควจณฺโฑ, ตถูปมาหํ ภิกฺขเว อิมํ ปุคฺคลํ วทามิ. อิเม โข ภิกฺขเว จตฺตาโร พลีพทฺทูปมา ปุคฺคลา สนฺโต สํวิชฺชมานา โลกสฺมินฺติ.  อฏฺฐมํ.}

\csromanmatikaanchor{mtk.158}
\csromantocmarkref{subsection}{9. รุกฺขสุตฺต}{mtk.158}
\bookmark[dest={mtk.158},level=2]{9. รุกฺขสุตฺต}
\csromanheader{2}{9. \textbf{รุกฺขสุตฺต}}
\proseitem{109}{จตฺตาโรเม ภิกฺขเว รุกฺขา. กตเม จตฺตาโร? เผคฺคุ เผคฺคุปริวาโร, เผคฺคุ สารปริวาโร, สาโร เผคฺคุปริวาโร, สาโร สารปริวาโร. อิเม โข ภิกฺขเว จตฺตาโร รุกฺขา. เอวเมวํ โข ภิกฺขเว จตฺตาโร รุกฺขูปมา ปุคฺคลา สนฺโต สํวิชฺชมานา โลกสฺมิํ. กตเม จตฺตาโร? เผคฺคุ เผคฺคุปริวาโร, เผคฺคุ สารปริวาโร, สาโร เผคฺคุปริวาโร, สาโร สารปริวาโร.}

\chapterheadpageread{(11) 1. วลาหกวคฺค}
\prose{\csromanfolio{425}กถญฺจ ภิกฺขเว ปุคฺคโล เผคฺคุ โหติ เผคฺคุปริวาโร? อิธ ภิกฺขเว เอกจฺโจ ปุคฺคโล ทุสฺสีโล โหติ ปาปธมฺโม, ปริสาปิสฺส โหติ ทุสฺสีลา ปาปธมฺมา. เอวํ โข ภิกฺขเว ปุคฺคโล เผคฺคุ โหติ เผคฺคุปริวาโร. เสยฺยถาปิ โส ภิกฺขเว รุกฺโข เผคฺคุ เผคฺคุปริวาโร, ตถูปมาหํ ภิกฺขเว อิมํ ปุคฺคลํ วทามิ.}

\prose{กถญฺจ ภิกฺขเว ปุคฺคโล เผคฺคุ โหติ สารปริวาโร? อิธ ภิกฺขเว เอกจฺโจ ปุคฺคโล ทุสฺสีโล โหติ ปาปธมฺโม, ปริสา จ ขฺวสฺส โหติ สีลวตี กลฺยาณธมฺมา. เอวํ โข ภิกฺขเว ปุคฺคโล เผคฺคุ โหติ สารปริวาโร. เสยฺยถาปิ โส ภิกฺขเว รุกฺโข เผคฺคุ สารปริวาโร, ตถูปมาหํ ภิกฺขเว อิมํ ปุคฺคลํ วทามิ.}

\prose{กถญฺจ ภิกฺขเว ปุคฺคโล สาโร โหติ เผคฺคุปริวาโร? อิธ ภิกฺขเว เอกจฺโจ ปุคฺคโล สีลวา โหติ กลฺยาณธมฺโม, ปริสา จ ขฺวสฺส โหติ ทุสฺสีลา ปาปธมฺมา. เอวํ โข ภิกฺขเว ปุคฺคโล สาโร โหติ เผคฺคุปริวาโร. เสยฺยถาปิ โส ภิกฺขเว รุกฺโข สาโร เผคฺคุปริวาโร, ตถูปมาหํ ภิกฺขเว อิมํ ปุคฺคลํ วทามิ.}

\prose{กถญฺจ ภิกฺขเว ปุคฺคโล สาโร โหติ สารปริวาโร? อิธ ภิกฺขเว เอกจฺโจ ปุคฺคโล สีลวา โหติ กลฺยาณธมฺโม, ปริสาปิสฺส โหติ สีลวตี กลฺยาณธมฺมา. เอวํ โข ภิกฺขเว ปุคฺคโล สาโร โหติ สารปริวาโร. เสยฺยถาปิ โส ภิกฺขเว รุกฺโข สาโร สารปริวาโร, ตถูปมาหํ ภิกฺขเว อิมํ ปุคฺคลํ วทามิ. อิเม โข ภิกฺขเว จตฺตาโร รุกฺขูปมา ปุคฺคลา สนฺโต สํวิชฺชมานา โลกสฺมินฺติ\csromansharedfootnote{f0389be4ba5240f8}{อภิ 3. 161 ปิฏฺเฐปิ.}.  นวมํ.}

\csromanmatikaanchor{mtk.159}
\csromantocmarkref{subsection}{10. อาสีวิสสุตฺต}{mtk.159}
\bookmark[dest={mtk.159},level=2]{10. อาสีวิสสุตฺต}
\csromanheader{2}{10. \textbf{อาสีวิสสุตฺต}}
\proseitem{110}{จตฺตาโรเม ภิกฺขเว อาสีวิสา\csromansharedfootnote{e246100a92553261}{อาสิวิสา (ก) อภิ 3. 155 ปิฏฺเฐปิ.}. กตเม จตฺตาโร? อาคตวิโส น โฆรวิโส, โฆรวิโส น อาคตวิโส, อาคตวิโส จ โฆรวิโส จ, เนวาคตวิโส น โฆรวิโส. อิเม โข ภิกฺขเว จตฺตาโร อาสีวิสา. เอวเมวํ โข ภิกฺขเว จตฺตาโร อาสีวิสูปมา ปุคฺคลา สนฺโต สํวิชฺชมานา โลกสฺมิํ. กตเม จตฺตาโร? อาคตวิโส น โฆรวิโส, โฆรวิโส น \csromanfolio{426}อาคตวิโส, อาคตวิโส จ โฆรวิโส จ, เนวาคตวิโส น โฆรวิโส.}

\prose{กถญฺจ ภิกฺขเว ปุคฺคโล อาคตวิโส โหติ น โฆรวิโส? อิธ ภิกฺขเว เอกจฺโจ ปุคฺคโล อภิณฺหํ กุชฺฌติ, โส จ ขฺวสฺส โกโธ\csromansharedfootnote{c1ce1a4328f60d17}{โกโป (ก) อภิ 3. 155 ปิฏฺเฐปิ.} น ทีฆรตฺตํ อนุเสติ. เอวํ โข ภิกฺขเว ปุคฺคโล อาคตวิโส โหติ น โฆรวิโส. เสยฺยถาปิ โส ภิกฺขเว อาสีวิโส อาคตวิโส น โฆรวิโส, ตถูปมาหํ ภิกฺขเว อิมํ ปุคฺคลํ วทามิ.}

\prose{กถญฺจ ภิกฺขเว ปุคฺคโล โฆรวิโส โหติ น อาคตวิโส? อิธ ภิกฺขเว เอกจฺโจ ปุคฺคโล น เหว โข อภิณฺหํ กุชฺฌติ, โส จ ขฺวสฺส โกโธ ทีฆรตฺตํ อนุเสติ. เอวํ โข ภิกฺขเว ปุคฺคโล โฆรวิโส โหติ น อาคตวิโส. เสยฺยถาปิ โส ภิกฺขเว อาสีวิโส โฆรวิโส น อาคตวิโส, ตถูปมาหํ ภิกฺขเว อิมํ ปุคฺคลํ วทามิ.}

\prose{กถญฺจ ภิกฺขเว ปุคฺคโล อาคตวิโส จ โหติ โฆรวิโส จ? อิธ ภิกฺขเว เอกจฺโจ ปุคฺคโล อภิณฺหํ กุชฺฌติ, โส จ ขฺวสฺส โกโธ ทีฆรตฺตํ อนุเสติ. เอวํ โข ภิกฺขเว ปุคฺคโล อาคตวิโส จ โหติ โฆรวิโส จ. เสยฺยถาปิ โส ภิกฺขเว อาสีวิโส อาคตวิโส จ โฆรวิโส จ, ตถูปมาหํ ภิกฺขเว อิมํ ปุคฺคลํ วทามิ.}

\prosewithclosermidruled{กถญฺจ ภิกฺขเว ปุคฺคโล เนวาคตวิโส โหติ น โฆรวิโส? อิธ ภิกฺขเว เอกจฺโจ ปุคฺคโล น เหว โข อภิณฺหํ กุชฺฌติ, โส จ ขฺวสฺส โกโธ น ทีฆรตฺตํ อนุเสติ. เอวํ โข ภิกฺขเว ปุคฺคโล เนวาคตวิโส โหติ น โฆรวิโส. เสยฺยถาปิ โส ภิกฺขเว อาสีวิโส เนวาคตวิโส น โฆรวิโส, ตถูปมาหํ ภิกฺขเว อิมํ ปุคฺคลํ วทามิ. อิเม โข ภิกฺขเว จตฺตาโร อาสีวิสูปมา ปุคฺคลา สนฺโต สํวิชฺชมานา โลกสฺมินฺติ.  ทสมํ.}{วลาหกวคฺโค ปฐโม.}
\csromancenter{\textbf{ตสฺสุทฺทานํ}}
\setlength{\csromangathapagewidth}{0pt}
\setlength{\csromangathawakwidth}{0pt}
\csromangathameasuregroup{เทฺว วลาหา กุมฺภ อุทก, \\ มูสิกา พลีพทฺธา รุกฺขา,}{\csromangathabat{เทฺว วลาหา กุมฺภ อุทก,}{รหทา เทฺว โหนฺติ อมฺพานิ.} \\ \csromangathabat{มูสิกา พลีพทฺธา รุกฺขา,}{อาสีวิเสน เต ทสาติ.}}
\csromangathasetleft{เทฺว วลาหา กุมฺภ อุทก, \\ มูสิกา พลีพทฺธา รุกฺขา,}
\csromangathagroup{\csromangathastanza{\csromangathabat{เทฺว วลาหา กุมฺภ อุทก,}{รหทา เทฺว โหนฺติ อมฺพานิ.} \\ \csromangathabat{มูสิกา พลีพทฺธา รุกฺขา,}{อาสีวิเสน เต ทสาติ.}}}

\chapterheadpageread{(12) 2. เกสิวคฺค \textbf{(12) 2. เกสิวคฺค}}
\csromanmatikaanchor{mtk.160}
\csromanmatikaanchor{mtk.161}
\csromantocmarknopage{section}{(12) 2. เกสิวคฺค}{mtk.160}
\bookmark[dest={mtk.160},level=1]{(12) 2. เกสิวคฺค}
\csromantocmarkref{subsection}{1. เกสิสุตฺต}{mtk.161}
\bookmark[dest={mtk.161},level=2]{1. เกสิสุตฺต}
\csromanheader{2}{1. \textbf{เกสิสุตฺต}}
\proseitem{111}{\csromanfolio{427}อถ โข เกสิ อสฺส\-ทมฺม\-สารถิ เยน ภควา เตนุปสงฺกมิ, อุปสงฺกมิตฺวา ภควนฺตํ อภิวาเทตฺวา เอกมนฺตํ นิสีทิ, เอกมนฺตํ นิสินฺนํ โข เกสิํ อสฺส\-ทมฺม\-สารถิํ ภควา เอตทโวจ “ตฺวํ โขสิ เกสิ ปญฺญาโต ‘อสฺส\-ทมฺม\-สารถี’ติ\csromansharedfootnote{9c79ef61576aa44c}{สญฺญาโต อสฺสทมฺมสารถิ (สี, สฺยา, กํ, อิ)}, กถํ ปน ตฺวํ เกสิ อสฺสทมฺมํ วิเนสี”ติ. อหํ โข ภนฺเต อสฺสทมฺมํ สเณฺหนปิ วิเนมิ, ผรุเสนปิ วิเนมิ, สณฺห\-ผรุเสนปิ วิเนมีติ. สเจ เต เกสิ อสฺสทมฺโม สเณฺหน วินยํ น อุเปติ, ผรุเสน วินยํ น อุเปติ, สณฺหผรุเสน วินยํ น อุเปติ, กินฺติ นํ กโรสีติ. สเจ เม ภนฺเต อสฺสทมฺโม สเณฺหน วินยํ น อุเปติ, ผรุเสน วินยํ น อุเปติ, สณฺหผรุเสน วินยํ น อุเปติ, หนามิ นํ ภนฺเต. ตํ กิสฺส เหตุ? “มา เม อาจริยกุลสฺส อวณฺโณ อโหสี”ติ.}

\prose{ภควา ปน ภนฺเต อนุตฺตโร ปุริส\-ทมฺม\-สารถิ, กถํ ปน ภนฺเต ภควา ปุริสทมฺมํ วิเนตีติ. อหํ โข เกสิ ปุริสทมฺมํ สเณฺหนปิ วิเนมิ, ผรุเสนปิ วิเนมิ, สณฺห\-ผรุเสนปิ วิเนมิ. ตตฺริทํ เกสิ สณฺหสฺมิํ\csromandash{}อิติ กายสุจริตํ อิติ กาย\-สุจริตสฺส วิปาโก, อิติ วจีสุจริตํ อิติ วจีสุจริตสฺส วิปาโก, อิติ มโนสุจริตํ อิติ มโน\-สุจริตสฺส วิปาโก, อิติ เทวา อิติ มนุสฺสาติ. ตตฺริทํ เกสิ ผรุสสฺมิํ\csromandash{}อิติ กายทุจฺจริตํ อิติ กาย\-ทุจฺจริตสฺส วิปาโก, อิติ วจีทุจฺจริตํ อิติ วจีทุจฺจริตสฺส วิปาโก, อิติ มโนทุจฺจริตํ อิติ มโน\-ทุจฺจริตสฺส วิปาโก, อิติ นิรโย อิติ ติรจฺฉานโยนิ อิติ เปตฺติวิสโยติ.}

\prose{ตตฺริทํ เกสิ สณฺห\-ผรุสสฺมิํ\csromandash{}อิติ กายสุจริตํ อิติ กาย\-สุจริตสฺส วิปาโก, อิติ กายทุจฺจริตํ อิติ กาย\-ทุจฺจริตสฺส วิปาโก, อิติ วจีสุจริตํ อิติ วจีสุจริตสฺส วิปาโก, อิติ วจีทุจฺจริตํ อิติ วจีทุจฺจริตสฺส วิปาโก, อิติ มโนสุจริตํ อิติ มโน\-สุจริตสฺส วิปาโก, อิติ มโนทุจฺจริตํ อิติ มโน\-ทุจฺจริตสฺส วิปาโก, อิติ เทวา อิติ มนุสฺสา อิติ นิรโย อิติ ติรจฺฉานโยนิ อิติ เปตฺติวิสโยติ.}

\prose{\csromanfolio{428}สเจ เต ภนฺเต ปุริสทมฺโม สเณฺหน วินยํ น อุเปติ, ผรุเสน วินยํ น อุเปติ, สณฺหผรุเสน วินยํ น อุเปติ, กินฺติ นํ ภควา กโรตีติ. สเจ เม เกสิ ปุริสทมฺโม สเณฺหน วินยํ น อุเปติ, ผรุเสน วินยํ น อุเปติ, สณฺหผรุเสน วินยํ น อุเปติ, หนามิ นํ เกสีติ. น โข ภนฺเต ภควโต ปาณาติปาโต กปฺปติ, อถ จ ปน ภควา เอวมาห “หนามิ นํ เกสี”ติ. สจฺจํ เกสิ, น ตถาคตสฺส ปาณาติปาโต กปฺปติ, อปิ จ โย ปุริสทมฺโม สเณฺหน วินยํ น อุเปติ, ผรุเสน วินยํ น อุเปติ, สณฺหผรุเสน วินยํ น อุเปติ. น ตํ ตถาคโต วตฺตพฺพํ อนุสาสิตพฺพํ มญฺญติ, นาปิ วิญฺญู สพฺรหฺมจารี วตฺตพฺพํ อนุสาสิตพฺพํ มญฺญนฺติ, วโธ เหโส เกสิ อริยสฺส วินเย, ยํ น ตถาคโต วตฺตพฺพํ อนุสาสิตพฺพํ มญฺญติ, นาปิ วิญฺญู สพฺรหฺมจารี วตฺตพฺพํ อนุสาสิตพฺพํ มญฺญนฺตีติ.}

\prose{โส หิ นูน ภนฺเต สุหโต โหติ ยํ น ตถาคโต วตฺตพฺพํ อนุสาสิตพฺพํ มญฺญติ, นาปิ วิญฺญู สพฺรหฺมจารี วตฺตพฺพํ อนุสาสิตพฺพํ มญฺญนฺตีติ. อภิกฺกนฺตํ ภนฺเต อภิกฺกนฺตํ ภนฺเต ฯเปฯ อุปาสกํ มํ ภนฺเต ภควา ธาเรตุ อชฺชตคฺเค ปาณุเปตํ สรณํ คตนฺติ.  ปฏฺฐมํ.}

\csromanmatikaanchor{mtk.162}
\csromantocmarkref{subsection}{2. ชวสุตฺต}{mtk.162}
\bookmark[dest={mtk.162},level=2]{2. ชวสุตฺต}
\csromanheader{2}{2. \textbf{ชวสุตฺต}}
\proseitem{112}{จตูหิ ภิกฺขเว องฺเคหิ สมนฺนาคโต รญฺโญ ภโทฺร อสฺสาชานีโย ราชารโห โหติ ราชโภคฺโค, รญฺโญ องฺคนฺเตว สงฺขํ คจฺฉติ. กตเมหิ จตูหิ? อชฺชเวน ชเวน ขนฺติยา โสรจฺเจน. อิเมหิ โข ภิกฺขเว จตูหิ องฺเคหิ สมนฺนาคโต รญฺโญ ภโทฺร อสฺสาชานีโย ราชารโห โหติ ราชโภคฺโค, รญฺโญ องฺคนฺเตว สงฺขํ คจฺฉติ.}

\prose{เอวเมวํ โข ภิกฺขเว จตูหิ ธมฺเมหิ สมนฺนาคโต ภิกฺขุ อาหุเนยฺโย โหติ ฯเปฯ อนุตฺตรํ ปุญฺญกฺเขตฺตํ โลกสฺส. กตเมหิ จตูหิ? อชฺชเวน ชเวน ขนฺติยา โสรจฺเจน. อิเมหิ โข ภิกฺขเว จตูหิ ธมฺเมหิ สมนฺนาคโต ภิกฺขุ อาหุเนยฺโย โหติ ฯเปฯ อนุตฺตรํ ปุญฺญกฺเขตฺตํ โลกสฺสาติ.  ทุติยํ.}

\csromanmatikaanchor{mtk.163}
\csromantocmarkref{subsection}{3. ปโตทสุตฺต}{mtk.163}
\bookmark[dest={mtk.163},level=2]{3. ปโตทสุตฺต}
\csromanmarkh{(12)}\csromanmarkhii{2. เกสิวคฺค 3. ปโตทสุตฺต}\csromanheadernomark{2}{(12) 2. เกสิวคฺค \textbf{3. ปโตทสุตฺต}}
\proseitem{113}{\csromanfolio{429}จตฺตาโรเม ภิกฺขเว ภทฺรา อสฺสาชานียา สนฺโต สํวิชฺชมานา โลกสฺมิํ. กตเม จตฺตาโร? อิธ ภิกฺขเว เอกจฺโจ ภโทฺร อสฺสาชานีโย ปโตทจฺฉายํ ทิสฺวา สํวิชฺชติ สํเวคํ อาปชฺชติ “กิํ นุ โข มํ อชฺช อสฺส\-ทมฺม\-สารถิ การณํ กาเรสฺสติ, กิมสฺสาหํ\csromansharedfootnote{ea739c4732d3f1f1}{กถมสฺสาหํ (ก)} ปฏิกโรมี”ติ. เอวรูโปปิ ภิกฺขเว อิเธกจฺโจ ภโทฺร อสฺสาชานีโย โหติ. อยํ ภิกฺขเว ปฐโม ภโทฺร อสฺสาชานีโย สนฺโต สํวิชฺชมาโน โลกสฺมิํ.}

\prose{ปุน จปรํ ภิกฺขเว อิเธกจฺโจ ภโทฺร อสฺสาชานีโย น เหว โข ปโตทจฺฉายํ ทิสฺวา สํวิชฺชติ สํเวคํ อาปชฺชติ, อปิ จ โข โลม\-เวธ\-วิทฺโธ สํวิชฺชติ สํเวคํ อาปชฺชติ. “กิํ นุ โข มํ อชฺช อสฺส\-ทมฺม\-สารถิ การณํ กาเรสฺสติ, กิมสฺสาหํ ปฏิกโรมี”ติ. เอวรูโปปิ ภิกฺขเว อิเธกจฺโจ ภโทฺร อสฺสาชานีโย โหติ. อยํ ภิกฺขเว ทุติโย ภโทฺร อสฺสาชานีโย สนฺโต สํวิชฺชมาโน โลกสฺมิํ.}

\prose{ปุน จปรํ ภิกฺขเว อิเธกจฺโจ ภโทฺร อสฺสาชานีโย น เหว โข ปโตทจฺฉายํ ทิสฺวา สํวิชฺชติ สํเวคํ อาปชฺชติ, นาปิ โลม\-เวธ\-วิทฺโธ สํวิชฺชติ สํเวคํ อาปชฺชติ, อปิ จ โข จมฺม\-เวธ\-วิทฺโธ สํวิชฺชติ สํเวคํ อาปชฺชติ “กิํ นุ โข มํ อชฺช อสฺส\-ทมฺม\-สารถิ การณํ กาเรสฺสติ, กิมสฺสาหํ ปฏิกโรมี”ติ. เอวรูโปปิ ภิกฺขเว อิเธกจฺโจ ภโทฺร อสฺสาชานีโย โหติ. อยํ ภิกฺขเว ตติโย ภโทฺร อสฺสาชานีโย สนฺโต สํวิชฺชมาโน โลกสฺมิํ.}

\prose{ปุน จปรํ ภิกฺขเว อิเธกจฺโจ ภโทฺร อสฺสาชานีโย น เหว โข ปโตทจฺฉายํ ทิสฺวา สํวิชฺชติ สํเวคํ อาปชฺชติ, นาปิ โลม\-เวธ\-วิทฺโธ สํวิชฺชติ สํเวคํ อาปชฺชติ, นาปิ จมฺม\-เวธ\-วิทฺโธ สํวิชฺชติ สํเวคํ อาปชฺชติ, อปิ จ โข อฏฺฐิ\-เวธ\-วิทฺโธ สํวิชฺชติ สํเวคํ อาปชฺชติ “กิํ นุ โข มํ อชฺช อสฺส\-ทมฺม\-สารถิ การณํ กาเรสฺสติ, กิมสฺสาหํ ปฏิกโรมี”ติ. เอวรูโปปิ ภิกฺขเว อิเธกจฺโจ ภโทฺร อสฺสาชานีโย โหติ, อยํ ภิกฺขเว จตุตฺโถ ภโทฺร อสฺสาชานีโย สนฺโต \csromanfolio{430}สํวิชฺชมาโน โลกสฺมิํ. อิเม โข ภิกฺขเว จตฺตาโร ภทฺรา อสฺสาชานียา สนฺโต สํวิชฺชมานา โลกสฺมิํ.}

\prose{เอวเมวํ โข ภิกฺขเว จตฺตาโรเม ภทฺรา ปุริสาชานียา สนฺโต สํวิชฺชมานา โลกสฺมิํ. กตเม จตฺตาโร? อิธ ภิกฺขเว เอกจฺโจ ภโทฺร ปุริสาชานีโย สุณาติ “อมุกสฺมิํ นาม คาเม วา นิคเม วา อิตฺถี วา ปุริโส วา ทุกฺขิโต วา กาลงฺกโต\csromansharedfootnote{41f80c016c431536}{กาลกโต (สี, สฺยา, กํ, อิ)} วา”ติ, โส เตน สํวิชฺชติ สํเวคํ อาปชฺชติ, สํวิคฺโค โยนิโส ปทหติ, ปหิตตฺโต กาเยน เจว ปรมสจฺจํ\csromansharedfootnote{ed5beec22604bc5d}{ปรมตฺถสจฺจํ (ก) ม 2. 145 ปิฏฺเฐ ปสฺสิตพฺพํ.} สจฺฉิกโรติ, ปญฺญาย จ อติวิชฺฌ ปสฺสติ. เสยฺยถาปิ โส ภิกฺขเว ภโทฺร อสฺสาชานีโย ปโตทจฺฉายํ ทิสฺวา สํวิชฺชติ สํเวคํ อาปชฺชติ, ตถูปมาหํ ภิกฺขเว อิมํ ภทฺรํ ปุริสาชานียํ วทามิ. เอวรูโปปิ ภิกฺขเว อิเธกจฺโจ ภโทฺร ปุริสาชานีโย โหติ. อยํ ภิกฺขเว ปฐโม ภโทฺร ปุริสาชานีโย สนฺโต สํวิชฺชมาโน โลกสฺมิํ.}

\prose{ปุน จปรํ ภิกฺขเว อิเธกจฺโจ ภโทฺร ปุริสาชานีโย น เหว โข สุณาติ “อมุกสฺมิํ นาม คาเม วา นิคเม วา อิตฺถี วา ปุริโส วา ทุกฺขิโต วา กาลงฺกโต วา”ติ, อปิ จ โข สามํ ปสฺสติ อิตฺถิํ วา ปุริสํ วา ทุกฺขิตํ วา กาลงฺกตํ วา. โส เตน สํวิชฺชติ สํเวคํ อาปชฺชติ, สํวิคฺโค โยนิโส ปทหติ, ปหิตตฺโต กาเยน เจว ปรมสจฺจํ สจฺฉิกโรติ, ปญฺญาย จ อติวิชฺฌ ปสฺสติ. เสยฺยถาปิ โส ภิกฺขเว ภโทฺร อสฺสาชานีโย โลม\-เวธ\-วิทฺโธ สํวิชฺชติ สํเวคํ อาปชฺชติ, ตถูปมาหํ ภิกฺขเว อิมํ ภทฺรํ ปุริสาชานียํ วทามิ. เอวรูโปปิ ภิกฺขเว อิเธกจฺโจ ภโทฺร ปุริสาชานีโย โหติ. อยํ ภิกฺขเว ทุติโย ภโทฺร ปุริสาชานีโย สนฺโต สํวิชฺชมาโน โลกสฺมิํ.}

\prose{ปุน จปรํ ภิกฺขเว อิเธกจฺโจ ภโทฺร ปุริสาชานิโย น เหว โข สุณาติ “อมุกสฺมิํ นาม คาเม วา นิคเม วา อิตฺถี วา ปุริโส วา ทุกฺขิโต วา กาลงฺกโต วา”ติ, นาปิ สามํ ปสฺสติ อิตฺถิํ วา ปุริสํ วา ทุกฺขิตํ วา กาลงฺกตํ วา, อปิ จ ขฺวสฺส ญาติ วา สาโลหิโต วา ทุกฺขิโต วา โหติ กาลงฺกโต วา. โส เตน สํวิชฺชติ สํเวคํ อาปชฺชติ, สํวิคฺโค โยนิโส ปทหติ, ปหิตตฺโต กาเยน เจว}

\vspace{\tipitakaparskip}
\csromancenter{(12) 2. เกสิวคฺค ปรมสจฺจํ สจฺฉิกโรติ, ปญฺญาย จ อติวิชฺฌ ปสฺสติ. เสยฺยถาปิ โส ภิกฺขเว ภโทฺร อสฺสาชานีโย จมฺม\-เวธ\-วิทฺโธ สํวิชฺชติ สํเวคํ อาปชฺชติ, ตถูปมาหํ ภิกฺขเว อิมํ ภทฺรํ ปุริสาชานียํ วทามิ. เอวรูโปปิ ภิกฺขเว อิเธกจฺโจ ภโทฺร ปุริสาชานีโย โหติ. อยํ ภิกฺขเว ตติโย ภโทฺร ปุริสาชานีโย สนฺโต สํวิชฺชมาโน โลกสฺมิํ.}
\prose{\csromanfolio{431}ปุน จปรํ ภิกฺขเว อิเธกจฺโจ ภโทฺร ปุริสาชานีโย น เหว โข สุณาติ “อมุกสฺมิํ นาม คาเม วา นิคเม วา อิตฺถี วา ปุริโส วา ทุกฺขิโต วา กาลงฺกโต วา”ติ, นาปิ สามํ ปสฺสติ อิตฺถิํ วา ปุริสํ วา ทุกฺขิตํ วา กาลงฺกตํ วา, นาปิสฺส ญาติ วา สาโลหิโต วา ทุกฺขิโต วา โหติ กาลงฺกโต วา, อปิ จ โข สามญฺเญว ผุฏฺโฐ โหติ สารีริกาหิ เวทนาหิ ทุกฺขาหิ ติพฺพาหิ\csromansharedfootnote{0f3ad6c08c2c9a0f}{ติปฺปาหิ (สี, อิ)} ขราหิ กฏุกาหิ อสาตาหิ อมนาปาหิ ปาณหราหิ. โส เตน สํวิชฺชติ สํเวคํ อาปชฺชติ, สํวิคฺโค โยนิโส ปทหติ, ปหิตตฺโต กาเยน เจว ปรมสจฺจํ สจฺฉิกโรติ, ปญฺญาย จ อติวิชฺฌ ปสฺสติ. เสยฺยถาปิ โส ภิกฺขเว ภโทฺร อสฺสาชานีโย อฏฺฐิ\-เวธ\-วิทฺโธ สํวิชฺชติ สํเวคํ อาปชฺชติ, ตถูปมาหํ ภิกฺขเว อิมํ ภทฺรํ ปุริสาชานียํ วทามิ. เอวรูโปปิ ภิกฺขเว อิเธกจฺโจ ภโทฺร ปุริสาชานีโย โหติ. อยํ ภิกฺขเว จตุตฺโถ ภโทฺร ปุริสาชานีโย สนฺโต สํวิชฺชมาโน โลกสฺมิํ. อิเม โข ภิกฺขเว จตฺตาโร ภทฺรา ปุริสาชานียา สนฺโต สํวิชฺชมานา โลกสฺมินฺติ.  ตติยํ.}

\csromanmatikaanchor{mtk.164}
\csromantocmarkref{subsection}{4. นาคสุตฺต}{mtk.164}
\bookmark[dest={mtk.164},level=2]{4. นาคสุตฺต}
\csromanheader{2}{4. \textbf{นาคสุตฺต}}
\proseitem{114}{จตูหิ ภิกฺขเว องฺเคหิ สมนฺนาคโต รญฺโญ นาโค ราชารโห โหติ ราชโภคฺโค, รญฺโญ องฺคนฺเตว สงฺขํ คจฺฉติ. กตเมหิ จตูหิ? อิธ ภิกฺขเว รญฺโญ นาโค โสตา จ โหติ หนฺตา จ ขนฺตา จ คนฺตา จ.}

\prose{กถญฺจ ภิกฺขเว รญฺโญ นาโค โสตา โหติ? อิธ ภิกฺขเว รญฺโญ นาโค ยเมนํ หตฺถิ\-ทมฺม\-สารถิ การณํ กาเรติ ยทิ วา กตปุพฺพํ, ยทิ วา อกตปุพฺพํ, ตํ อฏฺฐิํ กตฺวา\csromansharedfootnote{cf07f4161175797f}{อฏฺฐิกตฺวา (สี, สฺยา, กํ, อิ) อํ 2. 143 ปิฏฺเฐปิ.} มนสิ กตฺวา สพฺพเจตสา\csromansharedfootnote{a28120273788da69}{สพฺพํ เจตโส (สพฺพตฺถ)} \csromanfolio{432}สมนฺนาหริตฺวา โอหิตโสโต สุณาติ. เอวํ โข ภิกฺขเว รญฺโญ นาโค โสตา โหติ.}

\prose{กถญฺจ ภิกฺขเว รญฺโญ นาโค หนฺตา โหติ? อิธ ภิกฺขเว รญฺโญ นาโค สงฺคามคโต หตฺถิมฺปิ หนติ, หตฺถารุหมฺปิ หนติ, อสฺสมฺปิ หนติ, อสฺสารุหมฺปิ หนติ, รถมฺปิ หนติ, รถิกมฺปิ หนติ, ปตฺติกมฺปิ หนติ. เอวํ โข ภิกฺขเว รญฺโญ นาโค หนฺตา โหติ.}

\prose{กถญฺจ ภิกฺขเว รญฺโญ นาโค ขนฺตา โหติ? อิธ ภิกฺขเว รญฺโญ นาโค สงฺคามคโต ขโม โหติ สตฺติ\-ปฺปหารานํ อสิปฺปหารานํ อุสุปฺปหารานํ ผรสุ\-ปฺปหารานํ\csromansharedfootnote{fd58c5048b1330a1}{“ผรสุปฺปหารานนฺ”ติ อิทํ ปทํ สฺยามโปตฺถเก นตฺถิ. ม 3. 174 ปิฏฺเฐปิ ปสฺสิตพฺพํ.} เภริ\-ปณว\-สงฺข\-ติณว\-นินฺนาท\-สทฺทานํ. เอวํ โข ภิกฺขเว รญฺโญ นาโค ขนฺตา โหติ.}

\prose{กถญฺจ ภิกฺขเว รญฺโญ นาโค คนฺตา โหติ? อิธ ภิกฺขเว รญฺโญ นาโค ยเมนํ หตฺถิ\-ทมฺม\-สารถิ ทิสํ เปเสติ ยทิ วา คตปุพฺพํ, ยทิ วา อคตปุพฺพํ, ตํ ขิปฺปเมว คนฺตา โหติ. เอวํ โข ภิกฺขเว รญฺโญ นาโค คนฺตา โหติ. อิเมหิ โข ภิกฺขเว จตูหิ องฺเคหิ สมนฺนาคโต รญฺโญ นาโค ราชารโห โหติ ราชโภคฺโค, รญฺโญ องฺคนฺเตว สงฺขํ คจฺฉติ.}

\prose{เอวเมวํ โข ภิกฺขเว จตูหิ ธมฺเมหิ สมนฺนาคโต ภิกฺขุ อาหุเนยฺโย โหติ ฯเปฯ อนุตฺตรํ ปุญฺญกฺเขตฺตํ โลกสฺส. กตเมหิ จตูหิ? อิธ ภิกฺขเว ภิกฺขุ โสตา จ โหติ หนฺตา จ ขนฺตา จ คนฺตา จ.}

\prose{กถญฺจ ภิกฺขเว ภิกฺขุ โสตา โหติ? อิธ ภิกฺขเว ภิกฺขุ ตถาคต\-ปฺปเวทิเต ธมฺมวินเย เทสิยมาเน อฏฺฐิํ กตฺวา มนสิ กตฺวา สพฺพเจตสา สมนฺนาหริตฺวา โอหิตโสโต ธมฺมํ สุณาติ. เอวํ โข ภิกฺขเว ภิกฺขุ โสตา โหติ.}

\prose{กถญฺจ ภิกฺขเว ภิกฺขุ หนฺตา โหติ? อิธ ภิกฺขเว ภิกฺขุ อุปฺปนฺนํ กามวิตกฺกํ นาธิวาเสติ ปชหติ วิโนเทติ (หนติ)\csromansharedfootnote{fc79f4219f1ded52}{( ) นตฺถิ สีสฺยาอิโปตฺถเกสุ. อุปริ ปฏิปทาวคฺเค 470 ปิฏฺเฐ จตุตฺถสุตฺเต ปน “สเมตี”ติ ปทํ สพฺพตฺถปิ ทิสฺสติ.} พฺยนฺตีกโรติ อนภาวํ คเมติ. อุปฺปนฺนํ พฺยาปาท\-วิตกฺกํ ฯเปฯ. อุปฺปนฺนํ วิหิํสา\-วิตกฺกํ ฯเปฯ. อุปฺปนฺนุปฺปนฺเน ปาปเก อกุสเล ธมฺเม นาธิวาเสติ ปชหติ วิโนเทติ}

\vspace{\tipitakaparskip}
\csromancenter{(12) 2. เกสิวคฺค หนติ พฺยนฺตีกโรติ อนภาวํ คเมติ. เอวํ โข ภิกฺขเว ภิกฺขุ หนฺตา โหติ.}
\prose{\csromanfolio{433}กถญฺจ ภิกฺขเว ภิกฺขุ ขนฺตา โหติ? อิธ ภิกฺขเว ภิกฺขุ ขโม โหติ สีตสฺส อุณฺหสฺส ชิฆจฺฉาย ปิปาสาย ฑํส มกส วาตาตปสรีสปสมฺผสฺสานํ ทุรุตฺตานํ ทุราคตานํ วจนปถานํ, อุปฺปนฺนานํ สารีริกานํ เวทนานํ ทุกฺขานํ ติพฺพานํ ขรานํ กฏุกานํ อสาตานํ อมนาปานํ ปาณหรานํ อธิวาสก\-ชาติโก โหติ. เอวํ โข ภิกฺขเว ภิกฺขุ ขนฺตา โหติ.}

\prose{กถญฺจ ภิกฺขเว ภิกฺขุ คนฺตา โหติ? อิธ ภิกฺขเว ภิกฺขุ ยายํ ทิสา อคตปุพฺพา อิมินา ทีเฆน อทฺธุนา, ยทิทํ สพฺพ\-สงฺขาร\-สมโถ สพฺพู\-ปธิ\-ปฏินิสฺสคฺโค ตณฺหากฺขโย วิราโค นิโรโธ นิพฺพานํ, ตํ ขิปฺปญฺเญว คนฺตา โหติ. เอวํ โข ภิกฺขเว ภิกฺขุ คนฺตา โหติ. อิเมหิ โข ภิกฺขเว จตูหิ ธมฺเมหิ สมนฺนาคโต ภิกฺขุ อาหุเนยฺโย โหติ ฯเปฯ อนุตฺตรํ ปุญฺญกฺเขตฺตํ โลกสฺสาติ.  จตุตฺถํ.}

\csromanmatikaanchor{mtk.165}
\csromantocmarkref{subsection}{5. ฐานสุตฺต}{mtk.165}
\bookmark[dest={mtk.165},level=2]{5. ฐานสุตฺต}
\csromanheader{2}{5. \textbf{ฐานสุตฺต}}
\proseitem{115}{จตฺตาริมานิ ภิกฺขเว ฐานานิ. กตมานิ จตฺตาริ? อตฺถิ ภิกฺขเว ฐานํ อมนาปํ กาตุํ, ตญฺจ กยิรมานํ อนตฺถาย สํวตฺตติ. อตฺถิ ภิกฺขเว ฐานํ อมนาปํ กาตุํ, ตญฺจ กยิรมานํ อตฺถาย สํวตฺตติ. อตฺถิ ภิกฺขเว ฐานํ มนาปํ กาตุํ, ตญฺจ กยิรมานํ อนตฺถาย สํวตฺตติ. อตฺถิ ภิกฺขเว ฐานํ มนาปํ กาตุํ, ตญฺจ กยิรมานํ อตฺถาย สํวตฺตติ.}

\prose{ตตฺร ภิกฺขเว ยมิทํ\csromansharedfootnote{b13cc2deff52882a}{ยทิทํ (สฺยา, กํ, ก)} ฐานํ อมนาปํ กาตุํ, ตญฺจ กยิรมานํ อนตฺถาย สํวตฺตติ. อิทํ ภิกฺขเว ฐานํ อุภเยเนว น กตฺตพฺพํ มญฺญติ\csromansharedfootnote{30102dd7aae52117}{ปญฺญายติ (?)}\csromandash{}ยมฺปิทํ\csromansharedfootnote{935491e8bef34d5c}{ยทิทํ (ก) 4. ตํ (สี, อิ) สฺยามโปตฺถเก นตฺถิ.} ฐานํ อมนาปํ กาตุํ, อิมินาปิ นํ\csromansharedfootnote{f48e1654275f5546}{ตํ (อิ) สีสฺยาโปตฺถเกสุ นตฺถิ.} น กตฺตพฺพํ มญฺญติ. ยมฺปิทํ ฐานํ กยิรมานํ อนตฺถาย สํวตฺตติ, อิมินาปิ นํ\csromansharedfootnote{30a05007eddb39da}{ตํ (สี, อิ) สฺยามโปตฺถเก นตฺถิ.} น กตฺตพฺพํ มญฺญติ. อิทํ ภิกฺขเว ฐานํ อุภเยเนว น กตฺตพฺพํ มญฺญติ.}

\prose{ตตฺร ภิกฺขเว ยมิทํ ฐานํ อมนาปํ กาตุํ, ตญฺจ กยิรมานํ อตฺถาย สํวตฺตติ. อิมสฺมิํ ภิกฺขเว ฐาเน พาโล จ ปณฺฑิโต จ เวทิตพฺโพ \csromanfolio{434}ปุริสถาเม ปุริสวีริเย ปุริส\-ปรกฺกเม. น ภิกฺขเว พาโล อิติ ปฏิสญฺจิกฺขติ “กิญฺจาปิ โข อิทํ ฐานํ อมนาปํ กาตุํ, อถ จรหิทํ ฐานํ กยิรมานํ อตฺถาย สํวตฺตตี”ติ. โส ตํ ฐานํ น กโรติ, ตสฺส ตํ ฐานํ อกยิรมานํ อนตฺถาย สํวตฺตติ. ปณฺฑิโต จ โข ภิกฺขเว อิติ ปฏิสญฺจิกฺขติ “กิญฺจาปิ โข อิทํ ฐานํ อมนาปํ กาตุํ, อถ จรหิทํ ฐานํ กยิรมานํ อตฺถาย สํวตฺตตี”ติ. โส ตํ ฐานํ กโรติ, ตสฺส ตํ ฐานํ กยิรมานํ อตฺถาย สํวตฺตติ.}

\prose{ตตฺร ภิกฺขเว ยมิทํ\csromansharedfootnote{1ca9e85fa3e29f16}{ยทิทํ (สฺยา, กํ)} ฐานํ มนาปํ กาตุํ, ตญฺจ กยิรมานํ อนตฺถาย สํวตฺตติ. อิมสฺมิมฺปิ ภิกฺขเว ฐาเน พาโล จ ปณฺฑิโต จ เวทิตพฺโพ ปุริสถาเม ปุริสวีริเย ปุริส\-ปรกฺกเม. น ภิกฺขเว พาโล อิติ ปฏิสญฺจิกฺขติ “กิญฺจาปิ โข อิทํ ฐานํ มนาปํ กาตุํ, อถ จรหิทํ ฐานํ กยิรมานํ อนตฺถาย สํวตฺตตี”ติ. โส ตํ ฐานํ กโรติ, ตสฺส ตํ ฐานํ กยิรมานํ อนตฺถาย สํวตฺตติ. ปณฺฑิโต จ โข ภิกฺขเว อิติ ปฏิสญฺจิกฺขติ “กิญฺจาปิ โข อิทํ ฐานํ มนาปํ กาตุํ, อถ จรหิทํ ฐานํ กยิรมานํ อนตฺถาย สํวตฺตตี”ติ. โส ตํ ฐานํ น กโรติ, ตสฺส ตํ ฐานํ อกยิรมานํ อตฺถาย สํวตฺตติ.}

\prose{ตตฺร ภิกฺขเว ยมิทํ ฐานํ มนาปํ กาตุํ, ตญฺจ กยิรมานํ อตฺถาย สํวตฺตติ. อิทํ ภิกฺขเว ฐานํ อุภเยเนว กตฺตพฺพํ มญฺญติ\csromandash{}ยมฺปิทํ ฐานํ มนาปํ กาตุํ, อิมินาปิ นํ กตฺตพฺพํ มญฺญติ. ยมฺปิทํ ฐานํ กยิรมานํ อตฺถาย สํวตฺตติ, อิมินาปิ นํ กตฺตพฺพํ มญฺญติ. อิทํ ภิกฺขเว ฐานํ อุภเยเนว กตฺตพฺพํ มญฺญติ, อิมานิ โข ภิกฺขเว จตฺตาริ ฐานานีติ.}

\csromanmatikaanchor{mtk.166}
\csromantocmarkref{subsection}{6. อปฺปมาทสุตฺต}{mtk.166}
\bookmark[dest={mtk.166},level=2]{6. อปฺปมาทสุตฺต}
\csromanheader{2}{6. \textbf{อปฺปมาทสุตฺต}}
\proseitem{116}{จตูหิ ภิกฺขเว ฐาเนหิ อปฺปมาโท กรณีโย. กตเมหิ จตูหิ? กายทุจฺจริตํ ภิกฺขเว ปชหถ, กายสุจริตํ ภาเวถ, ตตฺถ จ มา ปมาทตฺถ. วจีทุจฺจริตํ ภิกฺขเว ปชหถ, วจีสุจริตํ ภาเวถ, ตตฺถ จ มา ปมาทตฺถ. มโนทุจฺจริตํ ภิกฺขเว ปชหถ, มโนสุจริตํ ภาเวถ, ตตฺถ จ มา ปมาทตฺถ. มิจฺฉาทิฏฺฐิํ ภิกฺขเว ปชหถ, สมฺมาทิฏฺฐิํ ภาเวถ, ตตฺถ จ มา ปมาทตฺถ.}

\csromanmarkh{(12)}\csromanmarkhii{2. เกสิวคฺค}\csromanheadernomark{1}{(12) 2. เกสิวคฺค}
\prose{\csromanfolio{435}ยโต โข ภิกฺขเว ภิกฺขุโน กายทุจฺจริตํ ปหีนํ โหติ, กายสุจริตํ ภาวิตํ. วจีทุจฺจริตํ ปหีนํ โหติ, วจีสุจริตํ ภาวิตํ. มโนทุจฺจริตํ ปหีนํ โหติ, มโนสุจริตํ ภาวิตํ. มิจฺฉาทิฏฺฐิ ปหีนา โหติ, สมฺมาทิฏฺฐิ ภาวิตา. โส น ภายติ สมฺปรายิกสฺส มรณสฺสาติ.  ฉฏฺฐํ.}

\csromanmatikaanchor{mtk.167}
\csromantocmarkref{subsection}{7. อารกฺขสุตฺต}{mtk.167}
\bookmark[dest={mtk.167},level=2]{7. อารกฺขสุตฺต}
\csromanheader{2}{7. \textbf{อารกฺขสุตฺต}}
\proseitem{117}{จตูสุ ภิกฺขเว ฐาเนสุ อตฺตรูเปน อปฺปมาโท สติ เจตโส อารกฺโข กรณีโย. กตเมสุ จตูสุ? “มา เม รชนีเยสุ ธมฺเมสุ จิตฺตํ รชฺชี”ติ อตฺตรูเปน อปฺปมาโท สติ เจตโส อารกฺโข กรณีโย. “มา เม โทสนีเยสุ ธมฺเมสุ จิตฺตํ ทุสฺสี”ติ อตฺตรูเปน อปฺปมาโท สติ เจตโส อารกฺโข กรณีโย. “มา เม โมหนีเยสุ ธมฺเมสุ จิตฺตํ มุยฺหี”ติ อตฺตรูเปน อปฺปมาโท สติ เจตโส อารกฺโข กรณีโย. “มา เม มทนีเยสุ ธมฺเมสุ จิตฺตํ มชฺชี”ติ อตฺตรูเปน อปฺปมาโท สติ เจตโส อารกฺโข กรณีโย.}

\prose{ยโต โข ภิกฺขเว ภิกฺขุโน รชนีเยสุ ธมฺเมสุ จิตฺตํ น รชฺชติ วีตราคตฺตา. โทสนีเยสุ ธมฺเมสุ จิตฺตํ น ทุสฺสติ วีตโทสตฺตา. โมหนีเยสุ ธมฺเมสุ จิตฺตํ น มุยฺหติ วีตโมหตฺตา. มทนีเยสุ ธมฺเมสุ จิตฺตํ น มชฺชติ วีตมทตฺตา. โส น ฉมฺภติ น กมฺปติ น เวธติ น สนฺตาสํ อาปชฺชติ. น จ ปน สมณ\-วจน\-เหตุปิ คจฺฉตีติ.  }

\csromanmatikaanchor{mtk.168}
\csromantocmarkref{subsection}{8. สํเวชนียสุตฺต}{mtk.168}
\bookmark[dest={mtk.168},level=2]{8. สํเวชนียสุตฺต}
\csromanheader{2}{8. สํเวชนีย\-สุตฺต}
\csromanmatikaanchor{mtk.169}
\csromantocmarkref{subsection}{9-10. ภยสุตฺตทฺวยํ}{mtk.169}
\bookmark[dest={mtk.169},level=2]{9-10. ภยสุตฺตทฺวยํ}
\proseitem{118}{จตฺตาริมานิ ภิกฺขเว สทฺธสฺส กุลปุตฺตสฺส ทสฺสนียานิ สํเวชนียานิ ฐานานิ. กตมานิ จตฺตาริ? “อิธ ตถาคโต ชาโต”ติ ภิกฺขเว สทฺธสฺส กุลปุตฺตสฺส ทสฺสนียํ สํเวชนียํ ฐานํ. “อิธ ตถาคโต อนุตฺตรํ สมฺมาสมฺโพธิํ อภิสมฺพุทฺโธ”ติ ภิกฺขเว สทฺธสฺส กุลปุตฺตสฺส ทสฺสนียํ สํเวชนียํ ฐานํ. “อิธ ตถาคโต อนุตฺตรํ ธมฺมจกฺกํ ปวตฺเตสี”ติ ภิกฺขเว สทฺธสฺส กุลปุตฺตสฺส ทสฺสนียํ สํเวชนียํ ฐานํ. “อิธ ตถาคโต อนุปาทิเสสาย นิพฺพานธาตุยา \csromanfolio{436}ปรินิพฺพุโต”ติ ภิกฺขเว สทฺธสฺส กุลปุตฺตสฺส ทสฺสนียํ สํเวชนียํ ฐานํ. อิมานิ โข ภิกฺขเว จตฺตาริ สทฺธสฺส กุลปุตฺตสฺส ทสฺสนียานิ สํเวชนียานิ ฐานานีติ.  อฏฺฐมํ.}

\csromanheader{2}{9. ปฐมภย\-สุตฺต}
\proseitem{119}{จตฺตาริมานิ ภิกฺขเว ภยานิ. กตมานิ จตฺตาริ? ชาติภยํ ชราภยํ พฺยาธิภยํ มรณภยํ. อิมานิ โข ภิกฺขเว จตฺตาริ ภยานีติ.  นวมํ.}

\csromanheader{2}{10. ทุติยภย\-สุตฺต}
\proseitemwithclosermidruled{120}{จตฺตาริมานิ ภิกฺขเว ภยานิ. กตมานิ จตฺตาริ? อคฺคิภยํ อุทกภยํ ราชภยํ โจรภยํ. อิมานิ โข ภิกฺขเว จตฺตาริ ภยานีติ.  ทสมํ.}{เกสิวคฺโค ทุติโย.}
\csromancenter{\textbf{ตสฺสุทฺทานํ}}
\setlength{\csromangathapagewidth}{0pt}
\setlength{\csromangathawakwidth}{0pt}
\csromangathameasuregroup{เกสิ ชโว ปโตโท จ, \\ อปฺปมาโท จ อารกฺโข,}{\csromangathabat{เกสิ ชโว ปโตโท จ,}{นาโค ฐาเนน ปญฺจมํ.} \\ \csromangathabat{อปฺปมาโท จ อารกฺโข,}{สํเวชนียญฺจ เทฺว ภยาติ.}}
\csromangathasetleft{เกสิ ชโว ปโตโท จ, \\ อปฺปมาโท จ อารกฺโข,}
\csromangathagroup{\csromangathastanza{\csromangathabat{เกสิ ชโว ปโตโท จ,}{นาโค ฐาเนน ปญฺจมํ.} \\ \csromangathabat{อปฺปมาโท จ อารกฺโข,}{สํเวชนียญฺจ เทฺว ภยาติ.}}}

\csromanmatikaanchor{mtk.170}
\csromantocmarknopage{section}{(13) 3. ภยวคฺค}{mtk.170}
\bookmark[dest={mtk.170},level=1]{(13) 3. ภยวคฺค}
\csromanmarkh{(13)}\csromanmarkhii{3. ภยวคฺค}\csromanheadernomark{1}{\textbf{(13) 3. ภยวคฺค}}
\csromanmatikaanchor{mtk.171}
\csromantocmarkref{subsection}{1. อตฺตานุวาทสุตฺต}{mtk.171}
\bookmark[dest={mtk.171},level=2]{1. อตฺตานุวาทสุตฺต}
\csromanheader{2}{1. อตฺตานุวาท\-สุตฺต}
\proseitem{121}{จตฺตาริมานิ ภิกฺขเว ภยานิ. กตมานิ จตฺตาริ? อตฺตานุวาทภยํ ปรานุวาทภยํ ทณฺฑภยํ ทุคฺคติภยํ.}

\prose{กตมญฺจ ภิกฺขเว อตฺตานุวาทภยํ? อิธ ภิกฺขเว เอกจฺโจ อิติ ปฏิสญฺจิกฺขติ “อหญฺเจว\csromansharedfootnote{cf278ffcd2fa33b2}{อหญฺเจ (?)} โข ปน กาเยน ทุจฺจริตํ จเรยฺยํ, วาจาย ทุจฺจริตํ จเรยฺยํ, มนสา ทุจฺจริตํ จเรยฺยํ. กิญฺจ ตํ ยํ มํ\csromansharedfootnote{37e84a2a51528e4a}{กิญฺจ ตํ มํ (สี), กิญฺจ มํ (สฺยา, กํ), กิญฺจ ตํ กมฺมํ (อิ, ก)} อตฺตา สีลโต น อุปวเทยฺยา”ติ. โส อตฺตานุวาท\-ภยสฺส ภีโต กายทุจฺจริตํ}

\vspace{\tipitakaparskip}
\csromancenter{(13) 3. ภยวคฺค ปหาย กายสุจริตํ ภาเวติ, วจีทุจฺจริตํ ปหาย วจีสุจริตํ ภาเวติ, มโนทุจฺจริตํ ปหาย มโนสุจริตํ ภาเวติ, สุทฺธํ อตฺตานํ ปริหรติ. อิทํ วุจฺจติ ภิกฺขเว อตฺตานุวาทภยํ.}
\prose{\csromanfolio{437}กตมญฺจ ภิกฺขเว ปรานุวาทภยํ? อิธ ภิกฺขเว เอกจฺโจ อิติ ปฏิสญฺจิกฺขติ “อหญฺเจว โข ปน กาเยน ทุจฺจริตํ จเรยฺยํ, วาจาย ทุจฺจริตํ จเรยฺยํ, มนสา ทุจฺจริตํ จเรยฺยํ. กิญฺจ ตํ ยํ มํ ปเร สีลโต น อุปวเทยฺยุนฺ”ติ. โส ปรานุวาท\-ภยสฺส ภีโต กายทุจฺจริตํ ปหาย กายสุจริตํ ภาเวติ, วจีทุจฺจริตํ ปหาย วจีสุจริตํ ภาเวติ, มโนทุจฺจริตํ ปหาย มโนสุจริตํ ภาเวติ, สุทฺธํ อตฺตานํ ปริหรติ. อิทํ วุจฺจติ ภิกฺขเว ปรานุวาทภยํ.}

\prose{กตมญฺจ ภิกฺขเว ทณฺฑภยํ? อิธ ภิกฺขเว เอกจฺโจ ปสฺสติ โจรํ อาคุจาริํ ราชาโน คเหตฺวา วิวิธา กมฺมการณา กาเรนฺเต กสาหิปิ ตาเฬนฺเต เวตฺเตหิปิ ตาเฬนฺเต อทฺธ\-ทณฺฑเกหิปิ ตาเฬนฺเต หตฺถมฺปิ ฉินฺทนฺเต ปาทมฺปิ ฉินฺทนฺเต หตฺถปาทมฺปิ ฉินฺทนฺเต กณฺณมฺปิ ฉินฺทนฺเต นาสมฺปิ ฉินฺทนฺเต กณฺณนาสมฺปิ ฉินฺทนฺเต พิลงฺค\-ถาลิกมฺปิ กโรนฺเต สงฺข\-มุณฺฑิกมฺปิ กโรนฺเต ราหุมุขมฺปิ กโรนฺเต โชติมาลิกมฺปิ กโรนฺเต หตฺถ\-ปชฺโชติกมฺปิ กโรนฺเต เอรก\-วตฺติกมฺปิ กโรนฺเต จีรก\-วาสิกมฺปิ กโรนฺเต เอเณยฺยกมฺปิ กโรนฺเต พลิสมํสิกมฺปิ กโรนฺเต กหาปณกมฺปิ กโรนฺเต ขารา\-ปตจฺฉิกมฺปิ กโรนฺเต ปลิฆ\-ปริวตฺติกมฺปิ กโรนฺเต ปลาล\-ปีฐกมฺปิ กโรนฺเต ตตฺเตนปิ เตเลน โอสิญฺจนฺเต สุนเขหิปิ ขาทาเปนฺเต ชีวนฺตมฺปิ สูเล อุตฺตาเสนฺเต อสินาปิ สีสํ ฉินฺทนฺเต.}

\prose{ตสฺส เอวํ โหติ “ยถารูปานํ โข ปาปกานํ กมฺมานํ เหตุ โจรํ อาคุจาริํ ราชาโน คเหตฺวา วิวิธา กมฺมการณา กาเรนฺติ, กสาหิปิ ตาเฬนฺติ ฯเปฯ อสินาปิ สีสํ ฉินฺทนฺติ, อหญฺเจว โข ปน เอวรูปํ ปาปกมฺมํ กเรยฺยํ, มมฺปิ ราชาโน คเหตฺวา เอวรูปา วิวิธา กมฺมการณา กาเรยฺยุํ, กสาหิปิ ตาเฬยฺยุํ, เวตฺเตหิปิ ตาเฬยฺยุํ, อทฺธ\-ทณฺฑเกหิปิ ตาเฬยฺยุํ, หตฺถมฺปิ ฉินฺเทยฺยุํ, ปาทมฺปิ ฉินฺเทยฺยุํ, หตฺถปาทมฺปิ ฉินฺเทยฺยุํ, กณฺณมฺปิ ฉินฺเทยฺยุํ, นาสมฺปิ ฉินฺเทยฺยุํ, กณฺณนาสมฺปิ ฉินฺเทยฺยุํ, พิลงฺค\-ถาลิกมฺปิ กเรยฺยุํ, สงฺข\-มุณฺฑิกมฺปิ กเรยฺยุํ, ราหุมุขมฺปิ กเรยฺยุํ, โชติมาลิกมฺปิ กเรยฺยุํ, หตฺถ\-ปชฺโชติกมฺปิ กเรยฺยุํ, เอรก\-วตฺติกมฺปิ กเรยฺยุํ, จีรก\-วาสิกมฺปิ กเรยฺยุํ, เอเณยฺยกมฺปิ กเรยฺยุํ, พลิสมํสิกมฺปิ กเรยฺยุํ, \csromanfolio{438}กหาปณกมฺปิ กเรยฺยุํ, ขารา\-ปตจฺฉิกมฺปิ กเรยฺยุํ, ปลิฆ\-ปริวตฺติกมฺปิ กเรยฺยุํ, ปลาล\-ปีฐกมฺปิ กเรยฺยุํ, ตตฺเตนปิ เตเลน โอสิญฺเจยฺยุํ, สุนเขหิปิ ขาทาเปยฺยุํ, ชีวนฺตมฺปิ สูเล อุตฺตาเสยฺยุํ, อสินาปิ สีสํ ฉินฺเตยฺยุนฺ”ติ. โส ทณฺฑภยสฺส ภีโต น ปเรสํ ปาภตํ วิลุมฺปนฺโต จรติ. กายทุจฺจริตํ ปหาย ฯเปฯ สุทฺธํ อตฺตานํ ปริหรติ. อิทํ วุจฺจติ ภิกฺขเว ทณฺฑภยํ.}

\prose{กตมญฺจ ภิกฺขเว ทุคฺคติภยํ? อิธ ภิกฺขเว เอกจฺโจ อิติ ปฏิสญฺจิกฺขติ “กาย\-ทุจฺจริตสฺส โข ปาปโก วิปาโก อภิสมฺปรายํ, วจีทุจฺจริตสฺส ปาปโก วิปาโก อภิสมฺปรายํ, มโน\-ทุจฺจริตสฺส ปาปโก วิปาโก อภิสมฺปรายํ, อหญฺเจว โข ปน กาเยน ทุจฺจริตํ จเรยฺยํ, วาจาย ทุจฺจริตํ จเรยฺยํ, มนสา ทุจฺจริตํ จเรยฺยํ, กิญฺจ ตํ ยาหํ น กายสฺส เภทา ปรํ มรณา อปายํ ทุคฺคติํ วินิปาตํ นิรยํ อุปปชฺเชยฺยนฺ”ติ. โส ทุคฺคติ\-ภยสฺส ภีโต กายทุจฺจริตํ ปหาย กายสุจริตํ ภาเวติ, วจีทุจฺจริตํ ปหาย วจีสุจริตํ ภาเวติ, มโนทุจฺจริตํ ปหาย มโนสุจริตํ ภาเวติ, สุทฺธํ อตฺตานํ ปริหรติ. อิทํ วุจฺจติ ภิกฺขเว ทุคฺคติภยํ. อิมานิ โข ภิกฺขเว จตฺตาริ ภยานีติ.}

\csromanmatikaanchor{mtk.172}
\csromantocmarkref{subsection}{2. อูมิภยสุตฺต}{mtk.172}
\bookmark[dest={mtk.172},level=2]{2. อูมิภยสุตฺต}
\csromanheader{2}{2. อู\textbf{มิภยสุตฺต}}
\proseitem{122}{จตฺตาริมานิ ภิกฺขเว ภยานิ อุทโก\-โรหนฺตสฺส\csromansharedfootnote{f5bcf879ce207450}{อุทโกโรหนฺเต (ม 2. 122 ปิฏฺเฐ)} ปาฏิกงฺขิตพฺพานิ. กตมานิ จตฺตาริ? อูมิภยํ กุมฺภีลภยํ อาวฏฺฏภยํ สุสุกาภยํ. อิมานิ โข ภิกฺขเว จตฺตาริ ภยานิ อุทโก\-โรหนฺตสฺส ปาฏิกงฺขิตพฺพานิ. เอวเมวํ โข ภิกฺขเว จตฺตาริ ภยานิ อิเธกจฺจสฺส กุลปุตฺตสฺส อิมสฺมิํ ธมฺมวินเย อคารสฺมา\csromansharedfootnote{66fde923d7d5257d}{สทฺธา อคารสฺมา (สี, สฺยา, กํ)} อนคาริยํ ปพฺพชิตสฺส\csromansharedfootnote{39c521afd7c8793b}{ปพฺพชโต (สี)} ปาฏิกงฺขิตพฺพานิ. กตมานิ จตฺตาริ? อูมิภยํ กุมฺภีลภยํ อาวฏฺฏภยํ สุสุกาภยํ.}

\prose{กตมญฺจ ภิกฺขเว อูมิภยํ? อิธ ภิกฺขเว เอกจฺโจ กุลปุตฺโต สทฺธา อคารสฺมา อนคาริยํ ปพฺพชิโต โหติ “โอติณฺโณมฺหิ ชาติยา ชราย มรเณน โสเกหิ ปริเทเวหิ ทุกฺเขหิ โสมนสฺเสหิ อุปายาเสหิ}

\vspace{\tipitakaparskip}
\csromancenter{(13) 3. ภยวคฺค ทุกฺโขติณฺโณ ทุกฺขปเรโต, อปฺเปว นาม อิมสฺส เกวลสฺส ทุกฺข\-กฺขนฺธสฺส อนฺตกิริยา ปญฺญาเยถา”ติ. ตเมนํ ตถา ปพฺพชิตํ สมานํ สพฺรหฺมจาริโน โอวทนฺติ อนุสาสนฺติ “เอวํ เต อภิกฺกมิตพฺพํ, เอวํ เต ปฏิกฺกมิตพฺพํ, เอวํ เต อาโลเกตพฺพํ เอวํ เต วิโลเกตพฺพํ, เอวํ เต สมิญฺชิตพฺพํ, เอวํ เต ปสาริตพฺพํ, เอวํ เต สงฺฆาฏิ\-ปตฺต\-จีวรํ ธาเรตพฺพนฺ”ติ. ตสฺส เอวํ โหติ “มยํ โข ปุพฺเพ อคาริยภูตา สมานา อญฺเญ โอวทามปิ อนุสาสามปิ. อิเม ปนมฺหากํ ปุตฺตมตฺตา มญฺเญ นตฺตมตฺตา มญฺเญ โอวทิตพฺพํ อนุสาสิตพฺพํ มญฺญนฺตี”ติ. โส กุปิโต อนตฺตมโน สิกฺขํ ปจฺจกฺขาย หีนายาวตฺตติ. อยํ วุจฺจติ ภิกฺขเว ภิกฺขุ อูมิภยสฺส ภีโต สิกฺขํ ปจฺจกฺขาย หีนายาวตฺโต. “อูมิภยนฺ”ติ โข ภิกฺขเว โกธูปายาสสฺเสตํ อธิวจนํ. อิทํ วุจฺจติ ภิกฺขเว อูมิภยํ.}
\prose{\csromanfolio{439}กตมญฺจ ภิกฺขเว กุมฺภีลภยํ? อิธ ภิกฺขเว เอกจฺโจ กุลปุตฺโต สทฺธา อคารสฺมา อนคาริยํ ปพฺพชิโต โหติ “โอติณฺโณมฺหิ ชาติยา ชราย มรเณน โสเกหิ ปริเทเวหิ ทุกฺเขหิ โทมนสฺเสหิ อุปายาเสหิ ทุกฺโขติณฺโณ ทุกฺขปเรโต, อปฺเปว นาม อิมสฺส เกวลสฺส ทุกฺข\-กฺขนฺธสฺส อนฺตกิริยา ปญฺญาเยถา”ติ. ตเมนํ ตถา ปพฺพชิตํ สมานํ สพฺรหฺมจาริโน โอวทนฺติ อนุสาสนฺติ “อิทํ เต ขาทิตพฺพํ, อิทํ เต น ขาทิตพฺพํ. อิทํ เต ภุญฺชิตพฺพํ, อิทํ เต น ภุญฺชิตพฺพํ. อิทํ เต สายิตพฺพํ, อิทํ เต น สายิตพฺพํ. อิทํ เต ปาตพฺพํ, อิทํ เต น ปาตพฺพํ. กปฺปิยํ เต ขาทิตพฺพํ, อกปฺปิยํ เต น ขาทิตพฺพํ. กปฺปิยํ เต ภุญฺชิตพฺพํ, อกปฺปิยํ เต น ภุญฺชิตพฺพํ. กปฺปิยํ เต สายิตพฺพํ, อกปฺปิยํ เต น สายิตพฺพํ. กปฺปิยํ เต ปาตพฺพํ, อกปฺปิยํ เต น ปาตพฺพํ. กาเล เต ขาทิตพฺพํ, วิกาเล เต น ขาทิตพฺพํ. กาเล เต ภุญฺชิตพฺพํ, วิกาเล เต น ภุญฺชิตพฺพํ. กาเล เต สายิตพฺพํ, วิกาเล เต น สายิตพฺพํ. กาเล เต ปาตพฺพํ, วิกาเล เต น ปาตพฺพนฺ”ติ. ตสฺส เอวํ โหติ “มยํ โข ปุพฺเพ อคาริยภูตา สมานา ยํ อิจฺฉาม ตํ ขาทาม, ยํ น อิจฺฉาม น ตํ ขาทาม. ยํ อิจฺฉาม ตํ ภุญฺชาม, ยํ น อิจฺฉาม น ตํ ภุญฺชาม. ยํ อิจฺฉาม ตํ สายาม, ยํ น อิจฺฉาม น ตํ สายาม. ยํ อิจฺฉาม ตํ ปิวาม, ยํ น อิจฺฉาม น ตํ ปิวาม. กปฺปิยมฺปิ ขาทาม, อกปฺปิยมฺปิ ขาทาม. กปฺปิยมฺปิ ภุญฺชาม, อกปฺปิยมฺปิ ภุญฺชาม. กปฺปิยมฺปิ สายาม, อกปฺปิยมฺปิ สายาม. กปฺปิยมฺปิ ปิวาม, อกปฺปิยมฺปิ ปิวาม. กาเลปิ ขาทาม, วิกาเลปิ ขาทาม. กาเลปิ ภุญฺชาม, วิกาเลปิ ภุญฺชาม. \csromanfolio{440}กาเลปิ สายาม, วิกาเลปิ สายาม. กาเลปิ ปิวาม, วิกาเลปิ ปิวาม. ยมฺปิ โน สทฺธา คหปติกา ทิวา วิกาเล ปณีตํ ขาทนียํ วา โภชนียํ วา เทนฺติ. ตตฺรปิเม\csromansharedfootnote{1d2216884594708a}{ตตฺถปิเม (ม 2. 124 ปิฏฺเฐ)} มุขาวรณํ มญฺเญ กโรนฺตี”ติ, โส กุปิโต อนตฺตมโน สิกฺขํ ปจฺจกฺขาย หีนายาวตฺตติ. อยํ วุจฺจติ ภิกฺขเว ภิกฺขุ กุมฺภีล\-ภยสฺส ภีโต สิกฺขํ ปจฺจกฺขาย หีนายาวตฺโต. “กุมฺภีลภยนฺ”ติ โข ภิกฺขเว โอทริกตฺตสฺเสตํ อธิวจนํ. อิทํ วุจฺจติ ภิกฺขเว กุมฺภีลภยํ.}

\prose{กตมญฺจ ภิกฺขเว อาวฏฺฏภยํ? อิธ ภิกฺขเว เอกจฺโจ กุลปุตฺโต สทฺธา อคารสฺมา อนคาริยํ ปพฺพชิโต โหติ “โอติณฺโณมฺหิ ชาติยา ชราย มรเณน โสเกหิ ปริเทเวหิ ทุกฺเขหิ โทมนสฺเสหิ อุปายาเสหิ ทุกฺโขติณฺโณ ทุกฺขปเรโต, อปฺเปว นาม อิมสฺส เกวลสฺส ทุกฺข\-กฺขนฺธสฺส อนฺตกิริยา ปญฺญาเยถา”ติ. โส เอวํ ปพฺพชิโต สมาโน ปุพฺพณฺหสมยํ นิวาเสตฺวา ปตฺต\-จีวรมา\-ทาย คามํ วา นิคมํ วา ปิณฺฑาย ปวิสติ อรกฺขิเตเนว กาเยน อรกฺขิตาย วาจาย อรกฺขิเตน จิตฺเตน อนุปฏฺฐิตาย สติยา อสํวุเตหิ อินฺทฺริเยหิ. โส ตตฺถ ปสฺสติ คหปติํ วา คหปติปุตฺตํ วา ปญฺจหิ กามคุเณหิ สมปฺปิตํ สมงฺคีภูตํ ปริจารย\-มานํ. ตสฺส เอวํ โหติ “มยํ โข ปุพฺเพ อคาริยภูตา สมานา ปญฺจหิ กามคุเณหิ สมปฺปิตา สมงฺคีภูตา ปริจาริมฺหา. สํวิชฺชนฺติ โข ปน เม กุเล โภคา, สกฺกา โภเค จ ภุญฺจิตุํ ปุญฺญานิ จ กาตุํ. ยํนูนาหํ สิกฺขํ ปจฺจกฺขาย หีนายาวตฺติตฺวา โภเค จ ภุญฺเชยฺยํ, ปุญฺญานิ จ กเรยฺยนฺ”ติ. โส สิกฺขํ ปจฺจกฺขาย หีนายาวตฺตติ. อยํ วุจฺจติ ภิกฺขเว ภิกฺขุ อาวฏฺฏภยสฺส ภีโต สิกฺขํ ปจฺจกฺขาย หีนายาวตฺโต. “อาวฏฺฏภยนฺ”ติ โข ภิกฺขเว ปญฺจนฺเนตํ กามคุณานํ อธิวจนํ. อิทํ วุจฺจติ ภิกฺขเว อาวฏฺฏภยํ.}

\prose{กตมญฺจ ภิกฺขเว สุสุกาภยํ, อิธ ภิกฺขเว เอกจฺโจ กุลปุตฺโต สทฺธา อคารสฺมา อนคาริยํ ปพฺพชิโต โหติ “โอติณฺโณมฺหิ ชาติยา ชราย มรเณน โสเกหิ ปริเทเวหิ ทุกฺเขหิ โทมนสฺเสหิ อุปายาเสหิ ทุกฺโขติณฺโณ ทุกฺขปเรโต, อปฺเปว นาม อิมสฺส เกวลสฺส ทุกฺข\-กฺขนฺธสฺส อนฺตกิริยา ปญฺญาเยถา”ติ. โส เอวํ ปพฺพชิโต สมาโน ปุพฺพณฺหสมยํ นิวาเสตฺวา ปตฺต\-จีวรมา\-ทาย คามํ วา}

\vspace{\tipitakaparskip}
\csromancenter{(13) 3. ภยวคฺค นิคมํ วา ปิณฺฑาย ปวิสติ อรกฺขิเตเนว กาเยน อรกฺขิตาย วาจาย อรกฺขิเตน จิตฺเตน อนุปฏฺฐิตาย สติยา อสํวุเตหิ อินฺทฺริเยหิ. โส ตตฺถ ปสฺสติ มาตุคามํ ทุนฺนิวตฺถํ วา ทุปฺปารุตํ วา, ตสฺส มาตุคามํ ทิสฺวา ทุนฺนิวตฺถํ วา ทุปฺปารุตํ วา ราโค จิตฺตํ อนุทฺธํเสติ, โส ราคา\-นุทฺธํสิเตน จิตฺเตน สิกฺขํ ปจฺจกฺขาย หีนายาวตฺตติ. อยํ วุจฺจติ ภิกฺขเว ภิกฺขุ สุสุกาภยสฺส ภีโต สิกฺขํ ปจฺจกฺขาย หีนายาวตฺโต. “สุสุกาภยนฺ”ติ โข ภิกฺขเว มาตุคามสฺเสตํ อธิวจนํ. อิทํ วุจฺจติ ภิกฺขเว สุสุกาภยํ. อิมานิ โข ภิกฺขเว จตฺตาริ ภยานิ อิเธกจฺจสฺส กุลปุตฺตสฺส อิมสฺมิํ ธมฺมวินเย อคารสฺมา อนคาริยํ ปพฺพชิตสฺส ปาฏิกงฺขิตพฺพานีติ.  ทุติยํ.}
\csromanheader{2}{3. ปฐม\-นานากรณ\-สุตฺต}
\csromanmatikaanchor{mtk.173}
\csromantocmarkref{subsection}{3-4. นานากรณสุตฺตทฺวยํ}{mtk.173}
\bookmark[dest={mtk.173},level=2]{3-4. นานากรณสุตฺตทฺวยํ}
\proseitem{123}{\csromanfolio{441}จตฺตาโรเม ภิกฺขเว ปุคฺคลา สนฺโต สํวิชฺชมานา โลกสฺมิํ. กตเม จตฺตาโร?\csromansymbolfootnote{*}{อภิ 4. 351 ปิฏฺเฐปิ.}อิธ ภิกฺขเว เอกจฺโจ ปุคฺคโล วิวิจฺเจว กาเมหิ วิวิจฺจ อกุสเลหิ ธมฺเมหิ สวิตกฺกํ สวิจารํ วิเวกชํ ปีติสุขํ ปฐมํ ฌานํ อุปสมฺปชฺช วิหรติ, โส ตทสฺสาเทติ, ตํ นิกาเมติ, เตน จ วิตฺติํ อาปชฺชติ. ตตฺถ ฐิโต ตทธิมุตฺโต ตพฺพหุล\-วิหารี อปริหีโน กาลํ กุรุมาโน พฺรหฺม\-กายิกานํ เทวานํ สหพฺยตํ อุปปชฺชติ. พฺรหฺม\-กายิกานํ ภิกฺขเว เทวานํ กปฺโป อายุปฺปมาณํ, ตตฺถ ปุถุชฺชโน ยาวตายุกํ ฐตฺวา ยาวตกํ เตสํ เทวานํ อายุปฺปมาณํ, ตํ สพฺพํ เขเปตฺวา นิรยมฺปิ คจฺฉติ, ติรจฺฉาน\-โยนิมฺปิ คจฺฉติ, เปตฺติวิสยมฺปิ คจฺฉติ. ภควโต ปน สาวโก ตตฺถ ยาวตายุกํ ฐตฺวา ยาวตกํ เตสํ เทวานํ อายุปฺปมาณํ, ตํ สพฺพํ เขเปตฺวา ตสฺมิํ เยว ภเว ปรินิพฺพายติ. อยํ โข ภิกฺขเว วิเสโส อยํ อธิปฺปยาโส อิทํ นานากรณํ สุตวโต อริยสาวกสฺส อสฺสุตวตา ปุถุชฺชเนน ยทิทํ คติยา อุปปตฺติยา สติ\csromansharedfootnote{64afadda7ceeb5ba}{ยทิทํ จุติยา อุปปตฺติยา จาติ (ก) 270 ปิฏฺเฐ ปสฺสิตพฺพํ.}.}

\prose{ปุน จปรํ ภิกฺขเว อิเธกจฺโจ ปุคฺคโล วิตกฺก\-วิจารานํ วูปสมา อชฺฌตฺตํ สมฺปสาทนํ เจตโส เอโกทิภาวํ อวิตกฺกํ อวิจารํ สมาธิชํ ปีติสุขํ ทุติยํ ฌานํ อุปสมฺปชฺช วิหรติ, โส ตทสฺสาเทติ, ตํ นิกาเมติ, เตน จ วิตฺติํ อาปชฺชติ. ตตฺถ ฐิโต ตทธิมุตฺโต \csromanfolio{442}ตพฺพหุล\-วิหารี อปริหีโน กาลํ กุรุมาโน อาภสฺสรานํ เทวานํ สหพฺยตํ อุปปชฺชติ. อาภสฺสรานํ ภิกฺขเว เทวานํ เทฺว กปฺปา อายุปฺปมาณํ. ตตฺถ ปุถุชฺชโน ยาวตายุกํ ฐตฺวา ยาวตกํ เตสํ เทวานํ อายุปฺปมาณํ, ตํ สพฺพํ เขเปตฺวา นิรยมฺปิ คจฺฉติ, ติรจฺฉาน\-โยนิมฺปิ คจฺฉติ, เปตฺติวิสยมฺปิ คจฺฉติ. ภควโต ปน สาวโก ตตฺถ ยาวตายุกํ ฐตฺวา ยาวตกํ เตสํ เทวานํ อายุปฺปมาณํ, ตํ สพฺพํ เขเปตฺวา ตสฺมิํ เยว ภเว ปรินิพฺพายติ. อยํ โข ภิกฺขเว วิเสโส อยํ อธิปฺปยาโส อิทํ นานากรณํ สุตวโต อริยสาวกสฺส อสฺสุตวตา ปุถุชฺชเนน ยทิทํ คติยา อุปปตฺติยา สติ.}

\prose{ปุน จปรํ ภิกฺขเว อิเธกจฺโจ ปุคฺคโล ปีติยา จ วิราคา อุเปกฺขโก จ วิหรติ, สโต จ สมฺปชาโน สุขญฺจ กาเยน ปฏิสํเวเทติ, ยํ ตํ อริยา อาจิกฺขนฺติ “อุเปกฺขโก สติมา สุขวิหารี”ติ, ตติยํ ฌานํ อุปสมฺปชฺช วิหรติ, โส ตทสฺสาเทติ, ตํ นิกาเมติ, เตน จ วิตฺติํ อาปชฺชติ. ตตฺถ ฐิโต ตทธิมุตฺโต ตพฺพหุล\-วิหารี อปริหีโน กาลํ กุรุมาโน สุภกิณฺหานํ เทวานํ สหพฺยตํ อุปปชฺชติ. สุภกิณฺหานํ ภิกฺขเว เทวานํ จตฺตาโร กปฺปา อายุปฺปมาณํ. ตตฺถ ปุถุชฺชโน ยาวตายุกํ ฐตฺวา ยาวตกํ เตสํ เทวานํ อายุปฺปมาณํ, ตํ สพฺพํ เขเปตฺวา นิรยมฺปิ คจฺฉติ, ติรจฺฉาน\-โยนิมฺปิ คจฺฉติ, เปตฺติวิสยมฺปิ คจฺฉติ. ภควโต ปน สาวโก ตตฺถ ยาวตายุกํ ฐตฺวา ยาวตกํ เตสํ เทวานํ อายุปฺปมาณํ, ตํ สพฺพํ เขเปตฺวา ตสฺมิํ เยว ภเว ปรินิพฺพายติ. อยํ โข ภิกฺขเว วิเสโส อยํ อธิปฺปยาโส อิทํ นานากรณํ สุตวโต อริยสาวกสฺส อสฺสุตวตา ปุถุชฺชเนน ยทิทํ คติยา อุปปตฺติยา สติ.}

\prose{ปุน จปรํ ภิกฺขเว อิเธกจฺโจ ปุคฺคโล สุขสฺส จ ปหานา ทุกฺขสฺส จ ปหานา ปุพฺเพว โสมนสฺส\-โทมนสฺสานํ อตฺถงฺคมา อทุกฺขมสุขํ อุเปกฺขา\-สติ\-ปาริสุทฺธิํ จตุตฺถํ ฌานํ อุปสมฺปชฺช วิหรติ, โส ตทสฺสาเทติ, ตํ นิกาเมติ, เตน จ วิตฺติํ อาปชฺชติ. ตตฺถ ฐิโต ตทธิมุตฺโต ตพฺพหุล\-วิหารี อปริหีโน กาลํ กุรุมาโน เวหปฺผลานํ เทวานํ สหพฺยตํ อุปปชฺชติ. เวหปฺผลานํ ภิกฺขเว เทวานํ ปญฺจ กปฺปสตานิ อายุปฺปมาณํ. ตตฺถ ปุถุชฺชโน ยาวตายุกํ ฐตฺวา ยาวตกํ เตสํ เทวานํ อายุปฺปมาณํ, ตํ สพฺพํ เขเปตฺวา นิรยมฺปิ คจฺฉติ, ติรจฺฉาน\-โยนิมฺปิ}

\vspace{\tipitakaparskip}
\csromancenter{(13) 3. ภยวคฺค คจฺฉติ, เปตฺติวิสยมฺปิ คจฺฉติ. ภควโต ปน สาวโก ตตฺถ ยาวตายุกํ ฐตฺวา ยาวตกํ เตสํ เทวานํ อายุปฺปมาณํ, ตํ สพฺพํ เขเปตฺวา ตสฺมิํ เยว ภเว ปรินิพฺพายติ. อยํ โข ภิกฺขเว วิเสโส อยํ อธิปฺปยาโส อิทํ นานากรณํ สุตวโต อริยสาวกสฺส อสฺสุตวตา ปุถุชฺชเนน ยทิทํ คติยา อุปปตฺติยา สติ. อิเม โข ภิกฺขเว จตฺตาโร ปุคฺคลา สนฺโต สํวิชฺชมานา โลกสฺมินฺติ.  ตติยํ.}
\csromanheader{2}{4. ทุติย\-นานากรณ\-สุตฺต}
\csromanmatikaanchor{mtk.174}
\csromantocmarkref{subsection}{5-6. เมตฺตาสุตฺตทฺวยํ}{mtk.174}
\bookmark[dest={mtk.174},level=2]{5-6. เมตฺตาสุตฺตทฺวยํ}
\proseitem{124}{\csromanfolio{443}จตฺตาโรเม ภิกฺขเว ปุคฺคลา สนฺโต สํวิชฺชมานา โลกสฺมิํ. กตเม จตฺตาโร? อิธ ภิกฺขเว เอกจฺโจ ปุคฺคโล วิวิจฺเจว กาเมหิ ฯเปฯ ปฐมํ ฌานํ อุปสมฺปชฺช วิหรติ. โส ยเทว ตตฺถ โหติ รูปคตํ เวทนาคตํ สญฺญาคตํ สงฺขารคตํ วิญฺญาณคตํ, เต ธมฺเม อนิจฺจโต ทุกฺขโต โรคโต คณฺฑโต สลฺลโต อฆโต อาพาธโต ปรโต ปโลกโต สุญฺญโต อนตฺตโต สมนุปสฺสติ. โส กายสฺส เภทา ปรํ มรณา สุทฺธาวาสานํ เทวานํ สหพฺยตํ อุปปชฺชติ. อยํ ภิกฺขเว อุปปตฺติ อสาธารณา ปุถุชฺชเนหิ.}

\prose{ปุน จปรํ ภิกฺขเว อิเธกจฺโจ ปุคฺคโล วิตกฺก\-วิจารานํ วูปสมา ฯเปฯ ทุติยํ ฌานํ ฯเปฯ ตติยํ ฌานํ ฯเปฯ จตุตฺถํ ฌานํ อุปสมฺปชฺช วิหรติ. โส ยเทว ตตฺถ โหติ รูปคตํ เวทนาคตํ สญฺญาคตํ สงฺขารคตํ วิญฺญาณคตํ, เต ธมฺเม อนิจฺจโต ทุกฺขโต โรคโต คณฺฑโต สลฺลโต อฆโต อาพาธโต ปรโต ปโลกโต สุญฺญโต อนตฺตโต สมนุปสฺสติ. โส กายสฺส เภทา ปรํ มรณา สุทฺธาวาสานํ เทวานํ สหพฺยตํ อุปปชฺชติ. อยํ ภิกฺขเว อุปปตฺติ อสาธารณา ปุถุชฺชเนหิ. อิเม โข ภิกฺขเว จตฺตาโร ปุคฺคลา สนฺโต สํวิชฺชมานา โลกสฺมินฺติ.  จตุตฺถํ.}

\csromanheader{2}{5. ปฐม\-เมตฺตา\-สุตฺต}
\proseitem{125}{จตฺตาโรเม ภิกฺขเว ปุคฺคลา สนฺโต สํวิชฺชมานา โลกสฺมิํ. กตเม จตฺตาโร? อิธ ภิกฺขเว เอกจฺโจ ปุคฺคโล เมตฺตา\-สหคเตน เจตสา เอกํ ทิสํ ผริตฺวา วิหรติ. ตถา ทุติยํ. ตถา ตติยํ. ตถา จตุตฺถํ. อิติ อุทฺธมโธ ติริยํ สพฺพธิ สพฺพตฺตตาย สพฺพาวนฺตํ \csromanfolio{444}โลกํ เมตฺตา\-สหคเตน เจตสา วิปุเลน มหคฺคเตน อปฺปมาเณน อเวเรน อพฺยาปชฺเชน ผริตฺวา วิหรติ, โส ตทสฺสาเทติ, ตํ นิกาเมติ, เตน จ วิตฺติํ อาปชฺชติ. ตตฺถ ฐิโต ตทธิมุตฺโต ตพฺพหุล\-วิหารี อปริหีโน กาลํ กุรุมาโน พฺรหฺม\-กายิกานํ เทวานํ สหพฺยตํ อุปปชฺชติ. พฺรหฺม\-กายิกานํ ภิกฺขเว เทวานํ กปฺโป อายุปฺปมาณํ. ตตฺถ ปุถุชฺชโน ยาวตายุกํ ฐตฺวา ยาวตกํ เตสํ เทวานํ อายุปฺปมาณํ, ตํ สพฺพํ เขเปตฺวา นิรยมฺปิ คจฺฉติ, ติรจฺฉาน\-โยนิมฺปิ คจฺฉติ, เปตฺติวิสยมฺปิ คจฺฉติ. ภควโต ปน สาวโก ตตฺถ ยาวตายุกํ ฐตฺวา ยาวตกํ เตสํ เทวานํ อายุปฺปมาณํ, ตํ สพฺพํ เขเปตฺวา ตสฺมิํเยว ภเว ปรินิพฺพายติ. อยํ โข ภิกฺขเว วิเสโส อยํ อธิปฺปยาโส อิทํ นานากรณํ สุตวโต อริยสาวกสฺส อสฺสุตวตา ปุถุชฺชเนน ยทิทํ คติยา อุปปตฺติยา สติ.}

\prose{ปุน จปรํ ภิกฺขเว อิเธกจฺโจ ปุคฺคโล กรุณา\-สหคเตน เจตสา ฯเปฯ มุทิตา\-สหคเตน เจตสา ฯเปฯ อุเปกฺขา\-สหคเตน เจตสา เอกํ ทิสํ ผริตฺวา วิหรติ. ตถา ทุติยํ. ตถา ตติยํ. ตถา จตุตฺถํ. อิติ อุทฺธมโธ ติริยํ สพฺพธิ สพฺพตฺตตาย สพฺพาวนฺตํ โลกํ อุเปกฺขา\-สหคเตน เจตสา วิปุเลน มหคฺคเตน อปฺปมาเณน อเวเรน อพฺยาปชฺเชน ผริตฺวา วิหรติ, โส ตทสฺสาเทติ, ตํ นิกาเมติ, เตน จ วิตฺติํ อาปชฺชติ. ตตฺถ ฐิโต ตทธิมุตฺโต ตพฺพหุล\-วิหารี อปริหีโน กาลํ กุรุมาโน อาภสฺสรานํ\csromansharedfootnote{0d965bc3af7d50bc}{สฺยามโปตฺถเก ปน กรุณาทโย ตโย วิหารา อาภสฺสราทีหิ ตีหิ วิสุํ วิสุํ โยเชตฺวา ปริปุณฺณเมว ทสฺสิตํ.} เทวานํ สหพฺยตํ อุปปชฺชติ. อาภสฺสรานํ ภิกฺขเว เทวานํ เทฺว กปฺปา อายุปฺปมาณํ ฯเปฯ สุภกิณฺหานํ\csromansharedfootnote{0d965bc3af7d50bc}{สฺยามโปตฺถเก ปน กรุณาทโย ตโย วิหารา อาภสฺสราทีหิ ตีหิ วิสุํ วิสุํ โยเชตฺวา ปริปุณฺณเมว ทสฺสิตํ.} เทวานํ สหพฺยตํ อุปปชฺชติ. สุภกิณฺหานํ ภิกฺขเว เทวานํ จตฺตาโร กปฺปา อายุปฺปมาณํ ฯเปฯ เวหปฺผลานํ\csromansharedfootnote{0d965bc3af7d50bc}{สฺยามโปตฺถเก ปน กรุณาทโย ตโย วิหารา อาภสฺสราทีหิ ตีหิ วิสุํ วิสุํ โยเชตฺวา ปริปุณฺณเมว ทสฺสิตํ.} เทวานํ สหพฺยตํ อุปปชฺชติ. เวหปฺผลานํ ภิกฺขเว เทวานํ ปญฺจ กปฺปสตานิ อายุปฺปมาณํ. ตตฺถ ปุถุชฺชโน ยาวตายุกํ ฐตฺวา ยาวตกํ เตสํ เทวานํ อายุปฺปมาณํ, ตํ สพฺพํ เขเปตฺวา นิรยมฺปิ คจฺฉติ, ติรจฺฉาน\-โยนิมฺปิ คจฺฉติ, เปตฺติวิสยมฺปิ คจฺฉติ. ภควโต ปน สาวโก ตตฺถ ยาวตายุกํ ฐตฺวา ยาวตกํ เตสํ เทวานํ อายุปฺปมาณํ, ตํ สพฺพํ เขเปตฺวา ตสฺมิํเยว ภเว ปรินิพฺพายติ. อยํ โข ภิกฺขเว}

\vspace{\tipitakaparskip}
\csromancenter{(13) 3. ภยวคฺค วิเสโส อยํ อธิปฺปยาโส อิทํ นานากรณํ สุตวโต อริยสาวกสฺส อสฺสุตวตา ปุถุชฺชเนน ยทิทํ คติยา อุปปตฺติยา สติ. อิเม โข ภิกฺขเว จตฺตาโร ปุคฺคลา สนฺโต สํวิชฺชมานา โลกสฺมินฺติ.  ปญฺจมํ.}
\csromanheader{2}{6. ทุติย\-เมตฺตา\-สุตฺต}
\proseitem{126}{\csromanfolio{445}จตฺตาโรเม ภิกฺขเว ปุคฺคลา สนฺโต สํวิชฺชมานา โลกสฺมิํ. กตเม จตฺตาโร? อิธ ภิกฺขเว เอกจฺโจ ปุคฺคโล เมตฺตา\-สหคเตน เจตสา เอกํ ทิสํ ผริตฺวา วิหรติ. ตถา ทุติยํ. ตถา ตติยํ. ตถา จตุตฺถํ. อิติ อุทฺธมโธ ติริยํ สพฺพธิ สพฺพตฺตตาย สพฺพาวนฺตํ โลกํ เมตฺตา\-สหคเตน เจตสา วิปุเลน มหคฺคเตน อปฺปมาเณน อเวเรน อพฺยาปชฺเชน ผริตฺวา วิหรติ. โส ยเทว ตตฺถ โหติ รูปคตํ เวทนาคตํ สญฺญาคตํ สงฺขารคตํ วิญฺญาณคตํ, เต ธมฺเม อนิจฺจโต ทุกฺขโต โรคโต คณฺฑโต สลฺลโต อฆโต อาพาธโต ปรโต ปโลกโต สุญฺญโต อนตฺตโต สมนุปสฺสติ. โส กายสฺส เภทา ปรํ มรณา สุทฺธาวาสานํ เทวานํ สหพฺยตํ อุปปชฺชติ. อยํ ภิกฺขเว อุปปตฺติ อสาธารณา ปุถุชฺชเนหิ.}

\prose{ปุน จปรํ ภิกฺขเว อิเธกจฺโจ ปุคฺคโล กรุณา ฯเปฯ มุทิตา ฯเปฯ อุเปกฺขา\-สหคเตน เจตสา เอกํ ทิสํ ผริตฺวา วิหรติ. ตถา ทุติยํ. ตถา ตติยํ. ตถา จตุตฺถํ. อิติ อุทฺธมโธ ติริยํ สพฺพธิ สพฺพตฺตตาย สพฺพาวนฺตํ โลกํ อุเปกฺขา\-สหคเตน เจตสา วิปุเลน มหคฺคเตน อปฺปมาเณน อเวเรน อพฺยาปชฺเชน ผริตฺวา วิหรติ. โส ยเทว ตตฺถ โหติ รูปคตํ เวทนาคตํ สญฺญาคตํ สงฺขารคตํ วิญฺญาณคตํ, เต ธมฺเม อนิจฺจโต ทุกฺขโต โรคโต คณฺฑโต สลฺลโต อฆโต อาพาธโต ปรโต ปโลกโต สุญฺญโต อนตฺตโต สมนุปสฺสติ. โส กายสฺส เภทา ปรํ มรณา สุทฺธาวาสานํ เทวานํ สหพฺยตํ อุปปชฺชติ. อยํ ภิกฺขเว อุปปตฺติ อสาธารณา ปุถุชฺชเนหิ. อิเม โข ภิกฺขเว จตฺตาโร ปุคฺคลา สนฺโต สํวิชฺชมานา โลกสฺมินฺติ.  ฉฏฺฐํ.}

\csromanmarkh{จตุกฺกนิปาต\-ปาฬิ}\csromanmarkhii{7. ปฐม\-ตถาคต\-อจฺฉริย\-สุตฺต}\csromanheadernomark{2}{จตุกฺกนิปาต\-ปาฬิ 7. ปฐม\-ตถาคต\-อจฺฉริย\-สุตฺต}
\csromanmatikaanchor{mtk.175}
\csromantocmarkref{subsection}{7-8. ตถาคตอจฺฉริยสุตฺตทฺวยํ}{mtk.175}
\bookmark[dest={mtk.175},level=2]{7-8. ตถาคตอจฺฉริยสุตฺตทฺวยํ}
\proseitem{127}{\csromanfolio{446}ตถาคตสฺส ภิกฺขเว อรหโต สมฺมา\-สมฺพุทฺธสฺส ปาตุภาวา จตฺตาโร อจฺฉริยา อพฺภุตา ธมฺมา\csromansharedfootnote{95b2d811d7757a52}{อพฺภุตธมฺมา (สฺยา, กํ, ก)} ปาตุภวนฺติ. กตเม จตฺตาโร? ยทา ภิกฺขเว โพธิสตฺโต ตุสิตา กายา จวิตฺวา สโต สมฺปชาโน มาตุกุจฺฉิํ โอกฺกมติ. อถ สเทวเก โลเก สมารเก สพฺรหฺมเก สสฺสมณ\-พฺราหฺมณิยา ปชาย สเทว\-มนุสฺสาย อปฺปมาโณ อุฬาโร โอภาโส ปาตุภวติ อติกฺกมฺเมว เทวานํ เทวานุภาวํ. ยาปิ ตา โลกนฺตริกา อฆา อสํวุตา อนฺธการา อนฺธ\-การ\-ติมิสา, ยตฺถปิเมสํ\csromansharedfootnote{97c9ffffa9266db8}{ยตฺถิเมสํ (สี, สฺยา, กํ)} จนฺทิม\-สูริยานํ เอวํ\-มหิทฺธิกานํ เอวํ\-มหา\-นุภาวานํ อาภา นานุโภนฺติ, ตตฺถปิ อปฺปมาโณ อุฬาโร โอภาโส ปาตุภวติ อติกฺกมฺเมว เทวานํ เทวานุภาวํ. เยปิ ตตฺถ สตฺตา อุปปนฺนา, เตปิ เตโนภาเสน อญฺญมญฺญํ สญฺชานนฺติ “อญฺเญปิ กิร โภ สนฺติ สตฺตา อิธูปปนฺนา”ติ. ตถาคตสฺส ภิกฺขเว อรหโต สมฺมา\-สมฺพุทฺธสฺส ปาตุภาวา อยํ ปฐโม อจฺฉริโย อพฺภุโต ธมฺโม\csromansharedfootnote{961a6e2a584abbdc}{อพฺภุตธมฺโม (สฺยา, กํ, ก)} ปาตุภวติ.}

\prose{ปุน จปรํ ภิกฺขเว ยทา โพธิสตฺโต สโต สมฺปชาโน มาตุกุจฺฉิมฺหา นิกฺขมติ. อถ สเทวเก โลเก สมารเก สพฺรหฺมเก สสฺสมณ\-พฺราหฺมณิยา ปชาย สเทว\-มนุสฺสาย อปฺปมาโณ อุฬาโร โอภาโส ปาตุภวติ อติกฺกมฺเมว เทวานํ เทวานุภาวํ. ยาปิ ตา โลกนฺตริกา อฆา อสํวุตา อนฺธการา อนฺธ\-การ\-ติมิสา, ยตฺถปิเมสํ จนฺทิม\-สูริยานํ เอวํ\-มหิทฺธิกานํ เอวํ\-มหา\-นุภาวานํ อาภา นานุโภนฺติ, ตตฺถปิ อปฺปมาโณ อุฬาโร โอภาโส ปาตุภวติ อติกฺกมฺเมว เทวานํ เทวานุภาวํ. เยปิ ตตฺถ สตฺตา อุปปนฺนา, เตปิ เตโนภาเสน อญฺญมญฺญํ สญฺชานนฺติ “อญฺเญปิ กิร โภ สนฺติ สตฺตา อิธูปปนฺนา”ติ. ตถาคตสฺส ภิกฺขเว อรหโต สมฺมา\-สมฺพุทฺธสฺส ปาตุภาวา อยํ ทุติโย อจฺฉริโย อพฺภุโต ธมฺโม ปาตุภวติ.}

\prose{ปุน จปรํ ภิกฺขเว ยทา ตถาคโต อนุตฺตรํ สมฺมาสมฺโพธิํ อภิสมฺ\-พุชฺฌติ. อถ สเทวเก โลเก สมารเก สพฺรหฺมเก สสฺสมณ\-พฺราหฺมณิยา ปชาย สเทว\-มนุสฺสาย อปฺปมาโณ อุฬาโร}

\vspace{\tipitakaparskip}
\csromancenter{(13) 3. ภยวคฺค โอภาโส ปาตุภวติ อติกฺกมฺเมว เทวานํ เทวานุภาวํ. ยาปิ ตา โลกนฺตริกา อฆา อสํวุตา อนฺธการา อนฺธ\-การ\-ติมิสา, ยตฺถปิเมสํ จนฺทิม\-สูริยานํ เอวํ\-มหิทฺธิกานํ เอวํ\-มหา\-นุภาวานํ อาภา นานุโภนฺติ, ตตฺถปิ อปฺปมาโณ อุฬาโร โอภาโส ปาตุภวติ อติกฺกมฺเมว เทวานํ เทวานุภาวํ. เยปิ ตตฺถ สตฺตา อุปปนฺนา, เตปิ เตโนภาเสน อญฺญมญฺญํ สญฺชานนฺติ “อญฺเญปิ กิร โภ สนฺติ สตฺตา อิธูปปนฺนา”ติ. ตถาคตสฺส ภิกฺขเว อรหโต สมฺมา\-สมฺพุทฺธสฺส ปาตุภาวา อยํ ตติโย อจฺฉริโย อพฺภุโต ธมฺโม ปาตุภวติ.}
\prose{\csromanfolio{447}ปุน จปรํ ภิกฺขเว ยทา ตถาคโต อนุตฺตรํ ธมฺมจกฺกํ ปวตฺเตติ. อถ สเทวเก โลเก สมารเก สพฺรหฺมเก สสฺสมณ\-พฺราหฺมณิยา ปชาย สเทว\-มนุสฺสาย อปฺปมาโณ อุฬาโร โอภาโส ปาตุภวติ อติกฺกมฺเมว เทวานํ เทวานุภาวํ. ยาปิ ตา โลกนฺตริกา อฆา อสํวุตา อนฺธการา อนฺธ\-การ\-ติมิสา, ยตฺถปิเมสํ จนฺทิม\-สูริยานํ เอวํ\-มหิทฺธิกานํ เอวํ\-มหา\-นุภาวานํ อาภา นานุโภนฺติ, ตตฺถปิ อปฺปมาโณ อุฬาโร โอภาโส ปาตุภวติ อติกฺกมฺเมว เทวานํ เทวานุภาวํ. เยปิ ตตฺถ สตฺตา อุปปนฺนา, เตปิ เตโนภาเสน อญฺญมญฺญํ สญฺชานนฺติ “อญฺเญปิ กิร โภ สนฺติ สตฺตา อิธูปปนฺนา”ติ. ตถาคตสฺส ภิกฺขเว อรหโต สมฺมา\-สมฺพุทฺธสฺส ปาตุภาวา อยํ จตุตฺโถ อจฺฉริโย อพฺภุโต ธมฺโม ปาตุภวติ. ตถาคตสฺส ภิกฺขเว อรหโต สมฺมา\-สมฺพุทฺธสฺส ปาตุภาวา อิเม จตฺตาโร อจฺฉริยา อพฺภุตา ธมฺมา ปาตุภวนฺตีติ.  สตฺตมํ.}

\prose{8. ทุติย\-ตถาคต\-อจฺฉริย\-สุตฺต ตถาคตสฺส ภิกฺขเว อรหโต สมฺมา\-สมฺพุทฺธสฺส ปาตุภาวา จตฺตาโร อจฺฉริยา อพฺภุตา ธมฺมา ปาตุภวนฺติ. กตเม จตฺตาโร? อาลยารามา\csromansharedfootnote{def813fc6c952340}{อาลยรามา (อญฺญสุตฺเตสุ)} ภิกฺขเว ปชา อาลยรตา อาลยสมฺมุทิตา. สา ตถาคเตน อนาลเย ธมฺเม เทสิยมาเน สุสฺสูสติ, โสตํ โอทหติ, อญฺญา จิตฺตํ อุปฏฺฐเปติ. ตถาคตสฺส ภิกฺขเว อรหโต สมฺมา\-สมฺพุทฺธสฺส ปาตุภาวา อยํ ปฐโม อจฺฉริโย อพฺภุโต ธมฺโม ปาตุภวติ.}

\csromanmatikaanchor{mtk.176}
\csromantocmarkref{subsection}{9-10. อานนฺทอจฺฉริยสุตฺตาทิ}{mtk.176}
\bookmark[dest={mtk.176},level=2]{9-10. อานนฺทอจฺฉริยสุตฺตาทิ}
\prose{\csromanfolio{448}มานารามา ภิกฺขเว ปชา มานรตา มานสมฺมุทิตา. สา ตถาคเตน มานวินเย ธมฺเม เทสิยมาเน สุสฺสูสติ, โสตํ โอทหติ, อญฺญา จิตฺตํ อุปฏฺฐเปติ. ตถาคตสฺส ภิกฺขเว อรหโต สมฺมา\-สมฺพุทฺธสฺส ปาตุภาวา อยํ ทุติโย อจฺฉริโย อพฺภุโต ธมฺโม ปาตุภวติ.}

\prose{อนุปสมารามา ภิกฺขเว ปชา อนุปสมรตา อนุปสม\-สมฺมุทิตา. สา ตถาคเตน โอปสมิเก ธมฺเม เทสิยมาเน สุสฺสูสติ, โสตํ โอทหติ, อญฺญา จิตฺตํ อุปฏฺฐเปติ. ตถาคตสฺส ภิกฺขเว อรหโต สมฺมา\-สมฺพุทฺธสฺส ปาตุภาวา อยํ ตติโย อจฺฉริโย อพฺภุโต ธมฺโม ปาตุภวติ.}

\prose{อวิชฺชาคตา ภิกฺขเว ปชา อณฺฑภูตา ปริโยนทฺธา. สา ตถาคเตน อวิชฺชาวินเย ธมฺเม เทสิยมาเน สุสฺสูสติ, โสตํ โอทหติ, อญฺญา จิตฺตํ อุปฏฺฐเปติ. ตถาคตสฺส ภิกฺขเว อรหโต สมฺมา\-สมฺพุทฺธสฺส ปาตุภาวา อยํ จตุตฺโถ อจฺฉริโย อพฺภุโต ธมฺโม ปาตุภวติ. ตถาคตสฺส ภิกฺขเว อรหโต สมฺมา\-สมฺพุทฺธสฺส ปาตุภาวา อิเม จตฺตาโร อจฺฉริยา อพฺภุตา ธมฺมา ปาตุภวนฺตีติ.  อฏฺฐมํ.}

\prose{9. อานนฺท\-อจฺฉริย\-สุตฺต}

\proseitem{129}{จตฺตาโรเม ภิกฺขเว อจฺฉริยา อพฺภุตา ธมฺมา อานนฺเท. กตเม จตฺตาโร? สเจ ภิกฺขเว ภิกฺขุปริสา อานนฺทํ ทสฺสนาย อุปสงฺกมติ, ทสฺสเนนปิ สา อตฺตมนา โหติ. ตตฺร เจ อานนฺโท ธมฺมํ ภาสติ, ภาสิเตนปิ สา อตฺตมนา โหติ. อติตฺตาว ภิกฺขเว ภิกฺขุปริสา โหติ, อถ อานนฺโท ตุณฺหี ภวติ.}

\prose{สเจ ภิกฺขเว ภิกฺขุนิ\-ปริสา อานนฺทํ ทสฺสนาย อุปสงฺกมติ, ทสฺสเนนปิ สา อตฺตมนา โหติ. ตตฺถ เจ อานนฺโท ธมฺมํ ภาสติ, ภาสิเตนปิ สา อตฺตมนา โหติ. อติตฺตาว ภิกฺขเว ภิกฺขุนิ\-ปริสา โหติ, อถ อานนฺโท ตุณฺหี ภวติ.}

\prose{สเจ ภิกฺขเว อุปาสกปริสา อานนฺทํ ทสฺสนาย อุปสงฺกมติ, ทสฺสเนนปิ สา อตฺตมนา โหติ. ตตฺร เจ อานนฺโท ธมฺมํ ภาสติ,}

\vspace{\tipitakaparskip}
\csromancenter{(13) 3. ภยวคฺค ภาสิเตนปิ สา อตฺตมนา โหติ. อติตฺตาว ภิกฺขเว อุปาสกปริสา โหติ, อถ อานนฺโท ตุณฺหี ภวติ.}
\prose{\csromanfolio{449}สเจ ภิกฺขเว อุปาสิกาปริสา อานนฺทํ ทสฺสนาย อุปสงฺกมติ, ทสฺสเนนปิ สา อตฺตมนา โหติ. ตตฺร เจ อานนฺโท ธมฺมํ ภาสติ, ภาสิเตนปิ สา อตฺตมนา โหติ. อติตฺตาว ภิกฺขเว อุปาสิกาปริสา โหติ, อถ อานนฺโท ตุณฺหี ภวติ. อิเม โข ภิกฺขเว จตฺตาโร อจฺฉริยา อพฺภุตา ธมฺมา อานนฺเทติ.  นวมํ.}

\prose{10. จกฺกวตฺติ\-อจฺฉริย\-สุตฺต}

\proseitem{130}{จตฺตาโรเม ภิกฺขเว อจฺฉริยา อพฺภุตา ธมฺมา รญฺเญ จกฺกวตฺติมฺหิ. กตเม จตฺตาโร? สเจ ภิกฺขเว ขตฺติยปริสา ราชานํ จกฺกวตฺติํ ทสฺสนาย อุปสงฺกมติ, ทสฺสเนนปิ สา อตฺตมนา โหติ. ตตฺร เจ ราชา จกฺกวตฺตี ภาสติ, ภาสิเตนปิ สา อตฺตมนา โหติ. อติตฺตาว ภิกฺขเว ขตฺติยปริสา โหติ, อถ ราชา จกฺกวตฺตี ตุณฺหี ภวติ.}

\prose{สเจ ภิกฺขเว พฺราหฺมณปริสา ราชานํ จกฺกวตฺติํ ทสฺสนาย อุปสงฺกมติ, ทสฺสเนนปิ สา อตฺตมนา โหติ. ตตฺร เจ ราชา จกฺกวตฺตี ภาสติ, ภาสิเตนปิ สา อตฺตมนา โหติ. อติตฺตาว ภิกฺขเว พฺราหฺมณปริสา โหติ, อถ ราชา จกฺกวตฺตี ตุณฺหี ภวติ.}

\prose{สเจ ภิกฺขเว คหปติปริสา ราชานํ จกฺกวตฺติํ ทสฺสนาย อุปสงฺกมติ, ทสฺสเนนปิ สา อตฺตมนา โหติ. ตตฺร เจ ราชา จกฺกวตฺตี ภาสติ, ภาสิเตนปิ สา อตฺตมนา โหติ. อติตฺตาว ภิกฺขเว คหปติปริสา โหติ, อถ ราชา จกฺกวตฺตี ตุณฺหี ภวติ.}

\prose{สเจ ภิกฺขเว สมณปริสา ราชานํ จกฺกวตฺติํ ทสฺสนาย อุปสงฺกมติ, ทสฺสเนนปิ สา อตฺตมนา โหติ. ตตฺร เจ ราชา จกฺกวตฺตี ภาสติ, ภาสิเตนปิ สา อตฺตมนา โหติ. อติตฺตาว ภิกฺขเว สมณปริสา โหติ, อถ ราชา จกฺกวตฺตี ตุณฺหี ภวติ. อิเม โข ภิกฺขเว จตฺตาโร อจฺฉริยา อพฺภุตา ธมฺมา รญฺเญ จกฺกวตฺติมฺหิ.}

\prose{\csromanfolio{450}เอวเมวํ โข ภิกฺขเว จตฺตาโร\csromansharedfootnote{1117a02e0906b839}{จกฺกาโรเม (ก)} อจฺฉริยา อพฺภุตา ธมฺมา อานนฺเท. กตเม จตฺตาโร? สเจ ภิกฺขเว ภิกฺขุปริสา อานนฺทํ ทสฺสนาย อุปสงฺกมติ, ทสฺสเนนปิ สา อตฺตมนา โหติ. ตตฺร เจ อานนฺโท ธมฺมํ ภาสติ, ภาสิเตนปิ สา อตฺตมนา โหติ. อติตฺตาว ภิกฺขเว ภิกฺขุปริสา โหติ, อถ อานนฺโท ตุณฺหี ภวติ.}

\prosewithclosermidruled{สเจ ภิกฺขเว ภิกฺขุนิ\-ปริสา ฯเปฯ. สเจ ภิกฺขเว อุปาสกปริสา ฯเปฯ. สเจ ภิกฺขเว อุปาสิกาปริสา อานนฺทํ ทสฺสนาย อุปสงฺกมติ, ทสฺสเนนปิ สา อตฺตมนา โหติ. ตตฺร เจ อานนฺโท ธมฺมํ ภาสติ, ภาสิเตนปิ สา อตฺตมนา โหติ. อติตฺตาว ภิกฺขเว อุปาสิกาปริสา โหติ, อถ อานนฺโท ตุณฺหี ภวติ. อิเม โข ภิกฺขเว จตฺตาโร อจฺฉริยา อพฺภุตา ธมฺมา อานนฺเทติ.  ทสมํ.}{ภยวคฺโค ตติโย.}
\csromancenter{\textbf{ตสฺสุทฺทานํ}}
\setlength{\csromangathapagewidth}{0pt}
\setlength{\csromangathawakwidth}{0pt}
\csromangathameasuregroup{อตฺตานุวาทอูมิ จ, \\ เมตฺตา เทฺว จ อจฺฉริยา,}{\csromangathabat{อตฺตานุวาทอูมิ จ,}{เทฺว จ นานา เทฺว จ โหนฺติ.} \\ \csromangathabat{เมตฺตา เทฺว จ อจฺฉริยา,}{อปรา จ ตถา ทุเวติ.}}
\csromangathasetleft{อตฺตานุวาทอูมิ จ, \\ เมตฺตา เทฺว จ อจฺฉริยา,}
\csromangathagroup{\csromangathastanza{\csromangathabat{อตฺตานุวาทอูมิ จ,}{เทฺว จ นานา เทฺว จ โหนฺติ.} \\ \csromangathabat{เมตฺตา เทฺว จ อจฺฉริยา,}{อปรา จ ตถา ทุเวติ.}}}

\csromanmatikaanchor{mtk.177}
\csromantocmarknopage{section}{(14) 4. ปุคฺคลวคฺค}{mtk.177}
\bookmark[dest={mtk.177},level=1]{(14) 4. ปุคฺคลวคฺค}
\csromanmarkh{(14)}\csromanmarkhii{4. ปุคฺคลวคฺค}\csromanheadernomark{1}{\textbf{(14) 4. ปุคฺคลวคฺค}}
\csromanmatikaanchor{mtk.178}
\csromantocmarkref{subsection}{1. สํโยชนสุตฺต}{mtk.178}
\bookmark[dest={mtk.178},level=2]{1. สํโยชนสุตฺต}
\csromanheader{2}{1. \textbf{สํโยชนสุตฺต}}
\proseitem{131}{จตฺตาโรเม ภิกฺขเว ปุคฺคลา สนฺโต สํวิชฺชมานา โลกสฺมิํ. กตเม จตฺตาโร? อิธ ภิกฺขเว เอกจฺจสฺส ปุคฺคลสฺส โอรมฺภาคิยานิ สํโยชนานิ อปฺปหีนานิ โหนฺติ, อุปปตฺติ\-ปฏิลาภิยานิ สํโยชนานิ อปฺปหีนานิ โหนฺติ, ภวปฏิลาภิยานิ สํโยชนานิ อปฺปหีนานิ โหนฺติ.}

\prose{อิธ ปน ภิกฺขเว เอกจฺจสฺส ปุคฺคลสฺส โอรมฺภาคิยานิ สํโยชนานิ ปหีนานิ โหนฺติ, อุปปตฺติ\-ปฏิลาภิยานิ สํโยชนานิ อปฺปหีนานิ โหนฺติ, ภวปฏิลาภิยานิ สํโยชนานิ อปฺปหีนานิ โหนฺติ.}

\chapterheadpageread{(14) 4. ปุคฺคลวคฺค}
\prose{\csromanfolio{451}อิธ ปน ภิกฺขเว เอกจฺจสฺส ปุคฺคลสฺส โอรมฺภาคิยานิ สํโยชนานิ ปหีนานิ โหนฺติ, อุปปตฺติ\-ปฏิลาภิยานิ สํโยชนานิ ปหีนานิ โหนฺติ, ภวปฏิลาภิยานิ สํโยชนานิ อปฺปหีนานิ โหนฺติ.}

\prose{อิธ ปน ภิกฺขเว เอกจฺจสฺส ปุคฺคลสฺส โอรมฺภาคิยานิ สํโยชนานิ ปหีนานิ โหนฺติ, อุปปตฺติ\-ปฏิลาภิยานิ สํโยชนานิ ปหีนานิ โหนฺติ, ภวปฏิลาภิยานิ สํโยชนานิ ปหีนานิ โหนฺติ.}

\prose{กตมสฺส ภิกฺขเว ปุคฺคลสฺส โอรมฺภาคิยานิ สํโยชนานิ อปฺปหีนานิ อุปปตฺติ\-ปฏิลาภิยานิ สํโยชนานิ อปฺปหีนานิ ภวปฏิลาภิยานิ สํโยชนานิ อปฺปหีนานิ? สกทาคามิสฺส. อิมสฺส โข ภิกฺขเว ปุคฺคลสฺส โอรมฺภาคิยานิ สํโยชนานิ อปฺปหีนานิ อุปปตฺติ\-ปฏิลาภิยานิ สํโยชนานิ อปฺปหีนานิ ภวปฏิลาภิยานิ สํโยชนานิ อปฺปหีนานิ.}

\prose{กตมสฺส ภิกฺขเว ปุคฺคลสฺส โอรมฺภาคิยานิ สํโยชนานิ ปหีนานิ อุปปตฺติ\-ปฏิลาภิยานิ สํโยชนานิ อปฺปหีนานิ ภวปฏิลาภิยานิ สํโยชนานิ อปฺปหีนานิ? อุทฺธํโสตสฺส อกนิฏฺฐ\-คามิโน. อิมสฺส โข ภิกฺขเว ปุคฺคลสฺส โอรมฺภาคิยานิ สํโยชนานิ ปหีนานิ อุปปตฺติ\-ปฏิลาภิยานิ สํโยชนานิ อปฺปหีนานิ ภวปฏิลาภิยานิ สํโยชนานิ อปฺปหีนานิ.}

\prose{กตมสฺส ภิกฺขเว ปุคฺคลสฺส โอรมฺภาคิยานิ สํโยชนานิ ปหีนานิ อุปปตฺติ\-ปฏิลาภิยานิ สํโยชนานิ ปหีนานิ ภวปฏิลาภิยานิ สํโยชนานิ อปฺปหีนานิ? อนฺตรา\-ปรินิพฺพายิสฺส. อิมสฺส โข ภิกฺขเว ปุคฺคลสฺส โอรมฺภาคิยานิ สํโยชนานิ ปหีนานิ อุปปตฺติ\-ปฏิลาภิยานิ สํโยชนานิ ปหีนานิ ภวปฏิลาภิยานิ สํโยชนานิ อปฺปหีนานิ.}

\prose{กตมสฺส ภิกฺขเว ปุคฺคลสฺส โอรมฺภาคิยานิ สํโยชนานิ ปหีนานิ อุปปตฺติ\-ปฏิลาภิยานิ สํโยชนานิ ปหีนานิ ภวปฏิลาภิยานิ สํโยชนานิ ปหีนานิ? อรหโต. อิมสฺส โข ภิกฺขเว ปุคฺคลสฺส โอรมฺภาคิยานิ สํโยชนานิ ปหีนานิ อุปปตฺติ\-ปฏิลาภิยานิ สํโยชนานิ ปหีนานิ ภวปฏิลาภิยานิ สํโยชนานิ ปหีนานิ. อิเม โข ภิกฺขเว จตฺตาโร ปุคฺคลา สนฺโต สํวิชฺชมานา โลกสฺมินฺติ.  ปฐมํ.}

\csromanpalirecto
\csromanmatikaanchor{mtk.179}
\csromantocmarkref{subsection}{2. ปฏิภานสุตฺต}{mtk.179}
\bookmark[dest={mtk.179},level=2]{2. ปฏิภานสุตฺต}
\csromanmarkh{จตุกฺกนิปาต\-ปาฬิ}\csromanmarkhii{2. ปฏิภานสุตฺต}\csromanheadernomark{2}{จตุกฺกนิปาต\-ปาฬิ 2. ปฏิภานสุตฺต}
\proseitem{132}{\csromanfolio{452}จตฺตาโรเม ภิกฺขเว ปุคฺคลา สนฺโต สํวิชฺชมานา โลกสฺมิํ. กตเม จตฺตาโร? ยุตฺต\-ปฺปฏิภาโน โน มุตฺต\-ปฺปฏิภาโน, มุตฺต\-ปฺปฏิภาโน โน ยุตฺต\-ปฺปฏิภาโน, ยุตฺต\-ปฺปฏิภาโน จ มุตฺต\-ปฺปฏิภาโน จ, เนว ยุตฺต\-ปฺปฏิภาโน น มุตฺต\-ปฺปฏิภาโน. อิเม โข ภิกฺขเว จตฺตาโร ปุคฺคลา สนฺโต สํวิชฺชมานา โลกสฺมินฺติ\csromansharedfootnote{0eaae954995ad28b}{อภิ 3. 147-8 ปิฏฺเฐสุปิ.}.  ทุติยํ.}

\csromanmatikaanchor{mtk.180}
\csromantocmarkref{subsection}{3. อุคฺฆฏิตญฺญูสุตฺต}{mtk.180}
\bookmark[dest={mtk.180},level=2]{3. อุคฺฆฏิตญฺญูสุตฺต}
\csromanheader{2}{3. อุคฺฆฏิตญฺญู\-สุตฺต}
\proseitem{133}{จตฺตาโรเม ภิกฺขเว ปุคฺคลา สนฺโต สํวิชฺชมานา โลกสฺมิํ. กตเม จตฺตาโร? อุคฺฆฏิตญฺญู วิปญฺจิตญฺญู เนยฺโย ปทปรโม. อิเม โข ภิกฺขเว จตฺตาโร ปุคฺคลา สนฺโต สํวิชฺชมานา โลกสฺมินฺติ\csromansharedfootnote{0eaae954995ad28b}{อภิ 3. 147-8 ปิฏฺเฐสุปิ.}. ตติยํ.}

\csromanmatikaanchor{mtk.181}
\csromantocmarkref{subsection}{4. อุฏฺฐานผลสุตฺต}{mtk.181}
\bookmark[dest={mtk.181},level=2]{4. อุฏฺฐานผลสุตฺต}
\csromanheader{2}{4. อุฏฺฐาน\-ผล\-สุตฺต}
\proseitem{134}{จตฺตาโรเม ภิกฺขเว ปุคฺคลา สนฺโต สํวิชฺชมานา โลกสฺมิํ. กตเม จตฺตาโร? อุฏฺฐาน\-ผลูปชีวี น กมฺม\-ผลูปชีวี, กมฺม\-ผลูปชีวี น อุฏฺฐาน\-ผลูปชีวี, อุฏฺฐาน\-ผลูปชีวี เจว กมฺม\-ผลูปชีวี จ, เนว อุฏฺฐาน\-ผลูปชีวี น กมฺม\-ผลูปชีวี. อิเม โข ภิกฺขเว จตฺตาโร ปุคฺคลา สนฺโต สํวิชฺชมานา โลกสฺมินฺติ\csromansharedfootnote{e31d13694865af0e}{อภิ 3. 158 ปิฏฺเฐปิ.}.  จตุตฺถํ.}

\csromanmatikaanchor{mtk.182}
\csromantocmarkref{subsection}{5. สาวชฺชสุตฺต}{mtk.182}
\bookmark[dest={mtk.182},level=2]{5. สาวชฺชสุตฺต}
\csromanheader{2}{5. \textbf{สาวชฺชสุตฺต}}
\proseitem{135}{จตฺตาโรเม ภิกฺขเว ปุคฺคลา\csromansharedfootnote{c30dfe92d7142058}{อภิ 3. 146 ปิฏฺเฐปิ.} สนฺโต สํวิชฺชมานา โลกสฺมิํ. กตเม จตฺตาโร? สาวชฺโช วชฺชพหุโล อปฺปวชฺโช อนวชฺโช.}

\prose{กถญฺจ ภิกฺขเว ปุคฺคโล สาวชฺโช โหติ? อิธ ภิกฺขเว เอกจฺโจ ปุคฺคโล สาวชฺเชน กายกมฺเมน สมนฺนาคโต โหติ, สาวชฺเชน วจีกมฺเมน สมนฺนาคโต โหติ, สาวชฺเชน มโนกมฺเมน สมนฺนาคโต โหติ. เอวํ โข ภิกฺขเว ปุคฺคโล สาวชฺโช โหติ.}

\prose{กถญฺจ ภิกฺขเว ปุคฺคโล วชฺชพหุโล โหติ? อิธ ภิกฺขเว เอกจฺโจ ปุคฺคโล สาวชฺเชน พหุลํ กายกมฺเมน สมนฺนาคโต โหติ อปฺปํ อนวชฺเชน, สาวชฺเชน พหุลํ วจีกมฺเมน สมนฺนาคโต โหติ อปฺปํ อนวชฺเชน, สาวชฺเชน พหุลํ มโนกมฺเมน สมนฺนาคโต โหติ อปฺปํ อนวชฺเชน. เอวํ โข ภิกฺขเว ปุคฺคโล วชฺชพหุโล โหติ.}

\chapterheadpageread{(14) 4. ปุคฺคลวคฺค}
\csromanmatikaanchor{mtk.183}
\csromantocmarkref{subsection}{6-7. สีลสุตฺตทฺวยํ}{mtk.183}
\bookmark[dest={mtk.183},level=2]{6-7. สีลสุตฺตทฺวยํ}
\prose{\csromanfolio{453}กถญฺจ ภิกฺขเว ปุคฺคโล อปฺปวชฺโช โหติ? อิธ ภิกฺขเว เอกจฺโจ ปุคฺคโล อนวชฺเชน พหุลํ กายกมฺเมน สมนฺนาคโต โหติ อปฺปํ สาวชฺเชน, อนวชฺเชน พหุลํ วจีกมฺเมน สมนฺนาคโต โหติ อปฺปํ สาวชฺเชน, อนวชฺเชน พหุลํ มโนกมฺเมน สมนฺนาคโต โหติ อปฺปํ สาวชฺเชน. เอวํ โข ภิกฺขเว ปุคฺคโล อปฺปวชฺโช โหติ.}

\prose{กถญฺจ ภิกฺขเว ปุคฺคโล อนวชฺโช โหติ? อิธ ภิกฺขเว เอกจฺโจ ปุคฺคโล อนวชฺเชน กายกมฺเมน สมนฺนาคโต โหติ, อนวชฺเชน วจีกมฺเมน สมนฺนาคโต โหติ, อนวชฺเชน มโนกมฺเมน สมนฺนาคโต โหติ. เอวํ โข ภิกฺขเว ปุคฺคโล อนวชฺโช โหติ. อิเม โข ภิกฺขเว จตฺตาโร ปุคฺคลา สนฺโต สํวิชฺชมานา โลกสฺมินฺติ.  ปญฺจมํ.}

\csromanheader{2}{6. ปฐม\-สีล\-สุตฺต}
\proseitem{136}{จตฺตาโรเม ภิกฺขเว ปุคฺคลา สนฺโต สํวิชฺชมานา โลกสฺมิํ. กตเม จตฺตาโร? อิธ ภิกฺขเว เอกจฺโจ ปุคฺคโล สีเลสุ น ปริปูรการี โหติ สมาธิสฺมิํ น ปริปูรการี ปญฺญาย น ปริปูรการี.}

\prose{อิธ ปน ภิกฺขเว เอกจฺโจ ปุคฺคโล สีเลสุ ปริปูรการี โหติ สมาธิสฺมิํ น ปริปูรการี ปญฺญาย น ปริปูรการี.}

\prose{อิธ ปน ภิกฺขเว เอกจฺโจ ปุคฺคโล สีเลสุ ปริปูรการี โหติ สมาธิสฺมิํ ปริปูรการี ปญฺญาย น ปริปูรการี.}

\prose{อิธ ปน ภิกฺขเว เอกจฺโจ ปุคฺคโล สีเลสุ ปริปูรการี โหติ สมาธิสฺมิํ ปริปูรการี ปญฺญาย ปริปูรการี. อิเม โข ภิกฺขเว จตฺตาโร ปุคฺคลา สนฺโต สํวิชฺชามานา โลกสฺมินฺติ.  ฉฏฺฐํ.}

\csromanheader{2}{7. ทุติย\-สีล\-สุตฺต}
\proseitem{137}{จตฺตาโรเม ภิกฺขเว ปุคฺคลา สนฺโต สํวิชฺชมานา โลกสฺมิํ. กตเม จตฺตาโร? อิธ ภิกฺขเว เอกจฺโจ ปุคฺคโล น สีลครุ โหติ น สีลาธิปเตยฺโย, น สมาธิครุ โหติ น สมาธา\-ธิปเตยฺโย, น ปญฺญาครุ โหติ น ปญฺญา\-ธิปเตยฺโย.}

\prose{\csromanfolio{454}อิธ ปน ภิกฺขเว เอกจฺโจ ปุคฺคโล สีลครุ โหติ สีลาธิปเตยฺโย, น สมาธิครุ โหติ น สมาธา\-ธิปเตยฺโย, น ปญฺญาครุ โหติ น ปญฺญา\-ธิปเตยฺโย.}

\prose{อิธ ปน ภิกฺขเว เอกจฺโจ ปุคฺคโล สีลครุ โหติ สีลาธิปเตยฺโย, สมาธิครุ โหติ สมาธา\-ธิปเตยฺโย, น ปญฺญาครุ โหติ น ปญฺญา\-ธิปเตยฺโย.}

\prose{อิธ ปน ภิกฺขเว เอกจฺโจ ปุคฺคโล สีลครุ โหติ สีลาธิปเตยฺโย, สมาธิครุ โหติ สมาธา\-ธิปเตยฺโย, ปญฺญาครุ โหติ ปญฺญา\-ธิปเตยฺโย. อิเม โข ภิกฺขเว จตฺตาโร ปุคฺคลา สนฺโต สํวิชฺชมานา โลกสฺมินฺติ.  สตฺตมํ.}

\csromanmatikaanchor{mtk.184}
\csromantocmarkref{subsection}{8. นิกฏฺฐสุตฺต}{mtk.184}
\bookmark[dest={mtk.184},level=2]{8. นิกฏฺฐสุตฺต}
\csromanheader{2}{8. \textbf{นิกฏฺฐสุตฺต}}
\proseitem{138}{จตฺตาโรเม ภิกฺขเว ปุคฺคลา สนฺโต สํวิชฺชมานา โลกสฺมิํ. กตเม จตฺตาโร? นิกฏฺฐกาโย อนิกฏฺฐจิตฺโต, อนิกฏฺฐกาโย นิกฏฺฐจิตฺโต, อนิกฏฺฐกาโย จ อนิกฏฺฐจิตฺโต จ, นิกฏฺฐกาโย จ นิกฏฺฐจิตฺโต จ.}

\prose{กถญฺจ ภิกฺขเว ปุคฺคโล นิกฏฺฐกาโย โหติ อนิกฏฺฐจิตฺโต? อิธ ภิกฺขเว เอกจฺโจ ปุคฺคโล อรญฺญ\-วนปตฺถานิ\csromansharedfootnote{d1a425e0757127c1}{อรญฺเญ วนปตฺถานิ (สี, อิ)} ปนฺตานิ เสนาสนานิ ปฏิเสวติ, โส ตตฺถ กามวิตกฺกมฺปิ วิตกฺเกติ, พฺยาปาท\-วิตกฺกมฺปิ วิตกฺเกติ, วิหิํสา\-วิตกฺกมฺปิ วิตกฺเกติ. เอวํ โข ภิกฺขเว ปุคฺคโล นิกฏฺฐกาโย โหติ อนิกฏฺฐจิตฺโต.}

\prose{กถญฺจ ภิกฺขเว ปุคฺคโล อนิกฏฺฐกาโย โหติ นิกฏฺฐจิตฺโต? อิธ ภิกฺขเว เอกจฺโจ ปุคฺคโล นเหว โข อรญฺญ\-วนปตฺถานิ ปนฺตานิ เสนาสนานิ ปฏิเสวติ, โส ตตฺถ เนกฺขมฺม\-วิตกฺกมฺปิ วิตกฺเกติ, อพฺยาปาทวิตกฺกมฺปิ วิตกฺเกติ, อวิหิํสา\-วิตกฺกมฺปิ วิตกฺเกติ. เอวํ โข ภิกฺขเว ปุคฺคโล อนิกฏฺฐกาโย โหติ นิกฏฺฐจิตฺโต.}

\prose{กถญฺจ ภิกฺขเว ปุคฺคโล อนิกฏฺฐกาโย จ โหติ อนิกฏฺฐจิตฺโต จ? อิธ ภิกฺขเว เอกจฺโจ ปุคฺคโล น เหว โข อรญฺญ\-วนปตฺถานิ ปนฺตานิ}

\vspace{\tipitakaparskip}
\csromancenter{(14) 4. ปุคฺคลวคฺค เสนาสนานิ ปฏิเสวติ, โส ตตฺถ กามวิตกฺกมฺปิ วิตกฺเกติ, พฺยาปาท\-วิตกฺกมฺปิ วิตกฺเกติ, วิหิํสา\-วิตกฺกมฺปิ วิตกฺเกติ. เอวํ โข ภิกฺขเว ปุคฺคโล อนิกฏฺฐกาโย จ โหติ อนิกฏฺฐจิตฺโต จ.}
\csromanmatikaanchor{mtk.185}
\csromantocmarkref{subsection}{9-10. ธมฺมกถิกสุตฺตาทิ}{mtk.185}
\bookmark[dest={mtk.185},level=2]{9-10. ธมฺมกถิกสุตฺตาทิ}
\prose{\csromanfolio{455}กถญฺจ ภิกฺขเว ปุคฺคโล นิกฏฺฐกาโย จ โหติ นิกฏฺฐจิตฺโต จ? อิธ ภิกฺขเว เอกจฺโจ ปุคฺคโล อรญฺญ\-วนปตฺถานิ ปนฺตานิ เสนาสนานิ ปฏิเสวติ, โส ตตฺถ เนกฺขมฺม\-วิตกฺกมฺปิ วิตกฺเกติ, อพฺยาปาทวิตกฺกมฺปิ วิตกฺเกติ, อวิหิํสา\-วิตกฺกมฺปิ วิตกฺเกติ. เอวํ โข ภิกฺกเว ปุคฺคโล นิกฏฺฐกาโย จ โหติ นิกฏฺฐจิตฺโต จ. อิเม โข ภิกฺขเว จตฺตาโร ปุคฺคลา สนฺโต สํวิชฺชมานา โลกสฺมินฺติ.  อฏฺฐมํ.}

\csromanheader{2}{9. ธมฺมกถิก\-สุตฺต}
\proseitem{139}{จตฺตาโรเม ภิกฺขเว ธมฺมกถิกา. กตเม จตฺตาโร? อิธ ภิกฺขเว เอกจฺโจ ธมฺมกถิโก อปฺปญฺจ ภาสติ อสหิตญฺจ, ปริสา จสฺส\csromansharedfootnote{43d2b821d8e0c8d6}{ปริสา จ (สี, สฺยา, กํ, อิ) อภิ 3. 148 ปิฏฺเฐปิ.} น กุสลา โหติ สหิตา\-สหิตสฺส. เอวรูโป ภิกฺขเว ธมฺมกถิโก เอวรูปาย ปริสาย ธมฺมกถิโกเตฺวว สงฺขํ คจฺฉติ.}

\prose{อิธ ปน ภิกฺขเว เอกจฺโจ ธมฺมกถิโก อปฺปญฺจ ภาสติ สหิตญฺจ, ปริสา จสฺส กุสลา โหติ สหิตา\-สหิตสฺส. เอวรูโป ภิกฺขเว ธมฺมกถิโก เอวรูปาย ปริสาย ธมฺมกถิโกเตฺวว สงฺขํ คจฺฉติ.}

\prose{อิธ ปน ภิกฺขเว เอกจฺโจ ธมฺมกถิโก พหุญฺจ ภาสติ อสหิตญฺจ, ปริสา จสฺส น กุสลา โหติ สหิตา\-สหิตสฺส. เอวรูโป ภิกฺขเว ธมฺมกถิโก เอวรูปาย ปริสาย ธมฺมกถิโกเตฺวว สงฺขํ คจฺฉติ.}

\prose{อิธ ปน ภิกฺขเว เอกจฺโจ ธมฺมกถิโก พหุญฺจ ภาสติ สหิตญฺจ, ปริสา จสฺส กุสลา โหติ สหิตา\-สหิตสฺส. เอวรูโป ภิกฺขเว ธมฺมกถิโก เอวรูปาย ปริสาย ธมฺมกถิโกเตฺวว สงฺขํ คจฺฉติ. อิเม โข ภิกฺขเว จตฺตาโร ธมฺมกถิกาติ.  นวมํ.}

\csromanheader{2}{10. \textbf{วาทีสุตฺต}}
\csromanmatikaanchor{mtk.186}
\csromanmatikaanchor{mtk.187}
\csromantocmarknopage{section}{(15) 5. อาภาวคฺค}{mtk.186}
\bookmark[dest={mtk.186},level=1]{(15) 5. อาภาวคฺค}
\csromantocmarkref{subsection}{1-5. อาภาสุตฺตาทีนิ}{mtk.187}
\bookmark[dest={mtk.187},level=2]{1-5. อาภาสุตฺตาทีนิ}
\proseitemwithclosermidruled{140}{จตฺตาโรเม ภิกฺขเว วาที. กตเม จตฺตาโร? อตฺถิ ภิกฺขเว วาที อตฺถโต ปริยาทานํ คจฺฉติ โน พฺยญฺชนโต, อตฺถิ ภิกฺขเว วาที \csromanfolio{456}พฺยญฺชนโต ปริยาทานํ คจฺฉติ โน อตฺถโต, อตฺถิ ภิกฺขเว วาที อตฺถโต จ พฺยญฺชนโต จ ปริยาทานํ คจฺฉติ, อตฺถิ ภิกฺขเว วาที เนวตฺถโต โน พฺยญฺชนโต ปริยาทานํ คจฺฉติ. อิเม โข ภิกฺขเว จตฺตาโร วาที. อฏฺฐานเมตํ ภิกฺขเว อนวกาโส ยํ จตูหิ ปฏิสมฺภิทาหิ สมนฺนาคโต\csromansharedfootnote{db0fa1fff094c1e2}{สมนฺนาคโต ภิกฺขุ (สี, สฺยา, กํ)} อตฺถโต วา พฺยญฺชนโต วา ปริยาทานํ คจฺเฉยฺยาติ.  ทสมํ.}{ปุคฺคลวคฺโค จตุตฺโถ.}
\csromancenter{\textbf{ตสฺสุทฺทานํ}}
\setlength{\csromangathapagewidth}{0pt}
\setlength{\csromangathawakwidth}{0pt}
\csromangathameasuregroup{สํโยชนํ ปฏิภาโน, \\ สาวชฺโช เทฺว จ สีลานิ,}{\csromangathabat{สํโยชนํ ปฏิภาโน,}{อุคฺฆฏิตญฺญุ อุฏฺฐานํ.} \\ \csromangathabat{สาวชฺโช เทฺว จ สีลานิ,}{นิกฏฺฐ ธมฺม วาที จาติ.}}
\csromangathasetleft{สํโยชนํ ปฏิภาโน, \\ สาวชฺโช เทฺว จ สีลานิ,}
\csromangathagroup{\csromangathastanza{\csromangathabat{สํโยชนํ ปฏิภาโน,}{อุคฺฆฏิตญฺญุ อุฏฺฐานํ.} \\ \csromangathabat{สาวชฺโช เทฺว จ สีลานิ,}{นิกฏฺฐ ธมฺม วาที จาติ.}}}

\csromanmarkh{(15)}\csromanmarkhii{5. อาภาวคฺค}\csromanheadernomark{1}{\textbf{(15) 5. อาภาวคฺค}}
\csromanheader{2}{1. \textbf{อาภาสุตฺต}}
\proseitem{141}{จตสฺโส อิมา ภิกฺขเว อาภา. กตมา จตสฺโส? จนฺทาภา สูริยาภา อคฺคาภา ปญฺญาภา. อิมา โข ภิกฺขเว จตสฺโส อาภา. เอตทคฺคํ ภิกฺขเว อิมาสํ จตุนฺนํ\csromansharedfootnote{350e443039a761b4}{จตสฺสนฺนํ (สฺยา, กํ) สทฺทนีติปทมาลา ปสฺสิตพฺพา.} อาภานํ ยทิทํ ปญฺญาภาติ.  ปฐมํ.}

\csromanheader{2}{2. \textbf{ปภาสุตฺต}}
\proseitem{142}{จตสฺโส อิมา ภิกฺขเว ปภา. กตมา จตสฺโส? จนฺทปฺปภา สูริยปฺปภา อคฺคิปฺปภา ปญฺญาปภา. อิมา โข ภิกฺขเว จตสฺโส ปภา. เอตทคฺคํ ภิกฺขเว อิมาสํ จตุนฺนํ ปภานํ ยทิทํ ปญฺญาปภาติ.  ทุติยํ.}

\csromanheader{2}{3. \textbf{อาโลกสุตฺต}}
\proseitem{143}{จตฺตาโรเม ภิกฺขเว อาโลกา. กตเม จตฺตาโร? จนฺทาโลโก สูริยาโลโก อคฺคาโลโก ปญฺญาโลโก. อิเม โข ภิกฺขเว จตฺตาโร อาโลกา. เอตทคฺคํ ภิกฺขเว อิเมสํ จตุนฺนํ อาโลกานํ ยทิทํ ปญฺญาโลโกติ.  ตติยํ.}

\chapterheadpageread{(15) 5. อาภาวคฺค \textbf{4. โอภาสสุตฺต}}
\csromanmatikaanchor{mtk.188}
\csromantocmarkref{subsection}{6-7. กาลสุตฺตทฺวยํ}{mtk.188}
\bookmark[dest={mtk.188},level=2]{6-7. กาลสุตฺตทฺวยํ}
\proseitem{144}{\csromanfolio{457}จตฺตาโรเม ภิกฺขเว โอภาสา. กตเม จตฺตาโร? จนฺโทภาโส สูริโยภาโส อคฺโคภาโส ปญฺโญภาโส. อิเม โข ภิกฺขเว จตฺตาโร โอภาสา. เอตทคฺคํ ภิกฺขเว อิเมสํ จตุนฺนํ โอภาสานํ ยทิทํ ปญฺโญภาโสติ.  จตุตฺถํ.}

\csromanheader{2}{5. \textbf{ปชฺโชตสุตฺต}}
\proseitem{145}{จตฺตาโรเม ภิกฺขเว ปชฺโชตา. กตเม จตฺตาโร? จนฺทปชฺโชโต สูริยปชฺโชโต อคฺคิปชฺโชโต ปญฺญาปชฺโชโต. อิเม โข ภิกฺขเว จตฺตาโร ปชฺโชตา. เอตทคฺคํ ภิกฺขเว อิเมสํ จตุนฺนํ ปชฺโชตานํ ยทิทํ ปญฺญาปชฺโชโตติ.  ปญฺจมํ.}

\csromanheader{2}{6. ปฐม\-กาล\-สุตฺต}
\proseitem{146}{จตฺตาโรเม ภิกฺขเว กาลา. กตเม จตฺตาโร? กาเลน ธมฺมสฺสวนํ, กาเลน ธมฺมสากจฺฉา, กาเลน สมฺมสนา\csromansharedfootnote{a1e9afab19da7331}{กาเลน สมโถ (สี, สฺยา, กํ, อิ)}, กาเลน วิปสฺสนา. อิเม โข ภิกฺขเว จตฺตาโร กาลาติ.  ฉฏฺฐํ.}

\csromanheader{2}{7. ทุติย\-กาล\-สุตฺต}
\proseitem{147}{จตฺตาโรเม ภิกฺขเว กาลา สมฺมา ภาวิยมานา สมฺมา อนุปริ\-วตฺติย\-มานา อนุปุพฺเพน อาสวานํ ขยํ ปาเปนฺติ. กตเม จตฺตาโร? กาเลน ธมฺมสฺสวนํ, กาเลน ธมฺมสากจฺฉา, กาเลน สมฺมสนา, กาเลน วิปสฺสนา. อิเม โข ภิกฺขเว จตฺตาโร กาลา สมฺมา ภาวิยมานา สมฺมา อนุปริ\-วตฺติย\-มานา อนุปุพฺเพน อาสวานํ ขยํ ปาเปนฺติ.}

\csromanmatikaanchor{mtk.189}
\csromantocmarkref{subsection}{8-10. ทุจฺจริตสุตฺตาทีนิ}{mtk.189}
\bookmark[dest={mtk.189},level=2]{8-10. ทุจฺจริตสุตฺตาทีนิ}
\prose{เสยฺยถาปิ ภิกฺขเว อุปริปพฺพเต ถุลฺล\-ผุสิตเก เทเว วสฺสนฺเต ตํ อุทกํ ยถานินฺนํ ปวตฺตมานํ ปพฺพต\-กนฺทร\-ปทร\-สาขา ปริปูเรติ, ปพฺพต\-กนฺทร\-ปทร\-สาขา ปริปูรา กุโสพฺเภ ปริปูเรนฺติ, กุโสพฺภา ปริปูรา มหาโสพฺเภ ปริปูเรนฺติ, มหาโสพฺภา ปริปูรา กุนฺนทิโย ปริปูเรนฺติ, \csromanfolio{458}กุนฺนทิโย ปริปูรา มหานทิโย ปริปูเรนฺติ, มหานทิโย ปริปูรา สมุทฺทํ\csromansharedfootnote{03367a49ccb9eadb}{สมุทฺทํ สาครํ (สี, อิ, ก), สมุทฺทสาครํ (สฺยา, กํ)} ปริปูเรนฺติ. เอวเมวํ โข ภิกฺขเว อิเม จตฺตาโร กาลา สมฺมา ภาวิยมานา สมฺมา อนุปริ\-วตฺติย\-มานา อนุปุพฺเพน อาสวานํ ขยํ ปาเปนฺตีติ.  สตฺตมํ.}

\csromanheader{2}{8. \textbf{ทุจฺจริตสุตฺต}}
\proseitem{148}{จตฺตาริมานิ ภิกฺขเว วจีทุจฺจริตานิ. กตมานิ จตฺตาริ? มุสาวาโท ปิสุณา วาจา ผรุสา วาจา สมฺผปฺปลาโป. อิมานิ โข ภิกฺขเว จตฺตาริ วจีทุจฺจริตานีติ.  อฏฺฐมํ.}

\csromanheader{2}{9. \textbf{สุจริตสุตฺต}}
\proseitem{149}{จตฺตาริมานิ ภิกฺขเว วจีสุจริตานิ. กตมานิ จตฺตาริ? สจฺจวาจา อปิสุณา วาจา สณฺหา วาจา มนฺตภาสา. อิมานิ โข ภิกฺขเว จตฺตาริ วจีสุจริตานีติ.  นวมํ.}

\csromanheader{2}{10. \textbf{สารสุตฺต}}
\proseitemwithclosermidruled{150}{จตฺตาโรเม ภิกฺขเว สารา. กตเม จตฺตาโร? สีลสาโร สมาธิสาโร ปญฺญาสาโร วิมุตฺติสาโร. อิเม โข ภิกฺขเว จตฺตาโร สาราติ.  ทสมํ.}{อาภาวคฺโค ปญฺจโม.}
\csromancenter{\textbf{ตสฺสุทฺทานํ}}
\setlength{\csromangathapagewidth}{0pt}
\setlength{\csromangathawakwidth}{0pt}
\csromangathameasuregroup{อาภา ปภา จ อาโลกา, \\ เทฺว กาลา จริตา เทฺว จ,}{\csromangathabat{อาภา ปภา จ อาโลกา,}{โอภาสา เจว ปชฺโชตา.} \\ \csromangathabat{เทฺว กาลา จริตา เทฺว จ,}{โหนฺติ สาเรน เต ทสาติ.} \\ \textbf{ตติโย ปณฺณาสโก สมตฺโต.}}
\csromangathasetleft{อาภา ปภา จ อาโลกา, \\ เทฺว กาลา จริตา เทฺว จ,}
\csromangathagroup{\csromangathastanza{\csromangathabat{อาภา ปภา จ อาโลกา,}{โอภาสา เจว ปชฺโชตา.} \\ \csromangathabat{เทฺว กาลา จริตา เทฺว จ,}{โหนฺติ สาเรน เต ทสาติ.} \\ \textbf{ตติโย ปณฺณาสโก สมตฺโต.}}}

\csromanmatikaanchor{mtk.190}
\csromantocmarknopage{chapter}{4. จตุตฺถปณฺณาสก}{mtk.190}
\bookmark[dest={mtk.190},level=0]{4. จตุตฺถปณฺณาสก}
\chapterheadpageread{4. จตุตฺถ\-ปณฺณาสก}
\csromanmatikaanchor{mtk.191}
\csromantocmarknopage{section}{(16) 1. อินฺทฺริยวคฺค}{mtk.191}
\bookmark[dest={mtk.191},level=1]{(16) 1. อินฺทฺริยวคฺค}
\csromanmarkh{(16)}\csromanmarkhii{1. อินฺทฺริยวคฺค}\csromanheadernomark{1}{\textbf{(16) 1. อินฺทฺริยวคฺค}}
\csromanmatikaanchor{mtk.192}
\csromantocmarkref{subsection}{1. อินฺทฺริยสุตฺต}{mtk.192}
\bookmark[dest={mtk.192},level=2]{1. อินฺทฺริยสุตฺต}
\csromanheader{2}{1. \textbf{อินฺทฺริยสุตฺต}}
\csromanmatikaanchor{mtk.193}
\csromantocmarkref{subsection}{2-5. สทฺธาพลสุตฺตาทีนิ}{mtk.193}
\bookmark[dest={mtk.193},level=2]{2-5. สทฺธาพลสุตฺตาทีนิ}
\proseitem{151}{\csromanfolio{459}จตฺตาริมานิ ภิกฺขเว อินฺทฺริยานิ. กตมานิ จตฺตาริ? สทฺธินฺทฺริยํ วีริยินฺทฺริยํ สตินฺทฺริยํ สมาธินฺทฺริยํ. อิมานิ โข ภิกฺขเว จตฺตาริ อินฺทฺริยานีติ.  ปฐมํ.}

\csromanheader{2}{2. สทฺธาพล\-สุตฺต}
\proseitem{152}{จตฺตาริมานิ ภิกฺขเว พลานิ. กตมานิ จตฺตาริ? สทฺธาพลํ วีริยพลํ สติพลํ สมาธิพลํ. อิมานิ โข ภิกฺขเว จตฺตาริ พลานีติ.  ทุติยํ.}

\csromanheader{2}{3. \textbf{ปญฺญาพลสุตฺต}}
\proseitem{153}{จตฺตาริมานิ ภิกฺขเว พลานิ. กตมานิ จตฺตาริ? ปญฺญาพลํ วีริยพลํ อนวชฺชพลํ สงฺคหพลํ\csromansharedfootnote{cfd5a8b222e60ab3}{สงฺคาหพลํ (สี, สฺยา, กํ, อิ)}. อิมานิ โข ภิกฺขเว จตฺตาริ พลานีติ.  ตติยํ.}

\csromanheader{2}{4. \textbf{สติพลสุตฺต}}
\proseitem{154}{จตฺตาริมานิ ภิกฺขเว พลานิ. กตมานิ จตฺตาริ? สติพลํ สมาธิพลํ อนวชฺชพลํ สงฺคหพลํ. อิมานิ โข ภิกฺขเว จตฺตาริ พลานีติ.  จตุตฺถํ.}

\csromanheader{2}{5. ปฏิสงฺขานพล\-สุตฺต}
\proseitem{155}{จตฺตาริมานิ ภิกฺขเว พลานิ. กตมานิ จตฺตาริ? ปฏิสงฺขานพลํ ภาวนาพลํ อนวชฺชพลํ สงฺคหพลํ. อิมานิ โข ภิกฺขเว จตฺตาริ พลานีติ.  ปญฺจมํ.}

\csromanpalirecto
\csromanmatikaanchor{mtk.194}
\csromantocmarkref{subsection}{6. กปฺปสุตฺต}{mtk.194}
\bookmark[dest={mtk.194},level=2]{6. กปฺปสุตฺต}
\csromanmarkh{จตุกฺกนิปาต\-ปาฬิ}\csromanmarkhii{6. กปฺปสุตฺต}\csromanheadernomark{2}{จตุกฺกนิปาต\-ปาฬิ 6. กปฺปสุตฺต}
\proseitem{156}{\csromanfolio{460}จตฺตาริมานิ ภิกฺขเว กปฺปสฺส อสงฺเขฺยยฺยานิ. กตมานิ จตฺตาริ? ยทา ภิกฺขเว กปฺโป สํวฏฺฏติ, ตํ น สุกรํ สงฺขาตุํ “เอตฺตกานิ วสฺสานี”ติ วา “เอตฺตกานิ วสฺสสตานี”ติ วา “เอตฺตกานิ วสฺสสหสฺสานี”ติ วา “เอตฺตกานิ วสฺสสตสหสฺสานี”ติ วา.}

\prose{ยทา ภิกฺขเว กปฺโป สํวฏฺโฏ ติฏฺฐติ, ตํ น สุกรํ สงฺขาตุํ “เอตฺตกานิ วสฺสานี”ติ วา “เอตฺตกานิ วสฺสสตานี”ติ วา “เอตฺตกานิ วสฺสสหสฺสานี”ติ วา “เอตฺตกานิ วสฺสสตสหสฺสานี”ติ วา.}

\prose{ยทา ภิกฺขเว กปฺโป วิวฏฺฏติ, ตํ น สุกรํ สงฺขาตุํ “เอตฺตกานิ วสฺสานี”ติ วา “เอตฺตกานิ วสฺสสตานี”ติ วา “เอตฺตกานิ วสฺสสหสฺสานี”ติ วา “เอตฺตกานิ วสฺสสตสหสฺสานี”ติ วา.}

\prose{ยทา ภิกฺขเว กปฺโป วิวฏฺโฏ ติฏฺฐติ, ตํ น สุกรํ สงฺขาตุํ “เอตฺตกานิ วสฺสานี”ติ วา “เอตฺตกานิ วสฺสสตานี”ติ วา “เอตฺตกานิ วสฺสสหสฺสานี”ติ วา “เอตฺตกานิ วสฺสสตสหสฺสานี”ติ วา. อิมานิ โข ภิกฺขเว จตฺตาริ กปฺปสฺส อสงฺเขฺยยฺยานีติ.  ฉฏฺฐํ.}

\csromanmatikaanchor{mtk.195}
\csromantocmarkref{subsection}{7. โรคสุตฺต}{mtk.195}
\bookmark[dest={mtk.195},level=2]{7. โรคสุตฺต}
\csromanheader{2}{7. \textbf{โรคสุตฺต}}
\proseitem{157}{เทฺวเม ภิกฺขเว โรคา. กตเม เทฺว? กายิโก จ โรโค เจตสิโก จ โรโค. ทิสฺสนฺติ ภิกฺขเว สตฺตา กายิเกน โรเคน เอกมฺปิ วสฺสํ อาโรคฺยํ ปฏิชานมานา เทฺวปิ วสฺสานิ อาโรคฺยํ ปฏิชานมานา ตีณิปิ วสฺสานิ อาโรคฺยํ ปฏิชานมานา จตฺตาริปิ วสฺสานิ อาโรคฺยํ ปฏิชานมานา ปญฺจปิ วสฺสานิ อาโรคฺยํ ปฏิชานมานา ทสปิ วสฺสานิ อาโรคฺยํ ปฏิชานมานา วีสติปิ วสฺสานิ อาโรคฺยํ ปฏิชานมานา ติํสมฺปิ วสฺสานิ อาโรคฺยํ ปฏิชานมานา จตฺตารีสมฺปิ วสฺสานิ อาโรคฺยํ ปฏิชานมานา ปญฺญาสมฺปิ วสฺสานิ อาโรคฺยํ ปฏิชานมานา วสฺสสตมฺปิ ภิยฺโยปิ อาโรคฺยํ ปฏิชานมานา. เต ภิกฺขเว สตฺตา สุทุลฺลภา\csromansharedfootnote{207e2676174afe7e}{ทุลฺลภา (สี, สฺยา, กํ, อิ)} โลกสฺมิํ, เย เจตสิเกน โรเคน มุหุตฺตมฺปิ อาโรคฺยํ ปฏิชานนฺติ อญฺญตฺร ขีณาสเวหิ.}

\chapterheadpageread{(16) 1. อินฺทฺริยวคฺค}
\prose{\csromanfolio{461}จตฺตาโรเม ภิกฺขเว ปพฺพชิตสฺส โรคา. กตเม จตฺตาโร? อิธ ภิกฺขเว ภิกฺขุ มหิจฺโฉ โหติ วิฆาตวา อสนฺตุฏฺโฐ อิตรีตร\-จีวร\-ปิณฺฑปาต\-เสนาสน\-คิลานปฺปจฺจย\-เภสชฺช\-ปริกฺขาเรน, โส มหิจฺโฉ สมาโน วิฆาตวา อสนฺตุฏฺโฐ อิตรีตร\-จีวร\-ปิณฺฑปาต\-เสนาสน\-คิลานปฺปจฺจย\-เภสชฺช\-ปริกฺขาเรน ปาปิกํ อิจฺฉํ ปณิทหติ อนวญฺญ\-ปฺปฏิลาภาย ลาภ\-สกฺการ\-สิโลกปฺปฏิลาภาย. โส อุฏฺฐหติ ฆฏติ วายมติ อนวญฺญ\-ปฺปฏิลาภาย ลาภสกฺกรสิโลกปฺปฏิลาภาย. โส สงฺขาย กุลานิ อุปสงฺกมติ, สงฺขาย นิสีทติ, สงฺขาย ธมฺมํ ภาสติ, สงฺขาย อุจฺจารปสฺสาวํ สนฺธาเรติ. อิเม โข ภิกฺขเว จตฺตาโร ปพฺพชิตสฺส โรคา.}

\prose{ตสฺมาติห ภิกฺขเว เอวํ สิกฺขิตพฺพํ “น มหิจฺฉา ภวิสฺสาม วิฆาตวนฺโต อสนฺตุฏฺฐา อิตรีตรจีวรปิณฺฑปาตเสนาสนคิลานปฺปจฺจย- เภสชฺช\-ปริกฺขาเรน, น ปาปิกํ อิจฺฉํ ปณิทหิสฺสาม อนวญฺญ\-ปฺปฏิลาภาย ลาภ\-สกฺการ\-สิโลกปฺปฏิลาภาย, น อุฏฺฐหิสฺสาม น ฆเฏสฺสาม น วายมิสฺสาม อนวญฺญ\-ปฺปฏิลาภาย ลาภ\-สกฺการ\-สิโลกปฺปฏิลาภาย, ขมา ภวิสฺสาม สีตสฺส อุณฺหสฺส ชิฆจฺฉาย ปิปาสาย ฑํสมกส\-วาตา\-ตป\-สรีสป\-สมฺผสฺสานํ ทุรุตฺตานํ ทุราคตานํ วจนปถานํ, อุปฺปนฺนานํ สารีริกานํ เวทนานํ ทุกฺขานํ ติพฺพานํ ขรานํ กฏุกานํ อสาตานํ อมนาปานํ ปาณหรานํ อธิวาสก\-ชาติกา ภวิสฺสามา”ติ. เอวํ หิ โว ภิกฺขเว สิกฺขิตพฺพนฺติ.  สตฺตมํ.}

\csromanmatikaanchor{mtk.196}
\csromantocmarkref{subsection}{8. ปริหานิสุตฺต}{mtk.196}
\bookmark[dest={mtk.196},level=2]{8. ปริหานิสุตฺต}
\csromanheader{2}{8. \textbf{ปริหานิสุตฺต}}
\proseitem{158}{ตตฺร โข อายสฺมา สาริปุตฺโต ภิกฺขู อามนฺเตสิ “อาวุโส ภิกฺขเว”ติ. “อาวุโส”ติ โข เต ภิกฺขู อายสฺมโต สาริปุตฺตสฺส ปจฺจสฺโสสุํ, อายสฺมา สาริปุตฺโต เอตทโวจ\csromandash{}}

\prose{โย หิ โกจิ อาวุโส ภิกฺขู วา ภิกฺขุนี วา จตฺตาโร ธมฺเม อตฺตนิ สมนุปสฺสติ, นิฏฺฐเมตฺถ คนฺตพฺพํ ปริหายามิ กุสเลหิ ธมฺเมหิ, ปริหานเมตํ วุตฺตํ ภควตา. กตเม จตฺตาโร? ราค\-เวปุลฺลตฺตํ\csromansharedfootnote{47bf7c59338953b8}{ราคเวปุลฺลตํ (สี, สฺยา, กํ, อิ)} โทส\-เวปุลฺลตฺตํ โมห\-เวปุลฺลตฺตํ คมฺภีเรสุ โข ปนสฺส ฐานาฐาเนสุ \csromanfolio{462}ปญฺญาจกฺขุ น กมติ. โย หิ โกจิ อาวุโส ภิกฺขุ วา ภิกฺขุนี วา อิเม จตฺตาโร ธมฺเม อตฺตนิ สมนุปสฺสติ, นิฏฺฐเมตฺถ คนฺตพฺพํ ปริหายามิ กุสเลหิ ธมฺเมหิ, ปริหานเมตํ วุตฺตํ ภควตา.}

\prose{โย หิ โกจิ อาวุโส ภิกฺขุ วา ภิกฺขุนี วา จตฺตาโร ธมฺเม อตฺตนิ สมนุปสฺสติ, นิฏฺฐเมตฺถ คนฺตพฺพํ น ปริหายามิ กุสเลหิ ธมฺเมหิ, อปริหานุเมตํ วุตฺตํ ภควตา. กตเม จตฺตาโร? ราคตนุตฺตํ\csromansharedfootnote{c48af97a35b49498}{ราคตนุตฺตนํ (ก)} โทสตนุตฺตํ โมหตนุตฺตํ คมฺภีเรสุ โข ปนสฺส ฐานาฐาเนสุ ปญฺญาจกฺขุ กมติ. โย หิ โกจิ อาวุโส ภิกฺขุ วา ภิกฺขุนี วา อิเม จตฺตาโร ธมฺเม อตฺตนิ สมนุปสฺสติ, นิฏฺฐเมตฺถ คนฺตพฺพํ น ปริหายามิ กุสเลหิ ธมฺเมหิ, อปริหานเมตํ วุตฺตํ ภควตาติ.  อฏฺฐมํ.}

\csromanmatikaanchor{mtk.197}
\csromantocmarkref{subsection}{9. ภิกฺขุนีสุตฺต}{mtk.197}
\bookmark[dest={mtk.197},level=2]{9. ภิกฺขุนีสุตฺต}
\csromanheader{2}{9. \textbf{ภิกฺขุนีสุตฺต}}
\proseitem{159}{เอวํ เม สุตํ\csromandash{}เอกํ สมยํ อายสฺมา อานนฺโท โกสมฺพิยํ วิหรติ โฆสิตาราเม. อถ โข อญฺญตรา ภิกฺขุนี อญฺญตรํ ปุริสํ อามนฺเตสิ “เอหิ ตฺวํ อมฺโภ ปุริส เยนยฺโย อานนฺโท เตนุปสงฺกม, อุปสงฺกมิตฺวา มม วจเนน อยฺยสฺส อานนฺทสฺส ปาเท สิรสา วนฺท ‘อิตฺถนฺนามา ภนฺเต ภิกฺขุนี อาพาธิกินี ทุกฺขิตา พาฬฺหคิลานา, สา อยฺยสฺส อานนฺทสฺส ปาเท สิรสา วนฺทตี’ติ, เอวญฺจ วเทหิ ‘สาธุ กิร ภนฺเต อยฺโย อานนฺโท เยน ภิกฺขุนุ\-ปสฺสโย เยน สา ภิกฺขุนี เตนุ\-ปสงฺกมตุ อนุกมฺปํ อุปาทายา’ติ”. “เอวํ อยฺเย”ติ โข โส ปุริโส ตสฺสา ภิกฺขุนิยา ปฏิสฺสุตฺวา เยนายสฺมา อานนฺโท เตนุปสงฺกมิ, อุปสงฺกมิตฺวา อายสฺมนฺตํ อานนฺทํ อภิวาเทตฺวา เอกมนฺตํ นิสีทิ, เอกมนฺตํ นิสินฺโน โข โส ปุริโส อายสฺมนฺตํ อานนฺทํ เอตทโวจ\csromandash{}}

\prose{“อิตฺถนฺนามา ภนฺเต ภิกฺขุนี อาพาธิกินี ทุกฺขิตา พาฬฺหคิลานา, สา อายสฺมโต อานนฺทสฺส ปาเท สิรสา วนฺทติ, เอวญฺจ วเทติ ‘สาธุ กิร ภนฺเต อายสฺมา อานนฺโท เยน ภิกฺขุนุ\-ปสฺสโย เยน สา ภิกฺขุนี เตนุ\-ปสงฺกมตุ อนุกมฺปํ อุปาทายา’ติ”. อธิวาเสสิ โข อายสฺมา อานนฺโท ตุณฺหีภาเวน.}

\chapterheadpageread{(16) 1. อินฺทฺริยวคฺค}
\prose{\csromanfolio{463}อถ โข อายสฺมา อานนฺโท นิวาเสตฺวา ปตฺต\-จีวรมา\-ทาย เยน ภิกฺขุนุ\-ปสฺสโย เยน สา ภิกฺขุนี เตนุปสงฺกมิ. อทฺทสา โข สา ภิกฺขุนี อายสฺมนฺตํ อานนฺทํ ทูรโตว อาคจฺฉนฺตํ ทิสฺวา สสีสํ ปารุปิตฺวา มญฺจเก นิปชฺชิ. อถ โข อายสฺมา อานนฺโท เยน สา ภิกฺขุนี เตนุปสงฺกมิ, อุปสงฺกมิตฺวา ปญฺญตฺเต อาสเน นิสีทิ, นิสชฺช โข อายสฺมา อานนฺโท ตํ ภิกฺขุนิํ เอตทโวจ\csromandash{}}

\prose{อาหารสมฺภูโต อยํ ภคินิ กาโย, อาหารํ นิสฺสาย อาหาโร ปหาตพฺโพ. ตณฺหาสมฺภูโต อยํ ภคินิ กาโย, ตณฺหํ นิสฺสาย ตณฺหา ปหาตพฺพา. มานสมฺภูโต อยํ ภคินิ กาโย, มานํ นิสฺสาย มาโน ปหาตพฺโพ. เมถุน\-สมฺภูโต อยํ ภคินิ กาโย, เมถุเน จ เสตุฆาโต วุตฺโต ภควตา.}

\prose{“อาหารสมฺภูโต อยํ ภคินิ กาโย, อาหารํ นิสฺสาย อาหาโร ปหาตพฺโพ”ติ อิติ โข ปเนตํ วุตฺตํ, กิญฺเจตํ ปฏิจฺจ วุตฺตํ, อิธ ภคินิ ภิกฺขุ ปฏิสงฺขา โยนิโส อาหารํ อาหาเรติ “เนว ทวาย น มทาย น มณฺฑนาย น วิภูสนาย ยาวเทว อิมสฺส กายสฺส ฐิติยา ยาปนาย วิหิํสู\-ปรติยา พฺรหฺมจริยา\-นุคฺคหาย อิติ ปุราณญฺจ เวทนํ ปฏิหงฺขามิ นวญฺจ เวทนํ น อุปฺปาเทสฺสามิ, ยาตฺรา จ เม ภวิสฺสติ อนวชฺชตา จ ผาสุวิหาโร จา”ติ. โส อปเรน สมเยน อาหารํ นิสฺสาย อาหารํ ปชหติ. “อาหารสมฺภูโต อยํ ภคินิ กาโย, อาหารํ นิสฺสาย อาหาโร ปหาตพฺโพ”ติ อิติ ยํ ตํ วุตฺตํ, อิทเมตํ ปฏิจฺจ วุตฺตํ.}

\prose{“ตณฺหาสมฺภูโต อยํ ภคินิ กาโย, ตณฺหํ นิสฺสาย ตณฺหา ปหาตพฺพา”ติ อิติ โข ปเนตํ วุตฺตํ, กิญฺเจตํ ปฏิจฺจ วุตฺตํ, อิธ ภคินิ ภิกฺขุ สุณาติ “อิตฺถนฺนาโม กิร ภิกฺขุ อาสวานํ ขยา อนาสวํ เจโตวิมุตฺติํ ปญฺญาวิมุตฺติํ ทิฏฺเฐว ธมฺเม สยํ อภิญฺญา สจฺฉิกตฺวา อุปสมฺปชฺช วิหรตี”ติ. ตสฺส เอวํ โหติ “กุทาสฺสุ นาม อหมฺปิ อาสวานํ ขยา อนาสวํ เจโตวิมุตฺติํ ปญฺญาวิมุตฺติํ ทิฏฺเฐว ธมฺเม สยํ อภิญฺญา สจฺฉิกตฺวา อุปสมฺปชฺช วิหริสฺสามี”ติ. โส อปเรน สมเยน ตณฺหํ นิสฺสาย ตณฺหํ ปชหติ. “ตณฺหาสมฺภูโต อยํ ภคินิ กาโย, ตณฺหํ นิสฺสาย ตณฺหา ปหาตพฺพา”ติ อิติ ยํ ตํ วุตฺตํ, อิทเมตํ ปฏิจฺจ วุตฺตํ.}

\prose{\csromanfolio{464}“มานสมฺภูโต อยํ ภคินิ กาโย, มานํ นิสฺสาย มาโน ปหาตพฺโพ”ติ อิติ โข ปเนตํ วุตฺตํ, กิญฺเจตํ ปฏิจฺจ วุตฺตํ. อิธ ภคินิ ภิกฺขุ สุณาติ “อิตฺถนฺนาโม กิร ภิกฺขุ อาสวานํ ขยา อนาสวํ เจโตวิมุตฺติํ ปญฺญาวิมุตฺติํ ทิฏฺเฐว ธมฺเม สยํ อภิญฺญา สจฺฉิกตฺวา อุปสมฺปชฺช วิหรตี”ติ. ตสฺส เอวํ โหติ “โส หิ นาม อายสฺมา อาสวานํ ขยา อนาสวํ เจโตวิมุตฺติํ ปญฺญาวิมุตฺติํ ทิฏฺเฐว ธมฺเม สยํ อภิญฺญา สจฺฉิกตฺวา อุปสมฺปชฺช วิหริสฺสติ, กิมงฺคํ\csromansharedfootnote{e945a5ae4617fdfa}{กิมงฺค (สี, อิ) อํ 2. 190; วิ 4. 339; สํ 3. 328 ปิฏฺเฐสุปิ.} ปนาหนฺ”ติ. โส อปเรน สมเยน มานํ นิสฺสาย มานํ ปชหติ. “มานสมฺภูโต อยํ ภคินิ กาโย, มานํ นิสฺสาย มาโน ปหาตพฺโพ”ติ อิติ ยํ ตํ วุตฺตํ, อิทเมตํ ปฏิจฺจ วุตฺตํ.}

\prose{เมถุน\-สมฺภูโต อยํ ภคินิ กาโย, เมถุเน จ เสตุฆาโต วุตฺโต ภควตาติ.}

\prose{อถ โข สา ภิกฺขูนี มญฺจกา วุฏฺฐหิตฺวา เอกํสํ อุตฺตราสงฺคํ กริตฺวา อายสฺมโต อานนฺทสฺส ปาเทสุ สิรสา นิปติตฺวา อายสฺมนฺตํ อานนฺทํ เอตทโวจ “อจฺจโย มํ ภนฺเต อจฺจคมา ยถาพาลํ ยถามูฬฺหํ ยถาอกุสลํ, ยาหํ เอวมกาสิํ, ตสฺสา เม ภนฺเต อยฺโย อานนฺโท อจฺจยํ อจฺจยโต ปฏิคฺคณฺหาตุ อายติํ สํวรายา”ติ. ตคฺฆ ตํ\csromansharedfootnote{7c70b4c9702c262e}{ตคฺฆ ตฺวํ (สี, อิ, ก)} ภคินิ อจฺจโย อจฺจคมา ยถาพาลํ ยถามูฬฺหํ ยถาอกุสลํ, ยา ตฺวํ เอวมกาสิ. ยโต จ โข ตฺวํ ภคินิ อจฺจยํ อจฺจยโต ทิสฺวา ยถาธมฺมํ ปฏิกโรสิ, ตํ เต มยํ ปฏิคฺคณฺหาม. วุทฺธิ เหสา ภคินิ อริยสฺส วินเย, โย อจฺจยํ อจฺจยโต ทิสฺวา ยถาธมฺมํ ปฏิกโรติ, อายติํ สํวรํ อาปชฺชตีติ.  นวมํ.}

\csromanmatikaanchor{mtk.198}
\csromantocmarkref{subsection}{10. สุคตวินยสุตฺต}{mtk.198}
\bookmark[dest={mtk.198},level=2]{10. สุคตวินยสุตฺต}
\csromanheader{2}{10. สุคตวินย\-สุตฺต}
\proseitem{160}{สุคโต วา ภิกฺขเว โลเก ติฏฺฐมาโน สุคตวินโย วา, ตทสฺส พหุชนหิตาย พหุชน\-สุขาย โลกานุกมฺปาย อตฺถาย หิตาย สุขาย เทวมนุสฺสานํ.}

\prose{กตโม จ ภิกฺขเว สุคโต? อิธ ภิกฺขเว ตถาคโต โลเก อุปฺปชฺชติ อรหํ สมฺมาสมฺพุทฺโธ วิชฺชา\-จรณ\-สมฺปนฺโน สุคโต โลกวิทู}

\vspace{\tipitakaparskip}
\csromancenter{(16) 1. อินฺทฺริยวคฺค อนุตฺตโร ปุริส\-ทมฺม\-สารถิ สตฺถา เทวมนุสฺสานํ พุทฺโธ ภควา. อยํ ภิกฺขเว สุคโต.}
\prose{\csromanfolio{465}กตโม จ ภิกฺขเว สุคตวินโย? โส ธมฺมํ เทเสติ อาทิกลฺยาณํ มชฺเฌกลฺยาณํ ปริโยสาน\-กลฺยาณํ สาตฺถํ สพฺยญฺชนํ เกวล\-ปริปุณฺณํ ปริสุทฺธํ พฺรหฺมจริยํ ปกาเสติ. อยํ ภิกฺขเว สุคตวินโย. เอวํ สุคโต วา ภิกฺขเว โลเก ติฏฺฐมาโน สุคตวินโย วา, ตทสฺส พหุชนหิตาย พหุชน\-สุขาย โลกานุกมฺปาย อตฺถาย หิตาย สุขาย เทว\-มนุสฺสา\-นนฺติ.}

\prose{จตฺตาโรเม ภิกฺขเว ธมฺมา สทฺธมฺมสฺส สมฺโมสาย อนฺตรธานาย สํวตฺตนฺติ. กตเม จตฺตาโร? อิธ ภิกฺขเว ภิกฺขู ทุคฺคหิตํ สุตฺตนฺตํ ปริยาปุณนฺติ ทุนฺนิกฺขิตฺเตหิ ปทพฺยญฺชเนหิ, ทุนฺนิกฺขิตฺตสฺส ภิกฺขเว ปทพฺยญฺชนสฺส อตฺโถปิ ทุนฺนโย โหติ. อยํ ภิกฺขเว ปฐโม ธมฺโม สทฺธมฺมสฺส สมฺโมสาย อนฺตรธานาย สํวตฺตติ.}

\prose{ปุน จปรํ ภิกฺขเว ภิกฺขู ทุพฺพจา โหนฺติ โทวจสฺส\-กรเณหิ ธมฺเมหิ สมนฺนาคตา อกฺขมา อปฺปทกฺขิณคฺคาหิโน อนุสาสนิํ. อยํ ภิกฺขเว ทุติโย ธมฺโม สทฺธมฺมสฺส สมฺโมสาย อนฺตรธานาย สํวตฺตติ.}

\prose{ปุน จปรํ ภิกฺขเว เย เต ภิกฺขู พหุสฺสุตา อาคตาคมา ธมฺมธรา วินยธรา มาติกาธรา, เต น สกฺกจฺจํ สุตฺตนฺตํ ปรํ วาเจนฺติ, เตสํ อจฺจเยน ฉินฺนมูลโก สุตฺตนฺโต โหติ อปฺปฏิสรโณ. อยํ ภิกฺขเว ตติโย ธมฺโม สทฺธมฺมสฺส สมฺโมสาย อนฺตรธานาย สํวตฺตติ.}

\prose{ปุน จปรํ ภิกฺขเว เถรา ภิกฺขู พาหุลิกา โหนฺติ สาถลิกา โอกฺกมเน ปุพฺพงฺคมา ปวิเวเก นิกฺขิตฺตธุรา, น วีริยํ อารภนฺติ อปฺปตฺตสฺส ปตฺติยา อนธิคตสฺสา อธิคมาย อสจฺฉิกตสฺส สจฺฉิกิริยาย. เตสํ ปจฺฉิมา ชนตา ทิฏฺฐานุคติํ อาปชฺชติ, สาปิ โหติ พหุลิกา สาถลิกา โอกฺกมเน ปุพฺพงฺคมา ปวิเวเก นิกฺขิตฺตธุรา, น วีริยํ อารภติ อปฺปตฺตสฺส ปตฺติยา อนธิคตสฺส อธิคมาย อสจฺฉิกตสฺส สจฺฉิกิริยาย. อยํ ภิกฺขเว จตุตฺโถ ธมฺโม สทฺธมฺมสฺส สมฺโมสาย อนฺตรธานาย สํวตฺตติ. อิเม โข ภิกฺขเว จตฺตาโร ธมฺมา สทฺธมฺมสฺส สมฺโมสาย อนฺตรธานาย สํวตฺตนฺตีติ.}

\prose{\csromanfolio{466}จตฺตาโรเม ภิกฺขเว ธมฺมา สทฺธมฺมสฺส ฐิติยา อสมฺโมสาย อนนฺตรธานาย สํวตฺตนฺติ. กตเม จตฺตาโร? อิธ ภิกฺขเว ภิกฺขู สุคฺคหิตํ สุตฺตนฺตํ ปริยาปุณนฺติ สุนิกฺขิตฺเตหิ ปทพฺยญฺชเนหิ, สุนิกฺขิตฺตสฺส ภิกฺขเว ปทพฺยญฺชนสฺส อตฺโถปิ สุนโย โหติ. อยํ ภิกฺขเว ปฐโม สทฺธมฺมสฺส ฐิติยา อสมฺโมสาย อนนฺตรธานาย สํวตฺตติ.}

\prose{ปุน จปรํ ภิกฺขเว ภิกฺขู สุวจา โหนฺติ โสวจสฺส\-กรเณหิ ธมฺเมหิ สมนฺนาคตา ขมา ปทกฺขิณ\-คฺคาหิโน อนุสาสนิํ. อยํ ภิกฺขเว ทุติโย ธมฺโม สทฺธมฺมสฺส ฐิติยา อสมฺโมสาย อนนฺตรธานาย สํวตฺตติ.}

\prose{ปุน จปรํ ภิกฺขเว เย เต ภิกฺขู พหุสฺสุตา อาคตาคมา ธมฺมธรา วินยธรา มาติกาธรา, เต สกฺกจฺจํ สุตฺตนฺตํ ปรํ วาเจนฺติ, เตสํ อจฺจเยน นจฺฉินฺน\-มูลโก\csromansharedfootnote{a853de95d64b820e}{อจฺฉินฺนมูลโก (สฺยา, กํ) อํ 2. 158 ปิฏฺเฐปิ.} สุตฺตนฺโต โหติ สปฺปฏิสรโณ. อยํ ภิกฺขเว ตติโย ธมฺโม สทฺธมฺมสฺส ฐิติยา อสมฺโมสาย อนนฺตรธานาย สํวตฺตติ.}

\prosewithclosermidruled{ปุน จปรํ ภิกฺขเว เถรา ภิกฺขู น พาหุลิกา โหนฺติ น สาถลิกา โอกฺกมเน นิกฺขิตฺตธุรา ปวิเวเก ปุพฺพงฺคมา, วีริยํ อารภนฺติ อปฺปตฺตสฺส ปตฺติยา อนธิคตสฺส อธิคมาย อสจฺฉิกตสฺส สจฺฉิกิริยาย, เตสํ ปจฺฉิมา ชนตา ทิฏฺฐานุคติํ อาปชฺชติ, สาปิ โหติ น พาหุลิกา น สาถลิกา โอกฺกมเน นิกฺขิตฺตธุรา ปวิเวเก ปุพฺพงฺคมา, วีริยํ อารภติ อปฺปตฺตสฺส ปตฺติยา อนธิคตสฺส อธิคมาย อสจฺฉิกตสฺส สจฺฉิกิริยาย. อยํ ภิกฺขเว จตุตฺโถ ธมฺโม สทฺธมฺมสฺส ฐิติยา อสมฺโมสาย อนนฺตรธานาย สํวตฺตติ. อิเม โข ภิกฺขเว จตฺตาโร ธมฺมา สทฺธมฺมสฺส ฐิติยา อสมฺโมสาย อนนฺตรธานาย สํวตฺตนฺตีติ.}{อินฺทฺริยวคฺโค ปฐโม.}
\csromancenter{\textbf{ตสฺสุทฺทานํ}}
\setlength{\csromangathapagewidth}{0pt}
\setlength{\csromangathawakwidth}{0pt}
\csromangathameasuregroup{อินฺทฺริยานิ สทฺธา ปญฺญา, \\ กปฺโป โรโค ปริหานิ,}{\csromangathabat{อินฺทฺริยานิ สทฺธา ปญฺญา,}{สติ สงฺขานปญฺจมํ.} \\ \csromangathabat{กปฺโป โรโค ปริหานิ,}{ภิกฺขุนี สุคเตน จาติ.}}
\csromangathasetleft{อินฺทฺริยานิ สทฺธา ปญฺญา, \\ กปฺโป โรโค ปริหานิ,}
\csromangathagroup{\csromangathastanza{\csromangathabat{อินฺทฺริยานิ สทฺธา ปญฺญา,}{สติ สงฺขานปญฺจมํ.} \\ \csromangathabat{กปฺโป โรโค ปริหานิ,}{ภิกฺขุนี สุคเตน จาติ.}}}

\chapterheadpageread{(17) 2. ปฏิปทาวคฺค \textbf{(17) 2. ปฏิปทาวคฺค}}
\csromanheader{2}{1. \textbf{สํขิตฺตสุตฺต}}
\csromanmatikaanchor{mtk.199}
\csromanmatikaanchor{mtk.200}
\csromantocmarknopage{section}{(17) 2. ปฏิปทาวคฺค}{mtk.199}
\bookmark[dest={mtk.199},level=1]{(17) 2. ปฏิปทาวคฺค}
\csromantocmarkref{subsection}{1-3. สํขิตฺตสุตฺตาทีนิ}{mtk.200}
\bookmark[dest={mtk.200},level=2]{1-3. สํขิตฺตสุตฺตาทีนิ}
\proseitem{161}{\csromanfolio{467}จตสฺโส อิมา ภิกฺขเว ปฏิปทา. กตมา จตสฺโส? ทุกฺขา ปฏิปทา ทนฺธาภิญฺญา, ทุกฺขา ปฏิปทา ขิปฺปาภิญฺญา, สุขา ปฏิปทา ทนฺธาภิญฺญา, สุขา ปฏิปทา ขิปฺปาภิญฺญา. อิมา โข ภิกฺขเว จตสฺโส ปฏิปทาติ.  ปฐมํ.}

\csromanheader{2}{2. \textbf{วิตฺถารสุตฺต}}
\proseitem{162}{จตสฺโส อิมา ภิกฺขเว ปฏิปทา. กตมา จตสฺโส? ทุกฺขา ปฏิปทา ทนฺธาภิญฺญา, ทุกฺขา ปฏิปทา ขิปฺปาภิญฺญา, สุขา ปฏิปทา ทนฺธาภิญฺญา, สุขา ปฏิปทา ขิปฺปาภิญฺญา.}

\prose{กตมา จ ภิกฺขเว ทุกฺขา ปฏิปทา ทนฺธาภิญฺญา? อิธ ภิกฺขเว เอกจฺโจ ปกติยาปิ ติพฺพ\-ราค\-ชาติโก โหติ, อภิกฺขณํ ราคชํ ทุกฺขํ โทมนสฺสํ ปฏิสํเวเทติ. ปกติยาปิ ติพฺพ\-โทส\-ชาติโก โหติ, อภิกฺขณํ โทสชํ ทุกฺขํ โทมนสฺสํ ปฏิสํเวเทติ. ปกติยาปิ ติพฺพ\-โมห\-ชาติโก โหติ, อภิกฺขณํ โมหชํ ทุกฺขํ โทมนสฺสํ ปฏิสํเวเทติ. ตสฺสิมานิ ปญฺจินฺทฺริยานิ มุทูนิ ปาตุภวนฺติ สทฺธินฺทฺริยํ วีริยินฺทฺริยํ สตินฺทฺริยํ สมาธินฺทฺริยํ ปญฺญินฺทฺริยํ. โส อิเมสํ ปญฺจนฺนํ อินฺทฺริยานํ มุทุตฺตา ทนฺธํ อานนฺตริยํ ปาปุณาติ อาสวานํ ขยาย. อยํ วุจฺจติ ภิกฺขเว ทุกฺขา ปฏิปทา ทนฺธาภิญฺญา.}

\prose{กตมา จ ภิกฺขเว ทุกฺขา ปฏิปทา ขิปฺปาภิญฺญา? อิธ ภิกฺขเว เอกจฺโจ ปกติยาปิ ติพฺพาราคชาติโก โหติ, อภิกฺขณํ ราคชํ ทุกฺขํ โทมนสฺสํ ปฏิสํเวเทติ. ปกติยาปิ ติพฺพ\-โทส\-ชาติโก โหติ, อภิกฺขณํ โทสชํ ทุกฺขํ โทมนสฺสํ ปฏิสํเวเทติ. ปกติยาปิ ติพฺพ\-โมห\-ชาติโก โหติ, อภิกฺขณํ โมหชํ ทุกฺขํ โทมนสฺสํ ปฏิสํเวเทติ. ตสฺสิมานิ ปญฺจินฺทฺริยานิ อธิมตฺตานิ ปาตุภวนฺติ สทฺธินฺทฺริยํ วีริยินฺทฺริยํ สตินฺทฺริยํ สมาธินฺทฺริยํ ปญฺญินฺทฺริยํ. โส อิเมสํ ปญฺจนฺนํ อินฺทฺริยานํ อธิมตฺตตฺตา ขิปฺปํ อานนฺตริยํ ปาปุณาติ อาสวานํ ขยาย. อยํ วุจฺจติ ภิกฺขเว ทุกฺขา ปฏิปทา ขิปฺปาภิญฺญา.}

\prose{กตมา จ ภิกฺขเว สุขา ปฏิปทา ทนฺธาภิญฺญา? อิธ ภิกฺขเว เอกจฺโจ ปกติยาปิ น ติพฺพ\-ราค\-ชาติโก โหติ, นาภิกฺขณํ ราคชํ ทุกฺขํ \csromanfolio{468}โทมนสฺสํ ปฏิสํเวเทติ. ปกติยาปิ น ติพฺพ\-โทส\-ชาติโก โหติ, นาภิกฺขณํ โทสชํ ทุกฺขํ โทมนสฺสํ ปฏิสํเวเทติ. ปกติยาปิ น ติพฺพ\-โมห\-ชาติโก โหติ, นาภิกฺขณํ โมหชํ ทุกฺขํ โทมนสฺสํ ปฏิสํเวเทติ. ตสฺสิมานิ ปญฺจินฺทฺริยานิ มุทูนิ ปาตุภวนฺติ สทฺธินฺทฺริยํ ฯเปฯ ปญฺญินฺทฺริยํ. โส อิเมสํ ปญฺจนฺนํ อินฺทฺริยานํ มุทุตฺตา ทนฺธํ อานนฺตริยํ ปาปุณาติ อาสวานํ ขยาย. อยํ วุจฺจติ ภิกฺขเว สุขา ปฏิปทา ทนฺธาภิญฺญา.}

\prose{กตมา จ ภิกฺขเว สุขา ปฏิปทา ขิปฺปาภิญฺญา? อิธ ภิกฺขเว เอกจฺโจ ปกติยาปิ น ติพฺพ\-ราค\-ชาติโก โหติ, นาภิกฺขณํ ราคชํ ทุกฺขํ โทมนสฺสํ ปฏิสํเวเทติ. ปกติยาปิ น ติพฺพ\-โทส\-ชาติโก โหติ, นาภิกฺขณํ โทสชํ ทุกฺขํ โทมนสฺสํ ปฏิสํเวเทติ. ปกติยาปิ น ติพฺพ\-โมห\-ชาติโก โหติ, นาภิกฺขณํ โมหชํ ทุกฺขํ โทมนสฺสํ ปฏิสํเวเทติ. ตสฺสิมานิ ปญฺจินฺทฺริยานิ อธิมตฺตานิ ปาตุภวนฺติ สทฺธินฺทฺริยํ วีริยินฺทฺริยํ สตินฺทฺริยํ สมาธินฺทฺริยํ ปญฺญินฺทฺริยํ. โส อิเมสํ ปญฺจนฺนํ อินฺทฺริยานํ อธิมตฺตตฺตา ขิปฺปํ อานนฺตริยํ ปาปุณาติ อาสวานํ ขยาย. อยํ วุจฺจติ ภิกฺขเว สุขา ปฏิปทา ขิปฺปาภิญฺญา. อิมา โข ภิกฺขเว จตสฺโส ปฏิปทาติ.  ทุติยํ.}

\csromanheader{2}{3. \textbf{อสุภสุตฺต}}
\proseitem{163}{จตสฺโส อิมา ภิกฺขเว ปฏิปทา. กตมา จตสฺโส? ทุกฺขา ปฏิปทา ทนฺธาภิญฺญา, ทุกฺขา ปฏิปทา ขิปฺปาภิญฺญา, สุขา ปฏิปทา ทนฺธาภิญฺญา, สุขา ปฏิปทา ขิปฺปาภิญฺญา.}

\prose{กตมา จ ภิกฺขเว ทุกฺขา ปฏิปทา ทนฺธาภิญฺญา? อิธ ภิกฺขเว ภิกฺขุ อสุภานุปสฺสี กาเย วิหรติ อาหาเร ปฏิกูลสญฺญี\csromansharedfootnote{34b7a5f1354b1da6}{ปฏิกฺกูลสญฺญี (สี, สฺยา, กํ, อิ)} สพฺพโลเก อนภิรติ\-สญฺญี\csromansharedfootnote{e4efc3faab3dbcba}{อนภิรตสญฺญี (สี, สฺยา, กํ, อิ)} สพฺพ\-สงฺขาเรสุ อนิจฺจานุปสฺสี, มรณสญฺญา โข ปนสฺส อชฺฌตฺตํ สูปฏฺฐิตา โหติ. โส อิมานิ ปญฺจ เสขพลานิ\csromansharedfootnote{3cccf4bf3dbe9ccf}{เสกฺเขพลานิ (สฺยา, กํ)} อุปนิสฺสาย วิหรติ สทฺธาพลํ หิริพลํ โอตฺตปฺปพลํ วีริยพลํ ปญฺญาพลํ. ตสฺสิมานิ ปญฺจินฺทฺริยานิ มุทูนิ ปาตุภวนฺติ สทฺธินฺทฺริยํ วีริยินฺทฺริยํ สตินฺทฺริยํ สมาธินฺทฺริยํ ปญฺญินฺทฺริยํ. โส อิเมสํ ปญฺจนฺนํ อินฺทฺริยานํ มุทุตฺตา ทนฺธํ อานนฺตริยํ ปาปุณาติ อาสวานํ ขยาย. อยํ วุจฺจติ ภิกฺขเว ทุกฺขา ปฏิปทา ทนฺธาภิญฺญา.}

\csromanmarkh{(17)}\csromanmarkhii{2. ปฏิปทาวคฺค}\csromanheadernomark{1}{(17) 2. ปฏิปทาวคฺค}
\prose{\csromanfolio{469}กตมา จ ภิกฺขเว ทุกฺขา ปฏิปทา ขิปฺปาภิญฺญา? อิธ ภิกฺขเว ภิกฺขุ อสุภานุปสฺสี กาเย วิหรติ อาหาเร ปฏิกูลสญฺญี สพฺพโลเก อนภิรติ\-สญฺญี สพฺพ\-สงฺขาเรสุ อนิจฺจานุปสฺสี, มรณสญฺญา โข ปนสฺส อชฺฌตฺตํ สูปฏฺฐิตา โหติ. โส อิมานิ ปญฺจ เสขพลานิ อุปนิสฺสาย วิหรติ สทฺธาพลํ ฯเปฯ ปญฺญาพลํ. ตสฺสิมานิ ปญฺจินฺทฺริยานิ อธิมตฺตานิ ปาตุภวนฺติ สทฺธินฺทฺริยํ ฯเปฯ ปญฺญินฺทฺริยํ. โส อิเมสํ ปญฺจนฺนํ อินฺทฺริยานํ อธิมตฺตตฺตา ขิปฺปํ อานนฺตริยํ ปาปุณาติ อาสวานํ ขยาย. อยํ วุจฺจติ ภิกฺขเว ทุกฺขา ปฏิปทา ขิปฺปาภิญฺญา.}

\prose{กตมา จ ภิกฺขเว สุขา ปฏิปทา ทนฺธาภิญฺญา? อิธ ภิกฺขเว ภิกฺขุ วิวิจฺเจว กาเมหิ วิวิจฺจ อกุสเลหิ ธมฺเมหิ สวิตกฺกํ สวิจารํ วิเวกชํ ปีติสุขํ ปฐมํ ฌานํ อุปสมฺปชฺช วิหรติ. วิตกฺก\-วิจารานํ วูปสมา อชฺฌตฺตํ สมฺปสาทนํ เจตโส เอโกทิภาวํ อวิตกฺกํ อวิจารํ สมาธิชํ ปีติสุขํ ทุติยํ ฌานํ อุปสมฺปชฺช วิหรติ. ปีติยา จ วิราคา อุเปกฺขโก จ วิหรติ สโต จ สมฺปชาโน สุขญฺจ กาเยน ปฏิสํเวเทติ, ยํ ตํ อริยา อาจิกฺขนฺติ “อุเปกฺขโก สติมา สุขวิหารี”ติ, ตติยํ ฌานํ อุปสมฺปชฺช วิหรติ. สุขสฺส จ ปหานา ทุกฺขสฺส จ ปหานา ปุพฺเพว โสมนสฺส\-โทมนสฺสานํ อตฺถงฺคมา อทุกฺขมสุขํ อุเปกฺขา\-สติ\-ปาริสุทฺธิํ จตุตฺถํ ฌานํ อุปสมฺปชฺช วิหรติ. โส อิมานิ ปญฺจ เสขพลานิ อุปนิสฺสาย วิหรติ สทฺธาพลํ ฯเปฯ ปญฺญาพลํ. ตสฺสิมานิ ปญฺจินฺทฺริยานิ มุทูนิ ปาตุภวนฺติ สทฺธินฺทฺริยํ ฯเปฯ ปญฺญินฺทฺริยํ. โส อิเมสํ ปญฺจนฺนํ อินฺทฺริยานํ มุทุตฺตา ทนฺธํ อานนฺตริยํ ปาปุณาติ อาสวานํ ขยาย. อยํ วุจฺจติ ภิกฺขเว สุขา ปฏิปทา ทนฺธาภิญฺญา.}

\prose{กตมา จ ภิกฺขเว สุขา ปฏิปทา ขิปฺปาภิญฺญา?\csromansymbolfootnote{*}{อภิ 4. 412 ปิฏฺเฐปิ.}อิธ ภิกฺขเว ภิกฺขุ วิวิจฺเจว กาเมหิ วิวิจฺจ อกุสเลหิ ธมฺเมหิ สวิตกฺกํ สวิจารํ วิเวกชํ ปีติสุขํ ปฐมํ ฌานํ อุปสมฺปชฺช วิหรติ ฯเปฯ ทุติยํ ฌานํ ฯเปฯ ตติยํ ฌานํ ฯเปฯ จตุตฺถํ ฌานํ อุปสมฺปชฺช วิหรติ. โส อิมานิ ปญฺจ เสขพลานิ อุปนิสฺสาย วิหรติ สทฺธาพลํ หิริพลํ โอตฺตปฺปพลํ วีริยพลํ ปญฺญาพลํ. ตสฺสิมานิ ปญฺจินฺทฺริยานิ อธิมตฺตานิ ปาตุภวนฺติ สทฺธินฺทฺริยํ วีริยินฺทฺริยํ สตินฺทฺริยํ สมาธินฺทฺริยํ ปญฺญินฺทฺริยํ. โส อิเมสํ ปญฺจนฺนํ อินฺทฺริยานํ อธิมตฺตตฺตา ขิปฺปํ อานนฺตริยํ ปาปุณาติ อาสวานํ ขยาย. อยํ วุจฺจติ ภิกฺขเว สุขา ปฏิปทา ขิปฺปาภิญฺญา. อิมา โข ภิกฺขเว จตสฺโส ปฏิปทาติ.  ตติยํ.}

\csromanpalirecto
\gambhira{จตุกฺกนิปาต\-ปาฬิ 4. ปฐม\-ขม\-สุตฺต}
\csromanmatikaanchor{mtk.201}
\csromantocmarkref{subsection}{4-5. ขมสุตฺตทฺวยํ}{mtk.201}
\bookmark[dest={mtk.201},level=2]{4-5. ขมสุตฺตทฺวยํ}
\proseitem{164}{\csromanfolio{470}จตสฺโส อิมา ภิกฺขเว ปฏิปทา. กตมา จตสฺโส? อกฺขมา ปฏิปทา, ขมา ปฏิปทา, ทมา ปฏิปทา, สมา ปฏิปทา. กตมา จ ภิกฺขเว อกฺขมา ปฏิปทา? อิธ ภิกฺขเว เอกจฺโจ\csromansharedfootnote{0c0b7b5ad47470a1}{เอกจฺโจ ปุคฺคโล (สี, สฺยา, กํ)} อกฺโกสนฺตํ ปจฺจกฺโกสติ, โรสนฺตํ ปฏิโรสติ, ภณฺฑนฺตํ ปฏิภณฺฑติ. อยํ วุจฺจติ ภิกฺขเว อกฺขมา ปฏิปทา.}

\prose{กตมา จ ภิกฺขเว ขมา ปฏิปทา? อิธ ภิกฺขเว เอกจฺโจ อกฺโกสนฺตํ น ปจฺจกฺโกสติ, โรสนฺตํ น ปฏิโรสติ, ภณฺฑนฺตํ น ปฏิภณฺฑติ. อยํ วุจฺจติ ภิกฺขเว ขมา ปฏิปทา.}

\prose{กตมา จ ภิกฺขเว ทมา ปฏิปทา? อิธ ภิกฺขเว ภิกฺขุ จกฺขุนา รูปํ ทิสฺวา น นิมิตฺตคฺคาหี โหติ นานุพฺยญฺชน\-คฺคาหี, ยตฺวาธิกรณเมนํ จกฺขุนฺทฺริยํ อสํวุตํ วิหรนฺตํ อภิชฺฌา\-โทมนสฺสา ปาปกา อกุสลา ธมฺมา อนฺวาสฺสเวยฺยุํ, ตสฺส สํวราย ปฏิปชฺชติ, รกฺขติ จกฺขุนฺทฺริยํ, จกฺขุนฺทฺริเย สํวรํ อาปชฺชติ. โสเตน สทฺทํ สุตฺวา. ฆาเนน คนฺธํ ฆายิตฺวา. ชิวฺหาย รสํ สายิตฺวา. กาเยน โผฏฺฐพฺพํ ผุสิตฺวา. มนสา ธมฺมํ วิญฺญาย น นิมิตฺตคฺคาหี โหติ นานุพฺยญฺชน\-คฺคาหี, ยตฺวาธิกรณเมนํ มนินฺทฺริยํ อสํวุตํ วิหรนฺตํ อภิชฺฌา\-โทมนสฺสา ปาปกา อกุสลา ธมฺมา อนฺวาสฺสเวยฺยุํ, ตสฺส สํวราย ปฏิปชฺชติ, รกฺขติ มนินฺทฺริยํ, มนินฺทฺริเย สํวรํ อาปชฺชติ. อยํ วุจฺจติ ภิกฺขเว ทมา ปฏิปทา.}

\prose{กตมา จ ภิกฺขเว สมา ปฏิปทา? อิธ ภิกฺขเว ภิกฺขุ อุปฺปนฺนํ กามวิตกฺกํ นาธิวาเสติ ปชหติ วิโนเทติ สเมติ พฺยนฺตีกโรติ อนภาวํ คเมติ. อุปฺปนฺนํ พฺยาปาท\-วิตกฺกํ ฯเปฯ. อุปฺปนฺนํ วิหิํสา\-วิตกฺกํ. อุปฺปนฺนุปฺปนฺเน ปาปเก อกุสเล ธมฺเม นาธิวาเสติ ปชหติ วิโนเทติ สเมติ พฺยนฺตีกโรติ อนภาวํ คเมติ. อยํ วุจฺจติ ภิกฺขเว สมา ปฏิปทา. อิมา โข ภิกฺขเว จตสฺโส ปฏิปทาติ.  จตุตฺถํ.}

\csromanheader{2}{5. ทุติย\-ขม\-สุตฺต}
\proseitem{165}{จตสฺโส อิมา ภิกฺขเว ปฏิปทา. กตมา จตสฺโส? อกฺขมา ปฏิปทา, ขมา ปฏิปทา, ทมา ปฏิปทา, สมา ปฏิปทา.}

\chapterheadpageread{(17) 2. ปฏิปทาวคฺค}
\prose{\csromanfolio{471}กตมา จ ภิกฺขเว อกฺขมา ปฏิปทา? อิธ ภิกฺขเว เอกจฺโจ อกฺขโม โหติ สีตสฺส อุณฺหสฺส ชิฆจฺฉาย ปิปาสาย ฑํสมกส\-วาตา\-ตป\-สรีสป\-สมฺผสฺสานํ ทุรุตฺตานํ ทุราคตานํ วจนปถานํ, อุปฺปนฺนานํ สารีริกานํ เวทนานํ ทุกฺขานํ ติพฺพานํ ขรานํ กฏุกานํ อสาตานํ อมนาปานํ ปาณหรานํ อนธิวาสก\-ชาติโก โหติ. อยํ วุจฺจติ ภิกฺขเว อกฺขมา ปฏิปทา.}

\prose{กตมา จ ภิกฺขเว ขมา ปฏิปทา? อิธ ภิกฺขเว เอกจฺโจ ขโม โหติ สีตสฺส อุณฺหสฺส ชิฆจฺฉาย ปิปาสาย ฑํสมกส\-วาตา\-ตป\-สรีสป\-สมฺผสฺสานํ ทุรุตฺตานํ ทุราคตานํ วจนปถานํ, อุปฺปนฺนานํ สารีริกานํ เวทนานํ ทุกฺขานํ ติพฺพานํ ขรานํ กฏุกานํ อสาตานํ อมนาปานํ ปาณหรานํ อธิวาสหชาติโก โหติ. อยํ วุจฺจติ ภิกฺขเว ขมา ปฏิปทา.}

\prose{กตมา จ ภิกฺขเว ทมา ปฏิปทา? อิธ ภิกฺขเว ภิกฺขุ จกฺขุนา รูปํ ทิสฺวา น นิมิตฺตคฺคาหี โหติ ฯเปฯ โสเตน สทฺทํ สุตฺวา. ฆาเนน คนฺธํ ฆายิตฺวา. ชิวฺหาย รสํ สายิตฺวา. กาเยน โผฏฺฐพฺพํ ผุสิตฺวา. มนสา ธมฺมํ วิญฺญาย น นิมิตฺตคฺคาหี โหติ นานุพฺยญฺชน\-คฺคาหี, ยตฺวาธิกรณเมนํ มนินฺทฺริยํ อสํวุตํ วิหรนฺตํ อภิชฺฌา\-โทมนสฺสา ปาปกา อกุสลา ธมฺมา อนฺวาสฺสเวยฺยุํ, ตสฺส สํวราย ปฏิปชฺชติ, รกฺขติ มนินฺทฺริยํ มนินฺทฺริเย สํวรํ อาปชฺชติ. อยํ วุจฺจติ ภิกฺขเว ทมา ปฏิปทา.}

\prose{กตมา จ ภิกฺขเว สมา ปฏิปทา? อิธ ภิกฺขเว ภิกฺขุ อุปฺปนฺนํ กามวิตกฺกํ นาธิวาเสติ ปชหติ วิโนเทติ สเมติ พฺยนฺตีกโรติ อนภาวํ คเมติ. อุปฺปนฺนํ พฺยาปาท\-วิตกฺกํ ฯเปฯ. อุปฺปนฺนํ วิหิํสา\-วิตกฺกํ. อุปฺปนฺนุปฺปนฺเน ปาปเก อกุสเล ธมฺเม นาธิวาเสติ ปชหติ วิโนเทติ สเมติ พฺยนฺตีกโรติ อนภาวํ คเมติ. อยํ วุจฺจติ ภิกฺขเว สมา ปฏิปทา. อิมา โข ภิกฺขเว จตสฺโส ปฏิปทาติ.  ปญฺจมํ.}

\csromanmatikaanchor{mtk.202}
\csromantocmarkref{subsection}{6. อุภยสุตฺต}{mtk.202}
\bookmark[dest={mtk.202},level=2]{6. อุภยสุตฺต}
\csromanheader{2}{6. \textbf{อุภยสุตฺต}}
\proseitem{166}{จตสฺโส อิมา ภิกฺขเว ปฏิปทา. กตมา จตสฺโส? ทุกฺขา ปฏิปทา ทนฺธาภิญฺญา, ทุกฺขา ปฏิปทา ขิปฺปาภิญฺญา, สุขา ปฏิปทา ทนฺธาภิญฺญา, สุขา ปฏิปทา ขิปฺปาภิญฺญา.}

\prose{\csromanfolio{472}ตตฺร ภิกฺขเว ยายํ ปฏิปทา ทุกฺขา ทนฺธาภิญฺญา, อยํ ภิกฺขเว ปฏิปทา อุภเยเนว หีนา อกฺขายติ. ยมฺปายํ ปฏิปทา ทุกฺขา, อิมินาปายํ หีนา อกฺขายติ. ยมฺปายํ ปฏิปทา ทนฺธา, อิมินาปายํ หีนา อกฺขายติ. อยํ ภิกฺขเว ปฏิปทา อุภเยเนว หีนา อกฺขายติ.}

\prose{ตตฺร ภิกฺขเว ยายํ ปฏิปทา ทุกฺขา ขิปฺปาภิญฺญา, อยํ ภิกฺขเว ปฏิปทา ทุกฺขตฺตา หีนา อกฺขายติ.}

\prose{ตตฺร ภิกฺขเว ยายํ ปฏิปทา สุขา ทนฺธาภิญฺญา, อยํ ภิกฺขเว ปฏิปทา ทนฺธตฺตา หีนา อกฺขายติ.}

\prose{ตตฺร ภิกฺขเว ยายํ ปฏิปทา สุขา ขิปฺปาภิญฺญา, อยํ ภิกฺขเว ปฏิปทา อุภเยเนว ปณีตา อกฺขายติ. ยมฺปายํ ปฏิปทา สุขา, อิมินาปายํ ปณีตา อกฺขายติ. ยมฺปายํ ปฏิปทา ขิปฺปา, อิมินาปายํ ปณีตา อกฺขายติ. อยํ ภิกฺขเว ปฏิปทา อุภเยเนว ปณีตา อกฺขายติ. อิมา โข ภิกฺขเว จตสฺโส ปฏิปทาติ.  ฉฏฺฐํ.}

\csromanmatikaanchor{mtk.203}
\csromantocmarkref{subsection}{7. มหาโมคฺคลฺลานสุตฺต}{mtk.203}
\bookmark[dest={mtk.203},level=2]{7. มหาโมคฺคลฺลานสุตฺต}
\csromanheader{2}{7. มหาโมคฺคลฺลาน\-สุตฺต}
\proseitem{167}{อถ โข อายสฺมา สาริปุตฺโต เยนายสฺมา มหาโมคฺคลฺลาโน เตนุปสงฺกมิ, อุปสงฺกมิตฺวา อายสฺมตา มหา\-โมคฺคลฺลาเนน สทฺธิํ สมฺโมทิ, สมฺโมทนียํ กถํ สารณียํ วีติสาเรตฺวา เอกมนฺตํ นิสีทิ, เอกมนฺตํ นิสินฺโน โข อายสฺมา สาริปุตฺโต อายสฺมนฺตํ มหา\-โมคฺคลฺลานํ เอตทโวจ\csromandash{}}

\prose{จตสฺโส อิมา อาวุโส โมคฺคลฺลาน ปฏิปทา. กตมา จตสฺโส? ทุกฺขา ปฏิปทา ทนฺธาภิญฺญา, ทุกฺขา ปฏิปทา ขิปฺปาภิญฺญา, สุขา ปฏิปทา ทนฺธาภิญฺญา, สุขา ปฏิปทา ขิปฺปาภิญฺญา. อิมา โข อาวุโส จตสฺโส ปฏิปทา. อิมาสํ อาวุโส จตุนฺนํ ปฏิปทานํ\csromansharedfootnote{c7c09d68df2474c5}{จตสฺสนฺนํ ปฏิปทานํ (สี, สฺยา, กํ)} กตมํ เต ปฏิปทํ อาคมฺม อนุปาทาย อาสเวหิ จิตฺตํ วิมุตฺตนฺติ.}

\prose{จตสฺโส อิมา อาวุโส สาริปุตฺต ปฏิปทา. กตมา จตสฺโส? ทุกฺขา ปฏิปทา ทนฺธาภิญฺญา, ทุกฺขา ปฏิปทา ขิปฺปาภิญฺญา, สุขา ปฏิปทา ทนฺธาภิญฺญา, สุขา ปฏิปทา ขิปฺปาภิญฺญา. อิมา โข อาวุโส จตสฺโส}

\vspace{\tipitakaparskip}
\csromancenter{(17) 2. ปฏิปทาวคฺค ปฏิปทา. อิมาสํ อาวุโส จตุนฺนํ ปฏิปทานํ ยายํ ปฏิปทา ทุกฺขา ขิปฺปาภิญฺญา, อิมํ เม ปฏิปทํ อาคมฺม อนุปาทาย อาสเวหิ จิตฺตํ วิมุตฺตนฺติ.}
\csromanmatikaanchor{mtk.204}
\csromantocmarkref{subsection}{8. สาริปุตฺตสุตฺต}{mtk.204}
\bookmark[dest={mtk.204},level=2]{8. สาริปุตฺตสุตฺต}
\csromanheader{2}{8. \textbf{สาริปุตฺตสุตฺต}}
\proseitem{168}{\csromanfolio{473}อถ โข อายสฺมา มหาโมคฺคลฺลาโน เยนายสฺมา สาริปุตฺโต เตนุปสงฺกมิ, อุปสงฺกมิตฺวา อายสฺมตา สาริปุตฺเตน สทฺธํ สมฺโมทิ, สมฺโมทนียํ กถํ สารณียํ วีติสาเรตฺวา เอกมนฺตํ นิสีทิ, เอกมนฺตํ นิสินฺโน โข อายสฺมา มหาโมคฺคลฺลโน อายสฺมนฺตํ สาริปุตฺตํ เอตทโวจ\csromandash{}}

\prose{จตสฺโส อิมา อาวุโส สาริปุตฺต ปฏิปทา. กตมา จตสฺโส? ทุกฺขา ปฏิปทา ทนฺธาภิญฺญา, ทุกฺขา ปฏิปทา ขิปฺปาภิญฺญา, สุขา ปฏิปทา ทนฺธาภิญฺญา, สุขา ปฏิปทา ขิปฺปาภิญฺญา. อิมา โข อาวุโส จตสฺโส ปฏิปทา. อิมาสํ อาวุโส จตุนฺนํ ปฏิปทานํ กตมํ เต ปฏิปทํ อาคมฺม อนุปาทาย อาสเวหิ จิตฺตํ วิมุตฺตนฺติ.}

\prose{จตสฺโส อิมา อาวุโส โมคฺคลฺลาน ปฏิปทา. กตมา จตสฺโส? ทุกฺขา ปฏิปทา ทนฺธาภิญฺญา, ทุกฺขา ปฏิปทา ขิปฺปาภิญฺญา, สุขา ปฏิปทา ทนฺธาภิญฺญา, สุขา ปฏิปทา ขิปฺปาภิญฺญา. อิมา โข อาวุโส จตสฺโส ปฏิปทา. อิมาสํ อาวุโส จตุนฺนํ ปฏิปทานํ ยายํ ปฏิปทา สุขา ขิปฺปาภิญฺญา, อิมํ เม ปฏิปทํ อาคมฺม อนุปาทาย อาสเวหิ จิตฺตํ วิมุตฺตนฺติ.  อฏฺฐมํ.}

\csromanmatikaanchor{mtk.205}
\csromantocmarkref{subsection}{9. สสงฺขารสุตฺต}{mtk.205}
\bookmark[dest={mtk.205},level=2]{9. สสงฺขารสุตฺต}
\csromanheader{2}{9. สสงฺขาร\-สุตฺต}
\proseitem{169}{จตฺตาโรเม ภิกฺขเว ปุคฺคลา สนฺโต สํวิชฺชมานา โลกสฺมิํ. กตเม จตฺตาโร? อิธ ภิกฺขเว เอกจฺโจ ปุคฺคโล ทิฏฺเฐว ธมฺเม สสงฺขาร\-ปรินิพฺพายี โหติ, อิธ ปน ภิกฺขเว เอกจฺโจ ปุคฺคโล กายสฺส เภทา สสงฺขาร\-ปรินิพฺพายี โหติ, อิธ ปน ภิกฺขเว เอกจฺโจ ปุคฺคโล ทิฏฺเฐว ธมฺเม อสงฺขารปรินิพฺพายี โหติ, อิธ ปน ภิกฺขเว เอกจฺโจ ปุคฺคโล กายสฺส เภทา อสงฺขารปรินิพฺพายี โหติ.}

\prose{กถญฺจ ภิกฺขเว ปุคฺคโล ทิฏฺเฐว ธมฺเม สสงฺขาร\-ปรินิพฺพายี โหติ? อิธ ภิกฺขเว ภิกฺขุ อสุภานุปสฺสี กาเย วิหรติ อาหาเร ปฏิกูลสญฺญี สพฺพโลเก อนภิรติ\-สญฺญี สพฺพ\-สงฺขาเรสุ อนิจฺจานุปสฺสี, \csromanfolio{474}มรณสญฺญา โข ปนสฺส อชฺฌตฺตํ สูปฏฺฐิตา โหติ. โส อิมานิ ปญฺจ เสขพลานิ อุปนิสฺสาย วิหรติ สทฺธาพลํ หิริพลํ โอตฺตปฺปพลํ วีริยพลํ ปญฺญาพลํ. ตสฺสิมานิ ปญฺจินฺทฺริยานิ อธิมตฺตานิ ปาตุภวนฺติ สทฺธินฺทฺริยํ วีริยินฺทฺริยํ สตินฺทฺริยํ สมาธินฺทฺริยํ ปญฺญินฺทฺริยํ. โส อิเมสํ ปญฺจนฺนํ อินฺทฺริยานํ อธิมตฺตตฺตา ทิฏฺเฐว ธมฺเม สสงฺขาร\-ปรินิพฺพายี โหติ. เอวํ โข ภิกฺขเว ปุคฺคโล ทิฏฺเฐว ธมฺเม สสงฺขาร\-ปรินิพฺพายี โหติ.}

\prose{กถญฺจ ภิกฺขเว ปุคฺคโล กายสฺส เภทา สสงฺขาร\-ปรินิพฺพายี โหติ? อิธ ภิกฺขเว ภิกฺขุ อสุภานุปสฺสี กาเย วิหรติ อาหาเร ปฏิกูลสญฺญี สพฺพโลเก อนภิรติ\-สญฺญี สพฺพ\-สงฺขาเรสุ อนิจฺจานุปสฺสี, มรณสญฺญา โข ปนสฺส อชฺฌตฺตํ สูปฏฺฐิตา โหติ. โส อิมานิ ปญฺจ เสขพลานิ อุปนิสฺสาย วิหรติ สทฺธาพลํ หิริพลํ โอตฺตปฺปพลํ วีริยพลํ ปญฺญาพลํ. ตสฺสิมานิ ปญฺจินฺทฺริยานิ มุทูนิ ปาตุภวนฺติ สทฺธินฺทฺริยํ วีริยินฺทฺริยํ สตินฺทฺริยํ สมาธินฺทฺริยํ ปญฺญินฺทฺริยํ. โส อิเมสํ ปญฺจนฺนํ อินฺทฺริยานํ มุทุตฺตา กายสฺส เภทา สสงฺขาร\-ปรินิพฺพายี โหติ. เอวํ โข ภิกฺขเว ปุคฺคโล กายสฺส เภทา สสงฺขาร\-ปรินิพฺพายี โหติ.}

\prose{กถญฺจ ภิกฺขเว ปุคฺคโล ทิฏฺเฐว ธมฺเม อสงฺขารปรินิพฺพายี โหติ? อิธ ภิกฺขเว ภิกฺขุ วิวิจฺเจว กาเมหิ ฯเปฯ ปฐมํ ฌานํ ฯเปฯ ทุติยํ ฌานํ ฯเปฯ ตติยํ ฌานํ ฯเปฯ จตุตฺถํ ฌานํ อุปสมฺปชฺช วิหรติ. โส อิมานิ ปญฺจ เสขพลานิ อุปนิสฺสาย วิหรติ สทฺธาพลํ ฯเปฯ ปญฺญาพลํ. ตสฺสิมานิ ปญฺจินฺทฺริยานิ อธิมตฺตานิ ปาตุภวนฺติ สทฺธินฺทฺริยํ ฯเปฯ ปญฺญินฺทฺริยํ. โส อิเมสํ ปญฺจนฺนํ อินฺทฺริยานํ อธิมตฺตตฺตา ทิฏฺเฐว ธมฺเม อสงฺขารปรินิพฺพายี โหติ. เอวํ โข ภิกฺขเว ปุคฺคโล ทิฏฺเฐว ธมฺเม อสงฺขารปรินิพฺพายี โหติ.}

\prose{กถญฺจ ภิกฺขเว ปุคฺคโล กายสฺส เภทา อสงฺขารปรินิพฺพายี โหติ? อิธ ภิกฺขเว ภิกฺขุ วิวิจฺเจว กาเมหิ ฯเปฯ ปฐมํ ฌานํ ฯเปฯ ทุติยํ ฌานํ ฯเปฯ ตติยํ ฌานํ ฯเปฯ จตุตฺถํ ฌานํ อุปสมฺปชฺช วิหรติ. โส อิมานิ ปญฺจ เสขพลานิ อุปนิสฺสาย วิหรติ สทฺธาพลํ หิริพลํ โอตฺตปฺปพลํ วีริยพลํ ปญฺญาพลํ. ตสฺสิมานิ ปญฺจินฺทฺริยานิ ฯเปฯ ปญฺญินฺทฺริยํ. โส อิเมสํ ปญฺจนฺนํ อินฺทฺริยานํ มุทุตฺตา กายสฺส เภทา อสงฺขารปรินิพฺพายี โหติ. เอวํ โข ภิกฺขเว ปุคฺคโล กายสฺส เภทา อสงฺขารปรินิพฺพายี โหติ. อิเม โข ภิกฺขเว จตฺตาโร ปุคฺคลา สนฺโต สํวิชฺชมานา โลกสฺมินฺติ.  นวมํ.}

\csromanmatikaanchor{mtk.206}
\csromantocmarkref{subsection}{10. ยุคนทฺธสุตฺต}{mtk.206}
\bookmark[dest={mtk.206},level=2]{10. ยุคนทฺธสุตฺต}
\csromanmarkh{(17)}\csromanmarkhii{2. ปฏิปทาวคฺค 10. ยุคนทฺธ\-สุตฺต}\csromanheadernomark{2}{(17) 2. ปฏิปทาวคฺค 10. ยุคนทฺธ\-สุตฺต}
\proseitem{170}{\csromanfolio{475}เอวํ เม สุตํ\csromandash{}เอกํ สมยํ อายสฺมา อานนฺโท โกสมฺพิยํ วิหรติ โฆสิตาราเม. ตตฺร โข อายสฺมา อานนฺโท ภิกฺขู อามนฺเตสิ “อาวุโส ภิกฺขเว”ติ. “อาวุโส”ติ โข เต ภิกฺขู อายสฺมโต อานนฺทสฺส ปจฺจสฺโสสุํ. อายสฺมา อานนฺโท เอตทโวจ\csromandash{}}

\prose{โย หิ โกจิ อาวุโส ภิกฺขุ วา ภิกฺขุนี วา มม สนฺติเก อรหตฺต\-ปฺปตฺติํ\csromansharedfootnote{24e57fe28150fca6}{อรหตฺตปฺปตฺตํ (ก) ขุ 9. 283 ปิฏฺเฐ ปฏิสมฺภิทามคฺเคปิ.} พฺยากโรติ, สพฺโพ โส จตูหิ มคฺเคหิ เอเตสํ วา อญฺญตเรน.}

\prose{กตเมหิ จตูหิ? อิธ อาวุโส ภิกฺขุ สมถ\-ปุพฺพงฺคมํ วิปสฺสนํ ภาเวติ, ตสฺส สมถ\-ปุพฺพงฺคมํ วิปสฺสนํ ภาวยโต มคฺโค สญฺชายติ, โส ตํ มคฺคํ อาเสวติ ภาเวติ พหุลีกโรติ. ตสฺส ตํ มคฺคํ อาเสวโต ภาวยโต พหุลีกโรโต สํโยชนานี ปหียนฺติ, อนุสยา พฺยนฺตีโหนฺติ.}

\prose{ปุน จปรํ อาวุโส ภิกฺขุ วิปสฺสนา\-ปุพฺพงฺคมํ สมถํ ภาเวติ, ตสฺส วิปสฺสนา\-ปุพฺพงฺคมํ สมถํ ภาวยโต มคฺโค สญฺชายติ, โส ตํ มคฺคํ อาเสวติ ภาวติ พหุลีกโรติ. ตสฺส ตํ มคฺคํ อาเสวโต ภาวยโต พหุลีกโรโต สํโยชนานิ ปหียนฺติ, อนุสยา พฺยนฺตีโหนฺติ.}

\prose{ปุน จปรํ อาวุโส ภิกฺขุ สมถ\-วิปสฺสนํ ยุคนทฺธํ ภาเวติ, ตสฺส สมถ\-วิปสฺสนํ ยุคนทฺธํ ภาวยโต มคฺโค สญฺชายติ, โส ตํ มคฺคํ อาเสวติ ภาเวติ พหุลีกโรติ. ตสฺส ตํ มคฺคํ อาเสวโต ภาวยโต พหุลีกโรโต สํโยชนานิ ปหียนฺติ, อนุสยา พฺยนฺตีโหนฺติ.}

\prose{ปุน จปรํ อาวุโส ภิกฺขุโน ธมฺมุ\-ทฺธจฺจ\-วิคฺคหิตํ มานสํ โหติ, โหติ โส อาวุโส สมโย ยํ ตํ จิตฺตํ อชฺฌตฺตเมว สนฺติฏฺฐติ สนฺนิสีทติ เอโกทิ โหติ สมาธิยติ, ตสฺส มคฺโค สญฺชายติ, โส ตํ มคฺคํ อาเสวติ ภาเวติ พหุลีกโรติ. ตสฺส ตํ มคฺคํ อาเสวโต \csromanfolio{476}ภาวยโต พหุลีกโรโต สํโยชนานิ ปหียนฺติ, อนุสยา พฺยนฺตีโหนฺติ.}

\prosewithclosermidruled{โย หิ โกจิ อาวุโส ภิกฺขุ วา ภิกฺขุนี วา มม สนฺติเก อรหตฺต\-ปฺปตฺติํ พฺยากโรติ, สพฺโพ โส อิเมหิ จตูหิ มคฺเคหิ เอเตสํ วา อญฺญตเรนาติ.}{ปฏิปทาวคฺโค ทุติโย.}
\csromancenter{\textbf{ตสฺสุทฺทานํ}}
\setlength{\csromangathapagewidth}{0pt}
\setlength{\csromangathawakwidth}{0pt}
\csromangathameasuregroup{สํขิตฺตํ วิตฺถาราสุภํ, \\ โมคฺคลฺลาโน สาริปุตฺโต,}{\csromangathabat{สํขิตฺตํ วิตฺถาราสุภํ,}{เทฺว ขมา อุภเยน จ.} \\ \csromangathabat{โมคฺคลฺลาโน สาริปุตฺโต,}{สสงฺขารํ ยุคนทฺเธน จาติ.}}
\csromangathasetleft{สํขิตฺตํ วิตฺถาราสุภํ, \\ โมคฺคลฺลาโน สาริปุตฺโต,}
\csromangathagroup{\csromangathastanza{\csromangathabat{สํขิตฺตํ วิตฺถาราสุภํ,}{เทฺว ขมา อุภเยน จ.} \\ \csromangathabat{โมคฺคลฺลาโน สาริปุตฺโต,}{สสงฺขารํ ยุคนทฺเธน จาติ.}}}

\csromanmatikaanchor{mtk.207}
\csromantocmarknopage{section}{(18) 3. สญฺเจตนิยวคฺค}{mtk.207}
\bookmark[dest={mtk.207},level=1]{(18) 3. สญฺเจตนิยวคฺค}
\csromanmarkh{(18)}\csromanmarkhii{3. สญฺเจตนิย\-วคฺค}\csromanheadernomark{1}{(18) 3. สญฺเจตนิย\-วคฺค}
\csromanmatikaanchor{mtk.208}
\csromantocmarkref{subsection}{1. เจตนาสุตฺต}{mtk.208}
\bookmark[dest={mtk.208},level=2]{1. เจตนาสุตฺต}
\csromanheader{2}{1. \textbf{เจตนาสุตฺต}}
\proseitem{171}{\csromansymbolfootnote{*}{อภิ 4. 290 ปิฏฺเฐปิ.}กาเย วา ภิกฺขเว สติ กาย\-สญฺเจตนา\-เหตุ อุปฺปชฺชติ อชฺฌตฺตํ สุขทุกฺขํ, วาจาย วา ภิกฺขเว สติ วจีสญฺเจตนา\-เหตุ อุปฺปชฺชติ อชฺฌตฺตํ สุขทุกฺขํ, มเน วา ภิกฺขเว สติ มโน\-สญฺเจตนา\-เหตุ อุปฺปชฺชติ อชฺฌตฺตํ สุขทุกฺขํ, อวิชฺชา\-ปจฺจยาว.}

\prose{สามํ วา ตํ ภิกฺขเว กายสงฺขารํ อภิสงฺขโรติ, ยํปจฺจยา’สฺส ตํ อุปฺปชฺชติ อชฺฌตฺตํ สุขทุกฺขํ. ปเร วา’สฺส\csromansharedfootnote{475643e488b5ee6c}{ปเร วา ตสฺส (ก)} ตํ ภิกฺขเว กายสงฺขารํ อภิส\-งฺขโรนฺติ, ยํปจฺจยา’สฺส ตํ อุปฺปชฺชติ อชฺฌตฺตํ สุขทุกฺขํ. สมฺปชาโน วา ตํ ภิกฺขเว กายสงฺขารํ อภิสงฺขโรติ, ยํปจฺจยา’สฺส ตํ อุปฺปชฺชติ อชฺฌตฺตํ สุขทุกฺขํ. อสมฺปชาโน วา ตํ ภิกฺขเว กายสงฺขารํ อภิสงฺขโรติ, ยํปจฺจยา’สฺส ตํ อุปฺปชฺชติ อชฺฌตฺตํ สุขทุกฺขํ.}

\prose{สามํ วา ตํ ภิกฺขเว วจีสงฺขารํ อภิสงฺขโรติ, ยํปจฺจยา’สฺส ตํ อุปฺปชฺชติ อชฺฌตฺตํ สุขทุกฺขํ. ปเร วา’สฺส ตํ ภิกฺขเว วจีสงฺขารํ อภิส\-งฺขโรนฺติ, ยํปจฺจยา’สฺส ตํ อุปฺปชฺชติ อชฺฌตฺตํ สุขทุกฺขํ. สมฺปชาโน วา ตํ ภิกฺขเว}

\vspace{\tipitakaparskip}
\csromancenter{(18) 3. สญฺเจตนิย\-วคฺค วจีสงฺขารํ อภิสงฺขโรติ, ยํปจฺจยาสฺส ตํ อุปฺปชฺชติ อชฺฌตฺตํ สุขทุกฺขํ. อสมฺปชาโน วา ตํ ภิกฺขเว วจีสงฺขารํ อภิสงฺขโรติ, ยํปจฺจยาสฺส ตํ อุปฺปชฺชติ อชฺฌตฺตํ สุขทุกฺขํ.}
\prose{\csromanfolio{477}สามํ วา ตํ ภิกฺขเว มโนสงฺขารํ อภิสงฺขโรติ, ยํปจฺจยา’สฺส ตํ อุปฺปชฺชติ อชฺฌตฺตํ สุขทุกฺขํ. ปเร วา’สฺส ตํ ภิกฺขเว มโนสงฺขารํ อภิส\-งฺขโรนฺติ, ยํปจฺจยา’สฺส ตํ อุปฺปชฺชติ อชฺฌตฺตํ สุขทุกฺขํ. สมฺปชาโน วา ตํ ภิกฺขเว มโนสงฺขารํ อภิสงฺขโรติ, ยํปจฺจยา’สฺส ตํ อุปฺปชฺชติ อชฺฌตฺตํ สุขทุกฺขํ. อสมฺปชาโน วา ตํ ภิกฺขเว มโนสงฺขารํ อภิสงฺขโรติ, ยํปจฺจยา’สฺส ตํ อุปฺปชฺชติ อชฺฌตฺตํ สุขทุกฺขํ.}

\prose{อิเมสุ ภิกฺขเว ธมฺเมสุ อวิชฺชา อนุปติตา, อวิชฺชายเตฺวว อเสส\-วิราค\-นิโรธา โส กาโย น โหติ, ยํปจฺจยา’สฺส ตํ อุปฺปชฺชติ อชฺฌตฺตํ สุขทุกฺขํ. สา วาจา น โหติ, ยํปจฺจยา’สฺส ตํ อุปฺปชฺชติ อชฺฌตฺตํ สุขทุกฺขํ. โส มโน น โหติ, ยํปจฺจยา’สฺส ตํ อุปฺปชฺชติ อชฺฌตฺตํ สุขทุกฺขํ. เขตฺตํ ตํ น โหติ ฯเปฯ วตฺถุ ตํ\csromansharedfootnote{f7c75958ba29c52a}{วตฺถุํ ตํ (สพฺพตฺถ)} น โหติ ฯเปฯ อยตนํ ตํ น โหติ ฯเปฯ อธิกรณํ ตํ น โหติ, ยํปจฺจยา’สฺส ตํ อุปฺปชฺชติ อชฺฌตฺตํ สุขทุกฺขนฺติ.}

\prose{จตฺตาโรเม ภิกฺขเว อตฺตภาว\-ปฏิลาภา. กตเม จตฺตาโร? อตฺถิ ภิกฺขเว อตฺตภาว\-ปฏิลาโภ, ยสฺมิํ อตฺตภาว\-ปฏิลาเภ อตฺตสญฺเจตนา กมติ โน ปรสญฺเจตนา. อตฺถิ ภิกฺขเว อตฺตภาว\-ปฏิลาโภ, ยสฺมิํ อตฺตภาว\-ปฏิลาเภ ปรสญฺเจตนา กมติ โน อตฺตสญฺเจตนา. อตฺถิ ภิกฺขเว อตฺตภาว\-ปฏิลาโภ, ยสฺมิํ อตฺตภาว\-ปฏิลาเภ อตฺตสญฺเจตนา จ กมติ ปรสญฺเจตนา จ. อตฺถิ ภิกฺขเว อตฺตภาว\-ปฏิลาโภ, ยสฺมิํ อตฺตภาว\-ปฏิลาเภ เนว\-ตฺต\-สญฺเจตนา กมติ โน ปรสญฺเจตนา. อิเม โข ภิกฺขเว จตฺตาโร อตฺตภาว\-ปฏิลาภาติ.}

\prose{เอวํ วุตฺเต อายสฺมา สาริปุตฺโต ภควนฺตํ เอตทโวจ\csromandash{}อิมสฺส โข อหํ ภนฺเต ภควตา สํขิตฺเตน ภาสิตสฺส เอวํ วิตฺถาเรน อตฺถํ อาชานามิ\csromandash{}ตตฺร ภนฺเต ยายํ อตฺตภาว\-ปฏิลาโภ, ยสฺมิํ อตฺตภาว\-ปฏิลาเภ อตฺตสญฺเจตนา กมติ โน ปรสญฺเจตนา, อตฺต\-สญฺเจตนา\-เหตุ เตสํ สตฺตานํ ตมฺหา กายา จุติ โหติ. ตตฺร ภนฺเต ยายํ \csromanfolio{478}อตฺตภาว\-ปฏิลาโภ, ยสฺมิํ อตฺตภาว\-ปฏิลาเภ ปรสญฺเจตนา กมติ โน อตฺตสญฺเจตนา, ปรสญฺเจตนา\-เหตุ เตสํ สตฺตานํ ตมฺหา กายา จุติ โหติ. ตตฺร ภนฺเต ยายํ อตฺตภาว\-ปฏิลาโภ, ยสฺมิํ อตฺตภาว\-ปฏิลาเภ อตฺตสญฺเจตนา จ กมติ ปรสญฺเจตนา จ. อตฺตสญฺเจตนา จ ปรสญฺเจตนา จ เหตุ เตสํ สตฺตานํ ตมฺหา กายา จุติ โหติ. ตตฺร ภนฺเต ยายํ อตฺตภาว\-ปฏิลาโภ, ยสฺมิํ อตฺตภาว\-ปฏิลาเภ เนว อตฺตสญฺเจตนา กมติ โน ปรสญฺเจตนา. กตเม เตน เทวา ทฏฺฐพฺพาติ. เนว\-สญฺญา\-นา\-สญฺญา\-ยตนู\-ปคา สาริปุตฺต เทวา เตน ทฏฺฐพฺพาติ.}

\prose{โก นุ โข ภนฺเต เหตุ โก ปจฺจโย, เยน มิเธกจฺเจ สตฺตา ตมฺหา กายา จุตา อาคามิโน โหนฺติ อาคนฺตาโร อิตฺถตฺตํ. โก ปน ภนฺเต เหตุ โก ปจฺจโย, เยน มิเธกจฺเจ สตฺตา ตมฺหา กายา จุตา อนาคามิโน โหนฺติ อนาคนฺตาโร อิตฺถตฺตนฺติ. อิธ สาริปุตฺต เอกจฺจสฺส ปุคฺคลสฺส โอรมฺภาคิยานิ สํโยชนานิ อปฺปหีนานิ โหนฺติ. โส ทิฏฺเฐว ธมฺเม เนว\-สญฺญา\-นา\-สญฺญา\-ยตนํ อุปสมฺปชฺช วิหรติ, โส ตทสฺสาเทติ ตํ นิกาเมติ, เตน จ วิตฺติํ อาปชฺชติ. ตตฺถ ฐิโต ตทธิมุตฺโต ตพฺพหุล\-วิหารี อปริหีโน กาลํ กุรุมาโน เนว\-สญฺญา\-นา\-สญฺญา\-ยตนู\-ปคานํ เทวานํ สหพฺยตํ อุปปชฺชติ. โส ตโต จุโต อาคามี โหติ อาคนฺตา อิตฺถตฺตํ.}

\prose{อิธ ปน สาริปุตฺต เอกจฺจสฺส ปุคฺคลสฺส โอรมฺภาคิยานิ สํโยชนานิ ปหีนานิ โหนฺติ. โส ทิฏฺเฐว ธมฺเม เนว\-สญฺญา\-นา\-สญฺญา\-ยตนํ อุปสมฺปชฺช วิหรติ. โส ตทสฺสาเทติ ตํ นิกาเมติ, เตน จ วิตฺติํ อาปชฺชติ. ตตฺถ ฐิโต ตทธิมุตฺโต ตพฺพหุล\-วิหารี อปริหีโน กาลํ กุรุมาโน เนว\-สญฺญา\-นา\-สญฺญา\-ยตนู\-ปคานํ เทวานํ สหพฺยตํ อุปปชฺชติ, โส ตโต จุโต อนาคามี โหติ อนาคนฺตา อิตฺถตฺตํ.}

\prose{อยํ โข สาริปุตฺต เหตุ อยํ ปจฺจโย, เยน มิเธกจฺเจ สตฺตา ตมฺหา กายา จุตา อาคามิโน โหนฺติ อาคนฺตาโร อิตฺถตฺตํ. อยํ ปน สาริปุตฺต เหตุ อยํ ปจฺจโย, เยน มิเธกจฺเจ สตฺตา ตมฺหา กายา จุตา อนาคามิโน โหนฺติ อนาคนฺตาโร อิตฺถตฺตนฺติ.  ปฐมํ.}

\csromanmatikaanchor{mtk.209}
\csromantocmarkref{subsection}{2. วิภตฺติสุตฺต}{mtk.209}
\bookmark[dest={mtk.209},level=2]{2. วิภตฺติสุตฺต}
\csromanmarkh{(18)}\csromanmarkhii{3. สญฺเจตนิย\-วคฺค 2. วิภตฺติสุตฺต}\csromanheadernomark{2}{(18) 3. สญฺเจตนิย\-วคฺค 2. วิภตฺติสุตฺต}
\proseitem{172}{\csromanfolio{479}ตตฺร โข อายสฺมา สาริปุตฺโต ภิกฺขู อามนฺเตสิ “อาวุโส ภิกฺขเว”ติ. “อาวุโส”ติ โข เต ภิกฺขู อายสฺมโต สาริปุตฺตสฺส ปจฺจสฺโสสุํ. อายสฺมา สาริปุตฺโต เอตทโวจ\csromandash{}}

\prose{อทฺธมาสู\-ปสมฺปนฺเนน เม อาวุโส อตฺถ\-ปฏิสมฺภิทา สจฺฉิกตา โอธิโส พฺยญฺชนโส, ตมหํ อเนก\-ปริยาเยน อาจิกฺขามิ เทเสมิ ปญฺญาเปมิ ปฏฺฐเปมิ วิวรามิ วิภชามิ อุตฺตานีกโรมิ. ยสฺส โข ปนสฺส กงฺขา วา วิมติ วา, โส มํ ปเญฺหน, อหํ เวยฺยากรเณน. สมฺมุขีภูโต โน สตฺถา, โย โน ธมฺมานํ สุกุสโล.}

\prose{อทฺธมาสู\-ปสมฺปนฺเนน เม อาวุโส ธมฺม\-ปฏิสมฺภิทา สจฺฉิกตา โอธิโส พฺยญฺชนโส, ตมหํ อเนก\-ปริยาเยน อาจิกฺขามิ เทเสมิ ปญฺญาเปมิ ปฏฺฐเปมิ วิวรามิ วิภชามิ อุตฺตานีกโรมิ. ยสฺส โข ปนสฺส กงฺขา วา วิมติ วา, โส มํ ปเญฺหน, อหํ เวยฺยากรเณน. สมฺมุขีภูโต โน สตฺถา, โย โน ธมฺมานํ สุกุสโล.}

\prose{อทฺธมาสู\-ปสมฺปนฺเนน เม อาวุโส นิรุตฺติ\-ปฏิสมฺภิทา สจฺฉิกตา โอธิโส พฺยญฺชนโส, ตมหํ อเนก\-ปริยาเยน อาจิกฺขามิ เทเสมิ ปญฺญาเปมิ ปฏฺฐเปมิ วิวรามิ วิภชามิ อุตฺตานีกโรมิ. ยสฺส โข ปนสฺส กงฺขา วา วิมติ วา, โส มํ ปเญฺหน, อหํ เวยฺยากรเณน. สมฺมุขีภูโต โน สตฺถา, โย โน ธมฺมานํ สุกุสโล.}

\prose{อทฺธมาสู\-ปสมฺปนฺเนน เม อาวุโส ปฏิภาน\-ปฏิสมฺภิทา สจฺฉิกตา โอธิโส พฺยญฺชนโส, ตมหํ อเนก\-ปริยาเยน อาจิกฺขามิ เทเสมิ ปญฺญาเปมิ ปฏฺฐเปมิ วิวรามิ วิภชามิ อุตฺตานีกโรมิ. ยสฺส โข ปนสฺส กงฺขา วา วิมติ วา, โส มํ ปเญฺหน, อหํ เวยฺยากรเณน. สมฺมุขีภูโต โน สตฺถา, โย โน ธมฺมานํ สุกุสโลติ.  ทุติยํ.}

\csromanmatikaanchor{mtk.210}
\csromantocmarkref{subsection}{3. มหาโกฏฺฐิกสุตฺต}{mtk.210}
\bookmark[dest={mtk.210},level=2]{3. มหาโกฏฺฐิกสุตฺต}
\csromanheader{2}{3. มหาโกฏฺฐิก\-สุตฺต}
\proseitem{173}{อถ โข อายสฺมา มหาโกฏฺฐิโก เยนายสฺมา สาริปุตฺโต เตนุปสงฺกมิ, อุปสงฺกมิตฺวา อายสฺมตา สาริปุตฺเตน สทฺธิํ สมฺโมทิ, สมฺโมทนียํ กถํ สารณียํ วีติสาเรตฺวา เอกมนฺตํ นิสีทิ, \csromanfolio{480}เอกมนฺตํ นิสินฺโน โข อายสฺมา มหาโกฏฺฐิโก อายสฺมนฺตํ สาริปุตฺตํ เอตทโวจ\csromandash{}}

\prose{“ฉนฺนํ อาวุโส ผสฺสา\-ยตนานํ อเสส\-วิราค\-นิโรธา อตฺถญฺญํ กิญฺจีติ. มา เหวํ อาวุโส. ฉนฺนํ อาวุโส ผสฺสา\-ยตนานํ อเสส\-วิราค\-นิโรธา นตฺถญฺญํ กิญฺจี”ติ. มา เหวํ อาวุโส. ฉนฺนํ อาวุโส ผสฺสา\-ยตนานํ อเสส\-วิราค\-นิโรธา อตฺถิ จ นตฺถิ จ อญฺญํ กิญฺจีติ. มา เหวํ อาวุโส. ฉนฺนํ อาวุโส ผสฺสา\-ยตนานํ อเสส\-วิราค\-นิโรธา เนวตฺถิ โน นตฺถญฺญํ กิญฺจีติ. มา เหวํ อาวุโส.}

\prose{“ฉนฺนํ อาวุโส ผสฺสา\-ยตนานํ อเสส\-วิราค\-นิโรธา อตฺถญฺญํ กิญฺจี”ติ อิติ ปุฏฺโฐ สมาโน “มา เหวํ อาวุโส”ติ วเทสิ. “ฉนฺนํ อาวุโส ผสฺสา\-ยตนานํ อเสส\-วิราค\-นิโรธา นตฺถญฺญํ กิญฺจี”ติ อิติ ปุฏฺโฐ สมาโน “มา เหวํ อาวุโส”ติ วเทสิ. “ฉนฺนํ อาวุโส ผสฺสา\-ยตนานํ อเสส\-วิราค\-นิโรธา อตฺถิ จ นตฺถิ จ อญฺญํ กิญฺจี”ติ อิติ ปุฏฺโฐ สมาโน “มา เหวํ อาวุโส”ติ วเทสิ. “ฉนฺนํ อาวุโส ผสฺสา\-ยตนานํ อเสส\-วิราค\-นิโรธา เนวตฺถิ โน นตฺถญฺญํ กิญฺจี”ติ อิติ ปุฏฺโฐ สมาโน “มา เหวํ อาวุโส”ติ วเทสิ. ยถา กถํ ปน อาวุโส อิมสฺส ภาสิตสฺส อตฺโถ ทฏฺฐพฺโพติ.}

\prose{“ฉนฺนํ อาวุโส ผสฺสา\-ยตนานํ อเสส\-วิราค\-นิโรธา อตฺถญฺญํ กิญฺจี”ติ อิติ วทํ อปฺปปญฺจํ ปปญฺเจติ. “ฉนฺนํ อาวุโส ผสฺสา\-ยตนานํ อเสส\-วิราค\-นิโรธา นตฺถญฺญํ กิญฺจี”ติ อิติ วทํ อปฺปปญฺจํ ปปญฺเจติ. “ฉนฺนํ อาวุโส ผสฺสา\-ยตนานํ อเสส\-วิราค\-นิโรธา อตฺถิ จ นตฺถิ จ อญฺญํ กิญฺจี”ติ อิติ วทํ อปฺปปญฺจํ ปปญฺเจติ. “ฉนฺนํ อาวุโส ผสฺสา\-ยตนานํ อเสส\-วิราค\-นิโรธา เนวตฺถิ โน นตฺถญฺญํ กิญฺจี”ติ อิติ วทํ อปฺปปญฺจํ ปปญฺเจติ. ยาวตา อาวุโส ฉนฺนํ ผสฺสา\-ยตนานํ คติ, ตาวตาปปญฺจสฺส คติ. ยาวตา ปปญฺจสฺส คติ, ตาวตา ฉนฺนํ ผสฺสา\-ยตนานํ คติ. ฉนฺนํ อาวุโส ผสฺสา\-ยตนานํ อเสส\-วิราค\-นิโรธา ปปญฺจนิโรโธ ปปญฺจ\-วูปสโมติ.  ตติยํ.}

\csromanmatikaanchor{mtk.211}
\csromantocmarkref{subsection}{4. อานนฺทสุตฺต}{mtk.211}
\bookmark[dest={mtk.211},level=2]{4. อานนฺทสุตฺต}
\csromanheader{2}{4. \textbf{อานนฺทสุตฺต}}
\proseitem{174}{อถ โข อายสฺมา อานนฺโท เยนายสฺมา มหาโกฏฺฐิโก เตนุปสงฺกมิ, อุปสงฺกมิตฺวา อายสฺมตา มหาโกฏฺฐิเกน}

\vspace{\tipitakaparskip}
\csromancenter{(18) 3. สญฺเจตนิย\-วคฺค สทฺธิํ สมฺโมทิ, สมฺโมทนียํ กถํ สารณียํ วีติสาเรตฺวา เอกมนฺตํ นิสีทิ, เอกมนฺตํ นิสินฺโน โข อายสฺมา อานนฺโท อายสฺมนฺตํ มหาโกฏฺฐิตํ เอตทโวจ\csromandash{}}
\prose{\csromanfolio{481}“ฉนฺนํ อาวุโส ผสฺสา\-ยตนานํ อเสส\-วิราค\-นิโรธา อตฺถญฺญํ กิญฺจี”ติ. มา เหวํ อาวุโส. ฉนฺนํ อาวุโส ผสฺสา\-ยตนานํ อเสส\-วิราค\-นิโรธา นตฺถญฺญํ กิญฺจีติ. มา เหวํ อาวุโส. ฉนฺนํ อาวุโส ผสฺสา\-ยตนานํ อเสส\-วิราค\-นิโรธา อตฺถิ จ นตฺถิ จ อญฺญํ กิญฺจีติ. มา เหวํ อาวุโส. ฉนฺนํ อาวุโส ผสฺสา\-ยตนานํ อเสส\-วิราค\-นิโรธา เนวตฺถิ โน นตฺถญฺญํ กิญฺจีติ. มา เหวํ อาวุโส.}

\prose{“ฉนฺนํ อาวุโส ผสฺสา\-ยตนานํ อเสส\-วิราค\-นิโรธา อตฺถญฺญํ กิญฺจี”ติ อิติ ปุฏฺโฐ สมาโน “มา เหวํ อาวุโส”ติ วเทสิ. “ฉนฺนํ อาวุโส ผสฺสา\-ยตนานํ อเสส\-วิราค\-นิโรธา นตฺถญฺญํ กิญฺจี”ติ อิติ ปุฏฺโฐ สมาโน “มา เหวํ อาวุโส”ติ วเทสิ. “ฉนฺนํ อาวุโส ผสฺสา\-ยตนานํ อเสส\-วิราค\-นิโรธา อตฺถิ จ นตฺถิ จ อญฺญํ กิญฺจี”ติ อิติ ปุฏฺโฐ สมาโน “มา เหวํ อาวุโส”ติ วเทสิ. “ฉนฺนํ อาวุโส ผสฺสา\-ยตนานํ อเสส\-วิราค\-นิโรธา เนวตฺถิ โน นตฺถญฺญํ กิญฺจี”ติ อิติ ปุฏฺโฐ สมาโน “มา เหวํ อาวุโส”ติ วเทสิ. ยถา กถํ ปนาวุโส อิมสฺส ภาสิตสฺส อตฺโถ ทฏฺฐพฺโพติ.}

\prose{“ฉนฺนํ อาวุโส ผสฺสา\-ยตนานํ อเสส\-วิราค\-นิโรธา อตฺถญฺญํ กิญฺจี”ติ อิติ วทํ อปฺปปญฺจํ ปปญฺเจติ. “ฉนฺนํ อาวุโส ผสฺสา\-ยตนานํ อเสส\-วิราค\-นิโรธา นตฺถญฺญํ กิญฺจี”ติ อิติ วทํ อปฺปปญฺจํ ปปญฺเจติ. “ฉนฺนํ อาวุโส ผสฺสา\-ยตนานํ อเสส\-วิราค\-นิโรธา อตฺถิ จ นตฺถิ จ อญฺญํ กิญฺจี”ติ อิติ วทํ อปฺปปญฺจํ ปปญฺเจติ. “ฉนฺนํ อาวุโส ผสฺสา\-ยตนานํ อเสส\-วิราค\-นิโรธา เนวตฺถิ โน นตฺถญฺญํ กิญฺจี”ติ อิติ วทํ อปฺปปญฺจํ ปปญฺเจติ. ยาวตา อาวุโส ฉนฺนํ ผสฺสา\-ยตนานํ คติ, ตาวตา ปปญฺจสฺส คติ. ยาวตา ปปญฺจสฺส คติ, ตาวตา ฉนฺนํ ผสฺสา\-ยตนานํ คติ. ฉนฺนํ อาวุโส ผสฺสา\-ยตนานํ อเสส\-วิราค\-นิโรธา ปปญฺจนิโรโธ ปปญฺจ\-วูปสโมติ.  จตุตฺถํ.}

\csromanpalirecto
\csromanmatikaanchor{mtk.212}
\csromantocmarkref{subsection}{5. อุปวาณสุตฺต}{mtk.212}
\bookmark[dest={mtk.212},level=2]{5. อุปวาณสุตฺต}
\csromanmarkh{จตุกฺกนิปาต\-ปาฬิ}\csromanmarkhii{5. อุปวาณสุตฺต}\csromanheadernomark{2}{จตุกฺกนิปาต\-ปาฬิ 5. อุปวาณสุตฺต}
\proseitem{175}{\csromanfolio{482}อถ โข อายสฺมา อุปวาโณ เยนายสฺมา สาริปุตฺโต เตนุปสงฺกมิ, อุปสงฺกมิตฺวา อายสฺมตา สาริปุตฺเตน สทฺธิํ สมฺโมทิ, สมฺโมทนียํ กถํ สารณียํ วีติสาเรตฺวา เอกมนฺตํ นิสีทิ, เอกมนฺตํ นิสินฺโน โข อายสฺมา อุปวาโณ อายสฺมนฺตํ สาริปุตฺตํ เอตทโวจ\csromandash{}“กิํ นุ โข อาวุโส สาริปุตฺต วิชฺชายนฺตกโร โหตี”ติ. โน หิทํ อาวุโส. กิํ ปนาวุโส สาริปุตฺต จรเณน\-นฺตกโร โหตีติ. โน หิทํ อาวุโส. กิํ ปนาวุโส สาริปุตฺต วิชฺชาจรเณน\-นฺตกโร โหตีติ. โน หิทํ อาวุโส. กิํ ปนาวุโส สาริปุตฺต อญฺญตฺร วิชฺชาจรเณน\-นฺตกโร โหตีติ. โน หิทํ อาวุโส.}

\prose{“กิํ นุ โข อาวุโส สาริปุตฺต วิชฺชายนฺตกโร โหตี”ติ อิติ ปุฏฺโฐ สมาโน “โน หิทํ อาวุโส”ติ วเทสิ. “กิํ ปนาวุโส สาริปุตฺต จรเณน\-นฺตกโร โหตี”ติ อิติ ปุฏฺโฐ สมาโน “โน หิทํ อาวุโส”ติ วเทสิ. “กิํ ปนาวุโส สาริปุตฺต วิชฺชาจรเณน\-นฺตกโร โหตี”ติ อิติ ปุฏฺโฐ สมาโน “โน หิทํ อาวุโส”ติ วเทสิ. “กิํ ปนาวุโส สาริปุตฺต อญฺญตฺร วิชฺชาจรเณน\-นฺตกโร โหตี”ติ อิติ ปุฏฺโฐ สมาโน “โน หิทํ อาวุโส”ติ วเทสิ. ยถา กถํ ปนาวุโส อนฺตกโร โหตีติ.}

\prose{วิชฺชาย เจ อาวุโส อนฺตกโร อภวิสฺส, สอุปาทาโนว สมาโน อนฺตกโร อภวิสฺส. จรเณน เจ อาวุโส อนฺตกโร อภวิสฺส, สอุปาทาโนว สมาโน อนฺตกโร อภวิสฺส. วิชฺชาจรเณน เจ อาวุโส อนฺตกโร อภวิสฺส, สอุปาทาโนว สมาโน อนฺตกโร อภวิสฺส. อญฺญตฺร วิชฺชาจรเณน เจ อาวุโส อนฺตกโร อภวิสฺส, ปุถุชฺชโน อนฺตกโร อภวิสฺส. ปุถุชฺชโน หิ อาวุโส อญฺญตฺร วิชฺชาจรเณน, จรณวิปนฺโน โข อาวุโส ยถาภูตํ น ชานาติ น ปสฺสติ. จรณสมฺปนฺโน ยถาภูตํ ชานาติ ปสฺสติ, ยถาภูตํ ชานํ ปสฺสํ อนฺตกโร โหตีติ.  ปญฺจมํ.}

\csromanmatikaanchor{mtk.213}
\csromantocmarkref{subsection}{6. อายาจนสุตฺต}{mtk.213}
\bookmark[dest={mtk.213},level=2]{6. อายาจนสุตฺต}
\csromanheader{2}{6. \textbf{อายาจนสุตฺต}}
\proseitem{176}{สทฺโธ ภิกฺขเว ภิกฺขุ เอวํ สมฺมา อายาจมาโน อายาเจยฺย “ตาทิโส โหมิ, ยาทิสา สาริปุตฺต\-โมคฺคลฺลานา”ติ\csromansharedfootnote{19dc6788d00ab0a6}{88 ปิฏฺเฐปิ อิทํ สุตฺตํ อาคตํ.}. เอสา}

\vspace{\tipitakaparskip}
\csromancenter{(18) 3. สญฺเจตนิย\-วคฺค ภิกฺขเว ตุลา เอตํ ปมาณํ มม สาวกานํ ภิกฺขูนํ ยทิทํ สาริปุตฺต\-โมคฺคลฺลานา.}
\prose{\csromanfolio{483}สทฺธา ภิกฺขเว ภิกฺขุนี เอวํ สมฺมา อายาจมานา อายาเจยฺย “ตาทิสา โหมิ, ยาทิสา เขมา จ ภิกฺขุนี อุปฺปลวณฺณา จา”ติ. เอสา ภิกฺขเว ตุลา เอตํ ปมาณํ มม สาวิกานํ ภิกฺขุนีนํ ยทิทํ เขมา จ ภิกฺขุนี อุปฺปลวณฺณา จ.}

\prose{สทฺโธ ภิกฺขเว อุปาสโก เอวํ สมฺมา อายาจมาโน อายาเจยฺย “ตาทิโส โหมิ, ยทิโส จิตฺโต จ คหปติ หตฺถโก จ อาฬวโก”ติ. เอสา ภิกฺขเว ตุลา เอตํ ปมาณํ มม สาวกานํ อุปาสกานํ ยทิทํ จิตฺโต จ คหปติ หตฺถโก จ อาฬวโก.}

\prose{สทฺธา ภิกฺขเว อุปาสิกา เอวํ สมฺมา อายาจมานา อายาเจยฺย “ตาทิสา โหมิ, ยาทิสา ขุชฺชุตฺตรา จ อุปาสิกา เวฬุกณฺฑกิยา จ นนฺทมาตา”ติ. เอสา ภิกฺขเว ตุลา เอตํ ปมาณํ มม สาวิกานํ อุปาสิกานํ ยทิทํ ขุชฺชุตฺตรา จ อุปาสิกา เวฬุกณฺฑกิยา จ นนฺทมาตาติ.  ฉฏฺฐํ.}

\csromanmatikaanchor{mtk.214}
\csromantocmarkref{subsection}{7. ราหุลสุตฺต}{mtk.214}
\bookmark[dest={mtk.214},level=2]{7. ราหุลสุตฺต}
\csromanheader{2}{7. \textbf{ราหุลสุตฺต}}
\proseitem{177}{อถ โข อายสฺมา ราหุโล เยน ภควา เตนุปสงฺกมิ, อุปสงฺกมิตฺวา ภควนฺตํ อภิวาเทตฺวา เอกมนฺตํ นิสีทิ, เอกมนฺตํ นิสินฺนํ โข อายสฺมนฺตํ ราหุลํ ภควา เอตทโวจ\csromandash{}}

\prose{ยา จ ราหุล อชฺฌตฺติกา ปถวีธาตุ ยา จ พาหิรา ปถวีธาตุ, ปถวีธาตุเรเวสา, ตํ “เนตํ มม, เนโสหมสฺมิ, น เมโส อตฺตา”ติ เอวเมตํ ยถาภูตํ สมฺมปฺปญฺญาย ทฏฺฐพฺพํ, เอวเมตํ ยถาภูตํ สมฺมปฺปญฺญาย ทิสฺวา ปถวีธาตุยา นิพฺพินฺทติ, ปถวีธาตุยา จิตฺตํ วิราเชติ.}

\prose{ยา จ ราหุล อชฺฌตฺติกา อาโปธาตุ ยา จ พาหิรา อาโปธาตุ, อาโปธาตุเรเวสา, ตํ “เนตํ มม, เนโสหมสฺมิ, น เมโส อตฺตา”ติ เอวเมตํ ยถาภูตํ สมฺมปฺปญฺญาย ทฏฺฐพฺพํ, เอวเมตํ ยถาภูตํ สมฺมปฺปญฺญาย ทิสฺวา อาโปธาตุยา นิพฺพินฺทติ, อาโปธาตุยา จิตฺตํ วิราเชติ.}

\prose{\csromanfolio{484}ยา จ ราหุล อชฺฌตฺติกา เตโชธาตุ ยา จ พาหิรา เตโชธาตุ, เตโชธาตุเรเวสา, ตํ “เนตํ มม, เนโสหมสฺมิ, น เมโส อตฺตา”ติ เอวเมตํ ยถาภูตํ สมฺมปฺปญฺญาย ทฏฺฐพฺพํ, เอวเมตํ ยถาภูตํ สมฺมปฺปญฺญาย ทิสฺวา เตโชธาตุยา นิพฺพินฺทติ, เตโชธาตุยา จิตฺตํ วิราเชติ.}

\prose{ยา จ ราหุล อชฺฌตฺติกา วาโยธาตุ ยา จ พาหิรา วาโยธาตุ วาโยธาตุเรเวสา, ตํ “เนตํ มม, เนโสหมสฺมิ, น เมโส อตฺตา”ติ เอวเมตํ ยถาภูตํ สมฺมปฺปญฺญาย ทฏฺฐพฺพํ, เอวเมตํ ยถาภูตํ สมฺมปฺปญฺญาย ทิสฺวา วาโยธาตุยา นิพฺพินฺทติ, วาโยธาตุยา จิตฺตํ วิราเชติ.}

\prose{ยโต โข ราหุล ภิกฺขุ อิมาสุ จตูสุ ธาตูสุ เนวตฺตานํ น อตฺตนิยํ สมนุปสฺสติ. อยํ วุจฺจติ ราหุล ภิกฺขุ อจฺเฉจฺฉิ ตณฺหํ, วิวตฺตยิ สํโยชนํ, สมฺมา มานาภิสมยา อนฺตมกาสิ ทุกฺขสฺสาติ.  }

\csromanmatikaanchor{mtk.215}
\csromantocmarkref{subsection}{8. ชมฺพาลีสุตฺต}{mtk.215}
\bookmark[dest={mtk.215},level=2]{8. ชมฺพาลีสุตฺต}
\csromanheader{2}{8. \textbf{ชมฺพาลีสุตฺต}}
\proseitem{178}{จตฺตาโรเม ภิกฺขเว ปุคฺคลา สนฺโต สํวิชฺชมานา โลกสฺมิํ. กตเม จตฺตาโร? อิธ ภิกฺขเว ภิกฺขุ อญฺญตรํ สนฺตํ เจโตวิมุตฺติํ อุปสมฺปชฺช วิหรติ, โส สกฺกาย\-นิโรธํ มนสิ กโรติ, ตสฺส สกฺกาย\-นิโรธํ มนสิ กโรโต สกฺกายนิโรเธ จิตฺตํ น ปกฺขนฺทติ นปฺปสีทติ น สนฺติฏฺฐติ นาธิมุจฺจติ. ตสฺส โข เอวํ\csromansharedfootnote{09b2e12a90aeb3db}{ตสฺส โข เอตํ (สี, สฺยา, กํ, อิ), เอวํ โข ตสฺส (?) “เอวํ หิ ตสฺสา ภิกฺขเว ชมฺพาลิยา”ติ ปาโฐ วิย.} ภิกฺขเว ภิกฺขุโน น สกฺกายนิโรโธ ปาฏิกงฺโข. เสยฺยถาปิ ภิกฺขเว ปุริโส เลปคเตน\csromansharedfootnote{831db0d61ba2b4c4}{ลสคเตน (สี, อิ)} หตฺเถน สาขํ คเณฺหยฺย, ตสฺส โส หตฺโถ สชฺเชยฺยปิ คเณฺหยฺยปิ\csromansharedfootnote{e4dd2a77318a0973}{คเยฺหยฺยปิ (?)} พชฺเฌยฺยปิ\csromansharedfootnote{f4e0dfdb2fc6f4af}{ขชฺเชยฺยปิ (สี)}. เอวเมวํ โข ภิกฺขเว ภิกฺขุ อญฺญตรํ สนฺตํ เจโตวิมุตฺติํ อุปสมฺปชฺช วิหรติ, โส สกฺกาย\-นิโรธํ มนสิ กโรติ. ตสฺส สกฺกาย\-นิโรธํ มนสิ กโรโต สกฺกายนิโรเธ จิตฺตํ น ปกฺขนฺทติ นปฺปสีทติ น สนฺติฏฺฐติ นาธิมุจฺจติ. ตสฺส โข เอวํ\csromansharedfootnote{09b2e12a90aeb3db}{ตสฺส โข เอตํ (สี, สฺยา, กํ, อิ), เอวํ โข ตสฺส (?) “เอวํ หิ ตสฺสา ภิกฺขเว ชมฺพาลิยา”ติ ปาโฐ วิย.} ภิกฺขเว ภิกฺขุโน น สกฺกายนิโรโธ ปาฏิกงฺโข.}

\chapterheadpageread{(18) 3. สญฺเจตนิย\-วคฺค}
\prose{\csromanfolio{485}อิธ ปน ภิกฺขเว ภิกฺขุ อญฺญตรํ สนฺตํ เจโตวิมุตฺติํ อุปสมฺปชฺช วิหรติ, โส สกฺกาย\-นิโรธํ มนสิ กโรติ, ตสฺส สกฺกาย\-นิโรธํ มนสิ กโรโต สกฺกายนิโรเธ จิตฺตํ ปกฺขนฺทติ ปสีทติ สนฺติฏฺฐติ อธิมุจฺจติ. ตสฺส โข เอวํ ภิกฺขเว ภิกฺขุโน สกฺกายนิโรโธ ปาฏิกงฺโข. เสยฺยถาปิ ภิกฺขเว ปุริโส สุทฺเธน หตฺเถน สาขํ คเณฺหยฺย, ตสฺส โส หตฺโถ เนว สชฺเชยฺย น คเณฺหยฺย น พชฺเฌยฺย. เอวเมวํ โข ภิกฺขเว ภิกฺขุ อญฺญตรํ สนฺตํ เจโตวิมุตฺติํ อุปสมฺปชฺช วิหรติ, โส สกฺกาย\-นิโรธํ มนสิ กโรติ, ตสฺส สกฺกาย\-นิโรธํ มนสิ กโรโต สกฺกายนิโรเธ จิตฺตํ ปกฺขนฺทติ ปสีทติ สนฺติฏฺฐติ อธิมุจฺจติ. ตสฺส โข เอวํ ภิกฺขเว ภิกฺขุโน สกฺกายนิโรโธ ปาฏิกงฺโข.}

\prose{อิธ ปน ภิกฺขเว ภิกฺขุ อญฺญตรํ สนฺตํ เจโตวิมุตฺติํ อุปสมฺปชฺช วิหรติ. โส อวิชฺชา\-ปฺปเภทํ มนสิ กโรติ, ตสฺส อวิชฺชา\-ปฺปเภทํ มนสิ กโรโต อวิชฺชาปฺปเภเท จิตฺตํ น ปกฺขนฺทติ นปฺปสีทติ น สนฺติฏฺฐติ นาธิมุจฺจติ. ตสฺส โข เอวํ ภิกฺขเว ภิกฺขุโน น อวิชฺชาปฺปเภโท ปาฏิกงฺโข. เสยฺยถาปิ ภิกฺขเว ชมฺพาลี อเนก\-วสฺส\-คณิกา, ตสฺสา ปุริโส ยานิ เจว อายมุขานิ, ตานิ ปิทเหยฺย. ยานิ จ อปายมุขานิ, ตานิ วิวเรยฺย. เทโว จ น สมฺมา ธารํ อนุปฺปเวจฺเฉยฺย. เอวํ หิ ตสฺสา ภิกฺขเว ชมฺพาลิยา น อาฬิปฺปเภโท ปาฏิกงฺโข. เอวเมวํ โข ภิกฺขเว ภิกฺขุ อญฺญตรํ สนฺตํ เจโตวิมุตฺติํ อุปสมฺปชฺช วิหรติ, โส อวิชฺชา\-ปฺปเภทํ มนสิ กโรติ, ตสฺส อวิชฺชา\-ปฺปเภทํ มนสิ กโรโต อวิชฺชาปฺปเภเท จิตฺตํ น ปกฺขนฺทติ นปฺปสีทติ น สนฺติฏฺฐติ นาธิมุจฺจติ. ตสฺส โข เอวํ ภิกฺขเว ภิกฺขุโน น อวิชฺชาปฺปเภโท ปาฏิกงฺโข.}

\prose{อิธ ปน ภิกฺขเว ภิกฺขุ อญฺญตรํ สนฺตํ เจโตวิมุตฺติํ อุปสมฺปชฺช วิหรติ, โส อวิชฺชา\-ปฺปเภทํ มนสิ กโรติ, ตสฺส อวิชฺชา\-ปฺปเภทํ มนสิ กโรโต อวิชฺชาปฺปเภเท จิตฺตํ ปกฺขนฺทติ ปสีทติ สนฺติฏฺฐติ อธิมุจฺจติ. ตสฺส โข เอวํ ภิกฺขเว ภิกฺขุโน อวิชฺชปฺปเภโท ปาฏิกงฺโข. เสยฺยถาปิ ภิกฺขเว ชมฺพาลี อเนก\-วสฺส\-คณิกา, ตสฺสา ปุริโส ยานิ เจว อายมุขานิ, ตานิ วิวเรยฺย. ยานิ จ อปายมุขานิ, ตานิ ปิทเหยฺย. เทโว จ สมฺมา ธารํ อนุปฺปเวจฺเฉยฺย. เอวํ หิ ตสฺสา ภิกฺขเว ชมฺพาลิยา อาฬิปฺปเภโท ปาฏิกงฺโข. เอวเมวํ โข ภิกฺขเว ภิกฺขุ}

\prose{\csromanfolio{486}อญฺญตรํ สนฺตํ เจโตวิมุตฺติํ อุปสมฺปชฺช วิหรติ, โส อวิชฺชา\-ปฺปเภทํ มนสิ กโรติ, ตสฺส อวิชฺชา\-ปฺปเภทํ มนสิ กโรโต อวิชฺชาปฺปเภเท จิตฺตํ ปกฺขนฺทติ ปสีทติ สนฺติฏฺฐติ อธิมุจฺจติ. ตสฺส โข เอวํ ภิกฺขเว ภิกฺขุโน อวิชฺชาปฺปเภโท ปาฏิกงฺโข. อิเม โข ภิกฺขเว จตฺตาโร ปุคฺคลา สนฺโต สํวิชฺชมานา โลกสฺมินฺติ.  อฏฺฐมํ.}

\csromanmatikaanchor{mtk.216}
\csromantocmarkref{subsection}{9. นิพฺพานสุตฺต}{mtk.216}
\bookmark[dest={mtk.216},level=2]{9. นิพฺพานสุตฺต}
\csromanheader{2}{9. \textbf{นิพฺพานสุตฺต}}
\proseitem{179}{อถ โข อายสฺมา อานนฺโท เยนายสฺมา สาริปุตฺโต เตนุปสงฺกมิ, อุปสงฺกมิตฺวา อายสฺมตา สาริปุตฺเตน สทฺธิํ สมฺโมทิ, สมฺโมทนียํ กถํ สารณียํ วีติสาเรตฺวา เอกมนฺตํ นิสีทิ, เอกมนฺตํ นิสินฺโน โข อายสฺมา อานนฺโท อายสฺมนฺตํ สาริปุตฺตํ เอตทโวจ “โก นุ โข อาวุโส สาริปุตฺต เหตุ โก ปจฺจโย, เยน มิเธกจฺเจ สตฺตา ทิฏฺเฐว ธมฺเม น ปรินิพฺพายนฺตี”ติ.}

\prose{อิธาวุโส อานนฺท สตฺตา “อิมา หานภาคิยา สญฺญา”ติ ยถาภูตํ นปฺปชานนฺติ, “อิมา ฐิติภาคิยา สญฺญา”ติ ยถาภูตํ นปฺปชานนฺติ, “อิมา วิเสสภาคิยา สญฺญา”ติ ยถาภูตํ นปฺปชานนฺติ, “อิมา นิพฺเพธ\-ภาคิยา สญฺญา”ติ ยถาภูตํ นปฺปชานนฺติ. อยํ โข อาวุโส อานนฺท เหตุ อยํ ปจฺจโย, เยน มิเธกจฺเจ สตฺตา ทิฏฺเฐว ธมฺเม น ปรินิพฺพายนฺตีติ.}

\prose{โก ปนาวุโส สาริปุตฺต เหตุ โก ปจฺจโย, เยน มิเธกจฺเจ สตฺตา ทิฏฺเฐว ธมฺเม ปรินิพฺพายนฺตีติ. อิธาวุโส อานนฺท สตฺตา “อิมา หานภาคิยา สญฺญา”ติ ยถาภูตํ ปชานนฺติ, “อิมา ฐิติภาคิยา สญฺญา”ติ ยถาภูตํ ปชานนฺติ, “อิมา วิเสสภาคิยา สญฺญา”ติ ยถาภูตํ ปชานนฺติ, “อิมา นิพฺเพธ\-ภาคิยา สญฺญา”ติ ยถาภูตํ ปชานนฺติ. อยํ โข อาวุโส อานนฺท เหตุ อยํ ปจฺจโย, เยน มิเธกจฺเจ สตฺตา ทิฏฺเฐว ธมฺเม ปรินิพฺพายนฺตีติ.  นวมํ.}

\csromanmatikaanchor{mtk.217}
\csromantocmarkref{subsection}{10. มหาปเทสสุตฺต}{mtk.217}
\bookmark[dest={mtk.217},level=2]{10. มหาปเทสสุตฺต}
\csromanheader{2}{10. มหาปเทส\-สุตฺต}
\proseitem{180}{เอกํ สมยํ ภควา โภคนคเร วิหรติ อานนฺทเจติเย\csromansharedfootnote{46edb0d4e78ea0f2}{อานนฺเท เจติเย (ที 2. 103 ปิฏฺเฐ) วิ 3. 348 ปิฏฺเฐ ปน อญฺญถา อาคตํ.}. ตตฺร โข ภควา ภิกฺขู อามนฺเตสิ “ภิกฺขโว”ติ. “ภทนฺเต”ติ เต}

\vspace{\tipitakaparskip}
\csromancenter{(18) 3. สญฺเจตนิย\-วคฺค ภิกฺขู ภควโต ปจฺจสฺโสสุํ. ภควา เอตทโวจ “จตฺตาโรเม ภิกฺขเว มหาปเทเส เทเสสฺสามิ, ตํ สุณาถ สาธุกํ มนสิ กโรถ, ภาสิสฺสามี”ติ. “เอวํ ภนฺเต”ติ โข เต ภิกฺขู ภควโต ปจฺจสฺโสสุํ. ภควา เอตทโวจ\csromandash{}}
\prose{\csromanfolio{487}กตเม ภิกฺขเว จตฺตาโร มหาปเทสา? อิธ ภิกฺขเว ภิกฺขุ เอวํ วเทยฺย “สมฺมุขา เมตํ อาวุโส ภควโต สุตํ, สมฺมุขา ปฏิคฺคหิตํ ‘อยํ ธมฺโม อยํ วินโย อิทํ สตฺถุสาสนนฺ’ติ”. ตสฺส ภิกฺขเว ภิกฺขุโน ภาสิตํ เนว อภินนฺทิตพฺพํ นปฺปฏิกฺโกสิตพฺพํ, อนภินนฺทิตฺวา อปฺปฏิกฺโกสิตฺวา ตานิ ปทพฺยญฺชนานิ สาธุกํ อุคฺคเหตฺวา สุตฺเต โอตาเรตพฺพานิ\csromansharedfootnote{4ad61fe77e912b03}{โอสาเรตพฺพานิ} วินเย สนฺทสฺเส\-ตพฺพานิ. ตานิ เจ สุตฺเต โอตาริยมานานิ\csromansharedfootnote{beeed3a46dfc10c3}{โอสาริยมานานิ} วินเย สนฺทสฺสิย\-มานานิ น เจว สุตฺเต โอตรนฺติ\csromansharedfootnote{ab30634926685478}{โอสรนฺติ (ที 2. 103)} น วินเย สนฺทิสฺสนฺติ, นิฏฺฐเมตฺถ คนฺตพฺพํ “อทฺธา อิทํ น เจว ตสฺส ภควโต วจนํ อรหโต สมฺมา\-สมฺพุทฺธสฺส, อิมสฺส จ ภิกฺขุโน ทุคฺคหิตนฺ”ติ. อิติ เหตํ\csromansharedfootnote{a0ca67b499d8e6fc}{อิติ หิทํ (สี, สฺยา, กํ, ก) 5-5. เอตฺตโก ปาโฐ ทีฆนิกาเย น ทิสฺสติ, เปยฺยาลมุเขน ทสฺสิโต ภเวยฺย.} ภิกฺขเว ฉฑฺเฑยฺยาถ.}

\prose{5อิธ ปน ภิกฺขเว ภิกฺขุ เอวํ วเทยฺย “สมฺมุขา เมตํ อาวุโส ภควโต สุตํ, สมฺมุขา ปฏิคฺคหิตํ ‘อยํ ธมฺโม อยํ วินโย อิทํ สตฺถุสาสนนฺ’ติ”. ตสฺส ภิกฺขเว ภิกฺขุโน ภาสิตํ เนว อภินนฺทิตพฺพํ นปฺปฏิกฺโกสิตพฺพํ, อนภินนฺทิตฺวา อปฺปฏิกฺโกสิตฺวา ตานิ ปทพฺยญฺชนานิ สาธุกํ อุคฺคเหตฺวา สุตฺเต โอตาเรตพฺพานิ วินเย สนฺทสฺเส\-ตพฺพานิ5. ตานิ เจ สุตฺเต โอตาริยมานานิ วินเย สนฺทสฺสิย\-มานานิ สุตฺเต เจว โอตรนฺติ วินเย จ สนฺทิสฺสนฺติ, นิฏฺฐเมตฺถ คนฺตพฺพํ “อทฺธา อิทํ ตสฺส ภควโต วจนํ อรหโต สมฺมา\-สมฺพุทฺธสฺส, อิมสฺส จ ภิกฺขุโน สุคฺคหิตนฺ”ติ. อิทํ ภิกฺขเว ปฐมํ มหาปเทสํ ธาเรยฺยาถ.}

\prose{อิธ ปน ภิกฺขเว ภิกฺขุ เอวํ วเทยฺย “อสุกสฺมิํ นาม อาวาเส สํโฆ วิหรติ สเถโร สปาโมกฺโข, ตสฺส เม สํฆสฺส สมฺมุขา สุตํ, สมฺมุขา ปฏิคฺคหิตํ ‘อยํ ธมฺโม อยํ วินโย อิทํ สตฺถุสาสนนฺ’ติ”. ตสฺส ภิกฺขเว ภิกฺขุโน ภาสิตํ เนว อภินนฺทิตพฺพํ นปฺปฏิกฺโกสิตพฺพํ, \csromanfolio{488}อนภินนฺทิตฺวา อปฺปฏิกฺโกสิตฺวา ตานิ ปทพฺยญฺชนานิ สาธุกํ อุคฺคเหตฺวา สุตฺเต โอตาเรตพฺพานิ วินเย สนฺทสฺเส\-ตพฺพานิ. ตานิ เจ สุตฺเต โอตาริยมานานิ วินเย สนฺทสฺสิย\-มานานิ น เจว สุตฺเต โอตรนฺติ น วินเย สนฺทิสฺสนฺติ, นิฏฺฐเมตฺถ คนฺตพฺพํ “อทฺธา อิทํ น เจว ตสฺส ภควโต วจนํ อรหโต สมฺมา\-สมฺพุทฺธสฺส, ตสฺส จ สํฆสฺส ทุคฺคหิตนฺ”ติ. อิติ เหตํ ภิกฺขเว ฉฑฺเฑยฺยาถ.}

\prose{อิธ ปน ภิกฺขเว ภิกฺขุ เอวํ วเทยฺย “อสุกสฺมิํ นาม อาวาเส สํโฆ วิหรติ สเถโร สปาโมกฺโข, ตสฺส เม สํฆสฺส สมฺมุขา สุตํ, สมฺมุขา ปฏิคฺคหิตํ ‘อยํ ธมฺโม อยํ วินโย อิทํ สตฺถุสาสนนฺ’ติ”. ตสฺส ภิกฺขเว ภิกฺขุโน ภาสิตํ เนว อภินนฺทิตพฺพํ นปฺปฏิกฺโกสิตพฺพํ. อนภินนฺทิตฺวา อปฺปฏิกฺโกสิตฺวา ตานิ ปทพฺยญฺชนานิ สาธุกํ อุคฺคเหตฺวา สุตฺเต โอตาเรตพฺพานิ วินเย สนฺทสฺเส\-ตพฺพานิ. ตานิ เจ สุตฺเต โอตาริยมานานิ วินเย สนฺทสฺสิย\-มานานิ สุตฺเต เจว โอตรนฺติ วินเย จ สนฺทิสฺสนฺติ, นิฏฺฐเมตฺถ คนฺตพฺพํ “อทฺธา อิทํ ตสฺส ภควโต วจนํ อรหโต สมฺมา\-สมฺพุทฺธสฺส, ตสฺส จ สํฆสฺส สุคฺคหิตนฺ”ติ. อิทํ ภิกฺขเว ทุติยํ มหาปเทสํ ธาเรยฺยาถ.}

\prose{อิธ ปน ภิกฺขเว ภิกฺขุ เอวํ วเทยฺย “อสุกสฺมิํ นาม อาวาเส สมฺพหุลา เถรา ภิกฺขู วิหรนฺติ พหุสฺสุตา อาคตาคมา ธมฺมธรา วินยธรา มาติกาธรา, เตสํ เม เถรานํ สมฺมุขา สุตํ, สมฺมุขา ปฏิคฺคหิตํ ‘อยํ ธมฺโม อยํ วินโย อิทํ สตฺถุสาสนนฺ’ติ”. ตสฺส ภิกฺขเว ภิกฺขุโน ภาสิตํ เนว อภินนฺทิตพฺพํ นปฺปฏิกฺโกสิตพฺพํ, อนภินนฺทิตฺวา อปฺปฏิกฺโกสิตฺวา ตานิ ปทพฺยญฺชนานิ สาธุกํ อุคฺคเหตฺวา สุตฺเต โอตาเรตพฺพานิ วินเย สนฺทสฺเส\-ตพฺพานิ. ตานิ เจ สุตฺเต โอตาริยมานานิ วินเย สนฺทสฺสิย\-มานานิ น เจว สุตฺเต โอตรนฺติ น วินเย สนฺทิสฺสนฺติ, นิฏฺฐเมตฺถ คนฺตพฺพํ “อทฺธา อิทํ น เจว ตสฺส ภควโต วจนํ อรหโต สมฺมา\-สมฺพุทฺธสฺส, เตสญฺจ เถรานํ ทุคฺคหิตนฺ”ติ. อิติ เหตํ ภิกฺขเว ฉฑฺเฑยฺยาถ.}

\prose{อิธ ปน ภิกฺขเว ภิกฺขุ เอวํ วเทยฺย “อสุกสฺมิํ นาม อาวาเส สมฺพหุลา เถรา ภิกฺขู วิหรนฺติ พหุสฺสุตา อาคตาคมา ธมฺมธรา วินยธรา มาติกาธรา, เตสํ เม เถรานํ สมฺมุขา สุตํ, สมฺมุขา ปฏิคฺคหิตํ ‘อยํ ธมฺโม อยํ วินโย อิทํ สตฺถุสาสนนฺ’ติ”. ตสฺส}

\vspace{\tipitakaparskip}
\csromancenter{(18) 3. สญฺเจตนิย\-วคฺค ภิกฺขเว ภิกฺขุโน ภาสิตํ เนว อภินนฺทิตพฺพํ นปฺปฏิกฺโกสิตพฺพํ, อนภินนฺทิตฺวา อปฺปฏิกฺโกสิตฺวา ตานิ ปทพฺยญฺชนานิ สาธุกํ อุคฺคเหตฺวา สุตฺเต โอตาเรตพฺพานิ วินเย สนฺทสฺเส\-ตพฺพานิ. ตานิ เจ สุตฺเต โอตาริยมานานิ วินเย สนฺทสฺสิย\-มานานิ สุตฺเต เจว โอตรนฺติ วินเย จ สนฺทิสฺสนฺติ, นิฏฺฐเมตฺถ คนฺตพฺพํ “อทฺธา อิทํ ตสฺส ภควโต วจนํ อรหโต สมฺมา\-สมฺพุทฺธสฺส, เตสญฺจ เถรานํ สุคฺคหิตนฺ”ติ. อิทํ ภิกฺขเว ตติยํ มหาปเทสํ ธาเรยฺยาถ.}
\prose{\csromanfolio{489}อิธ ปน ภิกฺขเว ภิกฺขุ เอวํ วเทยฺย “อสุกสฺมิํ นาม อาวาเส เอโก เถโร ภิกฺขุ วิหรติ พหุสฺสุโต อาคตาคโม ธมฺมธโร วินยธโร มาติกาธโร, ตสฺส เม เถรสฺส สมฺมุขา สุตํ, สมฺมุขา ปฏิคฺคหิตํ ‘อยํ ธมฺโม อยํ วินโย อิทํ สตฺถุสาสนนฺ’ติ”. ตสฺส ภิกฺขเว ภิกฺขุโน ภาสิตํ เนว อภินนฺทิตพฺพํ นปฺปฏิกฺโกสิตพฺพํ, อนภินนฺทิตฺวา อปฺปฏิกฺโกสิตฺวา ตานิ ปทพฺยญฺชนานิ สาธุกํ อุคฺคเหตฺวา สุตฺเต โอตาเรตพฺพานิ วินเย สนฺทสฺเส\-ตพฺพานิ. ตานิ เจ สุตฺเต โอตาริยมานานิ วินเย สนฺทสฺสิย\-มานานิ น เจว สุตฺเต โอตรนฺติ น วินเย สนฺทิสฺสนฺติ, นิฏฺฐเมตฺถ คนฺตพฺพํ “อทฺธา อิทํ น เจว ตสฺส ภควโต วจนํ อรหโต สมฺมา\-สมฺพุทฺธสฺส, ตสฺส จ เถรสฺส ทุคฺคหิตนฺ”ติ. อิติ เหตํ ภิกฺขเว ฉฑฺเฑยฺยาถ.}

\prosewithclosermidruled{อิธ ปน ภิกฺขเว ภิกฺขุ เอวํ วเทยฺย “อสุกสฺมิํ นาม อาวาเส เอโก เถโร ภิกฺขุ วิหรติ พหุสฺสุโต อาคตาคโม ธมฺมธโร วินยธโร มาติกาธโร, ตสฺส เม เถรสฺส สมฺมุขา สุตํ, สมฺมุขา ปฏิคฺคหิตํ ‘อยํ ธมฺโม อยํ วินโย อิทํ สตฺถุสาสนนฺ’ติ”. ตสฺส ภิกฺขเว ภิกฺขุโน ภาสิตํ เนว อภินนฺทิตพฺพํ นปฺปฏิกฺโกสิตพฺพํ, อนภินนฺทิตฺวา อปฺปฏิกฺโกสิตฺวา ตานิ ปทพฺยญฺชนานิ สาธุกํ อุคฺคเหตฺวา สุตฺเต โอตาเรตพฺพานิ วินเย สนฺทสฺเส\-ตพฺพานิ. ตานิ เจ สุตฺเต โอตาริยมานานิ วินเย สนฺทสฺสิย\-มานานิ สุตฺเต เจว โอตรนฺติ วินเย จ สนฺทิสฺสนฺติ, นิฏฺฐเมตฺถ คนฺตพฺพํ “อทฺธา อิทํ ตสฺส ภควโต วจนํ อรหโต สมฺมา\-สมฺพุทฺธสฺส, ตสฺส จ เถรสฺส สุคฺคหิตนฺ”ติ. อิทํ ภิกฺขเว จตุตฺถํ มหาปเทสํ ธาเรยฺยาถ. อิเม โข ภิกฺขเว จตฺตาโร มหาปเทสาติ.  ทสมํ.}{สญฺเจตนิย\-วคฺโค ตติโย.}
\csromanpalirecto
\gambhira{จตุกฺกนิปาต\-ปาฬิ ตสฺสุทฺทานํ}
\setlength{\csromangathapagewidth}{0pt}
\setlength{\csromangathawakwidth}{0pt}
\csromangathameasuregroup{เจตนา วิภตฺติ โกฏฺฐิโก, \\ อายาจน ราหุล ชมฺพาลี,}{\csromangathabat{เจตนา วิภตฺติ โกฏฺฐิโก,}{อานนฺโท อุปวาณปญฺจมํ.} \\ \csromangathabat{อายาจน ราหุล ชมฺพาลี,}{นิพฺพานํ มหาปเทเสนาติ.}}
\csromangathasetleft{เจตนา วิภตฺติ โกฏฺฐิโก, \\ อายาจน ราหุล ชมฺพาลี,}
\csromangathagroup{\csromangathastanza{\csromanfolio{490}\csromangathabat{เจตนา วิภตฺติ โกฏฺฐิโก,}{อานนฺโท อุปวาณปญฺจมํ.} \\ \csromangathabat{อายาจน ราหุล ชมฺพาลี,}{นิพฺพานํ มหาปเทเสนาติ.}}}

\csromanmatikaanchor{mtk.218}
\csromantocmarknopage{section}{(19) 4. พฺราหฺมณวคฺค}{mtk.218}
\bookmark[dest={mtk.218},level=1]{(19) 4. พฺราหฺมณวคฺค}
\csromanmarkh{(19)}\csromanmarkhii{4. พฺราหฺมณวคฺค}\csromanheadernomark{1}{\textbf{(19) 4. พฺราหฺมณวคฺค}}
\csromanmatikaanchor{mtk.219}
\csromantocmarkref{subsection}{1. โยธาชีวสุตฺต}{mtk.219}
\bookmark[dest={mtk.219},level=2]{1. โยธาชีวสุตฺต}
\csromanheader{2}{1. \textbf{โยธาชีวสุตฺต}}
\proseitem{181}{จตูหิ ภิกฺขเว องฺเคหิ สมนฺนาคโต โยธาชีโว ราชารโห โหติ ราชโภคฺโค, รญฺโญ องฺคนฺเตว สงฺขํ คจฺฉติ. กตเมหิ จตูหิ? อิธ ภิกฺขเว โยธาชีโว ฐานกุสโล จ โหติ, ทูเรปาตี จ, อกฺขณเวธี จ, มหโต จ กายสฺส ปทาเลตา. อิเมหิ โข ภิกฺขเว จตูหิ องฺเคหิ สมนฺนาคโต โยธาชีโว ราชารโห โหติ ราชโภคฺโค, รญฺโญ องฺคนฺเตว สงฺขํ คจฺฉติ. เอวเมวํ โข ภิกฺขเว จตูหิ ธมฺเมหิ สมนฺนาคโต ภิกฺขุ อาหุเนยฺโย โหติ ปาหุเนยฺโย ทกฺขิเณยฺโย อญฺชลิกรณีโย อนุตฺตรํ ปุญฺญกฺเขตฺตํ โลกสฺส. กตเมหิ จตูหิ? อิธ ภิกฺขเว ภิกฺขุ ฐานกุสโล จ โหติ, ทูเรปาตี จ, อกฺขณเวธี จ, มหโต จ กายสฺส ปทาเลตา.}

\prose{กถญฺจ ภิกฺขเว ภิกฺขุ ฐานกุสโล โหติ? อิธ ภิกฺขเว ภิกฺขุ สีลวา โหติ ฯเปฯ สมาทาย สิกฺขติ สิกฺขาปเทสุ. เอวํ โข ภิกฺขเว ภิกฺขุ ฐานกุสโล โหติ.}

\prose{กถญฺจ ภิกฺขเว ภิกฺขุ ทูเรปาตี โหติ? อิธ ภิกฺขเว ภิกฺขุ ยํ กิญฺจิ รูปํ อตีตา\-นาคต\-ปจฺจุปฺปนฺนํ อชฺฌตฺตํ วา พหิทฺธา วา โอฬาริกํ วา สุขุมํ วา หีนํ วา ปณีตํ วา ยํ ทูเร สนฺติเก วา, สพฺพํ รูปํ “เนตํ มม, เนโสหมสฺมิ, น เมโส อตฺตา”ติ เอวเมตํ ยถาภูตํ สมฺมปฺปญฺญาย ปสฺสติ. ยา กาจิ เวทนา. ยา กาจิ สญฺญา. เย เกจิ สงฺขารา. ยํ กิญฺจิ วิญฺญาณํ อตีตา\-นาคต\-ปจฺจุปฺปนฺนํ อชฺฌตฺตํ วา พหิทฺธา วา โอฬาริกํ วา สุขุมํ วา หีนํ วา ปณีตํ วา ยํ ทูเร สนฺติเก วา, สพฺพํ วิญฺญาณํ “เนตํ มม, เนโสหมสฺมิ, น เมโส อตฺตา”ติ เอวเมตํ ยถาภูตํ สมฺมปฺปญฺญาย ปสฺสติ. เอวํ โข ภิกฺขเว ภิกฺขุ ทูเรปาตี โหติ.}

\chapterheadpageread{(19) 4. พฺราหฺมณวคฺค}
\prose{\csromanfolio{491}กถญฺจ ภิกฺขเว ภิกฺขุ อกฺขณเวธี โหติ? อิธ ภิกฺขเว ภิกฺขุ “อิทํ ทุกฺขนฺ”ติ ยถาภูตํ ปชานาติ ฯเปฯ “อยํ ทุกฺขนิโรธ\-คามินี ปฏิปทา”ติ ยถาภูตํ ปชานาติ. เอวํ โข ภิกฺขเว ภิกฺขุ อกฺขณเวธี โหติ.}

\prose{กถญฺจ ภิกฺขเว ภิกฺขุ มหโต กายสฺส ปทาเลตา โหติ? อิธ ภิกฺขเว ภิกฺขุ มหนฺตํ อวิชฺชา\-กฺขนฺธํ ปทาเลตา. เอวํ โข ภิกฺขเว ภิกฺขุ มหโต กายสฺส ปทาเลตา โหติ. อิเมหิ โข ภิกฺขเว จตูหิ ธมฺเมหิ สมนฺนาคโต ภิกฺขุ อาหุเนยฺโย โหติ ฯเปฯ อนุตฺตรํ ปุญฺญกฺเขตฺตํ โลกสฺสาติ\csromansharedfootnote{568da7c755d763e4}{เหฏฺฐา 288 ปิฏฺเฐปิ.}.  ปฐมํ.}

\csromanmatikaanchor{mtk.220}
\csromantocmarkref{subsection}{2. ปาฏิโภคสุตฺต}{mtk.220}
\bookmark[dest={mtk.220},level=2]{2. ปาฏิโภคสุตฺต}
\csromanheader{2}{2. \textbf{ปาฏิโภคสุตฺต}}
\proseitem{182}{\csromansymbolfootnote{*}{อภิ 4. 332 ปิฏฺเฐปิ.}จตุนฺนํ ภิกฺขเว ธมฺมานํ นตฺถิ โกจิ ปาฏิโภโค สมโณ วา พฺราหฺมโณ วา เทโว วา มาโร วา พฺรหฺมา วา โกจิ วา โลกสฺมิํ.}

\prose{กตเมสํ จตุนฺนํ? ชราธมฺมํ “มา ชีรี”ติ นตฺถิ โกจิ ปาฏิโภโค สมโณ วา พฺราหฺมโณ วา เทโว วา มาโร วา พฺรหฺมา วา โกจิ วา โลกสฺมิํ. พฺยาธิธมฺมํ “มา พฺยาธิยี”ติ นตฺถิ โกจิ ปาฏิโภโค สมโณ วา พฺราหฺมโณ วา เทโว วา มาโร วา พฺราหฺมา วา โกจิ วา โลกสฺมิํ. มรณธมฺมํ “มา มียี”ติ นตฺถิ โกจิ ปาฏิโภโค สมโณ วา พฺราหฺมโณ วา เทโว วา มาโร วา พฺรหฺมา วา โกจิ วา โลกสฺมิํ. ยานิ โข ปน ตานิ ปุพฺเพ อตฺตนา กตานิ ปาปกานิ กมฺมานิ สํกิเลสิกานิ โปโนภวิกานิ สทรานิ ทุกฺขวิปากานิ อายติํ ชาติ\-ชรา\-มรณิ\-กานิ, เตสํ วิปาโก “มา นิพฺพตฺตี”ติ นตฺถิ โกจิ ปาฏิโภโค สมโณ วา พฺราหฺมโณ วา เทโว วา มาโร วา พฺรหฺมา วา โกจิ วา โลกสฺมิํ.}

\prose{อิเมสํ โข ภิกฺขเว จตุนฺนํ ธมฺมานํ นตฺถิ โกจิ ปาฏิโภโค สมโณ วา พฺราหฺมโณ วา เทโว วา มาโร วา พฺรหฺมา วา โกจิ วา โลกสฺมินฺติ.  ทุติยํ.}

\csromanmatikaanchor{mtk.221}
\csromantocmarkref{subsection}{3. สุตสุตฺต}{mtk.221}
\bookmark[dest={mtk.221},level=2]{3. สุตสุตฺต}
\csromanheader{2}{3. \textbf{สุตสุตฺต}}
\proseitem{183}{เอกํ สมยํ ภควา ราชคเห วิหรติ เวฬุวเน กลนฺทก\-นิวาเป. อถ โข วสฺสกาโร พฺราหฺมโณ มคธ\-มหามตฺโต เยน \csromanfolio{492}ภควา เตนุปสงฺกมิ, อุปสงฺกมิตฺวา ภควตา สทฺธิํ สมฺโมทิ, สมฺโมทนียํ กถํ สารณียํ วีติสาเรตฺวา เอกมนฺตํ นิสีทิ, เอกมนฺตํ นิสินฺโน โข วสฺสกาโร พฺราหฺมโณ มคธ\-มหามตฺโต ภควนฺตํ เอตทโวจ\csromandash{}}

\prose{อหํ หิ โภ โคตม เอวํวาที เอวํทิฏฺฐิ\csromansharedfootnote{7b9448bb472d99e7}{เอวํทิฏฺฐี (สพฺพตฺถ)} “โย โกจิ ทิฏฺฐํ ภาสติ ‘เอวํ เม ทิฏฺฐนฺ’ติ, นตฺถิ ตโต โทโส. โย โกจิ สุตํ ภาสติ ‘เอวํ เม สุตนฺ’ติ, นตฺถิ ตโต โทโส. โย โกจิ มุตํ ภาสติ ‘เอวํ เม มุตนฺ’ติ, นตฺถิ ตโต โทโส. โย โกจิ วิญฺญาตํ ภาสติ ‘เอวํ เม วิญฺญาตนฺ’ติ, นตฺถิ ตโต โทโส”ติ.}

\prose{นาหํ พฺราหฺมณ สพฺพํ ทิฏฺฐํ “ภาสิตพฺพนฺ”ติ วทามิ, น ปนาหํ พฺราหฺมณ สพฺพํ ทิฏฺฐํ “น ภาสิตพฺพนฺ”ติ วทามิ, นาหํ พฺราหฺมณ สพฺพํ สุตํ “ภาสิตพฺพนฺ”ติ วทามิ, น ปนาหํ พฺราหฺมณ สพฺพํ สุตํ “น ภาสิตพฺพนฺ”ติ วทามิ, นาหํ พฺราหฺมณ สพฺพํ มุตํ “ภาสิตพฺพนฺ”ติ วทามิ, น ปนาหํ พฺราหฺมณ สพฺพํ มุตํ “น ภาสิตพฺพนฺ”ติ วทามิ, นาหํ พฺราหฺมณ สพฺพํ วิญฺญาตํ “ภาสิตพฺพนฺ”ติ วทามิ, น ปนาหํ พฺราหฺมณ สพฺพํ วิญฺญาตํ “น ภาสิตพฺพนฺ”ติ วทามิ.}

\prose{ยํ หิ พฺราหฺมณ ทิฏฺฐํ ภาสโต อกุสลา ธมฺมา อภิวฑฺฒนฺติ, กุสลา ธมฺมา ปริหายนฺติ, เอวรูปํ ทิฏฺฐํ “น ภาสิตพฺพนฺ”ติ วทามิ. ยญฺจ ขฺวสฺส พฺราหฺมณ ทิฏฺฐํ อภาสโต กุสลา ธมฺมา ปริหายนฺติ, อกุสลา ธมฺมา อภิวฑฺฒนฺติ, เอวรูปํ ทิฏฺฐํ “ภาสิตพฺพนฺ”ติ วทามิ.}

\prose{ยํ หิ พฺราหฺมณ สุตํ ภาสโต อกุสลา ธมฺมา อภิวฑฺฒนฺติ, กุสลา ธมฺมา ปริหายนฺติ, เอวรูปํ สุตํ “น ภาสิตพฺพนฺ”ติ วทามิ. ยญฺจ ขฺวสฺส พฺราหฺมณ สุตํ อภาสโต กุสลา ธมฺมา ปริหายนฺติ, อกุสลา ธมฺมา อภิวฑฺฒนฺติ, เอวรูปํ สุตํ “ภาสิตพฺพนฺ”ติ วทามิ.}

\prose{ยํ หิ พฺราหฺมณ มุตํ ภาสโต อกุสลา ธมฺมา อภิวฑฺฒนฺติ, กุสลา ธมฺมา ปริหายนฺติ, เอวรูปํ มุตํ “น ภาสิตพฺพนฺ”ติ วทามิ, ยญฺจ ขฺวสฺส พฺราหฺมณ มุตํ อภาสโต กุสลา ธมฺมา ปริหายนฺติ, อกุสลา ธมฺมา อภิวฑฺฒนฺติ, เอวรูปํ มุตํ “ภาสิตพฺพนฺ”ติ วทามิ.}

\chapterheadpageread{(19) 4. พฺราหฺมณวคฺค}
\prose{\csromanfolio{493}ยํ หิ พฺราหฺมณ วิญฺญาตํ ภาสโต อกุสลา ธมฺมา อภิวฑฺฒนฺติ, กุสลา ธมฺมา ปริหายนฺติ, เอวรูปํ วิญฺญาตํ “น ภาสิตพฺพนฺ”ติ วทามิ, ยญฺจ ขฺวสฺส พฺราหฺมณ วิญฺญาตํ อภาสโต กุสลา ธมฺมา ปริหายนฺติ, อกุสลา ธมฺมา อภิวฑฺฒนฺติ เอวรูปํ วิญฺญาตํ “ภาสิตพฺพนฺ”ติ วทามีติ.}

\prose{อถ โข วสฺสกาโร พฺราหฺมโณ มคธ\-มหามตฺโต ภควโต ภาสิตํ อภินนฺทิตฺวา อนุโมทิตฺวา อุฏฺฐายาสนา ปกฺกามีติ.  ตติยํ.}

\csromanmatikaanchor{mtk.222}
\csromantocmarkref{subsection}{4. อภยสุตฺต}{mtk.222}
\bookmark[dest={mtk.222},level=2]{4. อภยสุตฺต}
\csromanheader{2}{4. \textbf{อภยสุตฺต}}
\proseitem{184}{อถ โข ชาณุสฺโสณิ พฺราหฺมโณ เยน ภควา เตนุปสงฺกมิ, อุปสงฺกมิตฺวา ภควตา สทฺธิํ สมฺโมทิ, สมฺโมทนียํ กถํ สารณียํ วีติสาเรตฺวา เอกมนฺตํ นิสีทิ, เอกมนฺตํ นิสีนฺโน โข ชาณุสฺโสณิ พฺราหฺมโณ ภควนฺตํ เอตทโวจ\csromandash{}}

\prose{อหํ หิ โภ โคตม เอวํวาที เอวํทิฏฺฐิ “นตฺถิ โย มรณธมฺโม สมาโน น ภายติ, น สนฺตาสํ อาปชฺชติ มรณสฺสา”ติ. อตฺถิ พฺราหฺมณ มรณธมฺโม สมาโน ภายติ, สนฺตาสํ อาปชฺชติ มรณสฺส. อตฺถิ ปน พฺราหฺมณ มรณธมฺโม สมาโน น ภายติ, น สนฺตาสํ อาปชฺชติ มรณสฺส.}

\prose{กตโม จ พฺราหฺมณ มรณธมฺโม สมาโน ภายติ, สนฺตาสํ อาปชฺชติ มรณสฺส? อิธ พฺราหฺมณ เอกจฺโจ กาเมสุ อวีตราโค โหติ อวิคตจฺฉนฺโท อวิคตเปโม อวิคตปิปาโส อวิคตปริฬาโห อวิคตตโณฺห, ตเมนํ อญฺญตโร คาโฬฺห โรคาตงฺโก ผุสติ, ตสฺส อญฺญตเรน คาเฬฺหน โรคาตงฺเกน ผุฏฺฐสฺส เอวํ โหติ “ปิยา วต มํ กามา ชหิสฺสนฺติ, ปิเย จาหํ กาเม ชหิสฺสามี”ติ. โส โสจติ กิลมติ ปริเทวติ อุรตฺตาฬิํ กนฺทติ สมฺโมหํ อาปชฺชติ. อยํ โข พฺราหฺมณ มรณธมฺโม สมาโน ภายติ, สนฺตาสํ อาปชฺชติ มรณสฺส.}

\prose{ปุน จปรํ พฺราหฺมณ อิเธกจฺโจ กาเย อวีตราโค โหติ อวิคตจฺฉนฺโท อวิคตเปโม อวิคตปิปาโส อวิคตปริฬาโห อวิคตตโณฺห, ตเมนํ อญฺญตโร คาโฬฺห โรคาตงฺโก ผุสติ, ตสฺส อญฺญตเรน คาเฬฺหน โรคาตงฺเกน ผุฏฺฐสฺส เอวํ โหติ “ปิโย \csromanfolio{494}วต มํ กาโย ชหิสฺสติ, ปิยญฺจาหํ กายํ ชหิสฺสามี”ติ. โส โสจติ กิลมติ ปริเทวติ อุรตฺตาฬิํ กนฺทติ สมฺโมหํ อาปชฺชติ. อยมฺปิ โข พฺราหฺมณ มรณธมฺโม สมาโน ภายติ, สนฺตาสํ อาปชฺชติ มรณสฺส.}

\prose{ปุน จปรํ พฺราหฺมณ อิเธกจฺโจ อกตกลฺยาโณ โหติ อกตกุสโล อกต\-ภีรุตฺตาโณ, กตปาโป กตลุทฺโท กตกิพฺพิโส. ตเมนํ อญฺญตโร คาโฬฺห โรคาตงฺโก ผุสติ, ตสฺส อญฺญตเรน คาเฬฺหน โรคาตงฺเกน ผุฏฺฐสฺส เอวํ โหติ “อกตํ วต เม กลฺยาณํ, อกตํ กุสลํ, อกตํ ภีรุตฺตาณํ, กตํ ปาปํ, กตํ ลุทฺทํ, กตํ กิพฺพิสํ. ยาวตา โภ อกต\-กลฺยาณานํ อกตกุสลานํ อกต\-ภีรุตฺตาณานํ กตปาปานํ กตลุทฺทานํ กตกิพฺพิสานํ คติ, ตํ คติํ เปจฺจ คจฺฉามี”ติ. โส โสจติ กิลมติ ปริเทวติ อุรตฺตาฬิํ กนฺทติ สมฺโมหํ อาปชฺชติ. อยมฺปิ โข พฺราหฺมณ มรณธมฺโม สมาโน ภายติ, สนฺตาสํ อาปชฺชติ มรณสฺส.}

\prose{ปุน จปรํ พฺราหฺมณ อิเธกจฺโจ กงฺขี โหติ วิจิกิจฺฉี อนิฏฺฐงฺคโต สทฺธมฺเม, ตเมนํ อญฺญตโร คาโฬฺห โรคาตงฺโก ผุสติ, ตสฺส อญฺญตเรน คาเฬฺหน โรคาตงฺเกน ผุฏฺฐสฺส เอวํ โหติ “กงฺขี วตมฺหิ วิจิกิจฺฉี อนิฏฺฐงฺคโต สทฺธมฺเม”ติ. โส โสจติ กิลมติ ปริเทวติ อุรตฺตาฬิํ กนฺทติ สมฺโมหํ อาปชฺชติ. อยมฺปิ โข พฺราหฺมณ มรณธมฺโม สมาโน ภายติ, สนฺตาสํ อาปชฺชติ มรณสฺส. อิเม โข พฺราหฺมณ จตฺตาโร มรณธมฺมา สมานา ภายนฺติ, สนฺตาสํ อาปชฺชนฺติ มรณสฺส.}

\prose{กตโม จ พฺราหฺมณ มรณธมฺโม สมาโน น ภายติ, น สนฺตาสํ อาปชฺชติ มรณสฺส? อิธ พฺราหฺมณ เอกจฺโจ กาเมสุ วีตราโค โหติ วิคตจฺฉนฺโท วิคตเปโม วิคตปิปาโส วิคตปริฬาโห วิคตตโณฺห, ตเมนํ อญฺญตโร คาโฬฺห โรคาตงฺโก ผุสติ, ตสฺส อญฺญตเรน คาเฬฺหน โรคาตงฺเกน ผุฏฺฐสฺส น เอวํ โหติ “ปิยา วต มํ กามา ชหิสฺสนฺติ, ปิเย จาหํ กาเม ชหิสฺสามี”ติ. โส น โสจติ น กิลมติ น ปริเทวติ น อุรตฺตาฬิํ กนฺทติ น สมฺโมหํ อาปชฺชติ. อยํ โข พฺราหฺมณ มรณธมฺโม สมาโน น ภายติ, น สนฺตาสํ อาปชฺชติ มรณสฺส.}

\chapterheadpageread{(19) 4. พฺราหฺมณวคฺค}
\prose{\csromanfolio{495}ปุน จปรํ พฺราหฺมณ อิเธกจฺโจ กาเย วีตราโค โหติ วิคตจฺฉนฺโท วิคตเปโม วิคตปิปาโส วิคตปริฬาโห วิคตตโณฺห, ตเมนํ อญฺญตโร คาโฬฺห โรคาตงฺโก ผุสติ, ตสฺส อญฺญตเรน คาเฬฺหน โรคาตงฺเกน ผุฏฺฐสฺส น เอวํ โหติ “ปิโย วต มํ กาโย ชหิสฺสติ, ปิยญฺจาหํ กายํ ชหิสฺสามี”ติ. โส น โสจติ น กิลมติ น ปริเทวติ น อุรตฺตาฬิํ กนฺทติ น สมฺโมหํ อาปชฺชติ. อยมฺปิ โข พฺราหฺมณ มรณธมฺโม สมาโน น ภายติ, น สนฺตาสํ อาปชฺชติ มรณสฺส.}

\prose{ปุน จปรํ พฺราหฺมณ อิเธกจฺโจ อกตปาโป โหติ อกตลุทฺโท อกตกิพฺพิโส, กตกลฺยาโณ กตกุสโล กตภีรุตฺตาโณ. ตเมนํ อญฺญตโร คาโฬฺห โรคาตงฺโก ผุสติ, ตสฺส อญฺญตเรน คาเฬฺหน โรคาตงฺเกน ผุฏฺฐสฺส เอวํ โหติ “อกตํ วต เม ปาปํ, อกตํ ลุทฺทํ, อกตํ กิพฺพิสํ, กตํ กลฺยาณํ, กตํ กุสลํ, กตํ ภีรุตฺตาณํ. ยาวตา โภ อกตปาปานํ อกตลุทฺทานํ อกตกิพฺพิสานํ กตกลฺยาณานํ กตกุสลานํ กตภีรุตฺตาณานํ คติ, ตํ คติํ เปจฺจ คจฺฉามี”ติ. โส น โสจติ น กิลมติ น ปริเทวติ น อุรตฺตาฬิํ กนฺทติ น สมฺโมหํ อาปชฺชติ. อยมฺปิ โข พฺราหฺมณ มรณธมฺโม สมาโน น ภายติ, น สนฺตาสํ อาปชฺชติ มรณสฺส.}

\prose{ปุน จปรํ พฺราหฺมณ อิเธกจฺโจ อกงฺขี โหติ อวิจิกิจฺฉี นิฏฺฐงฺคโต สทฺธมฺเม, ตเมนํ อญฺญตโร คาโฬฺห โรคาตงฺโก ผุสติ, ตสฺส อญฺญตเรน คาเฬฺหน โรคาตงฺเกน ผุฏฺฐสฺส เอวํ โหติ “อกงฺขี วตมฺหิ อวิจิกิจฺฉี นิฏฺฐงฺคโต สทฺธมฺเม”ติ. โส น โสจติ น กิลมติ น ปริเทวติ น อุรตฺตาฬิํ กนฺทติ น สมฺโมหํ อาปชฺชติ. อยมฺปิ โข พฺราหฺมณ มรณธมฺโม สมาโน น ภายติ, น สนฺตาสํ อาปชฺชติ มรณสฺส. อิเม โข พฺราหฺมณ จตฺตาโร มรณธมฺมา สมานา น ภายนฺติ, น สนฺตาสํ อาปชฺชนฺติ มรณสฺสาติ.}

\prose{อภิกฺกนฺตํ โภ โคตม, อภิกฺกนฺตํ โภ โคตม ฯเปฯ อุปาสกํ มํ ภวํ โคตโม ธาเรตุ อชฺชตคฺเค ปาณุเปตํ สรณํ คตนฺติ.  จตุตฺถํ.}

\csromanpalirecto
\csromanmatikaanchor{mtk.223}
\csromantocmarkref{subsection}{5. พฺราหฺมณสจฺจสุตฺต}{mtk.223}
\bookmark[dest={mtk.223},level=2]{5. พฺราหฺมณสจฺจสุตฺต}
\csromanmarkh{จตุกฺกนิปาต\-ปาฬิ}\csromanmarkhii{5. พฺราหฺมณสจฺจ\-สุตฺต}\csromanheadernomark{2}{จตุกฺกนิปาต\-ปาฬิ 5. พฺราหฺมณสจฺจ\-สุตฺต}
\proseitem{185}{\csromanfolio{496}เอกํ สมยํ ภควา ราชคเห วิหรติ คิชฺฌกูเฏ ปพฺพเต. เตน โข ปน สมเยน สมฺพหุลา อภิญฺญาตา อภิญฺญาตา ปริพฺพาชกา สิปฺปินิกาตีเร ปริพฺพาชกา\-ราเม ปฏิวสนฺติ. เสยฺยถิทํ? อนฺนภาโร วรธโร สกุลุทายี จ ปริพฺพาชโก อญฺเญ จ อภิญฺญาตา อภิญฺญาตา ปริพฺพาชกา. อถ โข ภควา สายนฺหสมยํ ปฏิสลฺลานา วุฏฺฐิโต เยน สิปฺปินิกาตีเร ปริพฺพาชกา\-ราโม เตนุปสงฺกมิ.}

\prose{เตน โข ปน สมเยน เตสํ อญฺญ\-ติตฺถิยานํ ปริพฺพาชกานํ สนฺนิสินฺนานํ สนฺนิปติตานํ อยมนฺตรา กถา อุทปาทิ “อิติปิ พฺราหฺมณ\-สจฺจานิ อิติปิ พฺราหฺมณสจฺจานี”ติ. อถ โข ภควา เยน เต ปริพฺพาชกา เตนุปสงฺกมิ, อุปสงฺกมิตฺวา ปญฺญตฺเต อาสเน นิสีทิ, นิสชฺช โข ภควา เต ปริพฺพาชเก เอตทโวจ\csromandash{}}

\prose{“กาย นุตฺถ ปริพฺพาชกา เอตรหิ กถาย สนฺนิสินฺนา, กา จ ปน โว อนฺตรากถา วิปฺปกตา”ติ. อิธ โภ โคตม อมฺหากํ สนฺนิสินฺนานํ สนฺนิปติตานํ อยม\-นฺตรากถา อุทปาทิ “อิติปิ พฺราหฺมณ\-สจฺจานิ อิติปิ พฺราหฺมณสจฺจานี”ติ.}

\prose{จตฺตาริมานิ ปริพฺพาชกา พฺราหฺมณ\-สจฺจานิ มยา สยํ อภิญฺญา สจฺฉิกตฺวา ปเวทิตานิ. กตมานิ จตฺตาริ? อิธ ปริพฺพาชกา พฺราหฺมโณ เอวมาห “สพฺเพ ปาณา อวชฺฌา”ติ, อิติ วทํ พฺราหฺมโณ สจฺจํ อาห โน มุสา. โส เตน น “สมโณ”ติ มญฺญติ, น “พฺราหฺมโณ”ติ มญฺญติ, น “เสยฺโยหมสฺมี”ติ มญฺญติ, น “สทิโสหมสฺมี”ติ มญฺญติ, น “หีโนหมสฺมี”ติ มญฺญติ. อปิ จ ยเทว ตตฺถ สจฺจํ, ตทภิญฺญาย ปาณานํเยว อนุทฺทยาย\csromansharedfootnote{4236e26b3ce130c1}{ตทภิญฺญาย อนุทยาย (ก)} อนุกมฺปาย ปฏิปนฺโน โหติ.}

\prose{ปุน จปรํ ปริพฺพาชกา พฺราหฺมโณ เอวมาห “สพฺเพ กามา อนิจฺจา ทุกฺขา วิปริณาม\-ธมฺมา”ติ, อิติ วทํ พฺราหฺมโณ สจฺจมาห โน มุสา. โส เตน น “สมโณ”ติ มญฺญติ, น “พฺราหฺมโณ”ติ มญฺญติ, น “เสยฺโยหมสฺมี”ติ มญฺญติ, น “สทิโสหมสฺมี”ติ มญฺญติ, น “หีโนหมสฺมี”ติ มญฺญติ. อปิ จ ยเทว ตตฺถ สจฺจํ, ตทภิญฺญาย กามานํเยว นิพฺพิทาย วิราคาย นิโรธาย ปฏิปนฺโน โหติ.}

\chapterheadpageread{(19) 4. พฺราหฺมณวคฺค}
\prose{\csromanfolio{497}ปุน จปรํ ปริพฺพาชกา พฺราหฺมโณ เอวมาห “สพฺเพ ภวา อนิจฺจา ฯเปฯ ตทภิญฺญาย ภวานํเยว นิพฺพิทาย วิราคาย นิโรธาย ปฏิปนฺโน โหติ.}

\prose{ปุน จปรํ ปริพฺพาชกา พฺราหฺมโณ เอวมาห “นาหํ กฺวจนิ\csromansharedfootnote{11b5067f60205336}{กฺวจน (สี, สฺยา)}, กสฺสจิ กิญฺจนตสฺมิํ, น จ มม กฺวจนิ, กตฺถจิ กิญฺจนตตฺถี”ติ, อิติ วทํ พฺราหฺมโณ สจฺจํ อาห โน มุสา. โส เตน น “สมโณ”ติ มญฺญติ, น พฺราหฺมโณ”ติ มญฺญติ, น “เสยฺโยหมสฺมี”ติ มญฺญติ, น “สทิโสหมสฺมี”ติ มญฺญติ, น “หีโนหมสฺมี”ติ มญฺญติ. อปิ จ ยเทว ตตฺถ สจฺจํ, ตทภิญฺญาย อากิญฺจญฺญํเยว ปฏิปทํ ปฏิปนฺโน โหติ. อิมานิ โข ปริพฺพาชกา จตฺตาริ พฺราหฺมณ\-สจฺจานิ มยา สยํ อภิญฺญา สจฺฉิกตฺวา ปเวทิตานีติ.  ปญฺจมํ.}

\csromanmatikaanchor{mtk.224}
\csromantocmarkref{subsection}{6. อุมฺมคฺคสุตฺต}{mtk.224}
\bookmark[dest={mtk.224},level=2]{6. อุมฺมคฺคสุตฺต}
\csromanheader{2}{6. \textbf{อุมฺมคฺคสุตฺต}}
\proseitem{186}{อถ โข อญฺญตโร ภิกฺขุ เยน ภควา เตนุปสงฺกมิ, อุปสงฺกมิตฺวา ภควนฺตํ อภิวาเทตฺวา เอกมนฺตํ นิสีทิ, เอกมนฺตํ นิสินฺโน โข โส ภิกฺขุ ภควนฺตํ เอตทโวจ “เกน นุ โข ภนฺเต โลโก นียติ, เกน โลโก ปริกสฺสติ, กสฺส จ อุปฺปนฺนสฺส วสํ คจฺฉตี”ติ.}

\prose{สาธุ สาธุ ภิกฺขุ ภทฺทโก โข เต ภิกฺขุ อุมฺมคฺโค\csromansharedfootnote{ede363d076a08fd7}{อุมฺมงฺโค (สฺยา, ก)} ภทฺทกํ ปฏิภานํ กลฺยาณี\csromansharedfootnote{acb1f720bd2d0695}{กลฺยาณา (ก)} ปริปุจฺฉา. เอวํ หิ ตฺวํ ภิกฺขุ ปุจฺฉสิ “เกน นุ โข ภนฺเต โลโก นียติ, เกน โลโก ปริกสฺสติ, กสฺส จ อุปฺปนฺนสฺส วสํ คจฺฉตี”ติ. เอวํ ภนฺเต. จิตฺเตน โข ภิกฺขุ โลโก นียติ, จิตฺเตน ปริกสฺสติ, จิตฺตสฺส อุปฺปนฺนสฺส วสํ คจฺฉตีติ.}

\prose{“สาธุ ภนฺเต”ติ โข โส ภิกฺขุ ภควโต ภาสิตํ อภินนฺทิตฺวา อนุโมทิตฺวา ภควนฺตํ อุตฺตริ ปญฺหํ อปุจฺฉิ “พหุสฺสุโต ธมฺมธโร, พหุสฺสุโต ธมฺมธโรติ ภนฺเต วุจฺจติ, กิตฺตาวตา นุ โข ภนฺเต พหุสฺสุโต ธมฺมธโร โหตี”ติ.}

\prose{สาธุ สาธุ ภิกฺขุ ภทฺทโก โข เต ภิกฺขุ อุมฺมคฺโค ภทฺทกํ ปฏิภานํ กลฺยาณี ปริปุจฺฉา. เอวํ หิ ตฺวํ ภิกฺขุ ปุจฺฉสิ “พหุสฺสุโต ธมฺมธโร, พหุสฺสุโต ธมฺมธโรติ ภนฺเต วุจฺจติ, กิตฺตาวตา นุ โข ภนฺเต \csromanfolio{498}พหุสฺสุโต ธมฺมธโร โหตี”ติ. เอวํ ภนฺเต. พหู โข ภิกฺขุ มยา ธมฺมา เทสิตา\csromansharedfootnote{16dc04346e819520}{พหุ โข ภิกฺขุ มยา ธมฺโม เทสิโต (ก)} สุตฺตํ เคยฺยํ เวยฺยากรณํ คาถา อุทานํ อิติวุตฺตกํ ชาตกํ อพฺภุตธมฺมํ เวทลฺลํ, จตุปฺปทาย เจปิ ภิกฺขุ คาถาย อตฺถมญฺญาย ธมฺมมญฺญาย ธมฺมา\-นุธมฺม\-ปฺปฏิปนฺโน โหติ, “พหุสฺสุโต ธมฺมธโร”ติ อลํ วจนายาติ.}

\prose{“สาธุ ภนฺเต”ติ โข โส ภิกฺขุ ภควโต ภาสิตํ อภินนฺทิตฺวา อนุโมทิตฺวา ภควนฺตํ อุตฺตริ ปญฺหํ อปุจฺฉิ “สุตวา นิพฺเพธิก\-ปญฺโญ, สุตวา นิพฺเพธิก\-ปญฺโญติ ภนฺเต วุจฺจติ, กิตฺตาวตา นุ โข ภนฺเต สุตวา นิพฺเพธิก\-ปญฺโญ โหตี”ติ.}

\prose{สาธุ สาธุ ภิกฺขุ ภทฺทโก โข เต ภิกฺขุ อุมฺมคฺโค ภทฺทกํ ปฏิภานํ กลฺยาณี ปริปุจฺฉา. เอวํ หิ ตฺวํ ภิกฺขุ ปุจฺฉสิ “สุตวา นิพฺเพธิก\-ปญฺโญ สุตวา นิพฺเพธิก\-ปญฺโญติ ภนฺเต วุจฺจติ, กิตฺตาวตา นุ โข ภนฺเต สุตวา นิพฺเพธิก\-ปญฺโญ โหตี”ติ. เอวํ ภนฺเต. อิธ ภิกฺขุ ภิกฺขุโน “อิทํ ทุกฺขนฺ”ติ สุตํ โหติ, ปญฺญาย จสฺส อตฺถํ อติวิชฺฌ ปสฺสติ. “อยํ ทุกฺขสมุทโย”ติ สุตํ โหติ, ปญฺญาย จสฺส อตฺถํ อติวิชฺฌ ปสฺสติ. “อยํ ทุกฺขนิโรโธ”ติ สุตํ โหติ, ปญฺญาย จสฺส อตฺถํ อติวิชฺฌ ปสฺสติ. “อยํ ทุกฺขนิโรธ\-คามินี ปฏิปทา”ติ สุตํ โหติ, ปญฺญาย จสฺส อตฺถํ อติวิชฺฌ ปสฺสติ. เอวํ โข ภิกฺขุ สุตวา นิพฺเพธิก\-ปญฺโญ โหตีติ.}

\prose{“สาธุ ภนฺเต”ติ โข โส ภิกฺขุ ภควโต ภาสิตํ อภินนฺทิตฺวา อนุโมทิตฺวา ภควนฺตํ อุตฺตริ ปญฺหํ อปุจฺฉิ “ปณฺฑิโต มหาปญฺโญ, ปณฺฑิโต มหาปญฺโญติ ภนฺเต วุจฺจติ, กิตฺตาวตา นุ โข ภนฺเต ปณฺฑิโต มหาปญฺโญ โหตี”ติ.}

\prose{สาธุ สาธุ ภิกฺขุ ภทฺทโก โข เต ภิกฺขุ อุมฺมคฺโค ภทฺทกํ ปฏิภานํ กลฺยาณี ปริปุจฺฉา. เอวํ หิ ตฺวํ ภิกฺขุ ปุจฺฉสิ “ปณฺฑิโต มหาปญฺโญ ปณฺฑิโต มหาปญฺโญติ ภนฺเต วุจฺจติ, กิตฺตาวตา นุ โข ภนฺเต ปณฺฑิโต มหาปญฺโญ โหตี”ติ. เอวํ ภนฺเต. อิธ ภิกฺขุ ปณฺฑิโต มหาปญฺโญ เนว\-ตฺตพฺยาพาธาย เจเตติ, น ปรพฺยาพาธาย เจเตติ,}

\vspace{\tipitakaparskip}
\csromancenter{(19) 4. พฺราหฺมณวคฺค น อุภย\-พฺยาพาธาย เจเตติ, อตฺตหิตปรหิตอุภยหิตสพฺพโลกหิภเมว จินฺตยมาโน จินฺเตติ. เอวํ โข ภิกฺขุ ปณฺฑิโต มหาปญฺโญ โหตีติ.  ฉฏฺฐํ.}
\csromanmatikaanchor{mtk.225}
\csromantocmarkref{subsection}{7. วสฺสการสุตฺต}{mtk.225}
\bookmark[dest={mtk.225},level=2]{7. วสฺสการสุตฺต}
\csromanheader{2}{7. \textbf{วสฺสการสุตฺต}}
\proseitem{187}{\csromanfolio{499}เอกํ สมยํ ภควา ราชคเห วิหรติ เวฬุวเน กลนฺทก\-นิวาเป. อถ โข วสฺสกาโร พฺราหฺมโณ มคธ\-มหามตฺโต เยน ภควา เตนุปสงฺกมิ, อุปสงฺกมิตฺวา ภควตา สทฺธิํ สมฺโมทิ, สมฺโมทนียํ กถํ สารณียํ วีติสาเรตฺวา เอกมนฺตํ นิสีทิ, เอกมนฺตํ นิสินฺโน โข วสฺสกาโร พฺราหฺมโณ มคธ\-มหามตฺโต ภควนฺตํ เอตทโวจ\csromandash{}}

\prose{ชาเนยฺย นุ โข โภ โคตม อสปฺปุริโส อสปฺปุริสํ “อสปฺปุริโส อยํ ภวนฺ”ติ, อฏฺฐานํ โข เอตํ พฺราหฺมณ อนวกาโส ยํ อสปฺปุริโส อสปฺปุริสํ ชาเนยฺย “อสปฺปุริโส อยํ ภวนฺ”ติ. ชาเนยฺย ปน โภ โคตม อสปฺปุริโส อปฺปุริสํ “สปฺปุริโส อยํ ภวนฺ”ติ, เอตมฺปิ โข พฺราหฺมณ อฏฺฐานํ อนวกาโส, ยํ อสปฺปุริโส อปฺปุริสํ ชาเนยฺย “สปฺปุริโส อยํ ภวนฺ”ติ. ชาเนยฺย นุ โข โภ โคตม สปฺปุริโส สปฺปุริสํ “สปฺปุริโส อยํ ภวนฺ”ติ. ฐานํ โข เอตํ พฺราหฺมณ วิชฺชติ, ยํ สปฺปุริโส สปฺปุริสํ ชาเนยฺย “สปฺปุริโส อยํ ภวนฺ”ติ. ชาเนยฺย ปน โภ โคตม สปฺปุริโส อสปฺปุริสํ “อสปฺปุริโส อยํ ภวนฺ”ติ, เอตมฺปิ โข พฺราหฺมณ ฐานํ วิชฺชติ, ยํ สปฺปุริโส อสปฺปุริสํ ชาเนยฺย “อสปฺปุริโส อยํ ภวนฺ”ติ.}

\prose{อจฺฉริยํ โภ โคตม, อพฺภุตํ โภ โคตม, ยาว สุภาสิตํ จิทํ โภตา โคตเมน “อฏฺฐานํ โข เอตํ พฺราหฺมณ อนวกาโส ยํ อสปฺปุริโส อสปฺปุริสํ ชาเนยฺย ‘อสปฺปุริโส อยํ ภวนฺ’ติ, เอตมฺปิ โข พฺราหฺมณ อฏฺฐานํ อนวกาโส, ยํ อสปฺปุริโส สปฺปุริสํ ชาเนยฺย ‘สปฺปุริโส อยํ ภวนฺ’ติ. ฐานํ โข เอตํ พฺราหฺมณ วิชฺชติ, ยํ สปฺปุริโส สปฺปุริสํ ชาเนยฺย ‘สปฺปุริโส อยํ ภวนฺ’ติ. เอตมฺปิ โข พฺราหฺมณ ฐานํ วิชฺชติ, ยํ สปฺปุริโส อสปฺปุริสํ ชาเนยฺย ‘อสปฺปุริโส อยํ ภวนฺ’ติ”.}

\prose{เอกมิทํ โภ โคตม สมยํ โตเทยฺยสฺส พฺราหฺมณสฺส ปริสติ ปรูปารมฺภํ วตฺเตนฺติ “พาโล อยํ ราชา เอเฬยฺโย สมเณ รามปุตฺเต อภิปฺปสนฺโน, สมเณ จ ปน รามปุตฺเต เอวรูปํ ปรม\-นิปจฺจ\-การํ กโรติ, \csromanfolio{500}ยทิทํ อภิวาทนํ ปจฺจุฏฺฐานํ อญฺชลิกมฺมํ สามีจิกมฺมนฺ”ติ. อิเมปิ รญฺโญ เอเฬยฺยสฺส ปริหารกา พาลา ยมโก โมคฺคลฺโล\csromansharedfootnote{3088832c27e17382}{ปุคฺคโล (ก)} อุคฺโค นาวินฺทกี คนฺธพฺโพ อคฺคิเวสฺโส, เย สมเณ รามปุตฺเต อภิปฺปสนฺนา, สมเณ จ ปน รามปุตฺเต เอวรูปํ ปรม\-นิปจฺจ\-การํ กโรนฺติ, ยทิทํ อภิวาทนํ ปจฺจุฏฺฐานํ อญฺชลิกมฺมํ สามีจิกมฺมนฺติ. ตฺยาสฺสุทํ โตเทยฺโย พฺราหฺมโณ อิมินา นเยน เนติ. ตํ กิํ มญฺญนฺติ โภนฺโต, ปณฺฑิโต ราชา เอเฬยฺโย กรณียา\-ธิกรณีเยสุ วจนียา\-ธิวจนีเยสุ อลม\-ตฺถทส\-ตเรหิ อลม\-ตฺถทส\-ตโรติ. เอวํ โภ ปณฺฑิโต ราชา เอเฬยฺโย กรณียา\-ธิกรณีเยสุ วจนียา\-ธิวจนีเยสุ อลม\-ตฺถทส\-ตเรหิ อลม\-ตฺถทส\-ตโรติ.}

\prose{ยสฺมา จ โข โภ สมโณ รามปุตฺโต รญฺญา เอเฬยฺเยน ปณฺฑิเตน ปณฺฑิตตโร กรณียา\-ธิกรณีเยสุ วจนียา\-ธิวจนีเยสุ, อลม\-ตฺถทส\-ตเรน อลม\-ตฺถทส\-ตโร, ตสฺมา ราชา เอเฬยฺโย สมเณ รามปุตฺเต อภิปฺปสนฺโน, สมเณ จ ปน รามปุตฺเต เอวรูปํ ปรม\-นิปจฺจ\-การํ กโรติ, ยทิทํ อภิวาทนํ ปจฺจุฏฺฐานํ อญฺชลิกมฺมํ สามีจิกมฺมํ.}

\prose{ตํ กิํ มญฺญนฺติ โภนฺโต, ปณฺฑิตา รญฺโญ เอเฬยฺยสฺส ปริหารกา ยมโก โมคฺคลฺโล อุคฺโค นาวินฺทกี คนฺธพฺโพ อคฺคิเวสฺโส กรณียา\-ธิกรณีเยสุ วจนียา\-ธิวจนีเยสุ อลม\-ตฺถทส\-ตเรหิ อลม\-ตฺถทส\-ตเรหิ อลม\-ตฺถทส\-ตราติ. เอวํ โภ ปณฺฑิตา รญฺโญ เอเฬยฺยสฺส ปริหารกา ยมโก โมคฺคลฺโล อุคฺโค นาวินฺทกี คนฺธพฺโพ อคฺคิเวสฺโส กรณียา\-ธิกรณีเยสุ วจนียา\-ธิวจนีเยสุ อลม\-ตฺถทส\-ตเรหิ อลม\-ตฺถทส\-ตราติ.}

\prose{ยสฺมา จ โข โภ สมโณ รามปุตฺโต รญฺโญ เอเฬยฺยสฺส ปริหารเกหิ ปณฺฑิเตหิ ปณฺฑิตตโร กรณียา\-ธิกรณีเยสุ วจนียา\-ธิวจนีเยสุ อลม\-ตฺถทส\-ตเรหิ อลม\-ตฺถทส\-ตโร, ตสฺมา รญฺโญ เอเฬยฺยสฺส ปริหารกา สมเณ รามปุตฺเต อภิปฺปสนฺนา, สมเณ จ ปน รามปุตฺเต เอวรูปํ ปรม\-นิปจฺจ\-การํ กโรนฺติ, ยทิทํ อภิวาทนํ ปจฺจุฏฺฐานํ อญฺชลิกมฺมํ สามีจิกมฺมนฺติ.}

\prose{อจฺฉริยํ โภ โคตม, อพฺภุตํ โภ โคตม, ยาว สุภาสิตํ จิทํ โภตา โคตเมน “อฏฺฐานํ โข เอตํ พฺราหฺมณ อนวกาโส, ยํ}

\vspace{\tipitakaparskip}
\csromancenter{(19) 4. พฺราหฺมณวคฺค อสปฺปุริโส อสปฺปุริสํ ชาเนยฺย ‘อสปฺปุริโส อยํ ภวนฺ’ติ. เอตมฺปิ โข พฺราหฺมณ อฏฺฐานํ อนวกาโส, ยํ อสปฺปุริโส สปฺปุริสํ ชาเนยฺย ‘สปฺปุริโส อยํ ภวนฺ’ติ. ฐานํ โข เอตํ พฺราหฺมณ วิชฺชติ, ยํ สปฺปุริโส สปฺปุริสํ ชาเนยฺย ‘สปฺปุริโส อยํ ภวนฺ’ติ. เอตมฺปิ โข พฺราหฺมณ ฐานํ วิชฺชติ, ยํ สปฺปุริโส อสปฺปุริสํ ชาเนยฺย ‘อสปฺปุริโส อยํ ภวนฺ’ติ. หนฺท จ ทานิ มยํ โภ โคตม คจฺฉาม พหุกิจฺจา มยํ พหุกรณียาติ. ยสฺสทานิ ตฺวํ พฺราหฺมณ กาลํ มญฺญสีติ. อถ โข วสฺสกาโร พฺราหฺมโณ มคธ\-มหามตฺโต ภควโต ภาสิตํ อภินนฺทิตฺวา อนุโมทิตฺวา อุฏฺฐายาสนา ปกฺกามีติ.  สตฺตมํ.}
\csromanmatikaanchor{mtk.226}
\csromantocmarkref{subsection}{8. อุปกสุตฺต}{mtk.226}
\bookmark[dest={mtk.226},level=2]{8. อุปกสุตฺต}
\csromanheader{2}{8. \textbf{อุปกสุตฺต}}
\proseitem{188}{\csromanfolio{501}เอกํ สมยํ ภควา ราชคเห วิหรติ คิชฺฌกูเฏ ปพฺพเต. อถ โข อุปโก มณฺฑิกาปุตฺโต เยน ภควา เตนุปสงฺกมิ, อุปสงฺกมิตฺวา ภควนฺตํ อภิวาเทตฺวา เอกมนฺตํ นิสีทิ, เอกมนฺตํ นิสินฺโน โข อุปโก มณฺฑิกาปุตฺโต ภควนฺตํ เอตทโวจ\csromandash{}}

\prose{อหํ หิ ภนฺเต เอวํวาที เอวํทิฏฺฐิ “โย โกจิ ปรูปารมฺภํ วตฺเตติ, ปรูปารมฺภํ วตฺเตนฺโต สพฺโพ โส\csromansharedfootnote{b2746a459db995b8}{สพฺพโส (สี, อิ)} น อุปปาเทติ, อนุปปาเทนฺโต คารโยฺห โหติ อุปวชฺโช”ติ. ปรูปารมฺภํ เจ อุปก วตฺเตติ, ปรูปารมฺภํ วตฺเตนฺโต น อุปปาเทติ, อนุปปาเทนฺโต คารโยฺห โหติ อุปวชฺโช. ตฺวํ โข อุปก ปรูปารมฺภํ วตฺเตสิ, ปรูปารมฺภํ วตฺเตนฺโต น อุปปาเทสิ, อนุปปาเทนฺโต คารโยฺห โหสิ อุปวชฺโชติ. เสยฺยถาปิ ภนฺเต อุมฺมุชฺชมานกํเยว มหตา ปาเสน พนฺเธยฺย, เอวเมวํ โข อหํ ภนฺเต อุมฺมุชฺชมานโกเยว ภควตา มหตา วาทปาเสน\csromansharedfootnote{8694b4a54529d985}{มหตา ปาเสน (ก)} พทฺโธติ.}

\prose{“อิทํ อกุสลนฺ”ติ โข อุปก มยา ปญฺญตฺตํ, ตตฺถ อปริมาณา ปทา อปริมาณา พฺยญฺชนา อปริมาณา ตถาคตสฺส ธมฺมเทสนา “อิติปิทํ อกุสลนฺ”ติ. “ตํ โข ปนิทํ อกุสลํ ปหาตพฺพนฺ”ติ โข อุปก มยา ปญฺญตฺตํ, ตตฺถ อปริมาณา ปทา อปริมาณา พฺยญฺชนา อปริมาณา ตถาคตสฺส ธมฺมเทสนา ‘อิติปิทํ อกุสลํ ปหาตพฺพนฺ”ติ.}

\prose{\csromanfolio{502}“อิทํ กุสลนฺ”ติ โข อุปก มยา ปญฺญตฺตํ, ตตฺถ อปริมาณา ปทา อปริมาณา พฺยาญฺชนา อปริมาณา ตถาคตสฺส ธมฺมเทสนา “อิติปิทํ กุสลนฺ”ติ, “ตํ โข ปนิทํ กุสลํ ภาเวตพฺพนฺ”ติ โข อุปก มยา ปญฺญตฺตํ, ตตฺถ อปริมาณา ปทา อปริมาณา พฺยญฺชนา อปริมาณา ตถาคตสฺส ธมฺมเทสนา “อิติปิทํ กุสลํ ภาเวตพฺพนฺ”ติ.}

\prose{อถ โข อุปโก มณฺฑิกาปุตฺโต ภควโต ภาสิตํ อภินนฺทิตฺวา อนุโมทิตฺวา อุฏฺฐายาสนา ภควนฺตํ อภิวาเทตฺวา ปทกฺขิณํ กตฺวา เยน ราชา มาคโธ อชาตสตฺตุ เวเทหิปุตฺโต เตนุปสงฺกมิ, อุปสงฺกมิตฺวา ยาวตโก อโหสิ ภควตา สทฺธิํ กถาสลฺลาโป, ตํ สพฺพํ รญฺโญ มาคธสฺส อชาตสตฺตุสฺส เวเทหิปุตฺตสฺส อาโรเจสิ.}

\prose{เอวํ วุตฺเต ราชา มาคโธ อชาตสตฺตุ เวเทหิปุตฺโต กุปิโต อนตฺตมโน อุปกํ มณฺฑิกาปุตฺตํ เอตทโวจ “ยาว ธํสี วตายํ โลณการทารโก ยาว มุขโร ยาว ปคพฺโพ, ยตฺร หิ นาม ตํ ภควนฺตํ อรหนฺตํ สมฺมา\-สมฺพุทฺธํ อาสาเทตพฺพํ มญฺญิสฺสติ, อเปหิ ตฺวํ อุปก วินสฺส, มา ตํ อทฺทสนฺ”ติ.  อฏฺฐมํ.}

\csromanmatikaanchor{mtk.227}
\csromantocmarkref{subsection}{9. สจฺฉิกรณียสุตฺต}{mtk.227}
\bookmark[dest={mtk.227},level=2]{9. สจฺฉิกรณียสุตฺต}
\csromanheader{2}{9. สจฺฉิกรณีย\-สุตฺต}
\proseitem{189}{จตฺตาโรเม ภิกฺขเว สจฺฉิกรณียา ธมฺมา. กตเม จตฺตาโร? อตฺถิ ภิกฺขเว ธมฺมา กาเยน สจฺฉิกรณียา, อตฺถิ ภิกฺขเว ธมฺมา สติยา สจฺฉิกรณียา, อตฺถิ ภิกฺขเว ธมฺมา จกฺขุนา สจฺฉิกรณียา, อตฺถิ ภิกฺขเว ธมฺมา ปญฺญาย สจฺฉิกรณียา. กตเม จ ภิกฺขเว ธมฺมา กาเยน สจฺฉิกรณียา? อฏฺฐ วิโมกฺขา ภิกฺขเว กาเยน สจฺฉิกรณียา.}

\prose{กตเม จ ภิกฺขเว ธมฺมา สติยา สจฺฉิกรณียา? ปุพฺเพนิวาโส ภิกฺขเว สติยา สจฺฉิกรณีโย.}

\prose{กตเม จ ภิกฺขเว ธมฺมา จกฺขุนา สจฺฉิกรณียา? สตฺตานํ จุตูปปาโต ภิกฺขเว จกฺขุนา สจฺฉิกรณีโย.}

\prose{กตเม จ ภิกฺขเว ธมฺมา ปญฺญาย สจฺฉิกรณียา? อาสวานํ ขโย ภิกฺขเว ปญฺญาย สจฺฉิกรณีโย. อิเม โข ภิกฺขเว จตฺตาโร สจฺฉิกรณียา ธมฺมาติ.  นวมํ.}

\csromanmatikaanchor{mtk.228}
\csromantocmarkref{subsection}{10. อุโปสถสุตฺต}{mtk.228}
\bookmark[dest={mtk.228},level=2]{10. อุโปสถสุตฺต}
\csromanmarkh{(19)}\csromanmarkhii{4. พฺราหฺมณวคฺค 10. อุโปสถสุตฺต}\csromanheadernomark{2}{(19) 4. พฺราหฺมณวคฺค \textbf{10. อุโปสถสุตฺต}}
\proseitem{190}{\csromanfolio{503}เอกํ สมยํ ภควา สาวตฺถิยํ วิหรติ ปุพฺพาราเม มิคาร\-มาตุ\-ปาสาเท. เตน โข ปน สมเยน ภควา ตทหุโปสเถ ภิกฺขุ\-สํฆ\-ปริวุโต นิสินฺโน โหติ. อถ โข ภควา ตุณฺหีภูตํ ตุณฺหีภูตํ ภิกฺขุสํฆํ อนุวิโลเกตฺวา ภิกฺขู อามนฺเตสิ\csromandash{}}

\prose{อปลาปายํ ภิกฺขเว ปริสา นิปฺปลาปายํ ภิกฺขเว ปริสา สุทฺธา สาเร ปติฏฺฐิตา. ตถารูโป อยํ ภิกฺขเว ภิกฺขุสํโฆ, ตถารูปายํ ภิกฺขเว ปริสา, ยถารูปา ปริสา ทุลฺลภา ทสฺสนายปิ โลกสฺมิํ. ตถารูโป อยํ ภิกฺขเว ภิกฺขุสํโฆ, ตถารูปายํ ภิกฺขเว ปริสา, ยถารูปา ปริสา อาหุเนยฺยา ปาหุเนยฺยา ทกฺขิเณยฺยา อญฺชลิกรณียา อนุตฺตรํ ปุญฺญกฺเขตฺตํ โลกสฺส. ตถารูโป อยํ ภิกฺขเว ภิกฺขุสํโฆ, ตถารูปายํ ภิกฺขเว ปริสา, ยถารูปาย ปริสาย อปฺปํ ทินฺนํ พหุ โหติ, พหุ ทินฺนํ พหุตรํ. ตถารูโป อยํ ภิกฺขเว ภิกฺขุสํโฆ, ตถารูปายํ ภิกฺขเว ปริสา, ยถารูปํ ปริสํ อลํ โยชน\-คณนานิปิ ทสฺสนาย คนฺตุํ อปิ ปุโฏเสนาปิ. ตถารูโป อยํ ภิกฺขเว ภิกฺขุสํโฆ. (ตถารูปายํ ภิกฺขเว ปริสา)\csromansharedfootnote{b9dab2ad278a4850}{( ) สีสฺยากํอิโปตฺถเกสุ นตฺถิ.}.}

\prose{สนฺติ ภิกฺขเว ภิกฺขู อิมสฺมิํ ภิกฺขุสํเฆ เทวปฺปตฺตา วิหรนฺติ, สนฺติ ภิกฺขเว ภิกฺขู อิมสฺมิํ ภิกฺขุสํเฆ พฺรหฺมปฺปตฺตา วิหรนฺติ, สนฺติ ภิกฺขเว ภิกฺขู อิมสฺมิํ ภิกฺขุสํเฆ อาเนญฺชปฺปตฺตา วิหรนฺติ, สนฺติ ภิกฺขเว ภิกฺขู อิมสฺมิํ ภิกฺขุสํเฆ อริยปฺปตฺตา วิหรนฺติ.}

\prose{กถญฺจ ภิกฺขเว ภิกฺขุ เทวปฺปตฺโต โหติ? อิธ ภิกฺขเว ภิกฺขุ วิวิจฺเจว กาเมหิ ฯเปฯ ปฐมํ ฌานํ อุปสมฺปชฺช วิหรติ, วิตกฺก\-วิจารานํ วูปสมา ฯเปฯ ทุติยํ ฌานํ. ตติยํ ฌานํ. จตุตฺถํ ฌานํ อุปสมฺปชฺช วิหรติ. เอวํ โข ภิกฺขเว ภิกฺขุ เทวปฺปตฺโต โหติ.}

\prose{กถญฺจ ภิกฺขเว ภิกฺขุ พฺรหฺมปฺปตฺโต โหติ? อิธ ภิกฺขเว ภิกฺขุ เมตฺตา\-สหคเตน เจตสา เอกํ ทิสํ ผริตฺวา วิหรติ. ตถา ทุติยํ. ตถา ตติยํ. ตถา จตุตฺถํ. อิติ อุทฺธมโธ ติริยํ สพฺพธิ สพฺพตฺตตาย สพฺพาวนฺตํ โลกํ เมตฺตา\-สหคเตน เจตสา วิปุเลน มหคฺคเตน อปฺปมาเณน อเวเรน อพฺยาปชฺเชน ผริตฺวา วิหรติ. กรุณา. มุทิตา. \csromanfolio{504}อุเปกฺขา\-สหคเตน เจตสา เอกํ ทิสํ ผริตฺวา วิหรติ. ตถา ทุติยํ. ตถา ตติยํ. ตถา จตุตฺถํ. อิติ อุทฺธมโธ ติริยํ สพฺพธิ สพฺพตฺตตาย สพฺพาวนฺตํ โลกํ อุเปกฺขา\-สหคเตน เจตสา วิปุเลน มหคฺคเตน อปฺปมาเณน อเวเรน อพฺยาปชฺเชน ผริตฺวา วิหรติ. เอวํ โข ภิกฺขเว ภิกฺขุ พฺราหฺมปฺปตฺโต โหติ.}

\prose{กถญฺจ ภิกฺขเว ภิกฺขุ อาเนญฺชปฺปตฺโต โหติ? อิธ ภิกฺขเว ภิกฺขุ สพฺพโส รูปสญฺญานํ สมติกฺกมา ปฏิฆ\-สญฺญานํ อตฺถงฺคมา นานตฺต\-สญฺญานํ อมนสิการา “อนนฺโต อากาโส”ติ อากาสา\-นญฺจา\-ยตนํ อุปสมฺปชฺช วิหรติ. สพฺพโส อากาสา\-นญฺจา\-ยตนํ สมติกฺกมฺม “อนนฺตํ วิญฺญาณนฺ”ติ วิญฺญาณญฺจา\-ยตนํ อุปสมฺปชฺช วิหรติ. สพฺพโส วิญฺญาณญฺจา\-ยตนํ สมติกฺกมฺม “นตฺถิ กิญฺจี”ติ อากิญฺจญฺญา\-ยตนํ อุปสมฺปชฺช วิหรติ. สพฺพโส อากิญฺจญฺญา\-ยตนํ สมติกฺกมฺม เนว\-สญฺญา\-นา\-สญฺญา\-ยตนํ อุปสมฺปชฺช วิหรติ. เอวํ โข ภิกฺขเว ภิกฺขุ อาเนญฺชปฺปตฺโต โหติ.}

\prosewithclosermidruled{กถญฺจ ภิกฺขเว ภิกฺขุ อริยปฺปตฺโต โหติ? อิธ ภิกฺขเว ภิกฺขุ “อิทํ ทุกฺขนฺ”ติ ยถาภูตํ ปชานาติ ฯเปฯ “อยํ ทุกฺขนิโรธ\-คามินี ปฏิปทา”ติ ยถาภูตํ ปชานาติ. เอวํ โข ภิกฺขเว ภิกฺขุ อริยปฺปตฺโต โหตีติ.  ทสมํ.}{พฺราหฺมณวคฺโค\csromansharedfootnote{338dd9d64fd7ee1e}{[มิสฺสินฺคฺ โนเต 0]} จตุตฺโถ.}
\csromancenter{\textbf{ตสฺสุทฺทานํ}}
\setlength{\csromangathapagewidth}{0pt}
\setlength{\csromangathawakwidth}{0pt}
\csromangathameasuregroup{โยธา ปาฏิโภคสุตํ, \\ อุมฺมคฺควสฺสกาโร,}{\csromangathabat{โยธา ปาฏิโภคสุตํ,}{อภยํ พฺราหฺมณสจฺเจน ปญฺจมํ.} \\ \csromangathabat{อุมฺมคฺควสฺสกาโร,}{อุปโก สจฺฉิกิริยา จ อุโปสโถติ.}}
\csromangathasetleft{โยธา ปาฏิโภคสุตํ, \\ อุมฺมคฺควสฺสกาโร,}
\csromangathagroup{\csromangathastanza{\csromangathabat{โยธา ปาฏิโภคสุตํ,}{อภยํ พฺราหฺมณสจฺเจน ปญฺจมํ.} \\ \csromangathabat{อุมฺมคฺควสฺสกาโร,}{อุปโก สจฺฉิกิริยา จ อุโปสโถติ.}}}

\csromanmatikaanchor{mtk.229}
\csromantocmarknopage{section}{(20) 5. มหาวคฺค}{mtk.229}
\bookmark[dest={mtk.229},level=1]{(20) 5. มหาวคฺค}
\csromanmarkh{(20)}\csromanmarkhii{5. มหาวคฺค}\csromanheadernomark{1}{\textbf{(20) 5. มหาวคฺค}}
\csromanmatikaanchor{mtk.230}
\csromantocmarkref{subsection}{1. โสตานุคตสุตฺต}{mtk.230}
\bookmark[dest={mtk.230},level=2]{1. โสตานุคตสุตฺต}
\csromanheader{2}{1. โสตานุคต\-สุตฺต}
\proseitem{191}{โสตานุคตานํ ภิกฺขเว ธมฺมานํ วจสา ปริจิตานํ มนสา\-นุเปกฺขิตานํ ทิฏฺฐิยา สุปฺปฏิวิทฺธานํ จตฺตาโร อานิสํสา ปาฏิกงฺขา.}

\vspace{\tipitakaparskip}
\csromancenter{(20) 5. มหาวคฺค กตเม จตฺตาโร? อิธ ภิกฺขเว ภิกฺขุ ธมฺมํ ปริยาปุณาติ สุตฺตํ เคยฺยํ เวยฺยากรณํ คาถํ อุทานํ อิติวุตฺตกํ ชาตกํ อพฺภุตธมฺมํ เวทลฺลํ, ตสฺส เต ธมฺมา โสตานุคตา โหนฺติ วจสา ปริจิตา มนสา\-นุเปกฺขิตา ทิฏฺฐิยา สุปฺปฏิวิทฺธา. โส มุฏฺฐสฺสติ\csromansharedfootnote{178c7ee3aa6535fe}{มุฏฺฐสฺสตี (สี)} กาลํ กุรุมาโน อญฺญตรํ เทวนิกายํ อุปปชฺชติ, ตสฺส ตตฺถ สุขิโน ธมฺมปทา ปฺลวนฺติ\csromansharedfootnote{d87bc4d4d79ab67b}{ปิลปนฺติ (สี, สฺยา, กํ, อิ)}. ทนฺโธ ภิกฺขเว สตุปฺปาโท, อถ โส สตฺโต ขิปฺปํเยว วิเสสคามี โหติ. โสตานุคตานํ ภิกฺขเว ธมฺมานํ วจสา ปริจิตานํ มนสา\-นุเปกฺขิตานํ ทิฏฺฐิยา สุปฺปฏิวิทฺธานํ อยํ ปฐโม อานิสํโส ปาฏิกงฺโข.}
\prose{\csromanfolio{505}ปุน จปรํ ภิกฺขเว ภิกฺขุ ธมฺมํ ปริยาปุณาติ สุตฺตํ เคยฺยํ เวยฺยากรณํ คาถํ อุทานํ อิติวุตฺตกํ ชาตกํ อพฺภุตธมฺมํ เวทลฺลํ, ตสฺส เต ธมฺมา โสตานุคตา โหนฺติ วจสา ปริจิตา มนสา\-นุเปกฺขิตา ทิฏฺฐิยา สุปฺปฏิวิทฺธา. โส มุฏฺฐสฺสติ กาลํ กุรุมาโน อญฺญตรํ เทวนิกายํ อุปปชฺชติ, ตสฺส ตตฺถ น เหว โข สุขิโน ธมฺมปทา ปฺลวนฺติ, อปิ จ โข ภิกฺขุ อิทฺธิมา เจโตวสิปฺปตฺโต เทวปริสายํ ธมฺมํ เทเสติ. ตสฺส เอวํ โหติ “อยํ วา โส ธมฺมวินโย, ยตฺถาหํ ปุพฺเพ พฺรหฺมจริยํ อจรินฺ”ติ. ทนฺโธ ภิกฺขเว สตุปฺปาโท, อถ โส สตฺโต ขิปฺปเมว วิเสสคามี โหติ. เสยฺยถาปิ ภิกฺขเว ปุริโส กุสโล เภริสทฺทสฺส, โส อทฺธาน\-มคฺค\-ปฺปฏิปนฺโน เภริสทฺทํ สุเณยฺย, ตสฺส น เหว โข อสฺส กงฺขา วา วิมติ วา “เภริสทฺโท นุ โข, น นุ โข เภริสทฺโท”ติ, อถ โข “เภริสทฺโท”เตฺวว นิฏฺฐํ คจฺเฉยฺย. เอวเมวํ โข ภิกฺขเว ภิกฺขุ ธมฺมํ ปริยาปุณาติ สุตฺตํ เคยฺยํ เวยฺยากรณํ คาถํ อุทานํ อิติวุตฺตกํ ชาตกํ อพฺภุตธมฺมํ เวทลฺลํ, ตสฺส เต ธมฺมา โสตานุคตา โหนฺติ วจสา ปริจิตา มนสา\-นุเปกฺขิตา ทิฏฺฐิยา สุปฺปฏิวิทฺธา. โส มุฏฺฐสฺสติ กาลํ กุรุมาโน อญฺญตรํ เทวนิกายํ อุปปชฺชติ, ตสฺส ตตฺถ น เหว โข สุขิโน ธมฺมปทา ปฺลวนฺติ, อปิ จ โข ภิกฺขุ อิทฺธิมา เจโตวสิปฺปตฺโต เทวปริสายํ ธมฺมํ เทเสติ. ตสฺส เอวํ โหติ “อยํ วา โส ธมฺมวินโย, ยตฺถายํ ปุพฺเพ พฺรหฺมจริยํ อจรินฺ”ติ. ทนฺโธ ภิกฺขเว สตุปฺปาโท, อถ โส สตฺโต ขิปฺปํเยว วิเสสคามี โหติ. โสตานุคตานํ ภิกฺขเว \csromanfolio{506}ธมฺมานํ วจสา ปริจิตานํ มนสา\-นุเปกฺขิตานํ ทิฏฺฐิยา สุปฺปฏิวิทฺธานํ อยํ ทุติโย อานิสํโส ปาฏิกงฺโข.}

\prose{ปุน จปรํ ภิกฺขเว ภิกฺขุ ธมฺมํ ปริยาปุณาติ สุตฺตํ เคยฺยํ เวยฺยากรณํ คาถํ อุทานํ อิติวุตฺตกํ ชาตกํ อพฺภุตธมฺมํ เวทลฺลํ, ตสฺส เต ธมฺมา โสตานุคตา โหนฺติ วจสา ปริจิตา มนสา\-นุเปกฺขิตา ทิฏฺฐิยา สุปฺปฏิวิทฺธา. โส มุฏฺฐสฺสติ กาลํ กุรุมาโน อญฺญตรํ เทวนิกายํ อุปปชฺชติ, ตสฺส ตตฺถ น เหว โข สุขิโน ธมฺมปทา ปฺลวนฺติ, นปิ ภิกฺขุ อิทฺธิมา เจโตวสิปฺปตฺโต เทวปริสายํ ธมฺมํ เทเสติ, อปิ จ โข เทวปุตฺโต เทวปริสายํ ธมฺมํ เทเสติ. ตสฺส เอวํ โหติ “อยํ วา โส ธมฺมวินโย, ยตฺถาหํ ปุพฺเพ พฺรหฺมจริยํ อจรินฺ”ติ. ทนฺโธ ภิกฺขเว สตุปฺปาโท, อถ โส สตฺโต ขิปฺปํเยว วิเสสคามี โหติ. เสยฺยถาปิ ภิกฺขเว ปุริโส กุสโล สงฺขสทฺทสฺส, โส อทฺธาน\-มคฺค\-ปฺปฏิปนฺโน สงฺขสทฺทํ สุเณยฺย, ตสฺส น เหว โข อสฺส กงฺขา วา วิมติ วา “สงฺขสทฺโท นุ โข, น นุ โข สงฺขสทฺโท”ติ. อถ โข “สงฺขสทฺโท”เตฺวว นิฏฺฐํ คจฺเฉยฺย. เอวเมวํ โข ภิกฺขเว ภิกฺขุ ธมฺมํ ปริยาปุณาติ สุตฺตํ เคยฺยํ เวยฺยากรณํ คาถํ อุทานํ อิติวุตฺตกํ ชาตกํ อพฺภุตธมฺมํ เวทลฺลํ, ตสฺส เต ธมฺมา โสตานุคตา โหนฺติ วจสา ปริจิตา มนสา\-นุเปกฺขิตา ทิฏฺฐิยา สุปฺปฏิวิทฺธา. โส มุฏฺฐสฺสติ กาลํ กุรุมาโน อญฺญตรํ เทวนิกายํ อุปปชฺชติ, ตสฺส ตตฺถ น เหว โข สุขิโน ธมฺมปทา ปฺลวนฺติ, นปิ ภิกฺขุ อิทฺธิมา เจโตวสิปฺปตฺโต เทวปริสายํ ธมฺมํ เทเสติ, อปิ จ โข เทวปุตฺโต เทวปริสายํ ธมฺมํ เทเสติ. ตสฺส เอวํ โหติ “อยํ วา โส ธมฺมวินโย, ยตฺถาหํ ปุพฺเพ พฺรหฺมจริยํ อจรินฺ”ติ. ทนฺโธ ภิกฺขเว สตุปฺปาโท, อถ โส สตฺโต ขิปฺปํเยว วิเสสคามี โหติ. โสตานุคตานํ ภิกฺขเว ธมฺมานํ วจสา ปริจิตานํ มนสา\-นุเปกฺขิตานํ ทิฏฺฐิยา สุปฺปฏิวิทฺธานํ อยํ ตติโย อานิสํโส ปาฏิกงฺโข.}

\prose{ปุน จปรํ ภิกฺขเว ภิกฺขุ ธมฺมํ ปริยาปุณาติ สุตฺตํ เคยฺยํ เวยฺยากรณํ คาถํ อุทานํ อิติวุตฺตกํ ชาตกํ อพฺภุตธมฺมํ เวทลฺลํ, ตสฺส เต ธมฺมา โสตานุคตา โหนฺติ วจสา ปริจิตา มนสา\-นุเปกฺขิตา ทิฏฺฐิยา สุปฺปฏิวิทฺธา. โส มุฏฺฐสฺสติ กาลํ กุรุมาโน อญฺญตรํ เทวนิกายํ อุปปชฺชติ, ตสฺส ตตฺถ น เหว โข สุขิโน ธมฺมปทา ปฺลวนฺติ, นปิ ภิกฺขุ อิทฺธิมา เจโตวสิปฺปตฺโต เทวปริสายํ ธมฺมํ เทเสติ, นปิ เทวปุตฺโต}

\vspace{\tipitakaparskip}
\csromancenter{(20) 5. มหาวคฺค เทวปริสายํ ธมฺมํ เทเสติ, อปิ จ โข โอปปาติโก โอปปาติกํ สาเรติ “สรสิ ตฺวํ มาริส, สรสิ ตฺวํ มาริส, ยตฺถ มยํ ปุพฺเพ พฺรหฺมจริยํ อจริมฺหา”ติ. โส เอวมาห “สรามิ มาริส, สรามิ มาริสา”ติ. ทนฺโธ ภิกฺขเว สตุปฺปาโท, อถ โส สตฺโต ขิปฺปํเยว วิเสสคามี โหติ. เสยฺยถาปิ ภิกฺขเว เทฺว สหายกา สหปํสุ\-กีฬิกา\csromansharedfootnote{d04858490fb21525}{สหปํสุกีฬกา (สฺยา, กํ)}, เต กทาจิ กรหจิ อญฺญมญฺญํ สมาคจฺเฉยฺยุํ, อญฺโญ ปน\csromansharedfootnote{82d3cec0c538acbd}{สมาคจฺเฉยฺยุํ, ตเมนํ (สี, สฺยา, กํ, อิ)} สหายโก สหายกํ เอวํ วเทยฺย “อิทมฺปิ สมฺม สรสิ, อิทมฺปิ สมฺม สรสี”ติ. โส เอวํ วเทยฺย “สรามิ สมฺม, สรามิ สมฺมา”ติ. เอวเมวํ โข ภิกฺขเว ภิกฺขุ ธมฺมํ ปริยาปุณาติ สุตฺตํ เคยฺยํ เวยฺยากรณํ คาถํ อุทานํ อิติวุตฺตกํ ชาตกํ อพฺภุตธมฺมํ เวทลฺลํ, ตสฺส เต ธมฺมา โสตานุคตา โหนฺติ วจสา ปริจิตา มนสา\-นุเปกฺขิตา ทิฏฺฐิยา สุปฺปฏิวิทฺธา. โส มุฏฺฐสฺสติ กาลํ กุรุมาโน อญฺญตรํ เทวนิกายํ อุปปชฺชติ, ตสฺส ตตฺถ น เหว โข สุขิโน ธมฺมปทา ปฺลวนฺติ, นปิ ภิกฺขุ อิทฺธิมา เจโตวสิปฺปตฺโต เทวปริสายํ ธมฺมํ เทเสติ, นปิ เทวปุตฺโต เทวปริสายํ ธมฺมํ เทเสติ, อปิ จ โข โอปปาติโก โอปปาติกํ สาเรติ “สรสิ ตฺวํ มาริส, สรสิ ตฺวํ มาริส, ยตฺถ มยํ ปุพฺเพ พฺรหฺมจริยํ อจริมฺหา”ติ. โส เอวมาห “สรามิ มาริส, สรามิ มาริสา”ติ. ทนฺโธ ภิกฺขเว สตุปฺปาโท, อถ โข โส สตฺโต ขิปฺปํเยว วิเสสคามี โหติ. โสตานุคตานํ ภิกฺขเว ธมฺมานํ วจสา ปริจิตานํ มนสา\-นุเปกฺขิตานํ ทิฏฺฐิยา สุปฺปฏิวิทฺธานํ อยํ จตุตฺโถ อานิสํโส ปาฏิกงฺโข. โสตานุคตานํ ภิกฺขเว ธมฺมานํ วจสา ปริจิตานํ มนสา\-นุเปกฺขิตานํ ทิฏฺฐิยา สุปฺปฏิวิทฺธานํ อิเม จตฺตาโร อานิสํสา ปาฏิกงฺขาติ.  ปฐมํ.}
\csromanmatikaanchor{mtk.231}
\csromantocmarkref{subsection}{2. ฐานสุตฺต}{mtk.231}
\bookmark[dest={mtk.231},level=2]{2. ฐานสุตฺต}
\csromanheader{2}{2. \textbf{ฐานสุตฺต}}
\proseitem{192}{\csromanfolio{507}จตฺตาริมานิ ภิกฺขเว ฐานานิ จตูหิ ฐาเนหิ เวทิตพฺพานิ. กตมานิ จตฺตาริ? สํวาเสน ภิกฺขเว สีลํ เวทิตพฺพํ, ตญฺจ โข ทีเฆน อทฺธุนา น อิตฺตรํ, มนสิกโรตา โน อมนสิกโรตา, ปญฺญวตา โน ทุปฺปญฺเญน. สํโวหาเรน ภิกฺขเว โสเจยฺยํ เวทิตพฺพํ, ตญฺจ โข ทีเฆน อทฺธุนา น อิตฺตรํ, มนสิกโรตา โน อมนสิกโรตา, ปญฺญวตา โน ทุปฺปญฺเญน. อาปทาสุ ภิกฺขเว ถาโม เวทิตพฺโพ, โส จ โข ทีเฆน \csromanfolio{508}อทฺธุนา น อิตฺตรํ, มนสิกโรตา โน อมนสิกโรตา, ปญฺญวตา โน ทุปฺปญฺเญน. สากจฺฉาย ภิกฺขเว ปญฺญา เวทิตพฺพา, สา จ โข ทีเฆน อทฺธุนา น อิตฺตรํ, มนสิกโรตา โน อมนสิกโรตา, ปญฺญวตา โน ทุปฺปญฺเญนาติ.}

\prose{\csromansymbolfootnote{*}{สํ 1. 78 ปิฏฺเฐปิ.}“สํวาเสน ภิกฺขเว สีลํ เวทิตพฺพํ, ตญฺจ โข ทีเฆน อทฺธุนา น อิตฺตรํ, มนสิกโรตา โน อมนสิกโรตา, ปญฺญวตา โน ทุปฺปญฺเญนา”ติ อิติ โข ปเนตํ วุตฺตํ, กิญฺเจตํ ปฏิจฺจ วุตฺตํ. อิธ ภิกฺขเว ปุคฺคโล ปุคฺคเลน สทฺธิํ สํวสมาโน เอวํ ชานาติ “ทีฆรตฺตํ โข อยมายสฺมา ขณฺฑการี ฉิทฺทการี สพลการี กมฺมาสการี น สนฺตตการี น สนฺตตวุตฺติ\csromansharedfootnote{278c5cc58d428e29}{สตตวุตฺติ (สฺยา, กํ)} สีเลสุ ทุสฺสีโล อยมายสฺมา, นายมายสฺมา สีลวา”ติ.}

\prose{อิธ ปน ภิกฺขเว ปุคฺคโล ปุคฺคเลน สทฺธิํ สํวสมาโน เอวํ ชานาติ “ทีฆรตฺตํ โข อยมายสฺมา อขณฺฑการี อจฺฉิทฺทการี อสพลการี อกมฺมาสการี สนฺตตการี สนฺตตวุตฺติ สีเลสุ, สีลวา อยมายสฺมา, นายมายสฺมา ทุสฺสีโล”ติ. “สํวาเสน ภิกฺขเว สีลํ เวทิตพฺพํ, ตญฺจ โข ทีเฆน อทฺธุนา น อิตฺตรํ, มนสิกโรตา โน อมนสิกโรตา, ปญฺญวตา โน ทุปฺปญฺเญนา”ติ อิติ ยํ ตํ วุตฺตํ, อิทเมตํ ปฏิจฺจ วุตฺตํ. (1)}

\prose{“สํโวหาเรน ภิกฺขเว โสเจยฺยํ เวทิตพฺพํ, ตญฺจ โข ทีเฆน อทฺธุนา น อิตฺตรํ, มนสิกโรตา โน อมนสิกโรตา, ปญฺญวตา โน ทุปฺปญฺเญนา”ติ อิติ โข ปเนตํ วุตฺตํ, กิญฺเจตํ ปฏิจฺจ วุตฺตํ. อิธ ภิกฺขเว ปุคฺคโล ปุคฺคเลน สทฺธิํ สํโวหรมาโน เอวํ ชานาติ “อญฺญถา โข อยมายสฺมา เอเกน เอโก โวหรติ, อญฺญถา ทฺวีหิ อญฺญถา ตีหิ อญฺญถา สมฺพหุเลหิ, โวกฺกมติ อยมายสฺมา ปุริมโวหารา ปจฺฉิม\-โวหารํ, อปริสุทฺธ\-โวหาโร อยมายสฺมา, นายมายสฺมา ปริสุทฺธ\-โวหาโร”ติ.}

\prose{อิธ ปน ภิกฺขเว ปุคฺคโล ปุคฺคเลน สทฺธิํ สํโวหรมาโน เอวํ ชานาติ “ยเถว โข อยมายสฺมา เอเกน เอโก โวหรติ, ตถา ทฺวีหิ ตถา ตีหิ ตถา สมฺพหุเลหิ, นายมายสฺมา โวกฺกมติ ปุริมโวหารา ปจฺฉิม\-โวหารํ. ปริสุทฺธ\-โวหาโร อยมายสฺมา, นายมายสฺมา}

\vspace{\tipitakaparskip}
\csromancenter{(20) 5. มหาวคฺค อปริสุทฺธ\-โวหาโร”ติ. “สํโวหาเรน ภิกฺขเว โสเจยฺยํ เวทิตพฺพํ, ตญฺจ โข ทีเฆน อทฺธุนา น อิตฺตรํ, มนสิกโรตา โน อมนสิกโรตา, ปญฺญวตา โน ทุปฺปญฺเญนา”ติ อิติ ยํ ตํ วุตฺตํ, อิทเมตํ ปฏิจฺจ วุตฺตํ. (2)}
\prose{\csromanfolio{509}“อาปทาสุ ภิกฺขเว ถาโม เวทิตพฺโพ, โส จ โข ทีเฆน อทฺธุนา น อิตฺตรํ, มนสิกโรตา โน อมนสิกโรตา, ปญฺญวตา โน ทุปฺปญฺเญนา”ติ อิติ โข ปเนตํ วุตฺตํ, กิญฺเจตํ ปฏิจฺจ วุตฺตํ. อิธ ภิกฺขเว เอกจฺโจ ญาติพฺยสเนน วา ผุฏฺโฐ สมาโน, โภคพฺยสเนน วา ผุฏฺโฐ สมาโน, โรคพฺยสเนน วา ผุฏฺโฐ สมาโน น อิติ ปฏิสญฺจิกฺขติ “ตถาภูโต โข อยํ โลกสนฺนิวาโส, ตถาภูโต อยํ อตฺตภาว\-ปฏิลาโภ, ยถาภูเต โลกสนฺนิวาเส ยถาภูเต อตฺตภาว\-ปฏิลาเภ อฏฺฐ โลกธมฺมา โลกํ อนุปริ\-วตฺตนฺติ, โลโก จ อฏฺฐ โลกธมฺเม อนุปริวตฺตติ\csromandash{}ลาโภ จ อลาโภ จ ยโส จ อยโส จ นินฺทา จ ปสํสา จ สุขญฺจ ทุกฺขญฺจา”ติ. โส ญาติพฺยสเนน วา ผุฏฺโฐ สมาโน, โภคพฺยสเนน วา ผุฏฺโฐ สมาโน, โรคพฺยสเนน วา ผุฏฺโฐ สมาโน โสจติ กิลมติ ปริเทวติ อุรตฺตาฬิํ กนฺทติ สมฺโมหํ อาปชฺชติ.}

\prose{อิธ ปน ภิกฺขเว เอกจฺโจ ญาติพฺยสเนน วา ผุฏฺโฐ สมาโน, โภคพฺยสเนน วา ผุฏฺโฐ สมาโน, โรคพฺยสเนน วา ผุฏฺโฐ สมาโน อิติ ปฏิสญฺจิกฺขติ “ตถาภูโต โข อยํ โลกสนฺนิวาโส, ตถาภูโต อยํ อตฺตภาว\-ปฏิลาโภ, ยถาภูเต โลกสนฺนิวาเส ยถาภูเต อตฺตภาว\-ปฏิลาเภ อฏฺฐ โลกธมฺมา โลกํ อนุปริ\-วตฺตนฺติ, โลโก จ อฏฺฐ โลกธมฺเม อนุปริวตฺตติ\csromandash{}ลาโภ จ อลาโภ จ ยโส จ อยโส จ นินฺทา จ ปสํสา จ สุขญฺจ ทุกฺขญฺจา”ติ. โส ญาติพฺยสเนน วา ผุฏฺโฐ สมาโน, โภคพฺยสเนน วา ผุฏฺโฐ สมาโน, โรคพฺยสเนน วา ผุฏฺโฐ สมาโน น โสจติ น กิลมติ น ปริเทวติ น อุรตฺตาฬิํ กนฺทติ น สมฺโมหํ อาปชฺชติ. “อาปทาสุ ภิกฺขเว ถาโม เวทิตพฺโพ, โส จ โข ทีเฆน อทฺธุนา น อิตฺตรํ, มนสิกโรตา โน อมนสิกโรตา, ปญฺญวตา โน ทุปฺปญฺเญนา”ติ อิติ ยํ ตํ วุตฺตํ, อิทเมตํ ปฏิจฺจ วุตฺตํ. (3)}

\prose{“สากจฺฉาย ภิกฺขเว ปญฺญา เวทิตพฺพา, สา จ โข ทีเฆน อทฺธุนา น อิตฺตรํ, มนสิกโรตา โน อมนสิกโรตา, ปญฺญวตา โน \csromanfolio{510}ทุปฺปญฺเญนา”ติ อิติ โข ปเนตํ วุตฺตํ, กิญฺเจตํ ปฏิจฺจ วุตฺตํ. อิธ ภิกฺขเว ปุคฺคโล ปุคฺคเลน สทฺธิํ สากจฺฉายมาโน เอวํ ชานาติ “ยถา โข อิมสฺส อายสฺมโต อุมฺมคฺโค ยถา จ อภินีหาโร ยถา จ ปญฺหา\-สมุทาหาโร, ทุปฺปญฺโญ อยมายสฺมา, นายมายสฺมา ปญฺญวา. ตํ กิสฺส เหตุ? ตถา หิ อยมายสฺมา น เจว คมฺภีรํ อตฺถปทํ อุทาหรติ สนฺตํ ปณีตํ อตกฺกาวจรํ นิปุณํ ปณฺฑิต\-เวทนียํ. ยญฺจ อยมายสฺมา ธมฺมํ ภาสติ, ตสฺส จ นปฺปฏิพโล สํขิตฺเตน วา วิตฺถาเรน วา อตฺถํ อาจิกฺขิตุํ เทเสตุํ ปญฺญาเปตุํ ปฏฺฐเปตุํ วิวริตุํ วิภชิตุํ อุตฺตานีกาตุํ. ทุปฺปญฺโญ อยมายสฺมา, นายมายสฺมา ปญฺญวา”ติ.}

\prose{เสยฺยถาปิ ภิกฺขเว จกฺขุมา ปุริโส อุทกรหทสฺส ตีเร ฐิโต ปสฺเสยฺย ปริตฺตํ มจฺฉํ อุมฺมุชฺชมานํ. ตสฺส เอวมสฺส “ยถา โข อิมสฺส มจฺฉสฺส อุมฺมคฺโค ยถา จ อูมิฆาโต ยถา จ เวคายิตตฺตํ, ปริตฺโต อยํ มจฺโฉ, นายํ มจฺโฉ มหนฺโต”ติ. เอวเมวํ โข ภิกฺขเว ปุคฺคโล ปุคฺคเลน สทฺธิํ สากจฺฉายมาโน เอวํ ชานาติ “ยถา โข อิมสฺส อายสฺมโต อุมฺมคฺโค ยถา จ อภินีหาโร ยถา จ ปญฺหา\-สมุทาหาโร, ทุปฺปญฺโญ อยมายสฺมา, นายมายสฺมา ปญฺญวา. ตํ กิสฺส เหตุ? ตถา หิ อยมายสฺมา น เจว คมฺภีรํ อตฺถปทํ อุทาหรติ สนฺตํ ปณีตํ อตกฺกาวจรํ นิปุณํ ปณฺฑิต\-เวทนียํ. ยญฺจ อยมายสฺมา ธมฺมํ ภาสติ, ตสฺส จ นปฏิพโล สํขิตฺเตน วา วิตฺถาเรน วา อตฺถํ อาจิกฺขิตุํ เทเสตุํ ปญฺญาเปตุํ ปฏฺฐเปตุํ วิวริตุํ วิภชิตุํ อุตฺตานีกาตุํ. ทุปฺปญฺโญ อยมายสฺมา, นายมายสฺมา ปญฺญวา”ติ.}

\prose{อิธ ปน ภิกฺขเว ปุคฺคโล ปุคฺคเลน สทฺธิํ สากจฺฉายมาโน เอวํ ชานาติ “ยถา โข อิมสฺส อายสฺมโต อุมฺมคฺโค ยถา จ อภินีหาโร ยถา จ ปญฺหา\-สมุทาหาโร, ปญฺญวา อยมายสฺมา, นายมายสฺมา ทุปฺปญฺโญ. ตํ กิสฺส เหตุ? ตถา หิ อยมายสฺมา คมฺภีรญฺเจว อตฺถปทํ อุทาหรติ สนฺตํ ปณีตํ อตกฺกาวจรํ นิปุณํ ปณฺฑิต\-เวทนียํ. ยญฺจ อยมายสฺมา ธมฺมํ ภาสติ, ตสฺส จ ปฏิพโล สํขิตฺเตน วา วิตฺถาเรน วา อตฺถํ อาจิกฺขิตุํ เทเสตุํ ปญฺญาเปตุํ ปฏฺฐเปตุํ วิวริตุํ วิภชิตุํ อุตฺตานีกาตุํ. ปญฺญวา อยมายสฺมา, นายมายสฺมา ทุปฺปญฺโญ”ติ.}

\chapterheadpageread{(20) 5. มหาวคฺค}
\prose{\csromanfolio{511}เสยฺยถาปิ ภิกฺขเว จกฺขุมา ปุริโส อุทกรหทสฺส ตีเร ฐิโต ปสฺเสยฺย มหนฺตํ มจฺฉํ อุมฺมุชฺชมานํ. ตสฺส เอวมสฺส “ยถา โข อิมสฺส มจฺฉสฺส อุมฺมคฺโค ยถา จ อูมิฆาโต ยถา จ เวคายิตตฺตํ, มหนฺโต อยํ มจฺโฉ, นายํ มจฺโฉ ปริตฺโต”ติ. เอวเมวํ โข ภิกฺขเว ปุคฺคโล ปุคฺคเลน สทฺธิํ สากจฺฉายมาโน เอวํ ชานาติ “ยถา โข อิมสฺส อายสฺมโต อุมฺมคฺโค ยถา จ อภินีหาโร ยถา จ ปญฺหา\-สมุทาหาโร, ปญฺญวา อยมายสฺมา, นายมายสฺมา ทุปฺปญฺโญ. ตํ กิสฺส เหตุ? ตถา หิ อยมายสฺมา คมฺภีรญฺเจว อตฺถปทํ อุทาหรติ สนฺตํ ปณีตํ อตกฺกาวจรํ นิปุณํ ปณฺฑิต\-เวทนียํ. ยญฺจ อยมายสฺมา ธมฺมํ ภาสติ, ตสฺส จ ปฏิพโล สํขิตฺเตน วา วิตฺถาเรน วา อตฺถํ อาจิกฺขิตุํ เทเสตุํ ปญฺญาเปตุํ ปฏฺฐเปตุํ วิวริตุํ วิภชิตุํ อุตฺตานีกาตุํ. ปญฺญวา อยมายสฺมา, นายมายสฺมา ทุปฺปญฺโญ”ติ. (4)}

\prose{“สากจฺฉาย ภิกฺขเว ปญฺญา เวทิตพฺพา, สา จ โข ทีเฆน อทฺธุนา น อิตฺตรํ, มนสิกโรตา โน อมนสิกโรตา, ปญฺญวตา โน ทุปฺปญฺเญนา”ติ อิติ ยํ ตํ วุตฺตํ, อิทเมตํ ปฏิจฺจ วุตฺตํ. อิมานิ โข ภิกฺขเว จตฺตาริ ฐานานิ อิเมหิ จตูหิ ฐาเนหิ เวทิตพฺพานีติ.  ทุติยํ.}

\csromanmatikaanchor{mtk.232}
\csromantocmarkref{subsection}{3. ภทฺทิยสุตฺต}{mtk.232}
\bookmark[dest={mtk.232},level=2]{3. ภทฺทิยสุตฺต}
\csromanheader{2}{3. \textbf{ภทฺทิยสุตฺต}}
\proseitem{193}{เอกํ สมยํ ภควา เวสาลิยํ วิหรติ มหาวเน กูฏาคาร\-สาลายํ. อถ โข ภทฺทิโย ลิจฺฉวิ เยน ภควา เตนุปสงฺกมิ, อุปสงฺกมิตฺวา ภควนฺตํ อภิวาเทตฺวา เอกมนฺตํ นิสีทิ, เอกมนฺตํ นิสินฺโน โข ภทฺทิโย ลิจฺฉวิ ภควนฺตํ เอตทโวจ\csromandash{}}

\prose{“สุตํ เมตํ ภนฺเต ‘มายาวี สมโณ โคตโม อาวฏฺฏนิํ มายํ\csromansharedfootnote{3aff2fab4efc006a}{อาวฏฺฏนีมายํ (สี), อาวฏฺฏนิมายํ (สฺยา, กํ, ก) ม 2. 38 ปิฏฺเฐปิ ปสฺสิตพฺพํ.} ชานาติ, ยาย อญฺญ\-ติตฺถิยานํ สาวเก อาวฏฺเฏตี’ติ. เย เต ภนฺเต เอวมาหํสุ ‘มายาวี สมโณ โคตโม อาวฏฺฏนิํ มายํ ชานาติ, ยาย อญฺญ\-ติตฺถิยานํ สาวเก อาวฏฺเฏตี’ติ. กจฺจิ เต ภนฺเต ภควโต วุตฺตวาทิโน, น จ ภควนฺตํ อภูเตน อพฺภาจิกฺขนฺติ, ธมฺมสฺส จ อนุธมฺมํ พฺยากโรนฺติ, น จ โกจิ สหธมฺมิโก วาทานุปาโต คารยฺหํ ฐานํ อาคจฺฉติ, อนพฺภกฺขาตุ\-กามา หิ มยํ ภนฺเต ภควนฺตนฺ”ติ.}

\prose{\csromanfolio{512}เอถ ตุเมฺห ภทฺทิย, มา อนุสฺสเวน, มา ปรมฺปราย, มา อิติกิราย, มา ปิฏก\-สมฺปทาเนน, มา ตกฺกเหตุ, มา นยเหตุ, มา อาการ\-ปริวิตกฺเกน, มา ทิฏฺฐิ\-นิชฺฌาน\-กฺขนฺติยา, มา ภพฺพรูปตาย, มา สมโณ โน ครูติ. ยทา ตุเมฺห ภทฺทิย อตฺตนาว ชาเนยฺยาถ “อิเม ธมฺมา อกุสลา, อิเม ธมฺมา สาวชฺชา, อิเม ธมฺมา วิญฺญุครหิตา, อิเม ธมฺมา สมตฺตา สมาทินฺนา อหิตาย ทุกฺขาย สํวตฺตนฺตี”ติ, อถ ตุเมฺห ภทฺทิย ปชเหยฺยาถ.}

\prose{ตํ กิํ มญฺญถ ภทฺทิย, โลโภ ปุริสสฺส อชฺฌตฺตํ อุปฺปชฺชมาโน อุปฺปชฺชติ หิตาย วา อหิตาย วาติ. อหิตาย ภนฺเต. ลุทฺโธ ปนายํ ภทฺทิย ปุริสปุคฺคโล โลเภน อภิภูโต ปริยาทินฺน\-จิตฺโต ปาณมฺปิ หนติ, อทินฺนมฺปิ อาทิยติ, ปรทารมฺปิ คจฺฉติ, มุสาปิ ภณติ, ปรมฺปิ ตถตฺตาย\csromansharedfootnote{99dfbfcbc3933f8e}{ตทตฺถาย (ก) เหฏฺฐา 191 ปิฏฺเฐปิ.} สมาทเปติ, ยํ’ส โหติ ทีฆรตฺตํ อหิตาย ทุกฺขายาติ. เอวํ ภนฺเต.}

\prose{ตํ กิํ มญฺญถ ภทฺทิย, โทโส ปุริสสฺส ฯเปฯ โมโห ปุริสสฺส ฯเปฯ สารมฺโภ ปุริสสฺส อชฺฌตฺตํ อุปฺปชฺชมาโน อุปฺปชฺชติ หิตาย วา อหิตาย วาติ. อหิตาย ภนฺเต. สารทฺโธ ปนายํ ภทฺทิย ปุริสปุคฺคโล สารมฺเภน อภิภูโต ปริยาทินฺน\-จิตฺโต ปาณมฺปิ หนติ, อทินฺนมฺปิ อาทิยติ, ปรทารมฺปิ คจฺฉติ, มุสาปิ ภณติ, ปรมฺปิ ตถตฺตาย สมาทเปติ, ยํ’ส โหติ ทีฆรตฺตํ อหิตาย ทุกฺขายาติ. เอวํ ภนฺเต.}

\prose{ตํ กิํ มญฺญถ ภทฺทิย, อิเม ธมฺมา กุสลา วา อกุสลา วาติ. อกุสลา ภนฺเต. สาวชฺชา วา อนวชฺชา วาติ. สาวชฺชา ภนฺเต. วิญฺญุครหิตา วา วิญฺญุปฺปสตฺถา วาติ. วิญฺญุครหิตา ภนฺเต. สมตฺตา สมาทินฺนา อหิตาย ทุกฺขาย สํวตฺตนฺติ โน วา กถํ วา เอตฺถ โหตีติ. สมตฺตา ภนฺเต สมาทินฺนา อหิตาย ทุกฺขาย สํวตฺตนฺติ, เอวํ โน เอตฺถ โหตีติ.}

\prose{อิติ โข ภทฺทิย ยํ ตํ เต อโวจุมฺหา “เอถ ตุเมฺห ภทฺทิย, มา อนุสฺสเวน, มา ปรมฺปราย, มา อิติกิราย, มา ปิฏก\-สมฺปทาเนน, มา ตกฺกเหตุ, มา นยเหตุ, มา อาการ\-ปริวิตกฺเกน, มา ทิฏฺฐิ\-นิชฺฌาน\-กฺขนฺติยา, มา ภพฺพรูปตาย, มา สมโณ โน ครูติ. ยทา ตุเมฺห ภทฺทิย อตฺตนาว}

\vspace{\tipitakaparskip}
\csromancenter{(20) 5. มหาวคฺค ชาเนยฺยาถ ‘อิเม ธมฺมา อกุสลา, อิเม ธมฺมา สาวชฺชา, อิเม ธมฺมา วิญฺญุครหิตา, อิเม ธมฺมา สมตฺตา สมาทินฺนา อหิตาย ทุกฺขาย สํวตฺตนฺตี’ติ, อถ ตุเมฺห ภทฺทิย ปชเหยฺยาถา”ติ อิติ ยํ ตํ วุตฺตํ, อิทเมตํ ปฏิจฺจ วุตฺตํ.}
\prose{\csromanfolio{513}เอถ ตุเมฺห ภทฺทิย, มา อนุสฺสเวน มา ปรมฺปราย มา อิติกิราย มา ปิฏก\-สมฺปทาเนน มา ตกฺกเหตุ มา นยเหตุ มา อาการ\-ปริวิตกฺเกน มา ทิฏฺฐิ\-นิชฺฌาน\-กฺขนฺติยา มา ภพฺพรูปตาย มา สมโณ โน ครูติ. ยทา ตุเมฺห ภทฺทิย อตฺตนาว ชาเนยฺยาถ “อิเม ธมฺมา กุสลา, อิเม ธมฺมา อนวชฺชา อิเม ธมฺมา วิญฺญุปฺปสตฺถา, อิเม ธมฺมา สมตฺตา สมาทินฺนา หิตาย สุขาย สํวตฺตนฺตี”ติ. อถ ตุเมฺห ภทฺทิย อุปสมฺปชฺช วิหเรยฺยาถาติ.}

\prose{ตํ กิํ มญฺญถ ภทฺทิย, อโลโภ ปุริสสฺส อชฺฌตฺตํ อุปฺปชฺชมาโน อุปฺปชฺชติ หิตาย วา อหิตาย วาติ. หิตาย ภนฺเต. อลุทฺโธ ปนายํ ภทฺทิย ปุริสปุคฺคโล โลเภน อนภิภูโต อปริยาทินฺนจิตฺโต เนว ปาณํ หนติ, น อทินฺนํ อาทิยติ, น ปรทารํ คจฺฉติ, น มุสา ภณติ, ปรมฺปิ ตถตฺตาย น สมาทเปติ, ยํ’ส โหติ ทีฆรตฺตํ หิตาย สุขายาติ. เอวํ ภนฺเต.}

\prose{ตํ กิํ มญฺญถ ภทฺทิย, อโทโส ปุริสสฺส ฯเปฯ อโมโห ปุริสสฺส ฯเปฯ อสารมฺโภ ปุริสสฺส อชฺฌตฺตํ อุปฺปชฺชมาโน อุปฺปชฺชติ หิตาย วา อหิตาย วาติ. หิตาย ภนฺเต. อสารทฺโธ ปนายํ ภทฺทิย ปุริสปุคฺคโล สารมฺเภน อนภิภูโต อปริยาทินฺนจิตฺโต เนว ปาณํ หนติ, น อทินฺนํ อาทิยติ, น ปรทารํ คจฺฉติ. น มุสา ภณติ, ปรมฺปิ ตถตฺตาย น สมาทเปติ, ยํ’ส โหติ ทีฆรตฺตํ หิตาย สุขายาติ. เอวํ ภนฺเต.}

\prose{ตํ กิํ มญฺญถ ภทฺทิย, อิเม ธมฺมา กุสลา วา อกุสลา วาติ. กุสลา ภนฺเต. สาวชฺชา วา อนวชฺชา วาติ. อนวชฺชา ภนฺเต. วิญฺญุครหิตา วา วิญฺญุปฺปสตฺถา วาติ. วิญฺญุปฺปสตฺถา ภนฺเต. สมตฺตา สมาทินฺนา หิตาย สุขาย สํวตฺตนฺติ โน วา, กถํ วา เอตฺถ โหตีติ. สมตฺตา ภนฺเต สมาทินฺนา หิตาย สุขาย สํวตฺตนฺติ, เอวํ โน เอตฺถ โหตีติ.}

\prose{\csromanfolio{514}อิติ โข ภทฺทิย ยํ ตํ เต อโวจุมฺหา “เอตฺถ ตุเมฺห ภทฺทิย, มา อนุสฺสเวน, มา ปรมฺปราย, มา อิติกิราย, มา ปิฏก\-สมฺปทาเนน, มา ตกฺกเหตุ, มา นยเหตุ, มา อาการ\-ปริวิตกฺเกน, มา ทิฏฺฐิ\-นิชฺฌาน\-กฺขนฺติยา, มา ภพฺพรูปตาย, มา สมโณ โน ครู”ติ. ยทา ตุเมฺห ภทฺทิย อตฺตนาว ชาเนยฺยาถ “อิเม ธมฺมา กุสลา, อิเม ธมฺมา อนวชฺชา, อิเม ธมฺมา วิญฺญุปฺปสตฺถา, อิเม ธมฺมา สมตฺตา สมาทินฺนา หิตาย สุขาย สํวตฺตนฺตี”ติ. อถ ตุเมฺห ภทฺทิย อุปสมฺปชฺช วิหเรยฺยาถาติ, อิติ ยํ ตํ วุตฺตํ, อิทเมตํ ปฏิจฺจ วุตฺตํ.}

\prose{เย โข เต ภทฺทิย โลเก สนฺโต สปฺปุริสา, เต สาวกํ เอวํ สมาทเปนฺติ เอหิ ตฺวํ อมฺโภ ปุริส โลภํ วิเนยฺย\csromansharedfootnote{21d90a0e5df99b09}{วิเนยฺย วิเนยฺย (สี, สฺยา, กํ)} วิหราหิ, โลภํ วิเนยฺย\csromansharedfootnote{21d90a0e5df99b09}{วิเนยฺย วิเนยฺย (สี, สฺยา, กํ)} วิหรนฺโต น โลภชํ กมฺมํ กริสฺสสิ กาเยน วาจาย มนสา. โทสํ วิเนยฺย วิหราหิ, โทสํ วิเนยฺย วิหรนฺโต น โทสชํ กมฺมํ กริสฺสสิ กาเยน วาจาย มนสา. โมหํ วิเนยฺย วิหราหิ, โมหํ วิเนยฺย วิหรนฺโต น โมหชํ กมฺมํ กริสฺสสิ กาเยน วาจาย มนสา. สารมฺภํ วิเนยฺย วิหราหิ, สารมฺภํ วิเนยฺย วิหรนฺโต น สารมฺภชํ กมฺมํ กริสฺสสิ กาเยน วาจาย มนสาติ.}

\prose{เอวํ วุตฺเต ภทฺทิโย ลิจฺฉวิ ภควนฺตํ เอตทโวจ “อภิกฺกนฺตํ ภนฺเต ฯเปฯ อุปาสกํ มํ ภนฺเต ภควา ธาเรตุ อชฺชตคฺเค ปาณุเปตํ สรณํ คตนฺ”ติ.}

\prose{อปิ นุ ตาหํ ภทฺทิย เอวํ อวจํ “เอหิ เม ตฺวํ ภทฺทิย สาวโก โหติ, อหํ สตฺถา ภวิสฺสามี”ติ. โน เหตํ ภนฺเต. เอวํวาทิํ โข มํ ภทฺทิย เอวมกฺขายิํ “เอเก สมณพฺราหฺมณา อสตา ตุจฺฉา มุสา อภูเตน อพฺภาจิกฺขนฺติ มายาวี สมโณ โคตโม อาวฏฺฏนิํ มายํ ชานาติ, ยาย อญฺญ\-ติตฺถิยานํ สาวเก อาวฏฺเฏตี”ติ. ภทฺทิกา ภนฺเต อาวฏฺฏนี มายา, กลฺยาณี ภนฺเต อาวฏฺฏนี มายา. ปิยา เม ภนฺเต ญาติสาโลหิตา อิมาย อาวฏฺฏนิยา อาวฏฺเฏยฺยุํ, ปิยานมฺปิ เม อสฺส ญาติ\-สาโลหิตานํ ทีฆรตฺตํ หิตาย สุขาย. สพฺเพ เจปิ ภนฺเต ขตฺติยา อิมาย อาวฏฺฏนิยา อาวฏฺเฏยฺยุํ, สพฺเพสมฺปิสฺส ขตฺติยานํ ทีฆรตฺตํ หิตาย สุขาย. สพฺเพ เจปิ ภนฺเต พฺราหฺมณา.}

\vspace{\tipitakaparskip}
\csromancenter{(20) 5. มหาวคฺค เวสฺสา. สุทฺทา อิมาย อาวฏฺฏนิยา อาวฏฺเฏยฺยุํ, สพฺเพสมฺปิสฺส สุทฺทานํ ทีฆรตฺตํ หิตาย สุขายาติ.}
\prose{\csromanfolio{515}เอวเมตํ ภทฺทิย เอวเมตํ ภทฺทิย, สพฺเพ เจปิ ภทฺทิย ขตฺติยา อิมาย อาวฏฺฏนิยา อาวฏฺเฏยฺยุํ อกุสลธมฺม\-ปฺปหานาย กุสลธมฺมู\-ปสมฺปทาย, สพฺเพสมฺปิสฺส ขตฺติยานํ ทีฆรตฺตํ หิตาย สุขาย. สพฺเพ เจปิ ภทฺทิย พฺราหฺมณา. เวสฺสา. สุทฺทา อาวฏฺเฏยฺยุํ อกุสลธมฺม\-ปฺปหานาย กุสลธมฺมู\-ปสมฺปทาย, สพฺเพสมฺปิสฺส สุทฺทานํ ทีฆรตฺตํ หิตาย สุขาย. สเทวโก เจปิ ภทฺทิย โลโก สมารโก สพฺรหฺมโก สสฺสมณ\-พฺราหฺมณี ปชา สเทวมนุสฺสา อิมาย อาวฏฺฏนิยา อาวฏฺเฏยฺยุํ\csromansharedfootnote{68c3fe7e08d3b07e}{อาวฏฺเฏยฺย (?)} อกุสลธมฺม\-ปฺปหานาย กุสลธมฺมู\-ปสมฺปทาย, สเทวกสฺสปิสฺส โลกสฺส สมารกสฺส สพฺรหฺมกสฺส สสฺสมณ\-พฺราหฺมณิยา ปชาย สเทว\-มนุสฺสาย ทีฆรตฺตํ หิตาย สุขาย. อิเม เจปิ ภทฺทิย มหาสาลา อิมาย อาวฏฺฏนิยา อาวฏฺเฏยฺยุํ อกุสลธมฺม\-ปฺปหานาย กุสลธมฺมู\-ปสมฺปทาย, อิเมสมฺปิสฺส มหาสาลานํ ทีฆรตฺตํ หิตาย สุขาย ( )\csromansharedfootnote{df35a417f5c9145e}{(สเจ เจเตยฺยุํ) (สี, สฺยา, กํ, อิ), (อาวฏฺเฏยฺยุํ) (ก) อํ 3. 88 ปิฏฺเฐปิ.}, โก ปน วาโท มนุสฺสภูตสฺสาติ.  ตติยํ.}

\csromanmatikaanchor{mtk.233}
\csromantocmarkref{subsection}{4. สามุคิยสุตฺต}{mtk.233}
\bookmark[dest={mtk.233},level=2]{4. สามุคิยสุตฺต}
\csromanheader{2}{4. \textbf{สามุคิยสุตฺต}}
\proseitem{194}{เอกํ สมยํ อายสฺมา อานนฺโท โกลิเยสุ วิหรติ สามุคํ นาม\csromansharedfootnote{cf68559f4e36050b}{สาปูคํ นาม (สี, สฺยา, กํ, อิ)} โกลิยานํ นิคโม, อถ โข สมฺพหุลา สามุคิยา โกลิยปุตฺตา เยนายสฺมา อานนฺโท เตนุ\-ปสงฺกมิํสุ, อุปสงฺกมิตฺวา อายสฺมนฺตํ อานนฺทํ อภิวาเทตฺวา เอกมนฺตํ นิสีทิํสุ, เอกมนฺตํ นิสินฺเน โข เต สามุคิเย โกลิยปุตฺเต อายสฺมา อานนฺโท เอตทโวจ\csromandash{}}

\prose{จตฺตาริมานิ พฺยคฺฆปชฺชา ปาริสุทฺธิ\-ปธานิย\-งฺคานิ เตน ภควตา ชานตา ปสฺสตา อรหตา สมฺมา\-สมฺพุทฺเธน สมฺมทกฺขาตานิ สตฺตานํ วิสุทฺธิยา โสก\-ปริเทวานํ\csromansharedfootnote{808c395c531b3b97}{โสกปริทฺทวานํ (สี)} สมติกฺกมาย ทุกฺข\-โทมนสฺสานํ อตฺถงฺคมาย ญายสฺส อธิคมาย นิพฺพานสฺส สจฺฉิกิริยาย. กตมานิ จตฺตาริ? สีลปาริสุทฺธิ\-ปธานิยงฺคํ จิตฺต\-ปาริสุทฺธิ\-ปธานิยงฺคํ ทิฏฺฐิ\-ปาริสุทฺธิ\-ปธานิยงฺคํ วิมุตฺติ\-ปาริสุทฺธิ\-ปธานิยงฺคํ.}

\prose{\csromanfolio{516}กตมญฺจ พฺยคฺฆปชฺชา สีลปาริสุทฺธิ\-ปธานิยงฺคํ? อิธ พฺยคฺฆปชฺชา ภิกฺขุ สีลวา โหติ ฯเปฯ สมาทาย สิกฺขติ สิกฺขาปเทสุ. อยํ วุจฺจติ พฺยคฺฆปชฺชา สีลปาริสุทฺธิ. “อิติ เอวรูปิํ สีลปาริสุทฺธิํ อปริปูรํ วา ปริปูเรสฺสามิ, ปริปูรํ วา ตตฺถ ตตฺถ ปญฺญาย อนุคฺคเหสฺสามี”ติ โย ตตฺถ ฉนฺโท จ วายาโม จ อุสฺสาโห จ อุสฺโสฬฺหี จ อปฺปฏิวานี จ สติ จ สมฺปชญฺญญฺจ. อิทํ วุจฺจติ พฺยคฺฆปชฺชา สีลปาริสุทฺธิ\-ปธานิยงฺคํ.}

\prose{กตมญฺจ พฺยคฺฆปชฺชา จิตฺต\-ปาริสุทฺธิ\-ปธานิยงฺคํ? อิธ พฺยคฺฆปชฺชา ภิกฺขุ วิวิจฺเจว กาเมหิ ฯเปฯ จตุตฺถํ ฌานํ อุปสมฺปชฺช วิหรติ. อยํ วุจฺจติ พฺยคฺฆปชฺชา จิตฺต\-ปาริสุทฺธิ. “อิติ เอวรูปิํ จิตฺต\-ปาริสุทฺธิํ อปริปูรํ วา ปริปูเรสฺสามิ, ปริปูรํ วา ตตฺถ ตตฺถ ปญฺญาย อนุคฺคเหสฺสามี”ติ โย ตตฺถ ฉนฺโท จ วายาโม จ อุสฺสาโห จ อุสฺโสฬฺหี จ อปฺปฏิวานี จ สติ จ สมฺปชญฺญญฺจ. อิทํ วุจฺจติ พฺยคฺฆปชฺชา จิตฺต\-ปาริสุทฺธิ\-ปธานิยงฺคํ.}

\prose{กตมญฺจ พฺยคฺฆปชฺชา ทิฏฺฐิ\-ปาริสุทฺธิ\-ปธานิยงฺคํ? อิธ พฺยคฺฆปชฺชา ภิกฺขุ “อิทํ ทุกฺขนฺ”ติ ยถาภูตํ ปชานาติ ฯเปฯ “อยํ ทุกฺขนิโรธ\-คามินี ปฏิปทา”ติ ยถาภูตํ ปชานาติ. อยํ วุจฺจติ พฺยคฺฆปชฺชา ทิฏฺฐิ\-ปาริสุทฺธิ. “อิติ เอวรูปิํ ทิฏฺฐิ\-ปาริสุทฺธิํ อปริปูรํ วา ฯเปฯ ตตฺถ ตตฺถ ปญฺญาย อนุคฺคเหสฺสามี”ติ โย ตตฺถ ฉนฺโท จ วายาโม จ อุสฺสาโห จ อุสฺโสฬฺหี จ อปฺปฏิวานี จ สติ จ สมฺปชญฺญญฺจ. อิทํ วุจฺจติ พฺยคฺฆปชฺชา ทิฏฺฐิ\-ปาริสุทฺธิ\-ปธานิยงฺคํ.}

\prose{กตมญฺจ พฺยคฺฆปชฺชา วิมุตฺติปาริสุกฺขิปธานิยงฺคํ? ส โข โส พฺยคฺฆปชฺชา อริยสาวโก อิมินา จ สีลปาริสุทฺธิ\-ปธานิย\-งฺเคน สมนฺนาคโต อิมินา จ จิตฺต\-ปาริสุทฺธิ\-ปธานิย\-งฺเคน สมนฺนาคโต อิมินา จ ทิฏฺฐิ\-ปาริสุทฺธิ\-ปธานิย\-งฺเคน สมนฺนาคโต รชนีเยสุ ธมฺเมสุ จิตฺตํ วิราเชติ, วิโมจนีเยสุ ธมฺเมสุ จิตฺตํ วิโมเจติ. โส รชนีเยสุ ธมฺเมสุ จิตฺตํ วิราเชตฺวา วิโมจนีเยสุ ธมฺเมสุ จิตฺตํ วิโมเจตฺวา สมฺมาวิมุตฺติํ ผุสติ. อยํ วุจฺจติ พฺยคฺฆปชฺชา วิมุตฺติ\-ปาริสุทฺธิ. “อิติ เอวรูปิํ วิมุตฺติ\-ปาริสุทฺธิํ อปริปูรํ วา ปริปูเรสฺสามิ, ปริปูรํ วา ตตฺถ ตตฺถ ปญฺญาย อนุคฺคเหสฺสามี”ติ โย ตตฺถ ฉนฺโท จ วายาโม จ อุสฺสาโห จ อุสฺโสฬฺหี จ อปฺปฏิวานี จ สติ จ สมฺปชญฺญญฺจ. อิทํ วุจฺจติ พฺยคฺฆปชฺชา วิมุตฺติ\-ปาริสุทฺธิ\-ปธานิยงฺคํ.}

\chapterheadpageread{(20) 5. มหาวคฺค}
\prose{\csromanfolio{517}อิมานิ โข พฺยคฺฆปชฺชา จตฺตาริ ปาริสุทฺธิ\-ปธานิย\-งฺคานิ เตน ภควตา ชานตา ปสฺสตา อรหตา สมฺมา\-สมฺพุทฺเธน สมฺมทกฺขาตานิ สตฺตานํ วิสุทฺธิยา โสก\-ปริเทวานํ สมติกฺกมาย ทุกฺข\-โทมนสฺสานํ อตฺถงฺคมาย ญายสฺส อธิคมาย นิพฺพานสฺส สจฺฉิกิริยายาติ.  จตุตฺถํ.}

\csromanmatikaanchor{mtk.234}
\csromantocmarkref{subsection}{5. วปฺปสุตฺต}{mtk.234}
\bookmark[dest={mtk.234},level=2]{5. วปฺปสุตฺต}
\csromanheader{2}{5. \textbf{วปฺปสุตฺต}}
\proseitem{195}{เอกํ สมยํ ภควา สกฺเกสุ วิหรติ กปิล\-วตฺถุสฺมิํ นิโคฺรธาราเม. อถ โข วปฺโป สกฺโก นิคณฺฐสาวโก เยนายสฺมา มหาโมคฺคลฺลาโน เตนุปสงฺกมิ, อุปสงฺกมิตฺวา อายสฺมนฺตํ มหา\-โมคฺคลฺลานํ อภิวาเทตฺวา เอกมนฺตํ นิสีทิ. เอกมนฺตํ นิสินฺนํ โข วปฺปํ สกฺกํ นิคณฺฐ\-สาวกํ อายสฺมา มหาโมคฺคลฺลาโน เอตทโวจ\csromandash{}}

\prose{อิธสฺส วปฺป กาเยน สํวุโต วาจาย สํวุโต มนสา สํวุโต อวิชฺชาวิราคา วิชฺชุปฺปาทา ปสฺสสิ โน ตฺวํ วปฺป ตํ ฐานํ, ยโตนิทานํ ปุริสํ ทุกฺขเวทนิยา อาสวา อสฺสเวยฺยุํ\csromansharedfootnote{7624d5d772aa834a}{อนฺวาสฺสเวยฺยุํ (ก)} อภิสมฺปรายนฺติ. ปสฺสามหํ ภนฺเต ตํ ฐานํ อิธสฺส ภนฺเต ปุพฺเพ ปาปกมฺมํ กตํ อวิปกฺก\-วิปากํ, ตโตนิทานํ ปุริสํ ทุกฺขเวทนิยา อาสวา อสฺสเวยฺยุํ อภิสมฺปรายนฺติ. อยญฺเจว โข ปน อายสฺมโต มหา\-โมคฺคลฺลานสฺส วปฺเปน สกฺเกน นิคณฺฐ\-สาวเกน สทฺธิํ อนฺตรากถา วิปฺปกตา โหติ.}

\prose{อถ โข ภควา สายนฺหสมยํ ปฏิสลฺลานา วุฏฺฐิโต เยน อุปฏฺฐานสาลา เตนุปสงฺกมิ, อุปสงฺกมิตฺวา ปญฺญตฺเต อาสเน นิสีทิ, นิสชฺช โข ภควา อายสฺมนฺตํ มหา\-โมคฺคลฺลานํ เอตทโวจ\csromandash{}}

\prose{กาย นุตฺถ โมคฺคลฺลาน เอตรหิ กถาย สนฺนิสินฺนา, กา จ ปน โว อนฺตรากถา วิปฺปกตาติ. อิธาหํ ภนฺเต วปฺปํ สกฺกํ นิคณฺฐ\-สาวกํ เอตทโวจํ “อิธสฺส วปฺป กาเยน สํวุโต วาจาย สํวุโต มนสา สํวุโต อวิชฺชาวิราคา วิชฺชุปฺปาทา ปสฺสสิ โน ตฺวํ วปฺป ตํ ฐานํ, ยโตนิทานํ ปุริสํ ทุกฺขเวทนิยา อาสวา อสฺสเวยฺยุํ อภิสมฺปรายนฺ”ติ. เอวํ วุตฺเต ภนฺเต วปฺโป สกฺโก นิคณฺฐสาวโก มํ เอตทโวจ “ปสฺสามหํ ภนฺเต ตํ ฐานํ อิธสฺส ภนฺเต ปุพฺเพ ปาปกมฺมํ กตํ อวิปกฺก\-วิปากํ, ตโตนิทานํ ปุริสํ ทุกฺขเวทนิยา อาสวา อสฺสเวยฺยุํ \csromanfolio{518}อภิสมฺปรายนฺ”ติ. อยํ โข โน ภนฺเต วปฺเปน สกฺเกน นิคณฺฐ\-สาวเกน สทฺธิํ อนฺตรากถา วิปฺปกตา, อถ ภควา อนุปฺปตฺโตติ.}

\prose{อถ โข ภควา วปฺปํ สกฺกํ นฺคณฺฐสาวกํ เอตทโวจ “สเจ เม ตฺวํ วปฺป อนุญฺเญยฺยญฺเจว อนุชาเนยฺยาสิ, ปฏิกฺโกสิตพฺพญฺจ ปฏิกฺโกเสยฺยาสิ, ยสฺส จ เม ภาสิตสฺส อตฺถํ น ชาเนยฺยาสิ, มเมเวตฺถ อุตฺตริ ปฏิปุจฺเฉยฺยาสิ ‘อิทํ ภนฺเต กถํ, อิมสฺส โก อตฺโถ’ติ, สิยา โน เอตฺถ กถาสลฺลาโป”ติ. อนุญฺเญยฺยญฺเจวาหํ ภนฺเต ภควโต อนุชานิสฺสามิ, ปฏิกฺโกสิตพฺพญฺจ ปฏิกฺโกสิสฺสามิ, ยสฺส จาหํ ภควโต ภาสิตสฺส อตฺถํ น ชานิสฺสามิ, ภควนฺตํเยเวตฺถ อุตฺตริ ปฏิปุจฺฉิสฺสามิ “อิทํ ภนฺเต กถํ, อิมสฺส โก อตฺโถ”ติ, โหตุ โน เอตฺถ กถาสลฺลาโปติ.}

\prose{ตํ กิํ มญฺญสิ วปฺป, เย กาย\-สมารมฺภ\-ปจฺจยา อุปฺปชฺชนฺติ อาสวา วิฆาต\-ปริฬาหา, กายสมารมฺภา ปฏิวิรตสฺส เอวํส เต อาสวา วิฆาต\-ปริฬาหา น โหนฺติ. โส นวญฺจ กมฺมํ น กโรติ, ปุราณญฺจ กมฺมํ ผุสฺส ผุสฺส พฺยนฺตีกโรติ สนฺทิฏฺฐิกา นิชฺชรา อกาลิกา เอหิปสฺสิกา โอปเนยฺยิกา ปจฺจตฺตํ เวทิตพฺพา วิญฺญูหิ\csromansharedfootnote{ea6f8efcce6ac777}{วิญฺญูหีติ (สี, อิ, ก) สํ 2. 520 ปิฏฺเฐ ปสฺสิตพฺพํ.}. ปสฺสสิ โน ตฺวํ วปฺป ตํ ฐานํ, ยโตนิทานํ ปุริสํ ทุกฺขเวทนิยา อาสวา อสฺสเวยฺยุํ อภิสมฺปรายนฺติ. โน เหตํ ภนฺเต.}

\prose{ตํ กิํ มญฺญสิ วปฺป, เย วจีสมารมฺภ\-ปจฺจยา อุปฺปชฺชนฺติ อาสวา วิฆาต\-ปริฬาหา, วจีสมารมฺภา ปฏิวิรตสฺส เอวํส เต อาสวา วิฆาต\-ปริฬาหา น โหนฺติ. โส นวญฺจ กมฺมํ น กโรติ, ปุราณญฺจ กมฺมํ ผุสฺส ผุสฺส พฺยนฺตีกโรติ สนฺทิฏฺฐิกา นิชฺชรา อกาลิกา เอหิปสฺสิกา โอปเนยฺยิกา ปจฺจตฺตํ เวทิตพฺพา วิญฺญูหิ. ปสฺสสิ โน ตฺวํ วปฺป ตํ ฐานํ, ยโตนิทานํ ปุริสํ ทุกฺขเวทนิยา อาสวา อสฺสเวยฺยุํ อภิสมฺปรายนฺติ. โน เหตํ ภนฺเต.}

\prose{ตํ กิํ มญฺญสิ วปฺป, เย มโน\-สมารมฺภ\-ปจฺจยา อุปฺปชฺชนฺติ อาสวา วิฆาต\-ปริฬาหา, มโนสมารมฺภา ปฏิวิรตสฺส เอวํส เต อาสวา วิฆาต\-ปริฬาหา น โหนฺติ. โส นวญฺจ กมฺมํ น กโรติ, ปุราณญฺจ กมฺมํ ผุสฺส ผุสฺส พฺยนฺตีกโรติ สนฺทิฏฺฐิกา นิชฺชรา อกาลิกา เอหิปสฺสิกา}

\vspace{\tipitakaparskip}
\csromancenter{(20) 5. มหาวคฺค โอปเนยฺยิกา ปจฺจตฺตํ เวทิตพฺพา วิญฺญูหิ. ปสฺสสิ โน ตฺวํ วปฺป ตํ ฐานํ, ยโตนิทานํ ปุริสํ ทุกฺขเวทนิยา อาสวา อสฺสเวยฺยุํ อภิสมฺปรายนฺติ. โน เหตํ ภนฺเต.}
\prose{\csromanfolio{519}ตํ กิํ มญฺญสิ วปฺป เย อวิชฺชาปจฺจยา อุปฺปชฺชนฺติ อาสวา วิฆาตปริฬฺหาหา, อวิชฺชาวิราคา วิชฺชุปฺปาทา เอวํส เต อาสวา วิฆาต\-ปริฬาหา น โหนฺติ. โส นวญฺจ กมฺมํ น กโรติ, ปุราณญฺจ กมฺมํ ผุสฺส ผุสฺส พฺยนฺตีกโรติ สนฺทิฏฺฐิกา นิชฺชรา อกาลิกา เอหิปสฺสิกา โอปเนยฺยิกา ปจฺจตฺตํ เวทิตพฺพา วิญฺญูหิ. ปสฺสสิ โน ตฺวํ วปฺป ตํ ฐานํ, ยโตนิทานํ ปุริสํ ทุกฺขเวทนิยา อาสวา อสฺสเวยฺยุํ อภิสมฺปรายนฺติ. โน เหตํ ภนฺเต.}

\prose{เอวํ สมฺมา วิมุตฺต\-จิตฺตสฺส โข วปฺป ภิกฺขุโน ฉ สตตวิหารา อธิคตา โหนฺติ. โส จกฺขุนา รูปํ ทิสฺวา เนว สุมโน โหติ น ทุมฺมโน, อุเปกฺขโก วิหรติ สโต สมฺปชาโน. โสเตน สทฺทํ สุตฺวา ฯเปฯ. ฆาเนน คนฺธํ ฆายิตฺวา ฯเปฯ. ชิวฺหาย รสํ สายิตฺวา ฯเปฯ. กาเยน โผฏฺฐพฺพํ ผุสิตฺวา ฯเปฯ. มนสา ธมฺมํ วิญฺญาย เนว สุมโน โหติ น ทุมฺมโน, อุเปกฺขโก วิหรติ สโต สมฺปชาโน, โส กาย\-ปริยนฺติกํ เวทนํ เวทิยมาโน “กาย\-ปริยนฺติกํ เวทนํ เวทิยามี”ติ ปชานาติ, ชีวิต\-ปริยนฺติกํ เวทนํ เวทิยมาโน “ชีวิต\-ปริยนฺติกํ เวทนํ เวทิยามี”ติ ปชานาติ, “กายสฺส เภทา อุทฺธํ ชีวิต\-ปริยาทานา อิเธว สพฺพ\-เวทยิตานิ อนภินนฺทิตานิ สีตี ภวิสฺสนฺตี”ติ ปชานาติ.}

\prose{เสยฺยถาปิ วปฺป ถูณํ ปฏิจฺจ ฉายา ปญฺญายติ, อถ ปุริโส อาคจฺเฉยฺย กุทาลปิฏกํ อาทาย. โส ตํ ถูณํ มูเล ฉินฺเทยฺย, มูเล ฉินฺทิตฺวา ปลิขเณยฺย, ปลิขณิตฺวา มูลานิ อุทฺธเรยฺย อนฺตมโส อุสีร\-นาฬิ\-มตฺตานิปิ\csromansharedfootnote{6609d2c0632db1a2}{อุสีรนาฬมตฺตานิปิ (สี)}. โส ตํ ถูณํ ขณฺฑาขณฺฑิกํ ฉินฺเทยฺย, ขณฺฑาขณฺฑิกํ เฉตฺวา ผาเลยฺย, ผาเลตฺวา สกลิกํ สกลิกํ กเรยฺย, สกลิกํ สกลิกํ กตฺวา วาตาตเป วิโสเสยฺย, วาตาตเป วิโสเสตฺวา อคฺคินา ฑเหยฺย, อคฺคินา ฑเหตฺวา มสิํ กเรยฺย, มสิํ กริตฺวา มหาวาเต วา โอผุเณยฺย, นทิยา วา สีฆโสตาย ปวาเหยฺย, \csromanfolio{520}เอวํ หิสฺส วปฺป ยา ถูณํ ปฏิจฺจ ฉายา, สา อุจฺฉินฺนมูลา ตาลาวตฺถุกตา อนภาวํกตา อายติํ อนุปฺปาทมฺมา.}

\prose{เอวเมวํ โข วปฺป เอวํ สมฺมา วิมุตฺต\-จิตฺตสฺส ภิกฺขุโน ฉ สตตวิหารา อธิคตา โหนฺติ, โส จกฺขุนา รูปํ ทิสฺวา เนว สุมโน โหติ น ทุมฺมโน, อุเปกฺขโก วิหรติ สโต สมฺปชาโน. โสเตน สทฺทํ สุตฺวา ฯเปฯ. ฆาเนน คนฺธํ ฆายิตฺวา ฯเปฯ. ชิวฺหาย รสํ สายิตฺวา ฯเปฯ. กาเยน โผฏฺฐพฺพํ ผุสิตฺวา ฯเปฯ. มนสา ธมฺมํ วิญฺญาย เนว สุมโน โหติ น ทุมฺมโน, อุเปกฺขโก วิหรติ สโต สมฺปชาโน. โส กาย\-ปริยนฺติกํ เวทนํ เวทิยมาโน “กาย\-ปริยนฺติกํ เวทนํ เวทิยามี”ติ ปชานาติ. ชีวิต\-ปริยนฺติกํ เวทนํ เวทิยมาโน “ชีวิต\-ปริยนฺติกํ เวทนํ เวทิยามี”ติ ปชานาติ, “กายสฺส เภทา อุทฺธํ ชีวิต\-ปริยาทานา อิเธว สพฺพ\-เวทยิตานิ อนภินนฺทิตานิ สีตี ภวิสฺสนฺตี”ติ ปชานาติ.}

\prose{เอวํ วุตฺเต วปฺโป สกฺโก นิคณฺฐสาวโก ภควนฺตํ เอตทโวจ “เสยฺยถาปิ ภนฺเต ปุริโส อุทยตฺถิโก อสฺสปณิยํ โปเสยฺย, โส อุทยญฺเจว นาธิคจฺเฉยฺย, อุตฺตริญฺจ กิลมถสฺส วิฆาตสฺส ภาคี อสฺส. เอวเมวํ โข อหํ ภนฺเต อุทยตฺถิโก พาเล นฺคณฺเฐ ปยิรุปาสิํ, สฺวาหํ อุทยญฺเจว นาธิคจฺฉิํ, อุตฺตริญฺจ กิลมถสฺส วิฆาตสฺส ภาคี อโหสิํ. เอสาหํ ภนฺเต อชฺชตคฺเค โย เม พาเลสุ นฺคณฺเฐสุ ปสาโท, ตํ มหาวาเต วา โอผุณามิ, นทิยา วา สีฆโสตาย ปวาเหมิ. อภิกฺกนฺตํ ภนฺเต ฯเปฯ อุปาสกํ มํ ภนฺเต ภควา ธาเรตุ อชฺชตคฺเค ปาณุเปตํ สรณํ คตนฺติ.  ปญฺจมํ.}

\csromanmatikaanchor{mtk.235}
\csromantocmarkref{subsection}{6. สาฬฺหสุตฺต}{mtk.235}
\bookmark[dest={mtk.235},level=2]{6. สาฬฺหสุตฺต}
\csromanheader{2}{6. \textbf{สาฬฺหสุตฺต}}
\proseitem{196}{เอกํ สมยํ ภควา เวสาลิยํ วิหรติ มหาวเน กูฏาคาร\-สาลายํ. อถ โข สาโฬฺห จ ลิจฺฉวิ, อภโย จ ลิจฺฉวิ เยน ภควา เตนุ\-ปสงฺกมิํสุ, อุปสงฺกมิตฺวา ภควนฺตํ อภิวาเทตฺวา เอกมนฺตํ นิสีทิํสุ, เอกมนฺตํ นิสินฺโน โข สาโฬฺห ลิจฺฉวิ ภควนฺตํ เอตทโวจ\csromandash{}}

\prose{“สนฺติ ภนฺเต เอเก สมณพฺราหฺมณา ทฺวเยน โอฆสฺส นิตฺถรณํ ปญฺญเปนฺติ สีลวิสุทฺธิ\-เหตุ จ ตโป\-ชิคุจฺฉา\-เหตุ จ, อิธ ภนฺเต ภควา กิมาหา”ติ.}

\chapterheadpageread{(20) 5. มหาวคฺค}
\prose{\csromanfolio{521}สีลวิสุทฺธิํ โข อหํ สาฬฺห “อญฺญตรํ สามญฺญงฺคนฺ”ติ วทามิ. เย เต สาฬฺห สมณพฺราหฺมณา ตโปชิ คุจฺฉาวาทา ตโปชิคุจฺฉา\-สารา ตโปชิคุจฺฉา\-อลฺลีนา วิหรนฺติ, อภพฺพา เต โอฆสฺส นิตฺถรณาย. เยปิ เต สาฬฺห สมณพฺราหฺมณา อปริสุทฺธ\-กายสมาจารา อปริสุทฺธ\-วจีสมาจารา อปริสุทฺธ\-มโนสมาจารา อปริสุทฺธา\-ชีวา, อภพฺพา เต ญาณทสฺสนาย อนุตฺตราย สมฺโพธาย.}

\prose{เสยฺยถาปิ สาฬฺห ปุริโส นทิํ ตริตุกาโม ติณฺหํ กุฐาริํ\csromansharedfootnote{ff32430e52b59506}{กุธาริํ (ก)} อาทาย วนํ ปวิเสยฺย, โส ตตฺถ ปสฺเสยฺย มหติํ สาลลฏฺฐิํ อุชุํ นวํ อกุกฺกุจฺจกชาตํ. ตเมนํ มูเล ฉินฺเทยฺย, มูเล เฉตฺวา อคฺเค ฉินฺเทยฺย, อคฺเค เฉตฺวา สาขาปลาสํ สุวิโสธิตํ วิโสเธยฺย, สาขาปลาสํ สุวิโสธิตํ วิโสเธตฺวา กุฐารีหิ ตจฺเฉยฺย, กุฐารีหิ ตจฺเฉตฺวา วาสีหิ ตจฺเฉยฺย, วาสีหิ ตจฺเฉตฺวา เลขณิยา ลิเขยฺย, เลขณิยา ลิขิตฺวา ปาสาณคุเฬน โธเวยฺย\csromansharedfootnote{0ce37e6a8c61676d}{โธเปยฺย (สี, สฺยา, กํ, อิ)}, ปาสาณคุเฬน โธเวตฺวา นทิํ ปตาเรยฺย.}

\prose{ตํ กิํ มญฺญสิ สาฬฺห, ภพฺโพ นุ โข โส ปุริโส นทิํ ตริตุนฺติ. โน เหตํ ภนฺเต. ตํ กิสฺส เหตุ? อสุ หิ ภนฺเต สาลลฏฺฐิ พหิทฺธา สุปริกมฺม\-กตา อนฺโต อวิสุทฺธา, ตสฺเสตํ ปาฏิกงฺขํ “สาลลฏฺฐิ สํสีทิสฺสติ, ปุริโส อนยพฺยสนํ อาปชฺชิสฺสตี”ติ.}

\prose{เอวเมวํ โข สาฬฺห เย เต สมณพฺราหฺมณา ตโปชิคุจฺฉา\-วาทา ตโปชิคุจฺฉา\-สารา ตโปชิคุจฺฉา\-อลฺลีนา วิหรนฺติ, อภพฺพา เต โอฆสฺส นิตฺถรณาย. เยปิ เต สาฬฺห สมณพฺราหฺมณา อปริสุทฺธ\-กายสมาจารา อปริสุทฺธ\-วจีสมาจารา อปริสุทฺธ\-มโนสมาจารา อปริสุทฺธา\-ชีวา, อภพฺพา เต ญาณทสฺสนาย อนุตฺตราย สมฺโพธาย.}

\prose{เย จ โข เต สาฬฺห สมณพฺราหฺมณา น ตโปชิคุจฺฉา\-วาทา น ตโปชิคุจฺฉา\-สารา น ตโปชิคุจฺฉา\-อลฺลีนา วิหรนฺติ, ภพฺพา เต โอฆสฺส นิตฺถรณาย. เยปิ เต สาฬฺห สมณพฺราหฺมณา ปริสุทฺธ\-กาย\-สมาจารา ปริสุทฺธ\-วจี\-สมาจารา ปริสุทฺธ\-มโน\-สมาจารา ปริสุทฺธาชีวา, ภพฺพา เต ญาณทสฺสนาย อนุตฺตราย สมฺโพธาย.}

\prose{\csromanfolio{522}เสยฺยถาปิ สาฬฺห ปุริโส นทิํ ตริตุกาโม ติณฺหํ กุฐาริํ อาทาย วนํ ปวิเสยฺย, โส ตตฺถ ปสฺเสยฺย มหติํ สาลลฏฺฐิํ อุชุํ นวํ อกุกฺกุจฺจกชาตํ, ตเมนํ มูเล ฉินฺเทยฺย, มูเล ฉินฺทิตฺวา อคฺเค ฉินฺเทยฺย, อคฺเค ฉินฺทิตฺวา สาขาปลาสํ สุวิโสธิตํ วิโสเธยฺย, สาขาปลาสํ สุวิโสธิตํ วิโสเธตฺวา กุฐารีหิ ตจฺเฉยฺย, กุฐารีหิ ตจฺเฉตฺวา วาสีหิ ตจฺเฉยฺย, วาสีหิ ตจฺเฉตฺวา นิขาทนํ อาทาย อนฺโต สุวิโสธิตํ วิโสเธยฺย, อนฺโต สุวิโสธิตํ วิโสเธตฺวา เลขณิยา ลิเขยฺย, เลขณิยา ลิขิตฺวา ปาสาณคุเฬน โธเวยฺย, ปาสาณคุเฬน โธเวตฺวา นาวํ กเรยฺย, นาวํ กตฺวา ผิยาริตฺตํ\csromansharedfootnote{11927c929e841540}{ปิยาริตฺตํ (สี, อิ)} พนฺเธยฺย, ผิยาริตฺตํ พนฺธิตฺวา นทิํ ปตาเรยฺย.}

\prose{ตํ กิํ มญฺญสิ สาฬฺห, ภพฺโพ นุ โข โส ปุริโส นทิํ ตริตุนฺติ. เอวํ ภนฺเต. ตํ กิสฺส เหตุ? อสุ หิ ภนฺเต สาลลฏฺฐิ พหิทฺธา สุปริกมฺม\-กตา อนฺโต สุวิสุทฺธา นาวากตา\csromansharedfootnote{146ba4efa452ffbb}{สุวิสุทฺธกตา (ก)} ผิยา\-ริตฺต\-พทฺธา, ตสฺเสตํ ปาฏิกงฺขํ “นาวา น สํสีทิสฺสติ, ปุริโส โสตฺถินา ปารํ คมิสฺสตี”ติ.}

\prose{เอวเมวํ โข สาฬฺห เย เต สมณพฺราหฺมณา น ตโปชิคุจฺฉา\-วาทา น ตโปชิคุจฺฉา\-สารา น ตโปชิคุจฺฉา\-อลฺลีนา วิหรนฺติ, ภพฺพา เต โอฆสฺส นิตฺถรณาย. เยปิ เต สาฬฺห สมณพฺราหฺมณา ปริสุทฺธ\-กาย\-สมาจารา ปริสุทฺธ\-วจี\-สมาจารา ปริสุทฺธ\-มโน\-สมาจารา ปริสุทฺธาชีวา, ภพฺพา เต ญาณทสฺสนาย อนุตฺตราย สมฺโพธาย. เสยฺยถาปิ สาฬฺห โยธาชีโว พหูนิ เจปิ กณฺฑจิตฺรกานิ ชานาติ. อถ โข โส ตีหิ ฐาเนหิ ราชารโห โหติ ราชโภคฺโค, รญฺโญ องฺคนฺเตว สงฺขํ คจฺฉติ. กตเมหิ ตีหิ? ทูเรปาตี จ อกฺขณเวธี จ มหโต จ กายสฺส ปทาเลตา.}

\prose{เสยฺยถาปิ สาฬฺห โยธาชีโว ทูเรปาตี, เอวเมวํ โข สาฬฺห อริยสาวโก สมฺมาสมาธิ โหติ. สมฺมาสมาธิ สาฬฺห อริยสาวโก ยํ กิญฺจิ รูปํ อตีตา\-นาคต\-ปจฺจุปฺปนฺนํ อชฺฌตฺตํ วา พหิทฺธา วา โอฬาริกํ วา สุขุมํ วา หีนํ วา ปณีตํ วา ยํ ทูเร สนฺติเก วา, สพฺพํ รูปํ “เนตํ มม, เนโสหมสฺมิ, น เมโส อตฺตา”ติ เอวเมตํ ยถาภูตํ สมฺมปฺปญฺญาย ปสฺสติ. ยา กาจิ เวทนา. ยา กาจิ สญฺญา.}

\vspace{\tipitakaparskip}
\csromancenter{(20) 5. มหาวคฺค เย เกจิ สงฺขารา. ยํ กิญฺจิ วิญฺญาณํ อตีตา\-นาคต\-ปจฺจุปฺปนฺนํ อชฺฌตฺตํ วา พหิทฺธา วา โอฬาริกํ วา สุขุมํ วา หีนํ วา ปณีตํ วา ยํ ทูเร สนฺติเก วา, สพฺพํ วิญฺญาณํ “เนตํ มม, เนโสหมสฺมิ, น เมโส อตฺตา”ติ เอวเมตํ ยถาภูตํ สมฺมปฺปญฺญาย ปสฺสติ.}
\prose{\csromanfolio{523}เสยฺยถาปิ สาฬฺห โยธาชีโว อกฺขณเวธี, เอวเมวํ โข สาฬฺห อริยสาวโก สมฺมาทิฏฺฐิ โหติ. สมฺมาทิฏฺฐิ สาฬฺห อริยสาวโก “อิทํ ทุกฺขนฺ”ติ ยถาภูตํ ปชานาติ ฯเปฯ “อยํ ทุกฺขนิโรธ\-คามินี ปฏิปทา”ติ ยถาภูตํ ปชานาติ.}

\prose{เสยฺยถาปิ สาฬฺห โยธาชีโว มหโต กายสา ปทาเลตา, เอวเมวํ โข สาฬฺห อริยสาวโก สมฺมาวิมุตฺติ โหติ. สมฺมาวิมุตฺติ สาฬฺห อริยสาวโก มหนฺตํ อวิชฺชา\-กฺขนฺธํ ปทาเลตีติ.  ฉฏฺฐํ.}

\csromanmatikaanchor{mtk.236}
\csromantocmarkref{subsection}{7. มลฺลิกาเทวีสุตฺต}{mtk.236}
\bookmark[dest={mtk.236},level=2]{7. มลฺลิกาเทวีสุตฺต}
\csromanheader{2}{7. มลฺลิกาเทวี\-สุตฺต}
\proseitem{197}{เอกํ สมยํ ภควา สาวตฺถิยํ วิหรติ เชตวเน อนาถ\-ปิณฺฑิกสฺส อาราเม. อถ โข มลฺลิกา เทวี เยน ภควา เตนุปสงฺกมิ, อุปสงฺกมิตฺวา ภควนฺตํ อภิวาเทตฺวา เอกมนฺตํ นิสีทิ, เอกมนฺตํ นิสินฺนา โข มลฺลิกา เทวี ภควนฺตํ เอตทโวจ\csromandash{}}

\prose{โก นุ โข ภนฺเต เหตุ โก ปจฺจโย, เยน มิเธกจฺโจ มาตุคาโม ทุพฺพณฺณา จ โหติ ทุรูปา สุปาปิกา\csromansharedfootnote{0a25994523b37266}{สฺยาอิโปตฺถเกสุ “ทุพฺพณฺโณ จ โหติ ทุรุโป สุปาปโก”ติ เอวมาทินา ปุลฺลิงฺคิกวเสน ทิสฺสติ.} ทสฺสนาย, ทลิทฺทา จ โหติ อปฺปสฺสกา อปฺปโภคา อปฺเปสกฺขา จ.}

\prose{โก ปน ภนฺเต เหตุ โก ปจฺจโย, เยน มิเธกจฺโจ มาตุคาโม ทุพฺพณฺณา จ โหติ ทุรูปา สุปาปิกา ทสฺสนาย, อฑฺฒา จ โหติ มหทฺธนา มหาโภคา มเหสกฺขา จ.}

\prose{โก นุ โข ภนฺเต เหตุ โก ปจฺจยา, เยน มิเธกจฺโจ มาตุคาโม อภิรูปา จ โหติ ทสฺสนียา ปาสาทิกา ปรมาย วณฺณ\-โปกฺขร\-ตาย สมนฺนาคตา, ทลิทฺทา จ โหติ อปฺปสฺสกา อปฺปโภคา อปฺเปสกฺขา จ.}

\prose{\csromanfolio{524}โก ปน ภนฺเต เหตุ โก ปจฺจโย, เยน มิเธกจฺโจ มาตุคาโม อภิรูปา จ โหติ ทสฺสนียา ปาสาทิกา ปรมาย วณฺณ\-โปกฺขร\-ตาย สมนฺนาคตา, อฑฺฒา จ โหติ มหทฺธนา มหาโภคา มเหสกฺขา จาติ.}

\prose{อิธ มลฺลิเก เอกจฺโจ มาตุคาโม โกธนา โหติ อุปายาสพหุลา, อปฺปมฺปิ วุตฺตา สมานา อภิสชฺชติ กุปฺปติ พฺยาปชฺชติ ปติตฺถียติ, โกปญฺจ โทสญฺจ อปฺปจฺจยญฺจ ปาตุกโรติ. สา\csromansharedfootnote{1d36d9a0dcdf31bb}{โส (สฺยา)} น ทาตา โหติ สมณสฺส วา พฺราหฺมณสฺส วา อนฺนํ ปานํ วตฺถํ ยานํ มาลา\-คนฺธ\-วิเลปนํ เสยฺยา\-วสถ\-ปทีเปยฺยํ. อิสฺสามนิกา\csromansharedfootnote{1b2bcef82850b1c4}{อิสฺสามนโก (สฺยา)} โข ปน โหติ ปรลาภ\-สกฺการ\-ครุการ\-มานน\-วนฺทน\-ปูชนาสุ อิสฺสติ อุปทุสฺสติ อิสฺสํ พนฺธติ. สา\csromansharedfootnote{1d36d9a0dcdf31bb}{โส (สฺยา)} เจ ตโต จุตา อิตฺถตฺตํ อาคจฺฉติ. สา\csromansharedfootnote{1d36d9a0dcdf31bb}{โส (สฺยา)} ยตฺถ ยตฺถ ปจฺจาชายติ, ทุพฺพณฺณา จ โหติ ทุรูปา สุปาปิกา ทสฺสนาย, ทลิทฺทา จ โหติ อปฺปสฺสกา อปฺปโภคา อปฺเปสกฺขา จ.}

\prose{อิธ ปน มลฺลิเก เอกจฺโจ มาตุคาโม โกธนา โหติ อุปายาสพหุลา, อปฺปมฺปิ วุตฺตา สมานา อภิสชฺชติ กุปฺปติ พฺยาปชฺชติ ปติตฺถียติ, โกปญฺจ โทสญฺจ อปฺปจฺจยญฺจ ปาตุกโรติ. สา ทาตา โหติ สมณสฺส วา พฺราหฺมณสฺส วา อนฺนํ ปานํ วตฺถํ ยานํ มาลา\-คนฺธ\-วิเลปนํ เสยฺยา\-วสถ\-ปทีเปยฺยํ. อนิสฺสามนิกา โข ปน โหติ ปรลาภ\-สกฺการ\-ครุการ\-มานน\-วนฺทน\-ปูชนาสุ น อิสฺสติ น อุปทุสฺสติ น อิสฺสํ พนฺธติ. สา เจ ตโต จุตา อิตฺถตฺตํ อาคจฺฉติ. สา ยตฺถ ยตฺถ ปจฺจาชายติ, ทุพฺพณฺณา จ โหติ ทุรูปา สุปาปิกา ทสฺสนาย, อฑฺฒา จ โหติ มหทฺธนา มหาโภคา มเหสกฺขา จ.}

\prose{อิธ ปน มลฺลิเก เอกจฺโจ มาตุคาโม อกฺโกธนา โหติ อนุปายาส\-พหุลา, พหุมฺปิ วุตฺตา สมานา นาภิสชฺชติ น กุปฺปติ น พฺยาปชฺชติ น ปติตฺถียติ, น โกปญฺจ โทสญฺจ อปฺปจฺจยญฺจ ปาตุกโรติ. สา น ทาตา โหติ สมณสฺส วา พฺราหฺมณสฺส วา อนฺนํ ปานํ วตฺถํ ยานํ มาลา\-คนฺธ\-วิเลปนํ เสยฺยา\-วสถ\-ปทีเปยฺยํ. อิสฺสามนิกา โข ปน โหติ ปรลาภ\-สกฺการ\-ครุการ\-มานน\-วนฺทน\-ปูชนาสุ อิสฺสติ อุปทุสฺสติ อิสฺสํ พนฺธติ. สา เจ ตโต จุตา อิตฺถตฺตํ อาคจฺฉติ. สา}

\vspace{\tipitakaparskip}
\csromancenter{(20) 5. มหาวคฺค ยตฺถ ยตฺถ ปจฺจาชายติ, อภิรูปา จ โหติ ทสฺสนียา ปาสาทิกา ปรมาย วณฺณ\-โปกฺขร\-ตาย สมนฺนาคตา, ทลิทฺทา จ โหติ อปฺปสฺสกา อปฺปโภคา อปฺเปสกฺขา จ.}
\prose{\csromanfolio{525}อิธ ปน มลฺลิเก เอกจฺโจ มาตุคาโม อกฺโกธนา โหติ อนุปายาส\-พหุลา, พหุมฺปิ วุตฺตา สมานา นาภิสชฺชติ น กุปฺปติ น พฺยาปชฺชติ น ปติตฺถียติ, น โกปญฺจ โทสญฺจ อปฺปจฺจยญฺจ ปาตุกโรติ. สา ทาตา โหติ สมณสฺส วา พฺราหฺมณสฺส วา อนฺนํ ปานํ วตฺถํ ยานํ มาลา\-คนฺธ\-วิเลปนํ เสยฺยา\-วสถ\-ปทีเปยฺยํ. อนิสฺสามนิกา โข ปน โหติ ปรลาภ\-สกฺการ\-ครุการ\-มานน\-วนฺทน\-ปูชนาสุ น อิสฺสติ น อุปทุสฺสติ น อิสฺสํ พนฺธติ. สา เจ ตโต จุตา อิตฺถตฺตํ อาคจฺฉติ. สา ยตฺถ ยตฺถ ปจฺจาชายติ, อภิรูปา จ โหติ ทสฺสนียา ปาสาทิกา ปรมาย วณฺณ\-โปกฺขร\-ตาย สมนฺนาคตา, อฑฺฒา จ โหติ มหทฺธนา มหาโภคา มเหสกฺขา จ.}

\prose{อยํ โข มลฺลิเก เหตุ อยํ ปจฺจโย, เยน มิเธกจฺโจ มาตุคาโม ทุพฺพณฺณา จ โหติ ทุรูปา สุปาปิกา ทสฺสนาย, ทลิทฺทา จ โหติ อปฺปสฺสกา อปฺปโภคา อปฺเปสกฺขา จ. อยํ ปน มลฺลิเก เหตุ อยํ ปจฺจโย, เยน มิเธกจฺโจ มาตุคาโม ทุพฺพณฺณา จ โหติ ทุรูปา สุปาปิกา ทสฺสนาย, อฑฺฒา จ โหติ มหทฺธนา มหาโภคา มเหสกฺขา จ. อยํ โข มลฺลิเก เหตุ อยํ ปจฺจโย, เยน มิเธกจฺโจ มาตุคาโม อภิรูปา จ โหติ ทสฺสนียา ปาสาทิกา ปรมาย วณฺณ\-โปกฺขร\-ตาย สมนฺนาคตา, ทลิทฺทา จ โหติ อปฺปสฺสกา อปฺปโภคา อปฺเปสกฺขา จ. อยํ ปน มลฺลิเก เหตุ อยํ ปจฺจโย, เยน มิเธกจฺโจ มาตุคาโม อภิรูปา จ โหติ ทสฺสนียา ปาสาทิกา ปรมาย วณฺณ\-โปกฺขร\-ตาย สมนฺนาคตา, อฑฺฒา จ โหติ มหทฺธนา มหาโภคา มเหสกฺขา จาติ.}

\prose{เอวํ วุตฺเต มลฺลิกา เทวี ภควนฺตํ เอตทโวจ\csromandash{}ยา นูนาหํ\csromansharedfootnote{9a7d86a1de7cb350}{สา นูนาหํ (สฺยา), ยํ นูนาหํ (ก)} ภนฺเต อญฺญํ ชาติํ\csromansharedfootnote{a85de30d02c2cefc}{อญฺญาย ชาติยา (สฺยา)} โกธนา อโหสิํ อุปายาสพหุลา อปฺปมฺปิ วุตฺตา}

\prose{\csromanfolio{526}สมานา อภิสชฺชิํ กุปฺปิํ พฺยาปชฺชิํ ปติตฺถียิํ, โกปญฺจ โทสญฺจ อปฺปจฺจยญฺจ ปาตฺวากาสิํ, สาหํ ภนฺเต เอตรหิ ทุพฺพณฺณา ทุรูปา สุปาปิกา ทสฺสนาย.}

\prose{ยา นูนาหํ ภนฺเต อญฺญํ ชาติํ ทาตา อโหสิํ สมณสฺส วา พฺราหฺมณสฺส วา อนฺนํ ปานํ วตฺถํ ยานํ มาลา\-คนฺธ\-วิเลปนํ เสยฺยา\-วสถ\-ปทีเปยฺยํ, สาหํ ภนฺเต เอตรหิ อฑฺฒา\csromansharedfootnote{4583d1fddc0db2bc}{อฑฺฒา จ (สี, อิ, ก)} มหทฺธนา มหาโภคา.}

\prose{ยา นูนาหํ ภนฺเต อญฺญํ ชาติํ อนิสฺสามนิกา อโหสิํ ปรลาภ\-สกฺการ\-ครุการ\-มานน\-วนฺทน\-ปูชนาสุ น อิสฺสิํ น อุปทุสฺสิํ น อิสฺสํ พนฺธิํ, สาหํ ภนฺเต เอตรหิ มเหสกฺขา. สนฺติ โข ปน ภนฺเต อิมสฺมิํ ราชกุเล ขตฺติยกญฺญาปิ พฺราหฺมณกญฺญาปิ คหปติกญฺญาปิ, ตาสาหํ อิสฺสรา\-ธิปจฺจํ กาเรมิ. เอสาหํ ภนฺเต อชฺชตคฺเค อกฺโกธนา ภวิสฺสามิ อนุปายาส\-พหุลา, พหุมฺปิ วุตฺตา สมานา นาภิสชฺชิสฺสามิ น กุปฺปิสฺสามิ น พฺยาปชฺชิสฺสามิ น ปติตฺถียิสฺสามิ, โกปญฺจ โทสญฺจ อปฺปจฺจยญฺจ น ปาตุกริสฺสามิ, ทสฺสามิ สมณสฺส วา พฺราหฺมณสฺส วา อนฺนํ ปานํ วตฺถํ ยานํ มาลา\-คนฺธ\-วิเลปนํ เสยฺยา\-วสถ\-ปทีเปยฺยํ, อนิสฺสามนิกา ภวิสฺสามิ ปรลาภ\-สกฺการ\-ครุการ\-มานน\-วนฺทน\-ปูชนาสุ น อิสฺสิสฺสามิ น อุปทุสฺสิสฺสามิ น อิสฺสํ พนฺธิสฺสามิ. อภิกฺกนฺตํ ภนฺเต ฯเปฯ อุปาสิกํ มํ ภนฺเต ภควา ธาเรตุ อชฺชตคฺเค ปาณุเปตํ สรณํ คตนฺติ.  }

\csromanmatikaanchor{mtk.237}
\csromantocmarkref{subsection}{8. อตฺตนฺตปสุตฺต}{mtk.237}
\bookmark[dest={mtk.237},level=2]{8. อตฺตนฺตปสุตฺต}
\csromanheader{2}{8. \textbf{อตฺตนฺตปสุตฺต}}
\proseitem{198}{จตฺตาโรเม ภิกฺขเว ปุคฺคลา สนฺโต สํวิชฺชมานา โลกสฺมิํ. กตเม จตฺตาโร? อิธ ภิกฺขเว เอกจฺโจ\csromansharedfootnote{1ac5e4063b3d02dd}{อภิ 3. 163; ม 2. 5; ที 3. 194 ปิฏฺเฐสุปิ, เหฏฺฐา 300 ปิฏฺเฐปิ.} ปุคฺคโล อตฺตนฺตโป โหติ อตฺต\-ปริตาปนา\-นุโยคม\-นุยุตฺโต, อิธ ปน ภิกฺขเว เอกจฺโจ ปุคฺคโล ปรนฺตโป โหติ ปรปริตาปนา\-นุโยคม\-นุยุตฺโต, อิธ ปน ภิกฺขเว เอกจฺโจ ปุคฺคโล อตฺตนฺตโป จ โหติ อตฺต\-ปริตาปนา\-นุโยคม\-นุยุตฺโต ปรนฺตโป จ ปรปริตาปนา\-นุโยคม\-นุยุตฺโต, อิธ ปน ภิกฺขเว เอกจฺโจ ปุคฺคโล เนวตฺตนฺตโป โหติ นาตฺตปริตาปนานุโยค}

\vspace{\tipitakaparskip}
\csromancenter{(20) 5. มหาวคฺค มนุยุตฺโต น ปรนฺตโป น ปรปริตาปนา\-นุโยคม\-นุยุตฺโต, โส เนว อตฺตนฺตโป น ปรนฺตโป ทิฏฺเฐว ธมฺเม นิจฺฉาโต นิพฺพุโต สีตีภูโต สุขปฺปฏิสํเวที พฺรหฺมภูเตน อตฺตนา วิหรติ.}
\prose{\csromanfolio{527}กถญฺจ ภิกฺขเว ปุคฺคโล อตฺตนฺตโป โหติ อตฺต\-ปริตาปนา\-นุโยคม\-นุยุตฺโต? อิธ ภิกฺขเว เอกจฺโจ อเจลโก โหติ มุตฺตาจาโร, หตฺถา\-ปเลขโน, นเอหิภทฺทนฺติโก, นติฏฺฐ\-ภทฺทนฺติโก, นาภิหฏํ, น อุทฺทิสฺสกตํ, น นิมนฺตนํ สาทิยติ. โส น กุมฺภิมุขา ปฏิคฺคณฺหาติ, น กโฬปิมุขา ปฏิคฺคณฺหาติ, น เอฬกมนฺตรํ, น ทณฺฑมนฺตรํ, น มุสลมนฺตรํ, น ทฺวินฺนํ ภุญฺชมานานํ, น คพฺภินิยา, น ปายมานาย, น ปุริส\-นฺตร\-คตาย, น สงฺกิตฺตีสุ, น ยตฺถ สา อุปฏฺฐิโต โหติ, น ยตฺถ มกฺขิกา สณฺฑ\-สณฺฑ\-จารินี, น มจฺฉํ, น มํสํ, น สุรํ, น เมรยํ, น ถุโสทกํ ปิวติ. โส เอกาคาริโก วา โหติ เอกาโลปิโก, ทฺวาคาริโก วา โหติ ทฺวาโลปิโก ฯเปฯ สตฺตาคาริโก วา โหติ สตฺตาโลปิโก, เอกิสฺสาปิ ทตฺติยา ยาเปติ, ทฺวีหิปิ ทตฺตีหิ ยาเปติ ฯเปฯ สตฺตหิปิ ทตฺตีหิ ยาเปติ, เอกาหิกมฺปิ อาหารํ อาหาเรติ, ทฺวาหิกมฺปิ อาหารํ อาหาเรติ ฯเปฯ สตฺตาหิกมฺปิ อาหารํ อาหาเรติ, อิติ เอวรูปํ อฑฺฒมาสิกมฺปิ ปริยาย\-ภตฺตโภชนานุโยคมนุยุตฺโต วิหรติ.}

\prose{โส สากภกฺโขปิ โหติ, สามาก\-ภกฺโขปิ โหติ, นีวาร\-ภกฺโขปิ โหติ, ททฺทุล\-ภกฺโขปิ โหติ, หฏภกฺโขปิ โหติ, กณภกฺโขปิ โหติ, อาจามภกฺโขปิ โหติ, ปิญฺญาก\-ภกฺโขปิ โหติ, ติณภกฺโขปิ โหติ, โคมย\-ภกฺโขปิ โหติ, วนมูล\-ผลา\-หาโรปิ ยาเปติ ปวตฺต\-ผล\-โภชี.}

\prose{โส สาณานิปิ ธาเรติ, มสาณานิปิ ธาเรติ, ฉวทุสฺสานิปิ ธาเรติ, ปํสุกูลานิปิ ธาเรติ, ติรีฐานิปิ ธาเรติ, อชินมฺปิ ธาเรติ, อชินกฺขิปมฺปิ ธาเรติ, กุสจีรมฺปิ ธาเรติ, วากจีรมฺปิ ธาเรติ, ผลกจีรมฺปิ ธาเรติ, เกสกมฺพลมฺปิ ธาเรติ, วาฬกมฺพลมฺปิ ธาเรติ, อุลูกปกฺขมฺปิ ธาเรติ, เกสมสฺสุ\-โลจโกปิ โหติ เกสมสฺสุ\-โลจนา\-นุโยคม\-นุยุตฺโต, อุพฺภฏฺฐโกปิ โหติ อาสน\-ปฺปฏิกฺขิตฺโต, อุกฺกุฏิโกปิ โหติ อุกฺกุฏิก\-ปฺปธานมนุยุตฺโต, กณฺฏกา\-ปสฺสยิ\-โกปิ โหติ กณฺฏกาปสฺสเย เสยฺยํ กปฺเปติ, สาย\-ตติยกมฺปิ อุทโก\-โรหนา\-นุโยคม\-นุยุตฺโต วิหรติ. อิติ เอวรูปํ อเนกวิหิตํ กายสฺส \csromanfolio{528}อาตาปน\-ปริตาปนา\-นุโยคม\-นุยุตฺโต วิหรติ. เอวํ โข ภิกฺขเว ปุคฺคโล อตฺตนฺตโป โหติ อตฺต\-ปริตาปนา\-นุโยคม\-นุยุตฺโต.}

\prose{กถญฺจ ภิกฺขเว ปุคฺคโล ปรนฺตโป โหติ ปรปริตาปนา\-นุโยคม\-นุยุตฺโต? อิธ ภิกฺขเว เอกจฺโจ ปุคฺคโล โอรมฺภิโก โหติ สูกริโก สากุณิโก มาควิโก ลุทฺโท มจฺฉฆาตโก โจโร โจรฆาตโก โคฆาตโก พนฺธนาคาริโก, เย วา ปนญฺเญปิ เกจิ กุรูรกมฺมนฺตา. เอวํ โข ภิกฺขเว ปุคฺคโล ปรนฺตโป โหติ ปรปริตาปนา\-นุโยคม\-นุยุตฺโต.}

\prose{กถญฺจ ภิกฺขเว ปุคฺคโล อตฺตนฺตโป จ โหติ อตฺต\-ปริตาปนา\-นุโยคม\-นุยุตฺโต ปรนฺตโป จ ปรปริตาปนา\-นุโยคม\-นุยุตฺโต? อิธ ภิกฺขเว เอกจฺโจ ปุคฺคโล ราชา วา โหติ ขตฺติโย มุทฺธาวสิตฺโต, พฺราหฺมโณ วา โหติ มหาสาโล, โส ปุรตฺถิเมน นครสฺส นวํ สนฺถาคารํ การาเปตฺวา เกสมสฺสุํ โอหาเรตฺวา ขราชินํ นิวาเสตฺวา สปฺปิเตเลน กายํ อพฺภญฺชิตฺวา มควิสาเณน ปิฏฺฐิํ กณฺฑุวมาโน นวํ สนฺถาคารํ ปวิสติ สทฺธิํ มเหสิยา พฺราหฺมเณน จ ปุโรหิเตน. โส ตตฺถ อนนฺตรหิตาย ภูมิยา หริตุ\-ปลิตฺตาย เสยฺยํ กปฺเปติ. เอกิสฺสาย คาวิยา สรูปวจฺฉาย ยํ เอกสฺมิํ ถเน ขีรํ โหติ, เตน ราชา ยาเปติ. ยํ ทุติยสฺมิํ ถเน ขีรํ โหติ, เตน มเหสี ยาเปติ. ยํ ตติยสฺมิํ ถเน ขีรํ โหติ, เตน พฺราหฺมโณ ปุโรหิโต ยาเปติ. ยํ จตุตฺถสฺมิํ ถเน ขีรํ โหติ, เตน อคฺคิํ ชุหติ\csromansharedfootnote{c41aab65c304ec16}{ชุหนฺติ (สี, อิ)}, อวเสเสน วจฺฉโก ยาเปติ. โส เอวมาห “เอตฺตกา อุสภา หญฺญนฺตุ ยญฺญตฺถาย, เอตฺตกา วจฺฉตรา หญฺญนฺตุ ยญฺญตฺถาย, เอตฺตกา วจฺฉตริโย หญฺญนฺตุ ยญฺญตฺถาย, เอตฺตกา อชา หญฺญนฺตุ ยญฺญตฺถาย, เอตฺตกา อุรพฺภา หญฺญนฺตุ ยญฺญตฺถาย, (เอตฺตกา อสฺสา หญฺญนฺตุ ยญฺญตฺถาย,)\csromansharedfootnote{31ca72165988f669}{( ) นตฺถิ สีสฺยากํอิโปตฺถเกสุ. 3. ปริกมฺมตฺถายาติ (ก)} เอตฺตกา รุกฺขา ฉิชฺชนฺตุ ยูปตฺถาย, เอตฺตกา ทพฺพา ลูยนฺตุ พริหิสตฺถายา”ติ\csromansharedfootnote{05199ee52c46abb1}{ปริกมฺมตฺถายาติ (ก)}. เยปิสฺส เต โหนฺติ ทาสาติ วา เปสฺสาติ วา กมฺมกราติ วา, เตปิ ทณฺฑตชฺชิตา ภยตชฺชิตา อสฺสุมุขา รุทมานา ปริกมฺมานิ กโรนฺติ. เอวํ โข ภิกฺขเว ปุคฺคโล อตฺตนฺตโป จ โหติ อตฺต\-ปริตาปนา\-นุโยคม\-นุยุตฺโต ปรนฺตโป จ ปรปริตาปนา\-นุโยคม\-นุยุตฺโต.}

\chapterheadpageread{(20) 5. มหาวคฺค}
\prose{\csromanfolio{529}กถญฺจ ภิกฺขเว ปุคฺคโล เนวตฺตนฺตโป โหติ นาตฺต\-ปริตาปนา\-นุโยคม\-นุยุตฺโต น ปรนฺตโปน ปรปริตาปนา\-นุโยคม\-นุยุตฺโต, โส อนตฺตนฺตโป อปรนฺตโป ทิฏฺเฐว ธมฺเม นิจฺฉาโต นิพฺพุโต สีตีภูโต สุขปฺปฏิสํเวที พฺรหฺมภูเตน อตฺตนา วิหรติ? อิธ ภิกฺขเว ตถาคโต โลเก อุปฺปชฺชติ อรหํ สมฺมาสมฺพุทฺโธ วิชฺชา\-จรณ\-สมฺปนฺโน สุคโต โลกวิทู อนุตฺตโร ปุริส\-ทมฺม\-สารถิ สตฺถา เทวมนุสฺสานํ พุทฺโธ ภควา. โส อิมํ โลกํ สเทวกํ สมารกํ สพฺรหฺมกํ สสฺสมณ\-พฺราหฺมณิํ ปชํ สเทวมนุสฺสํ สยํ อภิญฺญา สจฺฉิกตฺวา ปเวเทติ. โส ธมฺมํ เทเสติ อาทิกลฺยาณํ มชฺเฌกลฺยาณํ ปริโยสาน\-กลฺยาณํ สาตฺถํ สพฺยญฺชนํ เกวล\-ปริปุณฺณํ ปริสุทฺธํ พฺรหฺมจริยํ ปกาเสติ, ตํ ธมฺมํ สุณาติ คหปติ วา คหปติปุตฺโต วา อญฺญตรสฺมิํ วา กุเล ปจฺจาชาโต. โส ตํ ธมฺมํ สุตฺวา ตถาคเต สทฺธํ ปฏิลภติ, โส เตน สทฺธา\-ปฏิลาเภน สมนฺนาคโต อิติ ปฏิสญฺจิกฺขติ “สมฺพาโธ ฆราวาโส รชาปโถ อพฺโภกาโส ปพฺพชฺชา, นยิทํ สุกรํ อคารํ อชฺฌาวสตา เอกนฺต\-ปริปุณฺณํ เอกนฺต\-ปริสุทฺธํ สงฺขลิขิตํ พฺรหฺมจริยํ จริตุํ, ยํนูนาหํ เกสมสฺสุํ โอหาเรตฺวา กาสายานิ วตฺถานิ อจฺฉาเทตฺวา อคารสฺมา อนคาริยํ ปพฺพเชยฺยนฺ”ติ. โส อปเรน สมเยน อปฺปํ วา โภคกฺขนฺธํ ปหาย มหนฺตํ วา โภคกฺขนฺธํ ปหาย อปฺปํ วา ญาติปริวฏฺฏํ ปหาย มหนฺตํ วา ญาติปริวฏฺฏํ ปหาย เกสมสฺสุํ โอหาเรตฺวา กาสายานิ วตฺถานิ อจฺฉาเทตฺวา อคารสฺมา อนคาริยํ ปพฺพชติ.}

\prose{โส เอวํ ปพฺพชิโต สมาโน ภิกฺขูนํ สิกฺขา\-สาชีว\-สมาปนฺโน ปาณาติปาตํ ปหาย ปาณาติปาตา ปฏิวิรโต โหติ, นิหิตทณฺโฑ นิหิตสตฺโถ ลชฺชี ทยาปนฺโน สพฺพปาณภูต\-หิตา\-นุกมฺปี วิหรติ, อทินฺนาทานํ ปหาย อทินฺนาทานา ปฏิวิรโต โหติ ทินฺนาทายี ทินฺน\-ปาฏิกงฺขี อเถเนน สุจิภูเตน อตฺตนา วิหรติ, อพฺรหฺมจริยํ ปหาย พฺรหฺมจารี โหติ อาราจารี วิรโต อสทฺธมฺมา คามธมฺโม. มุสาวาทํ ปหาย มุสาวาทา ปฏิวิรโต โหติ สจฺจวาที สจฺจสนฺโธ เถโต ปจฺจยิโก อวิสํวาทโก โลกสฺส, ปิสุณํ วาจํ ปหาย ปิสุณาย วาจาย ปฏิวิรโต โหติ น อิโต สุตฺวา อมุตฺร อกฺขาตา อิเมสํ เภทาย, น อมุตฺร วา สุตฺวา อิเมสํ อกฺขาตา}

\prose{\csromanfolio{530}อมูสํ เภทาย, อิติ ภินฺนานํ วา สนฺธาตา, สหิตานํ วา อนุปฺปทาตา สมคฺคาราโม สมคฺครโต สมคฺคนนฺที สมคฺคกรณิํ วาจํ ภาสิตา โหติ, ผรุสํ วาจํ ปหาย ผรุสาย วาจาย ปฏิวิรโต โหติ ยา สา วาจา เนลา กณฺณสุขา เปมนียา หทยงฺคมา โปรี พหุชนกนฺตา พหุชนมนาปา, ตถารูปิํ วาจํ ภาสิตา โหติ, สมฺผปฺปลาปํ ปหาย สมฺผปฺปลาปา ปฏิวิรโต โหติ กาลวาที ภูตวาที อตฺถวาที ธมฺมวาที วินยวาที นิธานวติํ วาจํ ภาสิตา โหติ กาเลน สาปเทสํ ปริยนฺตวติํ อตฺถสํหิตํ.}

\prose{โส พีชคาม\-ภูตคาม\-สมารมฺภา ปฏิวิรโต โหติ, เอกภตฺติโก โหติ, รตฺตูปรโต วิรโต วิกาลโภชนา, นจฺจ\-คีต\-วาทิต\-วิสูก\-ทสฺสนา ปฏิวิรโต โหติ, มาลา\-คนฺธ\-วิเลปน\-ธารณ\-มณฺฑน\-วิภูสน\-ฏฺฐานา ปฏิวิรโต โหติ, อุจฺจาสยน\-มหาสยนา ปฏิวิรโต โหติ, ชาตรูป\-รชต\-ปฏิคฺคหณา ปฏิวิรโต โหติ, อามก\-ธญฺญ\-ปฏิคฺคหณา ปฏิวิรโต โหติ, อามก\-มํส\-ปฏิคฺคหณา ปฏิวิรโต โหติ, อิตฺถิ\-กุมาริก\-ปฏิคฺคหณา ปฏิวิรโต โหติ, ทาสิ\-ทาส\-ปฏิคฺคหณา ปฏิวิรโต โหติ, อเช\-ฬก\-ปฏิคฺคหณา ปฏิวิรโต โหติ, กุกฺกุฏ\-สูกร\-ปฏิคฺคหณา ปฏิวิรโต โหติ, หตฺถิ\-ควา\-สฺส\-วฬว\-ปฏิคฺคหณา ปฏิวิรโต โหติ, เขตฺตวตฺถุ\-ปฏิคฺคหณา ปฏิวิรโต โหติ, ทูเตยฺย\-ปหิณ\-คมนา\-นุโยคา ปฏิวิรโต โหติ, กยวิกฺกยา ปฏิวิรโต โหติ, ตุลา\-กูฏ\-กํส\-กูฏ\-มาน\-กูฏา ปฏิวิรโต โหติ, อุกฺโกฏน\-วญฺจน\-นิกติ\-สาจิโยคา ปฏิวิรโต โหติ, เฉทน\-วธ\-พนฺธน\-วิปราโมส\-อาโลป\-สหสาการา ปฏิวิรโต โหติ.}

\prose{โส สนฺตุฏฺโฐ โหติ กาย\-ปริหาริเกน จีวเรน กุจฺฉิ\-ปริหาริเกน ปิณฺฑปาเตน, โส เยน เยนว ปกฺกมติ, สมาทาเยว ปกฺกมติ. เสยฺยถาปิ นาม ปกฺขี สกุโณ เยน เยเนว เฑติ, สปตฺตภาโรว เฑติ. เอวเมวํ ภิกฺขุ สนฺตุฏฺโฐ โหติ กาย\-ปริหาริเกน จีวเรน กุจฺฉิ\-ปริหาริเกน ปิณฺฑปาเตน, โส เยน เยเนว ปกฺกมติ, สมาทาเยว ปกฺกมติ. โส อิมินา อริเยน สีลกฺขนฺเธน สมนฺนาคโต อชฺฌตฺตํ อนวชฺชสุขํ ปฏิสํเวเทติ.}

\prose{โส จกฺขุนา รูปํ ทิสฺวา น นิมิตฺตคฺคาหี โหติ นานุพฺยญฺชน\-คฺคาหี, ยตฺวาธิกรณเมนํ จกฺขุนฺทฺริยํ อสํวุตํ วิหรนฺตํ อภิชฺฌา\-โทมนสฺสา ปาปกา}

\vspace{\tipitakaparskip}
\csromancenter{(20) 5. มหาวคฺค อกุสลา ธมฺมา อนฺวาสฺสเวยฺยุํ. ตสฺส สํวราย ปฏิปชฺชติ, รกฺขติ จกฺขุนฺทฺริยํ, จกฺขุนฺทฺริเย สํวรํ อาปชฺชติ. โสเตน สทฺทํ สุตฺวา. ฆาเนน คนฺธํ ฆายิตฺวา. ชิวฺหาย รสํ สายิตฺวา. กาเยน โผฏฺฐพฺพํ ผุสิตฺวา. มนสา ธมฺมํ วิญฺญาย น นิมิตฺตคฺคาหี โหติ นานุพฺยญฺชน\-คฺคาหี, ยตฺวาธิกรณเมนํ มนินฺทฺริยํ อสํวุตํ วิหรนฺตํ อภิชฺฌา\-โทมนสฺสา ปาปกา อกุสลา ธมฺมา อนฺวาสฺสเวยฺยุํ. ตสฺส สํวราย ปฏิปชฺชติ, รกฺขติ มนินฺทฺริยํ, มนินฺทฺริเย สํวรํ อาปชฺชติ. โส อิมินา อริเยน อินฺทฺริย\-สํวเรน สมนฺนาคโต อชฺฌตฺตํ อพฺยาเสกสุขํ ปฏิสํเวเทติ.}
\prose{\csromanfolio{531}โส อภิกฺกนฺเต ปฏิกฺกนฺเต สมฺปชานการี โหติ, อาโลกิเต วิโลกิเต สมฺปชานการี โหติ, สมิญฺชิเต ปสาริเต สมฺปชานการี โหติ, สงฺฆาฏิ\-ปตฺต\-จีวร\-ธารเณ สมฺปชานการี โหติ, อสิเต ปีเต ขายิเต สายิเต สมฺปชานการี โหติ, อุจฺจาร\-ปสฺสาว\-กมฺเม สมฺปชานการี โหติ, คเต ฐิเต นิสินฺเน สุตฺเต ชาคริเต ภาสิเต ตุณฺหีภาเว สมฺปชานการี โหติ.}

\prose{โส อิมินา จ อริเยน สีลกฺขนฺเธน สมนฺนาคโต (อิมาย จ อริยาย สนฺตุฏฺฐิยา สมนฺนาคโต)\csromansharedfootnote{a456a47d6bed678d}{( ) นตฺถิ สีสฺยาโปตฺถเกสุ. ม 1. 239; ม 2. 10 ปิฏฺเฐสุ ปสฺสิตพฺพํ.} อิมินา จ อริเยน อินฺทฺริย\-สํวเรน สมนฺนาคโต อิมินา จ อริเยน สติสมฺปชญฺเญน สมนฺนาคโต\csromansharedfootnote{bff6f1aa45fbea6b}{สมนฺนาคโต. โส (ก)} วิวิตฺตํ เสนาสนํ ภชติ อรญฺญํ รุกฺขมูลํ ปพฺพตํ กนฺทรํ คิริคุหํ สุสานํ วนปฺปตฺถํ อพฺโภกาสํ ปลาลปุญฺชํ. โส ปจฺฉาภตฺตํ ปิณฺฑปาต\-ปฏิกฺกนฺโต นิสีทติ ปลฺลงฺกํ อาภุชิตฺวา อุชุํ กายํ ปณิธาย ปริมุขํ สติํ อุปฏฺฐเปตฺวา, โส อภิชฺฌํ โลเก ปหาย วิคตา\-ภิชฺเฌน เจตสา วิหรติ, อภิชฺฌาย จิตฺตํ ปริโสเธติ, พฺยาปาทปโทสํ ปหาย อพฺยาปนฺนจิตฺโต วิหรติ สพฺพปาณภูต\-หิตา\-นุกมฺปี, พฺยาปาทปโทสา จิตฺตํ ปริโสเธติ, ถินมิทฺธํ ปหาย วิคต\-ถินมิทฺโธ วิหรติ อาโลกสญฺญี สโต สมฺปชาโน, ถินมิทฺธา จิตฺตํ ปริโสเธติ, อุทฺธจฺจ\-กุกฺกุจฺจํ ปหาย อนุทฺธโต วิหรติ อชฺฌตฺตํ วูปสนฺตจิตฺโต, อุทฺธจฺจ\-กุกฺกุจฺจา จิตฺตํ ปริโสเธติ, วิจิกิจฺฉํ ปหาย ติณฺณ\-วิจิกิจฺโฉ วิหรติ อกถํกถี กุสเลสุ ธมฺเมสุ, วิจิกิจฺฉาย จิตฺตํ ปริโสเธติ, โส อิเม ปญฺจ นีวรเณ ปหาย เจตโส อุปกฺกิเลเส ปญฺญาย ทุพฺพลีกรเณ วิวิจฺเจว กาเมหิ ฯเปฯ จตุตฺถํ ฌานํ อุปสมฺปชฺช วิหรติ.}

\prose{\csromanfolio{532}โส เอวํ สมาหิเต จิตฺเต ปริสุทฺเธ ปริโยทาเต อนงฺคเณ วิคตู\-ปกฺกิเลเส มุทุภูเต กมฺมนิเย ฐิเต อาเนญฺชปฺปตฺเต ปุพฺเพ\-นิวาสา\-นุสฺสติ\-ญาณาย ฯเปฯ สตฺตานํ จุตูปปาต\-ญาณาย ฯเปฯ อาสวานํ ขยญาณาย จิตฺตํ อภินินฺนาเมติ, โส “อิทํ ทุกฺขนฺ”ติ ยถาภูตํ ปชานาติ, “อยํ ทุกฺขสมุทโย”ติ ยถาภูตํ ปชานาติ, “อยํ ทุกฺขนิโรโธ”ติ ยถาภูตํ ปชานาติ, “อยํ ทุกฺขนิโรธ\-คามินี ปฏิปทา”ติ ยถาภูตํ ปชานาติ, “อิเม อาสวา”ติ ยถาภูตํ ปชานาติ, “อยํ อาสวสมุทโย”ติ ยถาภูตํ ปชานาติ, “อยํ อาสวนิโรโธ”ติ ยถาภูตํ ปชานาติ, “อยํ อาสว\-นิโรธ\-คามินี ปฏิปทา”ติ ยถาภูตํ ปชานาติ.}

\prose{ตสฺส เอวํ ชานโต เอวํ ปสฺสโต กามาสวาปิ จิตฺตํ วิมุจฺจติ, ภวาสวาปิ จิตฺตํ วิมุจฺจติ, อวิชฺชาสวาปิ จิตฺตํ วิมุจฺจติ, วิมุตฺตสฺมิํ “วิมุตฺตมฺ”อิติ ญาณํ โหติ, “ขีณา ชาติ, วุสิตํ พฺรหฺมจริยํ, กตํ กรณียํ, นาปรํ อิตฺถตฺตายา”ติ ปชานาติ. เอวํ โข ภิกฺขเว ปุคฺคโล เนวตฺตนฺตโป โหติ นาตฺต\-ปริตาปนา\-นุโยคม\-นุยุตฺโต น ปรนฺตโป น ปรปริตาปนา\-นุโยคม\-นุยุตฺโต, โส น อตฺตนฺตโป น ปรนฺตโป ทิฏฺเฐว ธมฺเม นิจฺฉาโต นิพฺพุโต สีตีภูโต สุขปฺปฏิสํเวที พฺราหฺมภูเตน อตฺตนา วิหรติ, อิเม โข ภิกฺขเว จตฺตาโร ปุคฺคลา สนฺโต สํวิชฺชมานา โลกสฺมินฺติ.  อฏฺฐมํ.}

\csromanmatikaanchor{mtk.238}
\csromantocmarkref{subsection}{9. ตณฺหาสุตฺต}{mtk.238}
\bookmark[dest={mtk.238},level=2]{9. ตณฺหาสุตฺต}
\csromanheader{2}{9. \textbf{ตณฺหาสุตฺต}}
\proseitem{199}{ภควา เอตทโวจ “ตณฺหํ โว ภิกฺขเว เทเสสฺสามิ ชาลินิํ สริตํ วิสฏํ วิสตฺติกํ, ยาย อยํ โลโก อุทฺธสฺโต ปริโยนทฺโธ ตนฺตา\-กุลก\-ชาโต คุลา\-คุณฺฐิก\-ชาโต\csromansharedfootnote{7b5eee9799ca95c2}{กุลาคุณฺฐิกชาโต (สี, สฺยา, กํ, อิ) อํฏฺฐ 2. 386.} มุญฺช\-ปพฺพช\-ภูโต อปายํ ทุคฺคติํ วินิปาตํ สํสารํ นาติวตฺตติ, ตํ สุณาถ สาธุกํ มนสิ กโรถ, ภาสิสฺสามี”ติ. “เอวํ ภนฺเต”ติ โข เต ภิกฺขู ภควโต ปจฺจสฺโสสุํ. ภควา เอตทโวจ\csromandash{}}

\prose{กตมา จ สา ภิกฺขเว ตณฺหา ชาลินี สริตา วิสฏา วิสตฺติกา, ยาย อยํ โลโก อุทฺธสฺโต ปริโยนทฺโธ ตนฺตา\-กุลก\-ชาโต คุลา\-คุณฺฐิก\-ชาโต มุญฺช\-ปพฺพช\-ภูโต อปายํ ทุคฺคติํ วินิปาตํ สํสารํ}

\vspace{\tipitakaparskip}
\csromancenter{(20) 5. มหาวคฺค นาติวตฺตติ? อฏฺฐารส โข ปนิมานิ ภิกฺขเว ตณฺหา\-วิจริตานิ อชฺฌตฺติกสฺส อุปาทาย, อฏฺฐารส ตณฺหา\-วิจริตานิ พาหิรสฺส อุปาทาย.}
\prose{\csromanfolio{533}กตมานิ อฏฺฐารส ตณฺหา\-วิจริตานิ อชฺฌตฺติกสฺส อุปาทาย? “อสฺมี”ติ ภิกฺขเว สติ “อิตฺถสฺมี”ติ โหติ, “เอวํสฺมี”ติ\csromansharedfootnote{b3ee1b3ac07caced}{เอวมสฺมิ (สี), เอวสฺมิ (สฺยา, กํ, อิ) อภิ 2. 406 ปิฏฺเฐปิ ปสฺสิตพฺพํ.} โหติ, “อญฺญถาสฺมี”ติ โหติ, “อสสฺมี”ติ โหติ, “สตสฺมี”ติ โหติ, “สนฺ”ติ โหติ, “อิตฺถํ สนฺ”ติ โหติ, “เอวํ สนฺ”ติ โหติ, “อญฺญถา สนฺ”ติ โหติ, “อปิหํ\csromansharedfootnote{81392adb6e03ba0b}{อปิห (สี, อิ), อปิ (สฺยา, กํ)} สนฺ”ติ โหติ, “อปิหํ\csromansharedfootnote{387e461ed218d5b9}{อปิ (สี, สฺยา, กํ, อิ)} อิตฺถํ สนฺ”ติ โหติ. “อปิหํ\csromansharedfootnote{387e461ed218d5b9}{อปิ (สี, สฺยา, กํ, อิ)} เอวํ สนฺ”ติ โหติ, “อปิหํ\csromansharedfootnote{387e461ed218d5b9}{อปิ (สี, สฺยา, กํ, อิ)} อญฺญถา สนฺ”ติ โหติ, “ภวิสฺสนฺ”ติ โหติ, “อิตฺถํ ภวิสฺสนฺ”ติ โหติ, “เอวํ ภวิสฺสนฺ”ติ โหติ, “อญฺญถา ภวิสฺสนฺ”ติ โหติ. อิมานิ อฏฺฐารส ตณฺหา\-วิจริตานิ อชฺฌตฺติกสฺส อุปาทาย.}

\prose{กตมานิ อฏฺฐารส ตณฺหา\-วิจริตานิ พาหิรสฺส อุปาทาย? “อิมินาสฺมี”ติ ภิกฺขเว สติ “อิมินา อิตฺถสฺมี”ติ โหติ, “อิมินา เอวํสฺมี”ติ โหติ, “อิมินา อญฺญถาสฺมี”ติ โหติ, “อิมินา อสสฺมี”ติ โหติ, “อิมินา สตสฺมี”ติ โหติ, “อิมินา สนฺ”ติ โหติ, อิมินา “อิตฺถํ สนฺ”ติ โหติ, “อิมินา เอวํ สนฺ”ติ โหติ, “อิมินา อญฺญถา สนฺ”ติ โหติ, “อิมินา อปิหํ สนฺ”ติ โหติ, “อิมินา อปิหํ อิตฺถํ สนฺ”ติ โหติ, “อิมินา อปิหํ เอวํ สนฺ”ติ โหติ, อิมินา อปิหํ อญฺญถา สนฺ”ติ โหติ, “อิมินา ภวิสฺสนฺ”ติ โหติ, “อิมินา อิตฺถํ ภวิสฺสนฺ”ติ โหติ, “อิมินา เอวํ ภวิสฺสนฺ”ติ โหติ, “อิมินา อญฺญถา ภวิสฺสนฺ”ติ โหติ. อิมานิ อฏฺฐารส ตณฺหา\-วิจริตานิ พาหิรสฺส อุปาทาย.}

\prose{อิติ อฏฺฐารส ตณฺหา\-วิจริตานิ อชฺฌตฺติกสฺส อุปาทาย, อฏฺฐารส ตณฺหา\-วิจริตานิ พาหิรสฺส อุปาทาย. อิมานิ วุจฺจนฺติ ภิกฺขเว ฉตฺติํส ตณฺหา\-วิจริตานิ, อิติ เอวรูปานิ อตีตานิ ฉตฺติํส ตณฺหา\-วิจริตานิ, อนาคตานิ ฉตฺติํส ตณฺหา\-วิจริตานิ, ปจฺจุปฺปนฺนานิ ฉตฺติํส ตณฺหา\-วิจริตานิ, เอวํ อฏฺฐสตํ ตณฺหาวิจริตํ โหนฺติ.}

\prose{\csromanfolio{534}อยํ โข สา ภิกฺขเว ตณฺหา ชาลินี สริตา วิสฏา วิสตฺติกา, ยาย อยํ โลโก อุทฺธสฺโต ปริโยนทฺโธ ตนฺตา\-กุลก\-ชาโต คุณาคุณฺฐิกชาโต มุญฺช\-ปพฺพช\-ภูโต อปายํ ทุคฺคติํ วินิปาตํ สํสารํ นาติวตฺตตีติ.  นวมํ.}

\csromanmatikaanchor{mtk.239}
\csromantocmarkref{subsection}{10. เปมสุตฺต}{mtk.239}
\bookmark[dest={mtk.239},level=2]{10. เปมสุตฺต}
\csromanheader{2}{10. \textbf{เปมสุตฺต}}
\proseitem{200}{จตฺตาริมานิ ภิกฺขเว (เปมานิ)\csromansharedfootnote{070883e15d77b4d7}{( ) นตฺถิ สีสฺยากํอิโปตฺถเกสุ.} ชายนฺติ. กตมานิ จตฺตาริ? เปมา เปมํ ชายติ, เปมา โทโส ชายติ, โทสา เปมํ ชายติ, โทสา โทโส ชายติ.}

\prose{กถญฺจ ภิกฺขเว เปมา เปมํ ชายติ? อิธ ภิกฺขเว ปุคฺคโล ปุคฺคลสฺส อิฏฺโฐ โหติ กนฺโต มนาโป, ตํ ปเร อิฏฺเฐน กนฺเตน มนาเปน สมุทาจรนฺติ. ตสฺส เอวํ โหติ “โย โข มฺยายํ ปุคฺคโล อิฏฺโฐ กนฺโต มนาโป, ตํ ปเร อิฏฺเฐน กนฺเตน มนาเปน สมุทาจรนฺตี”ติ. โส เตสุ เปมํ ชเนติ. เอวํ โข ภิกฺขเว เปมา เปมํ ชายติ.}

\prose{กถญฺจ ภิกฺขเว เปมา โทโส ชายติ? อิธ ภิกฺขเว ปุคฺคโล ปุคฺคลสฺส อิฏฺโฐ โหติ กนฺโต มนาโป, ตํ ปเร อนิฏฺเฐน อกนฺเตน อมนาเปน สมุทาจรนฺติ. ตสฺส เอวํ โหติ “โย โข มฺยายํ ปุคฺคโล อิฏฺโฐ กนฺโต มนาโป, ตํ ปเร อนิฏฺเฐน อกนฺเตน อมนาเปน สมุทาจรนฺตี”ติ. โส เตสุ โทสํ ชเนติ. เอวํ โข ภิกฺขเว เปมา โทโส ชายติ.}

\prose{กถญฺจ ภิกฺขเว โทสา เปมํ ชายติ? อิธ ภิกฺขเว ปุคฺคโล ปุคฺคลสฺส อนิฏฺโฐ โหติ อกนฺโต อมนาโป, ตํ ปเร อนิฏฺเฐน อกนฺเตน อมนาเปน สมุทาจรนฺติ. ตสฺส เอวํ โหติ “โย โข มฺยายํ ปุคฺคโล อนิฏฺโฐ อกนฺโต อมนาโป, ตํ ปเร อนิฏฺเฐน อกนฺเตน อมนาเปน สมุทาจรนฺตี”ติ. โส เตสุ เปมํ ชเนติ. เอวํ โข ภิกฺขเว โทสา เปมํ ชายติ.}

\prose{กถญฺจ ภิกฺขเว โทสา โทโส ชายติ? อิธ ภิกฺขเว ปุคฺคโล ปุคฺคลสฺส อนิฏฺโฐ โหติ อกนฺโต อมนาโป, ตํ ปเร อิฏฺเฐน กนฺเตน มนาเปน สมุทาจรนฺติ. ตสฺส เอวํ โหติ “โย โข มฺยายํ ปุคฺคโล}

\vspace{\tipitakaparskip}
\csromancenter{(20) 5. มหาวคฺค อนิฏฺโฐ อกนฺโต อมนาโป, ตํ ปเร อิฏฺเฐน กนฺเตน มนาเปน สมุทาจรนฺตี”ติ. โส เตสุ โทสํ ชเนติ. เอวํ โข ภิกฺขเว โทสา โทโส ชายติ. อิมานิ โข ภิกฺขเว จตฺตาริ เปมานิ ชายนฺติ.}
\prose{\csromanfolio{535}ยสฺมิํ ภิกฺขเว สมเย ภิกฺขุ วิวิจฺเจว กาเมหิ ฯเปฯ ปฐมํ ฌานํ อุปสมฺปชฺช วิหรติ. ยมฺปิสฺส เปมา เปมํ ชายติ, ตมฺปิสฺส ตสฺมิํ สมเย น โหติ. โยปิสฺส เปมา โทโส ชายติ, โสปิสฺส ตสฺมิํ สมเย น โหติ. ตมฺปิสฺส โทสา เปมํ ชายติ, ตมฺปิสฺส ตสฺมิํ สมเย น โหติ. โยปิสฺส โทสา โทโส ชายติ, โสปิสฺส ตสฺมิํ สมเย น โหติ.}

\prose{ยสฺมิํ ภิกฺขเว สมเย ภิกฺขุ วิตกฺก\-วิจารานํ วูปสมา ฯเปฯ ทุติยํ ฌานํ ฯเปฯ ตติยํ ฌานํ ฯเปฯ จตุตฺถํ ฌานํ อุปสมฺปชฺช วิหรติ. ยมฺปิสฺส เปมา เปมํ ชายติ, ตมฺปิสฺส ตสฺมิํ สมเย น โหติ. โยปิสฺส เปมา โทโส ชายติ, โสปิสฺส ตสฺมิํ สมเย น โหติ. ยมฺปิสฺส โทสา เปมํ ชายติ, ตมฺปิสฺส ตสฺมิํ สมเย น โหติ. โยปิสฺส โทสา โทโส ชายติ, โสปิสฺส ตสฺมิํ สมเย น โหติ.}

\prose{ยสฺมิํ ภิกฺขเว สมเย ภิกฺขุ อาสวานํ ขยา อนาสวํ เจโตวิมุตฺติํ ปญฺญาวิมุตฺติํ ทิฏฺเฐว ธมฺเม สยํ อภิญฺญา สจฺฉิกตฺวา อุปสมฺปชฺช วิหรติ. ยมฺปิสฺส เปมา เปมํ ชายติ, ตมฺปิสฺส ปหีนํ โหติ อุจฺฉินฺนมูลํ ตาลา\-วตฺถุกตํ อนภาวํกตํ อายติํ อนุปฺปาท\-ธมฺมํ. โยปิสฺส เปมา โทโส ชายติ, โสปิสฺส ปหีโน โหติ อุจฺฉินฺนมูโล ตาลาวตฺถุกโต อนภาวํกโต อายติํ อนุปฺปาทธมฺโม. ยมฺปิสฺส โทสา เปมํ ชายติ, ตมฺปิสฺส ปหีนํ โหติ อุจฺฉินฺนมูลํ ตาลา\-วตฺถุกตํ อนภาวํกตํ อายติํ อนุปฺปาท\-ธมฺมํ. โยปิสฺส โทสา โทโส ชายติ, โสปิสฺส ปหีโน โหติ อุจฺฉินฺนมูโล ตาลาวตฺถุกโต อนภาวํกโต อายติํ อนุปฺปาทธมฺโม. อยํ วุจฺจติ ภิกฺขเว ภิกฺขุ เนว อุสฺเสเนติ น ปฏิเสเนติ\csromansharedfootnote{443c067bd37bbc00}{น ปฏิสฺเสเนติ (สี, อิ)} น ธูปายติ น ปชฺชลติ น สมฺปชฺฌายติ\csromansharedfootnote{0054b0e39e17e545}{น อปชฺฌายติ (สี), น ปชฺฌายติ (สฺยา, กํ, อิ)}.}

\prose{กถญฺจ ภิกฺขเว ภิกฺขุ อุสฺเสเนติ? อิธ ภิกฺขเว ภิกฺขุ รูปํ อตฺตโต สมนุปสฺสติ, รูปวนฺตํ วา อตฺตานํ, อตฺตนิ วา รูปํ, รูปสฺมิํ วา อตฺตานํ. เวทนํ \csromanfolio{536}อตฺตโต สมนุปสฺสติ, เวทนาวนฺตํ วา อตฺตานํ, อตฺตนิ วา เวทนํ, เวทนาย วา อตฺตานํ. สญฺญํ อตฺตโต สมนุปสฺสติ, สญฺญาวนฺตํ วา อตฺตานํ, อตฺตนิ วา สญฺญํ, สญฺญาย วา อตฺตานํ. สงฺขาเร อตฺตโต สมนุปสฺสติ, สงฺขารวนฺตํ วา อตฺตานํ, อตฺตนิ วา สงฺขาเร, สงฺขาเรสุ วา อตฺตานํ. วิญฺญาณํ อตฺตโต สมนุปสฺสติ, วิญฺญาณวนฺตํ วา อตฺตานํ, อตฺตนิ วา วิญฺญาณํ, วิญฺญาณสฺมิํ วา อตฺตานํ. เอวํ โข ภิกฺขเว ภิกฺขุ อุสฺเสเนติ.}

\prose{กถญฺจ ภิกฺขเว ภิกฺขุ น อุสฺเสเนติ? อิธ ภิกฺขเว ภิกฺขุ น รูปํ อตฺตโต สมนุปสฺสติ, น รูปวนฺตํ วา อตฺตานํ, น อตฺตนิ วา รูปํ, น รูปสฺมิํ วา อตฺตานํ. น เวทนํ อตฺตโต สมนุปสฺสติ, น เวทนาวนฺตํ วา อตฺตานํ, น อตฺตนิ วา เวทนํ, น เวทนาย วา อตฺตานํ. น สญฺญํ อตฺตโต สมนุปสฺสติ, น สญฺญาวนฺตํ วา อตฺตานํ, น อตฺตนิ วา สญฺญํ, น สญฺญาย วา อตฺตานํ. น สงฺขาเร อตฺตโต สมนุปสฺสติ, น สงฺขารวนฺตํ วา อตฺตานํ, น อตฺตนิ วา สงฺขาเร, น สงฺขาเรสุ วา อตฺตานํ. น วิญฺญาณํ อตฺตโต สมนุปสฺสติ, น วิญฺญาณวนฺตํ วา อตฺตานํ, น อตฺตนิ วา วิญฺญาณํ, น วิญฺญาณสฺมิํ วา อตฺตานํ. เอวํ โข ภิกฺขเว ภิกฺขุ น อุสฺเสเนติ.}

\prose{กถญฺจ ภิกฺขเว ภิกฺขุ ปฏิเสเนติ? อิธ ภิกฺขเว ภิกฺขุ อกฺโกสนฺตํ ปจฺจกฺโกสติ, โรสนฺตํ ปฏิโรสติ, ภณฺฑนฺตํ ปฏิภณฺฑติ. เอวํ โข ภิกฺขเว ภิกฺขุ ปฏิเสเนติ.}

\prose{กถญฺจ ภิกฺขเว ภิกฺขุ น ปฏิเสเนติ? อิธ ภิกฺขเว ภิกฺขุ อกฺโกสนฺตํ น ปจฺจกฺโกสติ, โรสนฺตํ น ปฏิโรสติ, ภณฺฑนฺตํ น ปฏิภณฺฑติ. เอวํ โข ภิกฺขเว ภิกฺขุ น ปฏิเสเนติ.}

\prose{กถญฺจ ภิกฺขเว ภิกฺขุ ธูปายติ? “อสฺมี”ติ ภิกฺขเว สติ “อิตฺถสฺมี”ติ โหติ, “เอวํสฺมี”ติ โหติ, “อญฺญถาสฺมี”ติ โหติ, “อสสฺมี”ติ โหติ, “สตสฺมี”ติ โหติ, “สนฺ”ติ โหติ, “อิตฺถํ สนฺ”ติ โหติ, “เอวํ สนฺ”ติ โหติ, “อญฺญถา สนฺ”ติ โหติ, “อปิหํ สนฺ”ติ โหติ, “อปิหํ อิตฺถํ สนฺ”ติ โหติ, “อปิหํ เอวํ สนฺ”ติ โหติ, “อปิหํ อญฺญถา สนฺ”ติ โหติ, “ภวิสฺสนฺ”ติ โหติ, “อิตฺถํ ภวิสฺสนฺ”ติ โหติ, “เอวํ ภวิสฺสนฺ”ติ โหติ, “อญฺญถา ภวิสฺสนฺ”ติ โหติ. เอวํ โข ภิกฺขเว ภิกฺขุ ธูปายติ.}

\chapterheadpageread{(20) 5. มหาวคฺค}
\prose{\csromanfolio{537}กถญฺจ ภิกฺขเว ภิกฺขุ น ธูปายติ? “อสฺมี”ติ ภิกฺขเว อสติ “อิตฺถสฺมี”ติ น โหติ, “เอวํสฺมี”ติ น โหติ, “อญฺญถาสฺมี”ติ น โหติ, “อสสฺมี”ติ น โหติ, “สตสฺมี”ติ น โหติ, “สนฺ”ติ น โหติ, “อิตฺถํ สนฺ”ติ น โหติ, “เอวํ สนฺ”ติ น โหติ, “อญฺญถา สนฺ”ติ น โหติ, “อปิหํ สนฺ”ติ น โหติ, “อปิหํ อิตฺถํ สนฺ”ติ น โหติ, “อปิหํ เอวํ สนฺ”ติ น โหติ, “อปิหํ อญฺญถา สนฺ”ติ น โหติ, “ภวิสฺสนฺ”ติ น โหติ, “อิตฺถํ ภวิสฺสนฺ”ติ น โหติ, “เอวํ ภวิสฺสนฺ”ติ น โหติ, “อญฺญถา ภวิสฺสนฺ”ติ น โหติ. เอวํ โข ภิกฺขเว ภิกฺขุ น ธูปายติ.}

\prose{กถญฺจ ภิกฺขเว ภิกฺขุ ปชฺชลติ? “อิมินา อสฺมี”ติ ภิกฺขเว สติ, “อิมินา อิตฺถสฺมี”ติ โหติ, “อิมินา เอวํสฺมี”ติ โหติ, “อิมินา อญฺญถาสฺมี”ติ โหติ, “อิมินา อสสฺมี”ติ โหติ, “อิมินา สตสฺมี”ติ โหติ, “อิมินา สนฺ”ติ โหติ, “อิมินา อิตฺถํ สนฺ”ติ โหติ, “อิมินา เอวํ สนฺ”ติ โหติ, “อิมินา อญฺญถา สนฺ”ติ โหติ, “อิมินา อปิหํ สนฺ”ติ โหติ, “อิมินา อปิหํ อิตฺถํ สนฺ”ติ โหติ, “อิมินา อปิหํ เอวํ สนฺ”ติ โหติ, “อิมินา อปิหํ อญฺญถา สนฺ”ติ โหติ, “อิมินา ภวิสฺสนฺ”ติ โหติ, “อิมินา อิตฺถํ ภวิสฺสนฺ”ติ โหติ, “อิมินา เอวํ ภวิสฺสนฺ”ติ โหติ, “อิมินา อญฺญถา ภวิสฺสนฺ”ติ โหติ. เอวํ โข ภิกฺขเว ภิกฺขุ ปชฺชลติ.}

\prose{กถญฺจ ภิกฺขเว ภิกฺขุ น ปชฺชลติ? “อิมินา อสฺมี”ติ ภิกฺขเว อสติ “อิมินา อิตฺถสฺมี”ติ น โหติ, “อิมินา เอวํสฺมี”ติ น โหติ, “อิมินา อญฺญถาสฺมี”ติ น โหติ, “อิมินา อสสฺมี”ติ น โหติ, “อิมินา สตสฺมี”ติ น โหติ, “อิมินา สนฺ”ติ น โหติ, “อิมินา อิตฺถํ สนฺ”ติ น โหติ, “อิมินา เอวํ สนฺ”ติ น โหติ, “อิมินา อญฺญถา สนฺ”ติ น โหติ, “อิมินา อปิหํ สนฺ”ติ น โหติ, “อิมินา อปิหํ อิตฺถํ สนฺ”ติ น โหติ, “อิมินา อปิหํ เอวํ สนฺ”ติ น โหติ, “อิมินา อปิหํ อญฺญถา สนฺ”ติ น โหติ, “อิมินา ภวิสฺสนฺ”ติ น โหติ, “อิมินา อิตฺถํ ภวิสฺสนฺ”ติ น โหติ, “อิมินา เอวํ ภวิสฺสนฺ”ติ น โหติ, “อิมินา อญฺญถา ภวิสฺสนฺ”ติ น โหติ. เอวํ โข ภิกฺขเว ภิกฺขุ น ปชฺชลติ.}

\prose{\csromanfolio{538}กถญฺจ ภิกฺขเว ภิกฺขุ สมฺปชฺฌายติ? อิธ ภิกฺขเว ภิกฺขุโน อสฺมิมาโน ปหีโน น โหติ อุจฺฉินฺนมูโล ตาลาวตฺถุกโต อนภาวํกโต อายติํ อนุปฺปาทธมฺโม. เอวํ โข ภิกฺขเว ภิกฺขุ สมฺปชฺฌายติ.}

\prosewithclosermidruled{กถญฺจ ภิกฺขเว ภิกฺขุ น สมฺปชฺฌายติ? อิธ ภิกฺขเว ภิกฺขุโน อสฺมิมาโน ปหีโน โหติ อุจฺฉินฺนมูโล ตาลาวตฺถุกโต อนภาวํกโต อายติํ อนุปฺปาทธมฺโม. เอวํ โข ภิกฺขเว ภิกฺขุ น สมฺป\-ชฺฌายตีติ.  ทสมํ.}{มหาวคฺโค ปญฺจโม.}
\csromancenter{\textbf{ตสฺสุทฺทานํ}}
\setlength{\csromangathapagewidth}{0pt}
\setlength{\csromangathawakwidth}{0pt}
\csromangathameasuregroup{โสตานุคตํ ฐานํ, \\ มลฺลิก อตฺตนฺตาโป,}{\csromangathabat{โสตานุคตํ ฐานํ,}{ภทฺทิย สามุคิย วปฺป สาฬฺหา จ.} \\ \csromangathabat{มลฺลิก อตฺตนฺตาโป,}{ตณฺหา เปเมน จ ทสา เตติ.} \\ \textbf{จตุตฺโถ มหาปณฺณาสโก สมตฺโต.}}
\csromangathasetleft{โสตานุคตํ ฐานํ, \\ มลฺลิก อตฺตนฺตาโป,}
\csromangathagroup{\csromangathastanza{\csromangathabat{โสตานุคตํ ฐานํ,}{ภทฺทิย สามุคิย วปฺป สาฬฺหา จ.} \\ \csromangathabat{มลฺลิก อตฺตนฺตาโป,}{ตณฺหา เปเมน จ ทสา เตติ.} \\ \textbf{จตุตฺโถ มหาปณฺณาสโก สมตฺโต.}}}

\csromanmatikaanchor{mtk.240}
\csromantocmarknopage{chapter}{5. ปญฺจมปณฺณาสก}{mtk.240}
\bookmark[dest={mtk.240},level=0]{5. ปญฺจมปณฺณาสก}
\chapterheadpageread{5. ปญฺจม\-ปณฺณาสก}
\csromanmatikaanchor{mtk.241}
\csromantocmarknopage{section}{(21) 1. สปฺปุริสวคฺค}{mtk.241}
\bookmark[dest={mtk.241},level=1]{(21) 1. สปฺปุริสวคฺค}
\csromanmarkh{(21)}\csromanmarkhii{1. สปฺปุริสวคฺค}\csromanheadernomark{1}{\textbf{(21) 1. สปฺปุริสวคฺค}}
\csromanmatikaanchor{mtk.242}
\csromantocmarkref{subsection}{1. สิกฺขาปทสุตฺต}{mtk.242}
\bookmark[dest={mtk.242},level=2]{1. สิกฺขาปทสุตฺต}
\csromanheader{2}{1. สิกฺขาปท\-สุตฺต}
\proseitem{201}{\csromanfolio{539}อสปฺปุริสญฺจ โว ภิกฺขเว เทเสสฺสามิ, อสปฺปุริเสน อสปฺปุริสตรญฺจ, สปฺปุริสญฺจ, สปฺปุริเสน สปฺปุริสตรญฺจ, ตํ สุณาถ สาธุกํ มนสิ กโรถ, ภาสิสฺสามีติ. “เอวํ ภนฺเต”ติ โข เต ภิกฺขู ภควโต ปจฺจสฺโสสุํ. ภควา เอตทโวจ\csromandash{}}

\prose{กตโม จ ภิกฺขเว อสปฺปุริโส? อิธ ภิกฺขเว เอกจฺโจ ปาณาติปาตี โหติ, อทินฺนาทายี โหติ, กาเมสุ\-มิจฺฉาจารี โหติ, มุสาวาที โหติ, สุรา\-เมรย\-มชฺช\-ปมาท\-ฏฺฐายี โหติ. อยํ วุจฺจติ ภิกฺขเว อสปฺปุริโส.}

\prose{กตโม จ ภิกฺขเว อสปฺปุริเสน อสปฺปุริสตโร? อิธ ภิกฺขเว เอกจฺโจ อตฺตนา จ ปาณาติปาตี โหติ, ปรญฺจ ปาณติปาเต สมาทเปติ. อตฺตนา จ อทินฺนาทายี โหติ, ปรญฺจ อทินฺนาทาเน สมาทเปติ. อตฺตนา จ กาเมสุ\-มิจฺฉาจารี โหติ, ปรญฺจ กาเมสุ\-มิจฺฉาจาเร สมาทเปติ. อตฺตนา จ มุสาวาที โหติ, ปรญฺจ มุสาวาเท สมาทเปติ. อตฺตนา จ สุรา\-เมรย\-มชฺช\-ปมาท\-ฏฺฐายี โหติ, ปรญฺจ สุรา\-เมรย\-มชฺช\-ปมาทฏฺฐาเน สมาทเปติ. อยํ วุจฺจติ ภิกฺขเว อสปฺปุริเสน อสปฺปุริสตโร.}

\prose{กตโม จ ภิกฺขเว สปฺปุริโส? อิธ ภิกฺขเว เอกจฺโจ ปาณาติปาตา ปฏิวิรโต โหติ, อทินฺนาทานา ปฏิวิรโต โหติ, กาเมสุ\-มิจฺฉาจารา ปฏิวิรโต โหติ, มุสาวาทา ปฏิวิรโต โหติ, สุรา\-เมรย\-มชฺช\-ปมาทฏฺฐานา ปฏิวิรโต โหติ. อยํ วุจฺจติ ภิกฺขเว สปฺปุริโส.}

\prose{กตโม จ ภิกฺขเว สปฺปุริเสน สปฺปุริสตโร? อิธ ภิกฺขเว เอกจฺโจ อตฺตนา จ ปาณาติปาตา ปฏิวิรโต โหติ, ปรญฺจ ปาณาติปาตา เวรมณิยา สมาทเปติ. อตฺตนา จ อทินฺนาทานา ปฏิวิรโต โหติ, ปรญฺจ อทินฺนาทานา เวรมณิยา สมาทเปติ. \csromanfolio{540}อตฺตนา จ กาเมสุ\-มิจฺฉาจารา ปฏิวิรโต โหติ, ปรญฺจ กาเมสุ\-มิจฺฉาจารา เวรมณิยา สมาทเปติ. อตฺตนา จ มุสาวาทา ปฏิวิรโต โหติ, ปรญฺจ มุสาวาทา เวรมณิยา สมาทเปติ. อตฺตนา จ สุรา\-เมรย\-มชฺช\-ปมาทฏฺฐานา ปฏิวิรโต โหติ, ปรญฺจ สุรา\-เมรย\-มชฺช\-ปมาทฏฺฐานา เวรมณิยา สมาทเปติ. อยํ วุจฺจติ ภิกฺขเว สปฺปุริเสน สปฺปุริสตโรติ*.  ปฏฺฐมํ.}

\csromanmatikaanchor{mtk.243}
\csromantocmarkref{subsection}{2. อสฺสทฺธสุตฺต}{mtk.243}
\bookmark[dest={mtk.243},level=2]{2. อสฺสทฺธสุตฺต}
\csromanheader{2}{2. \textbf{อสฺสทฺธสุตฺต}}
\proseitem{202}{อสปฺปุริสญฺจ โว ภิกฺขเว เทเสสฺสามิ, อสปฺปุริเสน อสปฺปุริสตรญฺจ, สปฺปุริสญฺจ, สปฺปุริเสน สปฺปุริสตรญฺจ, ตํ สุณาถ ฯเปฯ.}

\prose{กตโม จ ภิกฺขเว อสปฺปุริโส? อิธ ภิกฺขเว เอกจฺโจ อสฺสทฺโธ โหติ, อหิริโก โหติ, อโนตฺตปฺปี โหติ, อปฺปสฺสุโต โหติ, กุสีโต โหติ, มุฏฺฐสฺสติ โหติ, ทุปฺปญฺโญ โหติ. อยํ วุจฺจติ ภิกฺขเว อสปฺปุริโส.}

\prose{กตโม จ ภิกฺขเว อสปฺปุริเสน อสปฺปุริสตโร? อิธ ภิกฺขเว เอกจฺโจ อตฺตนา จ อสฺสทฺโธ โหติ, ปรญฺจ อสฺสทฺธิเย\csromansharedfootnote{800dc8ef3693e62e}{อสทฺธาย (ก)} สมาทเปติ. อตฺตนา จ อหิริโก โหติ, ปรญฺจ อหิริกตาย สมาทเปติ. อตฺตนา จ อโนตฺตปฺปี โหติ, ปรญฺจ อโนตฺตปฺเป สมาทเปติ. อตฺตนา จ อปฺปสฺสุโต โหติ, ปรญฺจ อปฺปสฺสุเต สมาทเปติ. อตฺตนา จ กุสีโต โหติ, ปรญฺจ โกสชฺเช สมาทเปติ. อตฺตนา จ มุฏฺฐสฺสติ โหติ, ปรญฺจ มุฏฺฐสฺสจฺเจ\csromansharedfootnote{d09672f024beb340}{มุฏฺฐสจฺเจ (สี, สฺยา, กํ, อิ)} สมาทเปติ. อตฺตนา จ ทุปฺปญฺโญ โหติ, ปรญฺจ ทุปฺปญฺญตาย สมาทเปติ. อยํ วุจฺจติ ภิกฺขเว อสปฺปุริเสน อสปฺปุริสตโร.}

\prose{กตโม จ ภิกฺขเว สปฺปุริโส? อิธ ภิกฺขเว เอกจฺโจ สทฺโธ โหติ, หิริมา โหติ, โอตฺตปฺปี โหติ, พหุสฺสุโต โหติ, อารทฺธวีริโย โหติ, สติมา โหติ, ปญฺญวา โหติ. อยํ วุจฺจติ ภิกฺขเว สปฺปุริโส.}

\prose{กตโม จ ภิกฺขเว สปฺปุริเสน สปฺปุริสตโร? อิธ ภิกฺขเว เอกจฺโจ อตฺตนา จ สทฺธาสมฺปนฺโน โหติ, ปรญฺจ สทฺธา\-สมฺปทาย}

\vspace{\tipitakaparskip}
\csromancenter{(21) 1. สปฺปุริสวคฺค สมาทเปติ. อตฺตนา จ หิริมา โหติ, ปรญฺจ หิริมตาย\csromansharedfootnote{cd23527b89423d63}{หิริสมฺปทาย (ก)} สมาทเปติ. อตฺตนา จ โอตฺตปฺปี โหติ, ปรญฺจ โอตฺตปฺเป สมาทเปติ. อตฺตนา จ พหุสฺสุโต โหติ, ปรญฺจ พาหุสจฺเจ สมาทเปติ. อตฺตนา จ อารทฺธวีริโย โหติ, ปรญฺจ วีริยารมฺเภ สมาทเปติ. อตฺตนา จ อุปฏฺฐิตสฺสติ โหติ, ปรญฺจ สติอุปฏฺฐาเน\csromansharedfootnote{0deeba9fd072b430}{สติปฏฺฐาเน (สี, สฺยา, กํ, อิ)} สมาทเปติ. อตฺตนา จ ปญฺญาสมฺปนฺโน โหติ, ปรญฺจ ปญฺญาสมฺปทาย สมาทเปติ. อยํ วุจฺจติ ภิกฺขเว สปฺปุริเสน สปฺปุริสตโรติ.  ทุติยํ.}
\csromanmatikaanchor{mtk.244}
\csromantocmarkref{subsection}{3. สตฺตกมฺมสุตฺต}{mtk.244}
\bookmark[dest={mtk.244},level=2]{3. สตฺตกมฺมสุตฺต}
\csromanheader{2}{3. สตฺตกมฺม\-สุตฺต}
\proseitem{203}{\csromanfolio{541}อสปฺปุริสญฺจ โว ภิกฺขเว เทเสสฺสามิ, อสปฺปุริเสน อสปฺปุริสตรญฺจ, สปฺปุริสญฺจ, สปฺปุริเสน สปฺปุริสตรญฺจ, ตํ สุณาถ ฯเปฯ.}

\prose{กตโม จ ภิกฺขเว อสปฺปุริโส? อิธ ภิกฺขเว เอกจฺโจ ปาณาติปาตี โหติ, อทินฺนาทายี โหติ, กาเมสุ\-มิจฺฉาจารี โหติ, มุสาวาที โหติ, ปิสุณวาโจ โหติ, ผรุสวาโจ โหติ, สมฺผปฺปลาปี โหติ. อยํ วุจฺจติ ภิกฺขเว อสปฺปุริโส.}

\prose{กตโม จ ภิกฺขเว อสปฺปุริเสน อสปฺปุริสตโร? อิธ ภิกฺขเว เอกจฺโจ อตฺตนา จ ปาณาติปาตี โหติ, ปรญฺจ ปาณติปาเต สมาทเปติ. อตฺตนา จ อทินฺนาทายี โหติ, ปรญฺจ อทินฺนาทาเน สมาทเปติ. อตฺตนา จ กาเมสุ\-มิจฺฉาจารี โหติ, ปรญฺจ กาเมสุ\-มิจฺฉาจาเร สมาทเปติ. อตฺตนา จ มุสาวาที โหติ, ปรญฺจ มุสาวาเท สมาทเปติ. อตฺตนา จ ปิสุณวาโจ โหติ, ปรญฺจ ปิสุณาย วาจาย สมาทเปติ. อตฺตนา จ ผรุสวาโจ โหติ, ปรญฺจ ผรุสาย วาจาย สมาทเปติ. อตฺตนา จ สมฺผปฺปลาปี โหติ, ปรญฺจ สมฺผปฺปลาเป สมาทเปติ. อยํ วุจฺจติ ภิกฺขเว อสปฺปุริเสน อสปฺปุริสตโร.}

\prose{กตโม จ ภิกฺขเว สปฺปุริโส? อิธ ภิกฺขเว เอกจฺโจ ปาณาติปาตา ปฏิวิรโต โหติ, อทินฺนาทานา ปฏิวิรโต โหติ, กาเมสุ\-มิจฺฉาจารา ปฏิวิรโต โหติ, มุสาวาทา ปฏิวิรโต โหติ, ปิสุณาย วาจาย ปฏิวิรโต โหติ, ผรุสาย วาจาย ปฏิวิรโต โหติ, สมฺผปฺปลาปา ปฏิวิรโต โหติ. อยํ วุจฺจติ ภิกฺขเว สปฺปุริโส.}

\prose{\csromanfolio{542}กตโม จ ภิกฺขเว สปฺปุริเสน สปฺปุริสตโร? อิธ ภิกฺขเว เอกจฺโจ อตฺตนา จ ปาณาติปาตา ปฏิวิรโต โหติ, ปรญฺจ ปาณาติปาตา เวรมณิยา สมาทเปติ. อตฺตนา จ อทินฺนาทานา ปฏิวิรโต โหติ, ปรญฺจ อทินฺนาทานา เวรมณิยา สมาทเปติ. อตฺตนา จ กาเมสุ\-มิจฺฉาจารา ปฏิวิรโต โหติ, ปรญฺจ กาเมสุ\-มิจฺฉาจารา เวรมณิยา สมาทเปติ. อตฺตนา จ มุสาวาทา ปฏิวิรโต โหติ, ปรญฺจ มุสาวาทา เวรมณิยา สมาทเปติ. อตฺตนา จ ปิสุณาย วาจาย ปฏิวิรโต โหติ, ปรญฺจ ปิสุณาย วาจาย เวรมณิยา สมาทเปติ. อตฺตนา จ ผรุสาย วาจาย ปฏิวิรโต โหติ, ปรญฺจ ผรุสาย วาจาย เวรมณิยา สมาทเปติ. อตฺตนา จ สมฺผปฺปลาปา ปฏิวิรโต โหติ, ปรญฺจ สมฺผปฺปลาปา เวรมณิยา สมาทเปติ. อยํ วุจฺจติ ภิกฺขเว สปฺปุริเสน สปฺปุริสตโรติ.  ตติยํ.}

\csromanmatikaanchor{mtk.245}
\csromantocmarkref{subsection}{4. ทสกมฺมสุตฺต}{mtk.245}
\bookmark[dest={mtk.245},level=2]{4. ทสกมฺมสุตฺต}
\csromanheader{2}{4. \textbf{ทสกมฺมสุตฺต}}
\proseitem{204}{อสปฺปุริสญฺจ โว ภิกฺขเว เทเสสฺสามิ, อสปฺปุริเสน อสปฺปุริสตรญฺจ, สปฺปุริสญฺจ, สปฺปุริเสน สปฺปุริสตรญฺจ, ตํ สุณาถ ฯเปฯ.}

\prose{กตโม จ ภิกฺขเว อสปฺปุริโส? อิธ ภิกฺขเว เอกจฺโจ ปาณาติปาตี โหติ, อทินฺนาทายี โหติ, กาเมสุ\-มิจฺฉาจารี โหติ, มุสาวาที โหติ, ปิสุณวาโจ โหติ, ผรุสวาโจ โหติ, สมฺผปฺปลาปี โหติ, อภิชฺฌาลุ โหติ, พฺยาปนฺนจิตฺโต โหติ, มิจฺฉาทิฏฺฐิโก โหติ. อยํ วุจฺจติ ภิกฺขเว อสปฺปุริโส.}

\prose{กตโม จ ภิกฺขเว อสปฺปุริเสน อสปฺปุริสตโร? อิธ ภิกฺขเว เอกจฺโจ อตฺตนา จ ปาณาติปาตี โหติ, ปรญฺจ ปาณาติปาเต สมาทเปติ ฯเปฯ อตฺตนา จ อภิชฺฌาลุ โหติ, ปรญฺจ อภิชฺฌาย สมาทเปติ. อตฺตนา จ พฺยาปนฺนจิตฺโต โหติ, ปรญฺจ พฺยาปาเท สมาทเปติ. อตฺตนา จ มิจฺฉาทิฏฺฐิโก โหติ, ปรญฺจ มิจฺฉาทิฏฺฐิยา สมาทเปติ. อยํ วุจฺจติ ภิกฺขเว อสปฺปุริเสน อสปฺปุริสตโร.}

\prose{กตโม จ ภิกฺขเว สปฺปุริโส? อิธ ภิกฺขเว เอกจฺโจ ปาณาติปาตา ปฏิวิรโต โหติ ฯเปฯ อนภิชฺฌาลุ โหติ, อพฺยาปนฺนจิตฺโต โหติ, สมฺมาทิฏฺฐิโก โหติ. อยํ วุจฺจติ ภิกฺขเว สปฺปุริโส.}

\chapterheadpageread{(21) 1. สปฺปุริสวคฺค}
\csromanmatikaanchor{mtk.246}
\csromantocmarkref{subsection}{5. อฏฺฐงฺคิกสุตฺต}{mtk.246}
\bookmark[dest={mtk.246},level=2]{5. อฏฺฐงฺคิกสุตฺต}
\prose{\csromanfolio{543}กตโม จ ภิกฺขเว สปฺปุริเสน สปฺปุริสตโร? อิธ ภิกฺขเว เอกจฺโจ อตฺตนา จ ปาณาติปาตา ปฏิวิรโต โหติ, ปรญฺจ ปาณาติปาตา เวรมณิยา สมาทเปติ ฯเปฯ อตฺตนา จ อนภิชฺฌาลุ โหติ, ปรญฺจ อนภิชฺฌาย สมาทเปติ. อตฺตนา จ อพฺยาปนฺนจิตฺโต โหติ, ปรญฺจ อพฺยาปาเท สมาทเปติ. อตฺตนา จ สมฺมาทิฏฺฐิโก โหติ, ปรญฺจ สมฺมาทิฏฺฐิยา สมาทเปติ. อยํ วุจฺจติ ภิกฺขเว สปฺปุริเสน สปฺปุริสตโรติ.  จตุตฺถํ.}

\csromanheader{2}{5. \textbf{อฏฺฐงฺคีกสุตฺต}}
\proseitem{205}{อสปฺปุริสญฺจ โว ภิกฺขเว เทเสสฺสามิ, อสปฺปุริเสน อสปฺปุริสตรญฺจ, สปฺปุริสญฺจ, สปฺปุริเสน สปฺปุริสตรญฺจ, ตํ สุณาถ ฯเปฯ.}

\prose{กตโม จ ภิกฺขเว อสปฺปุริโส? อิธ ภิกฺขเว เอกจฺโจ มิจฺฉาทิฏฺฐิโก โหติ, มิจฺฉาสงฺกปฺโป โหติ, มิจฺฉาวาโจ โหติ, มิจฺฉากมฺมนฺโต โหติ, มิจฺฉาอาชีโว โหติ, มิจฺฉาวายาโม โหติ, มิจฺฉาสติ โหติ, มิจฺฉาสมาธิ โหติ. อยํ วุจฺจติ ภิกฺขเว อสปฺปุริโส.}

\prose{กตโม จ ภิกฺขเว อสปฺปุริเสน อสปฺปุริสตโร? อิธ ภิกฺขเว เอกจฺโจ อตฺตนา จ มิจฺฉาทิฏฺฐิโก โหติ, ปรญฺจ มิจฺฉาทิฏฺฐิยา สมาทเปติ. อตฺตนา จ มิจฺฉาสงฺกปฺโป โหติ, ปรญฺจ มิจฺฉาสงฺกปฺเป สมาทเปติ. อตฺตนา จ มิจฺฉาวาโจ โหติ, ปรญฺจ มิจฺฉาวาจาย สมาทเปติ. อตฺตนา จ มิจฺฉากมฺมนฺโต โหติ, ปรญฺจ มิจฺฉากมฺมนฺเต สมาทเปติ. อตฺตนา จ มิจฺฉาอาชีโว โหติ, ปรญฺจ มิจฺฉาอาชีเว สมาทเปติ. อตฺตนา จ มิจฺฉาวายาโม โหติ, ปรญฺจ มิจฺฉาวายาเม สมาทเปติ. อตฺตนา จ มิจฺฉาสติ โหติ, ปรญฺจ มิจฺฉาสติยา สมาทเปติ. อตฺตนา จ มิจฺฉาสมาธิ โหติ, ปรญฺจ มิจฺฉา\-สมาธิมฺหิ สมาทเปติ. อยํ วุจฺจติ ภิกฺขเว อสปฺปุริเสน อสปฺปุริสตโร.}

\prose{กตโม จ ภิกฺขเว สปฺปุริโส? อิธ ภิกฺขเว เอกจฺโจ สมฺมาทิฏฺฐิโก โหติ, สมฺมาสงฺกปฺโป โหติ, สมฺมาวาโจ โหติ, สมฺมากมฺมนฺโต โหติ, สมฺมาอาชีโว โหติ, สมฺมาวายาโม โหติ, สมฺมาสติ โหติ, สมฺมาสมาธิ โหติ. อยํ วุจฺจติ ภิกฺขเว สปฺปุริโส.}

\prose{\csromanfolio{544}กตโม จ ภิกฺขเว สปฺปุริเสน สปฺปุริสตโร? อิธ ภิกฺขเว เอกจฺโจ อตฺตนา จ สมฺมาทิฏฺฐิโก โหติ, ปรญฺจ สมฺมาทิฏฺฐิยา สมาทเปติ. อตฺตนา จ สมฺมาสงฺกปฺโป โหติ, ปรญฺจ สมฺมาสงฺกปฺเป สมาทเปติ. อตฺตนา จ สมฺมาวาโจ โหติ, ปรญฺจ สมฺมาวาจาย สมาทเปติ. อตฺตนา จ สมฺมากมฺมนฺโต โหติ, ปรญฺจ สมฺมากมฺมนฺเต สมาทเปติ. อตฺตนา จ สมฺมาอาชีโว โหติ, ปรญฺจ สมฺมาอาชีเว สมาทเปติ. อตฺตนา จ สมฺมาวายาโม โหติ, ปรญฺจ สมฺมาวายาเม สมาทเปติ. อตฺตนา จ สมฺมาสติ โหติ, ปรญฺจ สมฺมาสติยา สมาทเปติ. อตฺตนา จ สมฺมาสมาธิ โหติ, ปรญฺจ สมฺมา\-สมาธิมฺหิ สมาทเปติ. อยํ วุจฺจติ ภิกฺขเว สปฺปุริเสน สปฺปุริสตโรติ.  ปญฺจมํ.}

\csromanmatikaanchor{mtk.247}
\csromantocmarkref{subsection}{6. ทสมคฺคสุตฺต}{mtk.247}
\bookmark[dest={mtk.247},level=2]{6. ทสมคฺคสุตฺต}
\csromanheader{2}{6. \textbf{ทสมคฺคสุตฺต}}
\proseitem{206}{อสปฺปุริสญฺจ โว ภิกฺขเว เทเสสฺสามิ, อสปฺปุริเสน อสปฺปุริสตรญฺจ, สปฺปุริสญฺจ, สปฺปุริเสน สปฺปุริสตรญฺจ, ตํ สุณาถ ฯเปฯ.}

\prose{กตโม จ ภิกฺขเว อสปฺปุริโส? อิธ ภิกฺขเว เอกจฺโจ มิจฺฉาทิฏฺฐิโก โหติ ฯเปฯ มิจฺฉาญาณี โหติ, มิจฺฉาวิมุตฺติ โหติ. อยํ วุจฺจติ ภิกฺขเว อสปฺปุริโส.}

\prose{กตโม จ ภิกฺขเว อสปฺปุริเสน อสปฺปุริสตโร? อิธ ภิกฺขเว เอกจฺโจ อตฺตนา จ มิจฺฉาทิฏฺฐิโก โหติ, ปรญฺจ มิจฺฉาทิฏฺฐิยา สมาทเปติ ฯเปฯ อตฺตนา จ มิจฺฉาญาณี โหติ, ปรญฺจ มิจฺฉาญาเณ สมาทเปติ. อตฺตนา จ มิจฺฉาวิมุตฺติ โหติ, ปรญฺจ มิจฺฉา\-วิมุตฺติยา สมาทเปติ. อยํ วุจฺจติ ภิกฺขเว อสปฺปุริเสน อสปฺปุริสตโร.}

\prose{กตโม จ ภิกฺขเว สปฺปุริโส? อิธ ภิกฺขเว เอกจฺโจ สมฺมาทิฏฺฐิโก โหติ ฯเปฯ สมฺมาญาณี โหติ, สมฺมาวิมุตฺติ โหติ. อยํ วุจฺจติ ภิกฺขเว สปฺปุริโส.}

\prose{กตโม จ ภิกฺขเว สปฺปุริเสน สปฺปุริสตโร? อิธ ภิกฺขเว เอกจฺโจ อตฺตนา จ สมฺมาทิฏฺฐิโก โหติ, ปรญฺจ สมฺมาทิฏฺฐิยา สมาทเปติ ฯเปฯ อตฺตนา จ สมฺมาญาณี โหติ, ปรญฺจ สมฺมาญาเณ สมาทเปติ. อตฺตนา จ สมฺมาวิมุตฺติ โหติ, ปรญฺจ สมฺมาวิมุตฺติยา สมาทเปติ. อยํ วุจฺจติ ภิกฺขเว สปฺปุริเสน สปฺปุริสตโรติ.  ฉฏฺฐํ.}

\chapterheadpageread{(21) 1. สปฺปุริสวคฺค 7. ปฐม\-ปาปธมฺม\-สุตฺต}
\csromanmatikaanchor{mtk.248}
\csromantocmarkref{subsection}{7-10. ปาปธมฺมสุตฺตานิ}{mtk.248}
\bookmark[dest={mtk.248},level=2]{7-10. ปาปธมฺมสุตฺตานิ}
\proseitem{207}{\csromanfolio{545}ปาปญฺจ โว ภิกฺขเว เทเสสฺสามิ, ปาเปน ปาปตรญฺจ, กลฺยาณญฺจ, กลฺยาเณน กลฺยาณตรญฺจ, ตํ สุณาถ ฯเปฯ.}

\prose{กตโม จ ภิกฺขเว ปาโป? อิธ ภิกฺขเว เอกจฺโจ ปาณาติปาตี โหติ ฯเปฯ มิจฺฉาทิฏฺฐิโก โหติ. อยํ วุจฺจติ ภิกฺขเว ปาโป.}

\prose{กตโม จ ภิกฺขเว ปาเปน ปาปตโร? อิธ ภิกฺขเว เอกจฺโจ อตฺตนา จ ปาณาติปาตี โหติ, ปรญฺจ ปาณาติปาเต สมาทเปติ ฯเปฯ อตฺตนา จ มิจฺฉาทิฏฺฐิโก โหติ, ปรญฺจ มิจฺฉาทิฏฺฐิยา สมาทเปติ. อยํ วุจฺจติ ภิกฺขเว ปาเปน ปาปตโร.}

\prose{กตโม จ ภิกฺขเว กลฺยาโณ? อิธ ภิกฺขเว เอกจฺโจ ปาณาติปาตา ปฏิวิรโต โหติ ฯเปฯ สมฺมาทิฏฺฐิโก โหติ. อยํ วุจฺจติ ภิกฺขเว กลฺยาโณ.}

\prose{กตโม จ ภิกฺขเว กลฺยาเณน กลฺยาณตโร? อิธ ภิกฺขเว เอกจฺโจ อตฺตนา จ ปาณาติปาตา ปฏิวิรโต โหติ, ปรญฺจ ปาณาติปาตา เวรมณิยา สมาทเปติ ฯเปฯ อตฺตนา จ สมฺมาทิฏฺฐิโก โหติ, ปรญฺจ สมฺมาทิฏฺฐิยา สมาทเปติ. อยํ วุจฺจติ ภิกฺขเว กลฺยาเณน กลฺยาณตโรติ.  สตฺตมํ.}

\csromanheader{2}{8. ทุติย\-ปาปธมฺม\-สุตฺต}
\proseitem{208}{ปาปญฺจ โว ภิกฺขเว เทเสสฺสามิ, ปาเปน ปาปตรญฺจ, กลฺยาณญฺจ, กลฺยาเณน กลฺยาณตรญฺจ, ตํ สุณาถ ฯเปฯ.}

\prose{กตโม จ ภิกฺขเว ปาโป? อิธ ภิกฺขเว เอกจฺโจ มิจฺฉาทิฏฺฐิโก โหติ ฯเปฯ มิจฺฉาญาณี โหติ, มิจฺฉาวิมุตฺติ โหติ. อยํ วุจฺจติ ภิกฺขเว ปาโป.}

\prose{กตโม จ ภิกฺขเว ปาเปน ปาปตโร? อิธ ภิกฺขเว เอกจฺโจ อตฺตนา จ มิจฺฉาทิฏฺฐิโก โหติ, ปรญฺจ มิจฺฉาทิฏฺฐิยา สมาทเปติ ฯเปฯ อตฺตนา จ มิจฺฉาญาณี โหติ, ปรญฺจ มิจฺฉาญาเณ สมาทเปติ. อตฺตนา จ มิจฺฉาวิมุตฺติ โหติ, ปรญฺจ มิจฺฉา\-วิมุตฺติยา สมาทเปติ. อยํ วุจฺจติ ภิกฺขเว ปาเปน ปาปตโร.}

\prose{\csromanfolio{546}กตโม จ ภิกฺขเว กลฺยาโณ? อิธ ภิกฺขเว เอกจฺโจ สมฺมาทิฏฺฐิโก โหติ ฯเปฯ สมฺมาญาณี โหติ, สมฺมาวิมุตฺติ โหติ. อยํ วุจฺจติ ภิกฺขเว กลฺยาโณ.}

\prose{กตโม จ ภิกฺขเว กลฺยาเณน กลฺยาณตโร? อิธ ภิกฺขเว เอกจฺโจ อตฺตนา จ สมฺมาทิฏฺฐิโก โหติ, ปรญฺจ สมฺมาทิฏฺฐิยา สมาทเปติ ฯเปฯ อตฺตนา จ สมฺมาญาณี โหติ, ปรญฺจ สมฺมาญาเณ สมาทเปติ. อตฺตนา จ สมฺมาวิมุตฺติ โหติ, ปรญฺจ สมฺมาวิมุตฺติยา สมาทเปติ. อยํ วุจฺจติ ภิกฺขเว กลฺยาเณน กลฺยาณตโรติ.  อฏฺฐมํ.}

\csromanheader{2}{9. ตติย\-ปาปธมฺม\-สุตฺต}
\proseitem{209}{ปาปธมฺมญฺจ โว ภิกฺขเว เทเสสฺสามิ, ปาปธมฺเมน ปาป\-ธมฺม\-ตรญฺจ, กลฺยาณธมฺมญฺจ, กลฺยาณ\-ธมฺเมน กลฺยาณ\-ธมฺม\-ตรญฺจ, ตํ สุณาถ ฯเปฯ.}

\prose{กตโม จ ภิกฺขเว ปาปธมฺโม? อิธ ภิกฺขเว เอกจฺโจ ปาณาติปาตี โหติ ฯเปฯ มิจฺฉาทิฏฺฐิโก โหติ. อยํ วุจฺจติ ภิกฺขเว ปาปธมฺโม.}

\prose{กตโม จ ภิกฺขเว ปาปธมฺเมน ปาปธมฺมตโร? อิธ ภิกฺขเว เอกจฺโจ อตฺตนา จ ปาณาติปาตี โหติ, ปรญฺจ ปาณาติปาเต สมาทเปติ ฯเปฯ อตฺตนา จ มิจฺฉาทิฏฺฐิโก โหติ, ปรญฺจ มิจฺฉาทิฏฺฐิยา สมาทเปติ. อยํ วุจฺจติ ภิกฺขเว ปาปธมฺเมน ปาปธมฺมตโร.}

\prose{กตโม จ ภิกฺขเว กลฺยาณธมฺโม? อิธ ภิกฺขเว เอกจฺโจ ปาณติปาตา ปฏิวิรโต โหติ ฯเปฯ สมฺมาทิฏฺฐิโก โหติ. อยํ วุจฺจติ ภิกฺขเว กลฺยาณธมฺโม.}

\prose{กตโม จ ภิกฺขเว กลฺยาณ\-ธมฺเมน กลฺยาณ\-ธมฺม\-ตโร? อิธ ภิกฺขเว เอกจฺโจ อตฺตนา จ ปาณาติปาตา ปฏิวิรโต โหติ, ปรญฺจ ปาณาติปาตา เวรมณิยา สมาทเปติ ฯเปฯ อตฺตนา จ สมฺมาทิฏฺฐิโก โหติ, ปรญฺจ สมฺมาทิฏฺฐิยา สมาทเปติ. อยํ วุจฺจติ ภิกฺขเว กลฺยาณ\-ธมฺเมน กลฺยาณ\-ธมฺม\-ตโรติ.  นวมํ.}

\csromanheader{2}{10. จตุตฺถ\-ปาปธมฺม\-สุตฺต}
\proseitem{210}{ปาปธมฺมญฺจ โว ภิกฺขเว เทเสสฺสามิ, ปาปธมฺเมน ปาป\-ธมฺม\-ตรญฺจ, กลฺยาณธมฺมญฺจ, กลฺยาณ\-ธมฺเมน กลฺยาณ\-ธมฺม\-ตรญฺจ, ตํ สุณาถ ฯเปฯ.}

\csromanmatikaanchor{mtk.249}
\csromantocmarknopage{section}{(22) 2. ปริสาวคฺค}{mtk.249}
\bookmark[dest={mtk.249},level=1]{(22) 2. ปริสาวคฺค}
\csromanmarkh{(22)}\csromanmarkhii{2. ปริสาวคฺค}\csromanheadernomark{1}{(22) 2. ปริสาวคฺค}
\prose{\csromanfolio{547}กตโม จ ภิกฺขเว ปาปธมฺโม? อิธ ภิกฺขเว เอกจฺโจ มิจฺฉาทิฏฺฐิโก โหติ ฯเปฯ มิจฺฉาญาณี โหติ, มิจฺฉาวิมุตฺติ โหติ. อยํ วุจฺจติ ภิกฺขเว ปาปธมฺโม.}

\prose{กตโม จ ภิกฺขเว ปาปธมฺเมน ปาปธมฺมตโร? อิธ ภิกฺขเว เอกจฺโจ อตฺตนา จ มิจฺฉาทิฏฺฐิโก โหติ, ปรญฺจ มิจฺฉาทิฏฺฐิยา สมาทเปติ ฯเปฯ อตฺตนา จ มิจฺฉาญาณี โหติ, ปรญฺจ มิจฺฉาญาเณ สมาทเปติ. อตฺตนา จ มิจฺฉาวิมุตฺติ โหติ, ปรญฺจ มิจฺฉา\-วิมุตฺติยา สมาทเปติ. อยํ วุจฺจติ ภิกฺขเว ปาปธมฺเมน ปาปธมฺมตโร.}

\prose{กตโม จ ภิกฺขเว กลฺยาณธมฺโม? อิธ ภิกฺขเว เอกจฺโจ สมฺมาทิฏฺฐิโก โหติ ฯเปฯ สมฺมาญาณี โหติ, สมฺมาวิมุตฺติ โหติ. อยํ วุจฺจติ ภิกฺขเว กลฺยาณธมฺโม.}

\prosewithclosermidruled{กตโม จ ภิกฺขเว กลฺยาณ\-ธมฺเมน กลฺยาณ\-ธมฺม\-ตโร? อิธ ภิกฺขเว เอกจฺโจ อตฺตนา จ สมฺมาทิฏฺฐิโก โหติ, ปรญฺจ สมฺมาทิฏฺฐิยา สมาทเปติ ฯเปฯ อตฺตนา จ สมฺมาญาณี โหติ, ปรญฺจ สมฺมาญาเณ สมาทเปติ. อตฺตนา จ สมฺมาวิมุตฺติ โหติ, ปรญฺจ สมฺมาวิมุตฺติยา สมาทเปติ. อยํ วุจฺจติ ภิกฺขเว กลฺยาณ\-ธมฺเมน กลฺยาณธมฺมกโรติ.  ทสมํ.}{สปฺปุริสวคฺโค ปฐโม.}
\csromancenter{\textbf{ตสฺสุทฺทานํ}}
\setlength{\csromangathapagewidth}{0pt}
\setlength{\csromangathawakwidth}{0pt}
\csromangathameasuregroup{สิกฺขาปทญฺจ อสฺสทฺธํ, \\ อฏฺฐงฺคิกญฺจ ทสมคฺคํ,}{\csromangathabat{สิกฺขาปทญฺจ อสฺสทฺธํ,}{สตฺตกมฺมํ อโถ จ ทสกมฺมํ.} \\ \csromangathabat{อฏฺฐงฺคิกญฺจ ทสมคฺคํ,}{เทฺว ปาปธมฺโม อปเร เทฺวติ.}}
\csromangathasetleft{สิกฺขาปทญฺจ อสฺสทฺธํ, \\ อฏฺฐงฺคิกญฺจ ทสมคฺคํ,}
\csromangathagroup{\csromangathastanza{\csromangathabat{สิกฺขาปทญฺจ อสฺสทฺธํ,}{สตฺตกมฺมํ อโถ จ ทสกมฺมํ.} \\ \csromangathabat{อฏฺฐงฺคิกญฺจ ทสมคฺคํ,}{เทฺว ปาปธมฺโม อปเร เทฺวติ.}}}

\chapterhead{(22) 2. ปริสาวคฺค}
\csromanmatikaanchor{mtk.250}
\csromantocmarkref{subsection}{1. ปริสาสุตฺต}{mtk.250}
\bookmark[dest={mtk.250},level=2]{1. ปริสาสุตฺต}
\csromanheader{2}{1. \textbf{ปริสาสุตฺต}}
\csromanmatikaanchor{mtk.251}
\csromanmatikaanchor{mtk.252}
\csromantocmarkref{subsection}{2. ทิฏฺฐิสุตฺต}{mtk.251}
\bookmark[dest={mtk.251},level=2]{2. ทิฏฺฐิสุตฺต}
\csromantocmarkref{subsection}{3-10. อกตญฺญุตาสุตฺตาทีนิ}{mtk.252}
\bookmark[dest={mtk.252},level=2]{3-10. อกตญฺญุตาสุตฺตาทีนิ}
\proseitem{211}{จตฺตาโรเม ภิกฺขเว ปริสทูสนา. กตเม จตฺตาโร? ภิกฺขุ ภิกฺขเว\csromansharedfootnote{caf34984a8c1a64c}{อิธ ภิกฺขเว ภิกฺขุ (อิ, ก)} ทุสฺสีโล ปาปธมฺโม ปริสทูสโน, ภิกฺขุนี ภิกฺขเว \csromanfolio{548}ทุสฺสีลา ปาปธมฺมา ปริสทูสนา, อุปาสโก ภิกฺขเว ทุสฺสีโล ปาปธมฺโม ปริสทูสโน, อุปาสิกา ภิกฺขเว ทุสฺสีลา ปาปธมฺมา ปริสทูสนา. อิเม โข ภิกฺขเว จตฺตาโร ปริสทูสนา.}

\prose{จตฺตาโรเม ภิกฺขเว ปริสโสภนา. กตเม จตฺตาโร? ภิกฺขุ ภิกฺขเว สีลวา กลฺยาณธมฺโม ปริสโสภโน, ภิกฺขุนี ภิกฺขเว สีลวตี กลฺยาณธมฺมา ปริสโสภนา, อุปาสโก ภิกฺขเว สีลวา กลฺยาณธมฺโม ปริสโสภโน, อุปาสิกา ภิกฺขเว สีลวตี กลฺยาณธมฺมา ปริสโสภนา. อิเม โข ภิกฺขเว จตฺตาโร ปริสโสภนาติ.  ปฐมํ.}

\csromanheader{2}{2. \textbf{ทิฏฺฐิสุตฺต}}
\proseitem{212}{จตูหิ ภิกฺขเว ธมฺเมหิ สมนฺนาคโต ยถาภตํ นิกฺขิตฺโต เอวํ นิรเย. กตเมหิ จตูหิ? กาย\-ทุจฺจริเตน วจีทุจฺจริเตน มโน\-ทุจฺจริเตน มิจฺฉาทิฏฺฐิยา. อิเมหิ โข ภิกฺขเว จตูหิ ธมฺเมหิ สมนฺนาคโต ยถาภตํ นิกฺขิตฺโต เอวํ นิรเย.}

\prose{จตูหิ ภิกฺขเว ธมฺเมหิ สมนฺนาคโต ยถาภตํ นิกฺขิตฺโต เอวํ สคฺเค. กตเมหิ จตูหิ? กายสุจริเตน วจีสุจริเตน มโนสุจริเตน สมฺมาทิฏฺฐิยา. อิเมหิ โข ภิกฺขเว จตูหิ ธมฺเมหิ สมนฺนาคโต ยถาภตํ นิกฺขิตฺโต เอวํ สคฺเคติ.  ทุติยํ.}

\csromanheader{2}{3. อกตญฺญุตา\-สุตฺต}
\proseitem{213}{จตูหิ ภิกฺขเว ธมฺเมหิ สมนฺนาคโต ยถาภตํ นิกฺขิตฺโต เอวํ นิรเย. กตเมหิ จตูหิ? กาย\-ทุจฺจริเตน, วจีทุจฺจริเตน, มโน\-ทุจฺจริเตน, อกตญฺญุตา อกตเวทิตา\csromansharedfootnote{84064cfdf1df9209}{อกตญฺญุตาอกตเวทิตาย (สี)}. อิเมหิ โข ภิกฺขเว จตูหิ ธมฺเมหิ สมนฺนาคโต ยถาภตํ นิกฺขิตฺโต เอวํ นิรเย.}

\prose{จตูหิ ภิกฺขเว ธมฺเมหิ สมนฺนาคโต ยถาภตํ นิกฺขิตฺโต เอวํ สคฺเค. กตเมหิ จตูหิ? กายสุจริเตน, วจีสุจริเตน, มโนสุจริเตน, กตญฺญุตา กตเวทิตา\csromansharedfootnote{7b39c4183cff366e}{กตญฺญุตากตเวทิตาย (สี)}. อิเมหิ โข ภิกฺขเว จตูหิ ธมฺเมหิ สมนฺนาคโต ยถาภตํ นิกฺขิตฺโต เอวํ สคฺเคติ.  ตติยํ.}

\chapterheadpageread{(22) 2. ปริสาวคฺค 4. ปาณาติปาตี\-สุตฺต}
\proseitem{214}{\csromanfolio{549}ฯเปฯ. ปาณาติปาตี โหติ, อทินฺนาทายี โหติ, กาเมสุ\-มิจฺฉาจารี โหติ, มุสาวาที โหติ ฯเปฯ ปาณาติปาตา ปฏิวิรโต โหติ, อทินฺนาทานา ปฏิวิรโต โหติ, กาเมสุ\-มิจฺฉาจารา ปฏิวิรโต โหติ, มุสาวาทา ปฏิวิรโต โหติ.  จตุตฺถํ.}

\csromanheader{2}{5. ปฐม\-มคฺค\-สุตฺต}
\proseitem{215}{ฯเปฯ. มิจฺฉาทิฏฺฐิโก โหติ, มิจฺฉาสงฺกปฺโป โหติ, มิจฺฉาวาโจ โหติ, มิจฺฉากมฺมนฺโต โหติ ฯเปฯ สมฺมาทิฏฺฐิโก โหติ, สมฺมาสงฺกปฺโป โหติ, สมฺมาวาโจ โหติ, สมฺมากมฺมนฺโต โหติ.  ปญฺจมํ.}

\csromanheader{2}{6. ทุติย\-มคฺค\-สุตฺต}
\proseitem{216}{ฯเปฯ. มิจฺฉาอาชีโว โหติ, มิจฺฉาวายาโม โหติ, มิจฺฉาสติ โหติ, มิจฺฉาสมาธิ โหติ ฯเปฯ สมฺมาอาชีโว โหติ, สมฺมาวายาโม โหติ, สมฺมาสติ โหติ, สมฺมาสมาธิ โหติ.  ฉฏฺฐํ.}

\csromanheader{2}{7. ปฐม\-โวหารปถ\-สุตฺต}
\proseitem{217}{ฯเปฯ. อทิฏฺเฐ ทิฏฺฐวาที โหติ, อสุเต สุตวาที โหติ, อมุเต มุตวาที โหติ, อวิญฺญาเต วิญฺญาตวาที โหติ ฯเปฯ อทิฏฺเฐ อทิฏฺฐวาที โหติ, อสุเต อสุตวาที โหติ, อมุเต อมุตวาที โหติ, อวิญฺญาเต อวิญฺญาตวาที โหติ.  สตฺตมํ.}

\csromanheader{2}{8. ทุติย\-โวหารปถ\-สุตฺต}
\proseitem{218}{ฯเปฯ. ทิฏฺเฐ อทิฏฺฐวาที โหติ, สุเต อสุตวาที โหติ, มุเต อมุตวาที โหติ, วิญฺญาเต อวิญฺญาตวาที โหติ ฯเปฯ ทิฏฺเฐ ทิฏฺฐวาที โหติ, สุเต สุตวาที โหติ, มุเต มุตวาที โหติ, วิญฺญาเต วิญฺญาตวาที โหติ.  อฏฺฐมํ.}

\csromanheader{2}{9. \textbf{อหิริกสุตฺต}}
\proseitem{219}{ฯเปฯ. อสฺสทฺโธ โหติ, ทุสฺสีโล โหติ, อหิริโก โหติ, อโนตฺตปฺปี โหติ ฯเปฯ สทฺโธ โหติ, สีลวา โหติ, หิริมา โหติ, โอตฺตปฺปี โหติ.  นวมํ.}

\csromanpalirecto
\gambhira{จตุกฺกนิปาต\-ปาฬิ 10. ทุสฺสีลสุตฺต}
\proseitem{220}{\csromanfolio{550}จตูหิ ภิกฺขเว ธมฺเมหิ สมนฺนาคโต ยถาภตํ นิกฺขิตฺโต เอวํ นิรเย. กตเมหิ จตูหิ? อสฺสทฺโธ โหติ, ทุสฺสีโล โหติ, กุสีโต โหติ, ทุปฺปญฺโญ โหติ. อิเมหิ โข ภิกฺขเว จตูหิ ธมฺเมหิ สมนฺนาคโต ยถาภตํ นิกฺขิตฺโต เอวํ นิรเย.}

\prosewithclosermidruled{จตูหิ ภิกฺขเว ธมฺเมหิ สมนฺนาคโต ยถาภตํ นิกฺขิตฺโต เอวํ สคฺเค. กตเมหิ จตูหิ? สทฺโธ โหติ, สีลวา โหติ, อารทฺธวีริโย โหติ, ปญฺญวา โหติ. อิเมหิ โข ภิกฺขเว จตูหิ ธมฺเมหิ สมนฺนาคโต ยถาภตํ นิกฺขิตฺโต เอวํ สคฺเคติ.  ทสมํ.}{ปริส\-โสภนวคฺโค\csromansharedfootnote{338dd9d64fd7ee1e}{[มิสฺสินฺคฺ โนเต 0]} ทุติโย.}
\csromancenter{\textbf{ตสฺสุทฺทานํ}}
\setlength{\csromangathapagewidth}{0pt}
\setlength{\csromangathawakwidth}{0pt}
\csromangathameasuregroup{ปริสา ทิฏฺฐิ อกตญฺญุตา, \\ เทฺว โวหารปถา วุตฺตา,}{\csromangathabat{ปริสา ทิฏฺฐิ อกตญฺญุตา,}{ปาณาติปาตาปิ เทฺว มคฺคา.} \\ \csromangathabat{เทฺว โวหารปถา วุตฺตา,}{อหิริกํ ทุปฺปญฺเญน จาติ.}}
\csromangathasetleft{ปริสา ทิฏฺฐิ อกตญฺญุตา, \\ เทฺว โวหารปถา วุตฺตา,}
\csromangathagroup{\csromangathastanza{\csromangathabat{ปริสา ทิฏฺฐิ อกตญฺญุตา,}{ปาณาติปาตาปิ เทฺว มคฺคา.} \\ \csromangathabat{เทฺว โวหารปถา วุตฺตา,}{อหิริกํ ทุปฺปญฺเญน จาติ.}}}

\csromanmatikaanchor{mtk.253}
\csromantocmarknopage{section}{(23) 3. ทุจฺจริตวคฺค}{mtk.253}
\bookmark[dest={mtk.253},level=1]{(23) 3. ทุจฺจริตวคฺค}
\csromanmarkh{(23)}\csromanmarkhii{3. ทุจฺจริตวคฺค}\csromanheadernomark{1}{\textbf{(23) 3. ทุจฺจริตวคฺค}}
\csromanmatikaanchor{mtk.254}
\csromantocmarkref{subsection}{1. ทุจฺจริตสุตฺต}{mtk.254}
\bookmark[dest={mtk.254},level=2]{1. ทุจฺจริตสุตฺต}
\csromanheader{2}{1. \textbf{ทุจฺจริตสุตฺต}}
\proseitem{221}{จตฺตาริมานิ ภิกฺขเว วจีทุจฺจริตานิ. กตมานิ จตฺตาริ? มุสาวาโท, ปิสุณา วาจา, ผรุสา วาจา, สมฺผปฺปลาโป. อิมานิ โข ภิกฺขเว จตฺตาริ วจีทุจฺจริตานิ. จตฺตาริมานิ ภิกฺขเว วจีสุจริตานิ. กตมานิ จตฺตาริ? สจฺจวาจา, อปิสุณา วาจา, สณฺหา วาจา, มนฺตวาจา\csromansharedfootnote{c3cb9f1365d6cdbe}{มนฺตา วาจา (สี, สฺยา, กํ, อิ)}. อิมานิ โข ภิกฺขเว จตฺตาริ วจีสุจรีตานีติ.  ปฐมํ.}

\csromanmatikaanchor{mtk.255}
\csromantocmarkref{subsection}{2. ทิฏฺฐิสุตฺต}{mtk.255}
\bookmark[dest={mtk.255},level=2]{2. ทิฏฺฐิสุตฺต}
\csromanheader{2}{2. \textbf{ทิฏฺฐิสุตฺต}}
\proseitem{222}{จตูหิ ภิกฺขเว ธมฺเมหิ สมนฺนาคโต พาโล อพฺยตฺโต อสปฺปุริโส ขตํ อุปหตํ อตฺตานํ ปริหรติ, สาวชฺโช จ โหติ}

\vspace{\tipitakaparskip}
\csromancenter{(23) 3. ทุจฺจริตวคฺค สานุวชฺโช วิญฺญูนํ, พหุญฺจ อปุญฺญํ ปสวติ. กตเมหิ จตูหิ? กาย\-ทุจฺจริเตน วจีทุจฺจริเตน มโน\-ทุจฺจริเตน มิจฺฉาทิฏฺฐิยา. อิเมหิ โข ภิกฺขเว จตูหิ ธมฺเมหิ สมนฺนาคโต พาโล อพฺยตฺโต อสปฺปุริโส ขตํ อุปหตํ อตฺตานํ ปริหรติ, สาวชฺโช จ โหติ สานุวชฺโช วิญฺญูนํ, พหุญฺจ อปุญฺญํ ปสวติ.}
\csromanmatikaanchor{mtk.256}
\csromantocmarkref{subsection}{3-10. อกตญฺญุตาสุตฺตาทีนิ}{mtk.256}
\bookmark[dest={mtk.256},level=2]{3-10. อกตญฺญุตาสุตฺตาทีนิ}
\prose{\csromanfolio{551}จตูหิ ภิกฺขเว ธมฺเมหิ สมนฺนาคโต ปณฺฑิโต วิยตฺโต สปฺปุริโส อกฺขตํ อนุปหตํ อตฺตานํ ปริหรติ, อนวชฺโช จ โหติ อนนุวชฺโช วิญฺญูนํ, พหุญฺจ ปุญฺญํ ปสวติ. กตเมหิ จตูหิ? กายสุจริเตน วจีสุจริเตน มโนสุจริเตน สมฺมาทิฏฺฐิยา. อิเมหิ โข ภิกฺขเว จตูหิ ธมฺเมหิ สมนฺนาคโต ปณฺฑิโต วิยตฺโต สปฺปุริโส อกฺขตํ อนุปหตํ อตฺตานํ ปริหรติ, อนวชฺโช จ โหติ อนนุวชฺโช วิญฺญูนํ, พหุญฺจ ปุญฺญํ ปสวตีติ.  ทุติยํ.}

\csromanheader{2}{3. อกตญฺญุตา\-สุตฺต}
\proseitem{223}{จตูหิ ภิกฺขเว ธมฺเมหิ สมนฺนาคโต พาโล อพฺยตฺโต อสปฺปุริโส ขตํ อุปหตํ อตฺตานํ ปริหรติ, สาวชฺโช จ โหติ สานุวชฺโช วิญฺญูนํ, พหุญฺจ อปุญฺญํ ปสวติ. กตเมหิ จตูหิ? กาย\-ทุจฺจริเตน, วจีทุจฺจริเตน, มโน\-ทุจฺจริเตน, อกตญฺญุตา อกตเวทิตา อิเมหิ ฯเปฯ กายสุจริเตน วจีสุจริเตน มโนสุจริเตน กตญฺญุตา กตเวทิตา ฯเปฯ.  ตติยํ.}

\csromanheader{2}{4. ปาณาติปาตี\-สุตฺต}
\proseitem{224}{ฯเปฯ. ปาณาติปาตี โหติ, อทินฺนาทายี โหติ, กาเมสุ\-มิจฺฉาจารี โหติ, มุสาวาที โหติ ฯเปฯ ปาณาติปาตา ปฏิวิรโต โหติ, อทินฺนาทานา ปฏิวิรโต โหติ, กาเมสุ\-มิจฺฉาจารา ปฏิวิรโต โหติ, มุสาวาทา ปฏิวิรโต โหติ ฯเปฯ.  จตุตฺถํ.}

\csromanheader{2}{5. ปฐม\-มคฺค\-สุตฺต}
\proseitem{225}{ฯเปฯ. มิจฺฉาทิฏฺฐิโก โหติ, มิจฺฉาสงฺกปฺโป โหติ, มิจฺฉาวาโจ โหติ, มิจฺฉากมฺมนฺโต โหติ ฯเปฯ สมฺมาทิฏฺฐิโก โหติ, สมฺมาสงฺกปฺโป โหติ, สมฺมาวาโจ โหติ, สมฺมากมฺมนฺโต โหติ ฯเปฯ.  ปญฺจมํ.}

\csromanpalirecto
\gambhira{จตุกฺกนิปาต\-ปาฬิ 6. ทุติย\-มคฺค\-สุตฺต}
\prose{\csromanfolio{552}226 ฯเปฯ. มิจฺฉาอาชีโว โหติ, มิจฺฉาวายาโม โหติ, มิจฺฉาสติ โหติ, มิจฺฉาสมาธิ โหติ ฯเปฯ สมฺมาอาชีโว โหติ, สมฺมาวายาโม โหติ, สมฺมาสติ โหติ, สมฺมาสมาธิ โหติ ฯเปฯ.  ฉฏฺฐํ.}

\csromanheader{2}{7. ปฐม\-โวหารปถ\-สุตฺต}
\proseitem{227}{ฯเปฯ. อทิฏฺเฐ ทิฏฺฐวาที โหติ, อสุเต สุตวาที โหติ, อมุเต มุตวาที โหติ, อวิญฺญาเต วิญฺญาตวาที โหติ ฯเปฯ อทิฏฺเฐ อทิฏฺฐวาที โหติ, อสุเต อสุตวาที โหติ, อมุเต อมุตวาที โหติ, อวิญฺญาเต อวิญฺญาตวาที โหติ ฯเปฯ.  สตฺตมํ.}

\csromanheader{2}{8. ทุติย\-โวหารปถ\-สุตฺต}
\proseitem{228}{ฯเปฯ. ทิฏฺเฐ อทิฏฺฐวาที โหติ, สุเต อสุตวาที โหติ, มุเต อมุตวาที โหติ, วิญฺญาเต อวิญฺญาตวาที โหติ ฯเปฯ ทิฏฺเฐ ทิฏฺฐวาที โหติ, สุเต สุตวาที โหติ, มุเต มุตวาที โหติ, วิญฺญาเต วิญฺญาตวาที โหติ ฯเปฯ.  อฏฺฐมํ.}

\csromanheader{2}{9. \textbf{อหิริกสุตฺต}}
\proseitem{229}{ฯเปฯ. อสฺสทฺโธ โหติ, ทุสฺสีโล โหติ, อหิริโก โหติ, อโนตฺตปฺปี โหติ ฯเปฯ สทฺโธ โหติ, สีลวา โหติ, หิริมา โหติ, โอตฺตปฺปี โหติ ฯเปฯ.  นวมํ.}

\csromanheader{2}{10. \textbf{ทุปฺปญฺญสุตฺต}}
\proseitem{230}{ฯเปฯ. อสฺสทฺโธ โหติ, ทุสฺสีโล โหติ, กุสีโต โหติ, ทุปฺปญฺโญ โหติ ฯเปฯ สทฺโธ โหติ, สีลวา โหติ, อารทฺธวีริโย โหติ, ปญฺญวา โหติ. อิเมหิ โข ภิกฺขเว จตูหิ ธมฺเมหิ สมนฺนาคโต ปณฺฑิโต วิยตฺโต สปฺปุริโส อกฺขตํ อนุปหตํ อตฺตานํ ปริหรติ, อนวชฺโช จ โหติ, อนนุวชฺโช วิญฺญูนํ, พหุญฺจ ปุญฺญํ ปสวตีติ.  ทสมํ.}

\csromanmatikaanchor{mtk.257}
\csromantocmarkref{subsection}{11. กวิสุตฺต}{mtk.257}
\bookmark[dest={mtk.257},level=2]{11. กวิสุตฺต}
\csromanmarkh{(24)}\csromanmarkhii{4. กมฺมวคฺค 11. กวิสุตฺต}\csromanheadernomark{2}{(24) 4. กมฺมวคฺค \textbf{11. กวิสุตฺต}}
\proseitemwithclosermidruled{231}{\csromanfolio{553}จตฺตาโรเม ภิกฺขเว กวี. กตเม จตฺตาโร? จินฺตากวิ สุตกวิ อตฺถกวิ ปฏิภานกวิ. อิเม โข ภิกฺขเว จตฺตาโร กวีติ.  เอกาทสมํ.}{ทุจฺจริตวคฺโค ตติโย.}
\csromancenter{\textbf{ตสฺสุทฺทานํ}}
\setlength{\csromangathapagewidth}{0pt}
\setlength{\csromangathawakwidth}{0pt}
\csromangathameasuregroup{ทุจฺจริตํ ทิฏฺฐิ อกตญฺญู จ, \\ เทฺว โวหารปถา วุตฺตา,}{\csromangathabat{ทุจฺจริตํ ทิฏฺฐิ อกตญฺญู จ,}{ปาณาติปาตาปิ เทฺว มคฺคา.} \\ \csromangathabat{เทฺว โวหารปถา วุตฺตา,}{อหิริกํ ทุปฺปญฺญกวินา จาติ.}}
\csromangathasetleft{ทุจฺจริตํ ทิฏฺฐิ อกตญฺญู จ, \\ เทฺว โวหารปถา วุตฺตา,}
\csromangathagroup{\csromangathastanza{\csromangathabat{ทุจฺจริตํ ทิฏฺฐิ อกตญฺญู จ,}{ปาณาติปาตาปิ เทฺว มคฺคา.} \\ \csromangathabat{เทฺว โวหารปถา วุตฺตา,}{อหิริกํ ทุปฺปญฺญกวินา จาติ.}}}

\csromanmatikaanchor{mtk.258}
\csromantocmarknopage{section}{(24) 4. กมฺมวคฺค}{mtk.258}
\bookmark[dest={mtk.258},level=1]{(24) 4. กมฺมวคฺค}
\csromanmarkh{(24)}\csromanmarkhii{4. กมฺมวคฺค}\csromanheadernomark{1}{(24) 4. กมฺมวคฺค}
\csromanmatikaanchor{mtk.259}
\csromantocmarkref{subsection}{1. สํขิตฺตสุตฺต}{mtk.259}
\bookmark[dest={mtk.259},level=2]{1. สํขิตฺตสุตฺต}
\csromanheader{2}{1. \textbf{สํขิตฺตสุตฺต}}
\proseitem{232}{จตฺตาริมานิ ภิกฺขเว กมฺมานิ มยา สยํ อภิญฺญา สจฺฉิกตฺวา ปเวทิตานิ. กตมานิ จตฺตาริ? อตฺถิ ภิกฺขเว กมฺมํ กณฺหํ กณฺหวิปากํ, อตฺถิ ภิกฺขเว กมฺมํ สุกฺกํ สุกฺกวิปากํ, อตฺถิ ภิกฺขเว กมฺมํ กณฺหสุกฺกํ กณฺห\-สุกฺก\-วิปากํ, อตฺถิ ภิกฺขเว กมฺมํ อกณฺหอสุกฺกํ\csromansharedfootnote{15175e2edfa0a222}{อกณฺหํ อสุกฺกํ (สี, สฺยา, อิ) (ที 3. 192; ม 2. 52 ปิฏฺเฐสุปิ)} อกณฺห\-อสุกฺก\-วิปากํ กมฺมกฺขยาย สํวตฺตติ. อิมานิ โข ภิกฺขเว จตฺตาริ กมฺมานิ มยา สยํ อภิญฺญา สจฺฉิกตฺวา ปเวทิตานีติ.  ปฐมํ.}

\csromanmatikaanchor{mtk.260}
\csromantocmarkref{subsection}{2. วิตฺถารสุตฺต}{mtk.260}
\bookmark[dest={mtk.260},level=2]{2. วิตฺถารสุตฺต}
\csromanheader{2}{2. \textbf{วิตฺถารสุตฺต}}
\proseitem{233}{จตฺตาริมานิ ภิกฺขเว กมฺมานิ มยา สยํ อภิญฺญา สจฺฉิกตฺวา ปเวทิตานิ. กตมานิ จตฺตาริ? อตฺถิ ภิกฺขเว กมฺมํ กณฺหํ กณฺหวิปากํ, อตฺถิ ภิกฺขเว กมฺมํ สุกฺกํ สุกฺกวิปากํ, อตฺถิ ภิกฺขเว กมฺมํ กณฺหสุกฺกํ กณฺห\-สุกฺก\-วิปากํ, อตฺถิ ภิกฺขเว กมฺมํ อกณฺหอสุกฺกํ อกณฺห\-สุกฺก\-วิปากํ, กมฺมกฺขยาย สํวตฺตติ.}

\prose{กตมญฺจ ภิกฺขเว กมฺมํ กณฺหํ กณฺหวิปากํ? อิธ ภิกฺขเว เอกจฺโจ สพฺยาพชฺฌํ\csromansharedfootnote{80f04476646516d9}{สพฺยาปชฺฌํ (สพฺพตฺถ)} กายสงฺขารํ อภิสงฺขโรติ, สพฺยาพชฺฌํ วจีสงฺขารํ อภิสงฺขโรติ, \csromanfolio{554}สพฺยาพชฺฌํ มโนสงฺขารํ อภิสงฺขโรติ. โส สพฺยาพชฺฌํ กายสงฺขารํ อภิส\-งฺขริ\-ตฺวา สพฺยาพชฺฌํ วจีสงฺขารํ อภิส\-งฺขริ\-ตฺวา สพฺยาพชฺฌํ มโนสงฺขารํ อภิส\-งฺขริ\-ตฺวา สพฺยาพชฺฌํ โลกํ อุปปชฺชติ, ตเมนํ สพฺยาพชฺฌํ โลกํ อุปปนฺนํ สมานํ สพฺยาพชฺฌา ผสฺสา ผุสนฺติ. โส สพฺยาพชฺเฌหิ ผสฺเสหิ ผุฏฺโฐ สมาโน สพฺยาพชฺฌํ เวทนํ เวทิยติ\csromansharedfootnote{61a1ab55b3e7e32d}{เวทยติ (ก) อํ 2. 360 ปิฏฺเฐปิ.} เอกนฺตทุกฺขํ, เสยฺยถาปิ สตฺตา เนรยิกา. อิทํ วุจฺจติ ภิกฺขเว กมฺมํ กณฺหํ กณฺหวิปากํ.}

\prose{กตมญฺจ ภิกฺขเว กมฺมํ สุกฺกํ สุกฺกวิปากํ? อิธ ภิกฺขเว เอกจฺโจ อพฺยาพชฺฌํ กายสงฺขารํ อภิสงฺขโรติ, อพฺยาพชฺฌํ วจีสงฺขารํ อภิสงฺขโรติ, อพฺยาพชฺฌํ มโนสงฺขารํ อภิสงฺขโรติ. โส อพฺยาพชฺฌํ กายสงฺขารํ อภิส\-งฺขริ\-ตฺวา อพฺยาพชฺฌํ วจีสงฺขารํ อภิส\-งฺขริ\-ตฺวา อพฺยาพชฺฌํ มโนสงฺขารํ อภิส\-งฺขริ\-ตฺวา อพฺยาพชฺฌํ โลกํ อุปปชฺชติ, ตเมนํ อพฺยาพชฺฌํ โลกํ อุปปนฺนํ สมานํ อพฺยาพชฺฌา ผสฺสา ผุสนฺติ. โส อพฺยาพชฺเฌหิ ผสฺเสหิ ผุฏฺโฐ สมาโน อพฺยาพชฺฌํ เวทนํ เวทิยติ เอกนฺตสุขํ, เสยฺยถาปิ เทวา สุภกิณฺหา. อิทํ วุจฺจติ ภิกฺขเว กมฺมํ สุกฺกํ สุกฺกวิปากํ.}

\prose{กตมญฺจ ภิกฺขเว กมฺมํ กณฺหสุกฺกํ กณฺห\-สุกฺก\-วิปากํ? อิธ ภิกฺขเว เอกจฺโจ สพฺยาพชฺฌมฺปิ อพฺยาพชฺฌมฺปิ กายสงฺขารํ อภิสงฺขโรติ, สพฺยาพชฺฌมฺปิ อพฺยาพชฺฌมฺปิ วจีสงฺขารํ อภิสงฺขโรติ, สพฺยาพชฺฌมฺปิ อพฺยาพชฺฌมฺปิ มโนสงฺขารํ อภิสงฺขโรติ. โส สพฺยาพชฺฌมฺปิ อพฺยาพชฺฌมฺปิ กายสงฺขารํ อภิส\-งฺขริ\-ตฺวา สพฺยาพชฺฌมฺปิ อพฺยาพชฺฌมฺปิ วจีสงฺขารํ อภิส\-งฺขริ\-ตฺวา สพฺยาพชฺฌมฺปิ อพฺยาพชฺฌมฺปิ มโนสงฺขารํ อภิส\-งฺขริ\-ตฺวา สพฺยาพชฺฌมฺปิ อพฺยาพชฺฌมฺปิ โลกํ อุปปชฺชติ, ตเมนํ สพฺยาพชฺฌมฺปิ อพฺยาพชฺฌมฺปิ โลกํ อุปปนฺนํ สมานํ สพฺยาพชฺฌาปิ อพฺยาพชฺฌาปิ ผสฺสา ผุสนฺติ. โส สพฺยาพชฺเฌหิปิ อพฺยาพชฺเฌหิปิ ผสฺเสหิ ผุฏฺโฐ สมาโน สพฺยาพชฺฌมฺปิ อพฺยาพชฺฌมฺปิ เวทนํ เวทิยติ โวกิณฺณสุขทุกฺขํ, เสยฺยถาปิ มนุสฺสา เอกจฺเจ จ เทวา เอกจฺเจ จ วินิปาติกา. อิทํ วุจฺจติ ภิกฺขเว กมฺมํ กณฺหสุกฺกํ กณฺห\-สุกฺก\-วิปากํ.}

\prose{กตมญฺจ ภิกฺขเว กมฺมํ อกณฺหอสุกฺกํ อกณฺห\-อสุกฺก\-วิปากํ กมฺมกฺขยาย สํวตฺตติ? ตตฺร ภิกฺขเว ยมิทํ กมฺมํ กณฺหํ กณฺหวิปากํ ตสฺส ปหานาย ยา เจตนา, ยมิทํ\csromansharedfootnote{da45c5f402c1b346}{ยมฺปิทํ (สี, สฺยา, กํ, อิ)} กมฺมํ สุกฺกํ สุกฺกวิปากํ ตสฺส ปหานาย ยา เจตนา, ยมิทํ\csromansharedfootnote{da45c5f402c1b346}{ยมฺปิทํ (สี, สฺยา, กํ, อิ)} กมฺมํ กณฺหสุกฺกํ กณฺห\-สุกฺก\-วิปากํ ตสฺส}

\vspace{\tipitakaparskip}
\csromancenter{(24) 4. กมฺมวคฺค ปหานาย ยา เจตนา. อิทํ วุจฺจติ ภิกฺขเว กมฺมํ อกณฺหอสุกฺกํ อกณฺห\-อสุกฺก\-วิปากํ กมฺมกฺขยาย สํวตฺตติ. อิมานิ โข ภิกฺขเว จตฺตาริ กมฺมานิ มยา สยํ อภิญฺญา สจฺฉิกตฺวา ปเวทิตานีติ.  ทุติยํ.}
\csromanmatikaanchor{mtk.261}
\csromantocmarkref{subsection}{3. โสณกายนสุตฺต}{mtk.261}
\bookmark[dest={mtk.261},level=2]{3. โสณกายนสุตฺต}
\csromanheader{2}{3. โสณกายน\-สุตฺต}
\proseitem{234}{\csromanfolio{555}อถ โข สิขา\-โมคฺคลฺลาโน พฺราหฺมโณ เยน ภควา เตนุปสงฺกมิ, อุปสงฺกมิตฺวา ภควตา สทฺธิํ สมฺโมทิ, สมฺโมทนียํ กถํ สารณียํ วีติสาเรตฺวา เอกมนฺตํ นิสีทิ, เอกมนฺตํ นิสินฺโน โข สิขา\-โมคฺคลฺลาโน พฺราหฺมโณ ภควนฺตํ เอตทโวจ\csromandash{}}

\prose{ปุริมานิ โภ โคตม ทิวสานิ ปุริมตรานิ โสณกายโน มาณโว เยนาหํ เตนุปสงฺกมิ, อุปสงฺกมิตฺวา มํ เอตทโวจ “สมโณ โคตโม สพฺพกมฺมานํ อกิริยํ ปญฺญเปติ, สพฺพกมฺมานํ โข ปน อกิริยํ ปญฺญเปนฺโต อุจฺเฉทํ อาห โลกสฺส. กมฺมสจฺจา’ยํ\csromansharedfootnote{32a3d9bf98ae456d}{กมฺมสจฺจายี (ก)} โภ โลโก กมฺมสมารมฺภ\-ฏฺฐายี”ติ.}

\prose{ทสฺสนมฺปิ โข อหํ พฺราหฺมณ โสณกายนสฺส มาณวสฺส นาภิชานามิ, กุโต ปเนวรูโป กถาสลฺลาโป. จตฺตาริมานิ พฺราหฺมณ กมฺมานิ มยา สยํ อภิญฺญา สจฺฉิกตฺวา ปเวทิตานิ. กตมานิ จตฺตาริ? อตฺถิ พฺราหฺมณ กมฺมํ กณฺหํ กณฺหวิปากํ, อตฺถิ พฺราหฺมณ กมฺมํ สุกฺกํ สุกฺกวิปากํ, อตฺถิ พฺราหฺมณ กมฺมํ กณฺหสุกฺกํ กณฺห\-สุกฺก\-วิปากํ, อตฺถิ พฺราหฺมณ กมฺมํ อกณฺหอสุกฺกํ อกณฺห\-อสุกฺก\-วิปากํ กมฺมกฺขยาย สํวตฺตติ.}

\prose{กตมญฺจ พฺราหฺมณ กมฺมํ กณฺหํ กณฺหวิปากํ? อิธ พฺราหฺมณ เอกจฺโจ สพฺยาพชฺฌํ กายสงฺขารํ อภิสงฺขโรติ, สพฺยาพชฺฌํ วจีสงฺขารํ อภิสงฺขโรติ, สพฺยาพชฺฌํ มโนสงฺขารํ อภิสงฺขโรติ. โส สพฺยาพชฺฌํ กายสงฺขารํ อภิส\-งฺขริ\-ตฺวา สพฺยาพชฺฌํ วจีสงฺขารํ อภิส\-งฺขริ\-ตฺวา สพฺยาพชฺฌํ มโนสงฺขารํ อภิส\-งฺขริ\-ตฺวา สพฺยาพชฺฌํ โลกํ อุปปชฺชติ, ตเมนํ สพฺยาพชฺฌํ โลกํ อุปปนฺนํ สมานํ สพฺยาพชฺฌา ผสฺสา ผุสนฺติ. โส สพฺยาพชฺเฌหิ ผสฺเสหิ ผุฏฺโฐ สมาโน สพฺยาพชฺฌํ เวทนํ เวทิยติ เอกนฺตทุกฺขํ, เสยฺยถาปิ สตฺตา เนรยิกา. อิทํ วุจฺจติ พฺราหฺมณ กมฺมํ กณฺหํ กณฺหวิปากํ.}

\csromanmatikaanchor{mtk.262}
\csromantocmarkref{subsection}{4-9. ปฐมสิกฺขาปทสุตฺตาทีนิ}{mtk.262}
\bookmark[dest={mtk.262},level=2]{4-9. ปฐมสิกฺขาปทสุตฺตาทีนิ}
\prose{\csromanfolio{556}กตมญฺจ พฺราหฺมณ กมฺมํ สุกฺกํ สุกฺกวิปากํ? อิธ พฺราหฺมณ เอกจฺโจ อพฺยาพชฺฌํ กายสงฺขารํ อภิสงฺขโรติ, อพฺยาพชฺฌํ วจีสงฺขารํ อภิสงฺขโรติ, อพฺยาพชฺฌํ มโนสงฺขารํ อภิสงฺขโรติ. โส อพฺยาพชฺฌํ กายสงฺขารํ อภิส\-งฺขริ\-ตฺวา อพฺยาพชฺฌํ วจีสงฺขารํ อภิส\-งฺขริ\-ตฺวา อพฺยาพชฺฌํ มโนสงฺขารํ อภิส\-งฺขริ\-ตฺวา อพฺยาพชฺฌํ โลกํ อุปปชฺชติ, ตเมนํ อพฺยาพชฺฌํ โลกํ อุปปนฺนํ สมานํ อพฺยาพชฺฌา ผสฺสา ผุสนฺติ. โส อพฺยาพชฺเฌหิ ผสฺเสหิ ผุฏฺโฐ สมาโน อพฺยาพชฺฌํ เวทนํ เวทิยติ เอกนฺตสุขํ, เสยฺยถาปิ เทวา สุภกิณฺหา. อิทํ วุจฺจติ พฺราหฺมณ กมฺมํ สุกฺกํ สุกฺกวิปากํ.}

\prose{กตมญฺจ พฺราหฺมณ กมฺมํ กณฺหสุกฺกํ กณฺห\-สุกฺก\-วิปากํ? อิธ พฺราหฺมณ เอกจฺโจ สพฺยาพชฺฌมฺปิ อพฺยาพชฺฌมฺปิ กายสงฺขารํ อภิสงฺขโรติ, สพฺยาพชฺฌมฺปิ อพฺยาพชฺฌมฺปิ วจีสงฺขารํ อภิสงฺขโรติ, สพฺยาพชฺฌมฺปิ อพฺยาพชฺฌมฺปิ มโนสงฺขารํ อภิสงฺขโรติ. โส สพฺยาพชฺฌมฺปิ อพฺยาพชฺฌมฺปิ กายสงฺขารํ อภิส\-งฺขริ\-ตฺวา สพฺยาพชฺฌมฺปิ อพฺยาพชฺฌมฺปิ วจีสงฺขารํ อภิส\-งฺขริ\-ตฺวา สพฺยาพชฺฌมฺปิ อพฺยาพชฺฌมฺปิ มโนสงฺขารํ อภิส\-งฺขริ\-ตฺวา สพฺยาพชฺฌมฺปิ อพฺยาพชฺฌมฺปิ โลกํ อุปปชฺชติ, ตเมนํ สพฺยาพชฺฌมฺปิ อพฺยาพชฺฌมฺปิ โลกํ อุปปนฺนํ สมานํ สพฺยาพชฺฌาปิ อพฺยาพชฺฌาปิ ผสฺสา ผุสนฺติ. โส สพฺยาพชฺเฌหิปิ อพฺยาพชฺเฌหิปิ ผสฺเสหิ ผุฏฺโฐ สมาโน สพฺยาพชฺฌมฺปิ อพฺยาพชฺฌมฺปิ เวทนํ เวทิยติ โวกิณฺณสุขทุกฺขํ, เสยฺยถาปิ มนุสฺสา เอกจฺเจ จ เทวา เอกจฺเจ จ วินิปาติกา. อิทํ วุจฺจติ พฺราหฺมณ กมฺมํ กณฺหสุกฺกํ กณฺห\-สุกฺก\-วิปากํ.}

\prose{กตมญฺจ พฺราหฺมณ กมฺมํ อกณฺหอสุกฺกํ อกณฺห\-อสุกฺก\-วิปากํ กมฺมกฺขยาย สํวตฺตติ? ตตฺร พฺราหฺมณ ยมิทํ กมฺมํ กณฺหํ กณฺหวิปากํ ตสฺส ปหานาย ยา เจตนา, ยมิทํ กมฺมํ สุกฺกํ สุกฺกวิปากํ ตสฺส ปหานาย ยา เจตนา, ยมิทํ กมฺมํ กณฺหสุกฺกํ กณฺห\-สุกฺก\-วิปากํ ตสฺส ปหานาย ยา เจตนา. อิทํ วุจฺจติ พฺราหฺมณ กมฺมํ อกณฺหอสุกฺกํ อกณฺห\-อสุกฺก\-วิปากํ กมฺมกฺขยาย สํวตฺตติ. อิมานิ โข พฺราหฺมณ จตฺตาริ กมฺมานิ มยา สยํ อภิญฺญา สจฺฉิกตฺวา ปเวทิตานีติ.  ตติยํ.}

\csromanheader{2}{4. ปฐม\-สิกฺขาปท\-สุตฺต}
\proseitem{235}{จตฺตาริมานิ ภิกฺขเว กมฺมานิ มยา สยํ อภิญฺญา สจฺฉิกตฺวา ปเวทิตานิ. กตมานิ จตฺตาริ? อตฺถิ ภิกฺขเว กมฺมํ กณฺหํ กณฺหวิปากํ, อตฺถิ ภิกฺขเว กมฺมํ สุกฺกํ สุกฺกวิปากํ, อตฺถิ ภิกฺขเว กมฺมํ กณฺหสุกฺกํ}

\vspace{\tipitakaparskip}
\csromancenter{(24) 4. กมฺมวคฺค กณฺห\-สุกฺก\-วิปากํ, อตฺถิ ภิกฺขเว กมฺมํ อกณฺหอสุกฺกํ อกณฺห\-อสุกฺก\-วิปากํ กมฺมกฺขยาย สํวตฺตติ. กตมญฺจ ภิกฺขเว กมฺมํ กณฺหํ กณฺหวิปากํ? อิธ ภิกฺขเว เอกจฺโจ ปาณาติปาตี โหติ, อทินฺนาทายี โหติ, กาเมสุ มิจฺฉาจารี โหติ, มุสาวาที โหติ, สุรา\-เมรย\-มชฺช\-ปมาท\-ฏฺฐายี โหติ. อิทํ วุจฺจติ ภิกฺขเว กมฺมํ กณฺหํ กณฺหวิปากํ.}
\prose{\csromanfolio{557}กตมญฺจ ภิกฺขเว กมฺมํ สุกฺกํ สุกฺกวิปากํ? อิธ ภิกฺขเว เอกจฺโจ ปาณาติปาตา ปฏิวิรโต โหติ, อทินฺนาทานา ปฏิวิรโต โหติ, กาเมสุ\-มิจฺฉาจารา ปฏิวิรโต โหติ, มุสาวาทา ปฏิวิรโต โหติ, สุรา\-เมรย\-มชฺช\-ปมาทฏฺฐานา ปฏิวิรโต โหติ. อิทํ วุจฺจติ ภิกฺขเว กมฺมํ สุกฺกํ สุกฺกวิปากํ.}

\prose{กตมญฺจ ภิกฺขเว กมฺมํ กณฺหสุกฺกํ กณฺห\-สุกฺก\-วิปากํ? อิธ ภิกฺขเว เอกจฺโจ สพฺยาพชฺฌมฺปิ อพฺยาพชฺฌมฺปิ กายสงฺขารํ อภิสงฺขโรติ ฯเปฯ. อิทํ วุจฺจติ ภิกฺขเว กมฺมํ กณฺหสุกฺกํ กณฺห\-สุกฺก\-วิปากํ.}

\prose{กตมญฺจ ภิกฺขเว กมฺมํ อกณฺหอสุกฺกํ อกณฺห\-อสุกฺก\-วิปากํ กมฺมกฺขยาย สํวตฺตติ? ตตฺร ภิกฺขเว ยมิทํ กมฺมํ กณฺหํ กณฺหวิปากํ ฯเปฯ. อิทํ วุจฺจติ ภิกฺขเว กมฺมํ อกณฺหอสุกฺกํ อกณฺห\-อสุกฺก\-วิปากํ กมฺมกฺขยาย สํวตฺตติ. อิมานิ โข ภิกฺขเว จตฺตาริ กมฺมานิ มยา สยํ อภิญฺญา สจฺฉิกตฺวา ปเวทิตานีติ.  จตุตฺถํ.}

\csromanheader{2}{5. ทุติย\-สิกฺขาปท\-สุตฺต}
\proseitem{236}{จตฺตาริมานิ ภิกฺขเว กมฺมานิ มยา สยํ อภิญฺญา สจฺฉิกตฺวา ปเวทิตานิ. กตมานิ จตฺตาริ? อตฺถิ ภิกฺขเว กมฺมํ กณฺหํ กณฺหวิปากํ, อตฺถิ ภิกฺขเว กมฺมํ สุกฺกํ สุกฺกวิปากํ, อตฺถิ ภิกฺขเว กมฺมํ กณฺหสุกฺกํ กณฺห\-สุกฺก\-วิปากํ, อตฺถิ ภิกฺขเว กมฺมํ อกณฺหอสุกฺกํ อกณฺห\-อสุกฺก\-วิปากํ กมฺมกฺขยาย สํวตฺตติ.}

\prose{กตมญฺจ ภิกฺขเว กมฺมํ กณฺหํ กณฺหวิปากํ? อิธ ภิกฺขเว เอกจฺเจน มาตา\csromansharedfootnote{f04f54715b4ef55f}{เอกจฺโจ มาตรํ (ก)} ชีวิตา โวโรปิตา โหติ, ปิตา\csromansharedfootnote{b07d13de766d1655}{ปิตรํ (ก)} ชีวิตา โวโรปิโต\csromansharedfootnote{e69a6a345b9f4144}{โวโรปิตา (ก)} โหติ, อรหํ\csromansharedfootnote{e0c23389a515895c}{อรหนฺตํ (ก)} ชีวิตา โวโรปิโต\csromansharedfootnote{e69a6a345b9f4144}{โวโรปิตา (ก)} โหติ, ตถาคตสฺส ทุฏฺเฐน จิตฺเตน \csromanfolio{558}โลหิตํ อุปฺปาทิตํ\csromansharedfootnote{c18d26f9aaf57b56}{อุปฺปาทิตา (ก)} โหติ, สํโฆ ภินฺโน โหติ. อิทํ วุจฺจติ ภิกฺขเว กมฺมํ กณฺหํ กณฺหวิปากํ.}

\prose{กตมญฺจ ภิกฺขเว กมฺมํ สุกฺกํ สุกฺกวิปากํ? อิธ ภิกฺขเว เอกจฺโจ ปาณาติปาตา ปฏิวิรโต โหติ, อทินฺนาทานา ปฏิวิรโต โหติ, กาเมสุ\-มิจฺฉาจารา ปฏิวิรโต โหติ, มุสาวาทา ปฏิวิรโต โหติ, ปิสุณาย วาจาย ปฏิวิรโต โหติ, ผรุสาย วาจาย ปฏิวิรโต โหติ, สมฺผปฺปลาปา ปฏิวิรโต โหติ, อนภิชฺฌาลุ โหติ, อพฺยาปนฺนจิตฺโต โหติ, สมฺมาทิฏฺฐิ\csromansharedfootnote{0e7345caaeba3557}{สมฺมาทิฏฺฐิโก (สี, สฺยา, กํ)} โหติ. อิทํ วุจฺจติ ภิกฺขเว กมฺมํ สุกฺกํ สุกฺกวิปากํ.}

\prose{กตมญฺจ ภิกฺขเว กมฺมํ กณฺหสุกฺกํ กณฺห\-สุกฺก\-วิปากํ? อิธ ภิกฺขเว เอกจฺโจ สพฺยาพชฺฌมฺปิ อพฺยาพชฺฌมฺปิ กายสงฺขารํ อภิสงฺขโรติ ฯเปฯ. อิทํ วุจฺจติ ภิกฺขเว กมฺมํ กณฺหสุกฺกํ กณฺห\-สุกฺก\-วิปากํ.}

\prose{กตมญฺจ ภิกฺขเว กมฺมํ อกณฺหอสุกฺกํ อกณฺห\-อสุกฺก\-วิปากํ กมฺมกฺขยาย สํวตฺตติ? ตตฺร ภิกฺขเว ยมิทํ กมฺมํ กณฺหํ กณฺหวิปากํ ฯเปฯ. อิทํ วุจฺจติ ภิกฺขเว กมฺมํ อกณฺหอสุกฺกํ อกณฺห\-อสุกฺก\-วิปากํ กมฺมกฺขยาย สํวตฺตติ. อิมานิ โข ภิกฺขเว จตฺตาริ กมฺมานิ มยา สยํ อภิญฺญา สจฺฉิกตฺวา ปเวทิตานีติ.  ปญฺจมํ.}

\csromanheader{2}{6. อริยมคฺค\-สุตฺต}
\proseitem{237}{จตฺตาริมานิ ภิกฺขเว กมฺมานิ มยา สยํ อภิญฺญา สจฺฉิกตฺวา ปเวทิตานิ. กตมานิ จตฺตาริ? อตฺถิ ภิกฺขเว กมฺมํ กณฺหํ กณฺหวิปากํ, อตฺถิ ภิกฺขเว กมฺมํ สุกฺกํ สุกฺกวิปากํ, อตฺถิ ภิกฺขเว กมฺมํ กณฺหสุกฺกํ กณฺห\-สุกฺก\-วิปากํ, อตฺถิ ภิกฺขเว กมฺมํ อกณฺหอสุกฺกํ อกณฺห\-อสุกฺก\-วิปากํ กมฺมกฺขยาย สํวตฺตติ.}

\prose{กตมญฺจ ภิกฺขเว กมฺมํ กณฺหํ กณฺหวิปากํ? อิธ ภิกฺขเว เอกจฺโจ สพฺยาพชฺฌํ กายสงฺขารํ อภิสงฺขโรติ ฯเปฯ. อิทํ วุจฺจติ ภิกฺขเว กมฺมํ กณฺหํ กณฺหวิปากํ.}

\chapterheadpageread{(24) 4. กมฺมวคฺค}
\prose{\csromanfolio{559}กตมญฺจ ภิกฺขเว กมฺมํ สุกฺกํ สุกฺกวิปากํ? อิธ ภิกฺขเว เอกจฺโจ อพฺยาพชฺฌํ กายสงฺขารํ อภิสงฺขโรติ ฯเปฯ. อิทํ วุจฺจติ ภิกฺขเว กมฺมํ สุกฺกํ สุกฺกวิปากํ.}

\prose{กตมญฺจ ภิกฺขเว กมฺมํ กณฺหสุตฺตํ กณฺห\-สุกฺก\-วิปากํ? อิธ ภิกฺขเว เอกจฺโจ สพฺยาพชฺฌมฺปิ อพฺยาพชฺฌมฺปิ กายสงฺขารํ อภิสงฺขโรติ ฯเปฯ. อิทํ วุจฺจติ ภิกฺขเว กมฺมํ กณฺหสุกฺกํ กณฺห\-สุกฺก\-วิปากํ.}

\prose{กตมญฺจ ภิกฺขเว กมฺมํ อกณฺหอสุกฺกํ อกณฺห\-อสุกฺก\-วิปากํ กมฺมกฺขยาย สํวตฺตติ. สมฺมาทิฏฺฐิ ฯเปฯ สมฺมาสมาธิ. อิทํ วุจฺจติ ภิกฺขเว กมฺมํ อกณฺหอสุกฺกํ อกณฺห\-อสุกฺก\-วิปากํ กมฺมกฺขยาย สํวตฺตติ. อิมานิ โข ภิกฺขเว จตฺตาริ กมฺมานิ มยา สยํ อภิญฺญา สจฺฉิกตฺวา ปเวทิตานีติ.  ฉฏฺฐํ.}

\csromanheader{2}{7. \textbf{โพชฺฌงฺคสุตฺต}}
\proseitem{238}{จตฺตาริมานิ ภิกฺขเว กมฺมานิ ฯเปฯ กณฺหํ กณฺหวิปากํ ฯเปฯ อิธ ภิกฺขเว เอกจฺโจ สพฺยาพชฺฌํ กายสงฺขารํ อภิสงฺขโรติ ฯเปฯ. อิทํ วุจฺจติ ภิกฺขเว กมฺมํ กณฺหํ กณฺหวิปากํ.}

\prose{กตมญฺจ ภิกฺขเว กมฺมํ สุกฺกํ สุกฺกวิปากํ? อิธ ภิกฺขเว เอกจฺโจ อพฺยาพชฺฌํ กายสงฺขารํ อภิสงฺขโรติ ฯเปฯ. อิทํ วุจฺจติ ภิกฺขเว กมฺมํ สุกฺกํ สุกฺกวิปากํ.}

\prose{กตมญฺจ ภิกฺขเว กมฺมํ กณฺหสุกฺกํ กณฺห\-สุกฺก\-วิปากํ? อิธ ภิกฺขเว เอกจฺโจ สพฺยาพชฺฌมฺปิ อพฺยาพชฺฌมฺปิ กายสงฺขารํ อภิสงฺขโรติ ฯเปฯ. อิทํ วุจฺจติ ภิกฺขเว กมฺมํ กณฺหสุกฺกํ กณฺห\-สุกฺก\-วิปากํ.}

\prose{กตมญฺจ ภิกฺขเว กมฺมํ อกณฺหอสุกฺกํ อกณฺห\-อสุกฺก\-วิปากํ กมฺมกฺขยาย สํวตฺตติ? สติสมฺโพชฺฌงฺโค ธมฺมวิจย\-สมฺโพชฺฌงฺโค วีริย\-สมฺโพชฺฌงฺโค ปีติสมฺโพชฺฌงฺโค ปสฺสทฺธิ\-สมฺโพชฺฌงฺโค สมาธิ\-สมฺโพชฺฌงฺโค อุเปกฺขา\-สมฺโพชฺฌงฺโค. อิทํ วุจฺจติ ภิกฺขเว กมฺมํ อกณฺหอสุกฺกํ อกณฺห\-อสุกฺก\-วิปากํ กมฺมกฺขยาย สํวตฺตติ. อิมานิ โข ภิกฺขเว จตฺตาริ กมฺมานิ มยา สยํ อภิญฺญา สจฺฉิกตฺวา ปเวทิตานีติ.  สตฺตมํ.}

\csromanpalirecto
\gambhira{จตุกฺกนิปาต\-ปาฬิ 8. สาวชฺชสุตฺต}
\proseitem{239}{\csromanfolio{560}จตูหิ ภิกฺขเว ธมฺเมหิ สมนฺนาคโต ยถาภตํ นิกฺขิตฺโต เอวํ นิรเย. กตเมหิ จตูหิ? สาวชฺเชน กายกมฺเมน, สาวชฺเชน วจีกมฺเมน, สาวชฺเชน มโนกมฺเมน, สาวชฺชาย ทิฏฺฐิยา. อิเมหิ โข ภิกฺขเว จตูหิ ธมฺเมหิ สมนฺนาคโต ยถาภตํ นิกฺขิตฺโต เอวํ นิรเย.}

\prose{จตูหิ ภิกฺขเว ธมฺเมหิ สมนฺนาคโต ยถาภตํ นิกฺขิตฺโต เอวํ สคฺเค. กตเมหิ จตูหิ? อนวชฺเชน กายกมฺเมน, อนวชฺเชน วจีกมฺเมน, อนวชฺเชน มโนกมฺเมน, อนวชฺชาย ทิฏฺฐิยา. อิเมหิ โข ภิกฺขเว จตูหิ ธมฺเมหิ สมนฺนาคโต ยถาภตํ นิกฺขิตฺโต เอวํ สคฺเคติ.  อฏฺฐมํ.}

\csromanheader{2}{9. อพฺยาพชฺฌ\-สุตฺต}
\proseitem{240}{จตูหิ ภิกฺขเว ธมฺเมหิ สมนฺนาคโต ยถาภตํ นิกฺขิตฺโต เอวํ นิรเย. กตเมหิ จตูหิ? สพฺยาพชฺเฌน กายกมฺเมน, สพฺยาพชฺเฌน วจีกมฺเมน, สพฺยาพชฺเฌน มโนกมฺเมน, สพฺยาพชฺฌาย ทิฏฺฐิยา. อิเมหิ โข ภิกฺขเว จตูหิ ธมฺเมหิ สมนฺนาคโต ยถาภตํ นิกฺขิตฺโต เอวํ นิรเย.}

\prose{จตูหิ ภิกฺขเว ธมฺเมหิ สมนฺนาคโต ยถาภตํ นิกฺขิตฺโต เอวํ สคฺเค. กตเมหิ จตูหิ? อพฺยาพชฺเฌน กายกมฺเมน, อพฺยาพชฺเฌน วจีกมฺเมน, อพฺยาพชฺเฌน มโนกมฺเมน, อพฺยาพชฺฌาย ทิฏฺฐิยา. อิเมหิ โข ภิกฺขเว จตูหิ ธมฺเมหิ สมนฺนาคโต ยถาภตํ นิกฺขิตฺโต เอวํ สคฺเคติ.  นวมํ.}

\csromanmatikaanchor{mtk.263}
\csromantocmarkref{subsection}{10. สมณสุตฺต}{mtk.263}
\bookmark[dest={mtk.263},level=2]{10. สมณสุตฺต}
\csromanheader{2}{10. \textbf{สมณสุตฺต}}
\proseitem{241}{“อิเธว ภิกฺขเว สมโณ, อิธ ทุติโย สมโณ, อิธ ตติโย สมโณ, อิธ จตุตฺโถ สมโณ, สุญฺญา ปรปฺปวาทา สมเณติ อญฺเญหี”ติ\csromansharedfootnote{92fbf1705c1d0d2b}{สมเณหิ อญฺเญติ (สี, อิ, ก) เอตฺถ อญฺเญหีติ สกาย ปฏิญฺญาย สจฺจาภิญฺเญหีติ อตฺโถ เวทิตพฺโพ. ที 2. 125; ม 1. 92 ปิฏฺเฐสุปิ.} เอวเมตํ ภิกฺขเว สมฺมา สีหนาทํ นทถ.}

\prose{กตโม จ ภิกฺขเว ปฐโม สมโณ? อิธ ภิกฺขเว ภิกฺขุ ติณฺณํ สํโยชนานํ ปริกฺขยา โสตาปนฺโน โหติ อวินิปาต\-ธมฺโม นิยโต สมฺโพธิ\-ปรายโณ. อยํ ภิกฺขเว ปฐโม สมโณ.}

\csromanmatikaanchor{mtk.264}
\csromanmatikaanchor{mtk.265}
\csromantocmarkref{subsection}{11. สปฺปุริสานิสํสสุตฺต}{mtk.264}
\bookmark[dest={mtk.264},level=2]{11. สปฺปุริสานิสํสสุตฺต}
\csromantocmarknopage{section}{(25) 5. อาปตฺติภยวคฺค}{mtk.265}
\bookmark[dest={mtk.265},level=1]{(25) 5. อาปตฺติภยวคฺค}
\csromanmarkh{(25)}\csromanmarkhii{5. อาปตฺติภย\-วคฺค}\csromanheadernomark{1}{(25) 5. อาปตฺติภย\-วคฺค}
\prose{\csromanfolio{561}กตโม จ ภิกฺขเว ทุติโย สมโณ? อิธ ภิกฺขเว ภิกฺขุ ติณฺณํ สํโยชนานํ ปริกฺขยา ราค\-โทส\-โมหานํ ตนุตฺตา สกทาคามี โหติ, สกิเทว อิมํ โลกํ อาคนฺตฺวา ทุกฺขสฺสนฺตํ กโรติ. อยํ ภิกฺขเว ทุติโย สมโณ.}

\prose{กตโม จ ภิกฺขเว ตติโย สมโณ? อิธ ภิกฺขเว ภิกฺขุ ปญฺจนฺนํ โอรมฺภาคิยานํ สํโยชนานํ ปริกฺขยา โอปปาติโก โหติ ตตฺถ ปรินิพฺพายี อนาวตฺติธมฺโม ตสฺมา โลกา. อยํ ภิกฺขเว ตติโย สมโณ.}

\prose{กตโม จ ภิกฺขเว จตุตฺโถ สมโณ? อิธ ภิกฺขเว ภิกฺขุ อาสวานํ ขยา อนาสวํ เจโตวิมุตฺติํ ปญฺญาวิมุตฺติํ ทิฏฺเฐว ธมฺเม สยํ อภิญฺญา สจฺฉิกตฺวา อุปสมฺปชฺช วิหรติ. อยํ ภิกฺขเว จตุตฺโถ สมโณ.}

\prose{อิเธว ภิกฺขเว ปฐโม สมโณ, อิธ ทุติโย สมโณ, อิธ ตติโย สมโณ, อิธ จตุตฺโถ สมโณ, สุญฺญา ปรปฺปวาทา สมเณภิ อญฺเญหีติ เอวเมตํ ภิกฺขเว สมฺมา สีหนาทํ นทถาติ.  ทสมํ.}

\csromanheader{2}{11. สปฺปุริสา\-นิสํส\-สุตฺต}
\proseitemwithclosermidruled{242}{สปฺปุริสํ ภิกฺขเว นิสฺสาย จตฺตาโร อานิสํสา ปาฏิกงฺขา. กตเม จตฺตาโร? อริเยน สีเลน วฑฺฒติ, อริเยน สมาธินา วฑฺฒติ, อริยาย ปญฺญาย วฑฺฒติ, อริยาย วิมุตฺติยา วฑฺฒติ. สปฺปุริสํ ภิกฺขเว นิสฺสาย อิเม จตฺตาโร อานิสํสา ปาฏิกงฺขาติ.  เอกาทสมํ.}{กมฺมวคฺโค จตุตฺโถ.}
\csromancenter{\textbf{ตสฺสุทฺทานํ}}
\setlength{\csromangathapagewidth}{0pt}
\setlength{\csromangathawakwidth}{0pt}
\csromangathameasuregroup{}{สํขิตฺต วิตฺถาร โสณกายน, \\ สิกฺขาปทํ อริยมคฺโค โพชฺฌงฺคํ. \\ สาวชฺชญฺเจว อพฺยาพชฺฌํ, \\ สมโณ จ สปฺปุริสานิสํโสติ.}
\csromangathameasurewak{สํขิตฺต วิตฺถาร โสณกายน, \\ สิกฺขาปทํ อริยมคฺโค โพชฺฌงฺคํ. \\ สาวชฺชญฺเจว อพฺยาพชฺฌํ, \\ สมโณ จ สปฺปุริสานิสํโสติ.}
\setlength{\csromangathaleftcol}{0pt}
\csromangathagroup{\csromangathastanza{\csromangathawak{สํขิตฺต วิตฺถาร โสณกายน, \\ สิกฺขาปทํ อริยมคฺโค โพชฺฌงฺคํ. \\ สาวชฺชญฺเจว อพฺยาพชฺฌํ, \\ สมโณ จ สปฺปุริสานิสํโสติ.}}}

\chapterheadpageread{จตุกฺกนิปาต\-ปาฬิ (25) 5. อาปตฺติภย\-วคฺค}
\csromanmatikaanchor{mtk.266}
\csromantocmarkref{subsection}{1. สํฆเภทกสุตฺต}{mtk.266}
\bookmark[dest={mtk.266},level=2]{1. สํฆเภทกสุตฺต}
\csromanheader{2}{1. สํฆเภทก\-สุตฺต}
\proseitem{243}{\csromanfolio{562}เอกํ สมยํ ภควา โกสมฺพิยํ วิหรติ โฆสิตาราเม. อถ โข อายสฺมา อานนฺโท เยน ภควา เตนุปสงฺกมิ, อุปสงฺกมิตฺวา ภควนฺตํ อภิวาเทตฺวา เอกมนฺตํ นิสีทิ. เอกมนฺตํ นิสินฺนํ โข อายสฺมนฺตํ อานนฺทํ ภควา เอตทโวจ “อปิ นุ ตํ อานนฺท อธิกรณํ วูปสนฺตนฺ”ติ. กุโต ตํ ภนฺเต อธิกรณํ วูปสมิสฺสติ\csromansharedfootnote{6f17765c6a60d8bd}{วูปสมฺมิสฺสติ (?)}, อายสฺมโต ภนฺเต อนุรุทฺธสฺส พาหิโย นาม สทฺธิวิหาริโก เกวลกปฺปํ สํฆเภทาย ฐิโต, ตตฺรายสฺมา อนุรุทฺโธ น เอกวาจิกมฺปิ ภณิตพฺพํ มญฺญตีติ.}

\prose{กทา ปนานนฺท อนุรุทฺโธ สํฆมชฺเฌ อธิกรเณสุ\csromansharedfootnote{5ed341d047c57b4f}{อธิกรเณสุ เตสุ (ก)} โวยุญฺชติ, นนุ อานนฺท ยานิ กานิจิ อธิกรณานิ อุปฺปชฺชนฺติ, สพฺพานิ ตานิ ตุเมฺห เจว วูปสเมถ สาริปุตฺต\-โมคฺคลฺลานา จ.}

\prose{จตฺตาโรเม อานนฺท อตฺถวเส สมฺปสฺสมาโน ปาปภิกฺขุ สํฆเภเทน นนฺทติ. กตเม จตฺตาโร? อิธานนฺท ปาปภิกฺขุ ทุสฺสีโล โหติ ปาปธมฺโม อสุจิ สงฺกสฺสร\-สมาจาโร ปฏิจฺฉนฺน\-กมฺมนฺโต อสฺสมโณ สมณปฏิญฺโญ อพฺรหฺมจารี พฺรหฺมจาริ\-ปฏิญฺโญ อนฺโตปูติ อวสฺสุโต กสมฺพุชาโต. ตสฺส เอวํ โหติ “สเจ โข มํ ภิกฺขู ชานิสฺสนฺติ ‘ทุสฺสีโล ปาปธมฺโม อสุจิ สงฺกสฺสร\-สมาจาโร ปฏิจฺฉนฺน\-กมฺมนฺโต อสฺสมโณ สมณปฏิญฺโญ อพฺรหฺมจารี พฺรหฺมจาริ\-ปฏิญฺโญ อนฺโตปูติ อวสฺสุโต กสมฺพุชาโต’ติ. สมคฺคา มํ สนฺตา นาเสสฺสนฺติ, วคฺคา ปน มํ น นาเสสฺสนฺตี”ติ. อิทํ อานนฺท ปฐมํ อตฺถวสํ สมฺปสฺสมาโน ปาปภิกฺขุ สํฆเภเทน นนฺทติ.}

\prose{ปุน จปรํ อานนฺท ปาปภิกฺขุ มิจฺฉาทิฏฺฐิโก โหติ อนฺตคฺคาหิกาย ทิฏฺฐิยา สมนฺนาคโต. ตสฺส เอวํ โหติ “สเจ โข มํ ภิกฺขู ชานิสฺสนฺติ ‘มิจฺฉาทิฏฺฐิโก อนฺตคฺคาหิกาย ทิฏฺฐิยา สมนฺนาคโต’ติ. สมคฺคา มํ สนฺตา นาเสสฺสนฺติ, วคฺคา ปน มํ น นาเสสฺสนฺตี”ติ. อิทํ อานนฺท ทุติยํ อตฺถวสํ สมฺปสฺสมาโน ปาปภิกฺขุ สํฆเภเทน นนฺทติ.}

\chapterheadpageread{(25) 5. อาปตฺติภย\-วคฺค}
\prose{\csromanfolio{563}ปุน จปรํ อานนฺท ปาปภิกฺขุ มิจฺฉาอาชีโว โหติ, มิจฺฉาอาชีเวน ชีวิกํ\csromansharedfootnote{c3b938a2536fef3c}{ชีวิตํ (สฺยา, กํ, อิ, ก)} กปฺเปติ. ตสฺส เอวํ โหติ “สเจ โข มํ ภิกฺขู ชานิสฺสนฺติ ‘มิจฺฉาอาชีโว มิจฺฉาอาชีเวน ชีวิกํ กปฺเปตี’ติ. สมคฺคา มํ สนฺตา นาเสสฺสนฺติ, วคฺคา ปน มํ น นาเสสฺสนฺตี”ติ. อิทํ อานนฺท ตติยํ อตฺถวสํ สมฺปสฺสมาโน ปาปภิกฺขุ สํฆเภเทน นนฺทติ.}

\prose{ปุน จปรํ อานนฺท ปาปภิกฺขุ ลาภกาโม โหติ สกฺการกาโม อนวญฺญตฺติ\-กาโม. ตสฺส เอวํ โหติ “สเจ โข มํ ภิกฺขู ชานิสฺสนฺติ ‘ลาภกาโม สกฺการกาโม อนวญฺญตฺติ\-กาโม’ติ. สมคฺคา มํ สนฺตา น สกฺกริสฺสนฺติ น ครุํ กริสฺสนฺติ น มาเนสฺสนฺติ น ปูเชสฺสนฺติ, วคฺคา ปน มํ สกฺกริสฺสนฺติ ครุํ กริสฺสนฺติ มาเนสฺสนฺติ ปูเชสฺสนฺตี”ติ. อิทํ อานนฺท จตุตฺถํ อตฺถวสํ สมฺปสฺสมาโน ปาปภิกฺขุ สํฆเภเทน นนฺทติ. อิเม โข อานนฺท จตฺตาโร อตฺถวเส สมฺปสฺสมาโน ปาปภิกฺขุ สํฆเภเทน นนฺทตีติ.  ปฐมํ.}

\csromanmatikaanchor{mtk.267}
\csromantocmarkref{subsection}{2. อาปตฺติภยสุตฺต}{mtk.267}
\bookmark[dest={mtk.267},level=2]{2. อาปตฺติภยสุตฺต}
\csromanheader{2}{2. อาปตฺติภย\-สุตฺต}
\proseitem{244}{จตฺตาริมานิ ภิกฺขเว อาปตฺติภยานิ. กตมานิ จตฺตาริ? เสยฺยถาปิ ภิกฺขเว โจรํ อาคุจาริํ คเหตฺวา รญฺโญ ทสฺเสยฺยุํ “อยํ เต เทว โจโร อาคุจารี, อิมสฺส เทโว ทณฺฑํ ปเณตู”ติ. ตเมนํ ราชา เอวํ วเทยฺย “คจฺฉถ โภ อิมํ ปุริสํ ทฬฺหาย รชฺชุยา ปจฺฉาพาหํ คาฬฺหพนฺธนํ พนฺธิตฺวา ขุรมุณฺฑํ กริตฺวา ขรสฺสเรน ปณเวน รถิกาย รถิกํ สิงฺฆาฏเกน สิงฺฆาฏกํ ปริเนตฺวา ทกฺขิเณน ทฺวาเรน นิกฺขาเมตฺวา ทกฺขิณโต นครสฺส สีสํ ฉินฺทถา”ติ. ตเมนํ รญฺโญ ปุริสา ทฬฺหาย รชฺชุยา ปจฺฉาพาหํ คาฬฺหพนฺธนํ พนฺธิตฺวา ขุรมุณฺฑํ กริตฺวา ขรสฺสเรน ปณเวน รถิกาย รถิกํ สิงฺฆาฏเกน สิงฺฆาฏกํ ปริเนตฺวา ทกฺขิเณน ทฺวาเรน นิกฺขาเมตฺวา ทกฺขิณโต นครสฺส สีสํ ฉินฺเทยฺยุํ. ตตฺร\-ญฺญตรสฺส ถลฏฺฐสฺส ปุริสสฺส เอวมสฺส “ปาปกํ วต โภ อยํ ปุริโส กมฺมํ อกาสิ คารยฺหํ สีสจฺเฉชฺชํ, ยตฺร หิ นาม รญฺโญ ปุริสา ทฬฺหาย รชฺชุยา ปจฺฉาพาหํ คาฬฺหพนฺธนํ พนฺธิตฺวา ขุรมุณฺฑํ กริตฺวา ขรสฺสเรน ปณเวน รถิกาย รถิกํ สิงฺฆาฏเกน สิงฺฆาฏกํ ปริเนตฺวา ทกฺขิเณน ทฺวาเรน นิกฺขาเมตฺวา ทกฺขิณโต นครสฺส สีสํ \csromanfolio{564}ฉินฺทิสฺสนฺติ, โส วตสฺสาหํ\csromansharedfootnote{59c7014c26d7ffbb}{โส วตสฺสายํ (สี)} เอวรูปํ ปาปกมฺมํ\csromansharedfootnote{65d3d041fe6c9a76}{ปาปํ กมฺมํ (สี, อิ)} น กเรยฺยํ\csromansharedfootnote{9b08542184c48a6a}{น กเรยฺย (สี) ที 1. 56 ปิฏฺเฐ ปาฬิยา ตทฏฺฐกถาย จ สํสนฺเทตพฺพํ.} คารยฺหํ สีสจฺเฉชฺชนฺ”ติ. เอวเมวํ โข ภิกฺขเว ยสฺส กสฺสจิ ภิกฺขุสฺส วา ภิกฺขุนิยา วา เอวํ ติพฺพา ภยสญฺญา ปจฺจุปฏฺฐิตา โหติ ปาราชิเกสุ ธมฺเมสุ. ตสฺเสตํ ปาฏิกงฺขํ “อนาปนฺโน วา ปาราชิกํ ธมฺมํ น อาปชฺชิสฺสติ, อาปนฺโน วา ปาราชิกํ ธมฺมํ ยถาธมฺมํ ปฏิกริสฺสติ”.}

\prose{เสยฺยถาปิ ภิกฺขเว ปุริโส กาฬวตฺถํ\csromansharedfootnote{b98864d317b3af5c}{กาฬกํ วตฺถํ (สี, สฺยา, กํ, อิ)} ปริธาย เกเส ปกิริตฺวา มุสลํ ขนฺเธ อาโรเปตฺวา มหาชนกายํ อุปสงฺกมิตฺวา เอวํ วเทยฺย “อหํ ภนฺเต ปาปกมฺมํ อกาสิํ คารยฺหํ โมสลฺลํ. เยน เม อายสฺมนฺโต อตฺตมนา โหนฺติ, ตํ กโรมี”ติ. ตตฺร\-ญฺญตรสฺส ถลฏฺฐสฺส ปุริสสฺส เอวมสฺส “ปาปกํ วต โภ อยํ ปุริโส กมฺมํ อกาสิ คารยฺหํ โมสลฺลํ, ยตฺร หิ นาม กาฬวตฺถํ ปริธาย เกเส ปกิริตฺวา มุสลํ ขนฺเธ อาโรเปตฺวา มหาชนกายํ อุปสงฺกมิตฺวา เอวํ วกฺขติ ‘อหํ ภนฺเต ปาปกมฺมํ อกาสิํ คารยฺหํ โมสลฺลํ, เยน เม อายสฺมนฺโต อตฺตมนา โหนฺติ, ตํ กโรมี’ติ. โส วตสฺสาหํ เอวรูปํ ปาปกมฺมํ น กเรยฺยํ คารยฺหํ โมสลฺลนฺ”ติ. เอวเมวํ โข ภิกฺขเว ยสฺส กสฺสจิ ภิกฺขุสฺส วา ภิกฺขุนิยา วา เอวํ ติพฺพา ภยสญฺญา ปจฺจุปฏฺฐิตา โหติ สํฆาทิเสเสสุ ธมฺเมสุ. ตสฺเสตํ ปาฏิกงฺขํ “อนาปนฺโน วา สํฆาทิเสสํ ธมฺมํ น อาปชฺชิสฺสติ, อาปนฺโน วา สํฆาทิเสสํ ธมฺมํ ยถาธมฺมํ ปฏิกริสฺสติ”.}

\prose{เสยฺยถาปิ ภิกฺขเว ปุริโส กาฬวตฺถํ ปริธาย เกเส ปกิริตฺวา ภสฺมปุฏํ\csromansharedfootnote{a4be77ef94a10d9a}{อสฺสปุฏํ (สี, สฺยา, กํ, อิ)} ขนฺเธ อาโรเปตฺวา มหาชนกายํ อุปสงฺกมิตฺวา เอวํ วเทยฺย ‘อหํ ภนฺเต ปาปกมฺมํ อกาสิํ คารยฺหํ ภสฺมปุฏํ, เยน เม อายสฺมนฺโต อตฺตมนา โหนฺติ, ตํ กโรมี”ติ. ตตฺร\-ญฺญตรสฺส ถลฏฺฐสฺส ปุริสสฺส เอวมสฺส “ปาปกํ วต โภ อยํ ปุริโส กมฺมํ อกาสิ คารยฺหํ ภสฺมปุฏํ, ยตฺร หิ นาม กาฬวตฺถํ ปริธาย เกเส ปกิริตฺวา ภสฺมปุฏํ ขนฺเธ อาโรเปตฺวา มหาชนกายํ อุปสงฺกมิตฺวา เอวํ วกฺขติ ‘อหํ ภนฺเต ปาปกมฺมํ อกาสิํ คารยฺหํ ภสฺมปุฏํ, เยน เม อายสฺมนฺโต อตฺตมนา โหนฺติ, ตํ กโรมี’ติ. โส วตสฺสาหํ เอวรูปํ ปาปกมฺมํ น กเรยฺยํ}

\vspace{\tipitakaparskip}
\csromancenter{(25) 5. อาปตฺติภย\-วคฺค คารยฺหํ ภสฺมปุฏนฺ”ติ. เอวเมวํ โข ภิกฺขเว ยสฺส กสฺสจิ ภิกฺขุสฺส วา ภิกฺขุนิยา วา เอวํ ติพฺพา ภยสญฺญา ปจฺจุปฏฺฐิตา โหติ ปาจิตฺติเยสุ ธมฺเมสุ. ตสฺเสตํ ปาฏิกงฺขํ “อนาปนฺโน วา ปาจิตฺติยํ ธมฺมํ น อาปชฺชิสฺสติ, อาปนฺโน วา ปาจิตฺติยํ ธมฺมํ ยถาธมฺมํ ปฏิกริสฺสติ”.}
\prose{\csromanfolio{565}เสยฺยถาปิ ภิกฺขเว ปุริโส กาฬวตฺถํ ปริธาย เกเส ปกิริตฺวา มหาชนกายํ อุปสงฺกมิตฺวา เอวํ วเทยฺย “อหํ ภนฺเต ปาปกมฺมํ อกาสิํ คารยฺหํ อุปวชฺชํ, เยน เม อายสฺมนฺโต อตฺตมนา โหนฺติ, ตํ กโรมี”ติ. ตตฺร\-ญฺญตรสฺส ถลฏฺฐสฺส ปุริสสฺส เอวมสฺส “ปาปกํ วต โภ อยํ ปุริโส กมฺมํ อกาสิ คารยฺหํ อุปวชฺชํ, ยตฺร หิ นาม กาฬวตฺถํ ปริธาย เกเส ปกิริตฺวา มหาชนกายํ อุปสงฺกมิตฺวา เอวํ วกฺขติ ‘อหํ ภนฺเต ปาปกมฺมํ อกาสิํ คารยฺหํ อุปวชฺชํ, เยน เม อายสฺมนฺโต อตฺตมนา โหนฺติ, ตํ กโรมี’ติ. โส วตสฺสาหํ เอวรูปํ ปาปกมฺมํ น กเรยฺยํ คารยฺหํ อุปวชฺชนฺ”ติ. เอวเมวํ โข ภิกฺขเว ยสฺส กสฺสจิ ภิกฺขุสฺส วา ภิกฺขุนิยา วา เอวํ ติพฺพา ภยสญฺญา ปจฺจุปฏฺฐิตา โหติ ปาฏิเทสนีเยสุ ธมฺเมสุ. ตสฺเสตํ ปาฏิกงฺขํ “อนาปนฺโน วา ปาฏิเทสนียํ ธมฺมํ น อาปชฺชิสฺสติ, อาปนฺโน วา ปาฏิเทสนียํ ธมฺมํ ยถาธมฺมํ ปฏิกริสฺสติ”. อิมานิ โข ภิกฺขเว จตฺตาริ อาปตฺติภยานีติ.  ทุติยํ.}

\csromanmatikaanchor{mtk.268}
\csromantocmarkref{subsection}{3. สิกฺขานิสํสสุตฺต}{mtk.268}
\bookmark[dest={mtk.268},level=2]{3. สิกฺขานิสํสสุตฺต}
\csromanheader{2}{3. สิกฺขานิสํส\-สุตฺต}
\proseitem{245}{สิกฺขานิสํสมิทํ ภิกฺขเว พฺรหฺมจริยํ วุสฺสติ ปญฺญุตฺตรํ วิมุตฺติสารํ สตา\-ธิปเตยฺยํ. กถญฺจ ภิกฺขเว สิกฺขานิสํสํ โหติ? อิธ ภิกฺขเว มยา สาวกานํ อาภิสมาจาริกา สิกฺขา ปญฺญตฺตา อปฺปสนฺนานํ ปสาทาย ปสนฺนานํ ภิยฺโยภาวาย. ยถา ยถา ภิกฺขเว มยา สาวกานํ อาภิสมาจาริกา สิกฺขา ปญฺญตฺตา อปฺปสนฺนานํ ปสาทาย ปสนฺนานํ ภิยฺโยภาวาย, ตถา ตถา โส ตสฺสา สิกฺขาย อขณฺฑการี โหติ อจฺฉิทฺทการี อสพลการี อกมฺมาสการี, สมาทาย สิกฺขติ สิกฺขาปเทสุ.}

\prose{ปุน จปรํ ภิกฺขเว มยา สาวกานํ อาทิพฺราหฺมจริยิกา สิกฺขา ปญฺญตฺตา สพฺพโส สมฺมา ทุกฺขกฺขยาย. ยถา ยถา ภิกฺขเว มยา สาวกานํ อาทิ\-พฺรหฺม\-จริยิกา สิกฺขา ปญฺญตฺตา สพฺพโส สมฺมา ทุกฺขกฺขยาย, ตถา ตถา โส ตสฺสา สิกฺขาย อขณฺฑการี \csromanfolio{566}โหติ อจฺฉิทฺทการี อสพลการี อกมฺมาสการี, สมาทาย สิกฺขติ สิกฺขาปเทสุ. เอวํ โข ภิกฺขเว สิกฺขานิสํสํ โหติ.}

\prose{กถญฺจ ภิกฺขเว ปญฺญุตฺตรํ โหติ? อิธ ภิกฺขเว มยา สาวกานํ ธมฺมา เทสิตา สพฺพโส สมฺมา ทุกฺขกฺขยาย. ยถา ยถา ภิกฺขเว มยา สาวกานํ ธมฺมา เทสิตา สพฺพโส สมฺมา ทุกฺขกฺขยาย, ตถา ตถาสฺส เต ธมฺมา ปญฺญาย สมเวกฺขิตา โหนฺติ. เอวํ โข ภิกฺขเว ปญฺญุตฺตรํ โหติ.}

\prose{กถญฺจ ภิกฺขเว วิมุตฺติสารํ โหติ? อิธ ภิกฺขเว มยา สาวกานํ ธมฺมา เทสิตา สพฺพโส สมฺมา ทุกฺขกฺขยาย. ยถา ยถา ภิกฺขเว มยา สาวกานํ ธมฺมา เทสิตา สพฺพโส สมฺมา ทุกฺขกฺขยาย, ตถา ตถาสฺส เต ธมฺมา วิมุตฺติยา ผุสิตา โหนฺติ. เอวํ โข ภิกฺขเว วิมุตฺติสารํ โหติ.}

\prose{กถญฺจ ภิกฺขเว สตา\-ธิปเตยฺยํ โหติ? “อิติ อปริปูรํ วา อาภิสมาจาริกํ สิกฺขํ ปริปูเรสฺสามิ, ปริปูรํ วา อาภิสมาจาริกํ สิกฺขํ ตตฺถ ตตฺถ ปญฺญาย อนุคฺคเหสฺสามี”ติ อชฺฌตฺตํเยว สติ สูปฏฺฐิตา โหติ. “อิติ อปริปูรํ วา อาทิ\-พฺรหฺม\-จริยิกํ สิกฺขํ ปริปูเรสฺสามิ, ปริปูรํ วา อาทิ\-พฺรหฺม\-จริยิกํ สิกฺขํ ตตฺถ ตตฺถ ปญฺญาย อนุคฺคเหสฺสามี”ติ อชฺฌตฺตํเยว สติ สูปฏฺฐิตา โหติ. “อิติ อสมเวกฺขิตํ วา ธมฺมํ ปญฺญาย สมเวกฺขิสฺสามิ, สมเวกฺขิตํ วา ธมฺมํ ตตฺถ ตตฺถ ปญฺญาย อนุคฺคเหสฺสามี”ติ อชฺฌตฺตํเยว สติ สูปฏฺฐิตา โหติ. “อิติ อผุสิตํ วา ธมฺมํ วิมุตฺติยา ผุสิสฺสามิ, ผุสิตํ วา ธมฺมํ ตตฺถ ตตฺถ ปญฺญาย อนุคฺคเหสฺสามี”ติ อชฺฌตฺตํเยว สติ สูปฏฺฐิตา โหติ. เอวํ โข ภิกฺขเว สตา\-ธิปเตยฺยํ โหติ. “สิกฺขานิสํสมิทํ ภิกฺขเว พฺรหฺมจริยํ วุสฺสติ ปญฺญุตฺตรํ วิมุตฺติสารํ สตาธิปเตยฺยนฺ”ติ อิติ ยํ ตํ วุตฺตํ, อิทเมตํ ปฏิจฺจ วุตฺตนฺติ.  ตติยํ.}

\csromanmatikaanchor{mtk.269}
\csromantocmarkref{subsection}{4. เสยฺยาสุตฺต}{mtk.269}
\bookmark[dest={mtk.269},level=2]{4. เสยฺยาสุตฺต}
\csromanheader{2}{4. \textbf{เสยฺยาสุตฺต}}
\proseitem{246}{จตสฺโส อิมา ภิกฺขเว เสยฺยา. กตมา จตสฺโส? เปตเสยฺยา กามโภคิเสยฺยา สีหเสยฺยา ตถาคตเสยฺยา.}

\vspace{\tipitakaparskip}
\csromancenter{(25) 5. อาปตฺติภย\-วคฺค กตมา จ ภิกฺขเว เปตเสยฺยา? เยภุยฺเยน ภิกฺขเว เปตา อุตฺตานา เสนฺติ. อยํ วุจฺจติ ภิกฺขเว เปตเสยฺยา.}
\csromanmatikaanchor{mtk.270}
\csromantocmarkref{subsection}{5. ถุปารหสุตฺต}{mtk.270}
\bookmark[dest={mtk.270},level=2]{5. ถุปารหสุตฺต}
\prose{\csromanfolio{567}กตมา จ ภิกฺขเว กามโภคิเสยฺยา? เยภุยฺเยน ภิกฺขเว กามโภคี วาเมน ปสฺเสน เสนฺติ. อยํ วุจฺจติ ภิกฺขเว กามโภคิเสยฺยา.}

\prose{กตมา จ ภิกฺขเว สีหเสยฺยา? สีโห ภิกฺขเว มิคราชา ทกฺขิเณน ปสฺเสน เสยฺยํ กปฺเปติ ปาเท ปาทํ อจฺจาธาย อนฺตร\-สตฺถิมฺหิ นงฺคุฏฺฐํ อนุปกฺขิปิตฺวา, โส ปฏิพุชฺฌิตฺวา ปุริมํ กายํ อพฺภุนฺนาเมตฺวา ปจฺฉิมํ กายํ อนุวิโลเกติ. สเจ ภิกฺขเว สีโห มิคราชา กิญฺจิ ปสฺสติ กายสฺส วิกฺขิตฺตํ วา วิสฏํ วา, เตน ภิกฺขเว สีโห มิคราชา อนตฺตมโน โหติ. สเจ ปน ภิกฺขเว สีโห มิคราชา น กิญฺจิ ปสฺสติ กายสฺส วิกฺขิตฺตํ วา วิสฏํ วา, เตน ภิกฺขเว สีโห มิคราชา อตฺตมโน โหติ. อยํ วุจฺจติ ภิกฺขเว สีหเสยฺยา.}

\prose{กตมา จ ภิกฺขเว ตถาคตเสยฺยา? อิธ ภิกฺขเว ตถาคโต วิวิจฺเจว กาเมหิ ฯเปฯ จตุตฺถํ ฌานํ อุปสมฺปชฺช วิหรติ. อยํ วุจฺจติ ภิกฺขเว ตถาคตเสยฺยา. อิมา โข ภิกฺขเว จตสฺโส เสยฺยาติ.  จตุตฺถํ.}

\csromanheader{2}{5. \textbf{ถูปารหสุตฺต}}
\proseitem{247}{จตฺตาโรเม ภิกฺขเว ถูปารหา. กตเม จตฺตาโร? ตถาคโต อรหํ สมฺมาสมฺพุทฺโธ ถูปารโห, ปจฺเจกพุทฺโธ ถูปารโห, ตถาคต\-สาวโก ถูปารโห, ราชา จกฺกวตฺตี ถูปารโห. อิเม โข ภิกฺขเว จตฺตาโร ถูปารหาติ.  ปญฺจมํ.}

\csromanmatikaanchor{mtk.271}
\csromantocmarkref{subsection}{6. ปญฺญาวุทฺธิสุตฺต}{mtk.271}
\bookmark[dest={mtk.271},level=2]{6. ปญฺญาวุทฺธิสุตฺต}
\csromanheader{2}{6. ปญฺญาวุทฺธิ\-สุตฺต}
\proseitem{248}{จตฺตาโรเม ภิกฺขเว ธมฺมา ปญฺญาวุทฺธิยา สํวตฺตนฺติ. กตเม จตฺตาโร? สปฺปุริส\-สํเสโว สทฺธมฺม\-สวนํ โยนิโส\-มนสิกาโร ธมฺมา\-นุธมฺม\-ปฺปฏิปตฺติ. อิเม โข ภิกฺขเว จตฺตาโร ธมฺมา ปญฺญาวุทฺธิยา สํวตฺตนฺตีติ.  ฉฏฺฐํ.}

\csromanpalirecto
\csromanmatikaanchor{mtk.272}
\csromantocmarkref{subsection}{7. พหุการสุตฺต}{mtk.272}
\bookmark[dest={mtk.272},level=2]{7. พหุการสุตฺต}
\csromanmarkh{จตุกฺกนิปาต\-ปาฬิ}\csromanmarkhii{7. พหุการสุตฺต}\csromanheadernomark{2}{จตุกฺกนิปาต\-ปาฬิ 7. พหุการสุตฺต}
\csromanmatikaanchor{mtk.273}
\csromantocmarkref{subsection}{8-11. โวหารสุตฺตานิ}{mtk.273}
\bookmark[dest={mtk.273},level=2]{8-11. โวหารสุตฺตานิ}
\proseitem{249}{\csromanfolio{568}จตฺตาโรเม ภิกฺขเว ธมฺมา มนุสฺส\-ภูตสฺส พหุการา โหนฺติ. กตเม จตฺตาโร? สปฺปุริส\-สํเสโว สทฺธมฺม\-สวนํ โยนิโส\-มนสิกาโร ธมฺมา\-นุธมฺม\-ปฺปฏิปตฺติ. อิเม โข ภิกฺขเว จตฺตาโร ธมฺมา มนุสฺส\-ภูตสฺส พหุการา โหนฺตีติ.  สตฺตมํ.}

\csromanheader{2}{8. ปฐม\-โวหาร\-สุตฺต}
\proseitem{250}{จตฺตาโรเม ภิกฺขเว อนริยโวหารา. กตเม จตฺตาโร? อทิฏฺเฐ ทิฏฺฐวาทิตา, อสุเต สุตวาทิตา, อมุเต มุตวาทิตา, อวิญฺญาเต วิญฺญาตวาทิตา. อิเม โข ภิกฺขเว จตฺตาโร อนริยโวหาราติ.  อฏฺฐมํ.}

\csromanheader{2}{9. ทุติย\-โวหาร\-สุตฺต}
\proseitem{251}{จตฺตาโรเม ภิกฺขเว อริยโวหารา. กตเม จตฺตาโร? อทิฏฺเฐ อทิฏฺฐวาทิตา, อสุเต อสุตวาทิตา, อมุเต อมุตวาทิตา, อวิญฺญาเต อวิญฺญาตวาทิตา. อิเม โข ภิกฺขเว จตฺตาโร อริยโวหาราติ.  นวมํ.}

\csromanheader{2}{10. ตติย\-โวหาร\-สุตฺต}
\proseitem{252}{จตฺตาโรเม ภิกฺขเว อนริยโวหารา. กตเม จตฺตาโร? ทิฏฺเฐ อทิฏฺฐวาทิตา, สุเต อสุตวาทิตา, มุเต อมุตวาทิตา, วิญฺญาเต อวิญฺญาตวาทิตา. อิเม โข ภิกฺขเว จตฺตาโร อนริยโวหาราติ.  ทสมํ.}

\csromanheader{2}{11. จตุตฺถ\-โวหาร\-สุตฺต}
\proseitemwithclosermidruled{253}{จตฺตาโรเม ภิกฺขเว อริยโวหารา. กตเม จตฺตาโร? ทิฏฺเฐ ทิฏฺฐวาทิตา, สุเต สุตวาทิตา, มุเต มุตวาทิตา, วิญฺญาเต วิญฺญาตวาทิตา. อิเม โข ภิกฺขเว จตฺตาโร อริยโวหาราติ.  เอกาทสมํ.}{อาปตฺติภย\-วคฺโค ปญฺจโม.}
\chapterheadpageread{(25) 5. อาปตฺติภย\-วคฺค ตสฺสุทฺทานํ}
\setlength{\csromangathapagewidth}{0pt}
\setlength{\csromangathawakwidth}{0pt}
\csromangathameasuregroup{เภทอาปตฺติ สิกฺขา จ, \\ ปญฺญาวุทฺธิ พหุการา,}{\csromangathabat{เภทอาปตฺติ สิกฺขา จ,}{เสยฺยา ถูปารเหน จ.} \\ \csromangathabat{ปญฺญาวุทฺธิ พหุการา,}{โวหารา จตุโร ฐิตาติ.}}
\csromangathasetleft{เภทอาปตฺติ สิกฺขา จ, \\ ปญฺญาวุทฺธิ พหุการา,}
\csromangathagroupwithclosermajor{\csromangathastanza{\csromanfolio{569}\csromangathabat{เภทอาปตฺติ สิกฺขา จ,}{เสยฺยา ถูปารเหน จ.} \\ \csromangathabat{ปญฺญาวุทฺธิ พหุการา,}{โวหารา จตุโร ฐิตาติ.}}}{\textbf{ปญฺจโม ปณฺณาสโก สมตฺโต.}}
\csromanpalirecto
\csromanmatikaanchor{mtk.274}
\csromantocmarknopage{section}{(26) 6. อภิญฺญาวคฺค}{mtk.274}
\bookmark[dest={mtk.274},level=1]{(26) 6. อภิญฺญาวคฺค}
\csromanmarkh{จตุกฺกนิปาต\-ปาฬิ (26)}\csromanmarkhii{6. อภิญฺญาวคฺค}\csromanheadernomark{1}{จตุกฺกนิปาต\-ปาฬิ (26) 6. อภิญฺญาวคฺค}
\csromanmatikaanchor{mtk.275}
\csromantocmarkref{subsection}{1. อภิญฺญาสุตฺต}{mtk.275}
\bookmark[dest={mtk.275},level=2]{1. อภิญฺญาสุตฺต}
\csromanheader{2}{1. \textbf{อภิญฺญาสุตฺต}}
\proseitem{254}{\csromanfolio{570}จตฺตาโรเม ภิกฺขเว ธมฺมา. กตเม จตฺตาโร? อตฺถิ ภิกฺขเว ธมฺมา อภิญฺญา ปริญฺเญยฺยา, อตฺถิ ภิกฺขเว ธมฺมา อภิญฺญา ปหาตพฺพา, อตฺถิ ภิกฺขเว ธมฺมา อภิญฺญา ภาเวตพฺพา, อตฺถิ ภิกฺขเว ธมฺมา อภิญฺญา สจฺฉิกาตพฺพา.}

\prose{กตเม จ ภิกฺขเว ธมฺมา อภิญฺญา ปริญฺเญยฺยา? ปญฺจุ\-ปาทาน\-กฺขนฺธา\csromansharedfootnote{4f9312ea79de2af6}{ปญฺจุปาทานกฺขนฺธาติสฺส วจนียํ (ก)}. อิเม วุจฺจนฺติ ภิกฺขเว ธมฺมา อภิญฺญา ปริญฺเญยฺยา.}

\prose{กตเม จ ภิกฺขเว ธมฺมา อภิญฺญา ปหาตพฺพา? อวิชฺชา จ ภวตณฺหา จ. อิเม วุจฺจนฺติ ภิกฺขเว ธมฺมา อภิญฺญา ปหาตพฺพา.}

\prose{กตเม จ ภิกฺขเว ธมฺมา อภิญฺญา ภาเวตพฺพา? สมโถ จ วิปสฺสนา จ. อิเม วุจฺจนฺติ ภิกฺขเว ธมฺมา อภิญฺญา ภาเวตพฺพา.}

\prose{กตเม จ ภิกฺขเว ธมฺมา อภิญฺญา สจฺฉิกาตพฺพา? วิชฺชา จ วิมุตฺติ จ. อิเม วุจฺจนฺติ ภิกฺขเว ธมฺมา อภิญฺญา สจฺฉิกาตพฺพา. อิเม โข ภิกฺขเว จตฺตาโร ธมฺมาติ.  ปฐมํ.}

\csromanmatikaanchor{mtk.276}
\csromantocmarkref{subsection}{2. ปริเยสนาสุตฺต}{mtk.276}
\bookmark[dest={mtk.276},level=2]{2. ปริเยสนาสุตฺต}
\csromanheader{2}{2. ปริเยสนา\-สุตฺต}
\proseitem{255}{จตสฺโส อิมา ภิกฺขเว อนริย\-ปริเยสนา. กตมา จตสฺโส? อิธ ภิกฺขเว เอกจฺโจ อตฺตนา ชราธมฺโม สมาโน ชรา\-ธมฺมํเยว ปริเยสติ, อตฺตนา พฺยาธิธมฺโม สมาโน พฺยาธิ\-ธมฺมํเยว ปริเยสติ, อตฺตนา มรณธมฺโม สมาโน มรณ\-ธมฺมํเยว ปริเยสติ, อตฺตนา สํกิเลส\-ธมฺโม สมาโน สํกิเลส\-ธมฺมํเยว ปริเยสติ. อิมา โข ภิกฺขเว จตสฺโส อนริย\-ปริเยสนา.}

\prose{จตสฺโส อิมา ภิกฺขเว อริย\-ปริเยสนา. กตมา จตสฺโส? อิธ ภิกฺขเว เอกจฺโจ อตฺตนา ชราธมฺโม สมาโน ชราธมฺเม อาทีนวํ วิทิตฺวา อชรํ อนุตฺตรํ โยคกฺเขมํ นิพฺพานํ ปริเยสติ, อตฺตนา พฺยาธิธมฺโม สมาโน พฺยาธิธมฺเม อาทีนวํ วิทิตฺวา อพฺยาธิํ อนุตฺตรํ}

\vspace{\tipitakaparskip}
\csromancenter{(26) 6. อภิญฺญาวคฺค โยคกฺเขมํ นิพฺพานํ ปริเยสติ, อตฺตนา มรณธมฺโม สมาโน มรณธมฺเม อาทีนวํ วิทิตฺวา อมตํ อนุตฺตรํ โยคกฺเขมํ นิพฺพานํ ปริเยสติ, อตฺตนา สํกิเลส\-ธมฺโม สมาโน สํกิเลส\-ธมฺเม อาทีนวํ วิทิตฺวา อสํกิลิฏฺฐํ อนุตฺตรํ โยคกฺเขมํ นิพฺพานํ ปริเยสติ. อิมา โข ภิกฺขเว จตสฺโส อริย\-ปริเยสนาติ.}
\csromanmatikaanchor{mtk.277}
\csromantocmarkref{subsection}{3. สงฺคหวตฺถุสุตฺต}{mtk.277}
\bookmark[dest={mtk.277},level=2]{3. สงฺคหวตฺถุสุตฺต}
\csromanheader{2}{3. สงฺคหวตฺถุ\-สุตฺต}
\proseitem{256}{\csromanfolio{571}จตฺตาริมานิ ภิกฺขเว สงฺคห\-วตฺถูนิ. กตมานิ จตฺตาริ? ทานํ เปยฺยวชฺชํ\csromansharedfootnote{d67598b9d9377d10}{ปิยวาจา (ก) เหฏฺฐา 341 ปิฏฺเฐปิ ปสฺสิตพฺพํ.} อตฺถจริยา สมานตฺตตา. อิมานิ โข ภิกฺขเว จตฺตาริ สงฺคหวตฺถูนีติ.}

\csromanmatikaanchor{mtk.278}
\csromantocmarkref{subsection}{4. มาลุกฺยปุตฺตสุตฺต}{mtk.278}
\bookmark[dest={mtk.278},level=2]{4. มาลุกฺยปุตฺตสุตฺต}
\csromanheader{2}{4. มาลุกฺยปุตฺต\-สุตฺต}
\proseitem{257}{อถ โข อายสฺมา มาลุกฺยปุตฺโต\csromansharedfootnote{ec854639dc76998b}{มาลุงฺกฺยปุตฺโต (สี, สฺยา, กํ, อิ)} เยน ภควา เตนุปสงฺกมิ, อุปสงฺกมิตฺวา ภควนฺตํ อภิวาเทตฺวา เอกมนฺตํ นิสีทิ, เอกมนฺตํ นิสินฺโน โข อายสฺมา มาลุกฺยปุตฺโต ภควนฺตํ เอตทโวจ\csromandash{}}

\prose{สาธุ เม ภนฺเต ภควา สํขิตฺเตน ธมฺมํ เทเสตุ, ยมหํ ภควโต ธมฺมํ สุตฺวา เอโก วูปกฏฺโฐ อปฺปมตฺโต อาตาปี ปหิตตฺโต วิหเรยฺยนฺติ. เอตฺถ อิทานิ มาลุกฺยปุตฺต กิํ ทหเร ภิกฺขู วกฺขาม. ยตฺร หิ นาม ตฺวํ ชิณฺโณ วุทฺโธ มหลฺลโก ตถาคตสฺส สํขิตฺเตน โอวาทํ ยาจสีติ. เทเสตุ เม ภนฺเต ภควา สํขิตฺเตน ธมฺมํ, เทเสตุ สุคโต สํขิตฺเตน ธมฺมํ, อปฺเปวนามาหํ ภควโต ภาสิตสฺส อตฺถํ อาชาเนยฺยํ, อปฺเปว นามาหํ ภควโต ภาสิตสฺส ทายาโท\csromansharedfootnote{186367fa18e1e607}{ภควโต สาวโก (ก)} อสฺสนฺติ.}

\csromanmatikaanchor{mtk.279}
\csromanmatikaanchor{mtk.280}
\csromantocmarkref{subsection}{5. กุลสุตฺต}{mtk.279}
\bookmark[dest={mtk.279},level=2]{5. กุลสุตฺต}
\csromantocmarkref{subsection}{6-7. อาชานียสุตฺตทฺวยํ}{mtk.280}
\bookmark[dest={mtk.280},level=2]{6-7. อาชานียสุตฺตทฺวยํ}
\prose{จตฺตาโรเม มาลุกฺยปุตฺต ตณฺหุปฺปาทา ยตฺถ ภิกฺขุโน ตณฺหา อุปฺปชฺชมานา อุปฺปชฺชติ. กตเม จตฺตาโร? จีวรเหตุ วา มาลุกฺยปุตฺต ภิกฺขุโน ตณฺหา อุปฺปชฺชมานา อุปฺปชฺชติ, ปิณฺฑปาตเหตุ วา มาลุกฺยปุตฺต ภิกฺขุโน ตณฺหา อุปฺปชฺชมานา อุปฺปชฺชติ, เสนาสนเหตุ วา มาลุกฺยปุตฺต ภิกฺขุโน ตณฺหา อุปฺปชฺชมานา อุปฺปชฺชติ, อิติภวา\-ภ\-ว\-เหตุ วา มาลุกฺยปุตฺต ภิกฺขุโน ตณฺหา อุปฺปชฺชมานา อุปฺปชฺชติ. อิเม โข มาลุกฺยปุตฺต จตฺตาโร ตณฺหุปฺปาทา ยตฺถ ภิกฺขุโน ตณฺหา อุปฺปชฺชมานา อุปฺปชฺชติ. ยโต โข \csromanfolio{572}มาลุกฺยปุตฺต ภิกฺขุโน ตณฺหา ปหีนา โหติ อุจฺฉินฺนมูลา ตาลาวตฺถุกตา อนภาวํกตา อายติํ อนุปฺปาทธมฺมา. อยํ วุจฺจติ มาลุกฺยปุตฺต ภิกฺขุ อจฺเฉจฺฉิ ตณฺหํ วิวตฺตยิ สํโยชนํ, สมฺมา มานาภิสมยา อนฺตมกาสิ ทุกฺขสฺสาติ.}

\prose{อถ โข อายสฺมา มาลุกฺยปุตฺโต ภควตา อิมินา โอวาเทน โอวทิโต อุฏฺฐายาสนา ภควนฺตํ อภิวาเทตฺวา ปทกฺขิณํ กตฺวา ปกฺกามิ. อถ โข อายสฺมา มาลุกฺยปุตฺโต เอโก วูปกฏฺโฐ อปฺปมตฺโต อาตาปี ปหิตตฺโต วิหรนฺโต นจิรสฺเสว ยสฺสตฺถาย กุลปุตฺตา สมฺมเทว อคารสฺมา อนคาริยํ ปพฺพชนฺติ, ตทนุตฺตรํ พฺรหฺมจริย\-ปริโยสานํ ทิฏฺเฐว ธมฺเม สยํ อภิญฺญา สจฺฉิกตฺวา อุปสมฺปชฺช วิหาสิ, “ขีณา ชาติ, วุสิตํ พฺรหฺมจริยํ, กตํ กรณียํ, นาปรํ อิตฺถตฺตายา”ติ อพฺภญฺญาสิ. อญฺญตโร จ ปนายสฺมา มาลุกฺยปุตฺโต อรหตํ อโหสีติ.  จตุตฺถํ.}

\csromanheader{2}{5. \textbf{กุลสุตฺต}}
\proseitem{258}{ยานิ กานิจิ ภิกฺขเว กุลานิ โภเคสุ มหตฺตํ ปตฺตานิ น จิรฏฺฐิติกานิ ภวนฺติ, สพฺพานิ ตานิ จตูหิ ฐาเนหิ เอเตสํ วา อญฺญตเรน. กตเมหิ จตูหิ? นฏฺฐํ น คเวสนฺติ, ชิณฺณํ น ปฏิสงฺขโรนฺติ, อปริมิต\-ปานโภชนา โหนฺติ, ทุสฺสีลํ อิตฺถิํ วา ปุริสํ วา อาธิปจฺเจ ฐเปนฺติ. ยานิ กานิจิ ภิกฺขเว กุลานิ โภเคสุ มหตฺตํ ปตฺตานิ น จิรฏฺฐิติกานิ ภวนฺติ, สพฺพานิ ตานิ อิเมหิ จตูหิ ฐาเนหิ เอเตสํ วา อญฺญตเรน.}

\prose{ยานิ กานิจิ ภิกฺขเว กุลานิ โภเคสุ มหตฺตํ ปตฺตานิ จิรฏฺฐิติกานิ ภวนฺติ, สพฺพานิ ตานิ จตูหิ ฐาเนหิ เอเตสํ วา อญฺญตเรน. กตเมหิ จตูหิ? นฏฺฐํ คเวสนฺติ, ชิณฺณํ ปฏิสงฺขโรนฺติ, ปริมิต\-ปานโภชนา โหนฺติ, สีลวนฺตํ อิตฺถิํ วา ปุริสํ วา อาธิปจฺเจ ฐเปนฺติ. ยานิ กานิจิ ภิกฺขเว กุลานิ โภเคสุ มหตฺตํ ปตฺตานิ จิรฏฺฐิติกานิ ภวนฺติ, สพฺพานิ ตานิ อิเมหิ จตูหิ ฐาเนหิ เอเตสํ วา อญฺญตเรนาติ.  ปญฺจมํ.}

\prose{6. ปฐม\-อาชานีย\-สุตฺต}

\proseitem{259}{จตูหิ ภิกฺขเว องฺเคหิ สมนฺนาคโต รญฺโญ ภโทฺร อสฺสาชานีโย ราชารโห โหติ ราชโภคฺโค, รญฺโญ องฺคนฺเตว}

\vspace{\tipitakaparskip}
\csromancenter{(26) 6. อภิญฺญาวคฺค สงฺขํ คจฺฉติ. กตเมหิ จตูหิ? อิธ ภิกฺขเว รญฺโญ ภโทฺร อสฺสาชานีโย วณฺณสมฺปนฺโน จ โหติ พลสมฺปนฺโน จ ชวสมฺปนฺโน จ อาโรหปริณาห\-สมฺปนฺโน จ. อิเมหิ โข ภิกฺขเว จตูหิ องฺเคหิ สมนฺนาคโต รญฺโญ ภโทฺร อสฺสาชานีโย ราชารโห โหติ ราชโภคฺโค, รญฺโญ องฺคนฺเตว สงฺขํ คจฺฉติ.}
\prose{\csromanfolio{573}เอวเมวํ โข ภิกฺขเว จตูหิ ธมฺเมหิ สมนฺนาคโต ภิกฺขุ อาหุเนยฺโย โหติ ฯเปฯ อนุตฺตรํ ปุญฺญกฺเขตฺตํ โลกสฺส. กตเมหิ จตูหิ? อิธ ภิกฺขเว ภิกฺขุ วณฺณสมฺปนฺโน จ โหติ พลสมฺปนฺโน จ ชวสมฺปนฺโน จ อาโรหปริณาห\-สมฺปนฺโน จ.}

\prose{กถญฺจ ภิกฺขเว ภิกฺขุ วณฺณสมฺปนฺโน โหติ? อิธ ภิกฺขเว ภิกฺขุ สีลวา โหติ ฯเปฯ สมาทาย สิกฺขติ สิกฺขาปเทสุ. เอวํ โข ภิกฺขเว ภิกฺขุ วณฺณสมฺปนฺโน โหติ.}

\prose{กถญฺจ ภิกฺขเว ภิกฺขุ พลสมฺปนฺโน โหติ? อิธ ภิกฺขเว ภิกฺขุ อารทฺธวีริโย วิหรติ อกุสลานํ ธมฺมานํ ปหานาย กุสลานํ ธมฺมานํ อุปสมฺปทาย ถามวา ทฬฺหปรกฺกโม อนิกฺขิตฺตธุโร กุสเลสุ ธมฺเมสุ. เอวํ โข ภิกฺขเว ภิกฺขุ พลสมฺปนฺโน โหติ.}

\prose{กถญฺจ ภิกฺขเว ภิกฺขุ ชวสมฺปนฺโน โหติ? อิธ ภิกฺขเว ภิกฺขุ “อิทํ ทุกฺขนฺ”ติ ยถาภูตํ ปชานาติ ฯเปฯ “อยํ ทุกฺขนิโรธ\-คามินี ปฏิปทา”ติ ยถาภูตํ ปชานาติ. เอวํ โข ภิกฺขเว ภิกฺขุ ชวสมฺปนฺโน โหติ.}

\prose{กถญฺจ ภิกฺขเว ภิกฺขุ อาโรหปริณาห\-สมฺปนฺโน โหติ? อิธ ภิกฺขเว ภิกฺขุ ลาภี โหติ จีวร\-ปิณฺฑปาต\-เสนาสน\-คิลานปจฺจย\-เภสชฺช\-ปริกฺขารานํ. เอวํ โข ภิกฺขเว ภิกฺขุ อาโรหปริณาห\-สมฺปนฺโน โหติ.}

\prose{อิเมหิ โข ภิกฺขเว จตูหิ ธมฺเมหิ สมนฺนาคโต ภิกฺขุ อาหุเนยฺโย โหติ ฯเปฯ อนุตฺตรํ ปุญฺญกฺเขตฺตํ โลกสฺสาติ.  ฉฏฺฐํ.}

\prose{7. ทุติย\-อาชานีย\-สุตฺต}

\proseitem{260}{จตูหิ ภิกฺขเว องฺเคหิ สมนฺนาคโต รญฺโญ ภโทฺร อสฺสาชานีโย ราชารโห โหติ ราชโภคฺโค, รญฺโญ องฺคนฺเตว \csromanfolio{574}สงฺขํ คจฺฉติ. กตเมหิ จตูหิ? อิธ ภิกฺขเว รญฺโญ ภโทฺร อสฺสาชานีโย วณฺณสมฺปนฺโน จ โหติ พลสมฺปนฺโน จ ชวสมฺปนฺโน จ อาโรหปริณาห\-สมฺปนฺโน จ. อิเมหิ โข ภิกฺขเว จตูหิ องฺเคหิ สมนฺนาคโต รญฺโญ ภโทฺร อสฺสาชานีโย ราชารโห โหติ ราชโภคฺโค, รญฺโญ องฺคนฺเตว สงฺขํ คจฺฉติ.}

\prose{เอวเมวํ โข ภิกฺขเว จตูหิ ธมฺเมหิ สมนฺนาคโต ภิกฺขุ อาหุเนยฺโย โหติ ฯเปฯ อนุตฺตรํ ปุญฺญกฺเขตฺตํ โลกสฺส. กตเมหิ จตูหิ? อิธ ภิกฺขเว ภิกฺขุ วณฺณสมฺปนฺโน จ โหติ พลสมฺปนฺโน จ ชวสมฺปนฺโน จ อาโรหปริณาห\-สมฺปนฺโน จ.}

\prose{กถญฺจ ภิกฺขเว ภิกฺขุ วณฺณสมฺปนฺโน โหติ? อิธ ภิกฺขเว ภิกฺขุ สีลวา โหติ ฯเปฯ สมาทาย สิกฺขติ สิกฺขาปเทสุ. เอวํ โข ภิกฺขเว ภิกฺขุ วณฺณสมฺปนฺโน โหติ.}

\prose{กถญฺจ ภิกฺขเว ภิกฺขุ พลสมฺปนฺโน โหติ? อิธ ภิกฺขเว ภิกฺขุ อารทฺธวีริโย วิหรติ อกุสลานํ ธมฺมานํ ปหานาย กุสลานํ ธมฺมานํ อุปสมฺปทาย ถามวา ทฬฺหปรกฺกโม อนิกฺขิตฺตธุโร กุสเลสุ ธมฺเมสุ. เอวํ โข ภิกฺขเว ภิกฺขุ พลสมฺปนฺโน โหติ.}

\prose{กถญฺจ ภิกฺขเว ภิกฺขุ ชวสมฺปนฺโน โหติ? อิธ ภิกฺขเว ภิกฺขุ อาสวานํ ขยา ฯเปฯ สจฺฉิกตฺวา อุปสมฺปชฺช วิหรติ. เอวํ โข ภิกฺขเว ภิกฺขุ ชวสมฺปนฺโน โหติ.}

\prose{กถญฺจ ภิกฺขเว ภิกฺขุ อาโรหปริณาห\-สมฺปนฺโน โหติ? อิธ ภิกฺขเว ภิกฺขุ ลาภี โหติ จีวร\-ปิณฺฑปาต\-เสนาสน\-คิลานปจฺจย\-เภสชฺช\-ปริกฺขารานํ. เอวํ โข ภิกฺขเว ภิกฺขุ อาโรหปริณาห\-สมฺปนฺโน โหติ.}

\prose{อิเมหิ โข ภิกฺขเว จตูหิ ธมฺเมหิ สมนฺนาคโต ภิกฺขุ อาหุเนยฺโย โหติ ฯเปฯ อนุตฺตรํ ปุญฺญกฺเขตฺตํ โลกสฺสาติ.  สตฺตมํ.}

\csromanmatikaanchor{mtk.281}
\csromantocmarkref{subsection}{8. พลสุตฺต}{mtk.281}
\bookmark[dest={mtk.281},level=2]{8. พลสุตฺต}
\csromanheader{2}{8. \textbf{พลสุตฺต}}
\proseitem{261}{จตฺตาริมานิ ภิกฺขเว พลานิ. กตมานิ จตฺตาริ? วีริยพลํ สติพลํ สมาธิพลํ ปญฺญาพลํ. อิมานิ โข ภิกฺขเว จตฺตาริ พลานีติ.  อฏฺฐมํ.}

\csromanmatikaanchor{mtk.282}
\csromantocmarkref{subsection}{9. อรญฺญสุตฺต}{mtk.282}
\bookmark[dest={mtk.282},level=2]{9. อรญฺญสุตฺต}
\csromanmarkh{(26)}\csromanmarkhii{6. อภิญฺญาวคฺค 9. อรญฺญสุตฺต}\csromanheadernomark{2}{(26) 6. อภิญฺญาวคฺค \textbf{9. อรญฺญสุตฺต}}
\proseitem{262}{\csromanfolio{575}จตูหิ ภิกฺขเว ธมฺเมหิ สมนฺนาคโต ภิกฺขุ นาลํ อรญฺญ\-วนปฺปตฺถานิ ปนฺตานิ เสนาสนานิ ปฏิเสวิตุํ. กตเมหิ จตูหิ? กามวิตกฺเกน, พฺยาปาท\-วิตกฺเกน, วิหิํสา\-วิตกฺเกน, ทุปฺปญฺโญ โหติ ชโฬ เอลมูโค. อิเมหิ โข ภิกฺขเว จตูหิ ธมฺเมหิ สมนฺนาคโต ภิกฺขุ นาลํ อรญฺญ\-วนปฺปตฺถานิ ปนฺตานิ เสนาสนานิ ปฏิเสวิตุํ.}

\prose{จตูหิ ภิกฺขเว ธมฺเมหิ สมนฺนาคโต ภิกฺขุ อลํ อรญฺญ\-วนปฺปตฺถานิ ปนฺตานิ เสนาสนานิ ปฏิเสวิตุํ. กตเมหิ จตูหิ? เนกฺขมฺม\-วิตกฺเกน, อพฺยาปาทวิตกฺเกน, อวิหิํสา\-วิตกฺเกน, ปญฺญวา โหติ อชโฬ อเนลมูโค. อิเมหิ โข ภิกฺขเว จตูหิ ธมฺเมหิ สมนฺนาคโต ภิกฺขุ อลํ อรญฺญ\-วนปฺปตฺถานิ ปนฺตานิ เสนาสนานิ ปฏิเสวิตุนฺติ.  นวมํ.}

\csromanmatikaanchor{mtk.283}
\csromantocmarkref{subsection}{10. กมฺมสุตฺต}{mtk.283}
\bookmark[dest={mtk.283},level=2]{10. กมฺมสุตฺต}
\csromanheader{2}{10. \textbf{กมฺมสุตฺต}}
\proseitem{263}{จตูหิ ภิกฺขเว ธมฺเมหิ สมนฺนาคโต พาโล อพฺยตฺโต อสปฺปุริโส ขตํ อุปหตํ อตฺตานํ ปริหรติ, สาวชฺโช จ โหติ สานุวชฺโช วิญฺญูนํ, พหุญฺจ อปุญฺญํ ปสวติ. กตเมหิ จตูหิ? สาวชฺเชน กายกมฺเมน, สาวชฺเชน วจีกมฺเมน, สาวชฺเชน มโนกมฺเมน, สาวชฺชาย ทิฏฺฐิยา. อิเมหิ โข ภิกฺขเว จตูหิ ธมฺเมหิ สมนฺนาคโต พาโล อพฺยตฺโต อสปฺปุริโส ขตํ อุปหตํ อตฺตานํ ปริหรติ, สาวชฺโช จ โหติ สานุวชฺโช วิญฺญูนํ, พหุญฺจ อปุญฺญํ ปสวติ.}

\prosewithclosermidruled{จตูหิ ภิกฺขเว ธมฺเมหิ สมนฺนาคโต ปณฺฑิโต วิยตฺโต สปฺปุริโส อกฺขตํ อนุปหตํ อตฺตานํ ปริหรติ, อนวชฺโช จ โหติ อนนุวชฺโช วิญฺญูนํ, พหุญฺจ ปุญฺญํ ปสวติ. กตเมหิ จตูหิ? อนวชฺเชน กายกมฺเมน, อนวชฺเชน วจีกมฺเมน, อนวชฺเชน มโนกมฺเมน, อนวชฺชาย ทิฏฺฐิยา. อิเมหิ โข ภิกฺขเว จตูหิ ธมฺเมหิ สมนฺนาคโต ปณฺฑิโต วิยตฺโต สปฺปุริโส อกฺขตํ อนุปหตํ อตฺตานํ ปริหรติ, อนวชฺโช จ โหติ อนนุวชฺโช วิญฺญูนํ, พหุญฺจ ปุญฺญํ ปสวตีติ.  ทสมํ.}{อภิญฺญาวคฺโค ฉฏฺโฐ.}
\csromanmatikaanchor{mtk.284}
\csromanmatikaanchor{mtk.285}
\csromantocmarknopage{section}{(27) 7. กมฺมปถวคฺค}{mtk.284}
\bookmark[dest={mtk.284},level=1]{(27) 7. กมฺมปถวคฺค}
\csromantocmarkref{subsection}{1-10. ปาณาติปาตีสุตฺตาทีนิ}{mtk.285}
\bookmark[dest={mtk.285},level=2]{1-10. ปาณาติปาตีสุตฺตาทีนิ}
\csromantocmarknopage{section}{(28) 8. ราคเปยฺยาล}{mtk.286}
\bookmark[dest={mtk.286},level=1]{(28) 8. ราคเปยฺยาล}
\setlength{\csromangathapagewidth}{0pt}
\setlength{\csromangathawakwidth}{0pt}
\csromangathameasuregroup{อภิญฺญา ปริเยสนา, \\ กุลํ เทฺว จ อาชานียา,}{\csromangathabat{อภิญฺญา ปริเยสนา,}{สงฺคหํ มาลุกฺยปุตฺโต.} \\ \csromangathabat{กุลํ เทฺว จ อาชานียา,}{พลํ อรญฺญกมฺมุนาติ.}}
\csromangathasetleft{อภิญฺญา ปริเยสนา, \\ กุลํ เทฺว จ อาชานียา,}
\csromangathagroup{\csromangathastanza{\csromanfolio{576}\csromangathabat{อภิญฺญา ปริเยสนา,}{สงฺคหํ มาลุกฺยปุตฺโต.} \\ \csromangathabat{กุลํ เทฺว จ อาชานียา,}{พลํ อรญฺญกมฺมุนาติ.}}}

\csromanmarkh{(27)}\csromanmarkhii{7. กมฺมปถ\-วคฺค}\csromanheadernomark{1}{(27) 7. กมฺมปถ\-วคฺค}
\csromanheader{2}{1. ปาณาติปาตี\-สุตฺต}
\proseitem{264}{จตูหิ ภิกฺขเว ธมฺเมหิ สมนฺนาคโต ยถาภตํ นิกฺขิตฺโต เอวํ นิรเย. กตเมหิ จตูหิ? อตฺตนา จ ปาณาติปาตี โหติ, ปรญฺจ ปาณาติปาเต สมาทเปติ, ปาณาติปาเต จ สมนุญฺโญ โหติ, ปาณาติปาตสฺส จ วณฺณํ ภาสติ. อิเมหิ โข ภิกฺขเว จตูหิ ธมฺเมหิ สมนฺนาคโต ยถาภตํ นิกฺขิตฺโต เอวํ นิรเย.}

\prose{จตูหิ ภิกฺขเว ธมฺเมหิ สมนฺนาคโต ยถาภตํ นิกฺขิตฺโต เอวํ สคฺเค. กตเมหิ จตูหิ? อตฺตนา จ ปาณาติปาตา ปฏิวิรโต โหติ, ปรญฺจ ปาณาติปาตา เวรมณิยา สมาทเปติ, ปาณาติปาตา เวรมณิยา จ สมนุญฺโญ โหติ, ปาณาติปาตา เวรมณิยา จ วณฺณํ ภาสติ. อิเมหิ โข ภิกฺขเว จตูหิ ธมฺเมหิ สมนฺนาคโต ยถาภตํ นิกฺขิตฺโต เอวํ สคฺเคติ.  ปฐมํ.}

\csromanheader{2}{2. อทินฺนาทายี\-สุตฺต}
\proseitem{265}{จตูหิ ภิกฺขเว ธมฺเมหิ สมนฺนาคโต ยถาภตํ นิกฺขิตฺโต เอวํ นิรเย. กตเมหิ จตูหิ? อตฺตนา จ อทินฺนาทายี โหติ, ปรญฺจ อทินฺนาทาเน สมาทเปติ, อทินฺนาทาเน จ สมนุญฺโญ โหติ, อทินฺนาทานสฺส จ วณฺณํ ภาสติ. อิเมหิ โข ฯเปฯ.}

\prose{อตฺตนา จ อทินฺนาทานา ปฏิวิรโต โหติ, ปรญฺจ อทินฺนาทานา เวรมณิยา สมาทเปติ, อทินฺนาทานา เวรมณิยา จ สมนุญฺโญ โหติ, อทินฺนาทานา เวรมณิยา จ วณฺณํ ภาสติ. อิเมหิ โข ภิกฺขเว.  ทุติยํ.}

\chapterheadpageread{(27) 7. กมฺมปถ\-วคฺค 3. มิจฺฉาจารี\-สุตฺต}
\proseitem{266}{\csromanfolio{577}ฯเปฯ. อตฺตนา จ กาเมสุ\-มิจฺฉาจารี โหติ, ปรญฺจ กาเมสุ\-มิจฺฉาจาเร สมาทเปติ, กาเมสุ\-มิจฺฉาจาเร จ สมนุญฺโญ โหติ, กาเมสุ\-มิจฺฉาจารสฺส จ วณฺณํ ภาสติ. อิเมหิ โข ฯเปฯ.}

\prose{อตฺตนา จ กาเมสุ\-มิจฺฉาจารา ปฏิวิรโต โหติ, ปรญฺจ กาเมสุ\-มิจฺฉาจารา เวรมณิยา สมาทเปติ, กาเมสุ\-มิจฺฉาจารา เวรมณิยา จ สมนุญฺโญ โหติ, กาเมสุ\-มิจฺฉาจารา เวรมณิยา จ วณฺณํ ภาสติ. อิเมหิ โข ฯเปฯ.  ตติยํ.}

\csromanheader{2}{4. \textbf{มุสาวาทีสุตฺต}}
\proseitem{267}{ฯเปฯ. อตฺตนา จ มุสาวาที โหติ, ปรญฺจ มุสาวาเท สมาทเปติ, มุสาวาเท จ สมนุญฺโญ โหติ, มุสาวาทสฺส จ วณฺณํ ภาสติ. อิเมหิ โข ฯเปฯ.}

\prose{อตฺตนา จ มุสาวาทา ปฏิวิรโต โหติ, ปรญฺจ มุสาวาทา เวรมณิยา สมาทเปติ, มุสาวาทา เวรมณิยา จ สมนุญฺโญ โหติ, มุสาวาทา เวรมณิยา จ วณฺณํ ภาสติ. อิเมหิ ฯเปฯ.  จตุตฺถํ.}

\csromanheader{2}{5. ปิสุณวาจา\-สุตฺต}
\proseitem{268}{ฯเปฯ. อตฺตนา จ ปิสุณวาโจ โหติ, ปรญฺจ ปิสุณาย วาจาย สมาทเปติ, ปิสุณาย วาจาย จ สมนุญฺโญ โหติ, ปิสุณาย วาจาย จ วณฺณํ ภาสติ. อิเมหิ ฯเปฯ.}

\prose{อตฺตนา จ ปิสุณาย วาจาย ปฏิวิรโต โหติ, ปรญฺจ ปิสุณาย วาจาย เวรมณิยา สมาทเปติ, ปิสุณาย วาจาย เวรมณิยา จ สมนุญฺโญ โหติ, ปิสุณาย วาจาย เวรมณิยา จ วณฺณํ ภาสติ. อิเมหิ ฯเปฯ.  ปญฺจมํ.}

\csromanheader{2}{6. ผรุสวาจา\-สุตฺต}
\proseitem{269}{ฯเปฯ. อตฺตนา จ ผรุสวาโจ โหติ, ปรญฺจ ผรุสาย วาจาย สมาทเปติ, ผรุสาย วาจาย จ สมนุญฺโญ โหติ, ผรุสาย วาจาย จ วณฺณํ ภาสติ ฯเปฯ. \csromanfolio{578}อตฺตนา จ ผรุสาย วาจาย ปฏิวิรโต โหติ, ปรญฺจ ผรุสาย วาจาย เวรมณิยา สมาทเปติ, ผรุสาย วาจาย เวรมณิยา จ สมนุญฺโญ โหติ, ผรุสาย วาจาย เวรมณิยา จ วณฺณํ ภาสติ. อิเมหิ โข ฯเปฯ.  ฉฏฺฐํ.}

\csromanheader{2}{7. สมฺผปฺปลาป\-สุตฺต}
\proseitem{270}{ฯเปฯ. อตฺตนา จ สมฺผปฺปลาปี โหติ, ปรญฺจ สมฺผปฺปลาเป สมาทเปติ, สมฺผปฺปลาเป จ สมนุญฺโญ โหติ, สมฺผ\-ปฺปลาปสฺส จ วณฺณํ ภาสติ. อิเมหิ ฯเปฯ.}

\prose{อตฺตนา จ สมฺผปฺปลาปา ปฏิวิรโต โหติ, ปรญฺจ สมฺผปฺปลาปา เวรมณิยา สมาทเปติ, สมฺผปฺปลาปา เวรมณิยา จ สมนุญฺโญ โหติ, สมฺผปฺปลาปา เวรมณิยา จ วณฺณํ ภาสติ. อิเมหิ โข ภิกฺขเว ฯเปฯ.  }

\csromanheader{2}{8. อภิชฺฌาลุ\-สุตฺต}
\proseitem{271}{ฯเปฯ. อตฺตนา จ อภิชฺฌาลุ โหติ, ปรญฺจ อภิชฺฌาย สมาทเปติ, อภิชฺฌาย จ สมนุญฺโญ โหติ, อภิชฺฌาย จ วณฺณํ ภาสติ ฯเปฯ.}

\prose{อตฺตนา จ อนภิชฺฌาลุ โหติ, ปรญฺจ อนภิชฺฌาย สมาทเปติ, อนภิชฺฌาย จ สมนุญฺโญ โหติ, อนภิชฺฌาย จ วณฺณํ ภาสติ. อิเมหิ โข ฯเปฯ.  อฏฺฐมํ.}

\csromanheader{2}{9. พฺยาปนฺนจิตฺต\-สุตฺต}
\proseitem{272}{ฯเปฯ. อตฺตนา จ พฺยาปนฺนจิตฺโต โหติ, ปรญฺจ พฺยาปาเท สมาทเปติ, พฺยาปาเท จ สมนุญฺโญ โหติ, พฺยาปาทสฺส จ วณฺณํ ภาสติ. อิเมหิ ฯเปฯ.}

\prose{อตฺตนา จ อพฺยาปนฺนจิตฺโต โหติ, ปรญฺจ อพฺยาปาเท สมาทเปติ, อพฺยาปาเท จ สมนุญฺโญ โหติ, อพฺยาปาทสฺส จ วณฺณํ ภาสติ. อิเมหิ ฯเปฯ.  นวมํ.}

\csromanheader{2}{10. มิจฺฉาทิฏฺฐิ\-สุตฺต}
\proseitem{273}{ฯเปฯ. อตฺตนา จ มิจฺฉาทิฏฺฐิโก โหติ, ปรญฺจ มิจฺฉาทิฏฺฐิยา สมาทเปติ, มิจฺฉาทิฏฺฐิยา จ สมนุญฺโญ โหติ, มิจฺฉาทิฏฺฐิยา จ วณฺณํ ภาสติ. อิเมหิ ฯเปฯ.}

\vspace{\tipitakaparskip}
\csromancenter{(28) 8. ราคเปยฺยาล}
\csromanmatikaanchor{mtk.287}
\csromantocmarkref{subsection}{1-510. สติปฏฺฐานสุตฺตาทีนิ}{mtk.287}
\bookmark[dest={mtk.287},level=2]{1-510. สติปฏฺฐานสุตฺตาทีนิ}
\prosewithclosermidruled{\csromanfolio{579}อตฺตนา จ สมฺมาทิฏฺฐิโก โหติ, ปรญฺจ สมฺมาทิฏฺฐิยา สมาทเปติ, สมฺมาทิฏฺฐิยา จ สมนุญฺโญ โหติ, สมฺมาทิฏฺฐิยา จ วณฺณํ ภาสติ. อิเมหิ โข ภิกฺขเว จตูหิ ธมฺเมหิ สมนฺนาคโต ยถาภตํ นิกฺขิตฺโต เอวํ สคฺเคติ.  ทสมํ.}{กมฺมปถ\-วคฺโค สตฺตโม.}
\csromancenter{(28) 8. ราคเปยฺยาล}
\csromanheader{2}{1. สติปฏฺฐาน\-สุตฺต}
\proseitem{274}{ราคสฺส ภิกฺขเว อภิญฺญาย จตฺตาโร ธมฺมา ภาเวตพฺพา. กตเม จตฺตาโร? อิธ ภิกฺขเว ภิกฺขุ กาเย กายานุปสฺสี วิหรติ อาตาปี สมฺปชาโน สติมา วิเนยฺย โลเก อภิชฺฌา\-โทมนสฺสํ. เวทนาสุ ฯเปฯ. จิตฺเต ฯเปฯ. ธมฺเมสุ ธมฺมานุปสฺสี วิหรติ อาตาปี สมฺปชาโน สติมา วิเนยฺย โลเก อภิชฺฌา\-โทมนสฺสํ. ราคสฺส ภิกฺขเว อภิญฺญาย อิเม จตฺตาโร ธมฺมา ภาเวตพฺพาติ.  ปฐมํ.}

\csromanheader{2}{2. สมฺมปฺปธาน\-สุตฺต}
\proseitem{275}{ราคสฺส ภิกฺขเว อภิญฺญาย จตฺตาโร ธมฺมา ภาเวตพฺพา. กตเม จตฺตาโร? อิธ ภิกฺขเว ภิกฺขุ อนุปฺปนฺนานํ ปาปกานํ อกุสลานํ ธมฺมานํ อนุปฺปาทาย ฉนฺทํ ชเนติ วายมติ, วีริยํ อารภติ, จิตฺตํ ปคฺคณฺหาติ ปทหติ. อุปฺปนฺนานํ ปาปกานํ อกุสลานํ ธมฺมานํ ปหานาย ฯเปฯ อนุปฺปนฺนานํ กุสลานํ ธมฺมานํ อุปฺปาทาย ฯเปฯ อุปฺปนฺนานํ กุสลานํ ธมฺมานํ ฐิติยา อสมฺโมสาย ภิยฺโยภาวาย เวปุลฺลาย ภาวนาย ปาริปูริยา ฉนฺทํ ชเนติ วายมติ, วีริยํ อารภติ, จิตฺตํ ปคฺคณฺหาติ ปทหติ. ราคสฺส ภิกฺขเว อภิญฺญาย อิเม จตฺตาโร ธมฺมา ภาเวตพฺพาติ.  ทุติยํ.}

\csromanheader{2}{3. \textbf{อิทฺธิปาทสุตฺต}}
\proseitem{276}{ราคสฺส ภิกฺขเว อภิญฺญาย จตฺตาโร ธมฺมา ภาเวตพฺพา. กตเม จตฺตาโร? อิธ ภิกฺขเว ภิกฺขุ ฉนฺท\-สมาธิ\-ปธาน\-สงฺขาร\-สมนฺนาคตํ \csromanfolio{580}อิทฺธิปาทํ ภาเวติ, วีริยสมาธิ ฯเปฯ จิตฺตสมาธิ ฯเปฯ วีมํสา\-สมาธิ\-ปธาน\-สงฺขาร\-สมนฺนาคตํ อิทฺธิปาทํ ภาเวติ. ราคสฺส ภิกฺขเว อภิญฺญาย อิเม จตฺตาโร ธมฺมา ภาเวตพฺพาติ.  ตติยํ.}

\csromanheader{2}{4-30. ปริญฺญาทิสุตฺต}
\proseitem{277-303}{ราคสฺส ภิกฺขเว ปริญฺญาย ฯเปฯ ปริกฺขยาย. ปหานาย. ขยาย. วยาย. วิราคาย. นิโรธาย. จาคาย. ปฏินิสฺสคฺคาย จตฺตาโร ธมฺมา ภาเวตพฺพา ฯเปฯ.  ติํสติมํ.}

\prose{31-510. โทสอภิญฺญาทิสุตฺต}

\proseitemwithcloser{304-783}{โทสสฺส ฯเปฯ. โมหสฺส. โกธสฺส. อุปนาหสฺส. มกฺขสฺส. ปฬาสสฺส. อิสฺสาย. มจฺฉริยสฺส. มายาย. สาเฐยฺยสฺส. ถมฺภสฺส. สารมฺภสฺส. มานสฺส. อติมานสฺส. มทสฺส. ปมาทสฺส อภิญฺญาย ปริญฺญาย ปริกฺขยาย ปหานาย ขยาย วยาย วิราคาย นิโรธาย จาคาย ปฏินิสฺสคฺคาย อิเม จตฺตาโร ธมฺมา ภาเวตพฺพาติ.  ทสุตฺตร\-ปญฺจสติมํ.}{ราคเปยฺยาลํ นิฏฺฐิตํ.}
\nitthitammajor{จตุกฺกนิปาต\-ปาฬิ นิฏฺฐิตา.}
\orphannote{ที 2. 92; ม 2. 205; อํ 3. 125 ปิฏฺเฐสุปิ.}

\orphannote{อภิ 3. 130 ปิฏฺเฐปิ.}

\orphannote{วิ 4. 304; สํ 1. 214 ปิฏฺเฐสุปิ.}

\orphannote{ขุ 2. 134, 152 ปิฏฺฐาทีสุปิ.}

\orphannote{อิมา อุทฺทานคาถาโย สีสฺยากํอิโปตฺถเกสุ น ทิสฺสนฺติ.}

\orphannote{สํ 2. 71 ปิฏฺเฐปิ.}

\orphannote{อภิ 3. 144 ปิฏฺเฐปิ.}

